\documentclass[11pt,a4paper]{article}
\usepackage[margin=25mm,headheight=15pt]{geometry}
\usepackage[T1]{fontenc}
\usepackage[utf8]{inputenc}
\usepackage{lmodern}
\usepackage{amsmath,amssymb,amsthm,amscd,mathtools,mathrsfs}
\usepackage{microtype,booktabs,longtable,array,xcolor,graphicx,enumitem,needspace,fancyhdr}
\usepackage[round,authoryear]{natbib}
\usepackage[unicode,colorlinks=true,linkcolor=blue!40!black,citecolor=blue!40!black,urlcolor=blue!40!black]{hyperref}
\usepackage{xurl}
\usepackage{listings}
\usepackage{tabularx}
\lstset{basicstyle=\ttfamily\small,breaklines=true,columns=fullflexible,keepspaces=true}
\usepackage{tikz}
\usepackage{tikz-cd}
\usetikzlibrary{arrows.meta,calc,patterns,positioning}
\usepackage{aliascnt}
\usepackage[nameinlink,noabbrev]{cleveref}
\allowdisplaybreaks[2]
\numberwithin{equation}{section}
\newcommand{\SsixRepositoryCommit}{83d76022bcd442278b760db9ff438afceed6c6ea}
\newcommand{\SsixPublishRecords}{https://github.com/KokunoYumeto/s6-complex-structure-workbench/tree/1fb9194767aef92abcc06071e3820c2f226b9512/pro_returns/S6_SERIES_PUBLISH_checkpoint_10_20261121T0612Z/}
\setcounter{MaxMatrixCols}{32}
\setlength{\emergencystretch}{3em}
\setlength{\parskip}{2pt}
\newtheorem{theorem}{Theorem}[section]
\newaliascnt{proposition}{theorem}
\newtheorem{proposition}[proposition]{Proposition}
\aliascntresetthe{proposition}
\crefname{proposition}{proposition}{propositions}
\newaliascnt{lemma}{theorem}
\newtheorem{lemma}[lemma]{Lemma}
\aliascntresetthe{lemma}
\crefname{lemma}{lemma}{lemmas}
\newaliascnt{corollary}{theorem}
\newtheorem{corollary}[corollary]{Corollary}
\aliascntresetthe{corollary}
\crefname{corollary}{corollary}{corollaries}
\newaliascnt{claim}{theorem}
\newtheorem{claim}[claim]{Claim}
\aliascntresetthe{claim}
\crefname{claim}{claim}{claims}
\newaliascnt{statement}{theorem}
\newtheorem{statement}[statement]{Statement}
\aliascntresetthe{statement}
\crefname{statement}{statement}{statements}
\theoremstyle{definition}
\newaliascnt{definition}{theorem}
\newtheorem{definition}[definition]{Definition}
\aliascntresetthe{definition}
\crefname{definition}{definition}{definitions}
\newaliascnt{remark}{theorem}
\newtheorem{remark}[remark]{Remark}
\aliascntresetthe{remark}
\crefname{remark}{remark}{remarks}
\newaliascnt{example}{theorem}
\newtheorem{example}[example]{Example}
\aliascntresetthe{example}
\crefname{example}{example}{examples}
\newaliascnt{problem}{theorem}
\newtheorem{problem}[problem]{Open problem}
\aliascntresetthe{problem}
\crefname{problem}{open problem}{open problems}
\newaliascnt{question}{theorem}
\newtheorem{question}[question]{Question}
\aliascntresetthe{question}
\crefname{question}{question}{questions}
\newaliascnt{auditrecord}{theorem}
\newtheorem{auditrecord}[auditrecord]{Verification record}
\aliascntresetthe{auditrecord}
\crefname{auditrecord}{verification record}{verification records}
\newaliascnt{test}{theorem}
\newtheorem{test}[test]{Verification}
\aliascntresetthe{test}
\crefname{test}{verification}{verifications}
\newaliascnt{conjecture}{theorem}
\newtheorem{conjecture}[conjecture]{Conjecture}
\aliascntresetthe{conjecture}
\crefname{conjecture}{conjecture}{conjectures}
\newcommand{\C}{\mathbb C}\newcommand{\R}{\mathbb R}\newcommand{\Z}{\mathbb Z}\newcommand{\Q}{\mathbb Q}\newcommand{\N}{\mathbb N}\newcommand{\F}{\mathbb F}
\newcommand{\CC}{\mathbb C}\newcommand{\RR}{\mathbb R}\newcommand{\ZZ}{\mathbb Z}\newcommand{\QQ}{\mathbb Q}\newcommand{\NN}{\mathbb N}
\newcommand{\PP}{\mathbb P}\newcommand{\Pp}{\mathbb P}\newcommand{\Pone}{\mathbb P^1}\newcommand{\HH}{\mathbb H}
\newcommand{\RP}{\mathbb{RP}}\newcommand{\CP}{\mathbb{CP}}
\newcommand{\OO}{\mathcal O}\newcommand{\cO}{\mathcal O}\newcommand{\cM}{\mathcal M}\newcommand{\cF}{\mathcal F}
\newcommand{\h}{\mathfrak h}\newcommand{\hz}{\mathfrak h_z}\newcommand{\Acon}{\mathfrak A}\newcommand{\wt}[1]{\widehat{#1}}
\newcommand{\red}{\mathrm{red}}\newcommand{\tor}{\mathrm{tor}}\newcommand{\dP}{\mathrm{dP}}
\newcommand{\gh}{\widehat\gamma}\newcommand{\uh}{\widehat u}\newcommand{\wh}{\widehat w}\newcommand{\dhat}{\widehat\delta}\renewcommand{\dh}{\widehat\delta}
\newcommand{\dd}{\mathrm d}\newcommand{\ii}{\mathrm i}\newcommand{\eps}{\varepsilon}\newcommand{\partt}{\mathrm{part}}\newcommand{\Bcirc}{B^\circ}\newcommand{\ddbar}{\partial\bar\partial}
\DeclareMathOperator{\Pic}{Pic}\DeclareMathOperator{\Alb}{Alb}\DeclareMathOperator{\Aut}{Aut}\DeclareMathOperator{\NS}{NS}
\DeclareMathOperator{\Hom}{Hom}\DeclareMathOperator{\End}{End}\DeclareMathOperator{\Ext}{Ext}\DeclareMathOperator{\Tor}{Tor}
\DeclareMathOperator{\Spec}{Spec}\DeclareMathOperator{\Proj}{Proj}\DeclareMathOperator{\Sym}{Sym}\DeclareMathOperator{\Diff}{Diff}
\DeclareMathOperator{\GL}{GL}\DeclareMathOperator{\SL}{SL}\DeclareMathOperator{\PSL}{PSL}\DeclareMathOperator{\Sp}{Sp}
\DeclareMathOperator{\im}{im}\DeclareMathOperator{\coker}{coker}\DeclareMathOperator{\rank}{rank}\DeclareMathOperator{\diag}{diag}
\DeclareMathOperator{\ord}{ord}\DeclareMathOperator{\vol}{vol}\DeclareMathOperator{\pr}{pr}\DeclareMathOperator{\id}{id}
\DeclareMathOperator{\divisor}{div}\DeclareMathOperator{\Tr}{Tr}\DeclareMathOperator{\Nm}{Nm}\DeclareMathOperator{\Sha}{Sha}
\DeclareMathOperator{\Imn}{Im}\DeclareMathOperator{\Ren}{Re}\DeclareMathOperator{\trdeg}{tr.deg}
\newcommand{\Uses}[1]{}
\newcommand{\Verified}{}\newcommand{\AuthorSelfCheck}{}\newcommand{\NeedsCheck}{}\newcommand{\Open}{\textbf{Open.}\ }

\newcommand{\RungZero}[1]{\par\noindent\textbf{Result and mechanism.} #1\par\smallskip}
\newcommand{\RungOne}{\par\Needspace{6\baselineskip}\noindent\textbf{Proof at the first rung.}\par\nobreak\smallskip}
\newcommand{\CurrentChecker}{GPT-6 Astra (OpenAI, ChatGPT Pro mode, effort Pro), 18 November 2026}
\newaliascnt{cited}{theorem}
\newtheorem{cited}[cited]{Cited result}
\aliascntresetthe{cited}
\crefname{cited}{cited result}{cited results}
\newaliascnt{refuted}{theorem}
\newtheorem{refuted}[refuted]{Correction}
\aliascntresetthe{refuted}
\crefname{refuted}{correction}{corrections}
\newaliascnt{computed}{theorem}
\newtheorem{computed}[computed]{Computation}
\aliascntresetthe{computed}
\crefname{computed}{computation}{computations}
\newaliascnt{openq}{theorem}
\newtheorem{openq}[openq]{Open question}
\aliascntresetthe{openq}
\crefname{openq}{open question}{open questions}

\providecommand{\rk}{\operatorname{rk}}
\providecommand{\Ann}{\operatorname{Ann}}
\providecommand{\Log}{\operatorname{Log}}
\providecommand{\vir}{\mathrm{vir}}

\usepackage{etoolbox}
\hypersetup{bookmarksnumbered=true,bookmarksdepth=2}
\setcounter{tocdepth}{2}
\makeatletter\renewcommand{\@pnumwidth}{2.5em}\renewcommand{\@tocrmarg}{3.5em}\renewcommand*{\l@subsection}{\@dottedtocline{2}{1.5em}{3.5em}}\makeatother
\newtheorem{maintheorem}{Theorem}
\renewcommand{\themaintheorem}{\Alph{maintheorem}}
\crefname{maintheorem}{theorem}{theorems}
\pretocmd{\section}{\Needspace{14\baselineskip}}{}{\PackageError{S6}{Section spacing hook failed}{}}
\pretocmd{\subsection}{\Needspace{9\baselineskip}}{}{\PackageError{S6}{Subsection spacing hook failed}{}}
\BeforeBeginEnvironment{theorem}{\Needspace{5\baselineskip}}
\BeforeBeginEnvironment{proposition}{\Needspace{5\baselineskip}}
\BeforeBeginEnvironment{lemma}{\Needspace{5\baselineskip}}
\BeforeBeginEnvironment{corollary}{\Needspace{5\baselineskip}}
\BeforeBeginEnvironment{claim}{\Needspace{5\baselineskip}}
\BeforeBeginEnvironment{statement}{\Needspace{5\baselineskip}}
\BeforeBeginEnvironment{definition}{\Needspace{5\baselineskip}}
\BeforeBeginEnvironment{remark}{\Needspace{5\baselineskip}}
\BeforeBeginEnvironment{example}{\Needspace{5\baselineskip}}
\BeforeBeginEnvironment{problem}{\Needspace{5\baselineskip}}
\BeforeBeginEnvironment{question}{\Needspace{5\baselineskip}}
\BeforeBeginEnvironment{auditrecord}{\Needspace{5\baselineskip}}
\BeforeBeginEnvironment{test}{\Needspace{5\baselineskip}}
\BeforeBeginEnvironment{conjecture}{\Needspace{5\baselineskip}}
\BeforeBeginEnvironment{cited}{\Needspace{5\baselineskip}}
\BeforeBeginEnvironment{refuted}{\Needspace{5\baselineskip}}
\BeforeBeginEnvironment{computed}{\Needspace{5\baselineskip}}
\BeforeBeginEnvironment{openq}{\Needspace{5\baselineskip}}
\BeforeBeginEnvironment{maintheorem}{\Needspace{5\baselineskip}}
\newcommand{\StatementRecord}[4]{{\upshape\footnotesize[tag~\texttt{#3}]}\enspace\ignorespaces}
\newcommand{\ExternalPaperRef}[3]{#1#2 (Paper~#3)}
\newcommand{\CompendiumRef}[2]{\ifmmode\text{\upshape S6\_09, Paper~#1, label~\nolinkurl{#2}}\else S6\_09, Paper~#1, label~\nolinkurl{#2}\fi}
\newcommand{\SeriesRefText}[1]{\ifcsname SRef@#1\endcsname\csname SRef@#1\endcsname\else\textbf{UNRESOLVED}\fi}
\newcommand{\SeriesRef}[1]{\ifmmode\text{\SeriesRefText{#1}}\else\SeriesRefText{#1}\fi}


\newcommand{\CurrentPaper}{I}
\expandafter\def\csname SRef@wb:pshift:eq:periods\endcsname{\hyperlink{tag-03CN}{I.\ref*{wb:pshift:eq:periods}}}
\expandafter\def\csname SRef@wb:pshift:eq:rank-one-map\endcsname{\hyperlink{tag-03CO}{I.\ref*{wb:pshift:eq:rank-one-map}}}
\expandafter\def\csname SRef@wb:pshift:eq:finite-data\endcsname{\hyperlink{tag-03C8}{I.\ref*{wb:pshift:eq:finite-data}}}
\expandafter\def\csname SRef@wb:pshift:eq:covariance\endcsname{\hyperlink{tag-03C2}{I.\ref*{wb:pshift:eq:covariance}}}
\expandafter\def\csname SRef@wb:pshift:eq:integer-shear\endcsname{\hyperlink{tag-03CF}{I.\ref*{wb:pshift:eq:integer-shear}}}
\expandafter\def\csname SRef@wb:pshift:eq:centralizer\endcsname{\hyperlink{tag-03BZ}{I.\ref*{wb:pshift:eq:centralizer}}}
\expandafter\def\csname SRef@wb:pshift:eq:period-morphism\endcsname{\hyperlink{tag-03CM}{I.\ref*{wb:pshift:eq:period-morphism}}}
\expandafter\def\csname SRef@wb:pshift:eq:norms\endcsname{\hyperlink{tag-03CK}{I.\ref*{wb:pshift:eq:norms}}}
\expandafter\def\csname SRef@wb:pshift:eq:finite-equation\endcsname{\hyperlink{tag-03C9}{I.\ref*{wb:pshift:eq:finite-equation}}}
\expandafter\def\csname SRef@wb:pshift:eq:finite-norm-obstruction\endcsname{\hyperlink{tag-03CB}{I.\ref*{wb:pshift:eq:finite-norm-obstruction}}}
\expandafter\def\csname SRef@wb:pshift:eq:finite-class\endcsname{\hyperlink{tag-03C7}{I.\ref*{wb:pshift:eq:finite-class}}}
\expandafter\def\csname SRef@wb:pshift:eq:local-corrections\endcsname{\hyperlink{tag-03CH}{I.\ref*{wb:pshift:eq:local-corrections}}}
\expandafter\def\csname SRef@wb:pshift:eq:local-matrix-identity\endcsname{\hyperlink{tag-03CI}{I.\ref*{wb:pshift:eq:local-matrix-identity}}}
\expandafter\def\csname SRef@wb:pshift:eq:finite-middle-sections\endcsname{\hyperlink{tag-03CA}{I.\ref*{wb:pshift:eq:finite-middle-sections}}}
\expandafter\def\csname SRef@wb:pshift:eq:local-affine-cocycle\endcsname{\hyperlink{tag-03CG}{I.\ref*{wb:pshift:eq:local-affine-cocycle}}}
\expandafter\def\csname SRef@wb:pshift:eq:actual-overlap-difference\endcsname{\hyperlink{tag-03BX}{I.\ref*{wb:pshift:eq:actual-overlap-difference}}}
\expandafter\def\csname SRef@wb:pshift:eq:cusp-periods\endcsname{\hyperlink{tag-03C5}{I.\ref*{wb:pshift:eq:cusp-periods}}}
\expandafter\def\csname SRef@wb:pshift:eq:cusp-actions-identical\endcsname{\hyperlink{tag-03C3}{I.\ref*{wb:pshift:eq:cusp-actions-identical}}}
\expandafter\def\csname SRef@wb:pshift:eq:cocycle-law\endcsname{\hyperlink{tag-03C0}{I.\ref*{wb:pshift:eq:cocycle-law}}}
\expandafter\def\csname SRef@wb:pshift:eq:cusp-cocycle\endcsname{\hyperlink{tag-03C4}{I.\ref*{wb:pshift:eq:cusp-cocycle}}}
\expandafter\def\csname SRef@wb:pshift:eq:cusp-translation\endcsname{\hyperlink{tag-03C6}{I.\ref*{wb:pshift:eq:cusp-translation}}}
\expandafter\def\csname SRef@wb:pshift:eq:affine-bundle-action\endcsname{\hyperlink{tag-03BY}{I.\ref*{wb:pshift:eq:affine-bundle-action}}}
\expandafter\def\csname SRef@wb:pshift:eq:vector-sheaf\endcsname{\hyperlink{tag-03CP}{I.\ref*{wb:pshift:eq:vector-sheaf}}}
\expandafter\def\csname SRef@wb:pshift:eq:no-analytic-obstruction\endcsname{\hyperlink{tag-03CJ}{I.\ref*{wb:pshift:eq:no-analytic-obstruction}}}
\expandafter\def\csname SRef@wb:pshift:eq:global-correction\endcsname{\hyperlink{tag-03CE}{I.\ref*{wb:pshift:eq:global-correction}}}
\expandafter\def\csname SRef@wb:pshift:eq:full-isomorphism\endcsname{\hyperlink{tag-03CD}{I.\ref*{wb:pshift:eq:full-isomorphism}}}
\expandafter\def\csname SRef@wb:pshift:eq:finite-overlap-identity\endcsname{\hyperlink{tag-03CC}{I.\ref*{wb:pshift:eq:finite-overlap-identity}}}
\expandafter\def\csname SRef@wb:pshift:eq:coherent-shift-action\endcsname{\hyperlink{tag-03C1}{I.\ref*{wb:pshift:eq:coherent-shift-action}}}
\expandafter\def\csname SRef@wb:pshift:eq:parameter-disc\endcsname{\hyperlink{tag-03CL}{I.\ref*{wb:pshift:eq:parameter-disc}}}
\expandafter\def\csname SRef@mapsI:quotient-eq-2\endcsname{\hyperlink{tag-02OZ}{I.\ref*{mapsI:quotient-eq-2}}}
\expandafter\def\csname SRef@mapsI:quotient-eq-3\endcsname{\hyperlink{tag-02P0}{I.\ref*{mapsI:quotient-eq-3}}}
\expandafter\def\csname SRef@mapsI:quotient-eq-4\endcsname{\hyperlink{tag-02P1}{I.\ref*{mapsI:quotient-eq-4}}}
\expandafter\def\csname SRef@mapsI:quotient-eq-5\endcsname{\hyperlink{tag-02P2}{I.\ref*{mapsI:quotient-eq-5}}}
\expandafter\def\csname SRef@mapsI:quotient-eq-7\endcsname{\hyperlink{tag-02P3}{I.\ref*{mapsI:quotient-eq-7}}}
\expandafter\def\csname SRef@app:rederive:fam-src-D.1.2\endcsname{\hyperlink{tag-000F}{I.\ref*{app:rederive:fam-src-D.1.2}}}
\expandafter\def\csname SRef@app:rederive:fam-src-D.1.6\endcsname{\hyperlink{tag-000G}{I.\ref*{app:rederive:fam-src-D.1.6}}}
\expandafter\def\csname SRef@app:rederive:fam-src-D.2.3\endcsname{\hyperlink{tag-000H}{I.\ref*{app:rederive:fam-src-D.2.3}}}
\expandafter\def\csname SRef@mapsI:real-eq-1\endcsname{\hyperlink{tag-02P4}{I.\ref*{mapsI:real-eq-1}}}
\expandafter\def\csname SRef@mapsI:real-eq-2\endcsname{\hyperlink{tag-02P4}{I.\ref*{mapsI:real-eq-2}}}
\expandafter\def\csname SRef@mapsI:real-eq-4\endcsname{\hyperlink{tag-02P5}{I.\ref*{mapsI:real-eq-4}}}
\expandafter\def\csname SRef@mapsI:real-eq-5\endcsname{\hyperlink{tag-02P6}{I.\ref*{mapsI:real-eq-5}}}
\expandafter\def\csname SRef@mapsI:real-eq-6\endcsname{\hyperlink{tag-02P6}{I.\ref*{mapsI:real-eq-6}}}
\expandafter\def\csname SRef@mapsI:real-eq-7\endcsname{\hyperlink{tag-02P7}{I.\ref*{mapsI:real-eq-7}}}
\expandafter\def\csname SRef@mapsI:real-eq-8\endcsname{\hyperlink{tag-02P8}{I.\ref*{mapsI:real-eq-8}}}
\expandafter\def\csname SRef@ari:open:norm\endcsname{\hyperlink{tag-003T}{I.\ref*{ari:open:norm}}}
\expandafter\def\csname SRef@push:src:C.3.5\endcsname{\hyperlink{tag-0322}{I.\ref*{push:src:C.3.5}}}
\expandafter\def\csname SRef@push:src:B.4.3\endcsname{\hyperlink{tag-030Z}{I.\ref*{push:src:B.4.3}}}
\expandafter\def\csname SRef@push:src:D.4.6\endcsname{\hyperlink{tag-032O}{I.\ref*{push:src:D.4.6}}}
\expandafter\def\csname SRef@inv:A1.0\endcsname{\hyperlink{tag-01RQ}{I.\ref*{inv:A1.0}}}
\expandafter\def\csname SRef@wb:lem:cusp-structure-sheaf-audit\endcsname{\hyperlink{tag-01RQ}{I.\ref*{wb:lem:cusp-structure-sheaf-audit}}}
\expandafter\def\csname SRef@wb:thm:structure-direct-images-audit\endcsname{\hyperlink{tag-01RQ}{I.\ref*{wb:thm:structure-direct-images-audit}}}
\expandafter\def\csname SRef@inv:C3.6\endcsname{\hyperlink{tag-01SJ}{I.\ref*{inv:C3.6}}}
\expandafter\def\csname SRef@push:src:E.3.1\endcsname{\hyperlink{tag-033I}{I.\ref*{push:src:E.3.1}}}
\expandafter\def\csname SRef@ari:C.5.1\endcsname{\hyperlink{tag-000U}{I.\ref*{ari:C.5.1}}}
\expandafter\def\csname SRef@ari:C.6.R4\endcsname{\hyperlink{tag-0013}{I.\ref*{ari:C.6.R4}}}
\expandafter\def\csname SRef@ari:C.4.1\endcsname{\hyperlink{tag-000P}{I.\ref*{ari:C.4.1}}}
\expandafter\def\csname SRef@ari:C.7.7\endcsname{\hyperlink{tag-001G}{I.\ref*{ari:C.7.7}}}
\expandafter\def\csname SRef@ari:C.7.8\endcsname{\hyperlink{tag-001H}{I.\ref*{ari:C.7.8}}}
\expandafter\def\csname SRef@push:src:A.4.5\endcsname{\hyperlink{tag-02ZV}{I.\ref*{push:src:A.4.5}}}
\expandafter\def\csname SRef@push:A45\endcsname{\hyperlink{tag-02ZV}{I.\ref*{push:A45}}}
\expandafter\def\csname SRef@push:src:B.3.5\endcsname{\hyperlink{tag-0309}{I.\ref*{push:src:B.3.5}}}
\expandafter\def\csname SRef@push:B35\endcsname{\hyperlink{tag-0309}{I.\ref*{push:B35}}}
\expandafter\def\csname SRef@push:src:B.4.1\endcsname{\hyperlink{tag-030D}{I.\ref*{push:src:B.4.1}}}
\expandafter\def\csname SRef@push:B41\endcsname{\hyperlink{tag-030D}{I.\ref*{push:B41}}}
\expandafter\def\csname SRef@push:src:B.2.2\endcsname{\hyperlink{tag-0303}{I.\ref*{push:src:B.2.2}}}
\expandafter\def\csname SRef@push:B22\endcsname{\hyperlink{tag-0303}{I.\ref*{push:B22}}}
\expandafter\def\csname SRef@push:src:B.5.1\endcsname{\hyperlink{tag-031A}{I.\ref*{push:src:B.5.1}}}
\expandafter\def\csname SRef@push:B51\endcsname{\hyperlink{tag-031A}{I.\ref*{push:B51}}}
\expandafter\def\csname SRef@push:src:B.2.3\endcsname{\hyperlink{tag-0304}{I.\ref*{push:src:B.2.3}}}
\expandafter\def\csname SRef@push:B23\endcsname{\hyperlink{tag-0304}{I.\ref*{push:B23}}}
\expandafter\def\csname SRef@push:src:D.3.2\endcsname{\hyperlink{tag-032I}{I.\ref*{push:src:D.3.2}}}
\expandafter\def\csname SRef@push:src:E.1.3\endcsname{\hyperlink{tag-032T}{I.\ref*{push:src:E.1.3}}}
\expandafter\def\csname SRef@fam:eq:cd-beta\endcsname{\hyperlink{tag-01P5}{I.\ref*{fam:eq:cd-beta}}}
\expandafter\def\csname SRef@ari:eq:ca-GHbounds\endcsname{\hyperlink{tag-0024}{I.\ref*{ari:eq:ca-GHbounds}}}
\expandafter\def\csname SRef@eq:P5L\endcsname{\hyperlink{tag-01OG}{I.\ref*{eq:P5L}}}
\expandafter\def\csname SRef@eq:P5floor\endcsname{\hyperlink{tag-01OI}{I.\ref*{eq:P5floor}}}
\expandafter\def\csname SRef@eq:P5normalization\endcsname{\hyperlink{tag-01OJ}{I.\ref*{eq:P5normalization}}}
\expandafter\def\csname SRef@eq:P5exponential\endcsname{\hyperlink{tag-01OH}{I.\ref*{eq:P5exponential}}}
\expandafter\def\csname SRef@eq:P5H2\endcsname{\hyperlink{tag-01OF}{I.\ref*{eq:P5H2}}}
\expandafter\def\csname SRef@ari:eq:cp-sigma\endcsname{\hyperlink{tag-003K}{I.\ref*{ari:eq:cp-sigma}}}
\expandafter\def\csname SRef@ari:eq:cp-klein\endcsname{\hyperlink{tag-003F}{I.\ref*{ari:eq:cp-klein}}}
\expandafter\def\csname SRef@ari:eq:cp-periods\endcsname{\hyperlink{tag-003G}{I.\ref*{ari:eq:cp-periods}}}
\expandafter\def\csname SRef@ari:eq:cp-v\endcsname{\hyperlink{tag-003L}{I.\ref*{ari:eq:cp-v}}}
\expandafter\def\csname SRef@ari:eq:cp-chi\endcsname{\hyperlink{tag-003E}{I.\ref*{ari:eq:cp-chi}}}
\expandafter\def\csname SRef@ari:eq:ca-derivatives\endcsname{\hyperlink{tag-002B}{I.\ref*{ari:eq:ca-derivatives}}}
\expandafter\def\csname SRef@ari:eq:cp-MTdim\endcsname{\hyperlink{tag-003D}{I.\ref*{ari:eq:cp-MTdim}}}
\expandafter\def\csname SRef@ari:eq:cp-positive-centre\endcsname{\hyperlink{tag-003H}{I.\ref*{ari:eq:cp-positive-centre}}}
\expandafter\def\csname SRef@ari:eq:cp-real-linear\endcsname{\hyperlink{tag-003I}{I.\ref*{ari:eq:cp-real-linear}}}
\expandafter\def\csname SRef@ari:eq:cp-relation\endcsname{\hyperlink{tag-003J}{I.\ref*{ari:eq:cp-relation}}}
\expandafter\def\csname SRef@ari:eq:ct-CMmult\endcsname{\hyperlink{tag-003Q}{I.\ref*{ari:eq:ct-CMmult}}}
\expandafter\def\csname SRef@ari:eq:ct-Gaussian\endcsname{\hyperlink{tag-003R}{I.\ref*{ari:eq:ct-Gaussian}}}
\expandafter\def\csname SRef@ari:eq:cd-slope\endcsname{\hyperlink{tag-0033}{I.\ref*{ari:eq:cd-slope}}}
\expandafter\def\csname SRef@ari:eq:cd-compatibility\endcsname{\hyperlink{tag-0032}{I.\ref*{ari:eq:cd-compatibility}}}
\expandafter\def\csname SRef@ari:eq:ce-ribet-r\endcsname{\hyperlink{tag-0036}{I.\ref*{ari:eq:ce-ribet-r}}}
\expandafter\def\csname SRef@ari:eq:ce-zeta\endcsname{\hyperlink{tag-0037}{I.\ref*{ari:eq:ce-zeta}}}
\expandafter\def\csname SRef@ari:eq:ca-delta-h\endcsname{\hyperlink{tag-0028}{I.\ref*{ari:eq:ca-delta-h}}}
\expandafter\def\csname SRef@ari:eq:ce-gap\endcsname{\hyperlink{tag-0035}{I.\ref*{ari:eq:ce-gap}}}
\expandafter\def\csname SRef@ari:eq:ce-explicit-witness\endcsname{\hyperlink{tag-0034}{I.\ref*{ari:eq:ce-explicit-witness}}}
\expandafter\def\csname SRef@ari:eq:cc-fixedfield\endcsname{\hyperlink{tag-002Q}{I.\ref*{ari:eq:cc-fixedfield}}}
\expandafter\def\csname SRef@ari:eq:cc-rankcriterion\endcsname{\hyperlink{tag-002X}{I.\ref*{ari:eq:cc-rankcriterion}}}
\expandafter\def\csname SRef@ari:eq:division-ratio-model\endcsname{\hyperlink{tag-003S}{I.\ref*{ari:eq:division-ratio-model}}}
\expandafter\def\csname SRef@ari:eq:ca-ladder\endcsname{\hyperlink{tag-002D}{I.\ref*{ari:eq:ca-ladder}}}
\expandafter\def\csname SRef@fam:eq:F-g\endcsname{\hyperlink{tag-01OT}{I.\ref*{fam:eq:F-g}}}
\expandafter\def\csname SRef@fam:eq:F-Gamma\endcsname{\hyperlink{tag-01OQ}{I.\ref*{fam:eq:F-Gamma}}}
\expandafter\def\csname SRef@fam:eq:F-law\endcsname{\hyperlink{tag-01OV}{I.\ref*{fam:eq:F-law}}}
\expandafter\def\csname SRef@fam:eq:F-SLpresentation\endcsname{\hyperlink{tag-01OS}{I.\ref*{fam:eq:F-SLpresentation}}}
\expandafter\def\csname SRef@fam:eq:F-H\endcsname{\hyperlink{tag-01OR}{I.\ref*{fam:eq:F-H}}}
\expandafter\def\csname SRef@fam:eq:F-qD\endcsname{\hyperlink{tag-01OW}{I.\ref*{fam:eq:F-qD}}}
\expandafter\def\csname SRef@fam:eq:cd-wp\endcsname{\hyperlink{tag-01P8}{I.\ref*{fam:eq:cd-wp}}}
\expandafter\def\csname SRef@fam:eq:F-tauaction\endcsname{\hyperlink{tag-01OX}{I.\ref*{fam:eq:F-tauaction}}}
\expandafter\def\csname SRef@fam:eq:F-j\endcsname{\hyperlink{tag-01OU}{I.\ref*{fam:eq:F-j}}}
\expandafter\def\csname SRef@fam:eq:cd-NS\endcsname{\hyperlink{tag-01P4}{I.\ref*{fam:eq:cd-NS}}}
\expandafter\def\csname SRef@fam:eq:cd-theta-sections\endcsname{\hyperlink{tag-01P7}{I.\ref*{fam:eq:cd-theta-sections}}}
\expandafter\def\csname SRef@mapsI:closing:eq:C-explicit-closing\endcsname{\hyperlink{tag-02OP}{I.\ref*{mapsI:closing:eq:C-explicit-closing}}}
\expandafter\def\csname SRef@mapsI:closing:eq:C-twelve-scalar\endcsname{\hyperlink{tag-02OQ}{I.\ref*{mapsI:closing:eq:C-twelve-scalar}}}
\expandafter\def\csname SRef@eq:d-period-normal\endcsname{\hyperlink{tag-01ON}{I.\ref*{eq:d-period-normal}}}
\expandafter\def\csname SRef@eq:d-limit-group\endcsname{\hyperlink{tag-01OM}{I.\ref*{eq:d-limit-group}}}
\expandafter\def\csname SRef@eq:d-dual-exact\endcsname{\hyperlink{tag-01OL}{I.\ref*{eq:d-dual-exact}}}
\expandafter\def\csname SRef@eq:P2H\endcsname{\hyperlink{tag-01O8}{I.\ref*{eq:P2H}}}
\expandafter\def\csname SRef@fam:eq:cd-norm\endcsname{\hyperlink{tag-01P6}{I.\ref*{fam:eq:cd-norm}}}
\expandafter\def\csname SRef@eq:P2zg\endcsname{\hyperlink{tag-01OC}{I.\ref*{eq:P2zg}}}
\expandafter\def\csname SRef@eq:P2F\endcsname{\hyperlink{tag-01O5}{I.\ref*{eq:P2F}}}
\expandafter\def\csname SRef@eq:P2FQ\endcsname{\hyperlink{tag-01O6}{I.\ref*{eq:P2FQ}}}
\expandafter\def\csname SRef@eq:P2FR\endcsname{\hyperlink{tag-01O7}{I.\ref*{eq:P2FR}}}
\expandafter\def\csname SRef@eq:P2chain\endcsname{\hyperlink{tag-01O9}{I.\ref*{eq:P2chain}}}
\expandafter\def\csname SRef@eq:P2split\endcsname{\hyperlink{tag-01OB}{I.\ref*{eq:P2split}}}
\expandafter\def\csname SRef@eq:P2ratio\endcsname{\hyperlink{tag-01OA}{I.\ref*{eq:P2ratio}}}
\expandafter\def\csname SRef@inv:eq:A-E-line\endcsname{\hyperlink{tag-01SZ}{I.\ref*{inv:eq:A-E-line}}}
\expandafter\def\csname SRef@inv:eq:A-H-images\endcsname{\hyperlink{tag-01T0}{I.\ref*{inv:eq:A-H-images}}}
\expandafter\def\csname SRef@inv:eq:A-N-images\endcsname{\hyperlink{tag-01T1}{I.\ref*{inv:eq:A-N-images}}}
\expandafter\def\csname SRef@ari:eq:cs-isogeny\endcsname{\hyperlink{tag-003O}{I.\ref*{ari:eq:cs-isogeny}}}
\expandafter\def\csname SRef@inv:eq:A-absolute-table\endcsname{\hyperlink{tag-01T2}{I.\ref*{inv:eq:A-absolute-table}}}
\expandafter\def\csname SRef@inv:eq:A-s-branch\endcsname{\hyperlink{tag-01T5}{I.\ref*{inv:eq:A-s-branch}}}
\expandafter\def\csname SRef@inv:eq:A-modular-K\endcsname{\hyperlink{tag-01T3}{I.\ref*{inv:eq:A-modular-K}}}
\expandafter\def\csname SRef@inv:eq:A-mu-value\endcsname{\hyperlink{tag-01T4}{I.\ref*{inv:eq:A-mu-value}}}
\expandafter\def\csname SRef@inv:eq:A-DI-branch\endcsname{\hyperlink{tag-01SY}{I.\ref*{inv:eq:A-DI-branch}}}
\expandafter\def\csname SRef@inv:eq:C-hodge-input\endcsname{\hyperlink{tag-01TF}{I.\ref*{inv:eq:C-hodge-input}}}
\expandafter\def\csname SRef@inv:eq:C-BC-A\endcsname{\hyperlink{tag-01T6}{I.\ref*{inv:eq:C-BC-A}}}
\expandafter\def\csname SRef@inv:eq:C-arrows\endcsname{\hyperlink{tag-01TA}{I.\ref*{inv:eq:C-arrows}}}
\expandafter\def\csname SRef@inv:eq:C-Pic-map\endcsname{\hyperlink{tag-01T8}{I.\ref*{inv:eq:C-Pic-map}}}
\expandafter\def\csname SRef@inv:eq:C-G-Green\endcsname{\hyperlink{tag-01T7}{I.\ref*{inv:eq:C-G-Green}}}
\expandafter\def\csname SRef@ari:eq:cs-division\endcsname{\hyperlink{tag-003N}{I.\ref*{ari:eq:cs-division}}}
\expandafter\def\csname SRef@inv:eq:C-factor12\endcsname{\hyperlink{tag-01TE}{I.\ref*{inv:eq:C-factor12}}}
\expandafter\def\csname SRef@inv:eq:C-collar-sum\endcsname{\hyperlink{tag-01TB}{I.\ref*{inv:eq:C-collar-sum}}}
\expandafter\def\csname SRef@inv:eq:C-delta\endcsname{\hyperlink{tag-01TC}{I.\ref*{inv:eq:C-delta}}}
\expandafter\def\csname SRef@inv:eq:C-Ricci\endcsname{\hyperlink{tag-01T9}{I.\ref*{inv:eq:C-Ricci}}}
\expandafter\def\csname SRef@inv:eq:C-eta-negative\endcsname{\hyperlink{tag-01TD}{I.\ref*{inv:eq:C-eta-negative}}}
\expandafter\def\csname SRef@ari:eq:cs-Li\endcsname{\hyperlink{tag-003M}{I.\ref*{ari:eq:cs-Li}}}
\expandafter\def\csname SRef@ari:eq:ch-global\endcsname{\hyperlink{tag-0039}{I.\ref*{ari:eq:ch-global}}}
\expandafter\def\csname SRef@ari:eq:ch-local-mult\endcsname{\hyperlink{tag-003B}{I.\ref*{ari:eq:ch-local-mult}}}
\expandafter\def\csname SRef@ari:eq:ch-local-exact\endcsname{\hyperlink{tag-003A}{I.\ref*{ari:eq:ch-local-exact}}}
\expandafter\def\csname SRef@ari:eq:ch-bound\endcsname{\hyperlink{tag-0038}{I.\ref*{ari:eq:ch-bound}}}
\expandafter\def\csname SRef@ari:eq:cs-torsionvalue\endcsname{\hyperlink{tag-003P}{I.\ref*{ari:eq:cs-torsionvalue}}}
\expandafter\def\csname SRef@ari:eq:ch-wp-bracket\endcsname{\hyperlink{tag-003C}{I.\ref*{ari:eq:ch-wp-bracket}}}
\expandafter\def\csname SRef@ari:eq:Cr-first\endcsname{\hyperlink{tag-0020}{I.\ref*{ari:eq:Cr-first}}}
\expandafter\def\csname SRef@ari:eq:Cr-second\endcsname{\hyperlink{tag-0022}{I.\ref*{ari:eq:Cr-second}}}
\expandafter\def\csname SRef@ari:eq:Cr-tails\endcsname{\hyperlink{tag-0023}{I.\ref*{ari:eq:Cr-tails}}}
\expandafter\def\csname SRef@ari:eq:Cr-rational-factor\endcsname{\hyperlink{tag-0021}{I.\ref*{ari:eq:Cr-rational-factor}}}
\expandafter\def\csname SRef@ari:eq:Cgz-central-integrality\endcsname{\hyperlink{tag-001X}{I.\ref*{ari:eq:Cgz-central-integrality}}}
\expandafter\def\csname SRef@ari:eq:Cgz-product\endcsname{\hyperlink{tag-001Z}{I.\ref*{ari:eq:Cgz-product}}}
\expandafter\def\csname SRef@ari:eq:Cgz-denominator\endcsname{\hyperlink{tag-001Y}{I.\ref*{ari:eq:Cgz-denominator}}}
\expandafter\def\csname SRef@ari:eq:ca-mahler\endcsname{\hyperlink{tag-002E}{I.\ref*{ari:eq:ca-mahler}}}
\expandafter\def\csname SRef@ari:eq:ca-wpmean\endcsname{\hyperlink{tag-002H}{I.\ref*{ari:eq:ca-wpmean}}}
\expandafter\def\csname SRef@ari:eq:ca-sup\endcsname{\hyperlink{tag-002G}{I.\ref*{ari:eq:ca-sup}}}
\expandafter\def\csname SRef@ari:eq:ca-nu\endcsname{\hyperlink{tag-002F}{I.\ref*{ari:eq:ca-nu}}}
\expandafter\def\csname SRef@ari:eq:ca-current\endcsname{\hyperlink{tag-0026}{I.\ref*{ari:eq:ca-current}}}
\expandafter\def\csname SRef@ari:eq:ca-density0\endcsname{\hyperlink{tag-0029}{I.\ref*{ari:eq:ca-density0}}}
\expandafter\def\csname SRef@ari:eq:ca-density1\endcsname{\hyperlink{tag-002A}{I.\ref*{ari:eq:ca-density1}}}
\expandafter\def\csname SRef@ari:eq:ca-cuspexp\endcsname{\hyperlink{tag-0027}{I.\ref*{ari:eq:ca-cuspexp}}}
\expandafter\def\csname SRef@ari:eq:cc-EZ\endcsname{\hyperlink{tag-002I}{I.\ref*{ari:eq:cc-EZ}}}
\expandafter\def\csname SRef@ari:eq:cc-radial\endcsname{\hyperlink{tag-002W}{I.\ref*{ari:eq:cc-radial}}}
\expandafter\def\csname SRef@ari:eq:cc-FR\endcsname{\hyperlink{tag-002J}{I.\ref*{ari:eq:cc-FR}}}
\expandafter\def\csname SRef@ari:eq:cc-cuspbound\endcsname{\hyperlink{tag-002O}{I.\ref*{ari:eq:cc-cuspbound}}}
\expandafter\def\csname SRef@ari:eq:cc-ode\endcsname{\hyperlink{tag-002T}{I.\ref*{ari:eq:cc-ode}}}
\expandafter\def\csname SRef@ari:eq:ca-exp\endcsname{\hyperlink{tag-002C}{I.\ref*{ari:eq:ca-exp}}}
\expandafter\def\csname SRef@ari:eq:cc-h\endcsname{\hyperlink{tag-002R}{I.\ref*{ari:eq:cc-h}}}
\expandafter\def\csname SRef@ari:eq:cc-Fseries\endcsname{\hyperlink{tag-002L}{I.\ref*{ari:eq:cc-Fseries}}}
\expandafter\def\csname SRef@ari:eq:cc-Ffactored\endcsname{\hyperlink{tag-002K}{I.\ref*{ari:eq:cc-Ffactored}}}
\expandafter\def\csname SRef@ari:eq:cc-norms\endcsname{\hyperlink{tag-002S}{I.\ref*{ari:eq:cc-norms}}}
\expandafter\def\csname SRef@ari:eq:cc-recurrence\endcsname{\hyperlink{tag-002Y}{I.\ref*{ari:eq:cc-recurrence}}}
\expandafter\def\csname SRef@ari:eq:cc-power\endcsname{\hyperlink{tag-002V}{I.\ref*{ari:eq:cc-power}}}
\expandafter\def\csname SRef@ari:eq:cc-constants\endcsname{\hyperlink{tag-002N}{I.\ref*{ari:eq:cc-constants}}}
\expandafter\def\csname SRef@ari:eq:cc-outer\endcsname{\hyperlink{tag-002U}{I.\ref*{ari:eq:cc-outer}}}
\expandafter\def\csname SRef@ari:eq:cc-tail60\endcsname{\hyperlink{tag-0031}{I.\ref*{ari:eq:cc-tail60}}}
\expandafter\def\csname SRef@ari:eq:cc-tail100\endcsname{\hyperlink{tag-0030}{I.\ref*{ari:eq:cc-tail100}}}
\expandafter\def\csname SRef@ari:eq:ca-H\endcsname{\hyperlink{tag-0025}{I.\ref*{ari:eq:ca-H}}}
\expandafter\def\csname SRef@ari:eq:cc-finite-table\endcsname{\hyperlink{tag-002P}{I.\ref*{ari:eq:cc-finite-table}}}
\expandafter\def\csname SRef@ari:eq:cc-scalar\endcsname{\hyperlink{tag-002Z}{I.\ref*{ari:eq:cc-scalar}}}
\expandafter\def\csname SRef@ari:eq:cc-P\endcsname{\hyperlink{tag-002M}{I.\ref*{ari:eq:cc-P}}}
\expandafter\def\csname SRef@fam:eq:ab-root-degree\endcsname{\hyperlink{tag-01P3}{I.\ref*{fam:eq:ab-root-degree}}}
\expandafter\def\csname SRef@fam:eq:ab-cusp-resolution\endcsname{\hyperlink{tag-01OY}{I.\ref*{fam:eq:ab-cusp-resolution}}}
\expandafter\def\csname SRef@fam:eq:ab-log-tangent\endcsname{\hyperlink{tag-01P0}{I.\ref*{fam:eq:ab-log-tangent}}}
\expandafter\def\csname SRef@fam:eq:ab-relative-dualizing\endcsname{\hyperlink{tag-01P2}{I.\ref*{fam:eq:ab-relative-dualizing}}}
\expandafter\def\csname SRef@fam:eq:ab-direct-images\endcsname{\hyperlink{tag-01OZ}{I.\ref*{fam:eq:ab-direct-images}}}
\expandafter\def\csname SRef@fam:eq:ab-normal-quotient\endcsname{\hyperlink{tag-01P1}{I.\ref*{fam:eq:ab-normal-quotient}}}
\expandafter\def\csname SRef@eq:P5tangent\endcsname{\hyperlink{tag-01OK}{I.\ref*{eq:P5tangent}}}
\expandafter\def\csname SRef@fig:I:orbifold\endcsname{\hyperlink{tag-01RL}{I.\ref*{fig:I:orbifold}}}
\expandafter\def\csname SRef@fig:I:parameter-disc\endcsname{\hyperlink{tag-01RM}{I.\ref*{fig:I:parameter-disc}}}
\expandafter\def\csname SRef@fig:I:monodromy-domain\endcsname{\hyperlink{tag-01RK}{I.\ref*{fig:I:monodromy-domain}}}
\expandafter\def\csname SRef@push:src:A.3.1\endcsname{\hyperlink{tag-02ZO}{I.\ref*{push:src:A.3.1}}}
\expandafter\def\csname SRef@push:A31\endcsname{\hyperlink{tag-02ZO}{I.\ref*{push:A31}}}
\expandafter\def\csname SRef@push:src:D.1.2\endcsname{\hyperlink{tag-032B}{I.\ref*{push:src:D.1.2}}}
\expandafter\def\csname SRef@push:src:D.2.1\endcsname{\hyperlink{tag-032D}{I.\ref*{push:src:D.2.1}}}
\expandafter\def\csname SRef@I:weierstrass-comparison\endcsname{\hyperlink{tag-0002}{I.\ref*{I:weierstrass-comparison}}}
\expandafter\def\csname SRef@wb:pshift:lem:centralizer\endcsname{\hyperlink{tag-03CQ}{I.\ref*{wb:pshift:lem:centralizer}}}
\expandafter\def\csname SRef@fam:src:F.1.2\endcsname{\hyperlink{tag-01RA}{I.\ref*{fam:src:F.1.2}}}
\expandafter\def\csname SRef@fam:st:F.1.2\endcsname{\hyperlink{tag-01RA}{I.\ref*{fam:st:F.1.2}}}
\expandafter\def\csname SRef@fam:src:F.2.1\endcsname{\hyperlink{tag-01RB}{I.\ref*{fam:src:F.2.1}}}
\expandafter\def\csname SRef@fam:st:F.2.1\endcsname{\hyperlink{tag-01RB}{I.\ref*{fam:st:F.2.1}}}
\expandafter\def\csname SRef@fam:src:F.2.3\endcsname{\hyperlink{tag-01RD}{I.\ref*{fam:src:F.2.3}}}
\expandafter\def\csname SRef@fam:st:F.2.3\endcsname{\hyperlink{tag-01RD}{I.\ref*{fam:st:F.2.3}}}
\expandafter\def\csname SRef@fam:src:F.3.2\endcsname{\hyperlink{tag-01RF}{I.\ref*{fam:src:F.3.2}}}
\expandafter\def\csname SRef@fam:st:F.3.2\endcsname{\hyperlink{tag-01RF}{I.\ref*{fam:st:F.3.2}}}
\expandafter\def\csname SRef@fam:src:F.3.3\endcsname{\hyperlink{tag-01RG}{I.\ref*{fam:src:F.3.3}}}
\expandafter\def\csname SRef@fam:st:F.3.3\endcsname{\hyperlink{tag-01RG}{I.\ref*{fam:st:F.3.3}}}
\expandafter\def\csname SRef@fam:src:D.1.1\endcsname{\hyperlink{tag-01QR}{I.\ref*{fam:src:D.1.1}}}
\expandafter\def\csname SRef@fam:st:D.1.1\endcsname{\hyperlink{tag-01QR}{I.\ref*{fam:st:D.1.1}}}
\expandafter\def\csname SRef@fam:src:C.3.1\endcsname{\hyperlink{tag-01QM}{I.\ref*{fam:src:C.3.1}}}
\expandafter\def\csname SRef@fam:st:C.3.1\endcsname{\hyperlink{tag-01QM}{I.\ref*{fam:st:C.3.1}}}
\expandafter\def\csname SRef@fam:src:C.3.2\endcsname{\hyperlink{tag-01QN}{I.\ref*{fam:src:C.3.2}}}
\expandafter\def\csname SRef@fam:st:C.3.2\endcsname{\hyperlink{tag-01QN}{I.\ref*{fam:st:C.3.2}}}
\expandafter\def\csname SRef@fam:src:C.3.4\endcsname{\hyperlink{tag-01QP}{I.\ref*{fam:src:C.3.4}}}
\expandafter\def\csname SRef@fam:st:C.3.4\endcsname{\hyperlink{tag-01QP}{I.\ref*{fam:st:C.3.4}}}
\expandafter\def\csname SRef@fam:src:D.1.2\endcsname{\hyperlink{tag-01QS}{I.\ref*{fam:src:D.1.2}}}
\expandafter\def\csname SRef@fam:st:D.1.2\endcsname{\hyperlink{tag-01QS}{I.\ref*{fam:st:D.1.2}}}
\expandafter\def\csname SRef@fam:src:D.1.4\endcsname{\hyperlink{tag-01QU}{I.\ref*{fam:src:D.1.4}}}
\expandafter\def\csname SRef@fam:st:D.1.4\endcsname{\hyperlink{tag-01QU}{I.\ref*{fam:st:D.1.4}}}
\expandafter\def\csname SRef@fam:src:D.1.5\endcsname{\hyperlink{tag-01QV}{I.\ref*{fam:src:D.1.5}}}
\expandafter\def\csname SRef@fam:st:D.1.5\endcsname{\hyperlink{tag-01QV}{I.\ref*{fam:st:D.1.5}}}
\expandafter\def\csname SRef@fam:src:D.1.7\endcsname{\hyperlink{tag-01QX}{I.\ref*{fam:src:D.1.7}}}
\expandafter\def\csname SRef@fam:st:D.1.7\endcsname{\hyperlink{tag-01QX}{I.\ref*{fam:st:D.1.7}}}
\expandafter\def\csname SRef@fam:src:D.2.1\endcsname{\hyperlink{tag-01QY}{I.\ref*{fam:src:D.2.1}}}
\expandafter\def\csname SRef@fam:st:D.2.1\endcsname{\hyperlink{tag-01QY}{I.\ref*{fam:st:D.2.1}}}
\expandafter\def\csname SRef@fam:src:D.2.2\endcsname{\hyperlink{tag-01QZ}{I.\ref*{fam:src:D.2.2}}}
\expandafter\def\csname SRef@fam:st:D.2.2\endcsname{\hyperlink{tag-01QZ}{I.\ref*{fam:st:D.2.2}}}
\expandafter\def\csname SRef@mapsI:src:A.3.1\endcsname{\hyperlink{tag-02PE}{I.\ref*{mapsI:src:A.3.1}}}
\expandafter\def\csname SRef@mapsI:src:A.3.4\endcsname{\hyperlink{tag-02PH}{I.\ref*{mapsI:src:A.3.4}}}
\expandafter\def\csname SRef@mapsI:real-eq-9\endcsname{\hyperlink{tag-02PH}{I.\ref*{mapsI:real-eq-9}}}
\expandafter\def\csname SRef@I:monodromy-fundamental-set\endcsname{\hyperlink{tag-0001}{I.\ref*{I:monodromy-fundamental-set}}}
\expandafter\def\csname SRef@fam:src:E.1.1\endcsname{\hyperlink{tag-01R1}{I.\ref*{fam:src:E.1.1}}}
\expandafter\def\csname SRef@fam:st:E.1.1\endcsname{\hyperlink{tag-01R1}{I.\ref*{fam:st:E.1.1}}}
\expandafter\def\csname SRef@fam:src:E.1.2\endcsname{\hyperlink{tag-01R2}{I.\ref*{fam:src:E.1.2}}}
\expandafter\def\csname SRef@fam:st:E.1.2\endcsname{\hyperlink{tag-01R2}{I.\ref*{fam:st:E.1.2}}}
\expandafter\def\csname SRef@fam:src:E.1.3\endcsname{\hyperlink{tag-01R3}{I.\ref*{fam:src:E.1.3}}}
\expandafter\def\csname SRef@fam:st:E.1.3\endcsname{\hyperlink{tag-01R3}{I.\ref*{fam:st:E.1.3}}}
\expandafter\def\csname SRef@fam:src:E.1.4\endcsname{\hyperlink{tag-01R4}{I.\ref*{fam:src:E.1.4}}}
\expandafter\def\csname SRef@fam:st:E.1.4\endcsname{\hyperlink{tag-01R4}{I.\ref*{fam:st:E.1.4}}}
\expandafter\def\csname SRef@fam:src:E.2.1\endcsname{\hyperlink{tag-01R6}{I.\ref*{fam:src:E.2.1}}}
\expandafter\def\csname SRef@fam:st:E.2.1\endcsname{\hyperlink{tag-01R6}{I.\ref*{fam:st:E.2.1}}}
\expandafter\def\csname SRef@push:src:A.1.1\endcsname{\hyperlink{tag-02ZH}{I.\ref*{push:src:A.1.1}}}
\expandafter\def\csname SRef@push:A11\endcsname{\hyperlink{tag-02ZH}{I.\ref*{push:A11}}}
\expandafter\def\csname SRef@push:src:A.4.1\endcsname{\hyperlink{tag-02ZR}{I.\ref*{push:src:A.4.1}}}
\expandafter\def\csname SRef@push:A41\endcsname{\hyperlink{tag-02ZR}{I.\ref*{push:A41}}}
\expandafter\def\csname SRef@push:src:A.4.2\endcsname{\hyperlink{tag-02ZS}{I.\ref*{push:src:A.4.2}}}
\expandafter\def\csname SRef@push:A42\endcsname{\hyperlink{tag-02ZS}{I.\ref*{push:A42}}}
\expandafter\def\csname SRef@push:src:A.4.3\endcsname{\hyperlink{tag-02ZT}{I.\ref*{push:src:A.4.3}}}
\expandafter\def\csname SRef@push:A43\endcsname{\hyperlink{tag-02ZT}{I.\ref*{push:A43}}}
\expandafter\def\csname SRef@push:src:A.4.4\endcsname{\hyperlink{tag-02ZU}{I.\ref*{push:src:A.4.4}}}
\expandafter\def\csname SRef@push:A44\endcsname{\hyperlink{tag-02ZU}{I.\ref*{push:A44}}}
\expandafter\def\csname SRef@push:src:A.4.6\endcsname{\hyperlink{tag-02ZW}{I.\ref*{push:src:A.4.6}}}
\expandafter\def\csname SRef@push:A46\endcsname{\hyperlink{tag-02ZW}{I.\ref*{push:A46}}}
\expandafter\def\csname SRef@push:src:A.4.9\endcsname{\hyperlink{tag-02ZZ}{I.\ref*{push:src:A.4.9}}}
\expandafter\def\csname SRef@push:A49\endcsname{\hyperlink{tag-02ZZ}{I.\ref*{push:A49}}}
\expandafter\def\csname SRef@push:src:A.2.1\endcsname{\hyperlink{tag-02ZJ}{I.\ref*{push:src:A.2.1}}}
\expandafter\def\csname SRef@push:A21\endcsname{\hyperlink{tag-02ZJ}{I.\ref*{push:A21}}}
\expandafter\def\csname SRef@push:src:A.2.5\endcsname{\hyperlink{tag-02ZN}{I.\ref*{push:src:A.2.5}}}
\expandafter\def\csname SRef@push:A25\endcsname{\hyperlink{tag-02ZN}{I.\ref*{push:A25}}}
\expandafter\def\csname SRef@push:src:A.3.3\endcsname{\hyperlink{tag-02ZQ}{I.\ref*{push:src:A.3.3}}}
\expandafter\def\csname SRef@push:A33\endcsname{\hyperlink{tag-02ZQ}{I.\ref*{push:A33}}}
\expandafter\def\csname SRef@push:src:B.1.1\endcsname{\hyperlink{tag-0300}{I.\ref*{push:src:B.1.1}}}
\expandafter\def\csname SRef@push:B11\endcsname{\hyperlink{tag-0300}{I.\ref*{push:B11}}}
\expandafter\def\csname SRef@push:src:B.3.6\endcsname{\hyperlink{tag-030A}{I.\ref*{push:src:B.3.6}}}
\expandafter\def\csname SRef@push:B36\endcsname{\hyperlink{tag-030A}{I.\ref*{push:B36}}}
\expandafter\def\csname SRef@push:src:B.3.7\endcsname{\hyperlink{tag-030B}{I.\ref*{push:src:B.3.7}}}
\expandafter\def\csname SRef@push:B37\endcsname{\hyperlink{tag-030B}{I.\ref*{push:B37}}}
\expandafter\def\csname SRef@push:src:B.3.8\endcsname{\hyperlink{tag-030C}{I.\ref*{push:src:B.3.8}}}
\expandafter\def\csname SRef@push:B38\endcsname{\hyperlink{tag-030C}{I.\ref*{push:B38}}}
\expandafter\def\csname SRef@push:src:B.4.4\endcsname{\hyperlink{tag-0314}{I.\ref*{push:src:B.4.4}}}
\expandafter\def\csname SRef@push:src:B.4.6\endcsname{\hyperlink{tag-0316}{I.\ref*{push:src:B.4.6}}}
\expandafter\def\csname SRef@push:B46\endcsname{\hyperlink{tag-0316}{I.\ref*{push:B46}}}
\expandafter\def\csname SRef@push:src:B.4.7\endcsname{\hyperlink{tag-0317}{I.\ref*{push:src:B.4.7}}}
\expandafter\def\csname SRef@push:src:B.4.8\endcsname{\hyperlink{tag-0318}{I.\ref*{push:src:B.4.8}}}
\expandafter\def\csname SRef@push:src:B.4.9\endcsname{\hyperlink{tag-0319}{I.\ref*{push:src:B.4.9}}}
\expandafter\def\csname SRef@push:src:B.4.10\endcsname{\hyperlink{tag-030E}{I.\ref*{push:src:B.4.10}}}
\expandafter\def\csname SRef@push:src:B.4.11\endcsname{\hyperlink{tag-030F}{I.\ref*{push:src:B.4.11}}}
\expandafter\def\csname SRef@push:src:B.4.12\endcsname{\hyperlink{tag-030G}{I.\ref*{push:src:B.4.12}}}
\expandafter\def\csname SRef@push:src:B.4.13\endcsname{\hyperlink{tag-030H}{I.\ref*{push:src:B.4.13}}}
\expandafter\def\csname SRef@push:src:B.4.14\endcsname{\hyperlink{tag-030I}{I.\ref*{push:src:B.4.14}}}
\expandafter\def\csname SRef@push:src:B.4.15\endcsname{\hyperlink{tag-030J}{I.\ref*{push:src:B.4.15}}}
\expandafter\def\csname SRef@push:B415\endcsname{\hyperlink{tag-030J}{I.\ref*{push:B415}}}
\expandafter\def\csname SRef@push:src:B.4.16\endcsname{\hyperlink{tag-030K}{I.\ref*{push:src:B.4.16}}}
\expandafter\def\csname SRef@push:B416\endcsname{\hyperlink{tag-030K}{I.\ref*{push:B416}}}
\expandafter\def\csname SRef@push:src:B.2.1\endcsname{\hyperlink{tag-0302}{I.\ref*{push:src:B.2.1}}}
\expandafter\def\csname SRef@push:B21\endcsname{\hyperlink{tag-0302}{I.\ref*{push:B21}}}
\expandafter\def\csname SRef@push:src:B.4.17\endcsname{\hyperlink{tag-030L}{I.\ref*{push:src:B.4.17}}}
\expandafter\def\csname SRef@push:src:B.4.18\endcsname{\hyperlink{tag-030M}{I.\ref*{push:src:B.4.18}}}
\expandafter\def\csname SRef@push:B418\endcsname{\hyperlink{tag-030M}{I.\ref*{push:B418}}}
\expandafter\def\csname SRef@push:src:B.4.19\endcsname{\hyperlink{tag-030N}{I.\ref*{push:src:B.4.19}}}
\expandafter\def\csname SRef@push:B419\endcsname{\hyperlink{tag-030N}{I.\ref*{push:B419}}}
\expandafter\def\csname SRef@push:src:B.4.20\endcsname{\hyperlink{tag-030P}{I.\ref*{push:src:B.4.20}}}
\expandafter\def\csname SRef@push:B420\endcsname{\hyperlink{tag-030P}{I.\ref*{push:B420}}}
\expandafter\def\csname SRef@push:src:B.4.21\endcsname{\hyperlink{tag-030Q}{I.\ref*{push:src:B.4.21}}}
\expandafter\def\csname SRef@push:src:B.4.22\endcsname{\hyperlink{tag-030R}{I.\ref*{push:src:B.4.22}}}
\expandafter\def\csname SRef@push:B422\endcsname{\hyperlink{tag-030R}{I.\ref*{push:B422}}}
\expandafter\def\csname SRef@push:src:B.4.23\endcsname{\hyperlink{tag-030S}{I.\ref*{push:src:B.4.23}}}
\expandafter\def\csname SRef@push:src:B.4.24\endcsname{\hyperlink{tag-030T}{I.\ref*{push:src:B.4.24}}}
\expandafter\def\csname SRef@push:B424\endcsname{\hyperlink{tag-030T}{I.\ref*{push:B424}}}
\expandafter\def\csname SRef@push:src:B.4.25\endcsname{\hyperlink{tag-030U}{I.\ref*{push:src:B.4.25}}}
\expandafter\def\csname SRef@push:B425\endcsname{\hyperlink{tag-030U}{I.\ref*{push:B425}}}
\expandafter\def\csname SRef@push:src:B.4.26\endcsname{\hyperlink{tag-030V}{I.\ref*{push:src:B.4.26}}}
\expandafter\def\csname SRef@push:B426\endcsname{\hyperlink{tag-030V}{I.\ref*{push:B426}}}
\expandafter\def\csname SRef@push:src:B.4.27\endcsname{\hyperlink{tag-030W}{I.\ref*{push:src:B.4.27}}}
\expandafter\def\csname SRef@push:B427\endcsname{\hyperlink{tag-030W}{I.\ref*{push:B427}}}
\expandafter\def\csname SRef@push:src:B.4.28\endcsname{\hyperlink{tag-030X}{I.\ref*{push:src:B.4.28}}}
\expandafter\def\csname SRef@push:src:B.4.29\endcsname{\hyperlink{tag-030Y}{I.\ref*{push:src:B.4.29}}}
\expandafter\def\csname SRef@push:B429\endcsname{\hyperlink{tag-030Y}{I.\ref*{push:B429}}}
\expandafter\def\csname SRef@push:src:B.4.30\endcsname{\hyperlink{tag-0310}{I.\ref*{push:src:B.4.30}}}
\expandafter\def\csname SRef@push:B430\endcsname{\hyperlink{tag-0310}{I.\ref*{push:B430}}}
\expandafter\def\csname SRef@push:src:B.4.31\endcsname{\hyperlink{tag-0311}{I.\ref*{push:src:B.4.31}}}
\expandafter\def\csname SRef@push:src:B.4.32\endcsname{\hyperlink{tag-0312}{I.\ref*{push:src:B.4.32}}}
\expandafter\def\csname SRef@push:B432\endcsname{\hyperlink{tag-0312}{I.\ref*{push:B432}}}
\expandafter\def\csname SRef@push:src:B.3.2\endcsname{\hyperlink{tag-0306}{I.\ref*{push:src:B.3.2}}}
\expandafter\def\csname SRef@push:B32\endcsname{\hyperlink{tag-0306}{I.\ref*{push:B32}}}
\expandafter\def\csname SRef@push:src:B.3.3\endcsname{\hyperlink{tag-0307}{I.\ref*{push:src:B.3.3}}}
\expandafter\def\csname SRef@push:B33\endcsname{\hyperlink{tag-0307}{I.\ref*{push:B33}}}
\expandafter\def\csname SRef@push:src:D.1.1\endcsname{\hyperlink{tag-032A}{I.\ref*{push:src:D.1.1}}}
\expandafter\def\csname SRef@push:src:D.4.1\endcsname{\hyperlink{tag-032J}{I.\ref*{push:src:D.4.1}}}
\expandafter\def\csname SRef@push:src:D.4.2\endcsname{\hyperlink{tag-032K}{I.\ref*{push:src:D.4.2}}}
\expandafter\def\csname SRef@push:src:D.4.3\endcsname{\hyperlink{tag-032L}{I.\ref*{push:src:D.4.3}}}
\expandafter\def\csname SRef@push:src:D.4.4\endcsname{\hyperlink{tag-032M}{I.\ref*{push:src:D.4.4}}}
\expandafter\def\csname SRef@push:src:D.4.5\endcsname{\hyperlink{tag-032N}{I.\ref*{push:src:D.4.5}}}
\expandafter\def\csname SRef@push:src:D.4.7\endcsname{\hyperlink{tag-032P}{I.\ref*{push:src:D.4.7}}}
\expandafter\def\csname SRef@push:src:D.4.8\endcsname{\hyperlink{tag-032Q}{I.\ref*{push:src:D.4.8}}}
\expandafter\def\csname SRef@push:src:D.2.2\endcsname{\hyperlink{tag-032E}{I.\ref*{push:src:D.2.2}}}
\expandafter\def\csname SRef@push:src:D.2.4\endcsname{\hyperlink{tag-032G}{I.\ref*{push:src:D.2.4}}}
\expandafter\def\csname SRef@inv:A3.2\endcsname{\hyperlink{tag-01RZ}{I.\ref*{inv:A3.2}}}
\expandafter\def\csname SRef@wb:lem:bcpr-eta-extension\endcsname{\hyperlink{tag-01RZ}{I.\ref*{wb:lem:bcpr-eta-extension}}}
\expandafter\def\csname SRef@inv:A4.1\endcsname{\hyperlink{tag-01S0}{I.\ref*{inv:A4.1}}}
\expandafter\def\csname SRef@inv:A4.2\endcsname{\hyperlink{tag-01S1}{I.\ref*{inv:A4.2}}}
\expandafter\def\csname SRef@inv:A1.1\endcsname{\hyperlink{tag-01RR}{I.\ref*{inv:A1.1}}}
\expandafter\def\csname SRef@inv:A1.2\endcsname{\hyperlink{tag-01RS}{I.\ref*{inv:A1.2}}}
\expandafter\def\csname SRef@wb:prop:log-relative-direct-images\endcsname{\hyperlink{tag-01RS}{I.\ref*{wb:prop:log-relative-direct-images}}}
\expandafter\def\csname SRef@inv:A1.3\endcsname{\hyperlink{tag-01RT}{I.\ref*{inv:A1.3}}}
\expandafter\def\csname SRef@wb:prop:absolute-cotangent-direct-images\endcsname{\hyperlink{tag-01RT}{I.\ref*{wb:prop:absolute-cotangent-direct-images}}}
\expandafter\def\csname SRef@inv:A2.1\endcsname{\hyperlink{tag-01RU}{I.\ref*{inv:A2.1}}}
\expandafter\def\csname SRef@wb:prop:log-ks-divisor\endcsname{\hyperlink{tag-01RU}{I.\ref*{wb:prop:log-ks-divisor}}}
\expandafter\def\csname SRef@wb:prop:period-derivative-determinant\endcsname{\hyperlink{tag-01RU}{I.\ref*{wb:prop:period-derivative-determinant}}}
\expandafter\def\csname SRef@inv:A2.2\endcsname{\hyperlink{tag-01RV}{I.\ref*{inv:A2.2}}}
\expandafter\def\csname SRef@inv:A2.3\endcsname{\hyperlink{tag-01RW}{I.\ref*{inv:A2.3}}}
\expandafter\def\csname SRef@inv:A2.4\endcsname{\hyperlink{tag-01RX}{I.\ref*{inv:A2.4}}}
\expandafter\def\csname SRef@wb:thm:h12-scalar\endcsname{\hyperlink{tag-01RX}{I.\ref*{wb:thm:h12-scalar}}}
\expandafter\def\csname SRef@inv:A3.1\endcsname{\hyperlink{tag-01RY}{I.\ref*{inv:A3.1}}}
\expandafter\def\csname SRef@wb:cor:exact-hodge-diamond\endcsname{\hyperlink{tag-01RY}{I.\ref*{wb:cor:exact-hodge-diamond}}}
\expandafter\def\csname SRef@wb:thm:parametric-hodge-diamond\endcsname{\hyperlink{tag-01RY}{I.\ref*{wb:thm:parametric-hodge-diamond}}}
\expandafter\def\csname SRef@wb:prop:workbench-holomorphic-form-vanishing\endcsname{\hyperlink{tag-01RY}{I.\ref*{wb:prop:workbench-holomorphic-form-vanishing}}}
\expandafter\def\csname SRef@inv:C1.1\endcsname{\hyperlink{tag-01SA}{I.\ref*{inv:C1.1}}}
\expandafter\def\csname SRef@wb:thm:frolicher-e2-candidate\endcsname{\hyperlink{tag-01SA}{I.\ref*{wb:thm:frolicher-e2-candidate}}}
\expandafter\def\csname SRef@inv:C3.7\endcsname{\hyperlink{tag-01SK}{I.\ref*{inv:C3.7}}}
\expandafter\def\csname SRef@inv:C1.2\endcsname{\hyperlink{tag-01SB}{I.\ref*{inv:C1.2}}}
\expandafter\def\csname SRef@inv:C2.1\endcsname{\hyperlink{tag-01SC}{I.\ref*{inv:C2.1}}}
\expandafter\def\csname SRef@wb:thm:exact-bott-chern-aeppli\endcsname{\hyperlink{tag-01SC}{I.\ref*{wb:thm:exact-bott-chern-aeppli}}}
\expandafter\def\csname SRef@inv:C2.2\endcsname{\hyperlink{tag-01SD}{I.\ref*{inv:C2.2}}}
\expandafter\def\csname SRef@inv:C3.1\endcsname{\hyperlink{tag-01SE}{I.\ref*{inv:C3.1}}}
\expandafter\def\csname SRef@wb:cor:picard-line-audit\endcsname{\hyperlink{tag-01SE}{I.\ref*{wb:cor:picard-line-audit}}}
\expandafter\def\csname SRef@inv:C3.2\endcsname{\hyperlink{tag-01SF}{I.\ref*{inv:C3.2}}}
\expandafter\def\csname SRef@wb:lem:bcpr-degree-one-input\endcsname{\hyperlink{tag-01SF}{I.\ref*{wb:lem:bcpr-degree-one-input}}}
\expandafter\def\csname SRef@inv:C3.3\endcsname{\hyperlink{tag-01SG}{I.\ref*{inv:C3.3}}}
\expandafter\def\csname SRef@wb:thm:bcpr-square-zero-algebra\endcsname{\hyperlink{tag-01SG}{I.\ref*{wb:thm:bcpr-square-zero-algebra}}}
\expandafter\def\csname SRef@wb:prop:bcpr-generators-are-bases\endcsname{\hyperlink{tag-01SG}{I.\ref*{wb:prop:bcpr-generators-are-bases}}}
\expandafter\def\csname SRef@inv:C3.4\endcsname{\hyperlink{tag-01SH}{I.\ref*{inv:C3.4}}}
\expandafter\def\csname SRef@inv:C3.5\endcsname{\hyperlink{tag-01SI}{I.\ref*{inv:C3.5}}}
\expandafter\def\csname SRef@inv:D4.1\endcsname{\hyperlink{tag-01SO}{I.\ref*{inv:D4.1}}}
\expandafter\def\csname SRef@inv:D4.2\endcsname{\hyperlink{tag-01SP}{I.\ref*{inv:D4.2}}}
\expandafter\def\csname SRef@inv:D4.3\endcsname{\hyperlink{tag-01SQ}{I.\ref*{inv:D4.3}}}
\expandafter\def\csname SRef@inv:D4.4\endcsname{\hyperlink{tag-01SR}{I.\ref*{inv:D4.4}}}
\expandafter\def\csname SRef@inv:D4.5\endcsname{\hyperlink{tag-01SS}{I.\ref*{inv:D4.5}}}
\expandafter\def\csname SRef@push:src:E.2.2\endcsname{\hyperlink{tag-0335}{I.\ref*{push:src:E.2.2}}}
\expandafter\def\csname SRef@push:src:E.2.3\endcsname{\hyperlink{tag-033B}{I.\ref*{push:src:E.2.3}}}
\expandafter\def\csname SRef@inv:B1.0\endcsname{\hyperlink{tag-01S2}{I.\ref*{inv:B1.0}}}
\expandafter\def\csname SRef@inv:B1.1\endcsname{\hyperlink{tag-01S3}{I.\ref*{inv:B1.1}}}
\expandafter\def\csname SRef@inv:B1.2\endcsname{\hyperlink{tag-01S4}{I.\ref*{inv:B1.2}}}
\expandafter\def\csname SRef@inv:B2.1\endcsname{\hyperlink{tag-01S5}{I.\ref*{inv:B2.1}}}
\expandafter\def\csname SRef@inv:B2.2\endcsname{\hyperlink{tag-01S6}{I.\ref*{inv:B2.2}}}
\expandafter\def\csname SRef@inv:B3.1\endcsname{\hyperlink{tag-01S7}{I.\ref*{inv:B3.1}}}
\expandafter\def\csname SRef@inv:B3.2\endcsname{\hyperlink{tag-01S8}{I.\ref*{inv:B3.2}}}
\expandafter\def\csname SRef@inv:B3.3\endcsname{\hyperlink{tag-01S9}{I.\ref*{inv:B3.3}}}
\expandafter\def\csname SRef@mapsI:src:A.1.1\endcsname{\hyperlink{tag-02PA}{I.\ref*{mapsI:src:A.1.1}}}
\expandafter\def\csname SRef@mapsI:quotient-eq-1\endcsname{\hyperlink{tag-02PA}{I.\ref*{mapsI:quotient-eq-1}}}
\expandafter\def\csname SRef@mapsI:src:A.1.2\endcsname{\hyperlink{tag-02PB}{I.\ref*{mapsI:src:A.1.2}}}
\expandafter\def\csname SRef@mapsI:src:A.1.3\endcsname{\hyperlink{tag-02PC}{I.\ref*{mapsI:src:A.1.3}}}
\expandafter\def\csname SRef@mapsI:quotient-eq-6\endcsname{\hyperlink{tag-02PC}{I.\ref*{mapsI:quotient-eq-6}}}
\expandafter\def\csname SRef@ari:C.2.1\endcsname{\hyperlink{tag-000L}{I.\ref*{ari:C.2.1}}}
\expandafter\def\csname SRef@ari:C.2.2\endcsname{\hyperlink{tag-000M}{I.\ref*{ari:C.2.2}}}
\expandafter\def\csname SRef@ari:C.3.1\endcsname{\hyperlink{tag-000N}{I.\ref*{ari:C.3.1}}}
\expandafter\def\csname SRef@ari:C.3.2\endcsname{\hyperlink{tag-000O}{I.\ref*{ari:C.3.2}}}
\expandafter\def\csname SRef@ari:C.6.1\endcsname{\hyperlink{tag-000V}{I.\ref*{ari:C.6.1}}}
\expandafter\def\csname SRef@ari:C.6.R5\endcsname{\hyperlink{tag-0014}{I.\ref*{ari:C.6.R5}}}
\expandafter\def\csname SRef@ari:C.6.T1\endcsname{\hyperlink{tag-0015}{I.\ref*{ari:C.6.T1}}}
\expandafter\def\csname SRef@ari:C.6.T2\endcsname{\hyperlink{tag-0016}{I.\ref*{ari:C.6.T2}}}
\expandafter\def\csname SRef@ari:C.8.6\endcsname{\hyperlink{tag-001T}{I.\ref*{ari:C.8.6}}}
\expandafter\def\csname SRef@ari:C.8.7\endcsname{\hyperlink{tag-001U}{I.\ref*{ari:C.8.7}}}
\expandafter\def\csname SRef@ari:C.8.8\endcsname{\hyperlink{tag-001V}{I.\ref*{ari:C.8.8}}}
\expandafter\def\csname SRef@ari:C.8.9\endcsname{\hyperlink{tag-001W}{I.\ref*{ari:C.8.9}}}
\expandafter\def\csname SRef@ari:C.8.10\endcsname{\hyperlink{tag-001K}{I.\ref*{ari:C.8.10}}}
\expandafter\def\csname SRef@ari:C.8.11\endcsname{\hyperlink{tag-001L}{I.\ref*{ari:C.8.11}}}
\expandafter\def\csname SRef@ari:C.6.2\endcsname{\hyperlink{tag-000W}{I.\ref*{ari:C.6.2}}}
\expandafter\def\csname SRef@ari:C.6.3\endcsname{\hyperlink{tag-000X}{I.\ref*{ari:C.6.3}}}
\expandafter\def\csname SRef@ari:C.6.4\endcsname{\hyperlink{tag-000Y}{I.\ref*{ari:C.6.4}}}
\expandafter\def\csname SRef@ari:C.6.5\endcsname{\hyperlink{tag-000Z}{I.\ref*{ari:C.6.5}}}
\expandafter\def\csname SRef@ari:C.6.R1\endcsname{\hyperlink{tag-0010}{I.\ref*{ari:C.6.R1}}}
\expandafter\def\csname SRef@ari:C.6.R2\endcsname{\hyperlink{tag-0011}{I.\ref*{ari:C.6.R2}}}
\expandafter\def\csname SRef@ari:C.6.R3\endcsname{\hyperlink{tag-0012}{I.\ref*{ari:C.6.R3}}}
\expandafter\def\csname SRef@ari:C.8.12\endcsname{\hyperlink{tag-001M}{I.\ref*{ari:C.8.12}}}
\expandafter\def\csname SRef@ari:C.8.13\endcsname{\hyperlink{tag-001N}{I.\ref*{ari:C.8.13}}}
\expandafter\def\csname SRef@ari:C.8.14\endcsname{\hyperlink{tag-001O}{I.\ref*{ari:C.8.14}}}
\expandafter\def\csname SRef@ari:C.8.15\endcsname{\hyperlink{tag-001P}{I.\ref*{ari:C.8.15}}}
\expandafter\def\csname SRef@fam:src:B.1.1\endcsname{\hyperlink{tag-01Q7}{I.\ref*{fam:src:B.1.1}}}
\expandafter\def\csname SRef@fam:st:B.1.1\endcsname{\hyperlink{tag-01Q7}{I.\ref*{fam:st:B.1.1}}}
\expandafter\def\csname SRef@fam:src:B.1.2\endcsname{\hyperlink{tag-01Q8}{I.\ref*{fam:src:B.1.2}}}
\expandafter\def\csname SRef@fam:st:B.1.2\endcsname{\hyperlink{tag-01Q8}{I.\ref*{fam:st:B.1.2}}}
\expandafter\def\csname SRef@fam:src:B.1.3\endcsname{\hyperlink{tag-01Q9}{I.\ref*{fam:src:B.1.3}}}
\expandafter\def\csname SRef@fam:st:B.1.3\endcsname{\hyperlink{tag-01Q9}{I.\ref*{fam:st:B.1.3}}}
\expandafter\def\csname SRef@fam:src:B.1.4\endcsname{\hyperlink{tag-01QA}{I.\ref*{fam:src:B.1.4}}}
\expandafter\def\csname SRef@fam:st:B.1.4\endcsname{\hyperlink{tag-01QA}{I.\ref*{fam:st:B.1.4}}}
\expandafter\def\csname SRef@fam:src:C.1.1\endcsname{\hyperlink{tag-01QH}{I.\ref*{fam:src:C.1.1}}}
\expandafter\def\csname SRef@fam:st:C.1.1\endcsname{\hyperlink{tag-01QH}{I.\ref*{fam:st:C.1.1}}}
\expandafter\def\csname SRef@fam:src:C.1.2\endcsname{\hyperlink{tag-01QI}{I.\ref*{fam:src:C.1.2}}}
\expandafter\def\csname SRef@fam:st:C.1.2\endcsname{\hyperlink{tag-01QI}{I.\ref*{fam:st:C.1.2}}}
\expandafter\def\csname SRef@fam:src:C.2.1\endcsname{\hyperlink{tag-01QJ}{I.\ref*{fam:src:C.2.1}}}
\expandafter\def\csname SRef@fam:st:C.2.1\endcsname{\hyperlink{tag-01QJ}{I.\ref*{fam:st:C.2.1}}}
\expandafter\def\csname SRef@fam:src:C.2.2\endcsname{\hyperlink{tag-01QK}{I.\ref*{fam:src:C.2.2}}}
\expandafter\def\csname SRef@fam:st:C.2.2\endcsname{\hyperlink{tag-01QK}{I.\ref*{fam:st:C.2.2}}}
\expandafter\def\csname SRef@ari:C.1.1\endcsname{\hyperlink{tag-000I}{I.\ref*{ari:C.1.1}}}
\expandafter\def\csname SRef@ari:C.1.2\endcsname{\hyperlink{tag-000J}{I.\ref*{ari:C.1.2}}}
\expandafter\def\csname SRef@ari:C.1.3\endcsname{\hyperlink{tag-000K}{I.\ref*{ari:C.1.3}}}
\expandafter\def\csname SRef@ari:C.4.2\endcsname{\hyperlink{tag-000Q}{I.\ref*{ari:C.4.2}}}
\expandafter\def\csname SRef@ari:C.4.3\endcsname{\hyperlink{tag-000R}{I.\ref*{ari:C.4.3}}}
\expandafter\def\csname SRef@ari:C.4.4\endcsname{\hyperlink{tag-000S}{I.\ref*{ari:C.4.4}}}
\expandafter\def\csname SRef@ari:C.4.5\endcsname{\hyperlink{tag-000T}{I.\ref*{ari:C.4.5}}}
\expandafter\def\csname SRef@ari:C.8.1\endcsname{\hyperlink{tag-001J}{I.\ref*{ari:C.8.1}}}
\expandafter\def\csname SRef@ari:C.8.2\endcsname{\hyperlink{tag-001Q}{I.\ref*{ari:C.8.2}}}
\expandafter\def\csname SRef@ari:C.8.3\endcsname{\hyperlink{tag-001R}{I.\ref*{ari:C.8.3}}}
\expandafter\def\csname SRef@ari:C.8.4\endcsname{\hyperlink{tag-001S}{I.\ref*{ari:C.8.4}}}
\expandafter\def\csname SRef@fam:src:A.1.1\endcsname{\hyperlink{tag-01PZ}{I.\ref*{fam:src:A.1.1}}}
\expandafter\def\csname SRef@fam:st:A.1.1\endcsname{\hyperlink{tag-01PZ}{I.\ref*{fam:st:A.1.1}}}
\expandafter\def\csname SRef@fam:src:A.1.2\endcsname{\hyperlink{tag-01Q0}{I.\ref*{fam:src:A.1.2}}}
\expandafter\def\csname SRef@fam:st:A.1.2\endcsname{\hyperlink{tag-01Q0}{I.\ref*{fam:st:A.1.2}}}
\expandafter\def\csname SRef@fam:src:A.2.1\endcsname{\hyperlink{tag-01Q1}{I.\ref*{fam:src:A.2.1}}}
\expandafter\def\csname SRef@fam:st:A.2.1\endcsname{\hyperlink{tag-01Q1}{I.\ref*{fam:st:A.2.1}}}
\expandafter\def\csname SRef@fam:src:A.3.1\endcsname{\hyperlink{tag-01Q2}{I.\ref*{fam:src:A.3.1}}}
\expandafter\def\csname SRef@fam:st:A.3.1\endcsname{\hyperlink{tag-01Q2}{I.\ref*{fam:st:A.3.1}}}
\expandafter\def\csname SRef@fam:src:A.4.1\endcsname{\hyperlink{tag-01Q3}{I.\ref*{fam:src:A.4.1}}}
\expandafter\def\csname SRef@fam:st:A.4.1\endcsname{\hyperlink{tag-01Q3}{I.\ref*{fam:st:A.4.1}}}
\expandafter\def\csname SRef@fam:src:A.4.2\endcsname{\hyperlink{tag-01Q4}{I.\ref*{fam:src:A.4.2}}}
\expandafter\def\csname SRef@fam:st:A.4.2\endcsname{\hyperlink{tag-01Q4}{I.\ref*{fam:st:A.4.2}}}
\expandafter\def\csname SRef@fam:src:A.4.3\endcsname{\hyperlink{tag-01Q5}{I.\ref*{fam:src:A.4.3}}}
\expandafter\def\csname SRef@fam:st:A.4.3\endcsname{\hyperlink{tag-01Q5}{I.\ref*{fam:st:A.4.3}}}
\expandafter\def\csname SRef@fam:src:B.2.1\endcsname{\hyperlink{tag-01QB}{I.\ref*{fam:src:B.2.1}}}
\expandafter\def\csname SRef@fam:st:B.2.1\endcsname{\hyperlink{tag-01QB}{I.\ref*{fam:st:B.2.1}}}
\expandafter\def\csname SRef@fam:src:B.2.2\endcsname{\hyperlink{tag-01QC}{I.\ref*{fam:src:B.2.2}}}
\expandafter\def\csname SRef@fam:st:B.2.2\endcsname{\hyperlink{tag-01QC}{I.\ref*{fam:st:B.2.2}}}
\expandafter\def\csname SRef@fam:src:B.2.3\endcsname{\hyperlink{tag-01QD}{I.\ref*{fam:src:B.2.3}}}
\expandafter\def\csname SRef@fam:st:B.2.3\endcsname{\hyperlink{tag-01QD}{I.\ref*{fam:st:B.2.3}}}
\expandafter\def\csname SRef@fam:src:B.3.1\endcsname{\hyperlink{tag-01QE}{I.\ref*{fam:src:B.3.1}}}
\expandafter\def\csname SRef@fam:st:B.3.1\endcsname{\hyperlink{tag-01QE}{I.\ref*{fam:st:B.3.1}}}
\expandafter\def\csname SRef@fam:src:B.4.1\endcsname{\hyperlink{tag-01QF}{I.\ref*{fam:src:B.4.1}}}
\expandafter\def\csname SRef@fam:st:B.4.1\endcsname{\hyperlink{tag-01QF}{I.\ref*{fam:st:B.4.1}}}
\expandafter\def\csname SRef@fam:src:B.4.2\endcsname{\hyperlink{tag-01QG}{I.\ref*{fam:src:B.4.2}}}
\expandafter\def\csname SRef@fam:st:B.4.2\endcsname{\hyperlink{tag-01QG}{I.\ref*{fam:st:B.4.2}}}
\expandafter\def\csname SRef@push:src:E.1.1\endcsname{\hyperlink{tag-032R}{I.\ref*{push:src:E.1.1}}}
\expandafter\def\csname SRef@inv:D1.1\endcsname{\hyperlink{tag-01SM}{I.\ref*{inv:D1.1}}}
\expandafter\def\csname SRef@inv:D1.2\endcsname{\hyperlink{tag-01SN}{I.\ref*{inv:D1.2}}}
\expandafter\def\csname SRef@ari:C.7.1\endcsname{\hyperlink{tag-0017}{I.\ref*{ari:C.7.1}}}
\expandafter\def\csname SRef@ari:C.7.11\endcsname{\hyperlink{tag-0019}{I.\ref*{ari:C.7.11}}}
\expandafter\def\csname SRef@ari:C.7.12\endcsname{\hyperlink{tag-001A}{I.\ref*{ari:C.7.12}}}
\expandafter\def\csname SRef@ari:C.7.2\endcsname{\hyperlink{tag-001B}{I.\ref*{ari:C.7.2}}}
\expandafter\def\csname SRef@ari:C.7.3\endcsname{\hyperlink{tag-001C}{I.\ref*{ari:C.7.3}}}
\expandafter\def\csname SRef@ari:C.7.4\endcsname{\hyperlink{tag-001D}{I.\ref*{ari:C.7.4}}}
\expandafter\def\csname SRef@ari:C.7.6\endcsname{\hyperlink{tag-001F}{I.\ref*{ari:C.7.6}}}
\expandafter\def\csname SRef@ari:C.7.9\endcsname{\hyperlink{tag-001I}{I.\ref*{ari:C.7.9}}}
\expandafter\def\csname SRef@ari:C.7.10\endcsname{\hyperlink{tag-0018}{I.\ref*{ari:C.7.10}}}
\expandafter\def\csname SRef@inv:C3.8\endcsname{\hyperlink{tag-01SL}{I.\ref*{inv:C3.8}}}
\expandafter\def\csname SRef@fam:src:C.3.3\endcsname{\hyperlink{tag-01QO}{I.\ref*{fam:src:C.3.3}}}
\expandafter\def\csname SRef@fam:st:C.3.3\endcsname{\hyperlink{tag-01QO}{I.\ref*{fam:st:C.3.3}}}
\expandafter\def\csname SRef@fam:src:E.1.5\endcsname{\hyperlink{tag-01R5}{I.\ref*{fam:src:E.1.5}}}
\expandafter\def\csname SRef@fam:st:E.1.5\endcsname{\hyperlink{tag-01R5}{I.\ref*{fam:st:E.1.5}}}
\expandafter\def\csname SRef@wb:rem:bcpr-literal-gauge\endcsname{\hyperlink{tag-03CT}{I.\ref*{wb:rem:bcpr-literal-gauge}}}
\expandafter\def\csname SRef@wb:rem:bcpr-global-gauge-boundary\endcsname{\hyperlink{tag-03CS}{I.\ref*{wb:rem:bcpr-global-gauge-boundary}}}
\expandafter\def\csname SRef@push:src:A.2.3\endcsname{\hyperlink{tag-02ZL}{I.\ref*{push:src:A.2.3}}}
\expandafter\def\csname SRef@push:A23\endcsname{\hyperlink{tag-02ZL}{I.\ref*{push:A23}}}
\expandafter\def\csname SRef@fam:section:F\endcsname{\hyperlink{tag-01PV}{I.\ref*{fam:section:F}}}
\expandafter\def\csname SRef@fam:section:D\endcsname{\hyperlink{tag-01PP}{I.\ref*{fam:section:D}}}
\expandafter\def\csname SRef@mapsI:real-forms\endcsname{\hyperlink{tag-02P9}{I.\ref*{mapsI:real-forms}}}
\expandafter\def\csname SRef@fam:section:E\endcsname{\hyperlink{tag-01PS}{I.\ref*{fam:section:E}}}
\expandafter\def\csname SRef@push:section:A\endcsname{\hyperlink{tag-02YZ}{I.\ref*{push:section:A}}}
\expandafter\def\csname SRef@push:section:B\endcsname{\hyperlink{tag-02Z4}{I.\ref*{push:section:B}}}
\expandafter\def\csname SRef@push:section:D\endcsname{\hyperlink{tag-02Z8}{I.\ref*{push:section:D}}}
\expandafter\def\csname SRef@inv:sec:partA\endcsname{\hyperlink{tag-01TH}{I.\ref*{inv:sec:partA}}}
\expandafter\def\csname SRef@inv:sec:partC\endcsname{\hyperlink{tag-01TI}{I.\ref*{inv:sec:partC}}}
\expandafter\def\csname SRef@def:common-conventions\endcsname{\hyperlink{tag-01O1}{I.\ref*{def:common-conventions}}}
\expandafter\def\csname SRef@mapsI:circle-quotient\endcsname{\hyperlink{tag-02OO}{I.\ref*{mapsI:circle-quotient}}}
\expandafter\def\csname SRef@inv:notation-dictionary\endcsname{\hyperlink{tag-01TG}{I.\ref*{inv:notation-dictionary}}}
\expandafter\def\csname SRef@fam:section:C\endcsname{\hyperlink{tag-01PK}{I.\ref*{fam:section:C}}}
\expandafter\def\csname SRef@fam:section:A\endcsname{\hyperlink{tag-01P9}{I.\ref*{fam:section:A}}}
\expandafter\def\csname SRef@fam:section:B\endcsname{\hyperlink{tag-01PF}{I.\ref*{fam:section:B}}}
\expandafter\def\csname SRef@push:section:E\endcsname{\hyperlink{tag-02ZD}{I.\ref*{push:section:E}}}
\expandafter\def\csname SRef@fam:section:F.1\endcsname{\hyperlink{tag-01PW}{I.\ref*{fam:section:F.1}}}
\expandafter\def\csname SRef@fam:section:F.2\endcsname{\hyperlink{tag-01PX}{I.\ref*{fam:section:F.2}}}
\expandafter\def\csname SRef@fam:section:F.3\endcsname{\hyperlink{tag-01PY}{I.\ref*{fam:section:F.3}}}
\expandafter\def\csname SRef@fam:section:D.1\endcsname{\hyperlink{tag-01PQ}{I.\ref*{fam:section:D.1}}}
\expandafter\def\csname SRef@fam:section:D.2\endcsname{\hyperlink{tag-01PR}{I.\ref*{fam:section:D.2}}}
\expandafter\def\csname SRef@fam:section:C.3\endcsname{\hyperlink{tag-01PN}{I.\ref*{fam:section:C.3}}}
\expandafter\def\csname SRef@fam:section:C.4\endcsname{\hyperlink{tag-01PO}{I.\ref*{fam:section:C.4}}}
\expandafter\def\csname SRef@fam:section:E.1\endcsname{\hyperlink{tag-01PT}{I.\ref*{fam:section:E.1}}}
\expandafter\def\csname SRef@fam:section:E.2\endcsname{\hyperlink{tag-01PU}{I.\ref*{fam:section:E.2}}}
\expandafter\def\csname SRef@push:section:A.1\endcsname{\hyperlink{tag-02Z0}{I.\ref*{push:section:A.1}}}
\expandafter\def\csname SRef@push:section:A.2\endcsname{\hyperlink{tag-02Z1}{I.\ref*{push:section:A.2}}}
\expandafter\def\csname SRef@push:section:A.3\endcsname{\hyperlink{tag-02Z2}{I.\ref*{push:section:A.3}}}
\expandafter\def\csname SRef@push:section:A.4\endcsname{\hyperlink{tag-02Z3}{I.\ref*{push:section:A.4}}}
\expandafter\def\csname SRef@push:section:B.1\endcsname{\hyperlink{tag-02Z5}{I.\ref*{push:section:B.1}}}
\expandafter\def\csname SRef@push:section:B.2\endcsname{\hyperlink{tag-02Z6}{I.\ref*{push:section:B.2}}}
\expandafter\def\csname SRef@push:section:B.3\endcsname{\hyperlink{tag-02Z7}{I.\ref*{push:section:B.3}}}
\expandafter\def\csname SRef@push:section:D.1\endcsname{\hyperlink{tag-02Z9}{I.\ref*{push:section:D.1}}}
\expandafter\def\csname SRef@push:section:D.2\endcsname{\hyperlink{tag-02ZA}{I.\ref*{push:section:D.2}}}
\expandafter\def\csname SRef@push:section:D.3\endcsname{\hyperlink{tag-02ZB}{I.\ref*{push:section:D.3}}}
\expandafter\def\csname SRef@push:section:D.4\endcsname{\hyperlink{tag-02ZC}{I.\ref*{push:section:D.4}}}
\expandafter\def\csname SRef@push:section:E.2\endcsname{\hyperlink{tag-02ZF}{I.\ref*{push:section:E.2}}}
\expandafter\def\csname SRef@push:section:E.3\endcsname{\hyperlink{tag-02ZG}{I.\ref*{push:section:E.3}}}
\expandafter\def\csname SRef@sec:e4\endcsname{\hyperlink{tag-034N}{I.\ref*{sec:e4}}}
\expandafter\def\csname SRef@push:section:E.4\endcsname{\hyperlink{tag-034N}{I.\ref*{push:section:E.4}}}
\expandafter\def\csname SRef@fam:section:B.1\endcsname{\hyperlink{tag-01PG}{I.\ref*{fam:section:B.1}}}
\expandafter\def\csname SRef@fam:section:C.1\endcsname{\hyperlink{tag-01PL}{I.\ref*{fam:section:C.1}}}
\expandafter\def\csname SRef@fam:section:C.2\endcsname{\hyperlink{tag-01PM}{I.\ref*{fam:section:C.2}}}
\expandafter\def\csname SRef@fam:section:A.1\endcsname{\hyperlink{tag-01PA}{I.\ref*{fam:section:A.1}}}
\expandafter\def\csname SRef@fam:section:A.2\endcsname{\hyperlink{tag-01PB}{I.\ref*{fam:section:A.2}}}
\expandafter\def\csname SRef@fam:section:A.3\endcsname{\hyperlink{tag-01PC}{I.\ref*{fam:section:A.3}}}
\expandafter\def\csname SRef@fam:section:A.4\endcsname{\hyperlink{tag-01PD}{I.\ref*{fam:section:A.4}}}
\expandafter\def\csname SRef@fam:section:A.5\endcsname{\hyperlink{tag-01PE}{I.\ref*{fam:section:A.5}}}
\expandafter\def\csname SRef@fam:section:B.2\endcsname{\hyperlink{tag-01PH}{I.\ref*{fam:section:B.2}}}
\expandafter\def\csname SRef@fam:section:B.3\endcsname{\hyperlink{tag-01PI}{I.\ref*{fam:section:B.3}}}
\expandafter\def\csname SRef@fam:section:B.4\endcsname{\hyperlink{tag-01PJ}{I.\ref*{fam:section:B.4}}}
\expandafter\def\csname SRef@push:section:E.1\endcsname{\hyperlink{tag-02ZE}{I.\ref*{push:section:E.1}}}
\expandafter\def\csname SRef@wb:sec:canonical-ring-recovery\endcsname{\hyperlink{tag-03CU}{I.\ref*{wb:sec:canonical-ring-recovery}}}
\expandafter\def\csname SRef@wb:sec:retained-period-deformation\endcsname{\hyperlink{tag-03CV}{I.\ref*{wb:sec:retained-period-deformation}}}
\expandafter\def\csname SRef@wb:pshift:sec:classification\endcsname{\hyperlink{tag-03CR}{I.\ref*{wb:pshift:sec:classification}}}
\expandafter\def\csname SRef@app:rederive:fam-src-A.5.1\endcsname{\hyperlink{tag-000C}{I.\ref*{app:rederive:fam-src-A.5.1}}}
\expandafter\def\csname SRef@app:rederive:fam-src-A.5.2\endcsname{\hyperlink{tag-000D}{I.\ref*{app:rederive:fam-src-A.5.2}}}
\expandafter\def\csname SRef@app:rederive:fam-src-B.4.1\endcsname{\hyperlink{tag-000E}{I.\ref*{app:rederive:fam-src-B.4.1}}}
\expandafter\def\csname SRef@fam:src:F.1.1\endcsname{\hyperlink{tag-01R9}{I.\ref*{fam:src:F.1.1}}}
\expandafter\def\csname SRef@fam:st:F.1.1\endcsname{\hyperlink{tag-01R9}{I.\ref*{fam:st:F.1.1}}}
\expandafter\def\csname SRef@cdp:st:D.1.14\endcsname{\hyperlink{tag-01R9}{I.\ref*{cdp:st:D.1.14}}}
\expandafter\def\csname SRef@fam:src:F.2.2\endcsname{\hyperlink{tag-01RC}{I.\ref*{fam:src:F.2.2}}}
\expandafter\def\csname SRef@fam:st:F.2.2\endcsname{\hyperlink{tag-01RC}{I.\ref*{fam:st:F.2.2}}}
\expandafter\def\csname SRef@fam:src:F.3.1\endcsname{\hyperlink{tag-01RE}{I.\ref*{fam:src:F.3.1}}}
\expandafter\def\csname SRef@fam:st:F.3.1\endcsname{\hyperlink{tag-01RE}{I.\ref*{fam:st:F.3.1}}}
\expandafter\def\csname SRef@fam:src:D.2.3\endcsname{\hyperlink{tag-01R0}{I.\ref*{fam:src:D.2.3}}}
\expandafter\def\csname SRef@fam:st:D.2.3\endcsname{\hyperlink{tag-01R0}{I.\ref*{fam:st:D.2.3}}}
\expandafter\def\csname SRef@fam:src:C.4.1\endcsname{\hyperlink{tag-01QQ}{I.\ref*{fam:src:C.4.1}}}
\expandafter\def\csname SRef@fam:st:C.4.1\endcsname{\hyperlink{tag-01QQ}{I.\ref*{fam:st:C.4.1}}}
\expandafter\def\csname SRef@fam:src:D.1.3\endcsname{\hyperlink{tag-01QT}{I.\ref*{fam:src:D.1.3}}}
\expandafter\def\csname SRef@fam:st:D.1.3\endcsname{\hyperlink{tag-01QT}{I.\ref*{fam:st:D.1.3}}}
\expandafter\def\csname SRef@fam:src:D.1.6\endcsname{\hyperlink{tag-01QW}{I.\ref*{fam:src:D.1.6}}}
\expandafter\def\csname SRef@fam:st:D.1.6\endcsname{\hyperlink{tag-01QW}{I.\ref*{fam:st:D.1.6}}}
\expandafter\def\csname SRef@mapsI:src:A.3.2\endcsname{\hyperlink{tag-02PF}{I.\ref*{mapsI:src:A.3.2}}}
\expandafter\def\csname SRef@mapsI:real-eq-3\endcsname{\hyperlink{tag-02PF}{I.\ref*{mapsI:real-eq-3}}}
\expandafter\def\csname SRef@mapsI:src:A.3.3\endcsname{\hyperlink{tag-02PG}{I.\ref*{mapsI:src:A.3.3}}}
\expandafter\def\csname SRef@push:src:A.4.7\endcsname{\hyperlink{tag-02ZX}{I.\ref*{push:src:A.4.7}}}
\expandafter\def\csname SRef@push:A47\endcsname{\hyperlink{tag-02ZX}{I.\ref*{push:A47}}}
\expandafter\def\csname SRef@push:src:A.4.8\endcsname{\hyperlink{tag-02ZY}{I.\ref*{push:src:A.4.8}}}
\expandafter\def\csname SRef@push:A48\endcsname{\hyperlink{tag-02ZY}{I.\ref*{push:A48}}}
\expandafter\def\csname SRef@push:src:A.1.2\endcsname{\hyperlink{tag-02ZI}{I.\ref*{push:src:A.1.2}}}
\expandafter\def\csname SRef@push:A12\endcsname{\hyperlink{tag-02ZI}{I.\ref*{push:A12}}}
\expandafter\def\csname SRef@push:src:A.2.2\endcsname{\hyperlink{tag-02ZK}{I.\ref*{push:src:A.2.2}}}
\expandafter\def\csname SRef@push:A22\endcsname{\hyperlink{tag-02ZK}{I.\ref*{push:A22}}}
\expandafter\def\csname SRef@push:src:A.2.4\endcsname{\hyperlink{tag-02ZM}{I.\ref*{push:src:A.2.4}}}
\expandafter\def\csname SRef@push:A24\endcsname{\hyperlink{tag-02ZM}{I.\ref*{push:A24}}}
\expandafter\def\csname SRef@push:src:A.3.2\endcsname{\hyperlink{tag-02ZP}{I.\ref*{push:src:A.3.2}}}
\expandafter\def\csname SRef@push:A32\endcsname{\hyperlink{tag-02ZP}{I.\ref*{push:A32}}}
\expandafter\def\csname SRef@push:src:B.4.2\endcsname{\hyperlink{tag-030O}{I.\ref*{push:src:B.4.2}}}
\expandafter\def\csname SRef@push:B42\endcsname{\hyperlink{tag-030O}{I.\ref*{push:B42}}}
\expandafter\def\csname SRef@push:src:B.4.5\endcsname{\hyperlink{tag-0315}{I.\ref*{push:src:B.4.5}}}
\expandafter\def\csname SRef@push:src:B.1.2\endcsname{\hyperlink{tag-0301}{I.\ref*{push:src:B.1.2}}}
\expandafter\def\csname SRef@push:B12\endcsname{\hyperlink{tag-0301}{I.\ref*{push:B12}}}
\expandafter\def\csname SRef@push:src:B.4.33\endcsname{\hyperlink{tag-0313}{I.\ref*{push:src:B.4.33}}}
\expandafter\def\csname SRef@push:B433\endcsname{\hyperlink{tag-0313}{I.\ref*{push:B433}}}
\expandafter\def\csname SRef@push:src:B.3.1\endcsname{\hyperlink{tag-0305}{I.\ref*{push:src:B.3.1}}}
\expandafter\def\csname SRef@push:B31\endcsname{\hyperlink{tag-0305}{I.\ref*{push:B31}}}
\expandafter\def\csname SRef@push:src:B.3.4\endcsname{\hyperlink{tag-0308}{I.\ref*{push:src:B.3.4}}}
\expandafter\def\csname SRef@push:B34\endcsname{\hyperlink{tag-0308}{I.\ref*{push:B34}}}
\expandafter\def\csname SRef@push:src:D.1.3\endcsname{\hyperlink{tag-032C}{I.\ref*{push:src:D.1.3}}}
\expandafter\def\csname SRef@push:src:D.2.3\endcsname{\hyperlink{tag-032F}{I.\ref*{push:src:D.2.3}}}
\expandafter\def\csname SRef@push:src:D.3.1\endcsname{\hyperlink{tag-032H}{I.\ref*{push:src:D.3.1}}}
\expandafter\def\csname SRef@night:P2.1\endcsname{\hyperlink{tag-02R9}{I.\ref*{night:P2.1}}}
\expandafter\def\csname SRef@stmt:P2.1\endcsname{\hyperlink{tag-02R9}{I.\ref*{stmt:P2.1}}}
\expandafter\def\csname SRef@night:P2.10\endcsname{\hyperlink{tag-02RA}{I.\ref*{night:P2.10}}}
\expandafter\def\csname SRef@stmt:P2.10\endcsname{\hyperlink{tag-02RA}{I.\ref*{stmt:P2.10}}}
\expandafter\def\csname SRef@night:P2.11\endcsname{\hyperlink{tag-02RB}{I.\ref*{night:P2.11}}}
\expandafter\def\csname SRef@stmt:P2.11\endcsname{\hyperlink{tag-02RB}{I.\ref*{stmt:P2.11}}}
\expandafter\def\csname SRef@night:P2.2\endcsname{\hyperlink{tag-02RC}{I.\ref*{night:P2.2}}}
\expandafter\def\csname SRef@stmt:P2.2\endcsname{\hyperlink{tag-02RC}{I.\ref*{stmt:P2.2}}}
\expandafter\def\csname SRef@night:P2.3\endcsname{\hyperlink{tag-02RD}{I.\ref*{night:P2.3}}}
\expandafter\def\csname SRef@stmt:P2.3\endcsname{\hyperlink{tag-02RD}{I.\ref*{stmt:P2.3}}}
\expandafter\def\csname SRef@night:P2.4\endcsname{\hyperlink{tag-02RE}{I.\ref*{night:P2.4}}}
\expandafter\def\csname SRef@stmt:P2.4\endcsname{\hyperlink{tag-02RE}{I.\ref*{stmt:P2.4}}}
\expandafter\def\csname SRef@night:P2.5\endcsname{\hyperlink{tag-02RF}{I.\ref*{night:P2.5}}}
\expandafter\def\csname SRef@stmt:P2.5\endcsname{\hyperlink{tag-02RF}{I.\ref*{stmt:P2.5}}}
\expandafter\def\csname SRef@night:P2.6\endcsname{\hyperlink{tag-02RG}{I.\ref*{night:P2.6}}}
\expandafter\def\csname SRef@stmt:P2.6\endcsname{\hyperlink{tag-02RG}{I.\ref*{stmt:P2.6}}}
\expandafter\def\csname SRef@night:P2.7\endcsname{\hyperlink{tag-02RH}{I.\ref*{night:P2.7}}}
\expandafter\def\csname SRef@stmt:P2.7\endcsname{\hyperlink{tag-02RH}{I.\ref*{stmt:P2.7}}}
\expandafter\def\csname SRef@night:P2.8\endcsname{\hyperlink{tag-02RI}{I.\ref*{night:P2.8}}}
\expandafter\def\csname SRef@stmt:P2.8\endcsname{\hyperlink{tag-02RI}{I.\ref*{stmt:P2.8}}}
\expandafter\def\csname SRef@night:P2.9\endcsname{\hyperlink{tag-02RJ}{I.\ref*{night:P2.9}}}
\expandafter\def\csname SRef@stmt:P2.9\endcsname{\hyperlink{tag-02RJ}{I.\ref*{stmt:P2.9}}}
\expandafter\def\csname SRef@push:src:E.2.1\endcsname{\hyperlink{tag-032U}{I.\ref*{push:src:E.2.1}}}
\expandafter\def\csname SRef@mapsI:src:A.2.1\endcsname{\hyperlink{tag-02PD}{I.\ref*{mapsI:src:A.2.1}}}
\expandafter\def\csname SRef@push:src:E.4.1\endcsname{\hyperlink{tag-033L}{I.\ref*{push:src:E.4.1}}}
\expandafter\def\csname SRef@push:src:E.4.2\endcsname{\hyperlink{tag-033M}{I.\ref*{push:src:E.4.2}}}
\expandafter\def\csname SRef@fam:src:C.2.3\endcsname{\hyperlink{tag-01QL}{I.\ref*{fam:src:C.2.3}}}
\expandafter\def\csname SRef@ari:C.8.5\endcsname{\hyperlink{tag-01QL}{I.\ref*{ari:C.8.5}}}
\expandafter\def\csname SRef@fam:st:C.2.3\endcsname{\hyperlink{tag-01QL}{I.\ref*{fam:st:C.2.3}}}
\expandafter\def\csname SRef@push:src:E.4.3\endcsname{\hyperlink{tag-01QL}{I.\ref*{push:src:E.4.3}}}
\expandafter\def\csname SRef@fam:src:A.5.1\endcsname{\hyperlink{tag-01Q6}{I.\ref*{fam:src:A.5.1}}}
\expandafter\def\csname SRef@wb:pshift:thm:classification\endcsname{\hyperlink{tag-01Q6}{I.\ref*{wb:pshift:thm:classification}}}
\expandafter\def\csname SRef@fam:st:A.5.1\endcsname{\hyperlink{tag-01Q6}{I.\ref*{fam:st:A.5.1}}}
\expandafter\def\csname SRef@fam:st:A.5.2\endcsname{\hyperlink{tag-01RH}{I.\ref*{fam:st:A.5.2}}}
\expandafter\def\csname SRef@fam:src:A.5.2\endcsname{\hyperlink{tag-01RH}{I.\ref*{fam:src:A.5.2}}}
\expandafter\def\csname SRef@inv:D1.3\endcsname{\hyperlink{tag-01RH}{I.\ref*{inv:D1.3}}}
\expandafter\def\csname SRef@inv:D2.1\endcsname{\hyperlink{tag-01RH}{I.\ref*{inv:D2.1}}}
\expandafter\def\csname SRef@inv:D3.1\endcsname{\hyperlink{tag-01RH}{I.\ref*{inv:D3.1}}}
\expandafter\def\csname SRef@cdp:st:E9\endcsname{\hyperlink{tag-01RH}{I.\ref*{cdp:st:E9}}}
\expandafter\def\csname SRef@push:src:E.1.2\endcsname{\hyperlink{tag-032S}{I.\ref*{push:src:E.1.2}}}
\expandafter\def\csname SRef@wb:thm:canonical-ring-recovery\endcsname{\hyperlink{tag-03DM}{I.\ref*{wb:thm:canonical-ring-recovery}}}
\expandafter\def\csname SRef@wb:thm:retained-period-deformation\endcsname{\hyperlink{tag-03DP}{I.\ref*{wb:thm:retained-period-deformation}}}
\expandafter\def\csname SRef@night:P5.1\endcsname{\hyperlink{tag-02YM}{I.\ref*{night:P5.1}}}
\expandafter\def\csname SRef@stmt:P5.1\endcsname{\hyperlink{tag-02YM}{I.\ref*{stmt:P5.1}}}
\expandafter\def\csname SRef@night:P5.2\endcsname{\hyperlink{tag-02YN}{I.\ref*{night:P5.2}}}
\expandafter\def\csname SRef@stmt:P5.2\endcsname{\hyperlink{tag-02YN}{I.\ref*{stmt:P5.2}}}
\expandafter\def\csname SRef@night:P5.3\endcsname{\hyperlink{tag-02YO}{I.\ref*{night:P5.3}}}
\expandafter\def\csname SRef@stmt:P5.3\endcsname{\hyperlink{tag-02YO}{I.\ref*{stmt:P5.3}}}
\expandafter\def\csname SRef@night:P5.4\endcsname{\hyperlink{tag-02YP}{I.\ref*{night:P5.4}}}
\expandafter\def\csname SRef@stmt:P5.4\endcsname{\hyperlink{tag-02YP}{I.\ref*{stmt:P5.4}}}
\expandafter\def\csname SRef@night:P5.5\endcsname{\hyperlink{tag-02YQ}{I.\ref*{night:P5.5}}}
\expandafter\def\csname SRef@stmt:P5.5\endcsname{\hyperlink{tag-02YQ}{I.\ref*{stmt:P5.5}}}
\expandafter\def\csname SRef@night:P5.6\endcsname{\hyperlink{tag-02YR}{I.\ref*{night:P5.6}}}
\expandafter\def\csname SRef@stmt:P5.6\endcsname{\hyperlink{tag-02YR}{I.\ref*{stmt:P5.6}}}
\expandafter\def\csname SRef@cdp:eq:A.RC.2\endcsname{\href{S6_SERIES_GRIND_II.pdf\#nameddest=tag-006V}{II.{A.RC.2} (p.~78)} }
\expandafter\def\csname SRef@cdp:eq:A.RC.3\endcsname{\href{S6_SERIES_GRIND_II.pdf\#nameddest=tag-006W}{II.{A.RC.3} (p.~78)} }
\expandafter\def\csname SRef@cdp:eq:A.RC.4\endcsname{\href{S6_SERIES_GRIND_II.pdf\#nameddest=tag-006X}{II.{A.RC.4} (p.~79)} }
\expandafter\def\csname SRef@cdp:eq:A.RC.5\endcsname{\href{S6_SERIES_GRIND_II.pdf\#nameddest=tag-006Y}{II.{A.RC.5} (p.~79)} }
\expandafter\def\csname SRef@cdp:eq:A.RC.6\endcsname{\href{S6_SERIES_GRIND_II.pdf\#nameddest=tag-006Z}{II.{A.RC.6} (p.~79)} }
\expandafter\def\csname SRef@cdp:eq:A.RC.7\endcsname{\href{S6_SERIES_GRIND_II.pdf\#nameddest=tag-0070}{II.{A.RC.7} (p.~80)} }
\expandafter\def\csname SRef@cdp:eq:A.Pi.1\endcsname{\href{S6_SERIES_GRIND_II.pdf\#nameddest=tag-006O}{II.{A.Pi.1} (p.~81)} }
\expandafter\def\csname SRef@cdp:eq:A.Pi.2\endcsname{\href{S6_SERIES_GRIND_II.pdf\#nameddest=tag-006P}{II.{A.Pi.2} (p.~81)} }
\expandafter\def\csname SRef@cdp:eq:A.Pi.3\endcsname{\href{S6_SERIES_GRIND_II.pdf\#nameddest=tag-006Q}{II.{A.Pi.3} (p.~81)} }
\expandafter\def\csname SRef@cdp:eq:A.Pi.4\endcsname{\href{S6_SERIES_GRIND_II.pdf\#nameddest=tag-006R}{II.{A.Pi.4} (p.~82)} }
\expandafter\def\csname SRef@cdp:eq:A.CD.1\endcsname{\href{S6_SERIES_GRIND_II.pdf\#nameddest=tag-0058}{II.{A.CD.1} (p.~82)} }
\expandafter\def\csname SRef@cdp:eq:A.CD.2\endcsname{\href{S6_SERIES_GRIND_II.pdf\#nameddest=tag-0059}{II.{A.CD.2} (p.~82)} }
\expandafter\def\csname SRef@cdp:eq:A.CD.3\endcsname{\href{S6_SERIES_GRIND_II.pdf\#nameddest=tag-005A}{II.{A.CD.3} (p.~83)} }
\expandafter\def\csname SRef@cdp:eq:A.CD.4\endcsname{\href{S6_SERIES_GRIND_II.pdf\#nameddest=tag-005B}{II.{A.CD.4} (p.~83)} }
\expandafter\def\csname SRef@cdp:eq:A.Q1.1\endcsname{\href{S6_SERIES_GRIND_II.pdf\#nameddest=tag-006S}{II.{A.Q1.1} (p.~83)} }
\expandafter\def\csname SRef@cdp:eq:A.Q1.2\endcsname{\href{S6_SERIES_GRIND_II.pdf\#nameddest=tag-006T}{II.{A.Q1.2} (p.~84)} }
\expandafter\def\csname SRef@cdp:eq:A.J.1\endcsname{\href{S6_SERIES_GRIND_II.pdf\#nameddest=tag-005C}{II.{A.J.1} (p.~84)} }
\expandafter\def\csname SRef@cdp:eq:A.J.2\endcsname{\href{S6_SERIES_GRIND_II.pdf\#nameddest=tag-005N}{II.{A.J.2} (p.~85)} }
\expandafter\def\csname SRef@cdp:eq:A.J.3\endcsname{\href{S6_SERIES_GRIND_II.pdf\#nameddest=tag-005Y}{II.{A.J.3} (p.~85)} }
\expandafter\def\csname SRef@cdp:eq:A.J.4\endcsname{\href{S6_SERIES_GRIND_II.pdf\#nameddest=tag-006C}{II.{A.J.4} (p.~85)} }
\expandafter\def\csname SRef@cdp:eq:A.J.5\endcsname{\href{S6_SERIES_GRIND_II.pdf\#nameddest=tag-006E}{II.{A.J.5} (p.~86)} }
\expandafter\def\csname SRef@cdp:eq:A.J.6\endcsname{\href{S6_SERIES_GRIND_II.pdf\#nameddest=tag-006F}{II.{A.J.6} (p.~86)} }
\expandafter\def\csname SRef@cdp:eq:A.J.7\endcsname{\href{S6_SERIES_GRIND_II.pdf\#nameddest=tag-006G}{II.{A.J.7} (p.~86)} }
\expandafter\def\csname SRef@cdp:eq:A.J.8\endcsname{\href{S6_SERIES_GRIND_II.pdf\#nameddest=tag-006H}{II.{A.J.8} (p.~86)} }
\expandafter\def\csname SRef@cdp:eq:A.J.9\endcsname{\href{S6_SERIES_GRIND_II.pdf\#nameddest=tag-006I}{II.{A.J.9} (p.~86)} }
\expandafter\def\csname SRef@cdp:eq:A.J.10\endcsname{\href{S6_SERIES_GRIND_II.pdf\#nameddest=tag-005D}{II.{A.J.10} (p.~87)} }
\expandafter\def\csname SRef@cdp:eq:A.J.11\endcsname{\href{S6_SERIES_GRIND_II.pdf\#nameddest=tag-005E}{II.{A.J.11} (p.~88)} }
\expandafter\def\csname SRef@cdp:eq:A.J.12\endcsname{\href{S6_SERIES_GRIND_II.pdf\#nameddest=tag-005F}{II.{A.J.12} (p.~88)} }
\expandafter\def\csname SRef@cdp:eq:A.J.13\endcsname{\href{S6_SERIES_GRIND_II.pdf\#nameddest=tag-005G}{II.{A.J.13} (p.~88)} }
\expandafter\def\csname SRef@cdp:eq:A.J.14\endcsname{\href{S6_SERIES_GRIND_II.pdf\#nameddest=tag-005H}{II.{A.J.14} (p.~89)} }
\expandafter\def\csname SRef@cdp:eq:A.J.15\endcsname{\href{S6_SERIES_GRIND_II.pdf\#nameddest=tag-005I}{II.{A.J.15} (p.~89)} }
\expandafter\def\csname SRef@cdp:eq:A.J.16\endcsname{\href{S6_SERIES_GRIND_II.pdf\#nameddest=tag-005J}{II.{A.J.16} (p.~89)} }
\expandafter\def\csname SRef@cdp:eq:A.J.17\endcsname{\href{S6_SERIES_GRIND_II.pdf\#nameddest=tag-005K}{II.{A.J.17} (p.~89)} }
\expandafter\def\csname SRef@cdp:eq:A.J.18\endcsname{\href{S6_SERIES_GRIND_II.pdf\#nameddest=tag-005L}{II.{A.J.18} (p.~89)} }
\expandafter\def\csname SRef@cdp:eq:A.J.19\endcsname{\href{S6_SERIES_GRIND_II.pdf\#nameddest=tag-005M}{II.{A.J.19} (p.~90)} }
\expandafter\def\csname SRef@cdp:eq:A.J.20\endcsname{\href{S6_SERIES_GRIND_II.pdf\#nameddest=tag-005O}{II.{A.J.20} (p.~90)} }
\expandafter\def\csname SRef@cdp:eq:A.J.30\endcsname{\href{S6_SERIES_GRIND_II.pdf\#nameddest=tag-005Z}{II.{A.J.30} (p.~90)} }
\expandafter\def\csname SRef@cdp:eq:A.J.31\endcsname{\href{S6_SERIES_GRIND_II.pdf\#nameddest=tag-0060}{II.{A.J.31} (p.~91)} }
\expandafter\def\csname SRef@cdp:eq:A.J.32\endcsname{\href{S6_SERIES_GRIND_II.pdf\#nameddest=tag-0061}{II.{A.J.32} (p.~91)} }
\expandafter\def\csname SRef@cdp:eq:A.J.21\endcsname{\href{S6_SERIES_GRIND_II.pdf\#nameddest=tag-005P}{II.{A.J.21} (p.~92)} }
\expandafter\def\csname SRef@cdp:eq:A.J.22\endcsname{\href{S6_SERIES_GRIND_II.pdf\#nameddest=tag-005Q}{II.{A.J.22} (p.~92)} }
\expandafter\def\csname SRef@cdp:eq:A.J.23\endcsname{\href{S6_SERIES_GRIND_II.pdf\#nameddest=tag-005R}{II.{A.J.23} (p.~92)} }
\expandafter\def\csname SRef@cdp:eq:A.J.24\endcsname{\href{S6_SERIES_GRIND_II.pdf\#nameddest=tag-005S}{II.{A.J.24} (p.~93)} }
\expandafter\def\csname SRef@cdp:eq:A.J.25\endcsname{\href{S6_SERIES_GRIND_II.pdf\#nameddest=tag-005T}{II.{A.J.25} (p.~93)} }
\expandafter\def\csname SRef@cdp:eq:A.J.26\endcsname{\href{S6_SERIES_GRIND_II.pdf\#nameddest=tag-005U}{II.{A.J.26} (p.~93)} }
\expandafter\def\csname SRef@cdp:eq:A.J.27\endcsname{\href{S6_SERIES_GRIND_II.pdf\#nameddest=tag-005V}{II.{A.J.27} (p.~93)} }
\expandafter\def\csname SRef@cdp:eq:A.J.28\endcsname{\href{S6_SERIES_GRIND_II.pdf\#nameddest=tag-005W}{II.{A.J.28} (p.~94)} }
\expandafter\def\csname SRef@cdp:eq:A.J.29\endcsname{\href{S6_SERIES_GRIND_II.pdf\#nameddest=tag-005X}{II.{A.J.29} (p.~94)} }
\expandafter\def\csname SRef@cdp:eq:A.J.33\endcsname{\href{S6_SERIES_GRIND_II.pdf\#nameddest=tag-0062}{II.{A.J.33} (p.~95)} }
\expandafter\def\csname SRef@cdp:eq:A.J.34\endcsname{\href{S6_SERIES_GRIND_II.pdf\#nameddest=tag-0065}{II.{A.J.34} (p.~96)} }
\expandafter\def\csname SRef@cdp:eq:A.J.33a\endcsname{\href{S6_SERIES_GRIND_II.pdf\#nameddest=tag-0063}{II.{A.J.33a} (p.~96)} }
\expandafter\def\csname SRef@cdp:eq:A.J.33b\endcsname{\href{S6_SERIES_GRIND_II.pdf\#nameddest=tag-0064}{II.{A.J.33b} (p.~97)} }
\expandafter\def\csname SRef@cdp:eq:A.J.35\endcsname{\href{S6_SERIES_GRIND_II.pdf\#nameddest=tag-0066}{II.{A.J.35} (p.~97)} }
\expandafter\def\csname SRef@cdp:eq:A.J.36\endcsname{\href{S6_SERIES_GRIND_II.pdf\#nameddest=tag-0067}{II.{A.J.36} (p.~97)} }
\expandafter\def\csname SRef@cdp:eq:A.J.37\endcsname{\href{S6_SERIES_GRIND_II.pdf\#nameddest=tag-0068}{II.{A.J.37} (p.~97)} }
\expandafter\def\csname SRef@cdp:eq:A.J.38\endcsname{\href{S6_SERIES_GRIND_II.pdf\#nameddest=tag-0069}{II.{A.J.38} (p.~98)} }
\expandafter\def\csname SRef@cdp:eq:A.J.38a\endcsname{\href{S6_SERIES_GRIND_II.pdf\#nameddest=tag-006A}{II.{A.J.38a} (p.~98)} }
\expandafter\def\csname SRef@cdp:eq:A.J.39\endcsname{\href{S6_SERIES_GRIND_II.pdf\#nameddest=tag-006B}{II.{A.J.39} (p.~99)} }
\expandafter\def\csname SRef@cdp:eq:A.J.40\endcsname{\href{S6_SERIES_GRIND_II.pdf\#nameddest=tag-006D}{II.{A.J.40} (p.~100)} }
\expandafter\def\csname SRef@cdp:eq:A.T.1\endcsname{\href{S6_SERIES_GRIND_II.pdf\#nameddest=tag-007M}{II.{A.T.1} (p.~101)} }
\expandafter\def\csname SRef@cdp:eq:A.T.2\endcsname{\href{S6_SERIES_GRIND_II.pdf\#nameddest=tag-007X}{II.{A.T.2} (p.~102)} }
\expandafter\def\csname SRef@cdp:eq:A.T.3\endcsname{\href{S6_SERIES_GRIND_II.pdf\#nameddest=tag-0082}{II.{A.T.3} (p.~102)} }
\expandafter\def\csname SRef@cdp:eq:A.T.4\endcsname{\href{S6_SERIES_GRIND_II.pdf\#nameddest=tag-0083}{II.{A.T.4} (p.~102)} }
\expandafter\def\csname SRef@cdp:eq:A.T.5\endcsname{\href{S6_SERIES_GRIND_II.pdf\#nameddest=tag-0084}{II.{A.T.5} (p.~103)} }
\expandafter\def\csname SRef@cdp:eq:A.T.6\endcsname{\href{S6_SERIES_GRIND_II.pdf\#nameddest=tag-0085}{II.{A.T.6} (p.~103)} }
\expandafter\def\csname SRef@cdp:eq:A.T.7\endcsname{\href{S6_SERIES_GRIND_II.pdf\#nameddest=tag-0086}{II.{A.T.7} (p.~103)} }
\expandafter\def\csname SRef@cdp:eq:A.T.8\endcsname{\href{S6_SERIES_GRIND_II.pdf\#nameddest=tag-0087}{II.{A.T.8} (p.~103)} }
\expandafter\def\csname SRef@cdp:eq:A.T.9\endcsname{\href{S6_SERIES_GRIND_II.pdf\#nameddest=tag-0088}{II.{A.T.9} (p.~103)} }
\expandafter\def\csname SRef@cdp:eq:A.T.10\endcsname{\href{S6_SERIES_GRIND_II.pdf\#nameddest=tag-007N}{II.{A.T.10} (p.~103)} }
\expandafter\def\csname SRef@cdp:eq:A.T.11\endcsname{\href{S6_SERIES_GRIND_II.pdf\#nameddest=tag-007O}{II.{A.T.11} (p.~104)} }
\expandafter\def\csname SRef@cdp:eq:A.T.12\endcsname{\href{S6_SERIES_GRIND_II.pdf\#nameddest=tag-007P}{II.{A.T.12} (p.~104)} }
\expandafter\def\csname SRef@cdp:eq:A.T.13\endcsname{\href{S6_SERIES_GRIND_II.pdf\#nameddest=tag-007Q}{II.{A.T.13} (p.~104)} }
\expandafter\def\csname SRef@cdp:eq:A.T.14\endcsname{\href{S6_SERIES_GRIND_II.pdf\#nameddest=tag-007R}{II.{A.T.14} (p.~105)} }
\expandafter\def\csname SRef@cdp:eq:A.T.15\endcsname{\href{S6_SERIES_GRIND_II.pdf\#nameddest=tag-007S}{II.{A.T.15} (p.~105)} }
\expandafter\def\csname SRef@cdp:eq:A.T.16\endcsname{\href{S6_SERIES_GRIND_II.pdf\#nameddest=tag-007T}{II.{A.T.16} (p.~106)} }
\expandafter\def\csname SRef@cdp:eq:A.T.17\endcsname{\href{S6_SERIES_GRIND_II.pdf\#nameddest=tag-007U}{II.{A.T.17} (p.~106)} }
\expandafter\def\csname SRef@cdp:eq:A.T.18\endcsname{\href{S6_SERIES_GRIND_II.pdf\#nameddest=tag-007V}{II.{A.T.18} (p.~106)} }
\expandafter\def\csname SRef@cdp:eq:A.T.19\endcsname{\href{S6_SERIES_GRIND_II.pdf\#nameddest=tag-007W}{II.{A.T.19} (p.~106)} }
\expandafter\def\csname SRef@cdp:eq:A.T.20\endcsname{\href{S6_SERIES_GRIND_II.pdf\#nameddest=tag-007Y}{II.{A.T.20} (p.~106)} }
\expandafter\def\csname SRef@cdp:eq:A.T.21\endcsname{\href{S6_SERIES_GRIND_II.pdf\#nameddest=tag-007Z}{II.{A.T.21} (p.~107)} }
\expandafter\def\csname SRef@cdp:eq:A.T.22\endcsname{\href{S6_SERIES_GRIND_II.pdf\#nameddest=tag-0080}{II.{A.T.22} (p.~108)} }
\expandafter\def\csname SRef@cdp:eq:A.T.23\endcsname{\href{S6_SERIES_GRIND_II.pdf\#nameddest=tag-0081}{II.{A.T.23} (p.~108)} }
\expandafter\def\csname SRef@cdp:eq:A.U.1\endcsname{\href{S6_SERIES_GRIND_II.pdf\#nameddest=tag-0089}{II.{A.U.1} (p.~109)} }
\expandafter\def\csname SRef@cdp:eq:A.U.2\endcsname{\href{S6_SERIES_GRIND_II.pdf\#nameddest=tag-008A}{II.{A.U.2} (p.~109)} }
\expandafter\def\csname SRef@cdp:eq:A.U.3\endcsname{\href{S6_SERIES_GRIND_II.pdf\#nameddest=tag-008B}{II.{A.U.3} (p.~110)} }
\expandafter\def\csname SRef@cdp:eq:A.U.4\endcsname{\href{S6_SERIES_GRIND_II.pdf\#nameddest=tag-008C}{II.{A.U.4} (p.~110)} }
\expandafter\def\csname SRef@cdp:eq:A.U.5\endcsname{\href{S6_SERIES_GRIND_II.pdf\#nameddest=tag-008D}{II.{A.U.5} (p.~110)} }
\expandafter\def\csname SRef@cdp:eq:A.U.6\endcsname{\href{S6_SERIES_GRIND_II.pdf\#nameddest=tag-008E}{II.{A.U.6} (p.~111)} }
\expandafter\def\csname SRef@cdp:eq:A.U.7\endcsname{\href{S6_SERIES_GRIND_II.pdf\#nameddest=tag-008F}{II.{A.U.7} (p.~111)} }
\expandafter\def\csname SRef@cdp:eq:A.RT.1\endcsname{\href{S6_SERIES_GRIND_II.pdf\#nameddest=tag-0071}{II.{A.RT.1} (p.~112)} }
\expandafter\def\csname SRef@cdp:eq:A.RT.2\endcsname{\href{S6_SERIES_GRIND_II.pdf\#nameddest=tag-007C}{II.{A.RT.2} (p.~112)} }
\expandafter\def\csname SRef@cdp:eq:A.RT.3\endcsname{\href{S6_SERIES_GRIND_II.pdf\#nameddest=tag-007F}{II.{A.RT.3} (p.~112)} }
\expandafter\def\csname SRef@cdp:eq:A.RT.4\endcsname{\href{S6_SERIES_GRIND_II.pdf\#nameddest=tag-007G}{II.{A.RT.4} (p.~112)} }
\expandafter\def\csname SRef@cdp:eq:A.RT.5\endcsname{\href{S6_SERIES_GRIND_II.pdf\#nameddest=tag-007H}{II.{A.RT.5} (p.~112)} }
\expandafter\def\csname SRef@cdp:eq:A.RT.6\endcsname{\href{S6_SERIES_GRIND_II.pdf\#nameddest=tag-007I}{II.{A.RT.6} (p.~112)} }
\expandafter\def\csname SRef@cdp:eq:A.RT.7\endcsname{\href{S6_SERIES_GRIND_II.pdf\#nameddest=tag-007J}{II.{A.RT.7} (p.~113)} }
\expandafter\def\csname SRef@cdp:eq:A.RT.8\endcsname{\href{S6_SERIES_GRIND_II.pdf\#nameddest=tag-007K}{II.{A.RT.8} (p.~113)} }
\expandafter\def\csname SRef@cdp:eq:A.RT.9\endcsname{\href{S6_SERIES_GRIND_II.pdf\#nameddest=tag-007L}{II.{A.RT.9} (p.~113)} }
\expandafter\def\csname SRef@cdp:eq:A.RT.10\endcsname{\href{S6_SERIES_GRIND_II.pdf\#nameddest=tag-0072}{II.{A.RT.10} (p.~113)} }
\expandafter\def\csname SRef@cdp:eq:A.RT.11\endcsname{\href{S6_SERIES_GRIND_II.pdf\#nameddest=tag-0073}{II.{A.RT.11} (p.~115)} }
\expandafter\def\csname SRef@cdp:eq:A.RT.12\endcsname{\href{S6_SERIES_GRIND_II.pdf\#nameddest=tag-0074}{II.{A.RT.12} (p.~115)} }
\expandafter\def\csname SRef@cdp:eq:A.RT.13\endcsname{\href{S6_SERIES_GRIND_II.pdf\#nameddest=tag-0075}{II.{A.RT.13} (p.~115)} }
\expandafter\def\csname SRef@cdp:eq:A.RT.14\endcsname{\href{S6_SERIES_GRIND_II.pdf\#nameddest=tag-0076}{II.{A.RT.14} (p.~115)} }
\expandafter\def\csname SRef@cdp:eq:A.RT.15\endcsname{\href{S6_SERIES_GRIND_II.pdf\#nameddest=tag-0077}{II.{A.RT.15} (p.~116)} }
\expandafter\def\csname SRef@cdp:eq:A.RT.16\endcsname{\href{S6_SERIES_GRIND_II.pdf\#nameddest=tag-0078}{II.{A.RT.16} (p.~116)} }
\expandafter\def\csname SRef@cdp:eq:A.RT.17\endcsname{\href{S6_SERIES_GRIND_II.pdf\#nameddest=tag-0079}{II.{A.RT.17} (p.~116)} }
\expandafter\def\csname SRef@cdp:eq:A.RT.18\endcsname{\href{S6_SERIES_GRIND_II.pdf\#nameddest=tag-007A}{II.{A.RT.18} (p.~116)} }
\expandafter\def\csname SRef@cdp:eq:A.RT.19\endcsname{\href{S6_SERIES_GRIND_II.pdf\#nameddest=tag-007B}{II.{A.RT.19} (p.~116)} }
\expandafter\def\csname SRef@cdp:eq:A.RT.20\endcsname{\href{S6_SERIES_GRIND_II.pdf\#nameddest=tag-007D}{II.{A.RT.20} (p.~117)} }
\expandafter\def\csname SRef@cdp:eq:A.RT.21\endcsname{\href{S6_SERIES_GRIND_II.pdf\#nameddest=tag-007E}{II.{A.RT.21} (p.~117)} }
\expandafter\def\csname SRef@cdp:eq:A.BP.1\endcsname{\href{S6_SERIES_GRIND_II.pdf\#nameddest=tag-004W}{II.{A.BP.1} (p.~118)} }
\expandafter\def\csname SRef@cdp:eq:A.BP.4\endcsname{\href{S6_SERIES_GRIND_II.pdf\#nameddest=tag-0052}{II.{A.BP.4} (p.~118)} }
\expandafter\def\csname SRef@cdp:eq:A.BP.2\endcsname{\href{S6_SERIES_GRIND_II.pdf\#nameddest=tag-0050}{II.{A.BP.2} (p.~118)} }
\expandafter\def\csname SRef@cdp:eq:A.BP.3\endcsname{\href{S6_SERIES_GRIND_II.pdf\#nameddest=tag-0051}{II.{A.BP.3} (p.~118)} }
\expandafter\def\csname SRef@cdp:eq:A.BP.5\endcsname{\href{S6_SERIES_GRIND_II.pdf\#nameddest=tag-0053}{II.{A.BP.5} (p.~119)} }
\expandafter\def\csname SRef@cdp:eq:A.BP.6\endcsname{\href{S6_SERIES_GRIND_II.pdf\#nameddest=tag-0054}{II.{A.BP.6} (p.~119)} }
\expandafter\def\csname SRef@cdp:eq:A.BP.7\endcsname{\href{S6_SERIES_GRIND_II.pdf\#nameddest=tag-0055}{II.{A.BP.7} (p.~119)} }
\expandafter\def\csname SRef@cdp:eq:A.BP.8\endcsname{\href{S6_SERIES_GRIND_II.pdf\#nameddest=tag-0056}{II.{A.BP.8} (p.~119)} }
\expandafter\def\csname SRef@cdp:eq:A.BP.9\endcsname{\href{S6_SERIES_GRIND_II.pdf\#nameddest=tag-0057}{II.{A.BP.9} (p.~120)} }
\expandafter\def\csname SRef@cdp:eq:A.BP.10\endcsname{\href{S6_SERIES_GRIND_II.pdf\#nameddest=tag-004X}{II.{A.BP.10} (p.~120)} }
\expandafter\def\csname SRef@cdp:eq:A.BP.11\endcsname{\href{S6_SERIES_GRIND_II.pdf\#nameddest=tag-004Y}{II.{A.BP.11} (p.~120)} }
\expandafter\def\csname SRef@cdp:eq:A.BP.12\endcsname{\href{S6_SERIES_GRIND_II.pdf\#nameddest=tag-004Z}{II.{A.BP.12} (p.~120)} }
\expandafter\def\csname SRef@cdp:eq:C.NC.4\endcsname{\href{S6_SERIES_GRIND_II.pdf\#nameddest=tag-00S3}{II.{C.NC.4} (p.~121)} }
\expandafter\def\csname SRef@cdp:eq:C.NC.5\endcsname{\href{S6_SERIES_GRIND_II.pdf\#nameddest=tag-00S4}{II.{C.NC.5} (p.~121)} }
\expandafter\def\csname SRef@cdp:eq:C.NC.6\endcsname{\href{S6_SERIES_GRIND_II.pdf\#nameddest=tag-00S5}{II.{C.NC.6} (p.~121)} }
\expandafter\def\csname SRef@cdp:eq:C.NC.7\endcsname{\href{S6_SERIES_GRIND_II.pdf\#nameddest=tag-00S6}{II.{C.NC.7} (p.~121)} }
\expandafter\def\csname SRef@cdp:eq:C.NC.8\endcsname{\href{S6_SERIES_GRIND_II.pdf\#nameddest=tag-00S7}{II.{C.NC.8} (p.~122)} }
\expandafter\def\csname SRef@cdp:eq:C.NC.9\endcsname{\href{S6_SERIES_GRIND_II.pdf\#nameddest=tag-00S8}{II.{C.NC.9} (p.~123)} }
\expandafter\def\csname SRef@cdp:eq:C.NC.10\endcsname{\href{S6_SERIES_GRIND_II.pdf\#nameddest=tag-00S1}{II.{C.NC.10} (p.~124)} }
\expandafter\def\csname SRef@cdp:eq:C.NC.11\endcsname{\href{S6_SERIES_GRIND_II.pdf\#nameddest=tag-00S2}{II.{C.NC.11} (p.~124)} }
\expandafter\def\csname SRef@cdp:eq:C.J.1\endcsname{\href{S6_SERIES_GRIND_II.pdf\#nameddest=tag-00OV}{II.{C.J.1} (p.~125)} }
\expandafter\def\csname SRef@cdp:eq:C.J.2\endcsname{\href{S6_SERIES_GRIND_II.pdf\#nameddest=tag-00OW}{II.{C.J.2} (p.~125)} }
\expandafter\def\csname SRef@cdp:eq:C.J.3\endcsname{\href{S6_SERIES_GRIND_II.pdf\#nameddest=tag-00OX}{II.{C.J.3} (p.~126)} }
\expandafter\def\csname SRef@cdp:eq:D1\endcsname{\href{S6_SERIES_GRIND_II.pdf\#nameddest=tag-019T}{II.{D1} (p.~21)} }
\expandafter\def\csname SRef@cdp:eq:D5\endcsname{\href{S6_SERIES_GRIND_II.pdf\#nameddest=tag-019X}{II.{D5} (p.~24)} }
\expandafter\def\csname SRef@cdp:eq:D4\endcsname{\href{S6_SERIES_GRIND_II.pdf\#nameddest=tag-019W}{II.{D4} (p.~25)} }
\expandafter\def\csname SRef@cdp:eq:D2\endcsname{\href{S6_SERIES_GRIND_II.pdf\#nameddest=tag-019U}{II.{D2} (p.~28)} }
\expandafter\def\csname SRef@cdp:eq:D3\endcsname{\href{S6_SERIES_GRIND_II.pdf\#nameddest=tag-019V}{II.{D3} (p.~28)} }
\expandafter\def\csname SRef@cdp:eq:C.SK.1\endcsname{\href{S6_SERIES_GRIND_II.pdf\#nameddest=tag-011H}{II.{C.SK.1} (p.~36)} }
\expandafter\def\csname SRef@cdp:eq:C.SK.2\endcsname{\href{S6_SERIES_GRIND_II.pdf\#nameddest=tag-011I}{II.{C.SK.2} (p.~36)} }
\expandafter\def\csname SRef@a.-one-dimensional-first-jets\endcsname{\href{S6_SERIES_GRIND_II.pdf\#nameddest=tag-0005}{II.{{$\langle 1\rangle $1.}} (p.~243)} }
\expandafter\def\csname SRef@three-smooth-branches-the-complete-excess-two-classification\endcsname{\href{S6_SERIES_GRIND_II.pdf\#nameddest=tag-0005}{II.{{$\langle 1\rangle $1.}} (p.~243)} }
\expandafter\def\csname SRef@cdp:eq:C.SK.3\endcsname{\href{S6_SERIES_GRIND_II.pdf\#nameddest=tag-011J}{II.{C.SK.3} (p.~37)} }
\expandafter\def\csname SRef@b.-two-dimensional-first-jets\endcsname{\href{S6_SERIES_GRIND_II.pdf\#nameddest=tag-003V}{II.{{$\langle 1\rangle $1.}} (p.~243)} }
\expandafter\def\csname SRef@multiplicity-four-one-cusp-and-two-smooth-branches\endcsname{\href{S6_SERIES_GRIND_II.pdf\#nameddest=tag-02QH}{II.{{$\langle 1\rangle $1.}} (p.~244)} }
\expandafter\def\csname SRef@a.-a-nonzero-coefficient-of-the-cusp-cubic\endcsname{\href{S6_SERIES_GRIND_II.pdf\#nameddest=tag-0004}{II.{{$\langle 1\rangle $1.}} (p.~245)} }
\expandafter\def\csname SRef@b.-a-zero-coefficient-of-the-cusp-cubic\endcsname{\href{S6_SERIES_GRIND_II.pdf\#nameddest=tag-003U}{II.{{$\langle 1\rangle $1.}} (p.~245)} }
\expandafter\def\csname SRef@cdp:eq:C.SK.4\endcsname{\href{S6_SERIES_GRIND_II.pdf\#nameddest=tag-011K}{II.{C.SK.4} (p.~37)} }
\expandafter\def\csname SRef@every-nonordinary-triple-exceeds-the-budget\endcsname{\href{S6_SERIES_GRIND_II.pdf\#nameddest=tag-01OO}{II.{{$\langle 1\rangle $1.}} (p.~246)} }
\expandafter\def\csname SRef@cdp:eq:C.SK.5\endcsname{\href{S6_SERIES_GRIND_II.pdf\#nameddest=tag-011L}{II.{C.SK.5} (p.~37)} }
\expandafter\def\csname SRef@cdp:eq:C.15\endcsname{\href{S6_SERIES_GRIND_II.pdf\#nameddest=tag-008T}{II.{C.15} (p.~255)} }
\expandafter\def\csname SRef@cdp:eq:C.16\endcsname{\href{S6_SERIES_GRIND_II.pdf\#nameddest=tag-008U}{II.{C.16} (p.~257)} }
\expandafter\def\csname SRef@cdp:eq:C.17\endcsname{\href{S6_SERIES_GRIND_II.pdf\#nameddest=tag-008V}{II.{C.17} (p.~257)} }
\expandafter\def\csname SRef@cdp:eq:C.18\endcsname{\href{S6_SERIES_GRIND_II.pdf\#nameddest=tag-008W}{II.{C.18} (p.~258)} }
\expandafter\def\csname SRef@cdp:eq:C.SK.6\endcsname{\href{S6_SERIES_GRIND_II.pdf\#nameddest=tag-011M}{II.{C.SK.6} (p.~37)} }
\expandafter\def\csname SRef@cdp:eq:C.19\endcsname{\href{S6_SERIES_GRIND_II.pdf\#nameddest=tag-008X}{II.{C.19} (p.~258)} }
\expandafter\def\csname SRef@cdp:eq:C.20\endcsname{\href{S6_SERIES_GRIND_II.pdf\#nameddest=tag-008Z}{II.{C.20} (p.~259)} }
\expandafter\def\csname SRef@cdp:eq:C.21\endcsname{\href{S6_SERIES_GRIND_II.pdf\#nameddest=tag-0090}{II.{C.21} (p.~260)} }
\expandafter\def\csname SRef@cdp:eq:C.22\endcsname{\href{S6_SERIES_GRIND_II.pdf\#nameddest=tag-0091}{II.{C.22} (p.~261)} }
\expandafter\def\csname SRef@cdp:eq:C.22b\endcsname{\href{S6_SERIES_GRIND_II.pdf\#nameddest=tag-0093}{II.{C.22b} (p.~261)} }
\expandafter\def\csname SRef@cdp:eq:C.22a\endcsname{\href{S6_SERIES_GRIND_II.pdf\#nameddest=tag-0092}{II.{C.22a} (p.~263)} }
\expandafter\def\csname SRef@cdp:eq:C.VC.1\endcsname{\href{S6_SERIES_GRIND_II.pdf\#nameddest=tag-015G}{II.{C.VC.1} (p.~263)} }
\expandafter\def\csname SRef@cdp:eq:C.VC.2\endcsname{\href{S6_SERIES_GRIND_II.pdf\#nameddest=tag-015H}{II.{C.VC.2} (p.~263)} }
\expandafter\def\csname SRef@cdp:eq:C.VC.3\endcsname{\href{S6_SERIES_GRIND_II.pdf\#nameddest=tag-015I}{II.{C.VC.3} (p.~264)} }
\expandafter\def\csname SRef@cdp:eq:C.VC.4\endcsname{\href{S6_SERIES_GRIND_II.pdf\#nameddest=tag-015J}{II.{C.VC.4} (p.~264)} }
\expandafter\def\csname SRef@cdp:eq:C.PK.1\endcsname{\href{S6_SERIES_GRIND_II.pdf\#nameddest=tag-00V8}{II.{C.PK.1} (p.~38)} }
\expandafter\def\csname SRef@cdp:eq:C.TEST.1\endcsname{\href{S6_SERIES_GRIND_II.pdf\#nameddest=tag-012O}{II.{C.TEST.1} (p.~265)} }
\expandafter\def\csname SRef@cdp:eq:C.TEST.2\endcsname{\href{S6_SERIES_GRIND_II.pdf\#nameddest=tag-012P}{II.{C.TEST.2} (p.~265)} }
\expandafter\def\csname SRef@cdp:eq:C.3\endcsname{\href{S6_SERIES_GRIND_II.pdf\#nameddest=tag-0094}{II.{C.3} (p.~266)} }
\expandafter\def\csname SRef@cdp:eq:C.4\endcsname{\href{S6_SERIES_GRIND_II.pdf\#nameddest=tag-0095}{II.{C.4} (p.~266)} }
\expandafter\def\csname SRef@cdp:eq:C.5\endcsname{\href{S6_SERIES_GRIND_II.pdf\#nameddest=tag-0096}{II.{C.5} (p.~267)} }
\expandafter\def\csname SRef@cdp:eq:C.6\endcsname{\href{S6_SERIES_GRIND_II.pdf\#nameddest=tag-0097}{II.{C.6} (p.~267)} }
\expandafter\def\csname SRef@cdp:eq:C.7\endcsname{\href{S6_SERIES_GRIND_II.pdf\#nameddest=tag-0098}{II.{C.7} (p.~267)} }
\expandafter\def\csname SRef@cdp:eq:C.8\endcsname{\href{S6_SERIES_GRIND_II.pdf\#nameddest=tag-0099}{II.{C.8} (p.~268)} }
\expandafter\def\csname SRef@cdp:eq:C.PK.2\endcsname{\href{S6_SERIES_GRIND_II.pdf\#nameddest=tag-00V9}{II.{C.PK.2} (p.~38)} }
\expandafter\def\csname SRef@cdp:eq:C.9\endcsname{\href{S6_SERIES_GRIND_II.pdf\#nameddest=tag-009A}{II.{C.9} (p.~269)} }
\expandafter\def\csname SRef@cdp:eq:C.10\endcsname{\href{S6_SERIES_GRIND_II.pdf\#nameddest=tag-008O}{II.{C.10} (p.~269)} }
\expandafter\def\csname SRef@cdp:eq:C.11\endcsname{\href{S6_SERIES_GRIND_II.pdf\#nameddest=tag-008P}{II.{C.11} (p.~269)} }
\expandafter\def\csname SRef@cdp:eq:C.12\endcsname{\href{S6_SERIES_GRIND_II.pdf\#nameddest=tag-008Q}{II.{C.12} (p.~269)} }
\expandafter\def\csname SRef@cdp:eq:C.13\endcsname{\href{S6_SERIES_GRIND_II.pdf\#nameddest=tag-008R}{II.{C.13} (p.~270)} }
\expandafter\def\csname SRef@cdp:eq:C.14\endcsname{\href{S6_SERIES_GRIND_II.pdf\#nameddest=tag-008S}{II.{C.14} (p.~271)} }
\expandafter\def\csname SRef@cdp:eq:D.18\endcsname{\href{S6_SERIES_GRIND_II.pdf\#nameddest=tag-0195}{II.{D.18} (p.~274)} }
\expandafter\def\csname SRef@cdp:eq:D.19\endcsname{\href{S6_SERIES_GRIND_II.pdf\#nameddest=tag-0196}{II.{D.19} (p.~275)} }
\expandafter\def\csname SRef@cdp:eq:D.S.1\endcsname{\href{S6_SERIES_GRIND_II.pdf\#nameddest=tag-019O}{II.{D.S.1} (p.~276)} }
\expandafter\def\csname SRef@cdp:eq:C.PK.3\endcsname{\href{S6_SERIES_GRIND_II.pdf\#nameddest=tag-00VA}{II.{C.PK.3} (p.~38)} }
\expandafter\def\csname SRef@cdp:eq:D.S.2\endcsname{\href{S6_SERIES_GRIND_II.pdf\#nameddest=tag-019P}{II.{D.S.2} (p.~276)} }
\expandafter\def\csname SRef@cdp:eq:D.S.3\endcsname{\href{S6_SERIES_GRIND_II.pdf\#nameddest=tag-019Q}{II.{D.S.3} (p.~277)} }
\expandafter\def\csname SRef@cdp:eq:D.S.4\endcsname{\href{S6_SERIES_GRIND_II.pdf\#nameddest=tag-019R}{II.{D.S.4} (p.~277)} }
\expandafter\def\csname SRef@cdp:eq:D.S.5\endcsname{\href{S6_SERIES_GRIND_II.pdf\#nameddest=tag-019S}{II.{D.S.5} (p.~277)} }
\expandafter\def\csname SRef@cdp:eq:D.20\endcsname{\href{S6_SERIES_GRIND_II.pdf\#nameddest=tag-0197}{II.{D.20} (p.~278)} }
\expandafter\def\csname SRef@cdp:eq:D.21\endcsname{\href{S6_SERIES_GRIND_II.pdf\#nameddest=tag-0198}{II.{D.21} (p.~279)} }
\expandafter\def\csname SRef@cdp:eq:D.22\endcsname{\href{S6_SERIES_GRIND_II.pdf\#nameddest=tag-0199}{II.{D.22} (p.~280)} }
\expandafter\def\csname SRef@cdp:eq:D.23\endcsname{\href{S6_SERIES_GRIND_II.pdf\#nameddest=tag-019A}{II.{D.23} (p.~280)} }
\expandafter\def\csname SRef@cdp:eq:D.24\endcsname{\href{S6_SERIES_GRIND_II.pdf\#nameddest=tag-019B}{II.{D.24} (p.~280)} }
\expandafter\def\csname SRef@cdp:eq:C.PK.4\endcsname{\href{S6_SERIES_GRIND_II.pdf\#nameddest=tag-00VB}{II.{C.PK.4} (p.~39)} }
\expandafter\def\csname SRef@cdp:eq:D.25\endcsname{\href{S6_SERIES_GRIND_II.pdf\#nameddest=tag-019C}{II.{D.25} (p.~280)} }
\expandafter\def\csname SRef@cdp:eq:D.26\endcsname{\href{S6_SERIES_GRIND_II.pdf\#nameddest=tag-019D}{II.{D.26} (p.~281)} }
\expandafter\def\csname SRef@cdp:eq:D.27\endcsname{\href{S6_SERIES_GRIND_II.pdf\#nameddest=tag-019E}{II.{D.27} (p.~281)} }
\expandafter\def\csname SRef@cdp:eq:D.28\endcsname{\href{S6_SERIES_GRIND_II.pdf\#nameddest=tag-019F}{II.{D.28} (p.~281)} }
\expandafter\def\csname SRef@cdp:eq:D.29\endcsname{\href{S6_SERIES_GRIND_II.pdf\#nameddest=tag-019G}{II.{D.29} (p.~282)} }
\expandafter\def\csname SRef@cdp:eq:D.30\endcsname{\href{S6_SERIES_GRIND_II.pdf\#nameddest=tag-019H}{II.{D.30} (p.~282)} }
\expandafter\def\csname SRef@cdp:eq:D.31\endcsname{\href{S6_SERIES_GRIND_II.pdf\#nameddest=tag-019I}{II.{D.31} (p.~282)} }
\expandafter\def\csname SRef@cdp:eq:D.32\endcsname{\href{S6_SERIES_GRIND_II.pdf\#nameddest=tag-019J}{II.{D.32} (p.~282)} }
\expandafter\def\csname SRef@cdp:eq:D.33\endcsname{\href{S6_SERIES_GRIND_II.pdf\#nameddest=tag-019K}{II.{D.33} (p.~283)} }
\expandafter\def\csname SRef@cdp:eq:C.PK.5\endcsname{\href{S6_SERIES_GRIND_II.pdf\#nameddest=tag-00VC}{II.{C.PK.5} (p.~40)} }
\expandafter\def\csname SRef@cdp:eq:D.34\endcsname{\href{S6_SERIES_GRIND_II.pdf\#nameddest=tag-019L}{II.{D.34} (p.~283)} }
\expandafter\def\csname SRef@cdp:eq:D.35\endcsname{\href{S6_SERIES_GRIND_II.pdf\#nameddest=tag-019M}{II.{D.35} (p.~283)} }
\expandafter\def\csname SRef@cdp:eq:D.36\endcsname{\href{S6_SERIES_GRIND_II.pdf\#nameddest=tag-019N}{II.{D.36} (p.~283)} }
\expandafter\def\csname SRef@cdp:eq:E.16.1\endcsname{\href{S6_SERIES_GRIND_II.pdf\#nameddest=tag-01A2}{II.{E.16.1} (p.~304)} }
\expandafter\def\csname SRef@cdp:eq:E.16.2\endcsname{\href{S6_SERIES_GRIND_II.pdf\#nameddest=tag-01A3}{II.{E.16.2} (p.~304)} }
\expandafter\def\csname SRef@cdp:eq:E.16.3\endcsname{\href{S6_SERIES_GRIND_II.pdf\#nameddest=tag-01A4}{II.{E.16.3} (p.~305)} }
\expandafter\def\csname SRef@cdp:eq:C.PK.6\endcsname{\href{S6_SERIES_GRIND_II.pdf\#nameddest=tag-00VD}{II.{C.PK.6} (p.~40)} }
\expandafter\def\csname SRef@cdp:eq:C.PK.7\endcsname{\href{S6_SERIES_GRIND_II.pdf\#nameddest=tag-00VE}{II.{C.PK.7} (p.~40)} }
\expandafter\def\csname SRef@cdp:eq:C.PK.8\endcsname{\href{S6_SERIES_GRIND_II.pdf\#nameddest=tag-00VF}{II.{C.PK.8} (p.~40)} }
\expandafter\def\csname SRef@cdp:eq:C.IN.1\endcsname{\href{S6_SERIES_GRIND_II.pdf\#nameddest=tag-00NK}{II.{C.IN.1} (p.~41)} }
\expandafter\def\csname SRef@invII:eq:E.2\endcsname{\href{S6_SERIES_GRIND_II.pdf\#nameddest=tag-01TR}{II.{E.2} (p.~354)} }
\expandafter\def\csname SRef@cdp:eq:C.IN.2\endcsname{\href{S6_SERIES_GRIND_II.pdf\#nameddest=tag-00NL}{II.{C.IN.2} (p.~41)} }
\expandafter\def\csname SRef@invII:eq:E.3\endcsname{\href{S6_SERIES_GRIND_II.pdf\#nameddest=tag-01TS}{II.{E.3} (p.~354)} }
\expandafter\def\csname SRef@invII:eq:E.4a\endcsname{\href{S6_SERIES_GRIND_II.pdf\#nameddest=tag-01TT}{II.{E.4a} (p.~354)} }
\expandafter\def\csname SRef@invII:eq:E.5\endcsname{\href{S6_SERIES_GRIND_II.pdf\#nameddest=tag-01TU}{II.{E.5} (p.~354)} }
\expandafter\def\csname SRef@invII:eq:E.6\endcsname{\href{S6_SERIES_GRIND_II.pdf\#nameddest=tag-01TV}{II.{E.6} (p.~356)} }
\expandafter\def\csname SRef@inv:wb:eq:9.19\endcsname{\href{S6_SERIES_GRIND_II.pdf\#nameddest=tag-01TJ}{II.{WBAPII.9.19} (p.~362)} }
\expandafter\def\csname SRef@inv:wb:eq:9.20\endcsname{\href{S6_SERIES_GRIND_II.pdf\#nameddest=tag-01TK}{II.{WBAPII.9.20} (p.~362)} }
\expandafter\def\csname SRef@inv:wb:eq:9.21\endcsname{\href{S6_SERIES_GRIND_II.pdf\#nameddest=tag-01TL}{II.{WBAPII.9.21} (p.~362)} }
\expandafter\def\csname SRef@inv:wb:eq:9.22\endcsname{\href{S6_SERIES_GRIND_II.pdf\#nameddest=tag-01TM}{II.{WBAPII.9.22} (p.~362)} }
\expandafter\def\csname SRef@inv:wb:eq:9.23\endcsname{\href{S6_SERIES_GRIND_II.pdf\#nameddest=tag-01TN}{II.{WBAPII.9.23} (p.~362)} }
\expandafter\def\csname SRef@inv:wb:eq:9.24\endcsname{\href{S6_SERIES_GRIND_II.pdf\#nameddest=tag-01TO}{II.{WBAPII.9.24} (p.~362)} }
\expandafter\def\csname SRef@cdp:eq:C.IN.3\endcsname{\href{S6_SERIES_GRIND_II.pdf\#nameddest=tag-00NM}{II.{C.IN.3} (p.~41)} }
\expandafter\def\csname SRef@inv:wb:eq:9.25\endcsname{\href{S6_SERIES_GRIND_II.pdf\#nameddest=tag-01TP}{II.{WBAPII.9.25} (p.~362)} }
\expandafter\def\csname SRef@wbII:eq:CN:2.2\endcsname{\href{S6_SERIES_GRIND_II.pdf\#nameddest=tag-03DW}{II.{W.145} (p.~364)} }
\expandafter\def\csname SRef@wbII:eq:W:5.3\endcsname{\href{S6_SERIES_GRIND_II.pdf\#nameddest=tag-03HS}{II.{W.3} (p.~365)} }
\expandafter\def\csname SRef@wbII:eq:W:5.4\endcsname{\href{S6_SERIES_GRIND_II.pdf\#nameddest=tag-03HX}{II.{W.4} (p.~365)} }
\expandafter\def\csname SRef@wbII:eq:W:5.5\endcsname{\href{S6_SERIES_GRIND_II.pdf\#nameddest=tag-03HY}{II.{W.5} (p.~366)} }
\expandafter\def\csname SRef@wbII:eq:W:5.6\endcsname{\href{S6_SERIES_GRIND_II.pdf\#nameddest=tag-03HZ}{II.{W.6} (p.~366)} }
\expandafter\def\csname SRef@cdp:eq:C.IN.4\endcsname{\href{S6_SERIES_GRIND_II.pdf\#nameddest=tag-00NN}{II.{C.IN.4} (p.~42)} }
\expandafter\def\csname SRef@wbII:eq:W:5.7\endcsname{\href{S6_SERIES_GRIND_II.pdf\#nameddest=tag-03I0}{II.{W.7} (p.~366)} }
\expandafter\def\csname SRef@wbII:eq:W:5.8\endcsname{\href{S6_SERIES_GRIND_II.pdf\#nameddest=tag-03I1}{II.{W.8} (p.~368)} }
\expandafter\def\csname SRef@wbII:eq:W:5.9\endcsname{\href{S6_SERIES_GRIND_II.pdf\#nameddest=tag-03I2}{II.{W.9} (p.~368)} }
\expandafter\def\csname SRef@wbII:eq:W:5.10\endcsname{\href{S6_SERIES_GRIND_II.pdf\#nameddest=tag-03EA}{II.{W.10} (p.~368)} }
\expandafter\def\csname SRef@wbII:eq:CN:4.3\endcsname{\href{S6_SERIES_GRIND_II.pdf\#nameddest=tag-03DX}{II.{W.160} (p.~369)} }
\expandafter\def\csname SRef@wbII:eq:CN:4.4\endcsname{\href{S6_SERIES_GRIND_II.pdf\#nameddest=tag-03DY}{II.{W.161} (p.~369)} }
\expandafter\def\csname SRef@wbII:eq:CN:4.5a\endcsname{\href{S6_SERIES_GRIND_II.pdf\#nameddest=tag-03E0}{II.{W.162} (p.~370)} }
\expandafter\def\csname SRef@cdp:eq:C.IN.5\endcsname{\href{S6_SERIES_GRIND_II.pdf\#nameddest=tag-00NO}{II.{C.IN.5} (p.~42)} }
\expandafter\def\csname SRef@wbII:eq:CN:4.5\endcsname{\href{S6_SERIES_GRIND_II.pdf\#nameddest=tag-03DZ}{II.{W.163} (p.~370)} }
\expandafter\def\csname SRef@wbII:eq:CN:4.6\endcsname{\href{S6_SERIES_GRIND_II.pdf\#nameddest=tag-03E1}{II.{W.164} (p.~370)} }
\expandafter\def\csname SRef@wbII:eq:CN:4.7\endcsname{\href{S6_SERIES_GRIND_II.pdf\#nameddest=tag-03E2}{II.{W.165} (p.~370)} }
\expandafter\def\csname SRef@wbII:eq:CN:4.8\endcsname{\href{S6_SERIES_GRIND_II.pdf\#nameddest=tag-03E3}{II.{W.166} (p.~371)} }
\expandafter\def\csname SRef@wbII:eq:CN:4.9\endcsname{\href{S6_SERIES_GRIND_II.pdf\#nameddest=tag-03E4}{II.{W.167} (p.~371)} }
\expandafter\def\csname SRef@wbII:eq:W:5.11\endcsname{\href{S6_SERIES_GRIND_II.pdf\#nameddest=tag-03EB}{II.{W.11} (p.~371)} }
\expandafter\def\csname SRef@wbII:eq:W:5.15a\endcsname{\href{S6_SERIES_GRIND_II.pdf\#nameddest=tag-03EC}{II.{W.12} (p.~372)} }
\expandafter\def\csname SRef@wbII:eq:W:5.16\endcsname{\href{S6_SERIES_GRIND_II.pdf\#nameddest=tag-03ED}{II.{W.13} (p.~374)} }
\expandafter\def\csname SRef@wbII:eq:W:5.17\endcsname{\href{S6_SERIES_GRIND_II.pdf\#nameddest=tag-03EE}{II.{W.14} (p.~374)} }
\expandafter\def\csname SRef@cdp:eq:C.IN.6\endcsname{\href{S6_SERIES_GRIND_II.pdf\#nameddest=tag-00NP}{II.{C.IN.6} (p.~43)} }
\expandafter\def\csname SRef@wbII:eq:W:5.17-V\endcsname{\href{S6_SERIES_GRIND_II.pdf\#nameddest=tag-03FU}{II.{W.15} (p.~374)} }
\expandafter\def\csname SRef@wbII:eq:W:5.17-T\endcsname{\href{S6_SERIES_GRIND_II.pdf\#nameddest=tag-03FS}{II.{W.16} (p.~374)} }
\expandafter\def\csname SRef@wbII:eq:W:5.17-U\endcsname{\href{S6_SERIES_GRIND_II.pdf\#nameddest=tag-03FT}{II.{W.17} (p.~374)} }
\expandafter\def\csname SRef@wbII:eq:W:5.17a\endcsname{\href{S6_SERIES_GRIND_II.pdf\#nameddest=tag-03FV}{II.{W.18} (p.~376)} }
\expandafter\def\csname SRef@wbII:eq:W:5.17b\endcsname{\href{S6_SERIES_GRIND_II.pdf\#nameddest=tag-03GL}{II.{W.19} (p.~376)} }
\expandafter\def\csname SRef@wbII:eq:W:5.17c\endcsname{\href{S6_SERIES_GRIND_II.pdf\#nameddest=tag-03GM}{II.{W.20} (p.~376)} }
\expandafter\def\csname SRef@wbII:eq:W:5.17d\endcsname{\href{S6_SERIES_GRIND_II.pdf\#nameddest=tag-03GN}{II.{W.21} (p.~377)} }
\expandafter\def\csname SRef@wbII:eq:W:5.17e\endcsname{\href{S6_SERIES_GRIND_II.pdf\#nameddest=tag-03GO}{II.{W.22} (p.~377)} }
\expandafter\def\csname SRef@wbII:eq:W:5.17g\endcsname{\href{S6_SERIES_GRIND_II.pdf\#nameddest=tag-03GQ}{II.{W.23} (p.~377)} }
\expandafter\def\csname SRef@wbII:eq:W:5.17h\endcsname{\href{S6_SERIES_GRIND_II.pdf\#nameddest=tag-03GR}{II.{W.24} (p.~377)} }
\expandafter\def\csname SRef@cdp:eq:C.PK.9\endcsname{\href{S6_SERIES_GRIND_II.pdf\#nameddest=tag-00VG}{II.{C.PK.9} (p.~44)} }
\expandafter\def\csname SRef@wbII:eq:W:5.17i\endcsname{\href{S6_SERIES_GRIND_II.pdf\#nameddest=tag-03GS}{II.{W.25} (p.~377)} }
\expandafter\def\csname SRef@wbII:eq:W:5.17j'\endcsname{\href{S6_SERIES_GRIND_II.pdf\#nameddest=tag-03GU}{II.{W.26} (p.~378)} }
\expandafter\def\csname SRef@wbII:eq:W:5.17j\endcsname{\href{S6_SERIES_GRIND_II.pdf\#nameddest=tag-03GT}{II.{W.27} (p.~378)} }
\expandafter\def\csname SRef@wbII:eq:W:5.17k\endcsname{\href{S6_SERIES_GRIND_II.pdf\#nameddest=tag-03GV}{II.{W.28} (p.~378)} }
\expandafter\def\csname SRef@wbII:eq:W:5.17l\endcsname{\href{S6_SERIES_GRIND_II.pdf\#nameddest=tag-03GW}{II.{W.29} (p.~379)} }
\expandafter\def\csname SRef@wbII:eq:W:5.17l'\endcsname{\href{S6_SERIES_GRIND_II.pdf\#nameddest=tag-03GX}{II.{W.30} (p.~379)} }
\expandafter\def\csname SRef@wbII:eq:W:5.17l''\endcsname{\href{S6_SERIES_GRIND_II.pdf\#nameddest=tag-03GY}{II.{W.31} (p.~379)} }
\expandafter\def\csname SRef@wbII:eq:W:5.17l'''\endcsname{\href{S6_SERIES_GRIND_II.pdf\#nameddest=tag-03GZ}{II.{W.32} (p.~379)} }
\expandafter\def\csname SRef@wbII:eq:W:5.17m\endcsname{\href{S6_SERIES_GRIND_II.pdf\#nameddest=tag-03H0}{II.{W.33} (p.~379)} }
\expandafter\def\csname SRef@wbII:eq:W:5.17m'\endcsname{\href{S6_SERIES_GRIND_II.pdf\#nameddest=tag-03H1}{II.{W.34} (p.~379)} }
\expandafter\def\csname SRef@cdp:eq:C.H.1\endcsname{\href{S6_SERIES_GRIND_II.pdf\#nameddest=tag-00L8}{II.{C.H.1} (p.~46)} }
\expandafter\def\csname SRef@wbII:eq:W:5.17m''\endcsname{\href{S6_SERIES_GRIND_II.pdf\#nameddest=tag-03H2}{II.{W.35} (p.~380)} }
\expandafter\def\csname SRef@wbII:eq:W:5.17n\endcsname{\href{S6_SERIES_GRIND_II.pdf\#nameddest=tag-03H3}{II.{W.36} (p.~380)} }
\expandafter\def\csname SRef@wbII:eq:W:5.17o\endcsname{\href{S6_SERIES_GRIND_II.pdf\#nameddest=tag-03H4}{II.{W.37} (p.~380)} }
\expandafter\def\csname SRef@wbII:eq:W:5.17p\endcsname{\href{S6_SERIES_GRIND_II.pdf\#nameddest=tag-03H5}{II.{W.38} (p.~381)} }
\expandafter\def\csname SRef@wbII:eq:W:5.17p0\endcsname{\href{S6_SERIES_GRIND_II.pdf\#nameddest=tag-03H7}{II.{W.39} (p.~381)} }
\expandafter\def\csname SRef@wbII:eq:W:5.17p'\endcsname{\href{S6_SERIES_GRIND_II.pdf\#nameddest=tag-03H6}{II.{W.40} (p.~381)} }
\expandafter\def\csname SRef@wbII:eq:W:5.17q\endcsname{\href{S6_SERIES_GRIND_II.pdf\#nameddest=tag-03H8}{II.{W.41} (p.~381)} }
\expandafter\def\csname SRef@wbII:eq:W:5.17r\endcsname{\href{S6_SERIES_GRIND_II.pdf\#nameddest=tag-03H9}{II.{W.42} (p.~382)} }
\expandafter\def\csname SRef@wbII:eq:W:5.17s\endcsname{\href{S6_SERIES_GRIND_II.pdf\#nameddest=tag-03HA}{II.{W.43} (p.~382)} }
\expandafter\def\csname SRef@wbII:eq:W:5.17t\endcsname{\href{S6_SERIES_GRIND_II.pdf\#nameddest=tag-03HB}{II.{W.61} (p.~382)} }
\expandafter\def\csname SRef@cdp:eq:A.1\endcsname{\href{S6_SERIES_GRIND_II.pdf\#nameddest=tag-003Y}{II.{A.1} (p.~48)} }
\expandafter\def\csname SRef@wbII:eq:W:5.17u\endcsname{\href{S6_SERIES_GRIND_II.pdf\#nameddest=tag-03HC}{II.{W.62} (p.~382)} }
\expandafter\def\csname SRef@wbII:eq:W:5.17v\endcsname{\href{S6_SERIES_GRIND_II.pdf\#nameddest=tag-03HD}{II.{W.63} (p.~382)} }
\expandafter\def\csname SRef@wbII:eq:W:5.17w\endcsname{\href{S6_SERIES_GRIND_II.pdf\#nameddest=tag-03HE}{II.{W.64} (p.~383)} }
\expandafter\def\csname SRef@wbII:eq:W:5.17x\endcsname{\href{S6_SERIES_GRIND_II.pdf\#nameddest=tag-03HF}{II.{W.65} (p.~383)} }
\expandafter\def\csname SRef@wbII:eq:W:5.17y\endcsname{\href{S6_SERIES_GRIND_II.pdf\#nameddest=tag-03HG}{II.{W.66} (p.~383)} }
\expandafter\def\csname SRef@wbII:eq:W:5.17z\endcsname{\href{S6_SERIES_GRIND_II.pdf\#nameddest=tag-03HH}{II.{W.67} (p.~383)} }
\expandafter\def\csname SRef@wbII:eq:W:5.17aa\endcsname{\href{S6_SERIES_GRIND_II.pdf\#nameddest=tag-03FW}{II.{W.68} (p.~383)} }
\expandafter\def\csname SRef@wbII:eq:W:5.17ab\endcsname{\href{S6_SERIES_GRIND_II.pdf\#nameddest=tag-03FX}{II.{W.69} (p.~383)} }
\expandafter\def\csname SRef@wbII:eq:W:5.17ac\endcsname{\href{S6_SERIES_GRIND_II.pdf\#nameddest=tag-03FY}{II.{W.70} (p.~384)} }
\expandafter\def\csname SRef@wbII:eq:W:5.17ad\endcsname{\href{S6_SERIES_GRIND_II.pdf\#nameddest=tag-03FZ}{II.{W.71} (p.~384)} }
\expandafter\def\csname SRef@cdp:eq:A.2\endcsname{\href{S6_SERIES_GRIND_II.pdf\#nameddest=tag-004B}{II.{A.2} (p.~48)} }
\expandafter\def\csname SRef@wbII:eq:W:5.17ae\endcsname{\href{S6_SERIES_GRIND_II.pdf\#nameddest=tag-03G0}{II.{W.72} (p.~384)} }
\expandafter\def\csname SRef@wbII:eq:W:5.17af\endcsname{\href{S6_SERIES_GRIND_II.pdf\#nameddest=tag-03G1}{II.{W.73} (p.~384)} }
\expandafter\def\csname SRef@wbII:eq:W:5.17ag\endcsname{\href{S6_SERIES_GRIND_II.pdf\#nameddest=tag-03G2}{II.{W.74} (p.~384)} }
\expandafter\def\csname SRef@wbII:eq:W:5.17ah\endcsname{\href{S6_SERIES_GRIND_II.pdf\#nameddest=tag-03G3}{II.{W.75} (p.~384)} }
\expandafter\def\csname SRef@wbII:eq:W:5.17ai\endcsname{\href{S6_SERIES_GRIND_II.pdf\#nameddest=tag-03G4}{II.{W.76} (p.~384)} }
\expandafter\def\csname SRef@wbII:eq:W:5.17aj\endcsname{\href{S6_SERIES_GRIND_II.pdf\#nameddest=tag-03G5}{II.{W.77} (p.~384)} }
\expandafter\def\csname SRef@wbII:eq:W:5.17ak\endcsname{\href{S6_SERIES_GRIND_II.pdf\#nameddest=tag-03G6}{II.{W.78} (p.~384)} }
\expandafter\def\csname SRef@wbII:eq:W:5.17al\endcsname{\href{S6_SERIES_GRIND_II.pdf\#nameddest=tag-03G7}{II.{W.79} (p.~385)} }
\expandafter\def\csname SRef@cdp:eq:A.3\endcsname{\href{S6_SERIES_GRIND_II.pdf\#nameddest=tag-004M}{II.{A.3} (p.~49)} }
\expandafter\def\csname SRef@wbII:eq:W:5.17am\endcsname{\href{S6_SERIES_GRIND_II.pdf\#nameddest=tag-03G8}{II.{W.80} (p.~385)} }
\expandafter\def\csname SRef@wbII:eq:W:5.17an\endcsname{\href{S6_SERIES_GRIND_II.pdf\#nameddest=tag-03G9}{II.{W.81} (p.~385)} }
\expandafter\def\csname SRef@wbII:eq:W:5.17ao\endcsname{\href{S6_SERIES_GRIND_II.pdf\#nameddest=tag-03GA}{II.{W.82} (p.~385)} }
\expandafter\def\csname SRef@wbII:eq:W:5.17ap\endcsname{\href{S6_SERIES_GRIND_II.pdf\#nameddest=tag-03GB}{II.{W.83} (p.~385)} }
\expandafter\def\csname SRef@wbII:eq:W:5.17aq\endcsname{\href{S6_SERIES_GRIND_II.pdf\#nameddest=tag-03GC}{II.{W.84} (p.~385)} }
\expandafter\def\csname SRef@wbII:eq:W:5.17ar\endcsname{\href{S6_SERIES_GRIND_II.pdf\#nameddest=tag-03GD}{II.{W.85} (p.~385)} }
\expandafter\def\csname SRef@wbII:eq:W:5.17as\endcsname{\href{S6_SERIES_GRIND_II.pdf\#nameddest=tag-03GE}{II.{W.86} (p.~385)} }
\expandafter\def\csname SRef@wbII:eq:W:5.17at\endcsname{\href{S6_SERIES_GRIND_II.pdf\#nameddest=tag-03GF}{II.{W.87} (p.~386)} }
\expandafter\def\csname SRef@wbII:eq:W:5.17au\endcsname{\href{S6_SERIES_GRIND_II.pdf\#nameddest=tag-03GG}{II.{W.88} (p.~386)} }
\expandafter\def\csname SRef@wbII:eq:W:5.17av\endcsname{\href{S6_SERIES_GRIND_II.pdf\#nameddest=tag-03GH}{II.{W.89} (p.~386)} }
\expandafter\def\csname SRef@cdp:eq:A.4\endcsname{\href{S6_SERIES_GRIND_II.pdf\#nameddest=tag-004O}{II.{A.4} (p.~49)} }
\expandafter\def\csname SRef@wbII:eq:W:5.17aw\endcsname{\href{S6_SERIES_GRIND_II.pdf\#nameddest=tag-03GI}{II.{W.90} (p.~386)} }
\expandafter\def\csname SRef@wbII:eq:W:5.17ax\endcsname{\href{S6_SERIES_GRIND_II.pdf\#nameddest=tag-03GJ}{II.{W.91} (p.~386)} }
\expandafter\def\csname SRef@wbII:eq:W:5.17ay\endcsname{\href{S6_SERIES_GRIND_II.pdf\#nameddest=tag-03GK}{II.{W.92} (p.~386)} }
\expandafter\def\csname SRef@wbII:eq:W:5.17-A1\endcsname{\href{S6_SERIES_GRIND_II.pdf\#nameddest=tag-03EF}{II.{W.93} (p.~387)} }
\expandafter\def\csname SRef@wbII:eq:W:5.17-A2\endcsname{\href{S6_SERIES_GRIND_II.pdf\#nameddest=tag-03EY}{II.{W.94} (p.~388)} }
\expandafter\def\csname SRef@wbII:eq:W:5.17-A3\endcsname{\href{S6_SERIES_GRIND_II.pdf\#nameddest=tag-03F9}{II.{W.95} (p.~388)} }
\expandafter\def\csname SRef@wbII:eq:W:5.17-A4\endcsname{\href{S6_SERIES_GRIND_II.pdf\#nameddest=tag-03FM}{II.{W.96} (p.~388)} }
\expandafter\def\csname SRef@wbII:eq:W:5.17-A5\endcsname{\href{S6_SERIES_GRIND_II.pdf\#nameddest=tag-03FN}{II.{W.97} (p.~388)} }
\expandafter\def\csname SRef@wbII:eq:W:5.17-A6\endcsname{\href{S6_SERIES_GRIND_II.pdf\#nameddest=tag-03FO}{II.{W.98} (p.~388)} }
\expandafter\def\csname SRef@wbII:eq:W:5.17-A7\endcsname{\href{S6_SERIES_GRIND_II.pdf\#nameddest=tag-03FP}{II.{W.99} (p.~388)} }
\expandafter\def\csname SRef@wbII:eq:W:5.17-A8\endcsname{\href{S6_SERIES_GRIND_II.pdf\#nameddest=tag-03FQ}{II.{W.100} (p.~388)} }
\expandafter\def\csname SRef@wbII:eq:W:5.17-A9\endcsname{\href{S6_SERIES_GRIND_II.pdf\#nameddest=tag-03FR}{II.{W.101} (p.~388)} }
\expandafter\def\csname SRef@wbII:eq:W:5.17-A10\endcsname{\href{S6_SERIES_GRIND_II.pdf\#nameddest=tag-03EG}{II.{W.102} (p.~388)} }
\expandafter\def\csname SRef@wbII:eq:W:5.17-A11\endcsname{\href{S6_SERIES_GRIND_II.pdf\#nameddest=tag-03EH}{II.{W.103} (p.~388)} }
\expandafter\def\csname SRef@wbII:eq:W:5.17-A12\endcsname{\href{S6_SERIES_GRIND_II.pdf\#nameddest=tag-03EI}{II.{W.104} (p.~389)} }
\expandafter\def\csname SRef@wbII:eq:W:5.17-A13\endcsname{\href{S6_SERIES_GRIND_II.pdf\#nameddest=tag-03EJ}{II.{W.105} (p.~389)} }
\expandafter\def\csname SRef@wbII:eq:W:5.17-A14\endcsname{\href{S6_SERIES_GRIND_II.pdf\#nameddest=tag-03EK}{II.{W.106} (p.~389)} }
\expandafter\def\csname SRef@wbII:eq:W:5.17-A14a\endcsname{\href{S6_SERIES_GRIND_II.pdf\#nameddest=tag-03EL}{II.{W.107} (p.~390)} }
\expandafter\def\csname SRef@cdp:eq:A.5\endcsname{\href{S6_SERIES_GRIND_II.pdf\#nameddest=tag-004P}{II.{A.5} (p.~49)} }
\expandafter\def\csname SRef@wbII:eq:W:5.17-A14c\endcsname{\href{S6_SERIES_GRIND_II.pdf\#nameddest=tag-03EM}{II.{W.108} (p.~390)} }
\expandafter\def\csname SRef@wbII:eq:W:5.17-A14d\endcsname{\href{S6_SERIES_GRIND_II.pdf\#nameddest=tag-03EN}{II.{W.109} (p.~390)} }
\expandafter\def\csname SRef@wbII:eq:W:5.17-A14e\endcsname{\href{S6_SERIES_GRIND_II.pdf\#nameddest=tag-03EO}{II.{W.110} (p.~392)} }
\expandafter\def\csname SRef@wbII:eq:W:5.17-A14f\endcsname{\href{S6_SERIES_GRIND_II.pdf\#nameddest=tag-03EP}{II.{W.111} (p.~392)} }
\expandafter\def\csname SRef@wbII:eq:W:5.17-A15\endcsname{\href{S6_SERIES_GRIND_II.pdf\#nameddest=tag-03EQ}{II.{W.112} (p.~393)} }
\expandafter\def\csname SRef@wbII:eq:W:5.17-A16\endcsname{\href{S6_SERIES_GRIND_II.pdf\#nameddest=tag-03ER}{II.{W.113} (p.~393)} }
\expandafter\def\csname SRef@wbII:eq:W:5.17-A17\endcsname{\href{S6_SERIES_GRIND_II.pdf\#nameddest=tag-03ES}{II.{W.114} (p.~393)} }
\expandafter\def\csname SRef@wbII:eq:W:5.17-A18\endcsname{\href{S6_SERIES_GRIND_II.pdf\#nameddest=tag-03ET}{II.{W.115} (p.~393)} }
\expandafter\def\csname SRef@wbII:eq:W:5.17-A19\endcsname{\href{S6_SERIES_GRIND_II.pdf\#nameddest=tag-03EU}{II.{W.116} (p.~393)} }
\expandafter\def\csname SRef@wbII:eq:W:5.17-A19a\endcsname{\href{S6_SERIES_GRIND_II.pdf\#nameddest=tag-03EV}{II.{W.117} (p.~393)} }
\expandafter\def\csname SRef@cdp:eq:A.6\endcsname{\href{S6_SERIES_GRIND_II.pdf\#nameddest=tag-004Q}{II.{A.6} (p.~50)} }
\expandafter\def\csname SRef@wbII:eq:W:5.17-A19b\endcsname{\href{S6_SERIES_GRIND_II.pdf\#nameddest=tag-03EW}{II.{W.118} (p.~393)} }
\expandafter\def\csname SRef@wbII:eq:W:5.17-A19c\endcsname{\href{S6_SERIES_GRIND_II.pdf\#nameddest=tag-03EX}{II.{W.119} (p.~394)} }
\expandafter\def\csname SRef@wbII:eq:W:5.17-A20\endcsname{\href{S6_SERIES_GRIND_II.pdf\#nameddest=tag-03EZ}{II.{W.120} (p.~394)} }
\expandafter\def\csname SRef@wbII:eq:W:5.17-A21\endcsname{\href{S6_SERIES_GRIND_II.pdf\#nameddest=tag-03F0}{II.{W.121} (p.~394)} }
\expandafter\def\csname SRef@wbII:eq:W:5.17-A22\endcsname{\href{S6_SERIES_GRIND_II.pdf\#nameddest=tag-03F1}{II.{W.122} (p.~394)} }
\expandafter\def\csname SRef@wbII:eq:W:5.17-A23\endcsname{\href{S6_SERIES_GRIND_II.pdf\#nameddest=tag-03F2}{II.{W.123} (p.~395)} }
\expandafter\def\csname SRef@wbII:eq:W:5.17-A24\endcsname{\href{S6_SERIES_GRIND_II.pdf\#nameddest=tag-03F3}{II.{W.124} (p.~395)} }
\expandafter\def\csname SRef@wbII:eq:W:5.17-A25\endcsname{\href{S6_SERIES_GRIND_II.pdf\#nameddest=tag-03F4}{II.{W.125} (p.~395)} }
\expandafter\def\csname SRef@wbII:eq:W:5.17-A26\endcsname{\href{S6_SERIES_GRIND_II.pdf\#nameddest=tag-03F5}{II.{W.126} (p.~395)} }
\expandafter\def\csname SRef@wbII:eq:W:5.17-A27\endcsname{\href{S6_SERIES_GRIND_II.pdf\#nameddest=tag-03F6}{II.{W.127} (p.~395)} }
\expandafter\def\csname SRef@cdp:eq:A.7\endcsname{\href{S6_SERIES_GRIND_II.pdf\#nameddest=tag-004R}{II.{A.7} (p.~50)} }
\expandafter\def\csname SRef@wbII:eq:W:5.17-A28\endcsname{\href{S6_SERIES_GRIND_II.pdf\#nameddest=tag-03F7}{II.{W.128} (p.~395)} }
\expandafter\def\csname SRef@wbII:eq:W:5.17-A29\endcsname{\href{S6_SERIES_GRIND_II.pdf\#nameddest=tag-03F8}{II.{W.129} (p.~395)} }
\expandafter\def\csname SRef@wbII:eq:W:5.17-A30\endcsname{\href{S6_SERIES_GRIND_II.pdf\#nameddest=tag-03FA}{II.{W.130} (p.~395)} }
\expandafter\def\csname SRef@wbII:eq:W:5.17-A31\endcsname{\href{S6_SERIES_GRIND_II.pdf\#nameddest=tag-03FB}{II.{W.131} (p.~395)} }
\expandafter\def\csname SRef@wbII:eq:W:5.17-A32\endcsname{\href{S6_SERIES_GRIND_II.pdf\#nameddest=tag-03FC}{II.{W.132} (p.~396)} }
\expandafter\def\csname SRef@wbII:eq:W:5.17-A33\endcsname{\href{S6_SERIES_GRIND_II.pdf\#nameddest=tag-03FD}{II.{W.133} (p.~396)} }
\expandafter\def\csname SRef@wbII:eq:W:5.17-A34\endcsname{\href{S6_SERIES_GRIND_II.pdf\#nameddest=tag-03FE}{II.{W.134} (p.~396)} }
\expandafter\def\csname SRef@wbII:eq:W:5.17-A34a\endcsname{\href{S6_SERIES_GRIND_II.pdf\#nameddest=tag-03FF}{II.{W.135} (p.~396)} }
\expandafter\def\csname SRef@wbII:eq:W:5.17-A34b\endcsname{\href{S6_SERIES_GRIND_II.pdf\#nameddest=tag-03FG}{II.{W.136} (p.~397)} }
\expandafter\def\csname SRef@wbII:eq:W:5.17-A34c\endcsname{\href{S6_SERIES_GRIND_II.pdf\#nameddest=tag-03FH}{II.{W.137} (p.~397)} }
\expandafter\def\csname SRef@cdp:eq:A.8\endcsname{\href{S6_SERIES_GRIND_II.pdf\#nameddest=tag-004S}{II.{A.8} (p.~50)} }
\expandafter\def\csname SRef@wbII:eq:W:5.17-A35\endcsname{\href{S6_SERIES_GRIND_II.pdf\#nameddest=tag-03FI}{II.{W.138} (p.~398)} }
\expandafter\def\csname SRef@wbII:eq:W:5.17-A36\endcsname{\href{S6_SERIES_GRIND_II.pdf\#nameddest=tag-03FJ}{II.{W.139} (p.~398)} }
\expandafter\def\csname SRef@wbII:eq:W:5.17-A37\endcsname{\href{S6_SERIES_GRIND_II.pdf\#nameddest=tag-03FK}{II.{W.140} (p.~398)} }
\expandafter\def\csname SRef@wbII:eq:W:5.17-A38\endcsname{\href{S6_SERIES_GRIND_II.pdf\#nameddest=tag-03FL}{II.{W.141} (p.~398)} }
\expandafter\def\csname SRef@wbII:eq:W:5.17f\endcsname{\href{S6_SERIES_GRIND_II.pdf\#nameddest=tag-03GP}{II.{W.44} (p.~399)} }
\expandafter\def\csname SRef@wbII:eq:W:5.18\endcsname{\href{S6_SERIES_GRIND_II.pdf\#nameddest=tag-03HI}{II.{W.45} (p.~400)} }
\expandafter\def\csname SRef@wbII:eq:W:5.19\endcsname{\href{S6_SERIES_GRIND_II.pdf\#nameddest=tag-03HJ}{II.{W.46} (p.~400)} }
\expandafter\def\csname SRef@wbII:eq:W:5.20\endcsname{\href{S6_SERIES_GRIND_II.pdf\#nameddest=tag-03HK}{II.{W.47} (p.~400)} }
\expandafter\def\csname SRef@wbII:eq:CN:5.2\endcsname{\href{S6_SERIES_GRIND_II.pdf\#nameddest=tag-03E5}{II.{W.169} (p.~401)} }
\expandafter\def\csname SRef@cdp:eq:A.9\endcsname{\href{S6_SERIES_GRIND_II.pdf\#nameddest=tag-004T}{II.{A.9} (p.~51)} }
\expandafter\def\csname SRef@wbII:eq:W:5.23\endcsname{\href{S6_SERIES_GRIND_II.pdf\#nameddest=tag-03HL}{II.{W.50} (p.~402)} }
\expandafter\def\csname SRef@wbII:eq:W:5.24\endcsname{\href{S6_SERIES_GRIND_II.pdf\#nameddest=tag-03HM}{II.{W.51} (p.~402)} }
\expandafter\def\csname SRef@wbII:eq:W:5.25\endcsname{\href{S6_SERIES_GRIND_II.pdf\#nameddest=tag-03HN}{II.{W.52} (p.~403)} }
\expandafter\def\csname SRef@wbII:eq:W:5.26\endcsname{\href{S6_SERIES_GRIND_II.pdf\#nameddest=tag-03HO}{II.{W.53} (p.~403)} }
\expandafter\def\csname SRef@wbII:eq:W:5.27\endcsname{\href{S6_SERIES_GRIND_II.pdf\#nameddest=tag-03HP}{II.{W.54} (p.~403)} }
\expandafter\def\csname SRef@wbII:eq:W:5.28\endcsname{\href{S6_SERIES_GRIND_II.pdf\#nameddest=tag-03HQ}{II.{W.55} (p.~403)} }
\expandafter\def\csname SRef@wbII:eq:W:5.29\endcsname{\href{S6_SERIES_GRIND_II.pdf\#nameddest=tag-03HR}{II.{W.56} (p.~403)} }
\expandafter\def\csname SRef@wbII:eq:W:5.30\endcsname{\href{S6_SERIES_GRIND_II.pdf\#nameddest=tag-03HT}{II.{W.57} (p.~404)} }
\expandafter\def\csname SRef@wbII:eq:W:5.31\endcsname{\href{S6_SERIES_GRIND_II.pdf\#nameddest=tag-03HU}{II.{W.58} (p.~404)} }
\expandafter\def\csname SRef@wbII:eq:W:5.32\endcsname{\href{S6_SERIES_GRIND_II.pdf\#nameddest=tag-03HV}{II.{W.59} (p.~404)} }
\expandafter\def\csname SRef@cdp:eq:A.9a\endcsname{\href{S6_SERIES_GRIND_II.pdf\#nameddest=tag-004U}{II.{A.9a} (p.~51)} }
\expandafter\def\csname SRef@wbII:eq:W:5.33\endcsname{\href{S6_SERIES_GRIND_II.pdf\#nameddest=tag-03HW}{II.{W.60} (p.~405)} }
\expandafter\def\csname SRef@wbII:eq:CN:6.1\endcsname{\href{S6_SERIES_GRIND_II.pdf\#nameddest=tag-03E6}{II.{W.176} (p.~406)} }
\expandafter\def\csname SRef@wbII:eq:CN:6.2\endcsname{\href{S6_SERIES_GRIND_II.pdf\#nameddest=tag-03E7}{II.{W.177} (p.~406)} }
\expandafter\def\csname SRef@wbII:eq:CN:6.3\endcsname{\href{S6_SERIES_GRIND_II.pdf\#nameddest=tag-03E8}{II.{W.178} (p.~406)} }
\expandafter\def\csname SRef@wbII:eq:CN:6.4\endcsname{\href{S6_SERIES_GRIND_II.pdf\#nameddest=tag-03E9}{II.{W.179} (p.~406)} }
\expandafter\def\csname SRef@cdp:eq:A.9b\endcsname{\href{S6_SERIES_GRIND_II.pdf\#nameddest=tag-004V}{II.{A.9b} (p.~51)} }
\expandafter\def\csname SRef@cdp:eq:A.NL.1\endcsname{\href{S6_SERIES_GRIND_II.pdf\#nameddest=tag-006J}{II.{A.NL.1} (p.~52)} }
\expandafter\def\csname SRef@invII:eq:E.7\endcsname{\href{S6_SERIES_GRIND_II.pdf\#nameddest=tag-01TW}{II.{E.7} (p.~795)} }
\expandafter\def\csname SRef@cdp:eq:A.NL.2\endcsname{\href{S6_SERIES_GRIND_II.pdf\#nameddest=tag-006K}{II.{A.NL.2} (p.~52)} }
\expandafter\def\csname SRef@cdp:eq:A.NL.3\endcsname{\href{S6_SERIES_GRIND_II.pdf\#nameddest=tag-006L}{II.{A.NL.3} (p.~53)} }
\expandafter\def\csname SRef@cdp:eq:A.NL.4\endcsname{\href{S6_SERIES_GRIND_II.pdf\#nameddest=tag-006M}{II.{A.NL.4} (p.~53)} }
\expandafter\def\csname SRef@cdp:eq:A.NL.5\endcsname{\href{S6_SERIES_GRIND_II.pdf\#nameddest=tag-006N}{II.{A.NL.5} (p.~54)} }
\expandafter\def\csname SRef@cdp:eq:B.1\endcsname{\href{S6_SERIES_GRIND_II.pdf\#nameddest=tag-008G}{II.{B.1} (p.~54)} }
\expandafter\def\csname SRef@cdp:eq:B.2\endcsname{\href{S6_SERIES_GRIND_II.pdf\#nameddest=tag-008H}{II.{B.2} (p.~55)} }
\expandafter\def\csname SRef@cdp:eq:B.3\endcsname{\href{S6_SERIES_GRIND_II.pdf\#nameddest=tag-008I}{II.{B.3} (p.~56)} }
\expandafter\def\csname SRef@cdp:eq:B.4\endcsname{\href{S6_SERIES_GRIND_II.pdf\#nameddest=tag-008J}{II.{B.4} (p.~56)} }
\expandafter\def\csname SRef@cdp:eq:B.5\endcsname{\href{S6_SERIES_GRIND_II.pdf\#nameddest=tag-008K}{II.{B.5} (p.~57)} }
\expandafter\def\csname SRef@cdp:eq:B.6\endcsname{\href{S6_SERIES_GRIND_II.pdf\#nameddest=tag-008L}{II.{B.6} (p.~57)} }
\expandafter\def\csname SRef@cdp:eq:B.7\endcsname{\href{S6_SERIES_GRIND_II.pdf\#nameddest=tag-008M}{II.{B.7} (p.~57)} }
\expandafter\def\csname SRef@cdp:eq:A.10\endcsname{\href{S6_SERIES_GRIND_II.pdf\#nameddest=tag-003Z}{II.{A.10} (p.~58)} }
\expandafter\def\csname SRef@cdp:eq:A.11\endcsname{\href{S6_SERIES_GRIND_II.pdf\#nameddest=tag-0040}{II.{A.11} (p.~58)} }
\expandafter\def\csname SRef@cdp:eq:A.12\endcsname{\href{S6_SERIES_GRIND_II.pdf\#nameddest=tag-0041}{II.{A.12} (p.~59)} }
\expandafter\def\csname SRef@cdp:eq:A.13\endcsname{\href{S6_SERIES_GRIND_II.pdf\#nameddest=tag-0042}{II.{A.13} (p.~60)} }
\expandafter\def\csname SRef@cdp:eq:A.14\endcsname{\href{S6_SERIES_GRIND_II.pdf\#nameddest=tag-0043}{II.{A.14} (p.~60)} }
\expandafter\def\csname SRef@cdp:eq:A.15a\endcsname{\href{S6_SERIES_GRIND_II.pdf\#nameddest=tag-0045}{II.{A.15a} (p.~60)} }
\expandafter\def\csname SRef@cdp:eq:A.15b\endcsname{\href{S6_SERIES_GRIND_II.pdf\#nameddest=tag-0046}{II.{A.15b} (p.~62)} }
\expandafter\def\csname SRef@cdp:eq:A.15\endcsname{\href{S6_SERIES_GRIND_II.pdf\#nameddest=tag-0044}{II.{A.15} (p.~63)} }
\expandafter\def\csname SRef@cdp:eq:A.16\endcsname{\href{S6_SERIES_GRIND_II.pdf\#nameddest=tag-0047}{II.{A.16} (p.~65)} }
\expandafter\def\csname SRef@cdp:eq:A.17\endcsname{\href{S6_SERIES_GRIND_II.pdf\#nameddest=tag-0048}{II.{A.17} (p.~66)} }
\expandafter\def\csname SRef@cdp:eq:A.18\endcsname{\href{S6_SERIES_GRIND_II.pdf\#nameddest=tag-0049}{II.{A.18} (p.~67)} }
\expandafter\def\csname SRef@cdp:eq:A.19\endcsname{\href{S6_SERIES_GRIND_II.pdf\#nameddest=tag-004A}{II.{A.19} (p.~67)} }
\expandafter\def\csname SRef@cdp:eq:A.20\endcsname{\href{S6_SERIES_GRIND_II.pdf\#nameddest=tag-004C}{II.{A.20} (p.~68)} }
\expandafter\def\csname SRef@cdp:eq:A.21\endcsname{\href{S6_SERIES_GRIND_II.pdf\#nameddest=tag-004D}{II.{A.21} (p.~69)} }
\expandafter\def\csname SRef@cdp:eq:A.22\endcsname{\href{S6_SERIES_GRIND_II.pdf\#nameddest=tag-004E}{II.{A.22} (p.~69)} }
\expandafter\def\csname SRef@cdp:eq:A.23\endcsname{\href{S6_SERIES_GRIND_II.pdf\#nameddest=tag-004F}{II.{A.23} (p.~70)} }
\expandafter\def\csname SRef@cdp:eq:A.24\endcsname{\href{S6_SERIES_GRIND_II.pdf\#nameddest=tag-004G}{II.{A.24} (p.~71)} }
\expandafter\def\csname SRef@cdp:eq:A.25\endcsname{\href{S6_SERIES_GRIND_II.pdf\#nameddest=tag-004H}{II.{A.25} (p.~74)} }
\expandafter\def\csname SRef@cdp:eq:A.26\endcsname{\href{S6_SERIES_GRIND_II.pdf\#nameddest=tag-004I}{II.{A.26} (p.~74)} }
\expandafter\def\csname SRef@cdp:eq:A.27\endcsname{\href{S6_SERIES_GRIND_II.pdf\#nameddest=tag-004J}{II.{A.27} (p.~76)} }
\expandafter\def\csname SRef@cdp:eq:A.28\endcsname{\href{S6_SERIES_GRIND_II.pdf\#nameddest=tag-004K}{II.{A.28} (p.~76)} }
\expandafter\def\csname SRef@cdp:eq:A.29\endcsname{\href{S6_SERIES_GRIND_II.pdf\#nameddest=tag-004L}{II.{A.29} (p.~76)} }
\expandafter\def\csname SRef@cdp:eq:A.30\endcsname{\href{S6_SERIES_GRIND_II.pdf\#nameddest=tag-004N}{II.{A.30} (p.~77)} }
\expandafter\def\csname SRef@cdp:eq:A.RC.1\endcsname{\href{S6_SERIES_GRIND_II.pdf\#nameddest=tag-006U}{II.{A.RC.1} (p.~78)} }
\expandafter\def\csname SRef@push:src:E.2.4\endcsname{\href{S6_SERIES_GRIND_II.pdf\#nameddest=tag-033C}{II.13.1 (p.~309)} }
\expandafter\def\csname SRef@push:src:E.2.24\endcsname{\href{S6_SERIES_GRIND_II.pdf\#nameddest=tag-033A}{II.13.21 (p.~348)} }
\expandafter\def\csname SRef@wbII:cor:surface-top-base-change\endcsname{\href{S6_SERIES_GRIND_II.pdf\#nameddest=tag-03DV}{II.16.10 (p.~374)} }
\expandafter\def\csname SRef@wbII:cor:candidate-pic0-route\endcsname{\href{S6_SERIES_GRIND_II.pdf\#nameddest=tag-03DS}{II.16.22 (p.~399)} }
\expandafter\def\csname SRef@wbII:cor:candidate-conductor-actual\endcsname{\href{S6_SERIES_GRIND_II.pdf\#nameddest=tag-03DR}{II.16.24 (p.~401)} }
\expandafter\def\csname SRef@wbII:cor:corrected-cdp-relative\endcsname{\href{S6_SERIES_GRIND_II.pdf\#nameddest=tag-03DT}{II.16.28 (p.~403)} }
\expandafter\def\csname SRef@wbII:cor:relative\endcsname{\href{S6_SERIES_GRIND_II.pdf\#nameddest=tag-03DT}{II.16.28 (p.~403)} }
\expandafter\def\csname SRef@wbII:cor:normal-crossing-kernel\endcsname{\href{S6_SERIES_GRIND_II.pdf\#nameddest=tag-03DU}{II.16.3 (p.~367)} }
\expandafter\def\csname SRef@wbII:test:cdp-comparison\endcsname{\href{S6_SERIES_GRIND_II.pdf\#nameddest=tag-03IV}{II.16.30 (p.~405)} }
\expandafter\def\csname SRef@cdp:A:filtration1\endcsname{\href{S6_SERIES_GRIND_II.pdf\#nameddest=tag-003W}{II.1 (p.~48)} }
\expandafter\def\csname SRef@cdp:eq:E13-residue\endcsname{\href{S6_SERIES_GRIND_II.pdf\#nameddest=tag-01A6}{II.10 (p.~296)} }
\expandafter\def\csname SRef@cdp:eq:E14-elliptic-family\endcsname{\href{S6_SERIES_GRIND_II.pdf\#nameddest=tag-01A8}{II.11 (p.~297)} }
\expandafter\def\csname SRef@cdp:eq:E14-canonical\endcsname{\href{S6_SERIES_GRIND_II.pdf\#nameddest=tag-01A7}{II.12 (p.~298)} }
\expandafter\def\csname SRef@cdp:eq:E14-log-coho\endcsname{\href{S6_SERIES_GRIND_II.pdf\#nameddest=tag-01AA}{II.13 (p.~299)} }
\expandafter\def\csname SRef@cdp:eq:E14-jet-map\endcsname{\href{S6_SERIES_GRIND_II.pdf\#nameddest=tag-01A9}{II.14 (p.~299)} }
\expandafter\def\csname SRef@cdp:eq:E15-positive\endcsname{\href{S6_SERIES_GRIND_II.pdf\#nameddest=tag-01AB}{II.15 (p.~301)} }
\expandafter\def\csname SRef@cdp:eq:E15-products\endcsname{\href{S6_SERIES_GRIND_II.pdf\#nameddest=tag-01AC}{II.16 (p.~301)} }
\expandafter\def\csname SRef@wbII:eq:analytic-toric-open-directions\endcsname{\href{S6_SERIES_GRIND_II.pdf\#nameddest=tag-03I5}{II.17 (p.~396)} }
\expandafter\def\csname SRef@wbII:eq:analytic-residual-descent-factor\endcsname{\href{S6_SERIES_GRIND_II.pdf\#nameddest=tag-03I4}{II.18 (p.~397)} }
\expandafter\def\csname SRef@wbII:eq:analytic-factor-horizontal-winding\endcsname{\href{S6_SERIES_GRIND_II.pdf\#nameddest=tag-03I3}{II.19 (p.~397)} }
\expandafter\def\csname SRef@cdp:A:filtration2\endcsname{\href{S6_SERIES_GRIND_II.pdf\#nameddest=tag-003X}{II.2 (p.~48)} }
\expandafter\def\csname SRef@eq:P4chi-resolution\endcsname{\href{S6_SERIES_GRIND_II.pdf\#nameddest=tag-01OD}{II.3 (p.~141)} }
\expandafter\def\csname SRef@eq:P4exceptional-genus\endcsname{\href{S6_SERIES_GRIND_II.pdf\#nameddest=tag-01OE}{II.4 (p.~141)} }
\expandafter\def\csname SRef@cdp:eq:E-Leray-map\endcsname{\href{S6_SERIES_GRIND_II.pdf\#nameddest=tag-019Y}{II.5 (p.~285)} }
\expandafter\def\csname SRef@cdp:eq:E-surj-embedding\endcsname{\href{S6_SERIES_GRIND_II.pdf\#nameddest=tag-01A0}{II.6 (p.~286)} }
\expandafter\def\csname SRef@cdp:eq:E-log-period\endcsname{\href{S6_SERIES_GRIND_II.pdf\#nameddest=tag-019Z}{II.7 (p.~288)} }
\expandafter\def\csname SRef@cdp:eq:E-theta-general\endcsname{\href{S6_SERIES_GRIND_II.pdf\#nameddest=tag-01A1}{II.8 (p.~289)} }
\expandafter\def\csname SRef@cdp:eq:E13-connection\endcsname{\href{S6_SERIES_GRIND_II.pdf\#nameddest=tag-01A5}{II.9 (p.~294)} }
\expandafter\def\csname SRef@wbII:lem:torsion-identification\endcsname{\href{S6_SERIES_GRIND_II.pdf\#nameddest=tag-03IF}{II.16.1 (p.~363)} }
\expandafter\def\csname SRef@wbII:lem:picard-crossed-defect\endcsname{\href{S6_SERIES_GRIND_II.pdf\#nameddest=tag-03IC}{II.16.14 (p.~377)} }
\expandafter\def\csname SRef@wbII:lem:picard-finite-norm\endcsname{\href{S6_SERIES_GRIND_II.pdf\#nameddest=tag-03IE}{II.16.15 (p.~378)} }
\expandafter\def\csname SRef@wbII:lem:picard-cusp-defect\endcsname{\href{S6_SERIES_GRIND_II.pdf\#nameddest=tag-03ID}{II.16.16 (p.~380)} }
\expandafter\def\csname SRef@wbII:lem:analytic-support-specialization\endcsname{\href{S6_SERIES_GRIND_II.pdf\#nameddest=tag-03I6}{II.16.20 (p.~394)} }
\expandafter\def\csname SRef@wbII:lem:generic-endpoints\endcsname{\href{S6_SERIES_GRIND_II.pdf\#nameddest=tag-03IA}{II.16.25 (p.~401)} }
\expandafter\def\csname SRef@wbII:lem:cdp-generic-endpoints\endcsname{\href{S6_SERIES_GRIND_II.pdf\#nameddest=tag-03I7}{II.16.26 (p.~402)} }
\expandafter\def\csname SRef@wbII:lem:gluing-bundle\endcsname{\href{S6_SERIES_GRIND_II.pdf\#nameddest=tag-03IB}{II.16.4 (p.~369)} }
\expandafter\def\csname SRef@wbII:lem:extending-gluing\endcsname{\href{S6_SERIES_GRIND_II.pdf\#nameddest=tag-03I9}{II.16.5 (p.~370)} }
\expandafter\def\csname SRef@wbII:lem:curve-flatness\endcsname{\href{S6_SERIES_GRIND_II.pdf\#nameddest=tag-03I8}{II.16.8 (p.~373)} }
\expandafter\def\csname SRef@inv:E2.3\endcsname{\href{S6_SERIES_GRIND_II.pdf\#nameddest=tag-01SW}{II.26.1 (p.~795)} }
\expandafter\def\csname SRef@cdp:st:D.2.1\endcsname{\href{S6_SERIES_GRIND_II.pdf\#nameddest=tag-01G8}{II.11.1 (p.~272)} }
\expandafter\def\csname SRef@cdp:st:D.2.3\endcsname{\href{S6_SERIES_GRIND_II.pdf\#nameddest=tag-01GA}{II.11.2 (p.~272)} }
\expandafter\def\csname SRef@cdp:st:D.2.4\endcsname{\href{S6_SERIES_GRIND_II.pdf\#nameddest=tag-01GB}{II.11.3 (p.~275)} }
\expandafter\def\csname SRef@cdp:st:D.2.5\endcsname{\href{S6_SERIES_GRIND_II.pdf\#nameddest=tag-01GC}{II.11.4 (p.~276)} }
\expandafter\def\csname SRef@cdp:st:D.2.6\endcsname{\href{S6_SERIES_GRIND_II.pdf\#nameddest=tag-01GD}{II.11.5 (p.~278)} }
\expandafter\def\csname SRef@cdp:st:D.2.7\endcsname{\href{S6_SERIES_GRIND_II.pdf\#nameddest=tag-01GE}{II.11.6 (p.~279)} }
\expandafter\def\csname SRef@cdp:st:D.2.8\endcsname{\href{S6_SERIES_GRIND_II.pdf\#nameddest=tag-01GF}{II.11.7 (p.~283)} }
\expandafter\def\csname SRef@cdp:st:D.2.2\endcsname{\href{S6_SERIES_GRIND_II.pdf\#nameddest=tag-01G9}{II.11.8 (p.~284)} }
\expandafter\def\csname SRef@cdp:st:E5\endcsname{\href{S6_SERIES_GRIND_II.pdf\#nameddest=tag-01GU}{II.12.1 (p.~285)} }
\expandafter\def\csname SRef@cdp:st:E15\endcsname{\href{S6_SERIES_GRIND_II.pdf\#nameddest=tag-01GM}{II.12.10 (p.~300)} }
\expandafter\def\csname SRef@cdp:st:E16\endcsname{\href{S6_SERIES_GRIND_II.pdf\#nameddest=tag-01GN}{II.12.11 (p.~304)} }
\expandafter\def\csname SRef@cdp:st:E17\endcsname{\href{S6_SERIES_GRIND_II.pdf\#nameddest=tag-01GO}{II.12.12 (p.~305)} }
\expandafter\def\csname SRef@cdp:st:E18\endcsname{\href{S6_SERIES_GRIND_II.pdf\#nameddest=tag-01GP}{II.12.13 (p.~307)} }
\expandafter\def\csname SRef@cdp:st:E6\endcsname{\href{S6_SERIES_GRIND_II.pdf\#nameddest=tag-01GV}{II.12.2 (p.~286)} }
\expandafter\def\csname SRef@cdp:st:E7\endcsname{\href{S6_SERIES_GRIND_II.pdf\#nameddest=tag-01GW}{II.12.3 (p.~288)} }
\expandafter\def\csname SRef@cdp:st:E8\endcsname{\href{S6_SERIES_GRIND_II.pdf\#nameddest=tag-01GX}{II.12.4 (p.~289)} }
\expandafter\def\csname SRef@cdp:st:E10\endcsname{\href{S6_SERIES_GRIND_II.pdf\#nameddest=tag-01GH}{II.12.5 (p.~290)} }
\expandafter\def\csname SRef@cdp:st:E11\endcsname{\href{S6_SERIES_GRIND_II.pdf\#nameddest=tag-01GI}{II.12.6 (p.~292)} }
\expandafter\def\csname SRef@cdp:st:E12\endcsname{\href{S6_SERIES_GRIND_II.pdf\#nameddest=tag-01GJ}{II.12.7 (p.~293)} }
\expandafter\def\csname SRef@cdp:st:E13\endcsname{\href{S6_SERIES_GRIND_II.pdf\#nameddest=tag-01GK}{II.12.8 (p.~294)} }
\expandafter\def\csname SRef@cdp:st:E14\endcsname{\href{S6_SERIES_GRIND_II.pdf\#nameddest=tag-01GL}{II.12.9 (p.~297)} }
\expandafter\def\csname SRef@push:src:E.2.13\endcsname{\href{S6_SERIES_GRIND_II.pdf\#nameddest=tag-032Y}{II.13.10 (p.~325)} }
\expandafter\def\csname SRef@push:src:E.2.14\endcsname{\href{S6_SERIES_GRIND_II.pdf\#nameddest=tag-032Z}{II.13.11 (p.~326)} }
\expandafter\def\csname SRef@push:src:E.2.15\endcsname{\href{S6_SERIES_GRIND_II.pdf\#nameddest=tag-0330}{II.13.12 (p.~328)} }
\expandafter\def\csname SRef@push:src:E.2.16\endcsname{\href{S6_SERIES_GRIND_II.pdf\#nameddest=tag-0331}{II.13.13 (p.~330)} }
\expandafter\def\csname SRef@push:src:E.2.17\endcsname{\href{S6_SERIES_GRIND_II.pdf\#nameddest=tag-0332}{II.13.14 (p.~332)} }
\expandafter\def\csname SRef@push:src:E.2.18\endcsname{\href{S6_SERIES_GRIND_II.pdf\#nameddest=tag-0333}{II.13.15 (p.~334)} }
\expandafter\def\csname SRef@push:src:E.2.19\endcsname{\href{S6_SERIES_GRIND_II.pdf\#nameddest=tag-0334}{II.13.16 (p.~335)} }
\expandafter\def\csname SRef@push:src:E.2.20\endcsname{\href{S6_SERIES_GRIND_II.pdf\#nameddest=tag-0336}{II.13.17 (p.~338)} }
\expandafter\def\csname SRef@push:src:E.2.21\endcsname{\href{S6_SERIES_GRIND_II.pdf\#nameddest=tag-0337}{II.13.18 (p.~341)} }
\expandafter\def\csname SRef@push:src:E.2.22\endcsname{\href{S6_SERIES_GRIND_II.pdf\#nameddest=tag-0338}{II.13.19 (p.~343)} }
\expandafter\def\csname SRef@push:src:E.2.5\endcsname{\href{S6_SERIES_GRIND_II.pdf\#nameddest=tag-033D}{II.13.2 (p.~310)} }
\expandafter\def\csname SRef@push:src:E.5.1\endcsname{\href{S6_SERIES_GRIND_II.pdf\#nameddest=tag-033N}{II.13.22 (p.~349)} }
\expandafter\def\csname SRef@push:src:E.5.2\endcsname{\href{S6_SERIES_GRIND_II.pdf\#nameddest=tag-033O}{II.13.23 (p.~349)} }
\expandafter\def\csname SRef@push:src:E.2.6\endcsname{\href{S6_SERIES_GRIND_II.pdf\#nameddest=tag-033E}{II.13.3 (p.~311)} }
\expandafter\def\csname SRef@push:src:E.2.7\endcsname{\href{S6_SERIES_GRIND_II.pdf\#nameddest=tag-033F}{II.13.4 (p.~313)} }
\expandafter\def\csname SRef@push:src:E.2.8\endcsname{\href{S6_SERIES_GRIND_II.pdf\#nameddest=tag-033G}{II.13.5 (p.~316)} }
\expandafter\def\csname SRef@push:src:E.2.9\endcsname{\href{S6_SERIES_GRIND_II.pdf\#nameddest=tag-033H}{II.13.6 (p.~318)} }
\expandafter\def\csname SRef@push:src:E.2.10\endcsname{\href{S6_SERIES_GRIND_II.pdf\#nameddest=tag-032V}{II.13.7 (p.~320)} }
\expandafter\def\csname SRef@push:src:E.2.11\endcsname{\href{S6_SERIES_GRIND_II.pdf\#nameddest=tag-032W}{II.13.8 (p.~321)} }
\expandafter\def\csname SRef@push:src:E.2.12\endcsname{\href{S6_SERIES_GRIND_II.pdf\#nameddest=tag-032X}{II.13.9 (p.~323)} }
\expandafter\def\csname SRef@inv:E2.2\endcsname{\href{S6_SERIES_GRIND_II.pdf\#nameddest=tag-01SV}{II.14.1 (p.~354)} }
\expandafter\def\csname SRef@inv:E2.4\endcsname{\href{S6_SERIES_GRIND_II.pdf\#nameddest=tag-01SX}{II.14.3 (p.~361)} }
\expandafter\def\csname SRef@wbII:prop:pic0-pullback\endcsname{\href{S6_SERIES_GRIND_II.pdf\#nameddest=tag-03IO}{II.16.11 (p.~376)} }
\expandafter\def\csname SRef@wbII:prop:all-pic-restriction\endcsname{\href{S6_SERIES_GRIND_II.pdf\#nameddest=tag-03IG}{II.16.12 (p.~376)} }
\expandafter\def\csname SRef@wbII:prop:leray-edge-restriction\endcsname{\href{S6_SERIES_GRIND_II.pdf\#nameddest=tag-03IM}{II.16.13 (p.~376)} }
\expandafter\def\csname SRef@wbII:prop:direct-boundary-picard\endcsname{\href{S6_SERIES_GRIND_II.pdf\#nameddest=tag-03IL}{II.16.17 (p.~381)} }
\expandafter\def\csname SRef@wbII:prop:local-cusp-picard-obstruction\endcsname{\href{S6_SERIES_GRIND_II.pdf\#nameddest=tag-03IN}{II.16.18 (p.~383)} }
\expandafter\def\csname SRef@wbII:prop:analytic-ah-boundary\endcsname{\href{S6_SERIES_GRIND_II.pdf\#nameddest=tag-03IH}{II.16.19 (p.~387)} }
\expandafter\def\csname SRef@wbII:prop:relative-local\endcsname{\href{S6_SERIES_GRIND_II.pdf\#nameddest=tag-03IP}{II.16.31 (p.~406)} }
\expandafter\def\csname SRef@wbII:prop:cartier-duality-absolute\endcsname{\href{S6_SERIES_GRIND_II.pdf\#nameddest=tag-03IJ}{II.16.7 (p.~372)} }
\expandafter\def\csname SRef@wbII:prop:analytic-base-change\endcsname{\href{S6_SERIES_GRIND_II.pdf\#nameddest=tag-03II}{II.16.9 (p.~374)} }
\expandafter\def\csname SRef@cdp:st:C.1.1\endcsname{\href{S6_SERIES_GRIND_II.pdf\#nameddest=tag-01BN}{II.3.1 (p.~18)} }
\expandafter\def\csname SRef@cdp:st:C.1.3\endcsname{\href{S6_SERIES_GRIND_II.pdf\#nameddest=tag-01BP}{II.3.2 (p.~19)} }
\expandafter\def\csname SRef@cdp:sec:E\endcsname{\href{S6_SERIES_GRIND_II.pdf\#nameddest=tag-01BP}{II.3 (p.~19)} }
\expandafter\def\csname SRef@cdp:st:E1\endcsname{\href{S6_SERIES_GRIND_II.pdf\#nameddest=tag-01GG}{II.3.3 (p.~19)} }
\expandafter\def\csname SRef@cdp:st:D.1.2\endcsname{\href{S6_SERIES_GRIND_II.pdf\#nameddest=tag-01G0}{II.4.1 (p.~20)} }
\expandafter\def\csname SRef@cdp:st:D.1.4\endcsname{\href{S6_SERIES_GRIND_II.pdf\#nameddest=tag-01G2}{II.4.10 (p.~28)} }
\expandafter\def\csname SRef@cdp:st:D.1.13\endcsname{\href{S6_SERIES_GRIND_II.pdf\#nameddest=tag-01FY}{II.4.11 (p.~28)} }
\expandafter\def\csname SRef@cdp:st:D.1.15\endcsname{\href{S6_SERIES_GRIND_II.pdf\#nameddest=tag-01FZ}{II.4.12 (p.~30)} }
\expandafter\def\csname SRef@cdp:st:D.1.5\endcsname{\href{S6_SERIES_GRIND_II.pdf\#nameddest=tag-01G3}{II.4.2 (p.~21)} }
\expandafter\def\csname SRef@cdp:st:D.1.6\endcsname{\href{S6_SERIES_GRIND_II.pdf\#nameddest=tag-01G4}{II.4.3 (p.~22)} }
\expandafter\def\csname SRef@cdp:st:D.1.7\endcsname{\href{S6_SERIES_GRIND_II.pdf\#nameddest=tag-01G5}{II.4.4 (p.~23)} }
\expandafter\def\csname SRef@cdp:st:D.1.8\endcsname{\href{S6_SERIES_GRIND_II.pdf\#nameddest=tag-01G6}{II.4.5 (p.~23)} }
\expandafter\def\csname SRef@cdp:st:D.1.9\endcsname{\href{S6_SERIES_GRIND_II.pdf\#nameddest=tag-01G7}{II.4.6 (p.~24)} }
\expandafter\def\csname SRef@cdp:st:D.1.10\endcsname{\href{S6_SERIES_GRIND_II.pdf\#nameddest=tag-01FV}{II.4.7 (p.~25)} }
\expandafter\def\csname SRef@cdp:st:D.1.11\endcsname{\href{S6_SERIES_GRIND_II.pdf\#nameddest=tag-01FW}{II.4.8 (p.~25)} }
\expandafter\def\csname SRef@cdp:st:D.1.12\endcsname{\href{S6_SERIES_GRIND_II.pdf\#nameddest=tag-01FX}{II.4.9 (p.~27)} }
\expandafter\def\csname SRef@cdp:st:C.3.2\endcsname{\href{S6_SERIES_GRIND_II.pdf\#nameddest=tag-01F8}{II.5.1 (p.~31)} }
\expandafter\def\csname SRef@cdp:st:C.3.11\endcsname{\href{S6_SERIES_GRIND_II.pdf\#nameddest=tag-01F6}{II.5.10 (p.~44)} }
\expandafter\def\csname SRef@cdp:st:C.3.5\endcsname{\href{S6_SERIES_GRIND_II.pdf\#nameddest=tag-01FB}{II.5.11 (p.~45)} }
\expandafter\def\csname SRef@cdp:st:C.3.6\endcsname{\href{S6_SERIES_GRIND_II.pdf\#nameddest=tag-01FC}{II.5.12 (p.~46)} }
\expandafter\def\csname SRef@cdp:st:C.3.3\endcsname{\href{S6_SERIES_GRIND_II.pdf\#nameddest=tag-01F9}{II.5.2 (p.~32)} }
\expandafter\def\csname SRef@cdp:st:C.3.4\endcsname{\href{S6_SERIES_GRIND_II.pdf\#nameddest=tag-01FA}{II.5.3 (p.~33)} }
\expandafter\def\csname SRef@cdp:st:C.3.7\endcsname{\href{S6_SERIES_GRIND_II.pdf\#nameddest=tag-01FD}{II.5.4 (p.~33)} }
\expandafter\def\csname SRef@cdp:st:C.3.8\endcsname{\href{S6_SERIES_GRIND_II.pdf\#nameddest=tag-01FE}{II.5.5 (p.~34)} }
\expandafter\def\csname SRef@cdp:st:C.3.9\endcsname{\href{S6_SERIES_GRIND_II.pdf\#nameddest=tag-01FF}{II.5.6 (p.~35)} }
\expandafter\def\csname SRef@cdp:st:C.3.10\endcsname{\href{S6_SERIES_GRIND_II.pdf\#nameddest=tag-01F5}{II.5.7 (p.~38)} }
\expandafter\def\csname SRef@cdp:st:C.3.12\endcsname{\href{S6_SERIES_GRIND_II.pdf\#nameddest=tag-01F7}{II.5.8 (p.~41)} }
\expandafter\def\csname SRef@cdp:st:C.3.1\endcsname{\href{S6_SERIES_GRIND_II.pdf\#nameddest=tag-01F4}{II.5.9 (p.~43)} }
\expandafter\def\csname SRef@cdp:st:A.1.1\endcsname{\href{S6_SERIES_GRIND_II.pdf\#nameddest=tag-01AD}{II.6.1 (p.~48)} }
\expandafter\def\csname SRef@wbII:lem:conormal\endcsname{\href{S6_SERIES_GRIND_II.pdf\#nameddest=tag-01AD}{II.6.1 (p.~48)} }
\expandafter\def\csname SRef@wbII:prop:criterion\endcsname{\href{S6_SERIES_GRIND_II.pdf\#nameddest=tag-01AD}{II.6.1 (p.~48)} }
\expandafter\def\csname SRef@wbII:prop:filtration\endcsname{\href{S6_SERIES_GRIND_II.pdf\#nameddest=tag-01AD}{II.6.1 (p.~48)} }
\expandafter\def\csname SRef@wbII:lem:conormal-injective\endcsname{\href{S6_SERIES_GRIND_II.pdf\#nameddest=tag-01AD}{II.6.1 (p.~48)} }
\expandafter\def\csname SRef@wbII:prop:missing-cotangent-filtration\endcsname{\href{S6_SERIES_GRIND_II.pdf\#nameddest=tag-01AD}{II.6.1 (p.~48)} }
\expandafter\def\csname SRef@wbII:prop:missing-global-section-criterion\endcsname{\href{S6_SERIES_GRIND_II.pdf\#nameddest=tag-01AD}{II.6.1 (p.~48)} }
\expandafter\def\csname SRef@cdp:st:B3\endcsname{\href{S6_SERIES_GRIND_II.pdf\#nameddest=tag-01BM}{II.6.10 (p.~58)} }
\expandafter\def\csname SRef@cdp:st:A.3.1\endcsname{\href{S6_SERIES_GRIND_II.pdf\#nameddest=tag-01AL}{II.6.11 (p.~58)} }
\expandafter\def\csname SRef@cdp:st:A.3.2\endcsname{\href{S6_SERIES_GRIND_II.pdf\#nameddest=tag-01AM}{II.6.12 (p.~59)} }
\expandafter\def\csname SRef@cdp:st:A.3.3\endcsname{\href{S6_SERIES_GRIND_II.pdf\#nameddest=tag-01AN}{II.6.13 (p.~60)} }
\expandafter\def\csname SRef@cdp:st:A.3.4\endcsname{\href{S6_SERIES_GRIND_II.pdf\#nameddest=tag-01AO}{II.6.14 (p.~60)} }
\expandafter\def\csname SRef@cdp:st:A.3.5\endcsname{\href{S6_SERIES_GRIND_II.pdf\#nameddest=tag-01AP}{II.6.15 (p.~62)} }
\expandafter\def\csname SRef@cdp:st:A.4.1\endcsname{\href{S6_SERIES_GRIND_II.pdf\#nameddest=tag-01AQ}{II.6.16 (p.~63)} }
\expandafter\def\csname SRef@cdp:st:A.4.2\endcsname{\href{S6_SERIES_GRIND_II.pdf\#nameddest=tag-01AR}{II.6.17 (p.~64)} }
\expandafter\def\csname SRef@cdp:st:A.4.3\endcsname{\href{S6_SERIES_GRIND_II.pdf\#nameddest=tag-01AS}{II.6.18 (p.~64)} }
\expandafter\def\csname SRef@cdp:st:A.4.4\endcsname{\href{S6_SERIES_GRIND_II.pdf\#nameddest=tag-01AT}{II.6.19 (p.~65)} }
\expandafter\def\csname SRef@cdp:st:A.2.1\endcsname{\href{S6_SERIES_GRIND_II.pdf\#nameddest=tag-01AF}{II.6.2 (p.~48)} }
\expandafter\def\csname SRef@cdp:st:A.4.5\endcsname{\href{S6_SERIES_GRIND_II.pdf\#nameddest=tag-01AU}{II.6.20 (p.~66)} }
\expandafter\def\csname SRef@cdp:st:A.4.6\endcsname{\href{S6_SERIES_GRIND_II.pdf\#nameddest=tag-01AV}{II.6.21 (p.~67)} }
\expandafter\def\csname SRef@cdp:st:A.4.7\endcsname{\href{S6_SERIES_GRIND_II.pdf\#nameddest=tag-01AW}{II.6.22 (p.~68)} }
\expandafter\def\csname SRef@cdp:st:A.4.8\endcsname{\href{S6_SERIES_GRIND_II.pdf\#nameddest=tag-01AX}{II.6.23 (p.~69)} }
\expandafter\def\csname SRef@cdp:st:A.4.9\endcsname{\href{S6_SERIES_GRIND_II.pdf\#nameddest=tag-01AY}{II.6.24 (p.~70)} }
\expandafter\def\csname SRef@cdp:st:A.2.2\endcsname{\href{S6_SERIES_GRIND_II.pdf\#nameddest=tag-01AG}{II.6.3 (p.~49)} }
\expandafter\def\csname SRef@wbII:lem:conductor\endcsname{\href{S6_SERIES_GRIND_II.pdf\#nameddest=tag-01AG}{II.6.3 (p.~49)} }
\expandafter\def\csname SRef@wbII:lem:conductor-square\endcsname{\href{S6_SERIES_GRIND_II.pdf\#nameddest=tag-01AG}{II.6.3 (p.~49)} }
\expandafter\def\csname SRef@cdp:st:A.1.2\endcsname{\href{S6_SERIES_GRIND_II.pdf\#nameddest=tag-01AE}{II.6.4 (p.~50)} }
\expandafter\def\csname SRef@cdp:st:A.2.4\endcsname{\href{S6_SERIES_GRIND_II.pdf\#nameddest=tag-01AI}{II.6.5 (p.~50)} }
\expandafter\def\csname SRef@cdp:st:A.2.5\endcsname{\href{S6_SERIES_GRIND_II.pdf\#nameddest=tag-01AJ}{II.6.6 (p.~51)} }
\expandafter\def\csname SRef@cdp:st:A.2.6\endcsname{\href{S6_SERIES_GRIND_II.pdf\#nameddest=tag-01AK}{II.6.7 (p.~52)} }
\expandafter\def\csname SRef@cdp:st:B1\endcsname{\href{S6_SERIES_GRIND_II.pdf\#nameddest=tag-01BK}{II.6.8 (p.~54)} }
\expandafter\def\csname SRef@wbII:thm:nodal-test\endcsname{\href{S6_SERIES_GRIND_II.pdf\#nameddest=tag-01BK}{II.6.8 (p.~54)} }
\expandafter\def\csname SRef@wbII:prop:family-geometry\endcsname{\href{S6_SERIES_GRIND_II.pdf\#nameddest=tag-01BK}{II.6.8 (p.~54)} }
\expandafter\def\csname SRef@wbII:prop:nodal-counterexample\endcsname{\href{S6_SERIES_GRIND_II.pdf\#nameddest=tag-01BK}{II.6.8 (p.~54)} }
\expandafter\def\csname SRef@cdp:st:B2\endcsname{\href{S6_SERIES_GRIND_II.pdf\#nameddest=tag-01BL}{II.6.9 (p.~56)} }
\expandafter\def\csname SRef@cdp:st:A.5.1\endcsname{\href{S6_SERIES_GRIND_II.pdf\#nameddest=tag-01AZ}{II.7.1 (p.~72)} }
\expandafter\def\csname SRef@cdp:st:A.5.5\endcsname{\href{S6_SERIES_GRIND_II.pdf\#nameddest=tag-01BF}{II.7.10 (p.~84)} }
\expandafter\def\csname SRef@cdp:st:A.5.7\endcsname{\href{S6_SERIES_GRIND_II.pdf\#nameddest=tag-01BH}{II.7.11 (p.~88)} }
\expandafter\def\csname SRef@cdp:st:A.5.8\endcsname{\href{S6_SERIES_GRIND_II.pdf\#nameddest=tag-01BI}{II.7.12 (p.~89)} }
\expandafter\def\csname SRef@cdp:st:A.5.14\endcsname{\href{S6_SERIES_GRIND_II.pdf\#nameddest=tag-01B4}{II.7.13 (p.~90)} }
\expandafter\def\csname SRef@cdp:st:A.5.10\endcsname{\href{S6_SERIES_GRIND_II.pdf\#nameddest=tag-01B0}{II.7.14 (p.~92)} }
\expandafter\def\csname SRef@cdp:st:A.5.20\endcsname{\href{S6_SERIES_GRIND_II.pdf\#nameddest=tag-01BB}{II.7.15 (p.~94)} }
\expandafter\def\csname SRef@cdp:st:A.5.15\endcsname{\href{S6_SERIES_GRIND_II.pdf\#nameddest=tag-01B5}{II.7.16 (p.~95)} }
\expandafter\def\csname SRef@cdp:st:A.5.17\endcsname{\href{S6_SERIES_GRIND_II.pdf\#nameddest=tag-01B7}{II.7.17 (p.~98)} }
\expandafter\def\csname SRef@cdp:st:A.5.18\endcsname{\href{S6_SERIES_GRIND_II.pdf\#nameddest=tag-01B8}{II.7.18 (p.~101)} }
\expandafter\def\csname SRef@cdp:st:A.5.19\endcsname{\href{S6_SERIES_GRIND_II.pdf\#nameddest=tag-01B9}{II.7.19 (p.~108)} }
\expandafter\def\csname SRef@cdp:st:A.5.2\endcsname{\href{S6_SERIES_GRIND_II.pdf\#nameddest=tag-01BA}{II.7.2 (p.~73)} }
\expandafter\def\csname SRef@cdp:st:A.5.16\endcsname{\href{S6_SERIES_GRIND_II.pdf\#nameddest=tag-01B6}{II.7.20 (p.~111)} }
\expandafter\def\csname SRef@cdp:st:A.5.21\endcsname{\href{S6_SERIES_GRIND_II.pdf\#nameddest=tag-01BC}{II.7.21 (p.~117)} }
\expandafter\def\csname SRef@cdp:st:A.5.3\endcsname{\href{S6_SERIES_GRIND_II.pdf\#nameddest=tag-01BD}{II.7.3 (p.~74)} }
\expandafter\def\csname SRef@cdp:st:A.5.4\endcsname{\href{S6_SERIES_GRIND_II.pdf\#nameddest=tag-01BE}{II.7.4 (p.~76)} }
\expandafter\def\csname SRef@cdp:st:A.5.6\endcsname{\href{S6_SERIES_GRIND_II.pdf\#nameddest=tag-01BG}{II.7.5 (p.~77)} }
\expandafter\def\csname SRef@cdp:st:A.5.9\endcsname{\href{S6_SERIES_GRIND_II.pdf\#nameddest=tag-01BJ}{II.7.6 (p.~78)} }
\expandafter\def\csname SRef@cdp:st:A.5.11\endcsname{\href{S6_SERIES_GRIND_II.pdf\#nameddest=tag-01B1}{II.7.7 (p.~80)} }
\expandafter\def\csname SRef@cdp:st:A.5.12\endcsname{\href{S6_SERIES_GRIND_II.pdf\#nameddest=tag-01B2}{II.7.8 (p.~82)} }
\expandafter\def\csname SRef@cdp:st:A.5.13\endcsname{\href{S6_SERIES_GRIND_II.pdf\#nameddest=tag-01B3}{II.7.9 (p.~83)} }
\expandafter\def\csname SRef@cdp:st:C.2.NC.2\endcsname{\href{S6_SERIES_GRIND_II.pdf\#nameddest=tag-01C9}{II.8.1 (p.~121)} }
\expandafter\def\csname SRef@cdp:st:C.2.NC.3\endcsname{\href{S6_SERIES_GRIND_II.pdf\#nameddest=tag-01CA}{II.8.2 (p.~122)} }
\expandafter\def\csname SRef@cdp:st:C.2.NC.4\endcsname{\href{S6_SERIES_GRIND_II.pdf\#nameddest=tag-01CB}{II.8.3 (p.~123)} }
\expandafter\def\csname SRef@cdp:st:C.2.NC.5\endcsname{\href{S6_SERIES_GRIND_II.pdf\#nameddest=tag-01CC}{II.8.4 (p.~124)} }
\expandafter\def\csname SRef@cdp:st:C2d\endcsname{\href{S6_SERIES_GRIND_II.pdf\#nameddest=tag-01FJ}{II.8.5 (p.~125)} }
\expandafter\def\csname SRef@night:P4.1\endcsname{\href{S6_SERIES_GRIND_II.pdf\#nameddest=tag-02WG}{II.9.1 (p.~127)} }
\expandafter\def\csname SRef@stmt:P4.1\endcsname{\href{S6_SERIES_GRIND_II.pdf\#nameddest=tag-02WG}{II.9.1 (p.~127)} }
\expandafter\def\csname SRef@night:P4.10\endcsname{\href{S6_SERIES_GRIND_II.pdf\#nameddest=tag-02WH}{II.9.10 (p.~138)} }
\expandafter\def\csname SRef@stmt:P4.10\endcsname{\href{S6_SERIES_GRIND_II.pdf\#nameddest=tag-02WH}{II.9.10 (p.~138)} }
\expandafter\def\csname SRef@night:P4.11\endcsname{\href{S6_SERIES_GRIND_II.pdf\#nameddest=tag-02WI}{II.9.11 (p.~139)} }
\expandafter\def\csname SRef@stmt:P4.11\endcsname{\href{S6_SERIES_GRIND_II.pdf\#nameddest=tag-02WI}{II.9.11 (p.~139)} }
\expandafter\def\csname SRef@night:P4.12\endcsname{\href{S6_SERIES_GRIND_II.pdf\#nameddest=tag-02WJ}{II.9.12 (p.~140)} }
\expandafter\def\csname SRef@stmt:P4.12\endcsname{\href{S6_SERIES_GRIND_II.pdf\#nameddest=tag-02WJ}{II.9.12 (p.~140)} }
\expandafter\def\csname SRef@night:P4.13\endcsname{\href{S6_SERIES_GRIND_II.pdf\#nameddest=tag-02WK}{II.9.13 (p.~142)} }
\expandafter\def\csname SRef@stmt:P4.13\endcsname{\href{S6_SERIES_GRIND_II.pdf\#nameddest=tag-02WK}{II.9.13 (p.~142)} }
\expandafter\def\csname SRef@night:P4.14\endcsname{\href{S6_SERIES_GRIND_II.pdf\#nameddest=tag-02WL}{II.9.14 (p.~143)} }
\expandafter\def\csname SRef@stmt:P4.14\endcsname{\href{S6_SERIES_GRIND_II.pdf\#nameddest=tag-02WL}{II.9.14 (p.~143)} }
\expandafter\def\csname SRef@night:P4.15\endcsname{\href{S6_SERIES_GRIND_II.pdf\#nameddest=tag-02WM}{II.9.15 (p.~144)} }
\expandafter\def\csname SRef@stmt:P4.15\endcsname{\href{S6_SERIES_GRIND_II.pdf\#nameddest=tag-02WM}{II.9.15 (p.~144)} }
\expandafter\def\csname SRef@night:P4.16\endcsname{\href{S6_SERIES_GRIND_II.pdf\#nameddest=tag-02WN}{II.9.16 (p.~145)} }
\expandafter\def\csname SRef@stmt:P4.16\endcsname{\href{S6_SERIES_GRIND_II.pdf\#nameddest=tag-02WN}{II.9.16 (p.~145)} }
\expandafter\def\csname SRef@night:P4.17\endcsname{\href{S6_SERIES_GRIND_II.pdf\#nameddest=tag-02WO}{II.9.17 (p.~147)} }
\expandafter\def\csname SRef@stmt:P4.17\endcsname{\href{S6_SERIES_GRIND_II.pdf\#nameddest=tag-02WO}{II.9.17 (p.~147)} }
\expandafter\def\csname SRef@night:P4.18\endcsname{\href{S6_SERIES_GRIND_II.pdf\#nameddest=tag-02WP}{II.9.18 (p.~148)} }
\expandafter\def\csname SRef@stmt:P4.18\endcsname{\href{S6_SERIES_GRIND_II.pdf\#nameddest=tag-02WP}{II.9.18 (p.~148)} }
\expandafter\def\csname SRef@night:P4.19\endcsname{\href{S6_SERIES_GRIND_II.pdf\#nameddest=tag-02WQ}{II.9.19 (p.~149)} }
\expandafter\def\csname SRef@stmt:P4.19\endcsname{\href{S6_SERIES_GRIND_II.pdf\#nameddest=tag-02WQ}{II.9.19 (p.~149)} }
\expandafter\def\csname SRef@night:P4.2\endcsname{\href{S6_SERIES_GRIND_II.pdf\#nameddest=tag-02WR}{II.9.2 (p.~128)} }
\expandafter\def\csname SRef@stmt:P4.2\endcsname{\href{S6_SERIES_GRIND_II.pdf\#nameddest=tag-02WR}{II.9.2 (p.~128)} }
\expandafter\def\csname SRef@night:P4.20\endcsname{\href{S6_SERIES_GRIND_II.pdf\#nameddest=tag-02WS}{II.9.20 (p.~149)} }
\expandafter\def\csname SRef@stmt:P4.20\endcsname{\href{S6_SERIES_GRIND_II.pdf\#nameddest=tag-02WS}{II.9.20 (p.~149)} }
\expandafter\def\csname SRef@night:P4.21\endcsname{\href{S6_SERIES_GRIND_II.pdf\#nameddest=tag-02WT}{II.9.21 (p.~151)} }
\expandafter\def\csname SRef@stmt:P4.21\endcsname{\href{S6_SERIES_GRIND_II.pdf\#nameddest=tag-02WT}{II.9.21 (p.~151)} }
\expandafter\def\csname SRef@night:P4.22\endcsname{\href{S6_SERIES_GRIND_II.pdf\#nameddest=tag-02WU}{II.9.22 (p.~153)} }
\expandafter\def\csname SRef@stmt:P4.22\endcsname{\href{S6_SERIES_GRIND_II.pdf\#nameddest=tag-02WU}{II.9.22 (p.~153)} }
\expandafter\def\csname SRef@night:P4.23\endcsname{\href{S6_SERIES_GRIND_II.pdf\#nameddest=tag-02WV}{II.9.23 (p.~154)} }
\expandafter\def\csname SRef@stmt:P4.23\endcsname{\href{S6_SERIES_GRIND_II.pdf\#nameddest=tag-02WV}{II.9.23 (p.~154)} }
\expandafter\def\csname SRef@night:P4.24\endcsname{\href{S6_SERIES_GRIND_II.pdf\#nameddest=tag-02WW}{II.9.24 (p.~155)} }
\expandafter\def\csname SRef@stmt:P4.24\endcsname{\href{S6_SERIES_GRIND_II.pdf\#nameddest=tag-02WW}{II.9.24 (p.~155)} }
\expandafter\def\csname SRef@night:P4.25\endcsname{\href{S6_SERIES_GRIND_II.pdf\#nameddest=tag-02WX}{II.9.25 (p.~158)} }
\expandafter\def\csname SRef@stmt:P4.25\endcsname{\href{S6_SERIES_GRIND_II.pdf\#nameddest=tag-02WX}{II.9.25 (p.~158)} }
\expandafter\def\csname SRef@night:P4.26\endcsname{\href{S6_SERIES_GRIND_II.pdf\#nameddest=tag-02WY}{II.9.26 (p.~160)} }
\expandafter\def\csname SRef@stmt:P4.26\endcsname{\href{S6_SERIES_GRIND_II.pdf\#nameddest=tag-02WY}{II.9.26 (p.~160)} }
\expandafter\def\csname SRef@night:P4.27\endcsname{\href{S6_SERIES_GRIND_II.pdf\#nameddest=tag-02WZ}{II.9.27 (p.~162)} }
\expandafter\def\csname SRef@stmt:P4.27\endcsname{\href{S6_SERIES_GRIND_II.pdf\#nameddest=tag-02WZ}{II.9.27 (p.~162)} }
\expandafter\def\csname SRef@night:P4.28\endcsname{\href{S6_SERIES_GRIND_II.pdf\#nameddest=tag-02X0}{II.9.28 (p.~164)} }
\expandafter\def\csname SRef@stmt:P4.28\endcsname{\href{S6_SERIES_GRIND_II.pdf\#nameddest=tag-02X0}{II.9.28 (p.~164)} }
\expandafter\def\csname SRef@night:P4.29\endcsname{\href{S6_SERIES_GRIND_II.pdf\#nameddest=tag-02X1}{II.9.29 (p.~166)} }
\expandafter\def\csname SRef@stmt:P4.29\endcsname{\href{S6_SERIES_GRIND_II.pdf\#nameddest=tag-02X1}{II.9.29 (p.~166)} }
\expandafter\def\csname SRef@night:P4.3\endcsname{\href{S6_SERIES_GRIND_II.pdf\#nameddest=tag-02X2}{II.9.3 (p.~129)} }
\expandafter\def\csname SRef@stmt:P4.3\endcsname{\href{S6_SERIES_GRIND_II.pdf\#nameddest=tag-02X2}{II.9.3 (p.~129)} }
\expandafter\def\csname SRef@night:P4.30\endcsname{\href{S6_SERIES_GRIND_II.pdf\#nameddest=tag-02X3}{II.9.30 (p.~167)} }
\expandafter\def\csname SRef@stmt:P4.30\endcsname{\href{S6_SERIES_GRIND_II.pdf\#nameddest=tag-02X3}{II.9.30 (p.~167)} }
\expandafter\def\csname SRef@night:P4.31\endcsname{\href{S6_SERIES_GRIND_II.pdf\#nameddest=tag-02X4}{II.9.31 (p.~168)} }
\expandafter\def\csname SRef@stmt:P4.31\endcsname{\href{S6_SERIES_GRIND_II.pdf\#nameddest=tag-02X4}{II.9.31 (p.~168)} }
\expandafter\def\csname SRef@night:P4.32\endcsname{\href{S6_SERIES_GRIND_II.pdf\#nameddest=tag-02X5}{II.9.32 (p.~170)} }
\expandafter\def\csname SRef@stmt:P4.32\endcsname{\href{S6_SERIES_GRIND_II.pdf\#nameddest=tag-02X5}{II.9.32 (p.~170)} }
\expandafter\def\csname SRef@night:P4.33\endcsname{\href{S6_SERIES_GRIND_II.pdf\#nameddest=tag-02X6}{II.9.33 (p.~171)} }
\expandafter\def\csname SRef@stmt:P4.33\endcsname{\href{S6_SERIES_GRIND_II.pdf\#nameddest=tag-02X6}{II.9.33 (p.~171)} }
\expandafter\def\csname SRef@night:P4.34\endcsname{\href{S6_SERIES_GRIND_II.pdf\#nameddest=tag-02X7}{II.9.34 (p.~172)} }
\expandafter\def\csname SRef@stmt:P4.34\endcsname{\href{S6_SERIES_GRIND_II.pdf\#nameddest=tag-02X7}{II.9.34 (p.~172)} }
\expandafter\def\csname SRef@night:P4.35\endcsname{\href{S6_SERIES_GRIND_II.pdf\#nameddest=tag-02X8}{II.9.35 (p.~174)} }
\expandafter\def\csname SRef@stmt:P4.35\endcsname{\href{S6_SERIES_GRIND_II.pdf\#nameddest=tag-02X8}{II.9.35 (p.~174)} }
\expandafter\def\csname SRef@night:P4.36\endcsname{\href{S6_SERIES_GRIND_II.pdf\#nameddest=tag-02X9}{II.9.36 (p.~176)} }
\expandafter\def\csname SRef@stmt:P4.36\endcsname{\href{S6_SERIES_GRIND_II.pdf\#nameddest=tag-02X9}{II.9.36 (p.~176)} }
\expandafter\def\csname SRef@night:P4.37\endcsname{\href{S6_SERIES_GRIND_II.pdf\#nameddest=tag-02XA}{II.9.37 (p.~178)} }
\expandafter\def\csname SRef@stmt:P4.37\endcsname{\href{S6_SERIES_GRIND_II.pdf\#nameddest=tag-02XA}{II.9.37 (p.~178)} }
\expandafter\def\csname SRef@night:P4.38\endcsname{\href{S6_SERIES_GRIND_II.pdf\#nameddest=tag-02XB}{II.9.38 (p.~179)} }
\expandafter\def\csname SRef@stmt:P4.38\endcsname{\href{S6_SERIES_GRIND_II.pdf\#nameddest=tag-02XB}{II.9.38 (p.~179)} }
\expandafter\def\csname SRef@night:P4.39\endcsname{\href{S6_SERIES_GRIND_II.pdf\#nameddest=tag-02XC}{II.9.39 (p.~181)} }
\expandafter\def\csname SRef@stmt:P4.39\endcsname{\href{S6_SERIES_GRIND_II.pdf\#nameddest=tag-02XC}{II.9.39 (p.~181)} }
\expandafter\def\csname SRef@night:P4.4\endcsname{\href{S6_SERIES_GRIND_II.pdf\#nameddest=tag-02XD}{II.9.4 (p.~130)} }
\expandafter\def\csname SRef@stmt:P4.4\endcsname{\href{S6_SERIES_GRIND_II.pdf\#nameddest=tag-02XD}{II.9.4 (p.~130)} }
\expandafter\def\csname SRef@night:P4.40\endcsname{\href{S6_SERIES_GRIND_II.pdf\#nameddest=tag-02XE}{II.9.40 (p.~183)} }
\expandafter\def\csname SRef@stmt:P4.40\endcsname{\href{S6_SERIES_GRIND_II.pdf\#nameddest=tag-02XE}{II.9.40 (p.~183)} }
\expandafter\def\csname SRef@night:P4.41\endcsname{\href{S6_SERIES_GRIND_II.pdf\#nameddest=tag-02XF}{II.9.41 (p.~184)} }
\expandafter\def\csname SRef@stmt:P4.41\endcsname{\href{S6_SERIES_GRIND_II.pdf\#nameddest=tag-02XF}{II.9.41 (p.~184)} }
\expandafter\def\csname SRef@night:P4.42\endcsname{\href{S6_SERIES_GRIND_II.pdf\#nameddest=tag-02XG}{II.9.42 (p.~187)} }
\expandafter\def\csname SRef@stmt:P4.42\endcsname{\href{S6_SERIES_GRIND_II.pdf\#nameddest=tag-02XG}{II.9.42 (p.~187)} }
\expandafter\def\csname SRef@night:P4.43\endcsname{\href{S6_SERIES_GRIND_II.pdf\#nameddest=tag-02XH}{II.9.43 (p.~188)} }
\expandafter\def\csname SRef@stmt:P4.43\endcsname{\href{S6_SERIES_GRIND_II.pdf\#nameddest=tag-02XH}{II.9.43 (p.~188)} }
\expandafter\def\csname SRef@night:P4.44\endcsname{\href{S6_SERIES_GRIND_II.pdf\#nameddest=tag-02XI}{II.9.44 (p.~189)} }
\expandafter\def\csname SRef@stmt:P4.44\endcsname{\href{S6_SERIES_GRIND_II.pdf\#nameddest=tag-02XI}{II.9.44 (p.~189)} }
\expandafter\def\csname SRef@night:P4.45\endcsname{\href{S6_SERIES_GRIND_II.pdf\#nameddest=tag-02XJ}{II.9.45 (p.~190)} }
\expandafter\def\csname SRef@stmt:P4.45\endcsname{\href{S6_SERIES_GRIND_II.pdf\#nameddest=tag-02XJ}{II.9.45 (p.~190)} }
\expandafter\def\csname SRef@night:P4.46\endcsname{\href{S6_SERIES_GRIND_II.pdf\#nameddest=tag-02XK}{II.9.46 (p.~191)} }
\expandafter\def\csname SRef@stmt:P4.46\endcsname{\href{S6_SERIES_GRIND_II.pdf\#nameddest=tag-02XK}{II.9.46 (p.~191)} }
\expandafter\def\csname SRef@night:P4.47\endcsname{\href{S6_SERIES_GRIND_II.pdf\#nameddest=tag-02XL}{II.9.47 (p.~192)} }
\expandafter\def\csname SRef@stmt:P4.47\endcsname{\href{S6_SERIES_GRIND_II.pdf\#nameddest=tag-02XL}{II.9.47 (p.~192)} }
\expandafter\def\csname SRef@night:P4.48\endcsname{\href{S6_SERIES_GRIND_II.pdf\#nameddest=tag-02XM}{II.9.48 (p.~193)} }
\expandafter\def\csname SRef@stmt:P4.48\endcsname{\href{S6_SERIES_GRIND_II.pdf\#nameddest=tag-02XM}{II.9.48 (p.~193)} }
\expandafter\def\csname SRef@night:P4.49\endcsname{\href{S6_SERIES_GRIND_II.pdf\#nameddest=tag-02XN}{II.9.49 (p.~194)} }
\expandafter\def\csname SRef@stmt:P4.49\endcsname{\href{S6_SERIES_GRIND_II.pdf\#nameddest=tag-02XN}{II.9.49 (p.~194)} }
\expandafter\def\csname SRef@night:P4.5\endcsname{\href{S6_SERIES_GRIND_II.pdf\#nameddest=tag-02XO}{II.9.5 (p.~134)} }
\expandafter\def\csname SRef@stmt:P4.5\endcsname{\href{S6_SERIES_GRIND_II.pdf\#nameddest=tag-02XO}{II.9.5 (p.~134)} }
\expandafter\def\csname SRef@night:P4.50\endcsname{\href{S6_SERIES_GRIND_II.pdf\#nameddest=tag-02XP}{II.9.50 (p.~198)} }
\expandafter\def\csname SRef@stmt:P4.50\endcsname{\href{S6_SERIES_GRIND_II.pdf\#nameddest=tag-02XP}{II.9.50 (p.~198)} }
\expandafter\def\csname SRef@night:P4.51\endcsname{\href{S6_SERIES_GRIND_II.pdf\#nameddest=tag-02XQ}{II.9.51 (p.~198)} }
\expandafter\def\csname SRef@stmt:P4.51\endcsname{\href{S6_SERIES_GRIND_II.pdf\#nameddest=tag-02XQ}{II.9.51 (p.~198)} }
\expandafter\def\csname SRef@night:P4.52\endcsname{\href{S6_SERIES_GRIND_II.pdf\#nameddest=tag-02XR}{II.9.52 (p.~199)} }
\expandafter\def\csname SRef@stmt:P4.52\endcsname{\href{S6_SERIES_GRIND_II.pdf\#nameddest=tag-02XR}{II.9.52 (p.~199)} }
\expandafter\def\csname SRef@night:P4.53\endcsname{\href{S6_SERIES_GRIND_II.pdf\#nameddest=tag-02XS}{II.9.53 (p.~201)} }
\expandafter\def\csname SRef@stmt:P4.53\endcsname{\href{S6_SERIES_GRIND_II.pdf\#nameddest=tag-02XS}{II.9.53 (p.~201)} }
\expandafter\def\csname SRef@night:P4.54\endcsname{\href{S6_SERIES_GRIND_II.pdf\#nameddest=tag-02XT}{II.9.54 (p.~204)} }
\expandafter\def\csname SRef@stmt:P4.54\endcsname{\href{S6_SERIES_GRIND_II.pdf\#nameddest=tag-02XT}{II.9.54 (p.~204)} }
\expandafter\def\csname SRef@night:P4.55\endcsname{\href{S6_SERIES_GRIND_II.pdf\#nameddest=tag-02XU}{II.9.55 (p.~210)} }
\expandafter\def\csname SRef@stmt:P4.55\endcsname{\href{S6_SERIES_GRIND_II.pdf\#nameddest=tag-02XU}{II.9.55 (p.~210)} }
\expandafter\def\csname SRef@night:P4.56\endcsname{\href{S6_SERIES_GRIND_II.pdf\#nameddest=tag-02XV}{II.9.56 (p.~213)} }
\expandafter\def\csname SRef@stmt:P4.56\endcsname{\href{S6_SERIES_GRIND_II.pdf\#nameddest=tag-02XV}{II.9.56 (p.~213)} }
\expandafter\def\csname SRef@night:P4.57\endcsname{\href{S6_SERIES_GRIND_II.pdf\#nameddest=tag-02XW}{II.9.57 (p.~216)} }
\expandafter\def\csname SRef@stmt:P4.57\endcsname{\href{S6_SERIES_GRIND_II.pdf\#nameddest=tag-02XW}{II.9.57 (p.~216)} }
\expandafter\def\csname SRef@night:P4.58\endcsname{\href{S6_SERIES_GRIND_II.pdf\#nameddest=tag-02XX}{II.9.58 (p.~218)} }
\expandafter\def\csname SRef@stmt:P4.58\endcsname{\href{S6_SERIES_GRIND_II.pdf\#nameddest=tag-02XX}{II.9.58 (p.~218)} }
\expandafter\def\csname SRef@night:P4R.1\endcsname{\href{S6_SERIES_GRIND_II.pdf\#nameddest=tag-02Y2}{II.9.59 (p.~219)} }
\expandafter\def\csname SRef@stmt:P4R.1\endcsname{\href{S6_SERIES_GRIND_II.pdf\#nameddest=tag-02Y2}{II.9.59 (p.~219)} }
\expandafter\def\csname SRef@night:P4.6\endcsname{\href{S6_SERIES_GRIND_II.pdf\#nameddest=tag-02XY}{II.9.6 (p.~135)} }
\expandafter\def\csname SRef@stmt:P4.6\endcsname{\href{S6_SERIES_GRIND_II.pdf\#nameddest=tag-02XY}{II.9.6 (p.~135)} }
\expandafter\def\csname SRef@night:P4R.2\endcsname{\href{S6_SERIES_GRIND_II.pdf\#nameddest=tag-02YD}{II.9.60 (p.~220)} }
\expandafter\def\csname SRef@stmt:P4R.2\endcsname{\href{S6_SERIES_GRIND_II.pdf\#nameddest=tag-02YD}{II.9.60 (p.~220)} }
\expandafter\def\csname SRef@night:P4R.3\endcsname{\href{S6_SERIES_GRIND_II.pdf\#nameddest=tag-02YF}{II.9.61 (p.~221)} }
\expandafter\def\csname SRef@stmt:P4R.3\endcsname{\href{S6_SERIES_GRIND_II.pdf\#nameddest=tag-02YF}{II.9.61 (p.~221)} }
\expandafter\def\csname SRef@night:P4R.4\endcsname{\href{S6_SERIES_GRIND_II.pdf\#nameddest=tag-02YG}{II.9.62 (p.~223)} }
\expandafter\def\csname SRef@stmt:P4R.4\endcsname{\href{S6_SERIES_GRIND_II.pdf\#nameddest=tag-02YG}{II.9.62 (p.~223)} }
\expandafter\def\csname SRef@night:P4R.5\endcsname{\href{S6_SERIES_GRIND_II.pdf\#nameddest=tag-02YH}{II.9.63 (p.~224)} }
\expandafter\def\csname SRef@stmt:P4R.5\endcsname{\href{S6_SERIES_GRIND_II.pdf\#nameddest=tag-02YH}{II.9.63 (p.~224)} }
\expandafter\def\csname SRef@night:P4R.6\endcsname{\href{S6_SERIES_GRIND_II.pdf\#nameddest=tag-02YI}{II.9.64 (p.~225)} }
\expandafter\def\csname SRef@stmt:P4R.6\endcsname{\href{S6_SERIES_GRIND_II.pdf\#nameddest=tag-02YI}{II.9.64 (p.~225)} }
\expandafter\def\csname SRef@night:P4R.7\endcsname{\href{S6_SERIES_GRIND_II.pdf\#nameddest=tag-02YJ}{II.9.65 (p.~227)} }
\expandafter\def\csname SRef@stmt:P4R.7\endcsname{\href{S6_SERIES_GRIND_II.pdf\#nameddest=tag-02YJ}{II.9.65 (p.~227)} }
\expandafter\def\csname SRef@night:P4R.8\endcsname{\href{S6_SERIES_GRIND_II.pdf\#nameddest=tag-02YK}{II.9.66 (p.~227)} }
\expandafter\def\csname SRef@stmt:P4R.8\endcsname{\href{S6_SERIES_GRIND_II.pdf\#nameddest=tag-02YK}{II.9.66 (p.~227)} }
\expandafter\def\csname SRef@night:P4R.9\endcsname{\href{S6_SERIES_GRIND_II.pdf\#nameddest=tag-02YL}{II.9.67 (p.~228)} }
\expandafter\def\csname SRef@stmt:P4R.9\endcsname{\href{S6_SERIES_GRIND_II.pdf\#nameddest=tag-02YL}{II.9.67 (p.~228)} }
\expandafter\def\csname SRef@night:P4R.10\endcsname{\href{S6_SERIES_GRIND_II.pdf\#nameddest=tag-02Y3}{II.9.68 (p.~229)} }
\expandafter\def\csname SRef@stmt:P4R.10\endcsname{\href{S6_SERIES_GRIND_II.pdf\#nameddest=tag-02Y3}{II.9.68 (p.~229)} }
\expandafter\def\csname SRef@night:P4R.11\endcsname{\href{S6_SERIES_GRIND_II.pdf\#nameddest=tag-02Y4}{II.9.69 (p.~231)} }
\expandafter\def\csname SRef@stmt:P4R.11\endcsname{\href{S6_SERIES_GRIND_II.pdf\#nameddest=tag-02Y4}{II.9.69 (p.~231)} }
\expandafter\def\csname SRef@night:P4.7\endcsname{\href{S6_SERIES_GRIND_II.pdf\#nameddest=tag-02XZ}{II.9.7 (p.~136)} }
\expandafter\def\csname SRef@stmt:P4.7\endcsname{\href{S6_SERIES_GRIND_II.pdf\#nameddest=tag-02XZ}{II.9.7 (p.~136)} }
\expandafter\def\csname SRef@night:P4R.12\endcsname{\href{S6_SERIES_GRIND_II.pdf\#nameddest=tag-02Y5}{II.9.70 (p.~232)} }
\expandafter\def\csname SRef@stmt:P4R.12\endcsname{\href{S6_SERIES_GRIND_II.pdf\#nameddest=tag-02Y5}{II.9.70 (p.~232)} }
\expandafter\def\csname SRef@night:P4R.13\endcsname{\href{S6_SERIES_GRIND_II.pdf\#nameddest=tag-02Y6}{II.9.71 (p.~233)} }
\expandafter\def\csname SRef@stmt:P4R.13\endcsname{\href{S6_SERIES_GRIND_II.pdf\#nameddest=tag-02Y6}{II.9.71 (p.~233)} }
\expandafter\def\csname SRef@night:P4R.14\endcsname{\href{S6_SERIES_GRIND_II.pdf\#nameddest=tag-02Y7}{II.9.72 (p.~237)} }
\expandafter\def\csname SRef@stmt:P4R.14\endcsname{\href{S6_SERIES_GRIND_II.pdf\#nameddest=tag-02Y7}{II.9.72 (p.~237)} }
\expandafter\def\csname SRef@night:P4R.15\endcsname{\href{S6_SERIES_GRIND_II.pdf\#nameddest=tag-02Y8}{II.9.73 (p.~239)} }
\expandafter\def\csname SRef@stmt:P4R.15\endcsname{\href{S6_SERIES_GRIND_II.pdf\#nameddest=tag-02Y8}{II.9.73 (p.~239)} }
\expandafter\def\csname SRef@night:P4R.16\endcsname{\href{S6_SERIES_GRIND_II.pdf\#nameddest=tag-02Y9}{II.9.74 (p.~240)} }
\expandafter\def\csname SRef@stmt:P4R.16\endcsname{\href{S6_SERIES_GRIND_II.pdf\#nameddest=tag-02Y9}{II.9.74 (p.~240)} }
\expandafter\def\csname SRef@night:P4R.17\endcsname{\href{S6_SERIES_GRIND_II.pdf\#nameddest=tag-02YA}{II.9.75 (p.~242)} }
\expandafter\def\csname SRef@stmt:P4R.17\endcsname{\href{S6_SERIES_GRIND_II.pdf\#nameddest=tag-02YA}{II.9.75 (p.~242)} }
\expandafter\def\csname SRef@bounds-without-an-ordinary-square-assumption\endcsname{\href{S6_SERIES_GRIND_II.pdf\#nameddest=tag-02YA}{II.{{$\langle 1\rangle $1.}} (p.~242)} }
\expandafter\def\csname SRef@night:P4R.18\endcsname{\href{S6_SERIES_GRIND_II.pdf\#nameddest=tag-02YB}{II.9.76 (p.~246)} }
\expandafter\def\csname SRef@stmt:P4R.18\endcsname{\href{S6_SERIES_GRIND_II.pdf\#nameddest=tag-02YB}{II.9.76 (p.~246)} }
\expandafter\def\csname SRef@night:P4R.19\endcsname{\href{S6_SERIES_GRIND_II.pdf\#nameddest=tag-02YC}{II.9.77 (p.~249)} }
\expandafter\def\csname SRef@stmt:P4R.19\endcsname{\href{S6_SERIES_GRIND_II.pdf\#nameddest=tag-02YC}{II.9.77 (p.~249)} }
\expandafter\def\csname SRef@night:P4R.20\endcsname{\href{S6_SERIES_GRIND_II.pdf\#nameddest=tag-02YE}{II.9.78 (p.~252)} }
\expandafter\def\csname SRef@stmt:P4R.20\endcsname{\href{S6_SERIES_GRIND_II.pdf\#nameddest=tag-02YE}{II.9.78 (p.~252)} }
\expandafter\def\csname SRef@night:P4.8\endcsname{\href{S6_SERIES_GRIND_II.pdf\#nameddest=tag-02Y0}{II.9.8 (p.~136)} }
\expandafter\def\csname SRef@stmt:P4.8\endcsname{\href{S6_SERIES_GRIND_II.pdf\#nameddest=tag-02Y0}{II.9.8 (p.~136)} }
\expandafter\def\csname SRef@night:P4.9\endcsname{\href{S6_SERIES_GRIND_II.pdf\#nameddest=tag-02Y1}{II.9.9 (p.~138)} }
\expandafter\def\csname SRef@stmt:P4.9\endcsname{\href{S6_SERIES_GRIND_II.pdf\#nameddest=tag-02Y1}{II.9.9 (p.~138)} }
\expandafter\def\csname SRef@cdp:st:C4a\endcsname{\href{S6_SERIES_GRIND_II.pdf\#nameddest=tag-01FK}{II.10.1 (p.~255)} }
\expandafter\def\csname SRef@cdp:st:C5c\endcsname{\href{S6_SERIES_GRIND_II.pdf\#nameddest=tag-01FT}{II.10.10 (p.~269)} }
\expandafter\def\csname SRef@cdp:st:C4b\endcsname{\href{S6_SERIES_GRIND_II.pdf\#nameddest=tag-01FL}{II.10.2 (p.~257)} }
\expandafter\def\csname SRef@cdp:st:C4c\endcsname{\href{S6_SERIES_GRIND_II.pdf\#nameddest=tag-01FM}{II.10.3 (p.~258)} }
\expandafter\def\csname SRef@cdp:st:C4d\endcsname{\href{S6_SERIES_GRIND_II.pdf\#nameddest=tag-01FN}{II.10.4 (p.~259)} }
\expandafter\def\csname SRef@cdp:st:C4e\endcsname{\href{S6_SERIES_GRIND_II.pdf\#nameddest=tag-01FO}{II.10.5 (p.~260)} }
\expandafter\def\csname SRef@cdp:st:C4f\endcsname{\href{S6_SERIES_GRIND_II.pdf\#nameddest=tag-01FP}{II.10.6 (p.~263)} }
\expandafter\def\csname SRef@cdp:st:C4g\endcsname{\href{S6_SERIES_GRIND_II.pdf\#nameddest=tag-01FQ}{II.10.7 (p.~264)} }
\expandafter\def\csname SRef@cdp:st:C5a\endcsname{\href{S6_SERIES_GRIND_II.pdf\#nameddest=tag-01FR}{II.10.8 (p.~266)} }
\expandafter\def\csname SRef@cdp:st:C5b\endcsname{\href{S6_SERIES_GRIND_II.pdf\#nameddest=tag-01FS}{II.10.9 (p.~267)} }
\expandafter\def\csname SRef@night:P3K.36\endcsname{\href{S6_SERIES_GRIND_II.pdf\#nameddest=tag-02UH}{II.17.2 (p.~411)} }
\expandafter\def\csname SRef@stmt:P3K.36\endcsname{\href{S6_SERIES_GRIND_II.pdf\#nameddest=tag-02UH}{II.17.2 (p.~411)} }
\expandafter\def\csname SRef@inv:wb:uniform-nontorsion-provenance\endcsname{\href{S6_SERIES_GRIND_II.pdf\#nameddest=tag-01TQ}{II.15.2 (p.~363)} }
\expandafter\def\csname SRef@wbII:rem:scope\endcsname{\href{S6_SERIES_GRIND_II.pdf\#nameddest=tag-03IS}{II.16.3 (p.~371)} }
\expandafter\def\csname SRef@wbII:rem:local-cusp-picard-obstruction\endcsname{\href{S6_SERIES_GRIND_II.pdf\#nameddest=tag-03IQ}{II.16.8 (p.~387)} }
\expandafter\def\csname SRef@wbII:rem:relative-local-scope\endcsname{\href{S6_SERIES_GRIND_II.pdf\#nameddest=tag-03IR}{II.16.12 (p.~407)} }
\expandafter\def\csname SRef@sec:II:conventions\endcsname{\href{S6_SERIES_GRIND_II.pdf\#nameddest=tag-034E}{II.2 (p.~18)} }
\expandafter\def\csname SRef@sec:II:13\endcsname{\href{S6_SERIES_GRIND_II.pdf\#nameddest=tag-0349}{II.11 (p.~272)} }
\expandafter\def\csname SRef@sec:II:14\endcsname{\href{S6_SERIES_GRIND_II.pdf\#nameddest=tag-034A}{II.12 (p.~285)} }
\expandafter\def\csname SRef@sec:II:push-stability\endcsname{\href{S6_SERIES_GRIND_II.pdf\#nameddest=tag-034G}{II.13 (p.~309)} }
\expandafter\def\csname SRef@sec:II:combs\endcsname{\href{S6_SERIES_GRIND_II.pdf\#nameddest=tag-034D}{II.14 (p.~354)} }
\expandafter\def\csname SRef@sec:II:uniform-nontorsion\endcsname{\href{S6_SERIES_GRIND_II.pdf\#nameddest=tag-034H}{II.15 (p.~361)} }
\expandafter\def\csname SRef@app:II:programs\endcsname{\href{S6_SERIES_GRIND_II.pdf\#nameddest=tag-000B}{II.26 (p.~795)} }
\expandafter\def\csname SRef@sec:II:01\endcsname{\href{S6_SERIES_GRIND_II.pdf\#nameddest=tag-033X}{II.3 (p.~18)} }
\expandafter\def\csname SRef@sec:II:02\endcsname{\href{S6_SERIES_GRIND_II.pdf\#nameddest=tag-033Y}{II.4 (p.~20)} }
\expandafter\def\csname SRef@sec:II:03\endcsname{\href{S6_SERIES_GRIND_II.pdf\#nameddest=tag-033Z}{II.5 (p.~31)} }
\expandafter\def\csname SRef@sec:II:04\endcsname{\href{S6_SERIES_GRIND_II.pdf\#nameddest=tag-0340}{II.6 (p.~47)} }
\expandafter\def\csname SRef@sec:II:05\endcsname{\href{S6_SERIES_GRIND_II.pdf\#nameddest=tag-0341}{II.7 (p.~72)} }
\expandafter\def\csname SRef@sec:II:08\endcsname{\href{S6_SERIES_GRIND_II.pdf\#nameddest=tag-0344}{II.8 (p.~121)} }
\expandafter\def\csname SRef@sec:II:night-sign\endcsname{\href{S6_SERIES_GRIND_II.pdf\#nameddest=tag-034F}{II.9 (p.~127)} }
\expandafter\def\csname SRef@sec:II:12\endcsname{\href{S6_SERIES_GRIND_II.pdf\#nameddest=tag-0348}{II.10 (p.~255)} }
\expandafter\def\csname SRef@wbII:subsec:local-cusp-picard-obstruction\endcsname{\href{S6_SERIES_GRIND_II.pdf\#nameddest=tag-03IU}{II.16.8 (p.~382)} }
\expandafter\def\csname SRef@wbII:subsec:analytic-all-picard\endcsname{\href{S6_SERIES_GRIND_II.pdf\#nameddest=tag-03IT}{II.16.9 (p.~387)} }
\expandafter\def\csname SRef@push:src:E.2.23\endcsname{\href{S6_SERIES_GRIND_II.pdf\#nameddest=tag-0339}{II.13.20 (p.~347)} }
\expandafter\def\csname SRef@push:src:E.5.3\endcsname{\href{S6_SERIES_GRIND_II.pdf\#nameddest=tag-033P}{II.13.24 (p.~351)} }
\expandafter\def\csname SRef@wb:thm:uniform-nontorsion-audit\endcsname{\href{S6_SERIES_GRIND_II.pdf\#nameddest=tag-03DQ}{II.15.1 (p.~361)} }
\expandafter\def\csname SRef@wbII:prop:conductor\endcsname{\href{S6_SERIES_GRIND_II.pdf\#nameddest=tag-03IK}{II.16.2 (p.~366)} }
\expandafter\def\csname SRef@wbII:prop:conductor-differential\endcsname{\href{S6_SERIES_GRIND_II.pdf\#nameddest=tag-03IK}{II.16.2 (p.~366)} }
\expandafter\def\csname SRef@wbII:thm:analytic-all-picard\endcsname{\href{S6_SERIES_GRIND_II.pdf\#nameddest=tag-03IW}{II.16.21 (p.~398)} }
\expandafter\def\csname SRef@wbII:thm:candidate-conductor\endcsname{\href{S6_SERIES_GRIND_II.pdf\#nameddest=tag-03IX}{II.16.23 (p.~400)} }
\expandafter\def\csname SRef@wbII:thm:corrected-cdp-global\endcsname{\href{S6_SERIES_GRIND_II.pdf\#nameddest=tag-03J0}{II.16.27 (p.~402)} }
\expandafter\def\csname SRef@wbII:thm:corrected-global\endcsname{\href{S6_SERIES_GRIND_II.pdf\#nameddest=tag-03J0}{II.16.27 (p.~402)} }
\expandafter\def\csname SRef@wbII:thm:corrected-cdp-fibrewise\endcsname{\href{S6_SERIES_GRIND_II.pdf\#nameddest=tag-03IZ}{II.16.29 (p.~404)} }
\expandafter\def\csname SRef@wbII:thm:cartier-duality\endcsname{\href{S6_SERIES_GRIND_II.pdf\#nameddest=tag-03IY}{II.16.6 (p.~371)} }
\expandafter\def\csname SRef@ii:LR:covering\endcsname{\href{S6_SERIES_GRIND_II.pdf\#nameddest=tag-01RO}{II.17.1 (p.~407)} }
\expandafter\def\csname SRef@maps:eq:D.O9\endcsname{\href{S6_SERIES_GRIND_III.pdf\#nameddest=tag-020U}{III.{D.O9} (p.~85)} }
\expandafter\def\csname SRef@maps:eq:D.FI1\endcsname{\href{S6_SERIES_GRIND_III.pdf\#nameddest=tag-01YT}{III.{D.FI1} (p.~87)} }
\expandafter\def\csname SRef@maps:eq:D.FI2\endcsname{\href{S6_SERIES_GRIND_III.pdf\#nameddest=tag-01YU}{III.{D.FI2} (p.~88)} }
\expandafter\def\csname SRef@maps:eq:D.FI3\endcsname{\href{S6_SERIES_GRIND_III.pdf\#nameddest=tag-01YV}{III.{D.FI3} (p.~88)} }
\expandafter\def\csname SRef@maps:eq:D.FI4\endcsname{\href{S6_SERIES_GRIND_III.pdf\#nameddest=tag-01YW}{III.{D.FI4} (p.~90)} }
\expandafter\def\csname SRef@maps:eq:D.FI5\endcsname{\href{S6_SERIES_GRIND_III.pdf\#nameddest=tag-01YX}{III.{D.FI5} (p.~91)} }
\expandafter\def\csname SRef@maps:eq:D.FG1\endcsname{\href{S6_SERIES_GRIND_III.pdf\#nameddest=tag-01YQ}{III.{D.FG1} (p.~92)} }
\expandafter\def\csname SRef@maps:eq:D.FG2\endcsname{\href{S6_SERIES_GRIND_III.pdf\#nameddest=tag-01YR}{III.{D.FG2} (p.~92)} }
\expandafter\def\csname SRef@maps:eq:D.FG3\endcsname{\href{S6_SERIES_GRIND_III.pdf\#nameddest=tag-01YS}{III.{D.FG3} (p.~92)} }
\expandafter\def\csname SRef@maps:eq:D.IS1\endcsname{\href{S6_SERIES_GRIND_III.pdf\#nameddest=tag-0207}{III.{D.IS1} (p.~93)} }
\expandafter\def\csname SRef@maps:eq:D.IS2\endcsname{\href{S6_SERIES_GRIND_III.pdf\#nameddest=tag-0208}{III.{D.IS2} (p.~93)} }
\expandafter\def\csname SRef@maps:eq:D.SC1\endcsname{\href{S6_SERIES_GRIND_III.pdf\#nameddest=tag-023Z}{III.{D.SC1} (p.~94)} }
\expandafter\def\csname SRef@maps:eq:D.SC2\endcsname{\href{S6_SERIES_GRIND_III.pdf\#nameddest=tag-0240}{III.{D.SC2} (p.~95)} }
\expandafter\def\csname SRef@maps:eq:D.SC3\endcsname{\href{S6_SERIES_GRIND_III.pdf\#nameddest=tag-0241}{III.{D.SC3} (p.~95)} }
\expandafter\def\csname SRef@maps:eq:D.SC4\endcsname{\href{S6_SERIES_GRIND_III.pdf\#nameddest=tag-0242}{III.{D.SC4} (p.~95)} }
\expandafter\def\csname SRef@maps:eq:D.SC5\endcsname{\href{S6_SERIES_GRIND_III.pdf\#nameddest=tag-0243}{III.{D.SC5} (p.~95)} }
\expandafter\def\csname SRef@maps:eq:D.SC6\endcsname{\href{S6_SERIES_GRIND_III.pdf\#nameddest=tag-0244}{III.{D.SC6} (p.~95)} }
\expandafter\def\csname SRef@maps:eq:D.SC7\endcsname{\href{S6_SERIES_GRIND_III.pdf\#nameddest=tag-0245}{III.{D.SC7} (p.~97)} }
\expandafter\def\csname SRef@maps:eq:D.SC8\endcsname{\href{S6_SERIES_GRIND_III.pdf\#nameddest=tag-0246}{III.{D.SC8} (p.~97)} }
\expandafter\def\csname SRef@maps:eq:D.HY1\endcsname{\href{S6_SERIES_GRIND_III.pdf\#nameddest=tag-01ZN}{III.{D.HY1} (p.~98)} }
\expandafter\def\csname SRef@maps:eq:D.HY2\endcsname{\href{S6_SERIES_GRIND_III.pdf\#nameddest=tag-01ZO}{III.{D.HY2} (p.~98)} }
\expandafter\def\csname SRef@maps:eq:D.HY3\endcsname{\href{S6_SERIES_GRIND_III.pdf\#nameddest=tag-01ZP}{III.{D.HY3} (p.~99)} }
\expandafter\def\csname SRef@maps:eq:D.HY4\endcsname{\href{S6_SERIES_GRIND_III.pdf\#nameddest=tag-01ZQ}{III.{D.HY4} (p.~99)} }
\expandafter\def\csname SRef@maps:eq:D.HY5\endcsname{\href{S6_SERIES_GRIND_III.pdf\#nameddest=tag-01ZR}{III.{D.HY5} (p.~99)} }
\expandafter\def\csname SRef@maps:eq:D.HY6\endcsname{\href{S6_SERIES_GRIND_III.pdf\#nameddest=tag-01ZS}{III.{D.HY6} (p.~100)} }
\expandafter\def\csname SRef@maps:eq:D.HY7\endcsname{\href{S6_SERIES_GRIND_III.pdf\#nameddest=tag-01ZT}{III.{D.HY7} (p.~100)} }
\expandafter\def\csname SRef@maps:eq:D.ST1\endcsname{\href{S6_SERIES_GRIND_III.pdf\#nameddest=tag-024V}{III.{D.ST1} (p.~102)} }
\expandafter\def\csname SRef@maps:eq:D.ST2\endcsname{\href{S6_SERIES_GRIND_III.pdf\#nameddest=tag-0250}{III.{D.ST2} (p.~103)} }
\expandafter\def\csname SRef@maps:eq:D.ST3\endcsname{\href{S6_SERIES_GRIND_III.pdf\#nameddest=tag-0251}{III.{D.ST3} (p.~103)} }
\expandafter\def\csname SRef@maps:eq:D.ST4\endcsname{\href{S6_SERIES_GRIND_III.pdf\#nameddest=tag-0252}{III.{D.ST4} (p.~103)} }
\expandafter\def\csname SRef@maps:eq:D.ST5\endcsname{\href{S6_SERIES_GRIND_III.pdf\#nameddest=tag-0253}{III.{D.ST5} (p.~103)} }
\expandafter\def\csname SRef@maps:eq:D.ST6\endcsname{\href{S6_SERIES_GRIND_III.pdf\#nameddest=tag-0254}{III.{D.ST6} (p.~103)} }
\expandafter\def\csname SRef@maps:eq:D.ST7\endcsname{\href{S6_SERIES_GRIND_III.pdf\#nameddest=tag-0255}{III.{D.ST7} (p.~104)} }
\expandafter\def\csname SRef@maps:eq:D.ST8\endcsname{\href{S6_SERIES_GRIND_III.pdf\#nameddest=tag-0256}{III.{D.ST8} (p.~104)} }
\expandafter\def\csname SRef@maps:eq:D.ST9\endcsname{\href{S6_SERIES_GRIND_III.pdf\#nameddest=tag-0257}{III.{D.ST9} (p.~104)} }
\expandafter\def\csname SRef@maps:eq:D.ST10\endcsname{\href{S6_SERIES_GRIND_III.pdf\#nameddest=tag-024W}{III.{D.ST10} (p.~105)} }
\expandafter\def\csname SRef@maps:eq:D.ST11\endcsname{\href{S6_SERIES_GRIND_III.pdf\#nameddest=tag-024X}{III.{D.ST11} (p.~105)} }
\expandafter\def\csname SRef@maps:eq:D.ST12\endcsname{\href{S6_SERIES_GRIND_III.pdf\#nameddest=tag-024Y}{III.{D.ST12} (p.~105)} }
\expandafter\def\csname SRef@maps:eq:D.ST13\endcsname{\href{S6_SERIES_GRIND_III.pdf\#nameddest=tag-024Z}{III.{D.ST13} (p.~106)} }
\expandafter\def\csname SRef@maps:eq:D.UF1\endcsname{\href{S6_SERIES_GRIND_III.pdf\#nameddest=tag-025Z}{III.{D.UF1} (p.~107)} }
\expandafter\def\csname SRef@maps:eq:D.UF2\endcsname{\href{S6_SERIES_GRIND_III.pdf\#nameddest=tag-0261}{III.{D.UF2} (p.~107)} }
\expandafter\def\csname SRef@maps:eq:D.UF3\endcsname{\href{S6_SERIES_GRIND_III.pdf\#nameddest=tag-0262}{III.{D.UF3} (p.~107)} }
\expandafter\def\csname SRef@maps:eq:D.UF4\endcsname{\href{S6_SERIES_GRIND_III.pdf\#nameddest=tag-0263}{III.{D.UF4} (p.~107)} }
\expandafter\def\csname SRef@maps:eq:D.UF5\endcsname{\href{S6_SERIES_GRIND_III.pdf\#nameddest=tag-0264}{III.{D.UF5} (p.~108)} }
\expandafter\def\csname SRef@maps:eq:D.UF6\endcsname{\href{S6_SERIES_GRIND_III.pdf\#nameddest=tag-0265}{III.{D.UF6} (p.~108)} }
\expandafter\def\csname SRef@maps:eq:D.UF7\endcsname{\href{S6_SERIES_GRIND_III.pdf\#nameddest=tag-0266}{III.{D.UF7} (p.~108)} }
\expandafter\def\csname SRef@maps:eq:D.UF8\endcsname{\href{S6_SERIES_GRIND_III.pdf\#nameddest=tag-0267}{III.{D.UF8} (p.~109)} }
\expandafter\def\csname SRef@maps:eq:D.UF9\endcsname{\href{S6_SERIES_GRIND_III.pdf\#nameddest=tag-0268}{III.{D.UF9} (p.~109)} }
\expandafter\def\csname SRef@maps:eq:D.UF10\endcsname{\href{S6_SERIES_GRIND_III.pdf\#nameddest=tag-0260}{III.{D.UF10} (p.~109)} }
\expandafter\def\csname SRef@maps:eq:D.AF1\endcsname{\href{S6_SERIES_GRIND_III.pdf\#nameddest=tag-01XI}{III.{D.AF1} (p.~110)} }
\expandafter\def\csname SRef@maps:eq:D.AF2\endcsname{\href{S6_SERIES_GRIND_III.pdf\#nameddest=tag-01XJ}{III.{D.AF2} (p.~110)} }
\expandafter\def\csname SRef@maps:eq:D.AF3\endcsname{\href{S6_SERIES_GRIND_III.pdf\#nameddest=tag-01XK}{III.{D.AF3} (p.~111)} }
\expandafter\def\csname SRef@maps:eq:D.AF4\endcsname{\href{S6_SERIES_GRIND_III.pdf\#nameddest=tag-01XL}{III.{D.AF4} (p.~111)} }
\expandafter\def\csname SRef@maps:eq:D.AF5\endcsname{\href{S6_SERIES_GRIND_III.pdf\#nameddest=tag-01XM}{III.{D.AF5} (p.~111)} }
\expandafter\def\csname SRef@maps:eq:D.AF6\endcsname{\href{S6_SERIES_GRIND_III.pdf\#nameddest=tag-01XN}{III.{D.AF6} (p.~111)} }
\expandafter\def\csname SRef@maps:eq:D.AF7\endcsname{\href{S6_SERIES_GRIND_III.pdf\#nameddest=tag-01XO}{III.{D.AF7} (p.~111)} }
\expandafter\def\csname SRef@maps:eq:D.AF8\endcsname{\href{S6_SERIES_GRIND_III.pdf\#nameddest=tag-01XP}{III.{D.AF8} (p.~111)} }
\expandafter\def\csname SRef@maps:eq:D.AF9\endcsname{\href{S6_SERIES_GRIND_III.pdf\#nameddest=tag-01XQ}{III.{D.AF9} (p.~112)} }
\expandafter\def\csname SRef@maps:eq:D.OD1\endcsname{\href{S6_SERIES_GRIND_III.pdf\#nameddest=tag-0212}{III.{D.OD1} (p.~113)} }
\expandafter\def\csname SRef@maps:eq:D.OD2\endcsname{\href{S6_SERIES_GRIND_III.pdf\#nameddest=tag-0213}{III.{D.OD2} (p.~113)} }
\expandafter\def\csname SRef@maps:eq:D.OD3\endcsname{\href{S6_SERIES_GRIND_III.pdf\#nameddest=tag-0214}{III.{D.OD3} (p.~113)} }
\expandafter\def\csname SRef@maps:eq:D.OD4\endcsname{\href{S6_SERIES_GRIND_III.pdf\#nameddest=tag-0215}{III.{D.OD4} (p.~114)} }
\expandafter\def\csname SRef@maps:eq:D.OD5\endcsname{\href{S6_SERIES_GRIND_III.pdf\#nameddest=tag-0216}{III.{D.OD5} (p.~114)} }
\expandafter\def\csname SRef@maps:eq:D.OC1\endcsname{\href{S6_SERIES_GRIND_III.pdf\#nameddest=tag-020V}{III.{D.OC1} (p.~116)} }
\expandafter\def\csname SRef@maps:eq:D.OC2\endcsname{\href{S6_SERIES_GRIND_III.pdf\#nameddest=tag-020W}{III.{D.OC2} (p.~116)} }
\expandafter\def\csname SRef@maps:eq:D.OC3\endcsname{\href{S6_SERIES_GRIND_III.pdf\#nameddest=tag-020X}{III.{D.OC3} (p.~116)} }
\expandafter\def\csname SRef@maps:eq:D.OC4\endcsname{\href{S6_SERIES_GRIND_III.pdf\#nameddest=tag-020Y}{III.{D.OC4} (p.~117)} }
\expandafter\def\csname SRef@maps:eq:D.OC5\endcsname{\href{S6_SERIES_GRIND_III.pdf\#nameddest=tag-020Z}{III.{D.OC5} (p.~117)} }
\expandafter\def\csname SRef@maps:eq:D.OC6\endcsname{\href{S6_SERIES_GRIND_III.pdf\#nameddest=tag-0210}{III.{D.OC6} (p.~117)} }
\expandafter\def\csname SRef@maps:eq:D.OC7\endcsname{\href{S6_SERIES_GRIND_III.pdf\#nameddest=tag-0211}{III.{D.OC7} (p.~118)} }
\expandafter\def\csname SRef@maps:eq:D.OM1\endcsname{\href{S6_SERIES_GRIND_III.pdf\#nameddest=tag-0222}{III.{D.OM1} (p.~119)} }
\expandafter\def\csname SRef@maps:eq:D.OM2\endcsname{\href{S6_SERIES_GRIND_III.pdf\#nameddest=tag-0224}{III.{D.OM2} (p.~119)} }
\expandafter\def\csname SRef@maps:eq:D.OM3\endcsname{\href{S6_SERIES_GRIND_III.pdf\#nameddest=tag-0225}{III.{D.OM3} (p.~120)} }
\expandafter\def\csname SRef@maps:eq:D.OM4\endcsname{\href{S6_SERIES_GRIND_III.pdf\#nameddest=tag-0226}{III.{D.OM4} (p.~120)} }
\expandafter\def\csname SRef@maps:eq:D.OM5\endcsname{\href{S6_SERIES_GRIND_III.pdf\#nameddest=tag-0227}{III.{D.OM5} (p.~120)} }
\expandafter\def\csname SRef@maps:eq:D.OM6\endcsname{\href{S6_SERIES_GRIND_III.pdf\#nameddest=tag-0228}{III.{D.OM6} (p.~120)} }
\expandafter\def\csname SRef@maps:eq:D.OM7\endcsname{\href{S6_SERIES_GRIND_III.pdf\#nameddest=tag-0229}{III.{D.OM7} (p.~120)} }
\expandafter\def\csname SRef@maps:eq:D.OM8\endcsname{\href{S6_SERIES_GRIND_III.pdf\#nameddest=tag-022A}{III.{D.OM8} (p.~120)} }
\expandafter\def\csname SRef@maps:eq:D.OM9\endcsname{\href{S6_SERIES_GRIND_III.pdf\#nameddest=tag-022B}{III.{D.OM9} (p.~121)} }
\expandafter\def\csname SRef@maps:eq:B.1.1\endcsname{\href{S6_SERIES_GRIND_III.pdf\#nameddest=tag-01UV}{III.{B.1.1} (p.~14)} }
\expandafter\def\csname SRef@maps:eq:D.OM10\endcsname{\href{S6_SERIES_GRIND_III.pdf\#nameddest=tag-0223}{III.{D.OM10} (p.~121)} }
\expandafter\def\csname SRef@maps:eq:D.WE1\endcsname{\href{S6_SERIES_GRIND_III.pdf\#nameddest=tag-026O}{III.{D.WE1} (p.~122)} }
\expandafter\def\csname SRef@maps:eq:D.WE2\endcsname{\href{S6_SERIES_GRIND_III.pdf\#nameddest=tag-026U}{III.{D.WE2} (p.~122)} }
\expandafter\def\csname SRef@maps:eq:D.WE3\endcsname{\href{S6_SERIES_GRIND_III.pdf\#nameddest=tag-026V}{III.{D.WE3} (p.~122)} }
\expandafter\def\csname SRef@maps:eq:D.WE4\endcsname{\href{S6_SERIES_GRIND_III.pdf\#nameddest=tag-026W}{III.{D.WE4} (p.~122)} }
\expandafter\def\csname SRef@maps:eq:D.WE5\endcsname{\href{S6_SERIES_GRIND_III.pdf\#nameddest=tag-026X}{III.{D.WE5} (p.~122)} }
\expandafter\def\csname SRef@maps:eq:D.WE6\endcsname{\href{S6_SERIES_GRIND_III.pdf\#nameddest=tag-026Y}{III.{D.WE6} (p.~122)} }
\expandafter\def\csname SRef@maps:eq:D.WE7\endcsname{\href{S6_SERIES_GRIND_III.pdf\#nameddest=tag-026Z}{III.{D.WE7} (p.~123)} }
\expandafter\def\csname SRef@maps:eq:D.WE8\endcsname{\href{S6_SERIES_GRIND_III.pdf\#nameddest=tag-0270}{III.{D.WE8} (p.~123)} }
\expandafter\def\csname SRef@maps:eq:D.WE9\endcsname{\href{S6_SERIES_GRIND_III.pdf\#nameddest=tag-0271}{III.{D.WE9} (p.~123)} }
\expandafter\def\csname SRef@maps:eq:B.1.2\endcsname{\href{S6_SERIES_GRIND_III.pdf\#nameddest=tag-01UW}{III.{B.1.2} (p.~14)} }
\expandafter\def\csname SRef@maps:eq:D.WE10\endcsname{\href{S6_SERIES_GRIND_III.pdf\#nameddest=tag-026P}{III.{D.WE10} (p.~123)} }
\expandafter\def\csname SRef@maps:eq:D.WE11\endcsname{\href{S6_SERIES_GRIND_III.pdf\#nameddest=tag-026Q}{III.{D.WE11} (p.~124)} }
\expandafter\def\csname SRef@maps:eq:D.WE12\endcsname{\href{S6_SERIES_GRIND_III.pdf\#nameddest=tag-026R}{III.{D.WE12} (p.~125)} }
\expandafter\def\csname SRef@maps:eq:D.WE13\endcsname{\href{S6_SERIES_GRIND_III.pdf\#nameddest=tag-026S}{III.{D.WE13} (p.~125)} }
\expandafter\def\csname SRef@maps:eq:D.WE14\endcsname{\href{S6_SERIES_GRIND_III.pdf\#nameddest=tag-026T}{III.{D.WE14} (p.~126)} }
\expandafter\def\csname SRef@maps:eq:D.C21\endcsname{\href{S6_SERIES_GRIND_III.pdf\#nameddest=tag-01Y7}{III.{D.C21} (p.~126)} }
\expandafter\def\csname SRef@maps:eq:D.C22\endcsname{\href{S6_SERIES_GRIND_III.pdf\#nameddest=tag-01YF}{III.{D.C22} (p.~126)} }
\expandafter\def\csname SRef@maps:eq:D.C23\endcsname{\href{S6_SERIES_GRIND_III.pdf\#nameddest=tag-01YG}{III.{D.C23} (p.~127)} }
\expandafter\def\csname SRef@maps:eq:D.C24\endcsname{\href{S6_SERIES_GRIND_III.pdf\#nameddest=tag-01YH}{III.{D.C24} (p.~127)} }
\expandafter\def\csname SRef@maps:eq:D.C25\endcsname{\href{S6_SERIES_GRIND_III.pdf\#nameddest=tag-01YI}{III.{D.C25} (p.~127)} }
\expandafter\def\csname SRef@maps:eq:B.1.3\endcsname{\href{S6_SERIES_GRIND_III.pdf\#nameddest=tag-01UX}{III.{B.1.3} (p.~14)} }
\expandafter\def\csname SRef@maps:eq:D.C26\endcsname{\href{S6_SERIES_GRIND_III.pdf\#nameddest=tag-01YJ}{III.{D.C26} (p.~127)} }
\expandafter\def\csname SRef@maps:eq:D.C27\endcsname{\href{S6_SERIES_GRIND_III.pdf\#nameddest=tag-01YK}{III.{D.C27} (p.~127)} }
\expandafter\def\csname SRef@maps:eq:D.C28\endcsname{\href{S6_SERIES_GRIND_III.pdf\#nameddest=tag-01YL}{III.{D.C28} (p.~127)} }
\expandafter\def\csname SRef@maps:eq:D.C29\endcsname{\href{S6_SERIES_GRIND_III.pdf\#nameddest=tag-01YM}{III.{D.C29} (p.~128)} }
\expandafter\def\csname SRef@maps:eq:D.C210\endcsname{\href{S6_SERIES_GRIND_III.pdf\#nameddest=tag-01Y8}{III.{D.C210} (p.~128)} }
\expandafter\def\csname SRef@maps:eq:D.C211\endcsname{\href{S6_SERIES_GRIND_III.pdf\#nameddest=tag-01Y9}{III.{D.C211} (p.~128)} }
\expandafter\def\csname SRef@maps:eq:D.C212\endcsname{\href{S6_SERIES_GRIND_III.pdf\#nameddest=tag-01YA}{III.{D.C212} (p.~129)} }
\expandafter\def\csname SRef@maps:eq:D.C213\endcsname{\href{S6_SERIES_GRIND_III.pdf\#nameddest=tag-01YB}{III.{D.C213} (p.~129)} }
\expandafter\def\csname SRef@maps:eq:D.C214\endcsname{\href{S6_SERIES_GRIND_III.pdf\#nameddest=tag-01YC}{III.{D.C214} (p.~129)} }
\expandafter\def\csname SRef@maps:eq:D.C215\endcsname{\href{S6_SERIES_GRIND_III.pdf\#nameddest=tag-01YD}{III.{D.C215} (p.~129)} }
\expandafter\def\csname SRef@maps:eq:B.2.1\endcsname{\href{S6_SERIES_GRIND_III.pdf\#nameddest=tag-01UY}{III.{B.2.1} (p.~15)} }
\expandafter\def\csname SRef@maps:eq:D.C216\endcsname{\href{S6_SERIES_GRIND_III.pdf\#nameddest=tag-01YE}{III.{D.C216} (p.~130)} }
\expandafter\def\csname SRef@maps:eq:D.OP1\endcsname{\href{S6_SERIES_GRIND_III.pdf\#nameddest=tag-022C}{III.{D.OP1} (p.~131)} }
\expandafter\def\csname SRef@maps:eq:D.OP2\endcsname{\href{S6_SERIES_GRIND_III.pdf\#nameddest=tag-022F}{III.{D.OP2} (p.~131)} }
\expandafter\def\csname SRef@maps:eq:D.OP3\endcsname{\href{S6_SERIES_GRIND_III.pdf\#nameddest=tag-022G}{III.{D.OP3} (p.~131)} }
\expandafter\def\csname SRef@maps:eq:D.OP4\endcsname{\href{S6_SERIES_GRIND_III.pdf\#nameddest=tag-022H}{III.{D.OP4} (p.~131)} }
\expandafter\def\csname SRef@maps:eq:D.OP5\endcsname{\href{S6_SERIES_GRIND_III.pdf\#nameddest=tag-022I}{III.{D.OP5} (p.~132)} }
\expandafter\def\csname SRef@maps:eq:D.OP6\endcsname{\href{S6_SERIES_GRIND_III.pdf\#nameddest=tag-022K}{III.{D.OP6} (p.~132)} }
\expandafter\def\csname SRef@maps:eq:D.OP5a\endcsname{\href{S6_SERIES_GRIND_III.pdf\#nameddest=tag-022J}{III.{D.OP5a} (p.~132)} }
\expandafter\def\csname SRef@maps:eq:D.OP7\endcsname{\href{S6_SERIES_GRIND_III.pdf\#nameddest=tag-022L}{III.{D.OP7} (p.~133)} }
\expandafter\def\csname SRef@maps:eq:D.OP8\endcsname{\href{S6_SERIES_GRIND_III.pdf\#nameddest=tag-022M}{III.{D.OP8} (p.~133)} }
\expandafter\def\csname SRef@maps:eq:B.3.1\endcsname{\href{S6_SERIES_GRIND_III.pdf\#nameddest=tag-01UZ}{III.{B.3.1} (p.~16)} }
\expandafter\def\csname SRef@maps:eq:D.OP9\endcsname{\href{S6_SERIES_GRIND_III.pdf\#nameddest=tag-022N}{III.{D.OP9} (p.~133)} }
\expandafter\def\csname SRef@maps:eq:D.OP10\endcsname{\href{S6_SERIES_GRIND_III.pdf\#nameddest=tag-022D}{III.{D.OP10} (p.~133)} }
\expandafter\def\csname SRef@maps:eq:D.OP11\endcsname{\href{S6_SERIES_GRIND_III.pdf\#nameddest=tag-022E}{III.{D.OP11} (p.~133)} }
\expandafter\def\csname SRef@maps:eq:D.WP1\endcsname{\href{S6_SERIES_GRIND_III.pdf\#nameddest=tag-027O}{III.{D.WP1} (p.~134)} }
\expandafter\def\csname SRef@maps:eq:D.WP2\endcsname{\href{S6_SERIES_GRIND_III.pdf\#nameddest=tag-027S}{III.{D.WP2} (p.~134)} }
\expandafter\def\csname SRef@maps:eq:D.WP3\endcsname{\href{S6_SERIES_GRIND_III.pdf\#nameddest=tag-027T}{III.{D.WP3} (p.~134)} }
\expandafter\def\csname SRef@maps:eq:D.WP4\endcsname{\href{S6_SERIES_GRIND_III.pdf\#nameddest=tag-027U}{III.{D.WP4} (p.~134)} }
\expandafter\def\csname SRef@maps:eq:D.WP5\endcsname{\href{S6_SERIES_GRIND_III.pdf\#nameddest=tag-027V}{III.{D.WP5} (p.~135)} }
\expandafter\def\csname SRef@maps:eq:D.WP6\endcsname{\href{S6_SERIES_GRIND_III.pdf\#nameddest=tag-027W}{III.{D.WP6} (p.~135)} }
\expandafter\def\csname SRef@maps:eq:D.WP7\endcsname{\href{S6_SERIES_GRIND_III.pdf\#nameddest=tag-027X}{III.{D.WP7} (p.~135)} }
\expandafter\def\csname SRef@maps:eq:B.3.2\endcsname{\href{S6_SERIES_GRIND_III.pdf\#nameddest=tag-01V0}{III.{B.3.2} (p.~16)} }
\expandafter\def\csname SRef@maps:eq:D.WP8\endcsname{\href{S6_SERIES_GRIND_III.pdf\#nameddest=tag-027Y}{III.{D.WP8} (p.~135)} }
\expandafter\def\csname SRef@maps:eq:D.WP9\endcsname{\href{S6_SERIES_GRIND_III.pdf\#nameddest=tag-027Z}{III.{D.WP9} (p.~135)} }
\expandafter\def\csname SRef@maps:eq:D.WP10\endcsname{\href{S6_SERIES_GRIND_III.pdf\#nameddest=tag-027P}{III.{D.WP10} (p.~136)} }
\expandafter\def\csname SRef@maps:eq:D.WP11\endcsname{\href{S6_SERIES_GRIND_III.pdf\#nameddest=tag-027Q}{III.{D.WP11} (p.~137)} }
\expandafter\def\csname SRef@maps:eq:D.WP12\endcsname{\href{S6_SERIES_GRIND_III.pdf\#nameddest=tag-027R}{III.{D.WP12} (p.~138)} }
\expandafter\def\csname SRef@maps:eq:D.MP1\endcsname{\href{S6_SERIES_GRIND_III.pdf\#nameddest=tag-0209}{III.{D.MP1} (p.~138)} }
\expandafter\def\csname SRef@maps:eq:D.MP2\endcsname{\href{S6_SERIES_GRIND_III.pdf\#nameddest=tag-020B}{III.{D.MP2} (p.~139)} }
\expandafter\def\csname SRef@maps:eq:D.MP3\endcsname{\href{S6_SERIES_GRIND_III.pdf\#nameddest=tag-020C}{III.{D.MP3} (p.~139)} }
\expandafter\def\csname SRef@maps:eq:D.MP4\endcsname{\href{S6_SERIES_GRIND_III.pdf\#nameddest=tag-020D}{III.{D.MP4} (p.~139)} }
\expandafter\def\csname SRef@maps:eq:D.MP5\endcsname{\href{S6_SERIES_GRIND_III.pdf\#nameddest=tag-020E}{III.{D.MP5} (p.~140)} }
\expandafter\def\csname SRef@maps:eq:B.3.3\endcsname{\href{S6_SERIES_GRIND_III.pdf\#nameddest=tag-01V1}{III.{B.3.3} (p.~16)} }
\expandafter\def\csname SRef@maps:eq:D.MP6\endcsname{\href{S6_SERIES_GRIND_III.pdf\#nameddest=tag-020F}{III.{D.MP6} (p.~140)} }
\expandafter\def\csname SRef@maps:eq:D.MP7\endcsname{\href{S6_SERIES_GRIND_III.pdf\#nameddest=tag-020G}{III.{D.MP7} (p.~140)} }
\expandafter\def\csname SRef@maps:eq:D.MP8\endcsname{\href{S6_SERIES_GRIND_III.pdf\#nameddest=tag-020H}{III.{D.MP8} (p.~140)} }
\expandafter\def\csname SRef@maps:eq:D.MP9\endcsname{\href{S6_SERIES_GRIND_III.pdf\#nameddest=tag-020I}{III.{D.MP9} (p.~141)} }
\expandafter\def\csname SRef@maps:eq:D.MP10\endcsname{\href{S6_SERIES_GRIND_III.pdf\#nameddest=tag-020A}{III.{D.MP10} (p.~141)} }
\expandafter\def\csname SRef@maps:eq:D.T31\endcsname{\href{S6_SERIES_GRIND_III.pdf\#nameddest=tag-025C}{III.{D.T31} (p.~142)} }
\expandafter\def\csname SRef@maps:eq:D.T32\endcsname{\href{S6_SERIES_GRIND_III.pdf\#nameddest=tag-025D}{III.{D.T32} (p.~142)} }
\expandafter\def\csname SRef@maps:eq:D.T33\endcsname{\href{S6_SERIES_GRIND_III.pdf\#nameddest=tag-025E}{III.{D.T33} (p.~142)} }
\expandafter\def\csname SRef@maps:eq:D.T34\endcsname{\href{S6_SERIES_GRIND_III.pdf\#nameddest=tag-025F}{III.{D.T34} (p.~143)} }
\expandafter\def\csname SRef@maps:eq:D.T35\endcsname{\href{S6_SERIES_GRIND_III.pdf\#nameddest=tag-025G}{III.{D.T35} (p.~143)} }
\expandafter\def\csname SRef@maps:eq:B.3.4\endcsname{\href{S6_SERIES_GRIND_III.pdf\#nameddest=tag-01V2}{III.{B.3.4} (p.~16)} }
\expandafter\def\csname SRef@maps:eq:D.T36\endcsname{\href{S6_SERIES_GRIND_III.pdf\#nameddest=tag-025H}{III.{D.T36} (p.~143)} }
\expandafter\def\csname SRef@maps:eq:D.T37\endcsname{\href{S6_SERIES_GRIND_III.pdf\#nameddest=tag-025I}{III.{D.T37} (p.~144)} }
\expandafter\def\csname SRef@maps:eq:D.T38\endcsname{\href{S6_SERIES_GRIND_III.pdf\#nameddest=tag-025J}{III.{D.T38} (p.~144)} }
\expandafter\def\csname SRef@maps:eq:D.T39\endcsname{\href{S6_SERIES_GRIND_III.pdf\#nameddest=tag-025K}{III.{D.T39} (p.~144)} }
\expandafter\def\csname SRef@maps:eq:D.OW1\endcsname{\href{S6_SERIES_GRIND_III.pdf\#nameddest=tag-022W}{III.{D.OW1} (p.~145)} }
\expandafter\def\csname SRef@maps:eq:D.OW2\endcsname{\href{S6_SERIES_GRIND_III.pdf\#nameddest=tag-022X}{III.{D.OW2} (p.~145)} }
\expandafter\def\csname SRef@maps:eq:D.OW3\endcsname{\href{S6_SERIES_GRIND_III.pdf\#nameddest=tag-022Y}{III.{D.OW3} (p.~145)} }
\expandafter\def\csname SRef@maps:eq:D.OW4\endcsname{\href{S6_SERIES_GRIND_III.pdf\#nameddest=tag-022Z}{III.{D.OW4} (p.~145)} }
\expandafter\def\csname SRef@maps:eq:D.OW5\endcsname{\href{S6_SERIES_GRIND_III.pdf\#nameddest=tag-0230}{III.{D.OW5} (p.~145)} }
\expandafter\def\csname SRef@maps:eq:D.OW6\endcsname{\href{S6_SERIES_GRIND_III.pdf\#nameddest=tag-0231}{III.{D.OW6} (p.~146)} }
\expandafter\def\csname SRef@maps:eq:B.3.5\endcsname{\href{S6_SERIES_GRIND_III.pdf\#nameddest=tag-01V3}{III.{B.3.5} (p.~18)} }
\expandafter\def\csname SRef@maps:eq:D.OW7\endcsname{\href{S6_SERIES_GRIND_III.pdf\#nameddest=tag-0232}{III.{D.OW7} (p.~147)} }
\expandafter\def\csname SRef@maps:eq:D.BC1\endcsname{\href{S6_SERIES_GRIND_III.pdf\#nameddest=tag-01XR}{III.{D.BC1} (p.~148)} }
\expandafter\def\csname SRef@maps:eq:D.BC2\endcsname{\href{S6_SERIES_GRIND_III.pdf\#nameddest=tag-01XY}{III.{D.BC2} (p.~148)} }
\expandafter\def\csname SRef@maps:eq:D.BC3\endcsname{\href{S6_SERIES_GRIND_III.pdf\#nameddest=tag-01XZ}{III.{D.BC3} (p.~148)} }
\expandafter\def\csname SRef@maps:eq:D.BC4\endcsname{\href{S6_SERIES_GRIND_III.pdf\#nameddest=tag-01Y0}{III.{D.BC4} (p.~148)} }
\expandafter\def\csname SRef@maps:eq:D.BC5\endcsname{\href{S6_SERIES_GRIND_III.pdf\#nameddest=tag-01Y1}{III.{D.BC5} (p.~148)} }
\expandafter\def\csname SRef@maps:eq:D.BC6\endcsname{\href{S6_SERIES_GRIND_III.pdf\#nameddest=tag-01Y2}{III.{D.BC6} (p.~149)} }
\expandafter\def\csname SRef@maps:eq:D.BC7\endcsname{\href{S6_SERIES_GRIND_III.pdf\#nameddest=tag-01Y3}{III.{D.BC7} (p.~149)} }
\expandafter\def\csname SRef@maps:eq:D.BC8\endcsname{\href{S6_SERIES_GRIND_III.pdf\#nameddest=tag-01Y4}{III.{D.BC8} (p.~149)} }
\expandafter\def\csname SRef@maps:eq:D.BC9\endcsname{\href{S6_SERIES_GRIND_III.pdf\#nameddest=tag-01Y5}{III.{D.BC9} (p.~149)} }
\expandafter\def\csname SRef@maps:eq:B.4.1\endcsname{\href{S6_SERIES_GRIND_III.pdf\#nameddest=tag-01V4}{III.{B.4.1} (p.~19)} }
\expandafter\def\csname SRef@maps:eq:D.BC10\endcsname{\href{S6_SERIES_GRIND_III.pdf\#nameddest=tag-01XS}{III.{D.BC10} (p.~149)} }
\expandafter\def\csname SRef@maps:eq:D.BC11\endcsname{\href{S6_SERIES_GRIND_III.pdf\#nameddest=tag-01XT}{III.{D.BC11} (p.~149)} }
\expandafter\def\csname SRef@maps:eq:D.BC12\endcsname{\href{S6_SERIES_GRIND_III.pdf\#nameddest=tag-01XU}{III.{D.BC12} (p.~149)} }
\expandafter\def\csname SRef@maps:eq:D.BC13\endcsname{\href{S6_SERIES_GRIND_III.pdf\#nameddest=tag-01XV}{III.{D.BC13} (p.~150)} }
\expandafter\def\csname SRef@maps:eq:D.BC14\endcsname{\href{S6_SERIES_GRIND_III.pdf\#nameddest=tag-01XW}{III.{D.BC14} (p.~150)} }
\expandafter\def\csname SRef@maps:eq:D.BC15\endcsname{\href{S6_SERIES_GRIND_III.pdf\#nameddest=tag-01XX}{III.{D.BC15} (p.~150)} }
\expandafter\def\csname SRef@maps:eq:D.MT1\endcsname{\href{S6_SERIES_GRIND_III.pdf\#nameddest=tag-020J}{III.{D.MT1} (p.~151)} }
\expandafter\def\csname SRef@maps:eq:D.MT2\endcsname{\href{S6_SERIES_GRIND_III.pdf\#nameddest=tag-020K}{III.{D.MT2} (p.~152)} }
\expandafter\def\csname SRef@maps:eq:D.MT3\endcsname{\href{S6_SERIES_GRIND_III.pdf\#nameddest=tag-020L}{III.{D.MT3} (p.~153)} }
\expandafter\def\csname SRef@maps:eq:D.GC1\endcsname{\href{S6_SERIES_GRIND_III.pdf\#nameddest=tag-01YY}{III.{D.GC1} (p.~153)} }
\expandafter\def\csname SRef@maps:eq:B.4.2\endcsname{\href{S6_SERIES_GRIND_III.pdf\#nameddest=tag-01V5}{III.{B.4.2} (p.~19)} }
\expandafter\def\csname SRef@maps:eq:D.GC2\endcsname{\href{S6_SERIES_GRIND_III.pdf\#nameddest=tag-01Z4}{III.{D.GC2} (p.~153)} }
\expandafter\def\csname SRef@maps:eq:D.GC3\endcsname{\href{S6_SERIES_GRIND_III.pdf\#nameddest=tag-01Z5}{III.{D.GC3} (p.~154)} }
\expandafter\def\csname SRef@maps:eq:D.GC4\endcsname{\href{S6_SERIES_GRIND_III.pdf\#nameddest=tag-01Z6}{III.{D.GC4} (p.~154)} }
\expandafter\def\csname SRef@maps:eq:D.GC5\endcsname{\href{S6_SERIES_GRIND_III.pdf\#nameddest=tag-01Z7}{III.{D.GC5} (p.~154)} }
\expandafter\def\csname SRef@maps:eq:D.GC6\endcsname{\href{S6_SERIES_GRIND_III.pdf\#nameddest=tag-01Z8}{III.{D.GC6} (p.~154)} }
\expandafter\def\csname SRef@maps:eq:D.GC7\endcsname{\href{S6_SERIES_GRIND_III.pdf\#nameddest=tag-01Z9}{III.{D.GC7} (p.~154)} }
\expandafter\def\csname SRef@maps:eq:D.GC8\endcsname{\href{S6_SERIES_GRIND_III.pdf\#nameddest=tag-01ZA}{III.{D.GC8} (p.~155)} }
\expandafter\def\csname SRef@maps:eq:D.GC9\endcsname{\href{S6_SERIES_GRIND_III.pdf\#nameddest=tag-01ZB}{III.{D.GC9} (p.~155)} }
\expandafter\def\csname SRef@maps:eq:D.GC10\endcsname{\href{S6_SERIES_GRIND_III.pdf\#nameddest=tag-01YZ}{III.{D.GC10} (p.~155)} }
\expandafter\def\csname SRef@maps:eq:D.GC11\endcsname{\href{S6_SERIES_GRIND_III.pdf\#nameddest=tag-01Z0}{III.{D.GC11} (p.~155)} }
\expandafter\def\csname SRef@maps:eq:B.4.3\endcsname{\href{S6_SERIES_GRIND_III.pdf\#nameddest=tag-01V6}{III.{B.4.3} (p.~20)} }
\expandafter\def\csname SRef@maps:eq:D.GC12\endcsname{\href{S6_SERIES_GRIND_III.pdf\#nameddest=tag-01Z1}{III.{D.GC12} (p.~155)} }
\expandafter\def\csname SRef@maps:eq:D.GC13\endcsname{\href{S6_SERIES_GRIND_III.pdf\#nameddest=tag-01Z2}{III.{D.GC13} (p.~155)} }
\expandafter\def\csname SRef@maps:eq:D.GC14\endcsname{\href{S6_SERIES_GRIND_III.pdf\#nameddest=tag-01Z3}{III.{D.GC14} (p.~156)} }
\expandafter\def\csname SRef@maps:eq:D44.1\endcsname{\href{S6_SERIES_GRIND_III.pdf\#nameddest=tag-0282}{III.{D44.1} (p.~157)} }
\expandafter\def\csname SRef@maps:eq:D44.0a\endcsname{\href{S6_SERIES_GRIND_III.pdf\#nameddest=tag-0280}{III.{D44.0a} (p.~157)} }
\expandafter\def\csname SRef@maps:eq:D44.0b\endcsname{\href{S6_SERIES_GRIND_III.pdf\#nameddest=tag-0281}{III.{D44.0b} (p.~157)} }
\expandafter\def\csname SRef@maps:eq:D44.2\endcsname{\href{S6_SERIES_GRIND_III.pdf\#nameddest=tag-0288}{III.{D44.2} (p.~158)} }
\expandafter\def\csname SRef@maps:eq:D44.3\endcsname{\href{S6_SERIES_GRIND_III.pdf\#nameddest=tag-0289}{III.{D44.3} (p.~158)} }
\expandafter\def\csname SRef@maps:eq:D44.4\endcsname{\href{S6_SERIES_GRIND_III.pdf\#nameddest=tag-028A}{III.{D44.4} (p.~159)} }
\expandafter\def\csname SRef@maps:eq:D44.5\endcsname{\href{S6_SERIES_GRIND_III.pdf\#nameddest=tag-028B}{III.{D44.5} (p.~159)} }
\expandafter\def\csname SRef@maps:eq:B.5.1\endcsname{\href{S6_SERIES_GRIND_III.pdf\#nameddest=tag-01V7}{III.{B.5.1} (p.~20)} }
\expandafter\def\csname SRef@maps:eq:D44.6\endcsname{\href{S6_SERIES_GRIND_III.pdf\#nameddest=tag-028C}{III.{D44.6} (p.~159)} }
\expandafter\def\csname SRef@maps:eq:D44.7\endcsname{\href{S6_SERIES_GRIND_III.pdf\#nameddest=tag-028D}{III.{D44.7} (p.~159)} }
\expandafter\def\csname SRef@maps:eq:D44.8\endcsname{\href{S6_SERIES_GRIND_III.pdf\#nameddest=tag-028E}{III.{D44.8} (p.~160)} }
\expandafter\def\csname SRef@maps:eq:D44.9\endcsname{\href{S6_SERIES_GRIND_III.pdf\#nameddest=tag-028F}{III.{D44.9} (p.~160)} }
\expandafter\def\csname SRef@maps:eq:D44.10\endcsname{\href{S6_SERIES_GRIND_III.pdf\#nameddest=tag-0283}{III.{D44.10} (p.~160)} }
\expandafter\def\csname SRef@maps:eq:D44.11\endcsname{\href{S6_SERIES_GRIND_III.pdf\#nameddest=tag-0284}{III.{D44.11} (p.~160)} }
\expandafter\def\csname SRef@maps:eq:D44.12\endcsname{\href{S6_SERIES_GRIND_III.pdf\#nameddest=tag-0285}{III.{D44.12} (p.~160)} }
\expandafter\def\csname SRef@maps:eq:D44.13\endcsname{\href{S6_SERIES_GRIND_III.pdf\#nameddest=tag-0286}{III.{D44.13} (p.~161)} }
\expandafter\def\csname SRef@maps:eq:D44.14\endcsname{\href{S6_SERIES_GRIND_III.pdf\#nameddest=tag-0287}{III.{D44.14} (p.~161)} }
\expandafter\def\csname SRef@maps:eq:D45.1\endcsname{\href{S6_SERIES_GRIND_III.pdf\#nameddest=tag-028G}{III.{D45.1} (p.~162)} }
\expandafter\def\csname SRef@maps:eq:B.5.2\endcsname{\href{S6_SERIES_GRIND_III.pdf\#nameddest=tag-01VH}{III.{B.5.2} (p.~21)} }
\expandafter\def\csname SRef@maps:eq:D45.2\endcsname{\href{S6_SERIES_GRIND_III.pdf\#nameddest=tag-028R}{III.{D45.2} (p.~162)} }
\expandafter\def\csname SRef@maps:eq:D45.3\endcsname{\href{S6_SERIES_GRIND_III.pdf\#nameddest=tag-028V}{III.{D45.3} (p.~162)} }
\expandafter\def\csname SRef@maps:eq:D45.4\endcsname{\href{S6_SERIES_GRIND_III.pdf\#nameddest=tag-028W}{III.{D45.4} (p.~162)} }
\expandafter\def\csname SRef@maps:eq:D45.5\endcsname{\href{S6_SERIES_GRIND_III.pdf\#nameddest=tag-028X}{III.{D45.5} (p.~162)} }
\expandafter\def\csname SRef@maps:eq:D45.6\endcsname{\href{S6_SERIES_GRIND_III.pdf\#nameddest=tag-028Y}{III.{D45.6} (p.~163)} }
\expandafter\def\csname SRef@maps:eq:D45.7\endcsname{\href{S6_SERIES_GRIND_III.pdf\#nameddest=tag-028Z}{III.{D45.7} (p.~163)} }
\expandafter\def\csname SRef@maps:eq:D45.8\endcsname{\href{S6_SERIES_GRIND_III.pdf\#nameddest=tag-0290}{III.{D45.8} (p.~163)} }
\expandafter\def\csname SRef@maps:eq:D45.9\endcsname{\href{S6_SERIES_GRIND_III.pdf\#nameddest=tag-0291}{III.{D45.9} (p.~163)} }
\expandafter\def\csname SRef@maps:eq:D45.10\endcsname{\href{S6_SERIES_GRIND_III.pdf\#nameddest=tag-028H}{III.{D45.10} (p.~164)} }
\expandafter\def\csname SRef@maps:eq:D45.11\endcsname{\href{S6_SERIES_GRIND_III.pdf\#nameddest=tag-028I}{III.{D45.11} (p.~164)} }
\expandafter\def\csname SRef@maps:eq:B.5.3\endcsname{\href{S6_SERIES_GRIND_III.pdf\#nameddest=tag-01VI}{III.{B.5.3} (p.~21)} }
\expandafter\def\csname SRef@maps:eq:D45.12\endcsname{\href{S6_SERIES_GRIND_III.pdf\#nameddest=tag-028J}{III.{D45.12} (p.~164)} }
\expandafter\def\csname SRef@maps:eq:D45.13\endcsname{\href{S6_SERIES_GRIND_III.pdf\#nameddest=tag-028K}{III.{D45.13} (p.~164)} }
\expandafter\def\csname SRef@maps:eq:D45.14\endcsname{\href{S6_SERIES_GRIND_III.pdf\#nameddest=tag-028L}{III.{D45.14} (p.~164)} }
\expandafter\def\csname SRef@maps:eq:D45.15\endcsname{\href{S6_SERIES_GRIND_III.pdf\#nameddest=tag-028M}{III.{D45.15} (p.~164)} }
\expandafter\def\csname SRef@maps:eq:D45.16\endcsname{\href{S6_SERIES_GRIND_III.pdf\#nameddest=tag-028N}{III.{D45.16} (p.~165)} }
\expandafter\def\csname SRef@maps:eq:D45.17\endcsname{\href{S6_SERIES_GRIND_III.pdf\#nameddest=tag-028O}{III.{D45.17} (p.~165)} }
\expandafter\def\csname SRef@maps:eq:D45.18\endcsname{\href{S6_SERIES_GRIND_III.pdf\#nameddest=tag-028P}{III.{D45.18} (p.~165)} }
\expandafter\def\csname SRef@maps:eq:D45.19\endcsname{\href{S6_SERIES_GRIND_III.pdf\#nameddest=tag-028Q}{III.{D45.19} (p.~165)} }
\expandafter\def\csname SRef@maps:eq:D45.20\endcsname{\href{S6_SERIES_GRIND_III.pdf\#nameddest=tag-028S}{III.{D45.20} (p.~165)} }
\expandafter\def\csname SRef@maps:eq:D45.21\endcsname{\href{S6_SERIES_GRIND_III.pdf\#nameddest=tag-028T}{III.{D45.21} (p.~166)} }
\expandafter\def\csname SRef@maps:eq:B.5.4\endcsname{\href{S6_SERIES_GRIND_III.pdf\#nameddest=tag-01VJ}{III.{B.5.4} (p.~21)} }
\expandafter\def\csname SRef@maps:eq:D45.22\endcsname{\href{S6_SERIES_GRIND_III.pdf\#nameddest=tag-028U}{III.{D45.22} (p.~166)} }
\expandafter\def\csname SRef@maps:eq:D.GI1\endcsname{\href{S6_SERIES_GRIND_III.pdf\#nameddest=tag-01ZC}{III.{D.GI1} (p.~167)} }
\expandafter\def\csname SRef@maps:eq:D.GI2\endcsname{\href{S6_SERIES_GRIND_III.pdf\#nameddest=tag-01ZF}{III.{D.GI2} (p.~167)} }
\expandafter\def\csname SRef@maps:eq:D.GI3\endcsname{\href{S6_SERIES_GRIND_III.pdf\#nameddest=tag-01ZG}{III.{D.GI3} (p.~167)} }
\expandafter\def\csname SRef@maps:eq:D.GI4\endcsname{\href{S6_SERIES_GRIND_III.pdf\#nameddest=tag-01ZH}{III.{D.GI4} (p.~167)} }
\expandafter\def\csname SRef@maps:eq:D.GI5\endcsname{\href{S6_SERIES_GRIND_III.pdf\#nameddest=tag-01ZI}{III.{D.GI5} (p.~168)} }
\expandafter\def\csname SRef@maps:eq:D.GI6\endcsname{\href{S6_SERIES_GRIND_III.pdf\#nameddest=tag-01ZJ}{III.{D.GI6} (p.~168)} }
\expandafter\def\csname SRef@maps:eq:D.GI7\endcsname{\href{S6_SERIES_GRIND_III.pdf\#nameddest=tag-01ZK}{III.{D.GI7} (p.~168)} }
\expandafter\def\csname SRef@maps:eq:D.GI8\endcsname{\href{S6_SERIES_GRIND_III.pdf\#nameddest=tag-01ZL}{III.{D.GI8} (p.~168)} }
\expandafter\def\csname SRef@maps:eq:D.GI9\endcsname{\href{S6_SERIES_GRIND_III.pdf\#nameddest=tag-01ZM}{III.{D.GI9} (p.~170)} }
\expandafter\def\csname SRef@maps:eq:B.5.5\endcsname{\href{S6_SERIES_GRIND_III.pdf\#nameddest=tag-01VK}{III.{B.5.5} (p.~22)} }
\expandafter\def\csname SRef@maps:eq:D.GI10\endcsname{\href{S6_SERIES_GRIND_III.pdf\#nameddest=tag-01ZD}{III.{D.GI10} (p.~171)} }
\expandafter\def\csname SRef@maps:eq:D.GI11\endcsname{\href{S6_SERIES_GRIND_III.pdf\#nameddest=tag-01ZE}{III.{D.GI11} (p.~171)} }
\expandafter\def\csname SRef@maps:eq:D47.1\endcsname{\href{S6_SERIES_GRIND_III.pdf\#nameddest=tag-0292}{III.{D47.1} (p.~172)} }
\expandafter\def\csname SRef@maps:eq:D47.2\endcsname{\href{S6_SERIES_GRIND_III.pdf\#nameddest=tag-0297}{III.{D47.2} (p.~172)} }
\expandafter\def\csname SRef@maps:eq:D47.3\endcsname{\href{S6_SERIES_GRIND_III.pdf\#nameddest=tag-0298}{III.{D47.3} (p.~172)} }
\expandafter\def\csname SRef@maps:eq:D47.4\endcsname{\href{S6_SERIES_GRIND_III.pdf\#nameddest=tag-0299}{III.{D47.4} (p.~172)} }
\expandafter\def\csname SRef@maps:eq:D47.5\endcsname{\href{S6_SERIES_GRIND_III.pdf\#nameddest=tag-029A}{III.{D47.5} (p.~173)} }
\expandafter\def\csname SRef@maps:eq:D47.6\endcsname{\href{S6_SERIES_GRIND_III.pdf\#nameddest=tag-029B}{III.{D47.6} (p.~173)} }
\expandafter\def\csname SRef@maps:eq:D47.7\endcsname{\href{S6_SERIES_GRIND_III.pdf\#nameddest=tag-029C}{III.{D47.7} (p.~173)} }
\expandafter\def\csname SRef@maps:eq:D47.8\endcsname{\href{S6_SERIES_GRIND_III.pdf\#nameddest=tag-029D}{III.{D47.8} (p.~173)} }
\expandafter\def\csname SRef@maps:eq:B.5.6\endcsname{\href{S6_SERIES_GRIND_III.pdf\#nameddest=tag-01VL}{III.{B.5.6} (p.~22)} }
\expandafter\def\csname SRef@maps:eq:D47.9\endcsname{\href{S6_SERIES_GRIND_III.pdf\#nameddest=tag-029E}{III.{D47.9} (p.~173)} }
\expandafter\def\csname SRef@maps:eq:D47.10\endcsname{\href{S6_SERIES_GRIND_III.pdf\#nameddest=tag-0293}{III.{D47.10} (p.~174)} }
\expandafter\def\csname SRef@maps:eq:D47.11\endcsname{\href{S6_SERIES_GRIND_III.pdf\#nameddest=tag-0294}{III.{D47.11} (p.~174)} }
\expandafter\def\csname SRef@maps:eq:D47.12\endcsname{\href{S6_SERIES_GRIND_III.pdf\#nameddest=tag-0295}{III.{D47.12} (p.~174)} }
\expandafter\def\csname SRef@maps:eq:D47.13\endcsname{\href{S6_SERIES_GRIND_III.pdf\#nameddest=tag-0296}{III.{D47.13} (p.~174)} }
\expandafter\def\csname SRef@maps:eq:D.WI1\endcsname{\href{S6_SERIES_GRIND_III.pdf\#nameddest=tag-0272}{III.{D.WI1} (p.~175)} }
\expandafter\def\csname SRef@maps:eq:D.WI2\endcsname{\href{S6_SERIES_GRIND_III.pdf\#nameddest=tag-027D}{III.{D.WI2} (p.~175)} }
\expandafter\def\csname SRef@maps:eq:D.WI3\endcsname{\href{S6_SERIES_GRIND_III.pdf\#nameddest=tag-027H}{III.{D.WI3} (p.~175)} }
\expandafter\def\csname SRef@maps:eq:D.WI14\endcsname{\href{S6_SERIES_GRIND_III.pdf\#nameddest=tag-0277}{III.{D.WI14} (p.~175)} }
\expandafter\def\csname SRef@maps:eq:D.WI15\endcsname{\href{S6_SERIES_GRIND_III.pdf\#nameddest=tag-0278}{III.{D.WI15} (p.~175)} }
\expandafter\def\csname SRef@maps:eq:B.5.7\endcsname{\href{S6_SERIES_GRIND_III.pdf\#nameddest=tag-01VM}{III.{B.5.7} (p.~22)} }
\expandafter\def\csname SRef@maps:eq:D.WI16\endcsname{\href{S6_SERIES_GRIND_III.pdf\#nameddest=tag-0279}{III.{D.WI16} (p.~175)} }
\expandafter\def\csname SRef@maps:eq:D.WI4\endcsname{\href{S6_SERIES_GRIND_III.pdf\#nameddest=tag-027I}{III.{D.WI4} (p.~176)} }
\expandafter\def\csname SRef@maps:eq:D.WI5\endcsname{\href{S6_SERIES_GRIND_III.pdf\#nameddest=tag-027J}{III.{D.WI5} (p.~176)} }
\expandafter\def\csname SRef@maps:eq:D.WI6\endcsname{\href{S6_SERIES_GRIND_III.pdf\#nameddest=tag-027K}{III.{D.WI6} (p.~177)} }
\expandafter\def\csname SRef@maps:eq:D.WI7\endcsname{\href{S6_SERIES_GRIND_III.pdf\#nameddest=tag-027L}{III.{D.WI7} (p.~177)} }
\expandafter\def\csname SRef@maps:eq:D.WI8\endcsname{\href{S6_SERIES_GRIND_III.pdf\#nameddest=tag-027M}{III.{D.WI8} (p.~177)} }
\expandafter\def\csname SRef@maps:eq:D.WI9\endcsname{\href{S6_SERIES_GRIND_III.pdf\#nameddest=tag-027N}{III.{D.WI9} (p.~177)} }
\expandafter\def\csname SRef@maps:eq:D.WI10\endcsname{\href{S6_SERIES_GRIND_III.pdf\#nameddest=tag-0273}{III.{D.WI10} (p.~177)} }
\expandafter\def\csname SRef@maps:eq:D.WI11\endcsname{\href{S6_SERIES_GRIND_III.pdf\#nameddest=tag-0274}{III.{D.WI11} (p.~177)} }
\expandafter\def\csname SRef@maps:eq:D.WI12\endcsname{\href{S6_SERIES_GRIND_III.pdf\#nameddest=tag-0275}{III.{D.WI12} (p.~177)} }
\expandafter\def\csname SRef@maps:eq:B.5.8\endcsname{\href{S6_SERIES_GRIND_III.pdf\#nameddest=tag-01VN}{III.{B.5.8} (p.~23)} }
\expandafter\def\csname SRef@maps:eq:D.WI13\endcsname{\href{S6_SERIES_GRIND_III.pdf\#nameddest=tag-0276}{III.{D.WI13} (p.~178)} }
\expandafter\def\csname SRef@maps:eq:D.WI17\endcsname{\href{S6_SERIES_GRIND_III.pdf\#nameddest=tag-027A}{III.{D.WI17} (p.~178)} }
\expandafter\def\csname SRef@maps:eq:D.WI18\endcsname{\href{S6_SERIES_GRIND_III.pdf\#nameddest=tag-027B}{III.{D.WI18} (p.~178)} }
\expandafter\def\csname SRef@maps:eq:D.WI19\endcsname{\href{S6_SERIES_GRIND_III.pdf\#nameddest=tag-027C}{III.{D.WI19} (p.~178)} }
\expandafter\def\csname SRef@maps:eq:D.WI20\endcsname{\href{S6_SERIES_GRIND_III.pdf\#nameddest=tag-027E}{III.{D.WI20} (p.~179)} }
\expandafter\def\csname SRef@maps:eq:D.WI21\endcsname{\href{S6_SERIES_GRIND_III.pdf\#nameddest=tag-027F}{III.{D.WI21} (p.~179)} }
\expandafter\def\csname SRef@maps:eq:D.TP1\endcsname{\href{S6_SERIES_GRIND_III.pdf\#nameddest=tag-025R}{III.{D.TP1} (p.~179)} }
\expandafter\def\csname SRef@maps:eq:D.TP2\endcsname{\href{S6_SERIES_GRIND_III.pdf\#nameddest=tag-025S}{III.{D.TP2} (p.~180)} }
\expandafter\def\csname SRef@maps:eq:D.TP3\endcsname{\href{S6_SERIES_GRIND_III.pdf\#nameddest=tag-025T}{III.{D.TP3} (p.~180)} }
\expandafter\def\csname SRef@maps:eq:B.5.9\endcsname{\href{S6_SERIES_GRIND_III.pdf\#nameddest=tag-01VO}{III.{B.5.9} (p.~24)} }
\expandafter\def\csname SRef@maps:eq:D.TP4\endcsname{\href{S6_SERIES_GRIND_III.pdf\#nameddest=tag-025U}{III.{D.TP4} (p.~180)} }
\expandafter\def\csname SRef@maps:eq:D.TP5\endcsname{\href{S6_SERIES_GRIND_III.pdf\#nameddest=tag-025V}{III.{D.TP5} (p.~180)} }
\expandafter\def\csname SRef@maps:eq:D.TP6\endcsname{\href{S6_SERIES_GRIND_III.pdf\#nameddest=tag-025W}{III.{D.TP6} (p.~180)} }
\expandafter\def\csname SRef@maps:eq:D.TP7\endcsname{\href{S6_SERIES_GRIND_III.pdf\#nameddest=tag-025X}{III.{D.TP7} (p.~180)} }
\expandafter\def\csname SRef@maps:eq:D.TP8\endcsname{\href{S6_SERIES_GRIND_III.pdf\#nameddest=tag-025Y}{III.{D.TP8} (p.~180)} }
\expandafter\def\csname SRef@maps:eq:D.WI22\endcsname{\href{S6_SERIES_GRIND_III.pdf\#nameddest=tag-027G}{III.{D.WI22} (p.~181)} }
\expandafter\def\csname SRef@maps:eq:D.SL1\endcsname{\href{S6_SERIES_GRIND_III.pdf\#nameddest=tag-0247}{III.{D.SL1} (p.~182)} }
\expandafter\def\csname SRef@maps:eq:D.SL2\endcsname{\href{S6_SERIES_GRIND_III.pdf\#nameddest=tag-024I}{III.{D.SL2} (p.~182)} }
\expandafter\def\csname SRef@maps:eq:D.SL3\endcsname{\href{S6_SERIES_GRIND_III.pdf\#nameddest=tag-024O}{III.{D.SL3} (p.~182)} }
\expandafter\def\csname SRef@maps:eq:D.SL4\endcsname{\href{S6_SERIES_GRIND_III.pdf\#nameddest=tag-024P}{III.{D.SL4} (p.~182)} }
\expandafter\def\csname SRef@maps:eq:B.5.10\endcsname{\href{S6_SERIES_GRIND_III.pdf\#nameddest=tag-01V8}{III.{B.5.10} (p.~24)} }
\expandafter\def\csname SRef@maps:eq:D.SL5\endcsname{\href{S6_SERIES_GRIND_III.pdf\#nameddest=tag-024Q}{III.{D.SL5} (p.~182)} }
\expandafter\def\csname SRef@maps:eq:D.SL6\endcsname{\href{S6_SERIES_GRIND_III.pdf\#nameddest=tag-024R}{III.{D.SL6} (p.~183)} }
\expandafter\def\csname SRef@maps:eq:D.SL7\endcsname{\href{S6_SERIES_GRIND_III.pdf\#nameddest=tag-024S}{III.{D.SL7} (p.~183)} }
\expandafter\def\csname SRef@maps:eq:D.SL8\endcsname{\href{S6_SERIES_GRIND_III.pdf\#nameddest=tag-024T}{III.{D.SL8} (p.~183)} }
\expandafter\def\csname SRef@maps:eq:D.SL9\endcsname{\href{S6_SERIES_GRIND_III.pdf\#nameddest=tag-024U}{III.{D.SL9} (p.~183)} }
\expandafter\def\csname SRef@maps:eq:D.SL10\endcsname{\href{S6_SERIES_GRIND_III.pdf\#nameddest=tag-0248}{III.{D.SL10} (p.~183)} }
\expandafter\def\csname SRef@maps:eq:D.SL11\endcsname{\href{S6_SERIES_GRIND_III.pdf\#nameddest=tag-0249}{III.{D.SL11} (p.~183)} }
\expandafter\def\csname SRef@maps:eq:D.SL12\endcsname{\href{S6_SERIES_GRIND_III.pdf\#nameddest=tag-024A}{III.{D.SL12} (p.~184)} }
\expandafter\def\csname SRef@maps:eq:D.SL13\endcsname{\href{S6_SERIES_GRIND_III.pdf\#nameddest=tag-024B}{III.{D.SL13} (p.~184)} }
\expandafter\def\csname SRef@maps:eq:D.SL14\endcsname{\href{S6_SERIES_GRIND_III.pdf\#nameddest=tag-024C}{III.{D.SL14} (p.~184)} }
\expandafter\def\csname SRef@maps:eq:B.5.11\endcsname{\href{S6_SERIES_GRIND_III.pdf\#nameddest=tag-01V9}{III.{B.5.11} (p.~24)} }
\expandafter\def\csname SRef@maps:eq:D.SL15\endcsname{\href{S6_SERIES_GRIND_III.pdf\#nameddest=tag-024D}{III.{D.SL15} (p.~185)} }
\expandafter\def\csname SRef@maps:eq:D.SL16\endcsname{\href{S6_SERIES_GRIND_III.pdf\#nameddest=tag-024E}{III.{D.SL16} (p.~185)} }
\expandafter\def\csname SRef@maps:eq:D.SL17\endcsname{\href{S6_SERIES_GRIND_III.pdf\#nameddest=tag-024F}{III.{D.SL17} (p.~185)} }
\expandafter\def\csname SRef@maps:eq:D.SL18\endcsname{\href{S6_SERIES_GRIND_III.pdf\#nameddest=tag-024G}{III.{D.SL18} (p.~185)} }
\expandafter\def\csname SRef@maps:eq:D.SL19\endcsname{\href{S6_SERIES_GRIND_III.pdf\#nameddest=tag-024H}{III.{D.SL19} (p.~185)} }
\expandafter\def\csname SRef@maps:eq:D.SL20\endcsname{\href{S6_SERIES_GRIND_III.pdf\#nameddest=tag-024J}{III.{D.SL20} (p.~186)} }
\expandafter\def\csname SRef@maps:eq:D.SL21\endcsname{\href{S6_SERIES_GRIND_III.pdf\#nameddest=tag-024K}{III.{D.SL21} (p.~186)} }
\expandafter\def\csname SRef@maps:eq:D.SL22\endcsname{\href{S6_SERIES_GRIND_III.pdf\#nameddest=tag-024L}{III.{D.SL22} (p.~186)} }
\expandafter\def\csname SRef@maps:eq:D.SL24\endcsname{\href{S6_SERIES_GRIND_III.pdf\#nameddest=tag-024N}{III.{D.SL24} (p.~186)} }
\expandafter\def\csname SRef@maps:eq:D.SL23\endcsname{\href{S6_SERIES_GRIND_III.pdf\#nameddest=tag-024M}{III.{D.SL23} (p.~187)} }
\expandafter\def\csname SRef@maps:eq:B.5.12\endcsname{\href{S6_SERIES_GRIND_III.pdf\#nameddest=tag-01VA}{III.{B.5.12} (p.~24)} }
\expandafter\def\csname SRef@maps:eq:D.VC1\endcsname{\href{S6_SERIES_GRIND_III.pdf\#nameddest=tag-026E}{III.{D.VC1} (p.~187)} }
\expandafter\def\csname SRef@maps:eq:D.VC7\endcsname{\href{S6_SERIES_GRIND_III.pdf\#nameddest=tag-026L}{III.{D.VC7} (p.~188)} }
\expandafter\def\csname SRef@maps:eq:D.VC2\endcsname{\href{S6_SERIES_GRIND_III.pdf\#nameddest=tag-026G}{III.{D.VC2} (p.~188)} }
\expandafter\def\csname SRef@maps:eq:D.VC3\endcsname{\href{S6_SERIES_GRIND_III.pdf\#nameddest=tag-026H}{III.{D.VC3} (p.~188)} }
\expandafter\def\csname SRef@maps:eq:D.VC4\endcsname{\href{S6_SERIES_GRIND_III.pdf\#nameddest=tag-026I}{III.{D.VC4} (p.~188)} }
\expandafter\def\csname SRef@maps:eq:D.VC5\endcsname{\href{S6_SERIES_GRIND_III.pdf\#nameddest=tag-026J}{III.{D.VC5} (p.~188)} }
\expandafter\def\csname SRef@maps:eq:D.VC8\endcsname{\href{S6_SERIES_GRIND_III.pdf\#nameddest=tag-026M}{III.{D.VC8} (p.~189)} }
\expandafter\def\csname SRef@maps:eq:D.VC6\endcsname{\href{S6_SERIES_GRIND_III.pdf\#nameddest=tag-026K}{III.{D.VC6} (p.~189)} }
\expandafter\def\csname SRef@maps:eq:D.VC9\endcsname{\href{S6_SERIES_GRIND_III.pdf\#nameddest=tag-026N}{III.{D.VC9} (p.~190)} }
\expandafter\def\csname SRef@maps:eq:D.VC10\endcsname{\href{S6_SERIES_GRIND_III.pdf\#nameddest=tag-026F}{III.{D.VC10} (p.~190)} }
\expandafter\def\csname SRef@maps:eq:B.5.13\endcsname{\href{S6_SERIES_GRIND_III.pdf\#nameddest=tag-01VB}{III.{B.5.13} (p.~25)} }
\expandafter\def\csname SRef@maps:eq:D.IC1\endcsname{\href{S6_SERIES_GRIND_III.pdf\#nameddest=tag-01ZY}{III.{D.IC1} (p.~191)} }
\expandafter\def\csname SRef@maps:eq:D.IC2\endcsname{\href{S6_SERIES_GRIND_III.pdf\#nameddest=tag-01ZZ}{III.{D.IC2} (p.~191)} }
\expandafter\def\csname SRef@maps:eq:D.IC3\endcsname{\href{S6_SERIES_GRIND_III.pdf\#nameddest=tag-0200}{III.{D.IC3} (p.~191)} }
\expandafter\def\csname SRef@maps:eq:D.IC4\endcsname{\href{S6_SERIES_GRIND_III.pdf\#nameddest=tag-0201}{III.{D.IC4} (p.~191)} }
\expandafter\def\csname SRef@maps:eq:D.IC5\endcsname{\href{S6_SERIES_GRIND_III.pdf\#nameddest=tag-0202}{III.{D.IC5} (p.~191)} }
\expandafter\def\csname SRef@maps:eq:D.IC6\endcsname{\href{S6_SERIES_GRIND_III.pdf\#nameddest=tag-0203}{III.{D.IC6} (p.~191)} }
\expandafter\def\csname SRef@maps:eq:D.IC7\endcsname{\href{S6_SERIES_GRIND_III.pdf\#nameddest=tag-0204}{III.{D.IC7} (p.~191)} }
\expandafter\def\csname SRef@maps:eq:D.IC8\endcsname{\href{S6_SERIES_GRIND_III.pdf\#nameddest=tag-0205}{III.{D.IC8} (p.~192)} }
\expandafter\def\csname SRef@maps:eq:D.IC9\endcsname{\href{S6_SERIES_GRIND_III.pdf\#nameddest=tag-0206}{III.{D.IC9} (p.~192)} }
\expandafter\def\csname SRef@maps:eq:D.CI1\endcsname{\href{S6_SERIES_GRIND_III.pdf\#nameddest=tag-01YN}{III.{D.CI1} (p.~195)} }
\expandafter\def\csname SRef@maps:eq:B.5.14\endcsname{\href{S6_SERIES_GRIND_III.pdf\#nameddest=tag-01VC}{III.{B.5.14} (p.~25)} }
\expandafter\def\csname SRef@maps:eq:D.OS1\endcsname{\href{S6_SERIES_GRIND_III.pdf\#nameddest=tag-022O}{III.{D.OS1} (p.~196)} }
\expandafter\def\csname SRef@maps:eq:D.OS2\endcsname{\href{S6_SERIES_GRIND_III.pdf\#nameddest=tag-022P}{III.{D.OS2} (p.~196)} }
\expandafter\def\csname SRef@maps:eq:D.OS3\endcsname{\href{S6_SERIES_GRIND_III.pdf\#nameddest=tag-022Q}{III.{D.OS3} (p.~196)} }
\expandafter\def\csname SRef@maps:eq:D.OS4\endcsname{\href{S6_SERIES_GRIND_III.pdf\#nameddest=tag-022R}{III.{D.OS4} (p.~196)} }
\expandafter\def\csname SRef@maps:eq:D.PC1\endcsname{\href{S6_SERIES_GRIND_III.pdf\#nameddest=tag-0233}{III.{D.PC1} (p.~196)} }
\expandafter\def\csname SRef@maps:eq:D.PC2\endcsname{\href{S6_SERIES_GRIND_III.pdf\#nameddest=tag-0239}{III.{D.PC2} (p.~197)} }
\expandafter\def\csname SRef@maps:eq:D.PC3\endcsname{\href{S6_SERIES_GRIND_III.pdf\#nameddest=tag-023A}{III.{D.PC3} (p.~197)} }
\expandafter\def\csname SRef@maps:eq:D.PC4\endcsname{\href{S6_SERIES_GRIND_III.pdf\#nameddest=tag-023B}{III.{D.PC4} (p.~197)} }
\expandafter\def\csname SRef@maps:eq:D.PC5\endcsname{\href{S6_SERIES_GRIND_III.pdf\#nameddest=tag-023C}{III.{D.PC5} (p.~197)} }
\expandafter\def\csname SRef@maps:eq:D.PC6\endcsname{\href{S6_SERIES_GRIND_III.pdf\#nameddest=tag-023D}{III.{D.PC6} (p.~197)} }
\expandafter\def\csname SRef@maps:eq:B.5.15\endcsname{\href{S6_SERIES_GRIND_III.pdf\#nameddest=tag-01VD}{III.{B.5.15} (p.~25)} }
\expandafter\def\csname SRef@maps:eq:D.PC7\endcsname{\href{S6_SERIES_GRIND_III.pdf\#nameddest=tag-023E}{III.{D.PC7} (p.~197)} }
\expandafter\def\csname SRef@maps:eq:D.PC8\endcsname{\href{S6_SERIES_GRIND_III.pdf\#nameddest=tag-023F}{III.{D.PC8} (p.~197)} }
\expandafter\def\csname SRef@maps:eq:D.PC9\endcsname{\href{S6_SERIES_GRIND_III.pdf\#nameddest=tag-023G}{III.{D.PC9} (p.~198)} }
\expandafter\def\csname SRef@maps:eq:D.PC10\endcsname{\href{S6_SERIES_GRIND_III.pdf\#nameddest=tag-0234}{III.{D.PC10} (p.~198)} }
\expandafter\def\csname SRef@maps:eq:D.PC11\endcsname{\href{S6_SERIES_GRIND_III.pdf\#nameddest=tag-0235}{III.{D.PC11} (p.~198)} }
\expandafter\def\csname SRef@maps:eq:D.PC12\endcsname{\href{S6_SERIES_GRIND_III.pdf\#nameddest=tag-0236}{III.{D.PC12} (p.~198)} }
\expandafter\def\csname SRef@maps:eq:D.PC13\endcsname{\href{S6_SERIES_GRIND_III.pdf\#nameddest=tag-0237}{III.{D.PC13} (p.~198)} }
\expandafter\def\csname SRef@maps:eq:D.PC14\endcsname{\href{S6_SERIES_GRIND_III.pdf\#nameddest=tag-0238}{III.{D.PC14} (p.~198)} }
\expandafter\def\csname SRef@maps:eq:D.OS5\endcsname{\href{S6_SERIES_GRIND_III.pdf\#nameddest=tag-022S}{III.{D.OS5} (p.~199)} }
\expandafter\def\csname SRef@maps:eq:B.5.16\endcsname{\href{S6_SERIES_GRIND_III.pdf\#nameddest=tag-01VE}{III.{B.5.16} (p.~25)} }
\expandafter\def\csname SRef@maps:eq:D.OS6\endcsname{\href{S6_SERIES_GRIND_III.pdf\#nameddest=tag-022T}{III.{D.OS6} (p.~199)} }
\expandafter\def\csname SRef@maps:eq:D.OS7\endcsname{\href{S6_SERIES_GRIND_III.pdf\#nameddest=tag-022U}{III.{D.OS7} (p.~199)} }
\expandafter\def\csname SRef@maps:eq:D.OS8\endcsname{\href{S6_SERIES_GRIND_III.pdf\#nameddest=tag-022V}{III.{D.OS8} (p.~200)} }
\expandafter\def\csname SRef@maps:eq:D.OF1\endcsname{\href{S6_SERIES_GRIND_III.pdf\#nameddest=tag-0217}{III.{D.OF1} (p.~200)} }
\expandafter\def\csname SRef@maps:eq:D.OF2\endcsname{\href{S6_SERIES_GRIND_III.pdf\#nameddest=tag-021I}{III.{D.OF2} (p.~200)} }
\expandafter\def\csname SRef@maps:eq:D.OF3\endcsname{\href{S6_SERIES_GRIND_III.pdf\#nameddest=tag-021J}{III.{D.OF3} (p.~201)} }
\expandafter\def\csname SRef@maps:eq:D.OF4\endcsname{\href{S6_SERIES_GRIND_III.pdf\#nameddest=tag-021K}{III.{D.OF4} (p.~201)} }
\expandafter\def\csname SRef@maps:eq:D.OF5\endcsname{\href{S6_SERIES_GRIND_III.pdf\#nameddest=tag-021L}{III.{D.OF5} (p.~201)} }
\expandafter\def\csname SRef@maps:eq:B.5.17\endcsname{\href{S6_SERIES_GRIND_III.pdf\#nameddest=tag-01VF}{III.{B.5.17} (p.~25)} }
\expandafter\def\csname SRef@maps:eq:D.OF6\endcsname{\href{S6_SERIES_GRIND_III.pdf\#nameddest=tag-021M}{III.{D.OF6} (p.~201)} }
\expandafter\def\csname SRef@maps:eq:D.OF7\endcsname{\href{S6_SERIES_GRIND_III.pdf\#nameddest=tag-021N}{III.{D.OF7} (p.~201)} }
\expandafter\def\csname SRef@maps:eq:D.OF8\endcsname{\href{S6_SERIES_GRIND_III.pdf\#nameddest=tag-021O}{III.{D.OF8} (p.~201)} }
\expandafter\def\csname SRef@maps:eq:D.OF9\endcsname{\href{S6_SERIES_GRIND_III.pdf\#nameddest=tag-021P}{III.{D.OF9} (p.~202)} }
\expandafter\def\csname SRef@maps:eq:D.OF10\endcsname{\href{S6_SERIES_GRIND_III.pdf\#nameddest=tag-0218}{III.{D.OF10} (p.~202)} }
\expandafter\def\csname SRef@maps:eq:D.OF11\endcsname{\href{S6_SERIES_GRIND_III.pdf\#nameddest=tag-0219}{III.{D.OF11} (p.~202)} }
\expandafter\def\csname SRef@maps:eq:D.OF12\endcsname{\href{S6_SERIES_GRIND_III.pdf\#nameddest=tag-021A}{III.{D.OF12} (p.~202)} }
\expandafter\def\csname SRef@maps:eq:D.OF13\endcsname{\href{S6_SERIES_GRIND_III.pdf\#nameddest=tag-021B}{III.{D.OF13} (p.~202)} }
\expandafter\def\csname SRef@maps:eq:D.OF14\endcsname{\href{S6_SERIES_GRIND_III.pdf\#nameddest=tag-021C}{III.{D.OF14} (p.~202)} }
\expandafter\def\csname SRef@maps:eq:B.5.18\endcsname{\href{S6_SERIES_GRIND_III.pdf\#nameddest=tag-01VG}{III.{B.5.18} (p.~26)} }
\expandafter\def\csname SRef@maps:eq:D.OF15\endcsname{\href{S6_SERIES_GRIND_III.pdf\#nameddest=tag-021D}{III.{D.OF15} (p.~203)} }
\expandafter\def\csname SRef@maps:eq:D.OF16\endcsname{\href{S6_SERIES_GRIND_III.pdf\#nameddest=tag-021E}{III.{D.OF16} (p.~203)} }
\expandafter\def\csname SRef@maps:eq:D.OF17\endcsname{\href{S6_SERIES_GRIND_III.pdf\#nameddest=tag-021F}{III.{D.OF17} (p.~203)} }
\expandafter\def\csname SRef@maps:eq:D.OF18\endcsname{\href{S6_SERIES_GRIND_III.pdf\#nameddest=tag-021G}{III.{D.OF18} (p.~203)} }
\expandafter\def\csname SRef@maps:eq:D.OL1\endcsname{\href{S6_SERIES_GRIND_III.pdf\#nameddest=tag-021Q}{III.{D.OL1} (p.~204)} }
\expandafter\def\csname SRef@maps:eq:D.OL2\endcsname{\href{S6_SERIES_GRIND_III.pdf\#nameddest=tag-021U}{III.{D.OL2} (p.~204)} }
\expandafter\def\csname SRef@maps:eq:D.OL3\endcsname{\href{S6_SERIES_GRIND_III.pdf\#nameddest=tag-021V}{III.{D.OL3} (p.~204)} }
\expandafter\def\csname SRef@maps:eq:D.OL4\endcsname{\href{S6_SERIES_GRIND_III.pdf\#nameddest=tag-021W}{III.{D.OL4} (p.~204)} }
\expandafter\def\csname SRef@maps:eq:D.OL5\endcsname{\href{S6_SERIES_GRIND_III.pdf\#nameddest=tag-021X}{III.{D.OL5} (p.~204)} }
\expandafter\def\csname SRef@maps:eq:C.1.1\endcsname{\href{S6_SERIES_GRIND_III.pdf\#nameddest=tag-01XE}{III.{C.1.1} (p.~27)} }
\expandafter\def\csname SRef@maps:eq:D.OL6\endcsname{\href{S6_SERIES_GRIND_III.pdf\#nameddest=tag-021Y}{III.{D.OL6} (p.~204)} }
\expandafter\def\csname SRef@maps:eq:D.OL7\endcsname{\href{S6_SERIES_GRIND_III.pdf\#nameddest=tag-021Z}{III.{D.OL7} (p.~204)} }
\expandafter\def\csname SRef@maps:eq:D.OL8\endcsname{\href{S6_SERIES_GRIND_III.pdf\#nameddest=tag-0220}{III.{D.OL8} (p.~205)} }
\expandafter\def\csname SRef@maps:eq:D.OL9\endcsname{\href{S6_SERIES_GRIND_III.pdf\#nameddest=tag-0221}{III.{D.OL9} (p.~205)} }
\expandafter\def\csname SRef@maps:eq:D.OL10\endcsname{\href{S6_SERIES_GRIND_III.pdf\#nameddest=tag-021R}{III.{D.OL10} (p.~205)} }
\expandafter\def\csname SRef@maps:eq:D.OL11\endcsname{\href{S6_SERIES_GRIND_III.pdf\#nameddest=tag-021S}{III.{D.OL11} (p.~205)} }
\expandafter\def\csname SRef@maps:eq:D.OL12\endcsname{\href{S6_SERIES_GRIND_III.pdf\#nameddest=tag-021T}{III.{D.OL12} (p.~205)} }
\expandafter\def\csname SRef@maps:eq:D.OF19\endcsname{\href{S6_SERIES_GRIND_III.pdf\#nameddest=tag-021H}{III.{D.OF19} (p.~206)} }
\expandafter\def\csname SRef@maps:eq:C.1.2\endcsname{\href{S6_SERIES_GRIND_III.pdf\#nameddest=tag-01XF}{III.{C.1.2} (p.~28)} }
\expandafter\def\csname SRef@maps:eq:C.1.3\endcsname{\href{S6_SERIES_GRIND_III.pdf\#nameddest=tag-01XG}{III.{C.1.3} (p.~30)} }
\expandafter\def\csname SRef@maps:eq:E.1\endcsname{\href{S6_SERIES_GRIND_V.pdf\#nameddest=tag-029F}{V.{E.1} (p.~14)} }
\expandafter\def\csname SRef@maps:eq:E.2\endcsname{\href{S6_SERIES_GRIND_V.pdf\#nameddest=tag-029O}{V.{E.2} (p.~15)} }
\expandafter\def\csname SRef@maps:eq:C.1.4\endcsname{\href{S6_SERIES_GRIND_III.pdf\#nameddest=tag-01XH}{III.{C.1.4} (p.~30)} }
\expandafter\def\csname SRef@maps:eq:E.3\endcsname{\href{S6_SERIES_GRIND_V.pdf\#nameddest=tag-029Q}{V.{E.3} (p.~15)} }
\expandafter\def\csname SRef@maps:eq:E.4\endcsname{\href{S6_SERIES_GRIND_V.pdf\#nameddest=tag-029R}{V.{E.4} (p.~15)} }
\expandafter\def\csname SRef@maps:eq:E.5\endcsname{\href{S6_SERIES_GRIND_V.pdf\#nameddest=tag-029S}{V.{E.5} (p.~15)} }
\expandafter\def\csname SRef@maps:eq:E.6\endcsname{\href{S6_SERIES_GRIND_V.pdf\#nameddest=tag-029T}{V.{E.6} (p.~16)} }
\expandafter\def\csname SRef@maps:eq:E.7\endcsname{\href{S6_SERIES_GRIND_V.pdf\#nameddest=tag-029U}{V.{E.7} (p.~16)} }
\expandafter\def\csname SRef@maps:eq:E.8\endcsname{\href{S6_SERIES_GRIND_V.pdf\#nameddest=tag-029V}{V.{E.8} (p.~18)} }
\expandafter\def\csname SRef@maps:eq:E.9\endcsname{\href{S6_SERIES_GRIND_V.pdf\#nameddest=tag-029W}{V.{E.9} (p.~18)} }
\expandafter\def\csname SRef@maps:eq:E.10\endcsname{\href{S6_SERIES_GRIND_V.pdf\#nameddest=tag-029G}{V.{E.10} (p.~18)} }
\expandafter\def\csname SRef@maps:eq:E.11\endcsname{\href{S6_SERIES_GRIND_V.pdf\#nameddest=tag-029H}{V.{E.11} (p.~18)} }
\expandafter\def\csname SRef@maps:eq:E.I1\endcsname{\href{S6_SERIES_GRIND_V.pdf\#nameddest=tag-02B5}{V.{E.I1} (p.~19)} }
\expandafter\def\csname SRef@maps:eq:D.C1\endcsname{\href{S6_SERIES_GRIND_III.pdf\#nameddest=tag-01Y6}{III.{D.C1} (p.~59)} }
\expandafter\def\csname SRef@maps:eq:E.I2\endcsname{\href{S6_SERIES_GRIND_V.pdf\#nameddest=tag-02B6}{V.{E.I2} (p.~20)} }
\expandafter\def\csname SRef@maps:eq:E.I3\endcsname{\href{S6_SERIES_GRIND_V.pdf\#nameddest=tag-02B7}{V.{E.I3} (p.~20)} }
\expandafter\def\csname SRef@maps:eq:E.I4\endcsname{\href{S6_SERIES_GRIND_V.pdf\#nameddest=tag-02B8}{V.{E.I4} (p.~20)} }
\expandafter\def\csname SRef@maps:eq:E.I5\endcsname{\href{S6_SERIES_GRIND_V.pdf\#nameddest=tag-02B9}{V.{E.I5} (p.~21)} }
\expandafter\def\csname SRef@maps:eq:E.12\endcsname{\href{S6_SERIES_GRIND_V.pdf\#nameddest=tag-029I}{V.{E.12} (p.~21)} }
\expandafter\def\csname SRef@maps:eq:E.13\endcsname{\href{S6_SERIES_GRIND_V.pdf\#nameddest=tag-029J}{V.{E.13} (p.~21)} }
\expandafter\def\csname SRef@maps:eq:E.14\endcsname{\href{S6_SERIES_GRIND_V.pdf\#nameddest=tag-029K}{V.{E.14} (p.~22)} }
\expandafter\def\csname SRef@maps:eq:E.16\endcsname{\href{S6_SERIES_GRIND_V.pdf\#nameddest=tag-029M}{V.{E.16} (p.~23)} }
\expandafter\def\csname SRef@maps:eq:E.17\endcsname{\href{S6_SERIES_GRIND_V.pdf\#nameddest=tag-029N}{V.{E.17} (p.~23)} }
\expandafter\def\csname SRef@maps:eq:E.20\endcsname{\href{S6_SERIES_GRIND_V.pdf\#nameddest=tag-029P}{V.{E.20} (p.~25)} }
\expandafter\def\csname SRef@maps:eq:D.R1\endcsname{\href{S6_SERIES_GRIND_III.pdf\#nameddest=tag-023H}{III.{D.R1} (p.~64)} }
\expandafter\def\csname SRef@maps:eq:E.MV1\endcsname{\href{S6_SERIES_GRIND_V.pdf\#nameddest=tag-02BA}{V.{E.MV1} (p.~26)} }
\expandafter\def\csname SRef@maps:eq:E.MV2\endcsname{\href{S6_SERIES_GRIND_V.pdf\#nameddest=tag-02BB}{V.{E.MV2} (p.~27)} }
\expandafter\def\csname SRef@maps:eq:E.MV3\endcsname{\href{S6_SERIES_GRIND_V.pdf\#nameddest=tag-02BC}{V.{E.MV3} (p.~27)} }
\expandafter\def\csname SRef@maps:eq:E.MV4\endcsname{\href{S6_SERIES_GRIND_V.pdf\#nameddest=tag-02BD}{V.{E.MV4} (p.~27)} }
\expandafter\def\csname SRef@maps:eq:E.MV5\endcsname{\href{S6_SERIES_GRIND_V.pdf\#nameddest=tag-02BE}{V.{E.MV5} (p.~27)} }
\expandafter\def\csname SRef@maps:eq:E.MV6\endcsname{\href{S6_SERIES_GRIND_V.pdf\#nameddest=tag-02BF}{V.{E.MV6} (p.~28)} }
\expandafter\def\csname SRef@maps:eq:E.MV7\endcsname{\href{S6_SERIES_GRIND_V.pdf\#nameddest=tag-02BG}{V.{E.MV7} (p.~28)} }
\expandafter\def\csname SRef@maps:eq:E.MV8\endcsname{\href{S6_SERIES_GRIND_V.pdf\#nameddest=tag-02BH}{V.{E.MV8} (p.~28)} }
\expandafter\def\csname SRef@maps:eq:E.MV9\endcsname{\href{S6_SERIES_GRIND_V.pdf\#nameddest=tag-02BI}{V.{E.MV9} (p.~29)} }
\expandafter\def\csname SRef@maps:eq:E.B1\endcsname{\href{S6_SERIES_GRIND_V.pdf\#nameddest=tag-029X}{V.{E.B1} (p.~30)} }
\expandafter\def\csname SRef@maps:eq:D.R2\endcsname{\href{S6_SERIES_GRIND_III.pdf\#nameddest=tag-023I}{III.{D.R2} (p.~64)} }
\expandafter\def\csname SRef@maps:eq:E.B2\endcsname{\href{S6_SERIES_GRIND_V.pdf\#nameddest=tag-029Y}{V.{E.B2} (p.~30)} }
\expandafter\def\csname SRef@maps:eq:E.B3\endcsname{\href{S6_SERIES_GRIND_V.pdf\#nameddest=tag-029Z}{V.{E.B3} (p.~31)} }
\expandafter\def\csname SRef@maps:eq:E.B4\endcsname{\href{S6_SERIES_GRIND_V.pdf\#nameddest=tag-02A0}{V.{E.B4} (p.~31)} }
\expandafter\def\csname SRef@maps:eq:E.B5\endcsname{\href{S6_SERIES_GRIND_V.pdf\#nameddest=tag-02A1}{V.{E.B5} (p.~31)} }
\expandafter\def\csname SRef@maps:eq:E.B6\endcsname{\href{S6_SERIES_GRIND_V.pdf\#nameddest=tag-02A2}{V.{E.B6} (p.~32)} }
\expandafter\def\csname SRef@maps:eq:E.B7\endcsname{\href{S6_SERIES_GRIND_V.pdf\#nameddest=tag-02A3}{V.{E.B7} (p.~32)} }
\expandafter\def\csname SRef@maps:eq:E.S1\endcsname{\href{S6_SERIES_GRIND_V.pdf\#nameddest=tag-02CC}{V.{E.S1} (p.~33)} }
\expandafter\def\csname SRef@maps:eq:E.S2\endcsname{\href{S6_SERIES_GRIND_V.pdf\#nameddest=tag-02CN}{V.{E.S2} (p.~33)} }
\expandafter\def\csname SRef@maps:eq:E.S3\endcsname{\href{S6_SERIES_GRIND_V.pdf\#nameddest=tag-02CY}{V.{E.S3} (p.~34)} }
\expandafter\def\csname SRef@maps:eq:E.S4\endcsname{\href{S6_SERIES_GRIND_V.pdf\#nameddest=tag-02D9}{V.{E.S4} (p.~34)} }
\expandafter\def\csname SRef@maps:eq:D.R3\endcsname{\href{S6_SERIES_GRIND_III.pdf\#nameddest=tag-023J}{III.{D.R3} (p.~64)} }
\expandafter\def\csname SRef@maps:eq:E.S5\endcsname{\href{S6_SERIES_GRIND_V.pdf\#nameddest=tag-02DK}{V.{E.S5} (p.~34)} }
\expandafter\def\csname SRef@maps:eq:E.S12\endcsname{\href{S6_SERIES_GRIND_V.pdf\#nameddest=tag-02CF}{V.{E.S12} (p.~35)} }
\expandafter\def\csname SRef@maps:eq:E.S6\endcsname{\href{S6_SERIES_GRIND_V.pdf\#nameddest=tag-02DR}{V.{E.S6} (p.~35)} }
\expandafter\def\csname SRef@maps:eq:E.S7\endcsname{\href{S6_SERIES_GRIND_V.pdf\#nameddest=tag-02DS}{V.{E.S7} (p.~35)} }
\expandafter\def\csname SRef@maps:eq:E.S8\endcsname{\href{S6_SERIES_GRIND_V.pdf\#nameddest=tag-02DT}{V.{E.S8} (p.~36)} }
\expandafter\def\csname SRef@maps:eq:E.S9\endcsname{\href{S6_SERIES_GRIND_V.pdf\#nameddest=tag-02DU}{V.{E.S9} (p.~36)} }
\expandafter\def\csname SRef@maps:eq:E.S10\endcsname{\href{S6_SERIES_GRIND_V.pdf\#nameddest=tag-02CD}{V.{E.S10} (p.~36)} }
\expandafter\def\csname SRef@maps:eq:E.S11\endcsname{\href{S6_SERIES_GRIND_V.pdf\#nameddest=tag-02CE}{V.{E.S11} (p.~36)} }
\expandafter\def\csname SRef@maps:eq:E.S13\endcsname{\href{S6_SERIES_GRIND_V.pdf\#nameddest=tag-02CG}{V.{E.S13} (p.~37)} }
\expandafter\def\csname SRef@maps:eq:E.S14\endcsname{\href{S6_SERIES_GRIND_V.pdf\#nameddest=tag-02CH}{V.{E.S14} (p.~37)} }
\expandafter\def\csname SRef@maps:eq:D.R4\endcsname{\href{S6_SERIES_GRIND_III.pdf\#nameddest=tag-023K}{III.{D.R4} (p.~64)} }
\expandafter\def\csname SRef@maps:eq:E.S15\endcsname{\href{S6_SERIES_GRIND_V.pdf\#nameddest=tag-02CI}{V.{E.S15} (p.~37)} }
\expandafter\def\csname SRef@maps:eq:E.S16\endcsname{\href{S6_SERIES_GRIND_V.pdf\#nameddest=tag-02CJ}{V.{E.S16} (p.~37)} }
\expandafter\def\csname SRef@maps:eq:E.S20\endcsname{\href{S6_SERIES_GRIND_V.pdf\#nameddest=tag-02CO}{V.{E.S20} (p.~37)} }
\expandafter\def\csname SRef@maps:eq:E.S17\endcsname{\href{S6_SERIES_GRIND_V.pdf\#nameddest=tag-02CK}{V.{E.S17} (p.~37)} }
\expandafter\def\csname SRef@maps:eq:E.S18\endcsname{\href{S6_SERIES_GRIND_V.pdf\#nameddest=tag-02CL}{V.{E.S18} (p.~37)} }
\expandafter\def\csname SRef@maps:eq:E.S19\endcsname{\href{S6_SERIES_GRIND_V.pdf\#nameddest=tag-02CM}{V.{E.S19} (p.~38)} }
\expandafter\def\csname SRef@maps:eq:E.S21\endcsname{\href{S6_SERIES_GRIND_V.pdf\#nameddest=tag-02CP}{V.{E.S21} (p.~39)} }
\expandafter\def\csname SRef@maps:eq:E.S22\endcsname{\href{S6_SERIES_GRIND_V.pdf\#nameddest=tag-02CQ}{V.{E.S22} (p.~40)} }
\expandafter\def\csname SRef@maps:eq:E.S23\endcsname{\href{S6_SERIES_GRIND_V.pdf\#nameddest=tag-02CR}{V.{E.S23} (p.~40)} }
\expandafter\def\csname SRef@maps:eq:E.S24\endcsname{\href{S6_SERIES_GRIND_V.pdf\#nameddest=tag-02CS}{V.{E.S24} (p.~40)} }
\expandafter\def\csname SRef@maps:eq:D.S1\endcsname{\href{S6_SERIES_GRIND_III.pdf\#nameddest=tag-023L}{III.{D.S1} (p.~65)} }
\expandafter\def\csname SRef@maps:eq:E.S25\endcsname{\href{S6_SERIES_GRIND_V.pdf\#nameddest=tag-02CT}{V.{E.S25} (p.~40)} }
\expandafter\def\csname SRef@maps:eq:E.S26\endcsname{\href{S6_SERIES_GRIND_V.pdf\#nameddest=tag-02CU}{V.{E.S26} (p.~40)} }
\expandafter\def\csname SRef@maps:eq:E.S27\endcsname{\href{S6_SERIES_GRIND_V.pdf\#nameddest=tag-02CV}{V.{E.S27} (p.~41)} }
\expandafter\def\csname SRef@maps:eq:E.S28\endcsname{\href{S6_SERIES_GRIND_V.pdf\#nameddest=tag-02CW}{V.{E.S28} (p.~41)} }
\expandafter\def\csname SRef@maps:eq:E.S29\endcsname{\href{S6_SERIES_GRIND_V.pdf\#nameddest=tag-02CX}{V.{E.S29} (p.~42)} }
\expandafter\def\csname SRef@maps:eq:E.S30\endcsname{\href{S6_SERIES_GRIND_V.pdf\#nameddest=tag-02CZ}{V.{E.S30} (p.~42)} }
\expandafter\def\csname SRef@maps:eq:E.S31\endcsname{\href{S6_SERIES_GRIND_V.pdf\#nameddest=tag-02D0}{V.{E.S31} (p.~42)} }
\expandafter\def\csname SRef@maps:eq:E.S32\endcsname{\href{S6_SERIES_GRIND_V.pdf\#nameddest=tag-02D1}{V.{E.S32} (p.~42)} }
\expandafter\def\csname SRef@maps:eq:E.S33\endcsname{\href{S6_SERIES_GRIND_V.pdf\#nameddest=tag-02D2}{V.{E.S33} (p.~42)} }
\expandafter\def\csname SRef@maps:eq:E.S34\endcsname{\href{S6_SERIES_GRIND_V.pdf\#nameddest=tag-02D3}{V.{E.S34} (p.~43)} }
\expandafter\def\csname SRef@maps:eq:D.S2\endcsname{\href{S6_SERIES_GRIND_III.pdf\#nameddest=tag-023R}{III.{D.S2} (p.~65)} }
\expandafter\def\csname SRef@maps:eq:E.S35\endcsname{\href{S6_SERIES_GRIND_V.pdf\#nameddest=tag-02D4}{V.{E.S35} (p.~43)} }
\expandafter\def\csname SRef@maps:eq:E.S36\endcsname{\href{S6_SERIES_GRIND_V.pdf\#nameddest=tag-02D5}{V.{E.S36} (p.~43)} }
\expandafter\def\csname SRef@maps:eq:E.S37\endcsname{\href{S6_SERIES_GRIND_V.pdf\#nameddest=tag-02D6}{V.{E.S37} (p.~43)} }
\expandafter\def\csname SRef@maps:eq:E.S38\endcsname{\href{S6_SERIES_GRIND_V.pdf\#nameddest=tag-02D7}{V.{E.S38} (p.~43)} }
\expandafter\def\csname SRef@maps:eq:E.S39\endcsname{\href{S6_SERIES_GRIND_V.pdf\#nameddest=tag-02D8}{V.{E.S39} (p.~44)} }
\expandafter\def\csname SRef@maps:eq:E.S40\endcsname{\href{S6_SERIES_GRIND_V.pdf\#nameddest=tag-02DA}{V.{E.S40} (p.~44)} }
\expandafter\def\csname SRef@maps:eq:E.S41\endcsname{\href{S6_SERIES_GRIND_V.pdf\#nameddest=tag-02DB}{V.{E.S41} (p.~44)} }
\expandafter\def\csname SRef@maps:eq:E.S42\endcsname{\href{S6_SERIES_GRIND_V.pdf\#nameddest=tag-02DC}{V.{E.S42} (p.~44)} }
\expandafter\def\csname SRef@maps:eq:E.S43\endcsname{\href{S6_SERIES_GRIND_V.pdf\#nameddest=tag-02DD}{V.{E.S43} (p.~45)} }
\expandafter\def\csname SRef@maps:eq:E.S44\endcsname{\href{S6_SERIES_GRIND_V.pdf\#nameddest=tag-02DE}{V.{E.S44} (p.~45)} }
\expandafter\def\csname SRef@maps:eq:D.S3\endcsname{\href{S6_SERIES_GRIND_III.pdf\#nameddest=tag-023S}{III.{D.S3} (p.~65)} }
\expandafter\def\csname SRef@maps:eq:E.S45\endcsname{\href{S6_SERIES_GRIND_V.pdf\#nameddest=tag-02DF}{V.{E.S45} (p.~45)} }
\expandafter\def\csname SRef@maps:eq:E.S46\endcsname{\href{S6_SERIES_GRIND_V.pdf\#nameddest=tag-02DG}{V.{E.S46} (p.~45)} }
\expandafter\def\csname SRef@maps:eq:E.S47\endcsname{\href{S6_SERIES_GRIND_V.pdf\#nameddest=tag-02DH}{V.{E.S47} (p.~45)} }
\expandafter\def\csname SRef@maps:eq:E.S48\endcsname{\href{S6_SERIES_GRIND_V.pdf\#nameddest=tag-02DI}{V.{E.S48} (p.~46)} }
\expandafter\def\csname SRef@maps:eq:E.S49\endcsname{\href{S6_SERIES_GRIND_V.pdf\#nameddest=tag-02DJ}{V.{E.S49} (p.~46)} }
\expandafter\def\csname SRef@maps:eq:E.S50\endcsname{\href{S6_SERIES_GRIND_V.pdf\#nameddest=tag-02DL}{V.{E.S50} (p.~46)} }
\expandafter\def\csname SRef@maps:eq:E.S51\endcsname{\href{S6_SERIES_GRIND_V.pdf\#nameddest=tag-02DM}{V.{E.S51} (p.~46)} }
\expandafter\def\csname SRef@maps:eq:E.S52\endcsname{\href{S6_SERIES_GRIND_V.pdf\#nameddest=tag-02DN}{V.{E.S52} (p.~46)} }
\expandafter\def\csname SRef@maps:eq:E.S53\endcsname{\href{S6_SERIES_GRIND_V.pdf\#nameddest=tag-02DO}{V.{E.S53} (p.~46)} }
\expandafter\def\csname SRef@maps:eq:E.S54\endcsname{\href{S6_SERIES_GRIND_V.pdf\#nameddest=tag-02DP}{V.{E.S54} (p.~46)} }
\expandafter\def\csname SRef@maps:eq:D.S4\endcsname{\href{S6_SERIES_GRIND_III.pdf\#nameddest=tag-023T}{III.{D.S4} (p.~65)} }
\expandafter\def\csname SRef@maps:eq:E.S55\endcsname{\href{S6_SERIES_GRIND_V.pdf\#nameddest=tag-02DQ}{V.{E.S55} (p.~47)} }
\expandafter\def\csname SRef@maps:eq:E.C1\endcsname{\href{S6_SERIES_GRIND_V.pdf\#nameddest=tag-02A4}{V.{E.C1} (p.~48)} }
\expandafter\def\csname SRef@maps:eq:E.C2\endcsname{\href{S6_SERIES_GRIND_V.pdf\#nameddest=tag-02AG}{V.{E.C2} (p.~48)} }
\expandafter\def\csname SRef@maps:eq:E.C3\endcsname{\href{S6_SERIES_GRIND_V.pdf\#nameddest=tag-02AJ}{V.{E.C3} (p.~48)} }
\expandafter\def\csname SRef@maps:eq:E.C4\endcsname{\href{S6_SERIES_GRIND_V.pdf\#nameddest=tag-02AK}{V.{E.C4} (p.~48)} }
\expandafter\def\csname SRef@maps:eq:E.C19a\endcsname{\href{S6_SERIES_GRIND_V.pdf\#nameddest=tag-02AE}{V.{E.C19a} (p.~48)} }
\expandafter\def\csname SRef@maps:eq:E.C19b\endcsname{\href{S6_SERIES_GRIND_V.pdf\#nameddest=tag-02AF}{V.{E.C19b} (p.~48)} }
\expandafter\def\csname SRef@maps:eq:E.C5\endcsname{\href{S6_SERIES_GRIND_V.pdf\#nameddest=tag-02AL}{V.{E.C5} (p.~49)} }
\expandafter\def\csname SRef@maps:eq:E.C6\endcsname{\href{S6_SERIES_GRIND_V.pdf\#nameddest=tag-02AM}{V.{E.C6} (p.~49)} }
\expandafter\def\csname SRef@maps:eq:E.C7\endcsname{\href{S6_SERIES_GRIND_V.pdf\#nameddest=tag-02AN}{V.{E.C7} (p.~49)} }
\expandafter\def\csname SRef@maps:eq:D.S5\endcsname{\href{S6_SERIES_GRIND_III.pdf\#nameddest=tag-023U}{III.{D.S5} (p.~66)} }
\expandafter\def\csname SRef@maps:eq:E.C8\endcsname{\href{S6_SERIES_GRIND_V.pdf\#nameddest=tag-02AO}{V.{E.C8} (p.~49)} }
\expandafter\def\csname SRef@maps:eq:E.C9\endcsname{\href{S6_SERIES_GRIND_V.pdf\#nameddest=tag-02AP}{V.{E.C9} (p.~49)} }
\expandafter\def\csname SRef@maps:eq:E.C10\endcsname{\href{S6_SERIES_GRIND_V.pdf\#nameddest=tag-02A5}{V.{E.C10} (p.~49)} }
\expandafter\def\csname SRef@maps:eq:E.C11\endcsname{\href{S6_SERIES_GRIND_V.pdf\#nameddest=tag-02A6}{V.{E.C11} (p.~49)} }
\expandafter\def\csname SRef@maps:eq:E.C12\endcsname{\href{S6_SERIES_GRIND_V.pdf\#nameddest=tag-02A7}{V.{E.C12} (p.~50)} }
\expandafter\def\csname SRef@maps:eq:E.C20\endcsname{\href{S6_SERIES_GRIND_V.pdf\#nameddest=tag-02AH}{V.{E.C20} (p.~50)} }
\expandafter\def\csname SRef@maps:eq:E.C13\endcsname{\href{S6_SERIES_GRIND_V.pdf\#nameddest=tag-02A8}{V.{E.C13} (p.~51)} }
\expandafter\def\csname SRef@maps:eq:E.C14\endcsname{\href{S6_SERIES_GRIND_V.pdf\#nameddest=tag-02A9}{V.{E.C14} (p.~51)} }
\expandafter\def\csname SRef@maps:eq:E.C15\endcsname{\href{S6_SERIES_GRIND_V.pdf\#nameddest=tag-02AA}{V.{E.C15} (p.~51)} }
\expandafter\def\csname SRef@maps:eq:E.C16\endcsname{\href{S6_SERIES_GRIND_V.pdf\#nameddest=tag-02AB}{V.{E.C16} (p.~51)} }
\expandafter\def\csname SRef@maps:eq:D.S6\endcsname{\href{S6_SERIES_GRIND_III.pdf\#nameddest=tag-023V}{III.{D.S6} (p.~66)} }
\expandafter\def\csname SRef@maps:eq:E.C17\endcsname{\href{S6_SERIES_GRIND_V.pdf\#nameddest=tag-02AC}{V.{E.C17} (p.~51)} }
\expandafter\def\csname SRef@maps:eq:E.C18\endcsname{\href{S6_SERIES_GRIND_V.pdf\#nameddest=tag-02AD}{V.{E.C18} (p.~52)} }
\expandafter\def\csname SRef@maps:eq:E.C21\endcsname{\href{S6_SERIES_GRIND_V.pdf\#nameddest=tag-02AI}{V.{E.C21} (p.~52)} }
\expandafter\def\csname SRef@maps:eq:E.P1\endcsname{\href{S6_SERIES_GRIND_V.pdf\#nameddest=tag-02BJ}{V.{E.P1} (p.~52)} }
\expandafter\def\csname SRef@maps:eq:E.P2\endcsname{\href{S6_SERIES_GRIND_V.pdf\#nameddest=tag-02BS}{V.{E.P2} (p.~53)} }
\expandafter\def\csname SRef@maps:eq:E.P3\endcsname{\href{S6_SERIES_GRIND_V.pdf\#nameddest=tag-02BT}{V.{E.P3} (p.~53)} }
\expandafter\def\csname SRef@maps:eq:E.P4\endcsname{\href{S6_SERIES_GRIND_V.pdf\#nameddest=tag-02BU}{V.{E.P4} (p.~53)} }
\expandafter\def\csname SRef@maps:eq:E.P5\endcsname{\href{S6_SERIES_GRIND_V.pdf\#nameddest=tag-02BV}{V.{E.P5} (p.~53)} }
\expandafter\def\csname SRef@maps:eq:E.P6\endcsname{\href{S6_SERIES_GRIND_V.pdf\#nameddest=tag-02BW}{V.{E.P6} (p.~53)} }
\expandafter\def\csname SRef@maps:eq:E.P7\endcsname{\href{S6_SERIES_GRIND_V.pdf\#nameddest=tag-02BX}{V.{E.P7} (p.~53)} }
\expandafter\def\csname SRef@maps:eq:D.S7\endcsname{\href{S6_SERIES_GRIND_III.pdf\#nameddest=tag-023W}{III.{D.S7} (p.~67)} }
\expandafter\def\csname SRef@maps:eq:E.P8\endcsname{\href{S6_SERIES_GRIND_V.pdf\#nameddest=tag-02BY}{V.{E.P8} (p.~53)} }
\expandafter\def\csname SRef@maps:eq:E.P9\endcsname{\href{S6_SERIES_GRIND_V.pdf\#nameddest=tag-02BZ}{V.{E.P9} (p.~53)} }
\expandafter\def\csname SRef@maps:eq:E.P10\endcsname{\href{S6_SERIES_GRIND_V.pdf\#nameddest=tag-02BK}{V.{E.P10} (p.~54)} }
\expandafter\def\csname SRef@maps:eq:E.P11\endcsname{\href{S6_SERIES_GRIND_V.pdf\#nameddest=tag-02BL}{V.{E.P11} (p.~54)} }
\expandafter\def\csname SRef@maps:eq:E.P12\endcsname{\href{S6_SERIES_GRIND_V.pdf\#nameddest=tag-02BM}{V.{E.P12} (p.~54)} }
\expandafter\def\csname SRef@maps:eq:E.P13\endcsname{\href{S6_SERIES_GRIND_V.pdf\#nameddest=tag-02BN}{V.{E.P13} (p.~54)} }
\expandafter\def\csname SRef@maps:eq:E.P14\endcsname{\href{S6_SERIES_GRIND_V.pdf\#nameddest=tag-02BO}{V.{E.P14} (p.~55)} }
\expandafter\def\csname SRef@maps:eq:E.P15\endcsname{\href{S6_SERIES_GRIND_V.pdf\#nameddest=tag-02BP}{V.{E.P15} (p.~55)} }
\expandafter\def\csname SRef@maps:eq:E.P16\endcsname{\href{S6_SERIES_GRIND_V.pdf\#nameddest=tag-02BQ}{V.{E.P16} (p.~55)} }
\expandafter\def\csname SRef@maps:eq:E.P17\endcsname{\href{S6_SERIES_GRIND_V.pdf\#nameddest=tag-02BR}{V.{E.P17} (p.~56)} }
\expandafter\def\csname SRef@maps:eq:D.S8\endcsname{\href{S6_SERIES_GRIND_III.pdf\#nameddest=tag-023X}{III.{D.S8} (p.~67)} }
\expandafter\def\csname SRef@maps:eq:E.Q1\endcsname{\href{S6_SERIES_GRIND_V.pdf\#nameddest=tag-02C0}{V.{E.Q1} (p.~57)} }
\expandafter\def\csname SRef@maps:eq:E.Q2\endcsname{\href{S6_SERIES_GRIND_V.pdf\#nameddest=tag-02C4}{V.{E.Q2} (p.~57)} }
\expandafter\def\csname SRef@maps:eq:E.Q3\endcsname{\href{S6_SERIES_GRIND_V.pdf\#nameddest=tag-02C5}{V.{E.Q3} (p.~57)} }
\expandafter\def\csname SRef@maps:eq:E.Q4\endcsname{\href{S6_SERIES_GRIND_V.pdf\#nameddest=tag-02C6}{V.{E.Q4} (p.~57)} }
\expandafter\def\csname SRef@maps:eq:E.Q5\endcsname{\href{S6_SERIES_GRIND_V.pdf\#nameddest=tag-02C7}{V.{E.Q5} (p.~58)} }
\expandafter\def\csname SRef@maps:eq:E.Q6\endcsname{\href{S6_SERIES_GRIND_V.pdf\#nameddest=tag-02C8}{V.{E.Q6} (p.~58)} }
\expandafter\def\csname SRef@maps:eq:E.Q7\endcsname{\href{S6_SERIES_GRIND_V.pdf\#nameddest=tag-02C9}{V.{E.Q7} (p.~58)} }
\expandafter\def\csname SRef@maps:eq:E.Q8\endcsname{\href{S6_SERIES_GRIND_V.pdf\#nameddest=tag-02CA}{V.{E.Q8} (p.~58)} }
\expandafter\def\csname SRef@maps:eq:E.Q9\endcsname{\href{S6_SERIES_GRIND_V.pdf\#nameddest=tag-02CB}{V.{E.Q9} (p.~59)} }
\expandafter\def\csname SRef@maps:eq:E.Q10\endcsname{\href{S6_SERIES_GRIND_V.pdf\#nameddest=tag-02C1}{V.{E.Q10} (p.~59)} }
\expandafter\def\csname SRef@maps:eq:D.S9\endcsname{\href{S6_SERIES_GRIND_III.pdf\#nameddest=tag-023Y}{III.{D.S9} (p.~68)} }
\expandafter\def\csname SRef@maps:eq:E.Q11\endcsname{\href{S6_SERIES_GRIND_V.pdf\#nameddest=tag-02C2}{V.{E.Q11} (p.~59)} }
\expandafter\def\csname SRef@maps:eq:E.Q12\endcsname{\href{S6_SERIES_GRIND_V.pdf\#nameddest=tag-02C3}{V.{E.Q12} (p.~60)} }
\expandafter\def\csname SRef@maps:eq:E.T1\endcsname{\href{S6_SERIES_GRIND_V.pdf\#nameddest=tag-02DV}{V.{E.T1} (p.~60)} }
\expandafter\def\csname SRef@maps:eq:E.T2\endcsname{\href{S6_SERIES_GRIND_V.pdf\#nameddest=tag-02DW}{V.{E.T2} (p.~60)} }
\expandafter\def\csname SRef@maps:eq:E.T3\endcsname{\href{S6_SERIES_GRIND_V.pdf\#nameddest=tag-02DX}{V.{E.T3} (p.~61)} }
\expandafter\def\csname SRef@maps:eq:E.T4\endcsname{\href{S6_SERIES_GRIND_V.pdf\#nameddest=tag-02DY}{V.{E.T4} (p.~61)} }
\expandafter\def\csname SRef@maps:eq:E.T5\endcsname{\href{S6_SERIES_GRIND_V.pdf\#nameddest=tag-02DZ}{V.{E.T5} (p.~61)} }
\expandafter\def\csname SRef@maps:eq:E.T6\endcsname{\href{S6_SERIES_GRIND_V.pdf\#nameddest=tag-02E0}{V.{E.T6} (p.~62)} }
\expandafter\def\csname SRef@maps:eq:E.T7\endcsname{\href{S6_SERIES_GRIND_V.pdf\#nameddest=tag-02E1}{V.{E.T7} (p.~62)} }
\expandafter\def\csname SRef@maps:eq:E.T9\endcsname{\href{S6_SERIES_GRIND_V.pdf\#nameddest=tag-02E3}{V.{E.T9} (p.~62)} }
\expandafter\def\csname SRef@maps:eq:D.S10\endcsname{\href{S6_SERIES_GRIND_III.pdf\#nameddest=tag-023M}{III.{D.S10} (p.~68)} }
\expandafter\def\csname SRef@maps:eq:E.T8\endcsname{\href{S6_SERIES_GRIND_V.pdf\#nameddest=tag-02E2}{V.{E.T8} (p.~62)} }
\expandafter\def\csname SRef@maps:eq:E.U1\endcsname{\href{S6_SERIES_GRIND_V.pdf\#nameddest=tag-02E4}{V.{E.U1} (p.~63)} }
\expandafter\def\csname SRef@maps:eq:E.U2\endcsname{\href{S6_SERIES_GRIND_V.pdf\#nameddest=tag-02E7}{V.{E.U2} (p.~63)} }
\expandafter\def\csname SRef@maps:eq:E.U3\endcsname{\href{S6_SERIES_GRIND_V.pdf\#nameddest=tag-02E8}{V.{E.U3} (p.~63)} }
\expandafter\def\csname SRef@maps:eq:E.U4\endcsname{\href{S6_SERIES_GRIND_V.pdf\#nameddest=tag-02E9}{V.{E.U4} (p.~63)} }
\expandafter\def\csname SRef@maps:eq:E.U5\endcsname{\href{S6_SERIES_GRIND_V.pdf\#nameddest=tag-02EA}{V.{E.U5} (p.~63)} }
\expandafter\def\csname SRef@maps:eq:E.U6\endcsname{\href{S6_SERIES_GRIND_V.pdf\#nameddest=tag-02EB}{V.{E.U6} (p.~64)} }
\expandafter\def\csname SRef@maps:eq:E.U7\endcsname{\href{S6_SERIES_GRIND_V.pdf\#nameddest=tag-02EC}{V.{E.U7} (p.~64)} }
\expandafter\def\csname SRef@maps:eq:E.U8\endcsname{\href{S6_SERIES_GRIND_V.pdf\#nameddest=tag-02ED}{V.{E.U8} (p.~64)} }
\expandafter\def\csname SRef@maps:eq:E.U9\endcsname{\href{S6_SERIES_GRIND_V.pdf\#nameddest=tag-02EE}{V.{E.U9} (p.~64)} }
\expandafter\def\csname SRef@maps:eq:D.S11\endcsname{\href{S6_SERIES_GRIND_III.pdf\#nameddest=tag-023N}{III.{D.S11} (p.~69)} }
\expandafter\def\csname SRef@maps:eq:E.U10\endcsname{\href{S6_SERIES_GRIND_V.pdf\#nameddest=tag-02E5}{V.{E.U10} (p.~64)} }
\expandafter\def\csname SRef@maps:eq:E.U11\endcsname{\href{S6_SERIES_GRIND_V.pdf\#nameddest=tag-02E6}{V.{E.U11} (p.~64)} }
\expandafter\def\csname SRef@maps:eq:E.H0\endcsname{\href{S6_SERIES_GRIND_V.pdf\#nameddest=tag-02AQ}{V.{E.H0} (p.~65)} }
\expandafter\def\csname SRef@maps:eq:E.H1\endcsname{\href{S6_SERIES_GRIND_V.pdf\#nameddest=tag-02AR}{V.{E.H1} (p.~65)} }
\expandafter\def\csname SRef@maps:eq:E.H2\endcsname{\href{S6_SERIES_GRIND_V.pdf\#nameddest=tag-02AW}{V.{E.H2} (p.~66)} }
\expandafter\def\csname SRef@maps:eq:E.H3\endcsname{\href{S6_SERIES_GRIND_V.pdf\#nameddest=tag-02AX}{V.{E.H3} (p.~66)} }
\expandafter\def\csname SRef@maps:eq:E.H4\endcsname{\href{S6_SERIES_GRIND_V.pdf\#nameddest=tag-02AY}{V.{E.H4} (p.~66)} }
\expandafter\def\csname SRef@maps:eq:E.H5\endcsname{\href{S6_SERIES_GRIND_V.pdf\#nameddest=tag-02AZ}{V.{E.H5} (p.~66)} }
\expandafter\def\csname SRef@maps:eq:E.H6\endcsname{\href{S6_SERIES_GRIND_V.pdf\#nameddest=tag-02B0}{V.{E.H6} (p.~66)} }
\expandafter\def\csname SRef@maps:eq:E.H7\endcsname{\href{S6_SERIES_GRIND_V.pdf\#nameddest=tag-02B1}{V.{E.H7} (p.~66)} }
\expandafter\def\csname SRef@maps:eq:D.S12\endcsname{\href{S6_SERIES_GRIND_III.pdf\#nameddest=tag-023O}{III.{D.S12} (p.~69)} }
\expandafter\def\csname SRef@maps:eq:E.H8\endcsname{\href{S6_SERIES_GRIND_V.pdf\#nameddest=tag-02B2}{V.{E.H8} (p.~67)} }
\expandafter\def\csname SRef@maps:eq:E.H9\endcsname{\href{S6_SERIES_GRIND_V.pdf\#nameddest=tag-02B3}{V.{E.H9} (p.~67)} }
\expandafter\def\csname SRef@maps:eq:E.H9a\endcsname{\href{S6_SERIES_GRIND_V.pdf\#nameddest=tag-02B4}{V.{E.H9a} (p.~67)} }
\expandafter\def\csname SRef@maps:eq:E.H10\endcsname{\href{S6_SERIES_GRIND_V.pdf\#nameddest=tag-02AS}{V.{E.H10} (p.~68)} }
\expandafter\def\csname SRef@maps:eq:E.H11\endcsname{\href{S6_SERIES_GRIND_V.pdf\#nameddest=tag-02AT}{V.{E.H11} (p.~68)} }
\expandafter\def\csname SRef@maps:eq:E.H12\endcsname{\href{S6_SERIES_GRIND_V.pdf\#nameddest=tag-02AU}{V.{E.H12} (p.~68)} }
\expandafter\def\csname SRef@maps:eq:E.H13\endcsname{\href{S6_SERIES_GRIND_V.pdf\#nameddest=tag-02AV}{V.{E.H13} (p.~69)} }
\expandafter\def\csname SRef@maps:eq:D.S13\endcsname{\href{S6_SERIES_GRIND_III.pdf\#nameddest=tag-023P}{III.{D.S13} (p.~70)} }
\expandafter\def\csname SRef@maps:eq:D.S14\endcsname{\href{S6_SERIES_GRIND_III.pdf\#nameddest=tag-023Q}{III.{D.S14} (p.~70)} }
\expandafter\def\csname SRef@maps:eq:D.I1\endcsname{\href{S6_SERIES_GRIND_III.pdf\#nameddest=tag-01ZU}{III.{D.I1} (p.~71)} }
\expandafter\def\csname SRef@maps:eq:D.I2\endcsname{\href{S6_SERIES_GRIND_III.pdf\#nameddest=tag-01ZV}{III.{D.I2} (p.~71)} }
\expandafter\def\csname SRef@maps:eq:E.15\endcsname{\href{S6_SERIES_GRIND_III.pdf\#nameddest=tag-029L}{III.{E.15} (p.~317)} }
\expandafter\def\csname SRef@maps:eq:D.I3\endcsname{\href{S6_SERIES_GRIND_III.pdf\#nameddest=tag-01ZW}{III.{D.I3} (p.~71)} }
\expandafter\def\csname SRef@maps:eq:D.I4\endcsname{\href{S6_SERIES_GRIND_III.pdf\#nameddest=tag-01ZX}{III.{D.I4} (p.~72)} }
\expandafter\def\csname SRef@maps:eq:D.T1\endcsname{\href{S6_SERIES_GRIND_III.pdf\#nameddest=tag-0258}{III.{D.T1} (p.~73)} }
\expandafter\def\csname SRef@maps:eq:D.T2\endcsname{\href{S6_SERIES_GRIND_III.pdf\#nameddest=tag-025A}{III.{D.T2} (p.~74)} }
\expandafter\def\csname SRef@maps:eq:D.T3\endcsname{\href{S6_SERIES_GRIND_III.pdf\#nameddest=tag-025B}{III.{D.T3} (p.~74)} }
\expandafter\def\csname SRef@maps:eq:D.T4\endcsname{\href{S6_SERIES_GRIND_III.pdf\#nameddest=tag-025L}{III.{D.T4} (p.~74)} }
\expandafter\def\csname SRef@maps:eq:D.T5\endcsname{\href{S6_SERIES_GRIND_III.pdf\#nameddest=tag-025M}{III.{D.T5} (p.~75)} }
\expandafter\def\csname SRef@maps:eq:D.T6\endcsname{\href{S6_SERIES_GRIND_III.pdf\#nameddest=tag-025N}{III.{D.T6} (p.~75)} }
\expandafter\def\csname SRef@maps:eq:D.T7\endcsname{\href{S6_SERIES_GRIND_III.pdf\#nameddest=tag-025O}{III.{D.T7} (p.~75)} }
\expandafter\def\csname SRef@maps:eq:D.T8\endcsname{\href{S6_SERIES_GRIND_III.pdf\#nameddest=tag-025P}{III.{D.T8} (p.~75)} }
\expandafter\def\csname SRef@maps:eq:D.T9\endcsname{\href{S6_SERIES_GRIND_III.pdf\#nameddest=tag-025Q}{III.{D.T9} (p.~76)} }
\expandafter\def\csname SRef@maps:eq:D.T10\endcsname{\href{S6_SERIES_GRIND_III.pdf\#nameddest=tag-0259}{III.{D.T10} (p.~76)} }
\expandafter\def\csname SRef@maps:eq:D.V1\endcsname{\href{S6_SERIES_GRIND_III.pdf\#nameddest=tag-0269}{III.{D.V1} (p.~78)} }
\expandafter\def\csname SRef@maps:eq:D.V2\endcsname{\href{S6_SERIES_GRIND_III.pdf\#nameddest=tag-026A}{III.{D.V2} (p.~78)} }
\expandafter\def\csname SRef@maps:eq:D.V3\endcsname{\href{S6_SERIES_GRIND_III.pdf\#nameddest=tag-026B}{III.{D.V3} (p.~78)} }
\expandafter\def\csname SRef@maps:eq:D.V4\endcsname{\href{S6_SERIES_GRIND_III.pdf\#nameddest=tag-026C}{III.{D.V4} (p.~79)} }
\expandafter\def\csname SRef@maps:eq:D.V5\endcsname{\href{S6_SERIES_GRIND_III.pdf\#nameddest=tag-026D}{III.{D.V5} (p.~79)} }
\expandafter\def\csname SRef@maps:eq:D.F1\endcsname{\href{S6_SERIES_GRIND_III.pdf\#nameddest=tag-01YO}{III.{D.F1} (p.~81)} }
\expandafter\def\csname SRef@maps:eq:D.F2\endcsname{\href{S6_SERIES_GRIND_III.pdf\#nameddest=tag-01YP}{III.{D.F2} (p.~81)} }
\expandafter\def\csname SRef@maps:eq:D.O1\endcsname{\href{S6_SERIES_GRIND_III.pdf\#nameddest=tag-020M}{III.{D.O1} (p.~82)} }
\expandafter\def\csname SRef@maps:eq:D.O2\endcsname{\href{S6_SERIES_GRIND_III.pdf\#nameddest=tag-020N}{III.{D.O2} (p.~82)} }
\expandafter\def\csname SRef@maps:eq:D.O3\endcsname{\href{S6_SERIES_GRIND_III.pdf\#nameddest=tag-020O}{III.{D.O3} (p.~82)} }
\expandafter\def\csname SRef@maps:eq:D.O4\endcsname{\href{S6_SERIES_GRIND_III.pdf\#nameddest=tag-020P}{III.{D.O4} (p.~83)} }
\expandafter\def\csname SRef@maps:eq:D.O5\endcsname{\href{S6_SERIES_GRIND_III.pdf\#nameddest=tag-020Q}{III.{D.O5} (p.~83)} }
\expandafter\def\csname SRef@maps:eq:D.O6\endcsname{\href{S6_SERIES_GRIND_III.pdf\#nameddest=tag-020R}{III.{D.O6} (p.~83)} }
\expandafter\def\csname SRef@maps:eq:D.O7\endcsname{\href{S6_SERIES_GRIND_III.pdf\#nameddest=tag-020S}{III.{D.O7} (p.~84)} }
\expandafter\def\csname SRef@maps:eq:D.O8\endcsname{\href{S6_SERIES_GRIND_III.pdf\#nameddest=tag-020T}{III.{D.O8} (p.~84)} }
\expandafter\def\csname SRef@maps:src:D.2.2\endcsname{\href{S6_SERIES_GRIND_III.pdf\#nameddest=tag-02KU}{III.5.12 (p.~42)} }
\expandafter\def\csname SRef@maps:src:D.2.17\endcsname{\href{S6_SERIES_GRIND_III.pdf\#nameddest=tag-02KR}{III.5.26 (p.~77)} }
\expandafter\def\csname SRef@maps:src:D.2.28\endcsname{\href{S6_SERIES_GRIND_III.pdf\#nameddest=tag-02L3}{III.5.37 (p.~106)} }
\expandafter\def\csname SRef@wb:cusp-cone-S12-sphere-map-null\endcsname{\href{S6_SERIES_GRIND_III.pdf\#nameddest=tag-035Z}{III.7.14 (p.~270)} }
\expandafter\def\csname SRef@wb:cusp-cone-S24-sphere-map-null\endcsname{\href{S6_SERIES_GRIND_III.pdf\#nameddest=tag-036E}{III.7.15 (p.~270)} }
\expandafter\def\csname SRef@wb:cusp-cone-literal-R4-open-cells\endcsname{\href{S6_SERIES_GRIND_III.pdf\#nameddest=tag-036V}{III.7.16 (p.~271)} }
\expandafter\def\csname SRef@wb:cusp-cone-sharp-global-R4-scope\endcsname{\href{S6_SERIES_GRIND_III.pdf\#nameddest=tag-037J}{III.7.20 (p.~278)} }
\expandafter\def\csname SRef@wb:cusp-s12-all-higher-smashes-null\endcsname{\href{S6_SERIES_GRIND_III.pdf\#nameddest=tag-038T}{III.7.22 (p.~280)} }
\expandafter\def\csname SRef@wb:cusp_s12_carrier_homotopy:L567\endcsname{\href{S6_SERIES_GRIND_III.pdf\#nameddest=tag-03A6}{III.7.24 (p.~286)} }
\expandafter\def\csname SRef@wb:sphere-metric-unit-volume\endcsname{\href{S6_SERIES_GRIND_III.pdf\#nameddest=tag-03DH}{III.7.28 (p.~294)} }
\expandafter\def\csname SRef@wb:h24:ret:q2-detected-quotient\endcsname{\href{S6_SERIES_GRIND_III.pdf\#nameddest=tag-03BE}{III.7.33 (p.~306)} }
\expandafter\def\csname SRef@wb:cusp-cone-target-slice-embeddings\endcsname{\href{S6_SERIES_GRIND_III.pdf\#nameddest=tag-037N}{III.7.9 (p.~261)} }
\expandafter\def\csname SRef@wb:cusp-cover-target-action\endcsname{\href{S6_SERIES_GRIND_III.pdf\#nameddest=tag-0388}{III.1 (p.~238)} }
\expandafter\def\csname SRef@wb:cusp-cover-normalization-diagram\endcsname{\href{S6_SERIES_GRIND_III.pdf\#nameddest=tag-0383}{III.10 (p.~239)} }
\expandafter\def\csname SRef@wb:cusp-cone-products-not-spheres\endcsname{\href{S6_SERIES_GRIND_III.pdf\#nameddest=tag-0378}{III.100 (p.~263)} }
\expandafter\def\csname SRef@wb:cusp-cone-join-dimensions\endcsname{\href{S6_SERIES_GRIND_III.pdf\#nameddest=tag-036T}{III.101 (p.~263)} }
\expandafter\def\csname SRef@wb:cusp-cone-kernel-contained\endcsname{\href{S6_SERIES_GRIND_III.pdf\#nameddest=tag-036U}{III.102 (p.~263)} }
\expandafter\def\csname SRef@wb:cusp-cone-finite-index-orbit\endcsname{\href{S6_SERIES_GRIND_III.pdf\#nameddest=tag-036Q}{III.103 (p.~263)} }
\expandafter\def\csname SRef@wb:cusp-cone-fibre-H4-nonzero\endcsname{\href{S6_SERIES_GRIND_III.pdf\#nameddest=tag-036P}{III.104 (p.~264)} }
\expandafter\def\csname SRef@wb:cusp-cone-sphere-to-sphere-maps\endcsname{\href{S6_SERIES_GRIND_III.pdf\#nameddest=tag-037L}{III.105 (p.~266)} }
\expandafter\def\csname SRef@wb:cusp-cone-S12-S4-radial-map\endcsname{\href{S6_SERIES_GRIND_III.pdf\#nameddest=tag-035R}{III.106 (p.~267)} }
\expandafter\def\csname SRef@wb:cusp-cone-S24-S4-radial-map\endcsname{\href{S6_SERIES_GRIND_III.pdf\#nameddest=tag-0361}{III.107 (p.~267)} }
\expandafter\def\csname SRef@wb:cusp-cone-S12-exact-collapsed-locus\endcsname{\href{S6_SERIES_GRIND_III.pdf\#nameddest=tag-035T}{III.108 (p.~267)} }
\expandafter\def\csname SRef@wb:cusp-cone-S24-exact-collapsed-locus\endcsname{\href{S6_SERIES_GRIND_III.pdf\#nameddest=tag-0362}{III.109 (p.~267)} }
\expandafter\def\csname SRef@wb:cusp-cover-monodromy-conjugacy\endcsname{\href{S6_SERIES_GRIND_III.pdf\#nameddest=tag-0381}{III.11 (p.~239)} }
\expandafter\def\csname SRef@wb:cusp-cone-moving-S4-R4-limits\endcsname{\href{S6_SERIES_GRIND_III.pdf\#nameddest=tag-036Y}{III.110 (p.~267)} }
\expandafter\def\csname SRef@wb:cusp-cone-transported-skeleton-collapse\endcsname{\href{S6_SERIES_GRIND_III.pdf\#nameddest=tag-037O}{III.111 (p.~268)} }
\expandafter\def\csname SRef@wb:cusp-cone-distance-to-skeleton\endcsname{\href{S6_SERIES_GRIND_III.pdf\#nameddest=tag-036M}{III.112 (p.~268)} }
\expandafter\def\csname SRef@wb:cusp-cone-S12-onto-exact-ball\endcsname{\href{S6_SERIES_GRIND_III.pdf\#nameddest=tag-035W}{III.113 (p.~269)} }
\expandafter\def\csname SRef@wb:cusp-cone-S24-onto-exact-ball\endcsname{\href{S6_SERIES_GRIND_III.pdf\#nameddest=tag-0369}{III.114 (p.~270)} }
\expandafter\def\csname SRef@wb:cusp-cone-S24-null-maps\endcsname{\href{S6_SERIES_GRIND_III.pdf\#nameddest=tag-0367}{III.115 (p.~270)} }
\expandafter\def\csname SRef@wb:cusp-cone-S24-null-special-cases\endcsname{\href{S6_SERIES_GRIND_III.pdf\#nameddest=tag-0368}{III.116 (p.~270)} }
\expandafter\def\csname SRef@wb:cusp-cone-S24-smash-square-factorization\endcsname{\href{S6_SERIES_GRIND_III.pdf\#nameddest=tag-036C}{III.117 (p.~271)} }
\expandafter\def\csname SRef@wb:cusp-cone-S24-sixfold-suspension\endcsname{\href{S6_SERIES_GRIND_III.pdf\#nameddest=tag-036B}{III.118 (p.~271)} }
\expandafter\def\csname SRef@wb:cusp-cone-fourth-stable-stem\endcsname{\href{S6_SERIES_GRIND_III.pdf\#nameddest=tag-036R}{III.119 (p.~271)} }
\expandafter\def\csname SRef@wb:cusp-cover-maximal-chart-map\endcsname{\href{S6_SERIES_GRIND_III.pdf\#nameddest=tag-0380}{III.12 (p.~241)} }
\expandafter\def\csname SRef@wb:cusp-cone-U12-homeomorphism\endcsname{\href{S6_SERIES_GRIND_III.pdf\#nameddest=tag-036G}{III.121 (p.~271)} }
\expandafter\def\csname SRef@wb:cusp-cone-U24-homeomorphism\endcsname{\href{S6_SERIES_GRIND_III.pdf\#nameddest=tag-036H}{III.122 (p.~271)} }
\expandafter\def\csname SRef@wb:cusp-cone-D12-embedding\endcsname{\href{S6_SERIES_GRIND_III.pdf\#nameddest=tag-035G}{III.123 (p.~273)} }
\expandafter\def\csname SRef@wb:cusp-cone-D24-embedding\endcsname{\href{S6_SERIES_GRIND_III.pdf\#nameddest=tag-035I}{III.124 (p.~273)} }
\expandafter\def\csname SRef@wb:cusp-cone-D12-slice\endcsname{\href{S6_SERIES_GRIND_III.pdf\#nameddest=tag-035H}{III.125 (p.~273)} }
\expandafter\def\csname SRef@wb:cusp-cone-D24-slice\endcsname{\href{S6_SERIES_GRIND_III.pdf\#nameddest=tag-035J}{III.126 (p.~273)} }
\expandafter\def\csname SRef@wb:cusp-cone-carrier-local-homology\endcsname{\href{S6_SERIES_GRIND_III.pdf\#nameddest=tag-036K}{III.127 (p.~273)} }
\expandafter\def\csname SRef@wb:cusp-cone-A-smash-wedge\endcsname{\href{S6_SERIES_GRIND_III.pdf\#nameddest=tag-0357}{III.128 (p.~274)} }
\expandafter\def\csname SRef@wb:cusp-cone-C-smash-wedge\endcsname{\href{S6_SERIES_GRIND_III.pdf\#nameddest=tag-035F}{III.129 (p.~274)} }
\expandafter\def\csname SRef@wb:cusp-cone-E12-fibres\endcsname{\href{S6_SERIES_GRIND_III.pdf\#nameddest=tag-035K}{III.130 (p.~276)} }
\expandafter\def\csname SRef@wb:cusp-cone-E24-fibres\endcsname{\href{S6_SERIES_GRIND_III.pdf\#nameddest=tag-035L}{III.131 (p.~276)} }
\expandafter\def\csname SRef@wb:cusp-s12-smash-factorization\endcsname{\href{S6_SERIES_GRIND_III.pdf\#nameddest=tag-0390}{III.132 (p.~280)} }
\expandafter\def\csname SRef@wb:cusp-s12-toda-smash-formula\endcsname{\href{S6_SERIES_GRIND_III.pdf\#nameddest=tag-0392}{III.133 (p.~280)} }
\expandafter\def\csname SRef@wb:cusp-s12-three-stem-groups\endcsname{\href{S6_SERIES_GRIND_III.pdf\#nameddest=tag-0391}{III.134 (p.~280)} }
\expandafter\def\csname SRef@wb:cusp-s12-sixfold-suspension-class\endcsname{\href{S6_SERIES_GRIND_III.pdf\#nameddest=tag-038Z}{III.135 (p.~281)} }
\expandafter\def\csname SRef@wb:cusp-s12-regular-inverse-image\endcsname{\href{S6_SERIES_GRIND_III.pdf\#nameddest=tag-038Y}{III.136 (p.~281)} }
\expandafter\def\csname SRef@wb:cusp-null:eq:periods\endcsname{\href{S6_SERIES_GRIND_III.pdf\#nameddest=tag-038N}{III.137 (p.~282)} }
\expandafter\def\csname SRef@wb:cusp-null:eq:action\endcsname{\href{S6_SERIES_GRIND_III.pdf\#nameddest=tag-038E}{III.138 (p.~282)} }
\expandafter\def\csname SRef@wb:cusp-null:eq:g\endcsname{\href{S6_SERIES_GRIND_III.pdf\#nameddest=tag-038J}{III.139 (p.~282)} }
\expandafter\def\csname SRef@wb:cusp-null:eq:chart\endcsname{\href{S6_SERIES_GRIND_III.pdf\#nameddest=tag-038H}{III.140 (p.~282)} }
\expandafter\def\csname SRef@wb:cusp-null:eq:moduli\endcsname{\href{S6_SERIES_GRIND_III.pdf\#nameddest=tag-038L}{III.141 (p.~282)} }
\expandafter\def\csname SRef@wb:cusp-null:eq:invertibility\endcsname{\href{S6_SERIES_GRIND_III.pdf\#nameddest=tag-038K}{III.142 (p.~283)} }
\expandafter\def\csname SRef@wb:cusp-null:eq:straightening\endcsname{\href{S6_SERIES_GRIND_III.pdf\#nameddest=tag-038Q}{III.144 (p.~283)} }
\expandafter\def\csname SRef@wb:cusp-null:eq:torus\endcsname{\href{S6_SERIES_GRIND_III.pdf\#nameddest=tag-038R}{III.145 (p.~283)} }
\expandafter\def\csname SRef@wb:cusp-null:eq:spatial-bound\endcsname{\href{S6_SERIES_GRIND_III.pdf\#nameddest=tag-038P}{III.146 (p.~283)} }
\expandafter\def\csname SRef@wb:cusp-null:eq:nearby-moduli\endcsname{\href{S6_SERIES_GRIND_III.pdf\#nameddest=tag-038M}{III.147 (p.~284)} }
\expandafter\def\csname SRef@wb:cusp-null:eq:J\endcsname{\href{S6_SERIES_GRIND_III.pdf\#nameddest=tag-038B}{III.148 (p.~284)} }
\expandafter\def\csname SRef@wb:cusp-null:eq:actual-frame\endcsname{\href{S6_SERIES_GRIND_III.pdf\#nameddest=tag-038F}{III.149 (p.~284)} }
\expandafter\def\csname SRef@wb:cusp-cover-equivariance-computation\endcsname{\href{S6_SERIES_GRIND_III.pdf\#nameddest=tag-037W}{III.15 (p.~241)} }
\expandafter\def\csname SRef@wb:cusp-null:eq:J0\endcsname{\href{S6_SERIES_GRIND_III.pdf\#nameddest=tag-038C}{III.150 (p.~284)} }
\expandafter\def\csname SRef@wb:cusp-null:eq:P\endcsname{\href{S6_SERIES_GRIND_III.pdf\#nameddest=tag-038D}{III.151 (p.~284)} }
\expandafter\def\csname SRef@wb:cusp-null:eq:q\endcsname{\href{S6_SERIES_GRIND_III.pdf\#nameddest=tag-038O}{III.152 (p.~285)} }
\expandafter\def\csname SRef@wb:cusp-null:eq:bordism\endcsname{\href{S6_SERIES_GRIND_III.pdf\#nameddest=tag-038G}{III.153 (p.~285)} }
\expandafter\def\csname SRef@wb:cusp-null:eq:extended-frame\endcsname{\href{S6_SERIES_GRIND_III.pdf\#nameddest=tag-038I}{III.154 (p.~285)} }
\expandafter\def\csname SRef@wb:sphere-metric-normalized-Q-limit\endcsname{\href{S6_SERIES_GRIND_III.pdf\#nameddest=tag-03DA}{III.155 (p.~287)} }
\expandafter\def\csname SRef@wb:sphere-metric-Hausdorff-collapse\endcsname{\href{S6_SERIES_GRIND_III.pdf\#nameddest=tag-03D2}{III.158 (p.~288)} }
\expandafter\def\csname SRef@wb:sphere-metric-A-fixed-metric\endcsname{\href{S6_SERIES_GRIND_III.pdf\#nameddest=tag-03CW}{III.159 (p.~288)} }
\expandafter\def\csname SRef@wb:cusp-cover-basechange-local-ring\endcsname{\href{S6_SERIES_GRIND_III.pdf\#nameddest=tag-037S}{III.16 (p.~242)} }
\expandafter\def\csname SRef@wb:sphere-metric-BC-fixed-metric\endcsname{\href{S6_SERIES_GRIND_III.pdf\#nameddest=tag-03CY}{III.160 (p.~288)} }
\expandafter\def\csname SRef@wb:sphere-metric-A-slice-exact\endcsname{\href{S6_SERIES_GRIND_III.pdf\#nameddest=tag-03CX}{III.161 (p.~289)} }
\expandafter\def\csname SRef@wb:sphere-metric-BC-slice-exact\endcsname{\href{S6_SERIES_GRIND_III.pdf\#nameddest=tag-03CZ}{III.162 (p.~289)} }
\expandafter\def\csname SRef@wb:sphere-metric-slice-fixed-blowups\endcsname{\href{S6_SERIES_GRIND_III.pdf\#nameddest=tag-03DE}{III.164 (p.~289)} }
\expandafter\def\csname SRef@wb:sphere-metric-ambient-fixed-blowups\endcsname{\href{S6_SERIES_GRIND_III.pdf\#nameddest=tag-03D3}{III.166 (p.~289)} }
\expandafter\def\csname SRef@wb:sphere-metric-scaling-trichotomy-result\endcsname{\href{S6_SERIES_GRIND_III.pdf\#nameddest=tag-03DD}{III.167 (p.~291)} }
\expandafter\def\csname SRef@wb:sphere-metric-vertex-distance\endcsname{\href{S6_SERIES_GRIND_III.pdf\#nameddest=tag-03DI}{III.168 (p.~292)} }
\expandafter\def\csname SRef@wb:sphere-metric-intrinsic-ambient-agreement\endcsname{\href{S6_SERIES_GRIND_III.pdf\#nameddest=tag-03D7}{III.169 (p.~292)} }
\expandafter\def\csname SRef@wb:cusp-cover-invariant-ring\endcsname{\href{S6_SERIES_GRIND_III.pdf\#nameddest=tag-037Z}{III.17 (p.~242)} }
\expandafter\def\csname SRef@wb:sphere-metric-intrinsic-ambient-ball-comparison\endcsname{\href{S6_SERIES_GRIND_III.pdf\#nameddest=tag-03D8}{III.170 (p.~292)} }
\expandafter\def\csname SRef@wb:sphere-metric-arbitrary-unrolling-ambient-bubbles\endcsname{\href{S6_SERIES_GRIND_III.pdf\#nameddest=tag-03D4}{III.172 (p.~292)} }
\expandafter\def\csname SRef@wb:sphere-metric-outward-shortcut-exclusion\endcsname{\href{S6_SERIES_GRIND_III.pdf\#nameddest=tag-03DB}{III.173 (p.~293)} }
\expandafter\def\csname SRef@wb:sphere-metric-C-smash-local-homology\endcsname{\href{S6_SERIES_GRIND_III.pdf\#nameddest=tag-03D1}{III.174 (p.~295)} }
\expandafter\def\csname SRef@wb:sphere-metric-C-boundary-quotient-homeomorphism\endcsname{\href{S6_SERIES_GRIND_III.pdf\#nameddest=tag-03D0}{III.175 (p.~296)} }
\expandafter\def\csname SRef@wb:sphere-metric-boundary-link-local-shift\endcsname{\href{S6_SERIES_GRIND_III.pdf\#nameddest=tag-03D5}{III.176 (p.~296)} }
\expandafter\def\csname SRef@wb:h24:ret:eq:original-periods\endcsname{\href{S6_SERIES_GRIND_III.pdf\#nameddest=tag-03B4}{III.177 (p.~298)} }
\expandafter\def\csname SRef@wb:h24:ret:eq:A\endcsname{\href{S6_SERIES_GRIND_III.pdf\#nameddest=tag-03AF}{III.178 (p.~298)} }
\expandafter\def\csname SRef@wb:h24:ret:eq:deck\endcsname{\href{S6_SERIES_GRIND_III.pdf\#nameddest=tag-03AS}{III.179 (p.~298)} }
\expandafter\def\csname SRef@wb:h24:ret:eq:logx\endcsname{\href{S6_SERIES_GRIND_III.pdf\#nameddest=tag-03B0}{III.180 (p.~298)} }
\expandafter\def\csname SRef@wb:h24:ret:eq:invertibility\endcsname{\href{S6_SERIES_GRIND_III.pdf\#nameddest=tag-03AZ}{III.181 (p.~298)} }
\expandafter\def\csname SRef@wb:h24:ret:eq:native\endcsname{\href{S6_SERIES_GRIND_III.pdf\#nameddest=tag-03B1}{III.182 (p.~298)} }
\expandafter\def\csname SRef@wb:h24:ret:eq:two-skeleton\endcsname{\href{S6_SERIES_GRIND_III.pdf\#nameddest=tag-03B9}{III.183 (p.~298)} }
\expandafter\def\csname SRef@wb:h24:ret:eq:A-target\endcsname{\href{S6_SERIES_GRIND_III.pdf\#nameddest=tag-03AH}{III.184 (p.~299)} }
\expandafter\def\csname SRef@wb:h24:ret:eq:q2\endcsname{\href{S6_SERIES_GRIND_III.pdf\#nameddest=tag-03B5}{III.185 (p.~299)} }
\expandafter\def\csname SRef@wb:h24:ret:eq:H\endcsname{\href{S6_SERIES_GRIND_III.pdf\#nameddest=tag-03AN}{III.186 (p.~299)} }
\expandafter\def\csname SRef@wb:h24:ret:eq:four-cell\endcsname{\href{S6_SERIES_GRIND_III.pdf\#nameddest=tag-03AW}{III.187 (p.~299)} }
\expandafter\def\csname SRef@wb:h24:ret:eq:chosen-value\endcsname{\href{S6_SERIES_GRIND_III.pdf\#nameddest=tag-03AR}{III.188 (p.~299)} }
\expandafter\def\csname SRef@wb:h24:ret:eq:z-chart\endcsname{\href{S6_SERIES_GRIND_III.pdf\#nameddest=tag-03BB}{III.189 (p.~299)} }
\expandafter\def\csname SRef@wb:cusp-cover-triple-fibre-length\endcsname{\href{S6_SERIES_GRIND_III.pdf\#nameddest=tag-038A}{III.19 (p.~242)} }
\expandafter\def\csname SRef@wb:h24:ret:eq:units\endcsname{\href{S6_SERIES_GRIND_III.pdf\#nameddest=tag-03BA}{III.190 (p.~299)} }
\expandafter\def\csname SRef@wb:h24:ret:eq:F\endcsname{\href{S6_SERIES_GRIND_III.pdf\#nameddest=tag-03AL}{III.191 (p.~299)} }
\expandafter\def\csname SRef@wb:h24:ret:eq:Finv\endcsname{\href{S6_SERIES_GRIND_III.pdf\#nameddest=tag-03AM}{III.192 (p.~299)} }
\expandafter\def\csname SRef@wb:h24:ret:eq:M\endcsname{\href{S6_SERIES_GRIND_III.pdf\#nameddest=tag-03AP}{III.193 (p.~300)} }
\expandafter\def\csname SRef@wb:h24:ret:eq:exact-u\endcsname{\href{S6_SERIES_GRIND_III.pdf\#nameddest=tag-03AV}{III.194 (p.~300)} }
\expandafter\def\csname SRef@wb:h24:ret:eq:full-derivative\endcsname{\href{S6_SERIES_GRIND_III.pdf\#nameddest=tag-03AX}{III.195 (p.~300)} }
\expandafter\def\csname SRef@wb:h24:ret:eq:derivative-on-M\endcsname{\href{S6_SERIES_GRIND_III.pdf\#nameddest=tag-03AT}{III.196 (p.~300)} }
\expandafter\def\csname SRef@wb:h24:ret:eq:LQ\endcsname{\href{S6_SERIES_GRIND_III.pdf\#nameddest=tag-03AO}{III.197 (p.~301)} }
\expandafter\def\csname SRef@wb:h24:ret:eq:tangents\endcsname{\href{S6_SERIES_GRIND_III.pdf\#nameddest=tag-03B8}{III.198 (p.~301)} }
\expandafter\def\csname SRef@wb:h24:ret:eq:normal-first\endcsname{\href{S6_SERIES_GRIND_III.pdf\#nameddest=tag-03B2}{III.199 (p.~301)} }
\expandafter\def\csname SRef@wb:cusp-cover-base-change\endcsname{\href{S6_SERIES_GRIND_III.pdf\#nameddest=tag-037Q}{III.2 (p.~238)} }
\expandafter\def\csname SRef@wb:cusp-cover-double-fibre-length\endcsname{\href{S6_SERIES_GRIND_III.pdf\#nameddest=tag-037V}{III.20 (p.~242)} }
\expandafter\def\csname SRef@wb:h24:ret:eq:normal-u\endcsname{\href{S6_SERIES_GRIND_III.pdf\#nameddest=tag-03B3}{III.200 (p.~301)} }
\expandafter\def\csname SRef@wb:h24:ret:eq:differential-table\endcsname{\href{S6_SERIES_GRIND_III.pdf\#nameddest=tag-03AU}{III.201 (p.~301)} }
\expandafter\def\csname SRef@wb:h24:ret:eq:A-homotopy\endcsname{\href{S6_SERIES_GRIND_III.pdf\#nameddest=tag-03AG}{III.202 (p.~301)} }
\expandafter\def\csname SRef@wb:h24:ret:eq:sign\endcsname{\href{S6_SERIES_GRIND_III.pdf\#nameddest=tag-03B7}{III.203 (p.~302)} }
\expandafter\def\csname SRef@wb:h24:ret:eq:D\endcsname{\href{S6_SERIES_GRIND_III.pdf\#nameddest=tag-03AJ}{III.204 (p.~302)} }
\expandafter\def\csname SRef@wb:h24:ret:eq:full-frame\endcsname{\href{S6_SERIES_GRIND_III.pdf\#nameddest=tag-03AY}{III.205 (p.~302)} }
\expandafter\def\csname SRef@wb:h24:ret:eq:D-homotopy\endcsname{\href{S6_SERIES_GRIND_III.pdf\#nameddest=tag-03AK}{III.206 (p.~302)} }
\expandafter\def\csname SRef@wb:h24:ret:eq:quadratic\endcsname{\href{S6_SERIES_GRIND_III.pdf\#nameddest=tag-03B6}{III.207 (p.~303)} }
\expandafter\def\csname SRef@wb:h24:ret:eq:Arf\endcsname{\href{S6_SERIES_GRIND_III.pdf\#nameddest=tag-03AI}{III.208 (p.~303)} }
\expandafter\def\csname SRef@wb:h24:ret:eq:actual-Arf\endcsname{\href{S6_SERIES_GRIND_III.pdf\#nameddest=tag-03AQ}{III.209 (p.~304)} }
\expandafter\def\csname SRef@wb:cusp-cover-additive-fibre-map\endcsname{\href{S6_SERIES_GRIND_III.pdf\#nameddest=tag-037P}{III.21 (p.~243)} }
\expandafter\def\csname SRef@mapsD:repair:cocycle-relation\endcsname{\href{S6_SERIES_GRIND_III.pdf\#nameddest=tag-02OF}{III.210 (p.~310)} }
\expandafter\def\csname SRef@mapsD:repair:canonical-matrices\endcsname{\href{S6_SERIES_GRIND_III.pdf\#nameddest=tag-02OE}{III.211 (p.~312)} }
\expandafter\def\csname SRef@mapsD:repair:free-peripherals\endcsname{\href{S6_SERIES_GRIND_III.pdf\#nameddest=tag-02OH}{III.212 (p.~312)} }
\expandafter\def\csname SRef@mapsD:repair:50-weierstrass\endcsname{\href{S6_SERIES_GRIND_III.pdf\#nameddest=tag-02OD}{III.213 (p.~312)} }
\expandafter\def\csname SRef@mapsD:repair:height-formula\endcsname{\href{S6_SERIES_GRIND_III.pdf\#nameddest=tag-02OI}{III.214 (p.~315)} }
\expandafter\def\csname SRef@wb:cusp-spec-radial-notation\endcsname{\href{S6_SERIES_GRIND_III.pdf\#nameddest=tag-03A3}{III.22 (p.~244)} }
\expandafter\def\csname SRef@wb:cusp-spec-holomorphic-errors\endcsname{\href{S6_SERIES_GRIND_III.pdf\#nameddest=tag-039V}{III.23 (p.~244)} }
\expandafter\def\csname SRef@wb:cusp-spec-full-inverse-transpose\endcsname{\href{S6_SERIES_GRIND_III.pdf\#nameddest=tag-039O}{III.24 (p.~244)} }
\expandafter\def\csname SRef@wb:cusp-spec-global-circle-map\endcsname{\href{S6_SERIES_GRIND_III.pdf\#nameddest=tag-039R}{III.25 (p.~246)} }
\expandafter\def\csname SRef@wb:cusp-spec-global-circle-quotient\endcsname{\href{S6_SERIES_GRIND_III.pdf\#nameddest=tag-039S}{III.26 (p.~246)} }
\expandafter\def\csname SRef@wb:cusp-spec-global-circle-coordinate\endcsname{\href{S6_SERIES_GRIND_III.pdf\#nameddest=tag-039Q}{III.27 (p.~246)} }
\expandafter\def\csname SRef@wb:cusp-spec-global-kernel-family\endcsname{\href{S6_SERIES_GRIND_III.pdf\#nameddest=tag-039T}{III.28 (p.~246)} }
\expandafter\def\csname SRef@wb:cusp-spec-radial-circle-character\endcsname{\href{S6_SERIES_GRIND_III.pdf\#nameddest=tag-03A2}{III.29 (p.~247)} }
\expandafter\def\csname SRef@wb:cusp-cover-source-factor\endcsname{\href{S6_SERIES_GRIND_III.pdf\#nameddest=tag-0387}{III.3 (p.~238)} }
\expandafter\def\csname SRef@wb:cusp-spec-global-circle-circumference\endcsname{\href{S6_SERIES_GRIND_III.pdf\#nameddest=tag-039P}{III.30 (p.~247)} }
\expandafter\def\csname SRef@wb:cusp-spec-fourfold-fibre-product\endcsname{\href{S6_SERIES_GRIND_III.pdf\#nameddest=tag-039N}{III.31 (p.~247)} }
\expandafter\def\csname SRef@wb:cusp-spec-fourfold-circle-map\endcsname{\href{S6_SERIES_GRIND_III.pdf\#nameddest=tag-039M}{III.32 (p.~247)} }
\expandafter\def\csname SRef@wb:cusp-spec-four-circle-fibre\endcsname{\href{S6_SERIES_GRIND_III.pdf\#nameddest=tag-039J}{III.33 (p.~247)} }
\expandafter\def\csname SRef@wb:cusp-spec-four-circle-limit\endcsname{\href{S6_SERIES_GRIND_III.pdf\#nameddest=tag-039K}{III.34 (p.~247)} }
\expandafter\def\csname SRef@wb:cusp-spec-doubled-local-quotient\endcsname{\href{S6_SERIES_GRIND_III.pdf\#nameddest=tag-0398}{III.35 (p.~247)} }
\expandafter\def\csname SRef@wb:cusp-spec-finite-affine-action\endcsname{\href{S6_SERIES_GRIND_III.pdf\#nameddest=tag-039A}{III.36 (p.~248)} }
\expandafter\def\csname SRef@wb:cusp-spec-finite-character-shift\endcsname{\href{S6_SERIES_GRIND_III.pdf\#nameddest=tag-039B}{III.37 (p.~248)} }
\expandafter\def\csname SRef@wb:cusp-spec-finite-Seifert-target\endcsname{\href{S6_SERIES_GRIND_III.pdf\#nameddest=tag-0399}{III.38 (p.~248)} }
\expandafter\def\csname SRef@wb:cusp-spec-finite-fourfold-normalization\endcsname{\href{S6_SERIES_GRIND_III.pdf\#nameddest=tag-039D}{III.39 (p.~248)} }
\expandafter\def\csname SRef@wb:cusp-cover-source-action\endcsname{\href{S6_SERIES_GRIND_III.pdf\#nameddest=tag-0386}{III.4 (p.~238)} }
\expandafter\def\csname SRef@wb:cusp-spec-finite-fourfold-target\endcsname{\href{S6_SERIES_GRIND_III.pdf\#nameddest=tag-039E}{III.40 (p.~248)} }
\expandafter\def\csname SRef@wb:cusp-spec-finite-normalization-map\endcsname{\href{S6_SERIES_GRIND_III.pdf\#nameddest=tag-039F}{III.41 (p.~248)} }
\expandafter\def\csname SRef@wb:cusp-spec-product-not-sphere\endcsname{\href{S6_SERIES_GRIND_III.pdf\#nameddest=tag-03A1}{III.42 (p.~248)} }
\expandafter\def\csname SRef@wb:cusp-spec-four-product-not-sphere\endcsname{\href{S6_SERIES_GRIND_III.pdf\#nameddest=tag-039L}{III.43 (p.~248)} }
\expandafter\def\csname SRef@wb:cusp-spec-smash-spheres\endcsname{\href{S6_SERIES_GRIND_III.pdf\#nameddest=tag-03A4}{III.44 (p.~249)} }
\expandafter\def\csname SRef@wb:cusp-spec-target-smash-map\endcsname{\href{S6_SERIES_GRIND_III.pdf\#nameddest=tag-03A5}{III.45 (p.~249)} }
\expandafter\def\csname SRef@wb:cusp-spec-cusp-kernel\endcsname{\href{S6_SERIES_GRIND_III.pdf\#nameddest=tag-0394}{III.46 (p.~249)} }
\expandafter\def\csname SRef@wb:cusp-spec-invariant-flag-calculation\endcsname{\href{S6_SERIES_GRIND_III.pdf\#nameddest=tag-039X}{III.47 (p.~249)} }
\expandafter\def\csname SRef@wb:cusp-spec-invariant-character-lattice\endcsname{\href{S6_SERIES_GRIND_III.pdf\#nameddest=tag-039W}{III.48 (p.~249)} }
\expandafter\def\csname SRef@wb:cusp-spec-circle-dual-vector\endcsname{\href{S6_SERIES_GRIND_III.pdf\#nameddest=tag-0393}{III.49 (p.~250)} }
\expandafter\def\csname SRef@wb:cusp-cover-toric-power-map\endcsname{\href{S6_SERIES_GRIND_III.pdf\#nameddest=tag-0389}{III.5 (p.~239)} }
\expandafter\def\csname SRef@wb:cusp-spec-no-finite-cover-matrix\endcsname{\href{S6_SERIES_GRIND_III.pdf\#nameddest=tag-039Z}{III.50 (p.~251)} }
\expandafter\def\csname SRef@wb:cusp-spec-no-finite-cover-orbit\endcsname{\href{S6_SERIES_GRIND_III.pdf\#nameddest=tag-03A0}{III.51 (p.~251)} }
\expandafter\def\csname SRef@wb:cusp-spec-finite-target-conjugacy\endcsname{\href{S6_SERIES_GRIND_III.pdf\#nameddest=tag-039I}{III.52 (p.~251)} }
\expandafter\def\csname SRef@wb:cusp-spec-finite-target-cocycle\endcsname{\href{S6_SERIES_GRIND_III.pdf\#nameddest=tag-039H}{III.53 (p.~251)} }
\expandafter\def\csname SRef@wb:cusp-spec-finite-overlap-identity\endcsname{\href{S6_SERIES_GRIND_III.pdf\#nameddest=tag-039G}{III.54 (p.~251)} }
\expandafter\def\csname SRef@wb:cusp-spec-finite-fibre-product-branches\endcsname{\href{S6_SERIES_GRIND_III.pdf\#nameddest=tag-039C}{III.55 (p.~252)} }
\expandafter\def\csname SRef@wb:cusp-spec-cusp-nonextension-position\endcsname{\href{S6_SERIES_GRIND_III.pdf\#nameddest=tag-0397}{III.56 (p.~252)} }
\expandafter\def\csname SRef@wb:cusp-spec-cusp-nonextension-character\endcsname{\href{S6_SERIES_GRIND_III.pdf\#nameddest=tag-0395}{III.57 (p.~252)} }
\expandafter\def\csname SRef@wb:cusp-spec-cusp-nonextension-chart\endcsname{\href{S6_SERIES_GRIND_III.pdf\#nameddest=tag-0396}{III.58 (p.~253)} }
\expandafter\def\csname SRef@wb:cusp-cone-disc\endcsname{\href{S6_SERIES_GRIND_III.pdf\#nameddest=tag-036L}{III.59 (p.~255)} }
\expandafter\def\csname SRef@wb:cusp-cover-completed-map\endcsname{\href{S6_SERIES_GRIND_III.pdf\#nameddest=tag-037T}{III.6 (p.~239)} }
\expandafter\def\csname SRef@wb:cusp-cone-quotient-lattice\endcsname{\href{S6_SERIES_GRIND_III.pdf\#nameddest=tag-037F}{III.60 (p.~255)} }
\expandafter\def\csname SRef@wb:cusp-cone-native-family-map\endcsname{\href{S6_SERIES_GRIND_III.pdf\#nameddest=tag-0371}{III.61 (p.~255)} }
\expandafter\def\csname SRef@wb:cusp-cone-radial-isometry\endcsname{\href{S6_SERIES_GRIND_III.pdf\#nameddest=tag-037G}{III.62 (p.~255)} }
\expandafter\def\csname SRef@wb:cusp-cone-native-restriction\endcsname{\href{S6_SERIES_GRIND_III.pdf\#nameddest=tag-0372}{III.63 (p.~255)} }
\expandafter\def\csname SRef@wb:cusp-cone-native-systole\endcsname{\href{S6_SERIES_GRIND_III.pdf\#nameddest=tag-0373}{III.64 (p.~255)} }
\expandafter\def\csname SRef@wb:cusp-cone-monodromy-image\endcsname{\href{S6_SERIES_GRIND_III.pdf\#nameddest=tag-036X}{III.65 (p.~255)} }
\expandafter\def\csname SRef@wb:cusp-cone-A-definition\endcsname{\href{S6_SERIES_GRIND_III.pdf\#nameddest=tag-0354}{III.66 (p.~256)} }
\expandafter\def\csname SRef@wb:cusp-cone-A-suspension\endcsname{\href{S6_SERIES_GRIND_III.pdf\#nameddest=tag-0358}{III.67 (p.~256)} }
\expandafter\def\csname SRef@wb:cusp-cone-q2\endcsname{\href{S6_SERIES_GRIND_III.pdf\#nameddest=tag-037D}{III.68 (p.~256)} }
\expandafter\def\csname SRef@wb:cusp-cone-q2-formula\endcsname{\href{S6_SERIES_GRIND_III.pdf\#nameddest=tag-037E}{III.69 (p.~256)} }
\expandafter\def\csname SRef@wb:cusp-cover-direct-ramification\endcsname{\href{S6_SERIES_GRIND_III.pdf\#nameddest=tag-037U}{III.7 (p.~239)} }
\expandafter\def\csname SRef@wb:cusp-cone-S12-map\endcsname{\href{S6_SERIES_GRIND_III.pdf\#nameddest=tag-035V}{III.70 (p.~256)} }
\expandafter\def\csname SRef@wb:cusp-cone-S12-source\endcsname{\href{S6_SERIES_GRIND_III.pdf\#nameddest=tag-035Y}{III.71 (p.~256)} }
\expandafter\def\csname SRef@wb:cusp-cone-S12-fibre-map\endcsname{\href{S6_SERIES_GRIND_III.pdf\#nameddest=tag-035U}{III.72 (p.~256)} }
\expandafter\def\csname SRef@wb:cusp-cone-S12-R4-limit\endcsname{\href{S6_SERIES_GRIND_III.pdf\#nameddest=tag-035Q}{III.73 (p.~256)} }
\expandafter\def\csname SRef@wb:cusp-cone-A-homology\endcsname{\href{S6_SERIES_GRIND_III.pdf\#nameddest=tag-0355}{III.74 (p.~256)} }
\expandafter\def\csname SRef@wb:cusp-cone-A-smash-homology\endcsname{\href{S6_SERIES_GRIND_III.pdf\#nameddest=tag-0356}{III.75 (p.~257)} }
\expandafter\def\csname SRef@wb:cusp-cone-B-definition\endcsname{\href{S6_SERIES_GRIND_III.pdf\#nameddest=tag-0359}{III.76 (p.~258)} }
\expandafter\def\csname SRef@wb:cusp-cone-q1-local\endcsname{\href{S6_SERIES_GRIND_III.pdf\#nameddest=tag-037B}{III.77 (p.~258)} }
\expandafter\def\csname SRef@wb:cusp-cone-q1-local-formula\endcsname{\href{S6_SERIES_GRIND_III.pdf\#nameddest=tag-037C}{III.78 (p.~258)} }
\expandafter\def\csname SRef@wb:cusp-cone-S24-local-fibre-map\endcsname{\href{S6_SERIES_GRIND_III.pdf\#nameddest=tag-0364}{III.79 (p.~258)} }
\expandafter\def\csname SRef@wb:cusp-cover-base-changed-target\endcsname{\href{S6_SERIES_GRIND_III.pdf\#nameddest=tag-037R}{III.8 (p.~239)} }
\expandafter\def\csname SRef@wb:cusp-cone-B-homology\endcsname{\href{S6_SERIES_GRIND_III.pdf\#nameddest=tag-035A}{III.80 (p.~258)} }
\expandafter\def\csname SRef@wb:cusp-cone-B-smash-homology\endcsname{\href{S6_SERIES_GRIND_III.pdf\#nameddest=tag-035B}{III.81 (p.~258)} }
\expandafter\def\csname SRef@wb:cusp-cone-pair-of-pants\endcsname{\href{S6_SERIES_GRIND_III.pdf\#nameddest=tag-0377}{III.82 (p.~259)} }
\expandafter\def\csname SRef@wb:cusp-cone-M-definition\endcsname{\href{S6_SERIES_GRIND_III.pdf\#nameddest=tag-035P}{III.83 (p.~259)} }
\expandafter\def\csname SRef@wb:cusp-cone-M-boundary\endcsname{\href{S6_SERIES_GRIND_III.pdf\#nameddest=tag-035N}{III.84 (p.~259)} }
\expandafter\def\csname SRef@wb:cusp-cone-C-definition\endcsname{\href{S6_SERIES_GRIND_III.pdf\#nameddest=tag-035C}{III.85 (p.~259)} }
\expandafter\def\csname SRef@wb:cusp-cone-q1\endcsname{\href{S6_SERIES_GRIND_III.pdf\#nameddest=tag-0379}{III.86 (p.~259)} }
\expandafter\def\csname SRef@wb:cusp-cone-q1-cusp-formula\endcsname{\href{S6_SERIES_GRIND_III.pdf\#nameddest=tag-037A}{III.87 (p.~259)} }
\expandafter\def\csname SRef@wb:cusp-cone-S24-map\endcsname{\href{S6_SERIES_GRIND_III.pdf\#nameddest=tag-0366}{III.88 (p.~260)} }
\expandafter\def\csname SRef@wb:cusp-cone-S24-source\endcsname{\href{S6_SERIES_GRIND_III.pdf\#nameddest=tag-036D}{III.89 (p.~260)} }
\expandafter\def\csname SRef@wb:cusp-cover-relative-factorization\endcsname{\href{S6_SERIES_GRIND_III.pdf\#nameddest=tag-0384}{III.9 (p.~239)} }
\expandafter\def\csname SRef@wb:cusp-cone-S24-fibre-map\endcsname{\href{S6_SERIES_GRIND_III.pdf\#nameddest=tag-0363}{III.90 (p.~260)} }
\expandafter\def\csname SRef@wb:cusp-cone-S24-R4-limit\endcsname{\href{S6_SERIES_GRIND_III.pdf\#nameddest=tag-0360}{III.91 (p.~260)} }
\expandafter\def\csname SRef@wb:cusp-cone-C-homology\endcsname{\href{S6_SERIES_GRIND_III.pdf\#nameddest=tag-035D}{III.92 (p.~260)} }
\expandafter\def\csname SRef@wb:cusp-cone-C-smash-homology\endcsname{\href{S6_SERIES_GRIND_III.pdf\#nameddest=tag-035E}{III.93 (p.~260)} }
\expandafter\def\csname SRef@wb:cusp-cone-Seifert-meridians\endcsname{\href{S6_SERIES_GRIND_III.pdf\#nameddest=tag-036F}{III.94 (p.~261)} }
\expandafter\def\csname SRef@wb:cusp-cone-M-H1\endcsname{\href{S6_SERIES_GRIND_III.pdf\#nameddest=tag-035M}{III.95 (p.~261)} }
\expandafter\def\csname SRef@wb:cusp-cone-M-boundary-map\endcsname{\href{S6_SERIES_GRIND_III.pdf\#nameddest=tag-035O}{III.96 (p.~261)} }
\expandafter\def\csname SRef@wb:cusp-cone-ambient-collapse\endcsname{\href{S6_SERIES_GRIND_III.pdf\#nameddest=tag-036I}{III.97 (p.~262)} }
\expandafter\def\csname SRef@wb:cusp-cone-intrinsic-expansion\endcsname{\href{S6_SERIES_GRIND_III.pdf\#nameddest=tag-036S}{III.98 (p.~262)} }
\expandafter\def\csname SRef@wb:cusp-cone-rank-bound\endcsname{\href{S6_SERIES_GRIND_III.pdf\#nameddest=tag-037H}{III.99 (p.~262)} }
\expandafter\def\csname SRef@maps:src:B.1.1\endcsname{\href{S6_SERIES_GRIND_III.pdf\#nameddest=tag-02J4}{III.3.2 (p.~13)} }
\expandafter\def\csname SRef@maps:src:C.1.2\endcsname{\href{S6_SERIES_GRIND_III.pdf\#nameddest=tag-02K3}{III.4.2 (p.~28)} }
\expandafter\def\csname SRef@maps:src:D.1.10\endcsname{\href{S6_SERIES_GRIND_III.pdf\#nameddest=tag-02K8}{III.5.10 (p.~40)} }
\expandafter\def\csname SRef@wb:h24:ret:chart-lemma\endcsname{\href{S6_SERIES_GRIND_III.pdf\#nameddest=tag-03AE}{III.7.30 (p.~300)} }
\expandafter\def\csname SRef@wb:h24:ret:orientation-lemma\endcsname{\href{S6_SERIES_GRIND_III.pdf\#nameddest=tag-03BD}{III.7.31 (p.~301)} }
\expandafter\def\csname SRef@mapsIII:carrier-three-complex-wedge\endcsname{\href{S6_SERIES_GRIND_III.pdf\#nameddest=tag-02QC}{III.7.34 (p.~306)} }
\expandafter\def\csname SRef@mapsIII:carrier-extension-lemma\endcsname{\href{S6_SERIES_GRIND_III.pdf\#nameddest=tag-02Q8}{III.7.36 (p.~308)} }
\expandafter\def\csname SRef@wb:cusp-cone-native-descent\endcsname{\href{S6_SERIES_GRIND_III.pdf\#nameddest=tag-0370}{III.7.5 (p.~255)} }
\expandafter\def\csname SRef@maps:src:D.2.4\endcsname{\href{S6_SERIES_GRIND_III.pdf\#nameddest=tag-02LG}{III.10.1 (p.~319)} }
\expandafter\def\csname SRef@maps:src:A.2.2\endcsname{\href{S6_SERIES_GRIND_III.pdf\#nameddest=tag-02J3}{III.3.1 (p.~12)} }
\expandafter\def\csname SRef@maps:src:C.1.4\endcsname{\href{S6_SERIES_GRIND_III.pdf\#nameddest=tag-02K5}{III.4.4 (p.~30)} }
\expandafter\def\csname SRef@maps:src:D.2.3\endcsname{\href{S6_SERIES_GRIND_III.pdf\#nameddest=tag-02L5}{III.5.13 (p.~43)} }
\expandafter\def\csname SRef@maps:src:D.2.19\endcsname{\href{S6_SERIES_GRIND_III.pdf\#nameddest=tag-02KT}{III.5.28 (p.~80)} }
\expandafter\def\csname SRef@maps:src:D.1.5\endcsname{\href{S6_SERIES_GRIND_III.pdf\#nameddest=tag-02KE}{III.5.5 (p.~36)} }
\expandafter\def\csname SRef@maps:src:E.1.1\endcsname{\href{S6_SERIES_GRIND_V.pdf\#nameddest=tag-02M4}{V.2.1 (p.~14)} }
\expandafter\def\csname SRef@maps:src:E.2.7\endcsname{\href{S6_SERIES_GRIND_V.pdf\#nameddest=tag-02MR}{V.2.11 (p.~29)} }
\expandafter\def\csname SRef@maps:src:E.2.8\endcsname{\href{S6_SERIES_GRIND_V.pdf\#nameddest=tag-02MS}{V.2.12 (p.~30)} }
\expandafter\def\csname SRef@maps:src:E.2.10\endcsname{\href{S6_SERIES_GRIND_V.pdf\#nameddest=tag-02MA}{V.2.14 (p.~33)} }
\expandafter\def\csname SRef@maps:src:E.2.18\endcsname{\href{S6_SERIES_GRIND_V.pdf\#nameddest=tag-02MI}{V.2.22 (p.~56)} }
\expandafter\def\csname SRef@maps:src:E.2.19\endcsname{\href{S6_SERIES_GRIND_V.pdf\#nameddest=tag-02MJ}{V.2.23 (p.~60)} }
\expandafter\def\csname SRef@maps:src:E.2.20\endcsname{\href{S6_SERIES_GRIND_V.pdf\#nameddest=tag-02ML}{V.2.24 (p.~63)} }
\expandafter\def\csname SRef@maps:src:E.2.4\endcsname{\href{S6_SERIES_GRIND_V.pdf\#nameddest=tag-02MO}{V.2.8 (p.~24)} }
\expandafter\def\csname SRef@maps:src:E.2.5\endcsname{\href{S6_SERIES_GRIND_V.pdf\#nameddest=tag-02MP}{V.2.9 (p.~25)} }
\expandafter\def\csname SRef@wb:cusp-cone-metric-scope\endcsname{\href{S6_SERIES_GRIND_III.pdf\#nameddest=tag-036W}{III.7.10 (p.~262)} }
\expandafter\def\csname SRef@wb:cusp-cone-explicit-carriers\endcsname{\href{S6_SERIES_GRIND_III.pdf\#nameddest=tag-036O}{III.7.17 (p.~273)} }
\expandafter\def\csname SRef@mapsIII:carrier-real-periods\endcsname{\href{S6_SERIES_GRIND_III.pdf\#nameddest=tag-02QB}{III.7.2 (p.~244)} }
\expandafter\def\csname SRef@wb:sphere-metric-tangent-scope\endcsname{\href{S6_SERIES_GRIND_III.pdf\#nameddest=tag-03DF}{III.7.29 (p.~294)} }
\expandafter\def\csname SRef@wb:cusp-spec-native-quotient\endcsname{\href{S6_SERIES_GRIND_III.pdf\#nameddest=tag-039Y}{III.7.3 (p.~245)} }
\expandafter\def\csname SRef@wb:h24:ret:framing-proposition\endcsname{\href{S6_SERIES_GRIND_III.pdf\#nameddest=tag-03BC}{III.7.32 (p.~302)} }
\expandafter\def\csname SRef@mapsaudit:finite-degree-three\endcsname{\href{S6_SERIES_GRIND_III.pdf\#nameddest=tag-02QE}{III.7.35 (p.~307)} }
\expandafter\def\csname SRef@mapsIII:angle-forgotten-one-copy-null\endcsname{\href{S6_SERIES_GRIND_III.pdf\#nameddest=tag-02Q7}{III.7.37 (p.~308)} }
\expandafter\def\csname SRef@wb:cusp-cone-S24-local-route\endcsname{\href{S6_SERIES_GRIND_III.pdf\#nameddest=tag-0365}{III.7.7 (p.~258)} }
\expandafter\def\csname SRef@maps:src:D.3.4\endcsname{\href{S6_SERIES_GRIND_III.pdf\#nameddest=tag-02M3}{III.9.2 (p.~317)} }
\expandafter\def\csname SRef@qa:III:layout:12\endcsname{\href{S6_SERIES_GRIND_III.pdf\#nameddest=tag-033S}{III.3 (p.~13)} }
\expandafter\def\csname SRef@wb:cusp-cover-naive-action-failure\endcsname{\href{S6_SERIES_GRIND_III.pdf\#nameddest=tag-0382}{III.7.1 (p.~243)} }
\expandafter\def\csname SRef@wb:cusp-cover-global-scope\endcsname{\href{S6_SERIES_GRIND_III.pdf\#nameddest=tag-037Y}{III.7.1 (p.~243)} }
\expandafter\def\csname SRef@qa:III:layout:11\endcsname{\href{S6_SERIES_GRIND_III.pdf\#nameddest=tag-033R}{III.{{$\langle 1\rangle $1.}} (p.~271)} }
\expandafter\def\csname SRef@wb:cusp-cone-moving-S4-null-homotopy\endcsname{\href{S6_SERIES_GRIND_III.pdf\#nameddest=tag-036Z}{III.7.3.1 (p.~272)} }
\expandafter\def\csname SRef@wb:cusp-cone-sphere-map-homotopy-status\endcsname{\href{S6_SERIES_GRIND_III.pdf\#nameddest=tag-037K}{III.7.3.1 (p.~273)} }
\expandafter\def\csname SRef@wb:cusp-s12-one-copy-invariant\endcsname{\href{S6_SERIES_GRIND_III.pdf\#nameddest=tag-038X}{III.7.4 (p.~281)} }
\expandafter\def\csname SRef@wb:cusp-s12-nullity-scope\endcsname{\href{S6_SERIES_GRIND_III.pdf\#nameddest=tag-038W}{III.7.4 (p.~281)} }
\expandafter\def\csname SRef@qa:III:layout:1\endcsname{\href{S6_SERIES_GRIND_III.pdf\#nameddest=tag-033Q}{III.{{$\langle 1\rangle $5.}} (p.~283)} }
\expandafter\def\csname SRef@qa:III:layout:2\endcsname{\href{S6_SERIES_GRIND_III.pdf\#nameddest=tag-033T}{III.{{$\langle 1\rangle $6.}} (p.~316)} }
\expandafter\def\csname SRef@qa:III:layout:3\endcsname{\href{S6_SERIES_GRIND_III.pdf\#nameddest=tag-033U}{III.{{$\langle 1\rangle $6.}} (p.~316)} }
\expandafter\def\csname SRef@qa:III:layout:4\endcsname{\href{S6_SERIES_GRIND_III.pdf\#nameddest=tag-033V}{III.8 (p.~316)} }
\expandafter\def\csname SRef@sec:III:conventions\endcsname{\href{S6_SERIES_GRIND_III.pdf\#nameddest=tag-034J}{III.2 (p.~12)} }
\expandafter\def\csname SRef@sec:III:maps\endcsname{\href{S6_SERIES_GRIND_III.pdf\#nameddest=tag-034K}{III.3 (p.~12)} }
\expandafter\def\csname SRef@sec:III:os49-night\endcsname{\href{S6_SERIES_GRIND_III.pdf\#nameddest=tag-034L}{III.6 (p.~215)} }
\expandafter\def\csname SRef@sec:III:carriers\endcsname{\href{S6_SERIES_GRIND_III.pdf\#nameddest=tag-034I}{III.7 (p.~238)} }
\expandafter\def\csname SRef@mapsD:repair:modular-collisions\endcsname{\href{S6_SERIES_GRIND_III.pdf\#nameddest=tag-02OJ}{III.8 (p.~310)} }
\expandafter\def\csname SRef@wb:cusp-cover-section\endcsname{\href{S6_SERIES_GRIND_III.pdf\#nameddest=tag-0385}{III.7.1 (p.~238)} }
\expandafter\def\csname SRef@mapsIII:carrier-radial-prerequisites\endcsname{\href{S6_SERIES_GRIND_III.pdf\#nameddest=tag-02QA}{III.7.2 (p.~244)} }
\expandafter\def\csname SRef@wb:cusp-cone-section\endcsname{\href{S6_SERIES_GRIND_III.pdf\#nameddest=tag-037I}{III.7.3 (p.~255)} }
\expandafter\def\csname SRef@wb:cusp-s12-carrier-homotopy\endcsname{\href{S6_SERIES_GRIND_III.pdf\#nameddest=tag-038U}{III.7.4 (p.~279)} }
\expandafter\def\csname SRef@wb:sphere-metric-limit-note\endcsname{\href{S6_SERIES_GRIND_III.pdf\#nameddest=tag-03D9}{III.7.6 (p.~287)} }
\expandafter\def\csname SRef@wb:h24:ret:section\endcsname{\href{S6_SERIES_GRIND_III.pdf\#nameddest=tag-03BF}{III.7.7 (p.~298)} }
\expandafter\def\csname SRef@mapsIII:carrier-extension-section\endcsname{\href{S6_SERIES_GRIND_III.pdf\#nameddest=tag-02Q9}{III.7.8 (p.~308)} }
\expandafter\def\csname SRef@the-source-section-and-component-divisor\endcsname{\href{S6_SERIES_GRIND_III.pdf\#nameddest=tag-0351}{III.6.0.1 (p.~215)} }
\expandafter\def\csname SRef@the-page-including-the-finite-cap-line\endcsname{\href{S6_SERIES_GRIND_III.pdf\#nameddest=tag-034Z}{III.6.0.10 (p.~231)} }
\expandafter\def\csname SRef@the-actual-graph-line-and-return\endcsname{\href{S6_SERIES_GRIND_III.pdf\#nameddest=tag-034Q}{III.6.0.11 (p.~232)} }
\expandafter\def\csname SRef@all-stabilizers-on-the-actual-root-chart\endcsname{\href{S6_SERIES_GRIND_III.pdf\#nameddest=tag-0008}{III.6.0.12 (p.~233)} }
\expandafter\def\csname SRef@the-full-fundamental-group-presentation\endcsname{\href{S6_SERIES_GRIND_III.pdf\#nameddest=tag-034U}{III.6.0.13 (p.~234)} }
\expandafter\def\csname SRef@the-fixed-exterior-and-the-actual-boundary-circle\endcsname{\href{S6_SERIES_GRIND_III.pdf\#nameddest=tag-034T}{III.6.0.14 (p.~235)} }
\expandafter\def\csname SRef@an-integral-basis-of-the-boundary-two-cohomology\endcsname{\href{S6_SERIES_GRIND_III.pdf\#nameddest=tag-000A}{III.6.0.15 (p.~236)} }
\expandafter\def\csname SRef@determine-the-exterior-restriction-without-guessing-a-splitting\endcsname{\href{S6_SERIES_GRIND_III.pdf\#nameddest=tag-01O2}{III.6.0.16 (p.~236)} }
\expandafter\def\csname SRef@the-two-smith-invariants-and-every-homology-group\endcsname{\href{S6_SERIES_GRIND_III.pdf\#nameddest=tag-0353}{III.6.0.17 (p.~237)} }
\expandafter\def\csname SRef@actual-local-fillings-and-their-gluing\endcsname{\href{S6_SERIES_GRIND_III.pdf\#nameddest=tag-0007}{III.6.0.2 (p.~216)} }
\expandafter\def\csname SRef@marked-monodromy-and-its-completeness\endcsname{\href{S6_SERIES_GRIND_III.pdf\#nameddest=tag-02QG}{III.6.0.3 (p.~217)} }
\expandafter\def\csname SRef@integral-topology\endcsname{\href{S6_SERIES_GRIND_III.pdf\#nameddest=tag-01RP}{III.6.0.4 (p.~218)} }
\expandafter\def\csname SRef@the-exact-line-subfan\endcsname{\href{S6_SERIES_GRIND_III.pdf\#nameddest=tag-034S}{III.6.0.5 (p.~225)} }
\expandafter\def\csname SRef@a-page-carrying-the-actual-source-line\endcsname{\href{S6_SERIES_GRIND_III.pdf\#nameddest=tag-0003}{III.6.0.6 (p.~227)} }
\expandafter\def\csname SRef@the-full-graph-to-page-comparison-by-regular-bands\endcsname{\href{S6_SERIES_GRIND_III.pdf\#nameddest=tag-034V}{III.6.0.7 (p.~228)} }
\expandafter\def\csname SRef@characteristic-calculation-and-the-scope-of-the-result\endcsname{\href{S6_SERIES_GRIND_III.pdf\#nameddest=tag-01GY}{III.6.0.8 (p.~229)} }
\expandafter\def\csname SRef@the-actual-closing-map-and-its-origins\endcsname{\href{S6_SERIES_GRIND_III.pdf\#nameddest=tag-034P}{III.6.0.9 (p.~231)} }
\expandafter\def\csname SRef@wb:cusp-cone-optimal-strengthening\endcsname{\href{S6_SERIES_GRIND_III.pdf\#nameddest=tag-0376}{III.7.3.1 (p.~264)} }
\expandafter\def\csname SRef@maps:src:B.5.2\endcsname{\href{S6_SERIES_GRIND_III.pdf\#nameddest=tag-02JC}{III.3.10 (p.~23)} }
\expandafter\def\csname SRef@maps:src:B.1.2\endcsname{\href{S6_SERIES_GRIND_III.pdf\#nameddest=tag-02J5}{III.3.3 (p.~14)} }
\expandafter\def\csname SRef@maps:src:B.2.1\endcsname{\href{S6_SERIES_GRIND_III.pdf\#nameddest=tag-02J6}{III.3.4 (p.~15)} }
\expandafter\def\csname SRef@maps:src:B.3.1\endcsname{\href{S6_SERIES_GRIND_III.pdf\#nameddest=tag-02J7}{III.3.5 (p.~16)} }
\expandafter\def\csname SRef@maps:src:B.3.2\endcsname{\href{S6_SERIES_GRIND_III.pdf\#nameddest=tag-02J8}{III.3.6 (p.~17)} }
\expandafter\def\csname SRef@maps:src:B.3.3\endcsname{\href{S6_SERIES_GRIND_III.pdf\#nameddest=tag-02J9}{III.3.7 (p.~18)} }
\expandafter\def\csname SRef@maps:src:B.4.1\endcsname{\href{S6_SERIES_GRIND_III.pdf\#nameddest=tag-02JA}{III.3.8 (p.~18)} }
\expandafter\def\csname SRef@maps:src:B.5.1\endcsname{\href{S6_SERIES_GRIND_III.pdf\#nameddest=tag-02JB}{III.3.9 (p.~20)} }
\expandafter\def\csname SRef@maps:src:C.1.1\endcsname{\href{S6_SERIES_GRIND_III.pdf\#nameddest=tag-02K2}{III.4.1 (p.~27)} }
\expandafter\def\csname SRef@maps:src:C.1.3\endcsname{\href{S6_SERIES_GRIND_III.pdf\#nameddest=tag-02K4}{III.4.3 (p.~29)} }
\expandafter\def\csname SRef@maps:src:C.1.5\endcsname{\href{S6_SERIES_GRIND_III.pdf\#nameddest=tag-02K6}{III.4.5 (p.~30)} }
\expandafter\def\csname SRef@maps:src:D.1.1\endcsname{\href{S6_SERIES_GRIND_III.pdf\#nameddest=tag-02K7}{III.5.1 (p.~32)} }
\expandafter\def\csname SRef@maps:src:D.1.12\endcsname{\href{S6_SERIES_GRIND_III.pdf\#nameddest=tag-02KA}{III.5.11 (p.~41)} }
\expandafter\def\csname SRef@maps:src:D.2.5\endcsname{\href{S6_SERIES_GRIND_III.pdf\#nameddest=tag-02LR}{III.5.14 (p.~44)} }
\expandafter\def\csname SRef@maps:src:D.2.6\endcsname{\href{S6_SERIES_GRIND_III.pdf\#nameddest=tag-02LW}{III.5.15 (p.~46)} }
\expandafter\def\csname SRef@maps:src:D.2.7\endcsname{\href{S6_SERIES_GRIND_III.pdf\#nameddest=tag-02LX}{III.5.16 (p.~48)} }
\expandafter\def\csname SRef@maps:src:D.2.8\endcsname{\href{S6_SERIES_GRIND_III.pdf\#nameddest=tag-02LY}{III.5.17 (p.~50)} }
\expandafter\def\csname SRef@maps:src:D.2.9\endcsname{\href{S6_SERIES_GRIND_III.pdf\#nameddest=tag-02LZ}{III.5.18 (p.~53)} }
\expandafter\def\csname SRef@maps:src:D.2.10\endcsname{\href{S6_SERIES_GRIND_III.pdf\#nameddest=tag-02KK}{III.5.19 (p.~55)} }
\expandafter\def\csname SRef@maps:src:D.1.2\endcsname{\href{S6_SERIES_GRIND_III.pdf\#nameddest=tag-02KB}{III.5.2 (p.~32)} }
\expandafter\def\csname SRef@maps:src:D.2.11\endcsname{\href{S6_SERIES_GRIND_III.pdf\#nameddest=tag-02KL}{III.5.20 (p.~59)} }
\expandafter\def\csname SRef@maps:src:D.2.12\endcsname{\href{S6_SERIES_GRIND_III.pdf\#nameddest=tag-02KM}{III.5.21 (p.~62)} }
\expandafter\def\csname SRef@maps:src:D.2.13\endcsname{\href{S6_SERIES_GRIND_III.pdf\#nameddest=tag-02KN}{III.5.22 (p.~65)} }
\expandafter\def\csname SRef@maps:src:D.2.14\endcsname{\href{S6_SERIES_GRIND_III.pdf\#nameddest=tag-02KO}{III.5.23 (p.~68)} }
\expandafter\def\csname SRef@maps:src:D.2.15\endcsname{\href{S6_SERIES_GRIND_III.pdf\#nameddest=tag-02KP}{III.5.24 (p.~70)} }
\expandafter\def\csname SRef@maps:src:D.2.16\endcsname{\href{S6_SERIES_GRIND_III.pdf\#nameddest=tag-02KQ}{III.5.25 (p.~73)} }
\expandafter\def\csname SRef@maps:src:D.2.18\endcsname{\href{S6_SERIES_GRIND_III.pdf\#nameddest=tag-02KS}{III.5.27 (p.~78)} }
\expandafter\def\csname SRef@maps:src:D.2.20\endcsname{\href{S6_SERIES_GRIND_III.pdf\#nameddest=tag-02KV}{III.5.29 (p.~82)} }
\expandafter\def\csname SRef@maps:src:D.1.3\endcsname{\href{S6_SERIES_GRIND_III.pdf\#nameddest=tag-02KC}{III.5.3 (p.~33)} }
\expandafter\def\csname SRef@maps:src:D.2.21\endcsname{\href{S6_SERIES_GRIND_III.pdf\#nameddest=tag-02KW}{III.5.30 (p.~86)} }
\expandafter\def\csname SRef@maps:src:D.2.22\endcsname{\href{S6_SERIES_GRIND_III.pdf\#nameddest=tag-02KX}{III.5.31 (p.~87)} }
\expandafter\def\csname SRef@maps:src:D.2.23\endcsname{\href{S6_SERIES_GRIND_III.pdf\#nameddest=tag-02KY}{III.5.32 (p.~91)} }
\expandafter\def\csname SRef@maps:src:D.2.24\endcsname{\href{S6_SERIES_GRIND_III.pdf\#nameddest=tag-02KZ}{III.5.33 (p.~93)} }
\expandafter\def\csname SRef@maps:src:D.2.25\endcsname{\href{S6_SERIES_GRIND_III.pdf\#nameddest=tag-02L0}{III.5.34 (p.~94)} }
\expandafter\def\csname SRef@maps:src:D.2.26\endcsname{\href{S6_SERIES_GRIND_III.pdf\#nameddest=tag-02L1}{III.5.35 (p.~98)} }
\expandafter\def\csname SRef@maps:src:D.2.27\endcsname{\href{S6_SERIES_GRIND_III.pdf\#nameddest=tag-02L2}{III.5.36 (p.~102)} }
\expandafter\def\csname SRef@maps:src:D.2.29\endcsname{\href{S6_SERIES_GRIND_III.pdf\#nameddest=tag-02L4}{III.5.38 (p.~107)} }
\expandafter\def\csname SRef@maps:src:D.2.30\endcsname{\href{S6_SERIES_GRIND_III.pdf\#nameddest=tag-02L6}{III.5.39 (p.~110)} }
\expandafter\def\csname SRef@maps:src:D.1.4\endcsname{\href{S6_SERIES_GRIND_III.pdf\#nameddest=tag-02KD}{III.5.4 (p.~34)} }
\expandafter\def\csname SRef@maps:src:D.2.31\endcsname{\href{S6_SERIES_GRIND_III.pdf\#nameddest=tag-02L7}{III.5.40 (p.~113)} }
\expandafter\def\csname SRef@maps:src:D.2.32\endcsname{\href{S6_SERIES_GRIND_III.pdf\#nameddest=tag-02L8}{III.5.41 (p.~116)} }
\expandafter\def\csname SRef@maps:src:D.2.33\endcsname{\href{S6_SERIES_GRIND_III.pdf\#nameddest=tag-02L9}{III.5.42 (p.~119)} }
\expandafter\def\csname SRef@maps:src:D.2.34\endcsname{\href{S6_SERIES_GRIND_III.pdf\#nameddest=tag-02LA}{III.5.43 (p.~121)} }
\expandafter\def\csname SRef@maps:src:D.2.35\endcsname{\href{S6_SERIES_GRIND_III.pdf\#nameddest=tag-02LB}{III.5.44 (p.~126)} }
\expandafter\def\csname SRef@maps:src:D.2.36\endcsname{\href{S6_SERIES_GRIND_III.pdf\#nameddest=tag-02LC}{III.5.45 (p.~130)} }
\expandafter\def\csname SRef@maps:src:D.2.37\endcsname{\href{S6_SERIES_GRIND_III.pdf\#nameddest=tag-02LD}{III.5.46 (p.~134)} }
\expandafter\def\csname SRef@maps:src:D.2.38\endcsname{\href{S6_SERIES_GRIND_III.pdf\#nameddest=tag-02LE}{III.5.47 (p.~138)} }
\expandafter\def\csname SRef@maps:src:D.2.39\endcsname{\href{S6_SERIES_GRIND_III.pdf\#nameddest=tag-02LF}{III.5.48 (p.~141)} }
\expandafter\def\csname SRef@maps:src:D.2.40\endcsname{\href{S6_SERIES_GRIND_III.pdf\#nameddest=tag-02LH}{III.5.49 (p.~145)} }
\expandafter\def\csname SRef@maps:src:D.2.41\endcsname{\href{S6_SERIES_GRIND_III.pdf\#nameddest=tag-02LI}{III.5.50 (p.~148)} }
\expandafter\def\csname SRef@maps:src:D.2.42\endcsname{\href{S6_SERIES_GRIND_III.pdf\#nameddest=tag-02LJ}{III.5.51 (p.~151)} }
\expandafter\def\csname SRef@maps:src:D.2.43\endcsname{\href{S6_SERIES_GRIND_III.pdf\#nameddest=tag-02LK}{III.5.52 (p.~153)} }
\expandafter\def\csname SRef@maps:src:D.2.44\endcsname{\href{S6_SERIES_GRIND_III.pdf\#nameddest=tag-02LL}{III.5.53 (p.~156)} }
\expandafter\def\csname SRef@maps:src:D.2.45\endcsname{\href{S6_SERIES_GRIND_III.pdf\#nameddest=tag-02LM}{III.5.54 (p.~161)} }
\expandafter\def\csname SRef@maps:src:D.2.46\endcsname{\href{S6_SERIES_GRIND_III.pdf\#nameddest=tag-02LN}{III.5.55 (p.~167)} }
\expandafter\def\csname SRef@maps:src:D.2.47\endcsname{\href{S6_SERIES_GRIND_III.pdf\#nameddest=tag-02LO}{III.5.56 (p.~172)} }
\expandafter\def\csname SRef@maps:src:D.2.48\endcsname{\href{S6_SERIES_GRIND_III.pdf\#nameddest=tag-02LP}{III.5.57 (p.~175)} }
\expandafter\def\csname SRef@maps:src:D.2.49\endcsname{\href{S6_SERIES_GRIND_III.pdf\#nameddest=tag-02LQ}{III.5.58 (p.~182)} }
\expandafter\def\csname SRef@maps:src:D.2.50\endcsname{\href{S6_SERIES_GRIND_III.pdf\#nameddest=tag-02LS}{III.5.59 (p.~187)} }
\expandafter\def\csname SRef@maps:src:D.1.6\endcsname{\href{S6_SERIES_GRIND_III.pdf\#nameddest=tag-02KF}{III.5.6 (p.~36)} }
\expandafter\def\csname SRef@maps:src:D.2.51\endcsname{\href{S6_SERIES_GRIND_III.pdf\#nameddest=tag-02LT}{III.5.60 (p.~191)} }
\expandafter\def\csname SRef@maps:src:D.2.52\endcsname{\href{S6_SERIES_GRIND_III.pdf\#nameddest=tag-02LU}{III.5.61 (p.~194)} }
\expandafter\def\csname SRef@maps:src:D.2.53\endcsname{\href{S6_SERIES_GRIND_III.pdf\#nameddest=tag-02LV}{III.5.62 (p.~196)} }
\expandafter\def\csname SRef@maps:src:D.3.1\endcsname{\href{S6_SERIES_GRIND_III.pdf\#nameddest=tag-02M0}{III.5.63 (p.~211)} }
\expandafter\def\csname SRef@maps:src:D.3.3\endcsname{\href{S6_SERIES_GRIND_III.pdf\#nameddest=tag-02M2}{III.5.64 (p.~213)} }
\expandafter\def\csname SRef@maps:src:D.1.11\endcsname{\href{S6_SERIES_GRIND_III.pdf\#nameddest=tag-02K9}{III.5.65 (p.~214)} }
\expandafter\def\csname SRef@maps:src:D.1.7\endcsname{\href{S6_SERIES_GRIND_III.pdf\#nameddest=tag-02KG}{III.5.7 (p.~37)} }
\expandafter\def\csname SRef@maps:src:D.1.8\endcsname{\href{S6_SERIES_GRIND_III.pdf\#nameddest=tag-02KH}{III.5.8 (p.~38)} }
\expandafter\def\csname SRef@maps:src:D.1.9\endcsname{\href{S6_SERIES_GRIND_III.pdf\#nameddest=tag-02KI}{III.5.9 (p.~39)} }
\expandafter\def\csname SRef@night:P3F.26\endcsname{\href{S6_SERIES_GRIND_III.pdf\#nameddest=tag-02S2}{III.6.1 (p.~215)} }
\expandafter\def\csname SRef@stmt:P3F.26\endcsname{\href{S6_SERIES_GRIND_III.pdf\#nameddest=tag-02S2}{III.6.1 (p.~215)} }
\expandafter\def\csname SRef@night:P3F.27\endcsname{\href{S6_SERIES_GRIND_III.pdf\#nameddest=tag-02S3}{III.6.2 (p.~220)} }
\expandafter\def\csname SRef@stmt:P3F.27\endcsname{\href{S6_SERIES_GRIND_III.pdf\#nameddest=tag-02S3}{III.6.2 (p.~220)} }
\expandafter\def\csname SRef@night:P3F.28\endcsname{\href{S6_SERIES_GRIND_III.pdf\#nameddest=tag-02S4}{III.6.3 (p.~221)} }
\expandafter\def\csname SRef@stmt:P3F.28\endcsname{\href{S6_SERIES_GRIND_III.pdf\#nameddest=tag-02S4}{III.6.3 (p.~221)} }
\expandafter\def\csname SRef@night:P3F.29\endcsname{\href{S6_SERIES_GRIND_III.pdf\#nameddest=tag-02S5}{III.6.4 (p.~223)} }
\expandafter\def\csname SRef@stmt:P3F.29\endcsname{\href{S6_SERIES_GRIND_III.pdf\#nameddest=tag-02S5}{III.6.4 (p.~223)} }
\expandafter\def\csname SRef@night:P3F.31\endcsname{\href{S6_SERIES_GRIND_III.pdf\#nameddest=tag-02S8}{III.6.5 (p.~225)} }
\expandafter\def\csname SRef@stmt:P3F.31\endcsname{\href{S6_SERIES_GRIND_III.pdf\#nameddest=tag-02S8}{III.6.5 (p.~225)} }
\expandafter\def\csname SRef@night:P3F.32\endcsname{\href{S6_SERIES_GRIND_III.pdf\#nameddest=tag-02S9}{III.6.6 (p.~230)} }
\expandafter\def\csname SRef@stmt:P3F.32\endcsname{\href{S6_SERIES_GRIND_III.pdf\#nameddest=tag-02S9}{III.6.6 (p.~230)} }
\expandafter\def\csname SRef@night:P3F.33\endcsname{\href{S6_SERIES_GRIND_III.pdf\#nameddest=tag-02SA}{III.6.7 (p.~233)} }
\expandafter\def\csname SRef@stmt:P3F.33\endcsname{\href{S6_SERIES_GRIND_III.pdf\#nameddest=tag-02SA}{III.6.7 (p.~233)} }
\expandafter\def\csname SRef@night:P3F.34\endcsname{\href{S6_SERIES_GRIND_III.pdf\#nameddest=tag-02SB}{III.6.8 (p.~234)} }
\expandafter\def\csname SRef@stmt:P3F.34\endcsname{\href{S6_SERIES_GRIND_III.pdf\#nameddest=tag-02SB}{III.6.8 (p.~234)} }
\expandafter\def\csname SRef@maps:src:E.2.6\endcsname{\href{S6_SERIES_GRIND_V.pdf\#nameddest=tag-02MQ}{V.2.10 (p.~26)} }
\expandafter\def\csname SRef@maps:src:E.2.9\endcsname{\href{S6_SERIES_GRIND_V.pdf\#nameddest=tag-02MT}{V.2.13 (p.~31)} }
\expandafter\def\csname SRef@maps:src:E.2.11\endcsname{\href{S6_SERIES_GRIND_V.pdf\#nameddest=tag-02MB}{V.2.15 (p.~35)} }
\expandafter\def\csname SRef@maps:src:E.2.12\endcsname{\href{S6_SERIES_GRIND_V.pdf\#nameddest=tag-02MC}{V.2.16 (p.~37)} }
\expandafter\def\csname SRef@maps:src:E.2.13\endcsname{\href{S6_SERIES_GRIND_V.pdf\#nameddest=tag-02MD}{V.2.17 (p.~40)} }
\expandafter\def\csname SRef@maps:src:E.2.14\endcsname{\href{S6_SERIES_GRIND_V.pdf\#nameddest=tag-02ME}{V.2.18 (p.~42)} }
\expandafter\def\csname SRef@maps:src:E.2.15\endcsname{\href{S6_SERIES_GRIND_V.pdf\#nameddest=tag-02MF}{V.2.19 (p.~45)} }
\expandafter\def\csname SRef@maps:src:E.1.2\endcsname{\href{S6_SERIES_GRIND_V.pdf\#nameddest=tag-02M5}{V.2.2 (p.~14)} }
\expandafter\def\csname SRef@maps:src:E.2.16\endcsname{\href{S6_SERIES_GRIND_V.pdf\#nameddest=tag-02MG}{V.2.20 (p.~48)} }
\expandafter\def\csname SRef@maps:src:E.2.17\endcsname{\href{S6_SERIES_GRIND_V.pdf\#nameddest=tag-02MH}{V.2.21 (p.~52)} }
\expandafter\def\csname SRef@maps:src:E.2.21\endcsname{\href{S6_SERIES_GRIND_V.pdf\#nameddest=tag-02MM}{V.2.25 (p.~65)} }
\expandafter\def\csname SRef@maps:src:E.1.3\endcsname{\href{S6_SERIES_GRIND_V.pdf\#nameddest=tag-02M6}{V.2.3 (p.~16)} }
\expandafter\def\csname SRef@maps:src:E.1.4\endcsname{\href{S6_SERIES_GRIND_V.pdf\#nameddest=tag-02M7}{V.2.4 (p.~18)} }
\expandafter\def\csname SRef@maps:src:E.1.5\endcsname{\href{S6_SERIES_GRIND_V.pdf\#nameddest=tag-02M8}{V.2.5 (p.~19)} }
\expandafter\def\csname SRef@maps:src:E.2.1\endcsname{\href{S6_SERIES_GRIND_V.pdf\#nameddest=tag-02M9}{V.2.6 (p.~21)} }
\expandafter\def\csname SRef@maps:src:E.2.3\endcsname{\href{S6_SERIES_GRIND_V.pdf\#nameddest=tag-02MN}{V.2.7 (p.~23)} }
\expandafter\def\csname SRef@wb:cusp-cover-extension-theorem\endcsname{\href{S6_SERIES_GRIND_III.pdf\#nameddest=tag-037X}{III.7.1 (p.~238)} }
\expandafter\def\csname SRef@wb:cusp-cone-obstructions\endcsname{\href{S6_SERIES_GRIND_III.pdf\#nameddest=tag-0374}{III.7.11 (p.~262)} }
\expandafter\def\csname SRef@wb:cusp-cone-optimal-exponents\endcsname{\href{S6_SERIES_GRIND_III.pdf\#nameddest=tag-0375}{III.7.12 (p.~265)} }
\expandafter\def\csname SRef@wb:cusp-cone-angle-forgotten-sphere-carriers\endcsname{\href{S6_SERIES_GRIND_III.pdf\#nameddest=tag-036J}{III.7.13 (p.~266)} }
\expandafter\def\csname SRef@wb:cusp-cone-target-homotopy-types\endcsname{\href{S6_SERIES_GRIND_III.pdf\#nameddest=tag-037M}{III.7.18 (p.~274)} }
\expandafter\def\csname SRef@wb:cusp-cone-exact-decompactification-space\endcsname{\href{S6_SERIES_GRIND_III.pdf\#nameddest=tag-036N}{III.7.19 (p.~276)} }
\expandafter\def\csname SRef@wb:cusp-s12-full-carrier-null\endcsname{\href{S6_SERIES_GRIND_III.pdf\#nameddest=tag-038V}{III.7.21 (p.~280)} }
\expandafter\def\csname SRef@wb:cusp-null:thm\endcsname{\href{S6_SERIES_GRIND_III.pdf\#nameddest=tag-038S}{III.7.23 (p.~282)} }
\expandafter\def\csname SRef@wb:sphere-metric-uniform-collapse\endcsname{\href{S6_SERIES_GRIND_III.pdf\#nameddest=tag-03DG}{III.7.25 (p.~288)} }
\expandafter\def\csname SRef@wb:sphere-metric-fixed-realization\endcsname{\href{S6_SERIES_GRIND_III.pdf\#nameddest=tag-03D6}{III.7.26 (p.~288)} }
\expandafter\def\csname SRef@wb:sphere-metric-scaling-trichotomy\endcsname{\href{S6_SERIES_GRIND_III.pdf\#nameddest=tag-03DC}{III.7.27 (p.~291)} }
\expandafter\def\csname SRef@wb:cusp-spec-globalization-obstruction\endcsname{\href{S6_SERIES_GRIND_III.pdf\#nameddest=tag-039U}{III.7.4 (p.~246)} }
\expandafter\def\csname SRef@wb:cusp-cone-S12-route\endcsname{\href{S6_SERIES_GRIND_III.pdf\#nameddest=tag-035X}{III.7.6 (p.~256)} }
\expandafter\def\csname SRef@wb:cusp-cone-S24-route\endcsname{\href{S6_SERIES_GRIND_III.pdf\#nameddest=tag-036A}{III.7.8 (p.~259)} }
\expandafter\def\csname SRef@mapsD:repair:collision-realization\endcsname{\href{S6_SERIES_GRIND_III.pdf\#nameddest=tag-02OG}{III.8.1 (p.~310)} }
\expandafter\def\csname SRef@maps:src:D.3.2\endcsname{\href{S6_SERIES_GRIND_III.pdf\#nameddest=tag-02M1}{III.9.1 (p.~316)} }
\expandafter\def\csname SRef@maps:src:E.2.2\endcsname{\href{S6_SERIES_GRIND_III.pdf\#nameddest=tag-02MK}{III.9.3 (p.~317)} }
\expandafter\def\csname SRef@maps:eq:B.6.2\endcsname{\href{S6_SERIES_GRIND_IV.pdf\#nameddest=tag-01VZ}{IV.{B.6.2} (p.~68)} }
\expandafter\def\csname SRef@maps:eq:B.6.3\endcsname{\href{S6_SERIES_GRIND_IV.pdf\#nameddest=tag-01WA}{IV.{B.6.3} (p.~68)} }
\expandafter\def\csname SRef@maps:eq:B.6.4\endcsname{\href{S6_SERIES_GRIND_IV.pdf\#nameddest=tag-01WL}{IV.{B.6.4} (p.~69)} }
\expandafter\def\csname SRef@maps:eq:B.6.5\endcsname{\href{S6_SERIES_GRIND_IV.pdf\#nameddest=tag-01WR}{IV.{B.6.5} (p.~69)} }
\expandafter\def\csname SRef@maps:eq:B.6.6\endcsname{\href{S6_SERIES_GRIND_IV.pdf\#nameddest=tag-01WS}{IV.{B.6.6} (p.~69)} }
\expandafter\def\csname SRef@maps:eq:B.6.7\endcsname{\href{S6_SERIES_GRIND_IV.pdf\#nameddest=tag-01WT}{IV.{B.6.7} (p.~70)} }
\expandafter\def\csname SRef@maps:eq:B.6.8\endcsname{\href{S6_SERIES_GRIND_IV.pdf\#nameddest=tag-01WU}{IV.{B.6.8} (p.~70)} }
\expandafter\def\csname SRef@maps:eq:B.6.9\endcsname{\href{S6_SERIES_GRIND_IV.pdf\#nameddest=tag-01WV}{IV.{B.6.9} (p.~70)} }
\expandafter\def\csname SRef@maps:eq:B.6.10\endcsname{\href{S6_SERIES_GRIND_IV.pdf\#nameddest=tag-01VQ}{IV.{B.6.10} (p.~70)} }
\expandafter\def\csname SRef@maps:eq:B.6.11\endcsname{\href{S6_SERIES_GRIND_IV.pdf\#nameddest=tag-01VR}{IV.{B.6.11} (p.~72)} }
\expandafter\def\csname SRef@maps:eq:B.6.12\endcsname{\href{S6_SERIES_GRIND_IV.pdf\#nameddest=tag-01VS}{IV.{B.6.12} (p.~72)} }
\expandafter\def\csname SRef@maps:eq:B.6.13\endcsname{\href{S6_SERIES_GRIND_IV.pdf\#nameddest=tag-01VT}{IV.{B.6.13} (p.~73)} }
\expandafter\def\csname SRef@maps:eq:B.6.14\endcsname{\href{S6_SERIES_GRIND_IV.pdf\#nameddest=tag-01VU}{IV.{B.6.14} (p.~73)} }
\expandafter\def\csname SRef@maps:eq:B.6.15\endcsname{\href{S6_SERIES_GRIND_IV.pdf\#nameddest=tag-01VV}{IV.{B.6.15} (p.~75)} }
\expandafter\def\csname SRef@maps:eq:B.6.16\endcsname{\href{S6_SERIES_GRIND_IV.pdf\#nameddest=tag-01VW}{IV.{B.6.16} (p.~75)} }
\expandafter\def\csname SRef@maps:eq:B.6.17\endcsname{\href{S6_SERIES_GRIND_IV.pdf\#nameddest=tag-01VX}{IV.{B.6.17} (p.~76)} }
\expandafter\def\csname SRef@maps:eq:B.C1\endcsname{\href{S6_SERIES_GRIND_IV.pdf\#nameddest=tag-01WW}{IV.{B.C1} (p.~77)} }
\expandafter\def\csname SRef@maps:eq:B.C2\endcsname{\href{S6_SERIES_GRIND_IV.pdf\#nameddest=tag-01WX}{IV.{B.C2} (p.~77)} }
\expandafter\def\csname SRef@maps:eq:B.C3\endcsname{\href{S6_SERIES_GRIND_IV.pdf\#nameddest=tag-01WY}{IV.{B.C3} (p.~77)} }
\expandafter\def\csname SRef@maps:eq:B.C4\endcsname{\href{S6_SERIES_GRIND_IV.pdf\#nameddest=tag-01WZ}{IV.{B.C4} (p.~78)} }
\expandafter\def\csname SRef@maps:eq:B.P1\endcsname{\href{S6_SERIES_GRIND_IV.pdf\#nameddest=tag-01X0}{IV.{B.P1} (p.~78)} }
\expandafter\def\csname SRef@maps:eq:B.P2\endcsname{\href{S6_SERIES_GRIND_IV.pdf\#nameddest=tag-01X1}{IV.{B.P2} (p.~79)} }
\expandafter\def\csname SRef@maps:eq:B.P3\endcsname{\href{S6_SERIES_GRIND_IV.pdf\#nameddest=tag-01X2}{IV.{B.P3} (p.~79)} }
\expandafter\def\csname SRef@maps:eq:B.P4\endcsname{\href{S6_SERIES_GRIND_IV.pdf\#nameddest=tag-01X3}{IV.{B.P4} (p.~79)} }
\expandafter\def\csname SRef@maps:eq:B.6.19\endcsname{\href{S6_SERIES_GRIND_IV.pdf\#nameddest=tag-01VY}{IV.{B.6.19} (p.~80)} }
\expandafter\def\csname SRef@maps:eq:B.6.20\endcsname{\href{S6_SERIES_GRIND_IV.pdf\#nameddest=tag-01W0}{IV.{B.6.20} (p.~81)} }
\expandafter\def\csname SRef@maps:eq:B.6.21\endcsname{\href{S6_SERIES_GRIND_IV.pdf\#nameddest=tag-01W1}{IV.{B.6.21} (p.~83)} }
\expandafter\def\csname SRef@maps:eq:B.6.22\endcsname{\href{S6_SERIES_GRIND_IV.pdf\#nameddest=tag-01W2}{IV.{B.6.22} (p.~83)} }
\expandafter\def\csname SRef@maps:eq:B.6.23\endcsname{\href{S6_SERIES_GRIND_IV.pdf\#nameddest=tag-01W3}{IV.{B.6.23} (p.~83)} }
\expandafter\def\csname SRef@maps:eq:B.6.24\endcsname{\href{S6_SERIES_GRIND_IV.pdf\#nameddest=tag-01W4}{IV.{B.6.24} (p.~84)} }
\expandafter\def\csname SRef@maps:eq:B.6.25\endcsname{\href{S6_SERIES_GRIND_IV.pdf\#nameddest=tag-01W5}{IV.{B.6.25} (p.~84)} }
\expandafter\def\csname SRef@maps:eq:B.T1\endcsname{\href{S6_SERIES_GRIND_IV.pdf\#nameddest=tag-01XD}{IV.{B.T1} (p.~87)} }
\expandafter\def\csname SRef@maps:eq:B.6.26\endcsname{\href{S6_SERIES_GRIND_IV.pdf\#nameddest=tag-01W6}{IV.{B.6.26} (p.~87)} }
\expandafter\def\csname SRef@maps:eq:B.6.27\endcsname{\href{S6_SERIES_GRIND_IV.pdf\#nameddest=tag-01W7}{IV.{B.6.27} (p.~88)} }
\expandafter\def\csname SRef@maps:eq:B.6.28\endcsname{\href{S6_SERIES_GRIND_IV.pdf\#nameddest=tag-01W8}{IV.{B.6.28} (p.~89)} }
\expandafter\def\csname SRef@maps:eq:B.6.29\endcsname{\href{S6_SERIES_GRIND_IV.pdf\#nameddest=tag-01W9}{IV.{B.6.29} (p.~89)} }
\expandafter\def\csname SRef@maps:eq:B.6.30\endcsname{\href{S6_SERIES_GRIND_IV.pdf\#nameddest=tag-01WB}{IV.{B.6.30} (p.~89)} }
\expandafter\def\csname SRef@maps:eq:B.6.31\endcsname{\href{S6_SERIES_GRIND_IV.pdf\#nameddest=tag-01WC}{IV.{B.6.31} (p.~90)} }
\expandafter\def\csname SRef@maps:eq:B.Q1\endcsname{\href{S6_SERIES_GRIND_IV.pdf\#nameddest=tag-01X4}{IV.{B.Q1} (p.~91)} }
\expandafter\def\csname SRef@maps:eq:B.Q2\endcsname{\href{S6_SERIES_GRIND_IV.pdf\#nameddest=tag-01X5}{IV.{B.Q2} (p.~91)} }
\expandafter\def\csname SRef@maps:eq:B.Q3\endcsname{\href{S6_SERIES_GRIND_IV.pdf\#nameddest=tag-01X6}{IV.{B.Q3} (p.~91)} }
\expandafter\def\csname SRef@maps:eq:B.Q4\endcsname{\href{S6_SERIES_GRIND_IV.pdf\#nameddest=tag-01X7}{IV.{B.Q4} (p.~92)} }
\expandafter\def\csname SRef@maps:eq:B.Q5\endcsname{\href{S6_SERIES_GRIND_IV.pdf\#nameddest=tag-01X8}{IV.{B.Q5} (p.~92)} }
\expandafter\def\csname SRef@maps:eq:B.Q6\endcsname{\href{S6_SERIES_GRIND_IV.pdf\#nameddest=tag-01X9}{IV.{B.Q6} (p.~92)} }
\expandafter\def\csname SRef@maps:eq:B.Q7\endcsname{\href{S6_SERIES_GRIND_IV.pdf\#nameddest=tag-01XA}{IV.{B.Q7} (p.~92)} }
\expandafter\def\csname SRef@maps:eq:B.Q8\endcsname{\href{S6_SERIES_GRIND_IV.pdf\#nameddest=tag-01XB}{IV.{B.Q8} (p.~92)} }
\expandafter\def\csname SRef@maps:eq:B.Q9\endcsname{\href{S6_SERIES_GRIND_IV.pdf\#nameddest=tag-01XC}{IV.{B.Q9} (p.~93)} }
\expandafter\def\csname SRef@maps:eq:B.6.32\endcsname{\href{S6_SERIES_GRIND_IV.pdf\#nameddest=tag-01WD}{IV.{B.6.32} (p.~93)} }
\expandafter\def\csname SRef@maps:eq:B.6.33\endcsname{\href{S6_SERIES_GRIND_IV.pdf\#nameddest=tag-01WE}{IV.{B.6.33} (p.~94)} }
\expandafter\def\csname SRef@maps:eq:B.6.34\endcsname{\href{S6_SERIES_GRIND_IV.pdf\#nameddest=tag-01WF}{IV.{B.6.34} (p.~94)} }
\expandafter\def\csname SRef@famIII:eq:E9\endcsname{\href{S6_SERIES_GRIND_IV.pdf\#nameddest=tag-01RJ}{IV.{E9} (p.~20)} }
\expandafter\def\csname SRef@maps:eq:B.6.35\endcsname{\href{S6_SERIES_GRIND_IV.pdf\#nameddest=tag-01WG}{IV.{B.6.35} (p.~95)} }
\expandafter\def\csname SRef@maps:eq:B.6.36\endcsname{\href{S6_SERIES_GRIND_IV.pdf\#nameddest=tag-01WH}{IV.{B.6.36} (p.~95)} }
\expandafter\def\csname SRef@maps:eq:B.6.37\endcsname{\href{S6_SERIES_GRIND_IV.pdf\#nameddest=tag-01WI}{IV.{B.6.37} (p.~95)} }
\expandafter\def\csname SRef@maps:eq:B.6.38\endcsname{\href{S6_SERIES_GRIND_IV.pdf\#nameddest=tag-01WJ}{IV.{B.6.38} (p.~99)} }
\expandafter\def\csname SRef@maps:eq:B.6.39\endcsname{\href{S6_SERIES_GRIND_IV.pdf\#nameddest=tag-01WK}{IV.{B.6.39} (p.~100)} }
\expandafter\def\csname SRef@famIII:eq:E10\endcsname{\href{S6_SERIES_GRIND_IV.pdf\#nameddest=tag-01RI}{IV.{E10} (p.~20)} }
\expandafter\def\csname SRef@maps:eq:B.6.40\endcsname{\href{S6_SERIES_GRIND_IV.pdf\#nameddest=tag-01WM}{IV.{B.6.40} (p.~100)} }
\expandafter\def\csname SRef@maps:eq:B.6.41\endcsname{\href{S6_SERIES_GRIND_IV.pdf\#nameddest=tag-01WN}{IV.{B.6.41} (p.~100)} }
\expandafter\def\csname SRef@maps:eq:B.6.42\endcsname{\href{S6_SERIES_GRIND_IV.pdf\#nameddest=tag-01WO}{IV.{B.6.42} (p.~100)} }
\expandafter\def\csname SRef@maps:eq:B.6.43\endcsname{\href{S6_SERIES_GRIND_IV.pdf\#nameddest=tag-01WP}{IV.{B.6.43} (p.~103)} }
\expandafter\def\csname SRef@maps:eq:B.6.44\endcsname{\href{S6_SERIES_GRIND_IV.pdf\#nameddest=tag-01WQ}{IV.{B.6.44} (p.~103)} }
\expandafter\def\csname SRef@maps:eq:B.6.1\endcsname{\href{S6_SERIES_GRIND_IV.pdf\#nameddest=tag-01VP}{IV.{B.6.1} (p.~68)} }
\expandafter\def\csname SRef@push:src:C.3.15\endcsname{\href{S6_SERIES_GRIND_IV.pdf\#nameddest=tag-031M}{IV.2.16 (p.~40)} }
\expandafter\def\csname SRef@push:src:C.3.27\endcsname{\href{S6_SERIES_GRIND_IV.pdf\#nameddest=tag-031Z}{IV.2.28 (p.~63)} }
\expandafter\def\csname SRef@push:src:C.5.1\endcsname{\href{S6_SERIES_GRIND_IV.pdf\#nameddest=tag-0329}{IV.2.31 (p.~65)} }
\expandafter\def\csname SRef@maps:src:B.6.8\endcsname{\href{S6_SERIES_GRIND_IV.pdf\#nameddest=tag-02K0}{IV.2.39 (p.~76)} }
\expandafter\def\csname SRef@maps:src:B.6.25\endcsname{\href{S6_SERIES_GRIND_IV.pdf\#nameddest=tag-02JU}{IV.2.55 (p.~104)} }
\expandafter\def\csname SRef@push:eq:C12-u-flow\endcsname{\href{S6_SERIES_GRIND_IV.pdf\#nameddest=tag-02YV}{IV.1 (p.~33)} }
\expandafter\def\csname SRef@push:eq:C12-limiting-identification\endcsname{\href{S6_SERIES_GRIND_IV.pdf\#nameddest=tag-02YU}{IV.2 (p.~34)} }
\expandafter\def\csname SRef@push:eq:C13-real-overlap\endcsname{\href{S6_SERIES_GRIND_IV.pdf\#nameddest=tag-02YX}{IV.3 (p.~35)} }
\expandafter\def\csname SRef@push:eq:C13-weighted-map\endcsname{\href{S6_SERIES_GRIND_IV.pdf\#nameddest=tag-02YY}{IV.4 (p.~35)} }
\expandafter\def\csname SRef@push:eq:C13-cusp-gauge\endcsname{\href{S6_SERIES_GRIND_IV.pdf\#nameddest=tag-02YW}{IV.5 (p.~37)} }
\expandafter\def\csname SRef@push:src:C.3.9\endcsname{\href{S6_SERIES_GRIND_IV.pdf\#nameddest=tag-0326}{IV.2.10 (p.~29)} }
\expandafter\def\csname SRef@push:src:C.3.22\endcsname{\href{S6_SERIES_GRIND_IV.pdf\#nameddest=tag-031U}{IV.2.23 (p.~54)} }
\expandafter\def\csname SRef@push:src:C.3.8\endcsname{\href{S6_SERIES_GRIND_IV.pdf\#nameddest=tag-0325}{IV.2.9 (p.~28)} }
\expandafter\def\csname SRef@iv:rank_two_specialization\endcsname{\href{S6_SERIES_GRIND_IV.pdf\#nameddest=tag-01UU}{IV.3.20 (p.~131)} }
\expandafter\def\csname SRef@fam:src:E.2.2\endcsname{\href{S6_SERIES_GRIND_IV.pdf\#nameddest=tag-01R7}{IV.2.1 (p.~20)} }
\expandafter\def\csname SRef@push:src:C.3.10\endcsname{\href{S6_SERIES_GRIND_IV.pdf\#nameddest=tag-031H}{IV.2.11 (p.~30)} }
\expandafter\def\csname SRef@push:src:C.3.11\endcsname{\href{S6_SERIES_GRIND_IV.pdf\#nameddest=tag-031I}{IV.2.12 (p.~32)} }
\expandafter\def\csname SRef@push:src:C.3.12\endcsname{\href{S6_SERIES_GRIND_IV.pdf\#nameddest=tag-031J}{IV.2.13 (p.~32)} }
\expandafter\def\csname SRef@push:src:C.3.13\endcsname{\href{S6_SERIES_GRIND_IV.pdf\#nameddest=tag-031K}{IV.2.14 (p.~34)} }
\expandafter\def\csname SRef@push:src:C.3.14\endcsname{\href{S6_SERIES_GRIND_IV.pdf\#nameddest=tag-031L}{IV.2.15 (p.~37)} }
\expandafter\def\csname SRef@push:src:C.3.16\endcsname{\href{S6_SERIES_GRIND_IV.pdf\#nameddest=tag-031N}{IV.2.17 (p.~40)} }
\expandafter\def\csname SRef@push:src:C.3.17\endcsname{\href{S6_SERIES_GRIND_IV.pdf\#nameddest=tag-031O}{IV.2.18 (p.~41)} }
\expandafter\def\csname SRef@push:src:C.3.18\endcsname{\href{S6_SERIES_GRIND_IV.pdf\#nameddest=tag-031P}{IV.2.19 (p.~44)} }
\expandafter\def\csname SRef@push:src:C.3.19\endcsname{\href{S6_SERIES_GRIND_IV.pdf\#nameddest=tag-031Q}{IV.2.20 (p.~47)} }
\expandafter\def\csname SRef@push:src:C.3.20\endcsname{\href{S6_SERIES_GRIND_IV.pdf\#nameddest=tag-031S}{IV.2.21 (p.~49)} }
\expandafter\def\csname SRef@push:src:C.3.21\endcsname{\href{S6_SERIES_GRIND_IV.pdf\#nameddest=tag-031T}{IV.2.22 (p.~51)} }
\expandafter\def\csname SRef@push:src:C.3.23\endcsname{\href{S6_SERIES_GRIND_IV.pdf\#nameddest=tag-031V}{IV.2.24 (p.~55)} }
\expandafter\def\csname SRef@push:src:C.2.2\endcsname{\href{S6_SERIES_GRIND_IV.pdf\#nameddest=tag-031E}{IV.2.3 (p.~22)} }
\expandafter\def\csname SRef@push:src:C.3.25\endcsname{\href{S6_SERIES_GRIND_IV.pdf\#nameddest=tag-031X}{IV.2.26 (p.~58)} }
\expandafter\def\csname SRef@push:src:C.3.26\endcsname{\href{S6_SERIES_GRIND_IV.pdf\#nameddest=tag-031Y}{IV.2.27 (p.~61)} }
\expandafter\def\csname SRef@push:src:C.4.1\endcsname{\href{S6_SERIES_GRIND_IV.pdf\#nameddest=tag-0327}{IV.2.29 (p.~64)} }
\expandafter\def\csname SRef@push:src:C.4.2\endcsname{\href{S6_SERIES_GRIND_IV.pdf\#nameddest=tag-0328}{IV.2.30 (p.~65)} }
\expandafter\def\csname SRef@push:src:C.2.3\endcsname{\href{S6_SERIES_GRIND_IV.pdf\#nameddest=tag-031F}{IV.2.4 (p.~23)} }
\expandafter\def\csname SRef@maps:src:B.6.1\endcsname{\href{S6_SERIES_GRIND_IV.pdf\#nameddest=tag-02JD}{IV.2.32 (p.~66)} }
\expandafter\def\csname SRef@maps:src:B.6.3\endcsname{\href{S6_SERIES_GRIND_IV.pdf\#nameddest=tag-02JV}{IV.2.34 (p.~67)} }
\expandafter\def\csname SRef@maps:src:B.6.4\endcsname{\href{S6_SERIES_GRIND_IV.pdf\#nameddest=tag-02JW}{IV.2.35 (p.~69)} }
\expandafter\def\csname SRef@maps:src:B.6.5\endcsname{\href{S6_SERIES_GRIND_IV.pdf\#nameddest=tag-02JX}{IV.2.36 (p.~71)} }
\expandafter\def\csname SRef@maps:src:B.6.6\endcsname{\href{S6_SERIES_GRIND_IV.pdf\#nameddest=tag-02JY}{IV.2.37 (p.~73)} }
\expandafter\def\csname SRef@maps:src:B.6.7\endcsname{\href{S6_SERIES_GRIND_IV.pdf\#nameddest=tag-02JZ}{IV.2.38 (p.~75)} }
\expandafter\def\csname SRef@maps:src:B.6.10\endcsname{\href{S6_SERIES_GRIND_IV.pdf\#nameddest=tag-02JE}{IV.2.40 (p.~77)} }
\expandafter\def\csname SRef@maps:src:B.6.11\endcsname{\href{S6_SERIES_GRIND_IV.pdf\#nameddest=tag-02JF}{IV.2.41 (p.~78)} }
\expandafter\def\csname SRef@push:src:C.3.2\endcsname{\href{S6_SERIES_GRIND_IV.pdf\#nameddest=tag-031R}{IV.2.5 (p.~25)} }
\expandafter\def\csname SRef@maps:src:B.6.12\endcsname{\href{S6_SERIES_GRIND_IV.pdf\#nameddest=tag-02JG}{IV.2.42 (p.~80)} }
\expandafter\def\csname SRef@maps:src:B.6.13\endcsname{\href{S6_SERIES_GRIND_IV.pdf\#nameddest=tag-02JH}{IV.2.43 (p.~83)} }
\expandafter\def\csname SRef@maps:src:B.6.14\endcsname{\href{S6_SERIES_GRIND_IV.pdf\#nameddest=tag-02JI}{IV.2.44 (p.~86)} }
\expandafter\def\csname SRef@maps:src:B.6.15\endcsname{\href{S6_SERIES_GRIND_IV.pdf\#nameddest=tag-02JJ}{IV.2.45 (p.~87)} }
\expandafter\def\csname SRef@maps:src:B.6.16\endcsname{\href{S6_SERIES_GRIND_IV.pdf\#nameddest=tag-02JK}{IV.2.46 (p.~88)} }
\expandafter\def\csname SRef@maps:src:B.6.17\endcsname{\href{S6_SERIES_GRIND_IV.pdf\#nameddest=tag-02JL}{IV.2.47 (p.~91)} }
\expandafter\def\csname SRef@maps:src:B.6.18\endcsname{\href{S6_SERIES_GRIND_IV.pdf\#nameddest=tag-02JM}{IV.2.48 (p.~93)} }
\expandafter\def\csname SRef@maps:src:B.6.19\endcsname{\href{S6_SERIES_GRIND_IV.pdf\#nameddest=tag-02JN}{IV.2.49 (p.~94)} }
\expandafter\def\csname SRef@push:src:C.3.4\endcsname{\href{S6_SERIES_GRIND_IV.pdf\#nameddest=tag-0321}{IV.2.6 (p.~25)} }
\expandafter\def\csname SRef@maps:src:B.6.22\endcsname{\href{S6_SERIES_GRIND_IV.pdf\#nameddest=tag-02JR}{IV.2.52 (p.~99)} }
\expandafter\def\csname SRef@push:src:C.3.6\endcsname{\href{S6_SERIES_GRIND_IV.pdf\#nameddest=tag-0323}{IV.2.7 (p.~25)} }
\expandafter\def\csname SRef@push:src:C.3.7\endcsname{\href{S6_SERIES_GRIND_IV.pdf\#nameddest=tag-0324}{IV.2.8 (p.~26)} }
\expandafter\def\csname SRef@iv:correction:P3F.25\endcsname{\href{S6_SERIES_GRIND_IV.pdf\#nameddest=tag-01TX}{IV.8.1 (p.~933)} }
\expandafter\def\csname SRef@iv:correction:P3K.29\endcsname{\href{S6_SERIES_GRIND_IV.pdf\#nameddest=tag-01U6}{IV.8.10 (p.~937)} }
\expandafter\def\csname SRef@iv:correction:P3K.30\endcsname{\href{S6_SERIES_GRIND_IV.pdf\#nameddest=tag-01U7}{IV.8.11 (p.~938)} }
\expandafter\def\csname SRef@iv:correction:P3K.31\endcsname{\href{S6_SERIES_GRIND_IV.pdf\#nameddest=tag-01U8}{IV.8.12 (p.~938)} }
\expandafter\def\csname SRef@iv:correction:P3K.32\endcsname{\href{S6_SERIES_GRIND_IV.pdf\#nameddest=tag-01U9}{IV.8.13 (p.~939)} }
\expandafter\def\csname SRef@iv:correction:P3K.33\endcsname{\href{S6_SERIES_GRIND_IV.pdf\#nameddest=tag-01UA}{IV.8.14 (p.~939)} }
\expandafter\def\csname SRef@iv:correction:P3K.34\endcsname{\href{S6_SERIES_GRIND_IV.pdf\#nameddest=tag-01UB}{IV.8.15 (p.~940)} }
\expandafter\def\csname SRef@iv:correction:P3F.53\endcsname{\href{S6_SERIES_GRIND_IV.pdf\#nameddest=tag-01TY}{IV.8.16 (p.~940)} }
\expandafter\def\csname SRef@iv:correction:P3K.35\endcsname{\href{S6_SERIES_GRIND_IV.pdf\#nameddest=tag-01UC}{IV.8.17 (p.~940)} }
\expandafter\def\csname SRef@iv:correction:P3K.36\endcsname{\href{S6_SERIES_GRIND_IV.pdf\#nameddest=tag-01UD}{IV.8.18 (p.~941)} }
\expandafter\def\csname SRef@iv:correction:P3K.37\endcsname{\href{S6_SERIES_GRIND_IV.pdf\#nameddest=tag-01UE}{IV.8.19 (p.~941)} }
\expandafter\def\csname SRef@iv:correction:P3K.6\endcsname{\href{S6_SERIES_GRIND_IV.pdf\#nameddest=tag-01US}{IV.8.2 (p.~934)} }
\expandafter\def\csname SRef@iv:correction:P3K.40\endcsname{\href{S6_SERIES_GRIND_IV.pdf\#nameddest=tag-01UF}{IV.8.20 (p.~942)} }
\expandafter\def\csname SRef@iv:correction:P3K.41\endcsname{\href{S6_SERIES_GRIND_IV.pdf\#nameddest=tag-01UG}{IV.8.21 (p.~942)} }
\expandafter\def\csname SRef@iv:correction:P3K.43\endcsname{\href{S6_SERIES_GRIND_IV.pdf\#nameddest=tag-01UH}{IV.8.22 (p.~943)} }
\expandafter\def\csname SRef@iv:correction:P3K.44\endcsname{\href{S6_SERIES_GRIND_IV.pdf\#nameddest=tag-01UI}{IV.8.23 (p.~943)} }
\expandafter\def\csname SRef@iv:correction:P3K.45\endcsname{\href{S6_SERIES_GRIND_IV.pdf\#nameddest=tag-01UJ}{IV.8.24 (p.~944)} }
\expandafter\def\csname SRef@iv:correction:P3K.46\endcsname{\href{S6_SERIES_GRIND_IV.pdf\#nameddest=tag-01UK}{IV.8.25 (p.~944)} }
\expandafter\def\csname SRef@iv:correction:P3K.50\endcsname{\href{S6_SERIES_GRIND_IV.pdf\#nameddest=tag-01UL}{IV.8.26 (p.~944)} }
\expandafter\def\csname SRef@iv:correction:P3K.52\endcsname{\href{S6_SERIES_GRIND_IV.pdf\#nameddest=tag-01UM}{IV.8.27 (p.~944)} }
\expandafter\def\csname SRef@iv:correction:P3K.54\endcsname{\href{S6_SERIES_GRIND_IV.pdf\#nameddest=tag-01UN}{IV.8.28 (p.~945)} }
\expandafter\def\csname SRef@iv:correction:P3K.56\endcsname{\href{S6_SERIES_GRIND_IV.pdf\#nameddest=tag-01UO}{IV.8.29 (p.~945)} }
\expandafter\def\csname SRef@iv:correction:P3K.10\endcsname{\href{S6_SERIES_GRIND_IV.pdf\#nameddest=tag-01TZ}{IV.8.3 (p.~934)} }
\expandafter\def\csname SRef@iv:correction:P3K.57\endcsname{\href{S6_SERIES_GRIND_IV.pdf\#nameddest=tag-01UP}{IV.8.30 (p.~945)} }
\expandafter\def\csname SRef@iv:correction:P3K.58\endcsname{\href{S6_SERIES_GRIND_IV.pdf\#nameddest=tag-01UQ}{IV.8.31 (p.~946)} }
\expandafter\def\csname SRef@iv:correction:P3K.59\endcsname{\href{S6_SERIES_GRIND_IV.pdf\#nameddest=tag-01UR}{IV.8.32 (p.~946)} }
\expandafter\def\csname SRef@iv:correction:P3K.60\endcsname{\href{S6_SERIES_GRIND_IV.pdf\#nameddest=tag-01UT}{IV.8.33 (p.~947)} }
\expandafter\def\csname SRef@iv:correction:P3K.17\endcsname{\href{S6_SERIES_GRIND_IV.pdf\#nameddest=tag-01U0}{IV.8.4 (p.~935)} }
\expandafter\def\csname SRef@iv:correction:P3K.22\endcsname{\href{S6_SERIES_GRIND_IV.pdf\#nameddest=tag-01U1}{IV.8.5 (p.~935)} }
\expandafter\def\csname SRef@iv:correction:P3K.23\endcsname{\href{S6_SERIES_GRIND_IV.pdf\#nameddest=tag-01U2}{IV.8.6 (p.~936)} }
\expandafter\def\csname SRef@iv:correction:P3K.24\endcsname{\href{S6_SERIES_GRIND_IV.pdf\#nameddest=tag-01U3}{IV.8.7 (p.~936)} }
\expandafter\def\csname SRef@iv:correction:P3K.25\endcsname{\href{S6_SERIES_GRIND_IV.pdf\#nameddest=tag-01U4}{IV.8.8 (p.~937)} }
\expandafter\def\csname SRef@iv:correction:P3K.27\endcsname{\href{S6_SERIES_GRIND_IV.pdf\#nameddest=tag-01U5}{IV.8.9 (p.~937)} }
\expandafter\def\csname SRef@fam:src:E.2.3\endcsname{\href{S6_SERIES_GRIND_IV.pdf\#nameddest=tag-01R8}{IV.2.2 (p.~21)} }
\expandafter\def\csname SRef@construction-and-proof\endcsname{\href{S6_SERIES_GRIND_IV.pdf\#nameddest=tag-01GZ}{IV.3.0.1 (p.~127)} }
\expandafter\def\csname SRef@the-marked-double-and-its-bordism\endcsname{\href{S6_SERIES_GRIND_IV.pdf\#nameddest=tag-034Y}{IV.{{$\langle 1\rangle $1.}} (p.~144)} }
\expandafter\def\csname SRef@removal-of-the-cubic-coordinate-by-an-extending-map\endcsname{\href{S6_SERIES_GRIND_IV.pdf\#nameddest=tag-033W}{IV.{{$\langle 1\rangle $2.}} (p.~144)} }
\expandafter\def\csname SRef@the-surgery-and-its-bounding-eight-manifold\endcsname{\href{S6_SERIES_GRIND_IV.pdf\#nameddest=tag-0352}{IV.{{$\langle 1\rangle $3.}} (p.~144)} }
\expandafter\def\csname SRef@the-invariant-and-the-resulting-homomorphism\endcsname{\href{S6_SERIES_GRIND_IV.pdf\#nameddest=tag-034W}{IV.{{$\langle 1\rangle $4.}} (p.~145)} }
\expandafter\def\csname SRef@proof\endcsname{\href{S6_SERIES_GRIND_IV.pdf\#nameddest=tag-02YT}{IV.3.0.2 (p.~129)} }
\expandafter\def\csname SRef@an-explicit-rational-elliptic-surface-and-primitive-section\endcsname{\href{S6_SERIES_GRIND_IV.pdf\#nameddest=tag-0009}{IV.3.0.3 (p.~135)} }
\expandafter\def\csname SRef@the-linearization-and-every-local-filling\endcsname{\href{S6_SERIES_GRIND_IV.pdf\#nameddest=tag-034X}{IV.3.0.4 (p.~136)} }
\expandafter\def\csname SRef@actual-integral-monodromy-not-an-abstract-substitute-tuple\endcsname{\href{S6_SERIES_GRIND_IV.pdf\#nameddest=tag-0006}{IV.3.0.5 (p.~138)} }
\expandafter\def\csname SRef@fundamental-group-and-every-integral-cohomology-differential\endcsname{\href{S6_SERIES_GRIND_IV.pdf\#nameddest=tag-01RN}{IV.3.0.6 (p.~139)} }
\expandafter\def\csname SRef@the-punctured-holomorphic-family-and-its-closing-map\endcsname{\href{S6_SERIES_GRIND_IV.pdf\#nameddest=tag-0350}{IV.3.0.7 (p.~142)} }
\expandafter\def\csname SRef@the-actual-page-and-its-section-line\endcsname{\href{S6_SERIES_GRIND_IV.pdf\#nameddest=tag-034R}{IV.3.0.8 (p.~142)} }
\expandafter\def\csname SRef@executed-verification\endcsname{\href{S6_SERIES_GRIND_IV.pdf\#nameddest=tag-01OP}{IV.3.0.9 (p.~143)} }
\expandafter\def\csname SRef@night:P3K.75\endcsname{\href{S6_SERIES_GRIND_IV.pdf\#nameddest=tag-02VN}{IV.11.1 (p.~1392)} }
\expandafter\def\csname SRef@night:P3K.117\endcsname{\href{S6_SERIES_GRIND_IV.pdf\#nameddest=tag-02TO}{IV.10.7 (p.~1213)} }
\expandafter\def\csname SRef@night:P3K.118\endcsname{\href{S6_SERIES_GRIND_IV.pdf\#nameddest=tag-02TP}{IV.10.8 (p.~1250)} }
\expandafter\def\csname SRef@night:P3K.119\endcsname{\href{S6_SERIES_GRIND_IV.pdf\#nameddest=tag-02TQ}{IV.10.9 (p.~1278)} }
\expandafter\def\csname SRef@night:P3K.109\endcsname{\href{S6_SERIES_GRIND_IV.pdf\#nameddest=tag-02TF}{IV.10.1 (p.~1116)} }
\expandafter\def\csname SRef@night:P3K.110\endcsname{\href{S6_SERIES_GRIND_IV.pdf\#nameddest=tag-02TH}{IV.10.2 (p.~1127)} }
\expandafter\def\csname SRef@night:P3K.111\endcsname{\href{S6_SERIES_GRIND_IV.pdf\#nameddest=tag-02TI}{IV.10.10 (p.~1305)} }
\expandafter\def\csname SRef@night:P3K.112\endcsname{\href{S6_SERIES_GRIND_IV.pdf\#nameddest=tag-02TJ}{IV.10.11 (p.~1369)} }
\expandafter\def\csname SRef@night:P3K.113\endcsname{\href{S6_SERIES_GRIND_IV.pdf\#nameddest=tag-02TK}{IV.10.3 (p.~1135)} }
\expandafter\def\csname SRef@night:P3K.114\endcsname{\href{S6_SERIES_GRIND_IV.pdf\#nameddest=tag-02TL}{IV.10.4 (p.~1147)} }
\expandafter\def\csname SRef@night:P3K.115\endcsname{\href{S6_SERIES_GRIND_IV.pdf\#nameddest=tag-02TM}{IV.10.5 (p.~1177)} }
\expandafter\def\csname SRef@night:P3K.116\endcsname{\href{S6_SERIES_GRIND_IV.pdf\#nameddest=tag-02TN}{IV.10.6 (p.~1193)} }
\expandafter\def\csname SRef@push:src:C.3.24\endcsname{\href{S6_SERIES_GRIND_IV.pdf\#nameddest=tag-031W}{IV.2.25 (p.~57)} }
\expandafter\def\csname SRef@maps:src:B.6.2\endcsname{\href{S6_SERIES_GRIND_IV.pdf\#nameddest=tag-02JO}{IV.2.33 (p.~67)} }
\expandafter\def\csname SRef@maps:src:B.6.20\endcsname{\href{S6_SERIES_GRIND_IV.pdf\#nameddest=tag-02JP}{IV.2.50 (p.~96)} }
\expandafter\def\csname SRef@maps:src:B.6.21\endcsname{\href{S6_SERIES_GRIND_IV.pdf\#nameddest=tag-02JQ}{IV.2.51 (p.~99)} }
\expandafter\def\csname SRef@maps:src:B.6.23\endcsname{\href{S6_SERIES_GRIND_IV.pdf\#nameddest=tag-02JS}{IV.2.53 (p.~101)} }
\expandafter\def\csname SRef@maps:src:B.6.24\endcsname{\href{S6_SERIES_GRIND_IV.pdf\#nameddest=tag-02JT}{IV.2.54 (p.~102)} }
\expandafter\def\csname SRef@night:P3F.1\endcsname{\href{S6_SERIES_GRIND_IV.pdf\#nameddest=tag-02RK}{IV.3.1 (p.~104)} }
\expandafter\def\csname SRef@night:P3F.10\endcsname{\href{S6_SERIES_GRIND_IV.pdf\#nameddest=tag-02RL}{IV.3.10 (p.~117)} }
\expandafter\def\csname SRef@night:P3F.11\endcsname{\href{S6_SERIES_GRIND_IV.pdf\#nameddest=tag-02RM}{IV.3.11 (p.~119)} }
\expandafter\def\csname SRef@night:P3F.12\endcsname{\href{S6_SERIES_GRIND_IV.pdf\#nameddest=tag-02RN}{IV.3.12 (p.~121)} }
\expandafter\def\csname SRef@night:P3F.13\endcsname{\href{S6_SERIES_GRIND_IV.pdf\#nameddest=tag-02RO}{IV.3.13 (p.~123)} }
\expandafter\def\csname SRef@night:P3F.14\endcsname{\href{S6_SERIES_GRIND_IV.pdf\#nameddest=tag-02RP}{IV.3.14 (p.~124)} }
\expandafter\def\csname SRef@night:P3F.15\endcsname{\href{S6_SERIES_GRIND_IV.pdf\#nameddest=tag-02RQ}{IV.3.15 (p.~125)} }
\expandafter\def\csname SRef@night:P3F.16\endcsname{\href{S6_SERIES_GRIND_IV.pdf\#nameddest=tag-02RR}{IV.3.16 (p.~126)} }
\expandafter\def\csname SRef@night:P3F.17\endcsname{\href{S6_SERIES_GRIND_IV.pdf\#nameddest=tag-02RS}{IV.3.17 (p.~128)} }
\expandafter\def\csname SRef@night:P3F.18\endcsname{\href{S6_SERIES_GRIND_IV.pdf\#nameddest=tag-02RT}{IV.3.18 (p.~129)} }
\expandafter\def\csname SRef@night:P3F.19\endcsname{\href{S6_SERIES_GRIND_IV.pdf\#nameddest=tag-02RU}{IV.3.19 (p.~130)} }
\expandafter\def\csname SRef@night:P3F.2\endcsname{\href{S6_SERIES_GRIND_IV.pdf\#nameddest=tag-02RV}{IV.3.2 (p.~107)} }
\expandafter\def\csname SRef@night:P3F.20\endcsname{\href{S6_SERIES_GRIND_IV.pdf\#nameddest=tag-02RW}{IV.3.21 (p.~134)} }
\expandafter\def\csname SRef@night:P3F.21\endcsname{\href{S6_SERIES_GRIND_IV.pdf\#nameddest=tag-02RX}{IV.3.22 (p.~134)} }
\expandafter\def\csname SRef@night:P3F.22\endcsname{\href{S6_SERIES_GRIND_IV.pdf\#nameddest=tag-02RY}{IV.3.23 (p.~144)} }
\expandafter\def\csname SRef@night:P3F.23\endcsname{\href{S6_SERIES_GRIND_IV.pdf\#nameddest=tag-02RZ}{IV.3.24 (p.~145)} }
\expandafter\def\csname SRef@night:P3F.24\endcsname{\href{S6_SERIES_GRIND_IV.pdf\#nameddest=tag-02S0}{IV.3.25 (p.~147)} }
\expandafter\def\csname SRef@night:P3F.30\endcsname{\href{S6_SERIES_GRIND_IV.pdf\#nameddest=tag-02S7}{IV.3.26 (p.~148)} }
\expandafter\def\csname SRef@night:P3F.35\endcsname{\href{S6_SERIES_GRIND_IV.pdf\#nameddest=tag-02SC}{IV.3.28 (p.~152)} }
\expandafter\def\csname SRef@night:P3F.36\endcsname{\href{S6_SERIES_GRIND_IV.pdf\#nameddest=tag-02SD}{IV.3.29 (p.~154)} }
\expandafter\def\csname SRef@night:P3F.37\endcsname{\href{S6_SERIES_GRIND_IV.pdf\#nameddest=tag-02SE}{IV.3.30 (p.~159)} }
\expandafter\def\csname SRef@night:P3F.3\endcsname{\href{S6_SERIES_GRIND_IV.pdf\#nameddest=tag-02S6}{IV.3.3 (p.~108)} }
\expandafter\def\csname SRef@night:P3F.38\endcsname{\href{S6_SERIES_GRIND_IV.pdf\#nameddest=tag-02SF}{IV.3.31 (p.~164)} }
\expandafter\def\csname SRef@night:P3F.39\endcsname{\href{S6_SERIES_GRIND_IV.pdf\#nameddest=tag-02SG}{IV.3.32 (p.~169)} }
\expandafter\def\csname SRef@night:P3F.40\endcsname{\href{S6_SERIES_GRIND_IV.pdf\#nameddest=tag-02SI}{IV.3.33 (p.~172)} }
\expandafter\def\csname SRef@night:P3F.41\endcsname{\href{S6_SERIES_GRIND_IV.pdf\#nameddest=tag-02SJ}{IV.3.34 (p.~175)} }
\expandafter\def\csname SRef@night:P3F.42\endcsname{\href{S6_SERIES_GRIND_IV.pdf\#nameddest=tag-02SK}{IV.3.35 (p.~180)} }
\expandafter\def\csname SRef@night:P3F.43\endcsname{\href{S6_SERIES_GRIND_IV.pdf\#nameddest=tag-02SL}{IV.3.36 (p.~186)} }
\expandafter\def\csname SRef@night:P3F.44\endcsname{\href{S6_SERIES_GRIND_IV.pdf\#nameddest=tag-02SM}{IV.3.37 (p.~188)} }
\expandafter\def\csname SRef@night:P3F.45\endcsname{\href{S6_SERIES_GRIND_IV.pdf\#nameddest=tag-02SN}{IV.3.38 (p.~191)} }
\expandafter\def\csname SRef@night:P3F.46\endcsname{\href{S6_SERIES_GRIND_IV.pdf\#nameddest=tag-02SO}{IV.3.39 (p.~195)} }
\expandafter\def\csname SRef@night:P3F.47\endcsname{\href{S6_SERIES_GRIND_IV.pdf\#nameddest=tag-02SP}{IV.3.40 (p.~202)} }
\expandafter\def\csname SRef@night:P3F.4\endcsname{\href{S6_SERIES_GRIND_IV.pdf\#nameddest=tag-02SH}{IV.3.4 (p.~109)} }
\expandafter\def\csname SRef@night:P3F.48\endcsname{\href{S6_SERIES_GRIND_IV.pdf\#nameddest=tag-02SQ}{IV.3.41 (p.~207)} }
\expandafter\def\csname SRef@night:P3F.49\endcsname{\href{S6_SERIES_GRIND_IV.pdf\#nameddest=tag-02SR}{IV.3.42 (p.~210)} }
\expandafter\def\csname SRef@night:P3F.50\endcsname{\href{S6_SERIES_GRIND_IV.pdf\#nameddest=tag-02ST}{IV.3.43 (p.~212)} }
\expandafter\def\csname SRef@night:P3F.51\endcsname{\href{S6_SERIES_GRIND_IV.pdf\#nameddest=tag-02SU}{IV.3.44 (p.~216)} }
\expandafter\def\csname SRef@night:P3F.52\endcsname{\href{S6_SERIES_GRIND_IV.pdf\#nameddest=tag-02SV}{IV.3.45 (p.~222)} }
\expandafter\def\csname SRef@night:P3F.54\endcsname{\href{S6_SERIES_GRIND_IV.pdf\#nameddest=tag-02SX}{IV.3.46 (p.~225)} }
\expandafter\def\csname SRef@night:P3F.55\endcsname{\href{S6_SERIES_GRIND_IV.pdf\#nameddest=tag-02SY}{IV.3.47 (p.~229)} }
\expandafter\def\csname SRef@night:P3F.56\endcsname{\href{S6_SERIES_GRIND_IV.pdf\#nameddest=tag-02SZ}{IV.3.48 (p.~234)} }
\expandafter\def\csname SRef@night:P3F.5\endcsname{\href{S6_SERIES_GRIND_IV.pdf\#nameddest=tag-02SS}{IV.3.5 (p.~110)} }
\expandafter\def\csname SRef@night:P3F.6\endcsname{\href{S6_SERIES_GRIND_IV.pdf\#nameddest=tag-02T0}{IV.3.6 (p.~111)} }
\expandafter\def\csname SRef@night:P3F.7\endcsname{\href{S6_SERIES_GRIND_IV.pdf\#nameddest=tag-02T1}{IV.3.7 (p.~112)} }
\expandafter\def\csname SRef@night:P3F.8\endcsname{\href{S6_SERIES_GRIND_IV.pdf\#nameddest=tag-02T2}{IV.3.8 (p.~114)} }
\expandafter\def\csname SRef@night:P3F.9\endcsname{\href{S6_SERIES_GRIND_IV.pdf\#nameddest=tag-02T3}{IV.3.9 (p.~115)} }
\expandafter\def\csname SRef@night:P3R.1\endcsname{\href{S6_SERIES_GRIND_IV.pdf\#nameddest=tag-02WE}{IV.4.1 (p.~243)} }
\expandafter\def\csname SRef@night:P3K.9\endcsname{\href{S6_SERIES_GRIND_IV.pdf\#nameddest=tag-02W3}{IV.4.10 (p.~254)} }
\expandafter\def\csname SRef@night:P3K.11\endcsname{\href{S6_SERIES_GRIND_IV.pdf\#nameddest=tag-02TG}{IV.4.11 (p.~257)} }
\expandafter\def\csname SRef@night:P3K.12\endcsname{\href{S6_SERIES_GRIND_IV.pdf\#nameddest=tag-02TR}{IV.4.12 (p.~260)} }
\expandafter\def\csname SRef@night:P3K.13\endcsname{\href{S6_SERIES_GRIND_IV.pdf\#nameddest=tag-02TS}{IV.4.13 (p.~261)} }
\expandafter\def\csname SRef@night:P3K.14\endcsname{\href{S6_SERIES_GRIND_IV.pdf\#nameddest=tag-02TT}{IV.4.14 (p.~263)} }
\expandafter\def\csname SRef@night:P3K.15\endcsname{\href{S6_SERIES_GRIND_IV.pdf\#nameddest=tag-02TU}{IV.4.15 (p.~264)} }
\expandafter\def\csname SRef@night:P3K.16\endcsname{\href{S6_SERIES_GRIND_IV.pdf\#nameddest=tag-02TV}{IV.4.16 (p.~268)} }
\expandafter\def\csname SRef@night:P3K.18\endcsname{\href{S6_SERIES_GRIND_IV.pdf\#nameddest=tag-02TX}{IV.4.17 (p.~270)} }
\expandafter\def\csname SRef@night:P3K.19\endcsname{\href{S6_SERIES_GRIND_IV.pdf\#nameddest=tag-02TY}{IV.4.18 (p.~271)} }
\expandafter\def\csname SRef@night:P3K.20\endcsname{\href{S6_SERIES_GRIND_IV.pdf\#nameddest=tag-02U0}{IV.4.19 (p.~273)} }
\expandafter\def\csname SRef@night:P3R.2\endcsname{\href{S6_SERIES_GRIND_IV.pdf\#nameddest=tag-02WF}{IV.4.2 (p.~243)} }
\expandafter\def\csname SRef@night:P3K.21\endcsname{\href{S6_SERIES_GRIND_IV.pdf\#nameddest=tag-02U1}{IV.4.20 (p.~276)} }
\expandafter\def\csname SRef@night:P3K.26\endcsname{\href{S6_SERIES_GRIND_IV.pdf\#nameddest=tag-02U6}{IV.4.21 (p.~278)} }
\expandafter\def\csname SRef@night:P3K.28\endcsname{\href{S6_SERIES_GRIND_IV.pdf\#nameddest=tag-02U8}{IV.4.22 (p.~283)} }
\expandafter\def\csname SRef@night:P3K.1\endcsname{\href{S6_SERIES_GRIND_IV.pdf\#nameddest=tag-02T4}{IV.4.3 (p.~244)} }
\expandafter\def\csname SRef@night:P3K.2\endcsname{\href{S6_SERIES_GRIND_IV.pdf\#nameddest=tag-02TZ}{IV.4.4 (p.~245)} }
\expandafter\def\csname SRef@night:P3K.3\endcsname{\href{S6_SERIES_GRIND_IV.pdf\#nameddest=tag-02UA}{IV.4.5 (p.~246)} }
\expandafter\def\csname SRef@night:P3K.4\endcsname{\href{S6_SERIES_GRIND_IV.pdf\#nameddest=tag-02UK}{IV.4.6 (p.~247)} }
\expandafter\def\csname SRef@night:P3K.5\endcsname{\href{S6_SERIES_GRIND_IV.pdf\#nameddest=tag-02UV}{IV.4.7 (p.~248)} }
\expandafter\def\csname SRef@night:P3K.7\endcsname{\href{S6_SERIES_GRIND_IV.pdf\#nameddest=tag-02VH}{IV.4.8 (p.~249)} }
\expandafter\def\csname SRef@night:P3K.8\endcsname{\href{S6_SERIES_GRIND_IV.pdf\#nameddest=tag-02VS}{IV.4.9 (p.~251)} }
\expandafter\def\csname SRef@night:P3K.38\endcsname{\href{S6_SERIES_GRIND_IV.pdf\#nameddest=tag-02UJ}{IV.5.1 (p.~291)} }
\expandafter\def\csname SRef@night:P3K.62\endcsname{\href{S6_SERIES_GRIND_IV.pdf\#nameddest=tag-02V9}{IV.5.10 (p.~333)} }
\expandafter\def\csname SRef@night:P3K.63\endcsname{\href{S6_SERIES_GRIND_IV.pdf\#nameddest=tag-02VA}{IV.5.11 (p.~340)} }
\expandafter\def\csname SRef@night:P3K.64\endcsname{\href{S6_SERIES_GRIND_IV.pdf\#nameddest=tag-02VB}{IV.5.12 (p.~348)} }
\expandafter\def\csname SRef@night:P3K.65\endcsname{\href{S6_SERIES_GRIND_IV.pdf\#nameddest=tag-02VC}{IV.5.13 (p.~352)} }
\expandafter\def\csname SRef@night:P3K.66\endcsname{\href{S6_SERIES_GRIND_IV.pdf\#nameddest=tag-02VD}{IV.5.14 (p.~362)} }
\expandafter\def\csname SRef@night:P3K.67\endcsname{\href{S6_SERIES_GRIND_IV.pdf\#nameddest=tag-02VE}{IV.5.15 (p.~369)} }
\expandafter\def\csname SRef@night:P3K.42\endcsname{\href{S6_SERIES_GRIND_IV.pdf\#nameddest=tag-02UN}{IV.5.2 (p.~298)} }
\expandafter\def\csname SRef@night:P3K.47\endcsname{\href{S6_SERIES_GRIND_IV.pdf\#nameddest=tag-02US}{IV.5.3 (p.~304)} }
\expandafter\def\csname SRef@night:P3K.48\endcsname{\href{S6_SERIES_GRIND_IV.pdf\#nameddest=tag-02UT}{IV.5.4 (p.~307)} }
\expandafter\def\csname SRef@night:P3K.49\endcsname{\href{S6_SERIES_GRIND_IV.pdf\#nameddest=tag-02UU}{IV.5.5 (p.~311)} }
\expandafter\def\csname SRef@night:P3K.51\endcsname{\href{S6_SERIES_GRIND_IV.pdf\#nameddest=tag-02UX}{IV.5.6 (p.~318)} }
\expandafter\def\csname SRef@night:P3K.53\endcsname{\href{S6_SERIES_GRIND_IV.pdf\#nameddest=tag-02UZ}{IV.5.7 (p.~321)} }
\expandafter\def\csname SRef@night:P3K.55\endcsname{\href{S6_SERIES_GRIND_IV.pdf\#nameddest=tag-02V1}{IV.5.8 (p.~325)} }
\expandafter\def\csname SRef@night:P3K.61\endcsname{\href{S6_SERIES_GRIND_IV.pdf\#nameddest=tag-02V8}{IV.5.9 (p.~330)} }
\expandafter\def\csname SRef@night:P3K.68\endcsname{\href{S6_SERIES_GRIND_IV.pdf\#nameddest=tag-02VF}{IV.6.1 (p.~374)} }
\expandafter\def\csname SRef@night:P3K.69\endcsname{\href{S6_SERIES_GRIND_IV.pdf\#nameddest=tag-02VG}{IV.6.2 (p.~380)} }
\expandafter\def\csname SRef@night:P3K.70\endcsname{\href{S6_SERIES_GRIND_IV.pdf\#nameddest=tag-02VI}{IV.6.3 (p.~384)} }
\expandafter\def\csname SRef@night:P3K.71\endcsname{\href{S6_SERIES_GRIND_IV.pdf\#nameddest=tag-02VJ}{IV.6.4 (p.~389)} }
\expandafter\def\csname SRef@night:P3K.72\endcsname{\href{S6_SERIES_GRIND_IV.pdf\#nameddest=tag-02VK}{IV.6.5 (p.~393)} }
\expandafter\def\csname SRef@night:P3K.73\endcsname{\href{S6_SERIES_GRIND_IV.pdf\#nameddest=tag-02VL}{IV.6.6 (p.~400)} }
\expandafter\def\csname SRef@night:P3K.74\endcsname{\href{S6_SERIES_GRIND_IV.pdf\#nameddest=tag-02VM}{IV.6.7 (p.~405)} }
\expandafter\def\csname SRef@night:P3K.76\endcsname{\href{S6_SERIES_GRIND_IV.pdf\#nameddest=tag-02VO}{IV.7.1 (p.~407)} }
\expandafter\def\csname SRef@night:P3K.85\endcsname{\href{S6_SERIES_GRIND_IV.pdf\#nameddest=tag-02VY}{IV.7.10 (p.~494)} }
\expandafter\def\csname SRef@night:P3K.86\endcsname{\href{S6_SERIES_GRIND_IV.pdf\#nameddest=tag-02VZ}{IV.7.11 (p.~499)} }
\expandafter\def\csname SRef@night:P3K.87\endcsname{\href{S6_SERIES_GRIND_IV.pdf\#nameddest=tag-02W0}{IV.7.12 (p.~505)} }
\expandafter\def\csname SRef@night:P3K.88\endcsname{\href{S6_SERIES_GRIND_IV.pdf\#nameddest=tag-02W1}{IV.7.13 (p.~509)} }
\expandafter\def\csname SRef@night:P3K.89\endcsname{\href{S6_SERIES_GRIND_IV.pdf\#nameddest=tag-02W2}{IV.7.14 (p.~513)} }
\expandafter\def\csname SRef@night:P3K.90\endcsname{\href{S6_SERIES_GRIND_IV.pdf\#nameddest=tag-02W4}{IV.7.15 (p.~521)} }
\expandafter\def\csname SRef@night:P3K.91\endcsname{\href{S6_SERIES_GRIND_IV.pdf\#nameddest=tag-02W5}{IV.7.16 (p.~525)} }
\expandafter\def\csname SRef@night:P3K.92\endcsname{\href{S6_SERIES_GRIND_IV.pdf\#nameddest=tag-02W6}{IV.7.17 (p.~531)} }
\expandafter\def\csname SRef@night:P3K.93\endcsname{\href{S6_SERIES_GRIND_IV.pdf\#nameddest=tag-02W7}{IV.7.18 (p.~538)} }
\expandafter\def\csname SRef@night:P3K.94\endcsname{\href{S6_SERIES_GRIND_IV.pdf\#nameddest=tag-02W8}{IV.7.19 (p.~543)} }
\expandafter\def\csname SRef@night:P3K.77\endcsname{\href{S6_SERIES_GRIND_IV.pdf\#nameddest=tag-02VP}{IV.7.2 (p.~414)} }
\expandafter\def\csname SRef@night:P3K.95\endcsname{\href{S6_SERIES_GRIND_IV.pdf\#nameddest=tag-02W9}{IV.7.20 (p.~551)} }
\expandafter\def\csname SRef@night:P3K.96\endcsname{\href{S6_SERIES_GRIND_IV.pdf\#nameddest=tag-02WA}{IV.7.21 (p.~562)} }
\expandafter\def\csname SRef@night:P3K.97\endcsname{\href{S6_SERIES_GRIND_IV.pdf\#nameddest=tag-02WB}{IV.7.22 (p.~569)} }
\expandafter\def\csname SRef@night:P3K.98\endcsname{\href{S6_SERIES_GRIND_IV.pdf\#nameddest=tag-02WC}{IV.7.23 (p.~578)} }
\expandafter\def\csname SRef@night:P3K.99\endcsname{\href{S6_SERIES_GRIND_IV.pdf\#nameddest=tag-02WD}{IV.7.24 (p.~587)} }
\expandafter\def\csname SRef@night:P3K.100\endcsname{\href{S6_SERIES_GRIND_IV.pdf\#nameddest=tag-02T6}{IV.7.25 (p.~662)} }
\expandafter\def\csname SRef@night:P3K.101\endcsname{\href{S6_SERIES_GRIND_IV.pdf\#nameddest=tag-02T7}{IV.7.26 (p.~669)} }
\expandafter\def\csname SRef@night:P3K.102\endcsname{\href{S6_SERIES_GRIND_IV.pdf\#nameddest=tag-02T8}{IV.7.27 (p.~758)} }
\expandafter\def\csname SRef@night:P3K.103\endcsname{\href{S6_SERIES_GRIND_IV.pdf\#nameddest=tag-02T9}{IV.7.28 (p.~768)} }
\expandafter\def\csname SRef@night:P3K.104\endcsname{\href{S6_SERIES_GRIND_IV.pdf\#nameddest=tag-02TA}{IV.7.29 (p.~774)} }
\expandafter\def\csname SRef@night:P3K.78\endcsname{\href{S6_SERIES_GRIND_IV.pdf\#nameddest=tag-02VQ}{IV.7.3 (p.~420)} }
\expandafter\def\csname SRef@night:P3K.105\endcsname{\href{S6_SERIES_GRIND_IV.pdf\#nameddest=tag-02TB}{IV.7.30 (p.~785)} }
\expandafter\def\csname SRef@night:P3K.106\endcsname{\href{S6_SERIES_GRIND_IV.pdf\#nameddest=tag-02TC}{IV.7.31 (p.~798)} }
\expandafter\def\csname SRef@night:P3K.108\endcsname{\href{S6_SERIES_GRIND_IV.pdf\#nameddest=tag-02TE}{IV.7.32 (p.~861)} }
\expandafter\def\csname SRef@iv:i3_all_reduced_caps\endcsname{\href{S6_SERIES_GRIND_IV.pdf\#nameddest=tag-02TE}{IV.7.32 (p.~861)} }
\expandafter\def\csname SRef@night:P3K.107\endcsname{\href{S6_SERIES_GRIND_IV.pdf\#nameddest=tag-02TD}{IV.7.33 (p.~867)} }
\expandafter\def\csname SRef@night:P3K.79\endcsname{\href{S6_SERIES_GRIND_IV.pdf\#nameddest=tag-02VR}{IV.7.4 (p.~426)} }
\expandafter\def\csname SRef@night:P3K.80\endcsname{\href{S6_SERIES_GRIND_IV.pdf\#nameddest=tag-02VT}{IV.7.5 (p.~431)} }
\expandafter\def\csname SRef@night:P3K.81\endcsname{\href{S6_SERIES_GRIND_IV.pdf\#nameddest=tag-02VU}{IV.7.6 (p.~467)} }
\expandafter\def\csname SRef@night:P3K.82\endcsname{\href{S6_SERIES_GRIND_IV.pdf\#nameddest=tag-02VV}{IV.7.7 (p.~475)} }
\expandafter\def\csname SRef@night:P3K.83\endcsname{\href{S6_SERIES_GRIND_IV.pdf\#nameddest=tag-02VW}{IV.7.8 (p.~478)} }
\expandafter\def\csname SRef@night:P3K.84\endcsname{\href{S6_SERIES_GRIND_IV.pdf\#nameddest=tag-02VX}{IV.7.9 (p.~490)} }
\expandafter\def\csname SRef@night:P3F.25\endcsname{\href{S6_SERIES_GRIND_IV.pdf\#nameddest=tag-02S1}{IV.9.1 (p.~947)} }
\expandafter\def\csname SRef@night:P3K.29\endcsname{\href{S6_SERIES_GRIND_IV.pdf\#nameddest=tag-02U9}{IV.9.10 (p.~988)} }
\expandafter\def\csname SRef@night:P3K.30\endcsname{\href{S6_SERIES_GRIND_IV.pdf\#nameddest=tag-02UB}{IV.9.11 (p.~992)} }
\expandafter\def\csname SRef@night:P3K.31\endcsname{\href{S6_SERIES_GRIND_IV.pdf\#nameddest=tag-02UC}{IV.9.12 (p.~997)} }
\expandafter\def\csname SRef@night:P3K.32\endcsname{\href{S6_SERIES_GRIND_IV.pdf\#nameddest=tag-02UD}{IV.9.13 (p.~1000)} }
\expandafter\def\csname SRef@night:P3K.33\endcsname{\href{S6_SERIES_GRIND_IV.pdf\#nameddest=tag-02UE}{IV.9.14 (p.~1010)} }
\expandafter\def\csname SRef@night:P3K.34\endcsname{\href{S6_SERIES_GRIND_IV.pdf\#nameddest=tag-02UF}{IV.9.15 (p.~1015)} }
\expandafter\def\csname SRef@night:P3F.53\endcsname{\href{S6_SERIES_GRIND_IV.pdf\#nameddest=tag-02SW}{IV.9.16 (p.~1020)} }
\expandafter\def\csname SRef@night:P3K.35\endcsname{\href{S6_SERIES_GRIND_IV.pdf\#nameddest=tag-02UG}{IV.9.17 (p.~1025)} }
\expandafter\def\csname SRef@night:P3K.37\endcsname{\href{S6_SERIES_GRIND_IV.pdf\#nameddest=tag-02UI}{IV.9.18 (p.~1032)} }
\expandafter\def\csname SRef@night:P3K.40\endcsname{\href{S6_SERIES_GRIND_IV.pdf\#nameddest=tag-02UL}{IV.9.19 (p.~1036)} }
\expandafter\def\csname SRef@night:P3K.6\endcsname{\href{S6_SERIES_GRIND_IV.pdf\#nameddest=tag-02V6}{IV.9.2 (p.~950)} }
\expandafter\def\csname SRef@night:P3K.41\endcsname{\href{S6_SERIES_GRIND_IV.pdf\#nameddest=tag-02UM}{IV.9.20 (p.~1040)} }
\expandafter\def\csname SRef@night:P3K.43\endcsname{\href{S6_SERIES_GRIND_IV.pdf\#nameddest=tag-02UO}{IV.9.21 (p.~1047)} }
\expandafter\def\csname SRef@night:P3K.44\endcsname{\href{S6_SERIES_GRIND_IV.pdf\#nameddest=tag-02UP}{IV.9.22 (p.~1050)} }
\expandafter\def\csname SRef@night:P3K.45\endcsname{\href{S6_SERIES_GRIND_IV.pdf\#nameddest=tag-02UQ}{IV.9.23 (p.~1055)} }
\expandafter\def\csname SRef@night:P3K.46\endcsname{\href{S6_SERIES_GRIND_IV.pdf\#nameddest=tag-02UR}{IV.9.24 (p.~1058)} }
\expandafter\def\csname SRef@night:P3K.50\endcsname{\href{S6_SERIES_GRIND_IV.pdf\#nameddest=tag-02UW}{IV.9.25 (p.~1061)} }
\expandafter\def\csname SRef@night:P3K.52\endcsname{\href{S6_SERIES_GRIND_IV.pdf\#nameddest=tag-02UY}{IV.9.26 (p.~1064)} }
\expandafter\def\csname SRef@night:P3K.54\endcsname{\href{S6_SERIES_GRIND_IV.pdf\#nameddest=tag-02V0}{IV.9.27 (p.~1067)} }
\expandafter\def\csname SRef@night:P3K.56\endcsname{\href{S6_SERIES_GRIND_IV.pdf\#nameddest=tag-02V2}{IV.9.28 (p.~1069)} }
\expandafter\def\csname SRef@night:P3K.57\endcsname{\href{S6_SERIES_GRIND_IV.pdf\#nameddest=tag-02V3}{IV.9.29 (p.~1072)} }
\expandafter\def\csname SRef@night:P3K.10\endcsname{\href{S6_SERIES_GRIND_IV.pdf\#nameddest=tag-02T5}{IV.9.3 (p.~952)} }
\expandafter\def\csname SRef@night:P3K.58\endcsname{\href{S6_SERIES_GRIND_IV.pdf\#nameddest=tag-02V4}{IV.9.30 (p.~1095)} }
\expandafter\def\csname SRef@night:P3K.59\endcsname{\href{S6_SERIES_GRIND_IV.pdf\#nameddest=tag-02V5}{IV.9.31 (p.~1104)} }
\expandafter\def\csname SRef@night:P3K.60\endcsname{\href{S6_SERIES_GRIND_IV.pdf\#nameddest=tag-02V7}{IV.9.32 (p.~1109)} }
\expandafter\def\csname SRef@night:P3K.17\endcsname{\href{S6_SERIES_GRIND_IV.pdf\#nameddest=tag-02TW}{IV.9.4 (p.~961)} }
\expandafter\def\csname SRef@night:P3K.22\endcsname{\href{S6_SERIES_GRIND_IV.pdf\#nameddest=tag-02U2}{IV.9.5 (p.~965)} }
\expandafter\def\csname SRef@night:P3K.23\endcsname{\href{S6_SERIES_GRIND_IV.pdf\#nameddest=tag-02U3}{IV.9.6 (p.~971)} }
\expandafter\def\csname SRef@night:P3K.24\endcsname{\href{S6_SERIES_GRIND_IV.pdf\#nameddest=tag-02U4}{IV.9.7 (p.~975)} }
\expandafter\def\csname SRef@night:P3K.25\endcsname{\href{S6_SERIES_GRIND_IV.pdf\#nameddest=tag-02U5}{IV.9.8 (p.~980)} }
\expandafter\def\csname SRef@night:P3K.27\endcsname{\href{S6_SERIES_GRIND_IV.pdf\#nameddest=tag-02U7}{IV.9.9 (p.~983)} }
\expandafter\def\csname SRef@maps:eq:F.6\endcsname{\href{S6_SERIES_GRIND_V.pdf\#nameddest=tag-02FA}{V.{F.6} (p.~85)} }
\expandafter\def\csname SRef@maps:eq:F.L2\endcsname{\href{S6_SERIES_GRIND_V.pdf\#nameddest=tag-02GV}{V.{F.L2} (p.~174)} }
\expandafter\def\csname SRef@maps:eq:F.L3\endcsname{\href{S6_SERIES_GRIND_V.pdf\#nameddest=tag-02GW}{V.{F.L3} (p.~174)} }
\expandafter\def\csname SRef@maps:eq:F.L4\endcsname{\href{S6_SERIES_GRIND_V.pdf\#nameddest=tag-02GX}{V.{F.L4} (p.~175)} }
\expandafter\def\csname SRef@maps:eq:F.L5\endcsname{\href{S6_SERIES_GRIND_V.pdf\#nameddest=tag-02GY}{V.{F.L5} (p.~175)} }
\expandafter\def\csname SRef@maps:eq:F.L6\endcsname{\href{S6_SERIES_GRIND_V.pdf\#nameddest=tag-02GZ}{V.{F.L6} (p.~175)} }
\expandafter\def\csname SRef@maps:eq:F.L7\endcsname{\href{S6_SERIES_GRIND_V.pdf\#nameddest=tag-02H0}{V.{F.L7} (p.~175)} }
\expandafter\def\csname SRef@maps:eq:F.L8\endcsname{\href{S6_SERIES_GRIND_V.pdf\#nameddest=tag-02H1}{V.{F.L8} (p.~175)} }
\expandafter\def\csname SRef@maps:eq:F.L9\endcsname{\href{S6_SERIES_GRIND_V.pdf\#nameddest=tag-02H2}{V.{F.L9} (p.~176)} }
\expandafter\def\csname SRef@maps:eq:F.L10\endcsname{\href{S6_SERIES_GRIND_V.pdf\#nameddest=tag-02GU}{V.{F.L10} (p.~176)} }
\expandafter\def\csname SRef@maps:eq:F.T1\endcsname{\href{S6_SERIES_GRIND_V.pdf\#nameddest=tag-02I7}{V.{F.T1} (p.~177)} }
\expandafter\def\csname SRef@maps:eq:F.T2\endcsname{\href{S6_SERIES_GRIND_V.pdf\#nameddest=tag-02IC}{V.{F.T2} (p.~177)} }
\expandafter\def\csname SRef@maps:eq:F.T3\endcsname{\href{S6_SERIES_GRIND_V.pdf\#nameddest=tag-02ID}{V.{F.T3} (p.~177)} }
\expandafter\def\csname SRef@maps:eq:F.T4\endcsname{\href{S6_SERIES_GRIND_V.pdf\#nameddest=tag-02IE}{V.{F.T4} (p.~178)} }
\expandafter\def\csname SRef@maps:eq:F.T5\endcsname{\href{S6_SERIES_GRIND_V.pdf\#nameddest=tag-02IF}{V.{F.T5} (p.~178)} }
\expandafter\def\csname SRef@maps:eq:F.T6\endcsname{\href{S6_SERIES_GRIND_V.pdf\#nameddest=tag-02IG}{V.{F.T6} (p.~178)} }
\expandafter\def\csname SRef@maps:eq:F.T7\endcsname{\href{S6_SERIES_GRIND_V.pdf\#nameddest=tag-02IH}{V.{F.T7} (p.~178)} }
\expandafter\def\csname SRef@maps:eq:F.T8\endcsname{\href{S6_SERIES_GRIND_V.pdf\#nameddest=tag-02II}{V.{F.T8} (p.~178)} }
\expandafter\def\csname SRef@maps:eq:F.T9\endcsname{\href{S6_SERIES_GRIND_V.pdf\#nameddest=tag-02IJ}{V.{F.T9} (p.~179)} }
\expandafter\def\csname SRef@maps:eq:F.T10\endcsname{\href{S6_SERIES_GRIND_V.pdf\#nameddest=tag-02I8}{V.{F.T10} (p.~179)} }
\expandafter\def\csname SRef@maps:eq:F.S1\endcsname{\href{S6_SERIES_GRIND_V.pdf\#nameddest=tag-02HX}{V.{F.S1} (p.~180)} }
\expandafter\def\csname SRef@maps:eq:F.S2\endcsname{\href{S6_SERIES_GRIND_V.pdf\#nameddest=tag-02HY}{V.{F.S2} (p.~180)} }
\expandafter\def\csname SRef@maps:eq:F.S3\endcsname{\href{S6_SERIES_GRIND_V.pdf\#nameddest=tag-02HZ}{V.{F.S3} (p.~180)} }
\expandafter\def\csname SRef@maps:eq:F.S4\endcsname{\href{S6_SERIES_GRIND_V.pdf\#nameddest=tag-02I0}{V.{F.S4} (p.~180)} }
\expandafter\def\csname SRef@maps:eq:F.S5\endcsname{\href{S6_SERIES_GRIND_V.pdf\#nameddest=tag-02I1}{V.{F.S5} (p.~180)} }
\expandafter\def\csname SRef@maps:eq:F.S6\endcsname{\href{S6_SERIES_GRIND_V.pdf\#nameddest=tag-02I2}{V.{F.S6} (p.~181)} }
\expandafter\def\csname SRef@maps:eq:F.S7\endcsname{\href{S6_SERIES_GRIND_V.pdf\#nameddest=tag-02I3}{V.{F.S7} (p.~181)} }
\expandafter\def\csname SRef@maps:eq:F.S8\endcsname{\href{S6_SERIES_GRIND_V.pdf\#nameddest=tag-02I4}{V.{F.S8} (p.~181)} }
\expandafter\def\csname SRef@maps:eq:F.S9\endcsname{\href{S6_SERIES_GRIND_V.pdf\#nameddest=tag-02I5}{V.{F.S9} (p.~181)} }
\expandafter\def\csname SRef@maps:eq:F.R10\endcsname{\href{S6_SERIES_GRIND_V.pdf\#nameddest=tag-02HK}{V.{F.R10} (p.~182)} }
\expandafter\def\csname SRef@maps:eq:F.R11\endcsname{\href{S6_SERIES_GRIND_V.pdf\#nameddest=tag-02HL}{V.{F.R11} (p.~182)} }
\expandafter\def\csname SRef@maps:eq:F.R12\endcsname{\href{S6_SERIES_GRIND_V.pdf\#nameddest=tag-02HM}{V.{F.R12} (p.~182)} }
\expandafter\def\csname SRef@maps:eq:F.R13\endcsname{\href{S6_SERIES_GRIND_V.pdf\#nameddest=tag-02HN}{V.{F.R13} (p.~182)} }
\expandafter\def\csname SRef@maps:eq:F.R14\endcsname{\href{S6_SERIES_GRIND_V.pdf\#nameddest=tag-02HO}{V.{F.R14} (p.~182)} }
\expandafter\def\csname SRef@maps:eq:F.X1\endcsname{\href{S6_SERIES_GRIND_V.pdf\#nameddest=tag-02IV}{V.{F.X1} (p.~183)} }
\expandafter\def\csname SRef@maps:eq:F.X2\endcsname{\href{S6_SERIES_GRIND_V.pdf\#nameddest=tag-02IW}{V.{F.X2} (p.~183)} }
\expandafter\def\csname SRef@maps:eq:F.X3\endcsname{\href{S6_SERIES_GRIND_V.pdf\#nameddest=tag-02IX}{V.{F.X3} (p.~184)} }
\expandafter\def\csname SRef@maps:eq:F.X4\endcsname{\href{S6_SERIES_GRIND_V.pdf\#nameddest=tag-02IY}{V.{F.X4} (p.~184)} }
\expandafter\def\csname SRef@maps:eq:F.X5\endcsname{\href{S6_SERIES_GRIND_V.pdf\#nameddest=tag-02IZ}{V.{F.X5} (p.~184)} }
\expandafter\def\csname SRef@maps:eq:F.T11\endcsname{\href{S6_SERIES_GRIND_V.pdf\#nameddest=tag-02I9}{V.{F.T11} (p.~185)} }
\expandafter\def\csname SRef@maps:eq:F.T12\endcsname{\href{S6_SERIES_GRIND_V.pdf\#nameddest=tag-02IA}{V.{F.T12} (p.~185)} }
\expandafter\def\csname SRef@maps:eq:F.T13\endcsname{\href{S6_SERIES_GRIND_V.pdf\#nameddest=tag-02IB}{V.{F.T13} (p.~185)} }
\expandafter\def\csname SRef@maps:eq:F.HC1\endcsname{\href{S6_SERIES_GRIND_V.pdf\#nameddest=tag-02G9}{V.{F.HC1} (p.~186)} }
\expandafter\def\csname SRef@maps:eq:F.HC2\endcsname{\href{S6_SERIES_GRIND_V.pdf\#nameddest=tag-02GE}{V.{F.HC2} (p.~186)} }
\expandafter\def\csname SRef@maps:eq:F.HC3\endcsname{\href{S6_SERIES_GRIND_V.pdf\#nameddest=tag-02GF}{V.{F.HC3} (p.~186)} }
\expandafter\def\csname SRef@maps:eq:F.HC4\endcsname{\href{S6_SERIES_GRIND_V.pdf\#nameddest=tag-02GG}{V.{F.HC4} (p.~187)} }
\expandafter\def\csname SRef@maps:eq:F.HC5\endcsname{\href{S6_SERIES_GRIND_V.pdf\#nameddest=tag-02GH}{V.{F.HC5} (p.~187)} }
\expandafter\def\csname SRef@maps:eq:F.HC6\endcsname{\href{S6_SERIES_GRIND_V.pdf\#nameddest=tag-02GI}{V.{F.HC6} (p.~187)} }
\expandafter\def\csname SRef@maps:eq:F.HC7\endcsname{\href{S6_SERIES_GRIND_V.pdf\#nameddest=tag-02GJ}{V.{F.HC7} (p.~188)} }
\expandafter\def\csname SRef@maps:eq:F.HC8\endcsname{\href{S6_SERIES_GRIND_V.pdf\#nameddest=tag-02GK}{V.{F.HC8} (p.~188)} }
\expandafter\def\csname SRef@maps:eq:F.HC9\endcsname{\href{S6_SERIES_GRIND_V.pdf\#nameddest=tag-02GL}{V.{F.HC9} (p.~189)} }
\expandafter\def\csname SRef@maps:eq:F.HC10\endcsname{\href{S6_SERIES_GRIND_V.pdf\#nameddest=tag-02GA}{V.{F.HC10} (p.~189)} }
\expandafter\def\csname SRef@maps:eq:F.HC11\endcsname{\href{S6_SERIES_GRIND_V.pdf\#nameddest=tag-02GB}{V.{F.HC11} (p.~189)} }
\expandafter\def\csname SRef@maps:eq:F.HC12\endcsname{\href{S6_SERIES_GRIND_V.pdf\#nameddest=tag-02GC}{V.{F.HC12} (p.~189)} }
\expandafter\def\csname SRef@maps:eq:F.HC13\endcsname{\href{S6_SERIES_GRIND_V.pdf\#nameddest=tag-02GD}{V.{F.HC13} (p.~190)} }
\expandafter\def\csname SRef@maps:eq:F.G1\endcsname{\href{S6_SERIES_GRIND_V.pdf\#nameddest=tag-02FO}{V.{F.G1} (p.~191)} }
\expandafter\def\csname SRef@maps:eq:F.G2\endcsname{\href{S6_SERIES_GRIND_V.pdf\#nameddest=tag-02FR}{V.{F.G2} (p.~191)} }
\expandafter\def\csname SRef@maps:eq:F.G3\endcsname{\href{S6_SERIES_GRIND_V.pdf\#nameddest=tag-02FS}{V.{F.G3} (p.~191)} }
\expandafter\def\csname SRef@maps:eq:F.G4\endcsname{\href{S6_SERIES_GRIND_V.pdf\#nameddest=tag-02FT}{V.{F.G4} (p.~191)} }
\expandafter\def\csname SRef@maps:eq:F.G5\endcsname{\href{S6_SERIES_GRIND_V.pdf\#nameddest=tag-02FU}{V.{F.G5} (p.~191)} }
\expandafter\def\csname SRef@maps:eq:F.G6\endcsname{\href{S6_SERIES_GRIND_V.pdf\#nameddest=tag-02FV}{V.{F.G6} (p.~191)} }
\expandafter\def\csname SRef@maps:eq:F.G7\endcsname{\href{S6_SERIES_GRIND_V.pdf\#nameddest=tag-02FW}{V.{F.G7} (p.~192)} }
\expandafter\def\csname SRef@maps:eq:F.G8\endcsname{\href{S6_SERIES_GRIND_V.pdf\#nameddest=tag-02FX}{V.{F.G8} (p.~192)} }
\expandafter\def\csname SRef@maps:eq:F.G9\endcsname{\href{S6_SERIES_GRIND_V.pdf\#nameddest=tag-02FY}{V.{F.G9} (p.~192)} }
\expandafter\def\csname SRef@maps:eq:F.G10\endcsname{\href{S6_SERIES_GRIND_V.pdf\#nameddest=tag-02FP}{V.{F.G10} (p.~192)} }
\expandafter\def\csname SRef@maps:eq:F.G11\endcsname{\href{S6_SERIES_GRIND_V.pdf\#nameddest=tag-02FQ}{V.{F.G11} (p.~194)} }
\expandafter\def\csname SRef@maps:eq:F.LT1\endcsname{\href{S6_SERIES_GRIND_V.pdf\#nameddest=tag-02H3}{V.{F.LT1} (p.~195)} }
\expandafter\def\csname SRef@maps:eq:F.LT2\endcsname{\href{S6_SERIES_GRIND_V.pdf\#nameddest=tag-02H4}{V.{F.LT2} (p.~195)} }
\expandafter\def\csname SRef@maps:eq:F.LT3\endcsname{\href{S6_SERIES_GRIND_V.pdf\#nameddest=tag-02H5}{V.{F.LT3} (p.~195)} }
\expandafter\def\csname SRef@maps:eq:F.LT4\endcsname{\href{S6_SERIES_GRIND_V.pdf\#nameddest=tag-02H6}{V.{F.LT4} (p.~195)} }
\expandafter\def\csname SRef@maps:eq:F.LT5\endcsname{\href{S6_SERIES_GRIND_V.pdf\#nameddest=tag-02H7}{V.{F.LT5} (p.~195)} }
\expandafter\def\csname SRef@maps:eq:F.LT6\endcsname{\href{S6_SERIES_GRIND_V.pdf\#nameddest=tag-02H8}{V.{F.LT6} (p.~196)} }
\expandafter\def\csname SRef@maps:eq:F.LT7\endcsname{\href{S6_SERIES_GRIND_V.pdf\#nameddest=tag-02H9}{V.{F.LT7} (p.~196)} }
\expandafter\def\csname SRef@maps:eq:F.LT8\endcsname{\href{S6_SERIES_GRIND_V.pdf\#nameddest=tag-02HA}{V.{F.LT8} (p.~196)} }
\expandafter\def\csname SRef@maps:eq:F.LT9\endcsname{\href{S6_SERIES_GRIND_V.pdf\#nameddest=tag-02HB}{V.{F.LT9} (p.~197)} }
\expandafter\def\csname SRef@maps:eq:F.CA1\endcsname{\href{S6_SERIES_GRIND_V.pdf\#nameddest=tag-02FE}{V.{F.CA1} (p.~197)} }
\expandafter\def\csname SRef@maps:eq:F.CA2\endcsname{\href{S6_SERIES_GRIND_V.pdf\#nameddest=tag-02FF}{V.{F.CA2} (p.~198)} }
\expandafter\def\csname SRef@maps:eq:F.CA3\endcsname{\href{S6_SERIES_GRIND_V.pdf\#nameddest=tag-02FG}{V.{F.CA3} (p.~198)} }
\expandafter\def\csname SRef@maps:eq:F.CA4\endcsname{\href{S6_SERIES_GRIND_V.pdf\#nameddest=tag-02FH}{V.{F.CA4} (p.~198)} }
\expandafter\def\csname SRef@maps:eq:F.CA5\endcsname{\href{S6_SERIES_GRIND_V.pdf\#nameddest=tag-02FI}{V.{F.CA5} (p.~198)} }
\expandafter\def\csname SRef@maps:eq:F.CK1\endcsname{\href{S6_SERIES_GRIND_V.pdf\#nameddest=tag-02FJ}{V.{F.CK1} (p.~200)} }
\expandafter\def\csname SRef@maps:eq:F.CK2\endcsname{\href{S6_SERIES_GRIND_V.pdf\#nameddest=tag-02FK}{V.{F.CK2} (p.~201)} }
\expandafter\def\csname SRef@maps:eq:F.GC1\endcsname{\href{S6_SERIES_GRIND_V.pdf\#nameddest=tag-02FZ}{V.{F.GC1} (p.~202)} }
\expandafter\def\csname SRef@maps:eq:F.GC2\endcsname{\href{S6_SERIES_GRIND_V.pdf\#nameddest=tag-02G0}{V.{F.GC2} (p.~203)} }
\expandafter\def\csname SRef@maps:eq:F.GC3\endcsname{\href{S6_SERIES_GRIND_V.pdf\#nameddest=tag-02G1}{V.{F.GC3} (p.~203)} }
\expandafter\def\csname SRef@maps:eq:F.GC4\endcsname{\href{S6_SERIES_GRIND_V.pdf\#nameddest=tag-02G2}{V.{F.GC4} (p.~203)} }
\expandafter\def\csname SRef@maps:eq:F.GC5\endcsname{\href{S6_SERIES_GRIND_V.pdf\#nameddest=tag-02G3}{V.{F.GC5} (p.~204)} }
\expandafter\def\csname SRef@maps:eq:F.GC5a\endcsname{\href{S6_SERIES_GRIND_V.pdf\#nameddest=tag-02G4}{V.{F.GC5a} (p.~204)} }
\expandafter\def\csname SRef@maps:eq:F.GC6\endcsname{\href{S6_SERIES_GRIND_V.pdf\#nameddest=tag-02G5}{V.{F.GC6} (p.~204)} }
\expandafter\def\csname SRef@maps:eq:F.GC7\endcsname{\href{S6_SERIES_GRIND_V.pdf\#nameddest=tag-02G6}{V.{F.GC7} (p.~205)} }
\expandafter\def\csname SRef@maps:eq:F.GC8\endcsname{\href{S6_SERIES_GRIND_V.pdf\#nameddest=tag-02G7}{V.{F.GC8} (p.~205)} }
\expandafter\def\csname SRef@maps:eq:F.GC9\endcsname{\href{S6_SERIES_GRIND_V.pdf\#nameddest=tag-02G8}{V.{F.GC9} (p.~205)} }
\expandafter\def\csname SRef@maps:eq:F.IC1\endcsname{\href{S6_SERIES_GRIND_V.pdf\#nameddest=tag-02GM}{V.{F.IC1} (p.~207)} }
\expandafter\def\csname SRef@maps:eq:F.IC2\endcsname{\href{S6_SERIES_GRIND_V.pdf\#nameddest=tag-02GN}{V.{F.IC2} (p.~207)} }
\expandafter\def\csname SRef@maps:eq:F.IC3\endcsname{\href{S6_SERIES_GRIND_V.pdf\#nameddest=tag-02GO}{V.{F.IC3} (p.~207)} }
\expandafter\def\csname SRef@maps:eq:F.CN1\endcsname{\href{S6_SERIES_GRIND_V.pdf\#nameddest=tag-02FL}{V.{F.CN1} (p.~209)} }
\expandafter\def\csname SRef@maps:eq:F.CN2\endcsname{\href{S6_SERIES_GRIND_V.pdf\#nameddest=tag-02FM}{V.{F.CN2} (p.~209)} }
\expandafter\def\csname SRef@maps:eq:F.CN3\endcsname{\href{S6_SERIES_GRIND_V.pdf\#nameddest=tag-02FN}{V.{F.CN3} (p.~210)} }
\expandafter\def\csname SRef@maps:eq:F.IV1\endcsname{\href{S6_SERIES_GRIND_V.pdf\#nameddest=tag-02GP}{V.{F.IV1} (p.~210)} }
\expandafter\def\csname SRef@maps:eq:F.IV2\endcsname{\href{S6_SERIES_GRIND_V.pdf\#nameddest=tag-02GQ}{V.{F.IV2} (p.~210)} }
\expandafter\def\csname SRef@maps:eq:F.IV3\endcsname{\href{S6_SERIES_GRIND_V.pdf\#nameddest=tag-02GR}{V.{F.IV3} (p.~211)} }
\expandafter\def\csname SRef@maps:eq:F.IV4\endcsname{\href{S6_SERIES_GRIND_V.pdf\#nameddest=tag-02GS}{V.{F.IV4} (p.~211)} }
\expandafter\def\csname SRef@maps:eq:F.SY1\endcsname{\href{S6_SERIES_GRIND_V.pdf\#nameddest=tag-02I6}{V.{F.SY1} (p.~212)} }
\expandafter\def\csname SRef@maps:eq:F.O1\endcsname{\href{S6_SERIES_GRIND_V.pdf\#nameddest=tag-02HC}{V.{F.O1} (p.~216)} }
\expandafter\def\csname SRef@maps:eq:F.O2\endcsname{\href{S6_SERIES_GRIND_V.pdf\#nameddest=tag-02HD}{V.{F.O2} (p.~216)} }
\expandafter\def\csname SRef@maps:eq:F.O3\endcsname{\href{S6_SERIES_GRIND_V.pdf\#nameddest=tag-02HE}{V.{F.O3} (p.~216)} }
\expandafter\def\csname SRef@maps:eq:F.O4\endcsname{\href{S6_SERIES_GRIND_V.pdf\#nameddest=tag-02HF}{V.{F.O4} (p.~217)} }
\expandafter\def\csname SRef@maps:eq:F.O5\endcsname{\href{S6_SERIES_GRIND_V.pdf\#nameddest=tag-02HG}{V.{F.O5} (p.~218)} }
\expandafter\def\csname SRef@maps:eq:F.O6\endcsname{\href{S6_SERIES_GRIND_V.pdf\#nameddest=tag-02HH}{V.{F.O6} (p.~219)} }
\expandafter\def\csname SRef@maps:eq:F.O7\endcsname{\href{S6_SERIES_GRIND_V.pdf\#nameddest=tag-02HI}{V.{F.O7} (p.~220)} }
\expandafter\def\csname SRef@maps:eq:F.7\endcsname{\href{S6_SERIES_GRIND_V.pdf\#nameddest=tag-02FB}{V.{F.7} (p.~104)} }
\expandafter\def\csname SRef@maps:eq:F.8\endcsname{\href{S6_SERIES_GRIND_V.pdf\#nameddest=tag-02FC}{V.{F.8} (p.~104)} }
\expandafter\def\csname SRef@maps:eq:F.9\endcsname{\href{S6_SERIES_GRIND_V.pdf\#nameddest=tag-02FD}{V.{F.9} (p.~105)} }
\expandafter\def\csname SRef@maps:eq:F.10\endcsname{\href{S6_SERIES_GRIND_V.pdf\#nameddest=tag-02EG}{V.{F.10} (p.~105)} }
\expandafter\def\csname SRef@maps:eq:F.11\endcsname{\href{S6_SERIES_GRIND_V.pdf\#nameddest=tag-02EH}{V.{F.11} (p.~106)} }
\expandafter\def\csname SRef@maps:eq:F.12\endcsname{\href{S6_SERIES_GRIND_V.pdf\#nameddest=tag-02EI}{V.{F.12} (p.~106)} }
\expandafter\def\csname SRef@maps:eq:F.13\endcsname{\href{S6_SERIES_GRIND_V.pdf\#nameddest=tag-02EJ}{V.{F.13} (p.~110)} }
\expandafter\def\csname SRef@maps:eq:F.14\endcsname{\href{S6_SERIES_GRIND_V.pdf\#nameddest=tag-02EK}{V.{F.14} (p.~110)} }
\expandafter\def\csname SRef@maps:eq:F.15\endcsname{\href{S6_SERIES_GRIND_V.pdf\#nameddest=tag-02EL}{V.{F.15} (p.~110)} }
\expandafter\def\csname SRef@maps:eq:F.16\endcsname{\href{S6_SERIES_GRIND_V.pdf\#nameddest=tag-02EM}{V.{F.16} (p.~110)} }
\expandafter\def\csname SRef@maps:eq:F.17\endcsname{\href{S6_SERIES_GRIND_V.pdf\#nameddest=tag-02EN}{V.{F.17} (p.~110)} }
\expandafter\def\csname SRef@maps:eq:F.18\endcsname{\href{S6_SERIES_GRIND_V.pdf\#nameddest=tag-02EO}{V.{F.18} (p.~110)} }
\expandafter\def\csname SRef@maps:eq:F.19\endcsname{\href{S6_SERIES_GRIND_V.pdf\#nameddest=tag-02EP}{V.{F.19} (p.~111)} }
\expandafter\def\csname SRef@maps:eq:F.20\endcsname{\href{S6_SERIES_GRIND_V.pdf\#nameddest=tag-02ER}{V.{F.20} (p.~111)} }
\expandafter\def\csname SRef@maps:eq:F.21\endcsname{\href{S6_SERIES_GRIND_V.pdf\#nameddest=tag-02ES}{V.{F.21} (p.~111)} }
\expandafter\def\csname SRef@maps:eq:F.22\endcsname{\href{S6_SERIES_GRIND_V.pdf\#nameddest=tag-02ET}{V.{F.22} (p.~111)} }
\expandafter\def\csname SRef@maps:eq:F.23\endcsname{\href{S6_SERIES_GRIND_V.pdf\#nameddest=tag-02EU}{V.{F.23} (p.~111)} }
\expandafter\def\csname SRef@maps:eq:F.24\endcsname{\href{S6_SERIES_GRIND_V.pdf\#nameddest=tag-02EV}{V.{F.24} (p.~115)} }
\expandafter\def\csname SRef@maps:eq:F.25\endcsname{\href{S6_SERIES_GRIND_V.pdf\#nameddest=tag-02EW}{V.{F.25} (p.~117)} }
\expandafter\def\csname SRef@maps:eq:F.3\endcsname{\href{S6_SERIES_GRIND_V.pdf\#nameddest=tag-02EZ}{V.{F.3} (p.~73)} }
\expandafter\def\csname SRef@maps:eq:F.27\endcsname{\href{S6_SERIES_GRIND_V.pdf\#nameddest=tag-02EX}{V.{F.27} (p.~147)} }
\expandafter\def\csname SRef@maps:eq:F.28\endcsname{\href{S6_SERIES_GRIND_V.pdf\#nameddest=tag-02EY}{V.{F.28} (p.~147)} }
\expandafter\def\csname SRef@maps:eq:F.4\endcsname{\href{S6_SERIES_GRIND_V.pdf\#nameddest=tag-02F7}{V.{F.4} (p.~75)} }
\expandafter\def\csname SRef@maps:eq:F.30\endcsname{\href{S6_SERIES_GRIND_V.pdf\#nameddest=tag-02F0}{V.{F.30} (p.~150)} }
\expandafter\def\csname SRef@maps:eq:F.31\endcsname{\href{S6_SERIES_GRIND_V.pdf\#nameddest=tag-02F1}{V.{F.31} (p.~150)} }
\expandafter\def\csname SRef@maps:eq:F.32\endcsname{\href{S6_SERIES_GRIND_V.pdf\#nameddest=tag-02F2}{V.{F.32} (p.~151)} }
\expandafter\def\csname SRef@maps:eq:F.33\endcsname{\href{S6_SERIES_GRIND_V.pdf\#nameddest=tag-02F3}{V.{F.33} (p.~152)} }
\expandafter\def\csname SRef@maps:eq:F.34\endcsname{\href{S6_SERIES_GRIND_V.pdf\#nameddest=tag-02F4}{V.{F.34} (p.~152)} }
\expandafter\def\csname SRef@maps:eq:F.36\endcsname{\href{S6_SERIES_GRIND_V.pdf\#nameddest=tag-02F5}{V.{F.36} (p.~155)} }
\expandafter\def\csname SRef@maps:eq:F.39\endcsname{\href{S6_SERIES_GRIND_V.pdf\#nameddest=tag-02F6}{V.{F.39} (p.~161)} }
\expandafter\def\csname SRef@maps:eq:F.40\endcsname{\href{S6_SERIES_GRIND_V.pdf\#nameddest=tag-02F8}{V.{F.40} (p.~162)} }
\expandafter\def\csname SRef@maps:eq:F.U1\endcsname{\href{S6_SERIES_GRIND_V.pdf\#nameddest=tag-02IK}{V.{F.U1} (p.~163)} }
\expandafter\def\csname SRef@maps:eq:F.5\endcsname{\href{S6_SERIES_GRIND_V.pdf\#nameddest=tag-02F9}{V.{F.5} (p.~76)} }
\expandafter\def\csname SRef@maps:eq:F.U2\endcsname{\href{S6_SERIES_GRIND_V.pdf\#nameddest=tag-02IO}{V.{F.U2} (p.~163)} }
\expandafter\def\csname SRef@maps:eq:F.U3\endcsname{\href{S6_SERIES_GRIND_V.pdf\#nameddest=tag-02IP}{V.{F.U3} (p.~163)} }
\expandafter\def\csname SRef@maps:eq:F.U4\endcsname{\href{S6_SERIES_GRIND_V.pdf\#nameddest=tag-02IQ}{V.{F.U4} (p.~164)} }
\expandafter\def\csname SRef@maps:eq:F.U5\endcsname{\href{S6_SERIES_GRIND_V.pdf\#nameddest=tag-02IR}{V.{F.U5} (p.~164)} }
\expandafter\def\csname SRef@maps:eq:F.U6\endcsname{\href{S6_SERIES_GRIND_V.pdf\#nameddest=tag-02IS}{V.{F.U6} (p.~164)} }
\expandafter\def\csname SRef@maps:eq:F.U7\endcsname{\href{S6_SERIES_GRIND_V.pdf\#nameddest=tag-02IT}{V.{F.U7} (p.~166)} }
\expandafter\def\csname SRef@maps:eq:F.U9\endcsname{\href{S6_SERIES_GRIND_V.pdf\#nameddest=tag-02IU}{V.{F.U9} (p.~166)} }
\expandafter\def\csname SRef@maps:eq:F.U10\endcsname{\href{S6_SERIES_GRIND_V.pdf\#nameddest=tag-02IL}{V.{F.U10} (p.~166)} }
\expandafter\def\csname SRef@maps:eq:F.U11\endcsname{\href{S6_SERIES_GRIND_V.pdf\#nameddest=tag-02IM}{V.{F.U11} (p.~167)} }
\expandafter\def\csname SRef@maps:eq:F.U12\endcsname{\href{S6_SERIES_GRIND_V.pdf\#nameddest=tag-02IN}{V.{F.U12} (p.~169)} }
\expandafter\def\csname SRef@maps:eq:FT.1\endcsname{\href{S6_SERIES_GRIND_V.pdf\#nameddest=tag-02J0}{V.{FT.1} (p.~81)} }
\expandafter\def\csname SRef@maps:eq:F.R1\endcsname{\href{S6_SERIES_GRIND_V.pdf\#nameddest=tag-02HJ}{V.{F.R1} (p.~170)} }
\expandafter\def\csname SRef@maps:eq:F.R2\endcsname{\href{S6_SERIES_GRIND_V.pdf\#nameddest=tag-02HP}{V.{F.R2} (p.~170)} }
\expandafter\def\csname SRef@maps:eq:F.R3\endcsname{\href{S6_SERIES_GRIND_V.pdf\#nameddest=tag-02HQ}{V.{F.R3} (p.~171)} }
\expandafter\def\csname SRef@maps:eq:F.R4\endcsname{\href{S6_SERIES_GRIND_V.pdf\#nameddest=tag-02HR}{V.{F.R4} (p.~171)} }
\expandafter\def\csname SRef@maps:eq:F.R5\endcsname{\href{S6_SERIES_GRIND_V.pdf\#nameddest=tag-02HS}{V.{F.R5} (p.~171)} }
\expandafter\def\csname SRef@maps:eq:F.R6\endcsname{\href{S6_SERIES_GRIND_V.pdf\#nameddest=tag-02HT}{V.{F.R6} (p.~171)} }
\expandafter\def\csname SRef@maps:eq:F.R7\endcsname{\href{S6_SERIES_GRIND_V.pdf\#nameddest=tag-02HU}{V.{F.R7} (p.~172)} }
\expandafter\def\csname SRef@maps:eq:F.R8\endcsname{\href{S6_SERIES_GRIND_V.pdf\#nameddest=tag-02HV}{V.{F.R8} (p.~172)} }
\expandafter\def\csname SRef@maps:eq:F.R9\endcsname{\href{S6_SERIES_GRIND_V.pdf\#nameddest=tag-02HW}{V.{F.R9} (p.~172)} }
\expandafter\def\csname SRef@maps:eq:F.L1\endcsname{\href{S6_SERIES_GRIND_V.pdf\#nameddest=tag-02GT}{V.{F.L1} (p.~174)} }
\expandafter\def\csname SRef@maps:mapsF:repair:cyclic-inclusion\endcsname{\href{S6_SERIES_GRIND_V.pdf\#nameddest=tag-02J1}{V.{{$\langle 1\rangle $6.}} (p.~199)} }
\expandafter\def\csname SRef@maps:mapsF:repair:log-cylinder-scope\endcsname{\href{S6_SERIES_GRIND_V.pdf\#nameddest=tag-02J2}{V.{{$\langle 1\rangle $5.}} (p.~201)} }
\expandafter\def\csname SRef@sec:fourfold\endcsname{\href{S6_SERIES_GRIND_V.pdf\#nameddest=tag-034O}{V.5 (p.~80)} }
\expandafter\def\csname SRef@maps:src:F.2.21\endcsname{\href{S6_SERIES_GRIND_V.pdf\#nameddest=tag-02NB}{V.11.15 (p.~173)} }
\expandafter\def\csname SRef@maps:src:F.2.30\endcsname{\href{S6_SERIES_GRIND_V.pdf\#nameddest=tag-02NL}{V.11.24 (p.~190)} }
\expandafter\def\csname SRef@maps:eq:F.1\endcsname{\href{S6_SERIES_GRIND_V.pdf\#nameddest=tag-02EF}{V.1 (p.~72)} }
\expandafter\def\csname SRef@maps:eq:F.2\endcsname{\href{S6_SERIES_GRIND_V.pdf\#nameddest=tag-02EQ}{V.2 (p.~72)} }
\expandafter\def\csname SRef@eq:P1Adegree\endcsname{\href{S6_SERIES_GRIND_V.pdf\#nameddest=tag-01O3}{V.3 (p.~134)} }
\expandafter\def\csname SRef@eq:P1Alogcotangent\endcsname{\href{S6_SERIES_GRIND_V.pdf\#nameddest=tag-01O4}{V.4 (p.~135)} }
\expandafter\def\csname SRef@maps:src:F.3.10\endcsname{\href{S6_SERIES_GRIND_V.pdf\#nameddest=tag-02O2}{V.12.5 (p.~218)} }
\expandafter\def\csname SRef@mapsaudit:order-four-torsion-transgression\endcsname{\href{S6_SERIES_GRIND_V.pdf\#nameddest=tag-02QF}{V.5.4 (p.~81)} }
\expandafter\def\csname SRef@maps:src:F.2.17\endcsname{\href{S6_SERIES_GRIND_V.pdf\#nameddest=tag-02N6}{V.11.11 (p.~166)} }
\expandafter\def\csname SRef@maps:src:F.2.8\endcsname{\href{S6_SERIES_GRIND_V.pdf\#nameddest=tag-02NZ}{V.11.2 (p.~147)} }
\expandafter\def\csname SRef@maps:src:F.2.9\endcsname{\href{S6_SERIES_GRIND_V.pdf\#nameddest=tag-02O0}{V.11.3 (p.~148)} }
\expandafter\def\csname SRef@maps:src:F.2.11\endcsname{\href{S6_SERIES_GRIND_V.pdf\#nameddest=tag-02N0}{V.11.5 (p.~151)} }
\expandafter\def\csname SRef@maps:src:F.2.15\endcsname{\href{S6_SERIES_GRIND_V.pdf\#nameddest=tag-02N4}{V.11.9 (p.~162)} }
\expandafter\def\csname SRef@maps:src:F.3.3\endcsname{\href{S6_SERIES_GRIND_V.pdf\#nameddest=tag-02O5}{V.12.1 (p.~215)} }
\expandafter\def\csname SRef@maps:src:F.1.4\endcsname{\href{S6_SERIES_GRIND_V.pdf\#nameddest=tag-02MX}{V.4.2 (p.~72)} }
\expandafter\def\csname SRef@maps:src:F.2.1\endcsname{\href{S6_SERIES_GRIND_V.pdf\#nameddest=tag-02MY}{V.4.3 (p.~73)} }
\expandafter\def\csname SRef@maps:src:F.2.2\endcsname{\href{S6_SERIES_GRIND_V.pdf\#nameddest=tag-02N9}{V.4.4 (p.~73)} }
\expandafter\def\csname SRef@maps:src:F.2.3\endcsname{\href{S6_SERIES_GRIND_V.pdf\#nameddest=tag-02NK}{V.4.5 (p.~74)} }
\expandafter\def\csname SRef@maps:src:F.3.2\endcsname{\href{S6_SERIES_GRIND_V.pdf\#nameddest=tag-02O4}{V.4.7 (p.~75)} }
\expandafter\def\csname SRef@maps:src:F.2.4\endcsname{\href{S6_SERIES_GRIND_V.pdf\#nameddest=tag-02NV}{V.4.8 (p.~75)} }
\expandafter\def\csname SRef@sec:V-known-inputs\endcsname{\href{S6_SERIES_GRIND_V.pdf\#nameddest=tag-034M}{V.3 (p.~72)} }
\expandafter\def\csname SRef@maps:src:F.2.7\endcsname{\href{S6_SERIES_GRIND_V.pdf\#nameddest=tag-02NY}{V.11.1 (p.~146)} }
\expandafter\def\csname SRef@maps:src:F.2.16\endcsname{\href{S6_SERIES_GRIND_V.pdf\#nameddest=tag-02N5}{V.11.10 (p.~163)} }
\expandafter\def\csname SRef@maps:src:F.2.18\endcsname{\href{S6_SERIES_GRIND_V.pdf\#nameddest=tag-02N7}{V.11.12 (p.~167)} }
\expandafter\def\csname SRef@maps:src:F.2.19\endcsname{\href{S6_SERIES_GRIND_V.pdf\#nameddest=tag-02N8}{V.11.13 (p.~169)} }
\expandafter\def\csname SRef@maps:src:F.2.20\endcsname{\href{S6_SERIES_GRIND_V.pdf\#nameddest=tag-02NA}{V.11.14 (p.~170)} }
\expandafter\def\csname SRef@maps:src:F.2.22\endcsname{\href{S6_SERIES_GRIND_V.pdf\#nameddest=tag-02NC}{V.11.16 (p.~174)} }
\expandafter\def\csname SRef@maps:src:F.2.23\endcsname{\href{S6_SERIES_GRIND_V.pdf\#nameddest=tag-02ND}{V.11.17 (p.~176)} }
\expandafter\def\csname SRef@maps:src:F.2.24\endcsname{\href{S6_SERIES_GRIND_V.pdf\#nameddest=tag-02NE}{V.11.18 (p.~177)} }
\expandafter\def\csname SRef@maps:src:F.2.25\endcsname{\href{S6_SERIES_GRIND_V.pdf\#nameddest=tag-02NF}{V.11.19 (p.~179)} }
\expandafter\def\csname SRef@maps:src:F.2.26\endcsname{\href{S6_SERIES_GRIND_V.pdf\#nameddest=tag-02NG}{V.11.20 (p.~181)} }
\expandafter\def\csname SRef@maps:src:F.2.27\endcsname{\href{S6_SERIES_GRIND_V.pdf\#nameddest=tag-02NH}{V.11.21 (p.~183)} }
\expandafter\def\csname SRef@maps:src:F.2.28\endcsname{\href{S6_SERIES_GRIND_V.pdf\#nameddest=tag-02NI}{V.11.22 (p.~185)} }
\expandafter\def\csname SRef@maps:src:F.2.29\endcsname{\href{S6_SERIES_GRIND_V.pdf\#nameddest=tag-02NJ}{V.11.23 (p.~186)} }
\expandafter\def\csname SRef@maps:src:F.2.31\endcsname{\href{S6_SERIES_GRIND_V.pdf\#nameddest=tag-02NM}{V.11.25 (p.~191)} }
\expandafter\def\csname SRef@maps:src:F.2.32\endcsname{\href{S6_SERIES_GRIND_V.pdf\#nameddest=tag-02NN}{V.11.26 (p.~195)} }
\expandafter\def\csname SRef@maps:src:F.2.33\endcsname{\href{S6_SERIES_GRIND_V.pdf\#nameddest=tag-02NO}{V.11.27 (p.~197)} }
\expandafter\def\csname SRef@maps:src:F.2.34\endcsname{\href{S6_SERIES_GRIND_V.pdf\#nameddest=tag-02NP}{V.11.28 (p.~200)} }
\expandafter\def\csname SRef@maps:src:F.2.35\endcsname{\href{S6_SERIES_GRIND_V.pdf\#nameddest=tag-02NQ}{V.11.29 (p.~202)} }
\expandafter\def\csname SRef@maps:src:F.2.36\endcsname{\href{S6_SERIES_GRIND_V.pdf\#nameddest=tag-02NR}{V.11.30 (p.~206)} }
\expandafter\def\csname SRef@maps:src:F.2.37\endcsname{\href{S6_SERIES_GRIND_V.pdf\#nameddest=tag-02NS}{V.11.31 (p.~208)} }
\expandafter\def\csname SRef@maps:src:F.2.38\endcsname{\href{S6_SERIES_GRIND_V.pdf\#nameddest=tag-02NT}{V.11.32 (p.~210)} }
\expandafter\def\csname SRef@maps:src:F.2.39\endcsname{\href{S6_SERIES_GRIND_V.pdf\#nameddest=tag-02NU}{V.11.33 (p.~212)} }
\expandafter\def\csname SRef@night:P1B.1\endcsname{\href{S6_SERIES_GRIND_V.pdf\#nameddest=tag-02R7}{V.11.34 (p.~213)} }
\expandafter\def\csname SRef@night:P1B.2\endcsname{\href{S6_SERIES_GRIND_V.pdf\#nameddest=tag-02R8}{V.11.35 (p.~214)} }
\expandafter\def\csname SRef@maps:src:F.2.10\endcsname{\href{S6_SERIES_GRIND_V.pdf\#nameddest=tag-02MZ}{V.11.4 (p.~150)} }
\expandafter\def\csname SRef@maps:src:F.2.12\endcsname{\href{S6_SERIES_GRIND_V.pdf\#nameddest=tag-02N1}{V.11.6 (p.~152)} }
\expandafter\def\csname SRef@maps:src:F.2.13\endcsname{\href{S6_SERIES_GRIND_V.pdf\#nameddest=tag-02N2}{V.11.7 (p.~155)} }
\expandafter\def\csname SRef@maps:src:F.2.14\endcsname{\href{S6_SERIES_GRIND_V.pdf\#nameddest=tag-02N3}{V.11.8 (p.~161)} }
\expandafter\def\csname SRef@maps:src:F.3.4\endcsname{\href{S6_SERIES_GRIND_V.pdf\#nameddest=tag-02O6}{V.12.2 (p.~215)} }
\expandafter\def\csname SRef@maps:src:F.3.8\endcsname{\href{S6_SERIES_GRIND_V.pdf\#nameddest=tag-02OA}{V.12.3 (p.~216)} }
\expandafter\def\csname SRef@maps:src:F.3.9\endcsname{\href{S6_SERIES_GRIND_V.pdf\#nameddest=tag-02OB}{V.12.4 (p.~217)} }
\expandafter\def\csname SRef@maps:src:F.3.11\endcsname{\href{S6_SERIES_GRIND_V.pdf\#nameddest=tag-02O3}{V.12.6 (p.~219)} }
\expandafter\def\csname SRef@maps:src:F.1.1\endcsname{\href{S6_SERIES_GRIND_V.pdf\#nameddest=tag-02MU}{V.4.1 (p.~72)} }
\expandafter\def\csname SRef@maps:src:F.3.1\endcsname{\href{S6_SERIES_GRIND_V.pdf\#nameddest=tag-02O1}{V.4.6 (p.~74)} }
\expandafter\def\csname SRef@maps:src:F.2.5\endcsname{\href{S6_SERIES_GRIND_V.pdf\#nameddest=tag-02NW}{V.5.1 (p.~76)} }
\expandafter\def\csname SRef@night:P1.1\endcsname{\href{S6_SERIES_GRIND_V.pdf\#nameddest=tag-02QI}{V.5.2 (p.~80)} }
\expandafter\def\csname SRef@night:P1.3\endcsname{\href{S6_SERIES_GRIND_V.pdf\#nameddest=tag-02QU}{V.5.3 (p.~80)} }
\expandafter\def\csname SRef@maps:src:F.2.6\endcsname{\href{S6_SERIES_GRIND_V.pdf\#nameddest=tag-02NX}{V.5.5 (p.~82)} }
\expandafter\def\csname SRef@night:P1.10\endcsname{\href{S6_SERIES_GRIND_V.pdf\#nameddest=tag-02QJ}{V.5.6 (p.~86)} }
\expandafter\def\csname SRef@night:P1.11\endcsname{\href{S6_SERIES_GRIND_V.pdf\#nameddest=tag-02QK}{V.5.7 (p.~87)} }
\expandafter\def\csname SRef@night:P1.12\endcsname{\href{S6_SERIES_GRIND_V.pdf\#nameddest=tag-02QL}{V.5.8 (p.~91)} }
\expandafter\def\csname SRef@night:P1.19\endcsname{\href{S6_SERIES_GRIND_V.pdf\#nameddest=tag-02QS}{V.5.9 (p.~94)} }
\expandafter\def\csname SRef@night:P1.4\endcsname{\href{S6_SERIES_GRIND_V.pdf\#nameddest=tag-02QV}{V.6.1 (p.~103)} }
\expandafter\def\csname SRef@night:P1.5\endcsname{\href{S6_SERIES_GRIND_V.pdf\#nameddest=tag-02QW}{V.6.2 (p.~104)} }
\expandafter\def\csname SRef@maps:src:F.3.5\endcsname{\href{S6_SERIES_GRIND_V.pdf\#nameddest=tag-02O7}{V.6.3 (p.~104)} }
\expandafter\def\csname SRef@night:P1.2\endcsname{\href{S6_SERIES_GRIND_V.pdf\#nameddest=tag-02QT}{V.6.4 (p.~106)} }
\expandafter\def\csname SRef@night:P1.7\endcsname{\href{S6_SERIES_GRIND_V.pdf\#nameddest=tag-02QY}{V.6.5 (p.~108)} }
\expandafter\def\csname SRef@maps:src:F.3.6\endcsname{\href{S6_SERIES_GRIND_V.pdf\#nameddest=tag-02O8}{V.6.6 (p.~109)} }
\expandafter\def\csname SRef@night:P1.8\endcsname{\href{S6_SERIES_GRIND_V.pdf\#nameddest=tag-02QZ}{V.6.7 (p.~112)} }
\expandafter\def\csname SRef@night:P1.9\endcsname{\href{S6_SERIES_GRIND_V.pdf\#nameddest=tag-02R0}{V.6.8 (p.~113)} }
\expandafter\def\csname SRef@maps:src:F.3.7\endcsname{\href{S6_SERIES_GRIND_V.pdf\#nameddest=tag-02O9}{V.6.9 (p.~114)} }
\expandafter\def\csname SRef@night:P1.13\endcsname{\href{S6_SERIES_GRIND_V.pdf\#nameddest=tag-02QM}{V.7.1 (p.~116)} }
\expandafter\def\csname SRef@maps:src:F.4.1\endcsname{\href{S6_SERIES_GRIND_V.pdf\#nameddest=tag-02OC}{V.7.2 (p.~117)} }
\expandafter\def\csname SRef@night:P1.15\endcsname{\href{S6_SERIES_GRIND_V.pdf\#nameddest=tag-02QO}{V.8.1 (p.~119)} }
\expandafter\def\csname SRef@night:P1.16\endcsname{\href{S6_SERIES_GRIND_V.pdf\#nameddest=tag-02QP}{V.8.2 (p.~121)} }
\expandafter\def\csname SRef@night:P1.18\endcsname{\href{S6_SERIES_GRIND_V.pdf\#nameddest=tag-02QR}{V.8.3 (p.~124)} }
\expandafter\def\csname SRef@night:P1A.1\endcsname{\href{S6_SERIES_GRIND_V.pdf\#nameddest=tag-02R1}{V.9.1 (p.~134)} }
\expandafter\def\csname SRef@night:P1A.2\endcsname{\href{S6_SERIES_GRIND_V.pdf\#nameddest=tag-02R2}{V.9.2 (p.~135)} }
\expandafter\def\csname SRef@night:P1A.3\endcsname{\href{S6_SERIES_GRIND_V.pdf\#nameddest=tag-02R3}{V.9.3 (p.~137)} }
\expandafter\def\csname SRef@night:P1A.4\endcsname{\href{S6_SERIES_GRIND_V.pdf\#nameddest=tag-02R4}{V.9.4 (p.~139)} }
\expandafter\def\csname SRef@night:P1.17\endcsname{\href{S6_SERIES_GRIND_V.pdf\#nameddest=tag-02QQ}{V.9.5 (p.~141)} }
\expandafter\def\csname SRef@night:P1A.5\endcsname{\href{S6_SERIES_GRIND_V.pdf\#nameddest=tag-02R5}{V.10.1 (p.~143)} }
\expandafter\def\csname SRef@night:P1A.6\endcsname{\href{S6_SERIES_GRIND_V.pdf\#nameddest=tag-02R6}{V.10.2 (p.~146)} }
\expandafter\def\csname SRef@cdp:eq:C.UN.1\endcsname{\href{S6_SERIES_GRIND_II.pdf\#nameddest=tag-0152}{II.{C.UN.1} (p.~413)} }
\expandafter\def\csname SRef@cdp:eq:C.FL.1\endcsname{\href{S6_SERIES_GRIND_II.pdf\#nameddest=tag-00KR}{II.{C.FL.1} (p.~443)} }
\expandafter\def\csname SRef@cdp:eq:C.QT.11\endcsname{\href{S6_SERIES_GRIND_II.pdf\#nameddest=tag-00WF}{II.{C.QT.11} (p.~708)} }
\expandafter\def\csname SRef@cdp:eq:C.QT.12\endcsname{\href{S6_SERIES_GRIND_II.pdf\#nameddest=tag-00WG}{II.{C.QT.12} (p.~708)} }
\expandafter\def\csname SRef@cdp:eq:C.QT.13\endcsname{\href{S6_SERIES_GRIND_II.pdf\#nameddest=tag-00WH}{II.{C.QT.13} (p.~708)} }
\expandafter\def\csname SRef@cdp:eq:C.WT.1\endcsname{\href{S6_SERIES_GRIND_II.pdf\#nameddest=tag-017O}{II.{C.WT.1} (p.~709)} }
\expandafter\def\csname SRef@cdp:eq:C.WT.2\endcsname{\href{S6_SERIES_GRIND_II.pdf\#nameddest=tag-017Z}{II.{C.WT.2} (p.~709)} }
\expandafter\def\csname SRef@cdp:eq:C.WT.3\endcsname{\href{S6_SERIES_GRIND_II.pdf\#nameddest=tag-018B}{II.{C.WT.3} (p.~709)} }
\expandafter\def\csname SRef@cdp:eq:C.WT.4\endcsname{\href{S6_SERIES_GRIND_II.pdf\#nameddest=tag-018M}{II.{C.WT.4} (p.~709)} }
\expandafter\def\csname SRef@cdp:eq:C.WT.5\endcsname{\href{S6_SERIES_GRIND_II.pdf\#nameddest=tag-018R}{II.{C.WT.5} (p.~709)} }
\expandafter\def\csname SRef@cdp:eq:C.WT.6\endcsname{\href{S6_SERIES_GRIND_II.pdf\#nameddest=tag-018S}{II.{C.WT.6} (p.~710)} }
\expandafter\def\csname SRef@cdp:eq:C.WT.7\endcsname{\href{S6_SERIES_GRIND_II.pdf\#nameddest=tag-018T}{II.{C.WT.7} (p.~710)} }
\expandafter\def\csname SRef@cdp:eq:C.FL.2\endcsname{\href{S6_SERIES_GRIND_II.pdf\#nameddest=tag-00KS}{II.{C.FL.2} (p.~443)} }
\expandafter\def\csname SRef@cdp:eq:C.WT.8\endcsname{\href{S6_SERIES_GRIND_II.pdf\#nameddest=tag-018U}{II.{C.WT.8} (p.~710)} }
\expandafter\def\csname SRef@cdp:eq:C.WT.9\endcsname{\href{S6_SERIES_GRIND_II.pdf\#nameddest=tag-018V}{II.{C.WT.9} (p.~710)} }
\expandafter\def\csname SRef@cdp:eq:C.WT.10\endcsname{\href{S6_SERIES_GRIND_II.pdf\#nameddest=tag-017P}{II.{C.WT.10} (p.~711)} }
\expandafter\def\csname SRef@cdp:eq:C.WT.11\endcsname{\href{S6_SERIES_GRIND_II.pdf\#nameddest=tag-017Q}{II.{C.WT.11} (p.~711)} }
\expandafter\def\csname SRef@cdp:eq:C.WT.12\endcsname{\href{S6_SERIES_GRIND_II.pdf\#nameddest=tag-017R}{II.{C.WT.12} (p.~711)} }
\expandafter\def\csname SRef@cdp:eq:C.WT.13\endcsname{\href{S6_SERIES_GRIND_II.pdf\#nameddest=tag-017S}{II.{C.WT.13} (p.~711)} }
\expandafter\def\csname SRef@cdp:eq:C.WT.14\endcsname{\href{S6_SERIES_GRIND_II.pdf\#nameddest=tag-017T}{II.{C.WT.14} (p.~712)} }
\expandafter\def\csname SRef@cdp:eq:C.WT.15\endcsname{\href{S6_SERIES_GRIND_II.pdf\#nameddest=tag-017U}{II.{C.WT.15} (p.~712)} }
\expandafter\def\csname SRef@cdp:eq:C.WT.16\endcsname{\href{S6_SERIES_GRIND_II.pdf\#nameddest=tag-017V}{II.{C.WT.16} (p.~712)} }
\expandafter\def\csname SRef@cdp:eq:C.WT.17\endcsname{\href{S6_SERIES_GRIND_II.pdf\#nameddest=tag-017W}{II.{C.WT.17} (p.~713)} }
\expandafter\def\csname SRef@cdp:eq:C.FL.3\endcsname{\href{S6_SERIES_GRIND_II.pdf\#nameddest=tag-00KT}{II.{C.FL.3} (p.~443)} }
\expandafter\def\csname SRef@cdp:eq:C.WT.18\endcsname{\href{S6_SERIES_GRIND_II.pdf\#nameddest=tag-017X}{II.{C.WT.18} (p.~713)} }
\expandafter\def\csname SRef@cdp:eq:C.WT.19\endcsname{\href{S6_SERIES_GRIND_II.pdf\#nameddest=tag-017Y}{II.{C.WT.19} (p.~713)} }
\expandafter\def\csname SRef@cdp:eq:C.WT.20\endcsname{\href{S6_SERIES_GRIND_II.pdf\#nameddest=tag-0180}{II.{C.WT.20} (p.~714)} }
\expandafter\def\csname SRef@cdp:eq:C.WT.21\endcsname{\href{S6_SERIES_GRIND_II.pdf\#nameddest=tag-0181}{II.{C.WT.21} (p.~714)} }
\expandafter\def\csname SRef@cdp:eq:C.WT.22\endcsname{\href{S6_SERIES_GRIND_II.pdf\#nameddest=tag-0182}{II.{C.WT.22} (p.~714)} }
\expandafter\def\csname SRef@cdp:eq:C.WT.23\endcsname{\href{S6_SERIES_GRIND_II.pdf\#nameddest=tag-0183}{II.{C.WT.23} (p.~714)} }
\expandafter\def\csname SRef@cdp:eq:C.WT.23a\endcsname{\href{S6_SERIES_GRIND_II.pdf\#nameddest=tag-0184}{II.{C.WT.23a} (p.~714)} }
\expandafter\def\csname SRef@cdp:eq:C.WT.24\endcsname{\href{S6_SERIES_GRIND_II.pdf\#nameddest=tag-0185}{II.{C.WT.24} (p.~714)} }
\expandafter\def\csname SRef@cdp:eq:C.WT.25\endcsname{\href{S6_SERIES_GRIND_II.pdf\#nameddest=tag-0186}{II.{C.WT.25} (p.~715)} }
\expandafter\def\csname SRef@cdp:eq:C.WT.26\endcsname{\href{S6_SERIES_GRIND_II.pdf\#nameddest=tag-0187}{II.{C.WT.26} (p.~715)} }
\expandafter\def\csname SRef@cdp:eq:C.FL.4\endcsname{\href{S6_SERIES_GRIND_II.pdf\#nameddest=tag-00KU}{II.{C.FL.4} (p.~443)} }
\expandafter\def\csname SRef@cdp:eq:C.WT.27\endcsname{\href{S6_SERIES_GRIND_II.pdf\#nameddest=tag-0188}{II.{C.WT.27} (p.~715)} }
\expandafter\def\csname SRef@cdp:eq:C.WT.28\endcsname{\href{S6_SERIES_GRIND_II.pdf\#nameddest=tag-0189}{II.{C.WT.28} (p.~715)} }
\expandafter\def\csname SRef@cdp:eq:C.WT.29\endcsname{\href{S6_SERIES_GRIND_II.pdf\#nameddest=tag-018A}{II.{C.WT.29} (p.~715)} }
\expandafter\def\csname SRef@cdp:eq:C.WT.30\endcsname{\href{S6_SERIES_GRIND_II.pdf\#nameddest=tag-018C}{II.{C.WT.30} (p.~715)} }
\expandafter\def\csname SRef@cdp:eq:C.WT.31\endcsname{\href{S6_SERIES_GRIND_II.pdf\#nameddest=tag-018D}{II.{C.WT.31} (p.~715)} }
\expandafter\def\csname SRef@cdp:eq:C.WT.32\endcsname{\href{S6_SERIES_GRIND_II.pdf\#nameddest=tag-018E}{II.{C.WT.32} (p.~716)} }
\expandafter\def\csname SRef@cdp:eq:C.WT.33\endcsname{\href{S6_SERIES_GRIND_II.pdf\#nameddest=tag-018F}{II.{C.WT.33} (p.~716)} }
\expandafter\def\csname SRef@cdp:eq:C.WT.34\endcsname{\href{S6_SERIES_GRIND_II.pdf\#nameddest=tag-018G}{II.{C.WT.34} (p.~716)} }
\expandafter\def\csname SRef@cdp:eq:C.WT.35\endcsname{\href{S6_SERIES_GRIND_II.pdf\#nameddest=tag-018H}{II.{C.WT.35} (p.~716)} }
\expandafter\def\csname SRef@cdp:eq:C.WT.36\endcsname{\href{S6_SERIES_GRIND_II.pdf\#nameddest=tag-018I}{II.{C.WT.36} (p.~716)} }
\expandafter\def\csname SRef@cdp:eq:C.FL.5\endcsname{\href{S6_SERIES_GRIND_II.pdf\#nameddest=tag-00KV}{II.{C.FL.5} (p.~444)} }
\expandafter\def\csname SRef@cdp:eq:C.WT.37\endcsname{\href{S6_SERIES_GRIND_II.pdf\#nameddest=tag-018J}{II.{C.WT.37} (p.~717)} }
\expandafter\def\csname SRef@cdp:eq:C.WT.38\endcsname{\href{S6_SERIES_GRIND_II.pdf\#nameddest=tag-018K}{II.{C.WT.38} (p.~717)} }
\expandafter\def\csname SRef@cdp:eq:C.WT.39\endcsname{\href{S6_SERIES_GRIND_II.pdf\#nameddest=tag-018L}{II.{C.WT.39} (p.~717)} }
\expandafter\def\csname SRef@cdp:eq:C.WT.40\endcsname{\href{S6_SERIES_GRIND_II.pdf\#nameddest=tag-018N}{II.{C.WT.40} (p.~717)} }
\expandafter\def\csname SRef@cdp:eq:C.WT.41\endcsname{\href{S6_SERIES_GRIND_II.pdf\#nameddest=tag-018O}{II.{C.WT.41} (p.~717)} }
\expandafter\def\csname SRef@cdp:eq:C.WT.42\endcsname{\href{S6_SERIES_GRIND_II.pdf\#nameddest=tag-018P}{II.{C.WT.42} (p.~718)} }
\expandafter\def\csname SRef@cdp:eq:C.WT.43\endcsname{\href{S6_SERIES_GRIND_II.pdf\#nameddest=tag-018Q}{II.{C.WT.43} (p.~718)} }
\expandafter\def\csname SRef@cdp:eq:C.NM.1\endcsname{\href{S6_SERIES_GRIND_II.pdf\#nameddest=tag-00SR}{II.{C.NM.1} (p.~719)} }
\expandafter\def\csname SRef@cdp:eq:C.NM.2\endcsname{\href{S6_SERIES_GRIND_II.pdf\#nameddest=tag-00ST}{II.{C.NM.2} (p.~719)} }
\expandafter\def\csname SRef@cdp:eq:C.NM.3\endcsname{\href{S6_SERIES_GRIND_II.pdf\#nameddest=tag-00SU}{II.{C.NM.3} (p.~719)} }
\expandafter\def\csname SRef@cdp:eq:C.PG.1\endcsname{\href{S6_SERIES_GRIND_II.pdf\#nameddest=tag-00UZ}{II.{C.PG.1} (p.~444)} }
\expandafter\def\csname SRef@cdp:eq:C.NM.4\endcsname{\href{S6_SERIES_GRIND_II.pdf\#nameddest=tag-00SV}{II.{C.NM.4} (p.~719)} }
\expandafter\def\csname SRef@cdp:eq:C.NM.5\endcsname{\href{S6_SERIES_GRIND_II.pdf\#nameddest=tag-00SW}{II.{C.NM.5} (p.~720)} }
\expandafter\def\csname SRef@cdp:eq:C.NM.6\endcsname{\href{S6_SERIES_GRIND_II.pdf\#nameddest=tag-00SX}{II.{C.NM.6} (p.~720)} }
\expandafter\def\csname SRef@cdp:eq:C.NM.7\endcsname{\href{S6_SERIES_GRIND_II.pdf\#nameddest=tag-00SY}{II.{C.NM.7} (p.~720)} }
\expandafter\def\csname SRef@cdp:eq:C.NM.8\endcsname{\href{S6_SERIES_GRIND_II.pdf\#nameddest=tag-00SZ}{II.{C.NM.8} (p.~720)} }
\expandafter\def\csname SRef@cdp:eq:C.NM.9\endcsname{\href{S6_SERIES_GRIND_II.pdf\#nameddest=tag-00T0}{II.{C.NM.9} (p.~721)} }
\expandafter\def\csname SRef@cdp:eq:C.NM.10\endcsname{\href{S6_SERIES_GRIND_II.pdf\#nameddest=tag-00SS}{II.{C.NM.10} (p.~721)} }
\expandafter\def\csname SRef@cdp:eq:C.NP.1\endcsname{\href{S6_SERIES_GRIND_II.pdf\#nameddest=tag-00T1}{II.{C.NP.1} (p.~722)} }
\expandafter\def\csname SRef@cdp:eq:C.NP.2\endcsname{\href{S6_SERIES_GRIND_II.pdf\#nameddest=tag-00T2}{II.{C.NP.2} (p.~722)} }
\expandafter\def\csname SRef@cdp:eq:C.NP.3\endcsname{\href{S6_SERIES_GRIND_II.pdf\#nameddest=tag-00T3}{II.{C.NP.3} (p.~723)} }
\expandafter\def\csname SRef@cdp:eq:C.PG.2\endcsname{\href{S6_SERIES_GRIND_II.pdf\#nameddest=tag-00V0}{II.{C.PG.2} (p.~444)} }
\expandafter\def\csname SRef@cdp:eq:C.NP.4\endcsname{\href{S6_SERIES_GRIND_II.pdf\#nameddest=tag-00T4}{II.{C.NP.4} (p.~723)} }
\expandafter\def\csname SRef@cdp:eq:C.NP.5\endcsname{\href{S6_SERIES_GRIND_II.pdf\#nameddest=tag-00T5}{II.{C.NP.5} (p.~723)} }
\expandafter\def\csname SRef@cdp:eq:C.NP.6\endcsname{\href{S6_SERIES_GRIND_II.pdf\#nameddest=tag-00T6}{II.{C.NP.6} (p.~724)} }
\expandafter\def\csname SRef@cdp:eq:C.NP.7\endcsname{\href{S6_SERIES_GRIND_II.pdf\#nameddest=tag-00T7}{II.{C.NP.7} (p.~724)} }
\expandafter\def\csname SRef@cdp:eq:C.NP.8\endcsname{\href{S6_SERIES_GRIND_II.pdf\#nameddest=tag-00T8}{II.{C.NP.8} (p.~725)} }
\expandafter\def\csname SRef@cdp:eq:C.NP.9\endcsname{\href{S6_SERIES_GRIND_II.pdf\#nameddest=tag-00T9}{II.{C.NP.9} (p.~725)} }
\expandafter\def\csname SRef@cdp:eq:C.TG.1\endcsname{\href{S6_SERIES_GRIND_II.pdf\#nameddest=tag-012Q}{II.{C.TG.1} (p.~726)} }
\expandafter\def\csname SRef@cdp:eq:C.TG.2\endcsname{\href{S6_SERIES_GRIND_II.pdf\#nameddest=tag-012X}{II.{C.TG.2} (p.~726)} }
\expandafter\def\csname SRef@cdp:eq:C.TG.3\endcsname{\href{S6_SERIES_GRIND_II.pdf\#nameddest=tag-012Y}{II.{C.TG.3} (p.~726)} }
\expandafter\def\csname SRef@cdp:eq:C.TG.4\endcsname{\href{S6_SERIES_GRIND_II.pdf\#nameddest=tag-012Z}{II.{C.TG.4} (p.~726)} }
\expandafter\def\csname SRef@cdp:eq:C.PG.3\endcsname{\href{S6_SERIES_GRIND_II.pdf\#nameddest=tag-00V1}{II.{C.PG.3} (p.~445)} }
\expandafter\def\csname SRef@cdp:eq:C.TG.5\endcsname{\href{S6_SERIES_GRIND_II.pdf\#nameddest=tag-0130}{II.{C.TG.5} (p.~726)} }
\expandafter\def\csname SRef@cdp:eq:C.TG.6\endcsname{\href{S6_SERIES_GRIND_II.pdf\#nameddest=tag-0131}{II.{C.TG.6} (p.~726)} }
\expandafter\def\csname SRef@cdp:eq:C.TG.7\endcsname{\href{S6_SERIES_GRIND_II.pdf\#nameddest=tag-0132}{II.{C.TG.7} (p.~727)} }
\expandafter\def\csname SRef@cdp:eq:C.TG.8\endcsname{\href{S6_SERIES_GRIND_II.pdf\#nameddest=tag-0133}{II.{C.TG.8} (p.~727)} }
\expandafter\def\csname SRef@cdp:eq:C.TG.9\endcsname{\href{S6_SERIES_GRIND_II.pdf\#nameddest=tag-0134}{II.{C.TG.9} (p.~728)} }
\expandafter\def\csname SRef@cdp:eq:C.TG.10\endcsname{\href{S6_SERIES_GRIND_II.pdf\#nameddest=tag-012R}{II.{C.TG.10} (p.~728)} }
\expandafter\def\csname SRef@cdp:eq:C.TG.11\endcsname{\href{S6_SERIES_GRIND_II.pdf\#nameddest=tag-012S}{II.{C.TG.11} (p.~728)} }
\expandafter\def\csname SRef@cdp:eq:C.TG.12\endcsname{\href{S6_SERIES_GRIND_II.pdf\#nameddest=tag-012T}{II.{C.TG.12} (p.~729)} }
\expandafter\def\csname SRef@cdp:eq:C.TG.13\endcsname{\href{S6_SERIES_GRIND_II.pdf\#nameddest=tag-012U}{II.{C.TG.13} (p.~729)} }
\expandafter\def\csname SRef@cdp:eq:C.TG.14\endcsname{\href{S6_SERIES_GRIND_II.pdf\#nameddest=tag-012V}{II.{C.TG.14} (p.~729)} }
\expandafter\def\csname SRef@cdp:eq:C.PG.4\endcsname{\href{S6_SERIES_GRIND_II.pdf\#nameddest=tag-00V2}{II.{C.PG.4} (p.~445)} }
\expandafter\def\csname SRef@cdp:eq:C.TG.15\endcsname{\href{S6_SERIES_GRIND_II.pdf\#nameddest=tag-012W}{II.{C.TG.15} (p.~729)} }
\expandafter\def\csname SRef@cdp:eq:C.TX.1\endcsname{\href{S6_SERIES_GRIND_II.pdf\#nameddest=tag-013Y}{II.{C.TX.1} (p.~730)} }
\expandafter\def\csname SRef@cdp:eq:C.TX.2\endcsname{\href{S6_SERIES_GRIND_II.pdf\#nameddest=tag-0144}{II.{C.TX.2} (p.~730)} }
\expandafter\def\csname SRef@cdp:eq:C.TX.3\endcsname{\href{S6_SERIES_GRIND_II.pdf\#nameddest=tag-0145}{II.{C.TX.3} (p.~730)} }
\expandafter\def\csname SRef@cdp:eq:C.TX.4\endcsname{\href{S6_SERIES_GRIND_II.pdf\#nameddest=tag-0146}{II.{C.TX.4} (p.~730)} }
\expandafter\def\csname SRef@cdp:eq:C.TX.5\endcsname{\href{S6_SERIES_GRIND_II.pdf\#nameddest=tag-0147}{II.{C.TX.5} (p.~731)} }
\expandafter\def\csname SRef@cdp:eq:C.TX.6\endcsname{\href{S6_SERIES_GRIND_II.pdf\#nameddest=tag-0148}{II.{C.TX.6} (p.~731)} }
\expandafter\def\csname SRef@cdp:eq:C.TX.7\endcsname{\href{S6_SERIES_GRIND_II.pdf\#nameddest=tag-0149}{II.{C.TX.7} (p.~731)} }
\expandafter\def\csname SRef@cdp:eq:C.TX.8\endcsname{\href{S6_SERIES_GRIND_II.pdf\#nameddest=tag-014A}{II.{C.TX.8} (p.~731)} }
\expandafter\def\csname SRef@cdp:eq:C.TX.9\endcsname{\href{S6_SERIES_GRIND_II.pdf\#nameddest=tag-014B}{II.{C.TX.9} (p.~731)} }
\expandafter\def\csname SRef@cdp:eq:C.PG.5\endcsname{\href{S6_SERIES_GRIND_II.pdf\#nameddest=tag-00V3}{II.{C.PG.5} (p.~445)} }
\expandafter\def\csname SRef@cdp:eq:C.TX.10\endcsname{\href{S6_SERIES_GRIND_II.pdf\#nameddest=tag-013Z}{II.{C.TX.10} (p.~731)} }
\expandafter\def\csname SRef@cdp:eq:C.TX.11\endcsname{\href{S6_SERIES_GRIND_II.pdf\#nameddest=tag-0140}{II.{C.TX.11} (p.~732)} }
\expandafter\def\csname SRef@cdp:eq:C.TX.12\endcsname{\href{S6_SERIES_GRIND_II.pdf\#nameddest=tag-0141}{II.{C.TX.12} (p.~732)} }
\expandafter\def\csname SRef@cdp:eq:C.TX.13\endcsname{\href{S6_SERIES_GRIND_II.pdf\#nameddest=tag-0142}{II.{C.TX.13} (p.~732)} }
\expandafter\def\csname SRef@cdp:eq:C.TX.14\endcsname{\href{S6_SERIES_GRIND_II.pdf\#nameddest=tag-0143}{II.{C.TX.14} (p.~732)} }
\expandafter\def\csname SRef@cdp:eq:C.WF.1\endcsname{\href{S6_SERIES_GRIND_II.pdf\#nameddest=tag-0165}{II.{C.WF.1} (p.~733)} }
\expandafter\def\csname SRef@cdp:eq:C.WF.2\endcsname{\href{S6_SERIES_GRIND_II.pdf\#nameddest=tag-016G}{II.{C.WF.2} (p.~733)} }
\expandafter\def\csname SRef@cdp:eq:C.WF.3\endcsname{\href{S6_SERIES_GRIND_II.pdf\#nameddest=tag-016Q}{II.{C.WF.3} (p.~733)} }
\expandafter\def\csname SRef@cdp:eq:C.WF.4\endcsname{\href{S6_SERIES_GRIND_II.pdf\#nameddest=tag-016R}{II.{C.WF.4} (p.~733)} }
\expandafter\def\csname SRef@cdp:eq:C.WF.5\endcsname{\href{S6_SERIES_GRIND_II.pdf\#nameddest=tag-016S}{II.{C.WF.5} (p.~733)} }
\expandafter\def\csname SRef@cdp:eq:C.UN.2\endcsname{\href{S6_SERIES_GRIND_II.pdf\#nameddest=tag-0153}{II.{C.UN.2} (p.~413)} }
\expandafter\def\csname SRef@cdp:eq:C.PG.6\endcsname{\href{S6_SERIES_GRIND_II.pdf\#nameddest=tag-00V4}{II.{C.PG.6} (p.~445)} }
\expandafter\def\csname SRef@cdp:eq:C.WF.6\endcsname{\href{S6_SERIES_GRIND_II.pdf\#nameddest=tag-016T}{II.{C.WF.6} (p.~734)} }
\expandafter\def\csname SRef@cdp:eq:C.WF.7\endcsname{\href{S6_SERIES_GRIND_II.pdf\#nameddest=tag-016U}{II.{C.WF.7} (p.~734)} }
\expandafter\def\csname SRef@cdp:eq:C.WF.8\endcsname{\href{S6_SERIES_GRIND_II.pdf\#nameddest=tag-016V}{II.{C.WF.8} (p.~734)} }
\expandafter\def\csname SRef@cdp:eq:C.WF.9\endcsname{\href{S6_SERIES_GRIND_II.pdf\#nameddest=tag-016W}{II.{C.WF.9} (p.~735)} }
\expandafter\def\csname SRef@cdp:eq:C.WF.10\endcsname{\href{S6_SERIES_GRIND_II.pdf\#nameddest=tag-0166}{II.{C.WF.10} (p.~735)} }
\expandafter\def\csname SRef@cdp:eq:C.WF.11\endcsname{\href{S6_SERIES_GRIND_II.pdf\#nameddest=tag-0167}{II.{C.WF.11} (p.~735)} }
\expandafter\def\csname SRef@cdp:eq:C.WF.12\endcsname{\href{S6_SERIES_GRIND_II.pdf\#nameddest=tag-0168}{II.{C.WF.12} (p.~735)} }
\expandafter\def\csname SRef@cdp:eq:C.WF.13\endcsname{\href{S6_SERIES_GRIND_II.pdf\#nameddest=tag-0169}{II.{C.WF.13} (p.~736)} }
\expandafter\def\csname SRef@cdp:eq:C.WF.14\endcsname{\href{S6_SERIES_GRIND_II.pdf\#nameddest=tag-016A}{II.{C.WF.14} (p.~736)} }
\expandafter\def\csname SRef@cdp:eq:C.WF.15\endcsname{\href{S6_SERIES_GRIND_II.pdf\#nameddest=tag-016B}{II.{C.WF.15} (p.~736)} }
\expandafter\def\csname SRef@cdp:eq:C.PG.7\endcsname{\href{S6_SERIES_GRIND_II.pdf\#nameddest=tag-00V5}{II.{C.PG.7} (p.~445)} }
\expandafter\def\csname SRef@cdp:eq:C.WF.16\endcsname{\href{S6_SERIES_GRIND_II.pdf\#nameddest=tag-016C}{II.{C.WF.16} (p.~736)} }
\expandafter\def\csname SRef@cdp:eq:C.WF.17\endcsname{\href{S6_SERIES_GRIND_II.pdf\#nameddest=tag-016D}{II.{C.WF.17} (p.~737)} }
\expandafter\def\csname SRef@cdp:eq:C.WF.18\endcsname{\href{S6_SERIES_GRIND_II.pdf\#nameddest=tag-016E}{II.{C.WF.18} (p.~737)} }
\expandafter\def\csname SRef@cdp:eq:C.WF.19\endcsname{\href{S6_SERIES_GRIND_II.pdf\#nameddest=tag-016F}{II.{C.WF.19} (p.~737)} }
\expandafter\def\csname SRef@cdp:eq:C.WF.20\endcsname{\href{S6_SERIES_GRIND_II.pdf\#nameddest=tag-016H}{II.{C.WF.20} (p.~737)} }
\expandafter\def\csname SRef@cdp:eq:C.WF.21\endcsname{\href{S6_SERIES_GRIND_II.pdf\#nameddest=tag-016I}{II.{C.WF.21} (p.~737)} }
\expandafter\def\csname SRef@cdp:eq:C.WF.22\endcsname{\href{S6_SERIES_GRIND_II.pdf\#nameddest=tag-016J}{II.{C.WF.22} (p.~737)} }
\expandafter\def\csname SRef@cdp:eq:C.WF.23\endcsname{\href{S6_SERIES_GRIND_II.pdf\#nameddest=tag-016K}{II.{C.WF.23} (p.~737)} }
\expandafter\def\csname SRef@cdp:eq:C.WF.24\endcsname{\href{S6_SERIES_GRIND_II.pdf\#nameddest=tag-016L}{II.{C.WF.24} (p.~737)} }
\expandafter\def\csname SRef@cdp:eq:C.WF.25\endcsname{\href{S6_SERIES_GRIND_II.pdf\#nameddest=tag-016M}{II.{C.WF.25} (p.~738)} }
\expandafter\def\csname SRef@cdp:eq:C.PG.8\endcsname{\href{S6_SERIES_GRIND_II.pdf\#nameddest=tag-00V6}{II.{C.PG.8} (p.~445)} }
\expandafter\def\csname SRef@cdp:eq:C.WF.26\endcsname{\href{S6_SERIES_GRIND_II.pdf\#nameddest=tag-016N}{II.{C.WF.26} (p.~738)} }
\expandafter\def\csname SRef@cdp:eq:C.WF.27\endcsname{\href{S6_SERIES_GRIND_II.pdf\#nameddest=tag-016O}{II.{C.WF.27} (p.~738)} }
\expandafter\def\csname SRef@cdp:eq:C.WF.28\endcsname{\href{S6_SERIES_GRIND_II.pdf\#nameddest=tag-016P}{II.{C.WF.28} (p.~738)} }
\expandafter\def\csname SRef@cdp:eq:C.RC.1\endcsname{\href{S6_SERIES_GRIND_II.pdf\#nameddest=tag-00XN}{II.{C.RC.1} (p.~739)} }
\expandafter\def\csname SRef@cdp:eq:C.RC.2\endcsname{\href{S6_SERIES_GRIND_II.pdf\#nameddest=tag-00XR}{II.{C.RC.2} (p.~739)} }
\expandafter\def\csname SRef@cdp:eq:C.RC.3\endcsname{\href{S6_SERIES_GRIND_II.pdf\#nameddest=tag-00XS}{II.{C.RC.3} (p.~739)} }
\expandafter\def\csname SRef@cdp:eq:C.RC.4\endcsname{\href{S6_SERIES_GRIND_II.pdf\#nameddest=tag-00XT}{II.{C.RC.4} (p.~740)} }
\expandafter\def\csname SRef@cdp:eq:C.RC.5\endcsname{\href{S6_SERIES_GRIND_II.pdf\#nameddest=tag-00XU}{II.{C.RC.5} (p.~740)} }
\expandafter\def\csname SRef@cdp:eq:C.RC.6\endcsname{\href{S6_SERIES_GRIND_II.pdf\#nameddest=tag-00XV}{II.{C.RC.6} (p.~740)} }
\expandafter\def\csname SRef@cdp:eq:C.RC.7\endcsname{\href{S6_SERIES_GRIND_II.pdf\#nameddest=tag-00XW}{II.{C.RC.7} (p.~740)} }
\expandafter\def\csname SRef@cdp:eq:C.PG.9\endcsname{\href{S6_SERIES_GRIND_II.pdf\#nameddest=tag-00V7}{II.{C.PG.9} (p.~446)} }
\expandafter\def\csname SRef@cdp:eq:C.RC.8\endcsname{\href{S6_SERIES_GRIND_II.pdf\#nameddest=tag-00XX}{II.{C.RC.8} (p.~740)} }
\expandafter\def\csname SRef@cdp:eq:C.RC.9\endcsname{\href{S6_SERIES_GRIND_II.pdf\#nameddest=tag-00XY}{II.{C.RC.9} (p.~741)} }
\expandafter\def\csname SRef@cdp:eq:C.RC.10\endcsname{\href{S6_SERIES_GRIND_II.pdf\#nameddest=tag-00XO}{II.{C.RC.10} (p.~741)} }
\expandafter\def\csname SRef@cdp:eq:C.RC.11\endcsname{\href{S6_SERIES_GRIND_II.pdf\#nameddest=tag-00XP}{II.{C.RC.11} (p.~742)} }
\expandafter\def\csname SRef@cdp:eq:C.RC.12\endcsname{\href{S6_SERIES_GRIND_II.pdf\#nameddest=tag-00XQ}{II.{C.RC.12} (p.~742)} }
\expandafter\def\csname SRef@cdp:eq:C.KG.1\endcsname{\href{S6_SERIES_GRIND_II.pdf\#nameddest=tag-00P8}{II.{C.KG.1} (p.~743)} }
\expandafter\def\csname SRef@cdp:eq:C.KG.2\endcsname{\href{S6_SERIES_GRIND_II.pdf\#nameddest=tag-00PA}{II.{C.KG.2} (p.~743)} }
\expandafter\def\csname SRef@cdp:eq:C.KG.3\endcsname{\href{S6_SERIES_GRIND_II.pdf\#nameddest=tag-00PB}{II.{C.KG.3} (p.~743)} }
\expandafter\def\csname SRef@cdp:eq:C.KG.4\endcsname{\href{S6_SERIES_GRIND_II.pdf\#nameddest=tag-00PC}{II.{C.KG.4} (p.~744)} }
\expandafter\def\csname SRef@cdp:eq:C.KG.5\endcsname{\href{S6_SERIES_GRIND_II.pdf\#nameddest=tag-00PD}{II.{C.KG.5} (p.~744)} }
\expandafter\def\csname SRef@cdp:eq:C.RG.1\endcsname{\href{S6_SERIES_GRIND_II.pdf\#nameddest=tag-00Z3}{II.{C.RG.1} (p.~446)} }
\expandafter\def\csname SRef@cdp:eq:C.KG.6\endcsname{\href{S6_SERIES_GRIND_II.pdf\#nameddest=tag-00PE}{II.{C.KG.6} (p.~744)} }
\expandafter\def\csname SRef@cdp:eq:C.KG.7\endcsname{\href{S6_SERIES_GRIND_II.pdf\#nameddest=tag-00PF}{II.{C.KG.7} (p.~745)} }
\expandafter\def\csname SRef@cdp:eq:C.KG.8\endcsname{\href{S6_SERIES_GRIND_II.pdf\#nameddest=tag-00PG}{II.{C.KG.8} (p.~745)} }
\expandafter\def\csname SRef@cdp:eq:C.KG.9\endcsname{\href{S6_SERIES_GRIND_II.pdf\#nameddest=tag-00PH}{II.{C.KG.9} (p.~747)} }
\expandafter\def\csname SRef@cdp:eq:C.KG.10\endcsname{\href{S6_SERIES_GRIND_II.pdf\#nameddest=tag-00P9}{II.{C.KG.10} (p.~747)} }
\expandafter\def\csname SRef@cdp:eq:C.HM.1\endcsname{\href{S6_SERIES_GRIND_II.pdf\#nameddest=tag-00MG}{II.{C.HM.1} (p.~748)} }
\expandafter\def\csname SRef@cdp:eq:C.HM.2\endcsname{\href{S6_SERIES_GRIND_II.pdf\#nameddest=tag-00MI}{II.{C.HM.2} (p.~748)} }
\expandafter\def\csname SRef@cdp:eq:C.HM.3\endcsname{\href{S6_SERIES_GRIND_II.pdf\#nameddest=tag-00MJ}{II.{C.HM.3} (p.~748)} }
\expandafter\def\csname SRef@cdp:eq:C.HM.4\endcsname{\href{S6_SERIES_GRIND_II.pdf\#nameddest=tag-00MK}{II.{C.HM.4} (p.~749)} }
\expandafter\def\csname SRef@cdp:eq:C.HM.5\endcsname{\href{S6_SERIES_GRIND_II.pdf\#nameddest=tag-00ML}{II.{C.HM.5} (p.~749)} }
\expandafter\def\csname SRef@cdp:eq:C.RG.2\endcsname{\href{S6_SERIES_GRIND_II.pdf\#nameddest=tag-00Z8}{II.{C.RG.2} (p.~446)} }
\expandafter\def\csname SRef@cdp:eq:C.HM.6\endcsname{\href{S6_SERIES_GRIND_II.pdf\#nameddest=tag-00MM}{II.{C.HM.6} (p.~749)} }
\expandafter\def\csname SRef@cdp:eq:C.HM.7\endcsname{\href{S6_SERIES_GRIND_II.pdf\#nameddest=tag-00MN}{II.{C.HM.7} (p.~749)} }
\expandafter\def\csname SRef@cdp:eq:C.HM.8\endcsname{\href{S6_SERIES_GRIND_II.pdf\#nameddest=tag-00MO}{II.{C.HM.8} (p.~750)} }
\expandafter\def\csname SRef@cdp:eq:C.HM.9\endcsname{\href{S6_SERIES_GRIND_II.pdf\#nameddest=tag-00MP}{II.{C.HM.9} (p.~750)} }
\expandafter\def\csname SRef@cdp:eq:C.HM.10\endcsname{\href{S6_SERIES_GRIND_II.pdf\#nameddest=tag-00MH}{II.{C.HM.10} (p.~751)} }
\expandafter\def\csname SRef@cdp:eq:C.BC.1\endcsname{\href{S6_SERIES_GRIND_II.pdf\#nameddest=tag-00AZ}{II.{C.BC.1} (p.~752)} }
\expandafter\def\csname SRef@cdp:eq:C.BC.2\endcsname{\href{S6_SERIES_GRIND_II.pdf\#nameddest=tag-00B4}{II.{C.BC.2} (p.~752)} }
\expandafter\def\csname SRef@cdp:eq:C.BC.3\endcsname{\href{S6_SERIES_GRIND_II.pdf\#nameddest=tag-00B5}{II.{C.BC.3} (p.~752)} }
\expandafter\def\csname SRef@cdp:eq:C.BC.4\endcsname{\href{S6_SERIES_GRIND_II.pdf\#nameddest=tag-00B6}{II.{C.BC.4} (p.~752)} }
\expandafter\def\csname SRef@cdp:eq:C.BC.5\endcsname{\href{S6_SERIES_GRIND_II.pdf\#nameddest=tag-00B7}{II.{C.BC.5} (p.~753)} }
\expandafter\def\csname SRef@cdp:eq:C.RG.3\endcsname{\href{S6_SERIES_GRIND_II.pdf\#nameddest=tag-00Z9}{II.{C.RG.3} (p.~447)} }
\expandafter\def\csname SRef@cdp:eq:C.BC.6\endcsname{\href{S6_SERIES_GRIND_II.pdf\#nameddest=tag-00B8}{II.{C.BC.6} (p.~753)} }
\expandafter\def\csname SRef@cdp:eq:C.BC.7\endcsname{\href{S6_SERIES_GRIND_II.pdf\#nameddest=tag-00B9}{II.{C.BC.7} (p.~753)} }
\expandafter\def\csname SRef@cdp:eq:C.BC.8\endcsname{\href{S6_SERIES_GRIND_II.pdf\#nameddest=tag-00BA}{II.{C.BC.8} (p.~753)} }
\expandafter\def\csname SRef@cdp:eq:C.BC.9\endcsname{\href{S6_SERIES_GRIND_II.pdf\#nameddest=tag-00BB}{II.{C.BC.9} (p.~753)} }
\expandafter\def\csname SRef@cdp:eq:C.BC.10\endcsname{\href{S6_SERIES_GRIND_II.pdf\#nameddest=tag-00B0}{II.{C.BC.10} (p.~754)} }
\expandafter\def\csname SRef@cdp:eq:C.BC.11\endcsname{\href{S6_SERIES_GRIND_II.pdf\#nameddest=tag-00B1}{II.{C.BC.11} (p.~754)} }
\expandafter\def\csname SRef@cdp:eq:C.BC.12\endcsname{\href{S6_SERIES_GRIND_II.pdf\#nameddest=tag-00B2}{II.{C.BC.12} (p.~754)} }
\expandafter\def\csname SRef@cdp:eq:C.BC.13\endcsname{\href{S6_SERIES_GRIND_II.pdf\#nameddest=tag-00B3}{II.{C.BC.13} (p.~754)} }
\expandafter\def\csname SRef@cdp:eq:C.EP.1\endcsname{\href{S6_SERIES_GRIND_II.pdf\#nameddest=tag-00I1}{II.{C.EP.1} (p.~755)} }
\expandafter\def\csname SRef@cdp:eq:C.EP.2\endcsname{\href{S6_SERIES_GRIND_II.pdf\#nameddest=tag-00I4}{II.{C.EP.2} (p.~755)} }
\expandafter\def\csname SRef@cdp:eq:C.RG.4\endcsname{\href{S6_SERIES_GRIND_II.pdf\#nameddest=tag-00ZA}{II.{C.RG.4} (p.~447)} }
\expandafter\def\csname SRef@cdp:eq:C.EP.3\endcsname{\href{S6_SERIES_GRIND_II.pdf\#nameddest=tag-00I5}{II.{C.EP.3} (p.~755)} }
\expandafter\def\csname SRef@cdp:eq:C.EP.4\endcsname{\href{S6_SERIES_GRIND_II.pdf\#nameddest=tag-00I6}{II.{C.EP.4} (p.~756)} }
\expandafter\def\csname SRef@cdp:eq:C.EP.5\endcsname{\href{S6_SERIES_GRIND_II.pdf\#nameddest=tag-00I7}{II.{C.EP.5} (p.~756)} }
\expandafter\def\csname SRef@cdp:eq:C.EP.6\endcsname{\href{S6_SERIES_GRIND_II.pdf\#nameddest=tag-00I8}{II.{C.EP.6} (p.~756)} }
\expandafter\def\csname SRef@cdp:eq:C.EP.7\endcsname{\href{S6_SERIES_GRIND_II.pdf\#nameddest=tag-00I9}{II.{C.EP.7} (p.~756)} }
\expandafter\def\csname SRef@cdp:eq:C.EP.8\endcsname{\href{S6_SERIES_GRIND_II.pdf\#nameddest=tag-00IA}{II.{C.EP.8} (p.~756)} }
\expandafter\def\csname SRef@cdp:eq:C.EP.9\endcsname{\href{S6_SERIES_GRIND_II.pdf\#nameddest=tag-00IB}{II.{C.EP.9} (p.~757)} }
\expandafter\def\csname SRef@cdp:eq:C.EP.10\endcsname{\href{S6_SERIES_GRIND_II.pdf\#nameddest=tag-00I2}{II.{C.EP.10} (p.~757)} }
\expandafter\def\csname SRef@cdp:eq:C.EP.11\endcsname{\href{S6_SERIES_GRIND_II.pdf\#nameddest=tag-00I3}{II.{C.EP.11} (p.~757)} }
\expandafter\def\csname SRef@cdp:eq:C.AP.1\endcsname{\href{S6_SERIES_GRIND_II.pdf\#nameddest=tag-009M}{II.{C.AP.1} (p.~757)} }
\expandafter\def\csname SRef@cdp:eq:C.RG.5\endcsname{\href{S6_SERIES_GRIND_II.pdf\#nameddest=tag-00ZB}{II.{C.RG.5} (p.~447)} }
\expandafter\def\csname SRef@cdp:eq:C.AP.2\endcsname{\href{S6_SERIES_GRIND_II.pdf\#nameddest=tag-009S}{II.{C.AP.2} (p.~758)} }
\expandafter\def\csname SRef@cdp:eq:C.AP.3\endcsname{\href{S6_SERIES_GRIND_II.pdf\#nameddest=tag-009T}{II.{C.AP.3} (p.~758)} }
\expandafter\def\csname SRef@cdp:eq:C.AP.4\endcsname{\href{S6_SERIES_GRIND_II.pdf\#nameddest=tag-009U}{II.{C.AP.4} (p.~758)} }
\expandafter\def\csname SRef@cdp:eq:C.AP.5\endcsname{\href{S6_SERIES_GRIND_II.pdf\#nameddest=tag-009V}{II.{C.AP.5} (p.~759)} }
\expandafter\def\csname SRef@cdp:eq:C.AP.6\endcsname{\href{S6_SERIES_GRIND_II.pdf\#nameddest=tag-009W}{II.{C.AP.6} (p.~759)} }
\expandafter\def\csname SRef@cdp:eq:C.AP.7\endcsname{\href{S6_SERIES_GRIND_II.pdf\#nameddest=tag-009X}{II.{C.AP.7} (p.~759)} }
\expandafter\def\csname SRef@cdp:eq:C.AP.8\endcsname{\href{S6_SERIES_GRIND_II.pdf\#nameddest=tag-009Y}{II.{C.AP.8} (p.~760)} }
\expandafter\def\csname SRef@cdp:eq:C.AP.9\endcsname{\href{S6_SERIES_GRIND_II.pdf\#nameddest=tag-009Z}{II.{C.AP.9} (p.~760)} }
\expandafter\def\csname SRef@cdp:eq:C.AP.10\endcsname{\href{S6_SERIES_GRIND_II.pdf\#nameddest=tag-009N}{II.{C.AP.10} (p.~760)} }
\expandafter\def\csname SRef@cdp:eq:C.AP.11\endcsname{\href{S6_SERIES_GRIND_II.pdf\#nameddest=tag-009O}{II.{C.AP.11} (p.~760)} }
\expandafter\def\csname SRef@cdp:eq:C.RG.6\endcsname{\href{S6_SERIES_GRIND_II.pdf\#nameddest=tag-00ZC}{II.{C.RG.6} (p.~448)} }
\expandafter\def\csname SRef@cdp:eq:C.AP.12\endcsname{\href{S6_SERIES_GRIND_II.pdf\#nameddest=tag-009P}{II.{C.AP.12} (p.~761)} }
\expandafter\def\csname SRef@cdp:eq:C.AP.13\endcsname{\href{S6_SERIES_GRIND_II.pdf\#nameddest=tag-009Q}{II.{C.AP.13} (p.~761)} }
\expandafter\def\csname SRef@cdp:eq:C.AP.14\endcsname{\href{S6_SERIES_GRIND_II.pdf\#nameddest=tag-009R}{II.{C.AP.14} (p.~761)} }
\expandafter\def\csname SRef@cdp:eq:C.ML.1\endcsname{\href{S6_SERIES_GRIND_II.pdf\#nameddest=tag-00QR}{II.{C.ML.1} (p.~762)} }
\expandafter\def\csname SRef@cdp:eq:C.ML.2\endcsname{\href{S6_SERIES_GRIND_II.pdf\#nameddest=tag-00QT}{II.{C.ML.2} (p.~762)} }
\expandafter\def\csname SRef@cdp:eq:C.ML.3\endcsname{\href{S6_SERIES_GRIND_II.pdf\#nameddest=tag-00QU}{II.{C.ML.3} (p.~763)} }
\expandafter\def\csname SRef@cdp:eq:C.ML.4\endcsname{\href{S6_SERIES_GRIND_II.pdf\#nameddest=tag-00QV}{II.{C.ML.4} (p.~763)} }
\expandafter\def\csname SRef@cdp:eq:C.ML.5\endcsname{\href{S6_SERIES_GRIND_II.pdf\#nameddest=tag-00QW}{II.{C.ML.5} (p.~763)} }
\expandafter\def\csname SRef@cdp:eq:C.ML.6\endcsname{\href{S6_SERIES_GRIND_II.pdf\#nameddest=tag-00QX}{II.{C.ML.6} (p.~763)} }
\expandafter\def\csname SRef@cdp:eq:C.ML.7\endcsname{\href{S6_SERIES_GRIND_II.pdf\#nameddest=tag-00QY}{II.{C.ML.7} (p.~764)} }
\expandafter\def\csname SRef@cdp:eq:C.UN.3\endcsname{\href{S6_SERIES_GRIND_II.pdf\#nameddest=tag-0154}{II.{C.UN.3} (p.~413)} }
\expandafter\def\csname SRef@cdp:eq:C.RG.7\endcsname{\href{S6_SERIES_GRIND_II.pdf\#nameddest=tag-00ZD}{II.{C.RG.7} (p.~448)} }
\expandafter\def\csname SRef@cdp:eq:C.ML.8\endcsname{\href{S6_SERIES_GRIND_II.pdf\#nameddest=tag-00QZ}{II.{C.ML.8} (p.~764)} }
\expandafter\def\csname SRef@cdp:eq:C.ML.9\endcsname{\href{S6_SERIES_GRIND_II.pdf\#nameddest=tag-00R0}{II.{C.ML.9} (p.~764)} }
\expandafter\def\csname SRef@cdp:eq:C.ML.10\endcsname{\href{S6_SERIES_GRIND_II.pdf\#nameddest=tag-00QS}{II.{C.ML.10} (p.~765)} }
\expandafter\def\csname SRef@cdp:eq:C.TI.1\endcsname{\href{S6_SERIES_GRIND_II.pdf\#nameddest=tag-0135}{II.{C.TI.1} (p.~766)} }
\expandafter\def\csname SRef@cdp:eq:C.TI.2\endcsname{\href{S6_SERIES_GRIND_II.pdf\#nameddest=tag-013B}{II.{C.TI.2} (p.~766)} }
\expandafter\def\csname SRef@cdp:eq:C.TI.3\endcsname{\href{S6_SERIES_GRIND_II.pdf\#nameddest=tag-013C}{II.{C.TI.3} (p.~766)} }
\expandafter\def\csname SRef@cdp:eq:C.TI.4\endcsname{\href{S6_SERIES_GRIND_II.pdf\#nameddest=tag-013D}{II.{C.TI.4} (p.~766)} }
\expandafter\def\csname SRef@cdp:eq:C.TI.5\endcsname{\href{S6_SERIES_GRIND_II.pdf\#nameddest=tag-013E}{II.{C.TI.5} (p.~766)} }
\expandafter\def\csname SRef@cdp:eq:C.TI.6\endcsname{\href{S6_SERIES_GRIND_II.pdf\#nameddest=tag-013F}{II.{C.TI.6} (p.~767)} }
\expandafter\def\csname SRef@cdp:eq:C.TI.7\endcsname{\href{S6_SERIES_GRIND_II.pdf\#nameddest=tag-013G}{II.{C.TI.7} (p.~767)} }
\expandafter\def\csname SRef@cdp:eq:C.RG.8\endcsname{\href{S6_SERIES_GRIND_II.pdf\#nameddest=tag-00ZE}{II.{C.RG.8} (p.~448)} }
\expandafter\def\csname SRef@cdp:eq:C.TI.8\endcsname{\href{S6_SERIES_GRIND_II.pdf\#nameddest=tag-013H}{II.{C.TI.8} (p.~767)} }
\expandafter\def\csname SRef@cdp:eq:C.TI.9\endcsname{\href{S6_SERIES_GRIND_II.pdf\#nameddest=tag-013I}{II.{C.TI.9} (p.~768)} }
\expandafter\def\csname SRef@cdp:eq:C.TI.10\endcsname{\href{S6_SERIES_GRIND_II.pdf\#nameddest=tag-0136}{II.{C.TI.10} (p.~768)} }
\expandafter\def\csname SRef@cdp:eq:C.TI.11\endcsname{\href{S6_SERIES_GRIND_II.pdf\#nameddest=tag-0137}{II.{C.TI.11} (p.~768)} }
\expandafter\def\csname SRef@cdp:eq:C.TI.12\endcsname{\href{S6_SERIES_GRIND_II.pdf\#nameddest=tag-0138}{II.{C.TI.12} (p.~768)} }
\expandafter\def\csname SRef@cdp:eq:C.TI.13\endcsname{\href{S6_SERIES_GRIND_II.pdf\#nameddest=tag-0139}{II.{C.TI.13} (p.~769)} }
\expandafter\def\csname SRef@cdp:eq:C.TI.14\endcsname{\href{S6_SERIES_GRIND_II.pdf\#nameddest=tag-013A}{II.{C.TI.14} (p.~769)} }
\expandafter\def\csname SRef@cdp:eq:C.TK.1\endcsname{\href{S6_SERIES_GRIND_II.pdf\#nameddest=tag-013J}{II.{C.TK.1} (p.~770)} }
\expandafter\def\csname SRef@cdp:eq:C.TK.2\endcsname{\href{S6_SERIES_GRIND_II.pdf\#nameddest=tag-013K}{II.{C.TK.2} (p.~770)} }
\expandafter\def\csname SRef@cdp:eq:C.TK.3\endcsname{\href{S6_SERIES_GRIND_II.pdf\#nameddest=tag-013L}{II.{C.TK.3} (p.~771)} }
\expandafter\def\csname SRef@cdp:eq:C.RG.9\endcsname{\href{S6_SERIES_GRIND_II.pdf\#nameddest=tag-00ZF}{II.{C.RG.9} (p.~448)} }
\expandafter\def\csname SRef@cdp:eq:C.TK.4\endcsname{\href{S6_SERIES_GRIND_II.pdf\#nameddest=tag-013M}{II.{C.TK.4} (p.~771)} }
\expandafter\def\csname SRef@cdp:eq:C.TK.5\endcsname{\href{S6_SERIES_GRIND_II.pdf\#nameddest=tag-013N}{II.{C.TK.5} (p.~771)} }
\expandafter\def\csname SRef@cdp:eq:C.TK.6\endcsname{\href{S6_SERIES_GRIND_II.pdf\#nameddest=tag-013O}{II.{C.TK.6} (p.~772)} }
\expandafter\def\csname SRef@cdp:eq:C.TK.7\endcsname{\href{S6_SERIES_GRIND_II.pdf\#nameddest=tag-013P}{II.{C.TK.7} (p.~772)} }
\expandafter\def\csname SRef@cdp:eq:C.RK.1\endcsname{\href{S6_SERIES_GRIND_II.pdf\#nameddest=tag-00ZQ}{II.{C.RK.1} (p.~773)} }
\expandafter\def\csname SRef@cdp:eq:C.RK.2\endcsname{\href{S6_SERIES_GRIND_II.pdf\#nameddest=tag-0101}{II.{C.RK.2} (p.~773)} }
\expandafter\def\csname SRef@cdp:eq:C.RK.3\endcsname{\href{S6_SERIES_GRIND_II.pdf\#nameddest=tag-0105}{II.{C.RK.3} (p.~773)} }
\expandafter\def\csname SRef@cdp:eq:C.RK.4\endcsname{\href{S6_SERIES_GRIND_II.pdf\#nameddest=tag-0106}{II.{C.RK.4} (p.~773)} }
\expandafter\def\csname SRef@cdp:eq:C.RK.5\endcsname{\href{S6_SERIES_GRIND_II.pdf\#nameddest=tag-0107}{II.{C.RK.5} (p.~773)} }
\expandafter\def\csname SRef@cdp:eq:C.RK.6\endcsname{\href{S6_SERIES_GRIND_II.pdf\#nameddest=tag-0108}{II.{C.RK.6} (p.~774)} }
\expandafter\def\csname SRef@cdp:eq:C.RG.10\endcsname{\href{S6_SERIES_GRIND_II.pdf\#nameddest=tag-00Z4}{II.{C.RG.10} (p.~448)} }
\expandafter\def\csname SRef@cdp:eq:C.RK.7\endcsname{\href{S6_SERIES_GRIND_II.pdf\#nameddest=tag-0109}{II.{C.RK.7} (p.~774)} }
\expandafter\def\csname SRef@cdp:eq:C.RK.8\endcsname{\href{S6_SERIES_GRIND_II.pdf\#nameddest=tag-010A}{II.{C.RK.8} (p.~774)} }
\expandafter\def\csname SRef@cdp:eq:C.RK.18\endcsname{\href{S6_SERIES_GRIND_II.pdf\#nameddest=tag-00ZZ}{II.{C.RK.18} (p.~774)} }
\expandafter\def\csname SRef@cdp:eq:C.RK.19\endcsname{\href{S6_SERIES_GRIND_II.pdf\#nameddest=tag-0100}{II.{C.RK.19} (p.~774)} }
\expandafter\def\csname SRef@cdp:eq:C.RK.20\endcsname{\href{S6_SERIES_GRIND_II.pdf\#nameddest=tag-0102}{II.{C.RK.20} (p.~774)} }
\expandafter\def\csname SRef@cdp:eq:C.RK.9\endcsname{\href{S6_SERIES_GRIND_II.pdf\#nameddest=tag-010B}{II.{C.RK.9} (p.~775)} }
\expandafter\def\csname SRef@cdp:eq:C.RK.10\endcsname{\href{S6_SERIES_GRIND_II.pdf\#nameddest=tag-00ZR}{II.{C.RK.10} (p.~775)} }
\expandafter\def\csname SRef@cdp:eq:C.RK.11\endcsname{\href{S6_SERIES_GRIND_II.pdf\#nameddest=tag-00ZS}{II.{C.RK.11} (p.~775)} }
\expandafter\def\csname SRef@cdp:eq:C.RK.12\endcsname{\href{S6_SERIES_GRIND_II.pdf\#nameddest=tag-00ZT}{II.{C.RK.12} (p.~776)} }
\expandafter\def\csname SRef@cdp:eq:C.RK.13\endcsname{\href{S6_SERIES_GRIND_II.pdf\#nameddest=tag-00ZU}{II.{C.RK.13} (p.~776)} }
\expandafter\def\csname SRef@cdp:eq:C.RG.11\endcsname{\href{S6_SERIES_GRIND_II.pdf\#nameddest=tag-00Z5}{II.{C.RG.11} (p.~448)} }
\expandafter\def\csname SRef@cdp:eq:C.RK.14\endcsname{\href{S6_SERIES_GRIND_II.pdf\#nameddest=tag-00ZV}{II.{C.RK.14} (p.~776)} }
\expandafter\def\csname SRef@cdp:eq:C.RK.15\endcsname{\href{S6_SERIES_GRIND_II.pdf\#nameddest=tag-00ZW}{II.{C.RK.15} (p.~776)} }
\expandafter\def\csname SRef@cdp:eq:C.RK.16\endcsname{\href{S6_SERIES_GRIND_II.pdf\#nameddest=tag-00ZX}{II.{C.RK.16} (p.~777)} }
\expandafter\def\csname SRef@cdp:eq:C.RK.17\endcsname{\href{S6_SERIES_GRIND_II.pdf\#nameddest=tag-00ZY}{II.{C.RK.17} (p.~777)} }
\expandafter\def\csname SRef@cdp:eq:C.RK.21\endcsname{\href{S6_SERIES_GRIND_II.pdf\#nameddest=tag-0103}{II.{C.RK.21} (p.~777)} }
\expandafter\def\csname SRef@cdp:eq:C.RK.22\endcsname{\href{S6_SERIES_GRIND_II.pdf\#nameddest=tag-0104}{II.{C.RK.22} (p.~778)} }
\expandafter\def\csname SRef@cdp:eq:C.NT.1\endcsname{\href{S6_SERIES_GRIND_II.pdf\#nameddest=tag-00TU}{II.{C.NT.1} (p.~779)} }
\expandafter\def\csname SRef@cdp:eq:C.NT.2\endcsname{\href{S6_SERIES_GRIND_II.pdf\#nameddest=tag-00TV}{II.{C.NT.2} (p.~779)} }
\expandafter\def\csname SRef@cdp:eq:C.NT.3\endcsname{\href{S6_SERIES_GRIND_II.pdf\#nameddest=tag-00TW}{II.{C.NT.3} (p.~779)} }
\expandafter\def\csname SRef@cdp:eq:C.NT.4\endcsname{\href{S6_SERIES_GRIND_II.pdf\#nameddest=tag-00TX}{II.{C.NT.4} (p.~779)} }
\expandafter\def\csname SRef@cdp:eq:C.RG.12\endcsname{\href{S6_SERIES_GRIND_II.pdf\#nameddest=tag-00Z6}{II.{C.RG.12} (p.~449)} }
\expandafter\def\csname SRef@cdp:eq:C.NT.5\endcsname{\href{S6_SERIES_GRIND_II.pdf\#nameddest=tag-00TY}{II.{C.NT.5} (p.~780)} }
\expandafter\def\csname SRef@cdp:eq:C.NT.6\endcsname{\href{S6_SERIES_GRIND_II.pdf\#nameddest=tag-00TZ}{II.{C.NT.6} (p.~780)} }
\expandafter\def\csname SRef@cdp:eq:C.NT.7\endcsname{\href{S6_SERIES_GRIND_II.pdf\#nameddest=tag-00U0}{II.{C.NT.7} (p.~780)} }
\expandafter\def\csname SRef@cdp:eq:C.NT.8\endcsname{\href{S6_SERIES_GRIND_II.pdf\#nameddest=tag-00U1}{II.{C.NT.8} (p.~781)} }
\expandafter\def\csname SRef@cdp:eq:C.GP.1\endcsname{\href{S6_SERIES_GRIND_II.pdf\#nameddest=tag-00KY}{II.{C.GP.1} (p.~782)} }
\expandafter\def\csname SRef@cdp:eq:C.GP.2\endcsname{\href{S6_SERIES_GRIND_II.pdf\#nameddest=tag-00L0}{II.{C.GP.2} (p.~783)} }
\expandafter\def\csname SRef@cdp:eq:C.GP.3\endcsname{\href{S6_SERIES_GRIND_II.pdf\#nameddest=tag-00L1}{II.{C.GP.3} (p.~783)} }
\expandafter\def\csname SRef@cdp:eq:C.GP.4\endcsname{\href{S6_SERIES_GRIND_II.pdf\#nameddest=tag-00L2}{II.{C.GP.4} (p.~783)} }
\expandafter\def\csname SRef@cdp:eq:C.GP.5\endcsname{\href{S6_SERIES_GRIND_II.pdf\#nameddest=tag-00L3}{II.{C.GP.5} (p.~783)} }
\expandafter\def\csname SRef@cdp:eq:C.GP.6\endcsname{\href{S6_SERIES_GRIND_II.pdf\#nameddest=tag-00L4}{II.{C.GP.6} (p.~784)} }
\expandafter\def\csname SRef@cdp:eq:C.RG.13\endcsname{\href{S6_SERIES_GRIND_II.pdf\#nameddest=tag-00Z7}{II.{C.RG.13} (p.~449)} }
\expandafter\def\csname SRef@cdp:eq:C.GP.7\endcsname{\href{S6_SERIES_GRIND_II.pdf\#nameddest=tag-00L5}{II.{C.GP.7} (p.~785)} }
\expandafter\def\csname SRef@cdp:eq:C.GP.8\endcsname{\href{S6_SERIES_GRIND_II.pdf\#nameddest=tag-00L6}{II.{C.GP.8} (p.~785)} }
\expandafter\def\csname SRef@cdp:eq:C.GP.9\endcsname{\href{S6_SERIES_GRIND_II.pdf\#nameddest=tag-00L7}{II.{C.GP.9} (p.~785)} }
\expandafter\def\csname SRef@cdp:eq:C.GP.10\endcsname{\href{S6_SERIES_GRIND_II.pdf\#nameddest=tag-00KZ}{II.{C.GP.10} (p.~786)} }
\expandafter\def\csname SRef@cdp:eq:C.AR.1\endcsname{\href{S6_SERIES_GRIND_II.pdf\#nameddest=tag-00A0}{II.{C.AR.1} (p.~788)} }
\expandafter\def\csname SRef@cdp:eq:C.AR.2\endcsname{\href{S6_SERIES_GRIND_II.pdf\#nameddest=tag-00A8}{II.{C.AR.2} (p.~788)} }
\expandafter\def\csname SRef@cdp:eq:C.AR.3\endcsname{\href{S6_SERIES_GRIND_II.pdf\#nameddest=tag-00A9}{II.{C.AR.3} (p.~788)} }
\expandafter\def\csname SRef@cdp:eq:C.AR.4\endcsname{\href{S6_SERIES_GRIND_II.pdf\#nameddest=tag-00AA}{II.{C.AR.4} (p.~789)} }
\expandafter\def\csname SRef@cdp:eq:C.AR.5\endcsname{\href{S6_SERIES_GRIND_II.pdf\#nameddest=tag-00AB}{II.{C.AR.5} (p.~789)} }
\expandafter\def\csname SRef@cdp:eq:C.CS.1\endcsname{\href{S6_SERIES_GRIND_II.pdf\#nameddest=tag-00DH}{II.{C.CS.1} (p.~450)} }
\expandafter\def\csname SRef@cdp:eq:C.AR.6\endcsname{\href{S6_SERIES_GRIND_II.pdf\#nameddest=tag-00AC}{II.{C.AR.6} (p.~789)} }
\expandafter\def\csname SRef@cdp:eq:C.AR.7\endcsname{\href{S6_SERIES_GRIND_II.pdf\#nameddest=tag-00AD}{II.{C.AR.7} (p.~789)} }
\expandafter\def\csname SRef@cdp:eq:C.AR.8\endcsname{\href{S6_SERIES_GRIND_II.pdf\#nameddest=tag-00AE}{II.{C.AR.8} (p.~789)} }
\expandafter\def\csname SRef@cdp:eq:C.AR.9\endcsname{\href{S6_SERIES_GRIND_II.pdf\#nameddest=tag-00AF}{II.{C.AR.9} (p.~789)} }
\expandafter\def\csname SRef@cdp:eq:C.AR.10\endcsname{\href{S6_SERIES_GRIND_II.pdf\#nameddest=tag-00A1}{II.{C.AR.10} (p.~789)} }
\expandafter\def\csname SRef@cdp:eq:C.AR.11\endcsname{\href{S6_SERIES_GRIND_II.pdf\#nameddest=tag-00A2}{II.{C.AR.11} (p.~789)} }
\expandafter\def\csname SRef@cdp:eq:C.AR.12\endcsname{\href{S6_SERIES_GRIND_II.pdf\#nameddest=tag-00A3}{II.{C.AR.12} (p.~790)} }
\expandafter\def\csname SRef@cdp:eq:C.AR.13\endcsname{\href{S6_SERIES_GRIND_II.pdf\#nameddest=tag-00A4}{II.{C.AR.13} (p.~790)} }
\expandafter\def\csname SRef@cdp:eq:C.AR.14\endcsname{\href{S6_SERIES_GRIND_II.pdf\#nameddest=tag-00A5}{II.{C.AR.14} (p.~791)} }
\expandafter\def\csname SRef@cdp:eq:C.AR.15\endcsname{\href{S6_SERIES_GRIND_II.pdf\#nameddest=tag-00A6}{II.{C.AR.15} (p.~791)} }
\expandafter\def\csname SRef@cdp:eq:C.CS.2\endcsname{\href{S6_SERIES_GRIND_II.pdf\#nameddest=tag-00DI}{II.{C.CS.2} (p.~451)} }
\expandafter\def\csname SRef@cdp:eq:C.AR.16\endcsname{\href{S6_SERIES_GRIND_II.pdf\#nameddest=tag-00A7}{II.{C.AR.16} (p.~791)} }
\expandafter\def\csname SRef@cdp:eq:C.G.1\endcsname{\href{S6_SERIES_GRIND_II.pdf\#nameddest=tag-00KW}{II.{C.G.1} (p.~793)} }
\expandafter\def\csname SRef@cdp:eq:C.G.2\endcsname{\href{S6_SERIES_GRIND_II.pdf\#nameddest=tag-00KX}{II.{C.G.2} (p.~794)} }
\expandafter\def\csname SRef@cdp:eq:C.CS.3\endcsname{\href{S6_SERIES_GRIND_II.pdf\#nameddest=tag-00DJ}{II.{C.CS.3} (p.~451)} }
\expandafter\def\csname SRef@cdp:eq:C.UN.4\endcsname{\href{S6_SERIES_GRIND_II.pdf\#nameddest=tag-0155}{II.{C.UN.4} (p.~414)} }
\expandafter\def\csname SRef@cdp:eq:C.CS.4\endcsname{\href{S6_SERIES_GRIND_II.pdf\#nameddest=tag-00DK}{II.{C.CS.4} (p.~451)} }
\expandafter\def\csname SRef@cdp:eq:C.CS.5\endcsname{\href{S6_SERIES_GRIND_II.pdf\#nameddest=tag-00DL}{II.{C.CS.5} (p.~452)} }
\expandafter\def\csname SRef@cdp:eq:C.CS.6\endcsname{\href{S6_SERIES_GRIND_II.pdf\#nameddest=tag-00DM}{II.{C.CS.6} (p.~452)} }
\expandafter\def\csname SRef@cdp:eq:C.EQ.1\endcsname{\href{S6_SERIES_GRIND_II.pdf\#nameddest=tag-00IC}{II.{C.EQ.1} (p.~453)} }
\expandafter\def\csname SRef@cdp:eq:C.EQ.2\endcsname{\href{S6_SERIES_GRIND_II.pdf\#nameddest=tag-00IF}{II.{C.EQ.2} (p.~453)} }
\expandafter\def\csname SRef@cdp:eq:C.EQ.3\endcsname{\href{S6_SERIES_GRIND_II.pdf\#nameddest=tag-00IG}{II.{C.EQ.3} (p.~454)} }
\expandafter\def\csname SRef@cdp:eq:C.EQ.4\endcsname{\href{S6_SERIES_GRIND_II.pdf\#nameddest=tag-00IH}{II.{C.EQ.4} (p.~454)} }
\expandafter\def\csname SRef@cdp:eq:C.EQ.5\endcsname{\href{S6_SERIES_GRIND_II.pdf\#nameddest=tag-00II}{II.{C.EQ.5} (p.~454)} }
\expandafter\def\csname SRef@cdp:eq:C.EQ.6\endcsname{\href{S6_SERIES_GRIND_II.pdf\#nameddest=tag-00IJ}{II.{C.EQ.6} (p.~455)} }
\expandafter\def\csname SRef@cdp:eq:C.EQ.7\endcsname{\href{S6_SERIES_GRIND_II.pdf\#nameddest=tag-00IK}{II.{C.EQ.7} (p.~455)} }
\expandafter\def\csname SRef@cdp:eq:C.UN.5\endcsname{\href{S6_SERIES_GRIND_II.pdf\#nameddest=tag-0156}{II.{C.UN.5} (p.~415)} }
\expandafter\def\csname SRef@cdp:eq:C.EQ.8\endcsname{\href{S6_SERIES_GRIND_II.pdf\#nameddest=tag-00IL}{II.{C.EQ.8} (p.~455)} }
\expandafter\def\csname SRef@cdp:eq:C.EQ.9\endcsname{\href{S6_SERIES_GRIND_II.pdf\#nameddest=tag-00IM}{II.{C.EQ.9} (p.~456)} }
\expandafter\def\csname SRef@cdp:eq:C.EQ.10\endcsname{\href{S6_SERIES_GRIND_II.pdf\#nameddest=tag-00ID}{II.{C.EQ.10} (p.~456)} }
\expandafter\def\csname SRef@cdp:eq:C.EQ.11\endcsname{\href{S6_SERIES_GRIND_II.pdf\#nameddest=tag-00IE}{II.{C.EQ.11} (p.~456)} }
\expandafter\def\csname SRef@cdp:eq:C.CL.1\endcsname{\href{S6_SERIES_GRIND_II.pdf\#nameddest=tag-00BW}{II.{C.CL.1} (p.~457)} }
\expandafter\def\csname SRef@cdp:eq:C.CL.2\endcsname{\href{S6_SERIES_GRIND_II.pdf\#nameddest=tag-00C1}{II.{C.CL.2} (p.~457)} }
\expandafter\def\csname SRef@cdp:eq:C.CL.3\endcsname{\href{S6_SERIES_GRIND_II.pdf\#nameddest=tag-00C2}{II.{C.CL.3} (p.~457)} }
\expandafter\def\csname SRef@cdp:eq:C.CL.4\endcsname{\href{S6_SERIES_GRIND_II.pdf\#nameddest=tag-00C3}{II.{C.CL.4} (p.~458)} }
\expandafter\def\csname SRef@cdp:eq:C.CL.5\endcsname{\href{S6_SERIES_GRIND_II.pdf\#nameddest=tag-00C4}{II.{C.CL.5} (p.~458)} }
\expandafter\def\csname SRef@cdp:eq:C.CL.6\endcsname{\href{S6_SERIES_GRIND_II.pdf\#nameddest=tag-00C5}{II.{C.CL.6} (p.~458)} }
\expandafter\def\csname SRef@cdp:eq:C.UN.6\endcsname{\href{S6_SERIES_GRIND_II.pdf\#nameddest=tag-0157}{II.{C.UN.6} (p.~415)} }
\expandafter\def\csname SRef@cdp:eq:C.CL.7\endcsname{\href{S6_SERIES_GRIND_II.pdf\#nameddest=tag-00C6}{II.{C.CL.7} (p.~458)} }
\expandafter\def\csname SRef@cdp:eq:C.CL.8\endcsname{\href{S6_SERIES_GRIND_II.pdf\#nameddest=tag-00C7}{II.{C.CL.8} (p.~458)} }
\expandafter\def\csname SRef@cdp:eq:C.CL.9\endcsname{\href{S6_SERIES_GRIND_II.pdf\#nameddest=tag-00C8}{II.{C.CL.9} (p.~459)} }
\expandafter\def\csname SRef@cdp:eq:C.CL.10\endcsname{\href{S6_SERIES_GRIND_II.pdf\#nameddest=tag-00BX}{II.{C.CL.10} (p.~460)} }
\expandafter\def\csname SRef@cdp:eq:C.CL.11\endcsname{\href{S6_SERIES_GRIND_II.pdf\#nameddest=tag-00BY}{II.{C.CL.11} (p.~460)} }
\expandafter\def\csname SRef@cdp:eq:C.CL.12\endcsname{\href{S6_SERIES_GRIND_II.pdf\#nameddest=tag-00BZ}{II.{C.CL.12} (p.~460)} }
\expandafter\def\csname SRef@cdp:eq:C.CL.13\endcsname{\href{S6_SERIES_GRIND_II.pdf\#nameddest=tag-00C0}{II.{C.CL.13} (p.~460)} }
\expandafter\def\csname SRef@cdp:eq:C.HB.1\endcsname{\href{S6_SERIES_GRIND_II.pdf\#nameddest=tag-00LF}{II.{C.HB.1} (p.~461)} }
\expandafter\def\csname SRef@cdp:eq:C.HB.2\endcsname{\href{S6_SERIES_GRIND_II.pdf\#nameddest=tag-00LG}{II.{C.HB.2} (p.~461)} }
\expandafter\def\csname SRef@cdp:eq:C.HB.3\endcsname{\href{S6_SERIES_GRIND_II.pdf\#nameddest=tag-00LH}{II.{C.HB.3} (p.~461)} }
\expandafter\def\csname SRef@cdp:eq:C.UC.1\endcsname{\href{S6_SERIES_GRIND_II.pdf\#nameddest=tag-014C}{II.{C.UC.1} (p.~416)} }
\expandafter\def\csname SRef@cdp:eq:C.HB.4\endcsname{\href{S6_SERIES_GRIND_II.pdf\#nameddest=tag-00LI}{II.{C.HB.4} (p.~461)} }
\expandafter\def\csname SRef@cdp:eq:C.HB.5\endcsname{\href{S6_SERIES_GRIND_II.pdf\#nameddest=tag-00LJ}{II.{C.HB.5} (p.~461)} }
\expandafter\def\csname SRef@cdp:eq:C.HB.6\endcsname{\href{S6_SERIES_GRIND_II.pdf\#nameddest=tag-00LK}{II.{C.HB.6} (p.~461)} }
\expandafter\def\csname SRef@cdp:eq:C.HB.7\endcsname{\href{S6_SERIES_GRIND_II.pdf\#nameddest=tag-00LL}{II.{C.HB.7} (p.~462)} }
\expandafter\def\csname SRef@cdp:eq:C.HB.8\endcsname{\href{S6_SERIES_GRIND_II.pdf\#nameddest=tag-00LM}{II.{C.HB.8} (p.~462)} }
\expandafter\def\csname SRef@cdp:eq:C.HB.9\endcsname{\href{S6_SERIES_GRIND_II.pdf\#nameddest=tag-00LN}{II.{C.HB.9} (p.~462)} }
\expandafter\def\csname SRef@cdp:eq:C.ZC.1\endcsname{\href{S6_SERIES_GRIND_II.pdf\#nameddest=tag-018W}{II.{C.ZC.1} (p.~463)} }
\expandafter\def\csname SRef@cdp:eq:C.ZC.2\endcsname{\href{S6_SERIES_GRIND_II.pdf\#nameddest=tag-018X}{II.{C.ZC.2} (p.~463)} }
\expandafter\def\csname SRef@cdp:eq:C.ZC.3\endcsname{\href{S6_SERIES_GRIND_II.pdf\#nameddest=tag-018Y}{II.{C.ZC.3} (p.~463)} }
\expandafter\def\csname SRef@cdp:eq:C.ZC.4\endcsname{\href{S6_SERIES_GRIND_II.pdf\#nameddest=tag-018Z}{II.{C.ZC.4} (p.~464)} }
\expandafter\def\csname SRef@cdp:eq:C.UC.2\endcsname{\href{S6_SERIES_GRIND_II.pdf\#nameddest=tag-014E}{II.{C.UC.2} (p.~416)} }
\expandafter\def\csname SRef@cdp:eq:C.ZC.5\endcsname{\href{S6_SERIES_GRIND_II.pdf\#nameddest=tag-0190}{II.{C.ZC.5} (p.~464)} }
\expandafter\def\csname SRef@cdp:eq:C.ZC.6\endcsname{\href{S6_SERIES_GRIND_II.pdf\#nameddest=tag-0191}{II.{C.ZC.6} (p.~465)} }
\expandafter\def\csname SRef@cdp:eq:C.ZC.7\endcsname{\href{S6_SERIES_GRIND_II.pdf\#nameddest=tag-0192}{II.{C.ZC.7} (p.~465)} }
\expandafter\def\csname SRef@cdp:eq:C.ZC.8\endcsname{\href{S6_SERIES_GRIND_II.pdf\#nameddest=tag-0193}{II.{C.ZC.8} (p.~465)} }
\expandafter\def\csname SRef@cdp:eq:C.ZC.9\endcsname{\href{S6_SERIES_GRIND_II.pdf\#nameddest=tag-0194}{II.{C.ZC.9} (p.~465)} }
\expandafter\def\csname SRef@cdp:eq:C.FG.1\endcsname{\href{S6_SERIES_GRIND_II.pdf\#nameddest=tag-00K8}{II.{C.FG.1} (p.~466)} }
\expandafter\def\csname SRef@cdp:eq:C.FG.2\endcsname{\href{S6_SERIES_GRIND_II.pdf\#nameddest=tag-00KJ}{II.{C.FG.2} (p.~466)} }
\expandafter\def\csname SRef@cdp:eq:C.FG.3\endcsname{\href{S6_SERIES_GRIND_II.pdf\#nameddest=tag-00KK}{II.{C.FG.3} (p.~466)} }
\expandafter\def\csname SRef@cdp:eq:C.FG.4\endcsname{\href{S6_SERIES_GRIND_II.pdf\#nameddest=tag-00KL}{II.{C.FG.4} (p.~466)} }
\expandafter\def\csname SRef@cdp:eq:C.FG.5\endcsname{\href{S6_SERIES_GRIND_II.pdf\#nameddest=tag-00KM}{II.{C.FG.5} (p.~466)} }
\expandafter\def\csname SRef@cdp:eq:C.UC.3\endcsname{\href{S6_SERIES_GRIND_II.pdf\#nameddest=tag-014F}{II.{C.UC.3} (p.~416)} }
\expandafter\def\csname SRef@cdp:eq:C.FG.6\endcsname{\href{S6_SERIES_GRIND_II.pdf\#nameddest=tag-00KN}{II.{C.FG.6} (p.~467)} }
\expandafter\def\csname SRef@cdp:eq:C.FG.7\endcsname{\href{S6_SERIES_GRIND_II.pdf\#nameddest=tag-00KO}{II.{C.FG.7} (p.~467)} }
\expandafter\def\csname SRef@cdp:eq:C.FG.8\endcsname{\href{S6_SERIES_GRIND_II.pdf\#nameddest=tag-00KP}{II.{C.FG.8} (p.~467)} }
\expandafter\def\csname SRef@cdp:eq:C.FG.9\endcsname{\href{S6_SERIES_GRIND_II.pdf\#nameddest=tag-00KQ}{II.{C.FG.9} (p.~467)} }
\expandafter\def\csname SRef@cdp:eq:C.FG.10\endcsname{\href{S6_SERIES_GRIND_II.pdf\#nameddest=tag-00K9}{II.{C.FG.10} (p.~468)} }
\expandafter\def\csname SRef@cdp:eq:C.FG.11\endcsname{\href{S6_SERIES_GRIND_II.pdf\#nameddest=tag-00KA}{II.{C.FG.11} (p.~468)} }
\expandafter\def\csname SRef@cdp:eq:C.FG.12\endcsname{\href{S6_SERIES_GRIND_II.pdf\#nameddest=tag-00KB}{II.{C.FG.12} (p.~468)} }
\expandafter\def\csname SRef@cdp:eq:C.FG.13\endcsname{\href{S6_SERIES_GRIND_II.pdf\#nameddest=tag-00KC}{II.{C.FG.13} (p.~469)} }
\expandafter\def\csname SRef@cdp:eq:C.FG.14\endcsname{\href{S6_SERIES_GRIND_II.pdf\#nameddest=tag-00KD}{II.{C.FG.14} (p.~469)} }
\expandafter\def\csname SRef@cdp:eq:C.FG.15\endcsname{\href{S6_SERIES_GRIND_II.pdf\#nameddest=tag-00KE}{II.{C.FG.15} (p.~469)} }
\expandafter\def\csname SRef@cdp:eq:C.UC.4\endcsname{\href{S6_SERIES_GRIND_II.pdf\#nameddest=tag-014G}{II.{C.UC.4} (p.~416)} }
\expandafter\def\csname SRef@cdp:eq:C.FG.16\endcsname{\href{S6_SERIES_GRIND_II.pdf\#nameddest=tag-00KF}{II.{C.FG.16} (p.~469)} }
\expandafter\def\csname SRef@cdp:eq:C.FG.17\endcsname{\href{S6_SERIES_GRIND_II.pdf\#nameddest=tag-00KG}{II.{C.FG.17} (p.~470)} }
\expandafter\def\csname SRef@cdp:eq:C.FG.18\endcsname{\href{S6_SERIES_GRIND_II.pdf\#nameddest=tag-00KH}{II.{C.FG.18} (p.~470)} }
\expandafter\def\csname SRef@cdp:eq:C.FG.19\endcsname{\href{S6_SERIES_GRIND_II.pdf\#nameddest=tag-00KI}{II.{C.FG.19} (p.~470)} }
\expandafter\def\csname SRef@cdp:eq:C.PF.1\endcsname{\href{S6_SERIES_GRIND_II.pdf\#nameddest=tag-00UN}{II.{C.PF.1} (p.~471)} }
\expandafter\def\csname SRef@cdp:eq:C.PF.2\endcsname{\href{S6_SERIES_GRIND_II.pdf\#nameddest=tag-00UR}{II.{C.PF.2} (p.~471)} }
\expandafter\def\csname SRef@cdp:eq:C.PF.3\endcsname{\href{S6_SERIES_GRIND_II.pdf\#nameddest=tag-00US}{II.{C.PF.3} (p.~471)} }
\expandafter\def\csname SRef@cdp:eq:C.PF.4\endcsname{\href{S6_SERIES_GRIND_II.pdf\#nameddest=tag-00UT}{II.{C.PF.4} (p.~471)} }
\expandafter\def\csname SRef@cdp:eq:C.PF.5\endcsname{\href{S6_SERIES_GRIND_II.pdf\#nameddest=tag-00UU}{II.{C.PF.5} (p.~472)} }
\expandafter\def\csname SRef@cdp:eq:C.PF.6\endcsname{\href{S6_SERIES_GRIND_II.pdf\#nameddest=tag-00UV}{II.{C.PF.6} (p.~472)} }
\expandafter\def\csname SRef@cdp:eq:C.UC.5\endcsname{\href{S6_SERIES_GRIND_II.pdf\#nameddest=tag-014H}{II.{C.UC.5} (p.~417)} }
\expandafter\def\csname SRef@cdp:eq:C.PF.7\endcsname{\href{S6_SERIES_GRIND_II.pdf\#nameddest=tag-00UW}{II.{C.PF.7} (p.~472)} }
\expandafter\def\csname SRef@cdp:eq:C.PF.8\endcsname{\href{S6_SERIES_GRIND_II.pdf\#nameddest=tag-00UX}{II.{C.PF.8} (p.~472)} }
\expandafter\def\csname SRef@cdp:eq:C.PF.9\endcsname{\href{S6_SERIES_GRIND_II.pdf\#nameddest=tag-00UY}{II.{C.PF.9} (p.~473)} }
\expandafter\def\csname SRef@cdp:eq:C.PF.10\endcsname{\href{S6_SERIES_GRIND_II.pdf\#nameddest=tag-00UO}{II.{C.PF.10} (p.~473)} }
\expandafter\def\csname SRef@cdp:eq:C.PF.11\endcsname{\href{S6_SERIES_GRIND_II.pdf\#nameddest=tag-00UP}{II.{C.PF.11} (p.~473)} }
\expandafter\def\csname SRef@cdp:eq:C.PF.12\endcsname{\href{S6_SERIES_GRIND_II.pdf\#nameddest=tag-00UQ}{II.{C.PF.12} (p.~473)} }
\expandafter\def\csname SRef@cdp:eq:C.UNN.1\endcsname{\href{S6_SERIES_GRIND_II.pdf\#nameddest=tag-015B}{II.{C.UNN.1} (p.~474)} }
\expandafter\def\csname SRef@cdp:eq:C.UNN.2\endcsname{\href{S6_SERIES_GRIND_II.pdf\#nameddest=tag-015C}{II.{C.UNN.2} (p.~475)} }
\expandafter\def\csname SRef@cdp:eq:C.CR.1\endcsname{\href{S6_SERIES_GRIND_II.pdf\#nameddest=tag-00CV}{II.{C.CR.1} (p.~476)} }
\expandafter\def\csname SRef@cdp:eq:C.CR.2\endcsname{\href{S6_SERIES_GRIND_II.pdf\#nameddest=tag-00D6}{II.{C.CR.2} (p.~476)} }
\expandafter\def\csname SRef@cdp:eq:C.UC.6\endcsname{\href{S6_SERIES_GRIND_II.pdf\#nameddest=tag-014I}{II.{C.UC.6} (p.~417)} }
\expandafter\def\csname SRef@cdp:eq:C.CR.3\endcsname{\href{S6_SERIES_GRIND_II.pdf\#nameddest=tag-00DA}{II.{C.CR.3} (p.~476)} }
\expandafter\def\csname SRef@cdp:eq:C.CR.4\endcsname{\href{S6_SERIES_GRIND_II.pdf\#nameddest=tag-00DB}{II.{C.CR.4} (p.~476)} }
\expandafter\def\csname SRef@cdp:eq:C.CR.5\endcsname{\href{S6_SERIES_GRIND_II.pdf\#nameddest=tag-00DC}{II.{C.CR.5} (p.~477)} }
\expandafter\def\csname SRef@cdp:eq:C.CR.6\endcsname{\href{S6_SERIES_GRIND_II.pdf\#nameddest=tag-00DD}{II.{C.CR.6} (p.~477)} }
\expandafter\def\csname SRef@cdp:eq:C.CR.7\endcsname{\href{S6_SERIES_GRIND_II.pdf\#nameddest=tag-00DE}{II.{C.CR.7} (p.~477)} }
\expandafter\def\csname SRef@cdp:eq:C.CR.8\endcsname{\href{S6_SERIES_GRIND_II.pdf\#nameddest=tag-00DF}{II.{C.CR.8} (p.~478)} }
\expandafter\def\csname SRef@cdp:eq:C.CR.9\endcsname{\href{S6_SERIES_GRIND_II.pdf\#nameddest=tag-00DG}{II.{C.CR.9} (p.~478)} }
\expandafter\def\csname SRef@cdp:eq:C.CR.10\endcsname{\href{S6_SERIES_GRIND_II.pdf\#nameddest=tag-00CW}{II.{C.CR.10} (p.~479)} }
\expandafter\def\csname SRef@cdp:eq:C.CR.11\endcsname{\href{S6_SERIES_GRIND_II.pdf\#nameddest=tag-00CX}{II.{C.CR.11} (p.~479)} }
\expandafter\def\csname SRef@cdp:eq:C.CR.12\endcsname{\href{S6_SERIES_GRIND_II.pdf\#nameddest=tag-00CY}{II.{C.CR.12} (p.~479)} }
\expandafter\def\csname SRef@cdp:eq:C.UC.7\endcsname{\href{S6_SERIES_GRIND_II.pdf\#nameddest=tag-014J}{II.{C.UC.7} (p.~417)} }
\expandafter\def\csname SRef@cdp:eq:C.CR.13\endcsname{\href{S6_SERIES_GRIND_II.pdf\#nameddest=tag-00CZ}{II.{C.CR.13} (p.~479)} }
\expandafter\def\csname SRef@cdp:eq:C.CR.14\endcsname{\href{S6_SERIES_GRIND_II.pdf\#nameddest=tag-00D0}{II.{C.CR.14} (p.~479)} }
\expandafter\def\csname SRef@cdp:eq:C.CR.15\endcsname{\href{S6_SERIES_GRIND_II.pdf\#nameddest=tag-00D1}{II.{C.CR.15} (p.~481)} }
\expandafter\def\csname SRef@cdp:eq:C.CR.16\endcsname{\href{S6_SERIES_GRIND_II.pdf\#nameddest=tag-00D2}{II.{C.CR.16} (p.~481)} }
\expandafter\def\csname SRef@cdp:eq:C.CR.17\endcsname{\href{S6_SERIES_GRIND_II.pdf\#nameddest=tag-00D3}{II.{C.CR.17} (p.~481)} }
\expandafter\def\csname SRef@cdp:eq:C.CR.18\endcsname{\href{S6_SERIES_GRIND_II.pdf\#nameddest=tag-00D4}{II.{C.CR.18} (p.~482)} }
\expandafter\def\csname SRef@cdp:eq:C.CR.19\endcsname{\href{S6_SERIES_GRIND_II.pdf\#nameddest=tag-00D5}{II.{C.CR.19} (p.~483)} }
\expandafter\def\csname SRef@cdp:eq:C.CR.20\endcsname{\href{S6_SERIES_GRIND_II.pdf\#nameddest=tag-00D7}{II.{C.CR.20} (p.~483)} }
\expandafter\def\csname SRef@cdp:eq:C.CR.21\endcsname{\href{S6_SERIES_GRIND_II.pdf\#nameddest=tag-00D8}{II.{C.CR.21} (p.~483)} }
\expandafter\def\csname SRef@cdp:eq:C.CR.22\endcsname{\href{S6_SERIES_GRIND_II.pdf\#nameddest=tag-00D9}{II.{C.CR.22} (p.~484)} }
\expandafter\def\csname SRef@cdp:eq:C.UC.8\endcsname{\href{S6_SERIES_GRIND_II.pdf\#nameddest=tag-014K}{II.{C.UC.8} (p.~417)} }
\expandafter\def\csname SRef@cdp:eq:C.UNNODE.1\endcsname{\href{S6_SERIES_GRIND_II.pdf\#nameddest=tag-015D}{II.{C.UNNODE.1} (p.~484)} }
\expandafter\def\csname SRef@cdp:eq:C.UNNODE.2\endcsname{\href{S6_SERIES_GRIND_II.pdf\#nameddest=tag-015E}{II.{C.UNNODE.2} (p.~485)} }
\expandafter\def\csname SRef@cdp:eq:C.UNNODE.3\endcsname{\href{S6_SERIES_GRIND_II.pdf\#nameddest=tag-015F}{II.{C.UNNODE.3} (p.~485)} }
\expandafter\def\csname SRef@cdp:eq:C.UNCR.1\endcsname{\href{S6_SERIES_GRIND_II.pdf\#nameddest=tag-0158}{II.{C.UNCR.1} (p.~486)} }
\expandafter\def\csname SRef@cdp:eq:C.UNCR.2\endcsname{\href{S6_SERIES_GRIND_II.pdf\#nameddest=tag-0159}{II.{C.UNCR.2} (p.~486)} }
\expandafter\def\csname SRef@cdp:eq:C.UNCR.3\endcsname{\href{S6_SERIES_GRIND_II.pdf\#nameddest=tag-015A}{II.{C.UNCR.3} (p.~486)} }
\expandafter\def\csname SRef@cdp:eq:C.IR.1\endcsname{\href{S6_SERIES_GRIND_II.pdf\#nameddest=tag-00OH}{II.{C.IR.1} (p.~487)} }
\expandafter\def\csname SRef@cdp:eq:C.IR.2\endcsname{\href{S6_SERIES_GRIND_II.pdf\#nameddest=tag-00OI}{II.{C.IR.2} (p.~487)} }
\expandafter\def\csname SRef@cdp:eq:C.IR.4\endcsname{\href{S6_SERIES_GRIND_II.pdf\#nameddest=tag-00OJ}{II.{C.IR.4} (p.~487)} }
\expandafter\def\csname SRef@cdp:eq:C.IR.5\endcsname{\href{S6_SERIES_GRIND_II.pdf\#nameddest=tag-00OK}{II.{C.IR.5} (p.~488)} }
\expandafter\def\csname SRef@cdp:eq:C.UC.9\endcsname{\href{S6_SERIES_GRIND_II.pdf\#nameddest=tag-014L}{II.{C.UC.9} (p.~418)} }
\expandafter\def\csname SRef@cdp:eq:C.IR.6\endcsname{\href{S6_SERIES_GRIND_II.pdf\#nameddest=tag-00OL}{II.{C.IR.6} (p.~489)} }
\expandafter\def\csname SRef@cdp:eq:C.IR.9\endcsname{\href{S6_SERIES_GRIND_II.pdf\#nameddest=tag-00OM}{II.{C.IR.9} (p.~489)} }
\expandafter\def\csname SRef@cdp:eq:C.VII.1\endcsname{\href{S6_SERIES_GRIND_II.pdf\#nameddest=tag-015K}{II.{C.VII.1} (p.~489)} }
\expandafter\def\csname SRef@cdp:eq:C.VII.2\endcsname{\href{S6_SERIES_GRIND_II.pdf\#nameddest=tag-015O}{II.{C.VII.2} (p.~490)} }
\expandafter\def\csname SRef@cdp:eq:C.VII.3\endcsname{\href{S6_SERIES_GRIND_II.pdf\#nameddest=tag-015P}{II.{C.VII.3} (p.~491)} }
\expandafter\def\csname SRef@cdp:eq:C.VII.4\endcsname{\href{S6_SERIES_GRIND_II.pdf\#nameddest=tag-015Q}{II.{C.VII.4} (p.~491)} }
\expandafter\def\csname SRef@cdp:eq:C.VII.5\endcsname{\href{S6_SERIES_GRIND_II.pdf\#nameddest=tag-015R}{II.{C.VII.5} (p.~491)} }
\expandafter\def\csname SRef@cdp:eq:C.VII.6\endcsname{\href{S6_SERIES_GRIND_II.pdf\#nameddest=tag-015S}{II.{C.VII.6} (p.~491)} }
\expandafter\def\csname SRef@cdp:eq:C.VII.7\endcsname{\href{S6_SERIES_GRIND_II.pdf\#nameddest=tag-015T}{II.{C.VII.7} (p.~492)} }
\expandafter\def\csname SRef@cdp:eq:C.VII.8\endcsname{\href{S6_SERIES_GRIND_II.pdf\#nameddest=tag-015U}{II.{C.VII.8} (p.~493)} }
\expandafter\def\csname SRef@cdp:eq:C.UC.10\endcsname{\href{S6_SERIES_GRIND_II.pdf\#nameddest=tag-014D}{II.{C.UC.10} (p.~418)} }
\expandafter\def\csname SRef@cdp:eq:C.VII.9\endcsname{\href{S6_SERIES_GRIND_II.pdf\#nameddest=tag-015V}{II.{C.VII.9} (p.~493)} }
\expandafter\def\csname SRef@cdp:eq:C.VII.10\endcsname{\href{S6_SERIES_GRIND_II.pdf\#nameddest=tag-015L}{II.{C.VII.10} (p.~493)} }
\expandafter\def\csname SRef@cdp:eq:C.VII.11\endcsname{\href{S6_SERIES_GRIND_II.pdf\#nameddest=tag-015M}{II.{C.VII.11} (p.~494)} }
\expandafter\def\csname SRef@cdp:eq:C.VII.12\endcsname{\href{S6_SERIES_GRIND_II.pdf\#nameddest=tag-015N}{II.{C.VII.12} (p.~494)} }
\expandafter\def\csname SRef@cdp:eq:C.RA.1\endcsname{\href{S6_SERIES_GRIND_II.pdf\#nameddest=tag-00WQ}{II.{C.RA.1} (p.~495)} }
\expandafter\def\csname SRef@cdp:eq:C.RA.2\endcsname{\href{S6_SERIES_GRIND_II.pdf\#nameddest=tag-00X1}{II.{C.RA.2} (p.~495)} }
\expandafter\def\csname SRef@cdp:eq:C.RA.3\endcsname{\href{S6_SERIES_GRIND_II.pdf\#nameddest=tag-00XC}{II.{C.RA.3} (p.~495)} }
\expandafter\def\csname SRef@cdp:eq:C.RA.4\endcsname{\href{S6_SERIES_GRIND_II.pdf\#nameddest=tag-00XH}{II.{C.RA.4} (p.~495)} }
\expandafter\def\csname SRef@cdp:eq:C.RA.5\endcsname{\href{S6_SERIES_GRIND_II.pdf\#nameddest=tag-00XI}{II.{C.RA.5} (p.~495)} }
\expandafter\def\csname SRef@cdp:eq:C.RA.6\endcsname{\href{S6_SERIES_GRIND_II.pdf\#nameddest=tag-00XJ}{II.{C.RA.6} (p.~495)} }
\expandafter\def\csname SRef@cdp:eq:C.VR.1\endcsname{\href{S6_SERIES_GRIND_II.pdf\#nameddest=tag-015W}{II.{C.VR.1} (p.~418)} }
\expandafter\def\csname SRef@cdp:eq:C.RA.7\endcsname{\href{S6_SERIES_GRIND_II.pdf\#nameddest=tag-00XK}{II.{C.RA.7} (p.~496)} }
\expandafter\def\csname SRef@cdp:eq:C.RA.8\endcsname{\href{S6_SERIES_GRIND_II.pdf\#nameddest=tag-00XL}{II.{C.RA.8} (p.~496)} }
\expandafter\def\csname SRef@cdp:eq:C.RA.9\endcsname{\href{S6_SERIES_GRIND_II.pdf\#nameddest=tag-00XM}{II.{C.RA.9} (p.~496)} }
\expandafter\def\csname SRef@cdp:eq:C.RA.10\endcsname{\href{S6_SERIES_GRIND_II.pdf\#nameddest=tag-00WR}{II.{C.RA.10} (p.~496)} }
\expandafter\def\csname SRef@cdp:eq:C.RA.11\endcsname{\href{S6_SERIES_GRIND_II.pdf\#nameddest=tag-00WS}{II.{C.RA.11} (p.~497)} }
\expandafter\def\csname SRef@cdp:eq:C.RA.12\endcsname{\href{S6_SERIES_GRIND_II.pdf\#nameddest=tag-00WT}{II.{C.RA.12} (p.~497)} }
\expandafter\def\csname SRef@cdp:eq:C.RA.13\endcsname{\href{S6_SERIES_GRIND_II.pdf\#nameddest=tag-00WU}{II.{C.RA.13} (p.~497)} }
\expandafter\def\csname SRef@cdp:eq:C.RA.14\endcsname{\href{S6_SERIES_GRIND_II.pdf\#nameddest=tag-00WV}{II.{C.RA.14} (p.~498)} }
\expandafter\def\csname SRef@cdp:eq:C.RA.15\endcsname{\href{S6_SERIES_GRIND_II.pdf\#nameddest=tag-00WW}{II.{C.RA.15} (p.~498)} }
\expandafter\def\csname SRef@cdp:eq:C.RA.16\endcsname{\href{S6_SERIES_GRIND_II.pdf\#nameddest=tag-00WX}{II.{C.RA.16} (p.~498)} }
\expandafter\def\csname SRef@cdp:eq:C.VR.2\endcsname{\href{S6_SERIES_GRIND_II.pdf\#nameddest=tag-015X}{II.{C.VR.2} (p.~419)} }
\expandafter\def\csname SRef@cdp:eq:C.RA.17\endcsname{\href{S6_SERIES_GRIND_II.pdf\#nameddest=tag-00WY}{II.{C.RA.17} (p.~498)} }
\expandafter\def\csname SRef@cdp:eq:C.RA.18\endcsname{\href{S6_SERIES_GRIND_II.pdf\#nameddest=tag-00WZ}{II.{C.RA.18} (p.~498)} }
\expandafter\def\csname SRef@cdp:eq:C.RA.19\endcsname{\href{S6_SERIES_GRIND_II.pdf\#nameddest=tag-00X0}{II.{C.RA.19} (p.~499)} }
\expandafter\def\csname SRef@cdp:eq:C.RA.20\endcsname{\href{S6_SERIES_GRIND_II.pdf\#nameddest=tag-00X2}{II.{C.RA.20} (p.~499)} }
\expandafter\def\csname SRef@cdp:eq:C.RA.21\endcsname{\href{S6_SERIES_GRIND_II.pdf\#nameddest=tag-00X3}{II.{C.RA.21} (p.~499)} }
\expandafter\def\csname SRef@cdp:eq:C.RA.22\endcsname{\href{S6_SERIES_GRIND_II.pdf\#nameddest=tag-00X4}{II.{C.RA.22} (p.~500)} }
\expandafter\def\csname SRef@cdp:eq:C.RA.23\endcsname{\href{S6_SERIES_GRIND_II.pdf\#nameddest=tag-00X5}{II.{C.RA.23} (p.~500)} }
\expandafter\def\csname SRef@cdp:eq:C.RA.24\endcsname{\href{S6_SERIES_GRIND_II.pdf\#nameddest=tag-00X6}{II.{C.RA.24} (p.~501)} }
\expandafter\def\csname SRef@cdp:eq:C.RA.25\endcsname{\href{S6_SERIES_GRIND_II.pdf\#nameddest=tag-00X7}{II.{C.RA.25} (p.~501)} }
\expandafter\def\csname SRef@cdp:eq:C.RA.26\endcsname{\href{S6_SERIES_GRIND_II.pdf\#nameddest=tag-00X8}{II.{C.RA.26} (p.~501)} }
\expandafter\def\csname SRef@cdp:eq:C.VR.3\endcsname{\href{S6_SERIES_GRIND_II.pdf\#nameddest=tag-015Y}{II.{C.VR.3} (p.~419)} }
\expandafter\def\csname SRef@cdp:eq:C.RA.27\endcsname{\href{S6_SERIES_GRIND_II.pdf\#nameddest=tag-00X9}{II.{C.RA.27} (p.~502)} }
\expandafter\def\csname SRef@cdp:eq:C.RA.28\endcsname{\href{S6_SERIES_GRIND_II.pdf\#nameddest=tag-00XA}{II.{C.RA.28} (p.~502)} }
\expandafter\def\csname SRef@cdp:eq:C.RA.29\endcsname{\href{S6_SERIES_GRIND_II.pdf\#nameddest=tag-00XB}{II.{C.RA.29} (p.~502)} }
\expandafter\def\csname SRef@cdp:eq:C.RA.30\endcsname{\href{S6_SERIES_GRIND_II.pdf\#nameddest=tag-00XD}{II.{C.RA.30} (p.~502)} }
\expandafter\def\csname SRef@cdp:eq:C.RA.31\endcsname{\href{S6_SERIES_GRIND_II.pdf\#nameddest=tag-00XE}{II.{C.RA.31} (p.~502)} }
\expandafter\def\csname SRef@cdp:eq:C.RA.31a\endcsname{\href{S6_SERIES_GRIND_II.pdf\#nameddest=tag-00XF}{II.{C.RA.31a} (p.~503)} }
\expandafter\def\csname SRef@cdp:eq:C.RA.32\endcsname{\href{S6_SERIES_GRIND_II.pdf\#nameddest=tag-00XG}{II.{C.RA.32} (p.~503)} }
\expandafter\def\csname SRef@cdp:eq:C.RD.1\endcsname{\href{S6_SERIES_GRIND_II.pdf\#nameddest=tag-00XZ}{II.{C.RD.1} (p.~503)} }
\expandafter\def\csname SRef@cdp:eq:C.RD.2\endcsname{\href{S6_SERIES_GRIND_II.pdf\#nameddest=tag-00Y9}{II.{C.RD.2} (p.~503)} }
\expandafter\def\csname SRef@cdp:eq:C.RD.3\endcsname{\href{S6_SERIES_GRIND_II.pdf\#nameddest=tag-00YK}{II.{C.RD.3} (p.~504)} }
\expandafter\def\csname SRef@cdp:eq:C.VR.4\endcsname{\href{S6_SERIES_GRIND_II.pdf\#nameddest=tag-015Z}{II.{C.VR.4} (p.~419)} }
\expandafter\def\csname SRef@cdp:eq:C.RD.4\endcsname{\href{S6_SERIES_GRIND_II.pdf\#nameddest=tag-00YT}{II.{C.RD.4} (p.~504)} }
\expandafter\def\csname SRef@cdp:eq:C.RD.11\endcsname{\href{S6_SERIES_GRIND_II.pdf\#nameddest=tag-00Y0}{II.{C.RD.11} (p.~504)} }
\expandafter\def\csname SRef@cdp:eq:C.RD.12\endcsname{\href{S6_SERIES_GRIND_II.pdf\#nameddest=tag-00Y1}{II.{C.RD.12} (p.~505)} }
\expandafter\def\csname SRef@cdp:eq:C.RD.13\endcsname{\href{S6_SERIES_GRIND_II.pdf\#nameddest=tag-00Y2}{II.{C.RD.13} (p.~505)} }
\expandafter\def\csname SRef@cdp:eq:C.RD.14\endcsname{\href{S6_SERIES_GRIND_II.pdf\#nameddest=tag-00Y3}{II.{C.RD.14} (p.~505)} }
\expandafter\def\csname SRef@cdp:eq:C.RD.15\endcsname{\href{S6_SERIES_GRIND_II.pdf\#nameddest=tag-00Y4}{II.{C.RD.15} (p.~506)} }
\expandafter\def\csname SRef@cdp:eq:C.RD.16\endcsname{\href{S6_SERIES_GRIND_II.pdf\#nameddest=tag-00Y5}{II.{C.RD.16} (p.~506)} }
\expandafter\def\csname SRef@cdp:eq:C.RD.17\endcsname{\href{S6_SERIES_GRIND_II.pdf\#nameddest=tag-00Y6}{II.{C.RD.17} (p.~506)} }
\expandafter\def\csname SRef@cdp:eq:C.RD.18\endcsname{\href{S6_SERIES_GRIND_II.pdf\#nameddest=tag-00Y7}{II.{C.RD.18} (p.~506)} }
\expandafter\def\csname SRef@cdp:eq:C.VR.5\endcsname{\href{S6_SERIES_GRIND_II.pdf\#nameddest=tag-0160}{II.{C.VR.5} (p.~419)} }
\expandafter\def\csname SRef@cdp:eq:C.RD.19\endcsname{\href{S6_SERIES_GRIND_II.pdf\#nameddest=tag-00Y8}{II.{C.RD.19} (p.~507)} }
\expandafter\def\csname SRef@cdp:eq:C.RD.20\endcsname{\href{S6_SERIES_GRIND_II.pdf\#nameddest=tag-00YA}{II.{C.RD.20} (p.~507)} }
\expandafter\def\csname SRef@cdp:eq:C.RD.21\endcsname{\href{S6_SERIES_GRIND_II.pdf\#nameddest=tag-00YB}{II.{C.RD.21} (p.~507)} }
\expandafter\def\csname SRef@cdp:eq:C.RD.22\endcsname{\href{S6_SERIES_GRIND_II.pdf\#nameddest=tag-00YC}{II.{C.RD.22} (p.~508)} }
\expandafter\def\csname SRef@cdp:eq:C.RD.23\endcsname{\href{S6_SERIES_GRIND_II.pdf\#nameddest=tag-00YD}{II.{C.RD.23} (p.~508)} }
\expandafter\def\csname SRef@cdp:eq:C.RD.24\endcsname{\href{S6_SERIES_GRIND_II.pdf\#nameddest=tag-00YE}{II.{C.RD.24} (p.~508)} }
\expandafter\def\csname SRef@cdp:eq:C.RD.25\endcsname{\href{S6_SERIES_GRIND_II.pdf\#nameddest=tag-00YF}{II.{C.RD.25} (p.~508)} }
\expandafter\def\csname SRef@cdp:eq:C.RD.26\endcsname{\href{S6_SERIES_GRIND_II.pdf\#nameddest=tag-00YG}{II.{C.RD.26} (p.~508)} }
\expandafter\def\csname SRef@cdp:eq:C.RD.27\endcsname{\href{S6_SERIES_GRIND_II.pdf\#nameddest=tag-00YH}{II.{C.RD.27} (p.~508)} }
\expandafter\def\csname SRef@cdp:eq:C.VR.6\endcsname{\href{S6_SERIES_GRIND_II.pdf\#nameddest=tag-0161}{II.{C.VR.6} (p.~419)} }
\expandafter\def\csname SRef@cdp:eq:C.RD.28\endcsname{\href{S6_SERIES_GRIND_II.pdf\#nameddest=tag-00YI}{II.{C.RD.28} (p.~509)} }
\expandafter\def\csname SRef@cdp:eq:C.RD.29\endcsname{\href{S6_SERIES_GRIND_II.pdf\#nameddest=tag-00YJ}{II.{C.RD.29} (p.~509)} }
\expandafter\def\csname SRef@cdp:eq:C.RD.30\endcsname{\href{S6_SERIES_GRIND_II.pdf\#nameddest=tag-00YL}{II.{C.RD.30} (p.~509)} }
\expandafter\def\csname SRef@cdp:eq:C.RD.31\endcsname{\href{S6_SERIES_GRIND_II.pdf\#nameddest=tag-00YM}{II.{C.RD.31} (p.~510)} }
\expandafter\def\csname SRef@cdp:eq:C.RD.32\endcsname{\href{S6_SERIES_GRIND_II.pdf\#nameddest=tag-00YN}{II.{C.RD.32} (p.~510)} }
\expandafter\def\csname SRef@cdp:eq:C.RD.33\endcsname{\href{S6_SERIES_GRIND_II.pdf\#nameddest=tag-00YO}{II.{C.RD.33} (p.~510)} }
\expandafter\def\csname SRef@cdp:eq:C.RD.34\endcsname{\href{S6_SERIES_GRIND_II.pdf\#nameddest=tag-00YP}{II.{C.RD.34} (p.~510)} }
\expandafter\def\csname SRef@cdp:eq:C.RD.35\endcsname{\href{S6_SERIES_GRIND_II.pdf\#nameddest=tag-00YQ}{II.{C.RD.35} (p.~510)} }
\expandafter\def\csname SRef@cdp:eq:C.VR.7\endcsname{\href{S6_SERIES_GRIND_II.pdf\#nameddest=tag-0162}{II.{C.VR.7} (p.~419)} }
\expandafter\def\csname SRef@cdp:eq:C.RD.36\endcsname{\href{S6_SERIES_GRIND_II.pdf\#nameddest=tag-00YR}{II.{C.RD.36} (p.~511)} }
\expandafter\def\csname SRef@cdp:eq:C.RD.37\endcsname{\href{S6_SERIES_GRIND_II.pdf\#nameddest=tag-00YS}{II.{C.RD.37} (p.~511)} }
\expandafter\def\csname SRef@cdp:eq:C.D4.1\endcsname{\href{S6_SERIES_GRIND_II.pdf\#nameddest=tag-00E1}{II.{C.D4.1} (p.~512)} }
\expandafter\def\csname SRef@cdp:eq:C.D4.2\endcsname{\href{S6_SERIES_GRIND_II.pdf\#nameddest=tag-00EC}{II.{C.D4.2} (p.~512)} }
\expandafter\def\csname SRef@cdp:eq:C.D4.3\endcsname{\href{S6_SERIES_GRIND_II.pdf\#nameddest=tag-00EE}{II.{C.D4.3} (p.~513)} }
\expandafter\def\csname SRef@cdp:eq:C.D4.4\endcsname{\href{S6_SERIES_GRIND_II.pdf\#nameddest=tag-00EF}{II.{C.D4.4} (p.~513)} }
\expandafter\def\csname SRef@cdp:eq:C.D4.5\endcsname{\href{S6_SERIES_GRIND_II.pdf\#nameddest=tag-00EG}{II.{C.D4.5} (p.~513)} }
\expandafter\def\csname SRef@cdp:eq:C.VR.8\endcsname{\href{S6_SERIES_GRIND_II.pdf\#nameddest=tag-0163}{II.{C.VR.8} (p.~420)} }
\expandafter\def\csname SRef@cdp:eq:C.D4.6\endcsname{\href{S6_SERIES_GRIND_II.pdf\#nameddest=tag-00EH}{II.{C.D4.6} (p.~513)} }
\expandafter\def\csname SRef@cdp:eq:C.D4.7\endcsname{\href{S6_SERIES_GRIND_II.pdf\#nameddest=tag-00EI}{II.{C.D4.7} (p.~513)} }
\expandafter\def\csname SRef@cdp:eq:C.D4.8\endcsname{\href{S6_SERIES_GRIND_II.pdf\#nameddest=tag-00EJ}{II.{C.D4.8} (p.~515)} }
\expandafter\def\csname SRef@cdp:eq:C.D4.9\endcsname{\href{S6_SERIES_GRIND_II.pdf\#nameddest=tag-00EK}{II.{C.D4.9} (p.~515)} }
\expandafter\def\csname SRef@cdp:eq:C.D4.10\endcsname{\href{S6_SERIES_GRIND_II.pdf\#nameddest=tag-00E2}{II.{C.D4.10} (p.~515)} }
\expandafter\def\csname SRef@cdp:eq:C.D4.11\endcsname{\href{S6_SERIES_GRIND_II.pdf\#nameddest=tag-00E3}{II.{C.D4.11} (p.~515)} }
\expandafter\def\csname SRef@cdp:eq:C.D4.12\endcsname{\href{S6_SERIES_GRIND_II.pdf\#nameddest=tag-00E4}{II.{C.D4.12} (p.~515)} }
\expandafter\def\csname SRef@cdp:eq:C.DD.1\endcsname{\href{S6_SERIES_GRIND_II.pdf\#nameddest=tag-00EL}{II.{C.DD.1} (p.~516)} }
\expandafter\def\csname SRef@cdp:eq:C.DD.2\endcsname{\href{S6_SERIES_GRIND_II.pdf\#nameddest=tag-00EW}{II.{C.DD.2} (p.~516)} }
\expandafter\def\csname SRef@cdp:eq:C.VR.9\endcsname{\href{S6_SERIES_GRIND_II.pdf\#nameddest=tag-0164}{II.{C.VR.9} (p.~420)} }
\expandafter\def\csname SRef@cdp:eq:C.DD.3\endcsname{\href{S6_SERIES_GRIND_II.pdf\#nameddest=tag-00EZ}{II.{C.DD.3} (p.~516)} }
\expandafter\def\csname SRef@cdp:eq:C.DD.4\endcsname{\href{S6_SERIES_GRIND_II.pdf\#nameddest=tag-00F0}{II.{C.DD.4} (p.~516)} }
\expandafter\def\csname SRef@cdp:eq:C.DD.5\endcsname{\href{S6_SERIES_GRIND_II.pdf\#nameddest=tag-00F1}{II.{C.DD.5} (p.~517)} }
\expandafter\def\csname SRef@cdp:eq:C.DD.6\endcsname{\href{S6_SERIES_GRIND_II.pdf\#nameddest=tag-00F2}{II.{C.DD.6} (p.~517)} }
\expandafter\def\csname SRef@cdp:eq:C.DD.7\endcsname{\href{S6_SERIES_GRIND_II.pdf\#nameddest=tag-00F3}{II.{C.DD.7} (p.~517)} }
\expandafter\def\csname SRef@cdp:eq:C.DD.8\endcsname{\href{S6_SERIES_GRIND_II.pdf\#nameddest=tag-00F4}{II.{C.DD.8} (p.~517)} }
\expandafter\def\csname SRef@cdp:eq:C.DD.9\endcsname{\href{S6_SERIES_GRIND_II.pdf\#nameddest=tag-00F5}{II.{C.DD.9} (p.~517)} }
\expandafter\def\csname SRef@cdp:eq:C.DD.10\endcsname{\href{S6_SERIES_GRIND_II.pdf\#nameddest=tag-00EM}{II.{C.DD.10} (p.~518)} }
\expandafter\def\csname SRef@cdp:eq:C.EC.1\endcsname{\href{S6_SERIES_GRIND_II.pdf\#nameddest=tag-00FO}{II.{C.EC.1} (p.~421)} }
\expandafter\def\csname SRef@cdp:eq:C.DD.11\endcsname{\href{S6_SERIES_GRIND_II.pdf\#nameddest=tag-00EN}{II.{C.DD.11} (p.~518)} }
\expandafter\def\csname SRef@cdp:eq:C.DD.12\endcsname{\href{S6_SERIES_GRIND_II.pdf\#nameddest=tag-00EO}{II.{C.DD.12} (p.~518)} }
\expandafter\def\csname SRef@cdp:eq:C.DD.13\endcsname{\href{S6_SERIES_GRIND_II.pdf\#nameddest=tag-00EP}{II.{C.DD.13} (p.~518)} }
\expandafter\def\csname SRef@cdp:eq:C.DD.14\endcsname{\href{S6_SERIES_GRIND_II.pdf\#nameddest=tag-00EQ}{II.{C.DD.14} (p.~518)} }
\expandafter\def\csname SRef@cdp:eq:C.DD.15\endcsname{\href{S6_SERIES_GRIND_II.pdf\#nameddest=tag-00ER}{II.{C.DD.15} (p.~518)} }
\expandafter\def\csname SRef@cdp:eq:C.DD.16\endcsname{\href{S6_SERIES_GRIND_II.pdf\#nameddest=tag-00ES}{II.{C.DD.16} (p.~519)} }
\expandafter\def\csname SRef@cdp:eq:C.EC.2\endcsname{\href{S6_SERIES_GRIND_II.pdf\#nameddest=tag-00FP}{II.{C.EC.2} (p.~422)} }
\expandafter\def\csname SRef@cdp:eq:C.DD.17\endcsname{\href{S6_SERIES_GRIND_II.pdf\#nameddest=tag-00ET}{II.{C.DD.17} (p.~519)} }
\expandafter\def\csname SRef@cdp:eq:C.DD.18\endcsname{\href{S6_SERIES_GRIND_II.pdf\#nameddest=tag-00EU}{II.{C.DD.18} (p.~519)} }
\expandafter\def\csname SRef@cdp:eq:C.DD.19\endcsname{\href{S6_SERIES_GRIND_II.pdf\#nameddest=tag-00EV}{II.{C.DD.19} (p.~519)} }
\expandafter\def\csname SRef@cdp:eq:C.DD.20\endcsname{\href{S6_SERIES_GRIND_II.pdf\#nameddest=tag-00EX}{II.{C.DD.20} (p.~519)} }
\expandafter\def\csname SRef@cdp:eq:C.DD.21\endcsname{\href{S6_SERIES_GRIND_II.pdf\#nameddest=tag-00EY}{II.{C.DD.21} (p.~520)} }
\expandafter\def\csname SRef@cdp:eq:C.D4.13\endcsname{\href{S6_SERIES_GRIND_II.pdf\#nameddest=tag-00E5}{II.{C.D4.13} (p.~520)} }
\expandafter\def\csname SRef@cdp:eq:C.D4.14\endcsname{\href{S6_SERIES_GRIND_II.pdf\#nameddest=tag-00E6}{II.{C.D4.14} (p.~521)} }
\expandafter\def\csname SRef@cdp:eq:C.D4.15\endcsname{\href{S6_SERIES_GRIND_II.pdf\#nameddest=tag-00E7}{II.{C.D4.15} (p.~521)} }
\expandafter\def\csname SRef@cdp:eq:C.SC.1\endcsname{\href{S6_SERIES_GRIND_II.pdf\#nameddest=tag-010L}{II.{C.SC.1} (p.~423)} }
\expandafter\def\csname SRef@cdp:eq:C.D4.16\endcsname{\href{S6_SERIES_GRIND_II.pdf\#nameddest=tag-00E8}{II.{C.D4.16} (p.~521)} }
\expandafter\def\csname SRef@cdp:eq:C.D4.17\endcsname{\href{S6_SERIES_GRIND_II.pdf\#nameddest=tag-00E9}{II.{C.D4.17} (p.~521)} }
\expandafter\def\csname SRef@cdp:eq:C.D4.18\endcsname{\href{S6_SERIES_GRIND_II.pdf\#nameddest=tag-00EA}{II.{C.D4.18} (p.~522)} }
\expandafter\def\csname SRef@cdp:eq:C.D4.19\endcsname{\href{S6_SERIES_GRIND_II.pdf\#nameddest=tag-00EB}{II.{C.D4.19} (p.~522)} }
\expandafter\def\csname SRef@cdp:eq:C.D4.20\endcsname{\href{S6_SERIES_GRIND_II.pdf\#nameddest=tag-00ED}{II.{C.D4.20} (p.~522)} }
\expandafter\def\csname SRef@cdp:eq:C.UI.1\endcsname{\href{S6_SERIES_GRIND_II.pdf\#nameddest=tag-014M}{II.{C.UI.1} (p.~522)} }
\expandafter\def\csname SRef@cdp:eq:C.UI.2\endcsname{\href{S6_SERIES_GRIND_II.pdf\#nameddest=tag-014U}{II.{C.UI.2} (p.~523)} }
\expandafter\def\csname SRef@cdp:eq:C.UI.3\endcsname{\href{S6_SERIES_GRIND_II.pdf\#nameddest=tag-014V}{II.{C.UI.3} (p.~523)} }
\expandafter\def\csname SRef@cdp:eq:C.UI.4\endcsname{\href{S6_SERIES_GRIND_II.pdf\#nameddest=tag-014W}{II.{C.UI.4} (p.~524)} }
\expandafter\def\csname SRef@cdp:eq:C.SC.2\endcsname{\href{S6_SERIES_GRIND_II.pdf\#nameddest=tag-010M}{II.{C.SC.2} (p.~423)} }
\expandafter\def\csname SRef@cdp:eq:C.UI.5\endcsname{\href{S6_SERIES_GRIND_II.pdf\#nameddest=tag-014X}{II.{C.UI.5} (p.~524)} }
\expandafter\def\csname SRef@cdp:eq:C.UI.6\endcsname{\href{S6_SERIES_GRIND_II.pdf\#nameddest=tag-014Y}{II.{C.UI.6} (p.~524)} }
\expandafter\def\csname SRef@cdp:eq:C.UI.7\endcsname{\href{S6_SERIES_GRIND_II.pdf\#nameddest=tag-014Z}{II.{C.UI.7} (p.~524)} }
\expandafter\def\csname SRef@cdp:eq:C.UI.8\endcsname{\href{S6_SERIES_GRIND_II.pdf\#nameddest=tag-0150}{II.{C.UI.8} (p.~525)} }
\expandafter\def\csname SRef@cdp:eq:C.UI.9\endcsname{\href{S6_SERIES_GRIND_II.pdf\#nameddest=tag-0151}{II.{C.UI.9} (p.~525)} }
\expandafter\def\csname SRef@cdp:eq:C.UI.10\endcsname{\href{S6_SERIES_GRIND_II.pdf\#nameddest=tag-014N}{II.{C.UI.10} (p.~526)} }
\expandafter\def\csname SRef@cdp:eq:C.UI.11\endcsname{\href{S6_SERIES_GRIND_II.pdf\#nameddest=tag-014O}{II.{C.UI.11} (p.~526)} }
\expandafter\def\csname SRef@cdp:eq:C.UI.12\endcsname{\href{S6_SERIES_GRIND_II.pdf\#nameddest=tag-014P}{II.{C.UI.12} (p.~526)} }
\expandafter\def\csname SRef@cdp:eq:C.UI.13\endcsname{\href{S6_SERIES_GRIND_II.pdf\#nameddest=tag-014Q}{II.{C.UI.13} (p.~527)} }
\expandafter\def\csname SRef@cdp:eq:C.UI.14\endcsname{\href{S6_SERIES_GRIND_II.pdf\#nameddest=tag-014R}{II.{C.UI.14} (p.~528)} }
\expandafter\def\csname SRef@cdp:eq:C.SC.3\endcsname{\href{S6_SERIES_GRIND_II.pdf\#nameddest=tag-010N}{II.{C.SC.3} (p.~423)} }
\expandafter\def\csname SRef@cdp:eq:C.UI.15\endcsname{\href{S6_SERIES_GRIND_II.pdf\#nameddest=tag-014S}{II.{C.UI.15} (p.~528)} }
\expandafter\def\csname SRef@cdp:eq:C.UI.16\endcsname{\href{S6_SERIES_GRIND_II.pdf\#nameddest=tag-014T}{II.{C.UI.16} (p.~528)} }
\expandafter\def\csname SRef@cdp:eq:C.IZ.1\endcsname{\href{S6_SERIES_GRIND_II.pdf\#nameddest=tag-00ON}{II.{C.IZ.1} (p.~529)} }
\expandafter\def\csname SRef@cdp:eq:C.IZ.2\endcsname{\href{S6_SERIES_GRIND_II.pdf\#nameddest=tag-00OO}{II.{C.IZ.2} (p.~529)} }
\expandafter\def\csname SRef@cdp:eq:C.IZ.3\endcsname{\href{S6_SERIES_GRIND_II.pdf\#nameddest=tag-00OP}{II.{C.IZ.3} (p.~530)} }
\expandafter\def\csname SRef@cdp:eq:C.IZ.4\endcsname{\href{S6_SERIES_GRIND_II.pdf\#nameddest=tag-00OQ}{II.{C.IZ.4} (p.~530)} }
\expandafter\def\csname SRef@cdp:eq:C.IZ.5\endcsname{\href{S6_SERIES_GRIND_II.pdf\#nameddest=tag-00OR}{II.{C.IZ.5} (p.~530)} }
\expandafter\def\csname SRef@cdp:eq:C.IZ.6\endcsname{\href{S6_SERIES_GRIND_II.pdf\#nameddest=tag-00OS}{II.{C.IZ.6} (p.~530)} }
\expandafter\def\csname SRef@cdp:eq:C.IZ.7\endcsname{\href{S6_SERIES_GRIND_II.pdf\#nameddest=tag-00OT}{II.{C.IZ.7} (p.~530)} }
\expandafter\def\csname SRef@cdp:eq:C.IZ.8\endcsname{\href{S6_SERIES_GRIND_II.pdf\#nameddest=tag-00OU}{II.{C.IZ.8} (p.~531)} }
\expandafter\def\csname SRef@cdp:eq:C.SC.4\endcsname{\href{S6_SERIES_GRIND_II.pdf\#nameddest=tag-010O}{II.{C.SC.4} (p.~423)} }
\expandafter\def\csname SRef@cdp:eq:C.INT.1\endcsname{\href{S6_SERIES_GRIND_II.pdf\#nameddest=tag-00NQ}{II.{C.INT.1} (p.~531)} }
\expandafter\def\csname SRef@cdp:eq:C.INT.2\endcsname{\href{S6_SERIES_GRIND_II.pdf\#nameddest=tag-00O1}{II.{C.INT.2} (p.~531)} }
\expandafter\def\csname SRef@cdp:eq:C.INT.3\endcsname{\href{S6_SERIES_GRIND_II.pdf\#nameddest=tag-00OA}{II.{C.INT.3} (p.~532)} }
\expandafter\def\csname SRef@cdp:eq:C.INT.4\endcsname{\href{S6_SERIES_GRIND_II.pdf\#nameddest=tag-00OB}{II.{C.INT.4} (p.~532)} }
\expandafter\def\csname SRef@cdp:eq:C.INT.5\endcsname{\href{S6_SERIES_GRIND_II.pdf\#nameddest=tag-00OC}{II.{C.INT.5} (p.~533)} }
\expandafter\def\csname SRef@cdp:eq:C.INT.6\endcsname{\href{S6_SERIES_GRIND_II.pdf\#nameddest=tag-00OD}{II.{C.INT.6} (p.~533)} }
\expandafter\def\csname SRef@cdp:eq:C.INT.7\endcsname{\href{S6_SERIES_GRIND_II.pdf\#nameddest=tag-00OE}{II.{C.INT.7} (p.~533)} }
\expandafter\def\csname SRef@cdp:eq:C.INT.8\endcsname{\href{S6_SERIES_GRIND_II.pdf\#nameddest=tag-00OF}{II.{C.INT.8} (p.~534)} }
\expandafter\def\csname SRef@cdp:eq:C.INT.9\endcsname{\href{S6_SERIES_GRIND_II.pdf\#nameddest=tag-00OG}{II.{C.INT.9} (p.~534)} }
\expandafter\def\csname SRef@cdp:eq:C.INT.10\endcsname{\href{S6_SERIES_GRIND_II.pdf\#nameddest=tag-00NR}{II.{C.INT.10} (p.~534)} }
\expandafter\def\csname SRef@cdp:eq:C.SC.5\endcsname{\href{S6_SERIES_GRIND_II.pdf\#nameddest=tag-010P}{II.{C.SC.5} (p.~424)} }
\expandafter\def\csname SRef@cdp:eq:C.INT.11\endcsname{\href{S6_SERIES_GRIND_II.pdf\#nameddest=tag-00NS}{II.{C.INT.11} (p.~535)} }
\expandafter\def\csname SRef@cdp:eq:C.INT.12\endcsname{\href{S6_SERIES_GRIND_II.pdf\#nameddest=tag-00NT}{II.{C.INT.12} (p.~535)} }
\expandafter\def\csname SRef@cdp:eq:C.INT.13\endcsname{\href{S6_SERIES_GRIND_II.pdf\#nameddest=tag-00NU}{II.{C.INT.13} (p.~535)} }
\expandafter\def\csname SRef@cdp:eq:C.INT.14\endcsname{\href{S6_SERIES_GRIND_II.pdf\#nameddest=tag-00NV}{II.{C.INT.14} (p.~535)} }
\expandafter\def\csname SRef@cdp:eq:C.INT.15\endcsname{\href{S6_SERIES_GRIND_II.pdf\#nameddest=tag-00NW}{II.{C.INT.15} (p.~535)} }
\expandafter\def\csname SRef@cdp:eq:C.INT.16\endcsname{\href{S6_SERIES_GRIND_II.pdf\#nameddest=tag-00NX}{II.{C.INT.16} (p.~536)} }
\expandafter\def\csname SRef@cdp:eq:C.INT.17\endcsname{\href{S6_SERIES_GRIND_II.pdf\#nameddest=tag-00NY}{II.{C.INT.17} (p.~536)} }
\expandafter\def\csname SRef@cdp:eq:C.INT.18\endcsname{\href{S6_SERIES_GRIND_II.pdf\#nameddest=tag-00NZ}{II.{C.INT.18} (p.~536)} }
\expandafter\def\csname SRef@cdp:eq:C.INT.19\endcsname{\href{S6_SERIES_GRIND_II.pdf\#nameddest=tag-00O0}{II.{C.INT.19} (p.~536)} }
\expandafter\def\csname SRef@cdp:eq:C.INT.20\endcsname{\href{S6_SERIES_GRIND_II.pdf\#nameddest=tag-00O2}{II.{C.INT.20} (p.~536)} }
\expandafter\def\csname SRef@cdp:eq:C.SC.6\endcsname{\href{S6_SERIES_GRIND_II.pdf\#nameddest=tag-010Q}{II.{C.SC.6} (p.~425)} }
\expandafter\def\csname SRef@cdp:eq:C.INT.21\endcsname{\href{S6_SERIES_GRIND_II.pdf\#nameddest=tag-00O3}{II.{C.INT.21} (p.~536)} }
\expandafter\def\csname SRef@cdp:eq:C.INT.22\endcsname{\href{S6_SERIES_GRIND_II.pdf\#nameddest=tag-00O4}{II.{C.INT.22} (p.~537)} }
\expandafter\def\csname SRef@cdp:eq:C.INT.23\endcsname{\href{S6_SERIES_GRIND_II.pdf\#nameddest=tag-00O5}{II.{C.INT.23} (p.~537)} }
\expandafter\def\csname SRef@cdp:eq:C.INT.24\endcsname{\href{S6_SERIES_GRIND_II.pdf\#nameddest=tag-00O6}{II.{C.INT.24} (p.~537)} }
\expandafter\def\csname SRef@cdp:eq:C.INT.25\endcsname{\href{S6_SERIES_GRIND_II.pdf\#nameddest=tag-00O7}{II.{C.INT.25} (p.~537)} }
\expandafter\def\csname SRef@cdp:eq:C.INT.26\endcsname{\href{S6_SERIES_GRIND_II.pdf\#nameddest=tag-00O8}{II.{C.INT.26} (p.~537)} }
\expandafter\def\csname SRef@cdp:eq:C.INT.27\endcsname{\href{S6_SERIES_GRIND_II.pdf\#nameddest=tag-00O9}{II.{C.INT.27} (p.~537)} }
\expandafter\def\csname SRef@cdp:eq:C.MH.1\endcsname{\href{S6_SERIES_GRIND_II.pdf\#nameddest=tag-00QA}{II.{C.MH.1} (p.~538)} }
\expandafter\def\csname SRef@cdp:eq:C.MH.2\endcsname{\href{S6_SERIES_GRIND_II.pdf\#nameddest=tag-00QJ}{II.{C.MH.2} (p.~538)} }
\expandafter\def\csname SRef@cdp:eq:C.MH.3\endcsname{\href{S6_SERIES_GRIND_II.pdf\#nameddest=tag-00QK}{II.{C.MH.3} (p.~539)} }
\expandafter\def\csname SRef@cdp:eq:C.CF.1\endcsname{\href{S6_SERIES_GRIND_II.pdf\#nameddest=tag-00BN}{II.{C.CF.1} (p.~425)} }
\expandafter\def\csname SRef@cdp:eq:C.MH.4\endcsname{\href{S6_SERIES_GRIND_II.pdf\#nameddest=tag-00QL}{II.{C.MH.4} (p.~540)} }
\expandafter\def\csname SRef@cdp:eq:C.MH.5\endcsname{\href{S6_SERIES_GRIND_II.pdf\#nameddest=tag-00QM}{II.{C.MH.5} (p.~540)} }
\expandafter\def\csname SRef@cdp:eq:C.MH.6\endcsname{\href{S6_SERIES_GRIND_II.pdf\#nameddest=tag-00QN}{II.{C.MH.6} (p.~540)} }
\expandafter\def\csname SRef@cdp:eq:C.MH.7\endcsname{\href{S6_SERIES_GRIND_II.pdf\#nameddest=tag-00QO}{II.{C.MH.7} (p.~540)} }
\expandafter\def\csname SRef@cdp:eq:C.MH.8\endcsname{\href{S6_SERIES_GRIND_II.pdf\#nameddest=tag-00QP}{II.{C.MH.8} (p.~541)} }
\expandafter\def\csname SRef@cdp:eq:C.MH.9\endcsname{\href{S6_SERIES_GRIND_II.pdf\#nameddest=tag-00QQ}{II.{C.MH.9} (p.~541)} }
\expandafter\def\csname SRef@cdp:eq:C.MH.10\endcsname{\href{S6_SERIES_GRIND_II.pdf\#nameddest=tag-00QB}{II.{C.MH.10} (p.~541)} }
\expandafter\def\csname SRef@cdp:eq:C.MH.11\endcsname{\href{S6_SERIES_GRIND_II.pdf\#nameddest=tag-00QC}{II.{C.MH.11} (p.~541)} }
\expandafter\def\csname SRef@cdp:eq:C.MH.12\endcsname{\href{S6_SERIES_GRIND_II.pdf\#nameddest=tag-00QD}{II.{C.MH.12} (p.~541)} }
\expandafter\def\csname SRef@cdp:eq:C.MH.13\endcsname{\href{S6_SERIES_GRIND_II.pdf\#nameddest=tag-00QE}{II.{C.MH.13} (p.~541)} }
\expandafter\def\csname SRef@cdp:eq:C.CF.2\endcsname{\href{S6_SERIES_GRIND_II.pdf\#nameddest=tag-00BO}{II.{C.CF.2} (p.~426)} }
\expandafter\def\csname SRef@cdp:eq:C.MH.14\endcsname{\href{S6_SERIES_GRIND_II.pdf\#nameddest=tag-00QF}{II.{C.MH.14} (p.~541)} }
\expandafter\def\csname SRef@cdp:eq:C.MH.15\endcsname{\href{S6_SERIES_GRIND_II.pdf\#nameddest=tag-00QG}{II.{C.MH.15} (p.~542)} }
\expandafter\def\csname SRef@cdp:eq:C.MH.16\endcsname{\href{S6_SERIES_GRIND_II.pdf\#nameddest=tag-00QH}{II.{C.MH.16} (p.~542)} }
\expandafter\def\csname SRef@cdp:eq:C.MH.17\endcsname{\href{S6_SERIES_GRIND_II.pdf\#nameddest=tag-00QI}{II.{C.MH.17} (p.~542)} }
\expandafter\def\csname SRef@cdp:eq:C.1\endcsname{\href{S6_SERIES_GRIND_II.pdf\#nameddest=tag-008N}{II.{C.1} (p.~543)} }
\expandafter\def\csname SRef@cdp:eq:C.2\endcsname{\href{S6_SERIES_GRIND_II.pdf\#nameddest=tag-008Y}{II.{C.2} (p.~544)} }
\expandafter\def\csname SRef@cdp:eq:C.MS.1\endcsname{\href{S6_SERIES_GRIND_II.pdf\#nameddest=tag-00RH}{II.{C.MS.1} (p.~546)} }
\expandafter\def\csname SRef@cdp:eq:C.MS.2\endcsname{\href{S6_SERIES_GRIND_II.pdf\#nameddest=tag-00RJ}{II.{C.MS.2} (p.~546)} }
\expandafter\def\csname SRef@cdp:eq:C.MS.3\endcsname{\href{S6_SERIES_GRIND_II.pdf\#nameddest=tag-00RK}{II.{C.MS.3} (p.~547)} }
\expandafter\def\csname SRef@cdp:eq:C.MS.10\endcsname{\href{S6_SERIES_GRIND_II.pdf\#nameddest=tag-00RI}{II.{C.MS.10} (p.~547)} }
\expandafter\def\csname SRef@cdp:eq:C.CF.3\endcsname{\href{S6_SERIES_GRIND_II.pdf\#nameddest=tag-00BP}{II.{C.CF.3} (p.~426)} }
\expandafter\def\csname SRef@cdp:eq:C.MS.4\endcsname{\href{S6_SERIES_GRIND_II.pdf\#nameddest=tag-00RL}{II.{C.MS.4} (p.~547)} }
\expandafter\def\csname SRef@cdp:eq:C.MS.5\endcsname{\href{S6_SERIES_GRIND_II.pdf\#nameddest=tag-00RM}{II.{C.MS.5} (p.~547)} }
\expandafter\def\csname SRef@cdp:eq:C.MS.6\endcsname{\href{S6_SERIES_GRIND_II.pdf\#nameddest=tag-00RN}{II.{C.MS.6} (p.~548)} }
\expandafter\def\csname SRef@cdp:eq:C.MS.7\endcsname{\href{S6_SERIES_GRIND_II.pdf\#nameddest=tag-00RO}{II.{C.MS.7} (p.~548)} }
\expandafter\def\csname SRef@cdp:eq:C.MS.8\endcsname{\href{S6_SERIES_GRIND_II.pdf\#nameddest=tag-00RP}{II.{C.MS.8} (p.~549)} }
\expandafter\def\csname SRef@cdp:eq:C.MS.9\endcsname{\href{S6_SERIES_GRIND_II.pdf\#nameddest=tag-00RQ}{II.{C.MS.9} (p.~549)} }
\expandafter\def\csname SRef@cdp:eq:C.K3.1\endcsname{\href{S6_SERIES_GRIND_II.pdf\#nameddest=tag-00OY}{II.{C.K3.1} (p.~551)} }
\expandafter\def\csname SRef@cdp:eq:C.K3.2\endcsname{\href{S6_SERIES_GRIND_II.pdf\#nameddest=tag-00P0}{II.{C.K3.2} (p.~551)} }
\expandafter\def\csname SRef@cdp:eq:C.K3.3\endcsname{\href{S6_SERIES_GRIND_II.pdf\#nameddest=tag-00P1}{II.{C.K3.3} (p.~551)} }
\expandafter\def\csname SRef@cdp:eq:C.K3.4\endcsname{\href{S6_SERIES_GRIND_II.pdf\#nameddest=tag-00P2}{II.{C.K3.4} (p.~551)} }
\expandafter\def\csname SRef@cdp:eq:C.CF.4\endcsname{\href{S6_SERIES_GRIND_II.pdf\#nameddest=tag-00BQ}{II.{C.CF.4} (p.~426)} }
\expandafter\def\csname SRef@cdp:eq:C.K3.5\endcsname{\href{S6_SERIES_GRIND_II.pdf\#nameddest=tag-00P3}{II.{C.K3.5} (p.~552)} }
\expandafter\def\csname SRef@cdp:eq:C.K3.6\endcsname{\href{S6_SERIES_GRIND_II.pdf\#nameddest=tag-00P4}{II.{C.K3.6} (p.~552)} }
\expandafter\def\csname SRef@cdp:eq:C.K3.7\endcsname{\href{S6_SERIES_GRIND_II.pdf\#nameddest=tag-00P5}{II.{C.K3.7} (p.~552)} }
\expandafter\def\csname SRef@cdp:eq:C.K3.8\endcsname{\href{S6_SERIES_GRIND_II.pdf\#nameddest=tag-00P6}{II.{C.K3.8} (p.~553)} }
\expandafter\def\csname SRef@cdp:eq:C.K3.9\endcsname{\href{S6_SERIES_GRIND_II.pdf\#nameddest=tag-00P7}{II.{C.K3.9} (p.~553)} }
\expandafter\def\csname SRef@cdp:eq:C.K3.10\endcsname{\href{S6_SERIES_GRIND_II.pdf\#nameddest=tag-00OZ}{II.{C.K3.10} (p.~553)} }
\expandafter\def\csname SRef@cdp:eq:C.KX.1\endcsname{\href{S6_SERIES_GRIND_II.pdf\#nameddest=tag-00PW}{II.{C.KX.1} (p.~554)} }
\expandafter\def\csname SRef@cdp:eq:C.KX.2\endcsname{\href{S6_SERIES_GRIND_II.pdf\#nameddest=tag-00Q2}{II.{C.KX.2} (p.~554)} }
\expandafter\def\csname SRef@cdp:eq:C.KX.3\endcsname{\href{S6_SERIES_GRIND_II.pdf\#nameddest=tag-00Q3}{II.{C.KX.3} (p.~554)} }
\expandafter\def\csname SRef@cdp:eq:C.KX.4\endcsname{\href{S6_SERIES_GRIND_II.pdf\#nameddest=tag-00Q4}{II.{C.KX.4} (p.~555)} }
\expandafter\def\csname SRef@cdp:eq:C.CF.5\endcsname{\href{S6_SERIES_GRIND_II.pdf\#nameddest=tag-00BR}{II.{C.CF.5} (p.~426)} }
\expandafter\def\csname SRef@cdp:eq:C.KX.5\endcsname{\href{S6_SERIES_GRIND_II.pdf\#nameddest=tag-00Q5}{II.{C.KX.5} (p.~555)} }
\expandafter\def\csname SRef@cdp:eq:C.KX.6\endcsname{\href{S6_SERIES_GRIND_II.pdf\#nameddest=tag-00Q6}{II.{C.KX.6} (p.~555)} }
\expandafter\def\csname SRef@cdp:eq:C.KX.7\endcsname{\href{S6_SERIES_GRIND_II.pdf\#nameddest=tag-00Q7}{II.{C.KX.7} (p.~555)} }
\expandafter\def\csname SRef@cdp:eq:C.KX.8\endcsname{\href{S6_SERIES_GRIND_II.pdf\#nameddest=tag-00Q8}{II.{C.KX.8} (p.~556)} }
\expandafter\def\csname SRef@cdp:eq:C.KX.9\endcsname{\href{S6_SERIES_GRIND_II.pdf\#nameddest=tag-00Q9}{II.{C.KX.9} (p.~556)} }
\expandafter\def\csname SRef@cdp:eq:C.KX.10\endcsname{\href{S6_SERIES_GRIND_II.pdf\#nameddest=tag-00PX}{II.{C.KX.10} (p.~556)} }
\expandafter\def\csname SRef@cdp:eq:C.KX.11\endcsname{\href{S6_SERIES_GRIND_II.pdf\#nameddest=tag-00PY}{II.{C.KX.11} (p.~556)} }
\expandafter\def\csname SRef@cdp:eq:C.KX.12\endcsname{\href{S6_SERIES_GRIND_II.pdf\#nameddest=tag-00PZ}{II.{C.KX.12} (p.~556)} }
\expandafter\def\csname SRef@cdp:eq:C.KX.13\endcsname{\href{S6_SERIES_GRIND_II.pdf\#nameddest=tag-00Q0}{II.{C.KX.13} (p.~557)} }
\expandafter\def\csname SRef@cdp:eq:C.KX.14\endcsname{\href{S6_SERIES_GRIND_II.pdf\#nameddest=tag-00Q1}{II.{C.KX.14} (p.~557)} }
\expandafter\def\csname SRef@cdp:eq:C.CF.6\endcsname{\href{S6_SERIES_GRIND_II.pdf\#nameddest=tag-00BS}{II.{C.CF.6} (p.~426)} }
\expandafter\def\csname SRef@cdp:eq:C.EL.1\endcsname{\href{S6_SERIES_GRIND_II.pdf\#nameddest=tag-00FY}{II.{C.EL.1} (p.~557)} }
\expandafter\def\csname SRef@cdp:eq:C.EL.2\endcsname{\href{S6_SERIES_GRIND_II.pdf\#nameddest=tag-00G1}{II.{C.EL.2} (p.~557)} }
\expandafter\def\csname SRef@cdp:eq:C.EL.3\endcsname{\href{S6_SERIES_GRIND_II.pdf\#nameddest=tag-00G2}{II.{C.EL.3} (p.~557)} }
\expandafter\def\csname SRef@cdp:eq:C.EL.4\endcsname{\href{S6_SERIES_GRIND_II.pdf\#nameddest=tag-00G3}{II.{C.EL.4} (p.~558)} }
\expandafter\def\csname SRef@cdp:eq:C.EL.5\endcsname{\href{S6_SERIES_GRIND_II.pdf\#nameddest=tag-00G4}{II.{C.EL.5} (p.~558)} }
\expandafter\def\csname SRef@cdp:eq:C.EL.6\endcsname{\href{S6_SERIES_GRIND_II.pdf\#nameddest=tag-00G5}{II.{C.EL.6} (p.~558)} }
\expandafter\def\csname SRef@cdp:eq:C.EL.7\endcsname{\href{S6_SERIES_GRIND_II.pdf\#nameddest=tag-00G6}{II.{C.EL.7} (p.~558)} }
\expandafter\def\csname SRef@cdp:eq:C.EL.8\endcsname{\href{S6_SERIES_GRIND_II.pdf\#nameddest=tag-00G7}{II.{C.EL.8} (p.~558)} }
\expandafter\def\csname SRef@cdp:eq:C.EL.9\endcsname{\href{S6_SERIES_GRIND_II.pdf\#nameddest=tag-00G8}{II.{C.EL.9} (p.~559)} }
\expandafter\def\csname SRef@cdp:eq:C.EL.10\endcsname{\href{S6_SERIES_GRIND_II.pdf\#nameddest=tag-00FZ}{II.{C.EL.10} (p.~560)} }
\expandafter\def\csname SRef@cdp:eq:C.CF.7\endcsname{\href{S6_SERIES_GRIND_II.pdf\#nameddest=tag-00BT}{II.{C.CF.7} (p.~427)} }
\expandafter\def\csname SRef@cdp:eq:C.EL.11\endcsname{\href{S6_SERIES_GRIND_II.pdf\#nameddest=tag-00G0}{II.{C.EL.11} (p.~560)} }
\expandafter\def\csname SRef@cdp:eq:C.CZ.1\endcsname{\href{S6_SERIES_GRIND_II.pdf\#nameddest=tag-00DN}{II.{C.CZ.1} (p.~561)} }
\expandafter\def\csname SRef@cdp:eq:C.CZ.2\endcsname{\href{S6_SERIES_GRIND_II.pdf\#nameddest=tag-00DT}{II.{C.CZ.2} (p.~561)} }
\expandafter\def\csname SRef@cdp:eq:C.CZ.3\endcsname{\href{S6_SERIES_GRIND_II.pdf\#nameddest=tag-00DU}{II.{C.CZ.3} (p.~561)} }
\expandafter\def\csname SRef@cdp:eq:C.CZ.4\endcsname{\href{S6_SERIES_GRIND_II.pdf\#nameddest=tag-00DV}{II.{C.CZ.4} (p.~561)} }
\expandafter\def\csname SRef@cdp:eq:C.CZ.5\endcsname{\href{S6_SERIES_GRIND_II.pdf\#nameddest=tag-00DW}{II.{C.CZ.5} (p.~562)} }
\expandafter\def\csname SRef@cdp:eq:C.CZ.6\endcsname{\href{S6_SERIES_GRIND_II.pdf\#nameddest=tag-00DX}{II.{C.CZ.6} (p.~562)} }
\expandafter\def\csname SRef@cdp:eq:C.CZ.7\endcsname{\href{S6_SERIES_GRIND_II.pdf\#nameddest=tag-00DY}{II.{C.CZ.7} (p.~562)} }
\expandafter\def\csname SRef@cdp:eq:C.CZ.8\endcsname{\href{S6_SERIES_GRIND_II.pdf\#nameddest=tag-00DZ}{II.{C.CZ.8} (p.~562)} }
\expandafter\def\csname SRef@cdp:eq:C.CZ.9\endcsname{\href{S6_SERIES_GRIND_II.pdf\#nameddest=tag-00E0}{II.{C.CZ.9} (p.~562)} }
\expandafter\def\csname SRef@cdp:eq:C.CF.8\endcsname{\href{S6_SERIES_GRIND_II.pdf\#nameddest=tag-00BU}{II.{C.CF.8} (p.~427)} }
\expandafter\def\csname SRef@cdp:eq:C.CZ.10\endcsname{\href{S6_SERIES_GRIND_II.pdf\#nameddest=tag-00DO}{II.{C.CZ.10} (p.~563)} }
\expandafter\def\csname SRef@cdp:eq:C.CZ.11\endcsname{\href{S6_SERIES_GRIND_II.pdf\#nameddest=tag-00DP}{II.{C.CZ.11} (p.~563)} }
\expandafter\def\csname SRef@cdp:eq:C.CZ.12\endcsname{\href{S6_SERIES_GRIND_II.pdf\#nameddest=tag-00DQ}{II.{C.CZ.12} (p.~564)} }
\expandafter\def\csname SRef@cdp:eq:C.CZ.13\endcsname{\href{S6_SERIES_GRIND_II.pdf\#nameddest=tag-00DR}{II.{C.CZ.13} (p.~564)} }
\expandafter\def\csname SRef@cdp:eq:C.CZ.14\endcsname{\href{S6_SERIES_GRIND_II.pdf\#nameddest=tag-00DS}{II.{C.CZ.14} (p.~564)} }
\expandafter\def\csname SRef@cdp:eq:C.MR.1\endcsname{\href{S6_SERIES_GRIND_II.pdf\#nameddest=tag-00R1}{II.{C.MR.1} (p.~566)} }
\expandafter\def\csname SRef@cdp:eq:C.MR.2\endcsname{\href{S6_SERIES_GRIND_II.pdf\#nameddest=tag-00R9}{II.{C.MR.2} (p.~566)} }
\expandafter\def\csname SRef@cdp:eq:C.MR.3\endcsname{\href{S6_SERIES_GRIND_II.pdf\#nameddest=tag-00RA}{II.{C.MR.3} (p.~566)} }
\expandafter\def\csname SRef@cdp:eq:C.MR.4\endcsname{\href{S6_SERIES_GRIND_II.pdf\#nameddest=tag-00RB}{II.{C.MR.4} (p.~566)} }
\expandafter\def\csname SRef@cdp:eq:C.MR.5\endcsname{\href{S6_SERIES_GRIND_II.pdf\#nameddest=tag-00RC}{II.{C.MR.5} (p.~567)} }
\expandafter\def\csname SRef@cdp:eq:C.CF.9\endcsname{\href{S6_SERIES_GRIND_II.pdf\#nameddest=tag-00BV}{II.{C.CF.9} (p.~427)} }
\expandafter\def\csname SRef@cdp:eq:C.MR.6\endcsname{\href{S6_SERIES_GRIND_II.pdf\#nameddest=tag-00RD}{II.{C.MR.6} (p.~567)} }
\expandafter\def\csname SRef@cdp:eq:C.MR.7\endcsname{\href{S6_SERIES_GRIND_II.pdf\#nameddest=tag-00RE}{II.{C.MR.7} (p.~567)} }
\expandafter\def\csname SRef@cdp:eq:C.MR.8\endcsname{\href{S6_SERIES_GRIND_II.pdf\#nameddest=tag-00RF}{II.{C.MR.8} (p.~568)} }
\expandafter\def\csname SRef@cdp:eq:C.MR.9\endcsname{\href{S6_SERIES_GRIND_II.pdf\#nameddest=tag-00RG}{II.{C.MR.9} (p.~568)} }
\expandafter\def\csname SRef@cdp:eq:C.MR.10\endcsname{\href{S6_SERIES_GRIND_II.pdf\#nameddest=tag-00R2}{II.{C.MR.10} (p.~569)} }
\expandafter\def\csname SRef@cdp:eq:C.MR.11\endcsname{\href{S6_SERIES_GRIND_II.pdf\#nameddest=tag-00R3}{II.{C.MR.11} (p.~569)} }
\expandafter\def\csname SRef@cdp:eq:C.MR.12\endcsname{\href{S6_SERIES_GRIND_II.pdf\#nameddest=tag-00R4}{II.{C.MR.12} (p.~570)} }
\expandafter\def\csname SRef@cdp:eq:C.MR.13\endcsname{\href{S6_SERIES_GRIND_II.pdf\#nameddest=tag-00R5}{II.{C.MR.13} (p.~570)} }
\expandafter\def\csname SRef@cdp:eq:C.MR.14\endcsname{\href{S6_SERIES_GRIND_II.pdf\#nameddest=tag-00R6}{II.{C.MR.14} (p.~570)} }
\expandafter\def\csname SRef@cdp:eq:C.MR.15\endcsname{\href{S6_SERIES_GRIND_II.pdf\#nameddest=tag-00R7}{II.{C.MR.15} (p.~571)} }
\expandafter\def\csname SRef@cdp:eq:C.CA.1\endcsname{\href{S6_SERIES_GRIND_II.pdf\#nameddest=tag-00BC}{II.{C.CA.1} (p.~428)} }
\expandafter\def\csname SRef@cdp:eq:C.MR.16\endcsname{\href{S6_SERIES_GRIND_II.pdf\#nameddest=tag-00R8}{II.{C.MR.16} (p.~572)} }
\expandafter\def\csname SRef@cdp:eq:C.RH.1\endcsname{\href{S6_SERIES_GRIND_II.pdf\#nameddest=tag-00ZG}{II.{C.RH.1} (p.~573)} }
\expandafter\def\csname SRef@cdp:eq:C.RH.2\endcsname{\href{S6_SERIES_GRIND_II.pdf\#nameddest=tag-00ZI}{II.{C.RH.2} (p.~573)} }
\expandafter\def\csname SRef@cdp:eq:C.RH.3\endcsname{\href{S6_SERIES_GRIND_II.pdf\#nameddest=tag-00ZJ}{II.{C.RH.3} (p.~573)} }
\expandafter\def\csname SRef@cdp:eq:C.RH.4\endcsname{\href{S6_SERIES_GRIND_II.pdf\#nameddest=tag-00ZK}{II.{C.RH.4} (p.~573)} }
\expandafter\def\csname SRef@cdp:eq:C.RH.5\endcsname{\href{S6_SERIES_GRIND_II.pdf\#nameddest=tag-00ZL}{II.{C.RH.5} (p.~573)} }
\expandafter\def\csname SRef@cdp:eq:C.RH.6\endcsname{\href{S6_SERIES_GRIND_II.pdf\#nameddest=tag-00ZM}{II.{C.RH.6} (p.~573)} }
\expandafter\def\csname SRef@cdp:eq:C.RH.7\endcsname{\href{S6_SERIES_GRIND_II.pdf\#nameddest=tag-00ZN}{II.{C.RH.7} (p.~574)} }
\expandafter\def\csname SRef@cdp:eq:C.RH.8\endcsname{\href{S6_SERIES_GRIND_II.pdf\#nameddest=tag-00ZO}{II.{C.RH.8} (p.~575)} }
\expandafter\def\csname SRef@cdp:eq:C.RH.9\endcsname{\href{S6_SERIES_GRIND_II.pdf\#nameddest=tag-00ZP}{II.{C.RH.9} (p.~575)} }
\expandafter\def\csname SRef@cdp:eq:C.CA.2\endcsname{\href{S6_SERIES_GRIND_II.pdf\#nameddest=tag-00BF}{II.{C.CA.2} (p.~428)} }
\expandafter\def\csname SRef@cdp:eq:C.RH.10\endcsname{\href{S6_SERIES_GRIND_II.pdf\#nameddest=tag-00ZH}{II.{C.RH.10} (p.~575)} }
\expandafter\def\csname SRef@cdp:eq:C.TC.1\endcsname{\href{S6_SERIES_GRIND_II.pdf\#nameddest=tag-012C}{II.{C.TC.1} (p.~576)} }
\expandafter\def\csname SRef@cdp:eq:C.TC.2\endcsname{\href{S6_SERIES_GRIND_II.pdf\#nameddest=tag-012G}{II.{C.TC.2} (p.~576)} }
\expandafter\def\csname SRef@cdp:eq:C.TC.3\endcsname{\href{S6_SERIES_GRIND_II.pdf\#nameddest=tag-012H}{II.{C.TC.3} (p.~576)} }
\expandafter\def\csname SRef@cdp:eq:C.TC.4\endcsname{\href{S6_SERIES_GRIND_II.pdf\#nameddest=tag-012I}{II.{C.TC.4} (p.~577)} }
\expandafter\def\csname SRef@cdp:eq:C.TC.5\endcsname{\href{S6_SERIES_GRIND_II.pdf\#nameddest=tag-012J}{II.{C.TC.5} (p.~577)} }
\expandafter\def\csname SRef@cdp:eq:C.TC.6\endcsname{\href{S6_SERIES_GRIND_II.pdf\#nameddest=tag-012K}{II.{C.TC.6} (p.~577)} }
\expandafter\def\csname SRef@cdp:eq:C.TC.7\endcsname{\href{S6_SERIES_GRIND_II.pdf\#nameddest=tag-012L}{II.{C.TC.7} (p.~577)} }
\expandafter\def\csname SRef@cdp:eq:C.TC.8\endcsname{\href{S6_SERIES_GRIND_II.pdf\#nameddest=tag-012M}{II.{C.TC.8} (p.~577)} }
\expandafter\def\csname SRef@cdp:eq:C.TC.9\endcsname{\href{S6_SERIES_GRIND_II.pdf\#nameddest=tag-012N}{II.{C.TC.9} (p.~578)} }
\expandafter\def\csname SRef@cdp:eq:C.CA.3\endcsname{\href{S6_SERIES_GRIND_II.pdf\#nameddest=tag-00BG}{II.{C.CA.3} (p.~429)} }
\expandafter\def\csname SRef@cdp:eq:C.TC.10\endcsname{\href{S6_SERIES_GRIND_II.pdf\#nameddest=tag-012D}{II.{C.TC.10} (p.~578)} }
\expandafter\def\csname SRef@cdp:eq:C.TC.11\endcsname{\href{S6_SERIES_GRIND_II.pdf\#nameddest=tag-012E}{II.{C.TC.11} (p.~578)} }
\expandafter\def\csname SRef@cdp:eq:C.TC.12\endcsname{\href{S6_SERIES_GRIND_II.pdf\#nameddest=tag-012F}{II.{C.TC.12} (p.~578)} }
\expandafter\def\csname SRef@cdp:eq:C.ER.1\endcsname{\href{S6_SERIES_GRIND_II.pdf\#nameddest=tag-00IN}{II.{C.ER.1} (p.~579)} }
\expandafter\def\csname SRef@cdp:eq:C.ER.2\endcsname{\href{S6_SERIES_GRIND_II.pdf\#nameddest=tag-00IX}{II.{C.ER.2} (p.~579)} }
\expandafter\def\csname SRef@cdp:eq:C.ER.3\endcsname{\href{S6_SERIES_GRIND_II.pdf\#nameddest=tag-00IY}{II.{C.ER.3} (p.~579)} }
\expandafter\def\csname SRef@cdp:eq:C.ER.4\endcsname{\href{S6_SERIES_GRIND_II.pdf\#nameddest=tag-00IZ}{II.{C.ER.4} (p.~579)} }
\expandafter\def\csname SRef@cdp:eq:C.ER.5\endcsname{\href{S6_SERIES_GRIND_II.pdf\#nameddest=tag-00J0}{II.{C.ER.5} (p.~579)} }
\expandafter\def\csname SRef@cdp:eq:C.ER.6\endcsname{\href{S6_SERIES_GRIND_II.pdf\#nameddest=tag-00J1}{II.{C.ER.6} (p.~579)} }
\expandafter\def\csname SRef@cdp:eq:C.ER.7\endcsname{\href{S6_SERIES_GRIND_II.pdf\#nameddest=tag-00J2}{II.{C.ER.7} (p.~579)} }
\expandafter\def\csname SRef@cdp:eq:C.CA.4\endcsname{\href{S6_SERIES_GRIND_II.pdf\#nameddest=tag-00BH}{II.{C.CA.4} (p.~429)} }
\expandafter\def\csname SRef@cdp:eq:C.ER.8\endcsname{\href{S6_SERIES_GRIND_II.pdf\#nameddest=tag-00J3}{II.{C.ER.8} (p.~580)} }
\expandafter\def\csname SRef@cdp:eq:C.ER.9\endcsname{\href{S6_SERIES_GRIND_II.pdf\#nameddest=tag-00J4}{II.{C.ER.9} (p.~580)} }
\expandafter\def\csname SRef@cdp:eq:C.ER.10\endcsname{\href{S6_SERIES_GRIND_II.pdf\#nameddest=tag-00IO}{II.{C.ER.10} (p.~580)} }
\expandafter\def\csname SRef@cdp:eq:C.ER.11\endcsname{\href{S6_SERIES_GRIND_II.pdf\#nameddest=tag-00IP}{II.{C.ER.11} (p.~580)} }
\expandafter\def\csname SRef@cdp:eq:C.ER.12\endcsname{\href{S6_SERIES_GRIND_II.pdf\#nameddest=tag-00IQ}{II.{C.ER.12} (p.~580)} }
\expandafter\def\csname SRef@cdp:eq:C.ER.13\endcsname{\href{S6_SERIES_GRIND_II.pdf\#nameddest=tag-00IR}{II.{C.ER.13} (p.~580)} }
\expandafter\def\csname SRef@cdp:eq:C.ER.14\endcsname{\href{S6_SERIES_GRIND_II.pdf\#nameddest=tag-00IS}{II.{C.ER.14} (p.~580)} }
\expandafter\def\csname SRef@cdp:eq:C.ER.15\endcsname{\href{S6_SERIES_GRIND_II.pdf\#nameddest=tag-00IT}{II.{C.ER.15} (p.~581)} }
\expandafter\def\csname SRef@cdp:eq:C.ER.16\endcsname{\href{S6_SERIES_GRIND_II.pdf\#nameddest=tag-00IU}{II.{C.ER.16} (p.~581)} }
\expandafter\def\csname SRef@cdp:eq:C.CA.5\endcsname{\href{S6_SERIES_GRIND_II.pdf\#nameddest=tag-00BI}{II.{C.CA.5} (p.~429)} }
\expandafter\def\csname SRef@cdp:eq:C.ER.17\endcsname{\href{S6_SERIES_GRIND_II.pdf\#nameddest=tag-00IV}{II.{C.ER.17} (p.~582)} }
\expandafter\def\csname SRef@cdp:eq:C.ER.18\endcsname{\href{S6_SERIES_GRIND_II.pdf\#nameddest=tag-00IW}{II.{C.ER.18} (p.~582)} }
\expandafter\def\csname SRef@cdp:eq:C.CP.1\endcsname{\href{S6_SERIES_GRIND_II.pdf\#nameddest=tag-00C9}{II.{C.CP.1} (p.~583)} }
\expandafter\def\csname SRef@cdp:eq:C.CP.2\endcsname{\href{S6_SERIES_GRIND_II.pdf\#nameddest=tag-00CB}{II.{C.CP.2} (p.~583)} }
\expandafter\def\csname SRef@cdp:eq:C.CP.3\endcsname{\href{S6_SERIES_GRIND_II.pdf\#nameddest=tag-00CC}{II.{C.CP.3} (p.~583)} }
\expandafter\def\csname SRef@cdp:eq:C.CP.4\endcsname{\href{S6_SERIES_GRIND_II.pdf\#nameddest=tag-00CD}{II.{C.CP.4} (p.~583)} }
\expandafter\def\csname SRef@cdp:eq:C.CP.5\endcsname{\href{S6_SERIES_GRIND_II.pdf\#nameddest=tag-00CE}{II.{C.CP.5} (p.~583)} }
\expandafter\def\csname SRef@cdp:eq:C.CP.6\endcsname{\href{S6_SERIES_GRIND_II.pdf\#nameddest=tag-00CF}{II.{C.CP.6} (p.~584)} }
\expandafter\def\csname SRef@cdp:eq:C.CP.7\endcsname{\href{S6_SERIES_GRIND_II.pdf\#nameddest=tag-00CG}{II.{C.CP.7} (p.~584)} }
\expandafter\def\csname SRef@cdp:eq:C.CA.6\endcsname{\href{S6_SERIES_GRIND_II.pdf\#nameddest=tag-00BJ}{II.{C.CA.6} (p.~429)} }
\expandafter\def\csname SRef@cdp:eq:C.CP.8\endcsname{\href{S6_SERIES_GRIND_II.pdf\#nameddest=tag-00CH}{II.{C.CP.8} (p.~584)} }
\expandafter\def\csname SRef@cdp:eq:C.CP.9\endcsname{\href{S6_SERIES_GRIND_II.pdf\#nameddest=tag-00CI}{II.{C.CP.9} (p.~584)} }
\expandafter\def\csname SRef@cdp:eq:C.CP.10\endcsname{\href{S6_SERIES_GRIND_II.pdf\#nameddest=tag-00CA}{II.{C.CP.10} (p.~585)} }
\expandafter\def\csname SRef@cdp:eq:C.KM.1\endcsname{\href{S6_SERIES_GRIND_II.pdf\#nameddest=tag-00PI}{II.{C.KM.1} (p.~586)} }
\expandafter\def\csname SRef@cdp:eq:C.KM.2\endcsname{\href{S6_SERIES_GRIND_II.pdf\#nameddest=tag-00PL}{II.{C.KM.2} (p.~586)} }
\expandafter\def\csname SRef@cdp:eq:C.KM.S1\endcsname{\href{S6_SERIES_GRIND_II.pdf\#nameddest=tag-00PT}{II.{C.KM.S1} (p.~586)} }
\expandafter\def\csname SRef@cdp:eq:C.KM.3\endcsname{\href{S6_SERIES_GRIND_II.pdf\#nameddest=tag-00PM}{II.{C.KM.3} (p.~587)} }
\expandafter\def\csname SRef@cdp:eq:C.KM.4\endcsname{\href{S6_SERIES_GRIND_II.pdf\#nameddest=tag-00PN}{II.{C.KM.4} (p.~587)} }
\expandafter\def\csname SRef@cdp:eq:C.KM.5\endcsname{\href{S6_SERIES_GRIND_II.pdf\#nameddest=tag-00PO}{II.{C.KM.5} (p.~587)} }
\expandafter\def\csname SRef@cdp:eq:C.KM.6\endcsname{\href{S6_SERIES_GRIND_II.pdf\#nameddest=tag-00PP}{II.{C.KM.6} (p.~587)} }
\expandafter\def\csname SRef@cdp:eq:C.CA.7\endcsname{\href{S6_SERIES_GRIND_II.pdf\#nameddest=tag-00BK}{II.{C.CA.7} (p.~430)} }
\expandafter\def\csname SRef@cdp:eq:C.KM.7\endcsname{\href{S6_SERIES_GRIND_II.pdf\#nameddest=tag-00PQ}{II.{C.KM.7} (p.~587)} }
\expandafter\def\csname SRef@cdp:eq:C.KM.8\endcsname{\href{S6_SERIES_GRIND_II.pdf\#nameddest=tag-00PR}{II.{C.KM.8} (p.~588)} }
\expandafter\def\csname SRef@cdp:eq:C.KM.9\endcsname{\href{S6_SERIES_GRIND_II.pdf\#nameddest=tag-00PS}{II.{C.KM.9} (p.~588)} }
\expandafter\def\csname SRef@cdp:eq:C.KM.10\endcsname{\href{S6_SERIES_GRIND_II.pdf\#nameddest=tag-00PJ}{II.{C.KM.10} (p.~588)} }
\expandafter\def\csname SRef@cdp:eq:C.KM.S2\endcsname{\href{S6_SERIES_GRIND_II.pdf\#nameddest=tag-00PU}{II.{C.KM.S2} (p.~589)} }
\expandafter\def\csname SRef@cdp:eq:C.KM.S3\endcsname{\href{S6_SERIES_GRIND_II.pdf\#nameddest=tag-00PV}{II.{C.KM.S3} (p.~589)} }
\expandafter\def\csname SRef@cdp:eq:C.KM.11\endcsname{\href{S6_SERIES_GRIND_II.pdf\#nameddest=tag-00PK}{II.{C.KM.11} (p.~589)} }
\expandafter\def\csname SRef@cdp:eq:C.PE.1\endcsname{\href{S6_SERIES_GRIND_II.pdf\#nameddest=tag-00U2}{II.{C.PE.1} (p.~590)} }
\expandafter\def\csname SRef@cdp:eq:C.PE.2\endcsname{\href{S6_SERIES_GRIND_II.pdf\#nameddest=tag-00UD}{II.{C.PE.2} (p.~590)} }
\expandafter\def\csname SRef@cdp:eq:C.PE.3\endcsname{\href{S6_SERIES_GRIND_II.pdf\#nameddest=tag-00UG}{II.{C.PE.3} (p.~590)} }
\expandafter\def\csname SRef@cdp:eq:C.CA.8\endcsname{\href{S6_SERIES_GRIND_II.pdf\#nameddest=tag-00BL}{II.{C.CA.8} (p.~430)} }
\expandafter\def\csname SRef@cdp:eq:C.PE.4\endcsname{\href{S6_SERIES_GRIND_II.pdf\#nameddest=tag-00UH}{II.{C.PE.4} (p.~590)} }
\expandafter\def\csname SRef@cdp:eq:C.PE.5\endcsname{\href{S6_SERIES_GRIND_II.pdf\#nameddest=tag-00UI}{II.{C.PE.5} (p.~590)} }
\expandafter\def\csname SRef@cdp:eq:C.PE.6\endcsname{\href{S6_SERIES_GRIND_II.pdf\#nameddest=tag-00UJ}{II.{C.PE.6} (p.~591)} }
\expandafter\def\csname SRef@cdp:eq:C.PE.7\endcsname{\href{S6_SERIES_GRIND_II.pdf\#nameddest=tag-00UK}{II.{C.PE.7} (p.~591)} }
\expandafter\def\csname SRef@cdp:eq:C.PE.8\endcsname{\href{S6_SERIES_GRIND_II.pdf\#nameddest=tag-00UL}{II.{C.PE.8} (p.~591)} }
\expandafter\def\csname SRef@cdp:eq:C.PE.9\endcsname{\href{S6_SERIES_GRIND_II.pdf\#nameddest=tag-00UM}{II.{C.PE.9} (p.~591)} }
\expandafter\def\csname SRef@cdp:eq:C.PE.10\endcsname{\href{S6_SERIES_GRIND_II.pdf\#nameddest=tag-00U3}{II.{C.PE.10} (p.~591)} }
\expandafter\def\csname SRef@cdp:eq:C.PE.11\endcsname{\href{S6_SERIES_GRIND_II.pdf\#nameddest=tag-00U4}{II.{C.PE.11} (p.~591)} }
\expandafter\def\csname SRef@cdp:eq:C.PE.12\endcsname{\href{S6_SERIES_GRIND_II.pdf\#nameddest=tag-00U5}{II.{C.PE.12} (p.~591)} }
\expandafter\def\csname SRef@cdp:eq:C.PE.13\endcsname{\href{S6_SERIES_GRIND_II.pdf\#nameddest=tag-00U6}{II.{C.PE.13} (p.~592)} }
\expandafter\def\csname SRef@cdp:eq:C.CA.9\endcsname{\href{S6_SERIES_GRIND_II.pdf\#nameddest=tag-00BM}{II.{C.CA.9} (p.~430)} }
\expandafter\def\csname SRef@cdp:eq:C.PE.14\endcsname{\href{S6_SERIES_GRIND_II.pdf\#nameddest=tag-00U7}{II.{C.PE.14} (p.~592)} }
\expandafter\def\csname SRef@cdp:eq:C.PE.15\endcsname{\href{S6_SERIES_GRIND_II.pdf\#nameddest=tag-00U8}{II.{C.PE.15} (p.~592)} }
\expandafter\def\csname SRef@cdp:eq:C.PE.16\endcsname{\href{S6_SERIES_GRIND_II.pdf\#nameddest=tag-00U9}{II.{C.PE.16} (p.~592)} }
\expandafter\def\csname SRef@cdp:eq:C.PE.17\endcsname{\href{S6_SERIES_GRIND_II.pdf\#nameddest=tag-00UA}{II.{C.PE.17} (p.~592)} }
\expandafter\def\csname SRef@cdp:eq:C.PE.18\endcsname{\href{S6_SERIES_GRIND_II.pdf\#nameddest=tag-00UB}{II.{C.PE.18} (p.~593)} }
\expandafter\def\csname SRef@cdp:eq:C.PE.19\endcsname{\href{S6_SERIES_GRIND_II.pdf\#nameddest=tag-00UC}{II.{C.PE.19} (p.~593)} }
\expandafter\def\csname SRef@cdp:eq:C.PE.20\endcsname{\href{S6_SERIES_GRIND_II.pdf\#nameddest=tag-00UE}{II.{C.PE.20} (p.~593)} }
\expandafter\def\csname SRef@cdp:eq:C.PE.21\endcsname{\href{S6_SERIES_GRIND_II.pdf\#nameddest=tag-00UF}{II.{C.PE.21} (p.~593)} }
\expandafter\def\csname SRef@cdp:eq:C.SKD.1\endcsname{\href{S6_SERIES_GRIND_II.pdf\#nameddest=tag-011N}{II.{C.SKD.1} (p.~594)} }
\expandafter\def\csname SRef@cdp:eq:C.SKD.2\endcsname{\href{S6_SERIES_GRIND_II.pdf\#nameddest=tag-011U}{II.{C.SKD.2} (p.~594)} }
\expandafter\def\csname SRef@cdp:eq:C.CA.10\endcsname{\href{S6_SERIES_GRIND_II.pdf\#nameddest=tag-00BD}{II.{C.CA.10} (p.~430)} }
\expandafter\def\csname SRef@cdp:eq:C.SKD.3\endcsname{\href{S6_SERIES_GRIND_II.pdf\#nameddest=tag-011V}{II.{C.SKD.3} (p.~594)} }
\expandafter\def\csname SRef@cdp:eq:C.SKD.4\endcsname{\href{S6_SERIES_GRIND_II.pdf\#nameddest=tag-011W}{II.{C.SKD.4} (p.~595)} }
\expandafter\def\csname SRef@cdp:eq:C.SKD.5\endcsname{\href{S6_SERIES_GRIND_II.pdf\#nameddest=tag-011X}{II.{C.SKD.5} (p.~595)} }
\expandafter\def\csname SRef@cdp:eq:C.SKD.6\endcsname{\href{S6_SERIES_GRIND_II.pdf\#nameddest=tag-011Y}{II.{C.SKD.6} (p.~595)} }
\expandafter\def\csname SRef@cdp:eq:C.SKD.7\endcsname{\href{S6_SERIES_GRIND_II.pdf\#nameddest=tag-011Z}{II.{C.SKD.7} (p.~595)} }
\expandafter\def\csname SRef@cdp:eq:C.SKD.8\endcsname{\href{S6_SERIES_GRIND_II.pdf\#nameddest=tag-0120}{II.{C.SKD.8} (p.~595)} }
\expandafter\def\csname SRef@cdp:eq:C.SKD.9\endcsname{\href{S6_SERIES_GRIND_II.pdf\#nameddest=tag-0121}{II.{C.SKD.9} (p.~596)} }
\expandafter\def\csname SRef@cdp:eq:C.SKD.10\endcsname{\href{S6_SERIES_GRIND_II.pdf\#nameddest=tag-011O}{II.{C.SKD.10} (p.~596)} }
\expandafter\def\csname SRef@cdp:eq:C.SKD.11\endcsname{\href{S6_SERIES_GRIND_II.pdf\#nameddest=tag-011P}{II.{C.SKD.11} (p.~596)} }
\expandafter\def\csname SRef@cdp:eq:C.SKD.14\endcsname{\href{S6_SERIES_GRIND_II.pdf\#nameddest=tag-011S}{II.{C.SKD.14} (p.~596)} }
\expandafter\def\csname SRef@cdp:eq:C.CA.11\endcsname{\href{S6_SERIES_GRIND_II.pdf\#nameddest=tag-00BE}{II.{C.CA.11} (p.~430)} }
\expandafter\def\csname SRef@cdp:eq:C.SKD.15\endcsname{\href{S6_SERIES_GRIND_II.pdf\#nameddest=tag-011T}{II.{C.SKD.15} (p.~597)} }
\expandafter\def\csname SRef@cdp:eq:C.SKD.12\endcsname{\href{S6_SERIES_GRIND_II.pdf\#nameddest=tag-011Q}{II.{C.SKD.12} (p.~597)} }
\expandafter\def\csname SRef@cdp:eq:C.SKD.13\endcsname{\href{S6_SERIES_GRIND_II.pdf\#nameddest=tag-011R}{II.{C.SKD.13} (p.~597)} }
\expandafter\def\csname SRef@cdp:eq:C.EN.1\endcsname{\href{S6_SERIES_GRIND_II.pdf\#nameddest=tag-00G9}{II.{C.EN.1} (p.~598)} }
\expandafter\def\csname SRef@cdp:eq:C.EN.18\endcsname{\href{S6_SERIES_GRIND_II.pdf\#nameddest=tag-00GI}{II.{C.EN.18} (p.~598)} }
\expandafter\def\csname SRef@cdp:eq:C.EN.2\endcsname{\href{S6_SERIES_GRIND_II.pdf\#nameddest=tag-00GJ}{II.{C.EN.2} (p.~598)} }
\expandafter\def\csname SRef@cdp:eq:C.EN.3\endcsname{\href{S6_SERIES_GRIND_II.pdf\#nameddest=tag-00GK}{II.{C.EN.3} (p.~598)} }
\expandafter\def\csname SRef@cdp:eq:C.EN.4\endcsname{\href{S6_SERIES_GRIND_II.pdf\#nameddest=tag-00GL}{II.{C.EN.4} (p.~598)} }
\expandafter\def\csname SRef@cdp:eq:C.EN.5\endcsname{\href{S6_SERIES_GRIND_II.pdf\#nameddest=tag-00GM}{II.{C.EN.5} (p.~598)} }
\expandafter\def\csname SRef@cdp:eq:C.EN.6\endcsname{\href{S6_SERIES_GRIND_II.pdf\#nameddest=tag-00GN}{II.{C.EN.6} (p.~598)} }
\expandafter\def\csname SRef@cdp:eq:C.HQ.1\endcsname{\href{S6_SERIES_GRIND_II.pdf\#nameddest=tag-00MQ}{II.{C.HQ.1} (p.~431)} }
\expandafter\def\csname SRef@cdp:eq:C.EN.7\endcsname{\href{S6_SERIES_GRIND_II.pdf\#nameddest=tag-00GO}{II.{C.EN.7} (p.~599)} }
\expandafter\def\csname SRef@cdp:eq:C.EN.8\endcsname{\href{S6_SERIES_GRIND_II.pdf\#nameddest=tag-00GP}{II.{C.EN.8} (p.~599)} }
\expandafter\def\csname SRef@cdp:eq:C.EN.9\endcsname{\href{S6_SERIES_GRIND_II.pdf\#nameddest=tag-00GQ}{II.{C.EN.9} (p.~599)} }
\expandafter\def\csname SRef@cdp:eq:C.EN.10\endcsname{\href{S6_SERIES_GRIND_II.pdf\#nameddest=tag-00GA}{II.{C.EN.10} (p.~599)} }
\expandafter\def\csname SRef@cdp:eq:C.EN.11\endcsname{\href{S6_SERIES_GRIND_II.pdf\#nameddest=tag-00GB}{II.{C.EN.11} (p.~599)} }
\expandafter\def\csname SRef@cdp:eq:C.EN.12\endcsname{\href{S6_SERIES_GRIND_II.pdf\#nameddest=tag-00GC}{II.{C.EN.12} (p.~599)} }
\expandafter\def\csname SRef@cdp:eq:C.EN.13\endcsname{\href{S6_SERIES_GRIND_II.pdf\#nameddest=tag-00GD}{II.{C.EN.13} (p.~599)} }
\expandafter\def\csname SRef@cdp:eq:C.EN.14\endcsname{\href{S6_SERIES_GRIND_II.pdf\#nameddest=tag-00GE}{II.{C.EN.14} (p.~600)} }
\expandafter\def\csname SRef@cdp:eq:C.EN.15\endcsname{\href{S6_SERIES_GRIND_II.pdf\#nameddest=tag-00GF}{II.{C.EN.15} (p.~600)} }
\expandafter\def\csname SRef@cdp:eq:C.EN.16\endcsname{\href{S6_SERIES_GRIND_II.pdf\#nameddest=tag-00GG}{II.{C.EN.16} (p.~600)} }
\expandafter\def\csname SRef@cdp:eq:C.HQ.2\endcsname{\href{S6_SERIES_GRIND_II.pdf\#nameddest=tag-00MS}{II.{C.HQ.2} (p.~431)} }
\expandafter\def\csname SRef@cdp:eq:C.EN.17\endcsname{\href{S6_SERIES_GRIND_II.pdf\#nameddest=tag-00GH}{II.{C.EN.17} (p.~600)} }
\expandafter\def\csname SRef@cdp:eq:C.EN.S1\endcsname{\href{S6_SERIES_GRIND_II.pdf\#nameddest=tag-00H7}{II.{C.EN.S1} (p.~601)} }
\expandafter\def\csname SRef@cdp:eq:C.EN.S2\endcsname{\href{S6_SERIES_GRIND_II.pdf\#nameddest=tag-00H8}{II.{C.EN.S2} (p.~601)} }
\expandafter\def\csname SRef@cdp:eq:C.EN.S3\endcsname{\href{S6_SERIES_GRIND_II.pdf\#nameddest=tag-00H9}{II.{C.EN.S3} (p.~601)} }
\expandafter\def\csname SRef@cdp:eq:C.EN.S4\endcsname{\href{S6_SERIES_GRIND_II.pdf\#nameddest=tag-00HA}{II.{C.EN.S4} (p.~602)} }
\expandafter\def\csname SRef@cdp:eq:C.EN.S5\endcsname{\href{S6_SERIES_GRIND_II.pdf\#nameddest=tag-00HB}{II.{C.EN.S5} (p.~602)} }
\expandafter\def\csname SRef@cdp:eq:C.EN.S6\endcsname{\href{S6_SERIES_GRIND_II.pdf\#nameddest=tag-00HC}{II.{C.EN.S6} (p.~602)} }
\expandafter\def\csname SRef@cdp:eq:C.EN.S7\endcsname{\href{S6_SERIES_GRIND_II.pdf\#nameddest=tag-00HD}{II.{C.EN.S7} (p.~603)} }
\expandafter\def\csname SRef@cdp:eq:C.EN.R1\endcsname{\href{S6_SERIES_GRIND_II.pdf\#nameddest=tag-00GV}{II.{C.EN.R1} (p.~603)} }
\expandafter\def\csname SRef@cdp:eq:C.EN.R2\endcsname{\href{S6_SERIES_GRIND_II.pdf\#nameddest=tag-00GZ}{II.{C.EN.R2} (p.~603)} }
\expandafter\def\csname SRef@cdp:eq:C.HQ.3\endcsname{\href{S6_SERIES_GRIND_II.pdf\#nameddest=tag-00MT}{II.{C.HQ.3} (p.~432)} }
\expandafter\def\csname SRef@cdp:eq:C.EN.R3\endcsname{\href{S6_SERIES_GRIND_II.pdf\#nameddest=tag-00H0}{II.{C.EN.R3} (p.~603)} }
\expandafter\def\csname SRef@cdp:eq:C.EN.R4\endcsname{\href{S6_SERIES_GRIND_II.pdf\#nameddest=tag-00H1}{II.{C.EN.R4} (p.~603)} }
\expandafter\def\csname SRef@cdp:eq:C.EN.R5\endcsname{\href{S6_SERIES_GRIND_II.pdf\#nameddest=tag-00H2}{II.{C.EN.R5} (p.~603)} }
\expandafter\def\csname SRef@cdp:eq:C.EN.R6\endcsname{\href{S6_SERIES_GRIND_II.pdf\#nameddest=tag-00H3}{II.{C.EN.R6} (p.~604)} }
\expandafter\def\csname SRef@cdp:eq:C.EN.R7\endcsname{\href{S6_SERIES_GRIND_II.pdf\#nameddest=tag-00H4}{II.{C.EN.R7} (p.~604)} }
\expandafter\def\csname SRef@cdp:eq:C.EN.R8\endcsname{\href{S6_SERIES_GRIND_II.pdf\#nameddest=tag-00H5}{II.{C.EN.R8} (p.~604)} }
\expandafter\def\csname SRef@cdp:eq:C.EN.R9\endcsname{\href{S6_SERIES_GRIND_II.pdf\#nameddest=tag-00H6}{II.{C.EN.R9} (p.~604)} }
\expandafter\def\csname SRef@cdp:eq:C.EN.R10\endcsname{\href{S6_SERIES_GRIND_II.pdf\#nameddest=tag-00GW}{II.{C.EN.R10} (p.~604)} }
\expandafter\def\csname SRef@cdp:eq:C.EN.R11\endcsname{\href{S6_SERIES_GRIND_II.pdf\#nameddest=tag-00GX}{II.{C.EN.R11} (p.~605)} }
\expandafter\def\csname SRef@cdp:eq:C.EN.R12\endcsname{\href{S6_SERIES_GRIND_II.pdf\#nameddest=tag-00GY}{II.{C.EN.R12} (p.~605)} }
\expandafter\def\csname SRef@cdp:eq:C.HQ.4\endcsname{\href{S6_SERIES_GRIND_II.pdf\#nameddest=tag-00MU}{II.{C.HQ.4} (p.~432)} }
\expandafter\def\csname SRef@cdp:eq:C.EN.G1\endcsname{\href{S6_SERIES_GRIND_II.pdf\#nameddest=tag-00GR}{II.{C.EN.G1} (p.~606)} }
\expandafter\def\csname SRef@cdp:eq:C.EN.G2\endcsname{\href{S6_SERIES_GRIND_II.pdf\#nameddest=tag-00GS}{II.{C.EN.G2} (p.~606)} }
\expandafter\def\csname SRef@cdp:eq:C.EN.G3\endcsname{\href{S6_SERIES_GRIND_II.pdf\#nameddest=tag-00GT}{II.{C.EN.G3} (p.~606)} }
\expandafter\def\csname SRef@cdp:eq:C.EN.G4\endcsname{\href{S6_SERIES_GRIND_II.pdf\#nameddest=tag-00GU}{II.{C.EN.G4} (p.~606)} }
\expandafter\def\csname SRef@cdp:eq:C.ET.1\endcsname{\href{S6_SERIES_GRIND_II.pdf\#nameddest=tag-00J5}{II.{C.ET.1} (p.~607)} }
\expandafter\def\csname SRef@cdp:eq:C.ET.2\endcsname{\href{S6_SERIES_GRIND_II.pdf\#nameddest=tag-00J9}{II.{C.ET.2} (p.~607)} }
\expandafter\def\csname SRef@cdp:eq:C.ET.3\endcsname{\href{S6_SERIES_GRIND_II.pdf\#nameddest=tag-00JA}{II.{C.ET.3} (p.~607)} }
\expandafter\def\csname SRef@cdp:eq:C.ET.4\endcsname{\href{S6_SERIES_GRIND_II.pdf\#nameddest=tag-00JB}{II.{C.ET.4} (p.~607)} }
\expandafter\def\csname SRef@cdp:eq:C.ET.5\endcsname{\href{S6_SERIES_GRIND_II.pdf\#nameddest=tag-00JC}{II.{C.ET.5} (p.~607)} }
\expandafter\def\csname SRef@cdp:eq:C.ET.6\endcsname{\href{S6_SERIES_GRIND_II.pdf\#nameddest=tag-00JD}{II.{C.ET.6} (p.~608)} }
\expandafter\def\csname SRef@cdp:eq:C.HQ.5\endcsname{\href{S6_SERIES_GRIND_II.pdf\#nameddest=tag-00MV}{II.{C.HQ.5} (p.~432)} }
\expandafter\def\csname SRef@cdp:eq:C.ET.7\endcsname{\href{S6_SERIES_GRIND_II.pdf\#nameddest=tag-00JE}{II.{C.ET.7} (p.~608)} }
\expandafter\def\csname SRef@cdp:eq:C.ET.8\endcsname{\href{S6_SERIES_GRIND_II.pdf\#nameddest=tag-00JF}{II.{C.ET.8} (p.~608)} }
\expandafter\def\csname SRef@cdp:eq:C.ET.9\endcsname{\href{S6_SERIES_GRIND_II.pdf\#nameddest=tag-00JG}{II.{C.ET.9} (p.~608)} }
\expandafter\def\csname SRef@cdp:eq:C.ET.10\endcsname{\href{S6_SERIES_GRIND_II.pdf\#nameddest=tag-00J6}{II.{C.ET.10} (p.~609)} }
\expandafter\def\csname SRef@cdp:eq:C.ET.11\endcsname{\href{S6_SERIES_GRIND_II.pdf\#nameddest=tag-00J7}{II.{C.ET.11} (p.~609)} }
\expandafter\def\csname SRef@cdp:eq:C.ET.12\endcsname{\href{S6_SERIES_GRIND_II.pdf\#nameddest=tag-00J8}{II.{C.ET.12} (p.~609)} }
\expandafter\def\csname SRef@cdp:eq:C.SF.1\endcsname{\href{S6_SERIES_GRIND_II.pdf\#nameddest=tag-010R}{II.{C.SF.1} (p.~611)} }
\expandafter\def\csname SRef@cdp:eq:C.SF.2\endcsname{\href{S6_SERIES_GRIND_II.pdf\#nameddest=tag-010W}{II.{C.SF.2} (p.~611)} }
\expandafter\def\csname SRef@cdp:eq:C.SF.3\endcsname{\href{S6_SERIES_GRIND_II.pdf\#nameddest=tag-010X}{II.{C.SF.3} (p.~611)} }
\expandafter\def\csname SRef@cdp:eq:C.SF.4\endcsname{\href{S6_SERIES_GRIND_II.pdf\#nameddest=tag-010Y}{II.{C.SF.4} (p.~611)} }
\expandafter\def\csname SRef@cdp:eq:C.HQ.6\endcsname{\href{S6_SERIES_GRIND_II.pdf\#nameddest=tag-00MW}{II.{C.HQ.6} (p.~432)} }
\expandafter\def\csname SRef@cdp:eq:C.SF.5\endcsname{\href{S6_SERIES_GRIND_II.pdf\#nameddest=tag-010Z}{II.{C.SF.5} (p.~611)} }
\expandafter\def\csname SRef@cdp:eq:C.SF.6\endcsname{\href{S6_SERIES_GRIND_II.pdf\#nameddest=tag-0110}{II.{C.SF.6} (p.~612)} }
\expandafter\def\csname SRef@cdp:eq:C.SF.7\endcsname{\href{S6_SERIES_GRIND_II.pdf\#nameddest=tag-0111}{II.{C.SF.7} (p.~612)} }
\expandafter\def\csname SRef@cdp:eq:C.SF.8\endcsname{\href{S6_SERIES_GRIND_II.pdf\#nameddest=tag-0112}{II.{C.SF.8} (p.~612)} }
\expandafter\def\csname SRef@cdp:eq:C.SF.9\endcsname{\href{S6_SERIES_GRIND_II.pdf\#nameddest=tag-0113}{II.{C.SF.9} (p.~612)} }
\expandafter\def\csname SRef@cdp:eq:C.SF.10\endcsname{\href{S6_SERIES_GRIND_II.pdf\#nameddest=tag-010S}{II.{C.SF.10} (p.~612)} }
\expandafter\def\csname SRef@cdp:eq:C.SF.11\endcsname{\href{S6_SERIES_GRIND_II.pdf\#nameddest=tag-010T}{II.{C.SF.11} (p.~613)} }
\expandafter\def\csname SRef@cdp:eq:C.SF.12\endcsname{\href{S6_SERIES_GRIND_II.pdf\#nameddest=tag-010U}{II.{C.SF.12} (p.~613)} }
\expandafter\def\csname SRef@cdp:eq:C.SF.13\endcsname{\href{S6_SERIES_GRIND_II.pdf\#nameddest=tag-010V}{II.{C.SF.13} (p.~613)} }
\expandafter\def\csname SRef@cdp:eq:C.EO.1\endcsname{\href{S6_SERIES_GRIND_II.pdf\#nameddest=tag-00HE}{II.{C.EO.1} (p.~614)} }
\expandafter\def\csname SRef@cdp:eq:C.HQ.7\endcsname{\href{S6_SERIES_GRIND_II.pdf\#nameddest=tag-00MX}{II.{C.HQ.7} (p.~432)} }
\expandafter\def\csname SRef@cdp:eq:C.EO.2\endcsname{\href{S6_SERIES_GRIND_II.pdf\#nameddest=tag-00HM}{II.{C.EO.2} (p.~614)} }
\expandafter\def\csname SRef@cdp:eq:C.EO.3\endcsname{\href{S6_SERIES_GRIND_II.pdf\#nameddest=tag-00HN}{II.{C.EO.3} (p.~614)} }
\expandafter\def\csname SRef@cdp:eq:C.EO.4\endcsname{\href{S6_SERIES_GRIND_II.pdf\#nameddest=tag-00HO}{II.{C.EO.4} (p.~614)} }
\expandafter\def\csname SRef@cdp:eq:C.EO.5\endcsname{\href{S6_SERIES_GRIND_II.pdf\#nameddest=tag-00HP}{II.{C.EO.5} (p.~615)} }
\expandafter\def\csname SRef@cdp:eq:C.EO.6\endcsname{\href{S6_SERIES_GRIND_II.pdf\#nameddest=tag-00HQ}{II.{C.EO.6} (p.~615)} }
\expandafter\def\csname SRef@cdp:eq:C.EO.H1\endcsname{\href{S6_SERIES_GRIND_II.pdf\#nameddest=tag-00HU}{II.{C.EO.H1} (p.~615)} }
\expandafter\def\csname SRef@cdp:eq:C.EO.H2\endcsname{\href{S6_SERIES_GRIND_II.pdf\#nameddest=tag-00HV}{II.{C.EO.H2} (p.~616)} }
\expandafter\def\csname SRef@cdp:eq:C.EO.H3\endcsname{\href{S6_SERIES_GRIND_II.pdf\#nameddest=tag-00HW}{II.{C.EO.H3} (p.~616)} }
\expandafter\def\csname SRef@cdp:eq:C.EO.H4\endcsname{\href{S6_SERIES_GRIND_II.pdf\#nameddest=tag-00HX}{II.{C.EO.H4} (p.~616)} }
\expandafter\def\csname SRef@cdp:eq:C.EO.H5\endcsname{\href{S6_SERIES_GRIND_II.pdf\#nameddest=tag-00HY}{II.{C.EO.H5} (p.~616)} }
\expandafter\def\csname SRef@cdp:eq:C.HQ.8\endcsname{\href{S6_SERIES_GRIND_II.pdf\#nameddest=tag-00MY}{II.{C.HQ.8} (p.~432)} }
\expandafter\def\csname SRef@cdp:eq:C.EO.H6\endcsname{\href{S6_SERIES_GRIND_II.pdf\#nameddest=tag-00HZ}{II.{C.EO.H6} (p.~617)} }
\expandafter\def\csname SRef@cdp:eq:C.EO.H7\endcsname{\href{S6_SERIES_GRIND_II.pdf\#nameddest=tag-00I0}{II.{C.EO.H7} (p.~617)} }
\expandafter\def\csname SRef@cdp:eq:C.EO.7\endcsname{\href{S6_SERIES_GRIND_II.pdf\#nameddest=tag-00HR}{II.{C.EO.7} (p.~617)} }
\expandafter\def\csname SRef@cdp:eq:C.EO.8\endcsname{\href{S6_SERIES_GRIND_II.pdf\#nameddest=tag-00HS}{II.{C.EO.8} (p.~617)} }
\expandafter\def\csname SRef@cdp:eq:C.EO.9\endcsname{\href{S6_SERIES_GRIND_II.pdf\#nameddest=tag-00HT}{II.{C.EO.9} (p.~617)} }
\expandafter\def\csname SRef@cdp:eq:C.EO.10\endcsname{\href{S6_SERIES_GRIND_II.pdf\#nameddest=tag-00HF}{II.{C.EO.10} (p.~618)} }
\expandafter\def\csname SRef@cdp:eq:C.EO.11\endcsname{\href{S6_SERIES_GRIND_II.pdf\#nameddest=tag-00HG}{II.{C.EO.11} (p.~618)} }
\expandafter\def\csname SRef@cdp:eq:C.EO.12\endcsname{\href{S6_SERIES_GRIND_II.pdf\#nameddest=tag-00HH}{II.{C.EO.12} (p.~618)} }
\expandafter\def\csname SRef@cdp:eq:C.EO.13\endcsname{\href{S6_SERIES_GRIND_II.pdf\#nameddest=tag-00HI}{II.{C.EO.13} (p.~618)} }
\expandafter\def\csname SRef@cdp:eq:C.EO.14\endcsname{\href{S6_SERIES_GRIND_II.pdf\#nameddest=tag-00HJ}{II.{C.EO.14} (p.~619)} }
\expandafter\def\csname SRef@cdp:eq:C.HQ.9\endcsname{\href{S6_SERIES_GRIND_II.pdf\#nameddest=tag-00MZ}{II.{C.HQ.9} (p.~433)} }
\expandafter\def\csname SRef@cdp:eq:C.EO.15\endcsname{\href{S6_SERIES_GRIND_II.pdf\#nameddest=tag-00HK}{II.{C.EO.15} (p.~619)} }
\expandafter\def\csname SRef@cdp:eq:C.EO.16\endcsname{\href{S6_SERIES_GRIND_II.pdf\#nameddest=tag-00HL}{II.{C.EO.16} (p.~620)} }
\expandafter\def\csname SRef@cdp:eq:C.EH.1\endcsname{\href{S6_SERIES_GRIND_II.pdf\#nameddest=tag-00FQ}{II.{C.EH.1} (p.~620)} }
\expandafter\def\csname SRef@cdp:eq:C.EH.2\endcsname{\href{S6_SERIES_GRIND_II.pdf\#nameddest=tag-00FR}{II.{C.EH.2} (p.~620)} }
\expandafter\def\csname SRef@cdp:eq:C.EH.3\endcsname{\href{S6_SERIES_GRIND_II.pdf\#nameddest=tag-00FS}{II.{C.EH.3} (p.~621)} }
\expandafter\def\csname SRef@cdp:eq:C.EH.4\endcsname{\href{S6_SERIES_GRIND_II.pdf\#nameddest=tag-00FT}{II.{C.EH.4} (p.~621)} }
\expandafter\def\csname SRef@cdp:eq:C.EH.5\endcsname{\href{S6_SERIES_GRIND_II.pdf\#nameddest=tag-00FU}{II.{C.EH.5} (p.~621)} }
\expandafter\def\csname SRef@cdp:eq:C.EH.6\endcsname{\href{S6_SERIES_GRIND_II.pdf\#nameddest=tag-00FV}{II.{C.EH.6} (p.~621)} }
\expandafter\def\csname SRef@cdp:eq:C.EH.7\endcsname{\href{S6_SERIES_GRIND_II.pdf\#nameddest=tag-00FW}{II.{C.EH.7} (p.~621)} }
\expandafter\def\csname SRef@cdp:eq:C.EH.8\endcsname{\href{S6_SERIES_GRIND_II.pdf\#nameddest=tag-00FX}{II.{C.EH.8} (p.~621)} }
\expandafter\def\csname SRef@cdp:eq:C.HQ.10\endcsname{\href{S6_SERIES_GRIND_II.pdf\#nameddest=tag-00MR}{II.{C.HQ.10} (p.~433)} }
\expandafter\def\csname SRef@cdp:eq:C.DH.1\endcsname{\href{S6_SERIES_GRIND_II.pdf\#nameddest=tag-00F6}{II.{C.DH.1} (p.~622)} }
\expandafter\def\csname SRef@cdp:eq:C.DH.2\endcsname{\href{S6_SERIES_GRIND_II.pdf\#nameddest=tag-00FG}{II.{C.DH.2} (p.~622)} }
\expandafter\def\csname SRef@cdp:eq:C.DH.3\endcsname{\href{S6_SERIES_GRIND_II.pdf\#nameddest=tag-00FH}{II.{C.DH.3} (p.~622)} }
\expandafter\def\csname SRef@cdp:eq:C.DH.4\endcsname{\href{S6_SERIES_GRIND_II.pdf\#nameddest=tag-00FI}{II.{C.DH.4} (p.~622)} }
\expandafter\def\csname SRef@cdp:eq:C.DH.5\endcsname{\href{S6_SERIES_GRIND_II.pdf\#nameddest=tag-00FJ}{II.{C.DH.5} (p.~623)} }
\expandafter\def\csname SRef@cdp:eq:C.DH.6\endcsname{\href{S6_SERIES_GRIND_II.pdf\#nameddest=tag-00FK}{II.{C.DH.6} (p.~623)} }
\expandafter\def\csname SRef@cdp:eq:C.DH.7\endcsname{\href{S6_SERIES_GRIND_II.pdf\#nameddest=tag-00FL}{II.{C.DH.7} (p.~623)} }
\expandafter\def\csname SRef@cdp:eq:C.DH.8\endcsname{\href{S6_SERIES_GRIND_II.pdf\#nameddest=tag-00FM}{II.{C.DH.8} (p.~623)} }
\expandafter\def\csname SRef@cdp:eq:C.DH.9\endcsname{\href{S6_SERIES_GRIND_II.pdf\#nameddest=tag-00FN}{II.{C.DH.9} (p.~623)} }
\expandafter\def\csname SRef@cdp:eq:C.DH.10\endcsname{\href{S6_SERIES_GRIND_II.pdf\#nameddest=tag-00F7}{II.{C.DH.10} (p.~624)} }
\expandafter\def\csname SRef@cdp:eq:C.FD.1\endcsname{\href{S6_SERIES_GRIND_II.pdf\#nameddest=tag-00K2}{II.{C.FD.1} (p.~433)} }
\expandafter\def\csname SRef@cdp:eq:C.DH.11\endcsname{\href{S6_SERIES_GRIND_II.pdf\#nameddest=tag-00F8}{II.{C.DH.11} (p.~624)} }
\expandafter\def\csname SRef@cdp:eq:C.DH.12\endcsname{\href{S6_SERIES_GRIND_II.pdf\#nameddest=tag-00F9}{II.{C.DH.12} (p.~624)} }
\expandafter\def\csname SRef@cdp:eq:C.DH.13\endcsname{\href{S6_SERIES_GRIND_II.pdf\#nameddest=tag-00FA}{II.{C.DH.13} (p.~625)} }
\expandafter\def\csname SRef@cdp:eq:C.DH.14\endcsname{\href{S6_SERIES_GRIND_II.pdf\#nameddest=tag-00FB}{II.{C.DH.14} (p.~625)} }
\expandafter\def\csname SRef@cdp:eq:C.DH.14a\endcsname{\href{S6_SERIES_GRIND_II.pdf\#nameddest=tag-00FC}{II.{C.DH.14a} (p.~625)} }
\expandafter\def\csname SRef@cdp:eq:C.DH.14b\endcsname{\href{S6_SERIES_GRIND_II.pdf\#nameddest=tag-00FD}{II.{C.DH.14b} (p.~625)} }
\expandafter\def\csname SRef@cdp:eq:C.DH.15\endcsname{\href{S6_SERIES_GRIND_II.pdf\#nameddest=tag-00FE}{II.{C.DH.15} (p.~625)} }
\expandafter\def\csname SRef@cdp:eq:C.DH.16\endcsname{\href{S6_SERIES_GRIND_II.pdf\#nameddest=tag-00FF}{II.{C.DH.16} (p.~626)} }
\expandafter\def\csname SRef@cdp:eq:C.NH.1\endcsname{\href{S6_SERIES_GRIND_II.pdf\#nameddest=tag-00S9}{II.{C.NH.1} (p.~626)} }
\expandafter\def\csname SRef@cdp:eq:C.NH.2\endcsname{\href{S6_SERIES_GRIND_II.pdf\#nameddest=tag-00SJ}{II.{C.NH.2} (p.~627)} }
\expandafter\def\csname SRef@cdp:eq:C.FD.2\endcsname{\href{S6_SERIES_GRIND_II.pdf\#nameddest=tag-00K3}{II.{C.FD.2} (p.~434)} }
\expandafter\def\csname SRef@cdp:eq:C.NH.3\endcsname{\href{S6_SERIES_GRIND_II.pdf\#nameddest=tag-00SK}{II.{C.NH.3} (p.~627)} }
\expandafter\def\csname SRef@cdp:eq:C.NH.4\endcsname{\href{S6_SERIES_GRIND_II.pdf\#nameddest=tag-00SL}{II.{C.NH.4} (p.~627)} }
\expandafter\def\csname SRef@cdp:eq:C.NH.5\endcsname{\href{S6_SERIES_GRIND_II.pdf\#nameddest=tag-00SM}{II.{C.NH.5} (p.~627)} }
\expandafter\def\csname SRef@cdp:eq:C.NH.6\endcsname{\href{S6_SERIES_GRIND_II.pdf\#nameddest=tag-00SN}{II.{C.NH.6} (p.~627)} }
\expandafter\def\csname SRef@cdp:eq:C.NH.7\endcsname{\href{S6_SERIES_GRIND_II.pdf\#nameddest=tag-00SO}{II.{C.NH.7} (p.~627)} }
\expandafter\def\csname SRef@cdp:eq:C.NH.8\endcsname{\href{S6_SERIES_GRIND_II.pdf\#nameddest=tag-00SP}{II.{C.NH.8} (p.~628)} }
\expandafter\def\csname SRef@cdp:eq:C.NH.9\endcsname{\href{S6_SERIES_GRIND_II.pdf\#nameddest=tag-00SQ}{II.{C.NH.9} (p.~628)} }
\expandafter\def\csname SRef@cdp:eq:C.NH.10\endcsname{\href{S6_SERIES_GRIND_II.pdf\#nameddest=tag-00SA}{II.{C.NH.10} (p.~628)} }
\expandafter\def\csname SRef@cdp:eq:C.NH.11\endcsname{\href{S6_SERIES_GRIND_II.pdf\#nameddest=tag-00SB}{II.{C.NH.11} (p.~628)} }
\expandafter\def\csname SRef@cdp:eq:C.NH.12\endcsname{\href{S6_SERIES_GRIND_II.pdf\#nameddest=tag-00SC}{II.{C.NH.12} (p.~629)} }
\expandafter\def\csname SRef@cdp:eq:C.FD.3\endcsname{\href{S6_SERIES_GRIND_II.pdf\#nameddest=tag-00K4}{II.{C.FD.3} (p.~434)} }
\expandafter\def\csname SRef@cdp:eq:C.NH.13\endcsname{\href{S6_SERIES_GRIND_II.pdf\#nameddest=tag-00SD}{II.{C.NH.13} (p.~629)} }
\expandafter\def\csname SRef@cdp:eq:C.NH.14\endcsname{\href{S6_SERIES_GRIND_II.pdf\#nameddest=tag-00SE}{II.{C.NH.14} (p.~629)} }
\expandafter\def\csname SRef@cdp:eq:C.NH.15\endcsname{\href{S6_SERIES_GRIND_II.pdf\#nameddest=tag-00SF}{II.{C.NH.15} (p.~629)} }
\expandafter\def\csname SRef@cdp:eq:C.NH.16\endcsname{\href{S6_SERIES_GRIND_II.pdf\#nameddest=tag-00SG}{II.{C.NH.16} (p.~629)} }
\expandafter\def\csname SRef@cdp:eq:C.NH.17\endcsname{\href{S6_SERIES_GRIND_II.pdf\#nameddest=tag-00SH}{II.{C.NH.17} (p.~629)} }
\expandafter\def\csname SRef@cdp:eq:C.NH.18\endcsname{\href{S6_SERIES_GRIND_II.pdf\#nameddest=tag-00SI}{II.{C.NH.18} (p.~630)} }
\expandafter\def\csname SRef@cdp:eq:C.SH.1\endcsname{\href{S6_SERIES_GRIND_II.pdf\#nameddest=tag-0114}{II.{C.SH.1} (p.~630)} }
\expandafter\def\csname SRef@cdp:eq:C.SH.2\endcsname{\href{S6_SERIES_GRIND_II.pdf\#nameddest=tag-0119}{II.{C.SH.2} (p.~631)} }
\expandafter\def\csname SRef@cdp:eq:C.SH.3\endcsname{\href{S6_SERIES_GRIND_II.pdf\#nameddest=tag-011A}{II.{C.SH.3} (p.~631)} }
\expandafter\def\csname SRef@cdp:eq:C.SH.4\endcsname{\href{S6_SERIES_GRIND_II.pdf\#nameddest=tag-011B}{II.{C.SH.4} (p.~631)} }
\expandafter\def\csname SRef@cdp:eq:C.FD.4\endcsname{\href{S6_SERIES_GRIND_II.pdf\#nameddest=tag-00K5}{II.{C.FD.4} (p.~434)} }
\expandafter\def\csname SRef@cdp:eq:C.SH.5\endcsname{\href{S6_SERIES_GRIND_II.pdf\#nameddest=tag-011C}{II.{C.SH.5} (p.~632)} }
\expandafter\def\csname SRef@cdp:eq:C.SH.6\endcsname{\href{S6_SERIES_GRIND_II.pdf\#nameddest=tag-011D}{II.{C.SH.6} (p.~632)} }
\expandafter\def\csname SRef@cdp:eq:C.SH.7\endcsname{\href{S6_SERIES_GRIND_II.pdf\#nameddest=tag-011E}{II.{C.SH.7} (p.~632)} }
\expandafter\def\csname SRef@cdp:eq:C.SH.8\endcsname{\href{S6_SERIES_GRIND_II.pdf\#nameddest=tag-011F}{II.{C.SH.8} (p.~632)} }
\expandafter\def\csname SRef@cdp:eq:C.SH.9\endcsname{\href{S6_SERIES_GRIND_II.pdf\#nameddest=tag-011G}{II.{C.SH.9} (p.~632)} }
\expandafter\def\csname SRef@cdp:eq:C.SH.10\endcsname{\href{S6_SERIES_GRIND_II.pdf\#nameddest=tag-0115}{II.{C.SH.10} (p.~632)} }
\expandafter\def\csname SRef@cdp:eq:C.SH.11\endcsname{\href{S6_SERIES_GRIND_II.pdf\#nameddest=tag-0116}{II.{C.SH.11} (p.~633)} }
\expandafter\def\csname SRef@cdp:eq:C.SH.12\endcsname{\href{S6_SERIES_GRIND_II.pdf\#nameddest=tag-0117}{II.{C.SH.12} (p.~633)} }
\expandafter\def\csname SRef@cdp:eq:C.SH.13\endcsname{\href{S6_SERIES_GRIND_II.pdf\#nameddest=tag-0118}{II.{C.SH.13} (p.~633)} }
\expandafter\def\csname SRef@cdp:eq:C.RU.1\endcsname{\href{S6_SERIES_GRIND_II.pdf\#nameddest=tag-010C}{II.{C.RU.1} (p.~634)} }
\expandafter\def\csname SRef@cdp:eq:C.FD.5\endcsname{\href{S6_SERIES_GRIND_II.pdf\#nameddest=tag-00K6}{II.{C.FD.5} (p.~434)} }
\expandafter\def\csname SRef@cdp:eq:C.RU.2\endcsname{\href{S6_SERIES_GRIND_II.pdf\#nameddest=tag-010D}{II.{C.RU.2} (p.~634)} }
\expandafter\def\csname SRef@cdp:eq:C.RU.3\endcsname{\href{S6_SERIES_GRIND_II.pdf\#nameddest=tag-010E}{II.{C.RU.3} (p.~634)} }
\expandafter\def\csname SRef@cdp:eq:C.RU.4\endcsname{\href{S6_SERIES_GRIND_II.pdf\#nameddest=tag-010F}{II.{C.RU.4} (p.~634)} }
\expandafter\def\csname SRef@cdp:eq:C.RU.5\endcsname{\href{S6_SERIES_GRIND_II.pdf\#nameddest=tag-010G}{II.{C.RU.5} (p.~634)} }
\expandafter\def\csname SRef@cdp:eq:C.RU.6\endcsname{\href{S6_SERIES_GRIND_II.pdf\#nameddest=tag-010H}{II.{C.RU.6} (p.~635)} }
\expandafter\def\csname SRef@cdp:eq:C.RU.7\endcsname{\href{S6_SERIES_GRIND_II.pdf\#nameddest=tag-010I}{II.{C.RU.7} (p.~635)} }
\expandafter\def\csname SRef@cdp:eq:C.RU.8\endcsname{\href{S6_SERIES_GRIND_II.pdf\#nameddest=tag-010J}{II.{C.RU.8} (p.~635)} }
\expandafter\def\csname SRef@cdp:eq:C.RU.9\endcsname{\href{S6_SERIES_GRIND_II.pdf\#nameddest=tag-010K}{II.{C.RU.9} (p.~636)} }
\expandafter\def\csname SRef@cdp:eq:C.HA.1\endcsname{\href{S6_SERIES_GRIND_II.pdf\#nameddest=tag-00L9}{II.{C.HA.1} (p.~636)} }
\expandafter\def\csname SRef@cdp:eq:C.HA.2\endcsname{\href{S6_SERIES_GRIND_II.pdf\#nameddest=tag-00LA}{II.{C.HA.2} (p.~636)} }
\expandafter\def\csname SRef@cdp:eq:C.FD.6\endcsname{\href{S6_SERIES_GRIND_II.pdf\#nameddest=tag-00K7}{II.{C.FD.6} (p.~434)} }
\expandafter\def\csname SRef@cdp:eq:C.HA.3\endcsname{\href{S6_SERIES_GRIND_II.pdf\#nameddest=tag-00LB}{II.{C.HA.3} (p.~637)} }
\expandafter\def\csname SRef@cdp:eq:C.HA.4\endcsname{\href{S6_SERIES_GRIND_II.pdf\#nameddest=tag-00LC}{II.{C.HA.4} (p.~637)} }
\expandafter\def\csname SRef@cdp:eq:C.HA.5\endcsname{\href{S6_SERIES_GRIND_II.pdf\#nameddest=tag-00LD}{II.{C.HA.5} (p.~638)} }
\expandafter\def\csname SRef@cdp:eq:C.HA.6\endcsname{\href{S6_SERIES_GRIND_II.pdf\#nameddest=tag-00LE}{II.{C.HA.6} (p.~638)} }
\expandafter\def\csname SRef@cdp:eq:C.HC.1\endcsname{\href{S6_SERIES_GRIND_II.pdf\#nameddest=tag-00LO}{II.{C.HC.1} (p.~638)} }
\expandafter\def\csname SRef@cdp:eq:C.HC.2\endcsname{\href{S6_SERIES_GRIND_II.pdf\#nameddest=tag-00LR}{II.{C.HC.2} (p.~638)} }
\expandafter\def\csname SRef@cdp:eq:C.HC.3\endcsname{\href{S6_SERIES_GRIND_II.pdf\#nameddest=tag-00LS}{II.{C.HC.3} (p.~639)} }
\expandafter\def\csname SRef@cdp:eq:C.HC.4\endcsname{\href{S6_SERIES_GRIND_II.pdf\#nameddest=tag-00LT}{II.{C.HC.4} (p.~639)} }
\expandafter\def\csname SRef@cdp:eq:C.HC.5\endcsname{\href{S6_SERIES_GRIND_II.pdf\#nameddest=tag-00LU}{II.{C.HC.5} (p.~640)} }
\expandafter\def\csname SRef@cdp:eq:C.HC.6\endcsname{\href{S6_SERIES_GRIND_II.pdf\#nameddest=tag-00LV}{II.{C.HC.6} (p.~640)} }
\expandafter\def\csname SRef@cdp:eq:C.RF.5\endcsname{\href{S6_SERIES_GRIND_II.pdf\#nameddest=tag-00YY}{II.{C.RF.5} (p.~435)} }
\expandafter\def\csname SRef@cdp:eq:C.HC.7\endcsname{\href{S6_SERIES_GRIND_II.pdf\#nameddest=tag-00LW}{II.{C.HC.7} (p.~640)} }
\expandafter\def\csname SRef@cdp:eq:C.HC.8\endcsname{\href{S6_SERIES_GRIND_II.pdf\#nameddest=tag-00LX}{II.{C.HC.8} (p.~641)} }
\expandafter\def\csname SRef@cdp:eq:C.HC.9\endcsname{\href{S6_SERIES_GRIND_II.pdf\#nameddest=tag-00LY}{II.{C.HC.9} (p.~641)} }
\expandafter\def\csname SRef@cdp:eq:C.HC.10\endcsname{\href{S6_SERIES_GRIND_II.pdf\#nameddest=tag-00LP}{II.{C.HC.10} (p.~641)} }
\expandafter\def\csname SRef@cdp:eq:C.HC.11\endcsname{\href{S6_SERIES_GRIND_II.pdf\#nameddest=tag-00LQ}{II.{C.HC.11} (p.~642)} }
\expandafter\def\csname SRef@cdp:eq:C.HG.1\endcsname{\href{S6_SERIES_GRIND_II.pdf\#nameddest=tag-00LZ}{II.{C.HG.1} (p.~642)} }
\expandafter\def\csname SRef@cdp:eq:C.HG.2\endcsname{\href{S6_SERIES_GRIND_II.pdf\#nameddest=tag-00M8}{II.{C.HG.2} (p.~643)} }
\expandafter\def\csname SRef@cdp:eq:C.HG.3\endcsname{\href{S6_SERIES_GRIND_II.pdf\#nameddest=tag-00M9}{II.{C.HG.3} (p.~643)} }
\expandafter\def\csname SRef@cdp:eq:C.HG.4\endcsname{\href{S6_SERIES_GRIND_II.pdf\#nameddest=tag-00MA}{II.{C.HG.4} (p.~643)} }
\expandafter\def\csname SRef@cdp:eq:C.HG.5\endcsname{\href{S6_SERIES_GRIND_II.pdf\#nameddest=tag-00MB}{II.{C.HG.5} (p.~643)} }
\expandafter\def\csname SRef@cdp:eq:C.RF.1\endcsname{\href{S6_SERIES_GRIND_II.pdf\#nameddest=tag-00YU}{II.{C.RF.1} (p.~435)} }
\expandafter\def\csname SRef@cdp:eq:C.HG.6\endcsname{\href{S6_SERIES_GRIND_II.pdf\#nameddest=tag-00MC}{II.{C.HG.6} (p.~643)} }
\expandafter\def\csname SRef@cdp:eq:C.HG.7\endcsname{\href{S6_SERIES_GRIND_II.pdf\#nameddest=tag-00MD}{II.{C.HG.7} (p.~643)} }
\expandafter\def\csname SRef@cdp:eq:C.HG.8\endcsname{\href{S6_SERIES_GRIND_II.pdf\#nameddest=tag-00ME}{II.{C.HG.8} (p.~644)} }
\expandafter\def\csname SRef@cdp:eq:C.HG.9\endcsname{\href{S6_SERIES_GRIND_II.pdf\#nameddest=tag-00MF}{II.{C.HG.9} (p.~644)} }
\expandafter\def\csname SRef@cdp:eq:C.HG.10\endcsname{\href{S6_SERIES_GRIND_II.pdf\#nameddest=tag-00M0}{II.{C.HG.10} (p.~644)} }
\expandafter\def\csname SRef@cdp:eq:C.HG.11\endcsname{\href{S6_SERIES_GRIND_II.pdf\#nameddest=tag-00M1}{II.{C.HG.11} (p.~644)} }
\expandafter\def\csname SRef@cdp:eq:C.HG.12\endcsname{\href{S6_SERIES_GRIND_II.pdf\#nameddest=tag-00M2}{II.{C.HG.12} (p.~644)} }
\expandafter\def\csname SRef@cdp:eq:C.HG.13\endcsname{\href{S6_SERIES_GRIND_II.pdf\#nameddest=tag-00M3}{II.{C.HG.13} (p.~645)} }
\expandafter\def\csname SRef@cdp:eq:C.HG.14\endcsname{\href{S6_SERIES_GRIND_II.pdf\#nameddest=tag-00M4}{II.{C.HG.14} (p.~645)} }
\expandafter\def\csname SRef@cdp:eq:C.HG.15\endcsname{\href{S6_SERIES_GRIND_II.pdf\#nameddest=tag-00M5}{II.{C.HG.15} (p.~646)} }
\expandafter\def\csname SRef@cdp:eq:C.RF.2\endcsname{\href{S6_SERIES_GRIND_II.pdf\#nameddest=tag-00YV}{II.{C.RF.2} (p.~436)} }
\expandafter\def\csname SRef@cdp:eq:C.HG.16\endcsname{\href{S6_SERIES_GRIND_II.pdf\#nameddest=tag-00M6}{II.{C.HG.16} (p.~646)} }
\expandafter\def\csname SRef@cdp:eq:C.HG.17\endcsname{\href{S6_SERIES_GRIND_II.pdf\#nameddest=tag-00M7}{II.{C.HG.17} (p.~646)} }
\expandafter\def\csname SRef@cdp:eq:C.NS.1\endcsname{\href{S6_SERIES_GRIND_II.pdf\#nameddest=tag-00TA}{II.{C.NS.1} (p.~647)} }
\expandafter\def\csname SRef@cdp:eq:C.NS.2\endcsname{\href{S6_SERIES_GRIND_II.pdf\#nameddest=tag-00TL}{II.{C.NS.2} (p.~647)} }
\expandafter\def\csname SRef@cdp:eq:C.NS.3\endcsname{\href{S6_SERIES_GRIND_II.pdf\#nameddest=tag-00TN}{II.{C.NS.3} (p.~647)} }
\expandafter\def\csname SRef@cdp:eq:C.NS.4\endcsname{\href{S6_SERIES_GRIND_II.pdf\#nameddest=tag-00TO}{II.{C.NS.4} (p.~648)} }
\expandafter\def\csname SRef@cdp:eq:C.NS.5\endcsname{\href{S6_SERIES_GRIND_II.pdf\#nameddest=tag-00TP}{II.{C.NS.5} (p.~648)} }
\expandafter\def\csname SRef@cdp:eq:C.NS.6\endcsname{\href{S6_SERIES_GRIND_II.pdf\#nameddest=tag-00TQ}{II.{C.NS.6} (p.~648)} }
\expandafter\def\csname SRef@cdp:eq:C.NS.7\endcsname{\href{S6_SERIES_GRIND_II.pdf\#nameddest=tag-00TR}{II.{C.NS.7} (p.~648)} }
\expandafter\def\csname SRef@cdp:eq:C.NS.8\endcsname{\href{S6_SERIES_GRIND_II.pdf\#nameddest=tag-00TS}{II.{C.NS.8} (p.~649)} }
\expandafter\def\csname SRef@cdp:eq:C.RF.3\endcsname{\href{S6_SERIES_GRIND_II.pdf\#nameddest=tag-00YW}{II.{C.RF.3} (p.~436)} }
\expandafter\def\csname SRef@cdp:eq:C.NS.9\endcsname{\href{S6_SERIES_GRIND_II.pdf\#nameddest=tag-00TT}{II.{C.NS.9} (p.~649)} }
\expandafter\def\csname SRef@cdp:eq:C.NS.10\endcsname{\href{S6_SERIES_GRIND_II.pdf\#nameddest=tag-00TB}{II.{C.NS.10} (p.~649)} }
\expandafter\def\csname SRef@cdp:eq:C.NS.11\endcsname{\href{S6_SERIES_GRIND_II.pdf\#nameddest=tag-00TC}{II.{C.NS.11} (p.~649)} }
\expandafter\def\csname SRef@cdp:eq:C.NS.12\endcsname{\href{S6_SERIES_GRIND_II.pdf\#nameddest=tag-00TD}{II.{C.NS.12} (p.~649)} }
\expandafter\def\csname SRef@cdp:eq:C.NS.13\endcsname{\href{S6_SERIES_GRIND_II.pdf\#nameddest=tag-00TE}{II.{C.NS.13} (p.~650)} }
\expandafter\def\csname SRef@cdp:eq:C.NS.14\endcsname{\href{S6_SERIES_GRIND_II.pdf\#nameddest=tag-00TF}{II.{C.NS.14} (p.~650)} }
\expandafter\def\csname SRef@cdp:eq:C.NS.15\endcsname{\href{S6_SERIES_GRIND_II.pdf\#nameddest=tag-00TG}{II.{C.NS.15} (p.~650)} }
\expandafter\def\csname SRef@cdp:eq:C.NS.16\endcsname{\href{S6_SERIES_GRIND_II.pdf\#nameddest=tag-00TH}{II.{C.NS.16} (p.~650)} }
\expandafter\def\csname SRef@cdp:eq:C.NS.17\endcsname{\href{S6_SERIES_GRIND_II.pdf\#nameddest=tag-00TI}{II.{C.NS.17} (p.~650)} }
\expandafter\def\csname SRef@cdp:eq:C.NS.18\endcsname{\href{S6_SERIES_GRIND_II.pdf\#nameddest=tag-00TJ}{II.{C.NS.18} (p.~650)} }
\expandafter\def\csname SRef@cdp:eq:C.RF.4\endcsname{\href{S6_SERIES_GRIND_II.pdf\#nameddest=tag-00YX}{II.{C.RF.4} (p.~437)} }
\expandafter\def\csname SRef@cdp:eq:C.NS.19\endcsname{\href{S6_SERIES_GRIND_II.pdf\#nameddest=tag-00TK}{II.{C.NS.19} (p.~651)} }
\expandafter\def\csname SRef@cdp:eq:C.NS.20\endcsname{\href{S6_SERIES_GRIND_II.pdf\#nameddest=tag-00TM}{II.{C.NS.20} (p.~651)} }
\expandafter\def\csname SRef@cdp:eq:C.SN.1\endcsname{\href{S6_SERIES_GRIND_II.pdf\#nameddest=tag-0122}{II.{C.SN.1} (p.~652)} }
\expandafter\def\csname SRef@cdp:eq:C.SN.2\endcsname{\href{S6_SERIES_GRIND_II.pdf\#nameddest=tag-0124}{II.{C.SN.2} (p.~652)} }
\expandafter\def\csname SRef@cdp:eq:C.SN.3\endcsname{\href{S6_SERIES_GRIND_II.pdf\#nameddest=tag-0125}{II.{C.SN.3} (p.~652)} }
\expandafter\def\csname SRef@cdp:eq:C.SN.4\endcsname{\href{S6_SERIES_GRIND_II.pdf\#nameddest=tag-0126}{II.{C.SN.4} (p.~652)} }
\expandafter\def\csname SRef@cdp:eq:C.SN.5\endcsname{\href{S6_SERIES_GRIND_II.pdf\#nameddest=tag-0127}{II.{C.SN.5} (p.~653)} }
\expandafter\def\csname SRef@cdp:eq:C.SN.6\endcsname{\href{S6_SERIES_GRIND_II.pdf\#nameddest=tag-0128}{II.{C.SN.6} (p.~653)} }
\expandafter\def\csname SRef@cdp:eq:C.SN.7\endcsname{\href{S6_SERIES_GRIND_II.pdf\#nameddest=tag-0129}{II.{C.SN.7} (p.~653)} }
\expandafter\def\csname SRef@cdp:eq:C.SN.8\endcsname{\href{S6_SERIES_GRIND_II.pdf\#nameddest=tag-012A}{II.{C.SN.8} (p.~653)} }
\expandafter\def\csname SRef@cdp:eq:C.RF.9\endcsname{\href{S6_SERIES_GRIND_II.pdf\#nameddest=tag-00Z2}{II.{C.RF.9} (p.~437)} }
\expandafter\def\csname SRef@cdp:eq:C.SN.9\endcsname{\href{S6_SERIES_GRIND_II.pdf\#nameddest=tag-012B}{II.{C.SN.9} (p.~653)} }
\expandafter\def\csname SRef@cdp:eq:C.SN.10\endcsname{\href{S6_SERIES_GRIND_II.pdf\#nameddest=tag-0123}{II.{C.SN.10} (p.~654)} }
\expandafter\def\csname SRef@cdp:eq:C.MX.1\endcsname{\href{S6_SERIES_GRIND_II.pdf\#nameddest=tag-00RR}{II.{C.MX.1} (p.~654)} }
\expandafter\def\csname SRef@cdp:eq:C.MX.2\endcsname{\href{S6_SERIES_GRIND_II.pdf\#nameddest=tag-00RT}{II.{C.MX.2} (p.~654)} }
\expandafter\def\csname SRef@cdp:eq:C.MX.3\endcsname{\href{S6_SERIES_GRIND_II.pdf\#nameddest=tag-00RU}{II.{C.MX.3} (p.~655)} }
\expandafter\def\csname SRef@cdp:eq:C.MX.4\endcsname{\href{S6_SERIES_GRIND_II.pdf\#nameddest=tag-00RV}{II.{C.MX.4} (p.~655)} }
\expandafter\def\csname SRef@cdp:eq:C.MX.5\endcsname{\href{S6_SERIES_GRIND_II.pdf\#nameddest=tag-00RW}{II.{C.MX.5} (p.~656)} }
\expandafter\def\csname SRef@cdp:eq:C.MX.6\endcsname{\href{S6_SERIES_GRIND_II.pdf\#nameddest=tag-00RX}{II.{C.MX.6} (p.~656)} }
\expandafter\def\csname SRef@cdp:eq:C.MX.7\endcsname{\href{S6_SERIES_GRIND_II.pdf\#nameddest=tag-00RY}{II.{C.MX.7} (p.~656)} }
\expandafter\def\csname SRef@cdp:eq:C.MX.8\endcsname{\href{S6_SERIES_GRIND_II.pdf\#nameddest=tag-00RZ}{II.{C.MX.8} (p.~656)} }
\expandafter\def\csname SRef@cdp:eq:C.RF.6\endcsname{\href{S6_SERIES_GRIND_II.pdf\#nameddest=tag-00YZ}{II.{C.RF.6} (p.~438)} }
\expandafter\def\csname SRef@cdp:eq:C.MX.9\endcsname{\href{S6_SERIES_GRIND_II.pdf\#nameddest=tag-00S0}{II.{C.MX.9} (p.~656)} }
\expandafter\def\csname SRef@cdp:eq:C.MX.10\endcsname{\href{S6_SERIES_GRIND_II.pdf\#nameddest=tag-00RS}{II.{C.MX.10} (p.~656)} }
\expandafter\def\csname SRef@cdp:eq:C.WH.0\endcsname{\href{S6_SERIES_GRIND_II.pdf\#nameddest=tag-016X}{II.{C.WH.0} (p.~658)} }
\expandafter\def\csname SRef@cdp:eq:C.WH.1\endcsname{\href{S6_SERIES_GRIND_II.pdf\#nameddest=tag-016Y}{II.{C.WH.1} (p.~658)} }
\expandafter\def\csname SRef@cdp:eq:C.WH.2\endcsname{\href{S6_SERIES_GRIND_II.pdf\#nameddest=tag-0177}{II.{C.WH.2} (p.~658)} }
\expandafter\def\csname SRef@cdp:eq:C.WH.3\endcsname{\href{S6_SERIES_GRIND_II.pdf\#nameddest=tag-0178}{II.{C.WH.3} (p.~659)} }
\expandafter\def\csname SRef@cdp:eq:C.WH.4\endcsname{\href{S6_SERIES_GRIND_II.pdf\#nameddest=tag-0179}{II.{C.WH.4} (p.~659)} }
\expandafter\def\csname SRef@cdp:eq:C.WH.5\endcsname{\href{S6_SERIES_GRIND_II.pdf\#nameddest=tag-017A}{II.{C.WH.5} (p.~660)} }
\expandafter\def\csname SRef@cdp:eq:C.WH.6\endcsname{\href{S6_SERIES_GRIND_II.pdf\#nameddest=tag-017B}{II.{C.WH.6} (p.~660)} }
\expandafter\def\csname SRef@cdp:eq:C.WH.7\endcsname{\href{S6_SERIES_GRIND_II.pdf\#nameddest=tag-017C}{II.{C.WH.7} (p.~661)} }
\expandafter\def\csname SRef@cdp:eq:C.RF.7\endcsname{\href{S6_SERIES_GRIND_II.pdf\#nameddest=tag-00Z0}{II.{C.RF.7} (p.~439)} }
\expandafter\def\csname SRef@cdp:eq:C.WH.8\endcsname{\href{S6_SERIES_GRIND_II.pdf\#nameddest=tag-017D}{II.{C.WH.8} (p.~661)} }
\expandafter\def\csname SRef@cdp:eq:C.WH.9\endcsname{\href{S6_SERIES_GRIND_II.pdf\#nameddest=tag-017E}{II.{C.WH.9} (p.~661)} }
\expandafter\def\csname SRef@cdp:eq:C.WH.10\endcsname{\href{S6_SERIES_GRIND_II.pdf\#nameddest=tag-016Z}{II.{C.WH.10} (p.~661)} }
\expandafter\def\csname SRef@cdp:eq:C.WH.11\endcsname{\href{S6_SERIES_GRIND_II.pdf\#nameddest=tag-0170}{II.{C.WH.11} (p.~662)} }
\expandafter\def\csname SRef@cdp:eq:C.WH.12\endcsname{\href{S6_SERIES_GRIND_II.pdf\#nameddest=tag-0171}{II.{C.WH.12} (p.~662)} }
\expandafter\def\csname SRef@cdp:eq:C.WH.13\endcsname{\href{S6_SERIES_GRIND_II.pdf\#nameddest=tag-0172}{II.{C.WH.13} (p.~663)} }
\expandafter\def\csname SRef@cdp:eq:C.WH.14\endcsname{\href{S6_SERIES_GRIND_II.pdf\#nameddest=tag-0173}{II.{C.WH.14} (p.~663)} }
\expandafter\def\csname SRef@cdp:eq:C.WH.15\endcsname{\href{S6_SERIES_GRIND_II.pdf\#nameddest=tag-0174}{II.{C.WH.15} (p.~663)} }
\expandafter\def\csname SRef@cdp:eq:C.WH.16\endcsname{\href{S6_SERIES_GRIND_II.pdf\#nameddest=tag-0175}{II.{C.WH.16} (p.~663)} }
\expandafter\def\csname SRef@cdp:eq:C.WH.17\endcsname{\href{S6_SERIES_GRIND_II.pdf\#nameddest=tag-0176}{II.{C.WH.17} (p.~663)} }
\expandafter\def\csname SRef@cdp:eq:C.RF.8\endcsname{\href{S6_SERIES_GRIND_II.pdf\#nameddest=tag-00Z1}{II.{C.RF.8} (p.~439)} }
\expandafter\def\csname SRef@cdp:eq:C.AS.1\endcsname{\href{S6_SERIES_GRIND_II.pdf\#nameddest=tag-00AG}{II.{C.AS.1} (p.~664)} }
\expandafter\def\csname SRef@cdp:eq:C.AS.2\endcsname{\href{S6_SERIES_GRIND_II.pdf\#nameddest=tag-00AR}{II.{C.AS.2} (p.~664)} }
\expandafter\def\csname SRef@cdp:eq:C.AS.3\endcsname{\href{S6_SERIES_GRIND_II.pdf\#nameddest=tag-00AS}{II.{C.AS.3} (p.~664)} }
\expandafter\def\csname SRef@cdp:eq:C.AS.4\endcsname{\href{S6_SERIES_GRIND_II.pdf\#nameddest=tag-00AT}{II.{C.AS.4} (p.~664)} }
\expandafter\def\csname SRef@cdp:eq:C.AS.5\endcsname{\href{S6_SERIES_GRIND_II.pdf\#nameddest=tag-00AU}{II.{C.AS.5} (p.~664)} }
\expandafter\def\csname SRef@cdp:eq:C.AS.6\endcsname{\href{S6_SERIES_GRIND_II.pdf\#nameddest=tag-00AV}{II.{C.AS.6} (p.~665)} }
\expandafter\def\csname SRef@cdp:eq:C.AS.7\endcsname{\href{S6_SERIES_GRIND_II.pdf\#nameddest=tag-00AW}{II.{C.AS.7} (p.~665)} }
\expandafter\def\csname SRef@cdp:eq:C.AS.8\endcsname{\href{S6_SERIES_GRIND_II.pdf\#nameddest=tag-00AX}{II.{C.AS.8} (p.~665)} }
\expandafter\def\csname SRef@cdp:eq:C.AS.9\endcsname{\href{S6_SERIES_GRIND_II.pdf\#nameddest=tag-00AY}{II.{C.AS.9} (p.~665)} }
\expandafter\def\csname SRef@cdp:eq:C.AS.10\endcsname{\href{S6_SERIES_GRIND_II.pdf\#nameddest=tag-00AH}{II.{C.AS.10} (p.~666)} }
\expandafter\def\csname SRef@cdp:eq:C.CQ.1\endcsname{\href{S6_SERIES_GRIND_II.pdf\#nameddest=tag-00CJ}{II.{C.CQ.1} (p.~440)} }
\expandafter\def\csname SRef@cdp:eq:C.AS.11\endcsname{\href{S6_SERIES_GRIND_II.pdf\#nameddest=tag-00AI}{II.{C.AS.11} (p.~666)} }
\expandafter\def\csname SRef@cdp:eq:C.AS.12\endcsname{\href{S6_SERIES_GRIND_II.pdf\#nameddest=tag-00AJ}{II.{C.AS.12} (p.~666)} }
\expandafter\def\csname SRef@cdp:eq:C.AS.13\endcsname{\href{S6_SERIES_GRIND_II.pdf\#nameddest=tag-00AK}{II.{C.AS.13} (p.~666)} }
\expandafter\def\csname SRef@cdp:eq:C.AS.14\endcsname{\href{S6_SERIES_GRIND_II.pdf\#nameddest=tag-00AL}{II.{C.AS.14} (p.~667)} }
\expandafter\def\csname SRef@cdp:eq:C.AS.15\endcsname{\href{S6_SERIES_GRIND_II.pdf\#nameddest=tag-00AM}{II.{C.AS.15} (p.~667)} }
\expandafter\def\csname SRef@cdp:eq:C.AS.16\endcsname{\href{S6_SERIES_GRIND_II.pdf\#nameddest=tag-00AN}{II.{C.AS.16} (p.~667)} }
\expandafter\def\csname SRef@cdp:eq:C.AS.17\endcsname{\href{S6_SERIES_GRIND_II.pdf\#nameddest=tag-00AO}{II.{C.AS.17} (p.~668)} }
\expandafter\def\csname SRef@cdp:eq:C.AS.18\endcsname{\href{S6_SERIES_GRIND_II.pdf\#nameddest=tag-00AP}{II.{C.AS.18} (p.~668)} }
\expandafter\def\csname SRef@cdp:eq:C.AS.19\endcsname{\href{S6_SERIES_GRIND_II.pdf\#nameddest=tag-00AQ}{II.{C.AS.19} (p.~668)} }
\expandafter\def\csname SRef@cdp:eq:C.HR.1\endcsname{\href{S6_SERIES_GRIND_II.pdf\#nameddest=tag-00N0}{II.{C.HR.1} (p.~669)} }
\expandafter\def\csname SRef@cdp:eq:C.CQ.2\endcsname{\href{S6_SERIES_GRIND_II.pdf\#nameddest=tag-00CN}{II.{C.CQ.2} (p.~440)} }
\expandafter\def\csname SRef@cdp:eq:C.HR.2\endcsname{\href{S6_SERIES_GRIND_II.pdf\#nameddest=tag-00N1}{II.{C.HR.2} (p.~669)} }
\expandafter\def\csname SRef@cdp:eq:C.HR.3\endcsname{\href{S6_SERIES_GRIND_II.pdf\#nameddest=tag-00N2}{II.{C.HR.3} (p.~669)} }
\expandafter\def\csname SRef@cdp:eq:C.HR.4\endcsname{\href{S6_SERIES_GRIND_II.pdf\#nameddest=tag-00N3}{II.{C.HR.4} (p.~670)} }
\expandafter\def\csname SRef@cdp:eq:C.HR.5\endcsname{\href{S6_SERIES_GRIND_II.pdf\#nameddest=tag-00N4}{II.{C.HR.5} (p.~670)} }
\expandafter\def\csname SRef@cdp:eq:C.HR.6\endcsname{\href{S6_SERIES_GRIND_II.pdf\#nameddest=tag-00N5}{II.{C.HR.6} (p.~670)} }
\expandafter\def\csname SRef@cdp:eq:C.HR.7\endcsname{\href{S6_SERIES_GRIND_II.pdf\#nameddest=tag-00N6}{II.{C.HR.7} (p.~671)} }
\expandafter\def\csname SRef@cdp:eq:C.HR.8\endcsname{\href{S6_SERIES_GRIND_II.pdf\#nameddest=tag-00N7}{II.{C.HR.8} (p.~671)} }
\expandafter\def\csname SRef@cdp:eq:C.TL.1\endcsname{\href{S6_SERIES_GRIND_II.pdf\#nameddest=tag-013Q}{II.{C.TL.1} (p.~673)} }
\expandafter\def\csname SRef@cdp:eq:C.TL.2\endcsname{\href{S6_SERIES_GRIND_II.pdf\#nameddest=tag-013R}{II.{C.TL.2} (p.~673)} }
\expandafter\def\csname SRef@cdp:eq:C.TL.3\endcsname{\href{S6_SERIES_GRIND_II.pdf\#nameddest=tag-013S}{II.{C.TL.3} (p.~673)} }
\expandafter\def\csname SRef@cdp:eq:C.CQ.3\endcsname{\href{S6_SERIES_GRIND_II.pdf\#nameddest=tag-00CO}{II.{C.CQ.3} (p.~440)} }
\expandafter\def\csname SRef@cdp:eq:C.TL.4\endcsname{\href{S6_SERIES_GRIND_II.pdf\#nameddest=tag-013T}{II.{C.TL.4} (p.~673)} }
\expandafter\def\csname SRef@cdp:eq:C.TL.5\endcsname{\href{S6_SERIES_GRIND_II.pdf\#nameddest=tag-013U}{II.{C.TL.5} (p.~673)} }
\expandafter\def\csname SRef@cdp:eq:C.TL.6\endcsname{\href{S6_SERIES_GRIND_II.pdf\#nameddest=tag-013V}{II.{C.TL.6} (p.~673)} }
\expandafter\def\csname SRef@cdp:eq:C.TL.7\endcsname{\href{S6_SERIES_GRIND_II.pdf\#nameddest=tag-013W}{II.{C.TL.7} (p.~674)} }
\expandafter\def\csname SRef@cdp:eq:C.TL.8\endcsname{\href{S6_SERIES_GRIND_II.pdf\#nameddest=tag-013X}{II.{C.TL.8} (p.~674)} }
\expandafter\def\csname SRef@cdp:eq:C.WM.1\endcsname{\href{S6_SERIES_GRIND_II.pdf\#nameddest=tag-017F}{II.{C.WM.1} (p.~675)} }
\expandafter\def\csname SRef@cdp:eq:C.WM.2\endcsname{\href{S6_SERIES_GRIND_II.pdf\#nameddest=tag-017G}{II.{C.WM.2} (p.~676)} }
\expandafter\def\csname SRef@cdp:eq:C.WM.3\endcsname{\href{S6_SERIES_GRIND_II.pdf\#nameddest=tag-017H}{II.{C.WM.3} (p.~676)} }
\expandafter\def\csname SRef@cdp:eq:C.WM.4\endcsname{\href{S6_SERIES_GRIND_II.pdf\#nameddest=tag-017I}{II.{C.WM.4} (p.~676)} }
\expandafter\def\csname SRef@cdp:eq:C.WM.5\endcsname{\href{S6_SERIES_GRIND_II.pdf\#nameddest=tag-017J}{II.{C.WM.5} (p.~676)} }
\expandafter\def\csname SRef@cdp:eq:C.CQ.4\endcsname{\href{S6_SERIES_GRIND_II.pdf\#nameddest=tag-00CP}{II.{C.CQ.4} (p.~440)} }
\expandafter\def\csname SRef@cdp:eq:C.WM.6\endcsname{\href{S6_SERIES_GRIND_II.pdf\#nameddest=tag-017K}{II.{C.WM.6} (p.~677)} }
\expandafter\def\csname SRef@cdp:eq:C.WM.7\endcsname{\href{S6_SERIES_GRIND_II.pdf\#nameddest=tag-017L}{II.{C.WM.7} (p.~677)} }
\expandafter\def\csname SRef@cdp:eq:C.WM.8\endcsname{\href{S6_SERIES_GRIND_II.pdf\#nameddest=tag-017M}{II.{C.WM.8} (p.~677)} }
\expandafter\def\csname SRef@cdp:eq:C.WM.9\endcsname{\href{S6_SERIES_GRIND_II.pdf\#nameddest=tag-017N}{II.{C.WM.9} (p.~677)} }
\expandafter\def\csname SRef@cdp:eq:C.AD.1\endcsname{\href{S6_SERIES_GRIND_II.pdf\#nameddest=tag-009B}{II.{C.AD.1} (p.~679)} }
\expandafter\def\csname SRef@cdp:eq:C.AD.2\endcsname{\href{S6_SERIES_GRIND_II.pdf\#nameddest=tag-009E}{II.{C.AD.2} (p.~680)} }
\expandafter\def\csname SRef@cdp:eq:C.AD.3\endcsname{\href{S6_SERIES_GRIND_II.pdf\#nameddest=tag-009F}{II.{C.AD.3} (p.~680)} }
\expandafter\def\csname SRef@cdp:eq:C.AD.4\endcsname{\href{S6_SERIES_GRIND_II.pdf\#nameddest=tag-009G}{II.{C.AD.4} (p.~680)} }
\expandafter\def\csname SRef@cdp:eq:C.AD.5\endcsname{\href{S6_SERIES_GRIND_II.pdf\#nameddest=tag-009H}{II.{C.AD.5} (p.~680)} }
\expandafter\def\csname SRef@cdp:eq:C.AD.6\endcsname{\href{S6_SERIES_GRIND_II.pdf\#nameddest=tag-009I}{II.{C.AD.6} (p.~680)} }
\expandafter\def\csname SRef@cdp:eq:C.CQ.5\endcsname{\href{S6_SERIES_GRIND_II.pdf\#nameddest=tag-00CQ}{II.{C.CQ.5} (p.~441)} }
\expandafter\def\csname SRef@cdp:eq:C.AD.7\endcsname{\href{S6_SERIES_GRIND_II.pdf\#nameddest=tag-009J}{II.{C.AD.7} (p.~681)} }
\expandafter\def\csname SRef@cdp:eq:C.AD.8\endcsname{\href{S6_SERIES_GRIND_II.pdf\#nameddest=tag-009K}{II.{C.AD.8} (p.~681)} }
\expandafter\def\csname SRef@cdp:eq:C.AD.9\endcsname{\href{S6_SERIES_GRIND_II.pdf\#nameddest=tag-009L}{II.{C.AD.9} (p.~681)} }
\expandafter\def\csname SRef@cdp:eq:C.AD.10\endcsname{\href{S6_SERIES_GRIND_II.pdf\#nameddest=tag-009C}{II.{C.AD.10} (p.~681)} }
\expandafter\def\csname SRef@cdp:eq:C.AD.11\endcsname{\href{S6_SERIES_GRIND_II.pdf\#nameddest=tag-009D}{II.{C.AD.11} (p.~682)} }
\expandafter\def\csname SRef@cdp:eq:C.HT.1\endcsname{\href{S6_SERIES_GRIND_II.pdf\#nameddest=tag-00N8}{II.{C.HT.1} (p.~682)} }
\expandafter\def\csname SRef@cdp:eq:C.HT.2\endcsname{\href{S6_SERIES_GRIND_II.pdf\#nameddest=tag-00NC}{II.{C.HT.2} (p.~683)} }
\expandafter\def\csname SRef@cdp:eq:C.HT.3\endcsname{\href{S6_SERIES_GRIND_II.pdf\#nameddest=tag-00ND}{II.{C.HT.3} (p.~683)} }
\expandafter\def\csname SRef@cdp:eq:C.HT.4\endcsname{\href{S6_SERIES_GRIND_II.pdf\#nameddest=tag-00NE}{II.{C.HT.4} (p.~684)} }
\expandafter\def\csname SRef@cdp:eq:C.HT.5\endcsname{\href{S6_SERIES_GRIND_II.pdf\#nameddest=tag-00NF}{II.{C.HT.5} (p.~684)} }
\expandafter\def\csname SRef@cdp:eq:C.CQ.6\endcsname{\href{S6_SERIES_GRIND_II.pdf\#nameddest=tag-00CR}{II.{C.CQ.6} (p.~441)} }
\expandafter\def\csname SRef@cdp:eq:C.HT.6\endcsname{\href{S6_SERIES_GRIND_II.pdf\#nameddest=tag-00NG}{II.{C.HT.6} (p.~684)} }
\expandafter\def\csname SRef@cdp:eq:C.HT.7\endcsname{\href{S6_SERIES_GRIND_II.pdf\#nameddest=tag-00NH}{II.{C.HT.7} (p.~684)} }
\expandafter\def\csname SRef@cdp:eq:C.HT.8\endcsname{\href{S6_SERIES_GRIND_II.pdf\#nameddest=tag-00NI}{II.{C.HT.8} (p.~684)} }
\expandafter\def\csname SRef@cdp:eq:C.HT.9\endcsname{\href{S6_SERIES_GRIND_II.pdf\#nameddest=tag-00NJ}{II.{C.HT.9} (p.~685)} }
\expandafter\def\csname SRef@cdp:eq:C.HT.10\endcsname{\href{S6_SERIES_GRIND_II.pdf\#nameddest=tag-00N9}{II.{C.HT.10} (p.~685)} }
\expandafter\def\csname SRef@cdp:eq:C.HT.11\endcsname{\href{S6_SERIES_GRIND_II.pdf\#nameddest=tag-00NA}{II.{C.HT.11} (p.~686)} }
\expandafter\def\csname SRef@cdp:eq:C.HT.12\endcsname{\href{S6_SERIES_GRIND_II.pdf\#nameddest=tag-00NB}{II.{C.HT.12} (p.~686)} }
\expandafter\def\csname SRef@cdp:eq:C.EW.1\endcsname{\href{S6_SERIES_GRIND_II.pdf\#nameddest=tag-00JH}{II.{C.EW.1} (p.~687)} }
\expandafter\def\csname SRef@cdp:eq:C.EW.2\endcsname{\href{S6_SERIES_GRIND_II.pdf\#nameddest=tag-00JS}{II.{C.EW.2} (p.~688)} }
\expandafter\def\csname SRef@cdp:eq:C.EW.3\endcsname{\href{S6_SERIES_GRIND_II.pdf\#nameddest=tag-00JV}{II.{C.EW.3} (p.~688)} }
\expandafter\def\csname SRef@cdp:eq:C.CQ.7\endcsname{\href{S6_SERIES_GRIND_II.pdf\#nameddest=tag-00CS}{II.{C.CQ.7} (p.~441)} }
\expandafter\def\csname SRef@cdp:eq:C.EW.4\endcsname{\href{S6_SERIES_GRIND_II.pdf\#nameddest=tag-00JW}{II.{C.EW.4} (p.~688)} }
\expandafter\def\csname SRef@cdp:eq:C.EW.5\endcsname{\href{S6_SERIES_GRIND_II.pdf\#nameddest=tag-00JX}{II.{C.EW.5} (p.~688)} }
\expandafter\def\csname SRef@cdp:eq:C.EW.6\endcsname{\href{S6_SERIES_GRIND_II.pdf\#nameddest=tag-00JY}{II.{C.EW.6} (p.~688)} }
\expandafter\def\csname SRef@cdp:eq:C.EW.7\endcsname{\href{S6_SERIES_GRIND_II.pdf\#nameddest=tag-00JZ}{II.{C.EW.7} (p.~689)} }
\expandafter\def\csname SRef@cdp:eq:C.EW.8\endcsname{\href{S6_SERIES_GRIND_II.pdf\#nameddest=tag-00K0}{II.{C.EW.8} (p.~689)} }
\expandafter\def\csname SRef@cdp:eq:C.EW.9\endcsname{\href{S6_SERIES_GRIND_II.pdf\#nameddest=tag-00K1}{II.{C.EW.9} (p.~689)} }
\expandafter\def\csname SRef@cdp:eq:C.EW.10\endcsname{\href{S6_SERIES_GRIND_II.pdf\#nameddest=tag-00JI}{II.{C.EW.10} (p.~689)} }
\expandafter\def\csname SRef@cdp:eq:C.EW.11\endcsname{\href{S6_SERIES_GRIND_II.pdf\#nameddest=tag-00JJ}{II.{C.EW.11} (p.~689)} }
\expandafter\def\csname SRef@cdp:eq:C.EW.12\endcsname{\href{S6_SERIES_GRIND_II.pdf\#nameddest=tag-00JK}{II.{C.EW.12} (p.~689)} }
\expandafter\def\csname SRef@cdp:eq:C.EW.13\endcsname{\href{S6_SERIES_GRIND_II.pdf\#nameddest=tag-00JL}{II.{C.EW.13} (p.~689)} }
\expandafter\def\csname SRef@cdp:eq:C.CQ.8\endcsname{\href{S6_SERIES_GRIND_II.pdf\#nameddest=tag-00CT}{II.{C.CQ.8} (p.~441)} }
\expandafter\def\csname SRef@cdp:eq:C.EW.14\endcsname{\href{S6_SERIES_GRIND_II.pdf\#nameddest=tag-00JM}{II.{C.EW.14} (p.~689)} }
\expandafter\def\csname SRef@cdp:eq:C.EW.15\endcsname{\href{S6_SERIES_GRIND_II.pdf\#nameddest=tag-00JN}{II.{C.EW.15} (p.~690)} }
\expandafter\def\csname SRef@cdp:eq:C.EW.16\endcsname{\href{S6_SERIES_GRIND_II.pdf\#nameddest=tag-00JO}{II.{C.EW.16} (p.~690)} }
\expandafter\def\csname SRef@cdp:eq:C.EW.17\endcsname{\href{S6_SERIES_GRIND_II.pdf\#nameddest=tag-00JP}{II.{C.EW.17} (p.~690)} }
\expandafter\def\csname SRef@cdp:eq:C.EW.18\endcsname{\href{S6_SERIES_GRIND_II.pdf\#nameddest=tag-00JQ}{II.{C.EW.18} (p.~691)} }
\expandafter\def\csname SRef@cdp:eq:C.EW.19\endcsname{\href{S6_SERIES_GRIND_II.pdf\#nameddest=tag-00JR}{II.{C.EW.19} (p.~691)} }
\expandafter\def\csname SRef@cdp:eq:C.EW.20\endcsname{\href{S6_SERIES_GRIND_II.pdf\#nameddest=tag-00JT}{II.{C.EW.20} (p.~691)} }
\expandafter\def\csname SRef@cdp:eq:C.EW.21\endcsname{\href{S6_SERIES_GRIND_II.pdf\#nameddest=tag-00JU}{II.{C.EW.21} (p.~691)} }
\expandafter\def\csname SRef@cdp:eq:C.PR.1\endcsname{\href{S6_SERIES_GRIND_II.pdf\#nameddest=tag-00VH}{II.{C.PR.1} (p.~693)} }
\expandafter\def\csname SRef@cdp:eq:C.PR.2\endcsname{\href{S6_SERIES_GRIND_II.pdf\#nameddest=tag-00VR}{II.{C.PR.2} (p.~693)} }
\expandafter\def\csname SRef@cdp:eq:C.CQ.9\endcsname{\href{S6_SERIES_GRIND_II.pdf\#nameddest=tag-00CU}{II.{C.CQ.9} (p.~441)} }
\expandafter\def\csname SRef@cdp:eq:C.PR.3\endcsname{\href{S6_SERIES_GRIND_II.pdf\#nameddest=tag-00VS}{II.{C.PR.3} (p.~693)} }
\expandafter\def\csname SRef@cdp:eq:C.PR.4\endcsname{\href{S6_SERIES_GRIND_II.pdf\#nameddest=tag-00VT}{II.{C.PR.4} (p.~693)} }
\expandafter\def\csname SRef@cdp:eq:C.PR.5\endcsname{\href{S6_SERIES_GRIND_II.pdf\#nameddest=tag-00VU}{II.{C.PR.5} (p.~693)} }
\expandafter\def\csname SRef@cdp:eq:C.PR.6\endcsname{\href{S6_SERIES_GRIND_II.pdf\#nameddest=tag-00VV}{II.{C.PR.6} (p.~694)} }
\expandafter\def\csname SRef@cdp:eq:C.PR.7\endcsname{\href{S6_SERIES_GRIND_II.pdf\#nameddest=tag-00VW}{II.{C.PR.7} (p.~694)} }
\expandafter\def\csname SRef@cdp:eq:C.PR.8\endcsname{\href{S6_SERIES_GRIND_II.pdf\#nameddest=tag-00VX}{II.{C.PR.8} (p.~695)} }
\expandafter\def\csname SRef@cdp:eq:C.PR.9\endcsname{\href{S6_SERIES_GRIND_II.pdf\#nameddest=tag-00VY}{II.{C.PR.9} (p.~695)} }
\expandafter\def\csname SRef@cdp:eq:C.PR.10\endcsname{\href{S6_SERIES_GRIND_II.pdf\#nameddest=tag-00VI}{II.{C.PR.10} (p.~695)} }
\expandafter\def\csname SRef@cdp:eq:C.PR.11\endcsname{\href{S6_SERIES_GRIND_II.pdf\#nameddest=tag-00VJ}{II.{C.PR.11} (p.~695)} }
\expandafter\def\csname SRef@cdp:eq:C.PR.12\endcsname{\href{S6_SERIES_GRIND_II.pdf\#nameddest=tag-00VK}{II.{C.PR.12} (p.~696)} }
\expandafter\def\csname SRef@cdp:eq:C.CQ.10\endcsname{\href{S6_SERIES_GRIND_II.pdf\#nameddest=tag-00CK}{II.{C.CQ.10} (p.~441)} }
\expandafter\def\csname SRef@cdp:eq:C.PR.13\endcsname{\href{S6_SERIES_GRIND_II.pdf\#nameddest=tag-00VL}{II.{C.PR.13} (p.~696)} }
\expandafter\def\csname SRef@cdp:eq:C.PR.14\endcsname{\href{S6_SERIES_GRIND_II.pdf\#nameddest=tag-00VM}{II.{C.PR.14} (p.~696)} }
\expandafter\def\csname SRef@cdp:eq:C.PR.15\endcsname{\href{S6_SERIES_GRIND_II.pdf\#nameddest=tag-00VN}{II.{C.PR.15} (p.~696)} }
\expandafter\def\csname SRef@cdp:eq:C.PR.16\endcsname{\href{S6_SERIES_GRIND_II.pdf\#nameddest=tag-00VO}{II.{C.PR.16} (p.~697)} }
\expandafter\def\csname SRef@cdp:eq:C.PR.17\endcsname{\href{S6_SERIES_GRIND_II.pdf\#nameddest=tag-00VP}{II.{C.PR.17} (p.~697)} }
\expandafter\def\csname SRef@cdp:eq:C.PR.18\endcsname{\href{S6_SERIES_GRIND_II.pdf\#nameddest=tag-00VQ}{II.{C.PR.18} (p.~697)} }
\expandafter\def\csname SRef@cdp:eq:C.PT.1\endcsname{\href{S6_SERIES_GRIND_II.pdf\#nameddest=tag-00VZ}{II.{C.PT.1} (p.~699)} }
\expandafter\def\csname SRef@cdp:eq:C.PT.2\endcsname{\href{S6_SERIES_GRIND_II.pdf\#nameddest=tag-00W5}{II.{C.PT.2} (p.~699)} }
\expandafter\def\csname SRef@cdp:eq:C.PT.3\endcsname{\href{S6_SERIES_GRIND_II.pdf\#nameddest=tag-00W6}{II.{C.PT.3} (p.~699)} }
\expandafter\def\csname SRef@cdp:eq:C.PT.4\endcsname{\href{S6_SERIES_GRIND_II.pdf\#nameddest=tag-00W7}{II.{C.PT.4} (p.~699)} }
\expandafter\def\csname SRef@cdp:eq:C.CQ.11\endcsname{\href{S6_SERIES_GRIND_II.pdf\#nameddest=tag-00CL}{II.{C.CQ.11} (p.~442)} }
\expandafter\def\csname SRef@cdp:eq:C.PT.5\endcsname{\href{S6_SERIES_GRIND_II.pdf\#nameddest=tag-00W8}{II.{C.PT.5} (p.~700)} }
\expandafter\def\csname SRef@cdp:eq:C.PT.6\endcsname{\href{S6_SERIES_GRIND_II.pdf\#nameddest=tag-00W9}{II.{C.PT.6} (p.~700)} }
\expandafter\def\csname SRef@cdp:eq:C.PT.7\endcsname{\href{S6_SERIES_GRIND_II.pdf\#nameddest=tag-00WA}{II.{C.PT.7} (p.~700)} }
\expandafter\def\csname SRef@cdp:eq:C.PT.8\endcsname{\href{S6_SERIES_GRIND_II.pdf\#nameddest=tag-00WB}{II.{C.PT.8} (p.~700)} }
\expandafter\def\csname SRef@cdp:eq:C.PT.9\endcsname{\href{S6_SERIES_GRIND_II.pdf\#nameddest=tag-00WC}{II.{C.PT.9} (p.~701)} }
\expandafter\def\csname SRef@cdp:eq:C.PT.10\endcsname{\href{S6_SERIES_GRIND_II.pdf\#nameddest=tag-00W0}{II.{C.PT.10} (p.~702)} }
\expandafter\def\csname SRef@cdp:eq:C.PT.11\endcsname{\href{S6_SERIES_GRIND_II.pdf\#nameddest=tag-00W1}{II.{C.PT.11} (p.~702)} }
\expandafter\def\csname SRef@cdp:eq:C.PT.12\endcsname{\href{S6_SERIES_GRIND_II.pdf\#nameddest=tag-00W2}{II.{C.PT.12} (p.~702)} }
\expandafter\def\csname SRef@cdp:eq:C.PT.13\endcsname{\href{S6_SERIES_GRIND_II.pdf\#nameddest=tag-00W3}{II.{C.PT.13} (p.~702)} }
\expandafter\def\csname SRef@cdp:eq:C.PT.14\endcsname{\href{S6_SERIES_GRIND_II.pdf\#nameddest=tag-00W4}{II.{C.PT.14} (p.~703)} }
\expandafter\def\csname SRef@cdp:eq:C.CQ.12\endcsname{\href{S6_SERIES_GRIND_II.pdf\#nameddest=tag-00CM}{II.{C.CQ.12} (p.~442)} }
\expandafter\def\csname SRef@cdp:eq:C.QT.1\endcsname{\href{S6_SERIES_GRIND_II.pdf\#nameddest=tag-00WD}{II.{C.QT.1} (p.~704)} }
\expandafter\def\csname SRef@cdp:eq:C.QT.2\endcsname{\href{S6_SERIES_GRIND_II.pdf\#nameddest=tag-00WI}{II.{C.QT.2} (p.~704)} }
\expandafter\def\csname SRef@cdp:eq:C.QT.3\endcsname{\href{S6_SERIES_GRIND_II.pdf\#nameddest=tag-00WJ}{II.{C.QT.3} (p.~705)} }
\expandafter\def\csname SRef@cdp:eq:C.QT.4\endcsname{\href{S6_SERIES_GRIND_II.pdf\#nameddest=tag-00WK}{II.{C.QT.4} (p.~705)} }
\expandafter\def\csname SRef@cdp:eq:C.QT.5\endcsname{\href{S6_SERIES_GRIND_II.pdf\#nameddest=tag-00WL}{II.{C.QT.5} (p.~705)} }
\expandafter\def\csname SRef@cdp:eq:C.QT.6\endcsname{\href{S6_SERIES_GRIND_II.pdf\#nameddest=tag-00WM}{II.{C.QT.6} (p.~705)} }
\expandafter\def\csname SRef@cdp:eq:C.QT.7\endcsname{\href{S6_SERIES_GRIND_II.pdf\#nameddest=tag-00WN}{II.{C.QT.7} (p.~706)} }
\expandafter\def\csname SRef@cdp:eq:C.QT.8\endcsname{\href{S6_SERIES_GRIND_II.pdf\#nameddest=tag-00WO}{II.{C.QT.8} (p.~706)} }
\expandafter\def\csname SRef@cdp:eq:C.QT.9\endcsname{\href{S6_SERIES_GRIND_II.pdf\#nameddest=tag-00WP}{II.{C.QT.9} (p.~706)} }
\expandafter\def\csname SRef@cdp:eq:C.QT.10\endcsname{\href{S6_SERIES_GRIND_II.pdf\#nameddest=tag-00WE}{II.{C.QT.10} (p.~706)} }
\expandafter\def\csname SRef@cdp:st:C.2.UN.1\endcsname{\href{S6_SERIES_GRIND_II.pdf\#nameddest=tag-01CM}{II.19.1 (p.~413)} }
\expandafter\def\csname SRef@cdp:st:C.2.UN.20\endcsname{\href{S6_SERIES_GRIND_II.pdf\#nameddest=tag-01CY}{II.19.10 (p.~433)} }
\expandafter\def\csname SRef@cdp:st:C.2.UN.22\endcsname{\href{S6_SERIES_GRIND_II.pdf\#nameddest=tag-01D0}{II.19.11 (p.~435)} }
\expandafter\def\csname SRef@cdp:st:C.2.UN.23\endcsname{\href{S6_SERIES_GRIND_II.pdf\#nameddest=tag-01D1}{II.19.12 (p.~439)} }
\expandafter\def\csname SRef@cdp:st:C.2.UN.26\endcsname{\href{S6_SERIES_GRIND_II.pdf\#nameddest=tag-01D4}{II.19.13 (p.~443)} }
\expandafter\def\csname SRef@cdp:st:C.2.UN.28\endcsname{\href{S6_SERIES_GRIND_II.pdf\#nameddest=tag-01D6}{II.19.14 (p.~444)} }
\expandafter\def\csname SRef@cdp:st:C.2.UN.29\endcsname{\href{S6_SERIES_GRIND_II.pdf\#nameddest=tag-01D7}{II.19.15 (p.~446)} }
\expandafter\def\csname SRef@cdp:st:C.2.UN.31\endcsname{\href{S6_SERIES_GRIND_II.pdf\#nameddest=tag-01DA}{II.19.16 (p.~450)} }
\expandafter\def\csname SRef@cdp:st:C.2.UN.33\endcsname{\href{S6_SERIES_GRIND_II.pdf\#nameddest=tag-01DC}{II.19.17 (p.~453)} }
\expandafter\def\csname SRef@cdp:st:C.2.UN.36\endcsname{\href{S6_SERIES_GRIND_II.pdf\#nameddest=tag-01DF}{II.19.18 (p.~457)} }
\expandafter\def\csname SRef@cdp:st:C.2.UN.38\endcsname{\href{S6_SERIES_GRIND_II.pdf\#nameddest=tag-01DH}{II.19.19 (p.~460)} }
\expandafter\def\csname SRef@cdp:st:C.2.UN.2\endcsname{\href{S6_SERIES_GRIND_II.pdf\#nameddest=tag-01CX}{II.19.2 (p.~414)} }
\expandafter\def\csname SRef@cdp:st:C.2.UN.39\endcsname{\href{S6_SERIES_GRIND_II.pdf\#nameddest=tag-01DI}{II.19.20 (p.~463)} }
\expandafter\def\csname SRef@cdp:st:C.2.UN.44\endcsname{\href{S6_SERIES_GRIND_II.pdf\#nameddest=tag-01DO}{II.19.21 (p.~466)} }
\expandafter\def\csname SRef@cdp:st:C.2.UN.47\endcsname{\href{S6_SERIES_GRIND_II.pdf\#nameddest=tag-01DR}{II.19.22 (p.~471)} }
\expandafter\def\csname SRef@cdp:st:C.2.UN.3\endcsname{\href{S6_SERIES_GRIND_II.pdf\#nameddest=tag-01D8}{II.19.3 (p.~415)} }
\expandafter\def\csname SRef@cdp:st:C.2.UN.4\endcsname{\href{S6_SERIES_GRIND_II.pdf\#nameddest=tag-01DJ}{II.19.4 (p.~418)} }
\expandafter\def\csname SRef@cdp:st:C.2.UN.14\endcsname{\href{S6_SERIES_GRIND_II.pdf\#nameddest=tag-01CR}{II.19.5 (p.~421)} }
\expandafter\def\csname SRef@cdp:st:C.2.UN.15\endcsname{\href{S6_SERIES_GRIND_II.pdf\#nameddest=tag-01CS}{II.19.6 (p.~423)} }
\expandafter\def\csname SRef@cdp:st:C.2.UN.16\endcsname{\href{S6_SERIES_GRIND_II.pdf\#nameddest=tag-01CT}{II.19.7 (p.~425)} }
\expandafter\def\csname SRef@cdp:st:C.2.UN.17\endcsname{\href{S6_SERIES_GRIND_II.pdf\#nameddest=tag-01CU}{II.19.8 (p.~428)} }
\expandafter\def\csname SRef@cdp:st:C.2.UN.19\endcsname{\href{S6_SERIES_GRIND_II.pdf\#nameddest=tag-01CW}{II.19.9 (p.~431)} }
\expandafter\def\csname SRef@cdp:st:C.2.UN.5\endcsname{\href{S6_SERIES_GRIND_II.pdf\#nameddest=tag-01DU}{II.20.1 (p.~474)} }
\expandafter\def\csname SRef@cdp:st:C.2.IRR.1\endcsname{\href{S6_SERIES_GRIND_II.pdf\#nameddest=tag-01C3}{II.20.10 (p.~487)} }
\expandafter\def\csname SRef@cdp:st:C.2.VII.1\endcsname{\href{S6_SERIES_GRIND_II.pdf\#nameddest=tag-01F0}{II.20.11 (p.~489)} }
\expandafter\def\csname SRef@cdp:st:C.2.VII.2\endcsname{\href{S6_SERIES_GRIND_II.pdf\#nameddest=tag-01F1}{II.20.12 (p.~491)} }
\expandafter\def\csname SRef@cdp:st:C.2.VII.3\endcsname{\href{S6_SERIES_GRIND_II.pdf\#nameddest=tag-01F2}{II.20.13 (p.~492)} }
\expandafter\def\csname SRef@cdp:st:C.2.VII.4\endcsname{\href{S6_SERIES_GRIND_II.pdf\#nameddest=tag-01F3}{II.20.14 (p.~493)} }
\expandafter\def\csname SRef@cdp:st:C.2.RAT.1\endcsname{\href{S6_SERIES_GRIND_II.pdf\#nameddest=tag-01CD}{II.20.15 (p.~494)} }
\expandafter\def\csname SRef@cdp:st:C.2.RAT.2\endcsname{\href{S6_SERIES_GRIND_II.pdf\#nameddest=tag-01CE}{II.20.16 (p.~498)} }
\expandafter\def\csname SRef@cdp:st:C.2.RAT.3\endcsname{\href{S6_SERIES_GRIND_II.pdf\#nameddest=tag-01CF}{II.20.17 (p.~499)} }
\expandafter\def\csname SRef@cdp:st:C.2.RAT.4\endcsname{\href{S6_SERIES_GRIND_II.pdf\#nameddest=tag-01CG}{II.20.18 (p.~501)} }
\expandafter\def\csname SRef@cdp:st:C.2.RAT.5\endcsname{\href{S6_SERIES_GRIND_II.pdf\#nameddest=tag-01CH}{II.20.19 (p.~503)} }
\expandafter\def\csname SRef@cdp:st:C.2.CR.1\endcsname{\href{S6_SERIES_GRIND_II.pdf\#nameddest=tag-01BQ}{II.20.2 (p.~476)} }
\expandafter\def\csname SRef@cdp:st:C.2.RAT.7\endcsname{\href{S6_SERIES_GRIND_II.pdf\#nameddest=tag-01CJ}{II.20.20 (p.~504)} }
\expandafter\def\csname SRef@cdp:st:C.2.RAT.8\endcsname{\href{S6_SERIES_GRIND_II.pdf\#nameddest=tag-01CK}{II.20.21 (p.~506)} }
\expandafter\def\csname SRef@cdp:st:C.2.RAT.9\endcsname{\href{S6_SERIES_GRIND_II.pdf\#nameddest=tag-01CL}{II.20.22 (p.~509)} }
\expandafter\def\csname SRef@cdp:st:C.2.D4.1\endcsname{\href{S6_SERIES_GRIND_II.pdf\#nameddest=tag-01BW}{II.20.23 (p.~511)} }
\expandafter\def\csname SRef@cdp:st:C.2.D4.2\endcsname{\href{S6_SERIES_GRIND_II.pdf\#nameddest=tag-01BX}{II.20.24 (p.~516)} }
\expandafter\def\csname SRef@cdp:st:C.2.D4.3\endcsname{\href{S6_SERIES_GRIND_II.pdf\#nameddest=tag-01BY}{II.20.25 (p.~520)} }
\expandafter\def\csname SRef@cdp:st:C.2.UN.7\endcsname{\href{S6_SERIES_GRIND_II.pdf\#nameddest=tag-01EG}{II.20.26 (p.~522)} }
\expandafter\def\csname SRef@cdp:st:C.2.UN.12\endcsname{\href{S6_SERIES_GRIND_II.pdf\#nameddest=tag-01CP}{II.20.27 (p.~528)} }
\expandafter\def\csname SRef@cdp:st:C.2.INT.1\endcsname{\href{S6_SERIES_GRIND_II.pdf\#nameddest=tag-01C0}{II.20.28 (p.~531)} }
\expandafter\def\csname SRef@cdp:st:C.2.INT.2\endcsname{\href{S6_SERIES_GRIND_II.pdf\#nameddest=tag-01C1}{II.20.29 (p.~533)} }
\expandafter\def\csname SRef@cdp:st:C.2.CR.2\endcsname{\href{S6_SERIES_GRIND_II.pdf\#nameddest=tag-01BR}{II.20.3 (p.~478)} }
\expandafter\def\csname SRef@cdp:st:C.2.INT.3\endcsname{\href{S6_SERIES_GRIND_II.pdf\#nameddest=tag-01C2}{II.20.30 (p.~535)} }
\expandafter\def\csname SRef@cdp:st:C.2.MH.1\endcsname{\href{S6_SERIES_GRIND_II.pdf\#nameddest=tag-01C4}{II.20.31 (p.~538)} }
\expandafter\def\csname SRef@cdp:st:C.2.MH.2\endcsname{\href{S6_SERIES_GRIND_II.pdf\#nameddest=tag-01C5}{II.20.32 (p.~540)} }
\expandafter\def\csname SRef@cdp:st:C.2.MH.3\endcsname{\href{S6_SERIES_GRIND_II.pdf\#nameddest=tag-01C6}{II.20.33 (p.~542)} }
\expandafter\def\csname SRef@cdp:st:C.2.MH.4\endcsname{\href{S6_SERIES_GRIND_II.pdf\#nameddest=tag-01C7}{II.20.34 (p.~542)} }
\expandafter\def\csname SRef@cdp:st:C.2.CR.3\endcsname{\href{S6_SERIES_GRIND_II.pdf\#nameddest=tag-01BS}{II.20.4 (p.~480)} }
\expandafter\def\csname SRef@cdp:st:C.2.CR.4\endcsname{\href{S6_SERIES_GRIND_II.pdf\#nameddest=tag-01BT}{II.20.5 (p.~481)} }
\expandafter\def\csname SRef@cdp:st:C.2.CR.5\endcsname{\href{S6_SERIES_GRIND_II.pdf\#nameddest=tag-01BU}{II.20.6 (p.~482)} }
\expandafter\def\csname SRef@cdp:st:C.2.CR.6\endcsname{\href{S6_SERIES_GRIND_II.pdf\#nameddest=tag-01BV}{II.20.7 (p.~483)} }
\expandafter\def\csname SRef@cdp:st:C.2.UN.9\endcsname{\href{S6_SERIES_GRIND_II.pdf\#nameddest=tag-01EZ}{II.20.8 (p.~484)} }
\expandafter\def\csname SRef@cdp:st:C.2.UN.10\endcsname{\href{S6_SERIES_GRIND_II.pdf\#nameddest=tag-01CN}{II.20.9 (p.~486)} }
\expandafter\def\csname SRef@cdp:st:C2a\endcsname{\href{S6_SERIES_GRIND_II.pdf\#nameddest=tag-01FG}{II.21.1 (p.~543)} }
\expandafter\def\csname SRef@cdp:st:C.2.UN.67\endcsname{\href{S6_SERIES_GRIND_II.pdf\#nameddest=tag-01ED}{II.21.10 (p.~573)} }
\expandafter\def\csname SRef@cdp:st:C.2.UN.68\endcsname{\href{S6_SERIES_GRIND_II.pdf\#nameddest=tag-01EE}{II.21.11 (p.~576)} }
\expandafter\def\csname SRef@cdp:st:C2b\endcsname{\href{S6_SERIES_GRIND_II.pdf\#nameddest=tag-01FH}{II.21.2 (p.~543)} }
\expandafter\def\csname SRef@cdp:st:C2c\endcsname{\href{S6_SERIES_GRIND_II.pdf\#nameddest=tag-01FI}{II.21.3 (p.~544)} }
\expandafter\def\csname SRef@cdp:st:C.2.UN.11\endcsname{\href{S6_SERIES_GRIND_II.pdf\#nameddest=tag-01CO}{II.21.4 (p.~546)} }
\expandafter\def\csname SRef@cdp:st:C.2.UN.24\endcsname{\href{S6_SERIES_GRIND_II.pdf\#nameddest=tag-01D2}{II.21.5 (p.~550)} }
\expandafter\def\csname SRef@cdp:st:C.2.UN.32\endcsname{\href{S6_SERIES_GRIND_II.pdf\#nameddest=tag-01DB}{II.21.6 (p.~554)} }
\expandafter\def\csname SRef@cdp:st:C.2.UN.64\endcsname{\href{S6_SERIES_GRIND_II.pdf\#nameddest=tag-01EA}{II.21.7 (p.~557)} }
\expandafter\def\csname SRef@cdp:st:C.2.UN.65\endcsname{\href{S6_SERIES_GRIND_II.pdf\#nameddest=tag-01EB}{II.21.8 (p.~561)} }
\expandafter\def\csname SRef@cdp:st:C.2.UN.66\endcsname{\href{S6_SERIES_GRIND_II.pdf\#nameddest=tag-01EC}{II.21.9 (p.~565)} }
\expandafter\def\csname SRef@cdp:st:C.2.UN.6\endcsname{\href{S6_SERIES_GRIND_II.pdf\#nameddest=tag-01E5}{II.22.1 (p.~579)} }
\expandafter\def\csname SRef@cdp:st:C.2.UN.35\endcsname{\href{S6_SERIES_GRIND_II.pdf\#nameddest=tag-01DE}{II.22.10 (p.~613)} }
\expandafter\def\csname SRef@cdp:st:C.2.UN.37\endcsname{\href{S6_SERIES_GRIND_II.pdf\#nameddest=tag-01DG}{II.22.11 (p.~620)} }
\expandafter\def\csname SRef@cdp:st:C.2.UN.41\endcsname{\href{S6_SERIES_GRIND_II.pdf\#nameddest=tag-01DL}{II.22.12 (p.~622)} }
\expandafter\def\csname SRef@cdp:st:C.2.UN.42\endcsname{\href{S6_SERIES_GRIND_II.pdf\#nameddest=tag-01DM}{II.22.13 (p.~626)} }
\expandafter\def\csname SRef@cdp:st:C.2.UN.43\endcsname{\href{S6_SERIES_GRIND_II.pdf\#nameddest=tag-01DN}{II.22.14 (p.~630)} }
\expandafter\def\csname SRef@cdp:st:C.2.UN.45\endcsname{\href{S6_SERIES_GRIND_II.pdf\#nameddest=tag-01DP}{II.22.15 (p.~634)} }
\expandafter\def\csname SRef@cdp:st:C.2.UN.8\endcsname{\href{S6_SERIES_GRIND_II.pdf\#nameddest=tag-01ER}{II.22.2 (p.~583)} }
\expandafter\def\csname SRef@cdp:st:C.2.UN.13\endcsname{\href{S6_SERIES_GRIND_II.pdf\#nameddest=tag-01CQ}{II.22.3 (p.~586)} }
\expandafter\def\csname SRef@cdp:st:C.2.UN.18\endcsname{\href{S6_SERIES_GRIND_II.pdf\#nameddest=tag-01CV}{II.22.4 (p.~590)} }
\expandafter\def\csname SRef@cdp:st:C.2.UN.21\endcsname{\href{S6_SERIES_GRIND_II.pdf\#nameddest=tag-01CZ}{II.22.5 (p.~594)} }
\expandafter\def\csname SRef@cdp:st:C.2.UN.25\endcsname{\href{S6_SERIES_GRIND_II.pdf\#nameddest=tag-01D3}{II.22.6 (p.~598)} }
\expandafter\def\csname SRef@cdp:st:C.2.UN.27\endcsname{\href{S6_SERIES_GRIND_II.pdf\#nameddest=tag-01D5}{II.22.7 (p.~603)} }
\expandafter\def\csname SRef@cdp:st:C.2.UN.30\endcsname{\href{S6_SERIES_GRIND_II.pdf\#nameddest=tag-01D9}{II.22.8 (p.~607)} }
\expandafter\def\csname SRef@cdp:st:C.2.UN.34\endcsname{\href{S6_SERIES_GRIND_II.pdf\#nameddest=tag-01DD}{II.22.9 (p.~610)} }
\expandafter\def\csname SRef@cdp:st:C.2.UN.40\endcsname{\href{S6_SERIES_GRIND_II.pdf\#nameddest=tag-01DK}{II.23.1 (p.~636)} }
\expandafter\def\csname SRef@cdp:st:C.2.UN.55\endcsname{\href{S6_SERIES_GRIND_II.pdf\#nameddest=tag-01E0}{II.23.10 (p.~672)} }
\expandafter\def\csname SRef@cdp:st:C.2.UN.56\endcsname{\href{S6_SERIES_GRIND_II.pdf\#nameddest=tag-01E1}{II.23.11 (p.~675)} }
\expandafter\def\csname SRef@cdp:st:C.2.UN.57\endcsname{\href{S6_SERIES_GRIND_II.pdf\#nameddest=tag-01E2}{II.23.12 (p.~679)} }
\expandafter\def\csname SRef@cdp:st:C.2.UN.58\endcsname{\href{S6_SERIES_GRIND_II.pdf\#nameddest=tag-01E3}{II.23.13 (p.~682)} }
\expandafter\def\csname SRef@cdp:st:C.2.UN.59\endcsname{\href{S6_SERIES_GRIND_II.pdf\#nameddest=tag-01E4}{II.23.14 (p.~686)} }
\expandafter\def\csname SRef@cdp:st:C.2.UN.60\endcsname{\href{S6_SERIES_GRIND_II.pdf\#nameddest=tag-01E6}{II.23.15 (p.~692)} }
\expandafter\def\csname SRef@cdp:st:C.2.UN.61\endcsname{\href{S6_SERIES_GRIND_II.pdf\#nameddest=tag-01E7}{II.23.16 (p.~698)} }
\expandafter\def\csname SRef@cdp:st:C.2.UN.62\endcsname{\href{S6_SERIES_GRIND_II.pdf\#nameddest=tag-01E8}{II.23.17 (p.~703)} }
\expandafter\def\csname SRef@cdp:st:C.2.UN.63\endcsname{\href{S6_SERIES_GRIND_II.pdf\#nameddest=tag-01E9}{II.23.18 (p.~708)} }
\expandafter\def\csname SRef@cdp:st:C.2.UN.46\endcsname{\href{S6_SERIES_GRIND_II.pdf\#nameddest=tag-01DQ}{II.23.2 (p.~638)} }
\expandafter\def\csname SRef@cdp:st:C.2.UN.48\endcsname{\href{S6_SERIES_GRIND_II.pdf\#nameddest=tag-01DS}{II.23.3 (p.~642)} }
\expandafter\def\csname SRef@cdp:st:C.2.UN.49\endcsname{\href{S6_SERIES_GRIND_II.pdf\#nameddest=tag-01DT}{II.23.4 (p.~646)} }
\expandafter\def\csname SRef@cdp:st:C.2.UN.50\endcsname{\href{S6_SERIES_GRIND_II.pdf\#nameddest=tag-01DV}{II.23.5 (p.~652)} }
\expandafter\def\csname SRef@cdp:st:C.2.UN.51\endcsname{\href{S6_SERIES_GRIND_II.pdf\#nameddest=tag-01DW}{II.23.6 (p.~654)} }
\expandafter\def\csname SRef@cdp:st:C.2.UN.52\endcsname{\href{S6_SERIES_GRIND_II.pdf\#nameddest=tag-01DX}{II.23.7 (p.~657)} }
\expandafter\def\csname SRef@cdp:st:C.2.UN.53\endcsname{\href{S6_SERIES_GRIND_II.pdf\#nameddest=tag-01DY}{II.23.8 (p.~664)} }
\expandafter\def\csname SRef@cdp:st:C.2.UN.54\endcsname{\href{S6_SERIES_GRIND_II.pdf\#nameddest=tag-01DZ}{II.23.9 (p.~669)} }
\expandafter\def\csname SRef@cdp:st:C.2.UN.69\endcsname{\href{S6_SERIES_GRIND_II.pdf\#nameddest=tag-01EF}{II.24.1 (p.~719)} }
\expandafter\def\csname SRef@cdp:st:C.2.UN.80\endcsname{\href{S6_SERIES_GRIND_II.pdf\#nameddest=tag-01ES}{II.24.10 (p.~755)} }
\expandafter\def\csname SRef@cdp:st:C.2.UN.81\endcsname{\href{S6_SERIES_GRIND_II.pdf\#nameddest=tag-01ET}{II.24.11 (p.~757)} }
\expandafter\def\csname SRef@cdp:st:C.2.UN.82\endcsname{\href{S6_SERIES_GRIND_II.pdf\#nameddest=tag-01EU}{II.24.12 (p.~762)} }
\expandafter\def\csname SRef@cdp:st:C.2.UN.83\endcsname{\href{S6_SERIES_GRIND_II.pdf\#nameddest=tag-01EV}{II.24.13 (p.~765)} }
\expandafter\def\csname SRef@cdp:st:C.2.UN.84\endcsname{\href{S6_SERIES_GRIND_II.pdf\#nameddest=tag-01EW}{II.24.14 (p.~770)} }
\expandafter\def\csname SRef@cdp:st:C.2.UN.78\endcsname{\href{S6_SERIES_GRIND_II.pdf\#nameddest=tag-01EP}{II.24.15 (p.~773)} }
\expandafter\def\csname SRef@cdp:st:C.2.UN.72\endcsname{\href{S6_SERIES_GRIND_II.pdf\#nameddest=tag-01EJ}{II.24.16 (p.~778)} }
\expandafter\def\csname SRef@cdp:st:C.2.UN.85\endcsname{\href{S6_SERIES_GRIND_II.pdf\#nameddest=tag-01EX}{II.24.17 (p.~782)} }
\expandafter\def\csname SRef@cdp:st:C.2.UN.86\endcsname{\href{S6_SERIES_GRIND_II.pdf\#nameddest=tag-01EY}{II.24.18 (p.~788)} }
\expandafter\def\csname SRef@cdp:st:C.2.UN.70\endcsname{\href{S6_SERIES_GRIND_II.pdf\#nameddest=tag-01EH}{II.24.2 (p.~722)} }
\expandafter\def\csname SRef@cdp:st:C.2.UN.71\endcsname{\href{S6_SERIES_GRIND_II.pdf\#nameddest=tag-01EI}{II.24.3 (p.~726)} }
\expandafter\def\csname SRef@cdp:st:C.2.UN.73\endcsname{\href{S6_SERIES_GRIND_II.pdf\#nameddest=tag-01EK}{II.24.4 (p.~730)} }
\expandafter\def\csname SRef@cdp:st:C.2.UN.76\endcsname{\href{S6_SERIES_GRIND_II.pdf\#nameddest=tag-01EN}{II.24.5 (p.~733)} }
\expandafter\def\csname SRef@cdp:st:C.2.UN.74\endcsname{\href{S6_SERIES_GRIND_II.pdf\#nameddest=tag-01EL}{II.24.6 (p.~739)} }
\expandafter\def\csname SRef@cdp:st:C.2.UN.75\endcsname{\href{S6_SERIES_GRIND_II.pdf\#nameddest=tag-01EM}{II.24.7 (p.~743)} }
\expandafter\def\csname SRef@cdp:st:C.2.UN.77\endcsname{\href{S6_SERIES_GRIND_II.pdf\#nameddest=tag-01EO}{II.24.8 (p.~747)} }
\expandafter\def\csname SRef@cdp:st:C.2.UN.79\endcsname{\href{S6_SERIES_GRIND_II.pdf\#nameddest=tag-01EQ}{II.24.9 (p.~752)} }
\expandafter\def\csname SRef@sec:II:06\endcsname{\href{S6_SERIES_GRIND_II.pdf\#nameddest=tag-0342}{II.19 (p.~413)} }
\expandafter\def\csname SRef@sec:II:07\endcsname{\href{S6_SERIES_GRIND_II.pdf\#nameddest=tag-0343}{II.20 (p.~474)} }
\expandafter\def\csname SRef@sec:II:all:08\endcsname{\href{S6_SERIES_GRIND_II.pdf\#nameddest=tag-034B}{II.21 (p.~543)} }
\expandafter\def\csname SRef@sec:II:09\endcsname{\href{S6_SERIES_GRIND_II.pdf\#nameddest=tag-0345}{II.22 (p.~579)} }
\expandafter\def\csname SRef@sec:II:10\endcsname{\href{S6_SERIES_GRIND_II.pdf\#nameddest=tag-0346}{II.23 (p.~636)} }
\expandafter\def\csname SRef@sec:II:11\endcsname{\href{S6_SERIES_GRIND_II.pdf\#nameddest=tag-0347}{II.24 (p.~719)} }
\expandafter\def\csname SRef@sec:II:all:12\endcsname{\href{S6_SERIES_GRIND_II.pdf\#nameddest=tag-034C}{II.25 (p.~792)} }
\expandafter\def\csname SRef@cdp:st:C.2.G.1\endcsname{\href{S6_SERIES_GRIND_II.pdf\#nameddest=tag-01BZ}{II.25.1 (p.~792)} }
\expandafter\def\csname SRef@inv:E1.1\endcsname{\href{S6_SERIES_GRIND_II.pdf\#nameddest=tag-01ST}{II.? (p.~?)} }
\expandafter\def\csname SRef@inv:E2.1\endcsname{\href{S6_SERIES_GRIND_II.pdf\#nameddest=tag-01SU}{II.? (p.~?)} }
\expandafter\def\csname SRef@push:src:C.1.1\endcsname{\href{S6_SERIES_GRIND_III.pdf\#nameddest=tag-031B}{III.? (p.~?)} }
\expandafter\def\csname SRef@push:src:C.1.2\endcsname{\href{S6_SERIES_GRIND_III.pdf\#nameddest=tag-031C}{III.? (p.~?)} }
\expandafter\def\csname SRef@push:src:C.2.1\endcsname{\href{S6_SERIES_GRIND_III.pdf\#nameddest=tag-031D}{III.? (p.~?)} }
\expandafter\def\csname SRef@push:src:C.3.1\endcsname{\href{S6_SERIES_GRIND_III.pdf\#nameddest=tag-031G}{III.? (p.~?)} }
\expandafter\def\csname SRef@push:src:C.3.3\endcsname{\href{S6_SERIES_GRIND_III.pdf\#nameddest=tag-0320}{III.? (p.~?)} }
\expandafter\def\csname SRef@push:src:E.3.2\endcsname{\hyperlink{tag-033J}{I.\ref*{push:src:E.3.2}}}
\expandafter\def\csname SRef@push:src:E.3.3\endcsname{\hyperlink{tag-033K}{I.\ref*{push:src:E.3.3}}}
\expandafter\def\csname SRef@cdp:st:A.2.3\endcsname{\href{S6_SERIES_GRIND_II.pdf\#nameddest=tag-01AH}{II.? (p.~?)} }
\expandafter\def\csname SRef@cdp:st:C.1.2\endcsname{\href{S6_SERIES_GRIND_II.pdf\#nameddest=tag-01BO}{II.? (p.~?)} }
\expandafter\def\csname SRef@cdp:st:C.2.NC.1\endcsname{\href{S6_SERIES_GRIND_II.pdf\#nameddest=tag-01C8}{II.? (p.~?)} }
\expandafter\def\csname SRef@cdp:st:C.2.RAT.6\endcsname{\href{S6_SERIES_GRIND_II.pdf\#nameddest=tag-01CI}{II.? (p.~?)} }
\expandafter\def\csname SRef@cdp:st:D.1.1\endcsname{\href{S6_SERIES_GRIND_II.pdf\#nameddest=tag-01FU}{II.? (p.~?)} }
\expandafter\def\csname SRef@cdp:st:D.1.3\endcsname{\href{S6_SERIES_GRIND_II.pdf\#nameddest=tag-01G1}{II.? (p.~?)} }
\expandafter\def\csname SRef@cdp:st:E2\endcsname{\href{S6_SERIES_GRIND_II.pdf\#nameddest=tag-01GR}{II.? (p.~?)} }
\expandafter\def\csname SRef@cdp:st:E3\endcsname{\href{S6_SERIES_GRIND_II.pdf\#nameddest=tag-01GS}{II.? (p.~?)} }
\expandafter\def\csname SRef@cdp:st:E4\endcsname{\href{S6_SERIES_GRIND_II.pdf\#nameddest=tag-01GT}{II.? (p.~?)} }
\expandafter\def\csname SRef@cdp:st:E19\endcsname{\href{S6_SERIES_GRIND_II.pdf\#nameddest=tag-01GQ}{II.? (p.~?)} }
\expandafter\def\csname SRef@maps:src:B.6.9\endcsname{\href{S6_SERIES_GRIND_IV.pdf\#nameddest=tag-02K1}{IV.? (p.~?)} }
\expandafter\def\csname SRef@maps:src:D.2.1\endcsname{\href{S6_SERIES_GRIND_III.pdf\#nameddest=tag-02KJ}{III.? (p.~?)} }
\expandafter\def\csname SRef@maps:src:F.1.2\endcsname{\href{S6_SERIES_GRIND_V.pdf\#nameddest=tag-02MV}{V.? (p.~?)} }
\expandafter\def\csname SRef@maps:src:F.1.3\endcsname{\href{S6_SERIES_GRIND_V.pdf\#nameddest=tag-02MW}{V.? (p.~?)} }
\expandafter\def\csname SRef@ari:C.7.5\endcsname{\hyperlink{tag-001E}{I.\ref*{ari:C.7.5}}}
\expandafter\def\csname SRef@wb:lem:triangle-generator-orientations\endcsname{\hyperlink{tag-03BM}{I.\ref*{wb:lem:triangle-generator-orientations}}}
\expandafter\def\csname SRef@wb:prop:period-input-audit\endcsname{\hyperlink{tag-03BS}{I.\ref*{wb:prop:period-input-audit}}}
\expandafter\def\csname SRef@wb:prop:period-exact-admissible-constants\endcsname{\hyperlink{tag-03BR}{I.\ref*{wb:prop:period-exact-admissible-constants}}}
\expandafter\def\csname SRef@wb:lem:toric-chart-control\endcsname{\hyperlink{tag-03BK}{I.\ref*{wb:lem:toric-chart-control}}}
\expandafter\def\csname SRef@wb:thm:candidate-cusp-geometry\endcsname{\hyperlink{tag-03DK}{I.\ref*{wb:thm:candidate-cusp-geometry}}}
\expandafter\def\csname SRef@wb:prop:candidate-cusp-overlap\endcsname{\hyperlink{tag-03BN}{I.\ref*{wb:prop:candidate-cusp-overlap}}}
\expandafter\def\csname SRef@wb:lem:middle-period-proper\endcsname{\hyperlink{tag-03BI}{I.\ref*{wb:lem:middle-period-proper}}}
\expandafter\def\csname SRef@wb:prop:candidate-finite-fillings\endcsname{\hyperlink{tag-03BO}{I.\ref*{wb:prop:candidate-finite-fillings}}}
\expandafter\def\csname SRef@wb:lem:finite-filling-stalk-map\endcsname{\hyperlink{tag-03BG}{I.\ref*{wb:lem:finite-filling-stalk-map}}}
\expandafter\def\csname SRef@wb:thm:candidate-global-geometry\endcsname{\hyperlink{tag-03DL}{I.\ref*{wb:thm:candidate-global-geometry}}}
\expandafter\def\csname SRef@wb:prop:monodromy-certificate\endcsname{\hyperlink{tag-03BQ}{I.\ref*{wb:prop:monodromy-certificate}}}
\expandafter\def\csname SRef@wb:thm:cusp-lattices\endcsname{\hyperlink{tag-03DO}{I.\ref*{wb:thm:cusp-lattices}}}
\expandafter\def\csname SRef@wb:thm:coinvariants\endcsname{\hyperlink{tag-03DN}{I.\ref*{wb:thm:coinvariants}}}
\expandafter\def\csname SRef@wb:prop:finite-indices\endcsname{\hyperlink{tag-03BP}{I.\ref*{wb:prop:finite-indices}}}
\expandafter\def\csname SRef@wb:prop:terminal-group\endcsname{\hyperlink{tag-03BV}{I.\ref*{wb:prop:terminal-group}}}
\expandafter\def\csname SRef@wb:prop:terminal-degree-two\endcsname{\hyperlink{tag-03BU}{I.\ref*{wb:prop:terminal-degree-two}}}
\expandafter\def\csname SRef@wb:prop:terminal-all-degree\endcsname{\hyperlink{tag-03BT}{I.\ref*{wb:prop:terminal-all-degree}}}
\expandafter\def\csname SRef@wb:ffcert:prop:integral-summands\endcsname{\hyperlink{tag-03A8}{I.\ref*{wb:ffcert:prop:integral-summands}}}
\expandafter\def\csname SRef@wb:ffcert:thm:product-covers\endcsname{\hyperlink{tag-03AD}{I.\ref*{wb:ffcert:thm:product-covers}}}
\expandafter\def\csname SRef@wb:ffcert:prop:quotient-charts\endcsname{\hyperlink{tag-03A9}{I.\ref*{wb:ffcert:prop:quotient-charts}}}
\expandafter\def\csname SRef@wb:ffcert:thm:normalization\endcsname{\hyperlink{tag-03AB}{I.\ref*{wb:ffcert:thm:normalization}}}
\expandafter\def\csname SRef@wb:ffcert:thm:peripheral-kernel\endcsname{\hyperlink{tag-03AC}{I.\ref*{wb:ffcert:thm:peripheral-kernel}}}
\expandafter\def\csname SRef@wb:ffcert:cor:covering-h1\endcsname{\hyperlink{tag-03A7}{I.\ref*{wb:ffcert:cor:covering-h1}}}
\expandafter\def\csname SRef@wb:ffcert:thm:normal-trivialization\endcsname{\hyperlink{tag-03AA}{I.\ref*{wb:ffcert:thm:normal-trivialization}}}
\expandafter\def\csname SRef@wb:lem:torus-algebraic-reduction-audit\endcsname{\hyperlink{tag-03BL}{I.\ref*{wb:lem:torus-algebraic-reduction-audit}}}
\expandafter\def\csname SRef@wb:prop:very-general-torus-audit\endcsname{\hyperlink{tag-03BW}{I.\ref*{wb:prop:very-general-torus-audit}}}
\expandafter\def\csname SRef@wb:lem:meromorphic-descent-audit\endcsname{\hyperlink{tag-03BH}{I.\ref*{wb:lem:meromorphic-descent-audit}}}
\expandafter\def\csname SRef@wb:thm:algebraic-dimension-audit\endcsname{\hyperlink{tag-03DJ}{I.\ref*{wb:thm:algebraic-dimension-audit}}}
\expandafter\def\csname SRef@wb:lem:semidirect-abelianization\endcsname{\hyperlink{tag-03BJ}{I.\ref*{wb:lem:semidirect-abelianization}}}
\expandafter\def\csname SRef@mapsI:top:thm:p0-integral-specialization\endcsname{\hyperlink{tag-02Q5}{I.\ref*{mapsI:top:thm:p0-integral-specialization}}}
\expandafter\def\csname SRef@mapsI:top:prop:leray-e2\endcsname{\hyperlink{tag-02PU}{I.\ref*{mapsI:top:prop:leray-e2}}}
\expandafter\def\csname SRef@mapsI:top:lem:first-transgression\endcsname{\hyperlink{tag-02PN}{I.\ref*{mapsI:top:lem:first-transgression}}}
\expandafter\def\csname SRef@mapsI:top:lem:higher-transgressions\endcsname{\hyperlink{tag-02PO}{I.\ref*{mapsI:top:lem:higher-transgressions}}}
\expandafter\def\csname SRef@mapsI:top:thm:h2-vanishing\endcsname{\hyperlink{tag-02Q3}{I.\ref*{mapsI:top:thm:h2-vanishing}}}
\expandafter\def\csname SRef@mapsI:top:prop:geometric-pi1\endcsname{\hyperlink{tag-02PT}{I.\ref*{mapsI:top:prop:geometric-pi1}}}
\expandafter\def\csname SRef@mapsI:top:cor:candidate-sphere\endcsname{\hyperlink{tag-02PK}{I.\ref*{mapsI:top:cor:candidate-sphere}}}
\expandafter\def\csname SRef@mapsI:top:prop:topological-spine\endcsname{\hyperlink{tag-02PV}{I.\ref*{mapsI:top:prop:topological-spine}}}
\expandafter\def\csname SRef@mapsI:top:lem:collapse-straightening\endcsname{\hyperlink{tag-02PM}{I.\ref*{mapsI:top:lem:collapse-straightening}}}
\expandafter\def\csname SRef@mapsI:top:lem:positive-collapse\endcsname{\hyperlink{tag-02PQ}{I.\ref*{mapsI:top:lem:positive-collapse}}}
\expandafter\def\csname SRef@mapsI:top:lem:phase-spread-collapse\endcsname{\hyperlink{tag-02PP}{I.\ref*{mapsI:top:lem:phase-spread-collapse}}}
\expandafter\def\csname SRef@mapsI:top:prop:verified-cellular-collapse\endcsname{\hyperlink{tag-02PW}{I.\ref*{mapsI:top:prop:verified-cellular-collapse}}}
\expandafter\def\csname SRef@mapsI:top:prop:verified-collapse-homology\endcsname{\hyperlink{tag-02PX}{I.\ref*{mapsI:top:prop:verified-collapse-homology}}}
\expandafter\def\csname SRef@mapsI:top:thm:collapse-equals-specialization\endcsname{\hyperlink{tag-02Q2}{I.\ref*{mapsI:top:thm:collapse-equals-specialization}}}
\expandafter\def\csname SRef@mapsI:top:remark:L600\endcsname{\hyperlink{tag-02Q0}{I.\ref*{mapsI:top:remark:L600}}}
\expandafter\def\csname SRef@mapsI:top:record:vankampen\endcsname{\hyperlink{tag-02PY}{I.\ref*{mapsI:top:record:vankampen}}}
\expandafter\def\csname SRef@mapsI:top:prop:actual-finite-band-maps\endcsname{\hyperlink{tag-02PR}{I.\ref*{mapsI:top:prop:actual-finite-band-maps}}}
\expandafter\def\csname SRef@mapsI:top:prop:cusp-mv-blocks\endcsname{\hyperlink{tag-02PS}{I.\ref*{mapsI:top:prop:cusp-mv-blocks}}}
\expandafter\def\csname SRef@mapsI:top:remark:L920\endcsname{\hyperlink{tag-02Q1}{I.\ref*{mapsI:top:remark:L920}}}
\expandafter\def\csname SRef@mapsI:top:thm:mv-leray-comparison\endcsname{\hyperlink{tag-02Q4}{I.\ref*{mapsI:top:thm:mv-leray-comparison}}}
\expandafter\def\csname SRef@mapsI:top:remark:L1113\endcsname{\hyperlink{tag-02PZ}{I.\ref*{mapsI:top:remark:L1113}}}
\expandafter\def\csname SRef@mapsI:top:lem:cohomological-closure\endcsname{\hyperlink{tag-02PL}{I.\ref*{mapsI:top:lem:cohomological-closure}}}
\expandafter\def\csname SRef@mapsI:top:thm:sphere-recognition\endcsname{\hyperlink{tag-02Q6}{I.\ref*{mapsI:top:thm:sphere-recognition}}}
\expandafter\def\csname SRef@mapsI:top:cor:abstract-s6\endcsname{\hyperlink{tag-02PI}{I.\ref*{mapsI:top:cor:abstract-s6}}}
\expandafter\def\csname SRef@mapsI:top:cor:candidate-mv-sphere\endcsname{\hyperlink{tag-02PJ}{I.\ref*{mapsI:top:cor:candidate-mv-sphere}}}
\expandafter\def\csname SRef@wb:cusp-cone-S12-almost-complex-obstruction\endcsname{\href{S6_SERIES_GRIND_III.pdf\#nameddest=tag-035S}{III.? (p.~?)} }
\expandafter\def\csname SRef@courseIII:almost-complex-structures-and-integrability:section-1\endcsname{\href{S6_SERIES_GRIND_III.pdf\#nameddest=tag-01I4}{III.? (p.~?)} }
\expandafter\def\csname SRef@courseIII:almost-complex-structures-and-integrability:section-2\endcsname{\href{S6_SERIES_GRIND_III.pdf\#nameddest=tag-01I7}{III.? (p.~?)} }
\expandafter\def\csname SRef@courseIII:almost-complex-structures-and-integrability:section-3\endcsname{\href{S6_SERIES_GRIND_III.pdf\#nameddest=tag-01I8}{III.? (p.~?)} }
\expandafter\def\csname SRef@courseIII:almost-complex-structures-and-integrability:section-4\endcsname{\href{S6_SERIES_GRIND_III.pdf\#nameddest=tag-01I9}{III.? (p.~?)} }
\expandafter\def\csname SRef@courseIII:almost-complex-structures-and-integrability:section-5\endcsname{\href{S6_SERIES_GRIND_III.pdf\#nameddest=tag-01IA}{III.? (p.~?)} }
\expandafter\def\csname SRef@courseIII:almost-complex-structures-and-integrability:section-6\endcsname{\href{S6_SERIES_GRIND_III.pdf\#nameddest=tag-01IB}{III.? (p.~?)} }
\expandafter\def\csname SRef@courseIII:almost-complex-structures-and-integrability:section-7\endcsname{\href{S6_SERIES_GRIND_III.pdf\#nameddest=tag-01IC}{III.? (p.~?)} }
\expandafter\def\csname SRef@courseIII:almost-complex-structures-and-integrability:section-8\endcsname{\href{S6_SERIES_GRIND_III.pdf\#nameddest=tag-01ID}{III.? (p.~?)} }
\expandafter\def\csname SRef@courseIII:almost-complex-structures-and-integrability:section-9\endcsname{\href{S6_SERIES_GRIND_III.pdf\#nameddest=tag-01IE}{III.? (p.~?)} }
\expandafter\def\csname SRef@courseIII:almost-complex-structures-and-integrability:section-10\endcsname{\href{S6_SERIES_GRIND_III.pdf\#nameddest=tag-01I5}{III.? (p.~?)} }
\expandafter\def\csname SRef@courseIII:almost-complex-structures-and-integrability:section-11\endcsname{\href{S6_SERIES_GRIND_III.pdf\#nameddest=tag-01I6}{III.? (p.~?)} }
\expandafter\def\csname SRef@mapsI:course:CG-S6-07-HOMOTOPY\endcsname{\hyperlink{tag-02OR}{I.\ref*{mapsI:course:CG-S6-07-HOMOTOPY}}}
\expandafter\def\csname SRef@courseIII:integral-homology-and-sphere-recognition:section\endcsname{\href{S6_SERIES_GRIND_III.pdf\#nameddest=tag-01K9}{III.? (p.~?)} }
\expandafter\def\csname SRef@courseIII:sphere-twists-and-the-original-smooth-six-sphere-comparison:twist-actual-tube\endcsname{\href{S6_SERIES_GRIND_III.pdf\#nameddest=tag-01MN}{III.? (p.~?)} }
\expandafter\def\csname SRef@courseIII:sphere-twists-and-the-original-smooth-six-sphere-comparison:twist-signed-action\endcsname{\href{S6_SERIES_GRIND_III.pdf\#nameddest=tag-01MT}{III.? (p.~?)} }
\expandafter\def\csname SRef@courseIII:sphere-twists-and-the-original-smooth-six-sphere-comparison:twist-integral-lagrangians\endcsname{\href{S6_SERIES_GRIND_III.pdf\#nameddest=tag-01MP}{III.? (p.~?)} }
\expandafter\def\csname SRef@courseIII:sphere-twists-and-the-original-smooth-six-sphere-comparison:twist-original-middle-level\endcsname{\href{S6_SERIES_GRIND_III.pdf\#nameddest=tag-01MR}{III.? (p.~?)} }
\expandafter\def\csname SRef@courseIII:sphere-twists-and-the-original-smooth-six-sphere-comparison:twist-original-gluing\endcsname{\href{S6_SERIES_GRIND_III.pdf\#nameddest=tag-01MQ}{III.? (p.~?)} }
\expandafter\def\csname SRef@courseIII:sphere-twists-and-the-original-smooth-six-sphere-comparison:twist-actual-cancellation\endcsname{\href{S6_SERIES_GRIND_III.pdf\#nameddest=tag-01MM}{III.? (p.~?)} }
\expandafter\def\csname SRef@courseIII:sphere-twists-and-the-original-smooth-six-sphere-comparison:twist-original-smooth-map\endcsname{\href{S6_SERIES_GRIND_III.pdf\#nameddest=tag-01MS}{III.? (p.~?)} }
\expandafter\def\csname SRef@courseIII:sphere-twists-and-the-original-smooth-six-sphere-comparison:twist-general-stable-framing\endcsname{\href{S6_SERIES_GRIND_III.pdf\#nameddest=tag-01MO}{III.? (p.~?)} }
\expandafter\def\csname SRef@courseIII:sphere-twists-and-the-original-smooth-six-sphere-comparison:twist-worked-exercises\endcsname{\href{S6_SERIES_GRIND_III.pdf\#nameddest=tag-01MU}{III.? (p.~?)} }
\expandafter\def\csname SRef@mapsaudit:almost-complex-homotopy\endcsname{\href{S6_SERIES_GRIND_III.pdf\#nameddest=tag-02QD}{III.? (p.~?)} }
\expandafter\def\csname SRef@mapsI:hodge:line\endcsname{\hyperlink{tag-02OV}{I.\ref*{mapsI:hodge:line}}}
\expandafter\def\csname SRef@mapsI:chain:lemma:1.1\endcsname{\hyperlink{tag-02OL}{I.\ref*{mapsI:chain:lemma:1.1}}}
\expandafter\def\csname SRef@mapsI:chain:lemma:2.1\endcsname{\hyperlink{tag-02OM}{I.\ref*{mapsI:chain:lemma:2.1}}}
\expandafter\def\csname SRef@mapsI:chain:corollary:2.2\endcsname{\hyperlink{tag-02OK}{I.\ref*{mapsI:chain:corollary:2.2}}}
\expandafter\def\csname SRef@mapsI:chain:lemma:4.1\endcsname{\hyperlink{tag-02ON}{I.\ref*{mapsI:chain:lemma:4.1}}}
\expandafter\def\csname SRef@mapsI:cw:hurewicz-theorem-result\endcsname{\hyperlink{tag-02OU}{I.\ref*{mapsI:cw:hurewicz-theorem-result}}}
\expandafter\def\csname SRef@mapsI:cw:exercise:5.1\endcsname{\hyperlink{tag-02OS}{I.\ref*{mapsI:cw:exercise:5.1}}}
\expandafter\def\csname SRef@mapsI:cw:exercise:5.2\endcsname{\hyperlink{tag-02OT}{I.\ref*{mapsI:cw:exercise:5.2}}}
\expandafter\def\csname SRef@mapsI:pf:theorem:5.1\endcsname{\hyperlink{tag-02OY}{I.\ref*{mapsI:pf:theorem:5.1}}}
\expandafter\def\csname SRef@mapsI:pf:exercise:6.1\endcsname{\hyperlink{tag-02OW}{I.\ref*{mapsI:pf:exercise:6.1}}}
\expandafter\def\csname SRef@mapsI:pf:exercise:6.2\endcsname{\hyperlink{tag-02OX}{I.\ref*{mapsI:pf:exercise:6.2}}}
\expandafter\def\csname SRef@courseIII:thom-euler-for-lesson11:CG-S6-CHAIN-LEMMA-5.1\endcsname{\href{S6_SERIES_GRIND_III.pdf\#nameddest=tag-01N3}{III.? (p.~?)} }
\expandafter\def\csname SRef@courseIII:thom-euler-for-lesson11:CG-S6-CHAIN-THEOREM-5.2\endcsname{\href{S6_SERIES_GRIND_III.pdf\#nameddest=tag-01N7}{III.? (p.~?)} }
\expandafter\def\csname SRef@courseIII:thom-euler-for-lesson11:CG-S6-CHAIN-PROPOSITION-6.1\endcsname{\href{S6_SERIES_GRIND_III.pdf\#nameddest=tag-01N4}{III.? (p.~?)} }
\expandafter\def\csname SRef@courseIII:thom-euler-for-lesson11:CG-S6-CHAIN-THEOREM-7.1\endcsname{\href{S6_SERIES_GRIND_III.pdf\#nameddest=tag-01N8}{III.? (p.~?)} }
\expandafter\def\csname SRef@courseIII:thom-euler-for-lesson11:CG-S6-CHAIN-PROPOSITION-8.1\endcsname{\href{S6_SERIES_GRIND_III.pdf\#nameddest=tag-01N5}{III.? (p.~?)} }
\expandafter\def\csname SRef@courseIII:thom-euler-for-lesson11:CG-S6-CHAIN-PROPOSITION-8.2\endcsname{\href{S6_SERIES_GRIND_III.pdf\#nameddest=tag-01N6}{III.? (p.~?)} }
\expandafter\def\csname SRef@courseIII:projective-splitting-for-lesson11:CG-S6-PROJECTIVE-1\endcsname{\href{S6_SERIES_GRIND_III.pdf\#nameddest=tag-01M7}{III.? (p.~?)} }
\expandafter\def\csname SRef@courseIII:projective-splitting-for-lesson11:CG-S6-PROJECTIVE-2\endcsname{\href{S6_SERIES_GRIND_III.pdf\#nameddest=tag-01M8}{III.? (p.~?)} }
\expandafter\def\csname SRef@courseIII:projective-splitting-for-lesson11:CG-S6-PROJECTIVE-3\endcsname{\href{S6_SERIES_GRIND_III.pdf\#nameddest=tag-01M9}{III.? (p.~?)} }
\expandafter\def\csname SRef@courseIII:projective-splitting-for-lesson11:CG-S6-PROJECTIVE-4\endcsname{\href{S6_SERIES_GRIND_III.pdf\#nameddest=tag-01MA}{III.? (p.~?)} }
\expandafter\def\csname SRef@courseIII:projective-splitting-for-lesson11:CG-S6-PROJECTIVE-5\endcsname{\href{S6_SERIES_GRIND_III.pdf\#nameddest=tag-01MB}{III.? (p.~?)} }
\expandafter\def\csname SRef@courseIII:integral-chern-classes-for-lesson11:CG-S6-CHERN-1\endcsname{\href{S6_SERIES_GRIND_III.pdf\#nameddest=tag-01K6}{III.? (p.~?)} }
\expandafter\def\csname SRef@courseIII:integral-chern-classes-for-lesson11:CG-S6-CHERN-2\endcsname{\href{S6_SERIES_GRIND_III.pdf\#nameddest=tag-01K7}{III.? (p.~?)} }
\expandafter\def\csname SRef@courseIII:integral-chern-classes-for-lesson11:CG-S6-CHERN-3\endcsname{\href{S6_SERIES_GRIND_III.pdf\#nameddest=tag-01K8}{III.? (p.~?)} }
\expandafter\def\csname SRef@courseIII:framed6:1\endcsname{\href{S6_SERIES_GRIND_III.pdf\#nameddest=tag-01J8}{III.? (p.~?)} }
\expandafter\def\csname SRef@courseIII:framed6:2\endcsname{\href{S6_SERIES_GRIND_III.pdf\#nameddest=tag-01JB}{III.? (p.~?)} }
\expandafter\def\csname SRef@courseIII:framed6:3\endcsname{\href{S6_SERIES_GRIND_III.pdf\#nameddest=tag-01JC}{III.? (p.~?)} }
\expandafter\def\csname SRef@courseIII:framed6:4\endcsname{\href{S6_SERIES_GRIND_III.pdf\#nameddest=tag-01JD}{III.? (p.~?)} }
\expandafter\def\csname SRef@courseIII:framed6:5\endcsname{\href{S6_SERIES_GRIND_III.pdf\#nameddest=tag-01JE}{III.? (p.~?)} }
\expandafter\def\csname SRef@courseIII:framed6:6\endcsname{\href{S6_SERIES_GRIND_III.pdf\#nameddest=tag-01JF}{III.? (p.~?)} }
\expandafter\def\csname SRef@courseIII:framed6:7\endcsname{\href{S6_SERIES_GRIND_III.pdf\#nameddest=tag-01JG}{III.? (p.~?)} }
\expandafter\def\csname SRef@courseIII:framed6:8\endcsname{\href{S6_SERIES_GRIND_III.pdf\#nameddest=tag-01JH}{III.? (p.~?)} }
\expandafter\def\csname SRef@courseIII:framed6:9\endcsname{\href{S6_SERIES_GRIND_III.pdf\#nameddest=tag-01JI}{III.? (p.~?)} }
\expandafter\def\csname SRef@courseIII:framed6:10.1\endcsname{\href{S6_SERIES_GRIND_III.pdf\#nameddest=tag-01J9}{III.? (p.~?)} }
\expandafter\def\csname SRef@courseIII:framed6:10.2\endcsname{\href{S6_SERIES_GRIND_III.pdf\#nameddest=tag-01JA}{III.? (p.~?)} }
\expandafter\def\csname SRef@courseIII:whitney-receiving:1\endcsname{\href{S6_SERIES_GRIND_III.pdf\#nameddest=tag-01NG}{III.? (p.~?)} }
\expandafter\def\csname SRef@courseIII:whitney-receiving:2\endcsname{\href{S6_SERIES_GRIND_III.pdf\#nameddest=tag-01NH}{III.? (p.~?)} }
\expandafter\def\csname SRef@courseIII:whitney-receiving:3\endcsname{\href{S6_SERIES_GRIND_III.pdf\#nameddest=tag-01NI}{III.? (p.~?)} }
\expandafter\def\csname SRef@courseIII:whitney-receiving:4\endcsname{\href{S6_SERIES_GRIND_III.pdf\#nameddest=tag-01NJ}{III.? (p.~?)} }
\expandafter\def\csname SRef@courseIII:whitney-receiving:5.1\endcsname{\href{S6_SERIES_GRIND_III.pdf\#nameddest=tag-01NK}{III.? (p.~?)} }
\expandafter\def\csname SRef@courseIII:whitney-receiving:5.2\endcsname{\href{S6_SERIES_GRIND_III.pdf\#nameddest=tag-01NL}{III.? (p.~?)} }
\expandafter\def\csname SRef@courseIII:higher-slides:1\endcsname{\href{S6_SERIES_GRIND_III.pdf\#nameddest=tag-01JY}{III.? (p.~?)} }
\expandafter\def\csname SRef@courseIII:higher-slides:2\endcsname{\href{S6_SERIES_GRIND_III.pdf\#nameddest=tag-01JZ}{III.? (p.~?)} }
\expandafter\def\csname SRef@courseIII:higher-slides:3\endcsname{\href{S6_SERIES_GRIND_III.pdf\#nameddest=tag-01K0}{III.? (p.~?)} }
\expandafter\def\csname SRef@courseIII:higher-slides:4\endcsname{\href{S6_SERIES_GRIND_III.pdf\#nameddest=tag-01K1}{III.? (p.~?)} }
\expandafter\def\csname SRef@courseIII:higher-slides:5\endcsname{\href{S6_SERIES_GRIND_III.pdf\#nameddest=tag-01K2}{III.? (p.~?)} }
\expandafter\def\csname SRef@courseIII:higher-slides:6\endcsname{\href{S6_SERIES_GRIND_III.pdf\#nameddest=tag-01K3}{III.? (p.~?)} }
\expandafter\def\csname SRef@courseIII:higher-slides:8.1\endcsname{\href{S6_SERIES_GRIND_III.pdf\#nameddest=tag-01K4}{III.? (p.~?)} }
\expandafter\def\csname SRef@courseIII:higher-slides:8.2\endcsname{\href{S6_SERIES_GRIND_III.pdf\#nameddest=tag-01K5}{III.? (p.~?)} }
\expandafter\def\csname SRef@courseIII:middle-cancel:1\endcsname{\href{S6_SERIES_GRIND_III.pdf\#nameddest=tag-01KT}{III.? (p.~?)} }
\expandafter\def\csname SRef@courseIII:middle-cancel:2\endcsname{\href{S6_SERIES_GRIND_III.pdf\#nameddest=tag-01KU}{III.? (p.~?)} }
\expandafter\def\csname SRef@courseIII:middle-cancel:3\endcsname{\href{S6_SERIES_GRIND_III.pdf\#nameddest=tag-01KV}{III.? (p.~?)} }
\expandafter\def\csname SRef@courseIII:middle-cancel:4\endcsname{\href{S6_SERIES_GRIND_III.pdf\#nameddest=tag-01KW}{III.? (p.~?)} }
\expandafter\def\csname SRef@courseIII:middle-cancel:5\endcsname{\href{S6_SERIES_GRIND_III.pdf\#nameddest=tag-01KX}{III.? (p.~?)} }
\expandafter\def\csname SRef@courseIII:middle-cancel:6\endcsname{\href{S6_SERIES_GRIND_III.pdf\#nameddest=tag-01KY}{III.? (p.~?)} }
\expandafter\def\csname SRef@courseIII:middle-cancel:7.1\endcsname{\href{S6_SERIES_GRIND_III.pdf\#nameddest=tag-01KZ}{III.? (p.~?)} }
\expandafter\def\csname SRef@courseIII:middle-cancel:7.2\endcsname{\href{S6_SERIES_GRIND_III.pdf\#nameddest=tag-01L0}{III.? (p.~?)} }
\expandafter\def\csname SRef@courseIII:morse-cancel:1\endcsname{\href{S6_SERIES_GRIND_III.pdf\#nameddest=tag-01L5}{III.? (p.~?)} }
\expandafter\def\csname SRef@courseIII:morse-cancel:2.1\endcsname{\href{S6_SERIES_GRIND_III.pdf\#nameddest=tag-01L6}{III.? (p.~?)} }
\expandafter\def\csname SRef@courseIII:morse-cancel:6\endcsname{\href{S6_SERIES_GRIND_III.pdf\#nameddest=tag-01L7}{III.? (p.~?)} }
\expandafter\def\csname SRef@courseIII:morse-cancel:7.1\endcsname{\href{S6_SERIES_GRIND_III.pdf\#nameddest=tag-01L8}{III.? (p.~?)} }
\expandafter\def\csname SRef@courseIII:morse-cancel:7.2\endcsname{\href{S6_SERIES_GRIND_III.pdf\#nameddest=tag-01L9}{III.? (p.~?)} }
\expandafter\def\csname SRef@courseIII:spin6:1\endcsname{\href{S6_SERIES_GRIND_III.pdf\#nameddest=tag-01MV}{III.? (p.~?)} }
\expandafter\def\csname SRef@courseIII:spin6:2\endcsname{\href{S6_SERIES_GRIND_III.pdf\#nameddest=tag-01MW}{III.? (p.~?)} }
\expandafter\def\csname SRef@courseIII:spin6:3\endcsname{\href{S6_SERIES_GRIND_III.pdf\#nameddest=tag-01MX}{III.? (p.~?)} }
\expandafter\def\csname SRef@courseIII:spin6:4\endcsname{\href{S6_SERIES_GRIND_III.pdf\#nameddest=tag-01MY}{III.? (p.~?)} }
\expandafter\def\csname SRef@courseIII:spin6:5\endcsname{\href{S6_SERIES_GRIND_III.pdf\#nameddest=tag-01MZ}{III.? (p.~?)} }
\expandafter\def\csname SRef@courseIII:spin6:6\endcsname{\href{S6_SERIES_GRIND_III.pdf\#nameddest=tag-01N0}{III.? (p.~?)} }
\expandafter\def\csname SRef@courseIII:spin6:7.1\endcsname{\href{S6_SERIES_GRIND_III.pdf\#nameddest=tag-01N1}{III.? (p.~?)} }
\expandafter\def\csname SRef@courseIII:spin6:7.2\endcsname{\href{S6_SERIES_GRIND_III.pdf\#nameddest=tag-01N2}{III.? (p.~?)} }
\expandafter\def\csname SRef@courseIII:integral-stable:1\endcsname{\href{S6_SERIES_GRIND_III.pdf\#nameddest=tag-01KA}{III.? (p.~?)} }
\expandafter\def\csname SRef@courseIII:integral-stable:2\endcsname{\href{S6_SERIES_GRIND_III.pdf\#nameddest=tag-01KB}{III.? (p.~?)} }
\expandafter\def\csname SRef@courseIII:integral-stable:3\endcsname{\href{S6_SERIES_GRIND_III.pdf\#nameddest=tag-01KC}{III.? (p.~?)} }
\expandafter\def\csname SRef@courseIII:integral-stable:4\endcsname{\href{S6_SERIES_GRIND_III.pdf\#nameddest=tag-01KD}{III.? (p.~?)} }
\expandafter\def\csname SRef@courseIII:integral-stable:5\endcsname{\href{S6_SERIES_GRIND_III.pdf\#nameddest=tag-01KE}{III.? (p.~?)} }
\expandafter\def\csname SRef@courseIII:integral-stable:6\endcsname{\href{S6_SERIES_GRIND_III.pdf\#nameddest=tag-01KF}{III.? (p.~?)} }
\expandafter\def\csname SRef@courseIII:integral-stable:7\endcsname{\href{S6_SERIES_GRIND_III.pdf\#nameddest=tag-01KG}{III.? (p.~?)} }
\expandafter\def\csname SRef@courseIII:integral-stable:8.1\endcsname{\href{S6_SERIES_GRIND_III.pdf\#nameddest=tag-01KH}{III.? (p.~?)} }
\expandafter\def\csname SRef@courseIII:integral-stable:8.2\endcsname{\href{S6_SERIES_GRIND_III.pdf\#nameddest=tag-01KI}{III.? (p.~?)} }
\expandafter\def\csname SRef@courseIII:pontryagin-thom:1\endcsname{\href{S6_SERIES_GRIND_III.pdf\#nameddest=tag-01LQ}{III.? (p.~?)} }
\expandafter\def\csname SRef@courseIII:pontryagin-thom:2\endcsname{\href{S6_SERIES_GRIND_III.pdf\#nameddest=tag-01LR}{III.? (p.~?)} }
\expandafter\def\csname SRef@courseIII:pontryagin-thom:3\endcsname{\href{S6_SERIES_GRIND_III.pdf\#nameddest=tag-01LS}{III.? (p.~?)} }
\expandafter\def\csname SRef@courseIII:pontryagin-thom:4\endcsname{\href{S6_SERIES_GRIND_III.pdf\#nameddest=tag-01LT}{III.? (p.~?)} }
\expandafter\def\csname SRef@courseIII:pontryagin-thom:5\endcsname{\href{S6_SERIES_GRIND_III.pdf\#nameddest=tag-01LU}{III.? (p.~?)} }
\expandafter\def\csname SRef@courseIII:pontryagin-thom:6\endcsname{\href{S6_SERIES_GRIND_III.pdf\#nameddest=tag-01LV}{III.? (p.~?)} }
\expandafter\def\csname SRef@courseIII:pontryagin-thom:7\endcsname{\href{S6_SERIES_GRIND_III.pdf\#nameddest=tag-01LW}{III.? (p.~?)} }
\expandafter\def\csname SRef@courseIII:pontryagin-thom:8\endcsname{\href{S6_SERIES_GRIND_III.pdf\#nameddest=tag-01LX}{III.? (p.~?)} }
\expandafter\def\csname SRef@courseIII:pontryagin-thom:9.1\endcsname{\href{S6_SERIES_GRIND_III.pdf\#nameddest=tag-01LY}{III.? (p.~?)} }
\expandafter\def\csname SRef@courseIII:pontryagin-thom:9.2\endcsname{\href{S6_SERIES_GRIND_III.pdf\#nameddest=tag-01LZ}{III.? (p.~?)} }
\expandafter\def\csname SRef@courseIII:relative-morse:1\endcsname{\href{S6_SERIES_GRIND_III.pdf\#nameddest=tag-01MG}{III.? (p.~?)} }
\expandafter\def\csname SRef@courseIII:relative-morse:2\endcsname{\href{S6_SERIES_GRIND_III.pdf\#nameddest=tag-01MH}{III.? (p.~?)} }
\expandafter\def\csname SRef@courseIII:relative-morse:3\endcsname{\href{S6_SERIES_GRIND_III.pdf\#nameddest=tag-01MI}{III.? (p.~?)} }
\expandafter\def\csname SRef@courseIII:relative-morse:4\endcsname{\href{S6_SERIES_GRIND_III.pdf\#nameddest=tag-01MJ}{III.? (p.~?)} }
\expandafter\def\csname SRef@courseIII:relative-morse:5.1\endcsname{\href{S6_SERIES_GRIND_III.pdf\#nameddest=tag-01MK}{III.? (p.~?)} }
\expandafter\def\csname SRef@courseIII:relative-morse:5.2\endcsname{\href{S6_SERIES_GRIND_III.pdf\#nameddest=tag-01ML}{III.? (p.~?)} }
\expandafter\def\csname SRef@courseIII:handle-arrangement:1\endcsname{\href{S6_SERIES_GRIND_III.pdf\#nameddest=tag-01JJ}{III.? (p.~?)} }
\expandafter\def\csname SRef@courseIII:handle-arrangement:2\endcsname{\href{S6_SERIES_GRIND_III.pdf\#nameddest=tag-01JK}{III.? (p.~?)} }
\expandafter\def\csname SRef@courseIII:handle-arrangement:3\endcsname{\href{S6_SERIES_GRIND_III.pdf\#nameddest=tag-01JL}{III.? (p.~?)} }
\expandafter\def\csname SRef@courseIII:handle-arrangement:4\endcsname{\href{S6_SERIES_GRIND_III.pdf\#nameddest=tag-01JM}{III.? (p.~?)} }
\expandafter\def\csname SRef@courseIII:handle-arrangement:5.1\endcsname{\href{S6_SERIES_GRIND_III.pdf\#nameddest=tag-01JN}{III.? (p.~?)} }
\expandafter\def\csname SRef@courseIII:handle-arrangement:5.2\endcsname{\href{S6_SERIES_GRIND_III.pdf\#nameddest=tag-01JO}{III.? (p.~?)} }
\expandafter\def\csname SRef@courseIII:zero-top:1\endcsname{\href{S6_SERIES_GRIND_III.pdf\#nameddest=tag-01NU}{III.? (p.~?)} }
\expandafter\def\csname SRef@courseIII:zero-top:2\endcsname{\href{S6_SERIES_GRIND_III.pdf\#nameddest=tag-01NV}{III.? (p.~?)} }
\expandafter\def\csname SRef@courseIII:zero-top:3\endcsname{\href{S6_SERIES_GRIND_III.pdf\#nameddest=tag-01NW}{III.? (p.~?)} }
\expandafter\def\csname SRef@courseIII:zero-top:4\endcsname{\href{S6_SERIES_GRIND_III.pdf\#nameddest=tag-01NX}{III.? (p.~?)} }
\expandafter\def\csname SRef@courseIII:zero-top:5\endcsname{\href{S6_SERIES_GRIND_III.pdf\#nameddest=tag-01NY}{III.? (p.~?)} }
\expandafter\def\csname SRef@courseIII:zero-top:6.1\endcsname{\href{S6_SERIES_GRIND_III.pdf\#nameddest=tag-01NZ}{III.? (p.~?)} }
\expandafter\def\csname SRef@courseIII:zero-top:6.2\endcsname{\href{S6_SERIES_GRIND_III.pdf\#nameddest=tag-01O0}{III.? (p.~?)} }
\expandafter\def\csname SRef@courseIII:trajectory-lowering:1\endcsname{\href{S6_SERIES_GRIND_III.pdf\#nameddest=tag-01N9}{III.? (p.~?)} }
\expandafter\def\csname SRef@courseIII:trajectory-lowering:2\endcsname{\href{S6_SERIES_GRIND_III.pdf\#nameddest=tag-01NA}{III.? (p.~?)} }
\expandafter\def\csname SRef@courseIII:trajectory-lowering:3\endcsname{\href{S6_SERIES_GRIND_III.pdf\#nameddest=tag-01NB}{III.? (p.~?)} }
\expandafter\def\csname SRef@courseIII:trajectory-lowering:4\endcsname{\href{S6_SERIES_GRIND_III.pdf\#nameddest=tag-01NC}{III.? (p.~?)} }
\expandafter\def\csname SRef@courseIII:trajectory-lowering:5\endcsname{\href{S6_SERIES_GRIND_III.pdf\#nameddest=tag-01ND}{III.? (p.~?)} }
\expandafter\def\csname SRef@courseIII:trajectory-lowering:6.1\endcsname{\href{S6_SERIES_GRIND_III.pdf\#nameddest=tag-01NE}{III.? (p.~?)} }
\expandafter\def\csname SRef@courseIII:trajectory-lowering:6.2\endcsname{\href{S6_SERIES_GRIND_III.pdf\#nameddest=tag-01NF}{III.? (p.~?)} }
\expandafter\def\csname SRef@courseIII:belt-complement:1\endcsname{\href{S6_SERIES_GRIND_III.pdf\#nameddest=tag-01IF}{III.? (p.~?)} }
\expandafter\def\csname SRef@courseIII:belt-complement:2\endcsname{\href{S6_SERIES_GRIND_III.pdf\#nameddest=tag-01IG}{III.? (p.~?)} }
\expandafter\def\csname SRef@courseIII:belt-complement:3\endcsname{\href{S6_SERIES_GRIND_III.pdf\#nameddest=tag-01IH}{III.? (p.~?)} }
\expandafter\def\csname SRef@courseIII:belt-complement:4\endcsname{\href{S6_SERIES_GRIND_III.pdf\#nameddest=tag-01II}{III.? (p.~?)} }
\expandafter\def\csname SRef@courseIII:belt-complement:5\endcsname{\href{S6_SERIES_GRIND_III.pdf\#nameddest=tag-01IJ}{III.? (p.~?)} }
\expandafter\def\csname SRef@courseIII:belt-complement:6\endcsname{\href{S6_SERIES_GRIND_III.pdf\#nameddest=tag-01IK}{III.? (p.~?)} }
\expandafter\def\csname SRef@courseIII:belt-complement:7\endcsname{\href{S6_SERIES_GRIND_III.pdf\#nameddest=tag-01IL}{III.? (p.~?)} }
\expandafter\def\csname SRef@courseIII:whitney-support:1\endcsname{\href{S6_SERIES_GRIND_III.pdf\#nameddest=tag-01NM}{III.? (p.~?)} }
\expandafter\def\csname SRef@courseIII:whitney-support:2\endcsname{\href{S6_SERIES_GRIND_III.pdf\#nameddest=tag-01NN}{III.? (p.~?)} }
\expandafter\def\csname SRef@courseIII:whitney-support:3\endcsname{\href{S6_SERIES_GRIND_III.pdf\#nameddest=tag-01NO}{III.? (p.~?)} }
\expandafter\def\csname SRef@courseIII:whitney-support:4\endcsname{\href{S6_SERIES_GRIND_III.pdf\#nameddest=tag-01NP}{III.? (p.~?)} }
\expandafter\def\csname SRef@courseIII:whitney-support:5\endcsname{\href{S6_SERIES_GRIND_III.pdf\#nameddest=tag-01NQ}{III.? (p.~?)} }
\expandafter\def\csname SRef@courseIII:whitney-support:6\endcsname{\href{S6_SERIES_GRIND_III.pdf\#nameddest=tag-01NR}{III.? (p.~?)} }
\expandafter\def\csname SRef@courseIII:whitney-support:7.1\endcsname{\href{S6_SERIES_GRIND_III.pdf\#nameddest=tag-01NS}{III.? (p.~?)} }
\expandafter\def\csname SRef@courseIII:whitney-support:7.2\endcsname{\href{S6_SERIES_GRIND_III.pdf\#nameddest=tag-01NT}{III.? (p.~?)} }
\expandafter\def\csname SRef@courseIII:handle-chains:1\endcsname{\href{S6_SERIES_GRIND_III.pdf\#nameddest=tag-01JP}{III.? (p.~?)} }
\expandafter\def\csname SRef@courseIII:handle-chains:2\endcsname{\href{S6_SERIES_GRIND_III.pdf\#nameddest=tag-01JQ}{III.? (p.~?)} }
\expandafter\def\csname SRef@courseIII:handle-chains:3\endcsname{\href{S6_SERIES_GRIND_III.pdf\#nameddest=tag-01JR}{III.? (p.~?)} }
\expandafter\def\csname SRef@courseIII:handle-chains:4\endcsname{\href{S6_SERIES_GRIND_III.pdf\#nameddest=tag-01JS}{III.? (p.~?)} }
\expandafter\def\csname SRef@courseIII:handle-chains:5\endcsname{\href{S6_SERIES_GRIND_III.pdf\#nameddest=tag-01JT}{III.? (p.~?)} }
\expandafter\def\csname SRef@courseIII:handle-chains:6.1\endcsname{\href{S6_SERIES_GRIND_III.pdf\#nameddest=tag-01JU}{III.? (p.~?)} }
\expandafter\def\csname SRef@courseIII:handle-chains:6.2\endcsname{\href{S6_SERIES_GRIND_III.pdf\#nameddest=tag-01JV}{III.? (p.~?)} }
\expandafter\def\csname SRef@courseIII:handle-chains:8.1\endcsname{\href{S6_SERIES_GRIND_III.pdf\#nameddest=tag-01JW}{III.? (p.~?)} }
\expandafter\def\csname SRef@courseIII:handle-chains:8.2\endcsname{\href{S6_SERIES_GRIND_III.pdf\#nameddest=tag-01JX}{III.? (p.~?)} }
\expandafter\def\csname SRef@courseIII:pair-creation:1\endcsname{\href{S6_SERIES_GRIND_III.pdf\#nameddest=tag-01LI}{III.? (p.~?)} }
\expandafter\def\csname SRef@courseIII:pair-creation:2\endcsname{\href{S6_SERIES_GRIND_III.pdf\#nameddest=tag-01LJ}{III.? (p.~?)} }
\expandafter\def\csname SRef@courseIII:pair-creation:3\endcsname{\href{S6_SERIES_GRIND_III.pdf\#nameddest=tag-01LK}{III.? (p.~?)} }
\expandafter\def\csname SRef@courseIII:pair-creation:4\endcsname{\href{S6_SERIES_GRIND_III.pdf\#nameddest=tag-01LL}{III.? (p.~?)} }
\expandafter\def\csname SRef@courseIII:pair-creation:5\endcsname{\href{S6_SERIES_GRIND_III.pdf\#nameddest=tag-01LM}{III.? (p.~?)} }
\expandafter\def\csname SRef@courseIII:pair-creation:6\endcsname{\href{S6_SERIES_GRIND_III.pdf\#nameddest=tag-01LN}{III.? (p.~?)} }
\expandafter\def\csname SRef@courseIII:pair-creation:7.1\endcsname{\href{S6_SERIES_GRIND_III.pdf\#nameddest=tag-01LO}{III.? (p.~?)} }
\expandafter\def\csname SRef@courseIII:pair-creation:7.2\endcsname{\href{S6_SERIES_GRIND_III.pdf\#nameddest=tag-01LP}{III.? (p.~?)} }
\expandafter\def\csname SRef@courseIII:low-index-slides:1\endcsname{\href{S6_SERIES_GRIND_III.pdf\#nameddest=tag-01KJ}{III.? (p.~?)} }
\expandafter\def\csname SRef@courseIII:low-index-slides:2\endcsname{\href{S6_SERIES_GRIND_III.pdf\#nameddest=tag-01KK}{III.? (p.~?)} }
\expandafter\def\csname SRef@courseIII:low-index-slides:3\endcsname{\href{S6_SERIES_GRIND_III.pdf\#nameddest=tag-01KL}{III.? (p.~?)} }
\expandafter\def\csname SRef@courseIII:low-index-slides:4\endcsname{\href{S6_SERIES_GRIND_III.pdf\#nameddest=tag-01KM}{III.? (p.~?)} }
\expandafter\def\csname SRef@courseIII:low-index-slides:5\endcsname{\href{S6_SERIES_GRIND_III.pdf\#nameddest=tag-01KN}{III.? (p.~?)} }
\expandafter\def\csname SRef@courseIII:low-index-slides:6\endcsname{\href{S6_SERIES_GRIND_III.pdf\#nameddest=tag-01KO}{III.? (p.~?)} }
\expandafter\def\csname SRef@courseIII:low-index-slides:7\endcsname{\href{S6_SERIES_GRIND_III.pdf\#nameddest=tag-01KP}{III.? (p.~?)} }
\expandafter\def\csname SRef@courseIII:low-index-slides:8\endcsname{\href{S6_SERIES_GRIND_III.pdf\#nameddest=tag-01KQ}{III.? (p.~?)} }
\expandafter\def\csname SRef@courseIII:low-index-slides:9.1\endcsname{\href{S6_SERIES_GRIND_III.pdf\#nameddest=tag-01KR}{III.? (p.~?)} }
\expandafter\def\csname SRef@courseIII:low-index-slides:9.2\endcsname{\href{S6_SERIES_GRIND_III.pdf\#nameddest=tag-01KS}{III.? (p.~?)} }
\expandafter\def\csname SRef@courseIII:adams-arithmetic:1\endcsname{\href{S6_SERIES_GRIND_III.pdf\#nameddest=tag-01HJ}{III.? (p.~?)} }
\expandafter\def\csname SRef@courseIII:adams-arithmetic:2\endcsname{\href{S6_SERIES_GRIND_III.pdf\#nameddest=tag-01HK}{III.? (p.~?)} }
\expandafter\def\csname SRef@courseIII:adams-arithmetic:3\endcsname{\href{S6_SERIES_GRIND_III.pdf\#nameddest=tag-01HL}{III.? (p.~?)} }
\expandafter\def\csname SRef@courseIII:adams-arithmetic:4\endcsname{\href{S6_SERIES_GRIND_III.pdf\#nameddest=tag-01HM}{III.? (p.~?)} }
\expandafter\def\csname SRef@courseIII:adams-arithmetic:5\endcsname{\href{S6_SERIES_GRIND_III.pdf\#nameddest=tag-01HN}{III.? (p.~?)} }
\expandafter\def\csname SRef@courseIII:adams-arithmetic:6\endcsname{\href{S6_SERIES_GRIND_III.pdf\#nameddest=tag-01HO}{III.? (p.~?)} }
\expandafter\def\csname SRef@courseIII:adams-arithmetic:7\endcsname{\href{S6_SERIES_GRIND_III.pdf\#nameddest=tag-01HP}{III.? (p.~?)} }
\expandafter\def\csname SRef@courseIII:adams-arithmetic:8\endcsname{\href{S6_SERIES_GRIND_III.pdf\#nameddest=tag-01HQ}{III.? (p.~?)} }
\expandafter\def\csname SRef@courseIII:adams-arithmetic:9\endcsname{\href{S6_SERIES_GRIND_III.pdf\#nameddest=tag-01HR}{III.? (p.~?)} }
\expandafter\def\csname SRef@courseIII:adams-tower:1\endcsname{\href{S6_SERIES_GRIND_III.pdf\#nameddest=tag-01HS}{III.? (p.~?)} }
\expandafter\def\csname SRef@courseIII:adams-tower:2\endcsname{\href{S6_SERIES_GRIND_III.pdf\#nameddest=tag-01HV}{III.? (p.~?)} }
\expandafter\def\csname SRef@courseIII:adams-tower:3\endcsname{\href{S6_SERIES_GRIND_III.pdf\#nameddest=tag-01HW}{III.? (p.~?)} }
\expandafter\def\csname SRef@courseIII:adams-tower:4\endcsname{\href{S6_SERIES_GRIND_III.pdf\#nameddest=tag-01HX}{III.? (p.~?)} }
\expandafter\def\csname SRef@courseIII:adams-tower:5\endcsname{\href{S6_SERIES_GRIND_III.pdf\#nameddest=tag-01HY}{III.? (p.~?)} }
\expandafter\def\csname SRef@courseIII:adams-tower:6\endcsname{\href{S6_SERIES_GRIND_III.pdf\#nameddest=tag-01HZ}{III.? (p.~?)} }
\expandafter\def\csname SRef@courseIII:adams-tower:7.1\endcsname{\href{S6_SERIES_GRIND_III.pdf\#nameddest=tag-01I0}{III.? (p.~?)} }
\expandafter\def\csname SRef@courseIII:adams-tower:7.2\endcsname{\href{S6_SERIES_GRIND_III.pdf\#nameddest=tag-01I1}{III.? (p.~?)} }
\expandafter\def\csname SRef@courseIII:adams-tower:8\endcsname{\href{S6_SERIES_GRIND_III.pdf\#nameddest=tag-01I2}{III.? (p.~?)} }
\expandafter\def\csname SRef@courseIII:adams-tower:9\endcsname{\href{S6_SERIES_GRIND_III.pdf\#nameddest=tag-01I3}{III.? (p.~?)} }
\expandafter\def\csname SRef@courseIII:adams-tower:10\endcsname{\href{S6_SERIES_GRIND_III.pdf\#nameddest=tag-01HT}{III.? (p.~?)} }
\expandafter\def\csname SRef@courseIII:adams-tower:11\endcsname{\href{S6_SERIES_GRIND_III.pdf\#nameddest=tag-01HU}{III.? (p.~?)} }
\expandafter\def\csname SRef@courseIII:formal-char2:1\endcsname{\href{S6_SERIES_GRIND_III.pdf\#nameddest=tag-01J2}{III.? (p.~?)} }
\expandafter\def\csname SRef@courseIII:formal-char2:2\endcsname{\href{S6_SERIES_GRIND_III.pdf\#nameddest=tag-01J3}{III.? (p.~?)} }
\expandafter\def\csname SRef@courseIII:formal-char2:3\endcsname{\href{S6_SERIES_GRIND_III.pdf\#nameddest=tag-01J4}{III.? (p.~?)} }
\expandafter\def\csname SRef@courseIII:formal-char2:4\endcsname{\href{S6_SERIES_GRIND_III.pdf\#nameddest=tag-01J5}{III.? (p.~?)} }
\expandafter\def\csname SRef@courseIII:formal-char2:4.1\endcsname{\href{S6_SERIES_GRIND_III.pdf\#nameddest=tag-01J6}{III.? (p.~?)} }
\expandafter\def\csname SRef@courseIII:formal-char2:4.2\endcsname{\href{S6_SERIES_GRIND_III.pdf\#nameddest=tag-01J7}{III.? (p.~?)} }
\expandafter\def\csname SRef@courseIII:projective-coop:1\endcsname{\href{S6_SERIES_GRIND_III.pdf\#nameddest=tag-01M0}{III.? (p.~?)} }
\expandafter\def\csname SRef@courseIII:projective-coop:2\endcsname{\href{S6_SERIES_GRIND_III.pdf\#nameddest=tag-01M1}{III.? (p.~?)} }
\expandafter\def\csname SRef@courseIII:projective-coop:3\endcsname{\href{S6_SERIES_GRIND_III.pdf\#nameddest=tag-01M2}{III.? (p.~?)} }
\expandafter\def\csname SRef@courseIII:projective-coop:4\endcsname{\href{S6_SERIES_GRIND_III.pdf\#nameddest=tag-01M3}{III.? (p.~?)} }
\expandafter\def\csname SRef@courseIII:projective-coop:5\endcsname{\href{S6_SERIES_GRIND_III.pdf\#nameddest=tag-01M4}{III.? (p.~?)} }
\expandafter\def\csname SRef@courseIII:projective-coop:6.1\endcsname{\href{S6_SERIES_GRIND_III.pdf\#nameddest=tag-01M5}{III.? (p.~?)} }
\expandafter\def\csname SRef@courseIII:projective-coop:6.2\endcsname{\href{S6_SERIES_GRIND_III.pdf\#nameddest=tag-01M6}{III.? (p.~?)} }
\expandafter\def\csname SRef@courseIII:actual-degree8:1\endcsname{\href{S6_SERIES_GRIND_III.pdf\#nameddest=tag-01H0}{III.? (p.~?)} }
\expandafter\def\csname SRef@courseIII:actual-degree8:2\endcsname{\href{S6_SERIES_GRIND_III.pdf\#nameddest=tag-01H1}{III.? (p.~?)} }
\expandafter\def\csname SRef@courseIII:actual-degree8:3\endcsname{\href{S6_SERIES_GRIND_III.pdf\#nameddest=tag-01H2}{III.? (p.~?)} }
\expandafter\def\csname SRef@courseIII:actual-degree8:4\endcsname{\href{S6_SERIES_GRIND_III.pdf\#nameddest=tag-01H3}{III.? (p.~?)} }
\expandafter\def\csname SRef@courseIII:actual-degree8:5\endcsname{\href{S6_SERIES_GRIND_III.pdf\#nameddest=tag-01H4}{III.? (p.~?)} }
\expandafter\def\csname SRef@courseIII:actual-degree8:6\endcsname{\href{S6_SERIES_GRIND_III.pdf\#nameddest=tag-01H5}{III.? (p.~?)} }
\expandafter\def\csname SRef@courseIII:actual-degree8:7\endcsname{\href{S6_SERIES_GRIND_III.pdf\#nameddest=tag-01H6}{III.? (p.~?)} }
\expandafter\def\csname SRef@courseIII:actual-degree8:8\endcsname{\href{S6_SERIES_GRIND_III.pdf\#nameddest=tag-01H7}{III.? (p.~?)} }
\expandafter\def\csname SRef@courseIII:actual-degree8:9.1\endcsname{\href{S6_SERIES_GRIND_III.pdf\#nameddest=tag-01H8}{III.? (p.~?)} }
\expandafter\def\csname SRef@courseIII:actual-degree8:9.2\endcsname{\href{S6_SERIES_GRIND_III.pdf\#nameddest=tag-01H9}{III.? (p.~?)} }
\expandafter\def\csname SRef@courseIII:pt-relative-embedding\endcsname{\href{S6_SERIES_GRIND_III.pdf\#nameddest=tag-01ME}{III.? (p.~?)} }
\expandafter\def\csname SRef@courseIII:pt-relative-normal-frame\endcsname{\href{S6_SERIES_GRIND_III.pdf\#nameddest=tag-01MF}{III.? (p.~?)} }
\expandafter\def\csname SRef@courseIII:pt-normal-inverse-comparison\endcsname{\href{S6_SERIES_GRIND_III.pdf\#nameddest=tag-01MD}{III.? (p.~?)} }
\expandafter\def\csname SRef@courseIII:pt-abstract-framed-comparison\endcsname{\href{S6_SERIES_GRIND_III.pdf\#nameddest=tag-01MC}{III.? (p.~?)} }
\expandafter\def\csname SRef@courseIII:actual-resolution:1\endcsname{\href{S6_SERIES_GRIND_III.pdf\#nameddest=tag-01HA}{III.? (p.~?)} }
\expandafter\def\csname SRef@courseIII:actual-resolution:2\endcsname{\href{S6_SERIES_GRIND_III.pdf\#nameddest=tag-01HB}{III.? (p.~?)} }
\expandafter\def\csname SRef@courseIII:actual-resolution:3\endcsname{\href{S6_SERIES_GRIND_III.pdf\#nameddest=tag-01HC}{III.? (p.~?)} }
\expandafter\def\csname SRef@courseIII:actual-resolution:4\endcsname{\href{S6_SERIES_GRIND_III.pdf\#nameddest=tag-01HD}{III.? (p.~?)} }
\expandafter\def\csname SRef@courseIII:actual-resolution:5\endcsname{\href{S6_SERIES_GRIND_III.pdf\#nameddest=tag-01HE}{III.? (p.~?)} }
\expandafter\def\csname SRef@courseIII:actual-resolution:6\endcsname{\href{S6_SERIES_GRIND_III.pdf\#nameddest=tag-01HF}{III.? (p.~?)} }
\expandafter\def\csname SRef@courseIII:actual-resolution:7\endcsname{\href{S6_SERIES_GRIND_III.pdf\#nameddest=tag-01HG}{III.? (p.~?)} }
\expandafter\def\csname SRef@courseIII:actual-resolution:8.1\endcsname{\href{S6_SERIES_GRIND_III.pdf\#nameddest=tag-01HH}{III.? (p.~?)} }
\expandafter\def\csname SRef@courseIII:actual-resolution:8.2\endcsname{\href{S6_SERIES_GRIND_III.pdf\#nameddest=tag-01HI}{III.? (p.~?)} }
\expandafter\def\csname SRef@courseIII:milnor-cobar:1.1\endcsname{\href{S6_SERIES_GRIND_III.pdf\#nameddest=tag-01L1}{III.? (p.~?)} }
\expandafter\def\csname SRef@courseIII:milnor-cobar:6.1\endcsname{\href{S6_SERIES_GRIND_III.pdf\#nameddest=tag-01L2}{III.? (p.~?)} }
\expandafter\def\csname SRef@courseIII:milnor-cobar:7.1\endcsname{\href{S6_SERIES_GRIND_III.pdf\#nameddest=tag-01L3}{III.? (p.~?)} }
\expandafter\def\csname SRef@courseIII:milnor-cobar:7.2\endcsname{\href{S6_SERIES_GRIND_III.pdf\#nameddest=tag-01L4}{III.? (p.~?)} }
\expandafter\def\csname SRef@courseIII:odd-coop:1\endcsname{\href{S6_SERIES_GRIND_III.pdf\#nameddest=tag-01LA}{III.? (p.~?)} }
\expandafter\def\csname SRef@courseIII:odd-coop:2\endcsname{\href{S6_SERIES_GRIND_III.pdf\#nameddest=tag-01LB}{III.? (p.~?)} }
\expandafter\def\csname SRef@courseIII:odd-coop:3\endcsname{\href{S6_SERIES_GRIND_III.pdf\#nameddest=tag-01LC}{III.? (p.~?)} }
\expandafter\def\csname SRef@courseIII:odd-coop:4\endcsname{\href{S6_SERIES_GRIND_III.pdf\#nameddest=tag-01LD}{III.? (p.~?)} }
\expandafter\def\csname SRef@courseIII:odd-coop:5\endcsname{\href{S6_SERIES_GRIND_III.pdf\#nameddest=tag-01LE}{III.? (p.~?)} }
\expandafter\def\csname SRef@courseIII:odd-coop:6\endcsname{\href{S6_SERIES_GRIND_III.pdf\#nameddest=tag-01LF}{III.? (p.~?)} }
\expandafter\def\csname SRef@courseIII:odd-coop:7.1\endcsname{\href{S6_SERIES_GRIND_III.pdf\#nameddest=tag-01LG}{III.? (p.~?)} }
\expandafter\def\csname SRef@courseIII:odd-coop:7.2\endcsname{\href{S6_SERIES_GRIND_III.pdf\#nameddest=tag-01LH}{III.? (p.~?)} }
\expandafter\def\csname SRef@courseIII:coefficient-rings:1\endcsname{\href{S6_SERIES_GRIND_III.pdf\#nameddest=tag-01IM}{III.? (p.~?)} }
\expandafter\def\csname SRef@courseIII:coefficient-rings:2\endcsname{\href{S6_SERIES_GRIND_III.pdf\#nameddest=tag-01IN}{III.? (p.~?)} }
\expandafter\def\csname SRef@courseIII:coefficient-rings:3\endcsname{\href{S6_SERIES_GRIND_III.pdf\#nameddest=tag-01IO}{III.? (p.~?)} }
\expandafter\def\csname SRef@courseIII:coefficient-rings:4\endcsname{\href{S6_SERIES_GRIND_III.pdf\#nameddest=tag-01IP}{III.? (p.~?)} }
\expandafter\def\csname SRef@courseIII:coefficient-rings:5\endcsname{\href{S6_SERIES_GRIND_III.pdf\#nameddest=tag-01IQ}{III.? (p.~?)} }
\expandafter\def\csname SRef@courseIII:coefficient-rings:6\endcsname{\href{S6_SERIES_GRIND_III.pdf\#nameddest=tag-01IR}{III.? (p.~?)} }
\expandafter\def\csname SRef@courseIII:coefficient-rings:7.1\endcsname{\href{S6_SERIES_GRIND_III.pdf\#nameddest=tag-01IS}{III.? (p.~?)} }
\expandafter\def\csname SRef@courseIII:coefficient-rings:7.2\endcsname{\href{S6_SERIES_GRIND_III.pdf\#nameddest=tag-01IT}{III.? (p.~?)} }
\expandafter\def\csname SRef@courseIII:duality:2.1\endcsname{\href{S6_SERIES_GRIND_III.pdf\#nameddest=tag-01IU}{III.? (p.~?)} }
\expandafter\def\csname SRef@courseIII:duality:2.2\endcsname{\href{S6_SERIES_GRIND_III.pdf\#nameddest=tag-01IV}{III.? (p.~?)} }
\expandafter\def\csname SRef@courseIII:duality:2.3\endcsname{\href{S6_SERIES_GRIND_III.pdf\#nameddest=tag-01IW}{III.? (p.~?)} }
\expandafter\def\csname SRef@courseIII:duality:3.1\endcsname{\href{S6_SERIES_GRIND_III.pdf\#nameddest=tag-01IX}{III.? (p.~?)} }
\expandafter\def\csname SRef@courseIII:duality:3.2\endcsname{\href{S6_SERIES_GRIND_III.pdf\#nameddest=tag-01IY}{III.? (p.~?)} }
\expandafter\def\csname SRef@courseIII:duality:3.3\endcsname{\href{S6_SERIES_GRIND_III.pdf\#nameddest=tag-01IZ}{III.? (p.~?)} }
\expandafter\def\csname SRef@courseIII:duality:4.1\endcsname{\href{S6_SERIES_GRIND_III.pdf\#nameddest=tag-01J0}{III.? (p.~?)} }
\expandafter\def\csname SRef@courseIII:duality:4.2\endcsname{\href{S6_SERIES_GRIND_III.pdf\#nameddest=tag-01J1}{III.? (p.~?)} }
\expandafter\def\csname SRef@night:P1.6\endcsname{\href{S6_SERIES_GRIND_V.pdf\#nameddest=tag-02QX}{V.? (p.~?)} }
\expandafter\def\csname SRef@night:P1.14\endcsname{\href{S6_SERIES_GRIND_V.pdf\#nameddest=tag-02QN}{V.? (p.~?)} }
\expandafter\def\csname SRef@night:P5.7\endcsname{\hyperlink{tag-02YS}{I.\ref*{night:P5.7}}}
\expandafter\def\csname SRef@ari:A.1.1\endcsname{\hyperlink{tag-03J1}{I.\ref*{ari:A.1.1}}}
\expandafter\def\csname SRef@ari:A.1.2\endcsname{\hyperlink{tag-03J2}{I.\ref*{ari:A.1.2}}}
\expandafter\def\csname SRef@ari:A.1.3\endcsname{\hyperlink{tag-03J3}{I.\ref*{ari:A.1.3}}}
\expandafter\def\csname SRef@ari:A.1.4\endcsname{\hyperlink{tag-03J4}{I.\ref*{ari:A.1.4}}}
\expandafter\def\csname SRef@ari:A.2.1\endcsname{\hyperlink{tag-03J5}{I.\ref*{ari:A.2.1}}}
\expandafter\def\csname SRef@ari:A.2.2\endcsname{\hyperlink{tag-03J6}{I.\ref*{ari:A.2.2}}}
\expandafter\def\csname SRef@ari:A.2.3\endcsname{\hyperlink{tag-03J7}{I.\ref*{ari:A.2.3}}}
\expandafter\def\csname SRef@ari:A.2.4\endcsname{\hyperlink{tag-03J8}{I.\ref*{ari:A.2.4}}}
\expandafter\def\csname SRef@ari:A.2.5\endcsname{\hyperlink{tag-03J9}{I.\ref*{ari:A.2.5}}}
\expandafter\def\csname SRef@ari:A.3.1\endcsname{\hyperlink{tag-03JA}{I.\ref*{ari:A.3.1}}}
\expandafter\def\csname SRef@ari:A.3.2\endcsname{\hyperlink{tag-03JB}{I.\ref*{ari:A.3.2}}}
\expandafter\def\csname SRef@ari:A.4.1\endcsname{\hyperlink{tag-03JC}{I.\ref*{ari:A.4.1}}}
\expandafter\def\csname SRef@ari:A.4.2\endcsname{\hyperlink{tag-03JD}{I.\ref*{ari:A.4.2}}}
\expandafter\def\csname SRef@ari:A.4.3\endcsname{\hyperlink{tag-03JE}{I.\ref*{ari:A.4.3}}}
\expandafter\def\csname SRef@ari:A.4.4\endcsname{\hyperlink{tag-03JF}{I.\ref*{ari:A.4.4}}}
\expandafter\def\csname SRef@ari:A.4.5\endcsname{\hyperlink{tag-03JG}{I.\ref*{ari:A.4.5}}}
\expandafter\def\csname SRef@ari:A.4.6\endcsname{\hyperlink{tag-03JH}{I.\ref*{ari:A.4.6}}}
\expandafter\def\csname SRef@ari:A.4.7\endcsname{\hyperlink{tag-03JI}{I.\ref*{ari:A.4.7}}}
\expandafter\def\csname SRef@ari:A.4.8\endcsname{\hyperlink{tag-03JJ}{I.\ref*{ari:A.4.8}}}
\expandafter\def\csname SRef@ari:A.5.1\endcsname{\hyperlink{tag-03JK}{I.\ref*{ari:A.5.1}}}
\expandafter\def\csname SRef@ari:A.5.2\endcsname{\hyperlink{tag-03JL}{I.\ref*{ari:A.5.2}}}
\expandafter\def\csname SRef@ari:A.5.3\endcsname{\hyperlink{tag-03JM}{I.\ref*{ari:A.5.3}}}
\expandafter\def\csname SRef@ari:A.5.4\endcsname{\hyperlink{tag-03JN}{I.\ref*{ari:A.5.4}}}
\expandafter\def\csname SRef@ari:A.5.5\endcsname{\hyperlink{tag-03JO}{I.\ref*{ari:A.5.5}}}
\expandafter\def\csname SRef@ari:A.5.6\endcsname{\hyperlink{tag-03JP}{I.\ref*{ari:A.5.6}}}
\expandafter\def\csname SRef@ari:A.5.8\endcsname{\hyperlink{tag-03JR}{I.\ref*{ari:A.5.8}}}
\expandafter\def\csname SRef@ari:A.5.7\endcsname{\hyperlink{tag-03JQ}{I.\ref*{ari:A.5.7}}}
\expandafter\def\csname SRef@series:I:main-theorems\endcsname{\hyperlink{tag-03JS}{I.\ref*{series:I:main-theorems}}}
\expandafter\def\csname SRef@series:I:A\endcsname{\hyperlink{tag-03JT}{I.\ref*{series:I:A}}}
\expandafter\def\csname SRef@series:I:B\endcsname{\hyperlink{tag-03JU}{I.\ref*{series:I:B}}}
\expandafter\def\csname SRef@series:I:C\endcsname{\hyperlink{tag-03JV}{I.\ref*{series:I:C}}}
\expandafter\def\csname SRef@series:I:D\endcsname{\hyperlink{tag-03JW}{I.\ref*{series:I:D}}}
\expandafter\def\csname SRef@series:I:E\endcsname{\hyperlink{tag-03JX}{I.\ref*{series:I:E}}}
\expandafter\def\csname SRef@series:I:F\endcsname{\hyperlink{tag-03JY}{I.\ref*{series:I:F}}}
\expandafter\def\csname SRef@series:I:G\endcsname{\hyperlink{tag-03JZ}{I.\ref*{series:I:G}}}
\expandafter\def\csname SRef@series:I:outlines\endcsname{\hyperlink{tag-03K0}{I.\ref*{series:I:outlines}}}
\expandafter\def\csname SRef@series:II:main-theorems\endcsname{\href{S6_SERIES_GRIND_II.pdf\#nameddest=tag-03K1}{II.1.2 (p.~5)} }
\expandafter\def\csname SRef@series:II:A\endcsname{\href{S6_SERIES_GRIND_II.pdf\#nameddest=tag-03K2}{II.A (p.~5)} }
\expandafter\def\csname SRef@series:II:B\endcsname{\href{S6_SERIES_GRIND_II.pdf\#nameddest=tag-03K3}{II.B (p.~6)} }
\expandafter\def\csname SRef@series:II:C\endcsname{\href{S6_SERIES_GRIND_II.pdf\#nameddest=tag-03K4}{II.C (p.~6)} }
\expandafter\def\csname SRef@series:II:D\endcsname{\href{S6_SERIES_GRIND_II.pdf\#nameddest=tag-03K5}{II.D (p.~7)} }
\expandafter\def\csname SRef@series:II:E\endcsname{\href{S6_SERIES_GRIND_II.pdf\#nameddest=tag-03K6}{II.E (p.~7)} }
\expandafter\def\csname SRef@series:II:F\endcsname{\href{S6_SERIES_GRIND_II.pdf\#nameddest=tag-03K7}{II.F (p.~8)} }
\expandafter\def\csname SRef@series:II:G\endcsname{\href{S6_SERIES_GRIND_II.pdf\#nameddest=tag-03K8}{II.G (p.~9)} }
\expandafter\def\csname SRef@series:II:outline\endcsname{\href{S6_SERIES_GRIND_II.pdf\#nameddest=tag-03K9}{II.1.3 (p.~13)} }
\expandafter\def\csname SRef@series:III:main-theorems\endcsname{\href{S6_SERIES_GRIND_III.pdf\#nameddest=tag-03KA}{III.1.2 (p.~3)} }
\expandafter\def\csname SRef@series:III:A\endcsname{\href{S6_SERIES_GRIND_III.pdf\#nameddest=tag-03KB}{III.A (p.~3)} }
\expandafter\def\csname SRef@series:III:B\endcsname{\href{S6_SERIES_GRIND_III.pdf\#nameddest=tag-03KC}{III.B (p.~3)} }
\expandafter\def\csname SRef@series:III:C\endcsname{\href{S6_SERIES_GRIND_III.pdf\#nameddest=tag-03KD}{III.C (p.~4)} }
\expandafter\def\csname SRef@series:III:D\endcsname{\href{S6_SERIES_GRIND_III.pdf\#nameddest=tag-03KE}{III.D (p.~4)} }
\expandafter\def\csname SRef@series:III:E\endcsname{\href{S6_SERIES_GRIND_III.pdf\#nameddest=tag-03KF}{III.E (p.~4)} }
\expandafter\def\csname SRef@series:III:F\endcsname{\href{S6_SERIES_GRIND_III.pdf\#nameddest=tag-03KG}{III.F (p.~5)} }
\expandafter\def\csname SRef@series:III:G\endcsname{\href{S6_SERIES_GRIND_III.pdf\#nameddest=tag-03KH}{III.G (p.~5)} }
\expandafter\def\csname SRef@series:IV:main-theorems\endcsname{\href{S6_SERIES_GRIND_IV.pdf\#nameddest=tag-03KI}{IV.1.2 (p.~3)} }
\expandafter\def\csname SRef@series:IV:A\endcsname{\href{S6_SERIES_GRIND_IV.pdf\#nameddest=tag-03KJ}{IV.A (p.~3)} }
\expandafter\def\csname SRef@series:IV:B\endcsname{\href{S6_SERIES_GRIND_IV.pdf\#nameddest=tag-03KK}{IV.B (p.~4)} }
\expandafter\def\csname SRef@series:IV:C\endcsname{\href{S6_SERIES_GRIND_IV.pdf\#nameddest=tag-03KL}{IV.C (p.~5)} }
\expandafter\def\csname SRef@series:IV:D\endcsname{\href{S6_SERIES_GRIND_IV.pdf\#nameddest=tag-03KM}{IV.D (p.~6)} }
\expandafter\def\csname SRef@series:IV:E\endcsname{\href{S6_SERIES_GRIND_IV.pdf\#nameddest=tag-03KN}{IV.E (p.~7)} }
\expandafter\def\csname SRef@series:IV:F\endcsname{\href{S6_SERIES_GRIND_IV.pdf\#nameddest=tag-03KO}{IV.F (p.~8)} }
\expandafter\def\csname SRef@series:IV:G\endcsname{\href{S6_SERIES_GRIND_IV.pdf\#nameddest=tag-03KP}{IV.G (p.~8)} }
\expandafter\def\csname SRef@series:V:main-theorems\endcsname{\href{S6_SERIES_GRIND_V.pdf\#nameddest=tag-03KQ}{V.1.2 (p.~3)} }
\expandafter\def\csname SRef@series:V:A\endcsname{\href{S6_SERIES_GRIND_V.pdf\#nameddest=tag-03KR}{V.A (p.~3)} }
\expandafter\def\csname SRef@series:V:B\endcsname{\href{S6_SERIES_GRIND_V.pdf\#nameddest=tag-03KS}{V.B (p.~3)} }
\expandafter\def\csname SRef@series:V:C\endcsname{\href{S6_SERIES_GRIND_V.pdf\#nameddest=tag-03KT}{V.C (p.~4)} }
\expandafter\def\csname SRef@series:V:D\endcsname{\href{S6_SERIES_GRIND_V.pdf\#nameddest=tag-03KU}{V.D (p.~4)} }
\expandafter\def\csname SRef@series:V:E\endcsname{\href{S6_SERIES_GRIND_V.pdf\#nameddest=tag-03KV}{V.E (p.~5)} }
\expandafter\def\csname SRef@series:V:F\endcsname{\href{S6_SERIES_GRIND_V.pdf\#nameddest=tag-03KW}{V.F (p.~5)} }
\expandafter\def\csname SRef@layout:V:fresh-gauges\endcsname{\href{S6_SERIES_GRIND_V.pdf\#nameddest=tag-03KX}{V.{{$\langle 2\rangle $1.}} (p.~126)} }
\expandafter\def\csname SRef@layout:V:fresh-minimum\endcsname{\href{S6_SERIES_GRIND_V.pdf\#nameddest=tag-03KY}{V.{{$\langle 1\rangle $4.}} (p.~130)} }
\expandafter\def\csname SRef@layout:V:fresh-bounds\endcsname{\href{S6_SERIES_GRIND_V.pdf\#nameddest=tag-03KZ}{V.{1.18.4} (p.~130)} }
\expandafter\def\csname SRef@layout:V:fresh-marking\endcsname{\href{S6_SERIES_GRIND_V.pdf\#nameddest=tag-03L0}{V.{{$\langle 1\rangle $6.}} (p.~132)} }
\expandafter\def\csname SRef@layout:V:fresh-splitting\endcsname{\href{S6_SERIES_GRIND_V.pdf\#nameddest=tag-03L1}{V.{{$\langle 1\rangle $7.}} (p.~133)} }
\expandafter\def\csname SRef@iv:source_transports\endcsname{\href{S6_SERIES_GRIND_IV.pdf\#nameddest=tag-03L2}{IV.3.27 (p.~150)} }
\expandafter\def\csname SRef@inv:section:coefficient-atlas\endcsname{\hyperlink{tag-03L3}{I.\ref*{inv:section:coefficient-atlas}}}
\expandafter\def\csname SRef@inv:eq:atlas-theta-factor\endcsname{\hyperlink{tag-03L4}{I.\ref*{inv:eq:atlas-theta-factor}}}
\expandafter\def\csname SRef@inv:eq:atlas-semiflat\endcsname{\hyperlink{tag-03L5}{I.\ref*{inv:eq:atlas-semiflat}}}
\expandafter\def\csname SRef@inv:eq:atlas-affine-lift\endcsname{\hyperlink{tag-03L6}{I.\ref*{inv:eq:atlas-affine-lift}}}
\expandafter\def\csname SRef@inv:eq:atlas-cusp-factor\endcsname{\hyperlink{tag-03L7}{I.\ref*{inv:eq:atlas-cusp-factor}}}
\expandafter\def\csname SRef@inv:eq:atlas-cusp-unit\endcsname{\hyperlink{tag-03L8}{I.\ref*{inv:eq:atlas-cusp-unit}}}
\expandafter\def\csname SRef@inv:eq:atlas-connection-inverse\endcsname{\hyperlink{tag-03L9}{I.\ref*{inv:eq:atlas-connection-inverse}}}
\expandafter\def\csname SRef@inv:eq:atlas-cusp-connection\endcsname{\hyperlink{tag-03LA}{I.\ref*{inv:eq:atlas-cusp-connection}}}
\expandafter\def\csname SRef@inv:eq:atlas-finite-connection\endcsname{\hyperlink{tag-03LB}{I.\ref*{inv:eq:atlas-finite-connection}}}
\expandafter\def\csname SRef@inv:eq:atlas-averaged-connection\endcsname{\hyperlink{tag-03LC}{I.\ref*{inv:eq:atlas-averaged-connection}}}
\expandafter\def\csname SRef@inv:eq:atlas-cocycle\endcsname{\hyperlink{tag-03LD}{I.\ref*{inv:eq:atlas-cocycle}}}
\expandafter\def\csname SRef@inv:eq:atlas-weights\endcsname{\hyperlink{tag-03LE}{I.\ref*{inv:eq:atlas-weights}}}
\expandafter\def\csname SRef@inv:eq:atlas-weight-derivative\endcsname{\hyperlink{tag-03LF}{I.\ref*{inv:eq:atlas-weight-derivative}}}
\expandafter\def\csname SRef@inv:eq:atlas-final-forms\endcsname{\hyperlink{tag-03LG}{I.\ref*{inv:eq:atlas-final-forms}}}
\expandafter\def\csname SRef@inv:eq:atlas-final-coefficients\endcsname{\hyperlink{tag-03LH}{I.\ref*{inv:eq:atlas-final-coefficients}}}
\expandafter\def\csname SRef@inv:C3.8:operator-comparison\endcsname{\hyperlink{tag-03LI}{I.\ref*{inv:C3.8:operator-comparison}}}
\expandafter\def\csname SRef@iv:open:K111:original_auxiliary_arrays\endcsname{\href{S6_SERIES_GRIND_IV.pdf\#nameddest=tag-03LJ}{IV.12.1 (p.~1396)} }
\expandafter\def\csname SRef@maps:proof:D.OS.refined\endcsname{\href{S6_SERIES_GRIND_III.pdf\#nameddest=tag-03LK}{III.5.2 (p.~207)} }
\expandafter\def\csname SRef@maps:eq:D.OR1\endcsname{\href{S6_SERIES_GRIND_III.pdf\#nameddest=tag-03LL}{III.{D.OR1} (p.~207)} }
\expandafter\def\csname SRef@maps:eq:D.OR2\endcsname{\href{S6_SERIES_GRIND_III.pdf\#nameddest=tag-03LM}{III.{D.OR2} (p.~207)} }
\expandafter\def\csname SRef@maps:eq:D.OR3\endcsname{\href{S6_SERIES_GRIND_III.pdf\#nameddest=tag-03LN}{III.{D.OR3} (p.~207)} }
\expandafter\def\csname SRef@maps:eq:D.OR4\endcsname{\href{S6_SERIES_GRIND_III.pdf\#nameddest=tag-03LO}{III.{D.OR4} (p.~207)} }
\expandafter\def\csname SRef@maps:eq:D.OR5\endcsname{\href{S6_SERIES_GRIND_III.pdf\#nameddest=tag-03LP}{III.{D.OR5} (p.~208)} }
\expandafter\def\csname SRef@maps:eq:D.OR6\endcsname{\href{S6_SERIES_GRIND_III.pdf\#nameddest=tag-03LQ}{III.{D.OR6} (p.~208)} }
\expandafter\def\csname SRef@maps:eq:D.OR7\endcsname{\href{S6_SERIES_GRIND_III.pdf\#nameddest=tag-03LR}{III.{D.OR7} (p.~208)} }
\expandafter\def\csname SRef@maps:eq:D.OR8\endcsname{\href{S6_SERIES_GRIND_III.pdf\#nameddest=tag-03LS}{III.{D.OR8} (p.~208)} }
\expandafter\def\csname SRef@maps:eq:D.OR9\endcsname{\href{S6_SERIES_GRIND_III.pdf\#nameddest=tag-03LT}{III.{D.OR9} (p.~208)} }
\expandafter\def\csname SRef@maps:eq:D.OR10\endcsname{\href{S6_SERIES_GRIND_III.pdf\#nameddest=tag-03LU}{III.{D.OR10} (p.~209)} }
\expandafter\def\csname SRef@maps:eq:D.OR11\endcsname{\href{S6_SERIES_GRIND_III.pdf\#nameddest=tag-03LV}{III.{D.OR11} (p.~209)} }
\expandafter\def\csname SRef@maps:eq:D.OR12\endcsname{\href{S6_SERIES_GRIND_III.pdf\#nameddest=tag-03LW}{III.{D.OR12} (p.~210)} }
\expandafter\def\csname SRef@maps:eq:D.OR13\endcsname{\href{S6_SERIES_GRIND_III.pdf\#nameddest=tag-03LX}{III.{D.OR13} (p.~210)} }
\expandafter\def\csname SRef@maps:eq:D.OR14\endcsname{\href{S6_SERIES_GRIND_III.pdf\#nameddest=tag-03LY}{III.{D.OR14} (p.~211)} }
\expandafter\def\csname SRef@inv:E2.5\endcsname{\href{S6_SERIES_GRIND_II.pdf\#nameddest=tag-03LZ}{II.14.2 (p.~356)} }
\expandafter\def\csname SRef@invII:eq:E12-node\endcsname{\href{S6_SERIES_GRIND_II.pdf\#nameddest=tag-03M0}{II.{E.12.1} (p.~358)} }
\expandafter\def\csname SRef@invII:eq:E12-jets\endcsname{\href{S6_SERIES_GRIND_II.pdf\#nameddest=tag-03M1}{II.{E.12.2} (p.~359)} }
\expandafter\def\csname SRef@invII:eq:E12-open\endcsname{\href{S6_SERIES_GRIND_II.pdf\#nameddest=tag-03M2}{II.{E.12.3} (p.~360)} }
\expandafter\def\csname SBack@I:weierstrass-comparison\endcsname{\hyperlink{step-I-A-1}{proves step 1 of Theorem A}}
\expandafter\def\csname SBack@fam:src:B.1.1\endcsname{\hyperlink{step-I-A-1}{proves step 1 of Theorem A}}
\expandafter\def\csname SBack@fam:src:C.1.1\endcsname{\hyperlink{step-I-A-2}{proves step 2 of Theorem A}}
\expandafter\def\csname SBack@ari:C.4.2\endcsname{\hyperlink{step-I-A-3}{proves step 3 of Theorem A}}
\expandafter\def\csname SBack@ari:C.4.3\endcsname{\hyperlink{step-I-A-3}{proves step 3 of Theorem A}}
\expandafter\def\csname SBack@ari:C.4.4\endcsname{\hyperlink{step-I-A-3}{proves step 3 of Theorem A}}
\expandafter\def\csname SBack@ari:C.8.1\endcsname{\hyperlink{step-I-A-4}{proves step 4 of Theorem A}}
\expandafter\def\csname SBack@ari:C.8.2\endcsname{\hyperlink{step-I-A-4}{proves step 4 of Theorem A}}
\expandafter\def\csname SBack@ari:C.8.3\endcsname{\hyperlink{step-I-A-4}{proves step 4 of Theorem A}}
\expandafter\def\csname SBack@ari:C.8.4\endcsname{\hyperlink{step-I-A-4}{proves step 4 of Theorem A}}
\expandafter\def\csname SBack@fam:src:C.2.3\endcsname{\hyperlink{step-I-A-4}{proves step 4 of Theorem A}}
\expandafter\def\csname SBack@ari:C.8.15\endcsname{\hyperlink{step-I-A-5}{proves step 5 of Theorem A}}
\expandafter\def\csname SBack@push:src:D.4.8\endcsname{\hyperlink{step-I-A-5}{proves step 5 of Theorem A}}
\expandafter\def\csname SBack@ari:C.6.1\endcsname{\hyperlink{step-I-A-6}{proves step 6 of Theorem A}}
\expandafter\def\csname SBack@ari:C.6.3\endcsname{\hyperlink{step-I-A-6}{proves step 6 of Theorem A}}
\expandafter\def\csname SBack@fam:src:A.1.1\endcsname{\hyperlink{step-I-B-1}{proves step 1 of Theorem B}}
\expandafter\def\csname SBack@fam:src:A.5.1\endcsname{\hyperlink{step-I-B-1}{proves step 1 of Theorem B}}
\expandafter\def\csname SBack@fam:src:B.2.1\endcsname{\hyperlink{step-I-B-2}{proves step 2 of Theorem B}}
\expandafter\def\csname SBack@fam:src:B.2.2\endcsname{\hyperlink{step-I-B-2}{proves step 2 of Theorem B}}
\expandafter\def\csname SBack@fam:src:B.2.3\endcsname{\hyperlink{step-I-B-2}{proves step 2 of Theorem B}}
\expandafter\def\csname SBack@fam:src:B.3.1\endcsname{\hyperlink{step-I-B-3}{proves step 3 of Theorem B}}
\expandafter\def\csname SBack@push:src:E.1.2\endcsname{\hyperlink{step-I-B-3}{proves step 3 of Theorem B}; \hyperlink{step-I-E-5}{proves step 5 of Theorem E}}
\expandafter\def\csname SBack@ari:C.2.1\endcsname{\hyperlink{step-I-B-4}{proves step 4 of Theorem B}}
\expandafter\def\csname SBack@fam:src:B.4.2\endcsname{\hyperlink{step-I-B-5}{proves step 5 of Theorem B}}
\expandafter\def\csname SBack@mapsI:src:A.3.1\endcsname{\hyperlink{step-I-B-5}{proves step 5 of Theorem B}}
\expandafter\def\csname SBack@mapsI:src:A.3.2\endcsname{\hyperlink{step-I-B-6}{proves step 6 of Theorem B}}
\expandafter\def\csname SBack@mapsI:src:A.3.3\endcsname{\hyperlink{step-I-B-6}{proves step 6 of Theorem B}}
\expandafter\def\csname SBack@fam:src:F.1.1\endcsname{\hyperlink{step-I-C-1}{proves step 1 of Theorem C}}
\expandafter\def\csname SBack@I:monodromy-fundamental-set\endcsname{\hyperlink{step-I-C-2}{proves step 2 of Theorem C}}
\expandafter\def\csname SBack@fam:src:F.2.2\endcsname{\hyperlink{step-I-C-3}{proves step 3 of Theorem C}}
\expandafter\def\csname SBack@fam:src:F.2.3\endcsname{\hyperlink{step-I-C-3}{proves step 3 of Theorem C}}
\expandafter\def\csname SBack@fam:src:F.3.1\endcsname{\hyperlink{step-I-C-4}{proves step 4 of Theorem C}}
\expandafter\def\csname SBack@fam:src:D.1.4\endcsname{\hyperlink{step-I-C-5}{proves step 5 of Theorem C}}
\expandafter\def\csname SBack@fam:src:D.1.6\endcsname{\hyperlink{step-I-C-5}{proves step 5 of Theorem C}}
\expandafter\def\csname SBack@inv:A2.2\endcsname{\hyperlink{step-I-D-1}{proves step 1 of Theorem D}}
\expandafter\def\csname SBack@inv:A2.4\endcsname{\hyperlink{step-I-D-1}{proves step 1 of Theorem D}}
\expandafter\def\csname SBack@inv:A1.0\endcsname{\hyperlink{step-I-D-2}{proves step 2 of Theorem D}}
\expandafter\def\csname SBack@inv:A3.1\endcsname{\hyperlink{step-I-D-2}{proves step 2 of Theorem D}}
\expandafter\def\csname SBack@inv:C1.1\endcsname{\hyperlink{step-I-D-3}{proves step 3 of Theorem D}}
\expandafter\def\csname SBack@inv:C2.1\endcsname{\hyperlink{step-I-D-3}{proves step 3 of Theorem D}}
\expandafter\def\csname SBack@inv:C3.3\endcsname{\hyperlink{step-I-D-3}{proves step 3 of Theorem D}}
\expandafter\def\csname SBack@inv:A4.1\endcsname{\hyperlink{step-I-D-4}{proves step 4 of Theorem D}}
\expandafter\def\csname SBack@inv:A4.2\endcsname{\hyperlink{step-I-D-4}{proves step 4 of Theorem D}}
\expandafter\def\csname SBack@inv:C3.8\endcsname{\hyperlink{step-I-D-5}{proves step 5 of Theorem D}}
\expandafter\def\csname SBack@inv:C3.2\endcsname{\hyperlink{step-I-E-1}{proves step 1 of Theorem E}; \hyperlink{step-I-E-4}{proves step 4 of Theorem E}}
\expandafter\def\csname SBack@inv:D4.1\endcsname{\hyperlink{step-I-E-1}{proves step 1 of Theorem E}}
\expandafter\def\csname SBack@inv:D4.2\endcsname{\hyperlink{step-I-E-2}{proves step 2 of Theorem E}}
\expandafter\def\csname SBack@inv:D4.3\endcsname{\hyperlink{step-I-E-2}{proves step 2 of Theorem E}}
\expandafter\def\csname SBack@inv:D4.4\endcsname{\hyperlink{step-I-E-3}{proves step 3 of Theorem E}}
\expandafter\def\csname SBack@inv:D4.5\endcsname{\hyperlink{step-I-E-3}{proves step 3 of Theorem E}; \hyperlink{step-I-E-4}{proves step 4 of Theorem E}}
\expandafter\def\csname SBack@inv:C3.1\endcsname{\hyperlink{step-I-E-4}{proves step 4 of Theorem E}}
\expandafter\def\csname SBack@push:src:E.2.1\endcsname{\hyperlink{step-I-E-5}{proves step 5 of Theorem E}; \hyperlink{step-I-E-6}{proves step 6 of Theorem E}}
\expandafter\def\csname SBack@fam:st:A.5.2\endcsname{\hyperlink{step-I-E-6}{proves step 6 of Theorem E}}
\expandafter\def\csname SBack@fam:src:B.4.1\endcsname{\hyperlink{step-I-E-6}{proves step 6 of Theorem E}}
\expandafter\def\csname SBack@night:P2.2\endcsname{\hyperlink{step-I-F-1}{proves step 1 of Theorem F}}
\expandafter\def\csname SBack@night:P2.3\endcsname{\hyperlink{step-I-F-1}{proves step 1 of Theorem F}}
\expandafter\def\csname SBack@night:P2.4\endcsname{\hyperlink{step-I-F-2}{proves step 2 of Theorem F}}
\expandafter\def\csname SBack@night:P2.5\endcsname{\hyperlink{step-I-F-3}{proves step 3 of Theorem F}}
\expandafter\def\csname SBack@night:P2.6\endcsname{\hyperlink{step-I-F-3}{proves step 3 of Theorem F}}
\expandafter\def\csname SBack@night:P2.7\endcsname{\hyperlink{step-I-F-4}{proves step 4 of Theorem F}; \hyperlink{step-I-F-5}{proves step 5 of Theorem F}}
\expandafter\def\csname SBack@night:P2.9\endcsname{\hyperlink{step-I-F-6}{proves step 6 of Theorem F}}
\expandafter\def\csname SBack@night:P2.10\endcsname{\hyperlink{step-I-F-7}{proves step 7 of Theorem F}}
\expandafter\def\csname SBack@night:P2.11\endcsname{\hyperlink{step-I-F-7}{proves step 7 of Theorem F}}
\expandafter\def\csname SBack@push:src:D.1.3\endcsname{\hyperlink{step-I-G-1}{proves step 1 of Theorem G}}
\expandafter\def\csname SBack@push:src:D.2.1\endcsname{\hyperlink{step-I-G-2}{proves step 2 of Theorem G}}
\expandafter\def\csname SBack@eq:d-dual-exact\endcsname{\hyperlink{step-I-G-2}{proves step 2 of Theorem G}}
\expandafter\def\csname SBack@push:src:D.3.1\endcsname{\hyperlink{step-I-G-3}{proves step 3 of Theorem G}; \hyperlink{step-I-G-4}{proves step 4 of Theorem G}}
\expandafter\def\csname SBack@push:src:D.2.3\endcsname{\hyperlink{step-I-G-5}{proves step 5 of Theorem G}}
\expandafter\def\csname SBack@mapsI:src:A.1.1\endcsname{\href{S6_SERIES_GRIND_III.pdf\#nameddest=step-III-A-1}{proves step 1 of Theorem III.A}}
\expandafter\def\csname SBack@mapsI:src:A.2.1\endcsname{\href{S6_SERIES_GRIND_III.pdf\#nameddest=step-III-A-2}{proves step 2 of Theorem III.A}}

\begin{document}
\title{The $(3,4,\infty)$ family of complex six-spheres\\[.5em]I. Periods, moduli and the two ends}
\author{AI-produced mathematical text}
\date{21 November 2026}
\maketitle
\begin{abstract}
We determine the periods, moduli and two limiting regimes of the
$(3,4,\infty)$ family of complex six-spheres. A theta expression for
its elliptic period gives the exact parameter disc $0<|w|<216$.
Integral affine and cusp comparison maps classify the discrete data,
parameter identifications and real structures, including the fixed
three-spheres. We compute the integral monodromy group, its maximal
domain of proper discontinuity and the algebraic dimension of every
smooth fibre. Logarithmic direct images and explicit relative
differential modules determine the Dolbeault, Bott--Chern, Aeppli
and tangent cohomology, all Picard twists and the deformation parameter.
Global Dolbeault representatives extend through the multiple fibres
and cusp. At the puncture, the corner monodromy obstructs every
proper Hausdorff analytic filling after any finite parameter cover;
the analogous obstruction holds for the coprime higher-multiplicity
sphere families. At the outer boundary, the intrinsic Albanese
invariant obstructs continuation, while the dual period lattices
determine the local Hausdorff reflection and all additive limits.
The arithmetic norm satisfies exact division and height identities;
we describe its torsion values, critical points and level curves.
\end{abstract}
\clearpage
\tableofcontents
\clearpage
\section{Results and their proofs in outline}\label{sec:results}
\subsection{The series and this paper}
This series studies the $(3,4,\infty)$ family of torus-fibred complex
structures on $S^6$ \citep{AC2026,Engel2026}, its periods, maps and smooth
monodromy, and the higher-dimensional complex manifolds built from it.
The table gives the main result of each paper and the papers it uses.
\begin{center}\small
\begin{tabularx}{\textwidth}{@{}lXl@{}}\toprule
Paper & Main result & Uses\\\midrule
I & Exact periods, moduli, fibre strata and the two ends & ---\\
II & Torus-fibred threefolds with $b_1=b_2=0$: the Campana--Demailly--Peternell argument and the sign theorem for all degenerations & I, IV\\
III & Maps from the six-spheres and the manifolds they give & I, II, IV, V\\
IV & Smooth monodromy in $\Theta_7$ & I, II, III, V\\
V & Complex fourfolds on $3\mathbb{HP}^2\#(-\mathbb{HP}^2)$ and a resolved fibre product & I, III, IV\\
\bottomrule\end{tabularx}\end{center}
\subsection{The main theorems}
\hypertarget{tag-03JS}{}\label{series:I:main-theorems}

We use the Claude--Alp\"oge construction with its specified periodic
cusp fan, clockwise finite root generators and affine translations.
Its base is $B=\Pp^1_t$, with special values $0,1,\infty$, and
$B^\circ=B\setminus\{0,1,\infty\}$.  The ordered homology lattice and
period matrix are
\[
 \Lambda=\Z\langle\widehat\gamma,\widehat u,
                  \widehat w,\widehat\delta\rangle,
 \qquad
 \Pi_c=\begin{pmatrix}6\mu&\tau&1&0\\
                       \beta_{\rm part}+c&\mu&0&1\end{pmatrix}.
\]
Write $d=\widehat\delta$, and use $\gamma,u,w,\delta$ for the dual
integral covectors when they occur in lattice forms.  The parameter
$w=e^{-\pi i c}$ and its square $u=w^2$ are distinguished by their
arguments or by the words ``parameter'' and ``covector.''  The
particular solution is fixed by
$(\beta_{\rm part}+\tau)(\infty)=0$; $\mu$ is the cusp-selected
zero of $\wp_\tau$.  A statement about $t=0$ or $t=1$ uses the
specified smooth finite root model.  Unless further data are
displayed, $X_c$ denotes the standard $(3,4)$ sphere member with
$v_1=\varepsilon=(1,2,-4,0)^{\mathsf t}$,
$v_2=-\varepsilon'=(-1,-3,3,0)^{\mathsf t}$ and the untwisted
cusp collar.  We write $e=(0,1)^{\mathsf t}=\Pi_c d$ for the
invariant vertical vector.  The broader admissible translation and higher
multiplicity data are stated in Theorems B and F.

\hypertarget{main-I-A}{}
\begin{maintheorem}[The period, its exact domain, and its arithmetic norm]
\hypertarget{tag-03JT}{}\label{series:I:A}
\StatementRecord{series:I:A}{proved}{03JT}{\SeriesRef{I:weierstrass-comparison}, \SeriesRef{fam:src:B.1.1}, \SeriesRef{fam:src:C.1.1}, \SeriesRef{fam:src:C.2.3}, \SeriesRef{ari:C.4.2}, \SeriesRef{ari:C.4.3}, \SeriesRef{ari:C.4.4}, \SeriesRef{ari:C.6.1}, \SeriesRef{ari:C.6.3}, \SeriesRef{ari:C.8.15}, \SeriesRef{push:src:D.4.8}, \SeriesRef{ari:C.8.1}, \SeriesRef{ari:C.8.2}, \SeriesRef{ari:C.8.3}, \SeriesRef{ari:C.8.4}}
The missing period and the normalized norm are
\[
 \beta_{\rm part}=\frac{6i}{\pi}\Log
 \frac{\vartheta_1(\mu,\tau)}{2\sqrt3\,\eta_D(\tau)},
 \qquad
 N(P)=\left|\frac{\vartheta_1(\mu,\tau)}{\eta_D(\tau)}\right|
           e^{-\pi(\Im\mu)^2/\Im\tau},
 \quad P=[\mu].
\]
The displayed logarithm is the branch fixed by the cusp condition.
Putting $L=\log N$ and
$D_c=\Im(\beta_{\rm part}+c)-6(\Im\mu)^2/\Im\tau$, one has
\[
 D_{
 \rm part}=\frac6\pi\log\frac{N}{2\sqrt3},\qquad
 \max N=\sqrt2\quad(t=1),\qquad
 \max D_{\rm part}=-\frac3\pi\log6.
\]
All ordinary fibres and both normalized finite root fibres are
compact complex tori exactly on
\[
 \Im c<H_*:=\frac3\pi\log6,
 \qquad 0<|w|<216.
\]
There is no coarse critical point of $L$ away from $t=1$;
$t=0$ is a noncritical $C^1$ point with $L_t(0)=7/384$.
For every angle $\theta$, $L(1+re^{i\theta})$ decreases strictly
from $\log\sqrt2$ to $-\infty$ as $r$ increases from zero.
Every level below the maximum is one Jordan curve strictly
star-shaped about $1$, real analytic except for a possible
$C^1$ point with nonzero tangent at $t=0$ on the level
$L=\tfrac14\log3$.  Its bounded complementary component is the strict
superlevel set; these sets form nested discs.

In any of the explicitly transported pointed Weierstrass models,
for $m\ge2$ and $mP\ne O$,
\[
 N(P)^{m^2}=\frac{N(mP)|\Delta_E|^{(m^2-1)/12}}{|\psi_m(P)|}.
\]
For torsion $P$ of order $n$ and $m\ge2$ with
$m\equiv\pm1\pmod n$, this gives
\[
 N(P)^{12}=\frac{|\Delta_E|}{|\psi_m(P)|^{12/(m^2-1)}}.
\]
In particular the five specified values are
\[
\begin{array}{c|c|c}
t&\operatorname{ord}(P)&N(P)^{12}\\\hline
0&3&27\\1&2&64\\32/5&4&10\\
64(47-27\sqrt3)/11&6&17+10\sqrt3\\
64(47+27\sqrt3)/11&6&10\sqrt3-17
\end{array}
\]
and the body gives the complete exact division-polynomial lists
for orders four through eight.  For every number field $K$ and
$t\in K\setminus\{0,1\}$, in the arithmetic model
$y^2=x^3-3t^3(t-1)x/4+t^4(t-1)^2/4$ with
$P=(0,t^2(t-1)/2)$,
\[
 \sum_{\sigma:K\hookrightarrow\C}\log N(P_{\sigma(t)})
 =\sum_{v\nmid\infty}n_v\lambda_v(P_t)
       -[K:\Q]\widehat h_O(P_t).
\]
The normalization is $\lambda_\sigma(P)=-\log N(P)$, with complex
embeddings counted separately.  For rational $t$, the positive
algebraic factor $A_t=\exp(\sum_p\lambda_p(P_t))$ gives
$N(P_t)=A_t e^{-\widehat h_O(P_t)}$.  At the eleven nonspecial
rational CM parameters $t=j/1728$ its values are
\[
\begin{array}{r|c@{\qquad}r|c}
j&A_t&j&A_t\\\hline
8000&\sqrt6&-3375&2\\
54000&2^{2/3}3^{1/4}&-32768&\sqrt3\\
287496&2^{3/4}&-884736&1\\
16581375&2&-12288000&3^{5/12}\\
-884736000&1&-147197952000&1\\
-262537412640768000&1&&
\end{array}
\]

\end{maintheorem}

\textit{Mechanism.} Theta transformation laws fix the additive constant, division identities control the norm, and the local-height product formula fixes its arithmetic normalization.
\textit{Literature status.} New period evaluations, norm values and level classification, using the elliptic section of \cite[Remark~3.12, p.~19]{AC2026}.
\hyperlink{outline-I-A}{Proof in outline}.

\hypertarget{main-I-B}{}
\begin{maintheorem}[All discrete data, the parameter, and real forms]
\hypertarget{tag-03JU}{}\label{series:I:B}
\StatementRecord{series:I:B}{proved}{03JU}{\SeriesRef{fam:src:A.5.1}, \SeriesRef{fam:src:B.2.1}, \SeriesRef{fam:src:B.2.2}, \SeriesRef{fam:src:B.2.3}, \SeriesRef{fam:src:B.3.1}, \SeriesRef{ari:C.2.1}, \SeriesRef{fam:src:B.4.2}, \SeriesRef{mapsI:src:A.3.1}, \SeriesRef{mapsI:src:A.3.2}, \SeriesRef{mapsI:src:A.3.3}, \SeriesRef{fam:src:A.1.1}, \SeriesRef{push:src:E.1.2}}
For the fixed period solution and cusp fan, take
\[
 w_1=r_1\varepsilon+k_1d,\qquad
 w_2=r_2\varepsilon'+k_2d,\qquad
 3\nmid r_1,\quad2\nmid r_2,
 \quad p=12l_0-4r_1-3r_2,
\]
where $l_0$ is the first integral winding of the full holomorphic
cusp collar in the stated positive convention.  Two such members
are biholomorphic exactly when
\[
 n=c'-c\in\Z,\qquad p'=ap\ (a\in\{1,-1\}),\qquad
 n\equiv k_2-k'_2\pmod2.
\]
Every biholomorphism is over $B$, and the map with sign $a$ has
fibre-linear part $aI_2$.  The complete invariants are
\[
 \bigl(|p|,e^{-\pi i(c+k_2)}\bigr),\qquad
 |p|\ge1,\quad\gcd(|p|,6)=1,
\]
and every displayed integer occurs.  The additional analytic
collar part and its second integral winding contribute no further
invariant.  For sphere members, $|p|=1$, a fixed datum has period
subgroup $2\Z$.  At each fixed admissible $u=w^2$ there are
exactly two sphere isomorphism classes, exchanged by a positive
circuit of $u$ about zero.

For positive $k,l$ with $\gcd(3k,4l)=1$, put $m_1=3k,m_2=4l$.
Use Viaclovsky's sphere data
$w_j=r_j\varepsilon_j+s_jd$, $\gcd(r_j,m_j)=1$,
$\varepsilon_1=\varepsilon$, $\varepsilon_2=\varepsilon'$, and
$p=m_1m_2l_0-m_2r_1-m_1r_2=\pm1$.
Two members are biholomorphic precisely when their multiplicity
pairs agree and, for some $a=\pm1$ and integer $n=c'-c$,
\[
 p'=ap,\quad r'_j\equiv ar_j\pmod{m_j},\quad
 nr'_1\equiv as_1-s'_1\pmod{k},\quad
 nr'_2\equiv as_2-s'_2\pmod{2l}.
\]
The fixed-data period subgroup is $2kl\Z$.  At fixed admissible
$u$ there are exactly $2kl$ classes for this multiplicity pair,
and a parameter circuit permutes them in one cycle.

For the standard family the reduced finite fibres have Albanese
nomes
\[
 q_2=iw/216,\qquad q_1=-w^2/48^3=(27/64)q_2^2.
\]
An antiholomorphic diffeomorphism $X_c\to X_{c'}$ exists exactly
when $c'-(1-\bar c)\in2\Z$, equivalently $w'=-\bar w$.
The members with real structures are precisely
$\Re c\in\tfrac12+\Z$.  On each,
\[
 \operatorname{Aut}^{\pm}(X_c)=\C^*\rtimes\Z/2,
 \qquad z\longmapsto1/\bar z,
\]
and there are two holomorphic-conjugacy classes of antiholomorphic
involutions.  For negative imaginary $w$ both fixed manifolds
are $S^3$, with the semifree circle action having the circle
$D_2^{\sigma_c}$ as fixed set.  For positive imaginary $w$ one
class is fixed-point free and the other has fixed manifold $S^3$
with the free Hopf circle action, up to the stated circle and
orientation conventions.

\end{maintheorem}

\textit{Mechanism.} Integral centralizers, finite norm maps and cusp winding determine the parameter equivalence, while the marked reflection and cap-meridian maps determine the real forms.
\textit{Literature status.} New classification for the displayed discrete data and real forms; the standard $2\mathbb Z$ equivalence is known \cite[Theorem~1.4, p.~3]{DuEtAl2026v1}.
\hyperlink{outline-I-B}{Proof in outline}.

\hypertarget{main-I-C}{}
\begin{maintheorem}[The monodromy quotient and every fibre's algebraic dimension]
\hypertarget{tag-03JV}{}\label{series:I:C}
\StatementRecord{series:I:C}{proved}{03JV}{\SeriesRef{fam:src:F.1.1}, \SeriesRef{I:monodromy-fundamental-set}, \SeriesRef{fam:src:F.2.2}, \SeriesRef{fam:src:F.2.3}, \SeriesRef{fam:src:F.3.1}, \SeriesRef{fam:src:D.1.4}, \SeriesRef{fam:src:D.1.6}}
Let $\psi=6\gamma\wedge\delta+u\wedge w$ on $\Lambda$, and
let $P_{d,\Z}$ fix $d$ as a vector in its integral symplectic
group.  Write $J_2=\left(\begin{smallmatrix}0&1\\-1&0\end{smallmatrix}\right)$.
Every member has unique block coordinates
\[
 g(B,x,t)=\begin{pmatrix}1&0&0\\6x&B&0\\
                   t&-x^{\mathsf t}J_2B&1\end{pmatrix},
 \quad B\in\operatorname{SL}_2(\Z),\ x\in\Z^2,\ t\in\Z.
\]
For $S=\left(\begin{smallmatrix}0&-1\\1&0\end{smallmatrix}\right)$,
$T=\left(\begin{smallmatrix}1&1\\0&1\end{smallmatrix}\right)$,
let $\alpha:\operatorname{SL}_2(\Z)\to\Z/12$ be the homomorphism
with $\alpha(S)=9$, $\alpha(T)=1$.  The monodromy group is
\[
 \Gamma=\{g(B,x,t):t+6(x_1+x_2+x_1x_2)-\alpha(B)=0\pmod{12}\}.
\]
It has index 12 in $P_{d,\Z}$ and index 24 in $P_\Z$, the group preserving
$\Z d\subset\ker\gamma\subset\Lambda$.  In the real vector stabilizer $P_d(\mathbb R)$,
a fundamental set for $\Gamma\backslash P_d(\mathbb R)$
is given by the classical modular region for
$\tau$, the Iwasawa angle $0\le\theta<\pi$,
$0\le x_1,x_2<1$ and $0\le t<12$, with the modular
lifts $g(S,0,9)$, $g(T,0,1)$, $g(-I,0,6)$ and the Heisenberg
face identifications
\[
 (x_1,x_2,t)\mapsto(x_1+1,x_2,t-6-6x_2),\quad
 (x_1,x_2,t)\mapsto(x_1,x_2+1,t-6+6x_1).
\]

On $\mathcal D=\operatorname{Sp}_4(\R)/\operatorname U(1,1)$ put
$q_d(J)=\psi(d,Jd)$.  The largest $\Gamma$-invariant open domain
of proper discontinuity is $\mathcal O=\{q_d\ne0\}$;
the same assertion holds for $P_\Z$-invariant open domains and
its action.  The null locus has a dense set of
points with infinite monodromy stabilizer.  Each sign region has
a Hausdorff complex analytic orbifold quotient.  In the family
chart, $\mathcal O_+$ has $\Im\tau>0$, $D_c<0$ and
$q_d=-6/D_c$.  For every admissible $c$ the period map is
injective and immersive on $B^\circ$ and extends to a proper
holomorphic embedding
\[
 \C=B\setminus\{\infty\}\lhook\joinrel\longrightarrow
 \Gamma\backslash\mathcal O_+.
\]

For every smooth fibre its N\'eron--Severi group is the integral
kernel in $\Z^6$ of
\[
 E_{12}-E_{13}\tau+(6E_{23}-E_{14})\mu+E_{24}\beta_c
             +E_{34}(6\mu^2-\tau\beta_c)=0.
\]
With its complex orientation,
$E^2=-2\operatorname{Pf}(E)$, where
$\operatorname{Pf}(E)=E_{12}E_{34}-E_{13}E_{24}+E_{14}E_{23}$.
Let $\mathcal E_c$ be the image on $B^\circ$ of the union of
these zero sets for $E\ne0$ with $\operatorname{Pf}(E)=0$,
and let $\mathcal A_c$ use $\operatorname{Pf}(E)<0$.
They are exactly the elliptic-quotient locus and abelian-surface
locus, respectively, and
\[
 a(F_t)=2\ (t\in\mathcal A_c),\quad
 a(F_t)=1\ (t\in\mathcal E_c\setminus\mathcal A_c),\quad
 a(F_t)=0\ (t\notin\mathcal E_c\cup\mathcal A_c).
\]
For each nonzero form $E$ in this integral kernel with zero
Pfaffian, put
$K_{\mathbb R}=\ker(E:\Lambda_{\mathbb R}\to\Lambda_{\mathbb R}^*)$
and $K=K_{\mathbb R}\cap\Lambda$. The maps
\[
 (\Pi_cK_{\mathbb R})/\Pi_cK\ \hookrightarrow F_t
 \longrightarrow\mathbb C^2/(\Pi_cK_{\mathbb R}+\Pi_c\Lambda)
\]
give every connected elliptic kernel and its quotient, including
kernels that mix the two displayed coordinates.  The invariant
class $\eta=6\gamma\wedge\delta+u\wedge w$ has square $-12$;
a fibre with $\operatorname{NS}=\Z\eta$ has no compact curves.

\end{maintheorem}

\textit{Mechanism.} The integral congruence kernel determines the parabolic quotient, and the wedge equation with its Hermitian form determines the fibre strata.
\textit{Literature status.} New integral monodromy and maximal proper domain, answering \cite[Remark~8.4, pp.~23--24]{Engel2026}.
\hyperlink{outline-I-C}{Proof in outline}.

\hypertarget{main-I-D}{}
\begin{maintheorem}[Cohomology and the ordinary relative differential modules]
\hypertarget{tag-03JW}{}\label{series:I:D}
\StatementRecord{series:I:D}{proved}{03JW}{\SeriesRef{inv:C3.8}, \SeriesRef{inv:A2.2}, \SeriesRef{inv:A3.1}, \SeriesRef{inv:C1.1}, \SeriesRef{inv:C2.1}, \SeriesRef{inv:C3.3}, \SeriesRef{inv:A4.1}, \SeriesRef{inv:A4.2}, \SeriesRef{inv:A2.4}, \SeriesRef{inv:A1.0}}
Every admissible standard sphere member has the following
Dolbeault, Bott--Chern and Aeppli dimensions, with columns
$p=0,1,2,3$ and rows $q=3,2,1,0$:
\[
\begin{array}{c|rrrr}q\backslash p&0&1&2&3\\\hline
3&0&0&0&1\\2&0&1&2&1\\1&1&2&1&0\\0&1&0&0&0
\end{array}_{\bar\partial}\quad
\begin{array}{c|rrrr}q\backslash p&0&1&2&3\\\hline
3&0&0&1&1\\2&0&1&2&1\\1&0&2&1&0\\0&1&0&0&0
\end{array}_{BC}\quad
\begin{array}{c|rrrr}q\backslash p&0&1&2&3\\\hline
3&0&0&0&1\\2&0&1&2&0\\1&1&2&1&0\\0&1&1&0&0
\end{array}_{A}.
\]
The Fr\"olicher sequence satisfies $E_1\ne E_2=E_\infty$;
the only nonzero $d_1$ have rank one at
$(p,q)=(0,1),(1,1),(1,2),(2,2)$.  The Bott--Chern algebra is
$\C\cdot1\oplus\mathcal N$ with $\mathcal N^2=0$.

For the ordinary relative differentials the local presentations
retain both the degree and the full fibre multiplicity.  For
$R=\C\{x,y,z\}$ and $t=xy$,
\[
 \Omega^1_{R/\C\{t\}}=R^3/R(y,x,0)\simeq(x,y)\oplus R,
 \quad\Omega^2\simeq(x,y)\oplus R/(x,y),
 \quad\Omega^3\simeq R/(x,y).
\]
For $t=xyz$, $\Omega^1=R^3/R(yz,xz,xy)$ is torsion-free of
rank two and not locally free on the double locus.  With
$I=(xy,xz,yz)$,
\[
 0\longrightarrow
 \frac{Ru\oplus Rv}{(xu,\,y(u-v),\,zv)}
 \longrightarrow\Omega^2\xrightarrow{\wedge dt}I
 \longrightarrow0,
 \qquad \Omega^3=R/I.
\]
The left term is exactly the torsion and is annihilated by $I$.
For $t=s^m$ in $R=\C\{s,y_1,y_2\}$, the three modules are
\[
 R^2\oplus R/(s^{m-1})\,ds,\qquad
 R\,dy_1\wedge dy_2\oplus(R/(s^{m-1}))^2,
 \qquad R/(s^{m-1})\,ds\wedge dy_1\wedge dy_2.
\]
Consequently the degree-one torsion of $\Omega^1_{X/B}$ is
supported exactly on $S_1\cup S_2$, while degree two also has
torsion along the cusp double locus.

There are specified global forms
$G\in\mathcal A^{0,1}(X)$ and $Q\in\mathcal A^{1,1}(X)$ satisfying
\[
 \bar\partial G=\bar\partial Q=0,\qquad
 [G]=g_\gamma,\qquad i_b^*[Q]=\eta,\qquad
 \exp(12[G])=f^*\mathcal O_{\mathbb P^1}(1).
\]
On the specified common refinement of the original atlas they are
$G=\bar\partial\sum w_nf_{a,n}$ and
$Q=\tfrac1{12}\bar\partial\sum w_nc_{a,n}$, with the locally finite
weights, holomorphic scalar cocycle and holomorphic one-form cocycle
of \eqref{inv:eq:atlas-final-forms}.  The elementary derivative
recursion specifies every coefficient through both multiple fibres
and every cusp stratum.  The form $Q$ is complex; a real
$\bar\partial$-closed form cannot have this fibre class.

\end{maintheorem}

\textit{Mechanism.} Root characters and logarithmic residues compute the connecting scalar and direct images, exterior quotients by $dt$ give the ordinary modules, and theta factors with cusp Cartier frames give the coefficient cocycles.
\textit{Literature status.} New evaluated cohomology and coefficient atlas; the ordinary-module calculation corrects \cite[Remark~9.21, p.~81]{AC2026}.
\hyperlink{outline-I-D}{Proof in outline}.

\hypertarget{main-I-E}{}
\begin{maintheorem}[Every Picard twist and the deformation parameter]
\hypertarget{tag-03JX}{}\label{series:I:E}
\StatementRecord{series:I:E}{proved}{03JX}{\SeriesRef{inv:D4.1}, \SeriesRef{inv:D4.2}, \SeriesRef{inv:D4.3}, \SeriesRef{inv:D4.4}, \SeriesRef{inv:D4.5}, \SeriesRef{fam:st:A.5.2}, \SeriesRef{fam:src:B.4.1}, \SeriesRef{push:src:E.2.1}, \SeriesRef{inv:C3.2}, \SeriesRef{inv:C3.1}, \SeriesRef{push:src:E.1.2}}
Use the primitive coordinate $M_z=\exp(z\mathfrak g)$ on
$\operatorname{Pic}(X)=\C$, so
$M_n=\mathcal O(n(S_1-S_2))$ for $n\in\Z$.
Let $r_1(n)\in\{0,1,2\}$ and $r_2(n)\in\{0,1,2,3\}$ be the
residues of $n$ modulo 3 and $3n$ modulo 4, respectively, and set
\label{layout:I:1}
\[
\begin{gathered}
d(n)=\frac{n-4r_1(n)-3r_2(n)}{12},\quad
 \\
U(n)=\max\{d(n)+1,0\},\quad V(n)=\max\{-d(n)-1,0\},
 \\
\quad\epsilon_n=\begin{cases}0&3\mid n,\\1&3\nmid n.\end{cases}
\end{gathered}
\]
Then
\[
\begin{aligned}
h^0(T_X\otimes M_n)&=U(n)+U(n-1),\\
h^1(T_X\otimes M_n)&=V(n)+V(n-1)+U(n)+U(n-1)+U(n+2),\\
h^2(T_X\otimes M_n)&=V(n)+V(n-1)+V(n+2)+U(n+6)+\epsilon_n,\\
h^3(T_X\otimes M_n)&=V(n+6).
\end{aligned}
\]
For $z\notin\Z$, let $\mathcal R_z$ be the set of distinct
regular base points for which integers $n,p,q,r$ satisfy
\[
 q\tau+r\mu=p,\qquad z=n-6q\mu-r\beta_c.
\]
It is finite; write $N_z=|\mathcal R_z|$.  Put
$\rho=e^{\pi i/3}$, $\beta_j=\beta_c(z_j)$ at the specified
order-three and order-four marked root centers.  Define
$\ell_1(z)=1+\epsilon_n$ at
$z=n+k(3\beta_1+4-2\rho)$, $k\ne0$, and zero elsewhere;
define $\ell_2(z)=1$ at
$z=n+k(2\beta_2+3-3i)$, $k\ne0$, $n\equiv1,2\pmod4$,
and zero elsewhere.  With
$J_z=2N_z+\ell_1(z)+\ell_2(z)$,
\[
 (h^0,h^1,h^2,h^3)(T_X\otimes M_z)=(0,J_z,1+J_z,0).
\]
The cohomology jump loci are locally discrete in this
Picard line, and the generic tuple is $(0,0,1,0)$.

The untwisted tangent tuple is $(1,1,1,0)$, the holomorphic
automorphism group of the standard member is $\C^*$, and its
Kuranishi germ is a smooth reduced line.  For every sphere member
with coprime multiplicities $(3k,4l)$ and the data of Theorem~\ref{series:I:B},
put $z_0=kl$, $d_0=3k+4l$.  On the analytic weighted projective root stack
$\mathfrak B_{k,l}=\Pp(3k,4l)$, the vertical root bundle $E$ and the root lift satisfy
\[
 E=\mathcal O\oplus\mathcal O(-z_0),\quad
 R^q\bar f_*\mathcal O_X=
 \mathcal O,\ \mathcal O\oplus\mathcal O(-z_0),\ \mathcal O(-z_0),
 \quad K_X=\bar f^*\mathcal O(z_0-d_0).
\]
Its tangent tuple is again $(1,1,1,0)$ and its Kuranishi germ is
a smooth reduced line.  For the groupoid derivative
$\kappa(\Phi)=(\partial_c\Phi)\Phi^{-1}$ the classes of
$\partial_c$ and $\partial_u$ are
$e\otimes[\gamma]$ and $-(2\pi i u)^{-1}e\otimes[\gamma]$,
respectively, and are nonzero.  Every stated fixed-data period
family is locally miniversal.

\end{maintheorem}

\textit{Mechanism.} Weighted section counts and local Koszul modules compute every Picard twist, while the nonzero Kodaira--Spencer class gives the coordinate on the Kuranishi line.
\textit{Literature status.} New all-twist and higher-multiplicity formulas; the standard untwisted deformation results are known \cite[Theorem~6.12, pp.~81--82; Proposition~6.16, p.~85]{DuEtAl2026v1}.
\hyperlink{outline-I-E}{Proof in outline}.

\hypertarget{main-I-F}{}
\begin{maintheorem}[The puncture and the obstruction to every proper filling]
\hypertarget{tag-03JY}{}\label{series:I:F}
\StatementRecord{series:I:F}{proved}{03JY}{\SeriesRef{night:P2.2}, \SeriesRef{night:P2.3}, \SeriesRef{night:P2.4}, \SeriesRef{night:P2.5}, \SeriesRef{night:P2.6}, \SeriesRef{night:P2.7}, \SeriesRef{night:P2.9}, \SeriesRef{night:P2.10}, \SeriesRef{night:P2.11}}
For an admissible $(3,4)$ logarithmic datum put
$b=4l_1+3l_2$, $p=12l_0-b$, $a=-p$, with
$3\nmid l_1$ and $l_2$ odd.  After $w=r^6$ there is a proper
holomorphic partial filling over
\[
 (\Pp^1\times\Delta_r)\setminus\{(\infty,0)\}
\]
with smooth total space and reduced normal-crossing parameter
central divisor.  Over a regular finite base point the central
fibre is a cycle of twelve ruled surfaces, with component cover
\[
 y^{12}=t^{4l_1}(t-1)^{3l_2}
\]
and rotations $4l_1,3l_2,-b$ modulo 12.
On $q=Q^{12m}$ and $r=R^n$, the difference between the 
cusp-extending and parameter-extending origins has deck
differences
\[
 am\gamma-6am^2\delta,\qquad
 an\gamma-6an(n-1)\delta,
\]
so the radial class modulo $\Z\widehat\delta$ is
$-pn\widehat\gamma$ in these fixed conventions.

The punctured family has no proper Hausdorff complex-analytic
filling at $w=0$ after any finite parameter cover.  The conclusion
permits arbitrary singularities, multiplicities and extra central
components, and holds for every admissible datum, including every
sphere member.  It also holds for every admissible datum of the
coprime $(3k,4l)$ construction.  In that case
\[
 N=12kl,\quad b=4l\,l_1+3k\,l_2,\quad a=b-Nl_0,
 \qquad \gcd(l_1,3k)=\gcd(l_2,4l)=1,
\]
after $w=r^{N/2}$ the corresponding partial model has $N$ components.
On $q=Q^{Nm}$ and $r=R^n$, the two displayed deck differences
have $N/2$ in place of 6.

For the untwisted elliptic-surface quotient $X(u)$, whose
logarithmic translations are zero, the origins have zero
corner winding.  That family admits a proper flat Hausdorff
analytic filling at $u=0$, obtained by the two-projective-line-bundle
plumbing, proper modification and finite quotient in the body.
The particular one-component central model proposed in
\cite[Conjecture~8.2, p.~22]{Engel2026} cannot be the central
fibre of a proper flat smoothing to these $X(u)$: its virtual
Euler characteristic is zero, whereas $e(X(u))=2$.

\end{maintheorem}

\textit{Mechanism.} The intrinsic canonical system recovers the marked base, finite-orbit vanishing kills horizontal meridians, and their constant ratio across the connected compact corner contradicts the nonzero cusp coefficient.
\textit{Literature status.} New filling obstruction and untwisted plumbing; the Euler calculation corrects the proposed central model in \cite[Conjecture~8.2, p.~22]{Engel2026}.
\hyperlink{outline-I-F}{Proof in outline}.

\hypertarget{main-I-G}{}
\begin{maintheorem}[The outer boundary and all marked additive limits]
\hypertarget{tag-03JZ}{}\label{series:I:G}
\StatementRecord{series:I:G}{proved}{03JZ}{\SeriesRef{push:src:D.1.3}, \SeriesRef{push:src:D.2.1}, \SeriesRef{push:src:D.2.3}, \SeriesRef{push:src:D.3.1}, \SeriesRef{eq:d-dual-exact}}
For every $w_0$ with $|w_0|=216$ and every connected open
neighbourhood $U$ of $w_0$, there is no flat proper holomorphic
map from a reduced complex space to $U$ whose fibre at each
$w\in U\cap\{0<|w|<216\}$ is biholomorphic to $X_w$.
The intrinsic divisor $2S_2$ and its Albanese nome $iw/216$
retain the invariant $J(iw/216)$, where
$J(q)=q^{-1}+744+196884q+\cdots$ has the unit circle as a
natural boundary.

Fix $c_*$ with $\Im c_*=H_*$, put $a=\Re c_*+3/2$, and use the
specified order-four root coordinate $h$ at $t=1$ and its fixed
complex-linear fibre coordinates $(\eta,\xi)$.  In the transformed
integral period basis the four vectors are
\[
 (1,1/2),\quad(\tau,1/2+h),\quad(0,1),
 \quad(6h,a+B(h)+c-c_*),
 \qquad B(h)=-3ih^2+O(h^4).
\]
For the total root quotient at fixed $c=c_*$, two central points
are nonseparated exactly when
\[
 \Re(\xi'-\xi)\in G_a=\tfrac12\Z+a\Z\quad(a\in\Q);
\]
every central pair is nonseparated for $a\notin\Q$.  Distinct
noncentral points and distinct base fibres are separated.
The Hausdorff reflection keeps the noncentral complex tori and
has central fibre $\R/G_a\simeq S^1$ for rational $a$, and one
point for irrational $a$.  For reduced $a=m/n$, $n>0$,
$G_a=\operatorname{lcm}(2,n)^{-1}\Z$.

For sequences $c_n\to c_*$ with $\Im c_n<H_*$ and $h_n\to0$,
write
\[
 \delta_n=-D_{c_n}(h_n)>0,\quad
 V_n=\delta_n^{-1}(6\Re h_n,6\Im h_n,\Re(c_n-c_*)),
\]
\[
 L_a=\{(p,q,r):r\in\Z,\ p+r/2\in\Z,
                      \ q+r/2\in\Z,\ ar\in\Z\}.
\]
It has rank three for rational $a$, and equals
$\Z^2\times\{0\}$ for irrational $a$.  Each convergent
subsequence of the closed period lattices has limit
\[
 \Gamma_\infty=\{(\eta,\xi):p\Re\eta+q\Im\eta+r\Re\xi
       +\ell(p,q,r)\Im\xi\in\Z\text{ for every }(p,q,r)\in K\},
\]
where $K\subset L_a$ is primitive,
$K=\{\lambda:\lambda\cdot V_n\text{ has a finite limit}\}$,
and $\ell(\lambda)=\lim\lambda\cdot V_n$.
For labels outside $K$ the absolute values tend to infinity.
Conversely every such primitive $K$ and homomorphism
$\ell:K\to\R$ occurs.  The resulting pointed additive quotient
limits are real tori of every dimension zero through three for
rational $a$, and zero through two for irrational $a$.

\end{maintheorem}

\textit{Mechanism.} The intrinsic Albanese invariant detects the outer boundary, and the exact dual-period equations with their quadratic cancellation determine nonseparation and every marked closed-group limit.
\textit{Literature status.} New boundary obstruction and marked limit formulas, using the annihilator duality of \cite[Theorem~3.1, p.~2017]{Cornulier2011}.
\hyperlink{outline-I-G}{Proof in outline}.

\subsection{The proofs in outline}
\hypertarget{tag-03K0}{}\label{series:I:outlines}

\hypertarget{outline-I-A}{}
\paragraph{Theorem~\ref{series:I:A}.}
\begin{enumerate}[label=\arabic*.]
\item\hypertarget{step-I-A-1}{} The pointed maps
$(x,y)\mapsto(2x,-2\sqrt2y)$ and
$(x_0,y_0)\mapsto(-x_0,y_0/(2i))$ transport the chosen section,
invariant differential and period lattice.  The quotient
exponential identifies the period model with the elliptic-surface
line-bundle quotient
(\hyperref[I:weierstrass-comparison]{\ref*{I:weierstrass-comparison}},
\hyperref[fam:src:B.1.1]{\ref*{fam:src:B.1.1}}).
\item\hypertarget{step-I-A-2}{} The two theta transformation
factors give the prescribed affine period laws; the cusp product
fixes the additive constant and hence the displayed logarithm
(\hyperref[fam:src:C.1.1]{\ref*{fam:src:C.1.1}}).
\item\hypertarget{step-I-A-3}{} The distribution identity gives
the exact division-polynomial norm formula, and $mP=\pm P$
cancels the remaining norm.  The stated polynomial roots and
values then determine every order-four-through-eight
specialization
(\hyperref[ari:C.4.2]{\ref*{ari:C.4.2}},
\hyperref[ari:C.4.3]{\ref*{ari:C.4.3}},
\hyperref[ari:C.4.4]{\ref*{ari:C.4.4}}).
\item\hypertarget{step-I-A-4}{} Exact polynomial bounds away
from the square elliptic point and the controlled local theta
expansion give $\max N=\sqrt2$.  The determinant identity and
its cusp limit then give exactly $\Im c<H_*$
(\hyperref[ari:C.8.1]{\ref*{ari:C.8.1}},
\hyperref[ari:C.8.2]{\ref*{ari:C.8.2}},
\hyperref[ari:C.8.3]{\ref*{ari:C.8.3}},
\hyperref[ari:C.8.4]{\ref*{ari:C.8.4}},
\hyperref[fam:src:C.2.3]{\ref*{fam:src:C.2.3}}).
\item\hypertarget{step-I-A-5}{} The normalized critical factor
satisfies $|1-X|<1$ and
$(t-1)L_t=X(X-2)/24$.  Its radial sign, the coarse derivative
at zero and the cusp limit give the complete critical and level
classification
(\hyperref[ari:C.8.15]{\ref*{ari:C.8.15}},
\hyperref[push:src:D.4.8]{\ref*{push:src:D.4.8}}).
\item\hypertarget{step-I-A-6}{} The local multiplication identity
evaluates every finite-place contribution on the stated minimal
models.  The product formula identifies the archimedean term and
gives the eleven exact CM factors
(\hyperref[ari:C.6.1]{\ref*{ari:C.6.1}},
\hyperref[ari:C.6.3]{\ref*{ari:C.6.3}}).
\end{enumerate}


\hypertarget{outline-I-B}{}
\paragraph{Theorem~\ref{series:I:B}.}
\begin{enumerate}[label=\arabic*.]
\item\hypertarget{step-I-B-1}{} The intrinsic canonical pencil
fixes the base.  Its period centralizer forces $c'-c\in\Z$ and
the fibre-linear part $aI_2$
(\hyperref[fam:src:A.1.1]{\ref*{fam:src:A.1.1}},
\hyperref[fam:src:A.5.1]{\ref*{fam:src:A.5.1}}).
\item\hypertarget{step-I-B-2}{} Solving the two finite norm
equations gives the parity condition; comparing the cusp
collar gives $p'=ap$.  The holomorphic correction solves the
remaining gluing, so these necessary equations are sufficient
(\hyperref[fam:src:B.2.1]{\ref*{fam:src:B.2.1}},
\hyperref[fam:src:B.2.2]{\ref*{fam:src:B.2.2}},
\hyperref[fam:src:B.2.3]{\ref*{fam:src:B.2.3}}).
\item\hypertarget{step-I-B-3}{} The higher root pullbacks multiply
the norm images by $k,l$, yielding exactly the congruences modulo
$k,2l$.  Their intersection is $2kl\Z$, and a positive $u$
circuit is the cyclic parameter translation
(\hyperref[fam:src:B.3.1]{\ref*{fam:src:B.3.1}},
\hyperref[push:src:E.1.2]{\ref*{push:src:E.1.2}}).
\item\hypertarget{step-I-B-4}{} The invariant-coordinate lattice,
including its affine half-period, and the two exact finite theta
values give the Albanese nomes
(\hyperref[ari:C.2.1]{\ref*{ari:C.2.1}}).
\item\hypertarget{step-I-B-5}{} The explicit reflected period
and affine lattice maps give $c'\equiv1-\bar c\pmod{2\Z}$.
Their square and holomorphic conjugacy identify the two
involution classes
(\hyperref[fam:src:B.4.2]{\ref*{fam:src:B.4.2}},
\hyperref[mapsI:src:A.3.1]{\ref*{mapsI:src:A.3.1}}).
\item\hypertarget{step-I-B-6}{} The fixed-lift congruences and
cusp units determine the nonempty pieces.  The exact cap-meridian
matrices close them to $S^3$ and determine the semifree or Hopf
circle action
(\hyperref[mapsI:src:A.3.2]{\ref*{mapsI:src:A.3.2}},
\hyperref[mapsI:src:A.3.3]{\ref*{mapsI:src:A.3.3}}).
\end{enumerate}


\hypertarget{outline-I-C}{}
\paragraph{Theorem~\ref{series:I:C}.}
\begin{enumerate}[label=\arabic*.]
\item\hypertarget{step-I-C-1}{} The parabolic multiplication law
and the polarization of $x_1+x_2+x_1x_2$ prove the character law.
The commutator and its conjugates generate its integral kernel,
giving the two exact indices
(\hyperref[fam:src:F.1.1]{\ref*{fam:src:F.1.1}}).
\item\hypertarget{step-I-C-2}{} The modular lifts first reduce
the Iwasawa factor, integral Heisenberg translations reduce $x$,
and the center reduces $t$.  Interior uniqueness and the displayed
face maps give the fundamental set
(\hyperref[I:monodromy-fundamental-set]{\ref*{I:monodromy-fundamental-set}}).
\item\hypertarget{step-I-C-3}{} Real parabolic stabilizers are
compact on $q_d\ne0$.  Rational parabolic conjugates of the
explicit null unipotent give a dense infinite-stabilizer locus,
which proves maximality
(\hyperref[fam:src:F.2.2]{\ref*{fam:src:F.2.2}},
\hyperref[fam:src:F.2.3]{\ref*{fam:src:F.2.3}}).
\item\hypertarget{step-I-C-4}{} The modular function $j/1728$
is the period map's left inverse, while the finite root charts
extend it at zero and one.  Its control on compact inverse
images proves properness on the finite $t$-line
(\hyperref[fam:src:F.3.1]{\ref*{fam:src:F.3.1}}).
\item\hypertarget{step-I-C-5}{} Wedge the period rows with an
integral alternating form to obtain the exact $(1,1)$ equation.
Its Hermitian determinant and the complex orientation give
$E^2=-2\operatorname{Pf}(E)$ and the quotient/projectivity tests
(\hyperref[fam:src:D.1.4]{\ref*{fam:src:D.1.4}},
\hyperref[fam:src:D.1.6]{\ref*{fam:src:D.1.6}}).
\end{enumerate}


\hypertarget{outline-I-D}{}
\paragraph{Theorem~\ref{series:I:D}.}
\begin{enumerate}[label=\arabic*.]
\item\hypertarget{step-I-D-1}{} Root characters and the cusp
normalization compute the logarithmic direct images.  The residue
connecting scalar is $-iC_4/4\ne0$, which fixes the
remaining Hodge integer at one
(\hyperref[inv:A2.2]{\ref*{inv:A2.2}},
\hyperref[inv:A2.4]{\ref*{inv:A2.4}}).
\item\hypertarget{step-I-D-2}{} Coherent Leray on the curve,
the structure-sheaf inputs and Serre duality give the complete
Dolbeault table, uniformly for every admissible parameter
(\hyperref[inv:A1.0]{\ref*{inv:A1.0}},
\hyperref[inv:A3.1]{\ref*{inv:A3.1}}).
\item\hypertarget{step-I-D-3}{} The middle rows and the
$S^6$ de Rham groups force the four rank-one differentials.
The double-complex decomposition then gives the Bott--Chern and
Aeppli tables; wedge products of the specified representatives
give the square-zero multiplication
(\hyperref[inv:C1.1]{\ref*{inv:C1.1}},
\hyperref[inv:C2.1]{\ref*{inv:C2.1}},
\hyperref[inv:C3.3]{\ref*{inv:C3.3}}).
\item\hypertarget{step-I-D-4}{} Quotienting by $dt$ gives the
degree-one modules.  The maps $(a,b,c)\mapsto(xa-zc,yb-zc)$
and $\Omega^2\xrightarrow{\wedge dt}I$ compute their kernels
and torsion exactly; finite \'etale descent gives the global
support assertion
(\hyperref[inv:A4.1]{\ref*{inv:A4.1}},
\hyperref[inv:A4.2]{\ref*{inv:A4.2}}).
\item\hypertarget{step-I-D-5}{} The theta factors with their original
finite affine phases and the extending cusp Cartier units define the
local twelfth-power lines on the four pieces.  The exact determinant divisor makes
the displayed holomorphic connection inverse valid on every overlap
annulus.  Its one-form cocycle and the specified coordinate weights
give $Q$; the scalar cocycle gives the primitive $G$ and its factor
twelve normalization
(\hyperref[inv:C3.8]{\ref*{inv:C3.8}}).
\end{enumerate}


\hypertarget{outline-I-E}{}
\paragraph{Theorem~\ref{series:I:E}.}
\begin{enumerate}[label=\arabic*.]
\item\hypertarget{step-I-E-1}{} The primitive Picard
normalization identifies the integer twists, and weighted
coarse degrees give their section counts and four dimensions
(\hyperref[inv:C3.2]{\ref*{inv:C3.2}},
\hyperref[inv:D4.1]{\ref*{inv:D4.1}}).
\item\hypertarget{step-I-E-2}{} The integral cocycle
modulo the two period rows gives the regular resonance equation
and both finite root lattices.  At the cusp the normalization
complex gives exactly the reduced degree-two skyscraper
(\hyperref[inv:D4.2]{\ref*{inv:D4.2}},
\hyperref[inv:D4.3]{\ref*{inv:D4.3}}).
\item\hypertarget{step-I-E-3}{} The joint order-three
coefficient $123ik/(64\pi)$ fixes the local normal form;
Smith modules and every stabilizer character give the complete
finite stalk lengths.  Leray then gives the formula for $J_z$
(\hyperref[inv:D4.4]{\ref*{inv:D4.4}},
\hyperref[inv:D4.5]{\ref*{inv:D4.5}}).
\item\hypertarget{step-I-E-4}{} The established
$\bar\partial$-closed form $G_\gamma=G_0/12$ defines the
Picard family by $\bar\partial-2\pi izG_\gamma$.
Proper coherent base change and a real nonintegral nonresonant
point make every jump locus locally discrete
(\hyperref[inv:C3.1]{\ref*{inv:C3.1}},
\hyperref[inv:C3.2]{\ref*{inv:C3.2}},
\hyperref[inv:D4.5]{\ref*{inv:D4.5}}).
\item\hypertarget{step-I-E-5}{} The root-forgetting map
$(x,y)\mapsto(x^l,y^k)$ with character
$\lambda\mapsto\lambda^{kl}$ pulls back the standard Hodge
line.  Its weighted direct images and logarithmic tangent
sequence give $(1,1,1,0)$ for every stated higher sphere member
(\hyperref[fam:src:B.3.1]{\ref*{fam:src:B.3.1}},
\hyperref[push:src:E.2.1]{\ref*{push:src:E.2.1}}).
\item\hypertarget{step-I-E-6}{} Differentiating the
period groupoid gives the nonzero class $e\otimes[\gamma]$.
The induced parameter map forces the one-variable Kuranishi
ring to be reduced and smooth.  The integral translation sheaf
gives the standard automorphism group
(\hyperref[fam:st:A.5.2]{\ref*{fam:st:A.5.2}},
\hyperref[fam:src:B.4.1]{\ref*{fam:src:B.4.1}}).
\end{enumerate}


\hypertarget{outline-I-F}{}
\paragraph{Theorem~\ref{series:I:F}.}
\begin{enumerate}[label=\arabic*.]
\item\hypertarget{step-I-F-1}{} The node charts
$X_i=W/r^i$, $Y_i=r^{i+1}/W$, $X_iY_i=r$ and their inverse
overlaps give the proper rank-one cycle.  The finite affine
actions rotate its components freely and produce the stated
cover. Subtraction of the two origins fixes both deck
differences
(\hyperref[night:P2.3]{\ref*{night:P2.3}}).
\item\hypertarget{step-I-F-2}{} Resolve a proper filling's
dominating component.  Its relative anticanonical direct image
recovers the pencil, and the three labelled critical-value
sections recover the marked base by their meromorphic cross-ratio
(\hyperref[night:P2.4]{\ref*{night:P2.4}}).
\item\hypertarget{step-I-F-3}{} Modulo $\widehat\delta$ the
cross-cocycles vanish. The three unipotent annihilator
equations have determinant 6; Stein factorization on a horizontal
central component then kills its normal-meridian class
(\hyperref[night:P2.5]{\ref*{night:P2.5}},
\hyperref[night:P2.6]{\ref*{night:P2.6}}).
\item\hypertarget{step-I-F-4}{} Principalization makes both
axes normal crossing.  The Stein finite factor is birational
over the normal bidisc, so the compact corner fibre is connected.
The lifts $Q_o=T_{-mb\widehat\gamma}Q$, $R_o=R$ commute;
local meridians therefore satisfy the constant-ratio equation
(\hyperref[night:P2.7]{\ref*{night:P2.7}}).
\item\hypertarget{step-I-F-5}{} The connected incidence graph
joins a horizontal ratio zero to the strict cusp ratio $a$.
Admissibility gives $\gcd(a,12)=1$, producing the contradiction
for every proper analytic filling and every finite cover
(\hyperref[night:P2.7]{\ref*{night:P2.7}}).
\item\hypertarget{step-I-F-6}{} For higher multiplicities,
$K_X^{-N}=f^*\mathcal O(4l+3k-kl)$ with nonzero degree.
The indicated canonical or anticanonical power and its three labelled
critical points recover the degree-one base parameter.  The same ratio
argument and $\gcd(a,N)=1$ give the full higher conclusion
(\hyperref[night:P2.9]{\ref*{night:P2.9}}).
\item\hypertarget{step-I-F-7}{} With zero logarithmic
translations the two line-bundle neck equations have zero
corner winding and give the proper flat plumbing after its
specified modification and finite quotient.  Normalization and
the double-locus smoothing formula give virtual Euler zero for
Engel's particular one-component proposal, excluding its
smoothing to $e=2$
(\hyperref[night:P2.10]{\ref*{night:P2.10}},
\hyperref[night:P2.11]{\ref*{night:P2.11}}).
\end{enumerate}


\hypertarget{outline-I-G}{}
\paragraph{Theorem~\ref{series:I:G}.}
\begin{enumerate}[label=\arabic*.]
\item\hypertarget{step-I-G-1}{} The canonical divisor recovers
$S_2$ intrinsically in a proper extension.  Its normalization
and Albanese invariant retain $J(iw/216)$, whose natural boundary
excludes crossing the outer circle
(\hyperref[push:src:D.1.3]{\ref*{push:src:D.1.3}}).
\item\hypertarget{step-I-G-2}{} The fixed complex-linear and
integral basis changes give the displayed four periods.
The local theta expansion gives $B(h)=-3ih^2+O(h^4)$ and
the exact inverse transpose gives every integral dual covector
(\hyperref[push:src:D.2.1]{\ref*{push:src:D.2.1}},
\hyperref[eq:d-dual-exact]{\ref*{eq:d-dual-exact}}).
\item\hypertarget{step-I-G-3}{} Bounded dual covectors have
constant integral labels after subselection; their labels form
$L_a$.  The quadratic real term cancels, leaving
$\lambda\cdot V_n+o(1)$.  The primitive surviving subgroup and
its limiting homomorphism give exactly $\Gamma_\infty$
(\hyperref[push:src:D.3.1]{\ref*{push:src:D.3.1}}).
\item\hypertarget{step-I-G-4}{} Choose $v$ orthogonal to $K$
and outside the annihilator hyperplanes of $L_a\setminus K$.
The vectors $V_n=nv+v_0$, with the positive scales $\delta_n$
chosen as in the proof, realize every primitive subgroup and limiting
homomorphism. Pontryagin--Chabauty duality converts their dual
graphs to the stated real-torus additive quotients
(\hyperref[push:src:D.3.1]{\ref*{push:src:D.3.1}}).
\item\hypertarget{step-I-G-5}{} At fixed boundary parameter,
irrational root rays make all noncentral dual labels escape.
The surviving rational label is
$\operatorname{lcm}(2,n)\Re\xi$; its circle-valued function
separates precisely the asserted classes.  Collapsing those
classes proves the universal Hausdorff-reflection property
(\hyperref[push:src:D.2.3]{\ref*{push:src:D.2.3}}).
\end{enumerate}


% RungZero block 1 belongs to the shared series overview.
% The coefficient-atlas portion of block 25 is included in Theorem~\ref{series:I:D};
% its arithmetic portion remains in the final open-problems part,
% with the explicit unresolved relative-period gate at ari:C.7.4.
\Needspace{8\baselineskip}
\subsection{Known and new}\label{pub:I:known-new}
\begingroup\small
\begin{longtable}{@{}p{.10\textwidth}p{.85\textwidth}@{}}
\toprule Theorem & Status and credited inputs\\\midrule\endhead
A & New explicit evaluations and sharp norm bounds from the known section \cite[Remark~3.12, p.~19]{AC2026}; the quotient and height inputs are \cite[Definition~2.3, p.~5; proof of Lemma~2.1, pp.~3--4]{Engel2026}, transported by \cref{I:weierstrass-comparison}.\\
B & New discrete-data and real-form extensions of the standard equivalence \cite[Theorem~1.4, p.~3]{DuEtAl2026v1}; the higher construction is \cite[Theorem~1.1, p.~1; Section~2.1, pp.~4--6]{ViaclovskyLog2026}, with $(k,l)=(1,1)$ giving the standard data.\\
C & New exact group, quotient and proper-domain answer to \cite[Remark~8.4, pp.~23--24]{Engel2026}; the fibre tests use the integral wedge equation of \cite[Lemma~9.2, p.~65]{AC2026} with the period row $\beta=\beta_c$.\\
D & New value of the Hodge integer, full cohomology and coefficient representatives extending \cite[Theorem~9.1(vi), p.~65; Proposition~9.20, p.~80]{AC2026}; the degree-one cusp module corrects its Remark~9.21, p.~81, and specifies the degree dependence of Remark~1.3, p.~5.\\
E & New Picard-twist and coprime-root calculations extending the known standard tangent, Kuranishi and automorphism results \cite[Theorems~1.3--1.4, p.~3; Theorem~6.12, pp.~81--82; Proposition~6.16, p.~85; Theorem~7.14, p.~95]{DuEtAl2026v1}, under the parameter map of \cref{night:P5.7}.\\
F & New marked nonfilling obstruction and explicit untwisted filling in the degeneration setting of \cite[Remark~8.1, p.~22]{Engel2026}; virtual Euler number zero excludes the particular central model in its Conjecture~8.2, p.~22.\\
G & New intrinsic boundary and marked additive-limit calculations; the map $w\mapsto iw/216$ supplies the modular boundary, and \cite[Theorem~3.1, p.~2017]{Cornulier2011} supplies closed-group duality under the real annihilator map.\\
\bottomrule\end{longtable}\endgroup

\subsection{Methodology}
This text is AI-produced; no part of it was written by a human.
GPT-5.6 Sol (OpenAI), in Codex at reasoning effort ultra, contributed
to the working edition of 24 August--9 September 2026. The pinned
repository's authorship record also credits GPT-6 Astra (OpenAI),
at reasoning effort ultra, from 5 September.
GPT-6 Astra (OpenAI), in ChatGPT Pro mode, carried out the research
of 12--16 November, the independent re-derivation of 16--18 November,
and the writing of 16--20 November 2026. GPT-6 Astra (OpenAI),
in ChatGPT Pro mode, also undertook the publication revision and
further checks of this paper on 20--21 November 2026, including
the continuation on 21 November.

Claude Fable 5.2 (Anthropic), at effort high and at effort max,
and Claude Opus 5.5 (Anthropic), at effort max, undertook readings,
intakes, checks and recomputations on 12--20 November 2026.
Their objections and calculations were compared with the statements,
proofs and cited sources, and the resulting corrections were
incorporated into the text. The records specify the passages and
calculations examined by each model.

For the statements printed in this paper, the record contains
767 examinations of proof steps, 10 independent
re-derivations, and 143 computations or formal certificates.
It also records 731 readings of the proofs and applications,
194 uses of background results without a fault found,
and 1 refutation associated with its corrected statement.
Each total counts recorded links from a check to a statement;
one calculation used for several statements can contribute several
links. Every outcome specifies its scope. Re-reading an edited proof
is recorded separately from deriving its steps before comparison.
Exact algebra, numerical enclosures and analytic arguments retain
their respective scopes, and the records identify the implementations
used for independent computational comparisons.
No human has checked the mathematics.

The check records and programs are linked from the
\href{\SsixPublishRecords}{project repository}.
The text, programs and check records are dedicated to the public domain under
\href{https://creativecommons.org/publicdomain/zero/1.0/}{CC0 1.0 Universal}.

\subsection{How the paper is organized}
The marked period and quotient coordinates lead to the exact norm,
its admissible domain, and the proper family. The intrinsic
classification is followed by discrete logarithmic data, the monodromy
quotient and the exceptional fibres. Sections~\ref{pub:I:cotangent}--\ref{pub:I:twists}
compute the cotangent, Bott--Chern and tangent cohomology; the complete
coefficient construction precedes its use in the Picard family.
The integral circle quotient and the real structures then lead to
the puncture and outer-boundary problems. Their local models,
component obstructions and topology are developed with the specified
marked comparison maps.

The remaining main-body sections treat other multiplicities and finite
covers, rational CM heights, the critical level curves, and arithmetic
periods. Appendices~\ref{pub:I:parameter-appendix}--\ref{pub:I:torsor-appendix}
contain the credited proofs of parameter equivalence, tangent
cohomology, local deformations and automorphisms.
The final part states the unresolved arithmetic questions.
Each arrow in the diagram identifies a particular supplying and
receiving body argument in the table below it.
\begin{figure}[htbp]\centering
\begin{tikzpicture}[>=stealth,every node/.style={font=\small}]
\node[draw,rounded corners,minimum width=2.2cm,minimum height=.65cm] (TA) at (0.0,0.0) {\hyperlink{main-I-A}{Theorem~\ref{series:I:A}}};
\node[draw,rounded corners,minimum width=2.2cm,minimum height=.65cm] (TB) at (3.5,0.0) {\hyperlink{main-I-B}{Theorem~\ref{series:I:B}}};
\node[draw,rounded corners,minimum width=2.2cm,minimum height=.65cm] (TC) at (7.0,0.0) {\hyperlink{main-I-C}{Theorem~\ref{series:I:C}}};
\node[draw,rounded corners,minimum width=2.2cm,minimum height=.65cm] (TD) at (0.0,-1.25) {\hyperlink{main-I-D}{Theorem~\ref{series:I:D}}};
\node[draw,rounded corners,minimum width=2.2cm,minimum height=.65cm] (TE) at (3.5,-1.25) {\hyperlink{main-I-E}{Theorem~\ref{series:I:E}}};
\node[draw,rounded corners,minimum width=2.2cm,minimum height=.65cm] (TF) at (7.0,-1.25) {\hyperlink{main-I-F}{Theorem~\ref{series:I:F}}};
\node[draw,rounded corners,minimum width=2.2cm,minimum height=.65cm] (TG) at (0.0,-2.5) {\hyperlink{main-I-G}{Theorem~\ref{series:I:G}}};
\draw[->,semithick] (TA) to (TB);
\draw[->,semithick] (TE) to (TB);
\draw[->,semithick] (TD) to (TE);
\draw[->,semithick] (TB.east) -- ++(.45,0) |- (TG.east);
\draw[->,semithick] (TG.west) -- ++(-.45,0) |- (TA.west);
\end{tikzpicture}
\caption{Reuse of the indicated body arguments among the main theorem groups.}
\par\medskip
\begingroup\small
\begin{tabularx}{\textwidth}{@{}p{.10\textwidth}X@{}}
\toprule Arrow & Supplying and receiving arguments\\\midrule
$A\to B$ & \SeriesRef{fam:src:C.2.3} $\longrightarrow$ \SeriesRef{ari:C.2.1}. The exact parameter domain is used for the intrinsic Albanese nome.\\
$E\to B$ & \SeriesRef{fam:src:B.4.1} $\longrightarrow$ \SeriesRef{mapsI:src:A.3.1}. The automorphism group acts in the conjugacy classification of real involutions.\\
$D\to E$ & \SeriesRef{inv:C1.1} $\longrightarrow$ \SeriesRef{inv:C3.1}. The differential and Bott--Chern dimensions determine the Picard curvature statement.\\
$B\to G$ & \SeriesRef{fam:src:A.1.1} $\longrightarrow$ \SeriesRef{push:src:D.1.3}. The intrinsic multiplicity-four divisor supplies the Albanese continuation obstruction.\\
$G\to A$ & \SeriesRef{push:src:D.2.1} $\longrightarrow$ \SeriesRef{ari:C.8.2}. The half-period derivative identity supplies the theta-heat calculation.\\
\bottomrule\end{tabularx}\endgroup\end{figure}

\subsubsection*{Further results}
\begin{itemize}[leftmargin=*]
\item \Cref{ari:C.4.2,ari:C.4.3,ari:C.4.4}: division polynomials give the exact torsion loci and their norm values.
\item \Cref{ari:C.6.3,ari:C.6.T2}: local heights give the rational CM factors, and two-isogeny descent determines the stated rank over $\mathbb Q(\sqrt3)$.
\item \Cref{mapsI:src:A.1.1,mapsI:src:A.2.1}: the integral circle quotient has an explicit cusp replacement and relative uniqueness.
\item \Cref{inv:C3.8}: holomorphic cocycles and a locally finite coordinate partition give global coefficient representatives.
\item \Cref{ari:C.8.15,push:src:D.4.8}: the theta norm has one maximum and strictly star-shaped lower level curves.
\end{itemize}

\clearpage
\section{The marked family and its coordinates}
\hypertarget{tag-01TG}{}\label{inv:notation-dictionary}

\begin{definition}[Base and integral marking]
\hypertarget{tag-01O1}{}\label{def:common-conventions}
\StatementRecord{def:common-conventions}{proved}{01O1}{}
Let $B=\mathbb P^1_t$, $p_1=0$, $p_2=1$, $p_0=\infty$,
$\Sigma=\{p_0,p_1,p_2\}$, and $B^\circ=B\setminus\Sigma$.
The full fibres of the marked family are $3S_1$, $4S_2$ and $W$;
$S_1,S_2,W$ are reduced, and $D=W+S_1+S_2$.
We orient base meridians counterclockwise in holomorphic local
coordinates and use clockwise root deck generators
$s_j(g_jz)=e^{-2\pi i/m_j}s_j(z)$, where $(m_1,m_2)=(3,4)$ and
$g_1g_2g_0=1$. The lift of a positive base loop ends at the inverse
deck translate, and its endpoint identification applies $g_j$.

The integral homology lattice and its dual basis are
\[
 \Lambda=\mathbb Z\langle\widehat\gamma,\widehat u,
                     \widehat w,\widehat\delta\rangle,
 \qquad (\gamma,u,w,\delta).
\]
The period matrix in coordinates $(\zeta_1,\zeta_2)$ is
\begin{equation}\hypertarget{tag-03CN}{}\label{wb:pshift:eq:periods}
 \Pi_c=\begin{pmatrix}6\mu&\tau&1&0\\
                    \beta_c&\mu&0&1\end{pmatrix},
 \qquad \beta_c=\beta^{\rm part}+c,
 \qquad (\beta^{\rm part}+\tau)(p_0)=0.
\end{equation}
Its rows are $\sigma_1=6\mu\gamma+\tau u+w$ and
$\sigma_2=\beta_c\gamma+\mu u+\delta$. The monodromy matrices are
\begin{equation}\hypertarget{tag-03C8}{}\label{wb:pshift:eq:finite-data}
 A_1=\begin{pmatrix}1&0&0&0\\6&0&1&0\\-6&-1&-1&0\\-2&1&0&1\end{pmatrix},
 \quad A_2=\begin{pmatrix}1&0&0&0\\0&0&-1&0\\-6&1&0&0\\3&0&1&1\end{pmatrix},
 \quad A_0=(A_1A_2)^{-1}.
\end{equation}
They act on columns in the displayed integral basis. Homology has
integral coefficients unless specified, and complex manifolds carry
their complex orientations. Reversing the two additive coordinates
acts by $J=\left(\begin{smallmatrix}0&1\\1&0\end{smallmatrix}\right)$:
the period and Hermitian matrices become $J\Pi_c$ and
$JHJ^{\mathsf t}$, respectively.
\end{definition}

\begin{definition}[The parameter and the finite translations]
\hypertarget{tag-PI01}{}\label{I:def:family-parameter}
\StatementRecord{I:def:family-parameter}{proved}{PI01}{}
Put
\begin{equation}
 D_{\rm part}=\Im\beta^{\rm part}
                  -6(\Im\mu)^2/\Im\tau,\qquad
 M=\sup D_{\rm part},\qquad
 \mathcal P=\{c\in\mathbb C:\Im c<-M\},\qquad w=e^{-\pi ic}.
\label{pub:I:eq:PD1}
\end{equation}
The scalar $w$ is the family parameter; in an exterior product,
$w$ denotes the marked covector of Definition~\ref{def:common-conventions}.
The standard affine translations are
$v_1=(1,2,-4,0)^{\mathsf t}$ and
$v_2=(-1,-3,3,0)^{\mathsf t}$, and the cusp collar is untwisted.
We write $u_{\rm Engel}=e^{-2\pi ic}=w^2$ for the multiplicative
parameter with the inverse period lift at $c=0$ as its reference.
The line-bundle identification and the effect of changing that reference
are proved in Proposition~\ref{fam:src:B.1.4}.
We use $\rho=e^{\pi i/3}$; the transformation $\tau\mapsto\tau+1$
carries the convention $e^{2\pi i/3}$ to $\rho$.
\end{definition}

\begin{proposition}[Integral changes of marking]
\hypertarget{tag-01S2}{}\label{inv:B1.0}
\StatementRecord{inv:B1.0}{proved}{01S2}{}
In Engel's common integral basis, our marked basis has column matrix
\[
 P_{\rm mark}=\begin{pmatrix}
 0&-2&-1&-2\\2&-1&0&-1\\0&2&1&3\\-1&0&0&0
 \end{pmatrix}.
\]
His base indices $(0,1,\infty)$ correspond to $(1,2,0)$ here.
The matrices satisfy $T_0P_{\rm mark}=P_{\rm mark}A_1$ and
$T_1P_{\rm mark}=P_{\rm mark}A_2$, and his distinguished circle is
$\delta_{\rm E}=P_{\rm mark}\widehat\delta$.
The positive-root affine vectors are
$a_0=-P_{\rm mark}v_1/3$ and $a_1=-P_{\rm mark}v_2/4$.
For reversed meridians the change of marking is
\[
 Q=P_{\rm mark}S,
 \quad S=\begin{pmatrix}-1&0&0&0\\-6&0&-1&0\\6&-1&0&0\\2&1&1&1\end{pmatrix},
 \quad Q=\begin{pmatrix}2&-1&0&-2\\2&-1&0&-1\\0&2&1&3\\1&0&0&0\end{pmatrix},
\]
and $T_0Q=QA_1^{-1}$, $T_1Q=QA_2^{-1}$.
\end{proposition}
\begin{proof}
The matrices in \cite[Theorem~5.1, p.~13; equation~(6.1), p.~16]{Engel2026}
in that common basis are
\[
 T_0=\begin{pmatrix}0&-1&-1&0\\1&-1&0&0\\0&0&1&0\\0&0&0&1\end{pmatrix},
 \qquad
 T_1=\begin{pmatrix}-1&1&-1&2\\-1&0&-1&1\\1&1&2&-1\\0&0&0&1\end{pmatrix}.
\]
Multiplication gives all four intertwining identities,
$\det P_{\rm mark}=-1$ and $\det S=1$. Thus these are integral
lattice isomorphisms with the stated actions. The last column of
$P_{\rm mark}$ is $(-2,-1,3,0)^{\mathsf t}$, the distinguished
circle in that basis. Since $A_jv_j=v_j$, inversion of the affine
deck generator sends its translation $v_j/m_j$ to $-v_j/m_j$.
Applying $P_{\rm mark}$ gives the displayed affine vectors.
\end{proof}

\begin{definition}[Root bundles and logarithmic differentials]
\hypertarget{tag-PI02}{}\label{I:def:root-bundles}
\StatementRecord{I:def:root-bundles}{proved}{PI02}{}
Let $\rho:\mathcal B\to B$ be the root curve of orders $(3,4)$,
with tautological root bundles $R_1,R_2$. Put $R=R_1R_2^{-1}$ and
$\mathcal O_{\mathcal B}(n)=R^n$.
Thus $R_1=R^4$, $R_2=R^3$, $\rho^*\mathcal O_B(1)=R^{12}$ and
$K_{\mathcal B}=R^{-7}$. Here $\mathcal O_B(n)$ has coarse degree $n$;
its pullback is $\mathcal O_{\mathcal B}(12n)$.
We use
\[
 \mathcal L=\Omega_B^1(\log\Sigma),\qquad
 \mathcal A=\Omega_X^1(\log D),\qquad
 \mathcal E=\mathcal A/f^*\mathcal L.
\]
Fibres in coherent base change are the full scheme fibres.
\end{definition}
\begin{proposition}[The root base and its explicit cohomology]\hypertarget{tag-01SM}{}\label{inv:D1.1}
\StatementRecord{inv:D1.1}{proved}{01SM}{the definitions and constructions stated in the proof}

Let $\rho:\mathcal B\to B$ be the analytic root curve with orders $3$ and
$4$ at $p_1,p_2$. Write $R_1,R_2$ for its tautological root line bundles.
Its Picard group is infinite cyclic, with generator
$R=R_1\otimes R_2^{-1}$, and
\[
 R_1=R^4,\quad R_2=R^3,\quad \rho^*\OO_B(1)=R^{12},
 \quad K_{\mathcal B}=R^{-7}.
\]
We abbreviate $R^n$ to $\OO_{\mathcal B}(n)$. Choose sections
$x\in H^0(R^3)$ and $y\in H^0(R^4)$ vanishing at $p_2$ and $p_1$,
respectively. For every integer $n$, put
\[
 r_1(n)\equiv n\pmod3,\quad 0\le r_1(n)<3,
 \qquad r_2(n)\equiv3n\pmod4,\quad0\le r_2(n)<4,
\]
\[
 d(n)=\frac{n-4r_1(n)-3r_2(n)}{12}.
\]
Then
\begin{equation}
 \hypertarget{tag-01P3}{}\label{fam:eq:ab-root-degree}
\rho_*\OO_{\mathcal B}(n)=\OO_B(d(n)),\qquad
 H^i(\mathcal B,\OO(n))=H^i(B,\OO(d(n))).                 
\end{equation}
Equivalently, $h^0(\OO(n))$ is the number of pairs
$(a,b)\in\Z_{\ge0}^2$ with $3a+4b=n$, and
$h^1(\OO(n))=h^0(\OO(-n-7))$.
\end{proposition}



\begin{proof}
The Picard group and graded cohomology are those of
\cite[Lemma~6.1, p.~74]{DuEtAl2026v1}; the following chart maps specify
the coarse invariant direct image used here.
\begin{enumerate}[label=\arabic*.,leftmargin=*]
\item On the root charts the group acts on a disc coordinate by a primitive
$m_j$th root of unity. The stabilizer character of any line bundle is
removed uniquely by tensoring with $R_1^{-r_1}R_2^{-r_2}$, with the
exponents in the displayed ranges. A line bundle with trivial stabilizer
characters descends: an equivariant nonzero local frame can be averaged,
and its invariant frame gives the descent transition function. It is
therefore the pullback of a unique line bundle on $B$. Consequently
\[
 \operatorname{Pic}(\mathcal B)=
 \langle H,R_1,R_2\mid 3R_1=H,\;4R_2=H\rangle=\Z R.
\]
The weighted quotient presentation is
\[
 \mathcal B\simeq[(\mathbb C^2\setminus\{0\})/\mathbb C^*],
 \quad\lambda(x,y)=(\lambda^3x,\lambda^4y),\qquad
 t=\frac{y^3}{y^3-x^4}.
\]
The root sections are $(\mathcal O(4),y)$ and $(\mathcal O(3),x)$.
They give isomorphisms on each root chart, and the same coarse
coordinate on their ordinary overlap makes those isomorphisms glue.
Additive notation in the Picard presentation means tensor product. The determinant of the coordinate change $t=s^m$ gives
$K_{\mathcal B}=\rho^*K_B\otimes R_1^2R_2^3=R^{-24+8+9}$.

\item An invariant section of $R_j^{r_j}$ on its chart is $s^{r_j}$ times a
holomorphic function of $t=s^{m_j}$. It follows, by the local invariant
comparison on overlaps, that
$\rho_*(\rho^*\OO(k)R_1^{r_1}R_2^{r_2})=\OO(k)$.
In characteristic zero taking finite-group invariants is exact: the
averaging idempotent splits the invariant submodule. Thus there are no
higher direct images for $\rho$, and \eqref{fam:eq:ab-root-degree} follows. This proves the
cohomology comparison with its canonical invariant map. Finally every monomial of degree $n$ is
$x^{r_2}y^{r_1}$ times a homogeneous polynomial of degree $d(n)$ in
$x^4,y^3$. This gives the stated count of sections. The negative-degree
\v Cech complex on $x\ne0$, $y\ne0$ has cokernel represented by
$x^{-a}y^{-b}$ with $a,b\ge1$ and $3a+4b=-n$; replacing $(a,b)$ by
$(a-1,b-1)$ proves the $h^1$ formula. \end{enumerate}
\end{proof}

\begin{definition}[Theta, period and height conventions]
\hypertarget{tag-PI03}{}\label{I:def:theta-height}
\StatementRecord{I:def:theta-height}{proved}{PI03}{}
The theta variable has periods $(1,\tau)$. With $q=e^{2\pi i\tau}$,
define
\[
 N(z,\tau)=\left|\frac{\vartheta_1(z,\tau)}{\eta(\tau)}\right|
                 e^{-\pi(\Im z)^2/\Im\tau},\quad
 \Delta_0=q\prod_{n\ge1}(1-q^n)^{24},\quad t=j(\tau)/1728.
\]
Then $t=E_4^3/(1728\Delta_0)$ and
$t-1=E_6^2/(1728\Delta_0)$. The arithmetic Weierstrass equation is
\[
 E_t:Y^2=X^3-\tfrac34t^3(t-1)X+\tfrac14t^4(t-1)^2,
 \qquad P_t=(0,t^2(t-1)/2),\qquad \omega=dX/(2Y).
\]
Its discriminant is $\Delta_E=27t^8(t-1)^3$.
The lattice coordinates are $X=\wp(u)$, $Y=\wp'(u)/2$;
for $\Lambda=\mathbb Z+\tau\mathbb Z$,
$g_2^3-27g_3^2=(2\pi)^{12}\eta(\tau)^{24}$.
We use the height of $(O)$,
$\widehat h_O(P)=\frac12\lim n^{-2}h(x(nP))$;
the surface self-pairing and rank-one BSD regulator are
$2\widehat h_O$. The real period is
$\Omega=\int_{E(\mathbb R)}|\omega|$, summed over both real
components when there are two. The archimedean local height is
$-\log N$. At finite places $|p|_v=p^{-1}$ and
$n_v=[K_v:\mathbb Q_p]$; complex embeddings are counted separately.

Quasi-periods satisfy $\zeta(z+\omega_j)=\zeta(z)+\eta_j$ and
$\omega_1\eta_2-\omega_2\eta_1=-2\pi i$; $\eta(u)$ denotes their
real-linear extension. In the critical-point calculation we use
$\chi=(\tau-i)/(\tau+i)$, $z=k\chi^2$, and
$\beta_{\rm cusp}=\log(5+2\sqrt6)/(2\pi)$.
\end{definition}

\begin{lemma}[The real period determinant and the fibre class]
\hypertarget{tag-PI04}{}\label{I:period-determinant}
\StatementRecord{I:period-determinant}{proved}{PI04}{}
Put $y=\Im\tau$, $m=\Im\mu$, $r=\Im\beta_c$, and
$D_c=r-6m^2/y=D_{\rm part}+\Im c$.
The real period determinant in the row order
$(\Re\zeta_1,\Im\zeta_1,\Re\zeta_2,\Im\zeta_2)$ is $yD_c$.
For $D_c<0$, the integral fibre class
$\eta_F=u\wedge w+6\gamma\wedge\delta$ satisfies
\[
 \int_{F,\operatorname{or}_{\mathbb C}}\eta_F^2=-12,
 \qquad
 \Gamma=\frac{i}{2D_c}
       \left(d\bar\zeta_2-\frac myd\bar\zeta_1\right)
       =\gamma^{0,1}.
\]
Moreover $\eta_F\wedge\Gamma\wedge d\zeta_2$ equals the marked
volume form $\gamma\wedge u\wedge w\wedge\delta$.
\end{lemma}
\begin{proof}
Writing $a=\Re\mu$, $x=\Re\tau$, $b=\Re\beta_c$, the real matrix is
\[
 \begin{pmatrix}6a&x&1&0\\6m&y&0&0\\b&a&0&1\\r&m&0&0\end{pmatrix}.
\]
Its determinant is $yr-6m^2=yD_c$. Also
$\eta_F^2=12\gamma\wedge u\wedge w\wedge\delta$.
The negative determinant reverses the marked orientation relative
to the complex orientation, giving the stated integral.
The imaginary parts of the two period rows give
$\gamma=(\Im d\zeta_2-(m/y)\Im d\zeta_1)/D_c$;
taking its $(0,1)$ part gives $\Gamma$.
Substitution of the same two rows into the exterior product gives
the final identity, whose coefficient is one.
\end{proof}

\begin{definition}[The Picard coordinate]
\hypertarget{tag-PI05}{}\label{I:def:Picard-coordinate}
\StatementRecord{I:def:Picard-coordinate}{proved}{PI05}{}
Let $\mathfrak g=g_\gamma$ be the coherent Leray class with fibre
image $\gamma^{0,1}$ and put $M_z=\exp(z\mathfrak g)$, where
$\exp$ is the sheaf map $z\mapsto e^{2\pi iz}$.
We normalize $\int_B\omega_{FS}=1$ and
$c_{1,BC}(L)=[iF_L/(2\pi)]$; the unnormalized Chern--Ricci class
is $2\pi c_{1,BC}(T_X)$.
\end{definition}

\begin{remark}[Special-fibre and Picard coordinates]
\hypertarget{tag-PI06}{}\label{I:special-coordinate-values}
\StatementRecord{I:special-coordinate-values}{proved}{PI06}{}
Propositions~\ref{ari:C.1.2}, \ref{fam:src:A.4.1} and
\ref{fam:src:B.1.4} give
\[
 \begin{array}{c|c|c|c}
 &\kappa&\text{upper-half-plane period}&\text{nome}\\\hline
 S_1&c+\frac12-\frac{3i}{2\pi}\log48&-\kappa_1&
 q_1=-w^2/48^3\\
 S_2&c+\frac32-\frac{3i}{\pi}\log6&-\kappa_2/2&
 q_2=iw/216.
 \end{array}
\]
In particular $q_1=(27/64)q_2^2$ and the full compact-family domain
is $0<|w|<216$, equivalently $0<|u_{\rm Engel}|<46656$.
Proposition~\ref{inv:C3.2} fixes the collar orientation and gives
$\exp(12g_\gamma)=f^*\mathcal O_B(-p)$, where
$p=12l_0-4l_1-3l_2$. For the standard data $p=-1$,
$\exp(12\mathfrak g)=f^*\mathcal O_B(1)$,
$\mathcal O_X(S_1-S_2)=\exp(\mathfrak g)$, and $[G_0]=12\mathfrak g$.
The prescribed Chern--Ricci form in Proposition~\ref{inv:C3.6}
is $\pi f^*\omega_{FS}$.
\end{remark}

\begin{lemma}[The exact comparison of the two Weierstrass equations]
\hypertarget{tag-0002}{}\label{I:weierstrass-comparison}
\StatementRecord{I:weierstrass-comparison}{proved}{0002}{\SeriesRef{ari:C.1.1}}
Put $C=t^3(t-1)$ and $D=t^4(t-1)^2$, with $t\ne0,1$.
The pointed map
\[
 \Phi:E_t^{\rm ar}\longrightarrow E_t^{\rm geom},\qquad
 (x,y)\longmapsto(2x,-2\sqrt2y)
\]
is an isomorphism from
$y^2=x^3-3Cx/4+D/4$ to $Y^2=X^3-3CX+2D$.
It sends $(0,t^2(t-1)/2)$ to
$(0,-\sqrt2t^2(t-1))$ and preserves the normalized theta norm
under the induced pointed elliptic-curve identification.
The discriminants satisfy $\Delta_{\rm geom}=64\Delta_{\rm ar}$.
\end{lemma}
\begin{proof}
\begin{enumerate}[label=\arabic*.]
\item Substitution gives
$Y^2=8y^2=8x^3-6Cx+2D=X^3-3CX+2D$.
The inverse is $(X,Y)\mapsto(X/2,-Y/(2\sqrt2))$ and both maps
extend over the unique point at infinity. The displayed section
identity follows by the same substitution.

The original equation in \cite[Remark~3.12, p.~19]{AC2026} is
\[
 y_0^2=4x_0^3-3Cx_0-D,
 \qquad P_0=(0,i t^2(t-1)).
\]
Its exact pointed map to $E_t^{\rm ar}$ is
$(x_0,y_0)\mapsto(-x_0,y_0/(2i))$, with inverse
$(x,y)\mapsto(-x,2iy)$; its composite to $E_t^{\rm geom}$ is
$(x_0,y_0)\mapsto(-2x_0,i\sqrt2y_0)$.
These maps carry the displayed sections to one another.  With
$\omega_0=dx_0/y_0$, the pullback of $dx/(2y)$ is $-i\omega_0$;
transporting the period basis and elliptic logarithm by this scalar
therefore gives the same normalized theta norm by the calculation below.

\item Set $a=-\sqrt2$. In Weierstrass uniformization the equations
are related by $X=a^2x$, $Y=a^3y$ and
$\Lambda_{\rm geom}=a^{-1}\Lambda_{\rm ar}$.
The induced elliptic logarithm map is $z\mapsto a^{-1}z$.
Writing $\Lambda_{\rm ar}=\omega(\Z+\tau\Z)$ shows directly
that the argument $z/\omega$ in the defining theta quotient is
unchanged. Thus the norms agree through this map.
\item The discriminant scales by $a^{12}=64$.
The divisor and leading term of the division polynomial give
$\psi_m(\Phi(Q))=a^{m^2-1}\psi_m(Q)$.
The two absolute scaling factors in
$|\Delta_E|^{(m^2-1)/12}/|\psi_m(Q)|$ cancel exactly.
\end{enumerate}
\end{proof}

\subsection{The elliptic surface and the period quotient}\hypertarget{tag-01PG}{}\label{fam:section:B.1}

\begin{proposition}[The period quotient and its marked elliptic point]\hypertarget{tag-01Q7}{}\label{fam:src:B.1.1}\label{fam:st:B.1.1}
\hypertarget{tag-01Q9}{}\label{fam:src:B.1.3}\label{fam:st:B.1.3}
\StatementRecord{fam:src:B.1.1}{proved}{01Q7}{\cref{I:weierstrass-comparison}; \cite[Remark~3.12, p.~19]{AC2026}}
\StatementRecord{fam:src:B.1.3}{proved}{01Q9}{\cref{I:weierstrass-comparison}; \cite[Remark~3.12, p.~19]{AC2026}}

Let $E_\tau=\mathbb C/(\mathbb Z+\tau\mathbb Z)$, $P=[\mu]$, and $L_\mu=\mathcal O_{E_\tau}(P-O)$. The holomorphic map
\[
 [\xi,\eta]\longmapsto[\xi,e^{2\pi i\eta}]
\]
identifies $\mathbb C^2/\Pi_c\Lambda$ with the quotient of $L_\mu^\times$ by the lift
$A(\xi,W)=(\xi+6\mu,e^{2\pi i\beta_c}W)$. The generator in \cite[Definition~2.3, p.~5]{Engel2026} is its inverse, with translation $-6P$ and multiplier $e^{-2\pi i\beta_c}$. The section $P=[\mu]$ is a generator of the Mordell--Weil group
of the stated pointed Weierstrass surface, and
\[
 \wp_\tau(\mu)=0.
\]
Under the pointed map in the proof it is one of
$P_*=(0,\sqrt2t^2(t-1))$ and $-P_*$.
\end{proposition}

\begin{proof}
\begin{enumerate}[label=\arabic*.,leftmargin=*]
\item First quotient by columns two, three and four. Exponentiation kills column four; columns three and two then give
\[
 (\xi,W)\sim(\xi+1,W)\sim(\xi+\tau,e^{2\pi i\mu}W).
\]
For the odd theta function with periods $1,\tau$, the quotient $\theta(\xi-\mu)/\theta(\xi)$ is a meromorphic section of this line bundle: its factors are $1,e^{2\pi i\mu}$ and its divisor is $P-O$. Sending its divisor section to this quotient is the asserted line-bundle isomorphism. Quotienting the remaining first column gives exactly $A$. On sufficiently small simply connected charts the logarithm of $W$ is a holomorphic inverse; changes of logarithm are column-four periods. Thus the displayed map is a biholomorphism of the two quotient complex spaces, with its inverse explicitly defined by local logarithms. Choosing $A^{-1}$ changes the generator of the same infinite cyclic group and gives precisely Engel's declared inverse-translation convention.

\item The pointed map to the elliptic surface is determined by the branch $f_{\rm sc}$ satisfying
\[
 f_{\rm sc}^4=\frac{12t^3(t-1)}{g_2(\tau)},\qquad
 f_{\rm sc}^6=-\frac{8t^4(t-1)^2}{g_3(\tau)},\qquad
 f_{\rm sc}/t\longrightarrow i\sqrt3/\pi.
\]
The identities $g_2^3/(g_2^3-27g_3^2)=t$ make these equations compatible. Their zero and pole orders give a holomorphic branch on the period cover; the indicated cusp sign fixes it. The map $\xi\mapsto(f_{\rm sc}^2\wp_\tau(\xi),\frac12f_{\rm sc}^3\wp'_\tau(\xi))$ is a pointed isomorphism to
$y^2=x^3-3t^3(t-1)x+2t^4(t-1)^2$. Its factors are $f_{\rm sc}(g_1z)=-f_{\rm sc}/\tau$, $f_{\rm sc}(g_2z)=f_{\rm sc}/\tau$; therefore this isomorphism descends over $B^\circ$.
It pulls $dx/(2y)$ back to $f_{\rm sc}^{-1}d\xi$.

\item Write $C=t^3(t-1)$ and $D=t^4(t-1)^2$.
In \cite[Remark~3.12, p.~19]{AC2026}, the surface
\[
 y_0^2=4x_0^3-3Cx_0-D
\]
has Mordell--Weil generator $P_0=(0,it^2(t-1))$, whose
normalized Abel--Jacobi coordinate is $-\mu$.
Its Weierstrass scale $r$ satisfies
\[
 r^4=3C/g_2(\tau),\qquad r^6=D/g_3(\tau),\qquad
 r^2=\frac{Dg_2(\tau)}{3Cg_3(\tau)}.
\]
Thus $x_0(P_0)=r^2\wp_\tau(-\mu)=0$.
Evenness of $\wp$ gives the asserted zero for the original
affine logarithm. Changing the sign of $r$ composes this
uniformization with $\xi\mapsto-\xi$, and leaves this zero
identity unchanged.
The exact pointed map of \cref{I:weierstrass-comparison},
\[
 (x_0,y_0)\longmapsto(-2x_0,i\sqrt2y_0),
\]
sends $P_0$ to $-P_*$ and $-P_0$ to $P_*$.
Both visible sections therefore generate the Mordell--Weil group
of the geometric equation. The preceding $f_{\rm sc}$ map sends
$[\mu]$ to one of these two sections, proving the generator claim.
The Shioda self-pairing is
$\langle P,P\rangle=2-4/3-1/2=1/6$
\cite[proof of Lemma~2.1, pp.~3--4]{Engel2026}; the two sign choices
have the same pairing. The chosen sign in our family is fixed by
the cusp calculation in \cref{fam:src:B.1.4}.

\item The nodal comparison needed for that sign is explicit.
Put $v=1/t$, $X=v^2x$, $Y=v^3y$; the geometric equation becomes
$Y^2=X^3-3(1-v)X+2(1-v)^2$.
At $v=0$, the chosen scale gives
\[
 X=1+\frac{12U}{(U-1)^2},\qquad
 Y=\frac{12\sqrt3\,U(U+1)}{(U-1)^3},\qquad
 U=e^{2\pi i\xi}.
\]
On the smooth locus away from $O$ its inverse is
$U=(Y/(X-1)+\sqrt3)/(Y/(X-1)-\sqrt3)$;
at $O$ a local parameter is $X/Y$.
The normally convergent multiplicative expansion of $\wp$
gives the corresponding holomorphic map for small $v$.
\end{enumerate}
\end{proof}

\begin{proposition}[The compactness radius]\hypertarget{tag-01Q8}{}\label{fam:src:B.1.2}\label{fam:st:B.1.2}
\StatementRecord{fam:src:B.1.2}{proved}{01Q8}{the definitions and constructions stated in the proof}

The exact set of nonzero parameters for which every ordinary fibre and both normalized finite root fibres of Engel's period-reference quotient are compact complex two-tori is
\[
 0<|u|<e^{-2\pi M}.
\]
This is the parameter domain for the compact family with both smooth finite root models. The condition on ordinary fibres alone at equality is determined by whether a maximizing point of $D_{\rm part}$ is ordinary; the maximizing locus is determined in \cref{ari:C.8.15}. For an independent reference, the radius is $|\lambda_0|e^{-2\pi M}$.
\end{proposition}


\begin{proof}
\begin{enumerate}[label=\arabic*.,leftmargin=*]
\item The function $N(\xi,W)=|W|\exp(2\pi\operatorname{Im}\mu\,\operatorname{Im}\xi/\operatorname{Im}\tau)$ is invariant under both line-bundle period identifications. The inverse lift multiplies it by $e^{2\pi D_c}$. More directly the real determinant of the four period columns is a nonzero sign times $(\operatorname{Im}\tau)D_c$. Thus the four columns generate a full lattice exactly when $D_c\ne0$, and then the quotient is a compact torus by a closed fundamental parallelepiped. At $D_c=0$ their real span has dimension three; a rank-four lattice in $\mathbb R^4$ cannot result, so the quotient cannot be a compact complex two-torus. The function $D_{\rm part}$ descends continuously to the connected base with the cusp removed and tends to $-\infty$ at the cusp. Its image is therefore $(-\infty,M]$: the maximum is attained, and the intermediate value theorem fills all smaller values. Consequently $D_c$ is nonzero at every base point exactly when $\operatorname{Im}c<-M$. Since $|u|=e^{2\pi\operatorname{Im}c}$, the assertion follows, with the indicated reference scaling.
\end{enumerate}
\end{proof}



\begin{proposition}[The reference linearization]\hypertarget{tag-01QA}{}\label{fam:src:B.1.4}\label{fam:st:B.1.4}
\StatementRecord{fam:src:B.1.4}{proved}{01QA}{\SeriesRef{fam:src:B.1.3}, \SeriesRef{fam:src:B.1.1}; \cite[Lemma~2.1, pp.~3--4]{Engel2026}; \cite[Corollary~2.2, p.~4]{Engel2026}}

The line-bundle comparison of the period family with Engel's family over $B^\circ$ transports its inverse period lift to a global section of $\operatorname{Hom}(t_{6P}^*M_S,M_S)$ with divisor $F_\infty$. Take this transported section for $c=0$ as the reference linearization. Then $\lambda_0=1$ exactly and $u=e^{-2\pi ic}=w^2$. For another independently fixed reference the precise constant is its scalar ratio with this section. With the cusp branch $e^{2\pi i\mu}\to-(5+2\sqrt6)$ and $f_{\rm sc}/t\to i\sqrt3/\pi$, the section is $P=(0,-\sqrt2t^2(t-1))$.
\end{proposition}


\begin{proof}
\begin{enumerate}[label=\arabic*.,leftmargin=*]
\item The zero identity $\wp_\tau(\mu)=0$, proved in \ref{fam:src:B.1.3}, identifies $P$ with the signed visible section. With this cusp branch and the specified scale, substitution of $U=-(5+2\sqrt6)$ in the nodal uniformization of \ref{fam:src:B.1.1} gives $Y=-\sqrt2$, so $P=(0,-\sqrt2t^2(t-1))$.

\item To compare the line bundles as families, use their common divisor $P-O$ and rigidify at the zero section. Put $F_\theta(\xi)=\theta(\xi-\mu)/\theta(\xi)$, $r_\theta=\theta(-\mu)/\theta'(0)$, and let $s_{P-O}$ be the canonical meromorphic divisor section. The map
\[
 [\xi,W]\longmapsto\frac{2r_\theta}{f_{\rm sc}}
        W\frac{s_{P-O}(\xi)}{F_\theta(\xi)}
\]
is a holomorphic line-bundle isomorphism, including at $P,O$ after cancellation of the equal divisors. At $O$ it sends the chosen frame to $dx/y$, since $F_\theta\sim r_\theta/\xi$ and $x/y\sim-\xi/f_{\rm sc}$ with the corresponding canonical divisor frame. Two such maps that agree at $O$ differ by a holomorphic scalar on each compact elliptic fibre and hence agree everywhere. This proves descent under the specified period identifications.

\item \citet[Lemma~2.1, pp.~3--4]{Engel2026} identifies $\operatorname{Hom}(t_{6P}^*M_S,M_S)$ with $\pi_S^*\mathcal O_B(1)$. The lemma applies because the preceding pointed cubic map and Mordell--Weil calculation identify his $(S,O,P)$ with ours, including the chosen sign. To determine which section is transported, it suffices to compute its order at one generic smooth point of each special fibre, since a section of this pulled-back line is a base section. At a finite centre write $6\mu=p+r\tau+\delta$ with $(p,r)=(4,-2),(3,-3)$ and $\delta/f_{\rm sc}$ holomorphic; this divisibility follows by linearizing the affine laws, whose homogeneous orders are two and one. Evaluate at $\xi=f_{\rm sc}z$ for generic fixed $z$. The factor $(2r_\theta/f_{\rm sc})/F_\theta(f_{\rm sc}z)$ tends to $2z$, a unit. Restoring the lattice shift $p+r\tau$ at the target changes the fibre multiplier to $e^{-2\pi i(\beta_c-r\mu)}$, whose finite limits are $e^{-2\pi i\kappa_1}$ and $e^{-2\pi i\kappa_2}$. The source and target theta ratios remain units. Therefore the base section has order zero at both finite points. At the cusp $f_{\rm sc}$ has a simple pole and $r_\theta$ has a nonzero limit; the two factors $2r_\theta/f_{\rm sc}$ cancel between source and target. The remaining multiplier is
$e^{-2\pi i\beta_c}=q e^{-2\pi i(\beta_c+\tau)}$, of exact order one. Thus the transported section has divisor $F_\infty$. \citet[Corollary~2.2, p.~4]{Engel2026} identifies it with a nonzero scalar multiple of his reference, since two base sections with the same divisor have constant ratio. Changing $\beta$ by $c$ multiplies it exactly by $e^{-2\pi ic}$. With our stated period-reference normalization, this scalar is one. For any independently chosen reference $\phi_{\rm ref}$, the scalar is instead the precisely defined ratio $\lambda_0=\phi_{\rm part}/\phi_{\rm ref}$, evaluated on any ordinary fibre. Scaling the reference by $a\in\mathbb C^*$ changes $\lambda_0$ to $\lambda_0/a$.
\end{enumerate}
\end{proof}


\section{Normalized periods and the real determinant}\hypertarget{tag-01PK}{}\label{fam:section:C}

\subsection{The normalized theta period}\hypertarget{tag-01PL}{}\label{fam:section:C.1}
\begin{definition}[Theta functions and the cusp branch]
\hypertarget{tag-IA21}{}\label{I:theta-series-cusp-branch}
\StatementRecord{I:theta-series-cusp-branch}{proved}{IA21}{}
The odd Jacobi function and the Dedekind function are
\[
 \begin{aligned}
 \vartheta_1(v,\tau)&=2\sum_{n\geq0}(-1)^n
 e^{\pi i\tau(n+1/2)^2}\sin((2n+1)\pi v),\\
 \eta_D(\tau)=\eta(\tau)&=
 e^{\pi i\tau/12}\prod_{n\geq1}(1-e^{2\pi in\tau}).
 \end{aligned}
\]
Choose the cusp lift reached from $\tau=iT$, $T>1$, on which
$\mu=1/2-iy(T)$ and $0<y(T)<T/2$. Put
$a=5+2\sqrt6$ and $\rho=e^{\pi i/3}$.
\end{definition}

\begin{proposition}[The normalized theta period]\hypertarget{tag-01QH}{}\label{fam:src:C.1.1}\label{fam:st:C.1.1}
\StatementRecord{fam:src:C.1.1}{proved}{01QH}{\SeriesRef{fam:src:B.1.3}}

For the period section identified in \cref{fam:src:B.1.3}, the particular solution
with \((\beta^{\rm part}+\tau)(p_0)=0\) is
\begin{equation}\hypertarget{tag-01P5}{}\label{fam:eq:cd-beta}
 \beta^{\rm part}=\frac{6i}{\pi}\operatorname{Log}
 \frac{\vartheta_1(\mu,\tau)}{2\sqrt3\,\eta_D(\tau)}.
\end{equation}
Here the logarithm is continued on the simply connected uniformizing
half-plane from its cusp expansion. Its special values are
\[
 \mu_0=\frac12-\frac{i\log a}{2\pi},\qquad
 \mu(z_1)=\frac{2-\rho}{3},\qquad \mu(z_2)=\frac{1-i}{2},
\]
\[
 \beta^{\rm part}(z_2)=i\left(\frac32-\frac3\pi\log6\right),
 \qquad \operatorname{Re}\beta^{\rm part}(z_1)=-\frac12.
\]
In particular \(e^{2\pi i\mu_0}=-a\).
\end{proposition}


\begin{proof}
\begin{enumerate}[label=\arabic*.,leftmargin=*]
\item \Cref{fam:src:B.1.3} identifies the elliptic point represented by \(\mu\) with
one of the sections \((0,\pm\sqrt2t^2(t-1))\). The uniformization map
\((v,\tau)\mapsto(f^2\wp_\tau(v),f^3\wp'_\tau(v)/2)\) therefore gives
\(\wp_\tau(\mu)=0\). No zero of \(\wp\) is a lattice point. Thus the numerator
of \eqref{fam:eq:cd-beta} is nonzero; the denominator is nonzero as well.
A holomorphic logarithm exists on the simply connected cover.

\item The cusp value follows from the normally convergent Fourier
expansion:
\begin{equation}\hypertarget{tag-01P8}{}\label{fam:eq:cd-wp}
 \frac{\wp_\tau(v)}{\pi^2}
 =\csc^2(\pi v)-\frac13+
 8\sum_{n\geq1}\frac{nq^n}{1-q^n}(1-\cos(2\pi nv)),
 \qquad q=e^{2\pi i\tau}.
\end{equation}
This identity follows by differentiating the logarithm of the Jacobi
product twice and using
\(E_2=1-24\sum nq^n/(1-q^n)\). Boundedness of the chosen elliptic logarithm
at the cusp permits passage to the limit. We obtain
\(\sin^2(\pi\mu_0)=3\), equivalently
\(u^2+10u+1=0\), \(u=e^{2\pi i\mu_0}\).
On the chosen imaginary-axis branch \(\operatorname{Im}\mu<0\), so \(u=-a\).
On the rectangular
elliptic curve with \(\tau=iT\), \(T>1\), \(\wp\) is real on the segment
\(1/2-iy\), \(0\leq y\leq T/2\). Its endpoint values are the largest and
middle real roots of its cubic, respectively positive and negative.
Its derivative has no zero in the open segment, since the zeros of
\(\wp'\) are the three nonzero half-periods. There is exactly one zero
there; this gives the indicated continuation. Consequently
\(2\sin(\pi\mu_0)=2\sqrt3\).

\item Put \(R=\vartheta_1/\eta_D\). Poisson summation for the defining Gaussian
series and its index shift give
\[
 \begin{aligned}
 R(v,\tau+1)&=e^{\pi i/6}R(v,\tau),\\
 R(v/\tau,-1/\tau)&=-i e^{\pi iv^2/\tau}R(v,\tau),\\
 R(1-v,\tau)&=R(v,\tau),\qquad R(v+1,\tau)=-R(v,\tau).
 \end{aligned}
\]
Inserting the transformed arguments gives
\[
 \begin{aligned}
 R(\mu\circ g_1,\tau\circ g_1)
 &=e^{-\pi i/3}e^{\pi i(1-\mu)^2/\tau}R(\mu,\tau),\\
 R(\mu\circ g_2,\tau\circ g_2)
 &=i e^{\pi i\mu^2/\tau}R(\mu,\tau).
 \end{aligned}
\]
Multiplication of the logarithms by \(6i/\pi\) gives exactly
\(2-6(1-\mu)^2/\tau\) and \(-3-6\mu^2/\tau\), modulo constants in
\(12\mathbb Z\). These constants vanish: each defect is constant on the
connected cover, and evaluation at the corresponding fixed point makes
both its two sides zero. The fixed-point equations themselves give
\(\mu(z_1)=(2-\rho)/3\) and \(\mu(z_2)=(1-i)/2\).
The Jacobi product has cusp expansion
\(R(\mu,\tau)=2q^{1/12}\sin(\pi\mu)(1+O(q))\).
Equation \eqref{fam:eq:cd-beta} consequently gives
\(\beta^{\rm part}+\tau\to0\). The difference from the specified particular
solution is an invariant holomorphic function extending across all three
points of \(B\), hence a constant, and the cusp normalization makes it zero.

\item At \(\tau=i\), the half-period translation formula gives
\(\vartheta_1((1-i)/2,i)=e^{\pi/4}\vartheta_3(0,i)\).
The identities \(2\eta_D^3=\vartheta_2\vartheta_3\vartheta_4\),
\(\vartheta_2(0,i)=\vartheta_4(0,i)=2^{-1/4}\vartheta_3(0,i)\)
give \(R=e^{\pi/4}\sqrt2\), proving the displayed value of \(\beta\).
At \(\rho\), \(q\) is negative real, \(\mu=1/2-i/(2\sqrt3)\), and the
product for \(R\) has phase \(\pi/12\), with all its remaining factors
positive real. This gives \(\operatorname{Re}\beta=-1/2\).
\end{enumerate}
\end{proof}



\subsection{The invariant norm}\hypertarget{tag-01PM}{}\label{fam:section:C.2}
\begin{proposition}[The norm and the real determinant]
\hypertarget{tag-01QJ}{}\label{fam:src:C.2.1}\label{fam:st:C.2.1}
\StatementRecord{fam:src:C.2.1}{proved}{01QJ}{\cref{fam:src:C.1.1,I:def:theta-height,I:period-determinant}}
For the conventions of \cref{I:def:theta-height,I:def:family-parameter},
\begin{equation}\hypertarget{tag-01P6}{}\label{fam:eq:cd-norm}
 D_{\rm part}=\frac6\pi\log\frac{N(\mu,\tau)}{2\sqrt3},
 \qquad \mathcal N(v,\tau)=N(v,\tau).
\end{equation}
The norm is preserved by integral translations and sign change of $v$,
and by the lattice isomorphism
\[
 [v]\longmapsto[v/(c\tau+d)],\qquad
 \tau\longmapsto\frac{a\tau+b}{c\tau+d},
 \qquad \begin{pmatrix}a&b\\c&d\end{pmatrix}\in\mathrm{SL}_2(\mathbb Z).
\]
At the cusp, $D_{\rm part}=-\Im\tau+o(1)$.
\end{proposition}
\begin{proof}
The imaginary part of \eqref{fam:eq:cd-beta} is
$(6/\pi)\log|\vartheta_1(\mu,\tau)/(2\sqrt3\eta_D(\tau))|$.
Subtracting $6(\Im\mu)^2/\Im\tau$ gives \eqref{fam:eq:cd-norm}.
Thus the two notations $\mathcal N$ and $N$ denote the identical
function of \cref{I:def:theta-height}.
The theta transformation formulas in \cref{fam:src:C.1.1} have
exactly the exponential moduli cancelled by this Gaussian; the
theta and eta square-root factors cancel in their quotient.
The displayed lattice map also sends $\wp_\tau(v)=0$ to the
zero equation on the transformed lattice, since $\wp$ has weight two.
Finally $\beta^{\rm part}+\tau\to0$ and the bounded cusp limit of
$\mu$ give the stated asymptotic directly from the real determinant.
\end{proof}






\section{Special values, local expansions and torsion}
\subsection{The normalized Green function and special arguments}
\begin{definition}[The logarithmic norm on a lattice]
\hypertarget{tag-IA22}{}\label{I:logarithmic-lattice-norm}
\StatementRecord{I:logarithmic-lattice-norm}{proved}{IA22}{}
Write $\rho=(1+i\sqrt3)/2$ and $T=\operatorname{Im}\tau$.
For the norm of \cref{I:def:theta-height}, put
\[
 U_\tau(z)=\log N(z,\tau)
 =\log\left|\frac{\vartheta_1(z,\tau)}{\eta(\tau)}\right|
       -\frac{\pi(\operatorname{Im}z)^2}{T}.
\]
The theta variable has periods $1,\tau$; a convention with
first period $\pi$ therefore uses the argument $\pi z$.
For $\Lambda=\omega(\mathbb Z+\tau\mathbb Z)$ and $z\notin\Lambda$
set $U_\Lambda(z)=U_\tau(z/\omega)$.
\end{definition}
\begin{proof}[Well-definedness]
Write $v=z/\omega$. An oriented basis change has
\[
 \tau'=\frac{a\tau+b}{c\tau+d},\qquad
 \omega'=\omega(c\tau+d),\qquad v'=\frac v{c\tau+d},
 \qquad ad-bc=1.
\]
The theta transformations
\cite[Equations~20.7.26 and 20.7.30]{DLMFTheta2026}, with
the first argument multiplied by $\pi$, and the eta law
\cite[Equations~23.18.5--23.18.6]{DLMFTheta2026} give the
two modular generators; the eta translation law also follows
directly from its product. Composition gives
\[
 \log\left|\frac{\vartheta_1(v',\tau')}{\eta(\tau')}\right|
 =\log\left|\frac{\vartheta_1(v,\tau)}{\eta(\tau)}\right|
   -\pi\Im\frac{cv^2}{c\tau+d}.
\]
The identity
\[
 \frac{(\Im v')^2}{\Im\tau'}
 =\frac{(\Im v)^2}{\Im\tau}
     -\Im\frac{cv^2}{c\tau+d}
\]
cancels this term, so $U_{\tau'}(v')=U_\tau(v)$.
This proves independence of the basis of $\Lambda$.
For a homothety $\lambda\ne0$, the same normalized argument gives
$U_{\lambda\Lambda}(\lambda z)=U_\Lambda(z)$.
\end{proof}

\begin{proposition}[Distribution under isogenies]\hypertarget{tag-000I}{}\label{ari:C.1.1}
\StatementRecord{ari:C.1.1}{proved}{000I}{the definitions and constructions stated in the proof}

Let $\Lambda,\Lambda'\subset\mathbb C$ be lattices and let
$\alpha\ne0$ satisfy $\alpha\Lambda\subset\Lambda'$. Write
$\phi:\mathbb C/\Lambda\to\mathbb C/\Lambda'$ for $z\mapsto\alpha z$,
$K=\alpha^{-1}\Lambda'/\Lambda$, and
$d=[\Lambda':\alpha\Lambda]$. Then
\begin{equation}\hypertarget{tag-003O}{}\label{ari:eq:cs-isogeny}
 U_{\Lambda'}(\alpha z)=\sum_{v\in K}U_\Lambda(z+v).
\end{equation}
The equality holds away from the zeros and as an equality of locally
integrable functions at them. If $\Delta_\Lambda=g_2(\Lambda)^3-
27g_3(\Lambda)^2$, then
\[
 \prod_{0\ne v\in K}N_\Lambda(v)
 =|\alpha|\left|\frac{\Delta_{\Lambda'}}{\Delta_\Lambda}\right|^{1/12}.
\]
In particular, for every nonzero endomorphism of one elliptic curve,
$\prod_{v\in\ker\alpha,\ v\ne0}N(v)=|\alpha|$.
\end{proposition}


\begin{proof}
\begin{enumerate}[label=\arabic*.,leftmargin=*]
\item The Jacobi product, with $q=e^{2\pi i\tau}$ and
$X=e^{2\pi iz}$, gives
\[
 \vartheta _1(z,\tau)/\eta(\tau)
 =i q^{1/12}X^{-1/2}(1-X)
     \prod_{m\ge1}(1-q^mX)(1-q^mX^{-1}),
\]
where a harmless fixed unit changes with the theta convention. For
$0<y<T$, integration of the logarithm over $z=x+iy$, $0\le x\le1$,
gives $-\pi T/6+\pi y-\pi y^2/T$. Each product factor has mean logarithm
zero by the absolutely convergent logarithmic series; the endpoint
singularities are integrable. Integrating in $y$ proves
$\int U_\Lambda\,d\mu_\Lambda=0$ for Haar probability measure.
Moreover
\[
 \Delta_{\rm Eucl}U_\Lambda
   =2\pi\delta_0-\frac{2\pi}{\operatorname{area}\Lambda},\qquad
 U_\Lambda(z)=\log|z|+\frac1{12}\log|\Delta_\Lambda|+o(1).
\]
The second formula follows from
$\vartheta _1'(0,\tau)=2\pi\eta(\tau)^3$ and
$\Delta_{\mathbb Z+\tau\mathbb Z}=(2\pi)^{12}\eta(\tau)^{24}$,
then scaling by $\omega$. The two sides of
\eqref{ari:eq:cs-isogeny} have the same logarithmic masses at $K$ and the
same constant Laplacian, since
$d=|\alpha|^2\operatorname{area}\Lambda/
\operatorname{area}\Lambda'$. Their difference is a smooth harmonic
function on the compact torus, hence constant. Both means vanish:
translations preserve Haar measure and an isogeny sends Haar
probability measure to Haar probability measure. The constant is zero.
Taking $z\to0$ and cancelling the logarithmic singularity gives the
product.
\end{enumerate}
\end{proof}
\begin{proposition}[Phases and the two periods]\hypertarget{tag-000J}{}\label{ari:C.1.2}
\StatementRecord{ari:C.1.2}{proved}{000J}{\SeriesRef{ari:C.1.1}}

With the branch continued from the specified cusp, the exact values are
\begin{align*}
 \frac{\vartheta _1((2-\rho)/3,\rho)}{\eta(\rho)}
   &=e^{\pi i/12}\,3^{1/4}e^{\pi/(6\sqrt3)},\\
 \frac{\vartheta _1((1-i)/2,i)}{\eta(i)}
   &=\sqrt2\,e^{\pi/4},\\
 \beta^{\rm part}(z_1)
   &=-\frac12+i\left(\frac1{\sqrt3}-\frac3{2\pi}\log48\right),\\
 \beta^{\rm part}(z_2)
   &=i\left(\frac32-\frac3\pi\log6\right).
\end{align*}
Consequently $N(z_1)=3^{1/4}$, $N(z_2)=\sqrt2$,
$D_{\rm part}(z_1)=-3\log48/(2\pi)$, and
$D_{\rm part}(z_2)=-3\log6/\pi$.
\end{proposition}


\begin{proof}
\begin{enumerate}[label=\arabic*.,leftmargin=*]
\item Multiplication by $1+\rho$ has matrix
$\left(\begin{smallmatrix}1&-1\\1&2\end{smallmatrix}\right)$,
its kernel is $\{0,\mu,-\mu\}$, and $(1+\rho)\mu=1$ for
$\mu=(2-\rho)/3$. \Cref{ari:C.1.1} gives $N(\mu)^2=\sqrt3$.
At $i$, multiplication by $1+i$
has kernel $\{0,(1-i)/2\}$ and absolute value $\sqrt2$; the same
identity gives $N((1-i)/2,i)=\sqrt2$.

\item At $\rho$, $\mu=1/2-i/(2\sqrt3)$ and $q=-e^{-\pi\sqrt3}$. In
\[
 R(z,\tau)=2e^{\pi i\tau/6}\sin(\pi z)
 \prod_{m\ge1}\bigl(1-2q^m\cos(2\pi z)+q^{2m}\bigr),
\]
$\sin(\pi\mu)$ is positive real. The factors are
$(1+q^m e^{\pi/\sqrt3})(1+q^m e^{-\pi/\sqrt3})$, both positive:
when $m$ is odd the larger subtracted number is at most
$e^{-2\pi/\sqrt3}<1$. Thus the phase is exactly $\pi/12$ modulo
$2\pi$. There is no winding ambiguity in the continued logarithm.
Indeed continue from the cusp along $\tau=1/2+iT$,
$T\ge\sqrt3/2$. Real conjugation of the rhombic lattice keeps the
chosen zero on $\mu=1/2-iy$, $0<y<T$. It cannot leave that component:
a collision of the two zeros would require $g_3=0$, which does not
occur on this ray. Indeed, writing $q=-r$, $0<r<1/200$,
the tail after the first term of the $E_6$ series has absolute value
at most $504\cdot32r^2/((1-r^2)(1-8r))<1/2$, while its
first term is positive. Thus $E_6>1/2$. At $y=0$ the point is a nonzero half-period, so a zero there
would force $g_3=0$. At $y=T$ it is the lattice point
$1-\tau$, where $\wp$ has a pole; a zero cannot reach that boundary. Each displayed product factor remains
positive because $e^{-2\pi T}e^{2\pi y}<1$. The cusp phase tends to
$\pi/12$, so the continued argument has that value. On the imaginary
ray $\tau=iT$, $T\ge1$, the same branch is
$\mu=1/2-iy$ with $0<y\le T/2$; every product factor is positive,
so its argument at $i$ is zero. The Gaussian factors are respectively
$e^{\pi/(6\sqrt3)}$ and $e^{\pi/4}$. Substitution in the definition of
$\beta^{\rm part}$ proves all four identities.
\end{enumerate}
\end{proof}

\begin{corollary}[Difference of the finite invariant periods]
\hypertarget{tag-01QI}{}\label{fam:src:C.1.2}\label{fam:st:C.1.2}
\hypertarget{tag-000K}{}\label{ari:C.1.3}
\StatementRecord{fam:src:C.1.2}{proved}{01QI}{\cref{ari:C.1.2}}
Put $\kappa_1^{\rm part}=\beta^{\rm part}(z_1)+2\mu(z_1)$ and
$\kappa_2^{\rm part}=\beta^{\rm part}(z_2)+3\mu(z_2)$. Then
\[
 \kappa_1^{\rm part}-\kappa_2^{\rm part}
   =-1-\frac{3i}{2\pi}\log\frac43,
 \qquad
 -\frac15<\Im(\kappa_1^{\rm part}-\kappa_2^{\rm part})
             <-\frac1{10}.
\]
\end{corollary}
\begin{proof}
Substitution of the exact values in \cref{ari:C.1.2} gives
$\kappa_1^{\rm part}=\frac12-\frac{3i}{2\pi}\log48$ and
$\kappa_2^{\rm part}=\frac32-\frac{3i}{\pi}\log6$.
Their difference is the displayed expression.
Integrating $1/(1+x)$ on $0\le x\le1/3$ gives
$1/4<\log(4/3)<1/3$; together with $3<\pi<15/4$,
these inequalities give the stated strict bounds.
\end{proof}


\subsection{Division identities and the entire torsion locus}
\begin{remark}[The arithmetic model in the division identities]
\hypertarget{tag-IA23}{}\label{I:division-model-convention}
\StatementRecord{I:division-model-convention}{proved}{IA23}{}
Use the model and point of \cref{I:def:theta-height}, written as
\[
 E_t:\quad y^2=x^3+A(t)x+B(t),\qquad
 A=-\tfrac34t^3(t-1),\quad B=\tfrac14t^4(t-1)^2,
 \quad P=(0,t^2(t-1)/2).
\]
Its discriminant is $\Delta_E(t)=27t^8(t-1)^3$.
Under $(x,y)\mapsto(\lambda^2x,\lambda^3y)$ the period lattice
and its logarithm change by $(\Lambda,z)\mapsto
(\lambda^{-1}\Lambda,\lambda^{-1}z)$ for the invariant differential.
The identity in \cref{I:logarithmic-lattice-norm} therefore preserves $N$.
\end{remark}

\begin{cited}[Siegel functions and local heights]\hypertarget{tag-000P}{}\label{ari:C.4.1}
\StatementRecord{ari:C.4.1}{proved}{000P}{\citet[Remark~2.2 and Proposition~2.3]{Shin2015}}

For $a=(a_1,a_2)\in\mathbb R^2\setminus\mathbb Z^2$ and
$v=a_1\tau+a_2$, define
\[
 g_a(\tau)=-e^{\pi ia_2(a_1-1)}q^{B_2(a_1)/2}
 (1-e^{2\pi iv})\prod_{r\ge1}
 (1-q^re^{2\pi iv})(1-q^re^{-2\pi iv}),
 \quad B_2(x)=x^2-x+\tfrac16.
\]
Then
\[
 g_a(\tau)=i e^{\pi ia_1v}
       \frac{\vartheta_1(v,\tau)}{\eta(\tau)},\qquad
 N(v,\tau)=|g_a(\tau)|.
\]
When $a\in n^{-1}\mathbb Z^2\setminus\mathbb Z^2$, $n\ge2$,
$g_a^{12n}$ is a modular unit for $\Gamma(n)$ with Fourier coefficients
in $\mathbb Q(\zeta_n)$ in the variable $q^{1/n}$. On every complex
elliptic curve the model-independent Néron local function for $(O)$ is
\[
 \lambda_\infty(Q)=-\log N(Q).
\]
Thus $D_{\rm part}=-(6/\pi)(\lambda_\infty(P)+\log(2\sqrt3))$.
\end{cited}
\begin{proof}
Dividing the defining products gives
$g_a=i e^{\pi ia_1v}\vartheta_1(v,\tau)/\eta(\tau)$.
Since $\Im v=a_1\Im\tau$, its absolute value is exactly $N(v,\tau)$.
For a nonintegral rational pair of denominator $n\ge2$,
\cite[Remark~2.2, p.~3, and Proposition~2.3, pp.~3--4]{Shin2015}
gives the modular-unit assertion. Its product has coefficients in
$\mathbb Q(\zeta_n)$ after taking the indicated power; the exponent
$6nB_2(r/n)$ belongs to $n^{-1}\mathbb Z$.

By \cref{ari:C.1.1}, $-\log N$ has mean zero, the logarithmic pole
at $O$, and the Laplacian of the normalized Green function of $(O)$.
The difference of two functions with these properties is harmonic
on the compact torus, and its mean makes the constant zero.
In the Weierstrass coordinate the normalization is
\[
 -\log N(z)-\frac12\log|x(z)|
       \longrightarrow-\frac1{12}\log|\Delta|.
\]
Under the exact homothety
$(\Lambda,z,x,\Delta)\mapsto(c\Lambda,cz,c^{-2}x,c^{-12}\Delta)$,
$N$ is unchanged and both sides of the last comparison change by
$\log|c|$. This proves the stated model transport and the formula
for $D_{\rm part}$.
\end{proof}


\begin{proposition}[The full division identity]\hypertarget{tag-000Q}{}\label{ari:C.4.2}
\StatementRecord{ari:C.4.2}{proved}{000Q}{\SeriesRef{ari:C.1.1}}

Let $\psi_m$ be the standard division polynomial of a smooth
short Weierstrass model, normalized by $\psi_1=1$, $\psi_2=2y$.
On a general equation
$y^2+a_1xy+a_3y=x^3+a_2x^2+a_4x+a_6$, the exact change
$Y=y+(a_1x+a_3)/2$, followed by
$X=x+(a_2+a_1^2/4)/3$, gives a short equation and
$\psi_2=2Y=2y+a_1x+a_3$.
The division functions below are transported by this isomorphism.
For $m\ge2$ and $Q,mQ\ne O$,
\begin{equation}\hypertarget{tag-003N}{}\label{ari:eq:cs-division}
 N(Q)^{m^2}
       =N(mQ)\frac{|\Delta_E|^{(m^2-1)/12}}{|\psi_m(Q)|}.
\end{equation}
For the displayed family this specializes to
\[
 N(P)^4=3^{3/4}|t-1|^{-1/4}N(2P),\qquad
 N(P)^9=16|t|^{-2/3}N(3P).
\]
If $P_t$ has exact order $n$ and $m\ge2$ satisfies
$m\equiv1$ or $-1\pmod n$, then
\begin{equation}\hypertarget{tag-003P}{}\label{ari:eq:cs-torsionvalue}
 N(P_t)^{12}
       =\frac{|\Delta(t)|}{|\psi_m(P_t)|^{12/(m^2-1)}}.
\end{equation}
The positive value on the right is independent of $m$.
\end{proposition}


\begin{proof}
\begin{enumerate}[label=\arabic*.,leftmargin=*]
\item Use the complex uniformization $x=\wp_\Lambda(z)$,
$y=\wp'_\Lambda(z)/2$. The rational function $\psi_m$ has divisor
$\sum_{0\ne Q\in E[m]}(Q)-(m^2-1)(O)$ and leading Laurent term
$(-1)^{m-1}m z^{1-m^2}$. Both the divisor and the leading Laurent coefficient follow from
$(-1)^{m-1}\sigma_\Lambda(mz)/\sigma_\Lambda(z)^{m^2}$.
Its lattice factors cancel, its zeros are exactly the nonzero
$m$-division points, and its expansion at zero is the displayed
one. The addition formula for $\sigma$ gives the standard
division-polynomial recurrence with initial values $1,2y$, so it is
the stated rational function.

\item The function
\[
 m^2U_\Lambda(z)-U_\Lambda(mz)
                    +\log|\psi_m(Q(z))|
\]
is harmonic after its logarithmic singularities cancel. Its
Laplacian has zero constant term as well, so it is constant.
The expansion in \ref{ari:C.1.1} and the leading term of $\psi_m$ make that
constant $(m^2-1)\log|\Delta_E|/12$. This proves
\eqref{ari:eq:cs-division}. Direct substitution of
\[
 \psi_2(P)=t^2(t-1),\qquad
 \psi_3(P)=-9t^6(t-1)^2/16
\]
gives the two special identities. For $m\equiv\pm1\pmod n$,
$mP=\pm P$ and $N$ is even. Cancelling $N(P)$ in
\eqref{ari:eq:cs-division} gives \eqref{ari:eq:cs-torsionvalue}.
Every right side is the same positive number $N(P)^{12}$, which
proves independence of $m$.
\end{enumerate}
\end{proof}

\begin{proposition}[Algebraicity and fields of division values]\hypertarget{tag-000R}{}\label{ari:C.4.3}
\StatementRecord{ari:C.4.3}{proved}{000R}{\SeriesRef{ari:eq:cs-torsionvalue}}

Let $t_0\in\mathbb C\setminus\{0,1\}$ and suppose $P_{t_0}$ has
exact order $n\ge2$. Then $t_0$ and $N(P_{t_0})$ are algebraic.
Put $K_0=\mathbb Q(t_0,\bar t_0)$,
\[
 W=\Delta(t_0)\Delta(\bar t_0)>0,\qquad
 V=\psi_{n+1}(P_{t_0})\psi_{n+1}(P_{\bar t_0})>0.
\]
With positive real radicals,
\[
 N(P_{t_0})^{12n}=W^{n/2}V^{-6/(n+2)}
 \ \in\ K_0\bigl(W^{1/2},V^{1/(n+2)}\bigr).
\]
More generally, if $m\equiv\pm1\pmod n$, then
\[
 \lambda_\infty(P_{t_0})
 =\frac{\log|\psi_m(P_{t_0})|}{m^2-1}
       -\frac1{12}\log|\Delta(t_0)|.
\]
These formulas determine the height and $N^{12}$ on the entire
torsion locus.
\end{proposition}


\begin{proof}
\begin{enumerate}[label=\arabic*.,leftmargin=*]
\item The division polynomial $\psi_n(P_t)$ is a nonzero rational
polynomial in $t$. To check nonvanishing for every $n$, a hypothetical
torsion identity would express the elliptic logarithm as
$a\tau+b$ with constant rational $a,b$. Its bounded cusp limit
forces $a=0$, but that limit has nonzero imaginary part, a
contradiction. Thus $t_0$ is a zero of a nonzero polynomial over
$\mathbb Q$. Formula \eqref{ari:eq:cs-torsionvalue} shows algebraicity
of $N$ using only positive rational powers of nonzero algebraic
absolute values. Taking $m=n+1$, raising to the $n$th power, and
using $|\Delta|=W^{1/2}$, $|\psi_{n+1}|=V^{1/2}$ gives the exact
field statement. The local-height formula is its logarithm.

\item Put $r=N^{12}$. Remove the factors
at $0,1$ from $\psi_n(P_t)$ and divide successively by the
corresponding polynomials for proper divisors of $n$; let $F_n$ be
the resulting primitive polynomial. At a point of exact order
$d\mid n$, the formal group gives
$z([n/d]Q)=(n/d)z(Q)+O(z(Q)^2)$, so the vanishing multiplicities for
$\psi_d$ and $\psi_n$ agree. Hence this cancellation leaves exactly
the exact-order-$n$ divisor, with its multiplicities. Set
$m=n+1$ and $e=m^2-1$. Clearing rational denominators in
\[
 \operatorname{Res}_S\left(
 F_n(S),\
 r^{2e}\{\psi_m(P_T)\psi_m(P_S)\}^{12}
              -\{\Delta(T)\Delta(S)\}^{e}\right)
\]
gives a polynomial in $\mathbb Q[T,r]$ that vanishes at
$(t_0,N(P_{t_0})^{12})$, since $S=\bar t_0$ is a common root.
It is nonzero: at an exact-order-$n$ root no $\psi_m$ vanishes,
so the leading coefficient in $r$ of the resultant is nonzero.
This is an explicit recurrence-and-resultant prescription for
every $n$, including nonreal torsion parameters.
\end{enumerate}
\end{proof}

\begin{proposition}[Five norm values and torsion through order eight]\hypertarget{tag-000S}{}\label{ari:C.4.4}
\StatementRecord{ari:C.4.4}{proved}{000S}{\SeriesRef{ari:eq:cs-torsionvalue}, \SeriesRef{ari:C.1.2}, \SeriesRef{ari:C.4.3}}

The exact special values are
\[
 \begin{array}{c|c|c}
 t&\operatorname{ord}(P)&N(P)^{12}\\ \hline
 0&3&27\\
 1&2&64\\
 32/5&4&10\\
 64(47-27\sqrt3)/11&6&17+10\sqrt3\\
 64(47+27\sqrt3)/11&6&10\sqrt3-17 .
 \end{array}
\]
At $t=0,1$ the orders refer to the smooth models after the specified
base changes. On the smooth base, $F_2=F_3=1$ and
\[
 \begin{aligned}
 F_4&=5t-32,\\
 F_5&=89t^2+4736t-4096,\\
 F_6&=11t^2-6016t+8192,\\
 F_7&=881t^4+4745344t^3-12079104t^2\\
    &\hspace{13mm}+24641536t-16777216,\\
 F_8&=49t^4-2614336t^3-7864320t^2\\
    &\hspace{13mm}+19398656t-8388608.
 \end{aligned}
\]
All these polynomials are squarefree. At each root of exact order
$n=4,5,6,7,8$, the positive number $r=N^{12}$ is respectively
determined by
\[
 \begin{array}{c|l}
 n&\text{equation for }r\\ \hline
 4&r=64/|t|\\
 5&r^5=432^5/\bigl((1-t)(5t-32)^4\bigr)\\
 6&r^2=1728^2/F_5(t)\\
 7&r^{35}=2^{192}3^{81}/
       \bigl(|t|^8|t-1|^3|F_6(t)|^{12}\bigr)\\
 8&r^4=1728^4/|F_7(t)|.
 \end{array}
\]
The order-five roots are real and below one; the order-six roots
are real and above one. The two complex order-seven roots are
included in the formula.
\end{proposition}


\begin{proof}
\begin{enumerate}[label=\arabic*.,leftmargin=*]
\item The division-polynomial recurrence, evaluated at $x=0$,
$y=t^2(t-1)/2$, gives
\[
 \begin{aligned}
 \psi_4(P)&=-t^{10}(t-1)^4F_4/32,\\
 \psi_5(P)&=t^{16}(t-1)^6F_5/4096,\\
 \psi_6(P)&=9t^{24}(t-1)^9F_6/65536,\\
 \psi_7(P)&=-t^{32}(t-1)^{12}F_7/16777216,\\
 \psi_8(P)&=-t^{42}(t-1)^{16}F_4F_8/268435456.
 \end{aligned}
\]
The recurrence below gives these identities over $\mathbb Q[t]$.
The Euclidean algorithm gives $\gcd(F_n,F_n')=1$ and excludes
common factors with the polynomials for proper orders. Thus the
listed factors give every point, with its exact order and multiplicity. Applying
\eqref{ari:eq:cs-torsionvalue} with $m=n-1$ gives every displayed
formula. The signs for $F_5,F_6$ follow directly from their quadratic
formula; in particular $F_5$ is positive at the two roots of $F_6$.
At order six, substituting
$t=64(47\mp27\sqrt3)/11$ gives the exact identities
\[
 F_5(t)(17\pm10\sqrt3)^2=1728^2.
\]
The positive roots are $17+10\sqrt3$ and $10\sqrt3-17$; their
product is $11$. The order-four value follows from
$64/(32/5)=10$, and the two central values are \ref{ari:C.1.2}.

\item The division-polynomial recurrence is
\[
 \begin{aligned}
 \psi_{2m+1}&=\psi_{m+2}\psi_m^3-\psi_{m-1}\psi_{m+1}^3,\\
 \psi_{2m}&=\frac{\psi_m}{\psi_2}
   (\psi_{m+2}\psi_{m-1}^2-\psi_{m-2}\psi_{m+1}^2).
 \end{aligned}
\]
Its initial values at $P$ are $\psi_0=0,\psi_1=1,\psi_2=2y$,
$\psi_3=-A^2$, $\psi_4=4y(-8B^2-A^3)$.
The resultant prescription of \cref{ari:C.4.3} gives a rational
polynomial relation for the nonreal order-seven values.
\end{enumerate}
\end{proof}

\begin{corollary}[The small-order torsion parameters]
\hypertarget{tag-000T}{}\label{ari:C.4.5}
\StatementRecord{ari:C.4.5}{proved}{000T}{\cref{ari:C.4.4}}
The numbers of smooth parameters of exact orders
$2,3,4,5,6,7,8$ are respectively
$0,0,1,2,2,4,4$. Thus there are thirteen parameters in total,
including the two nonreal conjugate parameters of order seven.
\end{corollary}
\begin{proof}
The polynomials in \cref{ari:C.4.4} are squarefree and their
roots have the indicated exact orders. Their degrees give the count.
\end{proof}


\section{Local expansions of the norm}

\begin{definition}[Local expansion constants and differential conventions]
\hypertarget{tag-IA03}{}\label{I:local-expansion-conventions}
\StatementRecord{I:local-expansion-conventions}{proved}{IA03}{}

Put $G=g_2(i)=\Gamma(1/4)^8/(16\pi^2)$ and
$D_*=-3\log6/\pi$.
Our theta convention is
$\vartheta_j(z,\tau)=\theta_j(\pi z\mid\tau)$, with lattice
$\mathbb Z+\tau\mathbb Z$.  The heat identity and modular
transformations are \cite[Equations 20.13.1, 20.7.26, 20.7.30
and 23.18.5]{DLMFTheta2026}.  The classical differential system,
in the notation $P=E_2$, $Q=E_4$, $R=E_6$, is
\[
 xP_x=(P^2-Q)/12,\qquad
 xQ_x=(PQ-R)/3,\qquad
 xR_x=(PR-Q^2)/2;
\]
see \cite[Equations (1.1)--(1.4), p.~1]{ErbanVanGorder2021}.
Here $x=e^{2\pi i\tau}$, so
$d/d\tau=2\pi i x\,d/dx$.
For $E=\pi^2E_2/3$, $g_2=4\pi^4E_4/3$ and
$g_3=8\pi^6E_6/27$, this gives the exact normalization
\[
 \begin{split}
 E'&=\frac{i}{2\pi}(E^2-g_2/12),\\
 g_2'&=\frac{i}{\pi}(2Eg_2-3g_3),\qquad
 g_3'=\frac{i}{\pi}(3Eg_3-g_2^2/6).
 \end{split}
\]
\end{definition}


\subsection{Differential identities}
\begin{proposition}[The two first-derivative identities]\hypertarget{tag-001L}{}\label{ari:C.8.11}
\StatementRecord{ari:C.8.11}{proved}{001L}{the definitions and constructions stated in the proof}

On an ordinary modulus chart write $L(\tau)=\log N(\mu(\tau),\tau)$ and put
\[
 a=\frac{\Im\mu}{T},\qquad
 W=\frac{\vartheta_1'(\mu,\tau)}{\vartheta_1(\mu,\tau)}
                         +2\pi ia,\qquad
 R=\frac{d\mu}{d\tau}-a.
\]
Then
\[
 L_\tau=\frac{W}{8\pi i}(W+4\pi iR),\qquad
 \det\frac{d}{d\tau}
 \begin{pmatrix}\tau&\mu\\\mu&\beta^{\rm part}/6\end{pmatrix}
       =\frac{(W+2\pi iR)^2}{4\pi^2}.
\]
Moreover
\[
 R=\frac{g_2}{6\pi i\wp'(\mu)}-\frac{W}{2\pi i},
 \qquad
 L_\tau=\frac{iW(3W\wp'(\mu)-2g_2)}
                         {24\pi\wp'(\mu)}.
\]
Multiplication by $6/\pi$ gives the critical-point factorization for
$D_{\rm part}$ in Proposition~\ref{push:src:D.4.4}.
\end{proposition}


\begin{proof}
\begin{enumerate}[label=\arabic*.,leftmargin=*]
\item Put $f=\vartheta_1'/\vartheta_1$ and
$E=\pi^2E_2/3$. The theta heat equation and
$\eta'/\eta=\pi iE_2/12$ give
\[
 \partial_\tau\log(\vartheta_1/\eta)
       =\frac{f_z+f^2+E}{4\pi i}.
\]
At $\wp(\mu)=0$, the identity $f_z=-\wp-E$ reduces
this to $f^2/(4\pi i)$. Thus, along the section,
\[
 \frac{d}{d\tau}\log(\vartheta_1(\mu,\tau)/\eta)
       =f\mu'+\frac{f^2}{4\pi i}.
\]
Differentiate the Gaussian $-\pi(\Im\mu)^2/T$ using
Wirtinger derivatives, retaining the factor $1/2$ in
the derivative of the real part of a holomorphic logarithm.
Substitution of $f=W-2\pi ia$ and $\mu'=R+a$ gives the
first formula. Since
$\beta^{\rm part}/6=(i/\pi)\log(\vartheta_1/(2\sqrt3\eta))$,
its derivative is $if\mu'/\pi+f^2/(4\pi^2)$.
Subtracting $\mu'^2$ gives
$(f+2\pi i\mu')^2/(4\pi^2)$, the second formula.

\item Finally differentiate $\wp_\tau(\mu(\tau))=0$ by the theta
heat equation. After substituting $\wp(\mu)=0$,
$\wp''(\mu)=-g_2/2$, and $f=W-2\pi ia$, the result is
$6\pi i\wp'R=g_2-3W\wp'$. This is the formula for $R$.
Inserting it into the first identity gives the last one.
The formulas involving $1/\wp'(\mu)$ apply where $\wp'(\mu)\ne0$.
The half-period case is treated in the root coordinate of
\cref{ari:C.3.1}.
\end{enumerate}
\end{proof}

\subsection{The square and equianharmonic points}

\begin{proposition}[The fourth-order expansion at the maximum]\hypertarget{tag-000N}{}\label{ari:C.3.1}
\StatementRecord{ari:C.3.1}{proved}{000N}{the definitions and constructions stated in the proof}

On either root branch at $\tau=i$, let
$h=\mu-(1-\tau)/2$ and $\delta=\tau-i$. Then
\begin{align*}
 \tau&=i-\frac{3\pi i}{2}h^2-\frac{9\pi^2i}{8}h^4+O(h^6),\\
 h^2&=\frac{2i}{3\pi}\delta+\frac1{3\pi}\delta^2+O(\delta^3),\\
 \mu&=\frac{1-\tau}{2}
  \mathbin{\pm}\left(\frac{2i\delta}{3\pi}\right)^{1/2}
          \left(1-\frac{i}{4}\delta+O(\delta^2)\right).
\end{align*}
The square root is taken holomorphically on the chosen root chart.
Furthermore
\begin{align}\hypertarget{tag-003M}{}\label{ari:eq:cs-Li}
 \log N={}&\tfrac12\log2-\frac\pi2|h|^2
   +\frac{3\pi^2}{8}|h|^4
   -\frac{3\pi^2}{4}|h|^2\operatorname{Re}(h^2)
   -\frac{5G}{192}\operatorname{Re}(h^4)+O(|h|^6),\\
 D_{\rm part}={}&D_* -\frac2\pi|\delta|
  +\frac1\pi\bigl(|\delta|^2+|\delta|\operatorname{Im}\delta\bigr)
  +\frac{5G}{72\pi^3}\operatorname{Re}(\delta^2)
  +O(|\delta|^3).\nonumber
\end{align}
In particular the leading terms agree with the boundary expansion in Lemma~\ref{push:src:D.2.1}.
\end{proposition}


\begin{proof}
\begin{enumerate}[label=\arabic*.,leftmargin=*]
\item Write $e(\tau)=\wp_\tau((1-\tau)/2)$ and
$E=\pi^2E_2/3$. The half-period Taylor series is
\[
 \wp_\tau((1-\tau)/2+h)
 =e+(3e^2-g_2/4)h^2+(3e^3-g_2e/4)h^4+O(h^6).
\]
The exact half-period coefficients used in Lemma~\ref{push:src:D.2.1} are
$e(i)=0$, $e'(i)=iG/(6\pi)$. This proves the claimed identity
$3e'(i)=(2i/\pi)G/4$. Its closed form follows from
$e=\pi^2(\vartheta_2^4-\vartheta_4^4)/3$ and the heat equation:
if $A=\vartheta_2(i)^4=\vartheta_4(i)^4$, then
$G=4\pi^4A^2$ and $e'(i)=2\pi^3iA^2/3$.
For the classical closed value, \cite[Equation 23.5.3]{DLMFTheta2026}
gives $e_1=\Gamma(1/4)^4/(32\pi\omega_1^2)$,
$e_2=0$ and $e_3=-e_1$.  The lattice identification is
$2\omega_1=1$, $2\omega_3=i$; hence
$G=-4(e_1e_2+e_2e_3+e_3e_1)=4e_1^2
=\Gamma(1/4)^8/(16\pi^2)$.

\item Ramanujan's identities give $g_2'(i)=2iG$. Differentiating
$e(-1/\tau)=\tau^2 e(\tau)$ twice at $i$ gives
$e''(i)=3i e'(i)=-G/(2\pi)$. Inserting
$\delta=A_2h^2+A_4h^4+O(h^6)$ into the displayed half-period series
therefore gives
\[
 A_2=-3\pi i/2,\qquad
 A_4=\frac{g_2'A_2/4-e''A_2^2/2}{e'}=-9\pi^2i/8.
\]
The implicit-function theorem applies because $e'(i)\ne0$;
evenness in $h$ gives the stated remainder. Series inversion gives
the formulas for $h^2$ and $\mu$.

\item Let $H(h,\tau)=\log(\vartheta_3(h,\tau)/\eta(\tau))$. The exact
identity from the half-period translation is
$\log N=\operatorname{Re}H-\pi(\operatorname{Im}h)^2/
\operatorname{Im}\tau$. At $(0,i)$ its relevant Taylor data are
\[
 \begin{gathered}
 H_0=\tfrac12\log2,\quad H_\tau=0,\quad
 [h^2]H=-\pi/2,\quad [h^4]H=G/48,\\
 \partial_\tau[h^2]H=-\pi i/4-iG/(16\pi),\qquad
 H_{\tau\tau}=-G/(24\pi^2).
 \end{gathered}
\]
Indeed $\partial_h^2\log\vartheta_3=-\wp(h+(1-\tau)/2)-E$;
the heat equation at $h=0$ gives $H_\tau=ie/(4\pi)$, and
$E'(i)=\pi i/2-iG/(24\pi)$. These identities prove every entry.
After substitution $\tau=\tau(h)$, the pure $h^4$ coefficient of
$H$ is $-3\pi^2/8-5G/192$. Expanding
$\operatorname{Im}\tau=1-(3\pi/2)\operatorname{Re}(h^2)+O(|h|^4)$
in the Gaussian gives \eqref{ari:eq:cs-Li}. Finally
\[
 |h|^2=\frac{2}{3\pi}|\delta|
          \left(1+\tfrac12\operatorname{Im}\delta+O(|\delta|^2)\right)
\]
and substitution into the quartic expression, followed by
multiplication by $6/\pi$, gives the formula for $D_{\rm part}$.

\end{enumerate}
\end{proof}

\begin{proposition}[The equianharmonic expansion]\hypertarget{tag-000O}{}\label{ari:C.3.2}
\StatementRecord{ari:C.3.2}{proved}{000O}{the definitions and constructions stated in the proof}

Set $T_0=\sqrt3/2$, $\mu_0=(2-\rho)/3$, and
$p=\wp'_\rho(\mu_0)$. Then $p\ne0$, $p^2=-g_3(\rho)$, and, for
$\delta=\tau-\rho$ on the order-three root chart,
\begin{align*}
 \mu&=\frac{2-\tau}{3}
      +\frac{3p}{16\pi^2}\delta^2
      +\frac{3ip}{32\pi^2T_0}\delta^3+O(\delta^4),\\
 t&=\frac{ip^2}{\pi^3}\delta^3
        \left(1+\frac{3i}{2T_0}\delta+O(\delta^2)\right),\\
 \log N&=\tfrac14\log3+\frac7{192}\operatorname{Re}t
       -\frac{9\pi}{512T_0|p|^{2/3}}|t|^{4/3}
       +O(|t|^{5/3}).
\end{align*}
Thus
\[
 D_{\rm part}(t)=-\frac3{2\pi}\log48
       +\frac7{32\pi}\operatorname{Re}t
       -\frac{27}{256T_0|p|^{2/3}}|t|^{4/3}
       +O(|t|^{5/3}).
\]
The point $t=0$ is neither a local maximum nor a local minimum.
\end{proposition}


\begin{proof}
\begin{enumerate}[label=\arabic*.,leftmargin=*]
\item Use the notation of \cref{ari:C.8.11}, with
$\ell=\partial_z\log\vartheta_1(\mu,\tau)$ and $Z=W$.
The differential identities proved there apply to the zero locus.
At $\rho$ the unit endomorphism fixing the
nonzero order-three point forces the gradient of $U_\rho$ there to
vanish, so $Z=0$. Explicitly the unit $\rho-1$ fixes $\mu_0$ modulo
the lattice and has no fixed nonzero real covector. Hence
$a=\mu'=-1/3$ and $\ell=2\pi i/3$.

\item Along the zero locus, a further heat-equation differentiation gives
\[
 \ell'=\frac{i\wp'}{4\pi}+\frac{iEg_2}{6\pi\wp'}.
\]
The Ramanujan derivatives at $\rho$ are
\[
 g_2'=\frac{3ip^2}{\pi},\quad
 g_2''=\frac{5i}{T_0}g_2',\quad
 \frac{(\wp')'}{\wp'}=\frac{3i}{2T_0},\quad E=\frac\pi{T_0}.
\]
They give $\mu''=3p/(8\pi^2)$,
$\mu'''=9ip/(16\pi^2T_0)$ and hence the asserted $\mu$ expansion.
Expansion of $t=g_2^3/(g_2^3-27g_3^2)$ gives the displayed $t$
formula. Since $\partial_\tau a=R/(2iT)$, one has
\[
 Z=\frac{ip}{4\pi}\delta+O(|\delta|^2),\qquad
 R=\frac{3p}{8\pi^2}\delta+O(|\delta|^2).
\]
The derivative formula now gives the harmonic cubic
$\operatorname{Re}(7ip^2\delta^3/(192\pi^3))=(7/192)
\operatorname{Re}t+O(|\delta|^4)$.

\item For the next term, direct differentiation of the Gaussian gives
\[
 \partial_\tau\partial_{\bar\tau}\log N
       =-\frac\pi{2T}|R|^2
       =-\frac{9|p|^2}{128\pi^3T_0}|\delta|^2+O(|\delta|^3).
\]
After the harmonic cubic has been written as a function of $t$, the
order-three stabilizer permits only a radial term in degree four:
the stabilizer map is $\gamma\tau=1-1/\tau$ and the zero
branch transforms as
$\mu(\gamma\tau)=(1-\mu(\tau))/\tau$.
The theta and eta transformations cited above cancel against the
Gaussian factor, so $N(\gamma\tau)=N(\tau)$.
Thus
in the linearizing coordinate
$(\tau-\rho)/(\tau-\bar\rho)$ invariant monomials have exponents
whose difference is divisible by three. The coefficient of
$|\delta|^4$ is therefore $-9|p|^2/(512\pi^3T_0)$, since
$\partial_\delta\partial_{\bar\delta}|\delta|^4=4|\delta|^2$.
Converting $|\delta|^4$ to $|t|^{4/3}$ proves the expansion, including
the remainder. Its nonzero linear term in $\operatorname{Re}t$
has both signs, which proves the last assertion. In particular
superharmonicity at this root point does not make it a local maximum.
\end{enumerate}
\end{proof}




\section{An analytic proof of the elliptic maximum}

\begin{remark}[Modular and arithmetic discriminants]
\hypertarget{tag-IA24}{}\label{I:modular-arithmetic-discriminants}
\StatementRecord{I:modular-arithmetic-discriminants}{proved}{IA24}{}
Use $q=e^{2\pi i\tau}$, $T=\Im\tau$, $\Delta_0$ and $t$
from \cref{I:def:theta-height}. Thus
$t=E_4^3/(1728\Delta_0)$ and
$t-1=E_6^2/(1728\Delta_0)$.
On the smooth rational model $t\ne0,1$, choose $\lambda$ with
\[
 \lambda^2=-\frac{t(t-1)g_2}{3g_3}.
\]
The exact uniformization is
$(x,y)=(\lambda^2\wp(z),\lambda^3\wp'(z)/2)$, where
$\wp$ has lattice $\mathbb Z+\tau\mathbb Z$.
Choose the sign of $\lambda$ so that the marked zero $\mu$
maps to $P_t$.
Indeed $g_2^3/g_3^2=27t/(t-1)$ gives
\[
 A=-\lambda^4g_2/4,\qquad B=-\lambda^6g_3/4,
 \qquad \Delta_E=\lambda^{12}(2\pi)^{12}\Delta_0.
\]
The invariant differential pulls back to $\lambda^{-1}dz$,
so its period lattice is $\lambda^{-1}(\mathbb Z+\tau\mathbb Z)$.
\end{remark}

\begin{proposition}[Three uniform bounds on the elliptic norm]\hypertarget{tag-001J}{}\label{ari:C.8.1}
\StatementRecord{ari:C.8.1}{proved}{001J}{\SeriesRef{ari:eq:ca-sup}, \SeriesRef{ari:eq:ca-exp}}

For every $z\in\C$ and every $\tau$ in the standard modular
fundamental region, we have
\[
 \begin{array}{c|ccc}
 \text{restriction}&T\le2&T\le7/5&T\le11/10\\ \hline
 N(z,\tau)&<17/10&<3/2&<143/100.
 \end{array}
\]
If $\wp_\tau(z)=0$ and $T\ge2$, then $N(z,\tau)<541/500$.
\end{proposition}


\begin{proof}
\begin{enumerate}[label=\arabic*.,leftmargin=*]
\item Taking absolute values in the theta series and completing the square
gives
\[
 e^{-\pi(\Im z)^2/T}|\vartheta_1(z,\tau)|
 \le\sum_{n\in\Z}
 e^{-\pi T(n+1/2+\Im z/T)^2}
 \le\sum_{n\in\Z}e^{-\pi Tn^2}.
\]
The last inequality follows from Poisson summation: the Fourier
coefficients of the periodized Gaussian are positive, so its maximum
is at an integer. The eta product and $n^2\ge1+3(n-1)$ for $n\ge1$
therefore imply, on $a\le T\le b$,
\begin{equation}\hypertarget{tag-002G}{}\label{ari:eq:ca-sup}
 N(z,\tau)\le
 \exp\!\left(\frac{r_a}{(1-r_a)^2}\right)
 \max\{F(a),F(b)\},\qquad
 F(T)=e^{\pi T/12}
 \left(1+\frac{2e^{-\pi T}}{1-e^{-3\pi T}}\right),
\end{equation}
where $r_a=e^{-2\pi a}$. Indeed,
$-\sum\log(1-r^n)\le r/(1-r)^2$. The function $F$ is convex,
as follows by expanding its denominator as a geometric series;
each summand is an exponential of a real linear function of $T$.
We use $6/7<\sqrt3/2$ and the following rational bounds:
\[
\begin{array}{c|@{\,}c@{\,}c@{\,}c@{\,}c@{\,}c}
T&6/7&1&11/10&7/5&2\\ \hline
F(T)&<14211/10000&<14116/10000&<14180/10000&<14782/10000&<16944/10000.
\end{array}
\]
The eta factors in \eqref{ari:eq:ca-sup} are less than $10047/10000$
for $a=6/7$ and less than $10019/10000$ for $a=1$.
Multiplying these fractions proves the asserted bounds, using
$[6/7,1]$ and $[1,2]$ separately for the first one.

\item The scalar estimates follow from $314159/100000<\pi<355/113$ and, for
$0\le x\le20$, the Taylor bounds
\begin{equation}\hypertarget{tag-002C}{}\label{ari:eq:ca-exp}
 S_{80}(x)\le e^x\le
 S_{80}(x)+\frac{x^{81}}{81!}\frac1{1-x/82},
 \qquad S_{80}(x)=\sum_{j=0}^{80}\frac{x^j}{j!}.
\end{equation}
Taking reciprocals bounds negative exponentials. The remainder
bound follows by comparison of consecutive terms. These inequalities
give each displayed rational bound.

\item For the cusp use the Fourier expansion of $\wp$.
Choose $|\Im z|\le T/2$ by translation, put $b_0=43/100$ and
$r_0=e^{-4\pi}$. The Fourier series
\[
 \frac{\wp_\tau(z)}{\pi^2}
 =\csc^2(\pi z)-\frac13
   +8\sum_{n\ge1}\frac{nq^n}{1-q^n}
          \bigl(1-\cos(2\pi nz)\bigr)
\]
shows, if $|\Im z|\ge b_0$, that
\begin{align*}
 &\left|\frac{\wp_\tau(z)}{\pi^2}+\frac13\right|\\
 &\le \frac1{\sinh^2(\pi b_0)}
 +8\left\{\frac{r_0}{(1-r_0)^3}
 +\frac1{2(1-r_0)}
 \left(\frac{e^{-2\pi}}{(1-e^{-2\pi})^2}
       +\frac{e^{-6\pi}}{(1-e^{-6\pi})^2}\right)\right\}
 <\frac13.
\end{align*}
The two terms in the last parentheses bound respectively
$r^n e^{2\pi n|\Im z|}$ and $r^n e^{-2\pi n|\Im z|}$.
This inequality excludes a zero. Hence a zero has
$|\Im z|<b_0$. The Jacobi product, $\log(1+x)\le x$, and
the monotonicity of $\cosh(\pi y)e^{-\pi y^2/T}$ on
$0\le y\le b_0$ give
\[
 N(z,\tau)\le
 2\cosh(\pi b_0)e^{-\pi/3-\pi b_0^2/2}
 \exp\!\left(
 \frac{2\cosh(2\pi b_0)r_0}{1-r_0}
 +\frac{r_0^2}{1-r_0^2}\right)<\frac{541}{500}.
\]
For the stated monotonicity, the concave function
$\tanh(\pi y)-y$ is zero at zero and positive at $b_0$.
For $T\ge2$ the right side before substituting $T=2$ decreases:
the derivative of $-\pi T/6-\pi b_0^2/T$ is negative and
both product errors decrease. The two strict scalar inequalities
follow from \eqref{ari:eq:ca-exp}. This proves the cusp assertion
on the whole cusp region.
\end{enumerate}
\end{proof}

\begin{proposition}[The maximum on an explicit local disc]\hypertarget{tag-001Q}{}\label{ari:C.8.2}\hypertarget{tag-01QK}{}\label{fam:src:C.2.2}\label{fam:st:C.2.2}
\StatementRecord{ari:C.8.2}{proved}{001Q}{\SeriesRef{ari:eq:ca-exp}, \SeriesRef{ari:eq:ca-derivatives}, \SeriesRef{ari:eq:ca-GHbounds}, \SeriesRef{ari:eq:ca-H}, \SeriesRef{ari:eq:ca-delta-h}}

If $|\tau-i|\le1/50$ and $\wp_\tau(z)=0$, then
$N(z,\tau)\le\sqrt2$, with equality exactly when $\tau=i$.
\end{proposition}


\begin{proof}
\begin{enumerate}[label=\arabic*.,leftmargin=*]
\item Write $\delta=\tau-i$, $h=z-(1-\tau)/2$ and
\[
 G(h,\tau)=\wp_\tau((1-\tau)/2+h),\qquad
 H(h,\tau)=\log\frac{\vartheta_3(h,\tau)}{\eta(\tau)}.
\]
We choose the holomorphic logarithm whose value at $(0,i)$ is
$\log\sqrt2$. The half-period theta identity gives
\begin{equation}\hypertarget{tag-0025}{}\label{ari:eq:ca-H}
 \log N(z,\tau)=\Re H(h,\tau)-\frac{\pi(\Im h)^2}{T}.
\end{equation}
Both $G$ and $H$ are even in $h$. On the polydisc
$|h|\le1/4$, $|\delta|\le1/10$, they satisfy
\begin{equation}\hypertarget{tag-0024}{}\label{ari:eq:ca-GHbounds}
 |G|<30,\qquad |H|<2/3.
\end{equation}
Here are sufficient series estimates, including their margins.
Put
\[
 A_k=\sum_{n\ge1}n^k e^{-9\pi n^2/10}\cosh(\pi n/2),
 \quad a_0=2A_0,\quad a_1=4\pi A_1,\quad a_2=8\pi^2A_2.
\]
Then $a_0<3/10$, $a_1<19/10$, $a_2<12$.
The first two summands suffice to verify these inequalities: for
$n\ge3$, each of the two exponentials in the hyperbolic cosine
has successive-term ratio at most $e^{-29\pi/5}$, and the
factor $n^k$ increases this ratio by at most $(4/3)^2$.
We use \eqref{ari:eq:ca-exp} for the first terms and this geometric
bound for the tail. With $r=e^{-9\pi/5}$, the identities
\[
 G=-\partial_h^2\log\vartheta_3-\frac{\pi^2}{3}E_2,
 \qquad
 |E_2|\le1+\frac{24r}{(1-r)^3}
\]
give
\[
 |G|\le\frac{a_2}{1-a_0}+\frac{a_1^2}{(1-a_0)^2}
       +\frac{\pi^2}{3}\left(1+\frac{24r}{(1-r)^3}\right)<30.
\]
Also
\[
 |H|\le-\log(1-a_0)+\frac{11\pi}{120}
                    +\frac{r}{(1-r)^2}<\frac23.
\]
The logarithm is well defined since $|\vartheta_3-1|<3/10$;
the displayed eta product fixes its branch. A geometric bound
for $-\log(1-x)=\sum_{n\ge1}x^n/n$ verifies the last estimate.

\item Let
\[
 e(\tau)=\frac{\pi^2}{3}(\vartheta_2(0,\tau)^4-
                              \vartheta_4(0,\tau)^4),\qquad
 E(\tau)=\frac{\pi^2}{3}E_2(\tau),\qquad
 g=\frac{g_2(i)}4.
\]
On $|\delta|\le1/50$ the following loose bounds hold:
\begin{equation}\hypertarget{tag-002B}{}\label{ari:eq:ca-derivatives}
 |e'|<14,\quad |e''|<60,\quad |g_2'|<500,
 \quad |E'|<2,\quad |E_6''|<100,
 \qquad 47<g<95/2.
\end{equation}
The scalar Fourier series give, at $T\ge49/50$, the following
upper bounds for $|\vartheta_j|,|\vartheta_j'|,|\vartheta_j''|$:
\[
\begin{array}{c|ccc}
j&0\text{ derivatives}&1\text{ derivative}&2\text{ derivatives}\\ \hline
2&929/1000&742/1000&670/1000\\
4&1093/1000&290/1000&911/1000.
\end{array}
\]
They follow by differentiating the absolutely convergent theta
series, taking absolute values, retaining two terms, and using
the same successive-term bound. Differentiating the fourth
powers proves the first two bounds in \eqref{ari:eq:ca-derivatives}.
For $r\le1/470$, the remaining derivative majorants are
\begin{align*}
 |g_2'|&\le640\pi^5\frac{r(1+11r+11r^2+r^3)}{(1-r)^7}<500,\\
 |E'|&\le16\pi^3\frac{r(1+r)}{(1-r)^5}<2,\\
 |E_6''|&\le2016\pi^2
 \frac{r(1+r)(1+120r+1191r^2+2416r^3+1191r^4+120r^5+r^6)}
 {(1-r)^{11}}<100.
\end{align*}
These follow by summing $\sum n^k r^n$ and bounding
$(1-r^n)^{-1}$ by $(1-r)^{-1}$. Substitution of $r=1/470$
and $\pi<355/113$ verifies the rational inequalities.
Finally $g=\pi^4E_4(i)/3$ and
$E_4=1+240\sum n^3q^n/(1-q^n)$ give $47<g<95/2$.
The half-period identity used in Lemma~\ref{push:src:D.2.1},
obtained by differentiating $e$ and using the theta heat equation, reads
\[
 e(i)=0,\qquad e'(i)=id,\qquad d=\frac{2g}{3\pi}>99/10.
\]

\item The differential equation $\wp''=6\wp^2-g_2/2$ determines the
first coefficients in the even expansion:
\[
 G(h,\tau)=e+a h^2+ea h^4+\mathcal R_G,
 \qquad a=3e^2-g_2/4.
\]
Consequently, on $|\delta|\le1/50$,
\[
 |e|\le7/25,\quad |a+g|\le\frac52+3(7/25)^2=\frac{3419}{1250},
 \quad |a|\le\frac{62794}{1250}.
\]
Cauchy's estimate from \eqref{ari:eq:ca-GHbounds} yields
\[
 |\mathcal R_G|\le
 \frac{30(4|h|)^6}{1-16|h|^2}.
\]
On $|h|=1/10$ the term $ah^2$ has modulus at least
$55331/125000>0.4426$, whereas the sum of the other terms is at most
\[
 \frac7{25}+\frac7{25}\frac{62794}{1250}\frac1{10000}
       +\frac{30(4/25)^3}{1-4/25}<0.4277.
\]
Rouch\'e's theorem puts two zeros, with multiplicity, in
$|h|<1/10$. The disc injects into the torus and $\wp$ has exactly
two zeros on the torus, so it contains both of them.

\item Write $G=e+h^2Q$. The preceding estimates give, for $|h|\le1/10$,
\[
 |Q+g|\le
 \frac{3419}{1250}+\frac7{25}\frac{62794}{1250}\frac1{100}
       +\frac{30(4/25)^3}{(1-4/25)/100}<\frac{88}{5}.
\]
Taylor's formula and $|e''|<60$ give
$|e|\ge(99/10-3/5)|\delta|=(93/10)|\delta|$.
At either zero, therefore,
\begin{equation}\hypertarget{tag-0028}{}\label{ari:eq:ca-delta-h}
 |\delta|<7|h|^2,
\end{equation}
because $(95/2+88/5)/(93/10)=7$.

\item Expand
$H=H_0(\tau)+H_2(\tau)h^2+H_4(\tau)h^4+\mathcal R_H$.
The heat equation and the identity for $G$ give
\[
 H_0'=\frac{ie}{4\pi},\quad
 H_2=-\frac{e+E}{2},\quad H_4=-\frac a{12}.
\]
Thus $H_0(i)=\log\sqrt2$, $H_0'(i)=0$,
$|H_0''|<6/5$, $H_2(i)=-\pi/2$, $|H_2'|<8$, and
$|H_4|<21/5$. Cauchy's estimate gives
\[
 |\mathcal R_H|\le
 \frac{(2/3)(4|h|)^6}{1-16|h|^2}
 \le\frac{2048}{63}|h|^4
 \quad (|h|\le1/10).
\]
Combining \eqref{ari:eq:ca-H} and \eqref{ari:eq:ca-delta-h}, we obtain
\begin{align*}
 \log N-\log\sqrt2
 &\le-\frac\pi2\Re(h^2)-\pi(\Im h)^2+C|h|^4\\
 &=-\frac\pi2|h|^2+C|h|^4,\\
 C&=\frac{147}{5}+56+\frac{21}{5}
                 +\frac{2048}{63}+\frac{1100}{49}<150.
\end{align*}
The final term of $C$ bounds the change of $1/T$ from $1$:
$7\pi/(1-1/50)<1100/49$. Since $150|h|^2<3/2<\pi/2$,
the result is strict unless $h=0$. Equation
\eqref{ari:eq:ca-delta-h} then gives $\delta=0$, where equality
holds by \cref{ari:C.1.2}.
\end{enumerate}
\end{proof}

\begin{proposition}[Comparison of the two local discs]\hypertarget{tag-001R}{}\label{ari:C.8.3}
\StatementRecord{ari:C.8.3}{proved}{001R}{\SeriesRef{ari:eq:ca-derivatives}, \SeriesRef{ari:C.8.2}}

Every complex number $t_0$ with $|t_0-1|\le1/500$ has a
representative $\tau$ with $|\tau-i|<1/50$ and $j(\tau)=1728t_0$.
Consequently every zero of $\wp$ for that modulus satisfies
$N\le\sqrt2$, with equality only for $t_0=1$.
\end{proposition}


\begin{proof}
\begin{enumerate}[label=\arabic*.,leftmargin=*]
\item At $i$, $E_6=0$ and the Ramanujan differential identity gives
$|E_6'(i)|=\pi E_4(i)^2>13/2$, since $E_4(i)>29/20$.
On $|\delta|=1/50$, \eqref{ari:eq:ca-derivatives} gives
\[
 |E_6(\tau)|>\frac{13}{100}-\frac1{50}=\frac{11}{100}.
\]
On the whole disc, $|q|<1/470$ and
$|\Delta_0|\le |q|\exp(24|q|/(1-|q|))<1/400$.
It follows on the boundary that
\[
 |t-1|=\frac{|E_6|^2}{1728|\Delta_0|}
 >\frac{121}{43200}>\frac1{500}.
\]
Rouch\'e applied to $E_6'(i)\delta$ shows that $E_6$ has
exactly one zero in the disc. Hence $t-1$ has two zeros there,
with multiplicity, and another application of Rouch\'e shows
that $t-t_0$ has two for each $|t_0-1|\le1/500$.
Thus the stated $t$-disc is covered. Elliptic
curves over $\C$ with the same $j$ are isomorphic. In terms of
lattices the isomorphism is a homothety; it carries the two
zeros of $\wp$ to the two zeros of the scaled $\wp$ and
preserves $N$. \Cref{ari:C.8.2} applies to either zero.
\end{enumerate}
\end{proof}

\begin{proposition}[The first three normalized duplication steps]\hypertarget{tag-001S}{}\label{ari:C.8.4}
\StatementRecord{ari:C.8.4}{proved}{001S}{\SeriesRef{ari:C.4.2}, \SeriesRef{ari:eq:cs-division}, \SeriesRef{ari:C.3.1}}

Put $u=t-1$. For $n=2^k$, $k=1,2,3$, and $nP\ne O$,
\begin{equation}\hypertarget{tag-002D}{}\label{ari:eq:ca-ladder}
 \left(\frac{N(P)}{\sqrt2}\right)^{n^2}
 =\frac{3^{3/4}}{2n}\,
       \frac{N(nP)|u|^{-1/4}}{|A_k(u)|},
\end{equation}
where
\begin{align*}
 A_1(u)&=1,\qquad A_2(u)=1-5u/27,\\
 A_3(u)&=1-\frac{217}{27}u-\frac{20486}{729}u^2
       +\frac{10910}{19683}u^3
       +\frac{161383}{177147}u^4
       -\frac{245}{14348907}u^5.
\end{align*}
On $|u|\le1/28$ we have $|A_3(u)|\ge1-10|u|$.
For every $k\ge1$, define
\[
 E_k(t)=\frac{|\Delta_E|^{(4^k-1)/12}|t-1|^{1/4}}
 {2^{4^k/2}|\psi_{2^k}(P)|}.
\]
Then $(N(P)/\sqrt2)^{4^k}=E_k(t)N(2^kP)|t-1|^{-1/4}$,
and the removable limiting value is
\[
 \lim_{t\to1}E_k(t)=\frac{3^{3/4}}{2^{k+1}}
                 =2^{1-k}\frac{3^{3/4}}4.
\]
\end{proposition}


\begin{proof}
\begin{enumerate}[label=\arabic*.,leftmargin=*]
\item Substitute the exact division polynomials from \cref{ari:C.4.2}
in \eqref{ari:eq:cs-division}. Equivalently, remove
$t^{2(n^2-1)/3}(t-1)^{n^2/4}$ from $\psi_n(P)$ and divide
the resulting polynomial by its value at $t=1$. For $n=2,4,8$
this gives the displayed polynomials and the constant
$3^{3/4}/(2n)$. These are rational polynomial operations using the recurrence in
\cref{ari:C.4.4}. The general identity is \cref{ari:C.4.2}.
For its limit write $\mu=(1-\tau)/2+h$ as in \ref{ari:C.3.1}. Since
$n=2^k$ is even, $n\mu$ is congruent to $nh$ modulo the
period lattice. Hence
\[
 N(nP)=2\pi|\eta(i)|^2 n|h|(1+o(1)),\qquad
 t-1=\tfrac94\pi^4E_4(i)h^4(1+o(1)).
\]
The identity $E_4(i)^3=1728\eta(i)^{24}$ then gives
$N(nP)|t-1|^{-1/4}\to2n/3^{3/4}$.
Since the left side of the general ladder tends to one,
this proves the stated limit for every $k$.
The triangle inequality gives
$|A_3(u)|\ge1-|u|M$, where
\[
 M\le\frac{81}{10}+\frac{29}{28}
                +\frac1{28^2}+\frac1{28^3}+\frac1{28^4}<10.
\]
Here we bounded respectively the five absolute coefficients by
$81/10,29,1,1,1$. The bound holds on the whole disc.
\end{enumerate}
\end{proof}

\begin{theorem}[Exact norm maximum and parameter discs]
\hypertarget{tag-01QL}{}\label{fam:src:C.2.3}\label{fam:st:C.2.3}
\label{ari:C.8.5}\label{push:src:E.4.3}
\StatementRecord{fam:src:C.2.3}{proved}{01QL}{\cref{ari:C.8.1,ari:C.8.2,ari:C.8.3,ari:C.8.4,fam:eq:cd-norm,I:period-determinant}}
At every zero of $\wp_\tau$,
\[
 N\le\sqrt2,
\]
with equality precisely in the modular orbit of $\tau=i$.
For the cusp normalization $\beta^{\rm part}+\tau\to0$ this gives
\[
 M=-\frac3\pi\log6,\qquad
 \operatorname*{arg\,max}_{B\setminus\{p_0\}}D_{\rm part}=\{1\}.
\]
The real period determinant is nonzero throughout the base exactly
when $\Im c<-M$. Consequently the parameter domains are
\[
 0<|w|<216,\qquad 0<|u_{\rm Engel}|<46656.
\]
Changing the reference to $u_{\lambda_0}=\lambda_0w^2$ multiplies
the latter radius by $|\lambda_0|$.
\end{theorem}

\begin{proof}
\begin{enumerate}[label=\arabic*.,leftmargin=*]
\item Reduce to the standard modular region. The portion $T\ge2$
is covered by \cref{ari:C.8.1}. Suppose $T\le2$ and put
$r=|t-1|$. For $r\ge9/10$, the step $n=2$ in
\eqref{ari:eq:ca-ladder} and $N(2P)<17/10$ give strict inequality,
because
\[
 \frac{27(17/10)^4}{256}<\frac9{10}.
\]
There is no $2$-torsion point of the section on the smooth base.

\item If $r\le9/10$, then $T<7/5$. Indeed, $T\ge7/5$ would give
$|q|<1/6500$ and
\[
 |E_4|>24/25,\qquad |\Delta_0|<1/6400,
 \qquad |t|>\frac{6400(24/25)^3}{1728}>19/10,
\]
contrary to $|t|\le19/10$. The estimates follow from
\[
 \sum_{n\ge1}\sigma_3(n)r^n
 \le\frac{r(1+4r+r^2)}{(1-r)^5},\qquad
 |\Delta_0|\le r\exp\frac{24r}{1-r}.
\]
For $1/28\le r\le9/10$, the $n=4$ step and
$N(4P)<3/2$ now apply. The function
$(1-5r/27)r^{1/4}$ is increasing up to $r=27/25$.
At its left endpoint its fourth power is
\[
 \frac1{28}\left(\frac{751}{756}\right)^4
 >\frac{2187}{65536}
 =\left(\frac{3^{3/4}}8\frac32\right)^4.
\]
This proves strict inequality on the whole annulus. The lower
bound for $A_2$ is positive there.

\item If $r\le1/28$, then $T<11/10$. For otherwise $|q|<1/1000$,
and the Lambert series with
\[
 \sum_{n\ge1}\sigma_5(n)r^n
 \le\frac{r(1+26r+66r^2+26r^3+r^4)}{(1-r)^7}
\]
gives $|E_6|>9/20$, $|\Delta_0|<1/900$, and hence
$|t-1|>900(9/20)^2/1728>1/28$.
On $1/500\le r\le1/28$, the $n=8$ step,
$N(8P)<143/100$, and $|A_3|\ge1-10r$ prove the claim.
Indeed $(1-10r)r^{1/4}$ has a single critical point in this
interval, a maximum at $r=1/50$, so its minimum is at an endpoint.
At both endpoints its fourth power exceeds $9/5000$, whereas
\[
 \left(\frac{3^{3/4}}{16}\frac{143}{100}\right)^4
 <\frac7{4000}<\frac9{5000}.
\]
Again the displayed positive bound excludes a division-polynomial
zero. The remaining disc $r\le1/500$ is covered by \cref{ari:C.8.3}. The smooth limiting fibre at $t=0$ has
$N=3^{1/4}<\sqrt2$; the formula also extends there by
continuity. The only equality occurs at $t=1$.

\item The norm identity \eqref{fam:eq:cd-norm} gives
$M=(6/\pi)\log(\sqrt2/(2\sqrt3))=-3\log6/\pi$.
The cusp asymptotic in \cref{fam:src:C.2.1} and continuity show
that $D_{\rm part}+\Im c$ has a zero for $\Im c\ge-M$:
use the maximum at $t=1$ and a path to the cusp.
For $\Im c<-M$ it is strictly negative everywhere.
By \cref{I:period-determinant}, this is exactly nonvanishing of
the real period determinant throughout the base.
Since $|w|=e^{\pi\Im c}$ and $u_{\rm Engel}=w^2$,
the radii are $e^{-\pi M}=6^3$ and $e^{-2\pi M}=6^6$.
The reference change is the displayed multiplication by $\lambda_0$.
\end{enumerate}
\end{proof}






\section{The family and its intrinsic classification}
\hypertarget{tag-01P9}{}\label{fam:section:A}

\begin{definition}[Affine covariance]
\hypertarget{tag-PI08}{}\label{I:def:affine-covariance}
\StatementRecord{I:def:affine-covariance}{proved}{PI08}{}
Use the standard data of Definition~\ref{I:def:family-parameter}
and write $\varepsilon=v_1$, $\varepsilon'=-v_2$,
$\epsilon_j=\gamma(v_j)$, so $(\epsilon_1,\epsilon_2)=(1,-1)$.
Put
\begin{equation}\hypertarget{tag-03CO}{}\label{wb:pshift:eq:rank-one-map}
 d=\widehat\delta,\qquad e=(0,1)^{\mathsf t}=\Pi_cd,
 \qquad E=d\otimes\gamma,\qquad E^2=0,\quad\Pi_cE=e\gamma.
\end{equation}
The fibre covariance is $R_g(z)\Pi_c(z)=\Pi_c(gz)A_g$, where
\begin{equation}\hypertarget{tag-03C2}{}\label{wb:pshift:eq:covariance}
 R_{g_1}=\begin{pmatrix}-1/\tau&0\\(1-\mu)/\tau&1\end{pmatrix},
 \qquad R_{g_2}=\begin{pmatrix}1/\tau&0\\-\mu/\tau&1\end{pmatrix},
 \qquad R_{g_0}=I_2.
\end{equation}
In the cusp coordinate $s$, the generator satisfies
$s(g_0z)=s(z)-1$. The matrices satisfy $R_ge=e$ and
$A_jE=EA_j=E$.
\end{definition}

\subsection{Canonical sections and the finite covers}\hypertarget{tag-01PA}{}\label{fam:section:A.1}
\hypertarget{tag-03CU}{}\label{wb:sec:canonical-ring-recovery}

\begin{proposition}[The canonical bundle and anticanonical pencil]
\hypertarget{tag-01PZ}{}\label{fam:src:A.1.1}\label{fam:st:A.1.1}\hypertarget{tag-03DM}{}\label{wb:thm:canonical-ring-recovery}
\StatementRecord{fam:src:A.1.1}{proved}{01PZ}{\cite[Lemma~9.10(ii), p.~70; Proposition~9.11, pp.~70--71]{AC2026}}

For the standard marked $(3,4,\infty)$ family, including every admissible
choice of its affine translations and cusp collar with the fixed fan,
the evaluation map
\[
 f^*(f_*\omega_{X/B})\longrightarrow\omega_{X/B}
\]
has zero divisor exactly $2S_2$. Consequently
\begin{equation}
 \omega_X\simeq f^*\mathcal O_B(-p_2)\otimes\mathcal O_X(2S_2)
       \simeq\mathcal O_X(-2S_2),\qquad
 \omega_X^{-2}\simeq f^*\mathcal O_B(p_2).
\label{pub:I:eq:CR1}
\end{equation}
For every integer $k$ the direct image of the vertical divisor is
\begin{equation}
 f_*\mathcal O_X(kS_2)
       =\mathcal O_B\bigl(\lfloor k/4\rfloor p_2\bigr)
\label{pub:I:eq:CR2}
\end{equation}
as subsheaves of the meromorphic-function sheaf, using pullback to identify
their sections. In particular, for each integer $m\geq0$,
\begin{equation}
 h^0(X,\omega_X^{-m})=\lfloor m/2\rfloor+1,
 \qquad h^0(X,\omega_X^m)=0\quad(m>0).
\label{pub:I:eq:CR3}
\end{equation}
There is a graded algebra isomorphism
\begin{equation}
 \mathbb C[U,V]\xrightarrow{\sim}
       \bigoplus_{m\geq0}H^0(X,\omega_X^{-m}),
       \qquad\deg U=1,\quad\deg V=2.
\label{pub:I:eq:CR4}
\end{equation}
The degree-two complete linear system has no base points and its map is
the map $f$, followed by a specified projective-coordinate map of $B$.
The degree-one system consists of the single divisor $2S_2$.
Every biholomorphism between two such threefolds induces the identity
on $B$.
\end{proposition}
\begin{proof}
\begin{enumerate}[label=\arabic*.,leftmargin=*]
\item \cite[Lemma~9.10(ii), p.~70, and Proposition~9.11,
pp.~70--71]{AC2026} give the evaluation divisor \(2S_2\) and
\(f_*\omega_{X/B}\simeq\mathcal O_B(p_0)\) for these same
marked fillings. Use the explicit base identifications
\[
 h=(t-1)^{-1},\quad \operatorname{div}(h)=p_0-p_2,\qquad
 \eta_B=dt/(t-1)^2,\quad\operatorname{div}(\eta_B)=-2p_2.
\]
These local canonical transitions hold for every admissible affine
translation and cusp collar with the fixed fan. A translation
$\zeta\mapsto\zeta+a(s)$ contributes only $ds$ terms to its
differential, which vanish after wedging with the base differential.
A cusp multiplier $x_j\mapsto h_j(t_c)x_j$ adds
$(h_j'/h_j)dt_c$ to $dx_j/x_j$, with the same cancellation.
Thus the determinant characters and evaluation divisor are unchanged.
Tensoring evaluation with \(f^*\eta_B\), and multiplying the
base line by \(f^*h\), gives \eqref{pub:I:eq:CR1}, since \(f^*p_2=4S_2\).

\item Connected compact fibres give $f_*\mathcal O_X=\mathcal O_B$ by \cite[Proposition~9.13, pp.~74--76]{AC2026}. A section of
\(\mathcal O_X(kS_2)\), viewed as a meromorphic function,
descends away from \(p_2\) to a function \(b(t)\).
On the finite cover \(t-1=s_2^4\), its pullback is
invariant under \(s_2\mapsto\zeta_2s_2\), so its Laurent
exponents are multiples of four. The exact order condition is
\[
 4\operatorname{ord}_{p_2}b+k\ge0
 \quad\Longleftrightarrow\quad
 \operatorname{ord}_{p_2}b\ge-\lfloor k/4\rfloor.
\]
Conversely every such \(b\) pulls back to an allowed section.
This is \eqref{pub:I:eq:CR2} as an equality of subsheaves, and applying it
to \(k=2m,-2m\) gives \eqref{pub:I:eq:CR3}.

\item Fix the isomorphism
\(\omega_X^{-1}\simeq\mathcal O_X(2S_2)\) from the first step.
Let \(u\) correspond to its canonical divisor section and
choose \(\ell\in H^0(B,\mathcal O_B(p_2))\) with
\(\ell(p_2)\ne0\). Let \(v\) correspond to \(f^*\ell\), so
\(u^2=f^*\sigma_{p_2}\) under the induced degree-two map.
The map \(U\mapsto u,V\mapsto v\) sends the degree-\(m\)
basis \(U^{m-2j}V^j\), \(0\le j\le\lfloor m/2\rfloor\),
to the pullback basis from \eqref{pub:I:eq:CR2}, multiplied by \(u\) when
\(m\) is odd. It is therefore the graded isomorphism \eqref{pub:I:eq:CR4}.
Replacing \(\ell\) by \(a\ell+b\sigma_{p_2}\) gives
\(v\mapsto av+bu^2\); multiplying the chosen line isomorphism
by \(c\ne0\) gives \((u,v)\mapsto(c^{-1}u,c^{-2}v)\).
The base-point-free pair \((u^2,v)\) induces exactly
\[
 x\longmapsto[\sigma_{p_2}(f(x)):\ell(f(x))].
\]
For \(\ell=A/(t-1)+B\), \(A\ne0\), the specified base map is
\(t\mapsto[t-1:A+B(t-1)]\). Degree one has the single
divisor \(2S_2\).
\item A biholomorphism preserves the anticanonical graded ring and
its complete degree-two system, so it induces a projective
automorphism of $B$. The reduced singular cusp fibre and the
finite fibres of multiplicities three and four distinguish the
three values: the pullback identity for a fibre divisor preserves
its multiplicity, and smoothness of its reduction is preserved
under the induced fibre isomorphism. Hence each of $0,1,\infty$
is fixed. A projective automorphism fixing these three points is
the identity.
\end{enumerate}
\end{proof}

\begin{remark}[Divisor classes in the Picard coordinate]
\hypertarget{tag-PI09}{}\label{I:canonical-Picard-classes}
\StatementRecord{I:canonical-Picard-classes}{proved}{PI09}{}

In the additive coordinate of \cref{I:def:Picard-coordinate},
the identification in \eqref{inv:eq:C-factor12} gives
\[
 \mathcal O_X(S_1)=L_4,\qquad \mathcal O_X(S_2)=L_3,
 \qquad \omega_X=L_{-6},\qquad \omega_X^{-2}=L_{12}.
\]
\end{remark}







\begin{proposition}[Normalized root covers and the finite fibres]\hypertarget{tag-01Q0}{}\label{fam:src:A.1.2}\label{fam:st:A.1.2}
\StatementRecord{fam:src:A.1.2}{proved}{01Q0}{the definitions and constructions stated in the proof}

The normalized base change $N_j\times_{D_j}\Delta_j\to\Delta_j$, $t_j=s_j^{m_j}$, is the original smooth torus family $J'_j$. The maps $E_1\times F_1\to S_1$ and $E_2\times F_2\to S_2$ described below are finite \'etale covers of degrees $9$ and $8$.
\end{proposition}


\begin{proof}
\begin{enumerate}[label=\arabic*.,leftmargin=*]
\item On $J'_j$, the deck generator is
$(s,\zeta)\mapsto(\zeta_js,R_j(s)\zeta+\Pi_c(g_js)v_j/m_j)$. It has exact order $m_j$. Its $r$th power changes the invariant circle coordinate $\gamma$ by $r\gamma(v_j)/m_j$; this is nonintegral for $0<r<m_j$, so the action is free. The quotient map is therefore a local biholomorphism. Its local base function is $s^{m_j}$ and its fibre multiplicity is exactly $m_j$. The map
\[
 J'_j\longrightarrow N_j\times_{D_j}\Delta_j,\qquad y\longmapsto(q_j(y),s_j(y))
\]
is finite, is an isomorphism over $s\ne0$, and has normal source. Locally at the central fibre the base-change ring has equation $z^{m_j}-s^{m_j}=0$. Its normalization is the disjoint union of branches $z=\zeta_j^rs$. The displayed map gives those branches in the covering charts. It is thus the normalization everywhere, with its specified root coordinate.

\item Put $\rho=e^{\pi i/3}$, $\mu_1=(2-\rho)/3$, $\mu_2=(1-i)/2$, $\kappa_1=\beta_c(z_1)+2\mu_1$, and $\kappa_2=\beta_c(z_2)+3\mu_2$. The fixed-point transformation laws for $\mu$ give the two displayed values directly. We have $\operatorname{Im}\kappa_j=D_c(z_j)<0$. Define
\[
 E_j=\mathbb C/(\mathbb Z+\kappa_j\mathbb Z),\quad
 F_1=\mathbb C/(\mathbb Z(\rho-2)-3\mathbb Z),\quad
 F_2=\mathbb C/(\mathbb Z(i-1)-2\mathbb Z).
\]
Let $e_1=(0,1,-2,0)^t,f_1=(0,0,-3,1)^t,e_2=(0,1,-1,0)^t,f_2=(0,0,-2,1)^t$. Direct solution of the fixed equations and inspection of the columns give
\[
 \Lambda^{A_1}=\mathbb Z\varepsilon+\mathbb Zd,\quad
 \Lambda^{A_2}=\mathbb Z\varepsilon'+\mathbb Zd,\quad
 (A_j-I)\Lambda=\mathbb Ze_j+\mathbb Zf_j.
\]
The image lattices are saturated: they are the integer solutions of $\gamma=0,2u+w+3\delta=0$ and $\gamma=0,u+w+2\delta=0$, respectively. The determinants of the fixed-plus-moving basis matrices are $-3,-2$.

\item The complex-linear maps $L_1(\xi,\eta)=(\eta,\xi-\eta/3)$ and $L_2(\xi,\eta)=(\eta,\xi-\eta/2)$ carry the product periods to those fixed and moving sublattices. They induce isogenies to the original central tori, with kernel generated respectively by $(2/3,\rho)$ and $(1/2,i)$, of exact orders three and two. For example $\widehat u=2d/3+e_1-2f_1/3=d/2+e_2-f_2/2$ gives both kernel representatives.

\item The affine deck maps lift to
\[
 (\xi,\eta)\longmapsto(\xi+\kappa_1/3,(\rho-1)\eta),\qquad
 (\xi,\eta)\longmapsto(\xi-\kappa_2/4,-i\eta).
\]
Substitution in $L_j$ gives exactly $R_jL_j+\Pi_cv_j/m_j$. They commute with their kernel translations: the moving differences are periods, since $(\rho-2)\rho=-1-\rho=-(\rho-2)-3$ can equally be checked in the lattice basis, and $(-i-1)i=1-i=-(i-1)$. More directly these assertions follow by the two-dimensional period matrices of the indicated rotations. A nonidentity element of either resulting product group translates the $E_j$ coordinate by $2a/3+b\kappa_1/3$ or $a/2-b\kappa_2/4$, respectively. Real independence of $1,\kappa_j$ proves that this translation vanishes only for the identity element. Hence the groups $C_3\times C_3$ and $C_2\times C_4$ act freely and have quotient precisely $S_j$.
\end{enumerate}
\end{proof}









\subsection{The proper holomorphic family and coherent descent}\hypertarget{tag-01PC}{}\label{fam:section:A.3}

\begin{proposition}[Coherent even shifts and the quotient family]\hypertarget{tag-01Q2}{}\label{fam:src:A.3.1}\label{fam:st:A.3.1}\label{I:coherent-descent-proof}
\StatementRecord{fam:src:A.3.1}{proved}{01Q2}{\SeriesRef{fam:src:A.2.1}}

For each $n\in2\mathbb Z$ there is a biholomorphism $F_{n,c}:X_c\to X_{c+n}$, holomorphic jointly in $c$, with $F_{2h,c+2k}F_{2k,c}=F_{2(h+k),c}$. The quotient family over $0<|w|<e^{-\pi M}$, $w=e^{-\pi ic}$, is proper and holomorphic, with pairwise non-biholomorphic fibres.

Before the quotient, every disc $Q$ with $\overline Q\subset\mathcal P$
has the proper holomorphic submersion and proper base map
\begin{equation}
 \pi:\mathcal X_Q\longrightarrow Q,\qquad
 F:\mathcal X_Q\longrightarrow B\times Q,\qquad
 \pi=\operatorname{pr}_Q\circ F.
\label{pub:I:eq:PD2}
\end{equation}
Their restrictions to $c$ are $X_c$ and $f_c$, and their three
distinguished values are the constant sections $0,1,\infty$.
\end{proposition}
\begin{proof}
\begin{enumerate}[label=\arabic*.,leftmargin=*]
\item \emph{The holomorphic family.}
Let $Q$ be a disc with $\overline Q\subset\mathcal P$.
On an ordinary lifted base chart use the period matrix
\eqref{wb:pshift:eq:periods}, jointly holomorphic in $(z,c)$.
Its real determinant is $\Im\tau\,D_c$, and its inverse is
uniformly bounded on compact base-parameter sets. Thus the lattice
action is free and properly discontinuous there. A fixed closed
real lattice parallelotope gives compact surjective representatives
over each compact subset of the base chart times $Q$.
The quotient is holomorphic and its projection to $Q$ is a submersion.
The covariance convention of \cref{I:def:affine-covariance} is
\[
 R_g(z)\Pi(z,c)=\Pi(gz,c)A_g,
 \qquad M_g=A_g^{-1},\qquad
 \Pi(gz,c)=R_g(z)\Pi(z,c)M_g.
\]
Consequently the base deck quotient is jointly holomorphic; its
proper discontinuity follows from that of the lifted base action.

At a finite root chart the family action is
\begin{equation}
 (z,c,\zeta)\longmapsto
 \left(g_jz,c,R_{g_j}(z)\zeta+
             \Pi(g_jz,c)\frac{v_j}{m_j}\right),
 \qquad (m_1,m_2)=(3,4).
\label{pub:I:eq:PD7}
\end{equation}
In the real marking it is $(z,c,u)\mapsto(g_jz,c,A_ju+v_j/m_j)$.
The invariant-covector freeness test in \cref{fam:src:A.1.2}
is independent of $c$, so these quotients are local biholomorphisms
and submersions to $Q$. The same parallelotope proves properness
over the finite base disc times $Q$. Their punctured connectors are
\begin{equation}
 \frac{\log s_j(z)}{2\pi i}\Pi(z,c)v_j.
\label{pub:I:eq:PD8}
\end{equation}
A logarithm change adds a period, so this is a jointly holomorphic
torus section. Writing $\sigma_j=(\log s_j)/(2\pi i)\,\Pi(z,c)v_j$, the
relations $A_jv_j=v_j$ and
$\log s_j(g_jz)=\log s_j(z)-2\pi i/m_j$ give
\[
 R_{g_j}(z)\sigma_j(z,c)-\sigma_j(g_jz,c)
       =\Pi(g_jz,c)v_j/m_j.
\]
Thus translation by $\sigma_j$ intertwines the finite affine action
with the middle linear action, with the specified root coordinate.

\item \emph{Uniform cusp charts and properness.}
On the fixed unimodular fan write the full cusp period matrix as
\begin{equation}
 C(t_c,c)=\begin{pmatrix}6\mu(t_c)&h(t_c)\\
 b^{\rm part}(t_c)-h(t_c)+c&\mu(t_c)\end{pmatrix},
 \qquad B_0=\begin{pmatrix}0&1\\-1&0\end{pmatrix}.
\label{pub:I:eq:PD9}
\end{equation}
Choose a closed cusp disc $|t_c|\le r_1$ on which these functions
are holomorphic, and set
\begin{equation}
 K_Q=\max_{|t_c|\le r_1,\ c\in\overline Q}
               \|-2\pi\Im C(t_c,c)\|,
 \qquad 0<\epsilon<\min(r_1,1),\quad
               |\log\epsilon|>2K_Q.
\label{pub:I:eq:PD10}
\end{equation}
The lattice acts on the toric family times $Q$ by
\begin{equation}
 (x,t_c,c)\longmapsto
 \bigl(e^{2\pi iC(t_c,c)\lambda}t_c^{B_0\lambda}x,t_c,c\bigr),
 \qquad\lambda\in\mathbb Z^2.
\label{pub:I:eq:PD11}
\end{equation}
The fan translation and nowhere-zero torus multiplier extend this
map across every toric stratum. With Euclidean operator norm,
$B_0^{\mathsf t}B_0=I$, and
\begin{equation}
 B_{t_c,c}=B_0+\frac{-2\pi\Im C(t_c,c)}{\log|t_c|},
 \quad \|B_{t_c,c}-B_0\|<\tfrac12,
 \quad \tfrac12\|\lambda\|<\|B_{t_c,c}\lambda\|
                    \le\tfrac32\|\lambda\|\quad(\lambda\ne0).
\label{pub:I:eq:PD12}
\end{equation}
These inequalities exclude stabilizers on the dense torus.
For a height-one triangle, use its three dual unimodular cone
coordinates $z_0,z_1,z_2$, with $t_c=z_0z_1z_2$.
In the open chart $|z_i|<2$, $|t_c|<\epsilon$, the normalized
logarithmic barycentric coefficients have sum one and are bounded
below by $-\log2/|\log\epsilon|$. Thus the logarithmic positions
in each fixed chart lie in a bounded set, uniformly in $c$.
For a fixed pair of charts, \eqref{pub:I:eq:PD12} bounds the finitely many
$\lambda$ which can translate one into the other. At the central
fibre a nonzero $B_0\lambda$ moves every bounded fan cell.
These charts cover the entire toric space, including its boundary.
On the dense torus the duality identity is
\[
 \log|z_i|=\log|t_c|\,\theta_i^\sigma(y),
\]
where $\theta_i^\sigma$ are the barycentric coordinates of the
height-one triangle. A triangle containing $y$ therefore gives
$|z_i|\le1$. For a boundary point in the orbit of the cone over a
cell $P$, the coordinates belonging to vertices of $P$ vanish.
The remaining characters are dual to the image cones of triangles
containing $P$ in the quotient lattice by $\operatorname{span}(P,1)$.
Those cones cover the quotient real space: lift a vector, add a
sufficiently small positive multiple of it to a relative interior
point of $(P,1)$, and rescale to height one. Local finiteness places
that point in a triangle containing $P$. Applied to the negative
logarithm of the orbit coordinates, this gives the same closed
unit-polydisc bound. Thus the open radius-two charts cover the
boundary as well. Every compact set has a finite subcover.
The chart-pair estimate applies to an intersection meeting the
boundary because that intersection is open and meets the dense
torus. It therefore proves freeness and proper discontinuity on
the entire toric family.

For properness, fix $\epsilon_1<\epsilon$ and let $K_0$ be the
finite union of closed unit polydiscs for height-one triangles
meeting the closed radius-$\sqrt2$ ball, intersected with
$|t_c|\le\epsilon_1$. For $t_c\ne0$, choose $\lambda$ with
$\|B_{t_c,c}^{-1}y+\lambda\|_\infty\le1/2$.
The translated logarithmic position has norm at most
$3/(2\sqrt2)<\sqrt2$. A triangle containing it has nonnegative
barycentric coordinates of sum one, so each corresponding cone
coordinate has modulus at most one. This gives a representative
in $K_0$. At $t_c=0$, translate the height-one component and use
the six closed toric bidiscs covering $E_0\simeq\mathrm dP_6$.
The multiplier $\exp(2\pi iC(0,c)B_0v)$ is included in this
translation; the bidiscs cover the resulting whole compact component.
For each compact $K\subset Q$, the compact set $K_0\times K$
therefore maps onto the full inverse image of
$\{|t_c|\le\epsilon_1\}\times K$. This proves properness.
The quotient charts are open subsets of a smooth toric variety
times $Q$, so projection to $Q$ is a submersion.

\item Exponentiating the logarithmic coordinates gives the cusp
connector with period identity $Z=sB_0+C(t_c,c)$.
Glue it and \eqref{pub:I:eq:PD8} to the middle family. The distinguished base
discs are disjoint. Points with distinct $(t,c)$ values separate
on the base; points with the same value lie in one Hausdorff local
piece. The resulting total space is therefore Hausdorff and smooth.
The local proper maps give a proper map $F$ to $B\times Q$:
cover a compact subset by finitely many relatively compact sets
inside the displayed base pieces. Compactness of $B$ proves
properness of $\pi=\operatorname{pr}_QF$.
The fibres and overlaps are the specified $X_c$ and $f_c$.
These local families agree on smaller cusp radii and glue over
$\mathcal P$.
\item \emph{Local maps.}
Let $n\in2\mathbb Z$. Define the following rational columns,
integral cocycles, and holomorphic translation vectors:
\begin{equation}
\begin{aligned}
 r_1&=-\frac n3(\widehat u+\widehat w),&
 \ell_1&=-n\widehat w,&
 b_1(z)&=\Pi_c(z)r_1=-\frac n3(\tau(z)+1,\mu(z))^{\mathsf t},\\
 r_2&= \frac n4(\widehat u+\widehat w),&
 \ell_2&=\frac n2\widehat u,&
 b_2(z)&=\Pi_c(z)r_2=\frac n4(\tau(z)+1,\mu(z))^{\mathsf t}.
\end{aligned}
\hypertarget{tag-03CH}{}\label{wb:pshift:eq:local-corrections}
\end{equation}
Direct multiplication of the complete matrices in \eqref{wb:pshift:eq:finite-data} gives
\begin{equation}
 (I-A_j)r_j=\frac{n\epsilon_j}{m_j}\widehat\delta+\ell_j.
 \hypertarget{tag-03CI}{}\label{wb:pshift:eq:local-matrix-identity}
\end{equation}
For $j=1$ the four entries are $(0,0,-n,n/3)$, and for $j=2$
they are $(0,n/2,0,-n/4)$. Since both $r_j$ and
$n\epsilon_j\widehat\delta/m_j+\ell_j$ lie in $\ker\gamma$,
the period covariance gives
\[
 b_j(g_jz)-R_{g_j}(z)b_j(z)
 =\Pi_{c+n}(g_jz)
       \left(\frac{n\epsilon_j}{m_j}\widehat\delta+\ell_j\right).
\]
This is the deck-equivariance identity for translation by $b_j$
between the finite fillings at $c$ and $c+n$. Its overlap with
the fixed middle identification is as follows.

Keep the marked linearizing coordinate, radius, and phase
$s_j(g_jz)=e^{-2\pi i/m_j}s_j(z)$. On a logarithm branch put
\begin{equation}
 w_j(z)=\frac{\log s_j(z)}{2\pi i},\qquad
 a_j(z)=b_j(z)+n\epsilon_j w_j(z)e.
 \hypertarget{tag-03CA}{}\label{wb:pshift:eq:finite-middle-sections}
\end{equation}
Changing the logarithm adds $n\epsilon_j\widehat\delta$ under the
period map, so $a_j$ is branch independent as a torus section.
For the specified negative generator, $w_j(g_jz)=w_j(z)-1/m_j$.
Equations \eqref{wb:pshift:eq:covariance} and \eqref{wb:pshift:eq:local-matrix-identity} consequently give the exact equality
\begin{equation}
 a_j(g_jz)=R_{g_j}(z)a_j(z)+\Pi_c(g_jz)\ell_j.
 \hypertarget{tag-03CG}{}\label{wb:pshift:eq:local-affine-cocycle}
\end{equation}
In fact the marked finite overlap maps are translations by
$\sigma_{j,c}=w_j\Pi_c v_j$, and
\begin{equation}
                  \sigma_{j,c+n}-\sigma_{j,c}
                     =n\epsilon_jw_je.
 \hypertarget{tag-03BX}{}\label{wb:pshift:eq:actual-overlap-difference}
\end{equation}
Thus $a_j$ is exactly the middle translation induced by the local
translation $b_j$, with the marked gluing maps retained.

At the cusp retain $s=\tau-h$, $e^{2\pi is}=t_c$, $s(g_0z)=s(z)-1$,
and the complete period matrix
\begin{equation}
 Z_c=sB_0+C_c(t_c),\quad
 B_0=\begin{pmatrix}0&1\\-1&0\end{pmatrix},\quad
 C_c=\begin{pmatrix}6\mu&h\\ b^{\mathrm{part}}-h+c&\mu\end{pmatrix}.
 \hypertarget{tag-03C5}{}\label{wb:pshift:eq:cusp-periods}
\end{equation}
For $\lambda=(\lambda_1,\lambda_2)^{\mathsf t}\in\mathbb Z^2$,
$C_{c+n}\lambda-C_c\lambda=(0,n\lambda_1)^{\mathsf t}$. Therefore
\begin{equation}
 \exp(2\pi iC_{c+n}\lambda)t_c^{B_0\lambda}
       =\exp(2\pi iC_c\lambda)t_c^{B_0\lambda}.
 \hypertarget{tag-03C3}{}\label{wb:pshift:eq:cusp-actions-identical}
\end{equation}
The entire cusp actions, including their nonconstant holomorphic units,
agree under the integer-shift identification.

The affine cocycle convention in \eqref{wb:pshift:eq:local-affine-cocycle} has multiplication law
\begin{equation}
                    \ell_{gh}=\ell_g+A_g\ell_h.
 \hypertarget{tag-03C0}{}\label{wb:pshift:eq:cocycle-law}
\end{equation}
Since $g_0=(g_1g_2)^{-1}$, it gives
\begin{equation}
\begin{aligned}
 \ell_1+A_1\ell_2
   &=-\frac{3n}{2}\widehat w+\frac n2\widehat\delta,\\
 \ell_0=-A_{g_0}(\ell_1+A_1\ell_2)
   &=\frac{3n}{2}\widehat w-\frac n2\widehat\delta.
\end{aligned}
\hypertarget{tag-03C4}{}\label{wb:pshift:eq:cusp-cocycle}
\end{equation}
Here $A_{g_0}$ fixes both toric lattice vectors, as is also immediate
from $A_{g_0}=\left(\begin{smallmatrix}I&0\\B_0&I\end{smallmatrix}\right)$.
Put
\begin{equation}
 q_n=(-3n/2,n/2)^{\mathsf t}\in\mathbb Z^2,
                  \qquad a_0(z)=s(z)q_n.
 \hypertarget{tag-03C6}{}\label{wb:pshift:eq:cusp-translation}
\end{equation}
Then $a_0(g_0z)-a_0(z)=-q_n=\Pi_c(g_0z)\ell_0$, exactly as
required by \eqref{wb:pshift:eq:local-affine-cocycle}. Exponentiating gives the map
$(x,t_c)\mapsto(t_c^{q_n}x,t_c)$ on the dense cusp torus.
It extends holomorphically and invertibly to the marked infinite fan:
the integral map
\[
 (y,y_3)\longmapsto(y+y_3q_n,y_3)
\]
translates each height-one triangle by the integer vector $q_n$ and
therefore permutes the cones, with inverse given by $-q_n$.
It commutes with the cusp quotient action in \eqref{wb:pshift:eq:cusp-actions-identical}. This proves extension across every cusp stratum.
\item \emph{The global additive correction.}
Let $\widetilde B^\circ$ be
the universal cover of the three-punctured base. The fundamental group
of $B^\circ$ is free on the marked two negative peripheral generators. Pull back
the marked $\Pi_c,R$ to this cover. Assign $\ell_1,\ell_2$ from
\eqref{wb:pshift:eq:local-corrections} to its free generators and extend by \eqref{wb:pshift:eq:cocycle-law}. Freeness makes
this an integral cocycle; its third peripheral value is
precisely \eqref{wb:pshift:eq:cusp-cocycle}. The maps
\begin{equation}
 (z,y)\longmapsto
       (gz,R_g(z)y+\Pi_c(gz)\ell_g)
 \hypertarget{tag-03BY}{}\label{wb:pshift:eq:affine-bundle-action}
\end{equation}
form an affine action. To check the action law, covariance gives
$R_g(hz)\Pi_c(hz)=\Pi_c(ghz)A_g$; substituting this equality
in the composition gives exactly \eqref{wb:pshift:eq:cocycle-law}.
The quotient is an affine holomorphic bundle over $B^\circ$, since
the deck action on the base is free and properly discontinuous and
one may use evenly covered base opens as affine charts.

On a finite punctured disc, \eqref{wb:pshift:eq:local-affine-cocycle} says that $a_j$ is a section of
this affine bundle. On the cusp punctured disc, \eqref{wb:pshift:eq:cusp-cocycle}--\eqref{wb:pshift:eq:cusp-translation} say
the same for $a_0$. Define a sheaf $\mathscr H_n$ on $B$ as follows.
Away from the special points it is the sheaf of sections of \eqref{wb:pshift:eq:affine-bundle-action}.
Near $p_j$, a section is required to have the form $a_j+\eta_j$,
where $\eta_j:\Delta_j\to\mathbb C^2$ is holomorphic and obeys
$\eta_j(g_jz)=R_{g_j}(z)\eta_j(z)$. Near the cusp it is required
to have the form $a_0+\eta_0$, with $\eta_0(t_c)$ holomorphic on
the full cusp disc. These are compatible definitions on the punctured
overlaps by the equal affine cocycles just proved. They also prove
local nonemptiness, including at all three special points.

The differences of these local sections form a vector sheaf
$\mathscr V$. Its ordinary sections have covariance $\eta(gz)=R_g\eta(z)$;
at the two finite points and the cusp its extensions are precisely
the holomorphic extensions specified above. Projection to the first
coordinate gives an exact sequence on the marked compact base
\begin{equation}
 0\longrightarrow\mathcal O_B e\longrightarrow\mathscr V
        \longrightarrow\mathcal L\longrightarrow0,
                  \qquad \mathcal L\simeq\mathcal O_B(-p_0).
 \hypertarget{tag-03CP}{}\label{wb:pshift:eq:vector-sheaf}
\end{equation}
The homogeneous first-coordinate line is determined by its transition factors. Indeed
the first diagonal entries in \eqref{wb:pshift:eq:covariance} are its factors $-1/\tau$ and
$1/\tau$, and at the cusp both extensions require ordinary
holomorphicity. The marked function
\[
 \Theta=\frac{(E_4\circ\tau)^2(E_6\circ\tau)^{1/2}}
                 {\Delta_{\mathrm{mod}}\circ\tau}
\]
has those exact factors, orders two and one in $s_1,s_2$, and a
simple pole in $t_c$. Dividing by $\Theta$ identifies the ordinary
and finite stalks with $\mathcal O_B$, and the cusp stalk with
$\mathfrak m_{p_0}$. This recovers the displayed $\mathcal L$
with its fixed cusp condition. 
For the finite stalks in \eqref{wb:pshift:eq:vector-sheaf}, use the full matrix $R_g$. Given an equivariant first component
$u(z)$, begin with the vector $(u(z),0)^{\mathsf t}$ and average it
under the finite group action on vector-valued holomorphic functions.
The resulting vector is equivariant, and its first component is
still $u$. Thus projection is surjective there. Its kernel is
exactly the invariant second-component functions, which are
holomorphic functions of $s_j^{m_j}=t_j$, multiplied by $e$.
The ordinary and cusp stalks are immediate. This also proves that
$\mathscr V$ is locally free of rank two: lift a local frame of
$\mathcal L$ and join it to $e$.

Both $H^1(B,\mathcal O_B)$ and
$H^1(B,\mathcal O_B(-p_0))$ vanish: on the two standard affine
charts of $\mathbb P^1$, split a holomorphic Laurent series on
their overlap into its nonnegative and negative powers. For
$\mathcal O_B$ these are the two coboundary contributions, and
for $\mathcal O_B(-p_0)$ its transition factor shifts the second
contribution by one, still leaving no missing exponent. The positive
part converges on the full first affine chart and the transformed
negative part on the full second affine chart. This proves the two
vanishings, equivalently the usual additive Cousin calculation on
these charts. The exact sequence \eqref{wb:pshift:eq:vector-sheaf} now gives
\begin{equation}
                             H^1(B,\mathscr V)=0.
 \hypertarget{tag-03CJ}{}\label{wb:pshift:eq:no-analytic-obstruction}
\end{equation}

Choose
local sections $h_\alpha$ of $\mathscr H_n$, taking the displayed
$a_j,a_0$ on sufficiently small special discs and affine charts
elsewhere. Their differences $d_{\alpha\beta}=h_\beta-h_\alpha$
are a \v Cech $1$-cocycle with values in $\mathscr V$. Project it to
$\mathcal L$ and solve that additive cocycle, using its just proved
$H^1$ vanishing. Lift the resulting local $\mathcal L$ sections to
$\mathscr V$ by the local lifting construction in \eqref{wb:pshift:eq:vector-sheaf}. Subtracting
their coboundary from $d$ leaves a scalar $\mathcal O_Be$ cocycle.
Solve that scalar additive cocycle by its $H^1$ vanishing. This
produces sections $v_\alpha$ with
$d_{\alpha\beta}=v_\beta-v_\alpha$. The local sections
$h_\alpha-v_\alpha$ agree and define a global section
$a\in\mathscr H_n(B)$. It gives holomorphic vectors $\eta_j,\eta_0$ with
\begin{equation}
 a=a_j+\eta_j\quad(j=1,2),\qquad a=a_0+\eta_0
                              \quad\hbox{at the cusp}.
 \hypertarget{tag-03CE}{}\label{wb:pshift:eq:global-correction}
\end{equation}
Any choices made in the additive solutions give maps with the same
asserted properties. Their differences as sections of this affine
bundle lie in $H^0(B,\mathscr V)=\mathbb C e$: this follows from
\eqref{wb:pshift:eq:vector-sheaf}, $H^0(B,\mathcal L)=0$, and $H^0(B,\mathcal O_B)=\mathbb C$.

Projecting $a$ modulo $\Pi_c\Lambda$ gives a holomorphic section of
the middle torus family because every $\ell_g$ in \eqref{wb:pshift:eq:affine-bundle-action} is integral.
The three local maps are therefore
\begin{equation}
\begin{aligned}
 F_{\mathrm{mid}}[z,\zeta]&=[z,\zeta+a(z)],\\
 F_j[z,\zeta]&=[z,\zeta+b_j(z)+\eta_j(z)]\quad(j=1,2),\\
 F_0[x,t_c]&=[\exp(2\pi i\eta_0(t_c))t_c^{q_n}x,t_c].
\end{aligned}
\hypertarget{tag-03CD}{}\label{wb:pshift:eq:full-isomorphism}
\end{equation}
The first descends by \eqref{wb:pshift:eq:affine-bundle-action} and equality of the integer-shift
lattices. The finite maps descend by
\eqref{wb:pshift:eq:local-matrix-identity}, the finite
deck-equivariance identity above, and the exact covariance of $\eta_j$. The cusp map descends by \eqref{wb:pshift:eq:cusp-actions-identical}; its
holomorphic-unit factor extends by torus multiplication on every
toric chart, and its integer factor extends by the fan map in
\eqref{wb:pshift:eq:cusp-translation}. Each map is a biholomorphism, with inverse given by the
negative translation or the inverse torus multiplier and fan map.

On a finite overlap, write the marked gluing as
$h_{j,c}(\zeta)=\zeta+\sigma_{j,c}$. Then \eqref{wb:pshift:eq:actual-overlap-difference}, \eqref{wb:pshift:eq:global-correction}, and
\eqref{wb:pshift:eq:full-isomorphism} give the equality of torus maps
\begin{equation}
                     h_{j,c+n}F_j=F_{\mathrm{mid}}h_{j,c}.
 \hypertarget{tag-03CC}{}\label{wb:pshift:eq:finite-overlap-identity}
\end{equation}
Both sides add
$\sigma_{j,c}+a=\sigma_{j,c+n}+b_j+\eta_j$.
On the cusp overlap the marked map is exponentiation of $\zeta$;
\eqref{wb:pshift:eq:cusp-translation} and \eqref{wb:pshift:eq:global-correction} give
$\exp(2\pi ia)=t_c^{q_n}\exp(2\pi i\eta_0)$, exactly the
third line of \eqref{wb:pshift:eq:full-isomorphism}. Thus that overlap also agrees, including
all nonconstant period units. The maps and their inverses glue to
a global biholomorphism $X_c\to X_{c+n}$ over $B$.

The choices can be made coherently for every parameter and every even
shift. Both initial cocycle columns in \eqref{wb:pshift:eq:local-corrections} have $\gamma=0$.
Because every $A_g$ preserves $\gamma$, \eqref{wb:pshift:eq:cocycle-law} shows that
$\gamma(\ell_g)=0$ for every word. Therefore the affine term
$\Pi_c(gz)\ell_g$ in \eqref{wb:pshift:eq:affine-bundle-action} is independent of $c$. The displayed
$a_j,a_0$, and the vector sheaf $\mathscr V$ are also independent
of $c$. Perform the additive construction once for $n=2$, keeping
one resulting section $a^{(2)}$. For $n=2k$ choose
$a^{(2k)}=k a^{(2)}$ and multiply its local correcting vectors by
$k$. All their affine cocycles and local formulas are then exactly
those for $n=2k$. The resulting maps satisfy the full group law
\begin{equation}
 F_{2h,c+2k}\circ F_{2k,c}=F_{2(h+k),c},\qquad
                 F_{0,c}=\mathrm{id}_{X_c}.
 \hypertarget{tag-03C1}{}\label{wb:pshift:eq:coherent-shift-action}
\end{equation}
This follows by adding the translation vectors in the first
two lines of \eqref{wb:pshift:eq:full-isomorphism}, and by multiplying their units and adding the
integer exponents in its third line. The formulas are jointly
holomorphic in $c$, since their corrections are independent of it.
\item The local families constructed above are proper and holomorphic.
The coherent $2\mathbb Z$ action is properly
discontinuous because its action on $\mathcal P$ is.
Local inverse logarithms for $w=e^{-\pi ic}$ identify its quotient
with the proper local families just proved.
Proposition~\ref{fam:src:A.2.1} gives pairwise non-biholomorphic fibres.
\end{enumerate}
\end{proof}

\subsection{Albanese invariants and the derivative of the deformation}\hypertarget{tag-01PD}{}\label{fam:section:A.4}

\begin{proposition}[Albanese curves and parameter recovery]\hypertarget{tag-01Q3}{}\label{fam:src:A.4.1}\label{fam:st:A.4.1}
\StatementRecord{fam:src:A.4.1}{proved}{01Q3}{\SeriesRef{fam:src:A.1.2}, \SeriesRef{fam:src:A.2.1}}

With $\kappa_j(c)=\kappa_j^{\rm part}+c$, the Albanese elliptic curves of the reduced finite fibres are
\[
 \operatorname{Alb}(S_1)=\mathbb C/(\mathbb Z+\kappa_1(c)\mathbb Z),\qquad
 \operatorname{Alb}(S_2)=\mathbb C/(\mathbb Z+\tfrac12\kappa_2(c)\mathbb Z).
\]
Writing $J(q)=q^{-1}+744+196884q+\cdots$, their invariants are
\[
 j_1(w)=J(e^{-2\pi i\kappa_1^{\rm part}}w^2),\qquad
 j_2(w)=J(e^{-\pi i\kappa_2^{\rm part}}w).
\]
These are convergent formulas because the displayed nome moduli are below one. Outside the explicitly defined countable exceptional set in the proof, equality of both Albanese invariants forces $c'-c\in2\mathbb Z$. Thus the generic Albanese invariants recover the parameter modulo its exact period subgroup.
\end{proposition}


\begin{proof}
\begin{enumerate}[label=\arabic*.,leftmargin=*]
\item In the product covers of \ref{fam:src:A.1.2} the invariant holomorphic one-forms are exactly the multiples of $d\xi$: the moving factor is multiplied by a primitive third or fourth root. Integration thus identifies the Albanese with the quotient of the fixed line by the fixed-coordinate projections of every deck translation. These are, respectively,
\[
 \mathbb Z+\kappa_1\mathbb Z+\tfrac23\mathbb Z+\tfrac{\kappa_1}3\mathbb Z
   =\tfrac13\mathbb Z+\tfrac{\kappa_1}3\mathbb Z,
\quad
 \mathbb Z+\kappa_2\mathbb Z+\tfrac12\mathbb Z+\tfrac{\kappa_2}4\mathbb Z
   =\tfrac12\mathbb Z+\tfrac{\kappa_2}4\mathbb Z.
\]
Multiplication by three and two gives the stated maps of complex tori. The two displayed lattices include both the product-isogeny kernels and the affine quotient translations. Using the positive-imaginary representatives $-\kappa_1$ and $-\kappa_2/2$ in the standard modular $j$-series gives the formulas. At sufficiently small nome $J$ is injective: $1/J(q)=q+O(q^2)$ has nonzero derivative at zero. Thus $j_2$ alone distinguishes parameters within a sufficiently small punctured subdisc when both parameters lie there. For the full generic assertion put $x=-\kappa_2(c)/2$ and $d_*=\kappa_1^{\rm part}-\kappa_2^{\rm part}$, so $-\kappa_1(c)=2x-d_*$. The definitions $\kappa_1^{\rm part}=\beta^{\rm part}(z_1)+2\mu_1$ and
$\kappa_2^{\rm part}=\beta^{\rm part}(z_2)+3\mu_2$, together with
\cref{ari:C.1.2}, give
\[
 \operatorname{Im}d_*=
 -\frac{3}{2\pi}\log48+\frac3\pi\log6
 =-\frac{3}{2\pi}\log\frac43<0.
\] Equality of the two elliptic invariants at $c,c'$ gives matrices $M,N\in\operatorname{PSL}_2(\mathbb Z)$ with $x'=M(x)$ and
\[
 N(2x-d_*)=2M(x)-d_*.
\]
Define the exceptional set to be the $c\in\mathcal P$ obtained from roots of these equations for all pairs for which the two rational functions are not identical. Each equation has at most two roots after clearing denominators, so this is a countable, explicitly specified set. For a pair giving an identity, let $C(x)=2x-d_*$. Then $N=CMC^{-1}$. The real matrix $M$ preserves both the real circle $\mathbb R\cup\{\infty\}$ and $C^{-1}(\mathbb R\cup\{\infty\})$, the parallel line with imaginary part $\operatorname{Im}d_*/2$ together with infinity. Their unique intersection is infinity, so $M$ fixes infinity. An element of $\operatorname{PSL}_2(\mathbb Z)$ fixing infinity is $x\mapsto x+b$, $b\in\mathbb Z$, and then $N(y)=y+2b$. It follows that $c'-c=-2b$. This proves parameter recovery from the two Albanese invariants outside the specified exceptional set. For biholomorphic total spaces, \cref{fam:src:A.2.1} gives the parameter congruence at every point.
\end{enumerate}
\end{proof}



\begin{remark}[Period variation and norm comparisons]\hypertarget{tag-01Q5}{}\label{fam:src:A.4.3}\label{fam:st:A.4.3}
\StatementRecord{fam:src:A.4.3}{proved}{01Q5}{\SeriesRef{fam:src:A.4.2}, \SeriesRef{fam:src:A.1.1}, \SeriesRef{fam:src:A.2.1}, \SeriesRef{fam:src:A.3.1}; \cite[Proposition 4.7]{ViaclovskyLog2026}; \cite[Lemmas 6.6--6.9 and Proposition 6.15]{Yang2026}}

\begin{enumerate}[label=\arabic*.,leftmargin=*]
\item On a fixed ordinary fibre, the three last periods remain fixed and span $\mathbb C^2$, whereas the first changes by $\delta e$. A marked complex-linear identification fixing those three periods must be the identity, so the marked period map is nonconstant. The set of integral markings of one torus is countable. A holomorphic nonconstant map from a disc has uncountable image; hence the unmarked fibre family is locally non-isotrivial. \ref{fam:src:A.1.1} transfers this to total spaces. This is the period argument of \cite[Proposition~4.7, pp.~14--15]{ViaclovskyLog2026}. A nonconstant classifying map gives positive dimension of its local target. It does not by itself determine the derivative at the base point, as the explicit holomorphic map $z\mapsto z^2$ has nonconstant image and derivative zero at zero. \ref{fam:src:A.4.2} computes the derivative here.

\item The comparison in \cite[Lemmas~6.6--6.7 and~6.9, pp.~31--32; Remark~6.8, p.~32; Proposition~6.15, pp.~37--38]{Yang2026} uses an integral commuting lattice map, an affine fibre lift, normalization of the root base change, and a norm-lattice defect. His Proposition 6.15 obtains $6/r$ from the first-column scaling and its order-six norm image. Here the first-column scaling is one and the two computed norm images impose parity, yielding two. These matrix substitutions identify the commuting maps, finite norm equations and collar corrections with the ones just displayed. The nonzero absolute Kodaira--Spencer conclusion in \cite[Theorem~6.18, p.~39]{Yang2026} is obtained for the present family by the explicit computation and intrinsic-pencil transfer in \cref{fam:src:A.4.2}.
\end{enumerate}
\end{remark}

\subsection{The standard isomorphism classes}\hypertarget{tag-01PE}{}\label{fam:section:A.5}

\begin{theorem}[Classification of the standard family]\hypertarget{tag-01Q6}{}\label{fam:src:A.5.1}\label{fam:st:A.5.1}\label{wb:pshift:thm:classification}\hypertarget{tag-000C}{}\label{app:rederive:fam-src-A.5.1}
\StatementRecord{fam:src:A.5.1}{proved}{01Q6}{\SeriesRef{app:rederive:fam-src-A.5.1}; \citet[Theorem~1.4, p.~3]{DuEtAl2026v1}}

Fix the original cusp-regular period solution, the clockwise peripheral marking, the prescribed periodic triangular cusp fan, the affine vectors $(\varepsilon,-\varepsilon')$, and the untwisted cusp collar. For every $c,c'\in\mathcal P$, the constructed compact complex threefolds satisfy
\[
 X_c\simeq X_{c'}\quad\Longleftrightarrow\quad c'-c\in2\mathbb Z.
\]
Every biholomorphism is over the identity of the original $\mathbb P^1$. The proper holomorphic family of their isomorphism classes has parameter $0<|w|<e^{-\pi M}$ and nonzero Kodaira--Spencer class at every point.
\end{theorem}
\begin{proof}
For the standard twist data, the classification of
\cite[Theorem~1.4, p.~3]{DuEtAl2026v1}, with the exact parameter
identification of \cref{night:P5.7}, is re-derived in
\cref{night:P5.6}. It gives $X_c\simeq X_{c'}$ exactly when
$c-c'\in2\mathbb Z$. The explicit maps on the regular, finite,
and cusp charts are those of \cref{fam:src:A.3.1}, displayed in
\eqref{wb:pshift:eq:full-isomorphism} and satisfying
\eqref{wb:pshift:eq:coherent-shift-action}.

As a complex parameter space for the isomorphism classes represented
by this marked family, the quotient $\mathcal P/(2\mathbb Z)$ has
the explicit coordinate
\begin{equation}
 w=\exp(-\pi ic),\qquad 0<|w|<\exp(-\pi M).
 \hypertarget{tag-03CL}{}\label{wb:pshift:eq:parameter-disc}
\end{equation}
Indeed $|w|=\exp(\pi\operatorname{Im}c)$, and two complex numbers
have the same value of this exponential exactly when they differ
by an even integer. Every point in the displayed punctured disc
has a logarithm giving a $c\in\mathcal P$. The exponential has
nonzero derivative, so these local inverse logarithms give the
claimed quotient complex structure. The integer-period quotient
seen by the middle family instead has coordinate $w^2$; it is
covered two-to-one by \eqref{wb:pshift:eq:parameter-disc}. Its two lifts are distinguished by
the exact order-two summand in \eqref{wb:pshift:eq:finite-class}. The resulting parameter space classifies the members of the marked family.
Proposition~\ref{fam:src:A.3.1} proves properness of this quotient
family and its joint holomorphic dependence. Its Kodaira--Spencer
class is nonzero by Proposition~\ref{fam:src:A.4.2}.
\end{proof}




\begin{remark}[A rank-one nodal model]
\hypertarget{tag-IP09}{}\label{I:remark:rank-one-nodal-model}
\StatementRecord{I:remark:rank-one-nodal-model}{proved}{IP09}{}
The family of \cref{fam:src:A.5.1,fam:src:A.5.2} uses the original period
system with multiplicities $(3,4)$. Its arbitrary affine and cusp data are
classified in \cref{fam:section:B.2}, and its higher multiplicities in
\cref{fam:section:B.3}.
For the elliptic surface $y^2=x^3+tx+1$ of
\cite[Section~2, pp.~3--5; Proposition~3.1, pp.~6--7;
Proposition~6.1, pp.~17--18]{ViaclovskyTwo2026}, fix the nodal value
$p_1$ with $4p_1^3+27=0$. In that paper's equation~(3.1), put
$\gamma_0=\gamma(0)$, $\beta_0=\beta(0)$ and
$c_0=\lambda\delta(0)$, with $0<|c_0|<1$.
The normalization of this rank-one filling is
\[
 \begin{gathered}
 E'=\mathbb C^*/\langle c_0\rangle,\qquad
 R=\mathbb P_{E'}(\mathcal O_{E'}\oplus L),\\
 R=(\mathbb P^1\times\mathbb C^*)/
       \langle(z,v)\mapsto(\gamma_0z,c_0v)\rangle,
 \end{gathered}
\]
and its boundary sections are identified by $v\mapsto\beta_0v$.
The invariant nodal volume form in equation~(6.2) of that paper
trivializes the ambient canonical line near this fibre. The branch-normal
map therefore has quotient $\nu_*K_R$ by the local coefficient map in
\eqref{fam:eq:ab-normal-quotient}.
For the ruling $\pi:R\to E'$, one has
$\pi_*\mathcal O_R=\mathcal O_{E'}$ and
$R^1\pi_*\mathcal O_R=0$; Leray and Serre duality give
$h^\bullet(K_R)=(0,1,1)$.
The corresponding summand of the connecting homomorphism is
\[
 H^1(R,K_R)=\mathbb C\longrightarrow H^2(T_X(-\log D)),
\]
where $D$ is the union of that model's three reduced nodal fibres.
Its defining boundary-translation \v Cech cocycle uses the full
holomorphic units in equation~(3.1), their central values
$(\gamma_0,c_0,\beta_0)$ and the vertical frames. The integral
comparison equations use that model's four monodromy matrices and
their collar defects. These are the explicit local and global inputs
for this model; the displayed cohomology calculation evaluates the
source of the connecting map.
\end{remark}

\section{Discrete data and the geometric parameter}\hypertarget{tag-01PF}{}\label{fam:section:B}

\subsection{All affine vectors and cusp sections}\hypertarget{tag-01PH}{}\label{fam:section:B.2}

\begin{proposition}[Classification with arbitrary affine and cusp data]\hypertarget{tag-01QB}{}\label{fam:src:B.2.1}\label{fam:st:B.2.1}
\StatementRecord{fam:src:B.2.1}{proved}{01QB}{\SeriesRef{fam:src:A.1.1}, \SeriesRef{fam:src:A.3.1}; \cite[Section~6.4, pp.~39--40]{AC2026}}

Let
\[
 w_1=r_1\varepsilon+k_1d,\qquad w_2=r_2\varepsilon'+k_2d,
 \qquad 3\nmid r_1,\quad 2\nmid r_2,
\]
and use an arbitrary holomorphic cusp collar section $\sigma$ in the construction of \cite[Section~6.4, pp.~39--40]{AC2026}. Write $l_0=\gamma(c(\sigma))$ in its positive cusp-winding convention, and set $p=12l_0-4r_1-3r_2$. Let primes denote a second such datum with the same period solution and cusp fan. Then
\[
 X(c;w_1,w_2,\sigma)\simeq X(c';w'_1,w'_2,\sigma')
\]
holds exactly for the following arithmetic conditions:
\[
 n=c'-c\in\mathbb Z,\qquad
 p'=a p\text{ for some }a\in\{1,-1\},\qquad
 n\equiv k_2-k'_2\pmod2.
\]
All maps are over $B$. A map with the indicated $a$ has fibre-linear part $aI_2$. In particular the complete invariants are
\[
 \bigl(|p|,\ e^{-\pi i(c+k_2)}\bigr),\qquad
 |p|\ge1,\quad\gcd(|p|,6)=1.
\]
Every value of this positive integer occurs. The full collar section enters the construction, but its additional analytic part and its second integral winding contribute no further invariant.
\end{proposition}


\begin{proof}
\begin{enumerate}[label=\arabic*.,leftmargin=*]
\item The determinant calculation in \ref{fam:src:A.1.1} was independent of translations, so the base is still intrinsic. The centralizer computation of \cref{fam:src:A.2.1} and its period equation \eqref{wb:pshift:eq:period-morphism} give $c'=c+n$ and fibre-linear part $aI_2$. Put $S_n=I+nE$ and compare all target periods in source coordinates. On the finite root cover the exact deck equation is
\begin{equation}
 B_j(g_jz)-R_jB_j(z)=\Pi_c(g_jz)\left(\frac{d_j}{m_j}+\ell_j\right),
 \quad d_j=S_nw'_j-a w_j,
 \quad N_j\ell_j=-d_j.
\label{pub:I:eq:B2-1}
\end{equation}
Necessity follows by the normalized-base-change lift already proved. Conversely the norm equation has a rational solution $r_j^*$ of
$(I-A_j)r_j^*=d_j/m_j+\ell_j$, since $\ker N_j$ over $\mathbb Q$ equals the image of $I-A_j$. We can choose $\gamma(r_j^*)=0$: the columns $\widehat u,\widehat w$ already span that image over $\mathbb Q$. Then $B_j=\Pi_cr_j^*$ is holomorphic on the whole root disc and solves the exact equation. Thus the norm condition supplies a biholomorphism on the complete root disc.

\item The computed norm images translate \eqref{pub:I:eq:B2-1} into
\begin{equation}
 r'_1\equiv ar_1\pmod3,\quad r'_2\equiv ar_2\pmod4,
 \quad ak_2-k'_2-nr'_2\equiv0\pmod2.
\label{pub:I:eq:B2-2}
\end{equation}
The integer $k_1$ is arbitrary. Since $r'_2$ is odd, the last condition is the displayed parity. On the finite collar the middle translation is
$A_j=B_j+(\log s_j)/(2\pi i)\,\Pi_cd_j$ and has integral defect $\ell_j$. Its $\gamma$ defect is
$\alpha_j=(ar_j-r'_j)/m_j$.

\item A holomorphic isomorphism between cusp neighbourhoods with fibre-linear part $aI_2$ has, after lifting near a generic smooth central point to the toric covers, the form
\begin{equation}
 (x,t_c)\longmapsto(t_c^q u(t_c)x^a,t_c),\qquad q\in\mathbb Z^2,
 \quad u(0)\in(\mathbb C^*)^2.
\label{pub:I:eq:B2-3}
\end{equation}
Here the assertion concerns this local lift near a generic central point. To prove it, take a source chart near the generic point of its component and a target chart near its image. The \'etale toric quotient permits a local lift over a simply connected neighbourhood. The torus characters in a chart at the height-one ray $v$ are $x_i=t_c^{v_i}X_i$, with $X_i$ invertible. On $t_c\ne0$ the affine fibre equation gives $x'_i/x_i^a$ independent of the $X$ coordinates. Its expression in the two charts is therefore $t_c^{v'_i-av_i}$ times a holomorphic unit of $t_c$, including zero. This proves \eqref{pub:I:eq:B2-3}. Taking logarithms shows that its winding is in $\Lambda_{\rm tor}=\mathbb Z\widehat w+\mathbb Zd$. Changing a lift adds $(I-A_0)\Lambda=\Lambda_{\rm tor}$, so the vanishing of its class in $\bar\Lambda=\Lambda/\Lambda_{\rm tor}$ is intrinsic. Thus every extending translation has the stated integral winding and holomorphic unit factor.

\item The map from the cusp piece to the middle family is $T_{-\sigma}G_0^{-1}$. Consequently the middle translation induced by \eqref{pub:I:eq:B2-3} is
\[
 A_0=sq+a\sigma-\sigma'+\eta_0,
\]
with $\eta_0$ holomorphic on the cusp disc after a logarithm of the unit. For the clockwise generator $s\mapsto s-1$, the class of its defect in $\bar\Lambda$ is $-a c(\sigma)+c(\sigma')$. On the other hand every middle integral cocycle satisfies
$\ell_0=-A_0(\ell_1+A_1\ell_2)$, whose $\gamma$ coordinate is $-\alpha_1-\alpha_2$. Equality of these coordinates is exactly
\[
 l'_0-a l_0=-(ar_1-r'_1)/3-(ar_2-r'_2)/4,
\]
or $p'=ap$. This proves the remaining necessity. Notice that $p\equiv-r_1\pmod3$ and $p\equiv r_2\pmod2$, so $p$ is nonzero and coprime to six.

\item For sufficiency, $p'=ap$ implies the first two congruences of \eqref{pub:I:eq:B2-2} by reduction modulo three and four. The parity supplies the remaining one. Choose integral solutions $\ell_1,\ell_2$ of \eqref{pub:I:eq:B2-1}. Their $\gamma$ coordinates already have the required values. Adding any element of $\ker N_1=(I-A_1)\Lambda$ preserves \eqref{pub:I:eq:B2-1}. In particular adding an arbitrary integer multiple of $e_1=(0,1,-2,0)^t$ changes the $\widehat u$ coordinate of $\ell_0$ by an arbitrary integer while leaving its $\gamma$ coordinate fixed. We can therefore arrange
$[\ell_0]=-a c(\sigma)+c(\sigma')$ in the entire rank-two group $\bar\Lambda$. Choose lifts of $\sigma,\sigma'$ on the cusp half-plane. Their remaining defect differs from $\ell_0$ by a toric lattice vector; an integral $q$ makes $sq+a\sigma-\sigma'$ have exactly defect $\ell_0$.

\item The finite translations and this cusp reference now define a locally nonempty affine torsor with linear difference sheaf $\mathscr V$ defined by \eqref{wb:pshift:eq:vector-sheaf} in the proof of \cref{fam:src:A.3.1}. The specified collar reference contains the full terms $a\sigma-\sigma'$, including any essential singularities of their chosen lifts at the puncture; differences of allowed extensions are holomorphic $\eta_0$. The proof of $H^1(B,\mathscr V)=0$ gives a global middle translation. The finite maps descend by \eqref{pub:I:eq:B2-1}, and the cusp map is $x\mapsto t_c^qu(t_c)x^a$. The fan is preserved by all integer translations and by $y\mapsto-y$; exponentials of the integer period shift agree. Hence this cusp map and its inverse extend over every stratum. The finite overlap is the equality defining $A_j$, and the cusp overlap is the equality defining $A_0$. Thus the maps glue globally.

\item Finally parity is equivalent to equality of the two displayed exponentials. Passing from $p$ to $-p$ is effected by $a=-1$ and the same proof; for fixed sign the congruences have just been shown sufficient. Every positive integer coprime to six is obtained by choosing $r_1\in\{1,2\}$ and odd $r_2$ in the appropriate residue classes modulo three and four, and then taking $l_0=(p+4r_1+3r_2)/12$. The construction in the preceding steps absorbs the analytic part and the second winding of the collar.
\end{enumerate}
\end{proof}

\begin{proposition}[The two classes over a fixed linearization parameter]\hypertarget{tag-01QC}{}\label{fam:src:B.2.2}\label{fam:st:B.2.2}
\StatementRecord{fam:src:B.2.2}{proved}{01QC}{\SeriesRef{fam:src:A.5.1}, \SeriesRef{fam:src:B.2.1}, \SeriesRef{fam:src:A.2.1}}

For multiplicities $(3,4)$ and $|p|=1$, every admissible discrete member is biholomorphic to one member of \ref{fam:src:A.5.1}, uniquely modulo $2\mathbb Z$. At fixed nonzero admissible Engel parameter $u$ there are exactly two isomorphism classes, exchanged by a positive circuit of $u$ about zero. In the convention that Engel's root generator is the inverse of our clockwise one and his circle class is $d$, the identity of the middle quotient sends
\[
 X_{c+1}\longleftrightarrow X(u;a_0-d/3,\ a_1+d/4).
\]
Reversing both declared root generators reverses these two displayed signs.
\end{proposition}


\begin{proof}
\begin{enumerate}[label=\arabic*.,leftmargin=*]
\item Take the canonical target data $(r'_1,r'_2,k'_1,k'_2,l'_0)=(1,-1,0,0,0)$, for which $p'=-1$. Choose $a$ so that $p'=ap$, and an integer $n\equiv k_2\pmod2$. \ref{fam:src:B.2.1} supplies the isomorphism to $X_{c+n}$. The uniqueness is \ref{fam:src:A.5.1}. At fixed $u=e^{-2\pi ic}$, its logarithms differ by integers; the invariant $e^{-\pi i(c+k_2)}$ takes exactly the two square roots of $u$. Both occur, already as $X_c,X_{c+1}$, and \ref{fam:src:A.2.1} distinguishes them.

\item Put $P_{\rm E}=P_{\rm mark}$, the integral isomorphism of
\cref{inv:B1.0}, and retain Engel's $Cf_4=e_1+e_4$.
Besides the intertwining identities in that proposition, multiplication gives
$P_{\rm E}d=\delta_0$, $\psi P_{\rm E}=-\gamma$,
$P_{\rm E}\varepsilon=-e_4$ and
$P_{\rm E}\varepsilon'=\delta_0-Cf_4$.
They also agree with the complex periods: at $z_1$ the complex-linear map
$(\zeta_1,\zeta_2)\mapsto(-\zeta_1/(\rho+1),\zeta_2+\zeta_1/3)$ carries $\Pi_cP_{\rm E}^{-1}$ to Engel's columns $(1,0),(\rho-1,0),((1+\rho)/3,1/3),(0,-\kappa_1)$. At $z_2$ his $Cf_4$ has preimage $d-\varepsilon'$ and fixed period $1-\kappa_2$. Thus the comparison fixes the positive circle class and the local complex structures, in addition to their monodromy matrices.

\item The inverse-generator convention of \cref{inv:B1.0} gives $a_{j-1}=-P_{\rm E}v_j/m_j$ in Engel's lattice marking (and its period image in the torus). Integer shift gives $\Pi_{c+1}v_j=\Pi_c(v_j+\gamma(v_j)d)$, with $\gamma(v_1)=1,\gamma(v_2)=-1$, proving the two signs. A positive $u$ circuit changes $\log u$ by $2\pi i$ and hence changes $c$ by $-1$, exchanging the square roots. This gives the claimed permutation with its orientation.
\end{enumerate}
\end{proof}

\begin{proposition}[The double cover of the linearization parameter]\hypertarget{tag-01QD}{}\label{fam:src:B.2.3}\label{fam:st:B.2.3}
\StatementRecord{fam:src:B.2.3}{proved}{01QD}{\SeriesRef{fam:src:A.3.1}, \SeriesRef{fam:src:A.2.1}, \SeriesRef{fam:src:B.2.1}, \SeriesRef{fam:src:A.4.1}; \cite[Remark 8.1]{Engel2026}}

Fix the sphere logarithmic data in \cite[Remark~8.1, p.~22]{Engel2026}.
The resulting family has a proper holomorphic realization on every
simply connected admissible $u$-open set and globally over the punctured
$w$-disc of \cref{fam:src:A.3.1}, with $u=w^2$.
The order-two class of $d$ in $\Lambda^{A_2}/N_2\Lambda$ obstructs
descent of this family to the punctured $u$-disc.
\end{proposition}


\begin{proof}
Local inverse logarithms of $u$ give the proper local families of
\cref{fam:src:A.3.1}; its coherent even shifts give their global
family over the punctured $w$-disc.
Descent through $u=w^2$ would furnish an isomorphism between the
two pulled-back fibres at $w$ and $-w$, namely $X_c$ and $X_{c+1}$.
\Cref{fam:src:A.2.1} excludes that isomorphism by the nonzero class
of $d$ in $\Lambda^{A_2}/N_2\Lambda$. This class has exact order two
on the normalized finite cover. The same obstruction applies to
all sphere logarithmic data by \cref{fam:src:B.2.1}, because $r_2$ is odd.
\end{proof}

\subsection{Higher multiplicities}\hypertarget{tag-01PI}{}\label{fam:section:B.3}

\begin{proposition}[Classification for higher multiplicities]\hypertarget{tag-01QE}{}\label{fam:src:B.3.1}\label{fam:st:B.3.1}
\StatementRecord{fam:src:B.3.1}{proved}{01QE}{\SeriesRef{fam:src:A.1.2}, \SeriesRef{fam:src:A.2.1}, \SeriesRef{fam:src:B.2.1}}

Let $k,l$ be positive integers with $\gcd(3k,4l)=1$, and put $m_1=3k,m_2=4l$. Consider Viaclovsky's sphere members expressed in the same clockwise marking, with $w_j=r_j\varepsilon_j+s_jd$, $\gcd(r_j,m_j)=1$, cusp winding $l_0$, and $p=m_1m_2l_0-m_2r_1-m_1r_2=\pm1$, where $\varepsilon_1=\varepsilon,\varepsilon_2=\varepsilon'$. Two such members are biholomorphic precisely when their multiplicity pair agrees and, for some $a\in\{1,-1\}$ and integer $n=c'-c$,
\[
 p'=ap,\quad r'_j\equiv ar_j\pmod{m_j},\quad
 nr'_1\equiv as_1-s'_1\pmod{k},\quad
 nr'_2\equiv as_2-s'_2\pmod{2l}.
\]
The congruences on $r'_j$ already follow from $p'=ap$. For fixed discrete vectors and cusp collar the identification subgroup is $2kl\mathbb Z$. For fixed admissible $u$ and this multiplicity pair, all the stated sphere translation data give exactly $2kl$ isomorphism classes, permuted in one cycle by a circuit of $u$ about zero.
\end{proposition}


\begin{proof}
\begin{enumerate}[label=\arabic*.,leftmargin=*]
\item The base is intrinsic for these sphere members. A biholomorphism $F:X\to X'$ gives $g=f'F$. Since $H^2(S^6,\mathbb R)=0$, the closed form $g^*\omega_{\rm FS}$ is exact. Restrict to a regular torus $T$ and wedge with a flat K\"ahler form on $T$. Stokes' theorem gives integral zero, while the integrand is nonnegative and strictly positive wherever $dg|_T\ne0$. Thus $g$ is constant on each regular fibre. The image of $(f,g)$ is a compact analytic curve whose projection to $B$ has degree one; its normalization gives a holomorphic $h:B\to B$ with $g=hf$. Apply the same construction to $F^{-1}$ to see that $h$ is an automorphism. The three homological local monodromy orders are three, four and infinity, independently of $k,l$; their intertwining maps under $F$ therefore fix each special point. Hence $h=1$. The nonreduced fibre multiplicities now force equal $k,l$. This is the base-preservation argument of \cite[Proposition~4.7, pp.~14--15]{ViaclovskyLog2026}, with the displayed maps.

\item The root model of degree $m_1$ is the pullback of the degree-three model by $s\mapsto s^k$, and likewise at the other point. Its linear action remains $A_1$ or $A_2$, while its root generator has order $m_j$. The primitive invariant coordinate $r_j/m_j$ proves freeness of every nonidentity power. The normalization and deck-commutation arguments of \ref{fam:src:A.1.2} and \ref{fam:src:A.2.1} apply to these root covers, preserving the chosen generator. Their norms are
\[
 N_{3k}=kN_1,\qquad N_{4l}=lN_2,
\]
so the fixed-coordinate norm images are $3k\mathbb Z\varepsilon+k\mathbb Zd$ and $4l\mathbb Z\varepsilon'+2l\mathbb Zd$. The norm equation \eqref{pub:I:eq:B2-1}, with $m_j$ replaced by these covering degrees, gives exactly the displayed congruences. On ordinary fibres the affine lift has simultaneous rational centralizer
$\mathbb Q I+\mathbb Q E$, as in \cref{fam:src:A.2.1}.
Comparison of the constant period columns gives fibre-linear part $aI$,
integer shift $n$ and lattice map $aS_n^{-1}$, where $S_n=I+nE$.
Functorial normalization of the root base change makes the lift commute
with the prescribed deck generator; summing its defect gives the norm
equation. The explicit base norms are
\[
 N_3v=3\gamma(v)\varepsilon+(2v_u+v_w+3v_\delta)d,\qquad
 N_4v=4\gamma(v)\varepsilon'+(2v_u+2v_w+4v_\delta)d,
\]
where $N_3=N_1$ and $N_4=N_2$ in the notation above.
At the cusp the character ratio of any extending map is $t_c^q$ times
a holomorphic unit, by \cref{fam:src:B.2.1}; its winding vanishes in
the quotient by the toric directions. The middle cocycle then forces
$p'=ap$.

\item The integer $p$ records the primitive sphere relation in this marking. The common torus coinvariants are detected by $\gamma$, since $(A_1-I)\Lambda+(A_2-I)\Lambda=\ker\gamma$ by their explicit saturated bases. The finite meridian relations are $x^{m_1}=c^{r_1}$, $y^{m_2}=c^{r_2}$, and the cusp relation is $xy=c^{l_0}$ with $c$ central. Eliminating $y$ gives an abelian relation matrix with determinant $-p$ and primitive first elementary divisor because $\gcd(m_1,m_2)=1$. Thus the sphere data have $p=\pm1$. Reducing $p'=ap$ modulo $m_1,m_2$ gives the two congruences for $r_j$.

\item For the local sufficiency maps put $\alpha=a$, $h=n$,
$S_h=I+hE$ and $d_j=S_hw'_j-\alpha w_j$.
Suppose $p'=\alpha p$ and
$hr'_1\equiv\alpha s_1-s'_1\pmod k$,
$hr'_2\equiv\alpha s_2-s'_2\pmod{2l}$.
Reduction of $p'=\alpha p$ modulo $m_1,m_2$ gives
$r'_j\equiv\alpha r_j\pmod{m_j}$.
We may choose integral solutions of $N_{m_j}\ell_j=-d_j$ explicitly:
\[
 \begin{split}
 \ell_1&=\left(\frac{\alpha r_1-r'_1}{m_1},0,
 -\frac{s'_1+hr'_1-\alpha s_1}{k},0\right)^t,\\
 \ell_2&=\left(\frac{\alpha r_2-r'_2}{m_2},
 -\frac{s'_2+hr'_2-\alpha s_2}{2l},0,0\right)^t.
 \end{split}
\]
Put $y_j=d_j/m_j+\ell_j$; these vectors belong to the rational
kernel of $N_{m_j}$ and have first coordinate zero.
The following rational vectors solve $(I-A_j)b_j=y_j$:
\[
 \begin{split}
 b_1&=(0,(2y_{1u}+y_{1w})/3,(y_{1w}-y_{1u})/3,0)^t,\\
 b_2&=(0,(y_{2u}-y_{2w})/2,(y_{2u}+y_{2w})/2,0)^t.
 \end{split}
\]
The last coordinates of the equations follow from the displayed
norm-kernel equations.  The maps
\[
 (s,\zeta)\longmapsto(s,\alpha\zeta+\Pi_c(s)b_j)
\]
are holomorphic maps of the normalized torus families over the
degree-$m_j$ root discs.  Their inverses are holomorphic, and period
covariance and $(I-A_j)b_j=d_j/m_j+\ell_j$ give exact commutation
with the prescribed affine deck generators, modulo the integral
period $\ell_j$.  They therefore descend to biholomorphisms of the
finite fillings, including the full multiple fibres.

\item On the punctured root collar the corresponding middle translation is
\[
 \mathcal A_j=\Pi_cb_j+\frac{\log s_j}{2\pi i}\Pi_cd_j.
\]
Its integral defect is $\ell_j$.  The middle cocycle relation fixes
$\ell_0=-A_0(\ell_1+A_1\ell_2)$.  In the quotient
$\overline\Lambda=\Lambda/(\mathbb Z\widehat w+\mathbb Z\widehat\delta)$,
the first-coordinate condition is exactly $p'=\alpha p$.
Adding an integral multiple of $(0,1,-2,0)^t\in\ker N_{m_1}$ to
$\ell_1$ changes the other quotient coordinate by a generator,
without changing this first coordinate.  We can consequently arrange
\[
 [\ell_0]=-\alpha c(\sigma)+c(\sigma')\quad\hbox{in }\overline\Lambda.
\]
An integral toric vector then supplies $q\in\mathbb Z^2$ such that
the cusp reference $sq+\alpha\sigma-\sigma'$ has precisely defect
$\ell_0$.  Here $s=(2\pi i)^{-1}\log t_c$ is the cusp logarithm;
the maps of cusp toric covers are
\[
 (x,t_c)\longmapsto(t_c^q u(t_c)x^\alpha,t_c),
 \qquad u(0)\in(\mathbb C^*)^2.
\]
Integer translations of the cocharacter lattice, its inversion,
and multiplication by a holomorphic unit preserve the specified
hexagonal fan and all its strata.  These maps and their inverses
extend holomorphically across the three double curves and both
triple points.

\item Differences of allowed
finite extensions are the equivariant holomorphic vertical vectors
on the complete root discs.  The extra subgroup of order $k$ at
the first point acts trivially on the vertical vector coefficients
and rotates $s$; its invariant holomorphic coefficients are exactly
those in $s^k$.  The remaining order-three equivariance is the
original one.  The argument at the second point uses $s^l$ and
order four.  These invariant-coefficient maps identify the coarse difference sheaf
with $\mathcal V_B$ for multiplicities $(3,4)$.  Its first-coordinate quotient is
$\mathcal O_B(-p_0)$: the modular frame has vanishing orders $2k,l$
on the two enlarged roots and a simple cusp pole.  Its fixed
second-coordinate subline is $\mathcal O_B$.  Finite averaging
provides local lifts, hence
\[
 0\longrightarrow\mathcal O_B\longrightarrow\mathcal V_B
 \longrightarrow\mathcal O_B(-p_0)\longrightarrow0.
\]
Both end terms have vanishing $H^1$ on $\mathbb P^1$, so
$H^1(B,\mathcal V_B)=0$.  The permitted local translations form a
torsor under this sheaf: on each overlap their defects agree, so
their difference is a single-valued section of the indicated
difference sheaf.  Vanishing of its \v Cech class gives a global
middle translation.  Its overlap equalities are exactly the two
finite formulas and the cusp formula just displayed.  Consequently
the local biholomorphisms glue, with their inverses, over $B$.

\item For identical data, $p\ne0$ excludes $a=-1$. The congruences reduce to $k\mid nr_1$ and $2l\mid nr_2$. The two primitivity assumptions imply $k\mid n$ and $2l\mid n$. Since $\gcd(k,2l)=1$, their least positive common multiple is $2kl$. For different sphere translation data, $r'_1$ is invertible modulo $k$, $r'_2$ modulo $2l$, and the Chinese remainder theorem gives exactly one integer shift class modulo $2kl$. The values of $s_1$ modulo $k$ and $s_2$ modulo $2l$ all occur without affecting freeness or $p$. Thus there are exactly $2kl$ classes at fixed $u$. The change $c\mapsto c-1$ on a positive $u$ circuit acts by a generator of this cyclic residue group.

\end{enumerate}
\end{proof}





\section{Classification in elliptic-surface coordinates}\hypertarget{tag-02ZD}{}\label{push:section:E}
\subsection{Higher logarithmic transformations}\hypertarget{tag-02ZE}{}\label{push:section:E.1}
\begin{definition}[Higher multiplicities and source coordinates]
\hypertarget{tag-IP10}{}\label{I:def:higher-multiplicity-coordinates}
\StatementRecord{I:def:higher-multiplicity-coordinates}{proved}{IP10}{}
Put $m=3k$, $n=4l$ with $\gcd(m,n)=1$.
Viaclovsky's indices $(0,1)$ correspond to the finite indices $(1,2)$
of \cref{def:common-conventions}. Lattice coordinates are integral
homology coordinates, and matrices act on columns.
\end{definition}

\begin{proposition}[The integral dictionary]\hypertarget{tag-032R}{}\label{push:src:E.1.1}
\StatementRecord{push:src:E.1.1}{proved}{032R}{\cite[Section~2, equations~(2.1)--(2.2), p.~3, and (2.7), (2.11)--(2.12), p.~4]{ViaclovskyLog2026}}

Let $(\delta,a,b,d)$ be the marking in
\cite[Section~2, equations~(2.1)--(2.2), p.~3, and (2.7), (2.11)--(2.12), p.~4]{ViaclovskyLog2026},
and put $u_0=d+2a+2b$, $u_1=d+3a+3b$.
The homomorphism from the marking of \cref{def:common-conventions} to this marking is
\[
 P_V=\begin{pmatrix}
 -2&1&0&1\\-6&0&-1&0\\0&-1&0&0\\-1&0&0&0
 \end{pmatrix},\qquad \det P_V=-1.
\]
It intertwines the two finite monodromies and satisfies
\[
 P_V\widehat\delta=\delta,\qquad
 P_V\varepsilon=-u_0,\qquad
 P_V\varepsilon'=\delta-u_1.
\]
With the counterclockwise deck generators of that paper, its data
$m_0=3k$, $m_1=4l$ and
$a_i=[v_i/m_i]$ with $v_i=\rho_i\delta+N_i u_i$ correspond to
\[
 w_1=N_0\varepsilon-\rho_0\widehat\delta,
 \qquad
 w_2=N_1\varepsilon'-(\rho_1+N_1)\widehat\delta.
\]
For the section-preserving cusp comparison used there, $l_0=0$ and
$p=-(nN_0+mN_1)$.  Thus its sphere condition is
$b:=nN_0+mN_1=\pm1$.
The brackets in $a_i$ mean the torsion point in the smooth
root torus, using its section-preserving origin and the
transported period vector.  Replacing $v_i$ by $v_i+m_i\lambda$, $\lambda\in\Lambda$,
leaves the class $[v_i/m_i]$ unchanged, while its chosen lift fixes the integral collar defect.
\end{proposition}


\begin{proof}
\begin{enumerate}[label=\arabic*.,leftmargin=*]
\item Use $P_E=P_{\rm mark}$ from \cref{inv:B1.0}. In Engel's marking,
the four columns of the change to Viaclovsky's basis are
\[
 Q_E=\begin{pmatrix}-2&1&0&-2\\-1&0&0&0\\3&-1&1&0\\0&0&0&1\end{pmatrix}.
\]
Direct integral multiplication gives $P_V=Q_E^{-1}P_E$ and
$V_0P_V=P_VA_1$, $V_1P_V=P_VA_2$, where $V_0,V_1$ are the
matrices in that paper's Lemma 2.1.  Multiplication by the three
specified invariant vectors gives the displayed identities.
Inverting a clockwise affine generator replaces its invariant
translation $w_j/m_j$ by $-w_j/m_j$.  Consequently the
counterclockwise translation vector is $-P_Vw_j$, which gives the
two formulas, including the term $N_1$ in the second coefficient.  Viaclovsky's reference section extends through
the smooth locus of the cusp; its comparison therefore has zero
class in the quotient lattice of toric directions.  This is precisely
$l_0=0$.  Substitution into the definition of $p$ proves its formula.

\end{enumerate}
\end{proof}

\begin{definition}[The residual cusp defect]
\hypertarget{tag-IP11}{}\label{I:def:residual-cusp-defect}
\StatementRecord{I:def:residual-cusp-defect}{proved}{IP11}{}
For a general cusp collar $\sigma$, write $c(\sigma)$ for
its residual translation defect in
$\overline\Lambda=\Lambda/(\mathbb Z\widehat w+
\mathbb Z\widehat\delta)$, using the convention
$\sigma(g_0z)-R_0\sigma(z)=\Pi_c(g_0z)\ell(\sigma)$
and $c(\sigma)=[\ell(\sigma)]$.
Set $l_0=\gamma(c(\sigma))$.  In the marked coordinates
$w_j=r_j\varepsilon_j+s_j\widehat\delta$, the integer is
$p=mn l_0-nr_1-mr_2$.
The section-preserving source data of \ref{push:src:E.1.1} have $c(\sigma)=0$.
\end{definition}

\begin{theorem}[Classification in Viaclovsky's coordinates]\hypertarget{tag-032S}{}\label{push:src:E.1.2}
\StatementRecord{push:src:E.1.2}{proved}{032S}{\SeriesRef{push:src:E.1.1}}

Use a local logarithm $c$ of the transported linearization,
$u=e^{-2\pi i c}$, and the same reference section and cusp comparison
on both sides.  Two sphere members with data
$(m,\rho_0,N_0;n,\rho_1,N_1;c)$ and primed data are biholomorphic
if and only if their multiplicities agree and there are
$\alpha\in\{1,-1\}$ and $h=c'-c\in\mathbb Z$ such that
\[
 b'=\alpha b,\qquad
 hN'_0\equiv\rho'_0-\alpha\rho_0\pmod{k},\qquad
 hN'_1\equiv\rho'_1-\alpha\rho_1\pmod{2l}.
\]
Every such biholomorphism covers the identity of $B$.
For fixed discrete data the period identification subgroup is
$2kl\mathbb Z$.  At fixed $u$ all sphere translation data give
exactly $2kl$ classes.  A positive $u$-circuit acts on them by a
single cycle.
\end{theorem}


\begin{proof}
Apply \cref{fam:src:B.3.1} through the integral isomorphism
$P_V$ of \cref{push:src:E.1.1}. Its affine and cusp coordinates are
\[
 r_1=N_0,\quad s_1=-\rho_0,\qquad
 r_2=N_1,\quad s_2=-(\rho_1+N_1),\qquad p=-b.
\]
Thus $p'=\alpha p$ is $b'=\alpha b$, and the first norm congruence is
$hN'_0\equiv\rho'_0-\alpha\rho_0\pmod{k}$.
The second becomes
\[
 hN'_1\equiv\rho'_1-\alpha\rho_1+N'_1-\alpha N_1\pmod{2l}.
\]
The relation $b'=\alpha b$, reduced modulo $4l$, gives
$N'_1\equiv\alpha N_1\pmod{4l}$, since $\gcd(3k,4l)=1$.
It removes the last two terms and yields the stated congruence.
The base-preservation assertion, identification subgroup and number of
classes transfer through this same integral map. It preserves
$u=e^{-2\pi ic}$, so a positive $u$-circuit is still $c\mapsto c-1$
and induces the single cycle proved in \cref{fam:src:B.3.1}.
\end{proof}



\begin{corollary}[Holomorphic dependence of the comparisons]\hypertarget{tag-032T}{}\label{push:src:E.1.3}
\StatementRecord{push:src:E.1.3}{proved}{032T}{\SeriesRef{push:src:E.1.2}, \SeriesRef{fam:src:B.4.1}}

For fixed discrete data in \ref{push:src:E.1.2} and a fixed admissible shift $h$,
the constructed comparisons are holomorphic isomorphisms of families
over a sufficiently small parameter disc.  In particular, for sphere
data with multiplicities $(3,4)$ the identification with the standard
family uses the parameter
\[
 W=e^{-\pi i(c+s_2)}=(-1)^{s_2}w.
\]
The smooth monodromy conjugacy class around the positive parameter
circle is unchanged by this identification of twist data.
\end{corollary}


\begin{proof}
\begin{enumerate}[label=\arabic*.,leftmargin=*]
\item The rational vectors $b_j$, integral defects $\ell_j$ and toric
exponents in the proof of \cref{fam:src:B.3.1}, transported through
\cref{push:src:E.1.1}, are independent of $c$.  All period matrices,
root comparisons, and cusp unit factors are holomorphic in $c$.
On a small Stein parameter disc $D$, the correction sheaf is
$\mathcal V_B\boxtimes\mathcal O_D$.  Its first cohomology on
$B\times D$ is zero: project to $D$, use the same two-chart
complex and $H^1(B,\mathcal V_B)=0$, and then Stein acyclicity.
Thus the affine correction can be chosen holomorphically in $c$.
For $(k,l)=(1,1)$ the only nonvacuous finite congruence is
$h\equiv s_2-s_2'\pmod2$, because the admissible $r_2'$ is odd.
Taking $s_2'=0$ gives the displayed parameter identification.
Two choices of these family comparisons differ fibrewise by an
automorphism of the standard member.  The automorphism calculation in \cref{fam:src:B.4.1} identifies
that group with $\mathbb C^*$;
it is connected.  On overlapping parameter arcs the transitions
are therefore isotopic to the identity.  Composing around a circle
changes the resulting diffeomorphism only by such an isotopy and
by conjugation at the chosen base fibre.  Multiplication of the
parameter by $(-1)^{s_2}$ preserves the positive orientation.
\end{enumerate}
\end{proof}


\section{The monodromy group and its period domain}\hypertarget{tag-01PV}{}\label{fam:section:F}


\begin{definition}[The period domain and proper discontinuity]\hypertarget{tag-IP01}{}\label{I:def:period-domain}
\StatementRecord{I:def:period-domain}{proved}{IP01}{}
We use the clockwise root deck generators and lattice basis of
Definition~\ref{def:common-conventions}. A group acts properly
discontinuously when, for every
compact subset $K$ of its space, only finitely many group elements $g$ satisfy
$gK\cap K\ne\varnothing$.  This definition permits finite stabilizers.  We
write $e=\widehat\delta$, $E=e\otimes\gamma$, and
\[
 J_2=\begin{pmatrix}0&1\\-1&0\end{pmatrix},\qquad
 \Omega=\begin{pmatrix}
 0&0&0&6\\0&0&1&0\\0&-1&0&0\\-6&0&0&0
 \end{pmatrix}.
\]
Thus the alternating form is
$\psi=u\wedge w+6\gamma\wedge\delta$, an integral alternating form of
type $(1,6)$ on $\Lambda$.  We use $\Omega$ for its matrix and retain $J$ for
a real complex structure.  The period domain $\mathcal D$ consists of the real
endomorphisms $J$ of $\Lambda_{\mathbb R}$ satisfying
$J^2=-I$, $J^{\mathsf t}\Omega J=\Omega$, and
$\operatorname{sign}(\Omega J)=(2,2)$.  It is
$\operatorname{Sp}_4(\mathbb R)/U(1,1)$, with its complex structure as an
open subset of the Lagrangian Grassmannian.  An integral symplectic matrix
acts by $J\mapsto gJg^{-1}$.
\end{definition}

\subsection{An exact arithmetic description}\hypertarget{tag-01PW}{}\label{fam:section:F.1}

\begin{theorem}[The integral monodromy group]\hypertarget{tag-01R9}{}\label{fam:src:F.1.1}\label{fam:st:F.1.1}\label{cdp:st:D.1.14}
\StatementRecord{fam:src:F.1.1}{proved}{01R9}{\SeriesRef{fam:eq:F-law}, \SeriesRef{fam:eq:F-SLpresentation}}

Let $P_{e,\mathbb Z}$ be the integral symplectic group fixing $e$ as a
vector, and let $P_{\mathbb Z}$ be the integral symplectic group preserving
$\mathbb Ze\subset\ker\gamma\subset\Lambda$.  Every element of
$P_{e,\mathbb Z}$ has a unique expression
\begin{equation}\hypertarget{tag-01OT}{}\label{fam:eq:F-g}
 g(B,x,t)=
 \begin{pmatrix}
 1&0&0\\6x&B&0\\t&-x^{\mathsf t}J_2B&1
 \end{pmatrix},\qquad
 B\in\operatorname{SL}_2(\mathbb Z),\quad x\in\mathbb Z^2,\quad t\in\mathbb Z,
\end{equation}
where the block sizes are $1,2,1$.  Put
\[
 S=\begin{pmatrix}0&-1\\1&0\end{pmatrix},\quad
 T=\begin{pmatrix}1&1\\0&1\end{pmatrix},\quad
 a=\begin{pmatrix}0&1\\-1&-1\end{pmatrix}.
\]
There is a homomorphism
$\alpha:\operatorname{SL}_2(\mathbb Z)\to\mathbb Z/12$ determined by
$\alpha(S)=9$ and $\alpha(T)=1$; in particular $\alpha(a)=4$.
With $Q(x)=x_1+x_2+x_1x_2$, the monodromy group is
\begin{equation}\hypertarget{tag-01OQ}{}\label{fam:eq:F-Gamma}
 \Gamma=\left\{g(B,x,t):t+6Q(x)-\alpha(B)=0\ \text{in }\mathbb Z/12\right\}.
\end{equation}
Consequently $[P_{e,\mathbb Z}:\Gamma]=12$ and
$P_{\mathbb Z}=\{\pm I\}\times P_{e,\mathbb Z}$, so
$[P_{\mathbb Z}:\Gamma]=24$.  Its image in the Levi quotient is exactly
$\operatorname{SL}_2(\mathbb Z)$, with scalar on $e$ equal to $1$.
\end{theorem}



\begin{proof}
\begin{enumerate}[label=\arabic*.,leftmargin=*]
\item Use the matrices $A_1,A_2$ of \eqref{wb:pshift:eq:finite-data}.
Multiplication gives $A_1^3=A_2^4=I$ and
$A_j^{\mathsf t}\Omega A_j=\Omega$.  In particular
$A_1=g(a,(1,-1)^{\mathsf t},-2)$ and
$A_2=g(S,(0,-1)^{\mathsf t},3)$.  Their images $a,S$ generate
$\operatorname{SL}_2(\mathbb Z)$ because $Sa=T$ and $S,T$ generate that
group by \cite[Theorem~1.1, p.~1]{ConradSL2Z}.

\item A matrix fixing $e$ and preserving $\Omega$ preserves
$e^{\perp}=\ker\gamma$.  Its first row is $(1,0,0,0)$, its last column is
$e$, its middle block $B$ is symplectic on
$e^{\perp}/\mathbb Re$, and its remaining middle column, say $v$, satisfies
$6r=-v^{\mathsf t}J_2B$.  Integrality of $r$ and unimodularity of $B$
give $v=6x$ with $x\in\mathbb Z^2$. The lower left entry $t$
varies freely in $\mathbb Z$.  This proves \eqref{fam:eq:F-g}.  Direct multiplication
gives
\begin{equation}\hypertarget{tag-01OV}{}\label{fam:eq:F-law}
 g(B,x,t)g(C,y,s)
 =g(BC,x+By,t+s-6x^{\mathsf t}J_2By).
\end{equation}

\item Since $a=-ST$, its image in $\operatorname{PSL}_2(\mathbb Z)$
is the image of the generator $R=ST$ in
\cite[Theorem~C.1, pp.~18--19]{ConradSL2Z}. That theorem gives
$\operatorname{PSL}_2(\mathbb Z)=\langle\bar a,\bar S\mid
\bar a^3=\bar S^2=1\rangle$. Lifting through the central kernel
$\{I,-I\}$ gives
\begin{equation}\hypertarget{tag-01OS}{}\label{fam:eq:F-SLpresentation}
 \operatorname{SL}_2(\mathbb Z)=
 \langle a,S\mid a^3=S^4=1,\ [a,S^2]=1\rangle.
\end{equation}
Indeed the group on the right modulo its central element $S^2$ is the
cited free product, and $S^2$ maps to the nonidentity matrix $-I$.
Its kernel in the displayed matrix group is consequently trivial.
The values $4,9$ respect all these relations, proving the assertion about
$\alpha$; $T=Sa$ has value $1$.

\item The quadratic function $Q$ modulo $2$ is zero at $(0,0)$ and one at all
three other points of $\mathbb F_2^2$.  Hence $Q(Bx)=Q(x)\pmod2$ for every
$B\in\operatorname{SL}_2(\mathbb Z)$.  Also
\[
 Q(x+y)-Q(x)-Q(y)\equiv x^{\mathsf t}J_2y\pmod2.
\]
Equation \eqref{fam:eq:F-law} now shows that
$h(g(B,x,t))=t+6Q(x)\pmod{12}$ is a homomorphism.
The homomorphism $\chi=h-\alpha\circ B$ is surjective, since
$\chi(g(I,0,1))=1$.  Its values on $A_1,A_2$ are zero.

\item We prove $\Gamma=\ker\chi$.  Let $N_{\mathbb Z}$ be the kernel of the map to the middle
block.  The kernel of
$\langle g_1,g_2\mid g_1^3=g_2^4=1\rangle\to
\operatorname{SL}_2(\mathbb Z)$ is normally generated by
$[g_1,g_2^2]$, by \eqref{fam:eq:F-SLpresentation}.  Its image in $\Gamma$ is
therefore normally generated by
\[
 K=[A_1,A_2^2]=g(I,(0,-1)^{\mathsf t},6).
\]
Put
\[
 U=A_2KA_2^{-1}=g(I,(1,0)^{\mathsf t},-6),\quad
 V=K^{-1}=g(I,(0,1)^{\mathsf t},-6),\quad Z=I+12E.
\]
Equation \eqref{fam:eq:F-law} gives $[U,V]=Z^{-1}$.  Thus $U,V,Z$ generate
\begin{equation}\hypertarget{tag-01OR}{}\label{fam:eq:F-H}
 H=\{g(I,x,t):t\in6\mathbb Z,\quad t/6\equiv Q(x)\pmod2\}.
\end{equation}
More explicitly,
$U^rV^sZ^k=g(I,(r,s)^{\mathsf t},-6(r+s+rs)+12k)$, so these expressions
give every member of the set.  Conjugation by $g(B,x,t)$ sends
$g(I,y,s)$ to $g(I,By,s-12x^{\mathsf t}J_2By)$; the invariance of $Q$
modulo $2$ proves that $H$ is normal in $P_{e,\mathbb Z}$.  The normal
closure of $K$ is contained in $H$ and contains $U,V,Z$, so it is exactly
$H$.  We have proved $\Gamma\cap N_{\mathbb Z}=H$.  Both $\Gamma$ and
$\ker\chi$ surject onto the same Levi group and have intersection $H$
with $N_{\mathbb Z}$.  Given $g\in\ker\chi$, choose $\gamma\in\Gamma$
with the same middle block.  Then $g\gamma^{-1}\in H$, and $g\in\Gamma$.
This proves \eqref{fam:eq:F-Gamma} and the index $12$.

\item Finally an integral automorphism preserving the primitive line
$\mathbb Ze$ sends $e$ to $e$ or $-e$.  Multiplication by the central
matrix $-I$ converts the latter case to the former.  The intersection of
$\{\pm I\}$ with $P_{e,\mathbb Z}$ is trivial.  This proves the direct
product and the index $24$.
\end{enumerate}
\end{proof}

\begin{proposition}[The center and the Zariski closure]\hypertarget{tag-01RA}{}\label{fam:src:F.1.2}\label{fam:st:F.1.2}
\StatementRecord{fam:src:F.1.2}{proved}{01RA}{\SeriesRef{fam:eq:F-Gamma}, \SeriesRef{fam:eq:F-law}}

The group in \eqref{fam:eq:F-Gamma} has center
$\langle Z\rangle\cong\mathbb Z$ and is the following explicit central
extension of $\mathbb Z^2\rtimes\operatorname{SL}_2(\mathbb Z)$ by
$\mathbb Z$.  Choose the representative $a(B)\in\{0,\ldots,11\}$ of
$\alpha(B)$.  Its underlying set consists of triples $(B,x,k)$, with
product
\[
 (B,x,k)(C,y,l)=(BC,x+By,k+l+\kappa((B,x),(C,y))),
\]
where
\begin{align*}
12\kappa((B,x),(C,y))={}&a(B)+a(C)-a(BC)-6Q(x)-6Q(y)\\
 &-6x^{\mathsf t}J_2By+6Q(x+By).
\end{align*}
Its real and rational Zariski closures are respectively the full real
and rational symplectic stabilizers of the vector $e$.  They have
dimension $6$, with three-dimensional Heisenberg unipotent radical and
Levi quotient $\operatorname{SL}_2$.
\end{proposition}



\begin{proof}
\begin{enumerate}[label=\arabic*.,leftmargin=*]
\item The numerator defining $\kappa$ is divisible by $12$, by the two
congruences proved above.  The bijection
\[
 (B,x,k)\longmapsto g(B,x,a(B)-6Q(x)+12k)
\]
intertwines the displayed multiplication with matrix multiplication by
\eqref{fam:eq:F-law}.  This proves associativity and identifies the central extension with the
displayed cocycle.  The kernel
of the homomorphism to $\mathbb Z^2\rtimes\operatorname{SL}_2(\mathbb Z)$
is $\langle Z\rangle$.  That affine group has trivial center: commuting
with all translations forces its linear part to be $I$, and commuting
with $S,T$ forces its translation part to vanish.  Thus the center of
$\Gamma$ is exactly the stated kernel.

\item The logarithms $U-I,V-I,Z-I$ are linearly independent elements of the
three-dimensional Heisenberg Lie algebra.  The square of each is zero.
A polynomial vanishing on all integral powers $I+n(U-I)$ vanishes on
the whole one-parameter line; the same holds for $V$ and $Z$.  Hence the
Zariski closure contains the full unipotent radical.  Its Levi image
contains $S,T$; the powers of $T$ and of $STS^{-1}$ similarly supply the
two opposite unipotent subgroups of $\operatorname{SL}_2$.  Their Lie
algebras and their bracket span $\mathfrak{sl}_2$.  They generate the
connected algebraic group $\operatorname{SL}_2$.  A subgroup containing
the entire unipotent kernel and surjecting onto this Levi quotient is
the entire vector stabilizer.  Its dimension is $3+3=6$.  This reasoning
uses equations over $\mathbb Q$, so proves the rational assertion as
well as the real assertion.
\end{enumerate}
\end{proof}

\subsection{Properness and the maximal open set}\hypertarget{tag-01PX}{}\label{fam:section:F.2}


\begin{definition}[The sign domains]\hypertarget{tag-IP02}{}\label{I:def:period-sign-domains}
\StatementRecord{I:def:period-sign-domains}{proved}{IP02}{}
Put
\[
 q_e:\mathcal D\longrightarrow\mathbb R,\qquad q_e(J)=\psi(e,Je),
 \qquad\mathcal O=\{J:q_e(J)\ne0\}.
\]
The function $q_e$ is invariant under $P_{\mathbb Z}$: replacing $e$ by
$-e$ changes both of its arguments.  Write $\mathcal O_+$ and
$\mathcal O_-$ for its two sign regions.
\end{definition}

\begin{proposition}[An infinite stabilizer]\hypertarget{tag-01RB}{}\label{fam:src:F.2.1}\label{fam:st:F.2.1}
\StatementRecord{fam:src:F.2.1}{proved}{01RB}{\SeriesRef{fam:eq:F-Gamma}}

The following point of $\mathcal D$ has an infinite $\Gamma$-stabilizer:
\[
 J_0=\begin{pmatrix}
 0&0&-1&0\\0&0&0&-6\\1&0&0&0\\0&1/6&0&0
 \end{pmatrix},\qquad
 R=\begin{pmatrix}
 1&0&0&0\\0&1&72&0\\0&0&1&0\\-12&0&0&1
 \end{pmatrix}\in\Gamma.
\]
Every power of $R$ fixes $J_0$, and all these powers are distinct.
\end{proposition}



\begin{proof}
\begin{enumerate}[label=\arabic*.,leftmargin=*]
\item Direct multiplication gives $J_0^2=-I$ and
$J_0^{\mathsf t}\Omega J_0=\Omega$.  Moreover
\[
 \Omega J_0=\begin{pmatrix}
 0&1&0&0\\1&0&0&0\\0&0&0&6\\0&0&6&0
 \end{pmatrix}
\]
has eigenvalues $1,-1,6,-6$, so $J_0\in\mathcal D$ and $q_e(J_0)=0$.
The matrix $R$ equals $g(T^{72},0,-12)$ and belongs to \eqref{fam:eq:F-Gamma},
because $\alpha(T^{72})=0\pmod{12}$.  Its commutator with $J_0$ is zero.
Writing $R=I+N$ gives $N\ne0$ and $N^2=0$, so $R^m=I+mN$ and the
powers are distinct.  The identity $N^2=0$ also gives
\((I+mN)(I+nN)=I+(m+n)N\). The compact singleton $\{J_0\}$ meets every one of its translates by
$R^m$, so the action on $\mathcal D$ is not properly discontinuous.
\end{enumerate}
\end{proof}

\begin{theorem}[The maximal domain of proper discontinuity]\hypertarget{tag-01RC}{}\label{fam:src:F.2.2}\label{fam:st:F.2.2}
\StatementRecord{fam:src:F.2.2}{proved}{01RC}{\SeriesRef{fam:eq:F-g}}

The set $\mathcal O$ is the largest $\Gamma$-invariant open subset of
$\mathcal D$ on which the action is properly discontinuous.  It is also
the largest such open subset for $P_{\mathbb Z}$.  Each sign region is
a complex manifold in the analytic topology, and its quotient by either
discrete group is a Hausdorff complex analytic orbifold.
\end{theorem}



\begin{proof}
\begin{enumerate}[label=\arabic*.,leftmargin=*]
\item Let $P_{\mathbb R}$ be the real symplectic group preserving the real line
$\mathbb Re$.  For a fixed
$J$, the form
$h_J(v,w)=\psi(v,Jw)+i\psi(v,w)$ is a nondegenerate Hermitian form of
signature $(1,1)$ on the complex vector space $(\Lambda_{\mathbb R},J)$,
with the convention conjugate linear in the first entry.  When
$q_e(J)\ne0$, the complex line $\mathbb Ce$ is nondegenerate and its
Hermitian orthogonal complement is a complex line of the opposite sign.
Normalize $e$ by a positive real scalar and choose a unit vector in this
complement.  Comparing these two Hermitian bases for two complex
structures with the same sign gives a real symplectic map preserving
$\mathbb Re$ and conjugating one complex structure to the other.  Thus
$P_{\mathbb R}$ is transitive on each sign region.

\item An element of the stabilizer of $J$ inside $P_{\mathbb R}$ is complex
linear and sends $e$ to $r e$ for a real nonzero $r$.  Preservation of
$q_e(J)$ gives $r^2=1$.  On the complementary complex line it acts by an
element of $U(1)$.  The stabilizer is consequently the compact group
$\{\pm1\}\times U(1)$.  Each sign region is therefore a homogeneous
space of $P_{\mathbb R}$ with compact stabilizer $L$.  Local sections of
$P_{\mathbb R}\to P_{\mathbb R}/L$, and a finite cover of a compact
subset $K$, give a compact subset $C$ of $P_{\mathbb R}$ mapping onto
$K$.  Every element carrying a point of $K$ back into $K$ belongs to
$CLC^{-1}$, which is compact.  This proves properness.  Restricting a
proper action to a closed discrete subgroup gives
the required finite-intersection condition.  Both $\Gamma$ and
$P_{\mathbb Z}$ are closed discrete matrix groups.  The two sign regions
are disjoint invariant open-and-closed subsets of $\mathcal O$, so the
same finite-intersection condition holds on their union.  Proper actions
by discrete groups have finite stabilizers and Hausdorff quotients; their
holomorphic local finite-group quotient charts supply the asserted
analytic orbifolds.

\item We prove maximality by exhibiting a dense set of points with infinite
stabilizer in the zero locus of $q_e$.
The real vector stabilizer $P_{e,\mathbb R}$ acts transitively on
$\mathcal D_0=\{J:q_e(J)=0\}$.  In the
Hermitian space just described, $e$ is a nonzero null vector.  Choose
$v$ with $h_J(e,v)=1$, and replace $v$ by
$v-\tfrac12h_J(v,v)e$.  This preserves $h_J(e,v)=1$ and makes $v$ null.
The ordered complex basis $(e,v)$ has Gram matrix
$\bigl(\begin{smallmatrix}0&1\\1&0\end{smallmatrix}\bigr)$.
The complex-linear map between two such bases fixes $e$ and preserves
$h$, hence its imaginary part $\psi$.  As a real map it belongs to
$P_{e,\mathbb R}$ and conjugates the two complex structures.

\item The subgroup $P_{e,\mathbb Q}$ is dense in $P_{e,\mathbb R}$.
For the unipotent coordinates in \eqref{fam:eq:F-g} this is density of
$\mathbb Q^3$ in $\mathbb R^3$.  For $\operatorname{SL}_2$ one may
approximate $a,b,c$ by rational numbers in the chart $a\ne0$ and set
$d=(1+bc)/a$; translates of that chart cover the group.  Hence
\[
 \{rJ_0r^{-1}:r\in P_{e,\mathbb Q}\}
\]
is dense in $\mathcal D_0$.
Put $R_0=g(T^6,0,-1)$, so $R_0^{12}=R$ and
$(R_0-I)^2=0$.  At $rJ_0r^{-1}$ the nonzero rational nilpotent
$L=r(R_0-I)r^{-1}$ commutes with the complex structure and satisfies
$L^2=0$.  Choose a positive integer $d$ clearing its denominators.
The matrix $I+dL$ and its inverse $I-dL$ are integral, symplectic, and
fix $e$; this follows either by conjugating the identity
$(I+s(R_0-I))^{\mathsf t}\Omega(I+s(R_0-I))=\Omega$ or by expanding it.
The twelfth power of $I+dL$ lies in $\Gamma$, because
$P_{e,\mathbb Z}/\Gamma\cong\mathbb Z/12$.  It has infinite order and
fixes $rJ_0r^{-1}$.  Thus points with infinite $\Gamma$-stabilizer are
dense in $\mathcal D_0$.

\item Every open neighborhood in $\mathcal D$ of any point of $\mathcal D_0$
therefore contains a point with infinite stabilizer.  An open set on
which the action is proper cannot meet $\mathcal D_0$.  This proves the
largest-open-set assertion for $\Gamma$.  The same points have infinite
$P_{\mathbb Z}$-stabilizer, so it proves that assertion for the full
parabolic as well.
\end{enumerate}
\end{proof}

\begin{proposition}[A global chart for positive periods]\hypertarget{tag-01RD}{}\label{fam:src:F.2.3}\label{fam:st:F.2.3}
\StatementRecord{fam:src:F.2.3}{proved}{01RD}{the definitions and constructions stated in the proof}

The positive sign region admits the global period chart
\[
 \mathcal O_+=\left\{
 (\tau,\mu,\beta)\in\mathbb C^3:
 \operatorname{Im}\tau>0,\quad
 \operatorname{Im}\beta-\frac{6(\operatorname{Im}\mu)^2}
 {\operatorname{Im}\tau}<0\right\},
\]
represented by the period matrix \eqref{wb:pshift:eq:periods}, with
$\beta_c$ replaced by the coordinate $\beta$. Put
$D=\Im\beta-6(\Im\mu)^2/\Im\tau$. In this chart,
\begin{equation}\hypertarget{tag-01OW}{}\label{fam:eq:F-qD}
 q_e(J)=-\frac6D.
\end{equation}
Every admissible period $\Pi_c(z)$ belongs to $\mathcal O_+$, for every
admissible $c$.  Thus this single maximal proper-action open set contains
all the period images in question.
\end{proposition}



\begin{proof}
\begin{enumerate}[label=\arabic*.,leftmargin=*]
\item Let $J\in\mathcal O_+$.  The two vectors $\widehat w,e$ are complex-linearly
independent in $(\Lambda_{\mathbb R},J)$.  Otherwise
$\widehat w=r e+sJe$ for real $r,s$.  Pairing with $e$ gives
$0=sq_e(J)$; hence $s=0$, contradicting the real linear independence
of $\widehat w,e$.  There is consequently a unique complex-linear
isomorphism $(\Lambda_{\mathbb R},J)\to\mathbb C^2$ sending these two
vectors to $(1,0)$ and $(0,1)$.  The corresponding period matrix has
right block $I$.  The equation
$\Pi\Omega^{-1}\Pi^{\mathsf t}=0$ says that its upper left entry is
six times its lower middle entry, giving exactly the displayed matrix.
This construction is the ordinary holomorphic Grassmannian chart,
so its coordinates and inverse are holomorphic.

\item In the real symplectic basis
$(\widehat u,\widehat\gamma/6,\widehat w,e)$ the normalized symmetric
period matrix is
$\bigl(\begin{smallmatrix}\tau&\mu\\\mu&\beta/6\end{smallmatrix}\bigr)$.
Its imaginary part has signature $(1,1)$, equivalently
$6(\operatorname{Im}\mu)^2-
\operatorname{Im}\tau\operatorname{Im}\beta>0$.
Solving $\Pi(Je)=i\Pi e=(0,i)$ in real coefficients gives
\[
 \gamma(Je)=
 \frac{\operatorname{Im}\tau}
 {\operatorname{Im}\tau\operatorname{Im}\beta-
 6(\operatorname{Im}\mu)^2}.
\]
Since $\psi(e,v)=-6\gamma(v)$, positivity of $q_e$ now gives
$\operatorname{Im}\tau>0$ and $D<0$, and proves \eqref{fam:eq:F-qD}.
Conversely these inequalities make the imaginary period matrix
invertible with the stated signature and give $q_e>0$.  The admissibility
inequality of Definition~\ref{I:def:family-parameter} is $D_c<0$
everywhere, so the displayed comparison map places every admissible
period in $\mathcal O_+$.
\end{enumerate}
\end{proof}

\subsection{The curve in the proper quotient}\hypertarget{tag-01PY}{}\label{fam:section:F.3}

\begin{theorem}[The embedded period curve]\hypertarget{tag-01RE}{}\label{fam:src:F.3.1}\label{fam:st:F.3.1}
\StatementRecord{fam:src:F.3.1}{proved}{01RE}{\SeriesRef{fam:eq:F-j}, \SeriesRef{fam:src:F.2.2}}

For every admissible $c$, the marked period map on the covering of
$B^\circ$ is a holomorphic immersion into $\mathcal O_+$.  The induced
map
\[
 p_c:B^\circ\longrightarrow\Gamma\backslash\mathcal O_+
\]
is injective and a holomorphic immersion.  It extends over $t=0,1$ to a
proper holomorphic embedding
\[
 \bar p_c:\mathbb C=B\setminus\{\infty\}
 \lhook\joinrel\longrightarrow\Gamma\backslash\mathcal O_+.
\]
The closure of $p_c(B^\circ)$ in this quotient is exactly
$\bar p_c(\mathbb C)$.
\end{theorem}



\begin{proof}
\begin{enumerate}[label=\arabic*.,leftmargin=*]
\item For $B=\bigl(\begin{smallmatrix}a&b\\c&d\end{smallmatrix}\bigr)$, changing
the marking by $g(B,x,t)$ replaces $\Pi$ by $\Pi g(B,x,t)^{-1}$ and then
renormalizes its right block.  On the quotient by the complex line of
$e$, the two periods are the row $(\tau,1)$.  Right multiplication by
$B^{-1}$ sends it to $(d\tau-c,a-b\tau)$.  Hence
\begin{equation}\hypertarget{tag-01OX}{}\label{fam:eq:F-tauaction}
 \tau\longmapsto\frac{d\tau-c}{a-b\tau}.
\end{equation}
The denominator never vanishes in the upper half plane.  The function
$j(\tau)$ is therefore a holomorphic invariant of this action and
descends to a holomorphic function $j_\Gamma$ on the proper quotient.
Here $j$ is the elliptic invariant, normalized by
$j(E)=1728\,4A^3/(4A^3+27B^2)$ for $y^2=x^3+Ax+B$; changing a lattice
basis preserves the isomorphism of elliptic curves and hence this
function.

\item The elliptic curve
whose period is $\tau$ has the Engel equation of
Lemma~\ref{I:weierstrass-comparison},
\[
 y^2=x^3-3t^3(t-1)x+2t^4(t-1)^2.
\]
Its discriminant is $1728t^8(t-1)^3$ and direct cancellation gives
\begin{equation}\hypertarget{tag-01OU}{}\label{fam:eq:F-j}
 j(\tau)=1728t\qquad(t\in B^\circ).
\end{equation}
In a local base coordinate $t$, the
derivative of the composite $j\circ\tau$ equals $1728$.  Thus the
derivative of the marked period map is nonzero.  The same identity after
passage to the quotient shows both injectivity and injectivity of the
tangent map of $p_c$: $j_\Gamma/1728$ is a holomorphic left inverse.

\item At the two elliptic points, the period functions are holomorphic in the
cyclic covering coordinates $s_1,s_2$, with $t=s_1^3$ and
$t-1=s_2^4$ in the root coordinates of Definition~\ref{def:common-conventions}.  Their
values are nondegenerate because $D_c<0$ there.  The maps from these
discs to $\mathcal O_+$ are equivariant for $A_1,A_2$.  Composing with
the quotient makes them invariant, so they descend to holomorphic maps
from the base discs.  This is the finite-quotient comparison map: a
holomorphic invariant germ on a disc is a convergent series in $s_j^{m_j}$,
which proves descent on local function rings.  The extensions agree
with $p_c$ off their centers and therefore with each other.  Identity
\eqref{fam:eq:F-j} extends by holomorphy, so the resulting $\bar p_c$ still
has holomorphic left inverse $j_\Gamma/1728$.

\item The quotient is Hausdorff by \ref{fam:src:F.2.2}.  For a compact subset $K$ of it,
$\bar p_c^{-1}(K)$ is closed in $\mathbb C$ and lies in the compact set
$j_\Gamma(K)/1728$.  Thus $\bar p_c$ is proper.  A holomorphic section of
a holomorphic map is locally a closed immersion: in local embeddings
into polydiscs, its coordinate along the base is $t$ and the remaining
coordinates are its holomorphic component functions, whose graph
defines the image.  Applying this to the left inverse shows that
$\bar p_c$ is a holomorphic embedding; properness makes its image
closed.  Since $\mathbb C\setminus\{0,1\}$ is dense in $\mathbb C$, its
image has the asserted closure.
\end{enumerate}
\end{proof}

\begin{proposition}[Finite limits and the cusp limit]\hypertarget{tag-01RF}{}\label{fam:src:F.3.2}\label{fam:st:F.3.2}
\StatementRecord{fam:src:F.3.2}{proved}{01RF}{\SeriesRef{fam:src:F.3.1}}

At the two finite punctures the period map has the holomorphic
finite-cover limits
\[
 \Pi_1=
 \begin{pmatrix}2(2-\rho)&\rho&1&0\\
 \beta_c(z_1)&(2-\rho)/3&0&1\end{pmatrix},\qquad
 \Pi_2=
 \begin{pmatrix}3(1-i)&i&1&0\\
 \beta_c(z_2)&(1-i)/2&0&1\end{pmatrix},
 \quad \rho=e^{\pi i/3}.
\]
Their monodromies are the given order-three and order-four matrices.
At the cusp the curve leaves every compact subset of
$\Gamma\backslash\mathcal O_+$.  On the prescribed cusp branch its
Hodge plane in the compact dual tends to
$\mathbb C\gamma+\mathbb Cu$.  The latter plane has complex dimension
two in common with its conjugate and therefore does not represent a
point of $\mathcal D$.
\end{proposition}



\begin{proof}
\begin{enumerate}[label=\arabic*.,leftmargin=*]
\item The finite values of $\tau,\mu$ follow by evaluating their transformation
laws at the fixed points: $\tau(z_1)=\rho$,
$\mu(z_1)=(2-\rho)/3$, $\tau(z_2)=i$,
$\mu(z_2)=(1-i)/2$.  Substitution in $\Pi$ gives the two matrices.
The finite covering comparison and its descent were proved in \ref{fam:src:F.3.1}.
The unipotent cusp transformation is
$\tau\mapsto\tau-1$, $\mu\mapsto\mu$, $\beta\mapsto\beta+1$.
On the cusp branch, $\operatorname{Im}\tau\to+\infty$,
$\mu\to\mu_0$, and $\beta_c+\tau\to b_0+c$, from the holomorphic
cusp expansions.  Consequently
\[
 D_c=-\operatorname{Im}\tau+O(1),\qquad
 q_e(\Pi_c)\longrightarrow0^+,
 \qquad j_\Gamma(p_c(t))=1728t\longrightarrow\infty.
\]
The last equality alone excludes a convergent subnet in any compact
subset of the proper quotient; the first two identities give the
position relative to the maximal proper-action open set.
The rows of $\Pi_c$, regarded as covectors, are
$6\mu\gamma+\tau u+w$ and $\beta_c\gamma+\mu u+\delta$.
Dividing each by $\tau$ gives limits $u$ and $-\gamma$.
Their exterior product has a nonzero limit, so their span converges
in the Grassmannian to the stated plane.  It is spanned by real
covectors, hence equals its conjugate.  A complex structure has
$H^{1,0}\cap H^{0,1}=0$; here the intersection is the entire
two-dimensional plane.  This explicit comparison proves the assertion
about the compact-dual limit.
\end{enumerate}
\end{proof}

\begin{proposition}[The determinant of the period derivative]\hypertarget{tag-01RG}{}\label{fam:src:F.3.3}\label{fam:st:F.3.3}
\StatementRecord{fam:src:F.3.3}{proved}{01RG}{\SeriesRef{fam:src:F.3.1}, \SeriesRef{fam:eq:F-j}}

Let $Z=\bigl(\begin{smallmatrix}6\mu&\tau\\\beta&\mu\end{smallmatrix}\bigr)$
be the left block of \eqref{wb:pshift:eq:periods}. Its derivative has
generically nonzero determinant, with a zero at $z_1$:
\[
 \det\dot Z=6\dot\mu^2-\dot\tau\dot\beta,\qquad
 \det\dot Z(z_1)=0.
\]
The period map is an immersion at every point of $B^\circ$ by
Theorem~\ref{fam:src:F.3.1}; its differential and $\det\dot Z$ are
related by the polynomial map $\dot Z\mapsto\det\dot Z$.
\end{proposition}



\begin{proof}
\begin{enumerate}[label=\arabic*.,leftmargin=*]
\item In the cusp coordinate $s$, the expansion
$Z=s\bigl(\begin{smallmatrix}0&1\\-1&0\end{smallmatrix}\bigr)+C(q)$,
$q=e^{2\pi is}$, gives
$\det(dZ/ds)=1+O(q)$, so the determinant is not the zero function.
At $z_1$ the fixed-point transformation law of $\mu$ gives
$\dot\mu=-\dot\tau/3$.  One may obtain this by differentiating
$\mu\circ g_1=(1-\mu)/\tau$, using
$g_1'(z_1)=-\rho$ and $\rho^2-\rho+1=0$.
The law
$\beta\circ g_1=\beta+2-6(1-\mu)^2/\tau$ then gives
$\dot\beta=2\dot\tau/3$.  Substitution yields
$6\dot\tau^2/9-2\dot\tau^2/3=0$.
The derivative of a map from a curve to the three-dimensional period
chart is a vector of three coordinate derivatives.  Identity
\eqref{fam:eq:F-j} proves that this vector is nonzero over $B^\circ$;
the determinant tests the rank of the associated two-by-two
Kodaira--Spencer matrix.  Thus the coordinate derivative is nonzero on $B^\circ$, while its
associated matrix has the displayed determinant.
\end{enumerate}
\end{proof}



\section{Exceptional fibres, Neron--Severi classes, and compact subspaces}\hypertarget{tag-01PP}{}\label{fam:section:D}
\subsection{Elliptic quotients and the torsion locus}\hypertarget{tag-01PQ}{}\label{fam:section:D.1}
\begin{proposition}[The distinguished elliptic projection]\hypertarget{tag-01QR}{}\label{fam:src:D.1.1}\label{fam:st:D.1.1}
\StatementRecord{fam:src:D.1.1}{proved}{01QR}{the definitions and constructions stated in the proof}

The specified linear projection \(\operatorname{pr}_1:\mathbb C^2\to
\mathbb C\) induces a holomorphic group homomorphism
\(F_t\to E_t=\mathbb C/(\mathbb Z\tau+\mathbb Z)\) exactly on
\(6\mu=A\tau+B\), \(A,B\in\mathbb Z\). It is then surjective with
connected kernel
\(\mathbb C/(\mathbb Z+(A\mu-\beta)\mathbb Z)\).
On the smooth base the two values of \(t\) are precisely
\[
 t_\pm=\frac{64(47\pm27\sqrt3)}{11}.
\]
\end{proposition}


\begin{proof}
\begin{enumerate}[label=\arabic*.,leftmargin=*]
\item The images of the four lattice generators under the linear projection
are \(6\mu,\tau,1,0\). This proves descent in both directions.
The integral basis change replacing \(\widehat\gamma\) with
\(\widehat\gamma-A\widehat u-B\widehat w\) changes its period to
\((0,\beta-A\mu)\); the remaining vertical generator is \((0,1)\).
The other two periods map onto the entire target lattice, so the kernel
is connected. Its second modulus has imaginary part
\(A\operatorname{Im}\mu-\operatorname{Im}\beta=-D>0\).

\item The geometric model and its marked section are transported by
$\Phi^{-1}(X,Y)=(X/2,-Y/(2\sqrt2))$ to the arithmetic model
below. The inverse map is
$\Phi(x,y)=(2x,-2\sqrt2y)$; substitution in both cubic
equations proves the isomorphism and its effect on the marked point
(\cref{I:weierstrass-comparison}). Thus the arithmetic equation is
\[
 Y^2=X^3-\frac34t^3(t-1)X+\frac14t^4(t-1)^2,
 \qquad P=(0,t^2(t-1)/2).
\]
Its sixth division polynomial at \(P\) is
\[
 \psi_6(P)=\frac9{65536}t^{24}(t-1)^9
 (11t^2-6016t+8192).
\]
For a smooth elliptic fibre, \(\psi_n(P)=0\) is exactly the condition
\(nP=O\), since the multiplication rational map has these and only these
poles in its affine coordinate. The two roots displayed above are simple
and avoid \(0,1\). Chord-and-tangent addition gives the rational function
\[
 \begin{aligned}
 x(6P)&=\frac{(89t^2+4736t-4096)Q(t)}
 {1296(11t^2-6016t+8192)^2},\\
 Q(t)&=881t^4+4745344t^3-12079104t^2\\
 &\hspace{2em}+24641536t-16777216.
 \end{aligned}
\]
The division-polynomial recurrence is given in
\cref{ari:C.4.4}.
\end{enumerate}
\end{proof}

\begin{proposition}[The torsion locus and its degrees]\hypertarget{tag-01QS}{}\label{fam:src:D.1.2}\label{fam:st:D.1.2}
\StatementRecord{fam:src:D.1.2}{proved}{01QS}{\SeriesRef{fam:src:C.2.1}, \SeriesRef{fam:src:D.1.1}; \cite[Lemma 2.1]{Engel2026}; \cite[Theorem 4.4, p.~3]{Schnell19}}

The locus on \(B^\circ\) where \(P(t)\) is torsion is countable and dense.
Every fibre over this locus is an abelian surface. The divisor on the
smooth base cut out by \(6kP=O\) has degree \(3k^2-1\), including local
intersection multiplicity, for every positive integer \(k\).
\end{proposition}

\begin{proof}
\begin{enumerate}[label=\arabic*.,leftmargin=*]
\item \hypertarget{tag-000F}{}\label{app:rederive:fam-src-D.1.2}
On a simply connected open subset write uniquely
\(\mu=a_0\tau+b_0\), \(a_0,b_0\) real. The section is torsion precisely
when \((a_0,b_0)\in\mathbb Q^2\). In a local complex parameter \(z\), direct
differentiation gives
\[
 \det\frac{\partial(a_0,b_0)}{\partial(\operatorname{Re}z,
 \operatorname{Im}z)}
 =-\frac{|\mu'(z)-a_0\tau'(z)|^2}{\operatorname{Im}\tau}.
\]
This determinant vanishes on no open set. On an open set of vanishing,
the two real differential equations imply \(da_0=db_0=0\), hence
\(\mu=a_0\tau+b_0\) with constant real coefficients. Analytic continuation
would give this identity on the cover. The cusp law
\((\tau,\mu)\circ g_0=(\tau-1,\mu)\) then forces \(a_0=0\), and the
\(g_2\)-law makes a constant \(\mu\) impossible. Thus every open subset
contains a point where this real map is a local diffeomorphism, and its
image contains rational pairs. This proves density. For each integer
triple \((n,A,B)\), \(n\ne0\), the function \(n\mu-A\tau-B\) is
holomorphic and not identically zero by the same argument. Its zeros are
discrete; a countable union proves countability.

\item For \(N\mu=A\tau+B\), \(A,B,N\in\mathbb Z\), \(N>0\), use the linear
coordinate change \(\zeta'_2=\zeta_2-(A/N)\zeta_1\).
The four integral vectors
\[
 N\widehat u-B\widehat\delta,\quad
 N\widehat w+A\widehat\delta,\quad
 N\widehat\gamma-6A\widehat u-6B\widehat w,\quad
 \widehat\delta
\]
span a sublattice of index \(N^3\). Its first two periods lie on
\(\zeta'_2=0\), and its last two lie on \(\zeta_1=0\).
Both intersections are genuine rank-two lattices because \(D<0\).
Their product quotient is therefore a product of elliptic curves, and
inclusion of this sublattice induces a finite surjective holomorphic
homomorphism onto \(F_t\). A positive integral form on the product can
be summed over its finite translation kernel and multiplied by the
index to descend an integral positive form to \(F_t\). The Riemann
polarization criterion proves projectivity.

\item For the degree statement, \cite[Lemma~2.1, pp.~3--4, with the height computation on p.~4]{Engel2026} gives \(\langle P,P\rangle=1/6\),
\(\chi(\mathcal O_S)=1\), the component groups at the two reducible fibres
have orders three and two, and the local corrections for \(P\) are
\(4/3\) and \(1/2\). The application to \(6kP\) is as follows.
The corrected divisor class of a section is its class minus \(O\) and
a fibre multiple, plus the unique rational combination of nonidentity
fibre components that makes it perpendicular to those components.
On restriction to the generic fibre this assignment is additive because
the divisor class \([Q]-[O]\) is the group law in \(\operatorname{Pic}^0\).
The perpendicular correction makes its self-intersection the negative of
the height, so the height is quadratic. The element \(6kP\) is in the
identity component at both reducible fibres; its correction terms vanish.
Hence
\[
 6k^2=(6k)^2\langle P,P\rangle
 =2+2((6kP)\mathbin{\cdot}O),
\]
which gives the formula. The rational functions for \(6P\) show
that it has finite affine coordinates at \(0,1\), so it is a nonzero point
of each additive identity component; multiplication by \(k\) remains
nonzero in \(\mathbb C\). At infinity its multiplicative coordinate is
\((-a)^{6k}=a^{6k}\ne1\). Thus all intersections counted by the height
occur on \(B^\circ\), so their total degree is \(3k^2-1\).
\end{enumerate}
\end{proof}

\Uses{the hypotheses, notation and definitions stated here; \cref{app:rederive:fam-src-D.1.2}}




\begin{theorem}[A second elliptic quotient]\hypertarget{tag-01QT}{}\label{fam:src:D.1.3}\label{fam:st:D.1.3}
\StatementRecord{fam:src:D.1.3}{proved}{01QT}{\SeriesRef{fam:src:D.1.1}, \SeriesRef{fam:src:C.2.1}}

At \(t=32/5\) there is an admissible parameter for which
\(L(\zeta_1,\zeta_2)=3\zeta_1+4\zeta_2\) induces a surjective
holomorphic group homomorphism \(F_t\to E_t\) with connected elliptic
kernel, while \(P(t)\) has exact order four. The distinguished
projection of \ref{fam:src:D.1.1} descends exactly at the order-dividing-six
locus specified there.
\end{theorem}


\begin{proof}
\begin{enumerate}[label=\arabic*.,leftmargin=*]
\item Set \(t=32/5\). The exact fourth division polynomial is
\[
 \psi_4(P)=-\frac1{32}t^{10}(t-1)^4(5t-32),
\]
whereas \(\psi_2(P)=t^2(t-1)\ne0\). Thus \(P\) has exact order four.
On the rectangular branch write \(\tau=iT\), \(T>1\), with
\(j(iT)=1728\cdot32/5\). The chosen zero has real part \(1/2\), negative
imaginary part of absolute value less than \(T/2\), and order four;
therefore \(\mu=(2-\tau)/4\). Choose the concrete parameter
\[
 c=\left(\frac98-\frac{101}{4}\right)\tau-\frac94
       -\beta^{\rm part}(t).
\]
It is globally admissible. Indeed the positive-real Jacobi product at this
\((\mu,\tau)\) gives
\(\operatorname{Im}\beta^{\rm part}\geq-T-(3/\pi)\log3\).
Consequently
\(\operatorname{Im}c\leq-185T/8+(3/\pi)\log3<-22<-M\), by \ref{fam:src:C.2.1}.
Thus the displayed parameter is admissible.

\item The holomorphic linear map \(L:\mathbb C^2\to\mathbb C\),
\(L(\zeta)=3\zeta_1+4\zeta_2\), takes the four periods to
\[
 -101\tau,\qquad2\tau+2,\qquad3,\qquad4.
\]
These generate exactly \(\mathbb Z\tau+\mathbb Z\): the last two generate
\(1\), and the first two then generate \(\tau\), since
\(\gcd(101,2)=1\). Thus \(L\) descends to a surjective holomorphic group
homomorphism \(F_{32/5}\to E_{32/5}\) with connected one-dimensional
kernel. The kernel is compact as a closed subset of a compact torus,
hence an elliptic curve. Nevertheless \(6P=2P\ne O\).
\end{enumerate}
\end{proof}

\begin{definition}[The Hermitian form of an integral class]\hypertarget{tag-IP03}{}\label{I:def:Hermitian-form}
\StatementRecord{I:def:Hermitian-form}{proved}{IP03}{}

At a smooth fibre the real vector-space isomorphism
$\Pi:\mathbb R^4\to\mathbb C^2$ transfers an alternating form $E$ to
$\widetilde E(v,w)=E(\Pi^{-1}v,\Pi^{-1}w)$. For a $J$-invariant
$\widetilde E$, put
\[
 H_E(v,w)=\widetilde E(v,Jw)-i\widetilde E(v,w).
\]
Here $H_E$ is complex-linear in the first argument. Direct substitution
gives $H_E(Jv,w)=iH_E(v,w)$ and
$H_E(w,v)=\overline{H_E(v,w)}$. Its real Chern form is
$-\operatorname{Im}H_E=\widetilde E$, since
$(i/2)(H_E(v,w)-H_E(w,v))=\widetilde E(v,w)$.
\end{definition}

\begin{proposition}[The N\'eron--Severi lattice]\hypertarget{tag-01QU}{}\label{fam:src:D.1.4}\label{fam:st:D.1.4}
\StatementRecord{fam:src:D.1.4}{proved}{01QU}{\SeriesRef{fam:src:C.3.1}, \SeriesRef{fam:src:D.1.6}}

For every smooth fibre and every admissible \(c\), its Neron--Severi lattice
is the following explicit subgroup of \(\mathbb Z^6\):
\begin{equation}\hypertarget{tag-01P4}{}\label{fam:eq:cd-NS}
 \left\{E:\ E_{12}-E_{13}\tau+(6E_{23}-E_{14})\mu+
 E_{24}\beta+E_{34}(6\mu^2-\tau\beta)=0\right\}.
\end{equation}
The complex orientation has intersection form
\(E^2=-2\operatorname{Pf}(E)\), where
\(\operatorname{Pf}(E)=E_{12}E_{34}-E_{13}E_{24}+E_{14}E_{23}\).
In particular the invariant class \(\eta=6\gamma\wedge\delta+u\wedge w\)
has \(\int_{F_t}\eta^2=-12\). Its Hermitian form is indefinite.
The fibres with lattice exactly \(\mathbb Z\eta\) have no compact curves
and algebraic dimension zero.
\end{proposition}


\begin{proof}
\begin{enumerate}[label=\arabic*.,leftmargin=*]
\item On the universal cover the holomorphic two-form is the wedge of the two
rows of \(\Pi\). Its six coefficients in the integral coframe are
\(6\mu^2-\tau\beta,-\beta,6\mu,-\mu,\tau,1\).
An integral real two-form is of type \((1,1)\) exactly when its wedge
with this form vanishes: its \((0,2)\) part is then zero, and reality
also annihilates its \((2,0)\) part. Expanding the wedge gives
\eqref{fam:eq:cd-NS}. The exponential sequence identifies these integral
\((1,1)\) classes with Chern classes of holomorphic line bundles, proving
the lattice assertion.

\item The ordered real period matrix in \cref{I:period-determinant} has determinant \(TD<0\). Thus the
integral orientation \((\widehat\gamma,\widehat u,\widehat w,
\widehat\delta)\) is opposite to the complex orientation. Algebraically
\(E\wedge E=2\operatorname{Pf}(E)\gamma\wedge u\wedge w\wedge\delta\),
which proves the sign in the displayed intersection formula. In
particular \cite[Lemma~9.4, pp.~66--67]{AC2026} gives the coefficient
identity \(\eta\wedge\eta=12\gamma\wedge u\wedge w\wedge\delta\).
The real period isomorphism above takes this integral orientation to the
negative of the complex orientation, yielding
\(\int_{F_t}\eta^2=-12\).

\item For the Hermitian form $H_E$ of Definition~\ref{I:def:Hermitian-form}, direct
substitution of the four periods gives $H_\eta=-2G^{-1}$, where
\[
 G=\begin{pmatrix}2T&2U\\2U&\operatorname{Im}\beta/3\end{pmatrix},
 \qquad \det G=2TD/3<0.
\]
Thus $H_\eta$ has one positive and one negative eigenvalue, and every
nonzero multiple remains indefinite. The integration
current of an effective divisor, averaged under translations of the
compact torus, is a nonzero semipositive constant \((1,1)\)-form in its
Chern class. It cannot be an indefinite multiple of \(\eta\).
Thus no curve exists when \(\operatorname{NS}=\mathbb Z\eta\).
A nonconstant meromorphic function would have a nonempty polar divisor
(after subtracting a constant if necessary), so the algebraic dimension
is zero.
\end{enumerate}
\end{proof}

\begin{proposition}[Generic and projective fibres]\hypertarget{tag-01QV}{}\label{fam:src:D.1.5}\label{fam:st:D.1.5}
\StatementRecord{fam:src:D.1.5}{proved}{01QV}{\SeriesRef{fam:eq:cd-NS}}

For a fixed \(c\), all but countably many smooth fibres have
\(\operatorname{NS}=\mathbb Z\eta\). Conversely, for every fixed smooth
\(t\), there are admissible parameters tending to the cusp of the
parameter disc for which \(F_t\) is an abelian surface, irrespective of
the order of \(P(t)\).
\end{proposition}


\begin{proof}
\begin{enumerate}[label=\arabic*.,leftmargin=*]
\item By \cite[Lemma~9.3, p.~66]{AC2026}, the functions
$1,\tau,\mu,\beta^{\rm part},6\mu^2-\tau\beta^{\rm part}$ are linearly
independent. For $\beta=\beta^{\rm part}+c$, the relation
$A+B\tau+C\mu+D\beta+E(6\mu^2-\tau\beta)=0$ has coefficient vector
\[
 (A+Dc,\ B-Ec,\ C,\ D,\ E)
\]
in that ordered list. This coefficient change is invertible, with inverse
$(A',B',C',D',E')\mapsto(A'-D'c,B'+E'c,C',D',E')$.
Thus every integral class outside $\mathbb Z\eta$ has a
nonidentically-zero holomorphic left side in \eqref{fam:eq:cd-NS}; its
zeros are discrete. Taking the countable union over $\mathbb Z^6$ proves
the assertion for each fixed $c$, as in the application
\cite[Proposition~9.5, p.~67]{AC2026}.

\item For the second assertion put
\(\beta_n=-3n^2\tau+12n\mu\),
\(c_n=\beta_n-\beta^{\rm part}(t)\). Its imaginary part tends to
minus infinity, so the already proved finite bound on \(M\) makes every
sufficiently large \(n\) admissible. The integral alternating matrix
with nonzero upper entries
\[
 E_{13}=-3n^2,\quad E_{14}=6n,\quad E_{23}=-n,\quad E_{24}=1
\]
satisfies \eqref{fam:eq:cd-NS} exactly at \(\beta_n\), and
\(\operatorname{Pf}(E)=-3n^2<0\). Its complex-oriented square is positive.
A nondegenerate Hermitian form in complex dimension two with positive
determinant is definite. Choosing its sign gives a positive integral
Riemann form, so \(F_t\) is projective. This supplies an explicit family
of projective fibres even at any chosen non-torsion value of the section.
\end{enumerate}
\end{proof}

\begin{theorem}[Elliptic quotients and algebraic dimension]\hypertarget{tag-01QW}{}\label{fam:src:D.1.6}\label{fam:st:D.1.6}
\StatementRecord{fam:src:D.1.6}{proved}{01QW}{\SeriesRef{fam:eq:cd-NS}, \SeriesRef{fam:src:D.1.4}; \cite[Theorem 4.1, p.~1, and Theorem 4.4, p.~3]{Schnell19}}

Fix any admissible $c$. For each integral alternating matrix $E$, let
$Z_E(c)$ be the zero set on the marking cover of the explicit holomorphic
function in \eqref{fam:eq:cd-NS}, with $\beta=\beta^{\rm part}+c$ given by
\eqref{fam:eq:cd-beta}. Let $\pi$ be the cover map to $B^\circ$, and define
\[
 \mathcal E_c=\pi\!\left(\bigcup_{E\ne0,\ \operatorname{Pf}(E)=0}Z_E(c)\right),
 \qquad
 \mathcal A_c=\pi\!\left(\bigcup_{\operatorname{Pf}(E)<0}Z_E(c)\right).
\]
These give the exact loci of all elliptic quotients and all abelian
surfaces, respectively. The algebraic dimension of every smooth fibre is
\[
 a(F_t)=
 \begin{cases}
 2,&t\in\mathcal A_c,\\
 1,&t\in\mathcal E_c\setminus\mathcal A_c,\\
 0,&t\notin\mathcal E_c\cup\mathcal A_c.
 \end{cases}
\]
Together with \eqref{fam:eq:cd-NS}, these classify every smooth fibre
by countable unions of explicit analytic zero sets. The saturated null
lattices in the proof determine all connected elliptic kernels and their
quotient maps.
\end{theorem}

\begin{proof}
\begin{enumerate}[label=\arabic*.,leftmargin=*]
\item \hypertarget{tag-000G}{}\label{app:rederive:fam-src-D.1.6}
We apply the Appell--Humbert correspondence and Riemann projectivity
criterion \cite[Theorem~4.1, p.~1, and Theorem~4.4, p.~3]{Schnell19}.
Equation \eqref{fam:eq:cd-NS} gives $J$-invariance, so
Definition~\ref{I:def:Hermitian-form} identifies the integral class $E$
with the real Chern form $-\operatorname{Im}H_E=\widetilde E$.
The Appell--Humbert integrality condition is therefore exactly the
integrality of $E$. The complex-oriented square computed in \ref{fam:src:D.1.4} shows that
$\operatorname{Pf}(E)<0$ is equivalent to definiteness of this Hermitian
form. Its sign can be chosen positive. The cited projectivity criterion
then gives precisely the asserted abelian locus, in both directions.

\item For nonzero $E$ with Pfaffian zero, its real rank is two. Its nullspace
$K_{\mathbb R}\subset\mathbb R^4$ is a rational plane because it is the
kernel of an integral matrix. $J$-invariance makes $\Pi K_{\mathbb R}$ a
complex line, and $K=K_{\mathbb R}\cap\mathbb Z^4$ is saturated. The maps
\[
 (\Pi K_{\mathbb R})/\Pi K\ \hookrightarrow F_t
 \ \longrightarrow\
 \mathbb C^2/(\Pi K_{\mathbb R}+\Pi\mathbb Z^4)
\]
are holomorphic group morphisms; the first is a closed embedded elliptic
curve and the second is its quotient onto an elliptic curve, with
connected kernel. A basis of the saturated $K$ extends to an integral basis of $\mathbb Z^4$.
The other two periods project to real-linearly independent vectors in the
one-dimensional complex quotient, because $\Pi$ is a real isomorphism.
They therefore form a discrete lattice there. Compactness follows by
taking the image of its compact real fundamental parallelogram. Conversely any elliptic quotient
pulls an integral positive area form back to a nonzero rank-two integral
$(1,1)$ form. This proves the characterization of $\mathcal E_c$.
Equivalently all its kernels can be found by saturated rank-two lattices
with a basis $v,w$ satisfying
$\det_{\mathbb C}(\Pi v,\Pi w)=0$ and real independence. This is a second
exact integer test, with explicit inclusion and quotient morphisms.

\item For the algebraic-dimension assertion, an effective divisor on a torus
has a semipositive Hermitian form: averaging its integration current under
translations gives that constant form. A nonzero semipositive form has
rank one or two. Rank two gives the abelian case. Rank one gives the
elliptic quotient just constructed. In the absence of a definite form,
there cannot be two different such kernel subtori. Indeed addition from
their product onto $F_t$ would be a surjective holomorphic homomorphism
with finite kernel, and the product polarization would descend after
multiplication by an integer, giving a definite form. Thus there is a
single kernel $K_t$ in this case. Every effective divisor is a sum of
its translates: its intersection number with $K_t$ is zero, whereas a
curve mapping nonconstantly to the elliptic quotient has strictly positive
intersection with a general quotient fibre. Every meromorphic function
therefore factors through this quotient. One can check the last statement
on a fibre avoiding its polar divisor, where it is a holomorphic function
on a compact elliptic curve and hence constant; analytic continuation
then gives the factorization. The quotient supplies nonconstant
meromorphic functions, so the algebraic dimension is exactly one.
If neither locus occurs, a nonconstant meromorphic function would have
a polar divisor of one of the two forbidden ranks. This proves dimension
zero. Finally monodromy acts by integral determinant-one changes of basis,
preserving \eqref{fam:eq:cd-NS} and its Pfaffian. Thus the stated unions descend
as the stated subsets of $B^\circ$.
\end{enumerate}
\end{proof}

\Uses{the hypotheses, notation and definitions stated here; \cref{fam:eq:cd-NS}; \cref{fam:eq:cd-beta}; \cref{app:rederive:fam-src-D.1.6}}




\begin{proposition}[An explicit locus of algebraic dimension one]\hypertarget{tag-01QX}{}\label{fam:src:D.1.7}\label{fam:st:D.1.7}
\StatementRecord{fam:src:D.1.7}{proved}{01QX}{\SeriesRef{fam:eq:cd-NS}, \SeriesRef{fam:src:D.1.6}, \SeriesRef{fam:src:C.2.1}, \SeriesRef{fam:src:D.1.2}}

There is an explicit nonempty locus of admissible members with ordinary
fibres of algebraic dimension one. On a sufficiently small punctured
cusp neighbourhood set
\[
 c(z)=-6\tau(z)-12\mu(z)-\beta^{\rm part}(z).
\]
Outside a countable set of $z$, the fibre of $X_{c(z)}$ over $\pi(z)$ has
Neron--Severi lattice $\mathbb Z\eta+\mathbb ZE$, where the nonzero upper
entries of $E$ are
$E_{13}=-6$, $E_{14}=-6$, $E_{23}=1$, $E_{24}=1$.
Its algebraic dimension is one. The section $P(\pi(z))$ can simultaneously
be required to be non-torsion.
\end{proposition}


\begin{proof}
\begin{enumerate}[label=\arabic*.,leftmargin=*]
\item At the specified parameter $\beta=-6\tau-12\mu$. Both $E$ and $\eta$
satisfy \eqref{fam:eq:cd-NS}, and exact multiplication gives
$\operatorname{Pf}(a\eta+bE)=6a^2$.
For any further integral class $F$, substitution in its period equation
gives a linear relation among $1,\tau,\mu,(\mu+\tau)^2$.
These four functions are linearly independent: applying $g_0$ to a
relation $A+B\tau+C\mu+D(\mu+\tau)^2=0$ and subtracting yields
$-B+D(1-2(\mu+\tau))=0$. A further cusp difference forces $D=0$,
and then $B=0$; nonconstancy of $\mu$ forces $A=C=0$.
Comparison of coefficients now gives
$F_{34}=F_{12}=0$, $F_{13}=-6F_{24}$, and
$6F_{23}-F_{14}=12F_{24}$, exactly
$F\in\mathbb Z\eta+\mathbb ZE$. Every other integral class therefore
has a discrete zero set along this period locus. Excluding their
countable union gives the asserted exact Neron--Severi lattice.
It contains the nonzero isotropic $E$ and no matrix of negative
Pfaffian; \ref{fam:src:D.1.6} gives algebraic dimension one.
Finally
$c=-5\tau-12\mu-(\beta^{\rm part}+\tau)$ has imaginary part tending
to minus infinity at the cusp. A sufficiently small cusp neighbourhood
is therefore admissible by \ref{fam:src:C.2.1}. The torsion locus of $P$ is countable
by \ref{fam:src:D.1.2}, so its exclusion still leaves such points in every open subset
of that neighbourhood.
\end{enumerate}
\end{proof}

\subsection{Compact curves and surfaces}\hypertarget{tag-01PR}{}\label{fam:section:D.2}
\begin{proposition}[Vertical compact subspaces]\hypertarget{tag-01QY}{}\label{fam:src:D.2.1}\label{fam:st:D.2.1}
\StatementRecord{fam:src:D.2.1}{proved}{01QY}{the definitions and constructions stated in the proof}

For a standard-twist member \(X_c\) with underlying manifold \(S^6\),
every compact irreducible complex curve lies in one fibre of \(f\), and
every compact irreducible complex surface is one of the reduced fibres.
\end{proposition}


\begin{proof}
\begin{enumerate}[label=\arabic*.,leftmargin=*]
\item The given diffeomorphism with \(S^6\) implies \(H^2(X_c;\mathbb Z)=0\).
For a compact irreducible curve \(C\), let
\(\nu:\widetilde C\to C\hookrightarrow X_c\) be its normalization map.
The holomorphic map \(f\nu\) has
\(\deg (f\nu)^*\mathcal O_{\mathbb P^1}(1)=0\), since its Chern class is
pulled back from the zero class on \(X_c\). A nonconstant holomorphic map
of compact Riemann surfaces has positive degree: a regular value has a
positive number of preimages, all with positive local multiplicity.
Therefore \(f\nu\), and hence \(f|_C\), is constant.

\item An irreducible surface \(Z\) in the smooth threefold is an effective
Cartier divisor. If its image were all of \(B\), a general smooth fibre
would meet it in a nonempty effective divisor \(D\) whose line bundle
is \(\mathcal O_{X_c}(Z)|_{F_t}\). Its integral Chern class is zero.
Choose a flat K\"ahler form \(\omega\) on the torus \(F_t\). The equality
\(\int_{F_t}c_1(D)\wedge\omega=\int_D\omega\) has left side zero and right
side positive, a contradiction. Thus \(Z\) lies in a single fibre.
Each reduced fibre is irreducible of dimension two, so containment forces
equality.
\end{enumerate}
\end{proof}

\begin{proposition}[Curves in the special fibres]\hypertarget{tag-01QZ}{}\label{fam:src:D.2.2}\label{fam:st:D.2.2}
\StatementRecord{fam:src:D.2.2}{proved}{01QZ}{\SeriesRef{fam:src:D.2.1}, \SeriesRef{fam:src:D.1.1}, \SeriesRef{fam:src:C.4.1}}

There is no rational curve in \(X_c\setminus W_c\). There are elliptic
curves outside all three special fibres, in particular in each of the
two smooth fibres \(F_{t_\pm}\). All curves in \(W_c\) are precisely the
images of irreducible algebraic curves in \(\mathrm{dP}_6\) under its
finite normalization map. Distinct curves have the same image only for
the three paired boundary lines. All curves in a bielliptic fibre
\(S_j\) are images of irreducible curves in its finite product cover.
\end{proposition}


\begin{proof}
\begin{enumerate}[label=\arabic*.,leftmargin=*]
\item By \ref{fam:src:D.2.1} every curve is vertical. A holomorphic map from
\(\mathbb P^1\) into a torus lifts to its vector-space universal cover,
since \(\mathbb P^1\) is simply connected, and its coordinate functions
are constant by compactness. Thus a smooth torus fibre contains no
rational curve. A map from \(\mathbb P^1\) into a bielliptic surface
lifts through its finite unramified torus cover by the same covering
criterion; the lifted map is constant. These arguments prove the first
assertion, including possibly singular rational images.
\ref{fam:src:D.1.1} gives an embedded connected elliptic kernel in each \(F_{t_\pm}\),
proving the second.

\item For the third assertion let \(\eta:\mathrm{dP}_6\to W_c\) be the finite
normalization morphism. Each component of the inverse image of a compact
curve is compact and one-dimensional and maps onto that curve, because
finite maps have zero-dimensional fibres. It is algebraic in the
projective surface \(\mathrm{dP}_6\). Conversely its image is a compact
analytic curve by properness. Away from the conductor, \(\eta\) is a
biholomorphism, so two distinct curves can have the same image only when
both lie in the conductor; its identifications are exactly the three
opposite pairs computed in \ref{fam:src:C.4.1}. This is also an explicit algebraic
parametrization: the curves of \(\mathrm{dP}_6\) are the three exceptional
lines and strict transforms of irreducible plane curves under the blow-up
of the three coordinate points. Their homogeneous equations and their
three multiplicities specify the strict transforms. The finite-product
cover statement for \(S_j\) follows from the identical proper-image and
finite-preimage argument applied to its unramified morphism
\(E_j\times F_j\to S_j\); curves in the product are algebraic.
\end{enumerate}
\end{proof}

\begin{theorem}[All compact curves]\hypertarget{tag-01R0}{}\label{fam:src:D.2.3}\label{fam:st:D.2.3}
\StatementRecord{fam:src:D.2.3}{proved}{01R0}{\SeriesRef{fam:src:D.1.6}, \SeriesRef{fam:src:D.2.1}, \SeriesRef{fam:src:D.2.2}; \cite[Theorem 4.1 and equation (4.3), pp.~1--2]{Schnell19}}

The following descriptions give all compact irreducible curves of $X_c$.
In $W_c$ and $S_j$ they are exactly the images described in \ref{fam:src:D.2.2}.
In a smooth fibre of algebraic dimension zero there are none. In a
smooth fibre of algebraic dimension one they are exactly the translates
of its unique elliptic kernel $K_t$. In an abelian fibre they are exactly
the translates of elliptic subtori, together with the irreducible zero
divisors of the positive Appell--Humbert theta sections described below.
These descriptions, \ref{fam:src:D.1.6}, and \ref{fam:src:D.2.1} constitute a classification including
every smooth exceptional fibre and every compact irreducible surface.
\end{theorem}

\begin{proof}
\begin{enumerate}[label=\arabic*.,leftmargin=*]
\item \hypertarget{tag-000H}{}\label{app:rederive:fam-src-D.2.3}
All curves are vertical by \ref{fam:src:D.2.1}, so the list of fibre types is exhaustive.
The dimension-zero and dimension-one assertions, including uniqueness,
were proved in \ref{fam:src:D.1.6}. Let now $F_t=\mathbb C^2/\Pi\Lambda$ be abelian.
For every positive definite Hermitian form $H$ whose imaginary part takes
integer values on $\Pi\Lambda$, and every function
$\alpha:\Pi\Lambda\to U(1)$ satisfying
\[
 \alpha(\lambda+\nu)=\alpha(\lambda)\alpha(\nu)
                 e^{\pi i\operatorname{Im}H(\lambda,\nu)},
\]
take the nonzero entire solutions $\theta:\mathbb C^2\to\mathbb C$ of
\begin{equation}\hypertarget{tag-01P7}{}\label{fam:eq:cd-theta-sections}
 \theta(\zeta+\lambda)=\alpha(\lambda)
 e^{\pi H(\zeta,\lambda)+\pi H(\lambda,\lambda)/2}\theta(\zeta).
\end{equation}
The cocycle identity follows by multiplying the displayed factors; its
integral imaginary cross-term has exponential one. It defines a
holomorphic line bundle by quotienting $\mathbb C^2\times\mathbb C$.
The invariant zero set of $\theta$ descends to a compact divisor on
$F_t$. Retain precisely those zero divisors that are reduced and
irreducible. This is an exact parametrization, with all periods and
multiplier equations given; multiplying $\theta$ by a nonzero scalar
does not change its divisor.

\item The correspondence \cite[Theorem 4.1, p.~1, and equation (4.3), p.~2]{Schnell19} proves
that every holomorphic line bundle and every section on this quotient
has this form. For the cohomology matrix the exact comparison is
$H_E=\widetilde E(\,\cdot\,,J\,\cdot\,)-i\widetilde E$ from Definition~\ref{I:def:Hermitian-form};
we choose the sign of $E$ making $H_E$ positive. Its imaginary part is
$-\widetilde E$, which gives the integral semicharacter equation displayed
above. The effective line-bundle criterion \cite[pp.~2--3]{Schnell19} shows that
a divisor class with positive Hermitian rank two gives exactly
these positive theta sections. A divisor of Hermitian rank one is pulled
back from its elliptic quotient and, when irreducible and reduced, is a
single fibre. Thus it is a translate of an elliptic subtorus. The null
lattice is computed by the explicit integer kernel in \ref{fam:src:D.1.6}. No other
rank is possible for an effective nonempty divisor, by positivity of its
averaged current. This proves exhaustiveness on the abelian fibre.

\item The same equations give the curves in the product covers of $S_j$:
compute them on its two-elliptic-curve lattice and take images under the
specified finite morphism. Conversely each curve downstairs is the image
of a component of its finite inverse image. In $W_c$ the analogous procedure uses the homogeneous polynomial
parametrization of the normalization in \ref{fam:src:D.2.2}, with the three boundary
pairings already identified. The maps used throughout are finite or
closed embeddings, so they preserve compactness and give one-dimensional
images.
\end{enumerate}
\end{proof}

\Uses{the hypotheses, notation and definitions stated here; \cref{fam:src:D.2.2}; \cref{fam:src:D.1.6}; \cref{fam:src:D.2.1}; \cref{app:rederive:fam-src-D.2.3}}






\subsection{The rank-three period quotient}\hypertarget{tag-01PN}{}\label{fam:section:C.3}
\begin{proposition}[Rank-three period quotients]\hypertarget{tag-01QM}{}\label{fam:src:C.3.1}\label{fam:st:C.3.1}
\StatementRecord{fam:src:C.3.1}{proved}{01QM}{the definitions and constructions stated in the proof}

At every point with \(D=0\), \(T>0\), the real-linear map
\(\Pi:\mathbb R^4\to\mathbb C^2\) has rank three. Write
\(\mu=a_0\tau+b_0\), \(a_0,b_0\in\mathbb R\), and put
\[
 k=(1,-6a_0,-6b_0,6a_0\operatorname{Re}\mu-\operatorname{Re}\beta).
\]
Its real kernel is \(\mathbb R k\). The subgroup \(\Pi\mathbb Z^4\) is
discrete exactly when \(k\in\mathbb Q^4\). In that case it is a
rank-three lattice in a real hyperplane and its quotient in \(\mathbb C^2\)
is a noncompact complex Lie group. In every other case the topological
quotient \(\mathbb C^2/\Pi\mathbb Z^4\) is non-Hausdorff.
\end{proposition}


\begin{proof}
\begin{enumerate}[label=\arabic*.,leftmargin=*]
\item The real matrix, with output order
\((\operatorname{Re}\zeta_1,\operatorname{Im}\zeta_1,
\operatorname{Re}\zeta_2,\operatorname{Im}\zeta_2)\), is
\[
 \begin{pmatrix}6\operatorname{Re}\mu&\operatorname{Re}\tau&1&0\\
 6\operatorname{Im}\mu&T&0&0\\
 \operatorname{Re}\beta&\operatorname{Re}\mu&0&1\\
 \operatorname{Im}\beta&\operatorname{Im}\mu&0&0\end{pmatrix}.
\]
Its last three columns are real-linearly independent because \(T>0\).
Its determinant is \(TD\), so at zero determinant the rank is exactly
three. Substitution gives the displayed kernel, including its signs.
If the image subgroup is discrete, it is a lattice of rank at most three
in its real span. The homomorphism \(\mathbb Z^4\to\Pi\mathbb Z^4\)
therefore has a nonzero kernel, forcing the displayed kernel vector with
first coordinate one to be rational. Conversely a rational line has a
primitive integral generator extendable to a basis of \(\mathbb Z^4\).
The remaining three basis images are a real basis of the image hyperplane
and hence generate a discrete lattice. Its quotient leaves one unbounded
real direction. Finally a non-discrete finitely generated subgroup of
this real hyperplane cannot be closed: a closed subgroup is a product of
a real vector group and a discrete group, and a nonzero real vector group
is uncountable. A quotient topological group is Hausdorff exactly when its
identity subgroup is closed. This proves the last assertion.
\end{enumerate}
\end{proof}

\begin{proposition}[The quotient at the parameter boundary]\hypertarget{tag-01QN}{}\label{fam:src:C.3.2}\label{fam:st:C.3.2}
\StatementRecord{fam:src:C.3.2}{proved}{01QN}{\SeriesRef{fam:src:C.1.1}, \SeriesRef{fam:src:C.2.3}}

The unique degenerate base point at the admissibility boundary is
\(t=1\). The boundary constants
\(c=x+i(3/\pi)\log6\) give
\[
 \ker_{\mathbb R}\Pi=\mathbb R(1,3,-3,-x-3/2).
\]
For rational $x+3/2=m/n$ in lowest terms, the quotient is a
finite holomorphic quotient of $E_i\times\mathbb C^*$, with a covering
of degree four. For irrational \(x+3/2\), the closure of the period subgroup is generated
by the real line \(\mathbb R(0,1)\) and the two periods
\((i,(1-i)/2)\), \((1,0)\). Its Hausdorff quotient is a real manifold
homeomorphic to \(\mathbb R\times(S^1)^2\).
\end{proposition}


\begin{proof}
\begin{enumerate}[label=\arabic*.,leftmargin=*]
\item \ref{fam:src:C.1.1} gives \(\operatorname{Im}\beta=3/2\) at this value of \(c\).
The vector \(\widehat\gamma+3\widehat u-3\widehat w\) maps to
\((0,x+3/2)\), while \(\widehat\delta\) maps to \((0,1)\).
For $x+3/2=m/n$ in lowest terms, change coordinates by
$(\zeta_1,\zeta_2)\mapsto(\zeta_1,\zeta_2+\zeta_1/2)$.
The vectors $2\widehat u-\widehat\delta$ and
$2\widehat w-\widehat\delta$ then have periods $(2i,0)$ and $(2,0)$.
Bezout's identity applied to $\widehat\gamma+3\widehat u-3\widehat w$
and $\widehat\delta$ gives the period $(0,1/n)$.
These three periods generate a sublattice of index four in the rank-three
image: its two remaining generators are represented by
$(i,1/2)$ and $(1,1/2)$, each of order two modulo this sublattice.
The product identification is the explicit map
$([\zeta_1],[\zeta_2])\mapsto([\zeta_1/2],e^{2\pi in\zeta_2})$
from $\mathbb C^2$ modulo these three periods to $E_i\times\mathbb C^*$. Lattice inclusion gives the claimed degree-four map.
For irrational $x+3/2$, its additive subgroup with $1$ is dense in the
indicated real line. The other two displayed periods are independent
modulo that line. This proves the closure and the quotient description.
By \ref{fam:src:C.2.3} this is the only degenerate torus cover at the admissibility
boundary; every other finite-base period matrix retains real rank four.
\end{enumerate}
\end{proof}


\begin{remark}[The local boundary obstruction]\hypertarget{tag-01QO}{}\label{fam:src:C.3.3}\label{fam:st:C.3.3}
\StatementRecord{fam:src:C.3.3}{proved}{01QO}{\SeriesRef{fam:src:C.3.2}, \SeriesRef{fam:src:C.3.4}, \SeriesRef{push:src:D.1.2}, \SeriesRef{push:src:D.1.3}}

The only degenerate base fibre at the parameter boundary lies over $t=1$.
By \ref{fam:src:C.3.2}, its rank-three period quotient is noncompact for
rational $x+3/2$ and non-Hausdorff for irrational $x+3/2$; neither quotient
can be a fibre of a proper family of complex spaces.
Proposition~\ref{fam:src:C.3.4} proves that the elliptic Albanese model
also has no holomorphic Weierstrass continuation across this boundary.
The obstruction to a general proper holomorphic extension is proved in
\cref{push:src:D.1.2,push:src:D.1.3} by recovering the intrinsic divisor.
\end{remark}

\begin{proposition}[The natural boundary of the Albanese invariant]\hypertarget{tag-01QP}{}\label{fam:src:C.3.4}\label{fam:st:C.3.4}
\StatementRecord{fam:src:C.3.4}{proved}{01QP}{\SeriesRef{fam:section:A.4}, \SeriesRef{fam:src:C.1.1}}

On its half-plane of definition, the $j$-invariant of
$\operatorname{Alb}(S_2)$ has no meromorphic continuation across any
point of the line
$\operatorname{Im}c=(3/\pi)\log6$. In particular the rank-three quotient
at $t=1$ cannot be replaced by a Weierstrass family of elliptic Albanese
curves whose coefficients are holomorphic in the same parameter across
that line.
\end{proposition}


\begin{proof}
\begin{enumerate}[label=\arabic*.,leftmargin=*]
\item The explicit Albanese lattice in \cref{fam:src:A.4.1} is
$\mathbb Z+(\kappa_2/2)\mathbb Z$, with
$\kappa_2=c+3/2-i(3/\pi)\log6$ by \ref{fam:src:C.1.1}. A positive-imaginary modulus is
$z=-\kappa_2/2$, so the invariant is $j(z)$. At a reduced rational
$r=a/b$ on the real boundary, choose an integral determinant-one modular
matrix with bottom row $(b,-a)$. Its value at $r+iy$ has imaginary part
$1/(b^2y)$. The convergent Fourier expansion $j=q^{-1}+744+O(q)$ therefore
gives growth of order $\exp(2\pi/(b^2y))$ as $y\downarrow0$.
A meromorphic function at $r$ is bounded by a fixed power of $1/y$,
so continuation there is impossible. A meromorphic continuation near an
irrational real point would also continue the same function near some
rational real point in that neighbourhood, which has just been excluded.
The affine change from $c$ to $z$ transfers this obstruction to the
stated line. A Weierstrass family with holomorphic coefficients has a
meromorphic $j$-invariant, expressed as the quotient of its cubic
invariants, so it cannot supply the indicated continuation. This
obstruction applies to continuation of the displayed elliptic Albanese
family.
\end{enumerate}
\end{proof}

\subsection{Classification of the non-normal cusp surface}\hypertarget{tag-01PO}{}\label{fam:section:C.4}
\begin{theorem}[Gluing classification of the cusp surface]\hypertarget{tag-01QQ}{}\label{fam:src:C.4.1}\label{fam:st:C.4.1}
\StatementRecord{fam:src:C.4.1}{proved}{01QQ}{\SeriesRef{fam:src:C.1.1}}

For the period normalization of \cref{fam:src:C.1.1}, the three opposite-side identification
constants of \(W_c\), in the characters specified below, are
\[
 (g_1,g_2,g_3)=(-a^{-1},a^6,-a^{-7}e^{-2\pi ic}).
\]
Two such compact complex surfaces are biholomorphic exactly when
\(c'-c\in\mathbb Z\). The vertical \(\mathbb C^*\)-action inherited from
translation in the \(\widehat\delta\)-direction fixes exactly the double
curve represented by the pair of edges of direction \(e_2\).
\end{theorem}


\begin{proof}
\begin{enumerate}[label=\arabic*.,leftmargin=*]
\item Keep the lattice of toric cocharacters \(\mathbb Z^2\), with
\(e_1=\widehat w\), \(e_2=\widehat\delta\), and cusp logarithm
\(s=(2\pi i)^{-1}\log t_c\). The cusp period matrix is
\([sB_0+C(t_c)\mid I]\), where
\[
 B_0=\begin{pmatrix}0&1\\-1&0\end{pmatrix},\qquad
 C(0)=\begin{pmatrix}6\mu_0&h_0\\b_0-h_0&\mu_0\end{pmatrix},
 \quad b_0=c.
\]
For \(v=e_1,e_2,e_1-e_2\), the gluing map from the side of direction
\(v\) to its opposite is the restriction of the toric action with lattice
label \(B_0v\). A character row \(r\) with \(rv=0\) pulls back by this
map to itself multiplied by
\(\exp(2\pi i\,rC(0)B_0v)\): the remaining power of \(t_c\) is
\(rB_0^2v=-rv=0\). Taking respectively
\(r=(0,1),(1,0),(1,1)\) gives exponents
\(-\mu_0,6\mu_0,-7\mu_0-c\). This proves the three constants, including
the cancellation of \(h_0\). These are morphisms of the whole projective
boundary lines, since multiplication by a nonzero constant extends over
both endpoints. The restrictions agree at both triple points, as the
ambient toric action supplies the same vertex identifications.

\item A biholomorphism of the two surfaces lifts uniquely to a biholomorphism
of their normalizations: integral closure of the reduced local rings is
functorial under an isomorphism. It preserves the inverse image of the
singular locus, the hexagon on \(\mathrm{dP}_6\). Its permutation of the
six components preserves their cyclic incidence, hence lies in the
order-twelve dihedral group. After composition with the corresponding
toric lattice automorphism, it fixes each boundary component. Each torus
character has its divisor supported on those components, so its pullback
divided by itself has divisor zero on the compact normalization and is
constant. Such an automorphism is therefore multiplication by an element
of \((\mathbb C^*)^2\). These torus multipliers cancel in the opposite-side
comparison because the same character is used on both sides.

\item The remaining six effective actions on the triple, with the central
lattice automorphism \(-I\) acting trivially, are described by the equations
\(g'_{k_j}{}^{\epsilon_j}=g_j\), where
\[
 ((k_1,\epsilon_1),(k_2,\epsilon_2),(k_3,\epsilon_3))
 \in\left\{
\begin{array}{lll}
((1,+),(2,+),(3,+)),&((1,-),(3,+),(2,+)),\\
((2,-),(3,-),(1,+)),&((2,+),(1,+),(3,-)),\\
((3,-),(2,-),(1,-)),&((3,+),(1,-),(2,-)).
\end{array}\right.
\]
Enumerate the $36$ ordered pairs of ray images of an integral basis,
retain those with determinant $\pm1$, and require preservation of the six
rays. This gives the $12$ fan automorphisms. For
$Av_j=\alpha v_k$ and $r_kA=\beta r_j$, the transported equation is
$(g'_k)^{\alpha\beta}=g_j$, yielding the six displayed actions.
Every nonidentity action requires one of
\(g_1^{-1}=g_1\), \(g_2^{-1}=g_2\), or equality of a positive real number
\(g_2^{\pm1}\) with a negative real number \(g_1^{\pm1}\).
All fail for \(a>1\). Therefore \(g'_3=g_3\), precisely
\(e^{-2\pi ic'}=e^{-2\pi ic}\). Conversely equality of the three constants
gives the identity gluing morphism on the normalization and hence a
biholomorphism of the quotients. This proves the unmarked assertion.

\item For the action of the subtorus with cocharacter \(e_2\), a one-dimensional
toric orbit on the normalization is pointwise fixed exactly when its
stabilizer line contains \(e_2\). Among the six rays this occurs exactly
for \(\pm e_2\). Their images form one double curve; all zero-dimensional
orbits map to its two endpoints. No dense two-dimensional orbit is fixed.
The quotient action has the same fixed set: a continuous orbit cannot
jump between the finitely many branches of a normalization fibre.
\end{enumerate}
\end{proof}


\section{Cotangent direct images and the period scalar}\label{pub:I:cotangent}
\hypertarget{tag-01TH}{}\label{inv:sec:partA}

\subsection{Root-atlas cohomology and logarithmic bundles}\hypertarget{tag-02ZG}{}\label{push:section:E.3}

\begin{proposition}[Root-atlas base change and relative duality]\hypertarget{tag-033I}{}\label{push:src:E.3.1}
\StatementRecord{push:src:E.3.1}{proved}{033I}{the normalized torus root charts and the periodic triangular cusp charts; \cite[Chapter III, Theorem 4.12(ii), pp.~134--135, and Theorem 4.10, pp.~132--133]{BS76}; \cite[introduction, equation (1), p.~261; Section 1, pp.~264--265; Proposition 2.1, p.~266]{RRV71}}

Let $k,l$ be positive integers with $\gcd(3k,4l)=1$.
For the root lift $\bar f:X\to\mathfrak B_{k,l}=\mathbb P(3k,4l)$
of the construction in \cref{fam:src:B.3.1}, the direct images
$R^q\bar f_*\mathcal O_X$, for $q=0,1,2$, are locally free of
ranks $1,2,1$. The canonical geometric-point maps
\[
 (R^q\bar f_*\mathcal O_X)\otimes k(b)
 \longrightarrow H^q(X_b,\mathcal O_{X_b})
\]
are isomorphisms, and the trace induces the canonical isomorphism
\[
 (R^2\bar f_*\mathcal O_X)^\vee
 \simeq\bar f_*\omega_{X/\mathfrak B_{k,l}}.
\]
Both comparisons commute with the analytic root-atlas groupoid.
Here the geometric fibre at a finite root point is the smooth torus
on the normalized root chart.
\end{proposition}

\begin{proof}
The fibre-divisor identities $f^*p_1=3kS_1$ and $f^*p_2=4lS_2$
define the root lift. On a full root chart its atlas pullback is
\[
 [J'_j/\mu_{m_j}]
 \mathbin{\times_{[\Delta_{s_j}/\mu_{m_j}]}}
 \Delta_{s_j}=J'_j,
 \qquad (m_1,m_2)=(3k,4l).
\]
The affine quotient action is free and the family $J'_j\to\Delta_{s_j}$
is the smooth proper torus family. The atlas map is \'etale.
At the cusp the local charts are $q=z_0z_1z_2$, with the
two-factor and smooth charts on the other strata. They have smooth
total space and pure two-dimensional Cartier fibres. A base parameter
is a non-zero-divisor, so the atlas family is flat. Properness is
properness of the compact torus and cusp fillings. These local
comparisons make the atlas families proper flat Gorenstein morphisms
of relative dimension two.

A smooth two-torus has structure-sheaf dimensions $(1,2,1)$.
Let $\nu:P\to W$ be the normalization, $P=\mathrm{dP}_6$, let
$C_h,C_v,C_d$ be the three double-curve normalizations, and let
$P_+,P_-$ be the two triple points. Each $C_e$ is $\mathbb P^1$.
Write $\rho_e^0,\rho_e^1:C_e\to P$ for its two paired boundary maps
and $j_e=\nu\rho_e^0=\nu\rho_e^1$. The sheaf sequence is
\begin{equation}\hypertarget{tag-01OY}{}\label{fam:eq:ab-cusp-resolution}
 0\longrightarrow\mathcal O_W\longrightarrow\nu_*\mathcal O_P
 \xrightarrow{\delta^0}\bigoplus_{e=h,v,d}j_{e*}\mathcal O_{C_e}
 \xrightarrow{\delta^1}\mathbb C_{P_+}\oplus\mathbb C_{P_-}
 \longrightarrow0,
\end{equation}
The differentials in this complex are
\begin{equation}\hypertarget{tag-01OJ}{}\label{eq:P5normalization}
\begin{aligned}
 \delta^0(s)_e&=(\rho_e^1)^*s-(\rho_e^0)^*s,\\
 \delta^1(t)&=\bigl((t_h-t_v+t_d)(P_+),
                        (t_h-t_v+t_d)(P_-)\bigr).
\end{aligned}
\end{equation}
Its exactness is local: for $xyz=0$ it is the alternating restriction
complex of the three coordinate planes, their pairwise lines and their
common point. Decomposing convergent power series by the variables
occurring in each monomial gives the augmented simplex complex, which
is exact. The one- and two-branch cases are the same calculation with
fewer planes. This verifies exactness for the boundary maps,
including the torus multipliers in their gluing.

All positive-degree structure-sheaf cohomology of $P,C_e$ vanishes.
Consequently \eqref{fam:eq:ab-cusp-resolution} computes cohomology by
\[
 \mathbb C\xrightarrow{0}\mathbb C^3
 \xrightarrow{\left(\begin{smallmatrix}1&-1&1\\1&-1&1\end{smallmatrix}\right)}
 \mathbb C^2.
\]
It gives $(h^0,h^1,h^2)(\mathcal O_W)=(1,2,1)$.
The branch-intersection resolution
is the augmented semisimplicial sheaf algebra with its restriction face
maps; its local exactness makes the augmentation a multiplicative
quasi-isomorphism. Its Alexander--Whitney product therefore computes
the sheaf-cohomology cup product. In the ordered triangles
$[0,e_1,e_2]$ and $[e_1,e_2,e_1+e_2]$ it is
\[
 \alpha\smile\beta=(\alpha_h\beta_d,\alpha_d\beta_h),
 \qquad [\alpha]\smile[\beta]
 =\alpha_h\beta_d-\alpha_d\beta_h
 \quad\hbox{in }\mathbb C^2/\mathbb C(1,1).
\]
The cocycles $(1,1,0)$ and $(0,1,1)$ pair to one. Thus
$\bigwedge^2H^1(\mathcal O_W)\to H^2(\mathcal O_W)$ is an
isomorphism.
Thus the three dimensions are constant on every
atlas chart, including its centre.

The constant-cohomology theorem
\cite[Chapter III, Theorem 4.12(ii), pp.~134--135]{BS76}
now gives local freeness, and the canonical analytic base-change
maps are those of
\cite[Chapter III, Theorem 4.10, pp.~132--133]{BS76}.
The hypotheses are properness, flatness of the structure sheaf,
a reduced smooth base, and the constant fibre dimensions just
computed. The constant-dimension formulation also appears in
\cite[Section 7, Satz 5, pp.~63--64]{Grauert60}, with
\cite[pp.~35--36]{Grauert1963ErratumD}.
On a double atlas overlap, both pullbacks identify with the direct
image of the same pulled-back family. Functoriality makes these
identifications satisfy the triple-overlap cocycle condition, so
they descend with their stabilizer representations.

For an atlas family $p:Y\to U$, proper analytic duality gives
\[
 Rp_*\omega_{Y/U}[2]
 \simeq R\mathcal Hom_U(Rp_*\mathcal O_Y,\mathcal O_U).
\]
This is the canonical trace identity of
\cite[introduction, equation (1), p.~261; Section 1, pp.~264--265;
Proposition 2.1, p.~266]{RRV71}.
The smooth-space dualizing complex is $K_Y[3]$ and the relative
Gorenstein complex is $\omega_{Y/U}[2]$; see
\cite[Section 3, Theorem 1, p.~81; Proposition 2 and the
complete-intersection calculation, p.~91]{RR70push}.
The locally free cohomology sheaves already obtained have no
positive sheaf Ext against $\mathcal O_U$. Degree $-2$ of the
trace identity is therefore
$p_*\omega_{Y/U}\simeq(R^2p_*\mathcal O_Y)^\vee$.
The trace commutes with the local biholomorphisms on the atlas
and has the same cocycle compatibility. It consequently descends
to the asserted isomorphism on $\mathfrak B_{k,l}$.
\end{proof}

\begin{proposition}[Vertical and logarithmic tangent bundles]\hypertarget{tag-01SN}{}\label{inv:D1.2}
\StatementRecord{inv:D1.2}{proved}{01SN}{\SeriesRef{inv:D1.1}}

There is a canonical representable proper flat holomorphic map
$\bar f:X\to\mathcal B$ with $f=\rho\bar f$. If
\[
 \mathcal T=T_X(-\log W),\qquad
 \mathcal V=\ker\bigl(\mathcal T\to
                   \bar f^*T_{\mathcal B}(-p_0)\bigr),
\]
then, for rank-two vector bundles $E,G$ on $\mathcal B$,
\begin{equation}
 \mathcal V=\bar f^*E,\quad E=\OO\oplus\OO(-1),\qquad
 G=R^1\bar f_*\OO_X=\OO\oplus\OO(-1),
\label{pub:I:eq:D-2}
\end{equation}
\begin{equation}
 \hypertarget{tag-01OZ}{}\label{fam:eq:ab-direct-images}
R^0\bar f_*\OO_X=\OO,\qquad R^2\bar f_*\OO_X=\OO(-1).
\end{equation}
and the following is an exact sequence of locally free sheaves on $X$:
\begin{equation}
 \hypertarget{tag-01P0}{}\label{fam:eq:ab-log-tangent}
0\longrightarrow\bar f^*E\longrightarrow\mathcal T
 \longrightarrow\bar f^*\OO(-5)\longrightarrow0.       
\end{equation}
The trivial summands in $E,G$ are generated by the vertical field $\xi$
and the constant integral cocycle $\gamma$, respectively.
\end{proposition}



\begin{proof}
These bundles and maps are the standard-data calculation of
\cite[Lemma~6.2, pp.~74--75; Lemma~6.7, p.~78; Proposition~6.6, p.~77; Proposition~6.10, pp.~79--80]{DuEtAl2026v1}.
\begin{enumerate}[label=\arabic*.,leftmargin=*]
\item The fibre-divisor root objects and their proper flat
atlas comparison are those of \cref{push:src:E.3.1} with $k=l=1$.
\item On a smooth torus chart a vertical field is a holomorphic linear
combination of the translation fields. Its coefficients are constant
on each compact fibre and therefore depend only on the base. The
vertical derivative of an affine transition is its linear part $R_g$;
the translating part, including the logarithmic regluing, has vertical
derivative the identity. At the cusp the two fields
$x_1\partial_{x_1},x_2\partial_{x_2}$ extend as logarithmic vector
fields in every toric chart. The deck maps multiply each $x_i$ by a
base-dependent unit and a power of the base parameter. Their vertical
logarithmic derivative fixes both coefficient directions. The
exponential overlap sends $\partial_{\zeta_i}$ to
$2\pi i x_i\partial_{x_i}$. These transitions depend only on the base
and define $E$; they prove the bundle identity
$\mathcal V=\bar f^*E$.

\item At
$t=xy$, the kernel of $df$ is freely generated by
$x\partial_x-y\partial_y$ and $\partial_z$. At $t=xyz$ it is generated
by $x\partial_x-y\partial_y$ and $x\partial_x-z\partial_z$.
Indeed the equation $yz a+xz b+xy c=0$ in the regular local ring forces
$x\mid a,y\mid b,z\mid c$; after division the only relation is that
the three coefficients sum to zero. The logarithmic map sends
$\sum a_i z_i\partial_{z_i}$ to $(\sum a_i)t\partial_t$ and is
surjective. At a finite root point it is the ordinary relative tangent
sequence of the smooth map to the root coordinate $s$. These comparisons prove \eqref{fam:eq:ab-log-tangent}, since
$T_{\mathcal B}(-p_0)=\OO(7-12)$.

\item The invariant vector $e_2$ defines a subbundle $\OO\xi\subset E$.
It remains nonzero as a vector of the logarithmic vertical bundle at
the cusp, even where its image in the ordinary tangent fibre is zero.
The quotient line has automorphy factors $-1/\tau$ and $1/\tau$ under
the order-three and order-four generators. The modular discriminant
has weight twelve and no zero in the upper half-plane. Therefore its
reciprocal is a meromorphic section of the twelfth tensor power of
this quotient line. On the cusp chart
$\tau=(\log q)/(2\pi i)+h(q)$ with $h$ holomorphic, so the product
expansion of the discriminant is $q$ times a holomorphic unit.
There are no other zeros or poles, including on the finite root
charts. The quotient line thus has twelfth power
$\OO(-p_0)=\OO(-12)$, and the torsion-free Picard group in \ref{inv:D1.1} makes
it $\OO(-1)$. Its extension by $\OO$ splits because
$H^1(\mathcal B,\OO(1))=0$.

\item Proposition~\ref{push:src:E.3.1} gives local freeness,
geometric-fibre comparison and trace duality on the root atlas.
The determinant of \eqref{fam:eq:ab-log-tangent}, using $W=\bar f^*p_0$, gives
\begin{equation}
 \hypertarget{tag-01P2}{}\label{fam:eq:ab-relative-dualizing}
\omega_{X/\mathcal B}=\bar f^*(\det E)^{-1}.
\end{equation}
Relative Serre duality, applied on that same atlas and descended by
its trace pairing, gives $R^2\bar f_*\OO_X=\det E=\OO(-1)$.
Cup product is an isomorphism $\det G\to R^2\bar f_*\OO_X$, since it
is so on every geometric fibre by the preceding calculation.

\item The class $\gamma$ is nowhere zero in $G$. On a smooth marked fibre,
$H^1(\OO)$ is the quotient of constant lattice cocycles by the two
period rows. If $\gamma$ were a combination of those rows, evaluation
on $\widehat w,\widehat\delta$ would make both coefficients zero,
contradicting $\gamma(\widehat\gamma)=1$. The same argument applies
on each finite root cover, and affine translations act trivially on
these classes. At the cusp $\gamma$ kills
$\Lambda_{\rm tor}=\mathbb Z\langle\widehat w,\widehat\delta\rangle$.
Local sections of the free toric quotient with transitions
$s_i=\Psi_{\lambda_{ij}}s_j$ give the holomorphic constant \v Cech
cocycle $\gamma(\lambda_{ij})$ on the whole neighbourhood.
For $B_0=\left(\begin{smallmatrix}0&1\\-1&0\end{smallmatrix}\right)$,
the edges $e_1,e_2,e_2-e_1$ have deck vectors
$(0,1),(-1,0),(-1,-1)$, so its central values are $(0,-1,-1)$.
Their cup product with $(1,0,-1)$ is one by
\eqref{fam:eq:ab-cusp-resolution} and its displayed product.
On the punctured overlap the exponential map kills exactly
$\Lambda_{\rm tor}$, giving the same full-lattice cocycle.
The Leray edge map over the Stein cusp disc is an isomorphism,
since its positive coherent cohomology vanishes. This gives
a holomorphic nonvanishing extension of $[\gamma]$.
Thus
$\OO\gamma\subset G$ is a subbundle; its quotient is $\OO(-1)$ by
the determinant just computed. Once again $H^1(\OO(1))=0$ splits the
extension. This proves every assertion in \eqref{pub:I:eq:D-2}.
\end{enumerate}
\end{proof}



\begin{definition}[Logarithmic cotangent notation]
\hypertarget{tag-PI17}{}\label{I:def:coarse-log-cotangent}
\StatementRecord{I:def:coarse-log-cotangent}{proved}{PI17}{}

For an admissible marked member, set
\[
 \omega_f=K_X\otimes f^*K_B^{-1},\qquad
 \mathcal L=\Omega_B^1(\log\Sigma)=\OO_B(1),\qquad
 \mathcal A=\Omega_X^1(\log D),\qquad
 \mathcal E=\mathcal A/f^*\mathcal L.
\]
All coherent cohomology groups have complex coefficients.
The symbols $\gamma,u$ in $R^1f_*\OO_X$ denote the images of
the marked integral classes under the constant-to-coherent
cohomology map. Representatives in the Dolbeault resolution are
specified separately by their displayed forms.
\end{definition}

\begin{cited}[Structure-sheaf direct images]\hypertarget{tag-01RQ}{}\label{inv:A1.0}\label{wb:thm:structure-direct-images-audit}\label{wb:lem:cusp-structure-sheaf-audit}
\StatementRecord{inv:A1.0}{proved}{01RQ}{\cite[Lemma~9.12, pp.~72--74; Proposition~9.13, pp.~74--76; p.~92]{AC2026}}

This is the structure-sheaf calculation of Claude and Alp\"oge: \cite[Lemma~9.12, pp.~72--74; Proposition~9.13, pp.~74--76, and the base-change clarification on p.~92]{AC2026}. The direct images $R^qf_*\OO_X$, for $q=0,1,2$, are
$\OO$, $G=\OO\gamma\oplus\OO(-1)$, and $\OO(-1)$.
The inclusion $\OO\gamma\subset G$ is a subbundle.  At the cusp,
$(\gamma,u)$ is a holomorphic frame, and cup product identifies
$\bigwedge^2G$ with $R^2f_*\OO_X$.  On the normalized root family, the cohomology pullback by the
cyclic generator has characters $(1,\zeta^{-a})$, where
$(m,a)=(3,2)$ or $(4,1)$ and $\zeta=e^{-2\pi i/m}$.
These assertions concern full scheme fibres, including $3S_1$ and $4S_2$.
At the cusp, the natural map $H^1(W;\mathbb C)\to H^1(W,\mathcal O_W)$ is an isomorphism. In degree two the comparison is surjective and is injective on $\ker[H^2(W;\mathbb C)\to H^2(\mathrm{dP}_6;\mathbb C)]$. The cup product $\bigwedge^2H^1(W,\mathcal O_W)\to H^2(W,\mathcal O_W)$ is an isomorphism. Every base-change map $(R^qf_*\mathcal O_X)\otimes k(b)\to H^q(X_b,\mathcal O_{X_b})$ is an isomorphism, including at the nonreduced fibres.
\end{cited}

\begin{proposition}[The root-to-coarse comparison maps]
\hypertarget{tag-PI10}{}\label{I:root-coarse-comparison}
\StatementRecord{I:root-coarse-comparison}{proved}{PI10}{\cref{inv:D1.1,inv:D1.2,push:src:E.3.1}}
Write $G_B=R^1f_*\mathcal O_X$ and
$G_{\mathcal B}=R^1\bar f_*\mathcal O_X$. Then
\[
 G_B=\rho_*G_{\mathcal B}.
\]
On a root disc $t=s^m$, take a frame $(e_1,e_2)$ in which the
covariant action is $(1,\zeta^a)$, where
$(m,a)=(3,2)$ or $(4,1)$ and $\zeta=e^{-2\pi i/m}$.
The invariant frame is $(e_1,s^ae_2)$, and the adjunction map is
\[
 \rho^*G_B\longrightarrow G_{\mathcal B},
 \qquad (u,v)\longmapsto(u,s^av).
\]
It is the map induced by normalization after the root base change.
The covariant action on cohomology is
$(g^{-1})^*:G_{\mathcal B,s}\to G_{\mathcal B,gs}$;
the pullback $g^*:G_{\mathcal B,gs}\to G_{\mathcal B,s}$ has
the inverse eigenvalues $(1,\zeta^{-a})$.

For the relative logarithmic cotangent bundle of
\cref{I:def:root-bundles}, the evaluation pairing gives
\[
 \mathcal E=\bar f^*(E^\vee),
\]
where $E$ is the vertical root bundle in \cref{inv:D1.2}.
\end{proposition}
\begin{proof}
The factorization $f=\rho\bar f$ and exactness of finite-group
invariants give $R^qf_*\mathcal O_X=
\rho_*R^q\bar f_*\mathcal O_X$.
In the chosen root frame, invariance of $h(s)e_2$ is exactly
$h(\zeta s)=\zeta^a h(s)$. Its convergent series is therefore
$s^aH(s^m)$. This gives both the invariant frame and the displayed
adjunction matrix. The normalization map agrees with this map on
the punctured disc by its marked torus identification, hence on
the whole disc by holomorphic continuation. At $s=0$ its second
component is zero; this specifies the comparison with the
normalized torus fibre while the coarse base-change theorem uses
the full scheme fibre. Inverting the covariant action gives the
displayed pullback eigenvalues.

On an ordinary chart the pairing is
$\langle d\zeta_i,\partial_{\zeta_j}\rangle=\delta_{ij}$.
At the cusp the exponential map sends
$\partial_{\zeta_i}$ to $2\pi i x_i\partial_{x_i}$ and
$d\zeta_i$ modulo the base differential to
$(2\pi i)^{-1}d\log x_i$, preserving that pairing.
On a finite root chart, $dt/t=m\,ds/s$ removes exactly the
base logarithmic direction from the quotient defining
$\mathcal E$. Thus these local perfect pairings agree under
the transition maps and give $\mathcal E=\bar f^*(E^\vee)$.
\end{proof}





\Uses{the hypotheses, notation and definitions stated here; \citet[Lemma~9.12 and Proposition~9.13]{AC2026}; \CompendiumRef{I}{app:rederive:inv-A1.0}}





\begin{proposition}[Local freeness and the determinant line]\hypertarget{tag-01RR}{}\label{inv:A1.1}
\StatementRecord{inv:A1.1}{proved}{01RR}{the definitions and constructions stated in the proof}

Both $\mathcal A$ and $\mathcal E$ are locally free, of ranks three and
two.  The vertical vector $\xi=\partial/\partial\zeta_2$ is nowhere
zero as a section of $\mathcal E^\vee$.  There is an exact sequence
of holomorphic vector bundles
\begin{equation}\hypertarget{tag-01SZ}{}\label{inv:eq:A-E-line}
0\longrightarrow N\longrightarrow\mathcal E
 \xrightarrow{\langle\cdot,\xi\rangle}\OO_X\longrightarrow0,
\qquad N=\det\mathcal E=\omega_f(-2S_1-3S_2).
\end{equation}
\end{proposition}


\begin{proof}
\begin{enumerate}[label=\arabic*.,leftmargin=*]
\item At a smooth point of a reduced cusp branch the assertion is the usual
coordinate calculation.  At a double point write $t_c=xy$ with third
coordinate $z$; then
$\mathcal A=R\,d\log x\oplus R\,d\log y\oplus R\,dz$ and its
base subline is $R(d\log x+d\log y)$.  At either triple point write
$t_c=xyz$; then $\mathcal A$ has the three logarithmic coordinate
basis vectors and its base subline is their sum.  In each case that
subline is a direct summand.  These calculations apply on every one
of the three double curves, including their endpoints.

\item At $S_j$ we use the specified free cyclic quotient
$q:J'_j\to N_j$, which is finite and \'etale on total spaces, and
$t_j=s^{m_j}$ on the base.  On $J'_j$ the logarithmic conormal is
$m_j\,ds/s$, and its coefficient is a unit over $\C$.  It is again a
direct summand of the logarithmic cotangent bundle.  The comparison
$q^*\mathcal A=\Omega_{J'_j}^1(\log(s=0))$ identifies the quotient
bundles and their exact sequence; faithful flat descent proves the
claim downstairs.

\item On the marked torus family $\xi$ is the constant second vertical
vector.  The linear monodromies fix it, and translations preserve it.
On a cusp toric chart it is a nonzero constant vector in the relative
logarithmic tangent lattice of rank two.  A toric coordinate change
$x_i\mapsto c_i(t_c)t_c^{b_i}x_i$ changes $d\log x_i$ by a multiple
of $d\log t_c$, so this constant relative vector descends and is
nowhere zero in that bundle, including where its image in the
ordinary tangent fibre vanishes. Hence evaluation is surjective and its kernel
is the determinant line.  Taking determinants of
$0\to f^*\mathcal L\to\mathcal A\to\mathcal E\to0$ gives
\[
\det\mathcal E=K_X(D)\otimes f^*\mathcal L^{-1}
=\omega_f\bigl(D-f^*\Sigma\bigr)
=\omega_f(-2S_1-3S_2),
\]
which proves \eqref{inv:eq:A-E-line}.
\end{enumerate}
\end{proof}

\begin{proposition}[The logarithmic direct images by vertical twisting and duality]\hypertarget{tag-01RS}{}\label{inv:A1.2}\label{wb:prop:log-relative-direct-images}
\StatementRecord{inv:A1.2}{proved}{01RS}{\SeriesRef{inv:eq:A-H-images}, \SeriesRef{inv:eq:A-E-line}, \SeriesRef{inv:eq:A-N-images}, \SeriesRef{inv:A1.0}}

One has, as holomorphic vector bundles on $B$,
\[
\begin{array}{c|ccc}
q&0&1&2\\\hline
R^qf_*\mathcal E&\OO\oplus\OO(-1)&
 \OO^2\oplus\OO(-1)^2&\OO\oplus\OO(-1).
\end{array}
\]
They commute with base change to every full scheme fibre.
\end{proposition}


\begin{proof}
\begin{enumerate}[label=\arabic*.,leftmargin=*]
\item Put $H=2S_1+3S_2$.  Sections of the line bundle $\OO_X(H)$
on the cover at $p_j$ allow poles of order at most $m_j-1$
in $s$.  A fixed character has invariant exponent zero and gains no
pole, whereas the variable character $\zeta^{-a}$ has least allowed
invariant exponent $a-m$.  Thus the inclusion
$R^1f_*\OO\to R^1f_*\OO(H)$ is the identity on $\OO\gamma$ and
changes $s^ab_1$ to $s^{a-m}b_1=t_j^{-1}s^ab_1$ on its quotient.
It makes exactly one upper modification at each finite point.
The same computation applies to $R^2$.  No modification is made at
the cusp or a regular fibre.  The result is
\begin{equation}\hypertarget{tag-01T0}{}\label{inv:eq:A-H-images}
f_*\OO(H)=\OO,\qquad
R^1f_*\OO(H)=\OO\oplus\OO(1),\qquad
R^2f_*\OO(H)=\OO(1).
\end{equation}
Here the middle extension splits because $H^1(\OO(-1))=0$.
The local invariant modules are free.  Tensoring them with
$\C\{t_j\}/(t_j)$ gives invariants of
$s^{-(m_j-1)}\C\{s\}/(s^m)$, with the indicated surviving character
basis vectors.  The Reynolds idempotent is linear over the base,
so this calculation also proves the comparison with the nonreduced
fibre.  At the cusp the line bundle is trivial.  Cohomology and base
change, applied successively from the top degree, proves the claimed
local freeness and comparison for \eqref{inv:eq:A-H-images}.

\item The relative duality map we use is the cup product followed
by the analytic trace, giving a morphism of bundles
\[
R^qf_*(\omega_f\otimes\OO(-H))\longrightarrow
\bigl(R^{2-q}f_*\OO(H)\bigr)^\vee.
\]
Every fibre is a compact Gorenstein surface: it is the Cartier fibre
in a smooth threefold.  Serre duality on that scheme makes the map
an isomorphism on its residue field, also for $m_jS_j$.  The constant
dimensions just computed imply local freeness on both sides, so
Nakayama's lemma makes the bundle map an isomorphism.  Therefore
\begin{equation}\hypertarget{tag-01T1}{}\label{inv:eq:A-N-images}
(f_*N,R^1f_*N,R^2f_*N)
=(\OO(-1),\OO\oplus\OO(-1),\OO).
\end{equation}

\item The connecting maps of \eqref{inv:eq:A-E-line} are zero on a regular
fibre: its relative cotangent is a trivial bundle and evaluation on
a constant nonzero vector splits.  Source and target direct images
in \eqref{inv:eq:A-N-images} and \cref{inv:A1.0} are vector bundles,
so these holomorphic connecting maps, zero on a dense open set, are
zero everywhere.  We obtain extensions in each degree of the listed
$R^qN$ by $R^q\OO$.  Every resulting pair of line-bundle degrees is
zero or minus one; their extension groups are
$H^1(\OO(-1)),H^1(\OO),H^1(\OO(1))$, all zero.  These extensions
give exactly the three asserted bundles.  Applying the already proved
fibre comparisons to the short exact sequences proves base change
for $R^q\mathcal E$ as well.  
\end{enumerate}
\end{proof}



\subsection{The elliptic logarithm and the period determinant}

\begin{definition}[The finite local coordinates]
\hypertarget{tag-PI07}{}\label{I:def:finite-coordinates}
\StatementRecord{I:def:finite-coordinates}{proved}{PI07}{}
At the order-four point choose $t-1=s^4$ so that on the positive
$t>1$ arc toward the cusp, $\tau=iy$ with $y>1$, and
\begin{equation}\hypertarget{tag-01T5}{}\label{inv:eq:A-s-branch}
\tau=i+\frac{i}{\pi\sqrt{E_4(i)}}s^2+O(s^4),\qquad
\operatorname{Im}\mu_s(0)>0.
\end{equation}
Replacing $s$ by $i^js$ changes $\mu_s(0)$ by $i^{-j}$.
For any other local coordinate $z$ the derivative is transported
by $\mu_z=\mu_s s_z$. At the order-three point use $r^3=t$,
with $r>0$ on the arc $0<t<1$ along which $\tau$ travels from
$\rho$ to $i$ on the unit circle.
\end{definition}

\begin{proposition}[The elliptic logarithm at the finite points]
\hypertarget{tag-01RW}{}\label{inv:A2.3}
\StatementRecord{inv:A2.3}{proved}{01RW}{\SeriesRef{fam:src:B.1.1}, \SeriesRef{fam:src:B.1.3}; \citet{DukeImamoglu2008}}
Let $[\mu]$ be the generator identified in
Proposition~\ref{fam:src:B.1.1} with
one of the two sections
$P_*=(0,\sqrt2t^2(t-1))$ and $-P_*$ of
$y^2=x^3-3t^3(t-1)x+2t^4(t-1)^2$.
Along the branch in \eqref{inv:eq:A-s-branch}, put $x_0=1-1/t$.
Then
\begin{equation}\hypertarget{tag-01SY}{}\label{inv:eq:A-DI-branch}
\mu=\frac{1-\tau}{2}
 +\frac{i\sqrt6}{3\pi}
 \frac{x_0^{1/4}\,{}_3F_2(1/3,2/3,1;3/4,5/4;x_0)}
 {{}_2F_1(1/12,5/12;1;1-x_0)}.
\end{equation}
The difference $\mu-(2-\tau)/3$ has exact order two at the
order-three point in the coordinate $r^3=t$.
\end{proposition}
\begin{proof}
The pointed Weierstrass isomorphism in
Proposition~\ref{fam:src:B.1.1} and the section calculation in
Proposition~\ref{fam:src:B.1.3} give $\wp(\mu,\tau)=0$.
In \cite[Theorem~1, equation~(5), author preprint p.~2]{DukeImamoglu2008},
write $z_0$ for the displayed zero. The exact affine map
$\mu=1-z_0$ preserves the zero identity, since
$\wp(1-z_0,\tau)=\wp(z_0,\tau)$, and gives
\eqref{inv:eq:A-DI-branch}. At $t=1+s^4$ its initial value is
$(1-i)/2$. The coefficient of $s$ is
\[
 \frac{i\sqrt6}{3\pi\,{}_2F_1(1/12,5/12;1;1)}
 =\frac{i\sqrt6}{3\pi E_4(i)^{1/4}}.
\]
It has positive imaginary part and therefore selects the branch of
\eqref{inv:eq:A-s-branch}. Continuation gives the chosen bounded
lift at the cusp.

The chord-and-tangent formulas applied to $P_*$ give
\[
 x(2P_*)=\frac{9t^2}{8},\qquad
 y(2P_*)=\frac{t^2(32-5t)}{16\sqrt2},\qquad
 x(3P_*)=\frac{16}{81}(t-1)(5t-32).
\]
The last expression tends to $512/81$ at $t=0$.
Let $F=f_{\rm sc}$ be the elliptic scale of
Proposition~\ref{fam:src:B.1.1}; it has exact order two in $r$.
For $\nu=\mu-(2-\tau)/3$, the point $3P$ has local logarithm
$3\nu$ modulo its period, and
$x(3P)=F^2\wp(3\nu,\tau)$.
The expansion $\wp(3\nu,\tau)=(3\nu)^{-2}+O(\nu^2)$
therefore proves that $\nu$ has exact order two. Its leading
coefficient is $\pm3F_2/(16\sqrt2)$, where
$F_2=\lim_{r\to0}F/r^2$; the sign is fixed by the same logarithm branch.
\end{proof}

\begin{proposition}[The complete raw period determinant]\hypertarget{tag-01RU}{}\label{inv:A2.1}\label{wb:prop:log-ks-divisor}\label{wb:prop:period-derivative-determinant}
\StatementRecord{inv:A2.1}{proved}{01RU}{the definitions and constructions stated in the proof}

On the marked upper half-plane,
\begin{equation}\hypertarget{tag-01T3}{}\label{inv:eq:A-modular-K}
6\mu_z^2-\tau_z\beta_z
=\tau_z^2\frac{E_4(\tau)^2}{E_6(\tau)}.
\end{equation}
The right side extends holomorphically through the apparent poles.
Its exact orders are two at $z_1$ and zero at $z_2$; it has no other
zero away from their translates.  The logarithmic connecting map
$\kappa_0:f_*\mathcal E\to\mathcal L\otimes G$ has divisor
$\operatorname{divisor}(\det\kappa_0)=p_1+p_2$.
\end{proposition}


\begin{proof}
\begin{enumerate}[label=\arabic*.,leftmargin=*]
\item Consider the holomorphic quadratic differential
$\mathscr Q=6(d\mu)^2-d\tau\,d\beta$.  Direct differentiation of
the two affine transformation laws in \cite[Definition~3.1, p.~14]{AC2026} gives
\[
g_j^*\mathscr Q=\tau^{-2}\mathscr Q\quad(j=1,2).
\]
In particular $g_2^2$ fixes $\tau$ and sends
\[
\mu\longmapsto1-\tau-\mu,\qquad
\beta\longmapsto\beta+6\tau+12\mu-6,
\]
which leaves $\mathscr Q$ exactly unchanged.
The kernel of
$\Delta=\Z/3*\Z/4\to\Z/3*\Z/2=\operatorname{PSL}_2(\Z)$
is the normal subgroup generated by $g_2^2$.  The squared multiplier
of $\mathscr Q$ is the ordinary modular multiplier
$(c\tau+d)^{-2}$, so it is one on every kernel element.  Two generic
points of the marked cover with the same $\tau$ have the same $t$
and differ by an element of $\Delta$; its modular image fixes a
generic point, hence is the identity.  It follows that
$k=\mathscr Q/(d\tau)^2$ is a single-valued meromorphic function of
$\tau$, transforming as a modular form of weight two.  

\item At $z_1$, write a linearizing parameter $r$ with
$g_1r=\zeta r$, $\zeta=e^{-2\pi i/3}$.  If $\mathscr Q=K_rdr^2$
has first nonzero coefficient of order $n$, its transformation law
at the fixed point reads $\zeta^{n+2}=\rho^{-2}=\zeta$.
Thus $n\equiv2\pmod3$, and $K_r$ vanishes to order at least two.
The map $\tau$ is unramified there by \cite[Theorem~3.4(i), p.~14; Proposition~3.6, pp.~16--17]{AC2026}.  At
$z_2$, $\tau-i$ has order two and $K_s$ is holomorphic; therefore
$k$ has at most a simple pole at $i$.  Away from these points
$\tau$ is unramified.  At the cusp $\mu$ and $\beta+\tau$ are
holomorphic in $q$, so $k=1+O(q)$.

\item The modular function $k E_6/E_4^2$ is therefore holomorphic at
$\rho$, at $i$, and at the cusp, as well as everywhere else.
The usual zero statements for $E_4,E_6$ (simple at $\rho,i$,
respectively) are the modular facts used in \cite[Section~3.1, pp.~15--17]{AC2026}.
This function descends to the compact modular curve $X(1)=\Pone$,
so it is constant.  Its cusp value is one.  This proves
\eqref{inv:eq:A-modular-K} and the exact two local orders.

\item For the coarse lattices, choose
cover eigenframes $e_0,e_1$ of $\mathcal E$ with characters
$1,\zeta^a$, and $b_0,b_1$ of $G$ with characters
$1,\zeta^{-a}$.  Their invariant bases are
$e_0,s^{m-a}e_1$ and $b_0,s^ab_1$.  The logarithmic base frame is
$\ell=dt_j/t_j=m\,ds/s$.  Changing determinants from these frames
to the marked ones gives
\[
\det\kappa_0=\epsilon(t_j)m^{-2}s^{m+2-2a}K_s,
\qquad \epsilon(0)\ne0.
\]
This equality follows by taking the two factors $s/m$ from replacing
$ds$ by $\ell$, the factor $s^{m-a}$ in the source, and the factor
$s^{-a}$ in the target.  Thus at $(m,a)=(3,2)$ the total order in
$s$ is $1+2=3$, and at $(4,1)$ it is $4+0=4$.
Both coarse zeros are simple.  At the cusp, the frame change has
no elementary divisor and $dZ/d\tau$ tends to
$\left(\begin{smallmatrix}0&1\\-1&0\end{smallmatrix}\right)$,
with determinant one.  Formula \eqref{inv:eq:A-modular-K} excludes any
further zero, proving the asserted divisor.
\end{enumerate}
\end{proof}

\begin{proposition}[The finite period coefficients]\hypertarget{tag-01RV}{}\label{inv:A2.2}
\StatementRecord{inv:A2.2}{proved}{01RV}{\SeriesRef{inv:eq:A-s-branch}, \SeriesRef{inv:eq:A-modular-K}, \SeriesRef{inv:A2.3}}

In the coordinate \eqref{inv:eq:A-s-branch},
\begin{equation}\hypertarget{tag-01T4}{}\label{inv:eq:A-mu-value}
\mu_s(0)=i\,C_4,
\qquad
C_4=\frac{\sqrt{2/3}}{\pi E_4(i)^{1/4}}
=\frac{4\sqrt\pi}{3^{3/4}\Gamma(1/4)^2}>0.
\end{equation}
For $r^3=t$, with $r>0$ on the arc $0<t<1$ and $\tau$ travelling
from $\rho$ to $i$ along the unit circle, set
\[
\lambda=\frac{3e^{5\pi i/6}}{2\pi E_6(\rho)^{1/3}}.
\]
Then $\tau=\rho+\lambda r+O(r^2)$ and
\[
K_r=-\frac{3i\lambda}{2\pi}r^2+O(r^3).
\]
The exact expressions have the rational enclosures
\[
\begin{aligned}
0.2366093140807852&<C_4<0.2366093140807853,\\
0.0801022110875382&<\operatorname{Re}K_2<0.0801022110875383,\\
0.1387410994022233&<\operatorname{Im}K_2<0.1387410994022234.
\end{aligned}
\]
\end{proposition}


\begin{proof}
\begin{enumerate}[label=\arabic*.,leftmargin=*]
\item Ramanujan's differential relations give
$E_6'(i)=-\pi i E_4(i)^2$ and
$E_4'(\rho)=-(2\pi i/3)E_6(\rho)$.
Using $j=1728E_4^3/(E_4^3-E_6^2)$, we obtain
\[
j(\tau)-1728=-1728\pi^2E_4(i)(\tau-i)^2+O((\tau-i)^3).
\]
Since $j=1728(1+s^4)$, the chosen branch is precisely
\eqref{inv:eq:A-s-branch}.  Formula \eqref{inv:eq:A-modular-K} now gives
\[
K_s(0)=-\frac4{\pi^2\sqrt{E_4(i)}}=6\mu_s(0)^2,
\]
because $\tau_s(0)=0$.  This proves nonvanishing and the first
formula in \eqref{inv:eq:A-mu-value}; the sign is fixed by the stated
coordinate convention.  Similarly
$j(\tau)=-512i\pi^3E_6(\rho)(\tau-\rho)^3+\cdots$.
This gives the stated $\lambda$, and substitution into
\eqref{inv:eq:A-modular-K} gives the coefficient $-3i\lambda/(2\pi)$.

\item The real period also has the integral expression
\[
\Omega_R=\frac{\sqrt2}{3^{1/4}}
\int_0^{\pi/2}(1-\tfrac12\sin^2\theta)^{-1/2}\,d\theta,
\qquad C_4=\frac{\sqrt2}{3\Omega_R}.
\]
It follows by the elliptic-integral substitution for the three roots
$\sqrt3,0,-\sqrt3$ of $Y^2=X^3-3X$ and the differential
$dX/(2Y)$. The beta-integral evaluation gives the gamma expression
in \eqref{inv:eq:A-mu-value}.

For the rational enclosures, evaluate the divisor-sum series of
$E_4(i)$ and $E_6(\rho)$ by interval arithmetic with $N=100$ terms.
The elementary inequalities $\sigma_3(n)\le n^4$,
$\sigma_5(n)\le n^6$, $e^{-2\pi}<1/500$ and
$e^{-\pi\sqrt3}<1/200$ give respective tail bounds
\[
250(N+1)^4 500^{-N-1},\qquad
1024(N+1)^6 200^{-N-1}\quad(N=100).
\]
The ratios of successive majorant terms are less than $1/480$
and $1/180$, respectively, which proves these tail estimates.
Outward rounding in the finite sums, together with these tail bounds,
gives the three stated enclosures. The preceding modular calculation
proves the exact nonvanishing and the local orders.
\end{enumerate}
\end{proof}



\begin{lemma}[Residues and the kernel of the logarithmic connecting map]
\hypertarget{tag-PI00}{}\label{pub:I:residue-kernel-rank}
\StatementRecord{pub:I:residue-kernel-rank}{proved}{PI00}{\SeriesRef{inv:A1.0}, \SeriesRef{inv:A1.1}, \SeriesRef{inv:A1.2}, \SeriesRef{inv:A2.1}}
For every admissible parameter, the logarithmic connecting map
$\kappa_1:R^1f_*\mathcal E\to\mathcal L\otimes R^2f_*\OO=\OO$
is surjective. Put
$\mathcal K=\ker\kappa_1$ and
$r=\operatorname{rank}H^0(\kappa_1)$. The residue maps give
\[
f_*\Omega_X^1=\Omega_B^1,\qquad
R^1f_*\Omega_X^1=\mathcal K,\qquad
R^2f_*\Omega_X^1=R^2f_*\mathcal E.
\]
Thus the first and third bundles are isomorphic to $\OO(-2)$ and
$\OO\oplus\OO(-1)$, respectively; the latter splitting is \cref{inv:A1.2}.
Moreover,
\[
h^0(\mathcal K)=2-r,\qquad h^1(\mathcal K)=1-r,\qquad
h^{1,2}=2-r,\qquad h^{1,1}=3-r.
\]
\end{lemma}
\begin{proof}
\begin{enumerate}[label=\arabic*.,leftmargin=*]
\item Let $d_{ij}(s)$ be the cover Kodaira--Spencer coefficients in the
eigenframes of \cref{inv:A2.1} after removing $ds$.
The invariant coarse matrix and its cup-product row have the forms
\[
[\kappa_0]=\begin{pmatrix}A&t_jB\\C&D_0\end{pmatrix},\qquad
[\kappa_1]=(-C,A,-D_0,B),
\]
in the ordered basis
$e_0b_0,s^ae_0b_1,s^{m-a}e_1b_0,e_1b_1$. Specifically
$A=sd_{00}/m$, $B=s^{1-a}d_{01}/m$,
$C=s^{1-a}d_{10}/m$, $D_0=s^{m+1-2a}d_{11}/m$.
These are holomorphic invariant functions; the row follows by
wedging the connecting class with $b_0,b_1$.
If the row had no unit at a finite point, every entry would vanish
modulo $t_j$ and $\det\kappa_0$ would vanish to order at least two,
contradicting \cref{inv:A2.1}. Away from the finite points, including
the cusp, $\kappa_0$ is invertible, which gives a unit in the
cup-product row. Thus $\kappa_1$ is surjective.

\item By \cref{inv:A2.1},
$\coker\kappa_0=\C_{p_1}\oplus\C_{p_2}$.
The logarithmic conormal sequence and \cref{inv:A1.2} give
\[
f_*\mathcal A=\OO(1),\quad
0\to T\to R^1f_*\mathcal A\to\mathcal K\to0,\quad
R^2f_*\mathcal A=\OO\oplus\OO(-1),\qquad
T=\C_{p_1}\oplus\C_{p_2}.
\]
Let $\nu:\widetilde D=\mathrm{dP}_6\sqcup S_1\sqcup S_2\to D$
be normalization. The local bases in \cref{inv:A1.1} give the exact
residue sequence
\[
0\to\Omega_X^1\to\mathcal A\to\nu_*\OO_{\widetilde D}\to0.
\]
Its three residue direct images are
$\C_{p_0}\oplus\C_{p_1}\oplus\C_{p_2}$,
$\C_{p_1}\oplus\C_{p_2}$, and zero: the normalization of $W$ is
rational, and a bielliptic surface has
$(h^0,h^1,h^2)(\OO)=(1,1,0)$.
In degree zero the map from $\OO(1)$ is local evaluation with the
nonzero multipliers $1,3,4$. Its kernel is $\OO(-2)$.
Under the base logarithmic inclusion this kernel is $\Omega_B^1$,
and its map to $f_*\Omega_X^1$ is pullback.

\item At $p_j$, in the invariant frame $(b_0,s^ab_1)$ of $G$, residue
of $(dt_j/t_j)(x b_0+y s^ab_1)$ in $H^1(S_j,\OO)$ is
$m_jx[b_0]$. The residue factors through the base logarithmic line
on the covering chart. Naturality of the connecting map gives
$\operatorname{res}_1\kappa_0=0$.
The determinant of $\kappa_0$ has a simple zero, so its reduction
has rank one, and its image is exactly the one-dimensional kernel
$\C s^ab_1$ of this residue. Hence
$\coker\kappa_0\to H^1(S_j,\OO)$ is an isomorphism.
Both torsion summands cancel, giving the stated canonical
identifications of $R^1f_*\Omega_X^1$ and $R^2f_*\Omega_X^1$.

\item The long exact cohomology sequence of
\[
0\to\mathcal K\to\OO^2\oplus\OO(-1)^2
 \xrightarrow{\kappa_1}\OO\to0
\]
gives $h^0(\mathcal K)=2-r$ and $h^1(\mathcal K)=1-r$.
Leray on the curve $B$ and the preceding direct images then give
\[
h^{1,2}=h^1(\mathcal K)+h^0(\OO\oplus\OO(-1))=2-r,
\qquad
h^{1,1}=h^1(\OO(-2))+h^0(\mathcal K)=3-r.
\]
\end{enumerate}
\end{proof}

\begin{proposition}[The global rank and the period scalar]\hypertarget{tag-01RX}{}\label{inv:A2.4}\label{wb:thm:h12-scalar}
\StatementRecord{inv:A2.4}{proved}{01RX}{\SeriesRef{pub:I:residue-kernel-rank}, \SeriesRef{inv:A2.1}, \SeriesRef{inv:A2.2}}

The map
$\kappa_1:R^1f_*\mathcal E\to\mathcal L\otimes R^2f_*\OO=\OO$
is surjective, and $H^0(\kappa_1):\C^2\to\C$ has rank one.
Let $h=a\,d\zeta_1+d\zeta_2$ be the unique global section of
$f_*\mathcal E$ with $h(\xi)=1$.  In an invariant target
frame normalized as below,
\[
\kappa_1(h\gamma)(p_2)=-\tfrac14\mu_s(0)=-\tfrac i4 C_4.
\]
The rank-one condition is equivalent to
$a\tau_z+\mu_z\not\equiv0$.  The same direct-image exact sequence expresses the two Hodge numbers
in terms of the global rank $r$ as
$h^{1,2}=2-r$ and $h^{1,1}=3-r$; here $r=1$.
\end{proposition}


\begin{proof}
\begin{enumerate}[label=\arabic*.,leftmargin=*]
\item Sheaf surjectivity is proved in
\cref{pub:I:residue-kernel-rank}. The lift $h$ exists and is unique because $f_*\mathcal E$ is an
extension of $\OO$ by $\OO(-1)$.  Its invariance under $g_2$ gives
$a(g_2z)=\tau a+\mu$, hence $a(z_2)=1/2$.
In the quotient frame $(\gamma,u)$ of $G$ the connecting map is
\[
\kappa_0(h)=\frac{s}{4}
\bigl((6a\mu_s+\beta_s)\gamma+(a\tau_s+\mu_s)u\bigr)\ell.
\]
Choose the variable eigenframe $b_1$ by Reynolds averaging so that
$\gamma\wedge b_1=v(s)\gamma\wedge u$, $v(0)=1$.
Then $s\,v(s)(\gamma\wedge u)\ell$ is the invariant
target frame. Cup product with $\gamma$ gives at the
central point $-(a\tau_s+\mu_s)/4=-\mu_s/4$, proving rank one.

\item To check the stated equivalence, there is a second global section
\[
q_E=\sigma_1\otimes u+6\sigma_2\otimes\gamma
\quad\text{of }R^1f_*\mathcal E.
\]
It is the negative of the relative $(1,1)$ image of the invariant
flat class $\eta=u\wedge w+6\gamma\wedge\delta$.
It extends at the cusp because the relative logarithmic forms and
$\gamma,u$ extend there.  At a finite point it is an invariant
holomorphic tensor on the smooth cover, hence belongs to the four
explicit invariant generators of \cref{pub:I:residue-kernel-rank}.  Its cup-product row is zero,
since the two terms are $6\mu_z$ and $-6\mu_z$.
It is independent of $h\gamma$ on every regular fibre (its
$\sigma_1\otimes u$ coefficient is nonzero), so these two sections
form a basis of the two-dimensional global section space.
Consequently the global map has rank one precisely when its value
on $h\gamma$, equivalently $a\tau_z+\mu_z$, is not identically
zero. The Hodge-number formulas follow from
\cref{pub:I:residue-kernel-rank}, and the nonzero scalar above gives
$r=1$.
\end{enumerate}
\end{proof}

\subsection{Absolute cotangent direct images}

\begin{proposition}[All absolute cotangent direct images]\hypertarget{tag-01RT}{}\label{inv:A1.3}\label{wb:prop:absolute-cotangent-direct-images}
\StatementRecord{inv:A1.3}{proved}{01RT}{\SeriesRef{pub:I:residue-kernel-rank}, \SeriesRef{inv:A2.4}}

For every admissible $c$, the six direct images are
\begin{equation}\hypertarget{tag-01T2}{}\label{inv:eq:A-absolute-table}
\begin{array}{c|cc}
q&R^qf_*\Omega_X^1&R^qf_*\Omega_X^2\\\hline
0&\OO(-2)&\OO(-2)\oplus\OO(-1)\\
1&\OO\oplus\OO(-1)^2&\OO(-2)\oplus\OO(-1)^2\\
2&\OO\oplus\OO(-1)&\OO.
\end{array}
\end{equation}
In particular every listed sheaf is torsion-free.
\end{proposition}


\begin{proof}
\begin{enumerate}[label=\arabic*.,leftmargin=*]
\item By \cref{pub:I:residue-kernel-rank}, the first-column direct
images are $\OO(-2)$, $\mathcal K=\ker\kappa_1$, and
$\OO\oplus\OO(-1)$. By \cref{inv:A2.4}, the global map
$H^0(\kappa_1)$ has rank one. Hence the bundle map
$\OO^2\oplus\OO(-1)^2\to\OO$ has a nonzero constant component
on a trivial summand. Eliminating all other components by bundle
automorphisms gives $\mathcal K=\OO\oplus\OO(-1)^2$.

\item Finally $\Omega_X^2=(\Omega_X^1)^\vee\otimes K_X$. Relative
duality, with the locally free direct images now established, gives
\[
R^qf_*\Omega_X^2
=\bigl(R^{2-q}f_*\Omega_X^1\bigr)^\vee\otimes K_B.
\]
Substitution of $K_B=\OO(-2)$ gives every entry of the second column.
\end{enumerate}
\end{proof}

\subsection{The Hodge table and its Leray generators}

\begin{proposition}[The complete Hodge table for every admissible parameter]\hypertarget{tag-01RY}{}\label{inv:A3.1}\label{wb:cor:exact-hodge-diamond}\label{wb:thm:parametric-hodge-diamond}\label{wb:prop:workbench-holomorphic-form-vanishing}
\StatementRecord{inv:A3.1}{proved}{01RY}{\SeriesRef{inv:eq:A-absolute-table}, \SeriesRef{inv:A1.0}}

For every $\operatorname{Im}c<-M$, the Hodge numbers are
\[
\begin{array}{c|rrrr}
q\backslash p&0&1&2&3\\\hline
3&0&0&0&1\\
2&0&1&2&1\\
1&1&2&1&0\\
0&1&0&0&0.
\end{array}
\]
\end{proposition}



The vanishing of holomorphic forms and the Hodge-diamond
pattern with one undetermined integer are due to
\citet[Proposition~9.20, p.~80; Theorem~9.1(vi), p.~65]{AC2026}.
The nonzero direct-image scalar in \cref{inv:A2.4} determines that integer.

\begin{proof}
\begin{enumerate}[label=\arabic*.,leftmargin=*]
\item Since $B$ is a curve, coherent Leray has no differential and gives
for $F=\Omega_X^p$ the exact sequences
$0\to H^1(B,R^{q-1}f_*F)\to H^q(X,F)\to H^0(B,R^qf_*F)\to0$.
Substitution of \eqref{inv:eq:A-absolute-table}, with
$h^0(\OO)=h^1(\OO(-2))=1$ and all cohomology of $\OO(-1)$ zero,
gives columns $p=1,2$ directly.  \Cref{inv:A1.0} gives column zero,
and Serre duality gives column three. Replacing $\beta$ by $\beta+c$ leaves
every derivative, finite eigencharacter, invariant lattice and
logarithmic residue used above unchanged.  The toric gluing units
change, but their relative logarithmic derivatives still differ only
by base logarithmic forms.  Hence the same computations apply at
each $c$ in the stated open parameter domain.
\end{enumerate}
\end{proof}

\begin{proposition}[Associated-graded generators and a Dolbeault lift]\hypertarget{tag-01RZ}{}\label{inv:A3.2}\label{wb:lem:bcpr-eta-extension}
\StatementRecord{inv:A3.2}{proved}{01RZ}{\SeriesRef{inv:A2.4}, \SeriesRef{inv:eq:A-E-line}, \SeriesRef{inv:A1.3}}

The Leray associated-graded spaces for the four middle groups are
\[
\begin{array}{c|cc}
&E_2^{1,q-1}&E_2^{0,q}\\\hline
H^{1,1}&H^1(\OO(-2))=\C\alpha&\C\eta_{\rm rel}\\
H^{1,2}&0&\C\upsilon\\
H^{2,1}&\C\upsilon^\vee&0\\
H^{2,2}&\C\eta_{\rm rel}^\vee&\C\alpha^\vee.
\end{array}
\]
Here $\alpha$ is the \v Cech class of $dt/t$ on the two standard affine
base charts, $\eta_{\rm rel}=-q_E$ is the invariant relative class
from \ref{inv:A2.4}, and $\upsilon$ is the image of the trace-dual of
$1\in f_*\OO(H)$ under $R^2f_*N\to R^2f_*\mathcal E$ and the
isomorphism $R^2f_*\Omega_X^1=R^2f_*\mathcal E$.
The superscript $\vee$ denotes the bases dual under the relative
trace and the base Serre pairing.  There exists a Dolbeault class
$Q\in H^1(X,\Omega_X^1)$ whose restriction to every regular fibre
is $\eta=u\wedge w+6\gamma\wedge\delta$.
\end{proposition}


\begin{proof}
\begin{enumerate}[label=\arabic*.,leftmargin=*]
\item The isomorphism $f_*\Omega_X^1=\Omega_B^1$ in the residue sequence
is the pullback map, so $\alpha$ gives its stated Leray
class; it can be represented by a nonzero normalization of
$f^*\omega_{\rm FS}$ in Dolbeault cohomology.  The section
$\eta_{\rm rel}$ belongs to the canonical kernel identification
$R^1f_*\Omega_X^1=\ker\kappa_1$ and spans its global sections.
Leray is surjective onto this space, so any preimage defines $Q$
with the asserted fibre restriction.  Its preimages form the affine space $Q+\C\alpha$.
The injection $R^2f_*N=\OO\to R^2f_*\mathcal E$ from
\eqref{inv:eq:A-E-line} spans its global sections, since the quotient
is $\OO(-1)$.  This specifies $\upsilon$ through the displayed injection.  Duality used in \ref{inv:A1.3} gives the remaining entries
and their specified dual pairings.  Choosing a lift of
$\eta_{\rm rel}$ and of $\alpha^\vee$ gives bases of the whole
middle cohomology groups.  The named trace pairings determine the associated-graded bases;
the choices of preimages determine the corresponding Leray splittings.
\end{enumerate}
\end{proof}

\subsection{Ordinary relative differentials and the correction at the cusp}

\begin{proposition}[Relative differentials in the three local models]\hypertarget{tag-01S0}{}\label{inv:A4.1}
\StatementRecord{inv:A4.1}{proved}{01S0}{the definitions and constructions stated in the proof}

Let $R=\C\{x,y,z\}$.  For $t=xy$,
\[
\Omega_{R/\C\{t\}}^1=R^3/R(y,x,0)
\simeq(x,y)\oplus R,
\]
and
\[
\Omega_{R/\C\{t\}}^2\simeq
(x,y)\oplus R/(x,y),\qquad
\Omega_{R/\C\{t\}}^3\simeq R/(x,y).
\]
For $t=xyz$,
$\Omega^1=R^3/R(yz,xz,xy)$ is torsion-free of rank two and is
not free at the triple point.  If $I=(xy,xz,yz)$, there is an exact
sequence
\[
0\to T\to\Omega^2\xrightarrow{\wedge dt} I\to0,
\qquad
T=\frac{R u\oplus R v}{(xu,\ y(u-v),\ zv)},
\quad IT=0,
\]
and $\Omega^3=R/I$.  Finally, for $t=s^m$ in
$R=\C\{s,y_1,y_2\}$,
\[
\begin{aligned}
\Omega^1&=R^2\oplus R/(s^{m-1})\,ds,\\
\Omega^2&=R\,dy_1\wedge dy_2
 \oplus\bigl(R/(s^{m-1})\bigr)^2,\\
\Omega^3&=R/(s^{m-1})\,ds\wedge dy_1\wedge dy_2.
\end{aligned}
\]
\end{proposition}


\begin{proof}
\begin{enumerate}[label=\arabic*.,leftmargin=*]
\item The degree-one presentation is the quotient by $dt$; higher degrees
are its exterior powers.  At a double point the transverse quotient
$R^2/R(y,x)$ maps isomorphically to $(x,y)$ by
$(a,b)\mapsto xa-yb$. Its kernel is exactly $R(y,x)$:
the equality $xa=yb$ gives $a=yh$ and $b=xh$ by relative primality.  In degree two,
with basis $A=dx\wedge dy,B=dx\wedge dz,C=dy\wedge dz$, the
relations are $xA=yA=0$, $yB+xC=0$.  This gives the displayed
torsion summand and the ideal $(x,y)$; wedging once more gives
$R/(x,y)$.

\item At the triple point the map from ordinary to logarithmic relative
one-forms sends $(a,b,c)$ to $(xa-zc,yb-zc)$.
Its kernel is precisely $R(yz,xz,xy)$: the two equations imply
$c=xyh$, $a=yzh$, $b=xzh$ by relative primality of the coordinates.
Thus it injects into $R^2$, proving torsion-freeness.  Its fibre at
the maximal ideal has dimension three although its rank is two,
so it is not free.  The same minimal-generator argument works at
the double locus.

\item In degree two, wedging with $dt$ is the row
$(xy,-xz,yz)$ in the basis $(A,B,C)$.
Its syzygies are generated freely by
$u=(z,y,0)$ and $v=(0,y,x)$.
To check generation, an equation $xyA-xzB+yzC=0$ implies
$x\mid C$, say $C=xC_1$; then $yA=z(B-yC_1)$ implies
$A=zA_1$, $B=y(A_1+C_1)$.  Thus the syzygy is $A_1u+C_1v$.
The three relations from $dt\wedge(dx,dy,dz)$ are
$-xu,y(u-v),zv$, proving the stated presentation of the kernel.
The ideal $I$ kills it by those relations.  Conversely the quotient
$I$ is torsion-free, so this kernel is exactly the torsion of
$\Omega^2$.  The presentation immediately gives $\Omega^3=R/I$.
For $t=s^m$ the sole degree-one
relation is $ms^{m-1}ds=0$; $m$ is a unit and taking exterior
powers gives the three formulas.
\end{enumerate}
\end{proof}

\begin{proposition}[The support of relative torsion]\hypertarget{tag-01S1}{}\label{inv:A4.2}
\StatementRecord{inv:A4.2}{proved}{01S1}{\SeriesRef{inv:A4.1}}

The torsion of $\Omega_{X/B}^1$ is supported precisely on
$S_1\cup S_2$. Its cusp modules inject into the free logarithmic
relative modules, and their minimal number of generators exceeds
their rank at the double locus. In degree two, the torsion also
has support on the cusp double locus.
This corrects \cite[Remark~9.21, p.~81]{AC2026}, where
degree-one torsion is asserted on $W\cup S_1\cup S_2$;
the support statement in \cite[Remark~1.3, p.~5]{AC2026}
therefore depends on the degree.
\end{proposition}


\begin{proof}
\begin{enumerate}[label=\arabic*.,leftmargin=*]
\item The double and triple charts in \cref{inv:A4.1} exhaust the
singular strata of $W$. Their injections into free logarithmic
modules prove absence of degree-one cusp torsion. Their degree-two
presentations have a nonzero torsion generator at every general
point of a double curve. At a multiple fibre the finite \'etale
comparison
\[
q^*\Omega^1_{N_j/B}=\Omega^1_{J'_j/B}
\]
is an isomorphism; its summand $R/(s^{m_j-1})\,ds$ is nonzero
for $m_j=3,4$. Faithful flatness descends both the torsion and its
support. There are no other critical fibres, and ordinary relative
differentials are locally free on their complement. Thus the
degree-one support is exactly $S_1\cup S_2$, whereas degree two
also has cusp torsion, as asserted.
\end{enumerate}
\end{proof}



\section{Fr\"olicher, Bott--Chern, Aeppli, and Hermitian metrics}\label{pub:I:hodge}
\hypertarget{tag-01TI}{}\label{inv:sec:partC}


\subsection{The differentials and their dependence on the Hodge number}

\begin{definition}[Cohomological conventions]
\hypertarget{tag-PI13}{}\label{I:def:hodge-conventions}
\StatementRecord{I:def:hodge-conventions}{proved}{PI13}{}
For the standard marked threefold, put $h=h^{1,2}(X)$.
Proposition~\ref{inv:A3.1} gives $h=1$. We write the dimensions
below as functions of $h$ in the common Hodge pattern
\begin{equation}\hypertarget{tag-01TF}{}\label{inv:eq:C-hodge-input}
\begin{array}{c|rrrr}
q\backslash p&0&1&2&3\\\hline
3&0&0&0&1\\
2&0&h&h+1&1\\
1&1&h+1&h&0\\
0&1&0&0&0.
\end{array}
\end{equation}
The numerical tables of this family are obtained by the evaluation
map $\mathbb Z[h]\to\mathbb Z$, $h\mapsto1$.
We use $d=\partial+\bar\partial$ on complex-valued smooth forms
and the complex orientation for integration.
\end{definition}

\begin{proposition}[The complete Fr\"olicher differential in terms of $h$]\hypertarget{tag-01SA}{}\label{inv:C1.1}\label{wb:thm:frolicher-e2-candidate}
\StatementRecord{inv:C1.1}{proved}{01SA}{\SeriesRef{inv:eq:C-hodge-input}}

For a compact complex threefold with \eqref{inv:eq:C-hodge-input} and the
complex de Rham cohomology of $S^6$, the Fr\"olicher spectral sequence
satisfies $E_1\ne E_2=E_\infty$.  Its only nonzero $d_1$ ranks are
\[
\rank d_1^{0,1}=1,\quad \rank d_1^{1,1}=h,\quad
\rank d_1^{1,2}=h,\quad \rank d_1^{2,2}=1.
\]
In particular its two middle rows are exact.
\end{proposition}


\begin{proof}
\begin{enumerate}[label=\arabic*.,leftmargin=*]
\item Here $E_1^{p,q}=H^q(X,\Omega_X^p)$ and
$d_r:E_r^{p,q}\longrightarrow E_r^{p+r,q-r+1}$ are complex-linear.
For $r=2$, any source in the two middle rows has target in one of
$(2,0),(3,0),(3,1)$ or outside the table, except the source $(0,2)$,
which is zero.  Sources in row $q=3$ outside $(3,3)$ are zero, and
$(3,3)$ has no target.  For $r=3$ the only possible targets inside the
table have sources $(0,2)$ or $(0,3)$, both zero.  For $r>3$ there is
no arrow inside the table.  A zero term remains zero on subsequent
pages.  Hence $E_2=E_\infty$.

\item The filtration on $H^k_{\mathrm{dR}}(X;\C)$ has no nonzero associated
graded term for $1\le k\le5$.  Every middle term of $E_2$ therefore
vanishes.  Consequently
\[
0\longrightarrow\C\longrightarrow\C^{h+1}
\longrightarrow\C^h\longrightarrow0,
\qquad
0\longrightarrow\C^h\longrightarrow\C^{h+1}
\longrightarrow\C\longrightarrow0
\]
are exact, giving the stated ranks.  The first rank is one, so $E_1$
cannot equal $E_2$.
\end{enumerate}
\end{proof}

\begin{proposition}[A basis on which all first differentials are specified]\hypertarget{tag-01SB}{}\label{inv:C1.2}
\StatementRecord{inv:C1.2}{proved}{01SB}{\SeriesRef{inv:A1.3}, \SeriesRef{inv:A3.2}, \SeriesRef{inv:C1.1}, \SeriesRef{inv:C3.7}}

Choose a nonzero $g\in H^{0,1}_{\bar\partial}(X)$, a basis
$q_1,\ldots,q_h$ of a complement to $\C d_1g$ in
$H^{1,1}_{\bar\partial}(X)$, a basis $v_1,\ldots,v_h$ of
$H^{1,2}_{\bar\partial}(X)$, and a class $e$ complementary to the
image of $d_1:H^{1,2}\to H^{2,2}$.  With
$a=d_1g$, $c_j=d_1q_j$, $d_j=d_1v_j$, and $t=d_1e$, the first
differential is
\[
g\mapsto a\mapsto0,\quad q_j\mapsto c_j\mapsto0,
\qquad v_j\mapsto d_j\mapsto0,\quad e\mapsto t\mapsto0.
\]
These are bases in the corresponding bidegrees.  If $h=1$ and a
class $q$ restricts to the invariant fibre class $\eta$ described below,
we may take $q_1=q$, $v_1=gq$, and $e=q^2$, after scaling $g$ if
necessary.  Propositions~\ref{inv:A1.3} and \ref{inv:A3.2} supply these data
for every member of the marked family.
\end{proposition}


\begin{proof}
\begin{enumerate}[label=\arabic*.,leftmargin=*]
\item The first assertions follow from the injectivity and surjectivity
proved in \ref{inv:C1.1}.  For the last assertion write
$\eta=u\wedge w+6\gamma\wedge\delta$ in the marked integral
cohomology of a regular fibre.  Its complex-oriented square has
integral $-12$.  Its product with the $(0,1)$ image of $\gamma$ is
nonzero: the explicit period calculation in \ref{inv:C3.7} gives its nonzero
pairing with $d\zeta_2$.  Thus $gq$ and $q^2$ have nonzero fibre
restrictions.  The images under $d_1$ restrict to zero on the fibre,
because the fibre is a complex torus and its invariant representatives
are $d$-closed.  Therefore $q$ complements $d_1g$, $gq$ spans the
one-dimensional $H^{1,2}$, and $q^2$ complements $d_1H^{1,2}$.
The lift in Proposition~\ref{inv:A3.2} maps to $\eta$ under the
Leray restriction, and hence supplies the class $q$ used here.
\end{enumerate}
\end{proof}

\subsection{A computation from the two definitions}

\begin{definition}[Bott--Chern and Aeppli cohomology]
\hypertarget{tag-PI14}{}\label{I:def:BC-Aeppli}
\StatementRecord{I:def:BC-Aeppli}{proved}{PI14}{}

On the double complex $A^{p,q}(X)$ of smooth forms, set
\[
H^{p,q}_{BC}(X)=
\frac{\ker\partial\cap\ker\bar\partial\subset A^{p,q}(X)}
{\operatorname{im}(\partial\bar\partial:A^{p-1,q-1}\to A^{p,q})},
\quad
H^{p,q}_{A}(X)=
\frac{\ker(\partial\bar\partial:A^{p,q}\to A^{p+1,q+1})}
{\partial A^{p-1,q}+\bar\partial A^{p,q-1}}.
\]
Complex conjugation sends bidegree $(p,q)$ to $(q,p)$.
\end{definition}

\begin{proposition}[Bott--Chern and Aeppli tables as functions of $h$]\hypertarget{tag-01SC}{}\label{inv:C2.1}\label{wb:thm:exact-bott-chern-aeppli}
\StatementRecord{inv:C2.1}{proved}{01SC}{\SeriesRef{inv:A3.1}, \SeriesRef{inv:C1.1}, \SeriesRef{inv:eq:C-arrows}; \citet[Theorem 3]{Stelzig2021}; \citet[Proposition 6(1)]{Stelzig2021}}

The complete dimension tables are
\begin{equation}\hypertarget{tag-01T6}{}\label{inv:eq:C-BC-A}
\begin{array}{c|rrrr}
q\backslash p&0&1&2&3\\\hline
3&0&0&1&1\\
2&0&h&2h&1\\
1&0&2&h&0\\
0&1&0&0&0
\end{array}_{BC}
\qquad
\begin{array}{c|rrrr}
q\backslash p&0&1&2&3\\\hline
3&0&0&0&1\\
2&0&h&2&0\\
1&1&2h&h&0\\
0&1&1&0&0.
\end{array}_{A}
\end{equation}
At the value $h=1$ of Proposition~\ref{inv:A3.1}, the totals are
$(1,0,2,2,2,2,1)$ and $(1,2,2,2,2,0,1)$, respectively.
\end{proposition}


\begin{proof}
\begin{enumerate}[label=\arabic*.,leftmargin=*]
\item Apply \cite[Theorem~3, p.~7; Proposition~6, p.~16;
Lemma~8, p.~17]{Stelzig2021} to the bounded-bidegree complex of
smooth forms. The theorem permits arbitrary cardinal dimensions.
Its odd zigzags contribute to total cohomology; the sphere Betti
groups and the single extreme bidegrees therefore leave only dots
at $(0,0)$ and $(3,3)$. By \cref{inv:C1.1} and conjugation,
both spectral sequences have no differential after $d_1$, so every
remaining zigzag is an arrow. The horizontal arrows are exactly
\begin{equation}\hypertarget{tag-01TA}{}\label{inv:eq:C-arrows}
\begin{array}{c|c|r}
\text{source}&\text{target}&\text{multiplicity}\\\hline
(0,1)&(1,1)&1\\
(1,1)&(2,1)&h\\
(1,2)&(2,2)&h\\
(2,2)&(3,2)&1.
\end{array}
\end{equation}
The vertical arrows are their reflections in $p=q$.

\item Under the cited decomposition, an arrow contributes its target
to Bott--Chern cohomology and its source to Aeppli cohomology;
a dot contributes to both and a square to neither. The induced
isomorphisms in Lemma~8 respect Definition~\ref{I:def:BC-Aeppli}.
Counting the displayed arrows and their reflections gives every
entry of \eqref{inv:eq:C-BC-A}. Kernels, images and these quotients
commute with direct sums, including the possible infinite sum of
squares.

\end{enumerate}
\end{proof}

\begin{proposition}[Comparison with McHugh and Angella]\hypertarget{tag-01SD}{}\label{inv:C2.2}
\StatementRecord{inv:C2.2}{proved}{01SD}{\SeriesRef{inv:eq:C-hodge-input}, \SeriesRef{inv:C2.1}, \SeriesRef{inv:C1.1}, \SeriesRef{inv:eq:C-BC-A}; \citet[arXiv v4, Theorem~6.3, pp.~20--23; Proposition~6.4 and (7), p.~23; \S7.1, pp.~24--25]{McHugh2018}; \citet[arXiv v2, \S4, pp.~13--15; Figure~5, p.~16]{Angella2018}}

The specialization of McHugh's Theorem 6.3, Proposition 6.4, equation
(7), and Section 7.1 to \eqref{inv:eq:C-hodge-input} agrees with \ref{inv:C2.1}.
The corresponding relations and Figure 5 in Angella's Section 4 also
agree through the displayed specialization formulas.
\end{proposition}


\begin{proof}
\begin{enumerate}[label=\arabic*.,leftmargin=*]
\item McHugh's quotient $K^{p,q}$ is zero in intermediate total degrees
on every arrow and square just used: a closed corner of an arrow is
the image of the incoming differential and a closed corner of a square
is in both images.  Thus $k^{p,q}=0$ for $0<p+q<6$.
Also $E_2^{3,1}=0$ by \ref{inv:C1.1} and $H^{1,1}(X;\R)=0$ because
$H^2(X;\R)=0$.  Reading the formulas in
\citet[arXiv v4, Theorem~6.3, pp.~20--23; Proposition~6.4 and (7), p.~23; \S7.1, pp.~24--25]{McHugh2018}
therefore gives
\[
h_{BC}^{1,1}=2h^{0,1}=2,\qquad
h_{BC}^{1,2}=h^{1,2}=h,\qquad
h_{BC}^{2,2}=-1+(h+1)+h=2h.
\]
The remaining entries in his table are precisely the boundary entries
of \eqref{inv:eq:C-BC-A}.  In
\citet[arXiv v2, \S4, pp.~13--15; Figure~5, p.~16]{Angella2018}, the relation
\[
h_{BC}^{2,2}=2h_{BC}^{2,1}-2h^{0,1}+2
\]
gives $2h$, and the boundary values reduce to those above because
$h^{2,0}=h^{0,2}=0$.  
\end{enumerate}
\end{proof}

\subsection{The algebra, line bundles, and metric consequences}

\begin{proposition}[The real Bott--Chern Chern-class map and a global $G$]\hypertarget{tag-01SE}{}\label{inv:C3.1}\label{wb:cor:picard-line-audit}
\StatementRecord{inv:C3.1}{proved}{01SE}{\SeriesRef{inv:C1.1}, \SeriesRef{inv:C2.1}, \SeriesRef{inv:A1.0}; \cite[Theorem~10.5(b), proof Steps~1 and~4, p.~88]{AC2026}}

The Picard-group identification and fibrewise injectivity below are due to Claude and Alp\"oge \cite[Theorem~10.5(b), proof Steps~1 and~4, p.~88]{AC2026}. Let $G$ be a $\bar\partial$-closed smooth $(0,1)$ representative of
a nonzero $g\in H^1(X,\OO_X)$.  Set
\[
\mathsf A=[\partial G]_{BC},\qquad
\mathsf B=[\bar\partial\bar G]_{BC}.
\]
These classes form a complex basis of $H^{1,1}_{BC}(X)$.
With the exponential $z\mapsto e^{2\pi iz}$ and the Chern convention
$c_{1,BC}=[iF/(2\pi)]$, the map of real Lie groups is
\begin{equation}\hypertarget{tag-01T8}{}\label{inv:eq:C-Pic-map}
\Pic(X)=\{\exp(zg):z\in\C\}\longrightarrow H^{1,1}_{BC}(X),
\qquad \exp(zg)\longmapsto z\mathsf A+\bar z\mathsf B.
\end{equation}
It is injective and its image is the real subspace fixed by conjugation.

There is an operator-explicit global representative.  Normalize the
Fubini--Study form by $\int_{\Pone}\omega_{FS}=1$, put
$\theta=f^*\omega_{FS}$, choose any smooth Riemannian metric on $X$,
and let $\mathcal G_d$ be its Green operator for the real de Rham
Laplacian.  Then
\begin{equation}\hypertarget{tag-01T7}{}\label{inv:eq:C-G-Green}
\alpha=d^*\mathcal G_d\theta,\qquad G_0=\alpha^{0,1}
\end{equation}
satisfy $d\alpha=\theta$, $\bar\partial G_0=0$,
$[G_0]\ne0$, and $\exp([G_0])=f^*\OO_{\Pone}(1)$.
Restriction $H^1(X,\mathcal O_X)\to H^1(X_b,\mathcal O_{X_b})$ is injective for every scheme-theoretic fibre; on a smooth fibre its image is $\mathbb C\gamma^{0,1}$.
\end{proposition}

\begin{proof}
\begin{enumerate}[label=\arabic*.,leftmargin=*]
\item By \ref{inv:C1.1}, the Dolbeault image of $\mathsf A$ is
$d_1g\ne0$. The class $\mathsf B=\overline{\mathsf A}$ is nonzero,
whereas its Dolbeault image is zero. Thus these two classes are
independent; the dimension calculation of \ref{inv:C2.1} makes
them a basis. In this basis conjugation sends $(a,b)$ to
$(\bar b,\bar a)$, so its fixed real subspace consists precisely
of $(z,\bar z)$.

\item The exponential sequence and $H^1(X;\Z)=H^2(X;\Z)=0$ give the
cited identification $\Pic(X)=H^1(X,\OO_X)=\C g$.
On a smooth trivial line with unit frame $s$, put
$\bar\partial_zs=-2\pi izG\otimes s$. Integrability follows from
$\bar\partial G=0$. On a cover with $G=\bar\partial\phi_i$,
the holomorphic frames $e_i=e^{2\pi iz\phi_i}s$ satisfy
$e_j=e^{2\pi iz(\phi_j-\phi_i)}e_i$; hence this is exactly the
line represented by the exponential of the \v Cech class $zg$.
Its unitary Chern connection and curvature are
\[
 \nabla_z=d-2\pi i(zG+\bar z\bar G),\qquad
 F_z=-2\pi i(z\partial G+\bar z\bar\partial\bar G).
\]
Multiplication by $i/(2\pi)$ proves \eqref{inv:eq:C-Pic-map}.
The two independent basis vectors prove injectivity, additivity
as real Lie groups and the assertion about the image.

\item The closed real form $\theta$ is exact since $H^2(X;\R)=0$.
It has zero harmonic projection. De Rham Hodge decomposition,
with $d\mathcal G_d=\mathcal G_dd$, therefore gives
$dd^*\mathcal G_d\theta=\theta$. The Green operator and $d^*$
preserve reality, so $\alpha$ is real and $d\alpha=\theta$.
The $(0,2)$ component yields $\bar\partial G_0=0$.
To prove nonvanishing, an equation $G_0=\bar\partial u$ would give
\[
 \theta=\partial\bar\partial(u-\bar u)
       =2i\partial\bar\partial\operatorname{Im}u.
\]
Semipositivity of $\theta$ would make $\operatorname{Im}u$
plurisubharmonic. The maximum principle on compact connected $X$
would make it constant, contradicting the nonzero pullback of
$\omega_{FS}$ on the regular locus of $f$.

\item Since $\alpha=G_0+\bar G_0$, the Chern formula already proved
sends $\exp([G_0])$ to $[\theta]_{BC}$, which is also
$c_{1,BC}(f^*\OO(1))$. Its injectivity proves the stated
line-bundle equality, with exactly the declared $2\pi i$
exponential and curvature conventions.

\item For the last assertion, the exact sequence of a full scheme
fibre is
\[
 0\longrightarrow f^*\OO(-1)\longrightarrow\OO_X
   \longrightarrow\OO_{X_b}\longrightarrow0.
\]
The cited direct images in \ref{inv:A1.0} give
$f_*f^*\OO(-1)=\OO(-1)$ and
$R^1f_*f^*\OO(-1)=\OO(-1)\oplus\OO(-2)$.
Both terms bounding $H^1(X,f^*\OO(-1))$ in the degree-one Leray
edge sequence vanish, so this group vanishes. Restriction on
$H^1(\OO)$ is injective. On a smooth fibre, the same edge map is
the evaluation of $H^0(B,G)=\C\gamma^{0,1}$ in $G_b$.
This identifies its image as claimed and reproduces exactly
the application of \cite[Theorem~10.5(b), Step~4, p.~88]{AC2026}.
\end{enumerate}
\end{proof}

\Uses{the hypotheses, notation and definitions stated here; \citet[Theorem~10.5]{AC2026}; \CompendiumRef{I}{app:rederive:inv-C3.1}}





\begin{proposition}[The factor $12$ with a fixed collar orientation]\hypertarget{tag-01SF}{}\label{inv:C3.2}\label{wb:lem:bcpr-degree-one-input}
\StatementRecord{inv:C3.2}{proved}{01SF}{\SeriesRef{inv:eq:C-collar-sum}, \SeriesRef{inv:C3.1}}

Use the marked cohomology class $\gamma$ dual to
$\hat\gamma$, the analytic peripheral generators and the lift relation of
\cite[\S7.5, pp.~54--57, especially (7.10), p.~55]{AC2026}, and put
$p=12l_0-4l_1-3l_2=\pm1$.  Let $g_\gamma$ be the unique coherent
Leray class whose fibre image is $\gamma^{0,1}$.  Then
\begin{equation}\hypertarget{tag-01TE}{}\label{inv:eq:C-factor12}
\exp(12g_\gamma)=f^*\OO_{\Pone}(-p).
\end{equation}
For the standard data $(l_0,l_1,l_2)=(0,1,-1)$, this is
$\exp(12g_\gamma)=f^*\OO(1)$.  For the opposite value $p=1$,
the exponent representing $f^*\OO(1)$ is $-12g_\gamma$.
\end{proposition}


\begin{proof}
\begin{enumerate}[label=\arabic*.,leftmargin=*]
\item A degree-one integral
class $k\gamma$ on a regular fibre extends over the cyclic filling of
order $m_j$ precisely when the affine generator $\hat g_j$ can be
assigned an integer value $n_j$ satisfying
$m_j n_j=k\gamma(v_j)=kl_j$.  Admissibility gives
$\gcd(l_j,m_j)=1$, so this is precisely $m_j\mid k$.
At the cusp $\gamma$ extends, as follows from its specialization
to $H^1(W;\Z)$ in \cite[Theorem~B.1, pp.~101--105; Proposition~7.12, pp.~51--53]{AC2026}.
The global monodromy invariants have rank one generated by $\gamma$.
It follows, by intersection of these three local images in the marked
local system, that
$H^0(B,R^1f_*\Z)=12\Z\gamma$.  This is the degree-one calculation of
\cite[Proposition~7.26, pp.~61--62]{AC2026}.
The closed embedding $F_{\mathrm{red}}\hookrightarrow F$ is a
homeomorphism on underlying spaces and therefore induces the
isomorphism on integral topological cohomology used here.
The coherent restriction map has target $H^1(F,\mathcal O_F)$.

\item Let $U=\bigsqcup_iD_i$ be small discs around the special points and
let $V$ be their pair-of-pants complement with collars.  Lift
$12\gamma$ to integral classes $\Theta_i$ on $f^{-1}(D_i)$ and
$\Theta_V$ on $f^{-1}(V)$.  Existence over $V$ follows from
$H^2(V;\Z)=0$ and the low-degree Leray sequence; over each disc it
follows from the local images just computed.  Write
$a_j=\Theta_V(\hat\rho_j)$.  On the two finite analytic peripheral generators
the differences $\Theta_i-\Theta_V$ have periods
\[
n_1=4l_1-a_1,\qquad n_2=3l_2-a_2.
\]
At the cusp, the specified lift relation
$\widetilde a_0=t_\mu\hat\rho_2^{-1}\hat\rho_1^{-1}$,
with $\gamma(\mu)=l_0$, and the null-homotopy of
$\widetilde a_0$ in the cusp filling give
$n_0=-12l_0+a_1+a_2$.  Hence
\begin{equation}\hypertarget{tag-01TB}{}\label{inv:eq:C-collar-sum}
n_0+n_1+n_2=-p.
\end{equation}
These boundary evaluations use the orientations in
\cite[Proposition~7.27, pp.~62--63]{AC2026}.

\item Represent the lifts by real closed one-forms.  On a collar their
difference is the pullback of a real closed base form $b_i$ with
clockwise period $n_i$, modulo an exact form.  This follows from the
degree-one Leray sequence over the annulus, since the difference has
zero fibre restriction and the annulus has no second cohomology.
Subtract exact terms on smaller collars, extending their primitives
with cutoffs, so the equality of forms holds wherever a partition of
unity changes.  For a real base cutoff $\rho_i$ equal to one toward
$D_i$ and zero toward $V$, the resulting glued real one-form $\alpha$
on $X$ has
\[
d\alpha=f^*\vartheta,\qquad
\vartheta=\sum_i d\rho_i\wedge b_i.
\]
The form $\vartheta$ is a smooth real base two-form; thus
$\bar\partial\alpha^{0,1}=0$.  Its fibre restriction is the
$(0,1)$ image of $12\gamma$, so the coherent Leray edge isomorphism
$H^1(X,\OO_X)\simeq H^0(B,R^1f_*\OO_X)$ identifies
$[\alpha^{0,1}]=12g_\gamma$.

\item In a positively oriented polar coordinate $(r,\varphi)$ on $D_i$,
$b_i=-n_i\,d\varphi/(2\pi)$ modulo an exact form.
Because $\rho_i$ decreases from one to zero as $r$ increases,
\[
\int_{D_i\cap V}d\rho_i\wedge b_i
=-n_i\int_0^1d\rho_i=n_i.
\]
The exact part does not change this integral, since $d\rho_i$ is
supported inside the collar.  Therefore
$\int_B\vartheta=-p$ by \eqref{inv:eq:C-collar-sum}.  \ref{inv:C3.1} identifies
the Bott--Chern class of $\exp(12g_\gamma)$ with $[f^*\vartheta]$.
On $\Pone$ two real $(1,1)$ forms of the same integral differ by
$i\partial\bar\partial$ of a real function (solve the ordinary
scalar Laplace equation).  Hence this class is the Chern class of
$f^*\OO(-p)$, and injectivity in \ref{inv:C3.1} proves
\eqref{inv:eq:C-factor12}.
\end{enumerate}
\end{proof}

\begin{definition}[The nine positive-degree Bott--Chern classes]
\hypertarget{tag-PI15}{}\label{I:def:BC-generators}
\StatementRecord{I:def:BC-generators}{proved}{PI15}{}

For $h=1$, choose $\bar\partial$-closed representatives $G,Q$
of $g,q$ as in Proposition~\ref{inv:C1.2}, and choose a smooth
top-degree form $\nu$ with $\int_X\nu=1$. Define
\begin{equation}
\begin{array}{c|c|c}
\text{class}&\text{representative}&\text{bidegree}\\\hline
\mathsf A&\partial G&(1,1)\\
\mathsf B&\bar\partial\bar G&(1,1)\\
\mathsf C&\partial Q&(2,1)\\
\overline{\mathsf C}&\bar\partial\bar Q&(1,2)\\
\mathsf D&\partial(G\wedge Q)&(2,2)\\
\overline{\mathsf D}&\bar\partial(\bar G\wedge\bar Q)&(2,2)\\
\mathsf E&\partial(Q^2)&(3,2)\\
\overline{\mathsf E}&\bar\partial(\bar Q^2)&(2,3)\\
\mathsf T&\nu&(3,3).
\end{array}
\label{I:eq:BC-generators}
\end{equation}
\end{definition}

\begin{proposition}[The complete Bott--Chern algebra, for every $h$]\hypertarget{tag-01SG}{}\label{inv:C3.3}\label{wb:thm:bcpr-square-zero-algebra}\label{wb:prop:bcpr-generators-are-bases}
\StatementRecord{inv:C3.3}{proved}{01SG}{\SeriesRef{inv:C2.1}, \SeriesRef{inv:C1.2}, \SeriesRef{inv:C3.1}, \SeriesRef{inv:C1.1}}

The positive-total-degree ideal $N\subset H^{\bullet,\bullet}_{BC}(X)$
satisfies
\[
H^{\bullet,\bullet}_{BC}(X)=\C\cdot1\oplus N,\qquad N^2=0.
\]
Its nonzero bidegrees and dimensions are those of \ref{inv:C2.1}.  Thus the
square-zero conclusion is independent of the evaluated value of $h$.
For $h=1$, the nine classes specified in \eqref{I:eq:BC-generators}, together with the unit, give bases of every nonzero Bott--Chern group.
\end{proposition}


\begin{proof}
\begin{enumerate}[label=\arabic*.,leftmargin=*]
\item Take global $\bar\partial$-closed representatives $G,Q_j,V_j,E$
of the classes $g,q_j,v_j,e$ of \ref{inv:C1.2}.  By the arrow computation in
\ref{inv:C2.1}, bases for the positive Bott--Chern groups consist of
\[
\partial G,\ \bar\partial\bar G;\qquad
\partial Q_j,\ \bar\partial\bar Q_j;\qquad
\partial V_j,\ \bar\partial\bar V_j;\qquad
\partial E,\ \bar\partial\bar E;\qquad \nu,
\]
where the pairs have bidegrees $(1,1)$, $(2,1),(1,2)$,
$(2,2)$, $(3,2),(2,3)$, respectively, and
$\int_X\nu=1$.  In particular the natural map
$H^{3,2}_{BC}\to H^{3,2}_{\bar\partial}$ is an isomorphism:
on its unique horizontal arrow the closed target maps nontrivially
to column cohomology.

\item Write $a=\partial G$, $b=\bar\partial\bar G$ and
$c_j=\partial Q_j$.  The graded Leibniz rule gives
\[
\partial\bar\partial(G\wedge\bar G)=-a\wedge b,
\quad
\partial\bar\partial(G\wedge\bar Q_j)
=-a\wedge\bar c_j,
\quad
\partial\bar\partial(\bar G\wedge Q_j)=b\wedge c_j.
\]
Consequently $\mathsf A\mathsf B=0$,
$\mathsf A\overline{\mathsf C_j}=0$, and
$\mathsf B\mathsf C_j=0$.

\item Let $f^*\OO(1)=\exp(\lambda g)$, where $\lambda\ne0$.
Its curvature may be chosen to be the pullback of a two-form on $B$,
whose square vanishes pointwise.  By \ref{inv:C3.1},
$\lambda^2\mathsf A^2+\bar\lambda^2\mathsf B^2=0$.
The Dolbeault image of $\mathsf B^2$ is zero.  Hence
$[a^2]_{\bar\partial}=0$.  But
$a^2=\partial(G\wedge a)$, and \ref{inv:C1.1} says that
$d_1:H^{1,2}\to H^{2,2}$ is injective.  Thus
$[G\wedge a]_{\bar\partial}=0$.  Choose a smooth $(1,1)$ form
$\Phi$ satisfying $G\wedge a=\bar\partial\Phi$.  Applying
$\partial$ proves $a^2=\partial\bar\partial\Phi$.
Conjugation gives $\mathsf B^2=0$ as well.

\item To treat $\mathsf A\mathsf C_j$, pair its Dolbeault image with the
generator $g\in H^{0,1}_{\bar\partial}$.  Since
$\bar\partial c_j=0$ and $G\wedge a=\bar\partial\Phi$,
\[
\int_X a\wedge c_j\wedge G
=-\int_X\bar\partial(\Phi\wedge c_j)=0.
\]
Dolbeault--Serre duality and the isomorphism
$H^{3,2}_{BC}\to H^{3,2}_{\bar\partial}$ give
$\mathsf A\mathsf C_j=0$.  Complex conjugation gives
$\mathsf B\overline{\mathsf C_j}=0$.

\item Every remaining product whose bidegree lies in $[0,3]^2$ lands in
$(3,3)$: it is a product of a $(1,1)$ class with a $(2,2)$ class,
or of a $(2,1)$ class with a $(1,2)$ class.  Each non-top generator
displayed above is $d$-exact, since its chosen primitive is
$\bar\partial$-closed, or is the conjugate of such a form.
The products therefore have integral zero.  Integration is an
isomorphism $H^{3,3}_{BC}\to\C$ by \ref{inv:C2.1} and the class of $\nu$,
so all these products vanish.  All other products exceed the complex
dimension in at least one bidegree.  This exhausts the possibilities.
\end{enumerate}
\end{proof}

\begin{proposition}[Non-K\"ahler degrees and comparison maps]\hypertarget{tag-01SH}{}\label{inv:C3.4}
\StatementRecord{inv:C3.4}{proved}{01SH}{\SeriesRef{inv:C2.1}}

For $k=0,\ldots,6$ the non-K\"ahler degrees are
\begin{equation}\hypertarget{tag-01TC}{}\label{inv:eq:C-delta}
(\Delta^k)=(0,2,2+2h,4h,2+2h,2,0).
\end{equation}
If $h\ge1$, the two maps
$\bigoplus_{p+q=k}H^{p,q}_{BC}\to H^k_{\mathrm{dR}}(X;\C)
\to\bigoplus_{p+q=k}H^{p,q}_A$
are both isomorphisms exactly for $k=0,6$.
For the weaker assertion that a $d$-exact pure $d$-closed form of
degree $k$ is $\partial\bar\partial$-exact, the valid degrees are
$k=0,1,6$, and also $k=3,4$ if $h=0$.
\end{proposition}


\begin{proof}
\begin{enumerate}[label=\arabic*.,leftmargin=*]
\item The totals from \ref{inv:C2.1} are
\[
(h^k_{BC})=(1,0,2,2h,2h,2,1),\qquad
(h^k_A)=(1,2,2h,2h,2,0,1).
\]
Subtracting twice the Betti numbers gives \eqref{inv:eq:C-delta}.
On the two endpoint dots both comparison maps are isomorphisms;
all intermediate de Rham groups vanish.  The first assertion follows
by inspecting when both outside groups vanish.  The weaker assertion
is exactly injectivity of the first comparison map.  In intermediate
degree this is equivalent to vanishing of the Bott--Chern group alone,
which gives the stated list.  In degree zero there is no nonzero
exact form, and in degree six integration identifies Bott--Chern and
de Rham cohomology.  In degree one the first comparison map is $0\to0$, while the
second is $0\to\mathbb C^2$; the formula gives $\Delta^1=2$.
\end{enumerate}
\end{proof}

\begin{proposition}[Obstructions from exact positive forms]\hypertarget{tag-01SI}{}\label{inv:C3.5}
\StatementRecord{inv:C3.5}{proved}{01SI}{the definitions and constructions stated in the proof}

The threefold $X$ admits no balanced, strongly Gauduchon, K\"ahler,
locally conformally K\"ahler, or Hermitian-symplectic metric.
\end{proposition}


\begin{proof}
\begin{enumerate}[label=\arabic*.,leftmargin=*]
\item Let $\theta=f^*\omega_{FS}$.  It is a nonzero semipositive exact
real $(1,1)$ form.  If $\omega$ were balanced, the real four-form
$\omega^2$ would be closed and
$\int_X\theta\wedge\omega^2>0$, because the integrand is positive
on the regular locus.  Exactness of $\theta$ and Stokes' theorem
give zero, a contradiction.

\item More generally, if $\omega$ were strongly Gauduchon, write
$\partial\omega^2=\bar\partial u$ with $u\in A^{3,1}$.
The real four-form $\Omega=\omega^2-u-\bar u$ is closed.
Bidegrees give $\theta\wedge\Omega=\theta\wedge\omega^2$,
so the same argument applies.  K\"ahler metrics are balanced.
If $d\omega=\ell\wedge\omega$ with a closed real Lee form
$\ell$, then $b_1=0$ gives $\ell=dv$, and
$d(e^{-v}\omega)=0$, excluding locally conformally K\"ahler metrics.
Finally any symplectic two-form on $X$ is exact because $b_2=0$.
Its cube has nonzero integral for its symplectic orientation, whereas
an exact cube has zero integral.  Thus $X$ is not symplectic and in
particular admits no Hermitian-symplectic metric.
\end{enumerate}
\end{proof}

\begin{cited}[Gauduchon metrics with specified Chern--Ricci form]\hypertarget{tag-01SJ}{}\label{inv:C3.6}
\StatementRecord{inv:C3.6}{proved}{01SJ}{\citet[Theorem~2]{Gauduchon1985}; \citet[Theorem~1.3]{STW2017}}

There exists a Gauduchon metric on $X$.  More precisely, for every
Gauduchon metric $\omega_0$ there is another Gauduchon metric
$\omega$ with $[\omega^2]_A=[\omega_0^2]_A$ and
\begin{equation}\hypertarget{tag-01T9}{}\label{inv:eq:C-Ricci}
\operatorname{Ric}^{\mathrm{Ch}}(\omega)=\pi f^*\omega_{FS}.
\end{equation}
This Ricci form is semipositive and has rank one at regular points
of $f$.  No Chern--Ricci-flat Hermitian metric exists on $X$.
\end{cited}
\begin{proof}
Gauduchon's conformal theorem
\cite[Theorem~2, pp.~132--133]{Gauduchon1985} supplies a Gauduchon
metric on the compact complex threefold. Fix any such $\omega_0$.
By \cref{fam:src:A.1.1} and $4S_2=f^*(p_2)$,
$K_X^{\otimes2}\simeq f^*\mathcal O_B(-1)$.
For $\theta=f^*\omega_{FS}$, with
$\int_{\mathbb P^1}\omega_{FS}=1$, this gives
\[
 c_{1,BC}(T_X)=\tfrac12[\theta]_{BC}.
\]
The convention
$\operatorname{Ric}^{\mathrm{Ch}}(\omega)=-i\partial\bar\partial
\log\det(g)$ represents $2\pi c_{1,BC}(T_X)$.
Thus $\pi\theta$ is the prescribed representative of that class.
Theorem~1.3 of \cite[arXiv v2, p.~3; conventions on p.~2]{STW2017}
gives a Gauduchon metric with this Chern--Ricci form and
$[\omega^2]_A=[\omega_0^2]_A$.

The pullback $\theta$ is semipositive and has rank one where
$f$ is a submersion. In particular
$\int_X\theta\wedge\omega_0^2>0$.
Every $\partial\bar\partial$-exact $(1,1)$ form pairs to zero
with the Gauduchon square, by integration by parts.
Consequently $[\theta]_{BC}\ne0$, and a zero Chern--Ricci form
would contradict the displayed Chern-class identity.
\end{proof}

\begin{proposition}[The fibre obstruction to SKT and astheno-K\"ahler metrics]\hypertarget{tag-01SK}{}\label{inv:C3.7}
\StatementRecord{inv:C3.7}{proved}{01SK}{\SeriesRef{inv:A3.2}, \SeriesRef{inv:C1.2}, \SeriesRef{inv:C2.1}, \SeriesRef{inv:eq:C-eta-negative}}

The threefold $X$ admits no pluriclosed (SKT) metric and no
astheno-K\"ahler metric.  The geometric input is the value $h=1$
and the global Dolbeault class $q$ in \ref{inv:A3.2}, whose restriction to a
regular fibre is $\eta=u\wedge w+6\gamma\wedge\delta$.
More generally, it is enough that the real image of
$H_A^{1,1}(X;\R)\to H_A^{1,1}(F;\R)$ be contained in
$\R\eta$ for one regular fibre $F$.
\end{proposition}


\begin{proof}
\begin{enumerate}[label=\arabic*.,leftmargin=*]
\item Put
$y=\operatorname{Im}\tau>0$, $m=\operatorname{Im}\mu$, and
$\Delta=\operatorname{Im}\beta-6m^2/y<0$.
For the marked period rows
$\sigma_1=6\mu\gamma+\tau u+w$ and
$\sigma_2=\beta\gamma+\mu u+\delta$ we have
$\eta=u\wedge\sigma_1+6\gamma\wedge\sigma_2$.
This proves that the real class $\eta$ has type $(1,1)$.
The determinant of the four real period columns, with rows ordered
as $\operatorname{Re}\zeta_1,\operatorname{Im}\zeta_1,
\operatorname{Re}\zeta_2,\operatorname{Im}\zeta_2$, is
$y\Delta<0$.  Since $\eta^2=12\gamma\wedge u\wedge w\wedge\delta$,
\begin{equation}\hypertarget{tag-01TD}{}\label{inv:eq:C-eta-negative}
\int_{F,\mathrm{or}_{\C}}\eta^2=-12.
\end{equation}
Equivalently the invariant representative of $\eta$ has Hermitian
coefficient matrix
\[
\frac1{y\Delta}
\begin{pmatrix}\operatorname{Im}\beta&-6m\\-6m&6y\end{pmatrix},
\qquad \det=\frac6{y\Delta}<0,
\]
so it has one positive and one negative eigenvalue.  

\item The explicit invariant fibre representative of $\gamma^{0,1}$ is
\begin{equation}
\Gamma=\frac{i}{2\Delta}
\left(d\bar\zeta_2-\frac myd\bar\zeta_1\right).
\label{I:eq:fibre-gamma}
\end{equation}
Its invariant $(1,1)$ representative is
\begin{equation}
Q_F=\frac1{2iy\Delta}\bigl(
\operatorname{Im}\beta\,d\zeta_1\wedge d\bar\zeta_1
-6m\,d\zeta_1\wedge d\bar\zeta_2
-6m\,d\zeta_2\wedge d\bar\zeta_1
+6y\,d\zeta_2\wedge d\bar\zeta_2\bigr).
\label{I:eq:fibre-eta-form}
\end{equation}
Direct exterior multiplication gives
\[
Q_F\wedge\Gamma=
\frac1{4y\Delta}
d\zeta_1\wedge d\bar\zeta_1\wedge d\bar\zeta_2,
\quad
\int_FQ_F\wedge\Gamma\wedge d\zeta_2=-1.
\]
This also proves the nonvanishing used in \ref{inv:C1.2}.

\item When $h=1$, \ref{inv:C2.1} identifies $H_A^{1,1}(X)$ with the sum of the
one horizontal source and its vertical conjugate.  A representative
$Q$ of $q$ is nonzero in the horizontal source, since its fibre
restriction is nonzero whereas $d_1g$ restricts to zero.  Thus
$[Q]_A,[\bar Q]_A$ form a basis.  Their restrictions to $F$ both
equal the real class $\eta$.  The real image of the restriction map
is therefore $\R\eta$.

\item If $\omega$ were a pluriclosed positive real $(1,1)$ form, its
restriction to $F$ would be pluriclosed and positive.  Average this
restriction under translations of the compact complex torus $F$.
The average $\omega_{\mathrm{av}}$ is invariant, closed, and
strictly positive.  Translation preserves the Aeppli class: on a
torus the $\partial\bar\partial$-lemma identifies Aeppli
cohomology with the invariant $(p,q)$ forms, so the average is the
unique invariant representative of that class.  This identification
can also be checked Fourier mode by Fourier mode: each nonzero mode
is a Koszul square for two independent type components of its real
nonzero covector and contributes zero Aeppli cohomology.  Thus
$[\omega_{\mathrm{av}}]=r\eta$ for some real $r$.
But positivity gives $\int_F\omega_{\mathrm{av}}^2>0$, whereas
\eqref{inv:eq:C-eta-negative} gives $-12r^2\le0$, a contradiction.
In dimension three the astheno-K\"ahler equation
$\partial\bar\partial\omega^{n-2}=0$ is exactly the SKT equation.
\end{enumerate}
\end{proof}





\section{Global coefficient representatives}
\hypertarget{tag-03L3}{}\label{inv:section:coefficient-atlas}

\subsection{A coefficient atlas for the two generators}

\begin{remark}[Sources and coefficient conventions]
\hypertarget{tag-PI16}{}\label{I:rem:coefficient-sources}
\StatementRecord{I:rem:coefficient-sources}{proved}{PI16}{}
The Green representative of \cref{inv:C3.1} and the harmonic
normalization of a lift of \cref{inv:A3.2} give operator descriptions.
The following construction specifies a holomorphic \v Cech
cocycle and every smooth coefficient in its Dolbeault representative.
It uses the original chart maps through all three special fibres.
The period matrix, two affine finite fillings, logarithmic collar maps
and untranslated cusp fan are the original construction
\cite[Sections~3--6, pp.~13--40]{AC2026}.  The same local pieces in the two
versions of \cite{DuEtAl2026v1,DuEtAl2026v2} are identified by the
parameter translation $a=c+\delta$ of \cref{night:P5.7}; in particular
$\Omega_a=\Pi_c$, using the same ordered lattice, vectors and fan.
The scalar cocycle below is precisely $f=-\eta_{\rm DGKX}$ on every
displayed generating arrow in
\cite[Section~6.4, scalar cocycle and Lemma~6.14 with proof, p.~84; conventions and generating arrows, pp.~82--83]{DuEtAl2026v1}.
The minus sign is the stated line-coordinate convention.  The
exponential quotient comparison with Engel's elliptic-surface
presentation is the map already given in
\cref{fam:src:B.1.1}.

Conversion of a holomorphic cocycle to a Dolbeault form by smooth
weights is the usual \v Cech--Dolbeault construction
\cite[Theorem~3.1 and proof, p.~10]{AbateBracciSuwaTovena2010Atiyah}.
On two charts its formula applied to $(0,0,c_{01}/12)$ is
$-\bar\partial w_0\wedge c_{01}/12
=\bar\partial(w_1c_{01})/12$, exactly the convention used here;
replace $c_{01}/12$ by $f_{01}$ for $G$.
\end{remark}

\begin{proposition}[A specified global coefficient atlas]
\hypertarget{tag-01SL}{}\label{inv:C3.8}
\StatementRecord{inv:C3.8}{proved}{01SL}{\SeriesRef{inv:A2.1}, \SeriesRef{inv:A3.2}, \SeriesRef{inv:C3.1}, \SeriesRef{inv:C3.2}, \SeriesRef{inv:C3.7}}
Use the standard marked family of \cref{fam:src:A.3.1},
with vectors $v_1=(1,2,-4,0)^t$, $v_2=(-1,-3,3,0)^t$,
the untranslated original cusp fan, and any
$c\in\mathcal P=\{\operatorname{Im}c<3\log(6)/\pi\}$.
There are specified smooth forms
\[
 G\in\mathcal A^{0,1}(X),\qquad Q\in\mathcal A^{1,1}(X),
 \qquad \bar\partial G=\bar\partial Q=0,
\]
whose classes satisfy
\[
 [G]=g_\gamma,\qquad
 i_b^*[Q]_{\bar\partial}=\eta,\qquad
 \exp(12[G])=f^*\mathcal O_{\mathbb P^1}(1).
\]
Here $i_b$ is any regular fibre inclusion and
$\eta=u\wedge w+6\gamma\wedge\delta$ in the original integral marking.
The formulas \eqref{inv:eq:atlas-final-forms} and
\eqref{inv:eq:atlas-final-coefficients}, with the elementary recursion
\eqref{inv:eq:atlas-weight-derivative}, give their coefficients on a
specified countable refinement of the original atlas.  Every sum is
locally finite, also at the multiple fibres, cusp double curves and
both cusp triple points.  The form $Q$ is a complex representative;
no real $\bar\partial$-closed representative can have this fibre class.
\end{proposition}

\begin{proof}\leavevmode\par


 
\begin{enumerate}[label=\arabic*.,leftmargin=0pt,labelsep=.5em]
\item \emph{The chart domains and their inverse maps.}
Use the ordinary, finite and cusp covering charts of
\cref{fam:src:A.1.2,fam:src:A.3.1}, with their specified affine,
toric and exponential transition maps.
More explicitly, write $t_1=t$, $t_2=t-1$, $q=1/t$, and retain
the fixed disjoint special discs of those constructions.  After
restricting their charts by inclusion, write their radii as
$r_1,r_2,\epsilon$, with $r_1,r_2<1/4$, and use
\[
 V_j=\{|t_j|<r_j\},\quad
 V_0=\{|q|<\epsilon\},\quad
 V_*=B\setminus\left(
 \{|t_1|\le r_1/2\}\cup\{|t_2|\le r_2/2\}
 \cup\{|q|\le\epsilon/2\}\right).
\]
Thus only $V_*\cap V_j$ are nonempty pairwise overlaps between
distinct pieces.  All three are annuli.  The finite coordinate domains
are the original torus covers over $|s_j|^{m_j}<r_j$, with
$t_j=s_j^{m_j}$ and $(m_1,m_2)=(3,4)$; their maps to $X$ are the
original free affine quotient maps.  The middle domains are the
ordinary lifted base opens over $V_*$, times $\mathbb C^2$.

For the cusp, retain the full holomorphic matrix $C(q)$ and a closed
disc $|q|\le r_0$ on which it is holomorphic.  The original radius
satisfies
\[
 0<\epsilon<\min(r_0,1),\qquad
 |\log\epsilon|>
 2\max_{|q|\le r_0}\|-2\pi\operatorname{Im}C(q)\|.
\]
These are the chart and properness bounds in
Proposition~\ref{fam:src:A.3.1}.  For every triangle
$\sigma$ of the original periodic triangulation, use
\[
 \mathcal U_\sigma(2,\epsilon)=
 \{(z_0,z_1,z_2): |z_i|<2,\ |z_0z_1z_2|<\epsilon\},
 \qquad q=z_0z_1z_2.
 \]
The toric characters are the basis dual to the three original
height-one ray generators.  Their integral exponent matrix is
unimodular, so its inverse gives every inverse toric overlap.
The toric orbit calculation in that proof shows that these domains
cover the cusp space before quotient, including its boundary.

All these maps are local biholomorphisms to $X$.  Their overlaps are
fibre products, retaining every deck branch.  A lattice map
$\zeta\mapsto\zeta+\Pi\lambda$ has inverse
$\zeta\mapsto\zeta-\Pi\lambda$.  A finite affine map
$(s,\zeta)\mapsto(\omega s,R(s)\zeta+\Pi(\omega s)v/m)$
has inverse
\[
 s=\omega^{-1}s',\qquad
 \zeta=R(s)^{-1}\bigl(\zeta'-\Pi(s')v/m\bigr).
\]
The cusp period map has inverse
$x\mapsto q^{-B_0n}\exp(-2\pi iC(q)n)x$.
The finite collar inverse subtracts its full translating section,
and the cusp collar inverse uses the branches
$\zeta=(2\pi i)^{-1}\log x$, whose differences are the
integral periods.  These maps specify the refinement and its local inverse branches
inside the given $X$.

 
\item \emph{The regular class and its sign.}
Put
\[
 \Pi=[Z\mid I_2],\quad
 Z=\begin{pmatrix}6\mu&\tau\\\beta&\mu\end{pmatrix},\quad
 T=\begin{pmatrix}\tau&6\mu\\\mu&\beta\end{pmatrix},\quad
 K=\begin{pmatrix}0&6\\1&0\end{pmatrix},\quad D=\operatorname{diag}(1,6).
\]
The exact real marking is
\[
 q_B=(r_u,r_\gamma)^t,\quad p_B=(r_w,r_\delta)^t,
 \qquad \zeta=p_B+Tq_B,
\]
with inverse $q_B=(\operatorname{Im}T)^{-1}\operatorname{Im}\zeta$,
$p_B=\operatorname{Re}\zeta-(\operatorname{Re}T)q_B$.
For $\lambda=(n_\gamma,n_u,m_w,m_\delta)$, write
$n=(n_\gamma,n_u)^t$ and define
\begin{equation}\hypertarget{tag-03L4}{}\label{inv:eq:atlas-theta-factor}
 a_\lambda(\zeta)=
 \exp\left(2\pi i\left[n^tK\zeta+\tfrac12 n^tKZn\right]\right).
\end{equation}
Symmetry of $KZ$ makes its logarithmic composition defect the integer
$n^tK m'$.  Thus these factors satisfy the multiplicative cocycle identity.

Let $Y=\operatorname{Im}T$, $S=DY^{-1}$ and
$\varphi=2\pi q_B^tDYq_B$.  The metric coefficient $e^\varphi$
obeys $h(\zeta+\Pi\lambda)|a_\lambda|^2=h(\zeta)$.
With $F=\bar\partial\partial\log h$ and $c_1=iF/(2\pi)$ its
Chern form is
\begin{equation}\hypertarget{tag-03L5}{}\label{inv:eq:atlas-semiflat}
 Q_{\rm sf}=-\frac{i}{2\pi}\partial\bar\partial\varphi
 =-\frac i2(d\zeta-dTq_B)^tS\wedge
             \overline{(d\zeta-dTq_B)}.
\end{equation}
Expanding under the exact real inverse gives
$dq_B^tD\wedge dp_B=dr_u\wedge dr_w+6dr_\gamma\wedge dr_\delta$.
In particular the fibre class is $\eta$, with its original sign.
If $y=\operatorname{Im}\tau$, $m=\operatorname{Im}\mu$,
$r=\operatorname{Im}\beta$ and $\Delta=r-6m^2/y$, then
\[
 S=\frac1{y\Delta}
 \begin{pmatrix}r&-6m\\-6m&6y\end{pmatrix},
\]
which gives the invariant fibre form \eqref{I:eq:fibre-eta-form}.

 
\item \emph{The line on the middle piece.}
Use the original matrices $A_1,A_2$ in \cref{fam:section:A}.
Their period transformation laws and complex-linear maps are
\[
\begin{array}{c|ccc}
j&\tau_j&\mu_j&\beta_j\\\hline
1&1-1/\tau&(1-\mu)/\tau&\beta-6(\mu-1)^2/\tau+2\\
2&-1/\tau&1+\mu/\tau&\beta-6\mu^2/\tau-3,
\end{array}
\]
\[
 R_1=\begin{pmatrix}-1/\tau&0\\(1-\mu)/\tau&1\end{pmatrix},
 \qquad R_2=\begin{pmatrix}1/\tau&0\\-\mu/\tau&1\end{pmatrix},
 \qquad R_j\Pi=\Pi_jA_j.
\]
For both generators put $F_j(\zeta)=-\zeta_1^2/(2\tau)$.
The differences
\[
 \frac{\log a_{A_j\lambda}(R_j\zeta)-\log a_\lambda(\zeta)}{2\pi i}
 -F_j(\zeta+\Pi\lambda)+F_j(\zeta)
\]
on the four original lattice generators are
$(24,0,1/2,0)$ and $(-9,0,0,0)$.  Their exponential ratios are
characters, so the four values determine the full lattice comparison.
Twelfth powers kill every defect.  The line $L_*=L_\eta^{12}$ has
the linear lifts
\[
 (\zeta,v)\longmapsto
 (R_j\zeta,e^{2\pi i\,12F_j(\zeta)}v).
\]
The Gaussian sums vanish on the order-three and order-four relations.
Applying generator two and then generator one gives
$(\tau,\mu,\beta,\zeta)\mapsto(\tau+1,\mu,\beta-1,\zeta)$
and Gaussian sum zero.  These identities give the middle line on the
specified quotient, with trivial fibre multiplier for that cusp logarithm
turn.

 
\item \emph{Both finite affine lifts, including their phases.}
Let $b_j=v_j/m_j$.  Use the same polynomial logarithm in
\eqref{inv:eq:atlas-theta-factor} to define $a_b$ for this rational
vector.  Since $12b_j$ is integral, the affine map
\[
 g_j(s,\zeta)=(e^{-2\pi i/m_j}s,
                R_j\zeta+\Pi_jb_j)
\]
has the following line lift:
\begin{equation}\hypertarget{tag-03L6}{}\label{inv:eq:atlas-affine-lift}
 k_j(s,\zeta)=\exp\left(2\pi i
 \left[12F_j(\zeta)+12\frac{\log a_{b_j}(R_j\zeta)}{2\pi i}
                  +\kappa_j\right]\right),
 \quad \kappa_1=\frac{32}{3},\quad\kappa_2=\frac{81}{8}.
\end{equation}
The factor $a_{b_j}$ uses the target period $Z_j$.
Direct composition gives $g_j^{m_j}\zeta=\zeta+\Pi v_j$.
Before adding $\kappa_j$, the logarithmic difference from
$a_{v_j}^{12}$ is $-32$ for $j=1$ and $-81/2$ for $j=2$.
The stated phases cancel these differences exactly.  Hence the power
of \eqref{inv:eq:atlas-affine-lift} is the lattice action,
and it defines a line $L_j$ on the original free finite quotient.
The functions are holomorphic through $s=0$: the original periods
are holomorphic there, $\tau\ne0$, and every displayed exponential
is a unit.  Thus this is a line over the full multiple-fibre
neighbourhood.  Its regular fibre Chern class is $12\eta$.

 
\item \emph{Cusp Cartier frames and all extending units.}
Write the original full period matrix as
\[
 Z=sB_0+C(q),\quad s=\frac{\log q}{2\pi i},\quad
 B_0=\begin{pmatrix}0&1\\-1&0\end{pmatrix}.
\]
The period $n\in\mathbb Z^2$ acts by
$x\mapsto q^{B_0n}\exp(2\pi iC(q)n)x$.
Use the factor
\begin{equation}\hypertarget{tag-03L7}{}\label{inv:eq:atlas-cusp-factor}
 A_n=x^{K^tn}q^{h(n)}\exp(\pi i n^tKC(q)n),
 \qquad h(n)=-3n_\gamma^2+\frac{n_u(n_u-1)}2.
\end{equation}
Both $KB_0=\operatorname{diag}(-6,1)$ and $KC(q)$ are symmetric.
Consequently multiplication proves the exact cocycle in
\eqref{inv:eq:atlas-cusp-factor}.  Its $q$ exponent is integral.
Its ratio to \eqref{inv:eq:atlas-theta-factor}, on a fixed punctured
fibre, is $q^{-n_u/2}$, independent of $x$; its fibre Chern class
is therefore $\eta$.

Interpolate $L(i,j)=i(i-1)/2-3j^2$ on the prescribed triangulation.
Both triangles in the square with lower corner $(i,j)$ have
\[
 m_{ij}=(i,-6j-3),\quad
 l_{ij}=-i(i+1)/2+3j(j+1),\quad f_{ij}=x^{m_{ij}}q^{l_{ij}}.
\]
These are Cartier frames.  The successive ratios $f_{i+1,j}/f_{ij}$
and $f_{i,j+1}/f_{ij}$ are $x_1q^{-i-1}$ and
$x_2^{-6}q^{6(j+1)}$; their exponents vanish on the corresponding
common face, so they are units on the toric overlap.  The two triangles
inside a square have ratio one.  In the line coordinate $w=z/f_{ij}$
the overlap multiplier is the reciprocal frame ratio.

For $v=B_0n$ the exact identity
$L(u+v)-L(u)=(K^tn)\cdot u+h(n)$ cancels every $x$ and $q$ power
in the period action on $w$.  The remaining factor is
\begin{equation}\hypertarget{tag-03L8}{}\label{inv:eq:atlas-cusp-unit}
 U_{n,ij}(q)=\exp\left(\pi i n^tKC(q)n
       -2\pi i\,m_{i+v_1,j+v_2}^{t}C(q)n\right).
\end{equation}
It is a holomorphic unit at zero.  Its cocycle follows either from
these exact cancellations or by direct multiplication.  Thus it
descends on the original free properly discontinuous toric quotient.
Let $L_0$ be its twelfth power.  On the punctured cusp the exact
ratio $L_0\otimes L_*^{-1}$ has factors $q^{-6n_u}$.

 
\item \emph{The inverse used on every annulus.}
For a line with normalized lattice factors $1$ on the two real
periods and $e^{2\pi i n_T^tl(t)}$ on the columns of $T$, the form
$A=a(t)^t d\zeta$ is a holomorphic connection when
\begin{equation}\hypertarget{tag-03L9}{}\label{inv:eq:atlas-connection-inverse}
 T'^t a=-2\pi i l',\qquad
 (T'^t)^{-1}=
 \frac1{\tau'\beta'-6\mu'^2}
 \begin{pmatrix}\beta'&-\mu'\\-6\mu'&\tau'\end{pmatrix}.
\end{equation}
Indeed translation by $Tn_T$ adds
$a^tT'n_T\,dt=-d\log(e^{2\pi in_T^tl})$.

Every use of this inverse is on one of the three annuli above.
By \cref{inv:A2.1},
$6\mu_z^2-\tau_z\beta_z=\tau_z^2E_4(\tau)^2/E_6(\tau)$,
with the apparent poles removed, has order two at the order-three
point, order zero at the order-four point, and no further zeros on
the regular locus.  A local change of base coordinate multiplies
the determinant by the square of a nonzero derivative.  Thus $\det T'$ is nonzero throughout the three annular domains
of the inverse.  In the cusp variable
its leading term is $1/(4\pi^2q^2)$, in agreement with the exact
logarithmic period matrix.

 
\item \emph{The cusp overlap connection.}
For $F_0=L_0\otimes L_*^{-1}$, use
$l=(-6\log q/(2\pi i),0)^t$.  Its connection is
\begin{equation}\hypertarget{tag-03LA}{}\label{inv:eq:atlas-cusp-connection}
 A_0=a_0^td\zeta,\qquad
 a_0=\frac6q(T'^t)^{-1}\binom10
 =\frac6{q(\tau'\beta'-6\mu'^2)}\binom{\beta'}{-6\mu'}.
\end{equation}
It is single-valued: a logarithm turn adds
$\operatorname{diag}(1,-1)$ to $T$ and changes neither $T'$ nor
$\zeta$.  In the cusp frame $f_{ij}^{12}$ the form is
$A_0+12d\log f_{ij}$.  This transformation retains the full
units \eqref{inv:eq:atlas-cusp-unit}.

 
\item \emph{The finite overlap connections and all logarithm turns.}
The collar map, as in
\cref{fam:src:A.3.1}, is
\[
 H(s,\zeta)=(s,\zeta+\sigma),\qquad
 \sigma=h\Pi v_j,\qquad h=\frac{\log s}{2\pi i}.
\]
The identities $h(g_js)=h(s)-1/m_j$ and $A_jv_j=v_j$ give the
exact square $Hg_j^{\rm aff}=g_j^{\rm lin}H$.
For $F=H^*L_*\otimes L_j^{-1}$ its lattice factors are
$e^{2\pi i\,12n_T^tD\sigma}$.  Put
$v_q=(v_u,v_\gamma)^t$ and $S=v_q^tD\Pi v_j$.  Then
\begin{equation}\hypertarget{tag-03LB}{}\label{inv:eq:atlas-finite-connection}
 A=a^td\zeta-12\pi i\,d(h^2S),\qquad
 a=-24\pi i(T'^t)^{-1}D\sigma'.
\end{equation}
The last differential is a base one-form.
Equation \eqref{inv:eq:atlas-connection-inverse} proves lattice descent.
For an integer logarithm turn $h\mapsto h+k$, the frame changes by
$a_{kv_j}(\zeta+\sigma)^{12}$.  The vertical change in
\eqref{inv:eq:atlas-finite-connection} is
$-24\pi i kDv_q$.  The base change is exactly
\[
 -12\pi i\,d((2kh+k^2)S)
 =-d\log a_{kv_j}(\zeta+\sigma)^{12}
      +24\pi i k v_q^tD\,d\zeta.
\]
Thus the complete connection descends across every logarithm turn,
with the displayed base term.

For the remaining finite cyclic descent its line multiplier is
\[
 k_F=\frac{e^{2\pi i\,12F_j(\zeta+\sigma)}}{k_j(s,\zeta)}.
\]
Let $k_F^{(r)}$ be the product over the first $r$ affine
iterates.  The finite sum
\begin{equation}\hypertarget{tag-03LC}{}\label{inv:eq:atlas-averaged-connection}
 A^{\rm inv}=\frac1{m_j}\sum_{r=0}^{m_j-1}
       \bigl((g_j^r)^*A+d\log k_F^{(r)}\bigr)
\end{equation}
is a connection on the same line.  The exact power relation in
\eqref{inv:eq:atlas-affine-lift} makes cyclic permutation of the
summands valid, so it descends to the original annulus.  For
$F_j=L_j\otimes L_*^{-1}$ use the dual connection $-A^{\rm inv}$.
This is the tensor-dual connection, with all transition factors
inverted and the displayed sign.  These holomorphic connections need not be flat:
$da$ can supply a $(2,0)$ curvature term.

 
\item \emph{The holomorphic one-form cocycle.}
Choose the specified line frames in each same-piece chart.  If
the line-coordinate map is $v_b=g_{ba}v_a$, put
\begin{equation}\hypertarget{tag-03LD}{}\label{inv:eq:atlas-cocycle}
 c_{ab}=\frac{i}{2\pi}d\log g_{ba}\quad\hbox{within a piece},
 \qquad
 c_{ab}=-\frac{i}{2\pi}A_{F_j,(a,b)}\quad
       (a\text{ middle},\ b\text{ in piece }j),
\end{equation}
and $c_{ba}=-c_{ab}$.
On a triple with two middle charts, the ratio frame of $F_j$
changes by the inverse middle transition, and its connection changes
by $+d\log g_*$.  On a triple with two special charts it changes
by the special transition, and its connection changes by
$-d\log g_j$.  These signs give
$c_{ab}+c_{bc}=c_{ac}$ in both cases.  The special pieces do not
intersect one another.  This proves the cocycle identity on every
component of every overlap.

The result is the holomorphic one-form cocycle $c_{ab}$.
Its construction uses the connection transformation law; the
connections may have the $(2,0)$ curvature $dA$ displayed above.

 
\item \emph{A completely specified smooth partition.}
In every original uniformizing domain enumerate all rational closed
Euclidean balls $B(c_n,R_n)$ compactly contained in that domain.
Use a fixed enumeration of the original charts by piece, cone label
and finite deck word.  Order the balls by integer height bounding
the chart index and all coordinate numerators and denominators,
then lexicographically, removing repeats.  Each height block is
finite.  The half-radius balls cover every chart.

Set
\[
 E(t)=\begin{cases}e^{-1/t}&t>0,\\0&t\le0,\end{cases}
 \quad \chi(t)=\frac{E(t)}{E(t)+E(1-t)},\quad
 b_n(u)=\chi\!\left(
       \frac{R_n^2-\|u-c_n\|^2}{3R_n^2/4}\right).
\]
Extend $b_n$ by zero outside its ball.  It is one on the half-radius
ball and is flat on its support boundary.  For the étale map
$\pi_n$, define
\begin{equation}\hypertarget{tag-03LE}{}\label{inv:eq:atlas-weights}
 \begin{aligned}
 S_n(x)&=\sum_{u\in\pi_n^{-1}(x)}b_n(u),\qquad
 a_n(x)=\chi(2S_n(x)),\\
 P_1&=1,\qquad P_{n+1}=P_n(1-a_n),\\
 R(s)&=\chi(2s)/s\quad(s>0),\qquad R(0)=0,\\
 w_n(u)&=P_n(\pi_nu)R(S_n(\pi_nu))b_n(u).
 \end{aligned}
\end{equation}
The function $R$ is smooth at zero because its numerator is flat.
Each trace $S_n$ is locally a finite smooth sum: its source support
is compact and the map is étale, hence has discrete fibres.  Every
point has a lift in an open half-radius ball.  For some $N$ one therefore
has $S_N\ge1$ on a neighbourhood, after shrinking that neighbourhood
if necessary.  There $a_N=1$, so all subsequent weights vanish.
Consequently the weights are locally finite.  The finite telescoping
product gives
\[
 \sum_n\sum_{u\in\pi_n^{-1}(x)}w_n(u)=1.
\]
This construction also handles original chart maps that are not
globally injective; it retains all their local inverse branches.

Every derivative in the final coefficients is
given by the following finite recursion.  In a local inverse branch
write $\mathcal D=\partial/\partial\bar z_l$ and
$v_n=(R_n^2-\|u-c_n\|^2)/(3R_n^2/4)$.  Then
\begin{equation}\hypertarget{tag-03LF}{}\label{inv:eq:atlas-weight-derivative}
 \begin{aligned}
 \mathcal D b_n&=-\frac4{3R_n^2}\chi'(v_n)
        \sum_j(u_j-c_{n,j})\mathcal D\bar u_j,\\
 \mathcal D S_n&=\sum_{u\in\pi_n^{-1}(x)}\mathcal D b_n,
 &\mathcal D a_n&=2\chi'(2S_n)\mathcal D S_n,\\
 \mathcal D P_1&=0,
 &\mathcal D P_{n+1}&=(1-a_n)\mathcal D P_n-P_n\mathcal D a_n,\\
 R'(s)&=\frac{2s\chi'(2s)-\chi(2s)}{s^2}\quad(s>0),
 &R'(0)&=0,\\
 \mathcal D w_n&=R(S_n)b_n\mathcal D P_n
       +P_nR'(S_n)b_n\mathcal D S_n+P_nR(S_n)\mathcal D b_n.
 \end{aligned}
\end{equation}
The derivatives $\mathcal D\bar u_j$ are the conjugate Jacobians
of the explicit inverse chart maps in step~1.  Finally
\[
 \chi'(t)=
 \frac{E'(t)E(1-t)+E(t)E'(1-t)}{(E(t)+E(1-t))^2},
 \qquad E'(t)=\begin{cases}e^{-1/t}/t^2&t>0,\\0&t\le0.\end{cases}
\]
Together with the inverse-chart Jacobians, this evaluates every
weight derivative in the coefficient formulas.

 
\item \emph{The primitive scalar cocycle.}
For a complex parameter $z$, use lattice factors
$e^{-2\pi iz\gamma(\lambda)}$ and middle linear factors one.
The finite affine factor is
$e^{-2\pi iz\gamma(v_j)/m_j}$; its $m_j$th power is the
period factor of $v_j$.  The cusp period factor is
$e^{-2\pi izn_\gamma}$, which extends in every toric chart.
The finite collar comparison is
$e^{-2\pi iz\gamma(v_j)h(s)}$.
Its changes under $h\mapsto h-1/m_j$ and $h\mapsto h+k$ are
exactly those required by the affine generator and the period $kv_j$.

Write every resulting transition as $e^{2\pi iz f_{ab}}$.
Its coefficients are respectively
$-\gamma(\lambda)$, $0$, $-\gamma(v_j)/m_j$,
$-n_\gamma$ and $-\gamma(v_j)h(s)$ on those arrows;
pure toric frame changes have coefficient zero.
Composition and inversion use addition and negation.
Since the identities hold for every complex $z$,
$f_{ab}+f_{bc}=f_{ac}$ on each overlap component.

 
\item \emph{The evaluated global forms and their normalization.}
In a chart $a$, with the indicated component of each fibre-product
overlap, put
\begin{equation}\hypertarget{tag-03LG}{}\label{inv:eq:atlas-final-forms}
 H_a=\sum_n\sum_u w_n(u)f_{a,n}(x,u),\quad
 B_a=\sum_n\sum_u w_n(u)c_{a,n}(x,u),\quad
 G=\bar\partial H_a,\quad Q=\frac1{12}\bar\partial B_a.
\end{equation}
The cocycles and sum of weights give
$H_a-H_b=f_{ab}$ and $B_a-B_b=c_{ab}$.
Their right sides are holomorphic.  Therefore these formulas agree
on every overlap and define global $\bar\partial$-closed forms.
Writing $c_{a,n}=\sum_kc_{a,n,k}dz_k$, their coefficients are
\begin{equation}\hypertarget{tag-03LH}{}\label{inv:eq:atlas-final-coefficients}
 G_{\bar l}=\sum_n\sum_u f_{a,n}(x,u)
                       \frac{\partial w_n(u)}{\partial\bar z_l},
 \qquad
 Q_{k\bar l}=-\frac1{12}\sum_n\sum_u c_{a,n,k}(x,u)
                       \frac{\partial w_n(u)}{\partial\bar z_l}.
\end{equation}
Every coefficient is a locally finite sum of the specified period,
Cartier, connection and elementary cutoff expressions.  The minus
sign in the second formula is from moving $d\bar z_l$ past $dz_k$.

On a regular fibre, the one-form cocycle is that of the line with
Chern class $12\eta$, with the Chern sign fixed in step~2.
Consequently $i_b^*[Q]=\eta$.  This is a statement about the exact
Dolbeault quotient; the resulting coefficient form need not be
pointwise translation invariant on the fibre.
For the scalar cocycle, $f_{ab}=-\Delta r_\gamma$, so its fibre
image is $\bar\partial r_\gamma$, explicitly
\[
 \frac{i}{2\Delta}
 \left(d\bar\zeta_2-\frac{\operatorname{Im}\mu}{\operatorname{Im}\tau}
                      d\bar\zeta_1\right).
\]
This is \eqref{I:eq:fibre-gamma}; injectivity in
\cref{inv:C3.1} gives $[G]=g_\gamma$.
At $z=1$ the scalar line has the finite root characters and
collar frames of $\mathcal O(S_1-S_2)$.  Its twelfth power is
\[
 \mathcal O(12S_1-12S_2)
 =f^*\mathcal O(4p_1-3p_2)\simeq f^*\mathcal O(1),
\]
where the last map is induced by multiplication by
$t^4/(t-1)^3$: its divisor is $4p_1-3p_2-p_0$, so it identifies
$\mathcal O(4p_1-3p_2)$ with the hyperplane line
$\mathcal O(p_0)=\mathcal O(1)$ in the original base coordinate.
This agrees with the original sign and factor in \cref{inv:C3.2}.

 
\item \emph{The fibre square of a closed real form.}
Suppose a real $(1,1)$ form with the stated fibre class were
$\bar\partial$ closed.  Reality would give $\partial Q=0$, hence
$dQ=0$.  Since $H^2(X;\mathbb R)=0$, write $Q=dA$.
On a regular compact fibre, Dolbeault equality with $\eta$ and
Stokes' theorem give $\int_F(i^*Q)^2=\int_F\eta^2=-12$.
But $(i^*Q)^2=d(i^*A\wedge i^*Q)$ has integral zero.
This contradicts the exact fibre square and proves the last assertion.


\end{enumerate}
\end{proof}

\Uses{the original marked family and full chart maps in
\cref{fam:section:A,fam:src:A.1.2,wb:thm:retained-period-deformation};
\CompendiumRef{I}{wb:lem:toric-chart-control};
\CompendiumRef{I}{wb:thm:candidate-cusp-geometry};
the exact determinant divisor in \cref{inv:A2.1};
the fibre class in \cref{inv:A3.2};
the Picard and primitive normalization in \cref{inv:C3.1,inv:C3.2}}

\subsection{Comparison with the Green and harmonic representatives}
\begin{proposition}[Green and harmonic representatives]
\hypertarget{tag-03LI}{}\label{inv:C3.8:operator-comparison}
\StatementRecord{inv:C3.8:operator-comparison}{proved}{03LI}{}


The real restriction map
\[
r_b:H_A^{1,1}(X;\R)\longrightarrow H_A^{1,1}(F_b;\R)
\]
has image $\R\eta$ by \ref{inv:A3.2} and \ref{inv:C3.7}.  Its source has real dimension
two and its target has real dimension four.  The map is induced by pullback of $\partial\bar\partial$-closed forms followed by the quotient by $\partial$- and $\bar\partial$-exact forms on $F_b$. Thus it specifies the class of each restricted representative through that exact quotient map.

The formula \eqref{inv:eq:C-G-Green} supplies a global $G$ after a metric
choice.  For the standard primitive marking, $G_\gamma=G_0/12$
represents $g_\gamma$ by \ref{inv:C3.2}.  There is a precise operator
construction of $Q_H$ as well.  Fix a Hermitian metric on $X$ and a
regular fibre $F$, identify $\eta$ with its real invariant form on
$F$, and choose an orthonormal basis $e_1,e_2$ of the Dolbeault-harmonic
$(1,1)$ forms for $\Delta_{\bar\partial}$.  Define
\[
\ell(e)=-\frac1{12}\int_F i_F^*e\wedge\eta,\qquad
Q_H=\frac{\overline{\ell(e_1)}e_1+
            \overline{\ell(e_2)}e_2}
          {|\ell(e_1)|^2+|\ell(e_2)|^2}.
\]
The denominator is positive by \ref{inv:A3.2} and
$\int_F\eta^2=-12$.  The expression is independent of the
orthonormal basis, is $\bar\partial$-closed, and has fibre
Dolbeault class $\eta$.  Indeed the image of fibre restriction on
$H^{1,1}_{\bar\partial}(X)$ is $\C\eta$, with kernel
$\C d_1g$; the displayed functional measures exactly the coefficient
of $\eta$, and the formula makes that coefficient one.

The exact comparison on cohomology is
\[
 [G]=[G_0]/12=g_\gamma,\qquad
 [Q_H]-[Q]\in\C d_1g.
\]
The second inclusion follows from the restriction kernel in
\cref{inv:A3.2}.  Thus the two formulas choose lifts of the same fibre
class; their difference is the specified one-dimensional Leray
ambiguity.  The harmonic normalization remains basis independent for
the fixed metric and regular fibre.
\end{proposition}
\begin{proof}
\begin{enumerate}[label=\arabic*.,leftmargin=*]
\item The operator formulas have the following normalization.
The Green formula and $\exp([G_0])=f^*\OO(1)$ were proved in
\ref{inv:C3.1}; the primitive marking in \ref{inv:C3.2} then gives
$[G_0]/12=g_\gamma$.

\item Restriction of a Dolbeault-harmonic representative is a closed
Dolbeault form on the regular fibre. By \ref{inv:A3.2}, its class
is $c\eta$. Integration against the closed $(1,1)$ form $\eta$
annihilates Dolbeault-exact changes by Stokes' theorem. Since
$\int_F\eta^2=-12$, the displayed functional gives $\ell(e)=c$.
The image is $\C\eta$, hence $\ell$ is nonzero on the
 two-dimensional harmonic space and its squared norm
$|\ell(e_1)|^2+|\ell(e_2)|^2$ is positive.

\item For another orthonormal basis $e'_j=\sum_k U_{jk}e_k$ with
$U$ unitary, put $l_k=\ell(e_k)$. Then $l'=Ul$ and
\[
 \sum_j\overline{l'_j}e'_j
 =\sum_{j,k,r}\overline{U_{jk}}\bar l_k U_{jr}e_r
 =\sum_k\bar l_ke_k,
 \qquad \sum_j|l'_j|^2=\sum_k|l_k|^2.
\]
Thus $Q_H$ is independent of the orthonormal basis. It is a linear
combination of harmonic forms and $\ell(Q_H)=1$, so its fibre class
is precisely $\eta$. These computations establish the operator
representatives and their restrictions. The evaluated coefficient
construction above provides another representative with the same
fibre normalization.
\end{enumerate}
\end{proof}

\subsection{Changes of the specified coefficient gauge}

\begin{remark}[Fibre restriction of the coefficient representative]
\hypertarget{tag-03CT}{}\label{wb:rem:bcpr-literal-gauge}
\StatementRecord{wb:rem:bcpr-literal-gauge}{proved}{03CT}{}

The formulas \eqref{inv:eq:atlas-final-forms} and
\eqref{inv:eq:atlas-final-coefficients} choose global smooth
forms through all three special fibres.  Their restriction statement
uses the pullback map followed by the Dolbeault quotient on a regular
fibre.  In particular, its translation-invariant value is the class
represented in \eqref{I:eq:fibre-eta-form}.
The specified coefficient form on that fibre can be related to the
invariant representative by a Dolbeault coboundary.  The cocycle and
weight formulas below describe the changes produced by these choices.
\end{remark}

\begin{remark}[Changes of coefficient gauge]
\hypertarget{tag-03CS}{}\label{wb:rem:bcpr-global-gauge-boundary}
\StatementRecord{wb:rem:bcpr-global-gauge-boundary}{proved}{03CS}{}

Use the same holomorphic cocycles $c_{ab},f_{ab}$ on a common
refinement, and two locally finite weight systems $w_n,w'_n$ with
sum one.  In a chart $a$ set
\[
 K_a=\sum_n\sum_u(w'_n-w_n)c_{a,n},\qquad
 U_a=\sum_n\sum_u(w'_n-w_n)f_{a,n}.
\]
The cocycle identities and $\sum(w'_n-w_n)=0$ give
$K_a-K_b=0$ and $U_a-U_b=0$.  They therefore define a global
smooth $(1,0)$ form $K$ and a global smooth function $U$, with
\[
 Q'-Q=\tfrac1{12}\bar\partial K,\qquad
 G'-G=\bar\partial U.
\]
For a holomorphic change of cocycles
$c'_{ab}=c_{ab}+r_a-r_b$ and
$f'_{ab}=f_{ab}+u_a-u_b$, where $r_a$ is a holomorphic one-form
and $u_a$ a holomorphic function, the same weights give
\[
 B'_a-B_a=r_a-\sum_n\sum_u w_n r_n,\qquad
 H'_a-H_a=u_a-\sum_n\sum_u w_n u_n.
\]
The two sums on the right are globally defined smooth forms of the
indicated degrees.  Consequently
\[
 Q'-Q=-\tfrac1{12}\bar\partial\!\sum_n\sum_u w_n r_n,
 \qquad
 G'-G=-\bar\partial\!\sum_n\sum_u w_n u_n.
\]
These formulas identify the Dolbeault coboundaries under the
specified gauge changes and preserve all class and product
calculations.  A change of coherent-Leray lift of the fibre class is
measured separately by $\C d_1g$, as in the operator comparison above.
\end{remark}

\clearpage
\section{Twisted tangent cohomology}\label{pub:I:twists}
\subsection{All integral twists and the generic nonintegral twist}

\begin{proposition}[Integral Picard twists, including their torsion]\hypertarget{tag-01SO}{}\label{inv:D4.1}
\StatementRecord{inv:D4.1}{proved}{01SO}{the definitions and constructions stated in the proof}

For $M_n=\OO_X(n(S_1-S_2))$, define
$u(n)=\max\{d(n)+1,0\}$, $v(n)=\max\{-d(n)-1,0\}$ and
$\epsilon_n=0$ if $3\mid n$, $\epsilon_n=1$ otherwise. Then
\[
 f_*(T_X\otimes M_n)=\OO(d(n))\oplus\OO(d(n-1)),
\]
\begin{equation}
 R^1f_*(T_X\otimes M_n)=
 \OO(d(n))\oplus\OO(d(n-1))\oplus\OO(d(n+2)),
\label{pub:I:eq:D-12}
\end{equation}
\[
 R^2f_*(T_X\otimes M_n)=
 \OO(d(n+6))\oplus\C_{p_1}^{\epsilon_n}.
\]
In particular its cusp stalk is free of rank one. The four global
dimensions, for every integer $n$, are
\begin{equation}\begin{aligned}
 h^0&=u(n)+u(n-1),\\
 h^1&=v(n)+v(n-1)+u(n)+u(n-1)+u(n+2),\\
 h^2&=v(n)+v(n-1)+v(n+2)+u(n+6)+\epsilon_n,\\
 h^3&=v(n+6).                                          
\end{aligned}
\label{pub:I:eq:D-13}
\end{equation}
\end{proposition}


\begin{proof}
\begin{enumerate}[label=\arabic*.,leftmargin=*]
\item The defining roots give $\bar f^*R=\OO_X(S_1-S_2)$, so the projection
formula applies to \eqref{pub:I:eq:D-11}, with its coherent torsion summand.
Tensoring \eqref{pub:I:eq:D-8} changes its generator character from $-1$ to $n-1$.
There is exactly one invariant coefficient among exponents $0,1$ if
$n-1$ is $0$ or $1$ modulo three, and none if it is $2$. Multiplication
by the coarse parameter $t=s^3$ is zero on this module. Its invariant
direct image is therefore exactly the displayed skyscraper, of
length $\epsilon_n$, not a higher-order thickening. The two other
special places contribute no torsion. Formula \eqref{fam:eq:ab-root-degree} proves \eqref{pub:I:eq:D-12},
and the two-column Leray sequence proves \eqref{pub:I:eq:D-13}. 
\end{enumerate}
\end{proof}

\begin{proposition}[The cusp length and the generic twist]\hypertarget{tag-01SP}{}\label{inv:D4.2}
\StatementRecord{inv:D4.2}{proved}{01SP}{\SeriesRef{inv:C3.2}, \SeriesRef{inv:D4.3}, \SeriesRef{inv:D1.2}}

Use the coordinate $M_z=\exp(z\mathfrak g)$ of \ref{inv:C3.2}. For every
$z\notin\Z$ the cusp contribution is exactly
\begin{equation}
 (R^0f_*(T_X\otimes M_z))_{p_0}=0,\quad
 (R^1f_*(T_X\otimes M_z))_{p_0}=0,\quad
 (R^2f_*(T_X\otimes M_z))_{p_0}=\C\{q\}/(q).
\label{pub:I:eq:D-14}
\end{equation}
Outside the countable exceptional set explicitly defined in \ref{inv:D4.3},
these are the only nonzero direct-image stalks. Hence the generic
twist has $(h^0,h^1,h^2,h^3)=(0,0,1,0)$, and the cusp nilpotence exponent $r$ equals one there.
\end{proposition}


\begin{proof}
\begin{enumerate}[label=\arabic*.,leftmargin=*]
\item The map $H^1(W,\C)\to H^1(W,\OO_W)$ is an isomorphism by the
normalization complex in \ref{inv:D1.2}. The restriction of $M_z$ has trivial
pullback to $P$ and descent character
$(\exp(2\pi iz),1)$ on the two generators of the dual torus. The
normalization complex with these coefficients is the cellular
cochain complex of that rank-one local system. It may be computed
on the product cell structure as
\[
 \C\xrightarrow{(a-1,0)^t}\C^2
 \xrightarrow{(0,a-1)}\C,
 \qquad a=\exp(2\pi iz).
\]
For $z\notin\Z$ this is acyclic. All normalization components remain
rational with trivial pulled-back twist, so this normalization resolution computes
$H^q(W,M_z|_W)$ through the displayed cellular complex. Analytic semicontinuity and base change make
$R^q\bar f_*M_z$ zero on a neighbourhood of $p_0$ for all $q$.
The projection formula and \eqref{fam:eq:ab-log-tangent} imply the same vanishing for
$\mathcal T\otimes M_z$. Tensoring \eqref{fam:eq:ab-normal-quotient} remains exact, and
$\nu^*M_z=\OO_P$ because $P$ is rational and the topological first
Chern class of $M_z$ is zero. Thus the degree-two direct image is
exactly the reduced skyscraper $H^2(P,K_P)\otimes\C_{p_0}$.
The reduced support gives annihilator $(q)$ and proves \eqref{pub:I:eq:D-14}.

\item If the twist is nontrivial on every smooth fibre and on both finite
root-cover central tori, all cohomology of its fibre restriction
vanishes at those places. Indeed a nontrivial topologically trivial
line bundle on a complex torus has zero cohomology: in a unitary flat
presentation its Dolbeault Laplacian has no zero Fourier mode.
Tensoring the tangent extension by that line preserves acyclicity.
Base change on the root charts and exact invariants then give zero
direct images away from $p_0$. \ref{inv:D4.3} proves that the excluded twists
form a countable set. Leray now gives the asserted four dimensions.
\end{enumerate}
\end{proof}

\begin{proposition}[The exceptional nonintegral twists]\hypertarget{tag-01SQ}{}\label{inv:D4.3}
\StatementRecord{inv:D4.3}{proved}{01SQ}{\SeriesRef{inv:D4.2}}

On a regular marked fibre, $M_z$ is trivial if and only if there are
integers $n,p,q,r$ such that
\begin{equation}
 q\tau(t)+r\mu(t)=p,\qquad
 z=n-6q\mu(t)-r\beta_c(t).
\label{pub:I:eq:D-15}
\end{equation}
For $z\notin\Z$ these resonances occur in a countable set of parameter
values $z$; every resonant $z$ is nonreal. On the finite root covers
the exceptional lattices are
\begin{equation}
 \Lambda_1=\Z+(3\beta_c(z_1)+4-2\rho)\Z,\qquad
 \Lambda_2=\Z+(2\beta_c(z_2)+3-3i)\Z.
\label{pub:I:eq:D-16}
\end{equation}
Take the union of $\Z$, these two lattices, and all values from \eqref{pub:I:eq:D-15}
over regular fibres; this is the countable exceptional set used in
\ref{inv:D4.2}. The set is larger than $\Z$: the regular value $t=32/5$
supplies nonintegral twists with $h^1(X,T_X\otimes M_z)\ge2$.
\end{proposition}


\begin{proof}
\begin{enumerate}[label=\arabic*.,leftmargin=*]
\item The exponential sequence on a torus identifies trivial line bundles
with integral constant cocycles modulo the two holomorphic period
rows. Write
$z\gamma=(n,p,q,r)+a(6\mu,\tau,1,0)+b(\beta_c,\mu,0,1)$.
The third and fourth coordinates give $a=-q,b=-r$; the first two
give exactly \eqref{pub:I:eq:D-15}. If $r=0$, the positivity of $\operatorname{Im}\tau$
forces $q=p=0$, leaving $z=n$. Otherwise the imaginary parts of
\eqref{pub:I:eq:D-15} give
\begin{equation}
 \operatorname{Im}\mu=-\frac qr\operatorname{Im}\tau,
 \qquad\operatorname{Im}z=-rD\ne0.
\label{pub:I:eq:D-17}
\end{equation}
For fixed $(p,q,r)$ with $(q,r)\ne(0,0)$, the first equation of
\eqref{pub:I:eq:D-15} has discrete zeros, since $1,\tau,\mu$ are linearly
independent. A discrete subset of this second-countable Riemann
surface is countable. There are countably many integer tuples and
markings, so the union of the resulting $z$ is countable.

\item At $z_1$, insert $\tau=\rho$, $\mu=(2-\rho)/3$ in the first
equation. Its solutions have $(q,r,p)=(k,3k,2k)$.
At $z_2$ they have $(q,r,p)=(k,2k,k)$.
The second equation gives \eqref{pub:I:eq:D-16}. These are discrete lattices because
their nonintegral generators have imaginary parts $3D$ and $2D$.

\item For the explicit regular resonance, use the section
$P=(0,\sqrt2t^2(t-1))$ of the stated elliptic surface. The
chord-and-tangent law gives
\[
 x(2P)=9t^2/8,\qquad
 y(2P)=t^2(32-5t)/(16\sqrt2).
\]
At $t=32/5$, the discriminant
$1728t^8(t-1)^3$ is nonzero, $P$ is not two-torsion, and $2P$ is a
nonzero two-torsion point. Thus $P$ has exact order four. The period
identification in \cref{I:weierstrass-comparison} supplies a nonzero relation
$q\tau+4\mu=p$ with integers $p,q$. Equation \eqref{pub:I:eq:D-15}, with $r=4$,
defines a twist $M_z$ for which the regular torus restriction is
trivial; \eqref{pub:I:eq:D-17} shows that it is nonintegral. The period derivative
there has nonzero determinant, so the tangent restriction has
$h^0(T_X|_F)=2$. The degree-zero base-change exact sequence gives
\[
 H^0(F,(T_X\otimes M_z)|_F)
   \simeq (R^1f_*(T_X\otimes M_z))_t[t],
\]
because $f_*(T_X\otimes M_z)=0$ for a twist nontrivial on the general
fibre. Its right side has dimension two. The degree-one direct image
is a finite-support sheaf, so its global sections have dimension at
least two, proving the bound on $h^1$. The integers $p,q$ are the marked elliptic-logarithm coordinates of
this order-four point; the displayed period relation determines the
corresponding resonance lattice.
\end{enumerate}
\end{proof}

\begin{proposition}[Local complexes and every noncusp torsion length]\hypertarget{tag-01SR}{}\label{inv:D4.4}
\StatementRecord{inv:D4.4}{proved}{01SR}{\SeriesRef{inv:D4.1}, \SeriesRef{inv:A2.3}}

Let $z\notin\Z$. At a regular resonance $t_*$ in \eqref{pub:I:eq:D-15}, with
$A=\C\{t-t_*\}$, the two stalks are
\begin{equation}
 (R^1f_*(T_X\otimes M_z))_{t_*}=(A/(t-t_*))^2,
 \qquad (R^2f_*(T_X\otimes M_z))_{t_*}=A/((t-t_*)^2).
\label{pub:I:eq:D-18}
\end{equation}
At $p_1$, write $z=n+k(3\beta_1+4-2\rho)$, $k\ne0$, when this is
possible. Both coarse stalks are killed by $t$ and have length
$1+\epsilon_n$, with $\epsilon_n$ as in \ref{inv:D4.1}.
At $p_2$, write $z=n+k(2\beta_2+3-3i)$, $k\ne0$, when possible.
Both coarse stalks are killed by $t-1$ and have length one for
$n\equiv1,2\pmod4$ and length zero for $n\equiv0,3\pmod4$.
If the relevant lattice representation does not exist, both stalks
at that finite fibre are zero.
\end{proposition}


\begin{proof}
\begin{enumerate}[label=\arabic*.,leftmargin=*]
\item On a simply connected disc of smooth
torus fibres, including a finite root cover, retain the real marking
and the holomorphic period matrix. Translation by $\Pi(s)v$, for
fixed $v\in\R^4$, is holomorphic on the total family. After
subtracting the integer cocycle of a resonance, the line bundle has
constant lattice multipliers. Multiplying the lifted real translation
by its constant factor of automorphy makes the lattice act as the
identity and gives a genuine compact real torus action on the line
bundle. The same action acts on the tangent bundle by differentiation.
It commutes with the full Dolbeault differential.

\item Averaging over this compact torus projects the relative Dolbeault
complex onto its invariant part. Choose the lift so that the central
trivial line has trivial character. On every nonzero Fourier character,
the relative Dolbeault symbol is exterior multiplication by a
nonzero covector; contraction by a dual covector is a contracting
homotopy. After shrinking the disc, these nonzero covectors have a
uniform positive lower bound for bounded Fourier frequency and grow
linearly for unbounded frequency, because the real period matrices
and their inverses stay bounded. The homotopies therefore preserve
smooth Fourier series, with all derivatives, and are uniform on
compact subsets of the disc. The tangent extension is triangular
with respect to its vertical and horizontal parts; its nilpotent
off-diagonal term is treated by the finite homological perturbation
series, so every nonzero Fourier complex is still contractible.
These homotopies exhibit averaging as a quasi-isomorphism.

\item The chain-level comparison uses the following frames. Write the normalized period matrix as $[Z(s),I]$
and use the real coordinates $\zeta=Z(s)x+y$, with
$x,y\in\R^2/\Z^2$. Set
\[
 G_j=(dx_j)^{0,1},\qquad e_i=\partial_{\zeta_i},\qquad
 h=\partial_s+\sum_i(Z'(s)x)_i e_i.
\]
These are invariant under the real translations. Since each $dx_j$
is closed on the total real product, $\bar\partial G_j=0$.
The forms $G_j$ frame the relative invariant $(0,1)$ forms because
$\operatorname{Im}Z$ is invertible, and
\[
 \bar\partial h=\sum_{i,j}Z'_{ij}(s)e_iG_j,
 \qquad (dy)^{0,1}=-Z(s)(dx)^{0,1}.
\]
For the smooth line frame defined by its flat character, the last
identity gives $\bar\partial\ell=2\pi i(a_1G_1+a_2G_2)\ell$,
with $a_1,a_2$ in \eqref{pub:I:eq:D-20}; there is no additional base term.
All displayed coefficients are holomorphic in $s$. The full
invariant complex is thus the ordinary base Dolbeault resolution
of the following finite complex of holomorphic bundles:
\[
 C^j=(V\oplus L)\otimes\bigwedge^jG,\quad0\le j\le2,
 \qquad
 d(v\otimes\eta)=v\otimes a\wedge\eta,
\]
\begin{equation}
 d(h\otimes\eta)=h\otimes a\wedge\eta+\kappa\wedge\eta.
\label{pub:I:eq:D-19}
\end{equation}
Here $V$ is the rank-two vertical bundle, $L$ the horizontal line,
$G=R^1\OO$ on the root chart, and $\kappa\in V\otimes G\otimes L^*$
is the tangent-extension cocycle. The matrix is obtained
by differentiating the period translations. In the raw period and
constant-cocycle frames it is
\begin{equation}
 \kappa=\begin{pmatrix}6\mu'&\tau'\\\beta_c'&\mu'\end{pmatrix},
 \qquad
 a=(z-n+6q\mu+r\beta_c,-p+q\tau+r\mu).
\label{pub:I:eq:D-20}
\end{equation}
Let $\lambda=2\pi i$. The Dolbeault exponential convention gives
the differential with Picard vector $\lambda a$ and tangent term
$\kappa$. Its explicit chain isomorphism to \eqref{pub:I:eq:D-19} is, in degree $j$,
\[
 F^j|_{V\otimes\bigwedge^jG}=\lambda^{-j}\operatorname{id},
 \qquad
 F^j|_{L\otimes\bigwedge^jG}=\lambda^{-j-1}\operatorname{id}.
\]
Direct substitution gives $F^{j+1}d_{\lambda a,\kappa}
=d_{a,\kappa}F^j$ on both summands. The inverse multiplies them
by $\lambda^j$ and $\lambda^{j+1}$, respectively. These constant
maps also commute with the finite cyclic action. The displayed formula also follows directly by
subtracting the holomorphic period rows from the constant lattice
cocycle. Together these formulas identify the chain map and its tangent extension.
For the nonzero Fourier characters, use the filtration by the number
of base $(0,1)$ factors: the relative contractions make the associated
graded complexes acyclic, and the two-step filtration then proves
acyclicity of the full complexes. The finite cyclic group
normalizes the real translation group, so averaging and this
comparison are equivariant; exact invariants compute the coarse
multiple-fibre stalks.

\item At a regular resonance, $a'=\kappa^t(q,r)^t$ and
$(q,r)\ne(0,0)$. The determinant of $\kappa$ is nonzero there by
\cref{inv:A2.1}. Therefore $a$ has a simple zero. Choose its nonzero leading
direction as the second frame of $G$ and use invertibility of
$\kappa$ to choose the two vertical frames. Holomorphic invertible
changes reduce \eqref{pub:I:eq:D-19} to
\begin{equation}
 a=(0,s),\qquad \kappa=\begin{pmatrix}0&1\\1&0\end{pmatrix}.
\label{pub:I:eq:D-21}
\end{equation}
The second cohomology is generated by $h\gamma\delta$ and is killed
by $s^2$; the relations eliminate the two vertical generators.
The first cohomology is generated by $e_1\delta,e_2\delta$, each
killed by $s$. These are equalities of modules over $\C\{s\}$
and prove \eqref{pub:I:eq:D-18}.

\item There is an additional joint coefficient to check at the order-three
point. Put $t=s^3$, $\nu=\mu-(2-\tau)/3$ and
$b=\beta_c+4\mu+(2/3)\tau-b(0)$. The order-three transformation
laws give $\nu=\nu_2s^2+O(s^3)$ and $b=b_3s^3+O(s^4)$.
The elliptic calculation in \ref{inv:A2.3} gives, with elliptic scale
$f(s)=F s^2+O(s^3)$,
\[
 \nu_2^2=9F^2/512.
\]
Use the original Weierstrass normalization. Since $x=f^2\wp$,
$y=f^3\wp'/2$, one has
$f^4g_2=12t^3(t-1)$ and $f^6g_3=-8t^4(t-1)^2$.
The identities $g_2=4\pi^4E_4/3$, $g_3=8\pi^6E_6/27$ and
$E_4'(\rho)=-2\pi iE_6(\rho)/3$ give
\[
 E_6(\rho)=-\frac{27}{\pi^6F^6},\qquad
 \tau(s)=\rho+\tau_1s+O(s^2),\quad \tau_1=\frac{i\pi F^2}{2}.
\]
The period determinant identity in \cref{inv:A2.1} gives
\[
 K(s)=6\mu'^2-\tau'\beta_c'=K_2s^2+O(s^3),
 \quad K_2=\frac{3F^2}{4}.
\]
Equivalently $K=6\nu'^2-\tau'b'$, so
$K_2=24\nu_2^2-3\tau_1b_3$ and $b_3=7i/(32\pi)$.

\item For a nonintegral resonance the relation is
$(q,r,p)=\pm(k,3k,2k)$, $k\ne0$; the sign only changes the
following nonzero scalar. After a constant change of the two
cocycle frames, its Picard section is $3k(\nu,b)$, and the
transpose period derivative becomes
\[
 \begin{pmatrix}\tau'&\nu'\\6\nu'&b'\end{pmatrix}.
\]
Eliminating the unit $\tau'$ leaves the joint transverse component
$3k(b-6\nu\nu'/\tau')$. Its coefficient of $s^3$ is
\begin{equation}
 -\frac{k(K_2+12\nu_2^2)}{\tau_1}
       =\frac{123ik}{64\pi}\ne0.
\label{pub:I:eq:D-22}
\end{equation}
This nonzero joint coefficient determines the two transverse
orders in the following Smith presentation.

\item Choose the second cocycle frame to be the Picard section divided by
$s^2$. Equivariance and \eqref{pub:I:eq:D-22} give a matrix with a unit in its
lower-right entry, first-order off-diagonal terms with nonzero
coefficients, and a second-order upper-left term. Subtracting a
holomorphic multiple of the second vertical row removes the last
term; the nonzero determinant coefficient and \eqref{pub:I:eq:D-22} keep both
off-diagonal coefficients units. Rescaling the equivariant frames
therefore reduces \eqref{pub:I:eq:D-19} to
\begin{equation}
 a=(0,s^2),\qquad \kappa=\begin{pmatrix}0&s\\s&1\end{pmatrix}.
\label{pub:I:eq:D-23}
\end{equation}
Its cohomology is
\begin{equation}
 H^1=\C\{s\}/(s^2)\oplus\C\{s\}/(s^3),\qquad
 H^2=\C\{s\}/(s^2)\oplus\C\{s\}/(s^3).
\label{pub:I:eq:D-24}
\end{equation}
For example in degree one put
$u=e_1\delta$, $v=e_2\gamma+s h\delta$; the remaining cycle
$e_2\delta$ is eliminated by the boundary
$s u+e_2\delta+s v$. The two surviving relations are
$s^2u=0,s^3v=0$. In degree two eliminate $e_2\gamma\delta$;
the remaining generators $e_1\gamma\delta,h\gamma\delta$
have annihilators $s^2,s^3$ respectively.

\item A frame has invariant exponent $w$ when its coefficient satisfies
$g(\zeta s)=\zeta^w g(s)$; its pullback-frame weight is the negative
of $w$. For the trivialized twist in \eqref{pub:I:eq:D-16}, these exponents are
$n$ at $p_1$ and $3n$ at $p_2$. This agrees with $M_n$ in \ref{inv:D4.1}
and follows by evaluating its integer cocycle on $v_1,v_2$.
At $p_1$ the two generators in $H^1$ of \eqref{pub:I:eq:D-24} have invariant
exponents $n+2,n+2$, and the two generators in $H^2$ have exponents
$n+2,n$. The length-three summand has one invariant coefficient
for every $n$; the length-two summand has one precisely when
$n\not\equiv0\pmod3$. The coarse parameter is $s^3$, so all these
invariants are killed by it. This proves the asserted order-three
stalks with the full nilpotent layers included.

\item At $p_2$ the determinant is a unit. The Picard section has a simple
zero, so \eqref{pub:I:eq:D-21} applies on the order-four cover. The invariant
exponents of the two degree-one generators are $3n+1,3n+2$;
the cyclic degree-two module of length two has generator exponent
$3n+2$. In both degrees its invariant dimension is one exactly for
$n\equiv1,2\pmod4$. The coarse parameter $s^4$ kills the modules.
This proves the stated order-four calculation. \end{enumerate}
\end{proof}

\begin{proposition}[The full twisted cohomology and its jump loci]\hypertarget{tag-01SS}{}\label{inv:D4.5}
\StatementRecord{inv:D4.5}{proved}{01SS}{\SeriesRef{inv:D4.1}, \SeriesRef{inv:D4.2}, \SeriesRef{inv:D4.4}, \SeriesRef{inv:C3.1}, \SeriesRef{inv:C3.2}}

For $z\notin\Z$, let $\mathcal R_z$ be the set of regular base
points satisfying \eqref{pub:I:eq:D-15}, counted as points, not as marked integer
tuples. It is finite. Put $N_z=|\mathcal R_z|$. Define
$\ell_1(z)=0$ off $\Lambda_1$, and
$\ell_1(z)=1+\epsilon_n$ if
$z=n+k(3\beta_1+4-2\rho)$, $k\ne0$.
Define $\ell_2(z)=1$ if
$z=n+k(2\beta_2+3-3i)$, $k\ne0$, with
$n\equiv1,2\pmod4$, and $\ell_2(z)=0$ otherwise.
Then, setting $J_z=2N_z+\ell_1(z)+\ell_2(z)$,
\begin{equation}
 \boxed{(h^0,h^1,h^2,h^3)(T_X\otimes M_z)=(0,J_z,1+J_z,0).}
\label{pub:I:eq:D-25}
\end{equation}
More precisely, with $\mathfrak m_t$ the maximal ideal at $t$,
\[
 R^1f_*(T_X\otimes M_z)=
   \bigoplus_{t\in\mathcal R_z}\C_t^2
      \oplus\C_{p_1}^{\ell_1(z)}\oplus\C_{p_2}^{\ell_2(z)},
\]
\begin{equation}
 R^2f_*(T_X\otimes M_z)=\C_{p_0}\oplus
   \bigoplus_{t\in\mathcal R_z}\OO_{B,t}/\mathfrak m_t^2
      \oplus\C_{p_1}^{\ell_1(z)}\oplus\C_{p_2}^{\ell_2(z)}.
\label{pub:I:eq:D-26}
\end{equation}
Together with \ref{inv:D4.1} these formulas cover every element of
$\operatorname{Pic}(X)=\C$. The cohomology jump loci are
locally discrete subsets of this complex line.
\end{proposition}


\begin{proof}
\begin{enumerate}[label=\arabic*.,leftmargin=*]
\item The support assertions and all stalks were proved in \ref{inv:D4.2}--\ref{inv:D4.4}.
The direct image in degree zero is zero for a nonintegral twist:
it has generic rank zero and is torsion-free, since multiplication
by a nonzero base function is injective on the bundle on the smooth
integral total space. The support of either higher coherent direct
image is a proper closed analytic subset of the compact curve,
and hence is finite. This proves finiteness of $\mathcal R_z$ as
well as the absence of an accumulation at the special values.
If two integer tuples at a marked fibre represented the same
trivialized line, their difference would be an integral real class
in $H^{1,0}$, which is zero. Thus there is no multiplicity from
duplicate tuples. Formula \eqref{pub:I:eq:D-26} follows by the decomposition of a
finite-support sheaf into its local modules. Its global sections
have dimension equal to its length and its $H^1$ is zero, so Leray
gives \eqref{pub:I:eq:D-25}.

\item The $\bar\partial$-closed representative
$G_\gamma=G_0/12$ of \cref{inv:C3.1,inv:C3.2} defines a holomorphic
line bundle on $X\times\C$ by the integrable operators
$\bar\partial-2\pi i zG_\gamma$. Proper
coherent base change for the projection to $\C$ makes each
cohomology-dimension threshold an analytic subset. A real noninteger
$z$ has no resonance by \eqref{pub:I:eq:D-17} and \eqref{pub:I:eq:D-16}, so the generic values are
attained. All possible exceptional values were shown to form a
countable set. Thus each threshold exceeding a generic
dimension is a proper analytic subset of a complex curve and is
locally discrete. The jump set is the union of the four thresholds
$h^0\ge1$, $h^1\ge1$, $h^2\ge2$ and $h^3\ge1$, hence is locally
discrete. The order-four congruences $n\equiv0,3\pmod4$
have zero contribution in \eqref{pub:I:eq:D-25}.

\item The jump loci can also be indexed algebraically on the elliptic surface. If a regular fibre
has torsion section $P$ of exact order $r_t$, choose its marked
relation $q_t\tau+r_t\mu=p_t$. Its full resonance lattice is
\[
 \Z+\bigl(6q_t\mu(t)+r_t\beta_c(t)\bigr)\Z.
\]
Every regular resonance arises this way. The underlying values
$t$ are exactly the non-special roots of the division polynomials
$\psi_N(P)$, $N\ge2$. The nonintegral lattice points contribute two
to $J_z$ for each distinct such fibre. The finite-point lattices
have the congruence weights already stated. The division-polynomial recursion gives the first nonconstant
regular factors $5t-32$, $89t^2+4736t-4096$ and
$11t^2-6016t+8192$. The displayed lattice formula sends every
such root to its resonance values in the fixed Picard coordinate.

\item For the index, $c_1(X)=c_2(X)=c_1(M)=0$ in integral
cohomology, so the degree-six Riemann--Roch integrand for the
twisted tangent bundle is $c_3/2$. Its integral is one. Both \eqref{pub:I:eq:D-13}
and \eqref{pub:I:eq:D-25} have alternating sum one. 
\end{enumerate}
\end{proof}








\begin{proposition}[Local separation and Kodaira dimension]\hypertarget{tag-0335}{}\label{push:src:E.2.2}
\StatementRecord{push:src:E.2.2}{proved}{0335}{\SeriesRef{push:src:E.1.2}, \SeriesRef{push:src:E.2.1}}

The multiplicity pair $(3k,4l)$ is constant on every sufficiently
small deformation of one of the preceding members.  Distinct
pairs therefore determine disjoint open loci among these
isomorphism classes.  Moreover
\[
 \kappa(X)=\begin{cases}-\infty,&kl<3k+4l,\\1,&kl>3k+4l.
 \end{cases}
\]
Equality $kl=3k+4l$ never occurs under $\gcd(3k,4l)=1$.
\end{proposition}


\begin{proof}
\begin{enumerate}[label=\arabic*.,leftmargin=*]
\item Local miniversality identifies every small deformation with the
fixed-data period family; \ref{push:src:E.1.2} makes its multiplicity pair an
intrinsic biholomorphic invariant.  The canonical formula of
\ref{push:src:E.2.1} and connected fibres identify its pluricanonical sections
with those of $\mathcal O(r(kl-3k-4l))$ on the root curve.
Negative weight has no sections.  Positive weight has linearly
growing section counts along multiples of $12kl$, giving
Kodaira dimension one.  Equality would imply $k\mid4l$;
coprimality forces $k=1$, and then $l=3+4l$ is impossible.
\end{enumerate}
\end{proof}

\begin{proposition}[Separation in irreducible analytic families]\hypertarget{tag-033B}{}\label{push:src:E.2.3}
\StatementRecord{push:src:E.2.3}{proved}{033B}{\SeriesRef{push:src:E.2.1}; \cite[Chapter~III, Theorem~4.1, p.~124]{BS76}}

Suppose a proper holomorphic submersion over a connected
irreducible reduced complex base has two fibres of the type
in \ref{push:src:E.2.1}.  Then their multiplicity pairs agree.
For these members the asymptotic number
\[
 s(X)=\lim_{r\to\infty}
 \frac{h^0(X,K_X^r)-h^0(X,K_X^{-r})}{r}
 =\frac1{12}-\frac1{3k}-\frac1{4l}
\]
recovers the ordered multiplicity pair uniquely.
\end{proposition}


\begin{proof}
\begin{enumerate}[label=\arabic*.,leftmargin=*]
\item For coprime $m,n$, the number of nonnegative solutions of
$mi+nj=N$ is $N/(mn)+O(1)$ as $N\to\infty$:
there is one permitted residue class of $j$ modulo $m$, and
its representatives lie between $0$ and $N/n$.
The canonical formula in \ref{push:src:E.2.1} proves the displayed limit,
with a negative value when the anticanonical counts grow.
Write $1/12-s(X)=a/b$ in lowest positive terms.
Since $\gcd(m+n,mn)=1$, we have $a=m+n$, $b=mn$.
Thus $m,n$ are the two roots of $T^2-aT+b$.
Their order is determined because only $m$ is divisible by
three and only $n$ is divisible by four.

\item For every fixed integer $r$, apply the finite-free-complex
theorem of \cite[Chapter~III, Theorem~4.1, p.~124]{BS76} to
the relative canonical power, which is an invertible flat
sheaf for a holomorphic submersion.  Its zeroth cohomology
dimension has analytic jumping loci.  On an irreducible
base any value constant on a nonempty open set is its
generic value: a proper analytic jumping locus has empty
interior.  Local miniversality in \ref{push:src:E.2.1} makes these dimensions
constant near either of the two given fibres, since all
nearby fibres have the same multiplicities and their
canonical formula is independent of the period parameter.
Consequently the two fibres have equal dimensions for
every positive and negative $r$.  Their limits $s(X)$
are equal, and the preceding arithmetic recovers the same
pair.  
\end{enumerate}
\end{proof}

% Compendium-only push:src:E.3.2: apparatus: historical/current dependency audit of the five credited DGKX version1 theorems; complete mathematical applications retained at their actual receiving results and in the credited-result appendix



% Compendium-only push:src:E.3.3: apparatus: exact scope comparison for the five DGKX version1 standard(3,4) results; the source-specific receiving results are separately retained



\section{The integral circle quotient and its cusp replacement}\hypertarget{tag-02OO}{}\label{mapsI:circle-quotient}


\begin{definition}[Circle-quotient conventions]\hypertarget{tag-IP06}{}\label{I:def:circle-quotient-conventions}
\StatementRecord{I:def:circle-quotient-conventions}{proved}{IP06}{}
We write $g,u,w,d$ for the ordered homology basis
$(\widehat\gamma,\widehat u,\widehat w,\widehat\delta)$ in this
section, and $\gamma$ for its first dual covector. The cusp coordinate is
$q=1/t$.  Positive target orientations are the complex orientation
on the base followed by the positive $\gamma$-circle orientation.
The complex orientation on $X$ fixes the source orientation.
\end{definition}

\begin{proposition}[The primitive integral quotient and its finite fillings]\hypertarget{tag-02PA}{}\label{mapsI:src:A.1.1}
\StatementRecord{mapsI:src:A.1.1}{proved}{02PA}{\cite{AC2026}}

Use the monodromy matrices of \eqref{wb:pshift:eq:finite-data} and the
standard affine vectors of Definition~\ref{I:def:family-parameter}. Put
\[
 M_0=(A_1A_2)^{-1}=
 \begin{pmatrix}1&0&0&0\\0&1&0&0\\0&1&1&0\\-1&0&0&1\end{pmatrix}.
\]
The common invariant vectors and covectors are $\mathbb Zd$ and
$\mathbb Z\gamma$, respectively.  The smallest invariant primitive
sublattice containing the cusp vanishing lattice is
$L=\mathbb Z\langle u,w,d\rangle=\ker\gamma$.
The formula
\begin{equation}
 [z,\Pi(z)x]\longmapsto([z],\gamma(x)\bmod\mathbb Z)
 \label{mapsI:quotient-eq-1}
\end{equation}
defines a real-analytic fibrewise group epimorphism
$J\longrightarrow B^\circ\times S^1$ with connected kernel $T_L$.
At the finite fillings it descends to a smooth proper $T^3$-bundle
over a Seifert solid torus, of multiplicities $3$ and $4$ and
translation residues $1$ and $-1$.
\end{proposition}



\begin{proof}
\begin{enumerate}[label=\arabic*.,leftmargin=*]
\item For the matrices of \eqref{wb:pshift:eq:finite-data}, direct integral
multiplication gives
$A_1^3=A_2^4=A_1A_2M_0=I$ and $A_jv_j=v_j$.
Solving the two fixed-vector equations gives $(0,0,0,d_0)$.
For a row vector $(a,b,c,e)$, the first fixed-covector equation
gives $b=2c,e=3c$ and the second gives $b=c,e=2c$;
therefore $b=c=e=0$.  Also
\[
 (M_0-I)(a,b,c,e)=(0,0,b,-a),\qquad
 A_1w=u-w.
\]
These identities prove the assertions about $L$ and the cusp
lattice $K=\mathbb Z\langle w,d\rangle$.  Their primitivity follows
from the displayed bases.  The flat real marking identifies the
deck action on a fibre with $x\mapsto A_jx$.  Since
$\gamma A_j=\gamma$, \eqref{mapsI:quotient-eq-1} descends.  Its kernel is connected
because $\gamma$ is primitive.

\item On a root disc the affine generator is
\[
 (s,x)\longmapsto
 (e^{-2\pi i/m}s,A_jx+v_j/m),\qquad (m,\kappa)=(3,1),(4,-1).
\]
It induces $(s,\theta)\mapsto
(e^{-2\pi i/m}s,\theta+\kappa/m)$ on the character quotient.
This action is free since $\gcd(\kappa,m)=1$.  Writing
$a=e^{2\pi i\theta}$ and choosing $r$ with $r\kappa\equiv1\pmod m$,
the invariant coordinates $s a^r$ and $a^m$ identify its quotient
with $D^2\times S^1$.  The base coordinate becomes
$(s a^r)^m(a^m)^{-r}$, and hence has multiplicity $m$ at its
central circle.  Before the finite quotient the map is a trivial
$T_L$-bundle.  The free actions on source and target are
equivariant for that map, so it descends locally to a smooth
$T_L$-bundle.  Properness follows from compactness of $T_L$.
The logarithmic translation used to glue the root chart to $J$
changes the character by its $\gamma$-component.  Applying
$\gamma$ to the gluing translation therefore gives the
target gluing, proving agreement on every overlap.
\end{enumerate}
\end{proof}


\begin{remark}[Comparison of the quotient characters]\hypertarget{tag-IP08}{}\label{I:rem:circle-character-comparison}
\StatementRecord{I:rem:circle-character-comparison}{proved}{IP08}{}

Engel's primitive invariant covector $\psi$
\cite[Lemma~3.8, p.~11]{Engel2026} is related to the present character by
\begin{equation}
 P_E=P_{\rm mark},\qquad \psi P_E=-\gamma,
 \hypertarget{tag-02OZ}{}\label{mapsI:quotient-eq-2}
\end{equation}
where $P_{\rm mark}$ is the integral isomorphism of
Proposition~\ref{inv:B1.0}. Thus the quotient characters agree after the
target-circle map $\theta\mapsto-\theta$. Inverting the root generators
gives the translation vectors $a_{j-1}=-P_Ev_j/m_j$, as in
Proposition~\ref{fam:src:B.2.2}.
\end{remark}

\begin{proposition}[Pointwise cusp limits and the fixed double curve]\hypertarget{tag-02PB}{}\label{mapsI:src:A.1.2}
\StatementRecord{mapsI:src:A.1.2}{proved}{02PB}{\SeriesRef{mapsI:quotient-eq-4}}

The literal character has no continuous extension at either
triple point or at generic points of two of the three double
curves.  It has the constant limit $0$ on the smooth locus of
$W$ and on the third double curve away from the triple points.
In the chart $q=z_0z_1z_2$, $x_1=z_1$, $x_2=z_2$, the extendible
double curve is $z_0=z_1=0$, $z_2\ne0$.  The fixed curve of the
$\mathbb C^*$-action in direction $d=e_2$ is instead
$z_0=z_2=0$, $z_1\ne0$, and is one of the nonextendible curves.
Consequently the extendible double curve and the fixed double curve are distinct.
\end{proposition}



\begin{proof}
\begin{enumerate}[label=\arabic*.,leftmargin=*]
\item Put $q=e^{-2\pi T}$ on a positive radial ray and use the exact
cusp periods
\[
 Z_T=iTB_0+C(q),\qquad
 B_0=\begin{pmatrix}0&1\\-1&0\end{pmatrix},
\]
where $C$ is bounded and holomorphic at $0$.  For a fixed
$a\in\mathbb R^2$ the representative
$x_T(a)=\exp(2\pi iZ_Ta)$ has character $a_1\bmod\mathbb Z$ and
logarithmic position
\begin{equation}
 \frac{\log|x_T(a)|}{\log|q|}
 =B_0a+T^{-1}\operatorname{Im}C(q)a.
 \hypertarget{tag-02P0}{}\label{mapsI:quotient-eq-3}
\end{equation}
The choices $a=(-1/4,1/4)$ and $a'=(-1/2,1/4)$ have logarithmic
positions tending to $(1/4,1/4)$ and $(1/4,1/2)$.
Both are interior to the triangle with vertices $0,e_1,e_2$,
so all three chart coordinates tend to $0$.  Their character
values are $3/4$ and $1/2$.  This proves nonextension at the
first triple point.  The other triangle type has vertices
$e_1,e_2,e_1+e_2$.  Its coordinates can be taken as
$x_1x_2/q,q/x_2,q/x_1$.  Positions $(3/4,3/4)$ and
$(3/4,1/2)$ lie in its interior, give distinct values of $-y_2$
modulo one, and converge to its origin.  Translation by the
lattice preserves triangle type; these are precisely the two
triple points of $W$.

\item For arbitrary convergent paths in the first chart, inversion
of the exact logarithmic period matrix gives
\begin{equation}
 \theta_\gamma=-\frac{\log|z_2|}{\log|q|}+o(1)\pmod{\mathbb Z}.
 \hypertarget{tag-02P1}{}\label{mapsI:quotient-eq-4}
\end{equation}
Here the error tends to zero whenever the logarithmic ratios
stay bounded; they do stay bounded near a fixed point of a
normal-crossing stratum because $q=z_0z_1z_2$ and the three
coordinates are bounded.  If $z_0,z_1\to0$ and $z_2\to c\ne0$,
the right side tends to zero.  If $z_0,z_2\to0$ and $z_1=c\ne0$,
the paths $z_2=q^a,z_0=q^{1-a}/c$, $0<a<1$, have limits
$-a\pmod1$ at the same double point.  The same construction
with $z_0=c$ proves nonextension on $z_1=z_2=0$.
At a smooth point of a component the ratio in \eqref{mapsI:quotient-eq-4} tends
to $0$ or $1$, giving zero modulo one.  Integer translates
of these charts give the stated assertion on all of $W$.

\item The torus action in direction $e_2=d$ is
\[
 \lambda\cdot(z_0,z_1,z_2)
       =(\lambda^{-1}z_0,z_1,\lambda z_2).
\]
Its fixed curve is $z_0=z_2=0$, while the character extends along
$z_0=z_1=0$ by the preceding limit formula. The displayed action and
limit formula give these two loci in the same toric chart.
\end{enumerate}
\end{proof}

\begin{lemma}[The cusp neighbourhood and its circle classes]\hypertarget{tag-IP07}{}\label{I:lem:cusp-fundamental-group}
\StatementRecord{I:lem:cusp-fundamental-group}{proved}{IP07}{}

For the cusp neighbourhood $N_0$ and the vanishing lattice
$K=\mathbb Z\langle w,d\rangle$,
\begin{equation}
 \pi_1(N_0)=\overline\Lambda=\Lambda/K.
 \hypertarget{tag-02P2}{}\label{mapsI:quotient-eq-5}
\end{equation}
The punctured-cusp map to this group kills $K$ and the base meridian.
Both long-period covectors extend in integral $H^1(N_0;\mathbb Z)$.
The literal character on an ordinary fibre has no extension to all of
$X\simeq S^6$.
\end{lemma}
\begin{proof}
The toric covering has rays $(v,1)$ for all $v\in\mathbb Z^2$.
Meridians about its toric divisors kill the whole cocharacter lattice:
$(0,1)$ is a ray and $(e_i,1)-(0,1)=(e_i,0)$. Hence the toric cover is
simply connected. Its free period deck group is
$\overline\Lambda=\Lambda/K$, proving the displayed isomorphism.
The base meridian bounds the disc given by a point of the dense orbit
on a central component. Both long-period covectors kill these
relations and therefore extend in integral $H^1$.
An extension of $\theta_\gamma$ to all of $X\simeq S^6$ would give a
nonzero homomorphism $H_1(F)\to\mathbb Z$ factoring through
$H_1(X)=0$, which is impossible.
\end{proof}

\begin{proposition}[The filled three-dimensional target]\hypertarget{tag-02PC}{}\label{mapsI:src:A.1.3}
\StatementRecord{mapsI:src:A.1.3}{proved}{02PC}{the definitions and constructions stated in the proof}

For the original data the quotient over the complement of
the cusp disc is the exterior of the $(3,4)$ torus knot,
with the specified Seifert structure.  For general admissible
$(l_0,l_1,l_2)$ the filled quotient is the lens space determined
by the following meridian slopes in one oriented integral
torus basis $(\rho_1,h)$:
\begin{equation}
 \mu_1=(3,-l_1),\qquad
 \mu_2=(-4,4l_0-l_2),\qquad
 \det(\mu_1,\mu_2)=p=12l_0-4l_1-3l_2.
 \label{mapsI:quotient-eq-6}
\end{equation}
In the lens-space convention that the second meridian is
$-q_0\mu_1+p\lambda_1$, where
$\det(\mu_1,\lambda_1)=1$, the lens parameter satisfies
$3q_0\equiv4\pmod p$.  In particular the target depends only
on the signed integer $p$, and is a three-sphere exactly
when $|p|=1$.
\end{proposition}



\begin{proof}
\begin{enumerate}[label=\arabic*.,leftmargin=*]
\item Let $P$ be the three-holed sphere, with oriented boundary
meridians $\rho_0+\rho_1+\rho_2=0$ and positive fibre $h$.
The two finite meridian relations are
$3\rho_1-l_1h=0$ and $4\rho_2-l_2h=0$; the cusp solid-torus
meridian is $\rho_0+l_0h$.  Eliminating $\rho_0$ gives
$\rho_2=l_0h-\rho_1$ on the annulus joining the finite solid
tori, which gives exactly \eqref{mapsI:quotient-eq-6}.  Since both slopes are
primitive, this is a genus-one Heegaard splitting.  Choose
$a,b\in\mathbb Z$ with $al_1+3b=1$ and put
$\lambda_1=(a,b)$.  Direct multiplication gives
\begin{equation}
 \mu_2=-q_0\mu_1+p\lambda_1,
 \quad q_0=a(4l_0-l_2)+4b,
 \quad 3q_0-ap=4.
 \hypertarget{tag-02P3}{}\label{mapsI:quotient-eq-7}
\end{equation}
This proves the claimed lens parameter in a convention fixed
by the displayed oriented meridians.  The fundamental
group of this gluing is $\mathbb Z/|p|$.  Admissibility implies
$\gcd(p,6)=1$, so the modular inverse of $3$ used above exists.

\item For $(l_0,l_1,l_2)=(0,1,-1)$ the two exceptional-fibre data
are $(3,1),(4,-1)$.  They are realized by the circle action
\[
 s\cdot(z_1,z_2)=(e^{4is}z_1,e^{3is}z_2)
 \quad\text{on } S^3\subset\mathbb C^2.
\]
Its invariant map is $[z_1^3:z_2^4]$; the local circle
translation residues are $4^{-1}=1\pmod3$ and
$3^{-1}=-1\pmod4$.  Removing a regular orbit gives the
same two Seifert solid tori and their product annulus as
above. A regular orbit is the $(3,4)$ torus knot under the coordinate
exchange $(z_1,z_2)\mapsto(z_2,z_1)$ when the opposite ordering is used.
The displayed orbit map and meridian slopes identify the manifold
with its peripheral structure.
\end{enumerate}
\end{proof}

\subsection{The cusp replacement and its relative uniqueness}

\begin{theorem}[The cusp replacement and its relative homotopy class]\hypertarget{tag-02PD}{}\label{mapsI:src:A.2.1}
\StatementRecord{mapsI:src:A.2.1}{proved}{02PD}{\SeriesRef{mapsI:quotient-eq-5}}

For every admissible member, there is a smooth map
$e_{\rm q}:X\to Q$ over $B$ whose restriction outside an arbitrarily
small cusp neighbourhood is the literal circle quotient.
For fixed quotient gluing, its homotopy class is independent
of the cusp replacement, provided the replacement maps the
cusp neighbourhood into its designated solid torus and is
kept fixed on its outer collar.  In particular there is a
well-defined class for the original member $X\simeq S^6$.
\end{theorem}



\begin{proof}
\begin{enumerate}[label=\arabic*.,leftmargin=*]
\item On the cusp boundary use the product coordinates
$(q,\theta_\gamma)$ after the specified collar shear.
By \eqref{mapsI:quotient-eq-5}, the boundary circle class extends to $N_0$.
Since $S^1=K(\mathbb Z,1)$, it has a continuous representative
$\theta':N_0\to S^1$.  Its restriction to the boundary is
homotopic to the given one: the two maps induce the same
homomorphism on $\pi_1$ of the boundary mapping torus.
The homotopy-extension property for a manifold with boundary
changes $\theta'$ on a collar so that it agrees exactly with
$\theta_\gamma$ there.  Relative smooth approximation gives
a smooth representative, still equal on a smaller collar.
Then $(f,\theta')$ is the required map into $D^2\times S^1$,
and agrees with the finite-filling maps on the overlap.
Both components of $(f,\theta')$ are smooth on $N_0$, so the
product map is smooth at the critical points of $f$ as well.

\item The inclusion $\partial N_0\to N_0$ surjects on $\pi_1$,
as \eqref{mapsI:quotient-eq-5} explicitly shows.  It follows that
$H^1(N_0;\mathbb Z)\to H^1(\partial N_0;\mathbb Z)$
is injective.  Since both spaces are connected, the long
exact sequence gives $H^1(N_0,\partial N_0;\mathbb Z)=0$.
Two circle replacements agreeing on the boundary therefore
have quotient whose relative $H^1$ class is zero, hence
have a relative homotopy.  Keeping the base coordinate
fixed gives the required homotopy of maps into the solid
torus.  Different choices of the collar size are compared
on a common smaller collar.  This proves the assertion.
\end{enumerate}
\end{proof}




\section{Real structures and their fixed three-spheres}\hypertarget{tag-02P9}{}\label{mapsI:real-forms}

\hypertarget{tag-01PJ}{}\label{fam:section:B.4}

\begin{definition}[Real-period data]\hypertarget{tag-IP04}{}\label{I:def:real-period-data}
\StatementRecord{I:def:real-period-data}{proved}{IP04}{}
Use the lattice, period matrix and clockwise root generators of
Definition~\ref{def:common-conventions}, and the affine vectors of
Definition~\ref{I:def:family-parameter}. In this section write
$d=\widehat\delta$, $g=\widehat\gamma$, $u=\widehat u$ and
$w_\Lambda=\widehat w$. An unadorned $w$ in a lattice expression denotes
$w_\Lambda$; the scalar parameter remains $w=e^{-\pi ic}$. Put
\[
 e=\Pi_c d=(0,1)^{\mathsf t},\qquad E=d\otimes\gamma,
 \qquad M_0=(A_1A_2)^{-1},\qquad D_c=D_{\rm part}+\Im c<0.
\]
The circle $\rho(s)$ is translation by $se$, $s\in\mathbb R/\mathbb Z$.
In cusp coordinates $q=1/t$ and $s=(2\pi i)^{-1}\log q$, define the
bounded matrix $C_c(q)$ by $Z_c=sB_0+C_c(q)$, where
$B_0=\left(\begin{smallmatrix}0&1\\-1&0\end{smallmatrix}\right)$.
Thus $\Pi_cg=(6\mu,-s+C_{c,21})^{\mathsf t}$.
The cusp double curves $D_1,D_2,D_{12}$ belong to the edge directions
$e_1,e_2,e_1-e_2$. The action $\rho$ fixes $D_2$, while the literal
character has a limit along $D_1$, as located by the explicit toric
coordinates in Proposition~\ref{mapsI:src:A.1.2}.
\end{definition}

\begin{proposition}[Conjugate members and antiholomorphic automorphisms]\hypertarget{tag-01QG}{}\label{fam:src:B.4.2}\label{fam:st:B.4.2}
\StatementRecord{fam:src:B.4.2}{proved}{01QG}{\SeriesRef{fam:src:B.4.1}, \SeriesRef{fam:src:A.3.1}, \SeriesRef{fam:src:B.2.1}, \SeriesRef{fam:src:A.5.1}}

Use the normalization $(\beta^{\rm part}+\tau)(p_0)=0$. There is an antiholomorphic diffeomorphism $X_c\to X_{c'}$ exactly when
\[
 c'-(1-\bar c)\in2\mathbb Z,
 \qquad\text{equivalently }w'=-\bar w.
\]
Consequently $\overline{X_c}\simeq X_{1-\bar c}$. The members admitting antiholomorphic self-maps are exactly those with $\operatorname{Re}c\in\tfrac12+\mathbb Z$, equivalently $w=-\bar w$. On each such member the group of holomorphic and antiholomorphic diffeomorphisms is
\[
 \operatorname{Aut}^{\pm}(X_c)\simeq\mathbb C^*\rtimes\mathbb Z/2,
 \qquad z\longmapsto1/\bar z
\]
for the action of the second factor. Its antiholomorphic maps form the nonidentity coset. For all other members the full group is the holomorphic automorphism group of Theorem~\ref{fam:src:B.4.1}.
\end{proposition}


\begin{proof}
\begin{enumerate}[label=\arabic*.,leftmargin=*]
\item Lift conjugation of the real orbifold base to the antiholomorphic involution $\iota$ of its universal cover whose cusp branch has $\tau(\iota z)=-\overline{\tau(z)}$. A lift exists by the universal covering property; it extends at the finite orbifold points, whose local equations are $s_j^{m_j}=t_j$. Choose the usual paths from a real point $t>1$, so reflection fixes the second centre and sends the first chosen lift to $g_2^{-1}z_1$. The induced involution of the deck group is
\[
 \phi(g_1)=g_2^{-1}g_1^{-1}g_2,\qquad
 \phi(g_2)=g_2^{-1},\qquad \phi(g_0)=g_0^{-1}.
\]
These identities also follow by reflecting the lifted peripheral paths; their product and their squares agree with $g_1g_2g_0=1$. The cusp branch fixes the lift uniquely.

\item Denote the affine period transformations by $G_1,G_2$ on the triple $(\tau,\mu,\beta)$ and put $K(\tau,\mu,\beta)=(-\bar\tau,1-\bar\mu,-\bar\beta)$. Direct rational substitution gives
\[
 KG_1K=G_2^{-1}G_1^{-1}G_2,\qquad KG_2K=G_2^{-1}.
\]
For example $G_2^{-1}(\tau,\mu,\beta)=(-1/\tau,(1-\mu)/\tau,\beta+3-6(1-\mu)^2/\tau)$ and
$G_1^{-1}(\tau,\mu,\beta)=(1/(1-\tau),1-\mu/(1-\tau),\beta-2+6\mu^2/(1-\tau))$. These formulas verify every constant in the two equalities. The reflected $\mu$ is cusp bounded and has the same affine laws, so their difference is a global section of $\mathcal O_B(-p_0)$ and vanishes. After this equality, the difference between the reflected and original $\beta^{\rm part}$ is a holomorphic invariant base function, hence constant; its cusp value is zero under the stated normalization. Therefore
\[
 \mu(\iota z)=1-\bar\mu(z),\qquad
 \beta^{\rm part}(\iota z)=-\overline{\beta^{\rm part}(z)}.
\]
The original $\tau$ identity follows from its $j$ equation and its specified cusp branch. In particular $D_{-\bar c}(\iota z)=D_c(z)$, so this comparison stays in the exact admissible domain.

\item Let
\begin{equation}
 H=\begin{pmatrix}1&0&0&0\\0&1&0&0\\-6&0&-1&0\\0&-1&0&-1\end{pmatrix}.
 \label{I:eq:real-reflection-matrix}
\end{equation}
Then $H^2=I$, $HA_2H=A_2^{-1}$, $HA_1H=A_2^{-1}A_1^{-1}A_2$, and
\[
 \Pi_{-\bar c}(\iota z)H=-\overline{\Pi_c(z)}.
\]
Thus $(z,\zeta)\mapsto(\iota z,-\bar\zeta)$ is a well-defined antiholomorphic torus map. It also descends under the middle deck action. Indeed covariance and the displayed lattice conjugacies imply $R_{\phi(g)}(\iota z)=\overline{R_g(z)}$ by multiplying the full real-invertible period matrices. Thus the reflection intertwines the middle deck actions.

\item At the first finite centre compose this map with the target deck identification $g_2$ to return the centre $g_2^{-1}z_1$ to $z_1$. Its lattice map becomes $A_2H$ and it conjugates the clockwise generator to the inverse generator. At the second centre its lattice map is $H$ with the same reversal. Consequently the target affine vectors are exactly
\begin{equation}
 w'_1=-A_2H\varepsilon=-\varepsilon+d,\qquad
 w'_2=-H(-\varepsilon')=\varepsilon'-3d.
 \label{I:eq:real-reflected-affine-data}
\end{equation}
The minus signs arise because the inverse of a generator with invariant translation $w'_j/m_j$ has translation $-w'_j/m_j$. This calculation verifies the full affine actions on the root covers. For the overlaps choose the root coordinates of the real local base coordinates $t$ and $t-1$ compatible with these local reflections. This choice has an explicit comparison with any original root coordinate: their ratio is a holomorphic unit, so changing it changes the collar by $(2\pi i)^{-1}\log h\,\Pi w_j$, a translation extending on the entire root family. Its inverse gives a holomorphic isomorphism of the two glued descriptions. Thus the extension property and final classification transfer to the original coordinates. Then $\log s'_j=\overline{\log s_j}$ on paired branches. Under $\zeta\mapsto-\bar\zeta$, conjugation of $1/(2\pi i)$ supplies the compensating sign, and the identity of the two logarithmic collar translations follows exactly. Two roots of the same fixed real coordinate differ by an $m_j$th root of unity; this changes the comparison only by a constant period-coordinate translation, and the additive correction of \ref{fam:src:A.3.1} glues those constants. The preceding holomorphic-unit comparison handles a different original local base coordinate. Its inverse has the same construction.

\item At the cusp choose $s=(2\pi i)^{-1}\log t_c$, so $s\mapsto-\bar s$. The full holomorphic unit matrix satisfies
\[
 C_{-\bar c}(\bar t_c)=-\overline{C_c(t_c)}+
       \begin{pmatrix}6&0\\0&1\end{pmatrix}.
\]
The integral extra matrix has trivial exponential. Hence $(x,t_c)\mapsto(\bar x,\bar t_c)$ intertwines the full cusp deck groups, with the same integral index $\lambda$. The toric fan is defined over $\mathbb Z$, so this map extends antiholomorphically on every toric chart and stratum. It preserves the untwisted cusp collar. We have therefore constructed an antiholomorphic map to the target datum with parameter $-\bar c$, $r'_1=-1,k'_1=1,r'_2=1,k'_2=-3,l'_0=0$, whose $p'$ is $+1$.

\item Apply \ref{fam:src:B.2.1} to normalize this target back to the canonical $p=-1$ datum. Take $a=-1$ and the odd shift $n=1$. This gives a holomorphic map to $X_{1-\bar c}$. Every other antiholomorphic map, composed with the inverse of this one, is a holomorphic comparison between canonical members. \ref{fam:src:A.5.1} gives exactly the stated congruence. Exponentiating it gives $w'=-\bar w$.

\item  For a real member, choose any resulting antiholomorphic self-map $J$. Its derivative on ordinary fibre coordinates is $+\text{complex conjugation}$: the initial reflection has derivative $-\text{complex conjugation}$, while the normalization just used has linear part $-I_2$. Hence it conjugates translation by $b\in\mathbb C/\mathbb Z$ to translation by $\bar b$. Its square is a holomorphic translation $b_0$. Commuting $J$ with $J^2$ shows $b_0=\bar b_0$ in $\mathbb C/\mathbb Z$. This equality forces the imaginary part of a representative of $b_0$ to vanish, so take $b_0\in\mathbb R$. Replacing $J$ by the composition of translation by $-b_0/2$ with $J$ changes its square to $b_0-b_0/2-b_0/2=0$. We have an antiholomorphic involution. Under $b\mapsto e^{2\pi ib}$, its conjugation action is $z\mapsto1/\bar z$. This proves the split group and specifies its antiholomorphic coset.
\end{enumerate}
\end{proof}

\begin{proposition}[Two conjugacy classes of real structures]\hypertarget{tag-02PE}{}\label{mapsI:src:A.3.1}
\StatementRecord{mapsI:src:A.3.1}{proved}{02PE}{\SeriesRef{fam:src:B.4.2}, \SeriesRef{fam:src:B.4.1}; \cite[Theorem~3.1, p.~48]{Illman1972}; \cite[Theorem~4.1(1), p.~10]{Knudson2004}}

For a real member the antiholomorphic involutions have exactly two
conjugacy classes under holomorphic automorphisms. For any fixed
involutive splitting $\sigma$ they are represented by $\sigma$ and
$\rho(1/2)\sigma$. Every nonempty fixed set is a closed connected
orientable totally real three-manifold with the mod-two homology of
$S^3$. The circle $\rho$ preserves each such fixed set and acts
semifreely on it.
\end{proposition}



\begin{proof}
\begin{enumerate}[label=\arabic*.,leftmargin=*]
\item Proposition~\ref{fam:src:B.4.2} identifies the full holomorphic group
with $\mathbb C^*$ and gives the conjugation action $z\mapsto1/\bar z$.
Thus the square of $z\sigma$ is $z/\bar z$, and it is an involution
exactly for $z\in\mathbb R^*$. Conjugation by $y\in\mathbb C^*$ changes
this coefficient to $|y|^2z$, giving exactly the two sign classes.

\item At a fixed point an antiholomorphic involution has a real three-dimensional
fixed tangent space, and its complexification is the full tangent space.
An invariant metric and its exponential map identify its fixed locus
locally with this real subspace. It is therefore a closed totally real
three-manifold. The smooth $\mathbb Z/2$ action admits an equivariant
triangulation by \cite[Theorem~3.1, p.~48]{Illman1972}; compactness and finite
dimension give the finite-dimensional $G$-CW hypothesis. Smith's theorem,
in the form of
\cite[Theorem~4.1(1), p.~10]{Knudson2004}, applies to the smooth finite action on
the sphere and gives the stated mod-two homology. A nonempty
three-dimensional fixed set is consequently connected, and its first
mod-two cohomology vanishes; its orientation class $w_1$ therefore
vanishes. On the unit circle $1/\bar z=z$, so $\sigma$ commutes with
$\rho$. Its stabilizers are inherited from the semifree action on $X$ and are
therefore either trivial or the whole circle.
\end{enumerate}
\end{proof}

\begin{definition}[Representatives of the real involutions]\hypertarget{tag-IP05}{}\label{I:def:real-representatives}
\StatementRecord{I:def:real-representatives}{proved}{IP05}{}

Put
\begin{equation}
 L_F=(2\pi i)^{-1}\log F,\qquad L_G=(2\pi i)^{-1}\log G,
 \qquad F^3=t,\quad G^4=t-1.
 \hypertarget{tag-02P4}{}\label{mapsI:real-eq-1}
\end{equation}
The roots are positive over $t>1$. We continue from that interval
through the upper half $t$-plane. Their logarithms are taken on the
orbifold uniformizing cover. Write $\iota$ for its reflection lift.
Every point of $-i\R_{>0}$ has a representative $c=1/2+iy$, and every
point of $i\R_{>0}$ has a representative $c=-1/2+iy$, with $y$ in the
admissible range. Set
\begin{equation}
\begin{array}{c|c|c}
 &\alpha&S(t)\\\hline
 c=1/2+iy&1/3&2L_F-4L_G-1/6\\
 c=-1/2+iy&1/2&-2L_G+1/4
\end{array}
\qquad b_c(z)=\alpha\Pi_c(z)g+S(t(z))e.
\label{mapsI:real-eq-2}
\end{equation}
\end{definition}

\begin{theorem}[Explicit real involutions and all fixed fibres]\hypertarget{tag-02PF}{}\label{mapsI:src:A.3.2}
\StatementRecord{mapsI:src:A.3.2}{proved}{02PF}{\SeriesRef{fam:src:B.4.2}, \SeriesRef{mapsI:real-eq-4}, \SeriesRef{mapsI:real-eq-2}, \SeriesRef{fam:src:A.3.1}, \SeriesRef{mapsI:real-eq-1}}

On the ordinary period charts the formulas
\begin{equation}
 \sigma_c[z,\zeta]=[\iota z,\bar\zeta+b_c(\iota z)],
 \qquad \sigma_c^- =\rho(1/2)\sigma_c
 \label{mapsI:real-eq-3}
\end{equation}
extend to antiholomorphic involutions of the compact threefold.
With these representatives their complete fixed-set data are as follows.
Here $2T^2$ means the disjoint union of two real two-tori; every displayed
circle in a double curve avoids its two triple-point endpoints.
\begin{center}\small
\begin{tabular}{cclccc}
\toprule
$w$&involution&image in $\RP^1$&on $S_1$&on $S_2$&on $W$\\
\midrule
$-ir$, $r>0$&$\sigma_c$&all $\RP^1$&$2T^2$&$2T^2$&$T^2\sqcup S^1_{D_2}\sqcup S^1_{D_{12}}$\\
$-ir$, $r>0$&$\sigma_c^-$&$[1,\infty]$&$\varnothing$&$T^2$&$S^1_{D_1}\sqcup S^1_{D_2}$\\
$ir$, $r>0$&$\sigma_c$&$\varnothing$&$\varnothing$&$\varnothing$&$\varnothing$\\
$ir$, $r>0$&$\sigma_c^-$&$[\infty,1]$ through $0$&$2T^2$&$T^2$&$S^1_{D_1}\sqcup S^1_{D_{12}}$\\
\bottomrule
\end{tabular}
\end{center}
The first interval is the closure of $t>1$; the second is the closure
of $t<1$. On every open real arc contained in the indicated image the
real regular fibre is $2T^2$. No triple point is real for these structures.
\end{theorem}



\begin{proof}
\begin{enumerate}[label=\arabic*.,leftmargin=*]
\item Start with the reflection $(z,\zeta)\mapsto(\iota z,-\bar\zeta)$
and the integral matrix $H$ of \eqref{I:eq:real-reflection-matrix}.
Proposition~\ref{fam:src:B.4.2} proves its period compatibility and gives
the target affine datum \eqref{I:eq:real-reflected-affine-data}.
The normalization from this
target at parameter $c_0=-\bar c$ to the canonical datum at $c_0+1$
has linear part $-I$ and translation
\begin{equation}
 b^{\rm norm}_{c_0}=\tfrac13\Pi_{c_0}g+(2L_F-4L_G)e.
 \hypertarget{tag-02P5}{}\label{mapsI:real-eq-4}
\end{equation}
Its period defects at the two finite generators are
$(0,-2,2,0)^t$ and $(0,0,2,0)^t$. To check them, the logarithmic increments
are $-2/3$ and $1$, and one multiplies the two displayed matrices $A_j$
by $g/3$. The finite difference vectors in the complex translation space are $2e$ and $-4e$, so subtracting
their logarithmic collars leaves holomorphic translations on both full
root discs. At the cusp
$2L_F-4L_G=s/3-(2\pi i)^{-1}\log(1-q)$ and
$\Pi g=(6\mu,-s+C_{21})^t$. The logarithmic terms cancel and give the
holomorphic toric unit
$(2\mu,C_{21}/3-(2\pi i)^{-1}\log(1-q))$.
Thus \eqref{mapsI:real-eq-4} extends over all three fillings, with its inverse.

\item For $c=1/2+iy$ compose this normalization with the initial reflection
and then with the constant translation $e/6$. This gives exactly the
first line of \eqref{mapsI:real-eq-2}. For $c=-1/2+iy$, additionally close the parameter
shift by the following explicit even-shift formula with $h=-2$,
\[
 b_{h,c_0}=-\frac h{12}\Pi_{c_0}g+
 \left(h(L_F-L_G)-\frac{h^2}{24}\right)e.
\]
A final constant translation $-e/4$ gives the second line of \eqref{mapsI:real-eq-2}.
The displayed translation represents the even shift of Proposition~\ref{fam:src:A.3.1}. Its finite-chart corrections are holomorphic. For its cusp expression, put $\ell=L_F-L_G=-s/12+\kappa_0(q)$ and
$C_{c,21}(q)=c+C_{0,21}(q)$, and set
$D_0(q)=C_{0,21}(q)-12\kappa_0(q)$.
The explicit even-shift translation becomes
\[
 b_{h,c}=
 \left(-h\mu/2,-h(c+D_0)/12-h^2/24\right)^t.
\]
The logarithmic terms have cancelled exactly. Exponentiating both
coordinates gives holomorphic nonzero torus multipliers on the whole cusp
disc; multiplication by these units extends over every toric stratum.
Compatibility with the period deck group is the integral identity
$\exp(2\pi i h\lambda_1)=1$ for $h\in2\mathbb Z$ and
$\lambda_1\in\mathbb Z$. The finite-chart defects are checked by the
same displayed $A_j$ and root logarithms used for the normalization.

\item Over $t>1$ we have
$\Re\mu=1/2$, $\Re\tau=\Re\beta^{\rm part}=0$ and the logarithms in
\eqref{mapsI:real-eq-1} are purely imaginary. The two squares are respectively the
integral period translations $2\widehat w$ and $3\widehat w$.
The square is a holomorphic automorphism and is consequently the
identity everywhere. This proves the involution claim.

\item  Put $a=\Re c$ and
$R=-H+2aE$. In the chosen branch the reflection identifications on
$t>1$, $0<t<1$, $t<0$ are respectively $1,g_2,g_0^{-1}$.
They follow by reflecting the chosen upper-half-plane paths; equivalently
the root arguments are $(0,0)$, $(0,\pi/4)$ and $(\pi/3,\pi/4)$.
The three fibre involution matrices are
\begin{equation}
 R_>=R,\qquad R_m=A_2^{-1}R,\qquad R_<=M_0R.
 \hypertarget{tag-02P6}{}\label{mapsI:real-eq-5}
\end{equation}
In marked real coordinates let $r_c(t)=\Pi_c(t)^{-1}b_c(t)$,
where $\Pi_c^{-1}$ is the real inverse. The fixed tori are exactly
\begin{equation}
 \{x\in\R^4/\Lambda:(I-R_\bullet)x=
                       r_c(t)+\epsilon d/2\pmod\Lambda\},
 \qquad\epsilon=0,1.
 \label{mapsI:real-eq-6}
\end{equation}
This specifies the cosets, including their parameter dependence. Their
identity subgroups are $T_{\langle w,d\rangle}$ on both outer arcs and
$T_{\langle u+w,d\rangle}$ on the middle arc. When nonempty there are
two cosets. Relative to any one solution their offsets are
$(g+u)/2$ on the right arc and $g/2$ on the other two arcs.

\item Existence of fixed points is decided by integral affine congruences. An affine involution
$x\mapsto R_\bullet x+\beta$ has a fixed point if and only if its
square translation $(I+R_\bullet)\beta$ belongs to
$(I+R_\bullet)\Lambda$. In \eqref{mapsI:real-eq-2} the square translation is
\[
 (I+R_\bullet)\alpha g+2\Re S\,d.
\]
For $c=1/2+iy$ the values $2\Re S$ on the three arcs are
$-1/3,-4/3,-2/3$; for $c=-1/2+iy$ they are $1/2,0,0$.
The resulting integral square vectors are respectively
$(2w,0,2w)$ and $(3w,d,3w)$. Adding $d/2$ to the involution adds $d$
to each square vector. Taking column Hermite normal forms of
$I+R_\bullet$ gives exactly the nonempty entries in the table.
All three $I-R_\bullet$ have nonzero Smith factors $1,2$.
This proves the coset counts and offsets. The Hermite and Smith
computations are listed in \cref{tab:I:programs}. The left-arc lift
$g_0^{-1}$ follows from the reflected upper-half-plane path above.

\item At $t=0$ the root degree is odd. Conjugation by the degree-three deck
group acts transitively on the three reflection lifts; the centralizer
of a reflection in that deck group is trivial.
Thus the quotient has exactly the same two real tori as the adjacent
nonempty real regular fibres. In the empty cases there are none. This
also follows by solving the limiting affine equations on the full
multiplicity-three root quotient.

\item At $t=1$ write the affine deck map as
$G(x)=A_2x+v_2/4$. The real root involution has matrix $R_>$ and
translation
\[
 b_-=(1/3,0,0,-1/6)^t,
 \qquad b_+=(1/2,0,0,1/4)^t
\]
for the two parameter signs. The real locus of the quotient is the
quotient of the fixed sets of $G^kh$, $0\leq k<4$. Its even and odd
lift parities are identified internally by $G$, while $G^2$ preserves
each parity. Direct integer congruences show that the nonempty lift
parities are: both for negative imaginary $w$ and $\epsilon=0$; only even for
negative imaginary $w$ and $\epsilon=1$; neither for positive imaginary $w$ and $\epsilon=0$;
only odd for positive imaginary $w$ and $\epsilon=1$. Every nonempty lift has
two tori, exchanged by $G^2$. Each parity therefore contributes one
torus downstairs. The congruence enumeration uses denominator $24$:
translation denominators divide $12$, and the nonzero Smith factors are
$1,2$, so every component has such a representative. The component
permutation is included in the exact enumeration listed in
\cref{tab:I:programs}. The normal
real line has fourth power of sign positive for the even parity and
negative for the odd parity. This establishes the two one-sided
endpoints at $t=1$ recorded in the table.

\item  Put
$A=C_{12}(0)$; the cusp normalization gives
\[
 \Re\mu(0)=1/2,\quad A\in i\R,\quad C_{21}(0)=c-A,
 \quad C_{11}(0)=6\mu(0),\quad C_{22}(0)=\mu(0).
\]
In a toric lift the involution is
$x\mapsto U/\bar x$, with
\begin{equation}
\begin{array}{c|cc}
 &b_1(0)&b_2(0)\\\hline
 w=-ir&2\mu(0)&C_{21}(0)/3-1/6\\
 w=ir&3\mu(0)&C_{21}(0)/2+1/4
\end{array}
\qquad U_j=e^{2\pi i b_j(0)}.
\hypertarget{tag-02P7}{}\label{mapsI:real-eq-7}
\end{equation}
Thus $U_2>0$ in both rows, whereas $U_1>0$ in the first and $U_1<0$
in the second. Replacing $\sigma_c$ by $\sigma_c^-$ negates $U_2$.
The smooth central orbit has a real torus precisely when both units
are positive, with equations $|x_j|^2=U_j$.

\item For an edge $v$, set $n=(-v_2,v_1)^t$. The period translate bringing
the inverted edge back is $B_0n=v$. Its curve coordinate $z=x^n$ has
involution
\begin{equation}
 z\longmapsto K_v/\bar z,
 \qquad K_v=U^n e^{2\pi i n^tC(0)n}.
 \hypertarget{tag-02P8}{}\label{mapsI:real-eq-8}
\end{equation}
The signs for $(D_1,D_2,D_{12})$ are $(-,+,+)$ for negative $w$
and $(-,-,-)$ for positive $w$. Negating $U_2$ reverses the first and
third signs, but not the second. A positive sign gives the single
circle $|z|^2=K_v$; a negative sign gives no fixed point. Central
inversion exchanges the two triangle classes of the periodic fan, so
it exchanges the two triple points. These calculations prove every
cusp entry. Finally a nonempty fixed three-manifold cannot be confined
to the special fibres: a holomorphic function constant on an open
piece of a totally real submanifold is locally constant. This also
proves that the positive-$w$ representative with no regular real
fibres is entirely fixed-point free.
\end{enumerate}
\end{proof}

\begin{theorem}[The real three-manifolds]\hypertarget{tag-02PG}{}\label{mapsI:src:A.3.3}
\StatementRecord{mapsI:src:A.3.3}{proved}{02PG}{\SeriesRef{mapsI:src:A.3.2}}

Every nonempty fixed set in Theorem~\ref{mapsI:src:A.3.2} is diffeomorphic to $S^3$.
Its fundamental group is trivial and its integral homology is
$(\Z,0,0,\Z)$. For negative imaginary $w$, the restricted circle action
is the standard semifree action on $S^3$ with fixed circle
$D_2^{\sigma_c}$. For positive imaginary $w$, the nonempty real form
has a free circle action isomorphic, up to circle and orientation
conventions, to the Hopf action.
\end{theorem}



\begin{proof}
\begin{enumerate}[label=\arabic*.,leftmargin=*]
\item For negative imaginary $w$ both fixed manifolds meet $D_2$ in one circle.
Their circle action is semifree, so the orbit space is a compact
orientable surface with boundary. The local quotient at a fixed circle
is the quotient of $S^1\times\C$ by scalar rotation on $\C$, hence is
a boundary collar. Over the surface interior the principal circle bundle
is trivial because the surface has boundary and $H^2=0$. Changes of
boundary trivialization extend over the attached normal discs by
rotation. Thus the three-manifold is
\[
 (B\times S^1)\ \cup_{\partial B\times S^1}\ (\partial B\times D^2).
\]
If $B$ has genus $g_B$ and $b_B$ boundary components, van Kampen or its
surface handle decomposition gives first homology
$\Z^{2g_B+b_B-1}$. Smith's theorem makes this rank zero. Since $b_B>0$,
we have $g_B=0,b_B=1$. The displayed manifold is then
$\partial(D^2\times D^2)=S^3$, with the indicated circle action. This
derives the smooth recognition from an explicit surface decomposition.

\item For positive imaginary $w$ the nonempty real manifold is free under
$\rho$. The two real tori on $t<1$ join at the single real torus over
$1$. At $0$ the odd root comparison is a torus product. At infinity
they are capped by the real neighbourhoods of $D_1$ and $D_{12}$.
These are solid tori: the primitive circle in the edge direction
collapses and the remaining real circle is their core. In the outer
real-fibre lattice $\Z\langle w,d\rangle$ the two cap meridians are
$w$ and $w-d$.

\item  Put
$B_1=A_2^{-1}A_1^{-1}A_2$ and $C=B_1^2$. The two real-arc involutions
satisfy $R_<R_m=B_1$, $R_mB_1R_m=B_1^{-1}$ and $B_1^3=I$.
Moving between the real root lifts requires
$B_1^kR_mB_1^{-k}=B_1R_m$, so $2k\equiv1\pmod3$ and $k=2$.
Thus the odd-root transport is $C$. It satisfies $R_<=CR_mC^{-1}$, and
$C(u+w)=-w+d$. The exchange across the multiplicity-four real torus
is $A_2^2$; after this odd-root transport it is
$CA_2^2C^{-1}$, whose restriction to $\langle w,d\rangle$ is
$w\mapsto-w$, $d\mapsto d$. Consequently the meridians, compared on
one boundary torus, have determinant of absolute value one. The
resulting genus-one Heegaard splitting is $S^3$: primitive meridians
intersecting once can be taken to the two coordinate circles by an
integral torus automorphism, giving the usual two solid tori of $S^3$.
The free circle $d$ then has Euler number of absolute value one, so
the displayed bundle over $S^2$ is the Hopf bundle. This proves the
fundamental-group and integral-homology assertions.
\end{enumerate}
\end{proof}

\begin{proposition}[Restrictions of the circle quotient]\hypertarget{tag-02PH}{}\label{mapsI:src:A.3.4}
\StatementRecord{mapsI:src:A.3.4}{proved}{02PH}{\SeriesRef{mapsI:src:A.2.1}, \SeriesRef{mapsI:real-eq-2}}

The map $e_{\rm q}$ of Theorem~\ref{mapsI:src:A.2.1} has a $\rho$-invariant
cusp replacement. Its restriction to
every nonempty real form is null-homotopic. On an ordinary real fibre
its image consists of the two quotient-circle values
\begin{equation}
 \theta=\tfrac12\gamma(r_c(t)),\qquad
 \theta=\tfrac12\gamma(r_c(t))+\tfrac12\pmod\Z,
 \quad
 \gamma(r_c(t))=\alpha+\frac{\Im S(t)}{D_c(t)},
 \label{mapsI:real-eq-9}
\end{equation}
with $\alpha,S$ from \eqref{mapsI:real-eq-2}. The circle action is along each real
two-torus and is annihilated by $e_{\rm q}$.
\end{proposition}



\begin{proof}
\begin{enumerate}[label=\arabic*.,leftmargin=*]
\item The literal quotient is circle invariant. In a cusp extension
$\theta':N_0\to S^1$, each orbit circle has degree zero: its class is
$d$, which vanishes in $\pi_1(N_0)$. The quotients
$\theta'(\rho(s)x)/\theta'(x)$ therefore have a continuous real lift
vanishing at $s=0$ and periodic in $s$. Averaging this real lift over
$s\in\R/\Z$ changes $\theta'$ to a circle-invariant map. The cocycle
identity proves invariance, and on an invariant outer collar the lift
is zero, so the replacement agrees exactly with the prescribed map.
The same construction is smooth. Thus $e_{\rm q}|_{\operatorname{Fix}\sigma}$
factors through its circle orbit space, which is $D^2$ or $S^2$ by the
preceding proof. Both have only two-dimensional homotopy type and
$S^3$ is two-connected, so this restriction is null-homotopic.
The first row of $I-R_\bullet$ is $(2,0,0,0)$, giving the two values
in \eqref{mapsI:real-eq-9}. The real inverse of the period matrix sends $S e$ to
first coordinate $\Im S/D_c$; this follows by subtracting
$(\Im\mu/\Im\tau)$ times its first imaginary-coordinate equation
from the second. Both real-fibre identity subgroups have zero
$\gamma$-coordinate, proving the remaining assertions.
\end{enumerate}
\end{proof}




\begin{proposition}[The coherent closing translation]\hypertarget{tag-0322}{}\label{push:src:C.3.5}
\StatementRecord{push:src:C.3.5}{proved}{0322}{\SeriesRef{fam:src:A.3.1}; \citet[Proposition~7.8, pp.~91--92, and Proposition~C.2, p.~102]{DuEtAl2026v1}}

Fix the standard twists and the root and cusp markings of Proposition~\ref{fam:src:A.3.1}.
There is a coherent choice of the even-shift maps
\(\widetilde F_{h,c}:X_c\to X_{c+h}\), \(h\in2\mathbb Z\),
with the following explicit ordinary-chart translation. Choose global
roots \(F^3=t\), \(G^4=t-1\) on the uniformizing cover and put
\(\ell=(2\pi i)^{-1}\log(F/G)\pmod{\mathbb Z}\), \(e=(0,1)^t\).
Then the translation is
\begin{equation}\hypertarget{tag-02OP}{}\label{mapsI:closing:eq:C-explicit-closing}
 b_{h,c}=-\frac h{12}\Pi_c\widehat\gamma
              +\left(h\ell-\frac{h^2}{24}\right)e.
\end{equation}
For \(h=12\), it is the scalar multiplication
\begin{equation}\hypertarget{tag-02OQ}{}\label{mapsI:closing:eq:C-twelve-scalar}
 K(t)=\frac{t^4}{(t-1)^3}
\end{equation}
on the ordinary \(\mathbb C^*\)-coordinate, after removing an integral
period. On the finite root charts its scalar representatives are
\(K(t)s_1^{-12}\) and \(K(t)s_2^{12}\), respectively, both units.
In a cusp coordinate \(q=t_c\), the scalar representative is
\(q^{-1}H(q)\), with \(H\) a holomorphic nonvanishing function.
The same map has a holomorphic \((\mathbb C^*)^2\)-unit representative
on the entire cusp filling. These two representatives differ by
exactly the \(\widehat\gamma\) cusp deck transformation.

For six positive \(w\)-loops the closing map is
\(\widetilde F_{12,c_0-12}\); thus the stated cusp exponent is \(-1\)
with $w=e^{-\pi i c}$ and positive loops in the $w$-plane. For parallel transport $P:X_{c_0}\to X_{c_0-12}$ along this lifted path, the resulting six-loop diffeomorphism is $\widetilde F_{12,c_0-12}\circ P$.
\end{proposition}

\begin{proof}
The even-shift translation and its quadratic correction are those of \citet[Proposition~7.8, pp.~91--92, and Proposition~C.2, p.~102]{DuEtAl2026v1}. Their representatives in the fixed root and cusp charts are as follows.
The global roots exist because the zero orders of \(t\) and \(t-1\)
on the simply connected uniformizing surface are divisible by three
and four. The local characters and their coprime orders give
\[
 \ell(g_1z)-\ell(z)=-\frac13,\qquad
 \ell(g_2z)-\ell(z)=\frac14,\qquad
 \ell(g_0z)-\ell(z)=\frac1{12}\pmod{\mathbb Z}.
\]
A change of logarithm adds an integral \(h\widehat\delta\) period,
so \eqref{mapsI:closing:eq:C-explicit-closing} is a well-defined torus translation.
Since \(\Pi_{c+h}=\Pi_c(I+hE)\), the two ordinary lattices agree;
the source-to-target lattice map of the identity linear part is
\(I-hE\).

For compatible logarithm lifts, covariance gives the affine defects
\[
 b_{h,c}(g_jz)-R_j(z)b_{h,c}(z)
       =\Pi_c(g_jz)\eta_j,\qquad
 \eta_1=\frac h2(\widehat u-\widehat w-\widehat\delta),\quad
 \eta_2=\frac h2(-\widehat w+\widehat\delta).
\]
Their \(\gamma\)-coordinates vanish, so the same right sides use
\(\Pi_{c+h}\). The vectors are integral for even \(h\). At the
cusp the defect is zero: \((A_0-I)\widehat\gamma=-\widehat\delta\)
cancels the increment of \(\ell\). These equations prove ordinary
and cusp descent with the stated orientation of the peripheral
loops.

For the finite comparison put \(w_j=(2\pi i)^{-1}\log s_j\).
The functions \(F/(s_1G)\) and \(Fs_2/G\) are invariant
holomorphic units on the respective full root discs. Choose
holomorphic logarithms \(\kappa_j\) of them. Then
\[
 \ell=\epsilon_jw_j+\kappa_j\pmod{\mathbb Z},\qquad
 \epsilon_1=1,\quad\epsilon_2=-1.
\]
The exact finite-to-ordinary comparison satisfies
\(\sigma_{j,c+h}-\sigma_{j,c}=h\epsilon_jw_je\).
Consequently the required finite-chart translation is
\[
 b^{(j)}_{h,c}
 =-\frac h{12}\Pi_c\widehat\gamma
       +\left(h\kappa_j-\frac{h^2}{24}\right)e.
\]
It is holomorphic on the full root disc and satisfies
\[
 b^{(j)}_{h,c}(g_jz)-R_j(z)b^{(j)}_{h,c}(z)
       -\frac{h\epsilon_j}{m_j}e
 =\Pi_c(g_jz)\eta_j.
\]
This is precisely equivariance for the source and target affine
finite actions, including the changed period column. The exact
identity
\(\sigma_{j,c+h}+b^{(j)}_{h,c}
=\sigma_{j,c}+b_{h,c}\) proves overlap agreement. 

At the cusp let \(q=e^{2\pi i s}\) and
\(C_{c,21}=\beta^{\rm part}+c+s\), which is holomorphic at
\(q=0\). The function
\[
 \widehat H=e^{2\pi i s/12}F/G
 \quad\hbox{satisfies}\quad
 \widehat H^{12}=q\frac{t^4}{(t-1)^3}.
\]
The right side is a unit because \(t\) has a simple pole at the
cusp. The character formula makes \(\widehat H\) single-valued;
it therefore extends as a nonvanishing holomorphic function. Write
\(\ell=-s/12+\kappa_0(q)\), with \(\kappa_0\) holomorphic.
Since
\(\Pi_c\widehat\gamma=(6\mu,-s+C_{c,21})^t\),
the logarithmic terms cancel in \eqref{mapsI:closing:eq:C-explicit-closing}:
\[
 \bar b_{h,c}(q)=
 \left(-\frac h2\mu(q),\,
 -\frac h{12}C_{c,21}(q)+h\kappa_0(q)-\frac{h^2}{24}\right)^t.
\]
Exponentiating gives a holomorphic unit in \((\mathbb C^*)^2\).
On every toric chart its action multiplies a character coordinate
by an integral character of this unit, so it and its inverse
extend across every double curve and both triple points. The
source and target cusp lattice multipliers differ by
\((1,e^{2\pi i h\lambda_1})=(1,1)\), for
\(\lambda\in\mathbb Z^2\). They are therefore identical. The
unit multiplication commutes with their toric shears and descends.
Exponentiation of the displayed formula gives exactly the ordinary
translation on the punctured cusp. 

The uncorrected translations have composition defect
\(-hk e/12\). The constants in
\eqref{mapsI:closing:eq:C-explicit-closing} cancel it, since
\((h+k)^2-h^2-k^2=2hk\). Hence
\[
 \widetilde F_{k,c+h}\widetilde F_{h,c}
       =\widetilde F_{h+k,c},\qquad
 \widetilde F_{0,c}=\mathrm{id}.
\]
Composing this even-shift map with the inverse of the fixed choice in Proposition~\ref{fam:src:A.3.1} gives an automorphism over the identity of the base. The full automorphism theorem, Theorem~\ref{night:P5.5}, identifies it with a scalar in $\mathbb C^*$. A path from that scalar to $1$ supplies an isotopy, so the two closing choices define the same smooth monodromy class.

For \(h=12\), remove the period
\(-\Pi_c\widehat\gamma\) from \eqref{mapsI:closing:eq:C-explicit-closing}.
The remaining translation is \((12\ell-6)e\). Its exponential
is \((F/G)^{12}=K(t)\), proving
\eqref{mapsI:closing:eq:C-twelve-scalar}. The finite corrections divide it by
\(s_1^{12}\) and multiply it by \(s_2^{12}\). With the original
uniformizers, these are units because
\(t/s_1^3\) and \((t-1)/s_2^4\) are units. 

Put \(H(q)=qK(t(q))=e^{24\pi i\kappa_0(q)}\). The cusp unit
for \(h=12\) is exactly
\[
 U_{12,c}(q)=
 \left(e^{-12\pi i\mu(q)},\,
       e^{-2\pi i C_{c,21}(q)}H(q)\right).
\]
The \(\widehat\gamma\) deck transformation has multiplier
\[
 \left(e^{12\pi i\mu(q)},\,
       e^{2\pi i C_{c,21}(q)}q^{-1}\right),
\]
because \(B_0(1,0)^t=(0,-1)^t\). Multiplying these two
expressions gives \((1,q^{-1}H)=(1,K)\). This proves the
claimed comparison, including its sign, on the entire cusp
quotient. The extending map is represented by $U_{12,c}$, and the displayed deck multiplier identifies it with the scalar representative on the punctured chart.

Finally \(w=e^{-\pi i c}\) makes six positive loops the path
\(c(r)=c_0-12r\); its endpoint is closed by the positive shift
\(h=12\). For parallel transport $P$ along that path the closed diffeomorphism is therefore $\widetilde F_{12,c_0-12}\circ P$. The overlap difference between the scalar logarithm and the two-coordinate unit logarithm is exactly $\Pi_c\widehat\gamma$. Scaling both logarithms by a real number $a$ changes this difference to $a\Pi_c\widehat\gamma$. For $0<a<1$, this vector is not a period: the real period map is an isomorphism and $a\widehat\gamma\notin\Lambda$. Thus the separately scaled chart maps fail the overlap equation at precisely this fractional period.
\end{proof}

\begin{figure}[htbp]
\centering
\begin{tikzpicture}[x=1cm,y=1cm,>=stealth,font=\small]
\draw[thick] (-5.3,0) .. controls (-1.8,-.35) and (1.8,-.35) .. (5.3,0);
\node at (0,-1.05) {$B=\mathbb P^1_t$};
\foreach \x/\lab in {-4/0,0/1,4/\infty}{
 \fill (\x,-.1) circle (2pt);
 \node[below=5pt] at (\x,-.1) {$\lab$};
 \draw[->] (\x,1.15)--(\x,.15);
}
\node[draw,rounded corners,align=center,text width=2.8cm,minimum height=1.2cm] at (-4,1.75)
 {$3S_1$\\free order-three\\root quotient};
\node[draw,rounded corners,align=center,text width=2.8cm,minimum height=1.2cm] at (0,1.75)
 {$4S_2$\\free order-four\\root quotient};
\node[draw,rounded corners,align=center,text width=2.8cm,minimum height=1.2cm] at (4,1.75)
 {$W$\\normalization $\mathrm{dP}_6$};
\draw[->] (-2,1.15)--(-2,.15);
\node[align=center] at (-2,2.85) {ordinary fibre\\$\mathbb C^2/\Pi_c\mathbb Z^4$};
\node[align=center] at (4,3.05) {three paired boundary curves\\two triple points};
\end{tikzpicture}
\caption{The fixed $(3,4,\infty)$ base data. The arrows denote the
restriction of $f$ to each indicated fibre, with its scheme
multiplicity. The ordinary period quotient, finite affine
quotients and the paired hexagon are the three local models
in the common conventions; their overlap maps are fixed there.}
\hypertarget{tag-01RL}{}\label{fig:I:orbifold}
\end{figure}

\begin{figure}[htbp]
\centering
\begin{tikzpicture}[x=1cm,y=1cm,>=stealth,font=\small]
\fill[blue!4] (0,0) circle (2.6);
\draw[thick,dashed] (0,0) circle (2.6);
\draw[->,gray] (-3.1,0)--(3.2,0) node[right] {$\Re w$};
\draw[->,gray] (0,-3.15)--(0,3.2) node[above] {$\Im w$};
\draw[very thick,blue!65!black] (0,-2.56)--(0,-.13);
\draw[very thick,blue!65!black] (0,.13)--(0,2.56);
\fill[white] (0,0) circle (.11);
\draw[thick] (0,0) circle (.11);
\draw[->] (.23,-.18)--(1.1,-.7);
\node[anchor=west,align=left] at (1.05,-.88) {puncture $w=0$\\corner winding $-pn$};
\node[align=center] at (-1.45,.8) {$0<|w|<216$};
\node[anchor=west,align=left] at (3.15,1.65)
 {outer circle $|w|=216$\\$|q_2|=1$\\$q_2=iw/216$};
\draw[->] (3.05,1.55)--(2.15,1.43);
\node[anchor=east,align=right] at (-.25,1.8)
 {real structures\\$w\in i\mathbb R$};
\node[anchor=west,align=left] at (3.15,-2.1)
 {all admissible members\\retain $3S_1,4S_2,W$};
\draw[->] (3.05,-1.95)--(1.85,-1.7);
\end{tikzpicture}
\caption{The exact parameter disc for the standard fixed data.
The coordinate map is $w=e^{-\pi i c}$ and the coarser coordinate
map is $u=w^2$. The imaginary diameter is the fixed locus of
$w\mapsto-\bar w$, which is the reflected-family map.
The two boundary obstructions use the meridian comparison at the
puncture and the intrinsic Albanese nome at the outer circle.
The special fibres are over $t=0,1,\infty$ in every member, as
indicated in Figure~\ref{fig:I:orbifold}.}
\hypertarget{tag-01RM}{}\label{fig:I:parameter-disc}
\end{figure}

\begin{figure}[htbp]
\centering
\begin{tikzpicture}[x=1cm,y=1cm,>=stealth,font=\small]
\begin{scope}[shift={(-3.3,0)}]
 \fill[blue!5] (-1,0)--(-1,2.8)--(1,2.8)--(1,0)
  arc[start angle=30,end angle=150,radius=1.1547];
 \draw[thick] (-1,0)--(-1,2.8);
 \draw[thick] (1,0)--(1,2.8);
 \draw[thick] (1,0) arc[start angle=30,end angle=150,radius=1.1547];
 \node at (0,1.3) {$\mathcal F$};
 \node[align=center] at (0,-1.05)
 {$|\Re\tau|\le\tfrac12,\ |\tau|\ge1$\\$0\le\theta<\pi$};
 \draw[->] (0,2.1)--(0,3.1) node[above] {cusp};
\end{scope}
\begin{scope}[shift={(1.0,.2)},x={(1.8cm,0cm)},y={(.75cm,.5cm)},z={(0cm,2.2cm)}]
 \fill[blue!4] (0,0,0)--(1,0,0)--(1,0,1)--(0,0,1)--cycle;
 \draw[thick] (0,0,0)--(1,0,0)--(1,1,0)--(1,1,1)--(0,1,1)--(0,0,1)--cycle;
 \draw[thick] (1,0,0)--(1,0,1)--(0,0,1);
 \draw[thick] (1,0,1)--(1,1,1);
 \draw[dashed] (0,0,0)--(0,1,0)--(1,1,0);
 \draw[dashed] (0,1,0)--(0,1,1);
 \node[below] at (.5,0,0) {$0\le x_1<1$};
 \node[right] at (1,.5,0) {$0\le x_2<1$};
 \node[left] at (0,0,.5) {$0\le t<12$};
\end{scope}
\node[align=center] at (1.95,-1.55)
 {$U:\ t\mapsto t-6-6x_2$\qquad
  $V:\ t\mapsto t-6+6x_1$\\
  $Z:\ t\mapsto t+12$};
\end{tikzpicture}
\caption{A fundamental set for the index-twelve monodromy kernel
in the real parabolic stabilizer of $d$.
The base is a modular fundamental set with its rotation coordinate;
the fibre is the indicated Heisenberg parallelepiped. The displayed
face maps include the central shears. The representative map and
the boundary identifications are proved in
\ref{I:monodromy-fundamental-set}; the index twelve occurs in the
central fibre coordinate.}
\hypertarget{tag-01RK}{}\label{fig:I:monodromy-domain}
\end{figure}

\begin{proposition}[The fundamental set of the monodromy kernel]
\hypertarget{tag-0001}{}\label{I:monodromy-fundamental-set}
\StatementRecord{I:monodromy-fundamental-set}{proved}{0001}{\SeriesRef{fam:src:F.1.1}, \SeriesRef{fam:src:F.1.2}}
In the real coordinates $g(B,x,t)$ of
\ref{fam:src:F.1.1}, a measurable fundamental set for
$\Gamma\backslash P_d(\mathbb R)$ consists of the matrices
\[
 g(B,x,t),\qquad
 B=n(\Re\tau)a(\Im\tau)k(\theta),\quad
 \tau\in\mathcal F,\quad 0\le\theta<\pi,
 \quad x_1,x_2\in[0,1),\quad t\in[0,12).
\]
Here
$n(u)=\left(\begin{smallmatrix}1&u\\0&1\end{smallmatrix}\right)$,
$a(v)=\operatorname{diag}(\sqrt v,1/\sqrt v)$,
and $k(\theta)$ is the rotation matrix. Paired modular edges
are assigned to one side, and their lower-dimensional
stabilizer identifications use the corresponding lifts in
$\Gamma$: for
$S=\left(\begin{smallmatrix}0&-1\\1&0\end{smallmatrix}\right)$
and $T=\left(\begin{smallmatrix}1&1\\0&1\end{smallmatrix}\right)$,
we use $\widetilde S=g(S,0,9)$,
$\widetilde T=g(T,0,1)$ and
$\widetilde{-I}=g(-I,0,6)$.
Their maps on the three coordinates are
$(B,x,t)\mapsto(AB,Ax,t+\alpha(A))$, with
$\alpha(S)=9$, $\alpha(T)=1$ and $\alpha(-I)=6$.
The Heisenberg face identifications are the
maps printed in Figure~\ref{fig:I:monodromy-domain}.
\end{proposition}
\begin{proof}
\begin{enumerate}[label=\arabic*.,leftmargin=*]
\item The projection $g(B,x,t)\mapsto B$ maps $\Gamma$
onto $\operatorname{SL}_2(\mathbb Z)$ by
\ref{fam:src:F.1.1}. For any $B\in\operatorname{SL}_2(\mathbb R)$,
the usual modular reduction gives an integral matrix $A$
with $AB=n(u)a(v)k(\theta)$ in the stated set. Choose a
lift $g(A,y,s)\in\Gamma$ and multiply on the left.
For the displayed edge generators, the character formula
$\chi(g(A,0,t))=t-\alpha(A)\pmod {12}$ proves that the
three specified lifts lie in $\Gamma$. Substitution in
the group law gives each displayed coordinate map.
This is the specified map from a general parabolic element
to one over the modular set.
\item Two choices of lift over the same $A$ differ by
$\Gamma\cap H$. The group calculation in
\ref{fam:src:F.1.1} identifies this kernel with
$\langle U,V,Z\rangle$, where
$U=g(I,(1,0),-6)$, $V=g(I,(0,1),-6)$ and
$Z=g(I,0,12)$.
\item The multiplication law sends $(x,t)$ under $U$
to $(x+(1,0),t-6-6x_2)$, and under $V$ to
$(x+(0,1),t-6+6x_1)$. Powers of these maps first put
$x_1$ and $x_2$ in $[0,1)$. A power of $Z$ then puts
$t$ in $[0,12)$. All three are left translations
in $\Gamma$, so they retain the original coset.
\item Within the interiors of these intervals equality
of two reduced representatives forces their kernel
element to have zero $x$-translation and then zero
central translation. Modular uniqueness gives the same
conclusion in the base interior. The usual edge and
elliptic identifications, lifted by the specified
surjection, account for the measure-zero boundary.
This proves the fundamental-set assertion and all
displayed Heisenberg identifications.
\end{enumerate}
\end{proof}

\subsection{The Albanese coordinates and real structures}
\begin{proposition}[Exact nomes and continuation radius]\hypertarget{tag-000L}{}\label{ari:C.2.1}
\StatementRecord{ari:C.2.1}{proved}{000L}{\SeriesRef{fam:src:A.4.1}, \SeriesRef{ari:C.1.2}, \SeriesRef{fam:src:C.2.3}, \SeriesRef{fam:src:A.5.1}; \cite{DuEtAl2026v1}}

For $w=e^{-\pi ic}$, the fixed branches give
\begin{align*}
 \kappa_1&=c+\tfrac12-\frac{3i}{2\pi}\log48,&
 q_1&=-w^2/48^3,\\
 \kappa_2&=c+\tfrac32-\frac{3i}{\pi}\log6,&
 q_2&=iw/216,\qquad q_1=(27/64)q_2^2.
\end{align*}
The family is parameterized by $0<|q_2|<1$ and exists precisely in
that domain. Writing $J(q)=q^{-1}+744+\cdots$, one has
\[
 j(\operatorname{Alb}S_1)=J(-w^2/48^3).
\]
This function is holomorphic on $0<|w|<48^{3/2}$, has a pole of order
two at zero, and has a natural boundary at $|w|=48^{3/2}$.
In particular it continues holomorphically through every point of the
family boundary $|w|=216$.
\end{proposition}


\begin{proof}
\begin{enumerate}[label=\arabic*.,leftmargin=*]
\item The precise Albanese lattices in \cref{fam:src:A.4.1} are verified by projecting all
translations on the product covers. At the order-three fibre the
projection is $(1/3)(\mathbb Z+\kappa_1\mathbb Z)$; multiplication by
three gives the first stated lattice. At the order-four fibre it is
$\tfrac12\mathbb Z+\tfrac14\kappa_2\mathbb Z$; multiplication by two
gives $\mathbb Z+(\kappa_2/2)\mathbb Z$. Thus the positive-imaginary
period representatives are $-\kappa_1$ and $-\kappa_2/2$.
Adding $2\mu_1=1-i/\sqrt3$ and
$3\mu_2=3/2-3i/2$ to \ref{ari:C.1.2} proves the two $\kappa$ identities,
including their real parts. Exponentiation gives the signs in $q_1,q_2$;
the ratio is $216^2/48^3=27/64$.

\item The identity $q_2=iw/216$ in Proposition~\ref{push:src:D.1.1} has
factor $1/2$ precisely because of the Albanese projection just computed.
Its natural-boundary proof uses modular invariance at rational boundary
points. The existence criterion and exact $2\mathbb Z$
identification are \cref{fam:src:C.2.3,fam:src:A.5.1}. The latter
is also \cite[Theorem~1.4, p.~3]{DuEtAl2026v1} under the explicit
parameter map of \cref{night:P5.7}: with
$\delta=\beta^{\rm part}-\beta_0$ and $a=c+\delta$, one has
$\Omega_a=\Pi_c$ and $a'-a=c'-c$. That comparison also identifies
the finite affine twists and the cusp collar.

\item The Fourier expansion for $J$ converges for $0<|q|<1$ and has a simple
pole at $q=0$. The inverse image under $q=-w^2/48^3$ is the claimed
punctured disc and the pole has order two. The inequality
$216<48^{3/2}$ is equivalent to $27<64$. At every nonzero point on the
larger boundary the map to $q$ is locally biholomorphic. A continuation
there would continue $J$ through the unit circle, contradicting the
rational-cusp argument proved in Proposition~\ref{push:src:D.1.1}.
\end{enumerate}
\end{proof}

\begin{corollary}[The real locus]
\hypertarget{tag-000M}{}\label{ari:C.2.2}
\StatementRecord{ari:C.2.2}{proved}{000M}{\SeriesRef{fam:src:B.4.2}}
An antiholomorphic equivalence $X_c\to X_{c'}$ exists exactly when
$c'-(1-\bar c)\in2\mathbb Z$. Consequently $X_c$ admits an
antiholomorphic involution exactly when $q_2\in\mathbb R\setminus\{0\}$.
\end{corollary}
\begin{proof}
The equivalence criterion and the involution are
\cref{fam:src:B.4.2}. Its real locus is $w=-\bar w$.
Under the exact coordinate change $q_2=iw/216$, this is
$q_2=\bar q_2$; the parameter domain excludes zero.
\end{proof}

\begin{remark}[Affine corrections for the real comparison]
\hypertarget{tag-IA20}{}\label{I:real-affine-corrections}
\StatementRecord{I:real-affine-corrections}{proved}{IA20}{}
The normalization in \cref{fam:src:B.4.2} has linear fibre map
$-I_2$, parameter shift $1$, and lattice map $-I+E$, where
$d=(0,0,0,1)^t$ and $E=d(1,0,0,0)$.
For the canonical vectors $v_1=\varepsilon$, $v_2=-\varepsilon'$
and the reflected vectors
$v'_1=-\varepsilon+d$, $v'_2=\varepsilon'-3d$, its integral
correcting cocycles and rational translations can be taken as
\[
 \ell_1=(0,0,-2,0)^t,\quad \ell_2=(0,2,0,0)^t,
 \qquad
 r_1=(0,-2/3,-2/3,0)^t,\quad r_2=(0,1,1,0)^t.
\]
Direct substitution in the matrices of \cref{fam:src:B.2.1} gives
\[
 N_1\ell_1=-2d,\qquad N_2\ell_2=4d,\qquad
 (I-A_j)r_j=\frac{v_j-(-I+E)v'_j}{m_j}+\ell_j.
\]
The cusp cocycle is
\[
 \ell_0=-A_0(\ell_1+A_1\ell_2)=(0,0,4,-2)^t.
\]
These are the explicit corrections for the root-disc and cusp
extensions in \cref{fam:src:B.2.1,fam:src:B.4.2}.
\end{remark}

\section{The uniform corner obstruction}

\subsection{Period transport and parameter descent}

\begin{theorem}[Parameter descent and cusp transport]\hypertarget{tag-02R9}{}\label{night:P2.1}\label{stmt:P2.1}
\StatementRecord{night:P2.1}{proved}{02R9}{the definitions and constructions stated in the proof}

For an admissible $(3,4)$ datum with integral twists $(l_0,l_1,l_2)$, where $3\nmid l_1$ and $l_2$ is odd, use $w=r^6$,
$c=-12v$, $r=e^{2\pi iv}$ and $q=e^{2\pi is}$, and put
\[
 b=4l_1+3l_2,\qquad p=12l_0-4l_1-3l_2,\qquad a=b-12l_0=-p.
\]
On the logarithmic cover of the cusp bidisc the period columns are
\[
 \gamma=\binom{6\mu}{C-s-12v},\quad
 u=\binom{s+h}{\mu},\quad w=\binom10,\quad\delta=\binom01,
\]
where $h,\mu,C$ are holomorphic at $q=0$. Their positive-loop column comparisons are
\[
 K_q=\begin{pmatrix}1&0&0&0\\0&1&0&0\\0&1&1&0\\-1&0&0&1\end{pmatrix},\qquad
 K_r=\begin{pmatrix}1&0&0&0\\0&1&0&0\\0&0&1&0\\-12&0&0&1\end{pmatrix}.
\]
The parameter descent on the ordinary Jacobian chart is
\begin{equation}\hypertarget{tag-01O8}{}\label{eq:P2H}
 R:(t,v,\zeta)\longmapsto(t,v+1,\zeta+H(t)\delta),\qquad
 H(t)=-\frac{4l_1\log t+3l_2\log(t-1)}{2\pi i}.
\end{equation}
In the cusp collar it extends across $q=0$ and has coefficient $a$ on $s\delta$.
\end{theorem}
\begin{proof}
\begin{enumerate}[label=\arabic*.,leftmargin=*]
\item We substitute $c=-12v$ in the period matrix \eqref{wb:pshift:eq:periods} and use its cusp expansions $\tau=s+h$, $\beta^{\rm part}=C-s$. Direct subtraction gives the two displayed matrices. On fixed vector coordinates the lattice conjugation matrices are $K_q^{-1}$ and $K_r^{-1}$. In the characters $Z,W$, addition of $u$ and $\gamma$ gives respectively
\[
 (Z,W)\mapsto(qe^{2\pi ih}Z,e^{-2\pi i\mu}W),\qquad
 (Z,W)\mapsto(e^{12\pi i\mu}Z,qr^{12}e^{-2\pi iC}W).
\]
These are holomorphic character identities, including the signs of both meridians.

\item The collar section carrying the primitive $\widehat\gamma$ class has lift
\begin{equation}\hypertarget{tag-01OC}{}\label{eq:P2zg}
 z_\gamma(s,v)=\binom{6s\mu}{-s(s+1)/2+s(C-12v)}.
\end{equation}
It satisfies $z_\gamma(s+1,v)-z_\gamma(s,v)=\gamma(s+1,v)$ and $z_\gamma(s,v+1)-z_\gamma(s,v)=-12s\delta$. We therefore use the cusp coordinate $\eta=\zeta+l_0z_\gamma$, retaining the collar integer.

\item At finite point $j$, the logarithmic gluing is translation by $(\log z_j)/(2\pi i)\Pi v_j$, with $t_j=z_j^{m_j}$ and $\widehat\gamma$ coefficient of $v_j$ equal to $l_j$. Changing $v$ by one changes this gluing by $-12l_j\log z_j/(2\pi i)\delta$. This is $-4l_1\log t/(2\pi i)\delta$ at $0$ and $-3l_2\log(t-1)/(2\pi i)\delta$ at $1$. The complementary summand of $H$ is holomorphic on each disc, so its translation extends. The finite affine action changes by $-12l_j/m_j$ times the integral period $\delta$, and is unchanged on the torus. Thus \eqref{eq:P2H} is the descent compatible with both finite caps. Branch changes of $H$ are integral $\delta$ periods.

\item At infinity,
\[
 H=bs+U(q),\qquad U(q)=-\frac{3l_2\log(1-q)}{2\pi i}.
\]
In $\eta$ the descent is translation by $(as+U)\delta$. Adding the lattice period $a\gamma$ rewrites this vector as $(6a\mu,U+aC-12av)$. It is holomorphic at $q=0$ for each $v$; its exponentials act by units on every toric chart. This proves the extension and its collar coefficient. The displayed formulas verify the collar coefficient by direct substitution.
\end{enumerate}
\end{proof}

\subsection{The partial filling and its origins}
\begin{theorem}[Partial filling and extending origins]\hypertarget{tag-02RD}{}\label{night:P2.3}\label{stmt:P2.3}
\StatementRecord{night:P2.3}{proved}{02RD}{the definitions and constructions stated in the proof}
\hypertarget{tag-02RC}{}\label{night:P2.2}\label{stmt:P2.2}
\StatementRecord{night:P2.2}{proved}{02RC}{the definitions and constructions stated in the proof}


There is a proper holomorphic partial filling over $(\Pp^1\times\Delta_r)\setminus\{(\infty,0)\}$ with smooth total space and reduced normal-crossing parameter central divisor. Over a regular finite base point its central fibre is a cycle of twelve ruled surfaces. Its component rotations at $0,1,\infty$ are $4l_1,3l_2,-b$ modulo $12$, with component cover
\[
 y^{12}=t^{4l_1}(t-1)^{3l_2}.
\]


On $q=Q^{12}$, extending origins have lifts
\begin{align*}
 o&=\binom{-b\mu s/2}{bs^2/24-bCs/12+l_0s/2+v(bs+U)},\\
 \sigma&=-l_0z_\gamma+\binom{6a\mu v}{vU+aCv-6av^2+6av}.
\end{align*}
The section $\sigma$ extends across the cusp for $r\ne0$, and $o$ extends across the parameter central fibre of the partial filling. Their difference is
\begin{equation}\hypertarget{tag-01O5}{}\label{eq:P2F}
 F=\sigma-o=a\binom{\mu(s+12v)/2}{C(s+12v)/12-(s+12v)^2/24+6v}.
\end{equation}
On $q=Q^{12m}$, $r=R^n$, with $m,n$ positive integers,
\begin{align}
 \Delta_QF&=am\gamma-6am^2\delta,\hypertarget{tag-01O6}{}\label{eq:P2FQ}\\
 \Delta_RF&=an\gamma-6an(n-1)\delta.\hypertarget{tag-01O7}{}\label{eq:P2FR}
\end{align}
The radial class modulo $\Z\widehat\delta$ is $-pn\widehat\gamma$.
\end{theorem}
\begin{proof}
\begin{enumerate}[label=\arabic*.,leftmargin=*]
\item On a simply connected finite-base chart subtract a local origin with $v$-difference $H\delta$, quotient by $\delta$, and put $W=e^{-2\pi i\eta_2}$, $z=\eta_1$. The remaining lattice transformations are
\begin{align*}
 (z,W)&\mapsto(z+1,W),\\
 (z,W)&\mapsto(z+\tau,e^{-2\pi i\mu}W),\\
 (z,W)&\mapsto(z+6\mu,d(t)r^{12}W),
\end{align*}
with $d$ a nonzero holomorphic unit. The first two give the complement of the zero section of a degree-zero line bundle over $\C/(\Z+\tau\Z)$. Use the infinite chain of node charts
\begin{equation}\hypertarget{tag-01O9}{}\label{eq:P2chain}
 X_i=W/r^i,\quad Y_i=r^{i+1}/W,\quad X_iY_i=r\qquad(i\in\Z).
\end{equation}
Adjacent charts are identified on $Y_i\ne0$ by $X_{i+1}=Y_i^{-1}$, $Y_{i+1}=rY_i$. The charts and their inverse overlaps give an infinite projective-line chain at $r=0$. Unit multiplication on $W$ multiplies $X_i,Y_i$ inversely, so all elliptic line-bundle transitions extend. The third period shifts the index by twelve, translates the elliptic coordinate by $6\mu$, and multiplies by $d$. Its quotient has the stated twelve-cycle.

\item The total node chart is smooth by eliminating $r=X_iY_i$, and its central divisor is reduced NC. On a compact smaller finite-base chart the elliptic lattices have uniformly bounded compact fundamental parallelograms. Choose middle annuli and node discs using upper and lower bounds for $d$ and the finitely many transition units. Twelve consecutive indices give a compact fundamental cover over a closed smaller parameter disc. The overlap formulas have closed graphs on that cover, and the positive index translation separates nonadjacent charts. Thus the quotient is Hausdorff and proper. At $r\ne0$ it is exactly the original torus quotient.

\item At root point $j$, addition of $v_j/m_j$ multiplies $W$ by a factor containing $r^{12l_j/m_j}$ and rotates the cycle by $12l_j/m_j$. Every nonidentity power moves the cycle, because $3\nmid l_1$ and $l_2$ is odd. The full finite action is free on the central divisor, including its intersection strata. Its quotient is smooth and preserves the reduced NC parameter divisor. Logarithmic gluing translations differ on overlaps by holomorphic units and integral chain shifts, which extend by the same formulas. The rotations give the cyclic cover; the third is their negative sum and equals $p$ modulo $12$.

\item The lift descent identities are
\begin{align*}
 \Delta_{s,12}o&=-b\gamma+6(b+l_0)\delta,&\Delta_{v,1}o&=H\delta,\\
 \Delta_{s,12}\sigma&=-12l_0\gamma+78l_0\delta,&\Delta_{v,1}\sigma&=H\delta+a\gamma.
\end{align*}
Every additional term is a lattice period. Both lifts therefore define holomorphic sections of the descended torus torsor. In the cusp chart $\sigma+l_0z_\gamma$ is holomorphic and bounded in $q$ and lies in the smooth toric open orbit at $q=0$. It extends there. Subtracting $o$ makes the vector coordinate invariant under the parameter deck. The zero coordinate extends in the node charts \eqref{eq:P2chain}. This constructs the required second origin.

\item Subtracting gives \eqref{eq:P2F}. Integral shifts of $s$ leave $\mu,C,U$ unchanged; expansion of the square gives \eqref{eq:P2FQ}--\eqref{eq:P2FR}. For $(l_0,l_1,l_2)=(0,1,-1)$, $p=-1$, $a=1$. Its radial class is $+n\widehat\gamma$. For every admissible twist $p\ne0$, so the radial class is nonzero in the free quotient lattice.
\item We glue this finite-base filling to the prescribed cusp filling over $r\ne0$ by the translation $\sigma-o$. The displayed origin formulas identify the overlap with the original family. Properness and the Hausdorff property are local on the displayed target, so the glued map has both. The local equations prove flatness over the parameter disc.
\end{enumerate}
\end{proof}

\subsection{Extending the marked base map}
\begin{theorem}[Extension of the marked base map]\hypertarget{tag-02RE}{}\label{night:P2.4}\label{stmt:P2.4}
\StatementRecord{night:P2.4}{proved}{02RE}{the definitions and constructions stated in the proof}
\hypertarget{tag-0314}{}\label{push:src:B.4.4}
\StatementRecord{push:src:B.4.4}{proved}{0314}{\SeriesRef{fam:src:A.1.1}, \SeriesRef{fam:src:A.1.2}; \citep[Section 6, Hauptsatz I]{Grauert60}}


For every admissible $(3,4)$ twist, the marked map $f:X_c\to\Pp^1$ extends meromorphically over every proper Hausdorff analytic filling of a finite parameter cover. For a smooth given total space, its graph is a proper projective modification carrying the marked map. In general, a smooth proper modification carries a holomorphic marked map to $\Pp^1$, and hence a proper map to $\Pp^1\times\Delta$.
\end{theorem}
\begin{proof}
\begin{enumerate}[label=\arabic*.,leftmargin=*]
\item Take the reduced irreducible component containing the punctured family. It is closed and proper; discarding other central components preserves that family. Resolve it isomorphically over the smooth punctured family. We use analytic resolution \cite[Theorem~1.6 and the following noncompact extension, p.~215]{BM97}: reduced complex analytic spaces admit proper modifications by locally finite sequences of smooth-centre blowups with smooth result. It can be chosen unchanged over the nonsingular locus. Properness makes the neighborhoods of the central fibre relatively compact, so the stated analytic hypotheses hold. Write $g:\mathcal Z\to\Delta_s$ for the smooth result. The following argument also applies directly to a smooth given total space.

\item We shrink the disc whenever necessary. Put
$L=K_{\mathcal Z/\Delta_s}^{-2}$. The relative canonical
bundle here is the line bundle
$K_{\mathcal Z}\otimes g^*K_{\Delta_s}^{-1}$; its
restriction to a nonzero fibre is that fibre's canonical bundle.
The precise \cref{fam:src:A.1.1}, \cref{fam:src:A.1.2} input is
$K_{X_s}^{-2}\simeq f_s^*\mathcal O_{\PP^1}(1)$,
with connected fibres and no further sections. Therefore
$h^0(X_s,L|_{X_s})=2$, and its complete linear system is
the given map $f_s$, up to a projective change of coordinate.
The three critical values are intrinsically distinguished by
their fibres: multiplicity three, multiplicity four, and reduced
normal crossings, respectively.

\item Grauert's proper direct-image theorem
\citep[\S6, Hauptsatz~I, p.~59]{Grauert60} makes
$\mathcal F=g_*L$ coherent. It is torsion-free: multiplication
by a nonzero base germ is injective on $L$, because $g$ is
nonconstant on the smooth connected total space, and hence is injective on
its direct image. A coherent torsion-free sheaf on the smooth
disc is locally free. Its rank is two. On the punctured disc,
constant $h^0=2$ and the analytic base-change theorem of Proposition~\ref{push:src:E.3.1} identify its fibres with the complete spaces of sections.
Choose a frame $(\sigma_0,\sigma_1)$ of $\mathcal F$.
Evaluation gives two global holomorphic sections of $L$ on
$\mathcal Z$, and thus a meromorphic map
$F_0=[\sigma_0:\sigma_1]$ to $\PP^1$.

\item Its graph is analytic and proper by the following local
construction. In a local
trivialization of $L$, write the sections as $a,b$. Remove their
common divisorial factor in the regular local ring of the smooth
space $\mathcal Z$. The equation
$v a-u b=0$ in $\mathcal Z\times\PP^1_{[u:v]}$
then has a unique component dominating $\mathcal Z$; it is
the closure of the graph where $(a,b)\ne(0,0)$. These local
components glue. Denote the resulting reduced graph by
$\mathcal G$. The projection $\mathcal G\to\mathcal Z$
is proper and projective, and is an isomorphism over
$\Delta_s^*$. The map
\[
 H:\mathcal G\longrightarrow\PP^1\times\Delta_s,
 \qquad H=(F_0,g),
\]
is proper. Locally $\mathcal G$ is a hypersurface in the smooth
space $\mathcal Z\times\PP^1$. The vanishing of the maximal
Jacobian minors for its equation and the two target coordinates
defines the closed analytic locus $C$ where $H$ is not a
submersion; singular points of $\mathcal G$ are included.

\item The image $H(C)$ is a closed analytic subset of the target.
One may deduce this precise proper-image assertion from the same
direct-image theorem: the coherent sheaf
$H_*\mathcal O_C$ has support exactly $H(C)$, since the
constant section $1$ is nonzero at every nonempty proper fibre.
Over $\Delta_s^*$ this image consists exactly of three distinct
critical-value sections. It cannot have a two-dimensional
component, since such an analytic component would contain an
open set of the connected two-dimensional target and hence meet
the punctured region in dimension two. Remove any component
contained in $\PP^1\times\{0\}$. The remaining reduced
analytic curve is proper with zero-dimensional fibres over the
disc, and is therefore finite over it.

\item Each of the three labelled sections on the punctured disc belongs
to a degree-one branch of that finite curve. Indeed their labels
are intrinsic and cannot be permuted by continuation. The
normalization of a degree-$d$ branch over a punctured disc has
cyclic monodromy of length $d$; a single-valued labelled section
therefore has $d=1$. The normalization extends it holomorphically
to a map from the full disc to $\PP^1$. Let these three maps be
$c_1(s),c_2(s),c_0(s)$, in the multiplicity-three,
multiplicity-four and cusp order. They may collide at $s=0$.

\item For $s\ne0$ the unique projective transformation sending
$(c_1,c_2,c_0)$ to $(0,1,\infty)$ has meromorphic coefficients
at $0$. In a common affine chart its value at $z$ is
\[
 A_s(z)=
 \frac{(z-c_1(s))(c_2(s)-c_0(s))}
      {(z-c_0(s))(c_2(s)-c_1(s))};
\]
homogeneous brackets give the same formula if a value is infinite.
The denominators are not identically zero, since the three values
are distinct on the punctured disc. Thus $A_s\circ F_0$ is
meromorphic on $\mathcal Z$. On every nonzero fibre it agrees
with the marked $f_s$, because both are the complete
bi-anticanonical map and have the same three labelled values.
Thus the two maps agree over the punctured disc.

\item Finally the graph of $A_s\circ F_0$ is constructed by the same
local two-section argument after clearing powers of the base
coordinate in its meromorphic coefficients. Its projection to
$\mathcal Z$ is the asserted proper projective modification,
and its second projection is the required holomorphic marked map.
Resolve this reduced graph isomorphically over the smooth punctured
family by the same analytic resolution theorem. Composing its
holomorphic marked map with the resolution gives the asserted
smooth proper modification.
\end{enumerate}
\end{proof}

\begin{theorem}[Vanishing of finite monodromy orbits]\hypertarget{tag-02RF}{}\label{night:P2.5}\label{stmt:P2.5}
\StatementRecord{night:P2.5}{proved}{02RF}{the definitions and constructions stated in the proof}

The global monodromy on $\Lambda_\Q/\Q\widehat\delta$ has no nonzero vector with finite orbit.
\end{theorem}
\begin{proof}
\begin{enumerate}[label=\arabic*.,leftmargin=*]
\item In coordinates $(g,u,w)$, a cusp unipotent has nilpotent part $N=E_{32}$. The two finite homology generators are
\[
 A_1=\begin{pmatrix}1&0&0&0\\6&0&1&0\\-6&-1&-1&0\\-2&1&0&1\end{pmatrix},\qquad
 A_2=\begin{pmatrix}1&0&0&0\\0&0&-1&0\\-6&1&0&0\\3&0&1&1\end{pmatrix}.
\]
A vector with finite orbit is fixed by a positive power of every conjugate of that unipotent. Since $(I+N)^j=I+jN$ in characteristic zero, it belongs to $\ker N$, $\ker(NA_1)$ and $\ker(NA_2)$. The equations are $u=0$, $6g+w=0$, $-w=0$, forcing zero. The coefficient matrix of these three equations has determinant $6$.
\end{enumerate}
\end{proof}

\begin{theorem}[Vanishing of horizontal meridians]\hypertarget{tag-02RG}{}\label{night:P2.6}\label{stmt:P2.6}
\StatementRecord{night:P2.6}{proved}{02RG}{the definitions and constructions stated in the proof}

In a smooth resolved proper filling carrying the marked base map, every parameter-central component dominating $\Pp^1$ has zero normal-meridian lattice class modulo $\Z\widehat\delta$, relative to the ordinary parameter origin.
\end{theorem}
\begin{proof}
\begin{enumerate}[label=\arabic*.,leftmargin=*]
\item Put $B^\circ=\Pp^1\setminus\{0,1,\infty\}$. After $r=R^n$, the quotient fundamental group of the punctured family over $B^\circ$ by the normal subgroup generated by $\widehat\delta$ is
\begin{equation}\hypertarget{tag-01OB}{}\label{eq:P2split}
 \bigl((\Lambda/\Z\widehat\delta)\rtimes\pi_1(B^\circ)\bigr)\times\Z_R.
\end{equation}
The finite Jacobian linear changes are independent of $c$ and fix $\widehat\delta$. The $R$ descent is the $\delta$ translation $H$, all its branch changes are integral $\delta$ periods, and its lattice conjugation differs from the identity by a $\delta$ multiple. The finite gluing changes under $c\mapsto c-12n$ also lie in $\Z\delta$. Thus every cross-cocycle disappears in exactly the quotient \eqref{eq:P2split}; this proves the presentation from the chart action.

\item Let $D$ dominate $B$. Its positive normal circle has parameter degree $e=\ord_D(R)>0$, and in \eqref{eq:P2split} is $(\lambda,e)$ with trivial base-loop coordinate. Changing a basing path by a fibre loop changes $\lambda$ by $(I-T_R^e)$ of that loop, which vanishes modulo $\delta$. Transport along $D$ therefore acts on $\lambda$ only by base monodromy.

\item Resolve the reduced compact component $D$ by a proper map
$\nu:\widetilde D\to D$ which is an isomorphism over $D_{\rm reg}$,
using the analytic resolution theorem in \ref{night:P2.4}.
Apply Stein factorization to $f|_D\circ\nu$:
\[
 \widetilde D\xrightarrow{p} C\xrightarrow{h} B,
 \qquad f|_D\circ\nu=h\circ p.
\]
Here $p$ has connected fibres, $p_*\mathcal O_{\widetilde D}=\mathcal O_C$,
and $h$ is finite, by \cite[Theorem~8.5, pp.~80--81]{Movasati24}.
The curve $C$ is normal: a locally meromorphic function integral over
$\mathcal O_C$ pulls back to a holomorphic function on the normal space
$\widetilde D$ and descends by the displayed direct-image equality.
Thus $C$ is a smooth compact curve.

Let $\Sigma\subset B$ consist of $0,1,\infty$, the branch values of $h$,
and the $h$-images of the critical values of $p$. This set is finite.
Put
\[
 B'=B\setminus\Sigma,\qquad C'=h^{-1}(B'),\qquad
 V=p^{-1}(C').
\]
Then $p:V\to C'$ is a proper submersion with connected fibres, and
$h:C'\to B'$ is a finite unramified cover. Path lifting followed by
a path in a fibre gives $\pi_1(V)\twoheadrightarrow\pi_1(C')$.
Consequently the image of $\pi_1(V)$ in $\pi_1(B')$ has index $\deg h$.

Let $E$ range over the other irreducible components of the reduced
parameter-central divisor, and set
\[
 A=V\cap\nu^{-1}\!\left(D_{\rm sing}\cup
                  \bigcup_{E\ne D}(D\cap E)\right),\qquad U=V\setminus A.
\]
This is a proper complex analytic subset $A$ of the smooth manifold $V$.
General position moves every based loop in $V$ off $A$, so
$\pi_1(U)\to\pi_1(V)$ is surjective. The inclusion
$B'\hookrightarrow B^\circ$ is also surjective on fundamental groups;
therefore the image of $\pi_1(U)$ in $\pi_1(B^\circ)$ has index at most
$\deg h$.

The restriction of $\nu$ identifies $U$ with its image in
$D_{\rm reg}$ away from the other central components and over $B'$.
Composing these loops with $\nu$ thus transports the normal circle
in the original smooth total space. Its complex orientation, lattice
class $\lambda$ and parameter degree $e=\ord_D(R)$ are those of step~2.
Hence a finite-index subgroup of the base group transports the
original normal meridian.
\item The orbit of $\lambda$ is finite. \ref{night:P2.5} kills it rationally, and the quotient lattice is free, so it vanishes integrally. This includes every multiplicity $e$. After a local finite root in the base, each horizontal component inherits this zero class from its global preimage.
\end{enumerate}
\end{proof}

\subsection{The commuting-meridian obstruction}
\begin{theorem}[Nonexistence of proper fillings]\hypertarget{tag-02RH}{}\label{night:P2.7}\label{stmt:P2.7}
\hypertarget{tag-0313}{}\label{push:src:B.4.33}
\label{push:B433}
\StatementRecord{push:src:B.4.33}{proved}{0313}{\SeriesRef{push:src:B.4.29}, \SeriesRef{push:src:B.4.32}, \SeriesRef{push:src:B.4.31}, \SeriesRef{push:src:B.4.4}, \SeriesRef{push:src:B.4.24}, \SeriesRef{push:src:B.4.30}, \SeriesRef{push:src:B.4.27}, \SeriesRef{push:src:B.4.19}, \SeriesRef{push:src:B.4.26}; \citet[Theorems~1.6 and 1.10, pp.~215--216, including the noncompact extension after Theorem~1.6]{BM97}}
\hypertarget{tag-030D}{}\label{push:src:B.4.1}\label{push:B41}
\StatementRecord{push:src:B.4.1}{proved}{030D}{\SeriesRef{push:src:B.4.33}, \SeriesRef{push:section:B.1}, \SeriesRef{push:section:B.2}, \SeriesRef{push:src:B.4.19}}
\hypertarget{tag-031A}{}\label{push:src:B.5.1}\label{push:B51}
\StatementRecord{push:src:B.5.1}{proved}{031A}{\SeriesRef{push:src:B.4.33}}
\StatementRecord{night:P2.7}{proved}{02RH}{the definitions and constructions stated in the proof}

For every admissible $(3,4)$ twist, the punctured family has no proper Hausdorff complex-analytic filling at $w=0$ after any finite parameter cover. This includes every sphere member and permits arbitrary singularities, multiplicities and extra central components.
\end{theorem}
\begin{proof}
\begin{enumerate}[label=\arabic*.,leftmargin=*]
\item A proper filling pulls back properly to a further cover $w=R^{6n}$. By \ref{night:P2.4} we resolve its dominating component and marked-map graph. Near infinity take the finite base change $q=Q^{12m}$, normalize the dominating component, and resolve it. Principalize the inverse image of the two axes using \cite[Theorem~1.10, p.~216, with the analytic category in (0.1)(2), p.~208]{BM97}. That theorem turns a coherent ideal on a smooth analytic space into a locally monomial ideal after proper smooth-centre blowups. The axis ideal is coherent and the relevant inverse image is compact over a smaller closed bidisc, so its hypotheses hold. The smooth torus family over $Q^*\times R^*$ is unchanged.

\item The proper map to the bidisc has connected fibres. Its Stein finite factor is generically one sheet because its general fibre is a connected torus. A finite birational map to the normal bidisc is an isomorphism: its finite local algebra is integral over, and contained in the fraction field of, the integrally closed local ring of that bidisc. In particular the corner fibre is connected.

\item In the axis complement, the original logarithmic deck lifts have
\begin{align*}
 A_Q&=I-12mE_{32}+12mE_{41},&
 A_R&=I+12nE_{41},&
 [Q,R]&=-12mnb\widehat\delta.
\end{align*}
The affine relation is the difference $H(s+12m)-H(s)=12mb$ through $n$ parameter turns. The lifts
\[
 Q_o=T_{-mb\widehat\gamma}Q,\qquad R_o=R
\]
commute. They give a specific semidirect product $\Lambda\rtimes\Z^2$; central $\delta$ corrections do not affect the following coefficients.

\item For every irreducible axis component $D_i$, put
\[
 \alpha_i=\ord_{D_i}(Q),\quad\beta_i=\ord_{D_i}(R),\quad
 \kappa_i=\text{the }\widehat\gamma\text{ coefficient of its lattice meridian}.
\]
Its positive meridian is $(\lambda_i,\alpha_i,\beta_i)$, with nonnegative orders and positive sum. Conjugation changes its lattice vector by a vector killed by the first-coordinate covector; thus $\kappa_i$ is independent of its basing path in the complement. Meeting components have commuting meridians, represented by coordinate circles in a smooth NC chart. The $\widehat\delta$ coordinate of their commutator gives
\begin{equation}\hypertarget{tag-01OA}{}\label{eq:P2ratio}
 (m\alpha_i+n\beta_i)\kappa_j=(m\alpha_j+n\beta_j)\kappa_i.
\end{equation}
Both nilpotents have square zero and their products vanish, so there are no higher terms. The $E_{32}$ contribution lies in the $\widehat w$ coordinate; $E_{41}$ gives exactly \eqref{eq:P2ratio}. The ratio $\kappa_i/(m\alpha_i+n\beta_i)$ is therefore constant along every intersection edge. Its denominator is positive for every axis component, including vertical exceptional components.

\item A parameter-central component dominating the original base exists, since the proper map onto $\Pp^1\times\Delta$ is surjective and finitely many vertical components cannot cover its central base. Its strict transform meets the corner by properness. Its ratio is zero by \ref{night:P2.6}.

\item The strict cusp component met by $\sigma$ also meets the corner. Its meridian in the smooth open orbit is the loop traced by the extending section. \ref{night:P2.2} gives $\widehat\gamma$ coefficient $-12ml_0$ in the original coordinates. Relative to $Q_o$ this becomes $-12ml_0+mb=ma$. Its orders are $(1,0)$, so its ratio is $a$. Further ramification multiplies the order and coefficient by the same positive factor, and proper transforms preserve a smooth open part of the branch.

\item The compact connected corner fibre is covered by its intersections with finitely many axis components. The incidence graph of this closed cover is connected: a partition with no edge would partition the fibre into disjoint nonempty closed subsets, both also open. A chain therefore joins the horizontal central component to the strict cusp component. Equation \eqref{eq:P2ratio} forces $a=0$.

\item Admissibility gives $3\nmid l_1$ and $l_2$ odd, whence $p$ is coprime to $12$ and nonzero. Since $a=-p$, this is a contradiction. Each comparison was proper and preserved the punctured family, so it applies to the original filling and every finite parameter cover.
\end{enumerate}
\end{proof}

\begin{theorem}[The logarithmic datum $(0,2,-1)$]\hypertarget{tag-02RI}{}\label{night:P2.8}\label{stmt:P2.8}
\StatementRecord{night:P2.8}{proved}{02RI}{the definitions and constructions stated in the proof}

For $(l_0,l_1,l_2)=(0,2,-1)$ we have $p=-5$, $a=b=5$, component cover $y^{12}=t^8/(t-1)^3$, and rotations $8,-3,-5$ modulo $12$. The descent and origins are
\begin{align*}
 H&=5s+U,& U&=\frac{3\log(1-q)}{2\pi i},\\
 o&=\binom{-5\mu s/2}{5s^2/24-5Cs/12+v(5s+U)},\\
 \sigma&=\binom{30\mu v}{vU+5Cv-30v^2+30v},& F&=5F_0,
\end{align*}
where $F_0$ is the vector in \eqref{eq:P2F} without $a$. Its two differences are $5m\gamma-30m^2\delta$ and $5n\gamma-30n(n-1)\delta$. No finite parameter cover admits a proper filling.
\end{theorem}
\begin{proof}
\begin{enumerate}[label=\arabic*.,leftmargin=*]
\item Substitution in \ref{night:P2.1}--\ref{night:P2.3} gives every displayed formula. The lattice and character laws remain those of \ref{night:P2.1} and the charts those of \eqref{eq:P2chain}. The corner has $[Q,R]=-60mn\widehat\delta$ and $Q_o=T_{-5m\widehat\gamma}Q$. Its horizontal ratio is zero and strict cusp ratio five. Thus \ref{night:P2.7} gives the contradiction for this member itself.
\end{enumerate}
\end{proof}

\subsection{Higher coprime multiplicities}
\begin{theorem}[The obstruction for coprime multiplicities]\hypertarget{tag-02RJ}{}\label{night:P2.9}\label{stmt:P2.9}
\StatementRecord{night:P2.9}{proved}{02RJ}{the definitions and constructions stated in the proof}

Every admissible logarithmic datum in the coprime $(3k,4l)$ family of Viaclovsky's logarithmic construction has no proper Hausdorff analytic filling at parameter zero after any finite cover.
\end{theorem}
\begin{proof}
\begin{enumerate}[label=\arabic*.,leftmargin=*]
\item Put $m_1=3k$, $m_2=4l$, $N=m_1m_2=12kl$, with $\gcd(m_1,m_2)=1$. Let $l_j$ be the $\widehat\gamma$ coefficient of the invariant finite twist. Admissibility is $\gcd(l_j,m_j)=1$. Retain $l_0$ and define $b=m_2l_1+m_1l_2$, $a=b-Nl_0$. On $w=r^{N/2}$, $c=-Nv$, and the gamma period is $(6\mu,C-s-Nv)$. The descent is
\[
 H=-\frac{m_2l_1\log t+m_1l_2\log(t-1)}{2\pi i}=bs+U,
 \qquad U=-\frac{m_1l_2\log(1-q)}{2\pi i}.
\]
The finite actions are unchanged by $c\mapsto c-N$, since $N/m_j$ is integral. The collar coefficient is $a$; adding $a\gamma$ gives a unit action on cusp charts. The rank-one filling has $N$ components with rotations $m_2l_1,m_1l_2,-b$. Admissibility makes the finite actions free on the cycle.

\item On $q=Q^N$ the origins are
\begin{align*}
 o&=\binom{-6b\mu s/N}{bs^2/(2N)-bCs/N+l_0s/2+v(bs+U)},\\
 \sigma&=-l_0z_\gamma+\binom{6a\mu v}{vU+aCv-Nav^2/2+Nav/2},
\end{align*}
with $z_\gamma=(6s\mu,-s(s+1)/2+s(C-Nv))$. Their difference is
\[
 aF_N=a\binom{6\mu(s+Nv)/N}{C(s+Nv)/N-(s+Nv)^2/(2N)+Nv/2}.
\]
On $q=Q^{Nm}$, $r=R^n$,
\begin{align*}
 \Delta_Q(aF_N)&=am\gamma-\tfrac N2am^2\delta,\\
 \Delta_R(aF_N)&=an\gamma-\tfrac N2an(n-1)\delta.
\end{align*}
All coefficients are integral because $12\mid N$. Expanding the two quadratic expressions gives the origin subtraction, the four descent identities and both covered differences.

\item Pulling the order-three and order-four tangent root charts through $k$ and $l$ powers changes their vertical determinant orders from $2,1$ to $2k,l$. The vertical determinant is therefore $f^*\OO(-1)\otimes\OO(2kS_1+lS_2)$. The determinant of the image of $df$ is $f^*\OO(2)\otimes\OO(-(3k-1)S_1-(4l-1)S_2)$; the cusp singular ideal has codimension two. Hence
\begin{align*}
 K_X^{-1}&=f^*\OO(1)\otimes\OO((1-k)S_1+(1-3l)S_2),\\
 K_X^{-N}&=f^*\OO(d),& d&=4l+3k-kl.
\end{align*}
The integer $d$ is nonzero. Otherwise $k\mid4l$, while $\gcd(k,4l)=1$, so $k=1$, giving $d=3l+3>0$.

\item Choose $K^{-N}$ when $d>0$ and $K^N$ when $d<0$; its restriction is $f^*\OO(|d|)$. On a resolved proper filling the corresponding relative direct image is coherent torsion free of rank $|d|+1$. Evaluation and graph resolution give a proper pluricanonical image. Over the field $K$ of meromorphic parameter germs, its image is a smooth rational normal curve. More explicitly, the kernels of multiplication on the direct-image modules, tensored with $K$, give its quadratic equations; on the punctured disc these are the Veronese equations and generate its ideal. The three labelled critical-value curves give distinct $K$-rational points. The coprime multiplicities are unequal, and the third point is the cusp.

\item Over an algebraic closure of $K$ there is a unique rational parameter taking those points to $0,1,\infty$. Uniqueness makes it invariant under every field automorphism over $K$, so its coefficients belong to $K$. They are meromorphic germs with finite principal parts. Composing and resolving the graph yields the marked base map required in \ref{night:P2.4}. This explicitly recovers a degree-one base parameter.

\item Modulo $\delta$, the global lattice and finite-orbit calculation are unchanged. The cross-cocycles still lie in $\Z\delta$, so \ref{night:P2.6} applies. At the corner,
\[
 A_Q=I-NmE_{32}+NmE_{41},\quad A_R=I+NnE_{41},\quad [Q,R]=-Nmnb\widehat\delta.
\]
The commuting lifts $Q_o=T_{-mb\widehat\gamma}Q$, $R_o=R$ give \eqref{eq:P2ratio}. Connectedness forces equality of the horizontal ratio zero and cusp ratio $a$. But $a\equiv m_2l_1\pmod{m_1}$ and $a\equiv m_1l_2\pmod{m_2}$, so $\gcd(a,N)=1$. This contradiction proves nonexistence. Proper pullback to the further cover preserves the punctured family and gives the implication back to every finite parameter cover.
\end{enumerate}
\end{proof}

\section{The boundary of the parameter disc}\hypertarget{tag-02Z8}{}\label{push:section:D}

\begin{definition}[Boundary coordinates and the Hausdorff reflection]
\hypertarget{tag-PT08}{}\label{I:def:boundary-coordinates}
\StatementRecord{I:def:boundary-coordinates}{proved}{PT08}{}
For the standard twists put $H_*=(3/\pi)\log6$, fix $c_*=x+iH_*$ and set $a=x+3/2$. All topological group quotients carry their quotient topology. A Hausdorff reflection is a continuous quotient universal for continuous maps to Hausdorff spaces.
\end{definition}

\subsection{The intrinsic obstruction to extension}\hypertarget{tag-02Z9}{}\label{push:section:D.1}

\begin{proposition}[The natural boundary in the multiplicative parameter]\hypertarget{tag-032A}{}\label{push:src:D.1.1}
\StatementRecord{push:src:D.1.1}{proved}{032A}{\SeriesRef{fam:src:A.4.1}, \SeriesRef{fam:src:C.1.1}, \SeriesRef{fam:src:C.2.3}, \SeriesRef{fam:src:C.3.4}}

For \(0<|w|<216\), the reduced multiplicity-four surface has
\[
 j(\operatorname{Alb}S_2(w))=J(iw/216),\qquad
 J(q)=q^{-1}+744+196884q+\cdots.
\]
This holomorphic function has no meromorphic continuation through any
point of \(|w|=216\).
\end{proposition}


\begin{proof}
\begin{enumerate}[label=\arabic*.,leftmargin=*]
\item The fixed-coordinate periods on the product cover, including the
isogeny kernel and the affine deck translation, generate
\(\frac12\mathbb Z+\frac{\kappa_2}{4}\mathbb Z\).
Multiplication by two therefore identifies the Albanese lattice with
\(\mathbb Z+(\kappa_2/2)\mathbb Z\). This verifies the precise step of
\cref{fam:src:A.4.1} used here; taking only the invariant eigenspace would omit the
half-period. The values in \cref{fam:src:C.1.1}, \cref{fam:src:C.2.3} and \(\mu(z_2)=(1-i)/2\) give
\[
 \kappa_2=c+\frac32-iH_*,\qquad
 e^{-\pi i\kappa_2}=e^{-3\pi i/2}e^{-\pi H_*}e^{-\pi ic}
 =\frac{iw}{216}.
\]
The positive-imaginary representative is \(-\kappa_2/2\), proving the
identity with the stated Fourier normalization.

\item Proposition~\ref{fam:src:C.3.4} proves the natural boundary in the $c$-coordinate for this exact Albanese lattice. The parameter map is $w=e^{-\pi i c}$, with nonzero derivative at every finite $c$. A local inverse logarithm at a point of $|w|=216$ maps that circle to $\operatorname{Im}c=\log216/\pi=H_*$. A meromorphic continuation in $w$ would pull back to the excluded continuation in $c$.
\end{enumerate}
\end{proof}

\begin{lemma}[The intrinsic anticanonical surface family]\hypertarget{tag-032B}{}\label{push:src:D.1.2}
\StatementRecord{push:src:D.1.2}{proved}{032B}{the definitions and constructions stated in the proof}

Let \(V\) be a connected Riemann surface and
\(\pi:\mathcal X\to V\) a flat proper holomorphic map of reduced
complex spaces. Suppose that over a nonempty open subset \(V_0\), the
fibres are smooth compact threefolds whose unique anticanonical divisor
is \(2S_v\), where \(S_v\) is a smooth connected surface with
\(b_1(S_v)=2\). Then there is an open subset \(V^{\rm good}\subset V\), with at
most countable complement, on which the reduced support of the
relative anticanonical divisor is a smooth proper surface family. Its
Albanese elliptic invariant is holomorphic.
\end{lemma}


\begin{proof}
\begin{enumerate}[label=\arabic*.,leftmargin=*]
\item Flatness and the smoothness of the fibres over \(V_0\) imply that
\(\pi\) is a submersion there: the local analytic smoothness criterion
applies because the base is smooth. The relative nonsmooth locus is a
closed analytic subset, by the Jacobian criterion. Its proper image is analytic by the coherent direct-image argument in Theorem~\ref{night:P2.4}, applied to the structure sheaf of this proper analytic locus.
It misses \(V_0\), hence is discrete. We remove it and work over the
connected remaining curve \(V'\). Removing a discrete subset from a
connected Riemann surface does not disconnect it. Thus
\(\mathcal X'\to V'\) is a smooth proper map.

\item Write \(L=K_{\mathcal X'/V'}^{-1}\). Grauert's direct-image theorem \cite[\S6, Hauptsatz~I, p.~59]{Grauert60} makes \(F=\pi_*L\) and \(R^1\pi_*L\) coherent. His erratum \cite[item~2, p.~36]{Grauert1963ErratumD} corrects the assertion in \S7, Satz~2 that the locus where the comparison map fails to be bijective is always closed. For the present smooth curve, the following exact sequence computes the canonical base-change map and its torsion obstruction.
For a local base parameter \(t\) vanishing at \(v\), smoothness and
invertibility of \(L\) give an exact sequence
\(0\to L\xrightarrow{t}L\to i_*L_v\to0\), where
\(i:X_v\hookrightarrow\mathcal X'\) is the closed embedding.
Its direct-image sequence gives, at \(v\),
\[
 0\longrightarrow F_v/tF_v\longrightarrow H^0(X_v,L_v)
 \longrightarrow (R^1\pi_*L)_v[t]\longrightarrow0.
\]
The last term denotes elements annihilated by \(t\). The support of
the torsion subsheaf of the coherent module \(R^1\pi_*L\) is closed
discrete, by the structure theorem for finitely generated modules
over the local discrete valuation rings of the smooth base curve.
Outside that support the canonical base-change map is an isomorphism.
Also \(F\) is torsion-free: multiplication by a nonzero base function
is injective on sections of \(L\). It is therefore locally free.
Its rank is one by choosing an interior point outside the discrete
torsion support.
Locally a frame gives a relative section \(\sigma\) of \(L\),
nowhere identically zero on a fibre. Its divisor is independent of the
frame. Let \(D\) be its reduced support. On \(V_0\), the divisor is
\(2S_v\) and \(D\) is exactly \(S_v\).

\item The map \(D\to V'\) is proper, since \(D\) is closed in a proper
space. There is no vertical divisor component: a vertical divisor in
the smooth total space would contain a whole irreducible fibre,
contrary to the nonzero restriction of \(\sigma\). Every horizontal
component has proper image all of the connected base. Since the
interior fibre is irreducible, there is exactly one horizontal
component, and its multiplicity in \(\operatorname{div}\sigma\) is
two. Thus \(\operatorname{div}\sigma=2D\) on this base. In the smooth
ambient space \(D\) is Cartier. It is flat over the curve because
its local rings are torsion-free over the base discrete valuation
ring. Restriction of the Cartier identity to an interior fibre gives
\(2D_v=2S_v\), whence \(D_v=S_v\) with its reduced structure.
Flatness and the smoothness of these fibres make \(D\to V'\) smooth
over \(V_0\). Its analytic nonsmooth locus therefore misses \(V_0\);
the proper mapping theorem once more makes its image discrete.
After deleting that image, \(D\to V^{\rm good}\) is a smooth proper
family of connected surfaces. Each deletion is closed discrete in the
then-current open curve. Their finite union is at most countable,
and the final good locus is open.

\item Ehresmann's theorem identifies the smooth fibres along every path;
in particular their first Betti number is two. \cite[Theorem~11, p.~297]{Buchdahl1999} then gives a K\"ahler metric on every fibre. Hence
\(h^{1,0}=1\), every holomorphic one-form is closed, and the integral
period map has a rank-two lattice image. Constant \(h^{1,0}\),
properness, and analytic base change supply a local relative
holomorphic one-form. Locally on a simply connected
base disc, a relative one-form and a basis of the free part of
\(H_1\) give two holomorphic periods. Their ratio lies in one
half-plane after orienting the basis. To see holomorphicity directly,
choose local primitives of the closed relative form on simply
connected fibre charts; their differences are holomorphic functions
of the base, and the finite sum around a transported cycle is its
period. The modular invariant of the ratio is independent of all
choices, so these local invariants glue to a holomorphic function on
\(V^{\rm good}\).
\end{enumerate}
\end{proof}

\begin{theorem}[Nonextension across the parameter boundary]\hypertarget{tag-032C}{}\label{push:src:D.1.3}
\StatementRecord{push:src:D.1.3}{proved}{032C}{\SeriesRef{fam:src:A.1.1}, \SeriesRef{fam:src:A.1.2}, \SeriesRef{fam:src:A.4.1}, \SeriesRef{push:src:D.1.2}, \SeriesRef{push:src:D.1.1}}

For every \(w_0\) with \(|w_0|=216\) and every connected open
neighbourhood \(U\) of \(w_0\), there is no flat proper holomorphic
map \(\mathcal X\to U\) from a reduced complex space whose fibres
over \(U\cap\{0<|w|<216\}\) are biholomorphic, fibre by fibre, to
\(X_w\). In particular an extension isomorphic to the given family
over that intersection does not exist.
\end{theorem}


\begin{proof}
\begin{enumerate}[label=\arabic*.,leftmargin=*]
\item We may restrict a putative extension to a small disc meeting the
circle, with closure away from zero. This is a holomorphic base
restriction; properness and flatness are preserved. The calculation
of the canonical evaluation in \cref{fam:src:A.1.1} gives
\(K_{X_w}=\mathcal O(-2S_2(w))\) and
\(h^0(K_{X_w}^{-1})=1\). Its sole divisor is \(2S_2(w)\), so it is
intrinsic under every fibrewise biholomorphism. The surface has
\(b_1=2\): on its finite product cover the rotation kills the moving
factor in invariant first cohomology and fixes the elliptic factor.
These are the exact \cref{fam:src:A.1.1}, \cref{fam:src:A.1.2} and \cref{fam:src:A.4.1} steps required in Lemma~\ref{push:src:D.1.2}.

\item The at most countable exceptional set in that lemma cannot contain
an arc of the circle. Choose a boundary point outside it and then a smaller disc
on which the reduced divisor is a smooth proper surface family. Its
holomorphic Albanese invariant agrees, pointwise on the inner part,
with \(J(iw/216)\). The pointwise agreement gives an analytic continuation through
that point, contrary to Proposition~\ref{push:src:D.1.1}.
\end{enumerate}
\end{proof}

\subsection{The local quotient and its Hausdorff reflection}\hypertarget{tag-02ZA}{}\label{push:section:D.2}

\begin{lemma}[The quadratic normal form at the boundary]\hypertarget{tag-032D}{}\label{push:src:D.2.1}
\StatementRecord{push:src:D.2.1}{proved}{032D}{\SeriesRef{fam:src:C.2.2}}

Near the order-four root centre there is a holomorphic coordinate
\[
 h=\mu-(1-\tau)/2
\]
with nonzero derivative in \(s_2\). In the fixed complex-linear fibre
coordinates \(\eta=\zeta_1\), \(\xi=\zeta_2+\zeta_1/2\), and the
integral ordered basis
\((\widehat w,\widehat u,\widehat\delta,
\widehat\gamma+3\widehat u-3\widehat w)\), the periods at \(c_*\) are
\begin{equation}\hypertarget{tag-01ON}{}\label{eq:d-period-normal}
 (1,\tfrac12),\quad(\tau,\tfrac12+h),\quad(0,1),\quad(6h,a+B(h)),
\end{equation}
where
\[
 \tau=i-\frac{3\pi i}{2}h^2+O(h^4),\qquad
 B(h)=-3ih^2+O(h^4),\qquad
 D_{c_*}=-3|h|^2+O(|h|^4).
\]
For \(c\) near \(c_*\), only the last \(\xi\)-period changes, by
\(c-c_*\).
\end{lemma}


\begin{proof}
\begin{enumerate}[label=\arabic*.,leftmargin=*]
\item The basis change has determinant of absolute value one. Substitution
into the original period matrix gives \eqref{eq:d-period-normal}
exactly, with
\(B=\beta+6h+3(1-\tau)/2-a\).
At the centre \(B(0)=0\). Put
\(G(h,\tau)=\wp_\tau((1-\tau)/2+h)\) and
\(H(h,\tau)=\log(\vartheta_3(h,\tau)/\eta_D(\tau))\).
The half-period identities and theta heat equation give
\[
 G(0,i)=0,\quad G_{h^2}(0,i)=-g_2(i)/4,\quad
 G_\tau(0,i)=\frac{i g_2(i)}{6\pi},
\]
\[
 H(0,i)=\log\sqrt2,\qquad H_\tau(0,i)=0,\qquad
 H_{h^2}(0,i)=-\pi/2.
\]
Here \(G_{h^2}\) is the power-series coefficient. For the derivative
constant, write \(A=\vartheta_2(0,i)^4=\vartheta_4(0,i)^4\).
Then \(g_2(i)=4\pi^4A^2\), while differentiating
\(e_2=(\pi^2/3)(\vartheta_2^4-\vartheta_4^4)\) with the heat
identities gives \(e_2'(i)=2\pi^3iA^2/3\).
The logarithmic derivatives used in this differentiation are
\[
 \frac{\partial_\tau\vartheta_2}{\vartheta_2}
  =\frac{\pi i}{12}(E_2+\vartheta_3^4+\vartheta_4^4),\qquad
 \frac{\partial_\tau\vartheta_4}{\vartheta_4}
  =\frac{\pi i}{12}(E_2-\vartheta_2^4-\vartheta_3^4).
\]
They follow by applying the theta heat equation to the half-period
identities. In particular \(G_\tau(0,i)\ne0\).
The holomorphic implicit-function theorem and evenness of \(G\) give
\(\tau=\tau(h)\), even in \(h\), with the displayed quadratic
coefficient. Since \(\tau-i\) has order two in the root coordinate,
\(h\) has order one.

\item The identity verified in \cref{fam:src:C.2.2} is
\[
 \log\mathcal N=\operatorname{Re}H-
       \frac{\pi(\operatorname{Im}h)^2}{\operatorname{Im}\tau}.
\]
The displayed coefficients, \(\tau-i=O(h^2)\), and evenness give
\(\log\mathcal N-\log\sqrt2=-(\pi/2)|h|^2+O(|h|^4)\).
Multiplication by \(6/\pi\), using the boundary normalization, gives
the formula for \(D_{c_*}\). Directly from the definition,
\[
 D_{c_*}=\operatorname{Im}B-
             6(\operatorname{Im}h)^2/\operatorname{Im}\tau.
\]
Consequently \(\operatorname{Im}B=-3\operatorname{Re}(h^2)
+O(|h|^4)\). Since \(B\) is holomorphic and \(B(0)=0\), its Taylor
coefficients force \(B=-3ih^2+O(h^4)\). 
\end{enumerate}
\end{proof}

\begin{proposition}[Nonseparation beyond the rank-drop line]\hypertarget{tag-032E}{}\label{push:src:D.2.2}
\StatementRecord{push:src:D.2.2}{proved}{032E}{\SeriesRef{push:src:D.2.1}}

At the boundary value \(a=0\), the central points with
\((\eta,\xi)=(0,0)\) and \((0,-3i)\) cannot be separated in the
root-family quotient. Their difference lies outside the sum of
\(\mathbb C(2,-1)^t\) and the central period group.
\end{proposition}


\begin{proof}
\begin{enumerate}[label=\arabic*.,leftmargin=*]
\item Choose \(h_n=1/n\), which defines root points by Lemma~\ref{push:src:D.2.1},
and take the integral lattice vector
\[
 \lambda_n=n^2(\widehat\gamma+3\widehat u-3\widehat w)
            -6n\widehat w+3n\widehat\delta.
\]
Its period is exactly \((0,n^2B(1/n))\), and this tends to
\((0,-3i)\). Thus the orbit of \((h_n,0)\) also has a representative
converging to \((0,0,-3i)\). The two central quotient points are
nonseparated. In these coordinates the line $\mathbb C(2,-1)^t$ is
\(\{\xi=0\}\), and every central period has real \(\xi\)-coordinate;
therefore their difference has nonzero imaginary second coordinate modulo that line and the central periods. The analytic error is
\(n^2O(n^{-4})\to0\), so the counterexample is for the family.
\end{enumerate}
\end{proof}

\begin{theorem}[The Hausdorff reflection of the boundary family]\hypertarget{tag-032F}{}\label{push:src:D.2.3}
\StatementRecord{push:src:D.2.3}{proved}{032F}{\SeriesRef{push:src:D.3.1}, \SeriesRef{eq:d-dual-exact}, \SeriesRef{push:src:D.2.2}}

Let \(\mathcal T=(\Delta_h\times\mathbb C^2)/\Lambda\) at the fixed
boundary parameter \(c_*\), after shrinking \(\Delta_h\). Two
central points represented by \((\eta,\xi)\) and
\((\eta',\xi')\) are nonseparated precisely when
\[
 \operatorname{Re}(\xi'-\xi)\in G_a
       :=\tfrac12\mathbb Z+a\mathbb Z\quad(a\in\mathbb Q),
\]
and every pair is nonseparated when \(a\notin\mathbb Q\).
There is no nonseparation between distinct base fibres, or between
distinct points of a noncentral fibre. The Hausdorff reflection
\(\mathcal T^{\mathrm H}\to\Delta_h\) has the original complex tori
over \(h\ne0\), and its central fibre is
\[
 \mathbb R/G_a\simeq S^1\quad(a\in\mathbb Q),\qquad
 \{*\}\quad(a\notin\mathbb Q).
\]
If \(a=m/n\) is in lowest terms, \(n>0\), then
\(G_a=\operatorname{lcm}(2,n)^{-1}\mathbb Z\).
\end{theorem}


\begin{proof}
\begin{enumerate}[label=\arabic*.,leftmargin=*]
\item The limiting-group formula in \ref{push:src:D.3.1} also
applies to the fixed boundary family with \(c=c_*\), \(h\ne0\).
Put \(d_0=\operatorname{lcm}(2,n)\) in the rational case. In the dual
lattice, the fixed integral covector whose central limit is
\(d_0\operatorname{Re}\xi\) has
\((p,q,r)=(0,0,d_0)\). The exact dual formula
\eqref{eq:d-dual-exact} below shows that it extends continuously, as a
real-linear functional \(\ell_h\) on \(\mathbb C^2\), to that limit:
its last coefficient is \(O(h^2)\), its other moving coefficients
are \(O(h)\). Therefore
\[
 (h,z)\longmapsto\exp(2\pi i\ell_h(z))
\]
is continuous, invariant under every period, and restricts at the
centre to \(\exp(2\pi i d_0\operatorname{Re}\xi)\).
Points with different indicated central classes are separated by this
map to the circle.

\item For the converse take a ray \(h=\epsilon\omega\),
\(\epsilon\downarrow0\), whose real and imaginary parts have
irrational ratio. More precisely choose \(\omega\) so that
\(p\operatorname{Re}\omega+q\operatorname{Im}\omega\ne0\)
for every nonzero pair \((p,q)\) in the countable half-integral
lattice in \ref{push:src:D.3.1}. The asymptotic dual formula shows that all such
covectors escape to infinity. For rational \(a\), the surviving dual
group is exactly \(\mathbb Z d_0\operatorname{Re}\xi\); for
irrational \(a\), it is zero. Pontryagin--Chabauty duality
\cite[Theorem~3.1, p.~2017; proof, p.~2018]{Cornulier2011} now gives the period-group limits
\[
 \{z:d_0\operatorname{Re}\xi(z)\in\mathbb Z\},\qquad
 \mathbb C^2,
\]
respectively. Every element of the indicated limit is the limit of
period vectors. Such a sequence supplies intersecting images
of every two prescribed central neighbourhoods, proving
nonseparation. This includes the quadratic-scale differences of
\ref{push:src:D.2.2}, not only first-order differences.

\item Collapse exactly these central classes and leave every other fibre
unchanged. The resulting quotient is Hausdorff. The base coordinate
separates different fibres. Within a noncentral fibre, continuous
functions supported over a compact subdisc away from zero separate
distinct points, because there the original action is properly
discontinuous and gives a smooth torus bundle. In the rational central
fibre the displayed circle-valued function separates different
classes. These functions descend through this quotient, so
it is Hausdorff. Every continuous map from \(\mathcal T\) to a
Hausdorff space identifies nonseparated pairs and therefore factors
uniquely and continuously through this quotient. This proves the
universal property, including its topology. The elementary greatest
common divisor calculation gives the stated value of \(G_a\).
\end{enumerate}
\end{proof}

\begin{proposition}[The finite boundary quotient and its leaves]\hypertarget{tag-032G}{}\label{push:src:D.2.4}
\StatementRecord{push:src:D.2.4}{proved}{032G}{\SeriesRef{fam:src:C.1.1}, \SeriesRef{fam:src:C.2.3}, \SeriesRef{fam:src:C.3.1}, \SeriesRef{fam:src:C.3.2}}

At the boundary, the construction over \(B\setminus\{1\}\) is a
complex manifold, since \cref{fam:src:C.1.1}, \cref{fam:src:C.2.3} makes every remaining finite period map
invertible and the cusp completion is unchanged.
The order-four deck action descends continuously to
\(\mathcal T^{\mathrm H}\). On its rational central circle it is
translation by \(-a/4\); on the irrational central point it is
trivial. Its finite quotient is Hausdorff with central fibre a circle
or a point, respectively. Neither quotient is a complex space over
the base compatible with the given holomorphic family off the centre.
The linear foliation by \(\mathbb C(2,-1)^t\) on the original central
torus cover has compact elliptic leaves. For rational \(a\) the
leaves are closed; for irrational \(a\) the closure of a leaf is the
set with its fixed value of \(\operatorname{Im}\xi\).
\end{proposition}


\begin{proof}
\begin{enumerate}[label=\arabic*.,leftmargin=*]
\item The deck transformation is a biholomorphism of the original quotient,
so the universal property extends it uniquely to a homeomorphism of
the Hausdorff reflection. The action on the fixed product coordinate
at the centre is \(\xi\mapsto\xi-\kappa_2/4=\xi-a/4\).
A finite group quotient of a Hausdorff space is Hausdorff: finite
orbits are separated by intersecting finitely many translated open
neighbourhoods. A finite rotation quotient of a circle is a circle.

\item On a closed small base disc every noncentral torus point has a
representative in the image of \([0,1]^4\) under its period matrix.
These representatives belong to a single compact subset of
\(\mathbb C^2\). At the centre their images meet every class of the
circle or the point. The image of that compact parameterized
parallelepiped is therefore the full inverse image of the closed
disc in the reflection. The reflected map is proper; the same holds
after the finite quotient. If it were holomorphic on a compatible
complex space, the central fibre would be analytic. A circle cannot
be a complex analytic space: a positive-dimensional analytic germ
has real dimension at least two, and a compact zero-dimensional
analytic space is finite. A central point is also impossible here:
upper semicontinuity of fibre dimension, or the local dimension
inequality at a point in the closure of the three-dimensional smooth
locus, forces a fibre dimension of at least two for a holomorphic
map to a curve. Thus this Hausdorff reflection cannot be glued as a
complex-space filling to the remaining constructed manifold.

\item For the leaf assertion, the moving line is \(\{\xi=0\}\).
At the centre the projection of the period group to the \(\xi\)-line
is \(G_a\subset\mathbb R\). Its intersection with the moving line
is a rank-two lattice in \(\mathbb C_\eta\): the periods are
\(p+qi\) with \((p+q)/2+ka+\ell=0\) for integers
\(p,q,k,\ell\). For irrational \(a\) this is the index-two lattice
\(p+q\equiv0\pmod2\); for rational \(a\) it is either that lattice
or \(\mathbb Z[i]\). Both give intrinsic elliptic leaf manifolds isomorphic to
\(E_i\). In the irrational case their maps into the non-Hausdorff
central quotient need not be topological embeddings; ``leaf'' here
means the image with this intrinsic immersed topology. When \(G_a\)
is discrete, the saturated leaf is closed.
When it is dense in \(\mathbb R\), its saturation has closure
\(\mathbb C_\eta\times(\operatorname{Re}\xi+\mathbb R)\), with
\(\operatorname{Im}\xi\) fixed. The natural map from this central quotient to the central fibre of $\mathcal T^{\mathrm H}$ is $[(\eta,\xi)]\mapsto\operatorname{Re}\xi\bmod G_a$ for rational $a$, and the constant map for irrational $a$, as in Theorem~\ref{push:src:D.2.3}. It is well defined because the real second coordinate of every central period belongs to $G_a$.
\end{enumerate}
\end{proof}

\subsection{All local limiting period groups}\hypertarget{tag-02ZB}{}\label{push:section:D.3}

\begin{theorem}[All marked limiting period groups]\hypertarget{tag-032H}{}\label{push:src:D.3.1}
\StatementRecord{push:src:D.3.1}{proved}{032H}{\SeriesRef{push:src:D.2.1}, \SeriesRef{eq:d-dual-exact}}

Let \(c_n\to c_*\) with \(\operatorname{Im}c_n<H_*\), and let
\(h_n\to0\). Set
\[
 \delta_n=-D_{c_n}(h_n)>0,\qquad
 V_n=\left(\frac{6\operatorname{Re}h_n}{\delta_n},
           \frac{6\operatorname{Im}h_n}{\delta_n},
           \frac{\operatorname{Re}(c_n-c_*)}{\delta_n}\right).
\]
Define the discrete group
\[
 L_a=\{(p,q,r)\in\mathbb R^3:
    r\in\mathbb Z,\ p+r/2\in\mathbb Z,
    \ q+r/2\in\mathbb Z,\ ar\in\mathbb Z\}.
\]
It has rank three for rational \(a\), and is
\(\mathbb Z^2\times\{0\}\) for irrational \(a\).
Every convergent subsequence of the closed lattices
\(\Gamma_n=\Pi_{c_n}(h_n)\mathbb Z^4\subset\mathbb C^2\) has the
following form. There is a primitive subgroup \(K\subset L_a\) and
a homomorphism \(\ell:K\to\mathbb R\) such that
\begin{equation}\hypertarget{tag-01OM}{}\label{eq:d-limit-group}
 \Gamma_\infty=
 \{(\eta,\xi):p\operatorname{Re}\eta+q\operatorname{Im}\eta
       +r\operatorname{Re}\xi+\ell(p,q,r)\operatorname{Im}\xi
       \in\mathbb Z\text{ for every }(p,q,r)\in K\}.
\end{equation}
Here
\[
 K=\{\lambda\in L_a:\lambda\cdot V_n
                      \text{ has a finite limit}\},\qquad
 \ell(\lambda)=\lim_n\lambda\cdot V_n,
\]
and for \(\lambda\notin K\) the absolute value tends to infinity.
Conversely every primitive \(K\subset L_a\) and every such
\(\ell\) occur for sequences with \(\operatorname{Im}c_n<H_*\).
Thus the possible pointed additive quotient limits are real tori of
any dimension from zero to three for rational \(a\), and from zero
to two for irrational \(a\), with the marked groups given explicitly
by \eqref{eq:d-limit-group}.
\end{theorem}


\begin{proof}
\begin{enumerate}[label=\arabic*.,leftmargin=*]
\item We use the fixed complex-linear coordinates of \ref{push:src:D.2.1}. Write
\(h=X+iY\), \(\tau=R+iT\), and
\(z=B(h)+c-c_*\). The real period matrix, with the basis of \ref{push:src:D.2.1},
is
\[
 A=\begin{pmatrix}
 1&R&0&6X\\
 0&T&0&6Y\\
 1/2&1/2+X&1&a+\operatorname{Re}z\\
 0&Y&0&\operatorname{Im}z
 \end{pmatrix},\qquad
 \det A=T D_c.
\]
For an integral dual label \((m_w,m_u,m_d,m_g)\), put
\(r=m_d\), \(p=m_w-r/2\), \(q=m_u-r/2\).
The real-linear covector \(A^{-t}(m_w,m_u,m_d,m_g)^t\) is
\begin{equation}\hypertarget{tag-01OL}{}\label{eq:d-dual-exact}
 \left(p,\frac{q-pR-rX-tY}{T},r,t\right),\quad
 t=\frac{m_g-ra-6pX-6Y(q-pR-rX)/T-r\operatorname{Re}z}{D_c}.
\end{equation}
Multiplication by \(A^t\) proves the formula.

\item A bounded sequence of dual vectors has bounded integral labels,
since \(A_n\) is bounded. After subselection its label is constant.
Its limit must satisfy \(m_g=ra\), because
\(A_0^t\) has that relation. The remaining labels are exactly
\(L_a\). Conversely consider a fixed label from \(L_a\).
Lemma~\ref{push:src:D.2.1} gives
\[
 T=1+O(|h|^2),\quad R=O(|h|^2),\quad
 \operatorname{Re}B=6XY+O(|h|^4),
\]
\[
 \delta_n=-\operatorname{Im}(c_n-c_*)
                       +3|h_n|^2+O(|h_n|^4).
\]
In particular \(\delta_n\ge C|h_n|^2\) for a fixed \(C>0\)
after shrinking the disc. The numerator in
\eqref{eq:d-dual-exact}, with \(m_g=ra\), is
\[
 -6pX-6qY-r\operatorname{Re}(c-c_*)
       +O(|h|^3)(|p|+|q|+|r|).
\]
The cancellation of the quadratic real term is essential. Division
by \(D_c=-\delta_n\) makes the error \(O(|h_n|)\). Thus
\[
 t_n=(p,q,r)\cdot V_n+o(1).
\]
Whenever this is bounded, the other three coefficients converge to
\((p,q,r)\). This proves that the dual limit, if it exists, is the
discrete graph
\(\{(p,q,r,\ell(p,q,r)):(p,q,r)\in K\}\).
There can be no further elements from varying labels, because a
bounded sequence of labels lies in a finite set. A convergent dual
group cannot contain two different limits with the same label: their
difference would have label zero, whose sole preimage under every
\(A_n^t\) is zero. Nor can a surviving label escape along a
subsequence, since an element of a Chabauty limit must be approximable
along the whole convergent sequence. Consequently every label either
has the indicated finite limit or escapes to infinity. Addition and
integer division show that \(K\) is a subgroup and is primitive.

\item By \cite[Theorem~3.1, p.~2017]{Cornulier2011}, the map from
closed subgroups of a finite-dimensional real vector space to their
Pontryagin annihilators is a homeomorphism. Applying it to the dual graph gives exactly
\eqref{eq:d-limit-group}. This theorem's hypotheses are met by the
real vector space underlying \(\mathbb C^2\) and by the closed full
lattices \(\Gamma_n\).

\item To prove realizability, extend \(\ell\) linearly to a covector
represented by a vector \(v_0\in\mathbb R^3\). Since \(K\) is
primitive, \(L_a\cap\operatorname{span}_{\mathbb R}K=K\).
Choose \(v\) orthogonal to that span and not orthogonal to any member
of \(L_a\setminus K\); the excluded proper hyperplanes are
countably many, and cannot cover the orthogonal space. If
\(K=L_a\), take \(v=0\). Put \(V_n=nv+v_0=(A_n',B_n',C_n')\).
Choose positive \(d_n\to0\) with
\(d_n(1+\|V_n\|^2)\to0\), and define
\[
 h_n=d_n(A_n'+iB_n')/6,\quad x_n=d_nC_n',\quad
 \epsilon_n=d_n+D_{c_*}(h_n)>0,
 \qquad c_n=c_*+x_n-i\epsilon_n.
\]
The inequality follows from \(D_{c_*}(h)=-3|h|^2+O(|h|^4)\).
These choices give \(-D_{c_n}(h_n)=d_n\) exactly, so their vectors
are the prescribed \(V_n\). On \(K\) the scalar products converge
to \(\ell\), and off \(K\) they escape. This constructs every
claimed limit inside the admissible parameter domain.

\item Finally the dual graph is a discrete free group of rank
\(\operatorname{rank}K\). Its annihilator quotient is a compact
connected real torus of that dimension. In the fixed Euclidean metric,
\(d_n([z],[z'])=\operatorname{dist}(z-z',\Gamma_n)\).
Chabauty convergence gives local uniform convergence of these
functions: minimizers relevant on a bounded set can be restricted to
a bounded set, and the two defining Chabauty approximation properties
then give both inequalities. This identifies the stated pointed
quotient limits. Their group structures and quotient metrics are the ones displayed above.
\end{enumerate}
\end{proof}

\begin{corollary}[Path dependence of the limiting torus]\hypertarget{tag-032I}{}\label{push:src:D.3.2}
\StatementRecord{push:src:D.3.2}{proved}{032I}{\SeriesRef{push:src:D.2.3}, \SeriesRef{push:src:D.3.1}, \SeriesRef{push:src:D.1.3}}

There is no path-independent limiting group at the boundary even
when \(c_n\to c_*\) from the admissible side. At \(h_n=0\) and
\(c_n=c_*-i\epsilon_n\), all of \(L_a\) survives, giving a
three-dimensional real torus if \(a\) is rational and a
two-dimensional real torus if \(a\) is irrational. Other admissible
paths give a point. The fixed-boundary Hausdorff reflection of \ref{push:src:D.2.3}
therefore cannot be identified with all pointed fibre limits.
\end{corollary}


\begin{proof}
\begin{enumerate}[label=\arabic*.,leftmargin=*]
\item The first sequence has \(V_n=0\), hence \(K=L_a\) and \(\ell=0\).
For the second, take \(K=0\) in the explicit realization in \ref{push:src:D.3.1}.
The limits have different real dimensions. Away from \(t=1\), the
periods stay invertible and the local complex families converge to
their ordinary value at \(c_*\); the ambiguity is at the specified
root centre. Theorem~\ref{push:src:D.1.3} excludes a flat proper holomorphic
extension across the parameter boundary.
\end{enumerate}
\end{proof}

\subsection{The supercritical locus}\hypertarget{tag-02ZC}{}\label{push:section:D.4}

\begin{proposition}[The small supercritical level curve]\hypertarget{tag-032J}{}\label{push:src:D.4.1}
\StatementRecord{push:src:D.4.1}{proved}{032J}{\SeriesRef{push:src:D.2.1}, \SeriesRef{fam:src:C.1.1}, \SeriesRef{fam:src:C.2.3}}

For \(\operatorname{Im}c=H_*+\epsilon\), with
\(0<\epsilon\ll1\), the zero locus \(D_c=0\) is one real-analytic
simple closed curve surrounding \(t=1\). In the coordinate of \ref{push:src:D.2.1},
\[
 |h|^2=\epsilon/3+O(\epsilon^2).
\]
In the fixed coarse coordinate it has the leading shape
\[
 |t-1|=\frac{\pi^4E_4(i)}4\epsilon^2
                         \bigl(1+O(\epsilon)\bigr).
\]
The error is uniform along the curve. The region containing \(t=1\)
has \(D_c>0\); the region containing the cusp has \(D_c<0\).
\end{proposition}


\begin{proof}
\begin{enumerate}[label=\arabic*.,leftmargin=*]
\item The Hessian of \(D_{c_*}(h)=-3|h|^2+O(|h|^4)\) is negative
at zero. Its radial derivative is \(-6|h|+O(|h|^3)\), uniformly
in angle. Consequently each sufficiently small positive level of
\(-D_{c_*}\) has exactly one radius in each direction and is a
real-analytic simple closed curve, with the displayed expansion.
The order-four deck action preserves it, acts freely on it, and the
coarse map \(t\) is a four-fold covering of its image. Since the
punctured root quotient is a punctured disc, that image is again a
simple closed curve. The strict global maximum in
Theorem~\ref{fam:src:C.2.3} gives a uniform gap on the
complement of any fixed root neighbourhood; the cusp tends to minus
infinity. Thus no other zero-locus component exists for small enough
\(\epsilon\).

\item Ramanujan's differential identities give
\(j''(i)=-2\pi^2j(i)E_4(i)\): indeed
\(j'=-2\pi i jE_6/E_4\) and
\(E_6'(i)=-\pi iE_4(i)^2\). Since \(j=1728t\), Lemma~\ref{push:src:D.2.1} gives
\[
 t-1=\frac{9\pi^4E_4(i)}4h^4+O(h^6).
\]
Insertion of the radial expansion proves the coarse-coordinate
formula, including its order of error. The sign statement follows
from the strict radial decrease near zero and from the cusp sign.
\end{enumerate}
\end{proof}

\begin{proposition}[The maximal open locus and the negative region]\hypertarget{tag-032K}{}\label{push:src:D.4.2}
\StatementRecord{push:src:D.4.2}{proved}{032K}{\SeriesRef{fam:src:C.3.1}, \SeriesRef{fam:src:C.3.2}, \SeriesRef{fam:src:C.1.1}, \SeriesRef{fam:src:C.2.3}, \SeriesRef{fam:src:A.1.1}, \SeriesRef{fam:src:A.1.2}}

For any height \(H=\operatorname{Im}c>H_*\), the zero locus is the
well-defined modular-norm level set
\[
 Z_H=\left\{t:\mathcal N(\mu(t),\tau(t))
                    =2\sqrt3\,e^{-\pi H/6}\right\},\qquad
 \wp_{\tau(t)}(\mu(t))=0.
\]
It is compact away from the cusp, and is a real-analytic set in the
finite root charts. At every regular point it is a smooth real curve.
The construction gives a complex manifold precisely over
\(B\setminus Z_H\), retaining the original Mumford fibre at the cusp.
The open sublevel set \(\{D_c<0\}\), with the cusp adjoined, is
connected. At every point of \(Z_H\), the period quotient is the
rank-three quotient of \cref{fam:src:C.3.1}, \cref{fam:src:C.3.2}; its exact Hausdorffization is described
below.
\end{proposition}


\begin{proof}
\begin{enumerate}[label=\arabic*.,leftmargin=*]
\item The formula for \(D\) in \cref{fam:src:C.1.1}, \cref{fam:src:C.2.3} gives the equation, and the theta
transformation law shows that it is independent of the modular and
period representatives. Compactness away from the cusp follows from
\(D_{\rm part}\to-\infty\). The analytic and regular-level
assertions follow in each root chart from the real-analytic implicit
function theorem. On the complement, the real period matrix is
invertible. On compact subsets of that complement its inverse is
bounded, which proves uniform discreteness and proper discontinuity
of the lattice action. The quotient is therefore a complex manifold
with compact torus fibres. The finite deck action is free for the
same invariant-coordinate reason as \cref{fam:src:A.1.1}, \cref{fam:src:A.1.2}: a nontrivial power translates
that circle by a noninteger fraction, regardless of the sign of
\(D\). Near the cusp \(D<0\), so its existing holomorphic Mumford
completion and its overlaps remain available.
At a zero with irrational kernel line, the fibre quotient is already
non-Hausdorff. At a zero with rational kernel choose a primitive
integral kernel vector. Its holomorphic period vanishes at that root
point but is not identically zero nearby. Its first nonzero Taylor
coefficient, multiplied by arbitrarily large integer multiples and
approached along suitable root directions, yields nonseparated
central points. Thus no neighbourhood of a zero in the naive quotient
is a Hausdorff complex manifold. This justifies the asserted maximal
open base locus.

\item For connectedness, in any holomorphic coordinate \(z\) with
\(T=\operatorname{Im}\tau>0\), differentiation gives
\[
 \partial_z\partial_{\bar z}D_{\rm part}
 =-\frac3T\left|\mu'(z)-
         \frac{\operatorname{Im}\mu}{T}\tau'(z)\right|^2\le0.
\]
Thus the function is superharmonic on the root charts. This property
and the minimum principle descend under the finite quotient; one
can equally apply the principle on a root cover of a putative
minimum. A component of the strict sublevel set not containing the
cusp would have compact closure with boundary at level \(-H\), and
would attain a smaller interior minimum. The strong minimum
principle excludes this. A sufficiently small cusp neighbourhood is
contained in the sublevel set, so all components reaching the cusp
are the same one. Hence the sublevel set is connected.
\end{enumerate}
\end{proof}

\begin{proposition}[Closed period groups at a rank-three fibre]\hypertarget{tag-032L}{}\label{push:src:D.4.3}
\StatementRecord{push:src:D.4.3}{proved}{032L}{\SeriesRef{fam:src:C.3.1}, \SeriesRef{fam:src:C.3.2}}

At a zero of \(D_c\), let \(A:\mathbb R^4\to\mathbb C^2\) be
the real period map and \(\ker A=\mathbb R k\), with the kernel
vector stated in \cref{fam:src:C.3.1}, \cref{fam:src:C.3.2}. Define
\[
 L_k=\{m\in\mathbb Z^4:m\cdot k=0\},\qquad r=\operatorname{rank}L_k,
 \qquad V_k=\bigcap_{m\in L_k}\ker m.
\]
Then
\[
 \overline{A\mathbb Z^4}=A(\mathbb Z^4+V_k),\qquad
 \mathbb C^2/\overline{A\mathbb Z^4}\simeq
                         \mathbb R\times(S^1)^r
\]
as real topological groups. The original quotient is a complex Lie
group precisely when \(k\) is rational, in which case \(r=3\).
Otherwise it is non-Hausdorff. Every Chabauty limit of nearby full
lattices contains \(\overline{A\mathbb Z^4}\); the quotient topology
at the fixed zero and a limit of full lattices need not coincide.
\end{proposition}


\begin{proof}
\begin{enumerate}[label=\arabic*.,leftmargin=*]
\item The last three period columns are real-linearly independent, so the
rank is exactly three, as in \cref{fam:src:C.3.1}, \cref{fam:src:C.3.2}. The annihilator characters of
\(A\mathbb Z^4\) are the real covectors \(\lambda\) with
\(A^t\lambda\in\mathbb Z^4\). Their integral labels belong exactly
to \(L_k\). The smallest rational subspace containing \(k\) is
\(V_k\); the closure of the image of \(\mathbb Z^4\) in
\(\mathbb R^4/\mathbb R k\) is its image after adding \(V_k\).
For an elementary verification, the two groups have the same
annihilator characters, and characters separate closed subgroups of
a real vector space by putting such a subgroup into the form
\(\mathbb R^d\times\mathbb Z^e\). This proves the closure formula.
Inside the three-dimensional image hyperplane its connected part
has dimension \(3-r\), and its discrete quotient has rank \(r\).
The one transverse real direction gives the displayed group.

\item The image is discrete exactly when the kernel line has a primitive
integral generator: extend such a generator to an integral basis for
the converse; for necessity, a discrete subgroup of a real
three-space has rank at most three, forcing a nonzero integral
kernel. A nondiscrete finitely generated subgroup cannot be closed,
since a nondiscrete closed subgroup contains an uncountable vector
group. Finally each fixed integral period converges to its limiting
period; a Chabauty limit is closed and therefore contains their
closure.
\end{enumerate}
\end{proof}

\begin{proposition}[The derivative along the moving Weierstrass zero]\hypertarget{tag-032M}{}\label{push:src:D.4.4}
\StatementRecord{push:src:D.4.4}{proved}{032M}{the definitions and constructions stated in the proof}

On an ordinary modulus chart where \(P=\wp'_\tau(\mu)\ne0\), put
\[
 L=\partial_\mu\log\vartheta_1(\mu,\tau),\qquad
 Z=L+2\pi i\frac{\operatorname{Im}\mu}{\operatorname{Im}\tau}.
\]
Then the exact real-analytic derivative is
\[
 \partial_\tau D_{\rm part}
       =\frac{iZ}{4\pi^2P}(3ZP-2g_2).
\]
Away also from the ramification of \(j\), the critical-point equation
on the coarse base is therefore \(Z=0\) or \(3ZP=2g_2\).
\end{proposition}


\begin{proof}
\begin{enumerate}[label=\arabic*.,leftmargin=*]
\item Put \(E=\pi^2E_2/3\). The theta heat equation,
\(\wp=-\partial_\mu L-E\), the differential equation
\(\wp''=6\wp^2-g_2/2\), and
\(4\pi i E'=-2E^2+g_2/6\) give, with \(\mu\) held fixed,
\[
 \partial_\tau\wp=
 \frac{2\wp^2-2E\wp+L\wp'-g_2/3}{2\pi i}.
\]
Indeed differentiating the heat equation twice in \(\mu\) gives
\((\wp''-2(\wp+E)^2+2L\wp')/(4\pi i)-E'\), which simplifies to
the displayed formula. Implicit differentiation of \(\wp(\mu)=0\)
therefore yields
\(\mu'=iL/(2\pi)-ig_2/(6\pi P)\).
At a zero of \(\wp\), the same heat equation and the eta derivative
give
\[
 \frac{d}{d\tau}\log(\vartheta_1(\mu,\tau)/\eta_D(\tau))
       =L\mu'+\frac{L^2}{4\pi i}.
\]
Differentiating \(D=\operatorname{Im}\beta-
6(\operatorname{Im}\mu)^2/\operatorname{Im}\tau\), and writing
\(b=\operatorname{Im}\mu/\operatorname{Im}\tau\), now gives
\[
 \partial_\tau D=\frac3\pi Z(\mu'-b)
                         -\frac{3i}{4\pi^2}Z^2,
 \qquad \mu'-b=\frac{iZ}{2\pi}-\frac{ig_2}{6\pi P}.
\]
Substitution proves the formula. Where \(dt/d\tau\ne0\), the chain
rule transfers vanishing of this derivative to the coarse base.

\end{enumerate}
\end{proof}

\begin{proposition}[The equianharmonic crossing]\hypertarget{tag-032N}{}\label{push:src:D.4.5}
\StatementRecord{push:src:D.4.5}{proved}{032N}{\SeriesRef{push:section:D.4}, \SeriesRef{fam:src:C.1.1}, \SeriesRef{fam:src:C.2.3}, \SeriesRef{push:src:D.4.4}}

At the equianharmonic point \(\tau=\rho=e^{\pi i/3}\), with the
original period representative \(\mu=(2-\rho)/3\), one has
\[
 \mathcal N(\mu,\rho)=3^{1/4},\qquad
 D_{\rm part}(z_1)=-\frac3{2\pi}\log48.
\]
Consequently the zero locus passes through the coarse point \(t=0\)
exactly at the height
\[
 H_0=\frac3{2\pi}\log48
 =H_*+\frac3{2\pi}\log\frac43.
\]
It has \(D_c(0)<0\) below this height and \(D_c(0)>0\) above it.
In the fixed coarse coordinate the function is \(C^1\) at zero and
\[
 D_{\rm part}(t)=D_{\rm part}(0)
       +\frac7{32\pi}\operatorname{Re}t+O(|t|^{4/3}),\qquad
 \partial_tD_{\rm part}(t)=\frac7{64\pi}+O(|t|^{1/3}).
\]
Thus at \(H=H_0\) the level set is a regular \(C^1\) arc through
zero, tangent to the imaginary axis.
\end{proposition}


\begin{proof}
\begin{enumerate}[label=\arabic*.,leftmargin=*]
\item For a fixed modulus write \(T=\operatorname{Im}\tau\) and
\[
 u(z)=\log\mathcal N(z,\tau)
 =\log\left|\frac{\vartheta_1(z,\tau)}{\eta_D(\tau)}\right|
       -\frac{\pi(\operatorname{Im}z)^2}{T}.
\]
The theta transformation law makes \(u\) periodic for
\(\Lambda=\mathbb Z+\mathbb Z\tau\). It is locally integrable,
with \(u(z)=\log|z|+c_\tau+o(1)\) at zero, and has no other
singularity on the torus.  In the Jacobi product, with
\(Q=e^{2\pi i\tau}\) and \(X=e^{2\pi iz}\), the quotient is
\[
 \frac{\vartheta_1(z,\tau)}{\eta_D(\tau)}
 =iQ^{1/12}X^{-1/2}(1-X)
       \prod_{m\ge1}(1-Q^mX)(1-Q^mX^{-1}),
\]
up to a constant of absolute value one. For
\(z=x+iy\), \(0<y<T\), every product factor has horizontal
mean logarithm zero, by the power series for
\(\log(1-ae^{2\pi ix})\), \(|a|<1\). The series converges
uniformly on closed substrips. Thus
\[
 \int_0^1u(x+iy)\,dx
       =-\frac{\pi T}{6}+\pi y-\frac{\pi y^2}{T}.
\]
The logarithmic singularities at the strip endpoints are
integrable; alternatively one integrates first on closed substrips
and passes to the limit. Integration in \(y\) gives
\(\int_{\mathbb C/\Lambda}u\,dA/T=0\). A tilted fundamental
parallelogram has the same horizontal averages, by period one.
This proves the zero-mean normalization directly.

\item For \(\tau=\rho\) put \(\alpha=1+\rho\). Since
\(\rho^2=\rho-1\), multiplication by \(\alpha\) on
\((1,\rho)\) has integral matrix
\[
 \begin{pmatrix}1&-1\\1&2\end{pmatrix},\qquad \det=3.
\]
Moreover \(\alpha\mu=1\), \(3\mu=2-\rho\in\Lambda\),
and \(\mu\notin\Lambda\). Its kernel on the torus is therefore
\(K=\{0,\mu,-\mu\}\), and \(|\alpha|^2=3\).
The distributional Euclidean Laplacian is
\[
 \Delta u=2\pi\delta_0-\frac{2\pi}{T}.
\]
Indeed the holomorphic theta factor contributes its simple zero,
and the Gaussian contributes the displayed constant. Pulling back
by the covering \(z\mapsto\alpha z\) shows that
\(u(\alpha z)\) and \(\sum_{v\in K}u(z+v)\) have the same
Laplacian: both have one logarithmic mass at every point of \(K\)
and constant term \(-6\pi/T\). Their difference is consequently
a periodic harmonic function. The logarithmic singularities cancel
locally, so the removable-singularity theorem gives an ordinary
harmonic function on the compact torus. It is constant.
Multiplication by \(\alpha\) and translation preserve normalized
Haar measure. The mean just computed therefore makes that constant
zero, proving the exact distribution identity
\[
 \mathcal N(\alpha z,\rho)
       =\prod_{v\in K}\mathcal N(z+v,\rho).
\]
As \(z\to0\), divide by \(\mathcal N(z,\rho)\). The left
ratio tends to \(|\alpha|=\sqrt3\). Since \(\mathcal N\) is
even, the other two factors tend to \(\mathcal N(\mu,\rho)^2\).
Positivity now gives \(\mathcal N(\mu,\rho)=3^{1/4}\).

\item Finally the exact \cref{fam:src:C.1.1}, \cref{fam:src:C.2.3} normalization is
\[
 D_{\rm part}=\frac6\pi
          \log\frac{\mathcal N(\mu,\tau)}{2\sqrt3}.
\]
Substitution gives \(-3\log48/(2\pi)\). The boundary height is
\(H_*=3\log6/\pi\), so the stated difference is
\(3\log(4/3)/(2\pi)>0\). Addition of \(\operatorname{Im}c\)
proves the crossing height and sign assertions.

\item For regularity in the coarse coordinate, multiplication by \(\rho^2\) is an
isometry of the torus and preserves the zero-mean Green norm just
used. Also \((\rho^2-1)\mu=\rho-1\in\Lambda\). Hence its
nontrivial rotation fixes \(\mu\), and the real gradient of
\(\log\mathcal N\) vanishes there. In the notation of \ref{push:src:D.4.4} this
says \(Z(\rho)=0\). Write \(P_0=\wp'_\rho(\mu)\ne0\),
\(E=\pi^2E_2/3\), and \(b=\operatorname{Im}\mu/T\).
The identities of \ref{push:src:D.4.4} give \(\mu'=b\) at \(\rho\), since
\(g_2(\rho)=0\) and \(L=-2\pi ib\). The heat equation, with
\(\mu\) held fixed, gives
\[
 L_\tau=\frac{-P-2EL}{4\pi i},\qquad L_\mu=-E.
\]
When differentiating \(Z\) along \(\mu(\tau)\), its two terms
involving \(E b\) cancel and the derivative of \(b\) vanishes at
\(\rho\). Thus
\[
 \partial_\tau Z(\rho)=\frac{iP_0}{4\pi},\qquad
 \partial_{\bar\tau}Z(\rho)=0.
\]
Ramanujan's identity and the Weierstrass equation yield
\[
 g_2'(\rho)=-\frac{3i}{\pi}g_3(\rho)
            =\frac{3i}{\pi}P_0^2.
\]
For \(v=\tau-\rho\), it follows that
\[
 3ZP-2g_2=-\frac{21iP_0^2}{4\pi}v+O(|v|^2),\qquad
 \partial_\tau D_{\rm part}
       =\frac{21iP_0^2}{64\pi^4}v^2+O(|v|^3).
\]
Here all remainders are real analytic in \(v,\bar v\).
The exact normalization \(t=j/1728\) gives
\[
 t=-\frac{i g_3(\rho)}{\pi^3}v^3+O(v^4)
   =\frac{iP_0^2}{\pi^3}v^3+O(v^4).
\]
For example this follows immediately from
\(t=E_4^3/(E_4^3-E_6^2)\),
\(E_4'(\rho)=-(2\pi i/3)E_6(\rho)\), and
\(g_3=8\pi^6E_6/27\). Dividing the derivative expansions on
\(v\ne0\) proves
\(\partial_tD_{\rm part}=7/(64\pi)+O(|t|^{1/3})\).
The conjugate derivative has the conjugate expansion. The resulting
real Taylor expansion extends the coarse function as a \(C^1\)
function with nonzero real gradient, giving the asserted local arc
by the implicit function theorem. In particular this is not a
singular level at \(t=0\), despite the vanishing modulus derivative
at the order-three ramification point.

\end{enumerate}
\end{proof}

\begin{remark}[The Green-function pullback along the moving zero]\hypertarget{tag-032O}{}\label{push:src:D.4.6}
\StatementRecord{push:src:D.4.6}{proved}{032O}{\SeriesRef{push:src:D.4.4}, \SeriesRef{push:src:D.4.5}, \SeriesRef{push:src:D.4.7}, \SeriesRef{push:src:D.4.8}; \citet[equation~(1.1), p.~912, and Lemma~2.3, p.~919]{LinWang2010}}
For a fixed modulus $\tau$, let
\[
 G_\tau(z)=-\frac1{2\pi}\log\mathcal N(z,\tau),\qquad
 Z(z,\tau)=\partial_z\log\vartheta_1(z,\tau)
                 +2\pi i\frac{\operatorname{Im}z}{\operatorname{Im}\tau}.
\]
The normalization proved in Proposition~\ref{push:src:D.4.5} gives the zero-mean Green function, and differentiation gives
\[
 \partial_zG_\tau(z)=-\frac{Z(z,\tau)}{4\pi}.
\]
This is the critical equation of \citet[equation~(1.1), p.~912, and Lemma~2.3, p.~919]{LinWang2010}. The section of the universal elliptic curve used here is
$\iota(\tau)=(\mu(\tau),\tau)$, with $\wp_\tau(\mu(\tau))=0$, and its exact pullback is
\[
 D_{\rm part}(\tau)=-12G_\tau(\mu(\tau))
                  -\frac6\pi\log(2\sqrt3).
\]
Along this section, Proposition~\ref{push:src:D.4.4} computes the full derivative, including the variation of $\tau$:
\[
 \partial_\tau D_{\rm part}
   =\frac{iZ}{4\pi^2P}(3ZP-2g_2),\qquad P=\wp'_\tau(\mu).
\]
Thus the two factors controlling the level curves on ordinary charts are the pullback of the fixed-modulus critical factor $Z$ and the additional factor $3ZP-2g_2$. Propositions~\ref{push:src:D.4.7} and~\ref{push:src:D.4.8} supply their nonvanishing and boundary estimates; Proposition~\ref{push:src:D.4.5} gives the coarse equianharmonic crossing.
\end{remark}

\begin{proposition}[Nonvanishing near the elliptic points and cusp]\hypertarget{tag-032P}{}\label{push:src:D.4.7}
\StatementRecord{push:src:D.4.7}{proved}{032P}{\SeriesRef{ari:C.8.15}, \SeriesRef{push:src:D.4.4}, \SeriesRef{push:src:D.4.5}, \SeriesRef{ari:C.1.1}}

The function $D_{\rm part}$ has no critical point on the punctured
root neighbourhood $0<|\tau-i|\leq3/20000$, no critical point on the
coarse neighbourhood represented by $|\tau-\rho|\leq1/50$, and no
critical point where $\Im\tau\geq2$. At $\rho$ the two factors
$Z$ and $F=3ZP-2g_2$ of \ref{push:src:D.4.4} vanish, but the coarse gradient is
nonzero, as in \ref{push:src:D.4.5}.
\end{proposition}

\begin{proof}
\begin{enumerate}[label=\arabic*.,leftmargin=*]
\item The two norms agree by their exact sigma--theta formula in
\ref{ari:C.1.1}. Hence, with $L=\log\mathcal N$,
\[
 D_{\rm part}=\frac6\pi\bigl(L-\log(2\sqrt3)\bigr).
\]
Theorem~\ref{ari:C.8.15} proves that the coarse differential
is nonzero away from $t=1$, including at $t=0$, where
$L_t(0)=7/384$. It gives
$(D_{\rm part})_t(0)=7/(64\pi)$. Its two exact root-critical
orbits are $i$ and $\rho$, with the factor vanishing specified
in \ref{push:src:D.4.4}, \ref{push:src:D.4.5}.

\item The root-to-coarse map is unramified off those elliptic orbits.
The stated small disc about $i$ contains no other elliptic
representative, and removing its centre therefore removes its
only critical point. The coarse disc about $\rho$ has the
nonzero differential just calculated. Finally
\[
 \Im\frac{ai+b}{ci+d}=\frac1{c^2+d^2}\le1,\qquad
 \Im\frac{a\rho+b}{c\rho+d}
 =\frac{\Im\rho}{c^2+cd+d^2}\le1.
\]
Thus no elliptic-orbit critical point lies at height at least two.
These facts prove every claimed domain of nonvanishing.
\end{enumerate}
\end{proof}



\begin{proposition}[The supercritical level classification]\hypertarget{tag-032Q}{}\label{push:src:D.4.8}
\StatementRecord{push:src:D.4.8}{proved}{032Q}{\SeriesRef{ari:C.8.15}, \SeriesRef{push:src:D.4.1}, \SeriesRef{push:src:D.4.2}, \SeriesRef{push:src:D.4.3}, \SeriesRef{push:src:D.4.5}, \SeriesRef{ari:C.1.1}}

For every \(H>H_*\), the degeneracy locus \(Z_H\) is exactly one
embedded \(C^1\) Jordan curve in the finite \(t\)-plane. It is
real analytic unless \(H=H_0=3\log48/(2\pi)\); at that height it
is real analytic away from \(t=0\) and has the regular crossing
described in \ref{push:src:D.4.5}. There are no singular coarse levels at any
supercritical height. The bounded complementary disc contains
\(t=1\) and has \(D_c>0\); the complementary disc in
\(\mathbb P^1\), containing the cusp, has \(D_c<0\). The point
\(t=0\) is in the latter disc for \(H<H_0\), on the curve for
\(H=H_0\), and in the former disc for \(H>H_0\).

The bounded discs are strictly nested as \(H\) increases, shrink
to \(\{1\}\) as \(H\downarrow H_*\), and exhaust the finite
plane as \(H\to\infty\). Their small-height boundary shape is
the expansion in \ref{push:src:D.4.1}. The maximal open base locus of the complex
manifold construction consists of these two complementary discs.
Approaching their common boundary, the covolume of the marked
root torus lattice is
\[
 \operatorname{covol}(\Pi_c\mathbb Z^4)
       =\operatorname{Im}\tau\,|D_c|\longrightarrow0.
\]
At the boundary the period map has rank three and the exact quotient
and Hausdorff reflection are those of \ref{push:src:D.4.3}.
\end{proposition}

\begin{proof}
\begin{enumerate}[label=\arabic*.,leftmargin=*]
\item Use the same exact affine relation and write
\[
 f=-D_{\rm part},\qquad
 Z_H=\{L=\log(2\sqrt3)-\pi H/6\}.
\]
The maximum $L(1)=\log\sqrt2$ corresponds exactly to
$H_*=3\log6/\pi$. Thus $H>H_*$ corresponds to a level
strictly below that maximum. Theorem~\ref{ari:C.8.15} proves
that each such level is one strictly star-shaped $C^1$ Jordan
curve about $1$, real analytic away from the possible point
$t=0$. Its nonzero coarse derivative there proves regularity
also on that exceptional level.

\item Since $L(0)=\tfrac14\log3$, its height is
$H_0=3\log48/(2\pi)$. The exact equation $D_c=H-f$ gives
the two signs, the position of $0$, and the bounded component
$\{f<H\}$ containing $1$. Strict radial monotonicity proves
strict nesting. The unique minimum $f(1)=H_*$ and properness
of $f$ off $1$ prove shrinking as $H\downarrow H_*$ and
exhaustion of every compact subset as $H\to\infty$.
The expansion at the minimum is \ref{push:src:D.4.1}.

\item In the ordered real coordinates and marked lattice basis of
\ref{push:src:D.4.3}, direct expansion of the period determinant
gives
\[
 \det_{\R}\Pi_c
   =\Im\tau\,\Im\beta_c-6(\Im\mu)^2
   =\Im\tau\,D_c.
\]
Its absolute value is the covolume and tends to zero at the
level curve. The same matrix retains three independent columns
there, so its rank is exactly three. Apply the exact quotient
and Hausdorff-reflection computation of \ref{push:src:D.4.3}.
The local proper-action criterion of \ref{push:src:D.4.2}
gives complex manifolds on both complementary regions and
excludes every point of their common boundary, including its
discrete fixed-fibre quotients. This proves the maximality and
all the local degeneration assertions.
\end{enumerate}
\end{proof}

\section{The limit at the puncture}\hypertarget{tag-01PS}{}\label{fam:section:E}
\subsection{The untwisted pinching and its local smoothings}\hypertarget{tag-01PT}{}\label{fam:section:E.1}
\hypertarget{tag-02Z0}{}\label{push:section:A.1}

\begin{definition}[The elliptic surface and the pinching datum]
\hypertarget{tag-PT01}{}\label{I:def:pinching-data}
\StatementRecord{I:def:pinching-data}{proved}{PT01}{}
$S$ is the smooth minimal rational elliptic surface with
Weierstrass equation
\[
 y^2=x^3-3t^3(t-1)x+2t^4(t-1)^2,
\]
$\pi:S\to B$ is its elliptic fibration, $P$ is the section used in the
period dictionary, $T=t_{6P}:S\to S$, and $L=M_S=\OO_S(P-O)$.
We use the convention that $\mathbb P_S(\OO\oplus L)$ contains the total
space of $L$, with fibre coordinate $a$, zero section $s_0$, and infinity
section $s_\infty$. The two normal bundles are respectively $L$ and
$L^{-1}$. We write $R=\mathbb P_S(\OO\oplus L)$.
\end{definition}

\begin{proposition}[The pinching and its smoothing line]\hypertarget{tag-01R1}{}\label{fam:src:E.1.1}\label{fam:st:E.1.1}
\StatementRecord{fam:src:E.1.1}{proved}{01R1}{\cite[Lemma 2.1]{Engel2026}}

The quotient
\[
 Q=R/\{s_0(x)\sim s_\infty(Tx):x\in S\}
\]
has a compact, reduced complex-space structure. Its normalization is
the finite holomorphic map $\nu:R\to Q$. Its singular locus is a copy
of $S$, and its local analytic equation along that locus is $ab=0$,
with two further coordinates on $S$. The induced map $Q\to B$ is
proper. Its sheaf of first-order local smoothing parameters is
\[
 \mathcal T^1_Q=\mathcal Ext^1_{\OO_Q}(\Omega_Q,\OO_Q)
       \simeq L\otimes T^*L^{-1}\simeq\pi^*\OO_B(1)
\]
on the double locus.
\end{proposition}



\begin{proof}
\begin{enumerate}[label=\arabic*.,leftmargin=*]
\item The two sections of $R$ are disjoint closed submanifolds. The
equivalence relation consists of their two graphs and the diagonal;
it is closed in $R\times R$. A quotient of a compact Hausdorff space
by a closed equivalence relation is Hausdorff: for distinct equivalence
classes, compactness gives disjoint saturated neighbourhoods by taking
finite subcovers of the complements of the relation. Thus the displayed
quotient is compact and Hausdorff.

\item For a coordinate neighbourhood $U\subset S$ and frames of $L|_U$ and
$L^{-1}|_{T(U)}$, the two tubular coordinate charts are $(x,a)$ at
$s_0(U)$ and $(Tx,b)$ at $s_\infty(T(U))$. We give their identification
the ring
\[
 \{(f(x,a),g(x,b)):f(x,0)=g(x,0)\}
   =\C\{x_1,x_2,a,b\}/(ab).
\]
Changes of frame are holomorphic and invertible on each branch and
have the same change of $x$ on their intersection. They therefore glue
these hypersurface charts to the inherited charts off the sections.
This constructs the asserted complex space. The direct sum of the two
regular branch rings is the integral closure: an integral element in
the total ring of fractions is integral on each smooth branch, and
conversely the two branch idempotents are integral. The resulting map
$\nu$ is proper with at most two points in each fibre, and its local
ring extension is finite. The map to $B$ is well defined because
$\pi T=\pi$; its properness follows from compactness of $Q$.

\item The conormal presentation for the hypersurface $ab=0$ gives
$\mathcal Ext^1(\Omega,\OO)=\OO/(a,b)$. A change of normal frames
$a'=g(x)a$, $b'=h(x)b$ changes the equation and its smoothing parameter
by $gh$. Consequently these local modules glue as
$N_{s_0/R}\otimes T^*N_{s_\infty/R}=L\otimes T^*L^{-1}$.
For this particular surface, Engel's Lemma 2.1 states
$\operatorname{Hom}(T^*L,L)=\pi^*\OO_B(1)$
\cite[Lemma~2.1, pp.~3--4]{Engel2026}.
Its identification applies here with $T=t_{6P}$ and $L=\OO(P-O)$. One can recover the degree in that lemma
directly: the bundle has multidegree zero on every fibre, since $T$
preserves every component; hence it is $\pi^*\OO_B(n)$. Its degree on
$O$ is
\[
 n=(O-P)\cdot 6P+(P-O)\cdot O=(2-2)+(0-(-1))=1.
\]
The section intersections used here are the explicit computations in
the proof of that lemma. This derives the displayed sheaf in the
normal-bundle convention of our pushout.
\end{enumerate}
\end{proof}

\begin{definition}[The smoothing section]
\hypertarget{tag-PT06}{}\label{I:def:smoothing-section}
\StatementRecord{I:def:smoothing-section}{proved}{PT06}{\SeriesRef{fam:src:E.1.1}; \cite[Corollary 2.2, p.~4]{Engel2026}}
Set $J=L\otimes T^*L^{-1}$. Proposition~\ref{fam:src:E.1.1} identifies $J$ with $\pi^*\mathcal O_B(1)$. We fix the section $\phi\in H^0(S,J)$ with divisor $S_\infty$ from \cite[Corollary 2.2, p.~4]{Engel2026}. Thus $\phi_x$ is the bundle map $L_{Tx}\to L_x$ used in the plumbing equation.
\end{definition}



\begin{proposition}[The nodal plumbing and its small modification]\hypertarget{tag-01R3}{}\label{fam:src:E.1.3}\label{fam:st:E.1.3}
\StatementRecord{fam:src:E.1.3}{proved}{01R3}{the definitions and constructions stated in the proof}

The first simultaneous local plumbing at the nodal point of $S_\infty$
has equation
\begin{equation}
                 xy=u zv,\qquad q=zv,
\label{pub:I:eq:E2}
\end{equation}
after a holomorphic change of one normal coordinate, where $q$ is a
local coordinate on $B$ at infinity. This hypersurface is flat over
the $u$-disc. For each $u\ne0$, its three-dimensional fibre has an
ordinary double point at $x=y=z=v=0$. The projective modification
defined by the blow-up of $(x,z)$ resolves that point for $u\ne0$.
Its two strict-transform charts are
\begin{equation}
       x=zr,\quad ry=uv;
       \qquad z=xa,\quad y=uav.
\label{pub:I:eq:E3}
\end{equation}
The modified family remains flat over $u$. At $u=0$, the normalization
of its central fibre contains the blow-up of one branch of \eqref{pub:I:eq:E2} along
the smooth curve $x=z=0$; hence this particular modification changes
the normalization of the unmodified pinching.
\end{proposition}



\begin{proof}
\begin{enumerate}[label=\arabic*.,leftmargin=*]
\item The surface $S$ is smooth and its nodal elliptic fibre has local
equation $q=zv$ at its node. The section $\phi$ has a simple zero
along that fibre. Thus $\phi=zv$ times an invertible holomorphic
function. Absorbing the function in one normal coordinate gives \eqref{pub:I:eq:E2}.
The convergent power-series ring $\C\{x,y,z,v,u\}$ is factorial,
and $u$ is coprime to $xy-uzv$. Hence $uh\in(xy-uzv)$ implies
$h\in(xy-uzv)$. The quotient has no $u$-torsion. At a nonzero
parameter the identical argument uses $u-u_0$. This proves flatness
over the smooth parameter curve.
For fixed $u\ne0$ the four first derivatives vanish only at the
origin, and the quadratic form $xy-uzv$ is nondegenerate. This proves
the ordinary-double-point assertion by its defining equation.

\item The two standard affine charts of the blow-up give \eqref{pub:I:eq:E3}, after
removing the common factors $z$ and $x$, respectively. For $u\ne0$
the first chart solves for $v$ and the second for $y$, so both are
smooth. The map is an isomorphism away from the origin in the fixed
$u$-fibre: where the ideal fails to be visibly principal, the
hypersurface equation and one of $y,v$ make it principal. The fibre
over the origin is $\mathbb P^1$. The modification is projective
over \eqref{pub:I:eq:E2}, since it is a closed subspace of its product with
$\mathbb P^1$. The first strict-transform equation is coprime to
$u-u_0$ in its convergent power-series ring for every $u_0$; the
second chart is smooth with free parameter $u$. Thus flatness over
$u$ persists.

\item At $u=0$ the first chart has $ry=0$. The branch $r=0$ maps
isomorphically to $x=0$ in the original central fibre. The branch
$y=0$, together with the second chart, is the ordinary blow-up of the
branch $y=0$ in coordinates $(x,z,v)$ along $(x,z)$. This description
follows from the two affine substitutions defining that blow-up and
identifies the normalizations.
For $u\ne0$, the base coordinate in the first chart is
$q=zry/u$ and in the second it is $q=xav$. Thus the modification also
exhibits the triple-coordinate strata that the cusp filling requires.

\end{enumerate}
\end{proof}

\begin{proposition}[The unsubdivided cusp fan]\hypertarget{tag-02ZH}{}\label{push:src:A.1.1}\label{push:A11}
\StatementRecord{push:src:A.1.1}{proved}{02ZH}{the definitions and constructions stated in the proof}

For $k,l\in\mathbb Z$, let $P_{k,l}$ be the cone defined by $h_q,h_u\geq0$, $kh_q\leq y_1\leq(k+1)h_q$ and $l(h_q+h_u)\leq y_2\leq(l+1)(h_q+h_u)$. These cones and their faces form a fan. Its two
period automorphisms on the cocharacter lattice are
\[
 G_1(y_1,y_2,h_q,h_u)=(y_1+h_q,y_2,h_q,h_u),\qquad
 G_2(y_1,y_2,h_q,h_u)=(y_1,y_2+h_q+h_u,h_q,h_u).
\]
The affine chart of $P_{0,0}$ is
\[
 U_{0,0}=\operatorname{Spec}\mathbb C[x,y,z,v,u]/(xy-uzv),
 \qquad q=zv.
\]
In particular the unsubdivided chart is not a smooth simultaneous
model of the cusp.
\end{proposition}


\begin{proof}
\begin{enumerate}[label=\arabic*.,leftmargin=*]
\item At fixed $h_q,h_u\geq0$, the cone $P_{k,l}$ cuts out the rectangle
\[
 kh_q\leq y_1\leq(k+1)h_q,\qquad
 l(h_q+h_u)\leq y_2\leq(l+1)(h_q+h_u).
\]
When $h_q=0$ the first interval is a point. These rectangles tile
their support, and their intersections are common faces, including
this collapsed boundary. This proves the fan assertion. The displayed
integral transformations preserve the fan and have integral inverses.

\item For a character $(a,b,c,d)$, the six inequalities of the dual cone
are $c,a+c,b+c,a+b+c,d,b+d\geq0$. The characters
\[
 x=(0,1,0,0),\quad y=(0,-1,1,1),\quad
 z=(1,0,0,0),\quad v=(-1,0,1,0),\quad u=(0,0,0,1)
\]
generate its integral monoid. Indeed if $b\geq0$ we subtract $b x$;
if $b<0$ we subtract $(-b)y$. The remaining character has second
coordinate zero and is a nonnegative integral combination of $z,v,u$.
The sole relation is $x+y=u+z+v$, and the character of the first
base coordinate is $z+v$. This proves the coordinate-ring assertion,
including normality. 
\end{enumerate}
\end{proof}

\begin{theorem}[The compact nodal plumbing]\hypertarget{tag-02ZI}{}\label{push:src:A.1.2}\label{push:A12}
\StatementRecord{push:src:A.1.2}{proved}{02ZI}{\SeriesRef{push:A11}; \cite{EGFS25}}
\hypertarget{tag-01R2}{}\label{fam:src:E.1.2}\label{fam:st:E.1.2}
\StatementRecord{fam:src:E.1.2}{proved}{01R2}{\SeriesRef{fam:src:E.1.1}; \cite{Engel2026}}

There is a reduced complex space $\mathcal X^\flat$ and a proper flat
holomorphic map
\[
 (f,\varpi):\mathcal X^\flat\longrightarrow B\times\Delta_u
\]
whose central threefold is $Q$. For $u\ne0$ its fibre is the
contraction of Engel's chosen $X(u)$ at one rational curve of the
Mumford fibre. This curve is the double curve corresponding to the
diagonal of the square in the chosen $A_2$ triangulation; in the
coordinates of the marked cusp fan its direction is $e_1-e_2$. Its normal bundle
in $X(u)$ is $\mathcal O_{\mathbb P^1}(-1)^{\oplus2}$. Its image is
the unique ordinary double point of $\mathcal X^\flat_u$.

Over $S^\circ=S\setminus S_\infty$ the total plumbing space is smooth. Its restriction over every relatively compact disc $U\Subset B\setminus\{\infty\}$ is proper and flat over $U\times\Delta_\epsilon$ for sufficiently small $\epsilon>0$. For $u\ne0$ it is the untwisted quotient by $u\phi$ covering $T^{-1}$. This includes discs containing $0$ or $1$, before the logarithmic transformations.
\end{theorem}


\begin{proof}
\begin{enumerate}[label=\arabic*.,leftmargin=*]
\item On the whole compact surface $S$, choose a
Hermitian metric on $L$, use the dual metric on $L^{-1}$, and choose
$0<r<1/4$. For sufficiently small $\epsilon>0$ we have
$|u|\|\phi_x\|<r^2$ for all $x\in S$ and $|u|<\epsilon$.
Take the core
\[
 C=\{(x,a,u):a\in L_x,\ r<\|a\|<r^{-1},\ |u|<\epsilon\}
\]
and the node space
\begin{equation}
 N=\{(x,a,b,u):a\in L_x,\ b\in L^{-1}_{Tx},\
 a\otimes b=u\phi_x,\ \|a\|,\|b\|<2r,\ |u|<\epsilon\}.
\label{pub:I:eq:E1}
\end{equation}
On $r<\|a\|<2r$ identify a point of $N$ with $(x,a,u)\in C$.
On $r<\|b\|<2r$ identify it with $(Tx,b^{-1},u)\in C$.
Both identifications are biholomorphisms between open complex spaces;
in the first collar $b=u\phi_x/a$, and in the second
$a=u\phi_xb^{-1}$. Their domains are disjoint because
$\|a\|\|b\|<r^2$. Changes of frames multiply the equation by
units and give the stated bundle-valued equation, so no choice of
frame enters the gluing.

\item This gluing is Hausdorff. At either end of a collar, exactly one
of its two charts contains the limiting circle: the core omits
$\|a\|=r$, and the node chart omits $\|a\|=2r$; the infinity
collar has the identical property in $b$. Thus the graph of the
identification has no unpaired pair of limiting points. The two
collars in each chart are separated, so there are no additional
chains of identifications. These observations give disjoint
saturated neighbourhoods of any two inequivalent points. Properness
is equally concrete. Over a compact subset of $B\times\Delta_\epsilon$,
the closed core $3r/2\leq\|a\|\leq2/(3r)$ together with the
closed node subspace $\|a\|,\|b\|\leq7r/4$ is compact and maps
onto the glued space. To see surjectivity, a point outside the smaller
core lies on a node collar, and a node point outside the smaller
node subspace lies on the smaller core. Compactness of $S$ supplies
the uniform bound used above. This proves properness of the joint
map to $(t,u)$.

\item The total node rings are hypersurface rings over the regular
two-dimensional local rings of $S$. They are Cohen--Macaulay.
Every fibre of the joint map has pure dimension two: first fix
$t$, and then impose $xy=u\phi$ on the two line coordinates over
the one-dimensional fibre of $S$. Consequently $t-t_0,u-u_0$ is a
regular sequence in every total local ring. The regular-sequence
flatness criterion over the smooth two-dimensional base proves
flatness. On $S^\circ$, the derivative of $ab-u\phi_x$ with respect to $u$ is $-\phi_x\ne0$, so the total node space is smooth. On the core the total space is an open subset of $L\times\Delta_u$
and is smooth; joint flatness follows from that of $S\to B$. This includes the nonreduced Kodaira fibres
at $0,1$, since the argument uses smoothness of $S$. At $u=0$ the node equation is $ab=0$,
and the two collar maps recover precisely the specified pinching
of the two sections of $R$.

\item For $u\ne0$ and $t\ne\infty$, the equation gives
$a_{Tx}=a_x/(u\phi_x)$. Hence the resulting quotient is exactly
$L^\times/\langle u\phi\rangle$, where the contracting generator
covers $T^{-1}$. Every orbit meets a core or a collar, by the
uniform contraction estimate; the displayed comparison identifies
exactly its orbit relation. It is therefore a biholomorphism.

\item At the nodal point of $S_\infty$, write $q=zv$. Absorbing a unit
in one line coordinate gives $xy=uzv$. At a smooth point of
$S_\infty$ the equation is $xy=uq$. For $u\ne0$ the latter
threefold is smooth; the former has one nondegenerate quadratic
singularity at the origin. There are no other singularities of a
nonzero $u$ fibre.

\item In a Tate covering of a neighbourhood of the $I_1$
fibre use the torus coordinates $(Z,W)$ of
\citet[Proposition~2.7, p.~6]{Engel2026}.
The two periods act as
\[
 (Z,W,q,u)\longmapsto(c_1qZ,c_2W,q,u),\qquad
 (Z,W,q,u)\longmapsto(c_3Z,c'_4uqW,q,u),
\]
with holomorphic units $c_i$. Their monomial parts are $G_1,G_2$.
The chart indexed by $(k,l)$ has
$z=Zq^{-k}$, $v=q^{k+1}/Z$, $x=W(qu)^{-l}$,
$y=qu/x$, with the unit changes given by the line
frames. It is the chart of Proposition~\ref{push:A11}.
The Tate cover unfolds the node of $S_\infty$ into its infinite
chain, and unfolding the two plumbing collars gives the additional
integer index $l$. The foregoing collar identifications are exactly
the two deck identifications. Thus, on an invariant neighbourhood
over a sufficiently small bidisc, the fan space maps by a
holomorphic covering to the plumbing, with deck group $\mathbb Z^2$.
Each core or node coordinate neighbourhood has disjoint lifts
indexed by the two unfolded integers; this proves freeness and
proper discontinuity, including the $\rho_l$ orbits. The first
period stabilizes such a ray but on its orbit translates the Tate
coordinate by $c_1q$, with $q\ne0$ and $|c_1q|<1$; it does not
fix a point. Properness of its quotient was proved by the compact
core and node representatives. This proves the analytic comparison of the fan quotient with the compact plumbing.

\item Subdividing the square at $h_u=0$ by the chosen diagonal gives
the small resolution of $xy=uzv$ for $u\ne0$. Put $Y=y/u$, so the equation is $xY=zv$.
The resolution associated with the diagonal $[e_1,e_2]$ has charts
\[
 (v,Y,a)\longmapsto(av,uY,aY,v),\qquad
 (x,z,b)\longmapsto(x,ubz,z,bx)
\]
in the coordinates $(x,y,z,v)$. On their overlap,
$b=a^{-1}$ and $(x,z)=a(v,Y)$. The character $a=x/v$ has
valuation $y_1+y_2-h_q$, so these are exactly the two triangles
separated by $[e_1,e_2]$. The exceptional curve is the zero
section in the two normal coordinates. Their transition is
multiplication by $a$, giving
$\mathcal O_{\mathbb P^1}(-1)\oplus\mathcal O_{\mathbb P^1}(-1)$. There is one
diagonal edge modulo the periods, so there is exactly one such
curve. Removing that edge contracts the paired opposite
$(-1)$-curves in the normalization $dP_6$ to the square model
$\mathbb P^1\times\mathbb P^1$. This identifies the contracted
curve in the selected toric model. The fixed double curve of
Proposition~\ref{mapsI:src:A.1.2} has direction $e_2$ in the same
marked lattice; the two primitive directions satisfy
$\det(e_1-e_2,e_2)=1$.
\end{enumerate}
\end{proof}

\begin{proposition}[The integral topology of the pinching]\hypertarget{tag-01R4}{}\label{fam:src:E.1.4}\label{fam:st:E.1.4}
\StatementRecord{fam:src:E.1.4}{proved}{01R4}{\SeriesRef{fam:src:E.1.1}; \cite[Lemma 2.1]{Engel2026}}

The compact pinching $Q$ of Proposition~\ref{fam:src:E.1.1} has
\[
 \pi_1(Q)=\Z,\qquad
 (H_i(Q;\Z))_{i=0}^6
   =(\Z,\Z,\Z^8,\Z^7,\Z^{10},0,\Z),\qquad\chi(Q)=12.
\]

\end{proposition}



\begin{proof}
\begin{enumerate}[label=\arabic*.,leftmargin=*]
\item The two section embeddings are cofibrations: a tubular neighbourhood
retracts onto each section, and the two neighbourhoods are disjoint.
The ordinary pushout is therefore homotopy equivalent to the space
formed from $R$ by attaching $S\times[0,1]$ along
$s_0$ and $s_\infty T$. This comparison follows by collapsing the
intervals in the mapping cylinder; the tubular retractions construct
its homotopy inverse. Both $S$ and $R$ are simply connected. Van
Kampen applied to this mapping-cylinder model gives one free loop
and no additional generator or relation, so $\pi_1(Q)=\Z$.
The same model has the exact sequence
\begin{equation}
 H_i(S)\xrightarrow{d_i=s_{0*}-s_{\infty*}T_*}H_i(R)
 \longrightarrow H_i(Q)\longrightarrow H_{i-1}(S).
\label{pub:I:eq:E4}
\end{equation}

\item Put $F=[S_t]$, $m=P-O$, and
$v_0=[6P]-[O]-3F\in H_2(S;\Z)$. The intersection computations in
\cite[Lemma~2.1, pp.~3--4]{Engel2026} give
\label{layout:I:2}
\begin{equation}
\begin{gathered}
F^2=Fv_0=Fm=0,\quad v_0^2=-6,\quad mv_0=-1,\quad m^2=-2,
 \\
\quad OF=1,\quad Ov_0=0,\quad Om=1.
\end{gathered}
\label{pub:I:eq:E5}
\end{equation}
The translation preserves all components of the two reducible
Kodaira fibres. Their root classes span a negative definite space of
rank seven. Its orthogonal complement is generated over $\Q$ by
$F,O,v_0$. Since $T_*O=O+v_0+3F$, preservation of the intersection
form gives, on this complement and hence on the entire lattice,
\begin{equation}
 T_*a=a+(aF)v_0-\bigl(av_0-3aF\bigr)F.
\label{pub:I:eq:E6}
\end{equation}
Its value on $v_0$ is $v_0+6F$: it has no $O$-component
and has zero intersections with the rank-seven root space, so it is
$\lambda v_0+aF$. The already established identity $T^*m=m-F$ gives
$mT_*v_0=(m-F)v_0=-1$, forcing $\lambda=1$. Orthogonality to
$T_*O$ then forces $a=6$. These
values on $F,O,v_0$ and the root space verify \eqref{pub:I:eq:E6} on a rational basis,
and both sides preserve the integral lattice. In particular
$T^*m=m-F$, in agreement with \ref{fam:src:E.1.1}.

\item The rational elliptic surface is the blow-up of $\mathbb P^2$ at nine
points, so its homology in even degrees is $\Z,\Z^{10},\Z$ and
vanishes in odd degrees; its intersection lattice is unimodular.
The projective-bundle theorem follows here by local trivialization
and the fibre hyperplane class: $H_2(R)=H_2(S)\oplus\Z f$ and
$H_4(R)=\Z\oplus H_2(S)$, with $f$ the projective-line fibre.
The hyperplane class that vanishes on $s_0$ restricts to $-m$ on
$s_\infty$. Therefore
\begin{equation}
 d_2(a)=(1-T_*)a+(mT_*a)f
       =-(aF)v_0+(av_0-3aF)F+(am-aF)f.
\label{pub:I:eq:E7}
\end{equation}
The span of $F,v_0,m$ is primitive. To check the only nontrivial
part of this assertion, split off the unimodular plane spanned by
$F,O$; $m-F$ and $v_0$ lie in its negative definite orthogonal
complement. Their Gram determinant is
$\det\left(\begin{smallmatrix}-2&-1\\-1&-6\end{smallmatrix}\right)=11$.
The square of an index of a larger integral lattice would divide
eleven, so the index is one. The span of $F,v_0$ is also primitive:
pairing with $O,m$ gives determinant $-1$. Unimodularity consequently
makes $a\mapsto(aF,av_0,am)$ surjective onto $\Z^3$. The coefficient
matrix in \eqref{pub:I:eq:E7} has determinant $-1$. Thus $d_2$ has primitive image
$\Z v_0\oplus\Z F\oplus\Z f$, rank three, kernel $\Z^7$,
and cokernel $\Z^8$.

\item The map $d_4$ sends $[S]$ to $s_0[S]-s_\infty[S]$, whose Poincar\'e
dual is $\pi_R^*m$. This is primitive and nonzero since $mO=1$.
Thus $d_4$ is injective with cokernel $\Z^{10}$. The map $d_0$ is
zero. Substitution in \eqref{pub:I:eq:E4} proves every asserted homology group.
The Euler characteristic is independently obtained by stratification:
$\chi(R)-\chi(S)=24-12=12$. 
\end{enumerate}
\end{proof}


\begin{remark}[The local modification and the compact central fibre]\hypertarget{tag-01R5}{}\label{fam:src:E.1.5}\label{fam:st:E.1.5}
\StatementRecord{fam:src:E.1.5}{proved}{01R5}{}
The blow-up charts of Proposition~\ref{fam:src:E.1.3} resolve the ordinary double point of $xy=uzv$ for $u\ne0$ and change one branch of the central normalization by the stated blow-up along $(x,z)$. The compact pinching $Q$ has the virtual Euler obstruction of Theorem~\ref{push:src:A.3.2}. The periodic modifications in Theorem~\ref{push:src:A.2.2} give the two smooth total spaces with central normalizations obtained by the explicitly described blow-up and small resolution.
\end{remark}

\subsection{Finite bielliptic limits and the smooth monodromy}\hypertarget{tag-01PU}{}\label{fam:section:E.2}

\begin{proposition}[Semistable limits of the bielliptic surfaces]\hypertarget{tag-01R6}{}\label{fam:src:E.2.1}\label{fam:st:E.2.1}
\StatementRecord{fam:src:E.2.1}{proved}{01R6}{the definitions and constructions stated in the proof}

After the auxiliary base change $w=s^6$, each of the two families of
\emph{reduced} special surfaces $S_j(w)$ has a proper flat
semistable filling by explicit finite free quotients. The central
surface for $j=1$ is a cycle of four ruled surfaces over an elliptic
curve; that for $j=2$ is a cycle of three. They have
\[
\begin{array}{c|ccccc|c}
 &H_0&H_1&H_2&H_3&H_4&\chi\\\hline
 j=1&\Z&\Z\oplus\Z/3&\Z^5&\Z^9&\Z^4&0\\
 j=2&\Z&\Z\oplus\Z/2&\Z^4&\Z^7&\Z^3&0
\end{array}
\]
with integral coefficients. Their fundamental groups are
$\Z^2\rtimes_{R_3}\Z$ and $\Z^2\rtimes_{R_4}\Z$, respectively,
where
\[
 R_3=\begin{pmatrix}0&-1\\1&-1\end{pmatrix},\qquad
 R_4=\begin{pmatrix}0&-1\\1&0\end{pmatrix}.
\]
\end{proposition}



\begin{proof}
\begin{enumerate}[label=\arabic*.,leftmargin=*]
\item The following product presentations specify the comparison
with the marked central surfaces. Let
$\rho=e^{\pi i/3}$,
\[
 E_j=\C/(\Z+\kappa_j\Z),\quad
 F_1=\C/(\Z(\rho-2)+\Z(-3)),\quad
 F_2=\C/(\Z(i-1)+\Z(-2)),
\]
where $\kappa_1=\beta_c(z_1)+2\mu(z_1)$ and
$\kappa_2=\beta_c(z_2)+3\mu(z_2)$ have negative imaginary parts.
The holomorphic maps to the central tori are induced by
$L_1(\xi,\eta)=(\eta,\xi-\eta/3)$ and
$L_2(\xi,\eta)=(\eta,\xi-\eta/2)$. Their kernels are generated by
\[
 k_1=(2/3,\rho),\qquad k_2=(1/2,i).
\]
The residual affine generators lift to
\begin{equation}
 h_1(\xi,\eta)=(\xi+\kappa_1/3,(\rho-1)\eta),\qquad
 h_2(\xi,\eta)=(\xi-\kappa_2/4,-i\eta).
\label{pub:I:eq:E8}
\end{equation}
These formulas are verified by substituting the four product periods
in $L_j$; the lattice indices are three and two. Moreover
$(\rho-1)\rho-\rho$ and $(-i)i-i$ are periods of the displayed
$F_j$, so the generators commute. They give free groups
$G_1=C_3\times C_3$ and $G_2=C_2\times C_4$. Indeed, equality of
the first coordinates of $k_j^a h_j^b(\xi,\eta)$ and $(\xi,\eta)$
forces $b=0$ from the coefficient of the nonreal period $\kappa_j$,
and then $a=0$ from the real period. Thus the composite quotient
maps are exactly the degree-nine and degree-eight covers of $S_j$,
with the stated signs.

\item Use the multiplicative coordinate $v=e^{-2\pi i\xi}$ on $E_j$.
Its Tate parameter is
\[
       q_j=e^{-2\pi i\kappa_j}=A_jw^2=A_js^{12},
       \qquad A_j=e^{-2\pi i\kappa_j^{\mathrm{part}}}\ne0.
\]
Here $A_j$ denotes the displayed scalar constant within this proof. Choose specifically
$A_j^{1/12}=\exp(-2\pi i\kappa_j^{\mathrm{part}}/12)$ and put
$r_j=A_j^{1/12}s$. This choice preserves the fractional period
translations in \eqref{pub:I:eq:E8} exactly. The regular Tate model for $q_j=r_j^{12}$
has a central twelve-gon. An infinite chain has node charts
$x_i y_i=r_j$; on a branch overlap one uses
$x_{i+1}=y_i^{-1}$ and $y_{i+1}=r_jy_i$. The coordinate
$r_j^i x_i$ on a general fibre is unchanged on overlaps.
The shift $i\mapsto i+12$ acts on that coordinate by $r_j^{12}$.
Quotienting a tubular neighbourhood of the chain by this shift gives
the proper twelve-gon family, and on $r_j\ne0$ it is
$\C^*/q_j^\Z$. Compactness follows from twelve projective components
and a compact annular fundamental domain on a general fibre;
the node charts give flatness and smoothness of its total space.

\item Multiplication by $q_1^{1/3}=r_1^4$ shifts the component index by four.
Multiplication by $q_2^{-1/4}=r_2^{-3}$ shifts it by minus three.
Hence the maps $h_j$ in \eqref{pub:I:eq:E8} extend, together with their rotations
on $F_j$. The maps $k_j$ multiply $v$ by roots of unity and translate
$F_j$, and extend to every projective component. Their relations
hold on the dense set $s\ne0$ and consequently on all charts.
The action is free on the central fibre as well. A nontrivial power
of $h_j$ moves every component and every node. A nontrivial power
of $k_j$ translates $F_j$ by a nonzero torsion point, including
above the nodes where the multiplicative action alone fixes a point.
The finite free quotient is therefore a smooth total space, proper
over $s$, with a reduced normal-crossing central surface. Flatness
descends along the finite \'etale quotient. The component orbits have
cardinality three and four, leaving four and three components.
Quotienting an individual $\mathbb P^1\times F_j$ by $k_j$ is a
ruled surface over $F_j/\langle\rho\rangle$ or
$F_j/\langle i\rangle$, respectively. Its two distinguished
sections have the same integral homology class: the corresponding
line bundle has finite order and degree zero on the elliptic base.

\item For the topology, unwrap the component cycle into an infinite chain
and then pass to the universal cover $\C$ of its elliptic base.
The resulting union of simply connected ruled pieces, with
simply connected intersections and a tree as incidence graph, is
simply connected by van Kampen. Its deck group is generated by
the enlarged lattice of translations on $F_j$ and one step round
the component cycle. The latter conjugates the lattice by the
order-three or order-four rotation. This gives the asserted
semidirect products; inverses of the displayed rotation matrices
give isomorphic groups by reversing the cycle generator.

\item For homology, the normalization has $n=4$ or $3$ components, each
with ranks $(1,2,2,2,1)$, and $n$ elliptic double curves, each with
ranks $(1,2,1)$. There are no triple points. The mapping-cylinder
sequence for the normalization and the double curves applies
integrally. In degree one the incidence matrix is a twisted cycle:
its cokernel is $\operatorname{coker}(R_j-I)$ and its kernel is
zero, because $R_j$ has no eigenvalue one. The degree-zero incidence
matrix contributes an additional $\Z$ to $H_1$. In degree two the
two section classes on a component agree, so the incidence matrix
is the ordinary cycle matrix on their classes. It has kernel
$\Z$ and primitive image of rank $n-1$. Thus $H_2=\Z^{n+1}$,
$H_3=\Z^{2n+1}$, and $H_4=\Z^n$. Finally
$\operatorname{coker}(R_3-I)=\Z/3$ and
$\operatorname{coker}(R_4-I)=\Z/2$. 
\end{enumerate}
\end{proof}








\section{The untwisted limit}\hypertarget{tag-02YZ}{}\label{push:section:A}






\subsection{Two minimal toroidal fillings and further modifications}\hypertarget{tag-02Z1}{}\label{push:section:A.2}

\begin{proposition}[The two minimal toroidal refinements]\hypertarget{tag-02ZJ}{}\label{push:src:A.2.1}\label{push:A21}
\StatementRecord{push:src:A.2.1}{proved}{02ZJ}{\SeriesRef{push:src:A.1.1}}

Among subdivisions with no new rays, the fan in \ref{push:src:A.1.1} has exactly
two smooth period-invariant subdivisions restricting to the chosen
$A_2$ triangulation at $h_u=0$. These are also all subdivisions
which are flat over the original $(q,u)$ bidisc, have reduced
$u$-central divisor, and retain that boundary triangulation.
Without the joint-flatness restriction there are infinitely many
smooth invariant subdivisions retaining that triangulation and
having reduced $u$-central divisor.
\end{proposition}


\begin{proof}
\begin{enumerate}[label=\arabic*.,leftmargin=*]
\item Write the six rays of $P_{0,0}$ as
$A=(0,0,1,0)$, $B=(1,0,1,0)$, $C=(0,1,1,0)$,
$D=(1,1,1,0)$, $E=(0,0,0,1)$, $F=(0,1,0,1)$.
The return's $A_2$ diagonal is $BC$. A triangular prism has six
triangulations without new vertices. With this bottom diagonal
the three possibilities are
\[
\begin{array}{c|c|c}
\text{maximal tetrahedra}&\text{left diagonal}&\text{right diagonal}\\\hline
ABCE,\ BCDE,\ CDEF&CE&DE\\
ABCE,\ BCDF,\ BCEF&CE&BF\\
ABCF,\ ABEF,\ BCDF&AF&BF.
\end{array}
\]
Periodicity identifies the left and right faces by $A\mapsto B$,
$C\mapsto D$, $E\mapsto E$, $F\mapsto F$; it excludes exactly
the middle row. Every listed maximal cone has determinant $\pm1$.
There are no missed triangulations: the prism has normalized volume
three, all its nondegenerate four-vertex tetrahedra have normalized
volume one, and enumeration of their triples with pairwise
face intersections gives exactly these six possibilities before
fixing the bottom. To test each pairwise intersection, write both
tetrahedra by their rational barycentric inequalities and enumerate
the vertices of the combined system. The intersection is their
common face exactly when each such vertex belongs to that face.

\item For the stronger assertion, a new ray with $h_q,h_u>0$ would
give a divisor mapping to the origin of the two-dimensional base.
Its fibre dimension would be three, contradicting joint flatness.
A ray with $h_u=0$ is excluded by the retained boundary
triangulation. On $h_q=0$, the fan support has $y_1=0$.
Reducedness requires the primitive ray to have $h_u=1$, so it
is already one of the $\rho_l$. Thus there are no new rays under
the stated three restrictions. Conversely the two triangulations
have all heights zero or one, and every cone maps onto a face of
the quadrant with saturated image. Their local monomial maps
are flat and their $u$ divisors reduced.

\item If joint flatness is omitted, star-subdivide the orbit of the
maximal cone $ABCE$ by its interior ray
$v_1=A+B+C+E=(1,1,3,1)$. Distinct translates have disjoint
interiors and distinct closed torus orbits. These are simultaneous
ordinary blow-ups at disjoint points, so the operation is
unambiguously period-invariant. Repeat in the cone $ABCv_{m-1}$,
where $v_0=E$. The new ray is
$v_m=v_{m-1}+A+B+C=(m,m,3m,1)$. Replacing any basis vector
by their sum preserves determinant one, and every new ray has
$h_u=1$. The boundary triangulations are unchanged. On the
compact quotient this blows up one smooth central point away
from the conductor at each step. There are therefore infinitely
many invariant smooth refinements, all reduced over $u=0$, with
increasing numbers of rays modulo the period group. They are
not jointly flat.
\end{enumerate}
\end{proof}

\begin{theorem}[The two minimal smooth fillings]\hypertarget{tag-02ZK}{}\label{push:src:A.2.2}\label{push:A22}
\StatementRecord{push:src:A.2.2}{proved}{02ZK}{\SeriesRef{push:src:A.2.1}, \SeriesRef{push:A12}, \SeriesRef{push:src:A.1.2}}

The two minimal subdivisions of \ref{push:src:A.2.1} give proper flat holomorphic
families $\mathcal X^+,\mathcal X^-\to\Delta_u$ with smooth total
spaces. Their joint maps to $B\times\Delta_u$ are flat. Over
$\Delta_u^*$ they are isomorphic as families to the specified
untwisted family $X(u)$. Each central fibre $Q^\pm$ is reduced,
irreducible, and has normal crossings, with smooth double locus
$S$ and no triple points.

Its normalization $R^\pm$ is obtained by blowing up $R$ along
$S_\infty$ in, respectively, one of its two sections and then
taking the small resolution at the unique ordinary double point
of that blow-up prescribed by the fixed bottom diagonal.
The two proper transforms of the sections are still copies of
$S$, glued by $T$. Their normal-bundle product is trivial;
in particular $\mathcal T^1_{Q^\pm}\simeq\mathcal O_S$.
\end{theorem}


\begin{proof}
\begin{enumerate}[label=\arabic*.,leftmargin=*]
\item Apply either fan subdivision over the cusp neighbourhood in
Theorem~\ref{push:A12}, and use the identity modification elsewhere.
Toric modifications are proper, and a subdivision is an
isomorphism off the strata it subdivides. Its restriction over
$q\ne0$ is the identity, so this gluing has no additional
transition choice. The period actions preserve the subdivisions,
and the covering description in \ref{push:src:A.1.2} descends the proper
modifications to the quotient. All maximal cones are regular.
Every local ring of the resulting fourfold is therefore regular.
The coordinate $u$ is the product of those affine coordinates
whose rays have $h_u=1$; there are at most two such rays in a
maximal cone. This proves reduced normal crossings and the absence
of triple points for the $u$-central fibre. The regular-sequence
argument, or the saturated face images in \ref{push:src:A.2.1}, proves joint
flatness. Properness follows from properness of the toric
modification and of $\mathcal X^\flat$.

\item  At a node of $S_\infty$, write the original normalization
as coordinates $(X,z,v)$, with section $X=0$ and centre
$(X,zv)$. Blowing up that ideal has charts
\[
 X=zvA;\qquad zv=XB.
\]
The first is smooth, and the second has one conifold at its
origin. A small resolution of the latter is, for one of the
bottom-diagonal conventions,
\[
 X=za,\quad v=aB;\qquad z=Xd,\quad B=dv.
\]
Interchanging $z,v$ gives the other diagonal convention. These
are holomorphic charts of a proper small modification: they are
the two usual charts of the blow-up of the appropriate rank-one
ideal on the conifold. The strict transform of $X=0$ is $A=0$
in the first chart and never passes through the conifold. It is
therefore unchanged as a surface. The other section is disjoint
from the centre and is unchanged as well. Over a smooth point
of $S_\infty$ the exceptional fibre is one $\mathbb P^1$;
over its node it is a chain of two $\mathbb P^1$'s, one being
the exceptional curve of the small resolution. This identifies
the normalization morphism, including its order of operations.

\item For example the corresponding total fourfold modification of
$xy=uzv$ can be written by first taking $x=ua$, $ay=zv$ and
then resolving the conifold in $(a,y,z,v)$, together with the
other chart $u=xb$, $y=bzv$. Its three smooth affine charts map
to $(x,y,z,v,u;q)$ by
\[
\begin{aligned}
 (u,z,\lambda,y)&\longmapsto(uz\lambda,y,z,\lambda y,u;z\lambda y),\\
 (u,a,\nu,v)&\longmapsto(ua,\nu v,a\nu,v,u;a\nu v),\\
 (x,b,z,v)&\longmapsto(x,bzv,z,v,xb;zv).
\end{aligned}
\]
The first two are the conifold-resolution charts
$a=z\lambda$, $v=\lambda y$ and $z=a\nu$, $y=\nu v$.
Substitution gives $xy=uzv$ and $q=zv$ in every chart.
Exchanging the chosen zero and infinity branches gives the other model.

\item The normal bundle of the section containing the centre changes
by tensoring with $\mathcal O_S(-S_\infty)$; the second normal
bundle is unchanged. Since the original product is $J$, the
new product is $J(-S_\infty)=\mathcal O_S$. The local
formula for $\mathcal T^1$ is the product of the two normal
bundles, with $T$ used on the infinity branch. This proves
d-semistability in this precise sense. The $\rho_l$ form one
orbit, so the central fibre is irreducible and its conductor is
exactly the two stated sections. Over $u\ne0$ both subdivisions
restrict to the fixed bottom triangulation, giving precisely
Engel's Mumford filling, simultaneously in $u$.

\item The birational comparison $\mathcal X^+\dashrightarrow
\mathcal X^-$ is the identity over the punctured $u$ disc and
off the exceptional loci over $(q,u)=(0,0)$. Locally it replaces
the two compatible side-face diagonals by the other two.
All its rays lie in the Gorenstein hyperplane $h_q+h_u=1$,
so both small modifications are crepant over
$\mathcal X^\flat$. Over a smooth point of the nodal curve
this is the ordinary Atiyah flop times a curve; at the node its
extension is exactly the displayed pair of prism subdivisions.
This specifies the birational map.
\end{enumerate}
\end{proof}

\subsection{An explicit two-bundle plumbing}
\begin{theorem}[A proper filling of the untwisted family]\hypertarget{tag-02RA}{}\label{night:P2.10}\label{stmt:P2.10}
\StatementRecord{night:P2.10}{proved}{02RA}{the definitions and constructions stated in the proof}

The untwisted family $X(u)$ has a proper flat Hausdorff analytic filling at $u=0$. Its corner winding is zero. The filling is obtained from two projective-line bundles by explicit plumbing, a proper modification and a finite quotient.
\end{theorem}
\begin{proof}
\begin{enumerate}[label=\arabic*.,leftmargin=*]
\item Let $S$ be the rational elliptic surface of Definition~\ref{I:def:pinching-data}, $M=\OO_S(P-O)$, $h=t_{6P}$, and $L=\Hom(h^*M,M)=\OO_S(F_\infty)$, with section $\phi$ vanishing on $F_\infty$. The absence of logarithmic twists gives $b=a=0$ in the corner character calculation. Thus the origins agree and their winding is zero. 

\item Take two copies $Y_0,Y_1$ of $\Pp_S(\OO\oplus M)$. The finite coordinate $z_i$ lies in $M$, its inverse $w_i$ in $M^{-1}$. The first neck, over $p\in S$, joins zero in $Y_0$ to infinity in $Y_1$ over $h(p)$ by $z_0w_1=r\phi_p$. The second joins zero in $Y_1$ to infinity in $Y_0$ over the same point by $z_1w_0=r$. These equations are typed in $M\otimes h^*M^{-1}=L$ and $M\otimes M^{-1}=\OO$.

\item Choose a Hermitian metric and $0<\epsilon<1/4$. Use middle annuli $\epsilon/2<|z_i|<2/\epsilon$ in $Y_i$, times the parameter disc, and neck hypersurfaces with $|z|,|w|<\epsilon$. Glue their ends on $\epsilon/2<|z|<\epsilon$ or $\epsilon/2<|w|<\epsilon$ by coordinate projection, using $h$ at the first infinity end. In holomorphic frames these are holomorphic maps, and they commute with bundle transitions because the product lines are exactly the ones displayed above. Choose $\rho$ so that $\rho\max(1,\|\phi\|)<\epsilon^2/16$. The product equation makes the two ends of a neck disjoint; the two ends in each middle chart are disjoint as well.

\item For $|r|\le\rho'<\rho$, a compact fundamental cover consists of closed middle annuli $3\epsilon/4\le|z_i|\le4/(3\epsilon)$ and closed neck pieces $|z|,|w|\le3\epsilon/4$. Every omitted neck end is represented in a middle annulus. The gluing relation on this finite union of compact bundles is the union of diagonals and closed graphs of the coordinate identifications. No additional graph accumulates, because the two neck ends are disjoint. The relation is closed, so its quotient is compact and Hausdorff. The open charts give the analytic structure. This proves properness over the smaller disc and the global Hausdorff property.

\item Over $S\setminus F_\infty$, traversing the two necks gives
\[
 (p,z_0)\longmapsto\left(h(p),\frac{z_0}{r^2\phi_p}\right).
\]
The inverse covers $h^{-1}$ and is multiplied by $u\phi$, with $u=r^2$. It is the precise generator of Engel's quotient on that open set.

\item At a smooth point of $F_\infty$ the first neck is $zw=rq$; at its node it is $zw=rxy$, with $q=xy$. The ideal of one globally distinguished central component is locally $(z,r)$. Because the two global components are distinct, this is a global coherent ideal. Blow it up. At the node its two charts are
\begin{align*}
 r=za,\quad w=axy;&\qquad (z,a,x,y)\text{ are smooth coordinates},\\
 z=rb,\quad bw=xy;&\qquad (r,b,w,x,y)\text{ with the displayed relation}.
\end{align*}
The remaining singularity is an ordinary threefold node times the parameter disc. Its singular section does not meet the first chart. Resolve $(b,y)$ in the second chart. Its two charts are
\begin{align*}
 b=yk,\quad x=kw;&\qquad (r,y,k,w),\\
 y=bj,\quad w=jx;&\qquad (r,b,j,x).
\end{align*}
They are smooth, with transition $j=k^{-1}$ and $(b,x)=k(y,w)$. Off the singular section the ideal is invertible, so this proper small resolution glues to the identity. The analogous blowup charts at smooth points of $F_\infty$ are already smooth. The resulting total space is smooth, and its parameter central divisor is reduced NC: in the first chart $r=za$, and in both small-resolution charts $r$ is a coordinate. Parameter torsion freeness proves flatness.

\item  For fixed $r\ne0$ the first blowup is an isomorphism. Put $W=rw$ and $u=r^2$; the first neck becomes $zW=u xy$. On its dense torus use $(x,z,q)$, with $y=q/x$ and $W=uq/z$. The four affine monomials have exponents
\[
 (1,0,0),\quad(-1,0,1),\quad(0,1,0),\quad(0,-1,1).
\]
The dual cone is $0\le\alpha,\beta\le\gamma$, the cone over the unit square at height one. The resolution of $(z,y)$, equal to that of $(b,y)$ when $r$ is a unit, divides it along $\alpha+\beta=\gamma$: the chart character $k=z/(ry)$ has valuation $\alpha+\beta-\gamma$. Its two cones, over the triangles with common edge $[e_1,e_2]$, are unimodular. The exact comparison with the coordinates $(X,Y,Z,V)$ in the node $XY=uZV$ of Theorem~\ref{push:src:A.1.2} is
\[
 (X,Y,Z,V)=(z,rw,x,y).
\]
In its first resolution chart the parameters are $(V,Y/u,A)=(y,w/r,rk)$, and in its second they are $(X,Z,B)=(rb,x,j/r)$. These give respectively
\[
 (X,Y,Z,V)=(rk y,rw,kw,y),\qquad
 (X,Y,Z,V)=(rb,rjx,x,bj),
\]
and $AB=kj=1$ on the overlap. Thus both chart maps and their transition agree with the selected resolution of Theorem~\ref{push:src:A.1.2}. In the two charts, $q=kyw$ and $q=bxj$, respectively. The Tate and linearization periods shift the two marked lattice directions by one; their unit factors act holomorphically on these charts. This is the chosen marked triangular filling supplied by Engel's general construction in \cite[Proposition~2.8, pp.~6--7]{Engel2026}. Together with the quotient comparison, the displayed map identifies the entire compact generic fibre with $X(u)$.

\item The involution $r\mapsto-r$ fixes the $Y_0$ fibre coordinate and negates the $Y_1$ coordinate. It preserves both equations and the global component ideal. In the nodal blowup chart it acts as $(b,w,x,y,r)\mapsto(-b,-w,x,y,-r)$ and preserves $(b,y)$. Its lifts on the two resolution charts are
\[
 (r,y,k,w)\longmapsto(-r,y,-k,-w),\qquad
 (r,b,j,x)\longmapsto(-r,-b,-j,x).
\]
The transition $j=k^{-1}$ intertwines these maps. Its finite quotient is proper and Hausdorff over $u=r^2$. Its invariant local rings are torsion free over the discrete valuation ring $\C\{u\}$ and hence flat. Away from zero the involution exchanges the two roots of $u$, while $W=rw$ and the $Y_0$ coordinate are invariant. The displayed marked comparison therefore identifies the quotient fibre with one copy of $X(u)$. This is a filling in the original parameter.
\end{enumerate}
\end{proof}


\begin{proposition}[A projective-space component in a point-star refinement]\hypertarget{tag-02ZL}{}\label{push:src:A.2.3}\label{push:A23}
\StatementRecord{push:src:A.2.3}{proved}{02ZL}{\SeriesRef{push:src:A.2.1}, \SeriesRef{push:src:A.2.4}, \SeriesRef{push:src:A.2.5}, \SeriesRef{fam:src:E.1.1}}
The first maximal-cone star of Proposition~\ref{push:src:A.2.1} contributes a central normalization component $\mathbb P^3$. There is no proper bimeromorphic holomorphic map from this component to $R$.
\end{proposition}
\begin{proof}
The star is the ordinary blow-up of a point in a smooth four-dimensional chart, with exceptional divisor $\mathbb P^3$. The projective-bundle basis gives $b_2(R)=b_2(S)+1=11$, while $b_2(\mathbb P^3)=1$. A proper bimeromorphic map $f:\mathbb P^3\to R$ has topological degree one, and the projection formula gives
\[
 \int_{\mathbb P^3}f^*a\,f^*b=\int_Rab
 \qquad(a\in H^2(R;\mathbb Q),\ b\in H^4(R;\mathbb Q)).
\]
Poincar\'e duality would make $f^*$ injective in degree two, contradicting these ranks.
\end{proof}

\begin{theorem}[Exhaustive finite fan data]\hypertarget{tag-02ZM}{}\label{push:src:A.2.4}\label{push:A24}
\StatementRecord{push:src:A.2.4}{proved}{02ZM}{\SeriesRef{push:src:A.2.1}, \SeriesRef{push:src:B.2.1}, \SeriesRef{push:src:A.1.2}; \cite[Section 2, p.6]{DMM16}}

The following finite-data procedure parametrizes all smooth
period-invariant subdivisions of the cusp fan after
$u=r^n$, retaining the prescribed bottom triangulation and
having finitely many cones modulo the periods. The minimal subdivisions
for $n=1,2,12$ are described in Proposition~\ref{push:src:A.2.1}
and Theorem~\ref{push:src:B.2.1}.
Each output defines a proper flat family over $\Delta_r$
with smooth total space, whose punctured restriction is the
given untwisted family pulled back by $u=r^n$.

For a bound $H\geq1$, retain the four bottom rays
$A,B,C,D$ and consider all primitive vectors
\begin{equation}
 \begin{gathered}
 v=(a,b,k,\ell),\qquad
 \ell\geq1,\quad k\geq0,\quad k+\ell\leq H,\\
 0\leq a\leq k,\qquad0\leq b\leq k+n\ell.
 \end{gathered}
\label{pub:I:eq:A-24a}
\end{equation}
Take any collection of four-element subsets of this finite
set satisfying the following exact tests:
\begin{enumerate}
\item Each four-column determinant has absolute value one.
\item After dividing each ray by $k+\ell$, the resulting
tetrahedra intersect pairwise in common faces, and their
volumes sum to $(n+2)/6$.
\item Their bottom boundary is the two triangles $ABC,BCD$.
Their boundary subdivisions on $a=0$ and $a=k$ agree under
$G_1$, and those on $b=0$ and $b=k+n\ell$ agree under $G_2$.
\end{enumerate}
All these are finite rational-arithmetic tests. Extend a
passing collection by the two periods, and discard duplicates
as $H$ increases. This is exhaustive. For a reduced central
divisor, retain only $\ell=1$ in \eqref{pub:I:eq:A-24a}.

The central prime components are indexed by period orbits of
the selected rays with $\ell>0$, and the multiplicity of the
component indexed by $v$ is exactly $\ell$. The normalization
charts of that component are the star fan of $v$ in
$N/\mathbb Zv$, with the inherited bundle units and the
stabilizer action. For $k>0$ the stabilizer is trivial and this
normalization is a smooth complete toric threefold. All
multiple loci, including possible triple and quadruple loci,
are specified by the positive-$\ell$ rays belonging to each
cone. In the reduced case
$\mathcal T^1\simeq\mathcal O_{\operatorname{Sing}D}$,
so every output is d-semistable.
\end{theorem}


\begin{proof}
\begin{enumerate}[label=\arabic*.,leftmargin=*]
\item Write the period shears as
\begin{align*}
 G_1(a,b,k,\ell)&=(a+k,b,k,\ell),\\
 G_2(a,b,k,\ell)&=(a,b+k+n\ell,k,\ell).
\end{align*}
One closed fundamental cone is
\[
 P_n=\{k,\ell\geq0,\ 0\leq a\leq k,
                    \ 0\leq b\leq k+n\ell\}.
\]
Its bottom is the original square; its top is the edge from
$(0,0,0,1)$ to $(0,n,0,1)$. The slice $k+\ell=1$, in
coordinates $(a,b,k)$, is the compact polytope
\[
 0\leq k\leq1,\qquad0\leq a\leq k,
 \qquad0\leq b\leq n-(n-1)k.
\]
Its volume is
$\int_0^1 k(n-(n-1)k)\,dk=(n+2)/6$.
A tetrahedron with ray columns $v_i$ has slice volume
\begin{equation}
 \frac{|\det(v_1,v_2,v_3,v_4)|}
      {6\prod_{i=1}^4(k_i+\ell_i)}.
\label{pub:I:eq:A-24b}
\end{equation}
This follows by applying the determinant-one coordinate
change from $(a,b,k,\ell)$ to $(a,b,k,k+\ell)$ and then
normalizing the last coordinate to one.

\item A common-face test is also completely explicit: write each
tetrahedron in barycentric halfspaces, intersect the two
finite lists of halfspaces, and enumerate the possible
vertices by solving triples of bounding equations. The
intersection is a common face precisely when all its vertices
are common tetrahedron vertices. Since the candidate
tetrahedra lie in $P_n$, pairwise face intersections and the
volume equality imply that they cover the slice. Otherwise
their finite closed union would omit a nonempty relatively
open set, which has positive volume. The boundary tests make
the translates agree on every shared face. This proves that
every passing collection gives the claimed invariant fan.

\item Conversely, intersect any invariant smooth refinement with
$P_n$. A smooth maximal cone has four primitive generators
with determinant $\pm1$, and its slice is one of the stated
tetrahedra. Finitely many cone orbits give finitely many rays
in this prism, hence some maximum height $H$. No additional
bottom rays are permitted by the boundary condition; every
other ray has the form \eqref{pub:I:eq:A-24a}. The resulting collection
passes all three tests. Thus the search at each height is finite,
and the procedure is exhaustive.

\item Properness of the quotient forces the finite-orbit condition. Each full-dimensional cone has a
closed torus orbit consisting of one point with $q=r=0$.
These points form a closed discrete set: an affine chart
of a full-dimensional cone contains exactly its own closed
torus orbit. After the properly discontinuous period
quotient the points are indexed by cone orbits and remain
closed and discrete. Infinitely many such orbits would
therefore give an infinite closed discrete subset of a
compact central fibre, which is impossible.

\item The fundamental orbit representatives of mixed rays are
\begin{equation}
 (a,b,k,\ell),\quad k,\ell>0,
 \quad0\leq a<k,\quad0\leq b<k+n\ell,
 \quad\gcd(a,b,k,\ell)=1.
\label{pub:I:eq:A-24c}
\end{equation}
For pure top rays they are $(0,b,0,\ell)$ with
$0\leq b<n\ell$ and $\gcd(b,\ell)=1$. A smooth top
subdivision necessarily includes each $(0,j,0,1)$,
$0\leq j\leq n$. Indeed an omitted such vector would lie
inside a regular two-cone; writing it in that cone's integral
basis would express its height one as the sum of two
positive integral heights, which is impossible.
In the reduced search there are exactly $n$ pure top orbits
and $\sum_{k=1}^{H-1}k(k+n)$ available mixed orbits.

\item Enumerating the primitive rays in \eqref{pub:I:eq:A-24a} and retaining the
four-element subsets whose column determinants have absolute value one
gives the following finite counts:
\[
\begin{array}{c|c|c|r|r}
n&H&\text{reduced}&\text{rays in }P_n&
       \text{unimodular four-subsets}\\\hline
1&1&\text{yes}&6&12\\
1&2&\text{yes}&12&267\\
1&2&\text{no}&13&357\\
2&2&\text{yes}&15&544\\
2&2&\text{no}&17&784\\
12&1&\text{yes}&17&100
\end{array}
\]
The table counts rays and unimodular four-subsets. The common-face,
volume and boundary tests above select the complete collections.
At height one these are the two and $n+1$ fans of
\cref{push:src:A.2.1,push:src:B.2.1}, respectively.

\item For existence of the corresponding compact family, start
with the proper nodal plumbing of \ref{push:src:A.1.2} and base-change it
by $u=r^n$. On the cusp cover apply the proper toric
modification defined by the refinement. Its support is
unchanged, which is the toric properness criterion
\cite[Section 2, p.6]{DMM16}, and each
original affine chart is replaced by a finite fan. On
$q\ne0$ use the subdivision of the top two-dimensional fan
to modify the chart $ab=r^n\varphi$ fibrewise over $S$.
Changes of trivialization multiply $a$ and $b$ by the given
nonvanishing line-bundle units; a torus action extends to
every toric fan chart, so these modifications glue. The same
top fan occurs on the overlap with the cusp construction.
The core outside the plumbing charts is unchanged. This defines
a proper modification of the base-changed compact nodal family.

\item The period maps lift because the fan is invariant. Their
lifted action is free: a fixed point would project to a fixed
point in the nodal covering. It is properly discontinuous:
compact sets project to compact sets, and any translating
element producing an intersection upstairs produces one
downstairs, where \ref{push:src:A.1.2} gives finiteness. Thus the quotient
exists as a Hausdorff complex space, and the equivariant
proper modification descends to a proper modification of
the compact nodal quotient. In a unimodular chart its
coordinates $z_i$ satisfy
\begin{equation}
 r=\text{unit}\cdot\prod_{i=1}^4 z_i^{\ell_i},
 \qquad q=\text{unit}\cdot\prod_{i=1}^4 z_i^{k_i}.
\label{pub:I:eq:A-24d}
\end{equation}
The total space is smooth. The nonconstant coordinate $r$
is a nonzerodivisor, so the morphism to the disc is flat.
Its central multiplicities are the indicated $\ell_i$;
it is reduced exactly when all selected positive heights
are one. The bottom fan gives the specified smooth
Mumford model over $r\ne0$, including its original period
identification. Thus the punctured-family assertion holds
simultaneously in the parameter.

\item The normalization charts are as follows. For a maximal cone with basis
$(v,v_2,v_3,v_4)$, let $(m_1,m_2,m_3,m_4)$ be the dual
integral basis. The chart of the normalized component is
$z_1=0$ with coordinates $z_i=\chi^{m_i}$, $i=2,3,4$.
On overlaps the transition matrix is the product of the two
integer basis matrices, and the original bundle units are
raised to those same integer exponents. This is precisely
the star fan in $N/\mathbb Zv$. When $k>0$, \eqref{pub:I:eq:A-24c} shows
that the period stabilizer of the ray is trivial. The ray
lies in the interior of the full support, so its star is
complete. Its normalized component is therefore the
asserted complete smooth toric threefold, lying over the
cusp. When $k=0$, the stabilizer is generated by $G_1$.
Away from $S_\infty$ the resulting component is the ruled
bundle $\mathbb P_S(\mathcal O\oplus L^\ell)$: a character
on its open ruled fibre is $W^\ell/r^b$, which transforms
by the $\ell$th power of the transition of $L$. The stated
star charts describe its replacement over $S_\infty$.
The gluing maps are the original period maps
on those same integer characters.

\item Intersections are read from common cones, with separate
local branches retained before the period quotient. In a
chart \eqref{pub:I:eq:A-24d}, the number of branches of the reduced
central divisor is the number of positive $\ell_i$.
When all these heights are one, the hypersurface
calculation gives
$\mathcal T^1_D=\mathcal O_D(D)|_{\operatorname{Sing}D}$.
The divisor is globally principal, cut out by $r$, so this
line on the singular locus is trivial. Finally, any two
output fans have an explicit common refinement by all
intersections of their cones. Its normal toric modification
gives a proper common model, and the resulting birational
map is the identity on the fixed lattice and on the
punctured family.
\end{enumerate}
\end{proof}

\begin{proposition}[The normalization and modification algorithm]\hypertarget{tag-02ZN}{}\label{push:src:A.2.5}\label{push:A25}
\StatementRecord{push:src:A.2.5}{proved}{02ZN}{\SeriesRef{push:src:A.2.4}, \SeriesRef{push:src:A.4.6}, \SeriesRef{push:src:A.4.8}, \SeriesRef{push:src:A.2.2}}

For every finite fan datum $\Sigma$ of Theorem~\ref{push:src:A.2.4}, the following algorithm describes the normalized components and their comparison with a reference model.

The star charts and integer transition matrices of \ref{push:src:A.2.4}
give every normalized component, its gluing maps, and
all multiple loci directly. There is also a terminating
finite-data algorithm that expresses its comparison
with a reference model as an explicit sequence of
blow-ups and blow-downs along specified smooth strata.
At each blow-up the algorithm gives the normalized
strict transforms, the exceptional projective bundle,
its multiplicity, and every new intersection.
The algorithm uses both blow-ups and blow-downs.

For a component of multiplicity $\ell$, the local
nonreduced central equation is the monomial in \eqref{pub:I:eq:A-24d}.
More generally, in a unimodular chart with central
equation $F=\prod_i z_i^{\ell_i}$, the cotangent-complex
first-order deformation sheaf is
\begin{equation}
 \mathcal T^1_D
  =\mathcal O_D\Big/
    \left(\ell_i z_i^{\ell_i-1}
          \prod_{j\ne i}z_j^{\ell_j}:\ell_i>0\right).
\label{pub:I:eq:A-25a}
\end{equation}
Here $\mathcal T^1_D$ means
$\mathcal Ext^1(\mathbb L_D,\mathcal O_D)$; in the
reduced case it also equals
$\mathcal Ext^1(\Omega_D,\mathcal O_D)$.
The reduced case is the trivial line on the singular
scheme, as in \ref{push:src:A.2.4}. For a nonreduced centre the displayed
Jacobian quotient is the appropriate deformation sheaf;
the stated d-semistability conclusion applies to the reduced case.
\end{proposition}


\begin{proof}
\begin{enumerate}[label=\arabic*.,leftmargin=*]
\item Proposition~\ref{push:src:A.2.3} proves that the point-star component $\mathbb P^3$ has no proper bimeromorphic map to $R$. In that example the full normalization is $\operatorname{Bl}_pR^+\sqcup\mathbb P^3$: the first component is the strict transform, and the two meet along the exceptional plane. These are the ordinary point-blow-up charts of the smooth total fourfold.

\item Retain all marking
data from \ref{push:src:A.2.4}: the lattice, the two period matrices,
the original line-bundle units, and the common
unsubdivided fan. Strictify the finite quotient cone
complex by the decreasing-dimensional barycentric
stars used in \ref{push:src:A.4.8}. Record these operations and do
the same for the reference model. Form their common
intersection complex. A regular common subdivision
can be found either by toric desingularization or by
enumerating rational subdivisions with increasing
primitive-ray bounds and testing their determinants
and common faces. A regular subdivision exists, so
the latter enumeration terminates.

\item On a finite strict regular complex there are finitely
many possible ordinary star blow-ups: choose a cone
of dimension at least two and insert the sum of its
primitive basis rays. There are also finitely many
possible inverse moves: choose a ray and test whether
its star is exactly the star created by such a
subdivision of a regular cone. Both tests are exact
integer calculations. Require that every intermediate
complex remain a subdivision of the marked original
complex. Record face identifications; equality of two
marked complexes is a finite rational comparison,
using the period orbit representatives in \eqref{pub:I:eq:A-24c} and
their pure top and pure bottom analogues.

\item Enumerate this locally finite graph of moves in
increasing path length, discarding repeated marked
complexes. Applied separately to the two proper
subdivision morphisms from the common regular
complex, the finite-complex weak-factorization theorem
used in \ref{push:src:A.4.8} guarantees a path, all of whose
intermediates are over the corresponding target.
Thus the search terminates. Concatenate the two paths
and the recorded strictification paths. This gives a
fully specified finite sequence once the input fan is
given. Temporary pure bottom subdivisions may occur;
they are allowed in a presentation of the birational
comparison, and are removed before the prescribed
endpoint is reached.

\item  If a cone has primitive basis vectors
$v_1,\ldots,v_c$, its centre is the closed stratum
\[
 C=\{z_1=\cdots=z_c=0\}
\]
in the marked toroidal chart. Strictness makes its
global closure smooth; the period translates defining
one centre are disjoint. Blow-up chart $i$ has
coordinates
\[
 z_i=t_i,\qquad z_j=t_i w_j\quad
 (1\leq j\leq c,\ j\ne i),
 \qquad z_j\quad(j>c).
\]
The changes between charts are these integral
monomials, with the original units raised to the
same powers. The exceptional divisor is
$\mathbb P_C(N_{C/Y})$, where $Y$ is the smooth
intermediate total space. If $C$ is central, its
central multiplicity is
\[
 \ell_E=\sum_{j=1}^c\ell_j.
\]
If all $\ell_j$ vanish, it is a horizontal exceptional
divisor and is not a new central component.

\item For an old normalized component, take the inverse
image of this centre on its marked branch and the
ordinary strict transform in these charts. It is
the blow-up along that smooth inverse-image centre;
when that centre is Cartier the blow-up is the
identity. Its intersections with the new exceptional
component are the appropriate coordinate-hyperplane
projective subbundles. Intersections of several old
branches are given by imposing the corresponding
coordinate equations simultaneously. Thus the same
chart prescription records double, triple and
quadruple loci, their normal lines, and their 
gluing maps. An inverse star is the inverse
contraction of this recorded configuration. This
specifies the centres and their order.

\item The algorithm starts with the explicit reference
normalizations of \ref{push:src:A.2.2} and the $n$-cycle reference
in \ref{push:src:A.4.8}. It is therefore a uniform geometric
presentation of every output by specified blow-ups,
blow-downs and projective bundles. Independently,
the direct star charts avoid running this search
when only a normalization chart or an intersection
is required. The mixed-ray components are complete
smooth toric threefolds; a pure top component
$(0,b,0,\ell)$ is bimeromorphic to
$\mathbb P_S(\mathcal O\oplus L^\ell)$.

\item Finally $D$ is globally the principal Cartier
divisor of the parameter $r$ in a smooth total
space. Its two-term cotangent complex is
$[\mathcal O_D\xrightarrow{dF}\Omega_Y|_D]$,
so its first-order deformation module is the
cokernel of the gradient. Equivalently, a lifted
hypersurface equation $F+\epsilon G$ is changed by
an infinitesimal coordinate change precisely by a
Jacobian term. This argument applies to nonreduced
monomials as well; an ordinary conormal sequence
need not be left-exact there. Absorbing the local unit, or reducing its
derivative modulo $F$, gives \eqref{pub:I:eq:A-25a}. A global
parameter makes the normal line trivial, so these
Jacobian quotients glue by the canonical
hypersurface presentation. In the reduced case
the Jacobian ideal defines the singular scheme,
giving the asserted trivial deformation line there.

\end{enumerate}
\end{proof}


\subsection{The virtual Euler obstruction and the conjecture}\hypertarget{tag-02Z2}{}\label{push:section:A.3}

\begin{lemma}[Analytic specialization of the virtual Euler number]\hypertarget{tag-02ZO}{}\label{push:src:A.3.1}\label{push:A31}
\StatementRecord{push:src:A.3.1}{proved}{02ZO}{\cite[Section~3.1, Lemmas~3.1--3.2, pp.~588--590]{Levy08}; \cite[Section~3, Theorems~1--2, pp.~463--464]{Loj64}; \cite[Theorem~1.6, p.~215; Remark~1.16, p.~221]{BM97}; \cite[Theorem~5.8 and Remark~5.9, p.~102]{FG10}}

Let $Z$ be a compact reduced local-complete-intersection complex
space of pure dimension $n$. In any proper flat holomorphic
family over a disc with central fibre $Z$ and smooth nearby
fibres, their Euler number is
\[
 e_{\mathrm{vir}}(Z)=\langle c_n(T_Z^{\mathrm{vir}}),[Z]\rangle.
\]
\end{lemma}


\begin{proof}
\begin{enumerate}[label=\arabic*.,leftmargin=*]
\item Flat deformation preserves the local-complete-intersection
property relative to the disc near $Z$. Locally this follows by
lifting an embedding and its regular sequence of equations;
Nakayama's lemma and flatness show that their lifts generate
the ideal and remain a regular sequence. Properness allows one
to shrink the disc so that these neighbourhoods cover the total
space. The relative cotangent complex is consequently perfect
of amplitude $[-1,0]$. Its dual defines a topological complex
$K$-class $V$ on that total space; local finite vector-bundle
complexes define the class by descent, or by stabilization and
partitions of unity; the analytic-to-topological comparison is
given in \cite[Section~3.1, Lemmas~3.1--3.2, pp.~588--590]{Levy08}.
Its restriction to a smooth fibre is its
tangent bundle, and its restriction to $Z$ is its intrinsic
virtual tangent class.

\item We compare the fundamental cycles on a smooth proper modification,
so that the ambient-manifold hypothesis of the triangulation theorem
is satisfied. Analytic resolution, in the complex-analytic scope of
\cite[Analytic category (0.1)(2), p.~208; Theorem~1.6 and its noncompact extension, p.~215; Remark~1.16, p.~221]{BM97}, gives
\[
 \rho:\widetilde{\mathcal X}\longrightarrow\mathcal X
\]
with smooth total space, isomorphic over the smooth locus. Only a
neighbourhood of the compact central fibre is needed; the locally
finite construction is finite there after shrinking the base.
The total space is reduced: a nilpotent local section lies in every
power of the base parameter by reducedness of the central fibre and
flatness, and hence vanishes by Krull intersection. Thus the stated
resolution applies to the dominating space.

\item Write the divisor of the pulled-back parameter as
$\sum_j m_jD_j$. The inverse image of a short radial segment
$[0,\epsilon]$ is a compact semianalytic subset of the smooth real
analytic manifold $\widetilde{\mathcal X}$. We triangulate it,
compatibly with its two endpoint fibres and the singular loci of the
$D_j$, by \cite[Section~3, Theorems~1--2, pp.~463--464]{Loj64}.
Over the interior of the segment its top-dimensional part is an
oriented real $(2n+1)$-manifold. Near a general point of $D_j$ the
parameter is $u z^{m_j}$, with $u$ a holomorphic unit; choosing a
local root of $u$ exhibits $m_j$ radial half-sheets, each with the
complex boundary orientation. The sum of their oriented top
simplices therefore has boundary
\[
 [\mathcal X_\epsilon]-\sum_jm_j[D_j].
\]
The singular and crossing sets have real dimension at most $2n-2$
and contribute no boundary simplex of dimension $2n$.

\item At every smooth point of the reduced central fibre the original
map is a submersion, by flatness and smoothness of that fibre germ.
The modification is consequently an isomorphism there and the
strict-transform multiplicity is one. Every other $D_j$ maps into
$\operatorname{Sing}Z$, of complex dimension at most $n-1$.
Its fundamental class has zero pushforward in dimension $2n$:
its map factors through that lower-dimensional analytic set.
Thus pushing the displayed chain forward gives exactly
$[\mathcal X_\epsilon]-[Z]$. This proves the required homology
identity in the original total space. Pairing the globally defined
$c_n(V)$ with these homologous cycles gives the equality of
Chern numbers. Gauss--Bonnet on the smooth compact complex
fibre gives its Euler number. This proves the analytic assertion
directly. The proper-scheme counterpart is
\cite[Theorem~5.8 and Remark~5.9, p.~102]{FG10}.
\end{enumerate}
\end{proof}

\begin{theorem}[Virtual Euler obstruction to the prescribed pinching]\hypertarget{tag-02ZP}{}\label{push:src:A.3.2}\label{push:A32}
\hypertarget{tag-02RB}{}\label{night:P2.11}\label{stmt:P2.11}
\StatementRecord{night:P2.11}{proved}{02RB}{the definitions and constructions stated in the proof}
\StatementRecord{push:src:A.3.2}{proved}{02ZP}{\SeriesRef{push:src:A.1.2}, \SeriesRef{push:src:A.2.2}, \SeriesRef{push:src:A.4.3}, \SeriesRef{push:A31}; \cite{Engel2026}}

For the pinching $Q$ of Proposition~\ref{fam:src:E.1.1},
one has $e_{\mathrm{vir}}(Q)=0$, whereas $e(X(u))=2$.
Every smooth proper-flat smoothing of $Q$ therefore has Euler
number zero. In particular there is no proper flat holomorphic
filling of Engel's untwisted family with central fibre exactly
$Q$. This disproves the central-fibre assertion in
\cite[Conjecture~8.2, p.~22]{Engel2026}. Theorem~\ref{push:src:A.1.2} constructs a nodal family with central fibre $Q$; Theorem~\ref{push:src:A.2.2} constructs the two smooth total spaces with central fibres $Q^\pm$.
\end{theorem}


\begin{proof}
\begin{enumerate}[label=\arabic*.,leftmargin=*]
\item Let $D_0,D_\infty\subset R$ be the conductor sections and
$\nu:R\to Q$ the normalization. On either branch of the
hypersurface $xy=0$, the derivative comparison with its
normalization gives in $K^0(R)$, where pullback of the
perfect tangent class means derived pullback,
\[
 \nu^*T_Q^{\mathrm{vir}}
 =T_R-i_{0*}(J|_{D_0})-i_{\infty*}(J|_{D_\infty}).
\]
Indeed in a smooth local ambient space the normal line $A$ of
the branch and the normal line of the whole hypersurface
differ by $D$: the latter is $A(D)$. Their difference in
$K$-theory is $-[i_*(A(D)|_D)]$, and $A(D)|_D$ is exactly
the product of the two branch normals, namely $\mathcal T^1_Q$.
These local comparisons are intrinsic and agree on overlaps.
On the infinity branch $J$ means its $T^{-1}$-pullback
from the double locus, canonically the same line bundle
because it comes from $B$. Thus $J$ extends from both sections by the pullback from $B$,
so $[i_*J]=[J]-[J(-D)]$.

\item Put $j=c_1(J)=\pi_R^*F$, $d_i=[D_i]$. Since the sections are
disjoint, $d_0d_\infty=0$, and since $F$ is a fibre, $j^2=0$.
Consequently
\[
 c(\nu^*T_Q^{\mathrm{vir}})
 =c(T_R)\frac{(1+j-d_0)(1+j-d_\infty)}{(1+j)^2}
 =c(T_R)\{1-d_0-d_\infty+j(d_0+d_\infty)\}.
\]
The degree-three part has integral
\[
 e(R)-\sum_i\int_{D_i}c_2(T_R|_{D_i})
       +\sum_i\int_{D_i}j\,c_1(T_R|_{D_i}).
\]
We have $e(R)=2e(S)=24$. On a section with normal class
$n_i$, $c_2(T_R|_{D_i})=c_2(S)+c_1(S)n_i$ and
$c_1(T_R|_{D_i})=c_1(S)+n_i$. The two normal classes on
the unglued normalization are $m,-m$, where $m=c_1(L)$.
Their sum is zero; $c_1(S)=F$ and $F^2=0$. The displayed
number is therefore $24-24+0=0$. Because
$\nu_*[R]=[Q]$, this is the intrinsic virtual number on $Q$.
In the projective-bundle presentation, the same calculation uses
$H(H+m)=0$, $d_0=H+m$, $d_\infty=H$.

\item The untwisted quotient away from its cusp has Euler number zero:
scalar multiplication supplies a free circle action, including
the two additive places before logarithmic transformation. Indeed
a circle scalar identified with an iterate of $u\phi$ must be
the identity, since every positive iterate strictly contracts
the Hermitian norm and every negative iterate expands it.
Compact fibres with a free circle action have Euler number zero.
The chosen Mumford fibre has two zero-dimensional torus orbits
and all other torus strata have Euler number zero, so
$e(X(u))=2$. Lemma~\ref{push:A31}
therefore excludes such a smoothing of $Q$. 

\end{enumerate}
\end{proof}

\begin{proposition}[An Euler-zero smoothing of the pinching]
\hypertarget{tag-PT02}{}\label{I:Euler-zero-pinching-smoothing}
\StatementRecord{I:Euler-zero-pinching-smoothing}{proved}{PT02}{}
The pinching $Q$ admits a proper flat smoothing with smooth fibres of Euler number zero, obtained by choosing the zero of its smoothing section on a smooth elliptic fibre of $S$.
\end{proposition}
\begin{proof}
Replace $\phi$ in the global construction of \ref{push:src:A.1.2}
by a section of $J$ vanishing on a smooth elliptic fibre of
$S$. The compact collar construction still applies. At its
zero divisor the fixed nonzero-parameter equation is $xy=uz$
with $dz\ne0$, so every nearby threefold is smooth. The
double-locus nearby Euler function is one on that elliptic
curve and zero off it; its integral is zero. Thus this proper flat
smoothing has Euler number zero, whereas the fibres of Engel's
family have Euler number two by Theorem~\ref{push:src:A.3.2}.
\end{proof}

\begin{proposition}[Infinitely many smooth fillings]\hypertarget{tag-02ZQ}{}\label{push:src:A.3.3}\label{push:A33}
\StatementRecord{push:src:A.3.3}{proved}{02ZQ}{\SeriesRef{push:section:A.3}, \SeriesRef{push:src:A.2.1}}

There are infinitely many fillings of the same punctured untwisted
family with smooth total space, reduced normal-crossings central
divisor and joint flatness over $B\times\Delta$. The two-model
classification in Proposition~\ref{push:src:A.2.1} concerns the minimal
toroidal refinements.
\end{proposition}


\begin{proof}
\begin{enumerate}[label=\arabic*.,leftmargin=*]
\item The following centres preserve joint flatness.
The section $2P$ is disjoint from $P$ and $O$. In the
Weierstrass equation its coordinates are
\[
 x(2P)=\frac{9t^2}{8},\qquad
 y(2P)=\frac{\sqrt2\,t^2(5t-32)}{32}.
\]
At $0$ it meets the nonidentity component different from
that met by $P$; at $1$ the displayed coordinates are finite
and different from those of $P,O$; at infinity the minimal
coordinates are finite and different as well. This also
checks directly the zero intersection numbers
$P\cdot2P=O\cdot2P=0$. The rational section of
$\mathcal O(P-O)$ therefore restricts to a nowhere-vanishing
holomorphic section on $2P\simeq B$. Its graph is a compact
curve $C\subset L^\times\subset R^\pm$ disjoint from both
conductor sections and from the modification centres.
It gives a curve in the smooth locus of $Q^\pm$ mapping
isomorphically to $B$.

\item Blow up $C$ in the smooth total fourfold $\mathcal X^\pm$.
Locally the coordinates are $(t,u,a,b)$ with
$C=(u,a,b)=0$. The blow-up charts replace $u$ by $u$ itself
or by $aU$ or $bU$. Thus the new central divisor is reduced
normal crossings, and the total space is smooth. All fibres
over $(t,u)$ still have pure dimension two; the regular base
criterion proves joint flatness. The morphism is proper and
is an isomorphism for $u\ne0$. It adds a
$\mathbb P^2$-bundle over $C$, hence a new central irreducible
component. To justify repetition, $N_{C/Q^\pm}$ is an extension
of $\mathcal O_B(-1)$ by $\mathcal O_B$ and splits, and
$H^1(N_{C/Q^\pm})=0$. The normal sequence in the total family
therefore splits off the trivial parameter direction. Its
section in the exceptional $\mathbb P^2$-bundle is disjoint
from the intersection with the strict central transform and
has the same normal sequence. We may repeat indefinitely.
Their differing component counts preclude
isomorphism as fillings.
\end{enumerate}
\end{proof}

\subsection{Integral topology and nearby cycles}\hypertarget{tag-02Z3}{}\label{push:section:A.4}

\begin{proposition}[The topology of the two modified pinchings]\hypertarget{tag-02ZR}{}\label{push:src:A.4.1}\label{push:A41}
\StatementRecord{push:src:A.4.1}{proved}{02ZR}{\SeriesRef{fam:src:E.1.1}, \SeriesRef{fam:src:E.1.2}, \SeriesRef{fam:src:E.1.3}, \SeriesRef{fam:src:E.1.4}, \SeriesRef{push:src:A.2.2}}

With integral coefficients, the spaces constructed above have
\[
\begin{array}{c|rrrrrrr|r}
&b_0&b_1&b_2&b_3&b_4&b_5&b_6&\chi\\\hline
R&1&0&11&0&11&0&1&24\\
R^\pm&1&0&12&0&12&0&1&26\\
Q&1&1&8&7&10&0&1&12\\
Q^\pm&1&1&9&7&11&0&1&14
\end{array}
\]
All these homology groups are torsion free. The fundamental
groups of $R,R^\pm$ are trivial and those of $Q,Q^\pm$ are
$\mathbb Z$. The dual complex of either $Q^\pm$ is a circle.
\end{proposition}


\begin{proof}
\begin{enumerate}[label=\arabic*.,leftmargin=*]
\item Proposition~\ref{fam:src:E.1.4} gives the integral groups and fundamental groups of $R,Q$, together with the normalization maps $d_i=s_{0*}-s_{\infty*}T_*$. We use exactly its lattice $\mathcal L=H_2(S;\mathbb Z)$, $F=[S_t]$, $m=P-O$, and set $v=v_0=[6P]-[O]-3F$. In these coordinates its map $a\mapsto(aF,av,am)$ is onto $\mathbb Z^3$, its $d_2$ has primitive image of rank three, and its $d_4$ is primitive and injective.

\item Now let $h:R^+\to R$ be the modification constructed in \ref{push:src:A.2.2}.
Its fibres are points off the nodal curve $C=S_\infty$, one
$\mathbb P^1$ over $C_{\mathrm{reg}}$, and a chain of two
$\mathbb P^1$'s over the node. The chart calculation identifies
the sheaf $R^2h_*\mathbb Z$ with
$\eta_*\mathbb Z_{\mathbb P^1}$, where
$\eta:\mathbb P^1\to C$ is the normalization. If $(a,b)$ are the cohomological
degrees on the two exceptional curves, generization to the
two branches is $a$ and $a+b$: on one branch the generic
exceptional fibre specializes to the first curve, on the
other it specializes to their sum. The matrix
$\left(\begin{smallmatrix}1&0\\1&1\end{smallmatrix}\right)$
is unimodular. These generizations follow by substituting
$X=za$, $v=aB$ in the blow-up charts, or by taking the two
limits of their exceptional fibres. Proper local neighbourhoods
retract onto these fibres by the positive toric contraction,
so the fibre calculation computes the higher direct images.
There is no other positive-degree direct image.

\item The Leray spectral sequence has only the row $H^*(R,\mathbb Z)$
and the row $H^{*-2}(\mathbb P^1,\mathbb Z)$; both are free
and in even degrees. It degenerates integrally for degree
reasons, giving the stated groups of $R^+$. The same local
neighbourhood retractions and Van Kampen show that $h$ induces
an isomorphism on fundamental groups, so $R^+$ is simply
connected. The case $R^-$ interchanges the branches.

\item For $Q^+$ use the normalization mapping cylinder once again.
The new degree-two difference map is the old $d_2$ with
one extra coordinate $aF$: the exceptional divisor restricts
to $F$ on the blown-up section and to zero on the other.
It has the same primitive rank-three image and the same
kernel; its cokernel rank increases by one. The degree-four
difference remains primitive, since its pushforward to $R$
already was primitive, and its cokernel likewise increases
by one. This proves the groups of $Q^+$. Interchanging the
sections changes the extra coordinate by a sign and does
not change any group. The dual complex has
one vertex and one self-edge, with no higher cells, and
Van Kampen again gives $\mathbb Z$.

\item In the modification $h$,
$C_{\mathrm{reg}}\simeq\mathbb C^*$ contributes zero and
the exceptional chain above the node increases Euler number
by $3-1=2$. Thus $\chi(R^\pm)=26$ and
$\chi(Q^\pm)=26-12=14$.
\end{enumerate}
\end{proof}

\begin{proposition}[Integral homology and monodromy of the smooth fibre]\hypertarget{tag-02ZS}{}\label{push:src:A.4.2}\label{push:A42}
\StatementRecord{push:src:A.4.2}{proved}{02ZS}{\SeriesRef{push:src:A.2.2}, \SeriesRef{push:src:A.4.1}}

The untwisted smooth fibre has torsion-free integral homology
with Betti numbers
\[
 (b_0,b_1,b_2,b_3,b_4,b_5,b_6)=(1,1,8,14,8,1,1).
\]
The positive-$u$ monodromy is the identity outside degree three.
In degree three there is an integral exact sequence
\[
 0\longrightarrow \mathcal L/\langle F,v,m\rangle
 \longrightarrow H_3(X(u);\mathbb Z)
 \longrightarrow K\longrightarrow0,\qquad
 K=\{a\in\mathcal L:aF=av=am=0\},
\]
in which $T_u-I$ is zero on the left term and, on the quotient,
is the map
\[
 K\longrightarrow \mathcal L/\langle F,v,m\rangle,\qquad
 a\longmapsto a\bmod\langle F,v,m\rangle.
\]
The sign convention orients the neck circle by increasing
$\arg x$ in $xy=u$; the boundary convention for a transported
chain is chosen so that a positive $u$ loop adds this circle
times its boundary cycle. With these conventions
$(T_u-I)^2=0$, its rank is six, and its nonzero Smith factors
are $1,1,1,1,1,11$.
\end{proposition}


\begin{proof}
\begin{enumerate}[label=\arabic*.,leftmargin=*]
\item We use the semistable model \ref{push:src:A.2.2}. Remove disjoint
open tubular neighbourhoods of the two conductor sections
from $R^+$, and denote the remaining compact manifold by
$Z$. Its two boundary circle bundles are identified by the
plumbing; let $E$ denote the first. The Euler class of
$E\to S$ is $n=m-F$. It is primitive. The Thom sequence
for $(R^+,Z)$ gives
\[
 H_5(Z)=\mathbb Z,\quad H_4(Z)=0,\quad
 H_3(Z)=\mathcal L/\langle F,m\rangle\simeq\mathbb Z^8,
 \quad H_2(Z)\simeq\mathbb Z^{10},\quad H_1(Z)=0.
\]
 If $H$ is
the projective-bundle class restricting to zero on $s_0$
and to $-m$ on $s_\infty$, and $A$ is the exceptional
divisor class, the two restriction maps on $H^2(R^+)$ are
\[
 (\alpha,h,a)\longmapsto\alpha+aF,\qquad
 (\alpha,h,a)\longmapsto\alpha-hm.
\]
Their combined image is primitive of rank twelve, and its
cokernel is $\mathcal L/\langle F,m\rangle$. In degree two,
intersection with the two sections is onto $\mathbb Z^2$:
a generic projective fibre meets both once, and an exceptional
projective fibre meets just the section containing the centre
once. In degree six the map is the diagonal
$\mathbb Z\to\mathbb Z^2$. These are precisely the Thom
maps giving the asserted groups.

\item The Gysin sequence of the primitive circle bundle gives
\[
 H_2(E)=n^\perp\subset\mathcal L,\qquad
 H_3(E)=\mathcal L/\mathbb Z n,\qquad
 H_1(E)=H_4(E)=0,\quad H_5(E)=\mathbb Z.
\]
The smooth fibre is obtained from $Z$ by the boundary
mapping cylinder. In its Mayer--Vietoris sequence the
degree-two difference map is the normalization difference
$d_2^+$ restricted to $n^\perp$. In the coefficient
coordinates $(s,t,h)=(aF,av,am)$ the constraint is $h=s$,
so this map has primitive rank two and kernel $K$.
The degree-three difference map is, modulo $F,m$,
$(T_*-I)a=(aF)v$, hence has primitive image $\mathbb Z v$.
These statements use the surjectivity of the three
intersection coordinates proved in \ref{push:src:A.4.1}. The degree-five difference is zero:
the plumbing reverses the normal circle orientation, and
the two boundary fundamental classes sum to zero in $Z$.
The degree-zero difference is zero as well. Exactness now
gives the claimed groups and the displayed degree-three
extension. Poincar\'e duality, or the same integral exact
sequence, confirms the absence of torsion in complementary
degrees.

\item We also compute monodromy directly from this decomposition.
In the neck $xy=u$, a positive circuit of $u$ rotates one
boundary normal coordinate once relative to the other.
A cycle in $Z$ is unchanged. A three-cycle whose cut-open
boundary represents $a\in K$ acquires the tube obtained
by sweeping $a$ once around that normal circle. The Gysin
description sends this tube to the class of $a$ in
$H_3(E)=\mathcal L/\mathbb Z n$, and then to its class in
$\mathcal L/\langle F,m,v\rangle$. This is exactly the
map in the statement. The same description in the other
degrees is zero, using the groups and difference maps just
computed. A tube has no cut-open boundary, which proves
$(T_u-I)^2=0$.

\item Finally $K\cap\langle F,m,v\rangle=\mathbb ZF$, since
the Gram matrix of $m-F,v$ has determinant eleven.
Hence the rank is six. In the unimodular $E_8$ summand,
the primitive rank-two lattice spanned by $m-F,v$ has
discriminant eleven; its orthogonal complement maps into
the dual quotient with cokernel of order eleven. Thus the nonzero Smith factors are five ones and eleven.
\end{enumerate}
\end{proof}

\begin{proposition}[The pointwise nearby Euler function]\hypertarget{tag-02ZT}{}\label{push:src:A.4.3}\label{push:A43}
\StatementRecord{push:src:A.4.3}{proved}{02ZT}{\SeriesRef{push:src:A.3.2}}

For $\mathcal X^\flat\to\Delta_u$ the nearby-fibre Euler
function on the double locus of $Q$ is zero off $S_\infty$
and one at every point of $S_\infty$, including its node.
Its integral is one. The smooth part of $S_\infty$ has
Euler number zero, so its integrated contribution is zero.
For either semistable family $\mathcal X^\pm$ the nearby
Euler function is one off the double locus and zero on
it, giving $\chi(X(u))=26-2\cdot12=2$.
\end{proposition}


\begin{proof}
\begin{enumerate}[label=\arabic*.,leftmargin=*]
\item Off the zero divisor of $\phi$, the transverse nearby
fibre of $xy=u\phi$ is an annulus and has Euler number
zero. At a smooth zero, the local equation is $xy=uz$;
for fixed $u\ne0$ one solves for $z$, and the small
local nearby space is contractible. Its Euler number
is one. At the node it is the quadratic cone
$xy=uzv$, also contractible, with Euler number one.
The smooth part of the nodal elliptic curve is
$\mathbb C^*$ and has Euler number zero; the node
contributes one. The core $Q\setminus S$ is a
$\mathbb C^*$-bundle and contributes zero.
This proves the first count.

\item In the resolved family the local map on the double locus
is simply $ab=u$ times two transverse coordinates,
including above the old node. Its nearby annulus has
Euler number zero everywhere on that locus. Away from
it, the family is a submersion and the local value is
one. Integrating gives the second count. If the nodal
threefold in the first family is smoothed rather than
small-resolved, its ordinary-double-point Milnor fibre
is homotopy equivalent to $S^3$ and has Euler number
zero; the replacement changes $1$ to $0$, whereas
the small resolution changes $1$ to $2$. These agree
with the intrinsic virtual number in \ref{push:src:A.3.2}.
\end{enumerate}
\end{proof}

\begin{proposition}[Non-K\"ahler models and the monodromy weight filtration]\hypertarget{tag-02ZU}{}\label{push:src:A.4.4}\label{push:A44}
\StatementRecord{push:src:A.4.4}{proved}{02ZU}{\SeriesRef{push:src:A.4.1}}

The smooth fibres of these untwisted models have $b_1=1$ and
admit no K\"ahler metric.
Moreover the exceptional small-resolution curve in
$R^\pm$ is null-homologous, so that normalization is
not K\"ahler either. The rational monodromy weight
filtration, defined algebraically from $N=T_u-I$ and
centred at degree three, has graded ranks $6,2,6$ in
weights $2,3,4$.
\end{proposition}


\begin{proof}
\begin{enumerate}[label=\arabic*.,leftmargin=*]
\item The odd first Betti number excludes a K\"ahler metric.
In the generization matrix of \ref{push:src:A.4.1} a global section of
$R^2h_*\mathbb Z$ has equal branch values, hence degree
zero on the small exceptional curve. All classes pulled
back from $R$ also have degree zero on it. The Leray
calculation exhausts $H^2(R^\pm)$, so the curve has
zero homology class; it cannot have positive area for
a K\"ahler form. Finally $N^2=0$, rank $N=6$, and
$\dim H_3=14$, so the centred filtration is
$W_2=\operatorname{im}N$, $W_3=\ker N$, $W_4=H_3$.
Its graded dimensions are the asserted ones.
\end{enumerate}
\end{proof}



\Uses{the hypotheses, notation and definitions stated here; \cref{push:src:A.2.4}; \cref{push:src:A.4.7}; \cref{push:src:A.4.9}; \cref{push:src:A.4.8}; \cref{push:src:A.4.4}}


\begin{proposition}[An infinite explicit subclass]\hypertarget{tag-02ZW}{}\label{push:src:A.4.6}\label{push:A46}
\StatementRecord{push:src:A.4.6}{proved}{02ZW}{\SeriesRef{push:src:A.2.1}, \SeriesRef{push:src:A.4.2}, \SeriesRef{push:src:A.4.1}}

Let $D_m$ be the central fibre after the $m$ maximal-cone
star subdivisions in \ref{push:src:A.2.1}, starting with $Q^+$, where
$m\geq0$. Then
\[
 \begin{gathered}
 \pi_1(D_m)=\mathbb Z,\qquad
 H_*(D_m;\mathbb Z)\text{ is free},\\
 (b_0,\ldots,b_6)=(1,1,9+m,7,11+m,0,1+m).
 \end{gathered}
\]
and $\chi(D_m)=14+3m$. For $m\geq1$ its normalization is
the disjoint union of $\operatorname{Bl}_pR^+$, $m-1$
copies of $\operatorname{Bl}_{p_i}\mathbb P^3$, and a final
$\mathbb P^3$. The original two sections are still glued
by $T$. The new components form a chain attached to the
original component. Consecutive components meet in one
$\mathbb P^2$, with no triple intersections. The dual
complex is the original circle with a chain of $m$ edges
attached. The nearby Euler function is one off the double
locus and zero on it; its integral remains two. The nearby
integral monodromy remains exactly that of \ref{push:src:A.4.2}.
\end{proposition}


\begin{proof}
\begin{enumerate}[label=\arabic*.,leftmargin=*]
\item The torus orbit of the cone $ABCE$ is a point on a single
local central branch, because precisely one of its four
rays has positive $u$-height. Hence it is a smooth point
of the central fibre, away from the conductor. A star at
the sum of the four basis rays is the ordinary blow-up
of that point in the smooth fourfold. The central divisor
is the union of its strict transform and the exceptional
$\mathbb P^3$, both with multiplicity one. Their
intersection is the exceptional $\mathbb P^2$ in the
blow-up of the central threefold, identified with a
hyperplane in $\mathbb P^3$. Their two normal lines there
are $\mathcal O(-1)$ and $\mathcal O(1)$.

\item The next distinguished cone $ABCv_1$ has only the new
central ray $v_1$. Its point lies on the new $\mathbb P^3$
away from its attaching hyperplane. Repeating therefore
blows up that component at a point away from the preceding
double locus and attaches the next $\mathbb P^3$. The
incoming hyperplane and outgoing exceptional plane are
disjoint in each middle $\operatorname{Bl}_{p_i}\mathbb P^3$.
This proves the complete normalization and gluing
description, the absence of triples and the dual complex.

\item For the integral calculation, blowing up a smooth point
of a compact complex threefold adds one free generator
in each of $H_2$ and $H_4$ and changes no other homology.
This also holds when the threefold has singularities away
from that point: remove a small ball and use the
Mayer--Vietoris sequence with boundary $S^5$; the
replacement retracts to $\mathbb P^2$. Now attach
$\mathbb P^3$ along that $\mathbb P^2$. In degrees
$0,2,4$ the inclusion into $\mathbb P^3$ is an integral
isomorphism, so the Mayer--Vietoris boundary map is
injective and primitive. It contributes no torsion or
odd homology. The new component contributes one free
$H_6$ generator, since $H_6(\mathbb P^2)=0$. Each step
therefore adds precisely one to $b_2,b_4,b_6$.
Van Kampen gives no change of fundamental group: the
blow-up ball replacement, the attaching plane and the
new projective component are simply connected.
Induction from \ref{push:src:A.4.1} proves the asserted groups.

\item Independently, the strict-transform blow-up adds two to
the central Euler number, the new $\mathbb P^3$ adds
four, and its intersection $\mathbb P^2$ subtracts
three. Thus each step adds three, agreeing with the
homology computation. The double locus is the disjoint
union of the original $S$ and $m$ planes, of Euler number
$12+3m$. Since the local equation is $u=x$ off this locus
and $u=xy$ on it, the nearby Euler integral is
$(14+3m)-(12+3m)=2$. All modifications are the identity
over $u\ne0$, so they do not alter the marked punctured
family or its homological monodromy.
\end{enumerate}
\end{proof}

\begin{theorem}[Euler number and fundamental group for every toroidal output]\hypertarget{tag-02ZX}{}\label{push:src:A.4.7}\label{push:A47}
\StatementRecord{push:src:A.4.7}{proved}{02ZX}{\SeriesRef{push:src:A.2.4}, \SeriesRef{push:src:B.4.2}; \cite[Theorem~1.6 and the following noncompact extension, p.~215]{BM97}}

Let $\Sigma$ be any output of \ref{push:src:A.2.4} for $u=r^n$, and let
$D_\Sigma$ denote the underlying reduced central space.
Let $k(\Sigma)$ be the number of selected pure top ray
orbits, and let $t_4(\Sigma)$ be the number of maximal
four-cones in one fundamental prism. Then
\[
 \chi(D_\Sigma)=11k(\Sigma)+t_4(\Sigma),
 \qquad \pi_1(D_\Sigma)=\mathbb Z,
 \qquad H_1(D_\Sigma;\mathbb Z)=\mathbb Z.
\]
These assertions include nonreduced central divisors.
At a point with precisely one local central branch of
multiplicity $\ell$, the nearby Euler function is $\ell$;
at a point with at least two local branches it is zero.
Its integral is always two. For every reduced output the
pure top count is $k(\Sigma)=n$, so
$\chi(D_\Sigma)=11n+t_4(\Sigma)$.
The full integral homology is given by Theorem~\ref{push:src:A.4.8}.
\end{theorem}


\begin{proof}
\begin{enumerate}[label=\arabic*.,leftmargin=*]
\item Over $B\setminus\{\infty\}$ the top fan gives a cycle
of $k(\Sigma)$ ruled bundles over
$S\setminus S_\infty$, glued along their zero and
infinity sections with the specified translations.
There are no triple intersections there. Its Euler
number is therefore
\[
 2k(\Sigma)\chi(S\setminus S_\infty)
       -k(\Sigma)\chi(S\setminus S_\infty)
 =11k(\Sigma),
\]
because $\chi(S)=12$ and $\chi(S_\infty)=1$.
This computation concerns the support of the divisor,
so the multiplicities do not enter.

\item At $q=r=0$ the strata are the torus orbits of cones whose
generators collectively have positive $q$-height and
positive $r$-height. Each such cone has trivial period
stabilizer. Indeed, an element $G_1^aG_2^b$ preserving
its finite set of primitive rays has a power fixing them
all. A ray of positive $q$-height then forces $a=0$;
any ray with positive $q$- or $r$-height forces $b=0$.
Consequently a stratum of cone dimension $d$ is
$(\mathbb C^*)^{4-d}$.
The induced gluing units only identify different period
translates of the stratum. Its compactly supported Euler
number is zero unless $d=4$, in which case it is one.
Every maximal cone has both kinds of height, and each has
one representative in the interior of the fundamental
prism. Thus the cusp contributes exactly $t_4(\Sigma)$.
Additivity proves the displayed central Euler formula.

\item The local total space is smooth, and its parameter is
$r=z_1^{\ell_1}\cdots z_s^{\ell_s}$ times a holomorphic
unit, with $\ell_i>0$. Its Milnor fibre has the homotopy
type of a finite disjoint union of $(S^1)^{s-1}$.
For $s=1$ it consists of $\ell_1$ contractible pieces;
for $s\geq2$ its Euler number is zero. This proves the
pointwise assertion. Off the cusp the one-branch parts
are $\mathbb C^*$-bundles over the appropriate parts of
$S\setminus S_\infty$ and have Euler number zero; the
double parts have nearby Euler value zero. At the cusp,
positive-dimensional torus strata again contribute zero.
A maximal cone with exactly one positive-$r$ ray must
have three bottom rays. Its bottom face is consequently
one of the two prescribed triangles $ABC$ and $BCD$.
Each triangle belongs to exactly one such maximal cone.
Its determinant is the $r$-height $\ell$ of the fourth
ray, because the bottom triangle is unimodular.
Smoothness forces $\ell=1$. Exactly these two points
contribute one each. Thus the nearby Euler integral is
two for every output, even when some other components
have multiplicity greater than one.

\item For the fundamental group, first choose the reference
model obtained by base-changing $\mathcal X^+$ by
$u=r^n$ and resolving the surface equation $xy=r^n$
along its smooth conductor $S$. Its normalized central
components are $R^+$ and $n-1$ projective-line bundles
over $S$. All are simply connected, and their double
surfaces are copies of the simply connected $S$.
The gluing graph is an $n$-cycle, interpreted as the
self-gluing circle when $n=1$. Van Kampen therefore
gives fundamental group $\mathbb Z$ for this reference
centre.

\item Any two fans of \ref{push:src:A.2.4} have the explicit proper common
normal modification formed by intersecting their cones.
Resolve this analytic space to a smooth common model;
analytic resolution, including the noncompact version
near the compact central set, is supplied by
\cite[Theorem~1.6 and the following noncompact extension, p.~215]{BM97}.

\item A proper bimeromorphic map $F:Y\to X$ between smooth
complex manifolds induces an isomorphism of fundamental
groups. To see this directly, choose an analytic
$A\subset X$ of complex codimension at least two such
that $F$ is an isomorphism outside $A$. General position
gives $\pi_1(X\setminus A)\simeq\pi_1(X)$, and gives
surjectivity of
$\pi_1(Y\setminus F^{-1}A)\to\pi_1(Y)$.
Under the common-open identification, composition with
$F_*$ is the displayed isomorphism to $\pi_1(X)$.
Thus the first map is also injective, and $F_*$ is an
isomorphism. Apply this argument to both maps from the
smooth common model, after restricting to a sufficiently
small parameter disc.

\item Finally, each smooth total space retracts to the support
of its central divisor, including in the nonreduced case.
The gradient proof in \ref{push:src:B.4.2} applies with the explicit
bound for $\rho=|r|^2$
\[
 \|\nabla\rho\|\geq c\rho^{1-1/(2L)},
 \qquad L=\ell_1+\cdots+\ell_s.
\]
On the complement of the divisor it follows from the
weighted arithmetic--geometric mean inequality applied
to $\sum_i\ell_i^2/|z_i|^2$. A holomorphic unit can be
absorbed into one coordinate locally. Compactness of
the central set gives a uniform exponent below one,
and the resulting finite remaining path-length bound
gives the continuous gradient retraction. Hence all
the centres and their smooth total neighbourhoods have
the same fundamental group as the reference centre.
Abelianization gives the stated $H_1$.

\item For the nonreduced conductor-star family with added pure top
rays $(0,1,0,j)$, $2\leq j\leq m+1$, the pure-top ray-orbit count increases by $m$
and the maximal-cone-orbit count increases by $m$. The formula therefore gives
$\chi(D_m)=14+12m$. Equivalently, each conductor blow-up adds the Euler
number $12$ of its centre.
\end{enumerate}
\end{proof}

\begin{theorem}[Integral homology of every toroidal output]\hypertarget{tag-02ZY}{}\label{push:src:A.4.8}\label{push:A48}
\StatementRecord{push:src:A.4.8}{proved}{02ZY}{\SeriesRef{push:src:A.2.4}, \SeriesRef{push:src:A.4.2}, \SeriesRef{push:src:A.4.1}, \SeriesRef{push:src:A.4.7}; \cite[Section 5.6, pp.565--566]{AKMW02AB}; \cite[Theorem~1.6.1, p.~543; Theorem~2.7.1, pp.~551--552]{AKMW02AB}; \cite[Section 2, p.2]{Matsuki00ErrAB}; \cite[Section~5.5, pp.~564--565]{AKMW02AB}}

Let $D_\Sigma$ be the reduced central space of an output of
\ref{push:src:A.2.4} for $u=r^n$. Write $k$ for its number of pure top
ray orbits, $m$ for its number of mixed ray orbits, and
$t_4$ for its number of maximal four-cone orbits. Then
every integral homology group is free, and
\[
 \begin{split}
 (b_0,b_1,b_2,b_3,b_4,b_5,b_6)
   ={}&(1,1,8+k+m,7,\\
       &\hspace{20mm}9k+t_4-2m-1,0,k+m).
 \end{split}
\]
This includes nonreduced central divisors: the homology
is that of their underlying support. The marked nearby
fibre is $X(r^n)$, and its integral monodromy is
$I+nN$ in degree three, where $N$ is the map in \ref{push:src:A.4.2};
it is the identity in every other degree. In particular
the six nonzero Smith factors of this degree-three
monodromy minus the identity are
$n,n,n,n,n,11n$.
\end{theorem}


\begin{proof}
\begin{enumerate}[label=\arabic*.,leftmargin=*]
\item We first compute a reference centre by its normalization
mapping cylinder. Base-change $\mathcal X^+$ by $u=r^n$
and resolve $xy=r^n$ along its smooth conductor $S$.
Put $N_0=\mathcal O_S(m-F)$, using the lattice notation
of \ref{push:src:A.4.1}. The normalization consists of $R^+$ and
$n-1$ copies of $\mathbb P_S(\mathcal O\oplus N_0)$,
in a cycle. The latter assertion follows directly on
the surface-resolution charts: if $x$ has transition
function $g$, the coordinates $x/r^i$ have the same
transition function, while their complementary
coordinates have transition function $g^{-1}$.
Only the closing identification is $T$.

\item In the integral mapping-cylinder sequence the degree-two
map has domain $\mathcal L^n$ and target of rank $11n+1$.
Project each intermediate ruled component to $S$.
The resulting $n-1$ difference blocks are integral
identity matrices. Eliminate them by integral row and
column operations. The remaining map from $\mathcal L$
is $d_2^+$ from \ref{push:src:A.4.1}, together with $n-1$ coordinates
$a\mapsto a(m-F)$. Those coordinates are integral
combinations of the three primitive coordinates of
$d_2^+$. Thus the original degree-two map has primitive
image of rank $10(n-1)+3$, kernel $\mathbb Z^7$, and
cokernel $\mathbb Z^{n+8}$.

\item In degree four, the same elimination uses the $n-1$
identity blocks on the fundamental classes of the bases.
The remaining map from $\mathbb Z$ contains $d_4^+$,
which is primitive by \ref{push:src:A.4.1}; the other entries are the
section differences $m-F$. Hence the degree-four map
is injective with primitive image of rank $n$, and
cokernel of rank $10n+1$. In degree zero it is the
incidence map of an $n$-cycle. In degree six the
normalization contributes one independent fundamental
class per component. Consequently the reference centre
has free integral homology and Betti tuple
\[
 (1,1,n+8,7,10n+1,0,n).
\]

\item On each total
space take the toroidal boundary
\[
 E=\bigl(\{r=0\}\cup f^{-1}(\infty)\bigr)_{\rm red}.
\]
In the cusp charts it is the toric boundary, and outside
the cusp its nontrivial toroidal charts are the plumbing
charts along $S$. The unit transition functions preserve
monomial ideals. \ref{push:src:A.2.4} gives a finite cone complex modulo
the two periods, with the possible self-identifications
of its faces explicitly retained.

\item Make this complex strict by barycentric star subdivisions,
in decreasing cone dimension. At the moment a stratum
is blown up, all of its deeper self-intersections have
already been separated. Its closure is therefore smooth,
and different translates in the chosen orbit are disjoint.
The star at the sum of the primitive basis generators is
the ordinary blow-up of that closure. Repeating over the
finitely many cone orbits gives smooth strict toroidal
embeddings. This is the geometric barycentric subdivision
used in \cite[Section 5.6, pp.565--566]{AKMW02AB}; its
analytic realization by the same monomial ideals is
described in \cite[Section~5.5, pp.~564--565]{AKMW02AB}. Pure $q$-strata
are included in this preparatory step and in the comparison of
the punctured $r$-families.

\item Intersect the two subdivided fans, and take a regular
subdivision, to obtain their explicit common strict
toroidal modification. Apply toroidal weak factorization
to each of its two subdivision maps. The needed theorem
is the finite polyhedral-complex factorization by
nonsingular star subdivisions and inverse subdivisions,
\cite[Theorem~1.6.1, p.~543; Theorem~2.7.1, pp.~551--552]{AKMW02AB}. The corrected
form for a morphism states that its centres are smooth
irreducible closed strata and that all intermediate
complexes map to the target
\cite[Section 2, p.2]{Matsuki00ErrAB}.

\item The complexes are finite, so the cited combinatorial
factorization applies. The analytic
subdivision--modification correspondence of
\cite[Section~5.5, pp.~564--565]{AKMW02AB} realizes each of its finitely
many stars by the indicated monomial blow-up.
The correspondence respects the face identifications
and the bundle units, so the local analytic blow-ups
glue. All intermediate morphisms are proper over the
target of the subdivision map, hence over a sufficiently
small $r$-disc. They remain smooth, and both parameters
remain monomials times units. This gives the required
finite chain of analytic blow-ups and blow-downs
between the total neighbourhoods.

\item  A
nontrivial centre has codimension at least two. For a
cone having a positive $q$-height ray, its period
stabilizer is trivial: a power of a stabilizing element
fixes each of its finitely many primitive rays; the
positive $q$-height first forces the exponent of $G_1$
to vanish and then the exponent of $G_2$ to vanish.
The same argument holds for pure bottom cones. The normalization of the closure of their torus stratum
is the toric variety of its star. If the closed stratum
is smooth, this finite birational normalization map is
an isomorphism.

\item For a mixed cone the centre is contained in $q=r=0$;
its star is complete, and it is a compact smooth toric
point, curve, or surface. It has free integral
cohomology in even degrees only. For a pure top centre,
the cone must have dimension two. Its stabilizer is
precisely $\langle G_1\rangle$, and the chart quotient
and the plumbing identification show that its dense
open part is the corresponding open part of $S$.
The compact smooth centre is therefore a surface
birational to the rational surface $S$. Equivalently,
its proper map to the unmodified top stratum is a
birational map of smooth surfaces. It is a smooth
rational surface and again has free even cohomology.
It need not remain isomorphic to $S$ after earlier
mixed modifications.

\item For a pure bottom centre the cone has dimension two
or three. Its star gives a smooth toric surface or
curve proper over the $r$-disc. The induced character
$r$ is primitive on its quotient lattice, since it
was primitive before quotienting by a smooth cone
on which it vanishes. Its generic fibre is therefore
connected, and is respectively a complete toric
$\mathbb P^1$ or a point. In the surface case the
ordered two-dimensional fan gives a chain of rational
curves as the central support, possibly with positive
multiplicities. In the curve case the support is one
point. Monomial gradient retraction, as in \ref{push:src:A.4.7},
shows that the ordinary integral cohomology of the
centre is free and even, with $H^0=\mathbb Z$ and
$H^4=0$. These centres are not compact; this is why
they must be distinguished from the first two types.

\item For any closed smooth complex centre $C\subset Y$ of
complex codimension $c$, compact or not, the integral
blow-up splitting is
\[
 H^j(\operatorname{Bl}_C Y;\mathbb Z)
 \simeq H^j(Y;\mathbb Z)\oplus
       \bigoplus_{a=1}^{c-1}H^{j-2a}(C;\mathbb Z).
\]
Tubular excision identifies the relative
groups by the integral Thom class; the projective-bundle
theorem for the exceptional divisor supplies the powers
of its hyperplane class. Their exceptional Gysin maps,
together with pullback from $Y$, give the displayed
map. The relation that restriction of the exceptional
class is minus the hyperplane class makes the relative
matrix triangular with unit diagonal. Proper
blow-down pushforward, defined by Borel--Moore duality,
splits the pullback summand. This proves the splitting
also for the noncompact centres just listed.

\item It follows that all odd cohomology groups and all
torsion are unchanged along the factorization. Each
connected centre adds one to $b_2$ and one to the number
of prime boundary components. Each compact centre
adds one to $b_6$ and one to the number of central
prime components: in degree six the only new group
is its top cohomology. A pure bottom centre adds
neither, because its $H^4$ (and, for a curve, its
$H^2$) vanishes. The reverse operations subtract
the same contributions.

\item Consequently $b_2$ minus the number of all boundary
ray orbits, and $b_6$ minus the number of central ray
orbits, are invariants of this comparison. Both
endpoints have exactly one pure bottom ray orbit.
The reference has $n$ central components and
$b_2=n+8$, so the arbitrary output has
\[
 b_2=8+k+m,\qquad b_6=k+m,
 \qquad (b_1,b_3,b_5)=(1,7,0).
\]
All groups are free by the integral splitting in both
directions. The total neighbourhoods retract onto
their central supports by \ref{push:src:A.4.7}, and those compact
analytic supports have finite triangulations. The
universal coefficient theorem therefore transfers
the same freeness and ranks to integral homology.
Finally the independently computed torus-stratum
Euler number $11k+t_4$ of \ref{push:src:A.4.7} gives
$b_4=9k+t_4-2m-1$.

\item Every output agrees with the marked family
$X(r^n)$ off $r=0$. A positive $r$-loop makes $n$
positive $u$-loops. Since $N^2=0$, its monodromy is
$(I+N)^n=I+nN$. Multiplication of the already computed
integral Smith matrix by $n$ gives the stated factors.

\end{enumerate}
\end{proof}

\begin{proposition}[Integral cohomology under central blowups]\hypertarget{tag-02ZZ}{}\label{push:src:A.4.9}\label{push:A49}
\StatementRecord{push:src:A.4.9}{proved}{02ZZ}{\SeriesRef{push:src:A.4.8}, \SeriesRef{push:src:A.2.5}, \SeriesRef{push:src:A.4.7}}

The homology formula of \ref{push:src:A.4.8} can also be computed
on the central spaces, by their normalization
resolutions and the strata of each recorded star.
More precisely, let $Y'\to Y$ be one ordinary
blow-up in the algorithm of \ref{push:src:A.2.5}, with closed
smooth centre $C$ of complex codimension $c$.
Let $D,D'$ be the reduced central supports and
$C_0=C\cap D$. In the derived category of integral
sheaves on $D$ there is an isomorphism
\begin{equation}
 R(\pi_D)_*\mathbb Z_{D'}\simeq
 \mathbb Z_D\oplus
 \bigoplus_{a=1}^{c-1}(i_0)_*\mathbb Z_{C_0}[-2a],
\label{pub:I:eq:A-49a}
\end{equation}
where $\pi_D:D'\to D$ and $i_0:C_0\hookrightarrow D$
are the induced continuous maps. In particular
\begin{equation}
 H^j(D';\mathbb Z)\simeq H^j(D;\mathbb Z)
 \oplus\bigoplus_{a=1}^{c-1}H^{j-2a}(C_0;\mathbb Z).
\label{pub:I:eq:A-49b}
\end{equation}
The maps in this splitting are integral, with unit
diagonal on the exceptional-fibre basis. Thus it
also determines the primitive boundary maps and
the absence of torsion in the normalization
mapping-cylinder computation along the entire
finite algorithm.
\end{proposition}


\begin{proof}
\begin{enumerate}[label=\arabic*.,leftmargin=*]
\item To construct the sheaf maps, let $E=\mathbb P_C(N_{C/Y})$
be the exceptional divisor, let $p:E\to C$, and put
$\xi=c_1(\mathcal O_E(1))$. Pullback from $Y$ gives
the degree-zero summand. For $1\leq a\leq c-1$,
pullback by $p$, multiplication by $\xi^{a-1}$,
and the exceptional Gysin morphism give a derived
map
\[
 i_*\mathbb Z_C[-2a]\longrightarrow
 R\pi_*\mathbb Z_{Y'},\qquad i:C\hookrightarrow Y.
\]
The Gysin map exists integrally for a closed complex
divisor, by its oriented tubular Thom class; neither
the centre nor the ambient space must be compact.
At a point of $C$, proper base change identifies
the stalk of the target with
$H^*(\mathbb P^{c-1};\mathbb Z)$.
Restriction of the exceptional normal line is
$\mathcal O_E(-1)$, so the displayed maps send
the generators to
\[
 1,-\xi,-\xi^2,\ldots,-\xi^{c-1}.
\]
This is an integral basis. Off $C$ the map is
the identity. It is therefore a stalkwise
quasi-isomorphism
\[
 \mathbb Z_Y\oplus
 \bigoplus_{a=1}^{c-1}i_*\mathbb Z_C[-2a]
 \xrightarrow{\ \simeq\ }R\pi_*\mathbb Z_{Y'}.
\]

\item Now restrict this sheaf morphism to $D$.
The central support upstairs is the whole inverse
image $\pi^{-1}D$, including when the centre is
horizontal. Proper base change applies to this
Cartesian square. The inverse-image functor for
integral sheaves is exact, and the inverse image
of $i_*\mathbb Z_C$ is
$(i_0)_*\mathbb Z_{C_0}$. This proves \eqref{pub:I:eq:A-49a}.
Taking derived global sections proves \eqref{pub:I:eq:A-49b}.
This calculation concerns the stratified map of central spaces.

\item To compare this splitting with the normalization mapping
cylinders, resolve
$\mathbb Z_{D'}$ by the alternating ordered-local-branch
complex: its first term is the constant sheaf on
the normalization, its next term is on the normalized
double-branch space, and so on. Its augmentation is
an integral quasi-isomorphism, because at a point
with $s$ branches it is the augmented simplex
complex on those branches. Use alternating integral cochains, or
equivalently the orientation local system of each
unordered branch set. Proper pushforward of this
resolution computes the left side of \eqref{pub:I:eq:A-49a}.
The same resolution for $D$, together with the
centre complex, computes its right side.

\item The local matrices can be read directly from the
blow-up charts in \ref{push:src:A.2.5}. If the centre lies on $k$
central branches among its $c$ defining coordinates,
the exceptional fibre is $\mathbb P^{c-1}$. The new
exceptional branch supplies that whole projective
space; the strict transform of old branch $i$
supplies its coordinate hyperplane $H_i$.
Intersections supply the corresponding coordinate
projective subspaces. In each even cohomological
degree the normalization differential is the
alternating restriction matrix with entries
$0,1,-1$. The cover consisting of the whole
projective space and its selected hyperplanes
has the elementary simplex contraction, leaving
one generator $\xi^a$ in each degree $2a$,
$0\leq a<c$. For a horizontal centre, every
central branch meeting its generic central
stratum supplies the whole exceptional fibre;
the same simplex contraction applies to these
copies. Higher boundary strata are treated by
the same coordinate restrictions. The global
Gysin maps just constructed glue the surviving
classes with unit coefficient.

\item Starting from the explicitly computed reference
normalization matrices in \ref{push:src:A.4.8}, \eqref{pub:I:eq:A-49a} therefore
computes the normalization complex along every
recorded blow-up or inverse blow-up. A compact
centre has $C_0=C$. A horizontal toric surface
centre retracts to its rational-chain $C_0$;
a horizontal curve centre has $C_0$ a point.
Their cohomology is free and even. Consequently
the central computation preserves odd cohomology
and torsion, adds one to $b_2$ for every new
boundary ray, and adds one to $b_6$ exactly for
a new central ray. It gives the same
$(b_1,b_3,b_5)=(1,7,0)$,
$b_2=8+k+m$ and $b_6=k+m$ as \ref{push:src:A.4.8}.

\item Independently, stratification of $\pi_D$ into
$D\setminus C_0$ and $C_0$, with fibres a point
and $\mathbb P^{c-1}$, gives
\[
 \chi(D')-\chi(D)=(c-1)\chi(C_0).
\]
The absolute torus-stratum count of \ref{push:src:A.4.7} is
$\chi(D)=11k+t_4$. Together with the central
group calculation this gives again
$b_4=9k+t_4-2m-1$. Thus the two computations
agree integrally and in Euler number.

\end{enumerate}
\end{proof}


\begin{remark}[Topology determined by the fan data]\hypertarget{tag-02ZV}{}\label{push:src:A.4.5}\label{push:A45}
\StatementRecord{push:src:A.4.5}{proved}{02ZV}{\SeriesRef{push:src:A.2.4}, \SeriesRef{push:src:A.4.4}, \SeriesRef{push:src:A.4.7}, \SeriesRef{push:src:A.4.8}, \SeriesRef{push:src:A.4.9}}
For an output of Theorem~\ref{push:src:A.2.4},
Theorem~\ref{push:src:A.4.7} computes the Euler number, fundamental
group and nearby Euler function, and Theorem~\ref{push:src:A.4.8}
gives all integral homology groups and the marked monodromy.
Proposition~\ref{push:src:A.4.9} identifies the integral maps in the
normalization calculation through its direct-image splitting and
unit exceptional-fibre diagonals. All these calculations use the
same marked cone data.
\end{remark}

\section{Local limits and the global component obstruction}\hypertarget{tag-02Z4}{}\label{push:section:B}
\subsection{The finite root neighbourhoods}\hypertarget{tag-02Z5}{}\label{push:section:B.1}

\begin{proposition}[Finite-root plumbing and ruled central fibres]\hypertarget{tag-0300}{}\label{push:src:B.1.1}\label{push:B11}
\StatementRecord{push:src:B.1.1}{proved}{0300}{\SeriesRef{fam:src:B.1.1}, \SeriesRef{fam:src:B.1.4}, \SeriesRef{push:src:A.1.2}, \SeriesRef{fam:src:E.2.1}}

The marked root family admits a proper flat holomorphic filling over
$\Delta_{s_j}\times\Delta_r$, where $u=r^{12}$, with smooth
total space and reduced normal-crossings $r$-central divisor.
Its central surface over an ordinary value of $s_j$ is a cycle
of twelve ruled surfaces over the smooth elliptic curve
$E'_{j,s_j}$. Their normalizations are projective bundles
associated with the degree-zero line bundle represented by
$P_{s_j}-O_{s_j}$, with the appropriate pullbacks by translations.
Whenever $P_s-O_s$ is non-torsion, no finite isogeny of the
elliptic base makes these ruled normalizations into products.
\end{proposition}


\begin{proof}
\begin{enumerate}[label=\arabic*.,leftmargin=*]
\item On the root disc the elliptic family $E'_j\to\Delta_{s_j}$ is
smooth. In the additive period coordinates $(\xi,\zeta_2)$,
put $W=\exp(-2\pi i\zeta_2)$. Quotienting by
$\ker\gamma\subset\Lambda$ gives a holomorphic
$\mathbb C^*$-bundle over this elliptic family. The period
$\widehat\gamma$ acts by
\[
 (\xi,W)\longmapsto
 (\xi+6\mu(s),a(s)uW),\qquad
 a(s)=\exp(-2\pi i\beta^{\rm part}(s)).
\]
This is the bundle and translation in \cref{fam:src:B.1.1}, \cref{fam:src:B.1.4}, with inverse fibre
coordinate; the period identities check its transition
functions directly. The factor $a$ is a holomorphic unit on
the entire root disc. Apply the compact-core plumbing of
\ref{push:src:A.1.2} relatively over a closed smaller root disc. Here the
smoothing section never vanishes, so the only local equation
at the conductor is $xy=u$ after a change of frame. This
gives a smooth total space before the base change.

\item After $u=r^{12}$, the local equation is $xy=r^{12}$,
times the two smooth coordinates on $E'_j$. Resolve its
$A_{11}$ surface singularity by its minimal toric resolution.
In lattice coordinates this inserts each primitive height-one
ray between the two boundary rays, at unit intervals. The
construction is canonical for changes $x\mapsto a x$,
$y\mapsto a^{-1}y$ by a holomorphic unit, so these local
resolutions glue. Its affine charts have the equation
$r=XY$, or $r=X$ at a smooth point of the central divisor,
times the elliptic and root coordinates. Thus the total space
is smooth, the divisor is reduced with normal crossings, and
the joint map to $(s_j,r)$ is flat. The resolution is proper,
and the plumbing was proper, so the joint map is proper.
The comparison over $r\ne0$ is the original period quotient,
since the resolution is an isomorphism there. There are
twelve components in the cycle because one application of
$\widehat\gamma$ shifts the unfolded component index by
twelve.

\item The ruled normalization of a component is obtained by adding
the zero and infinity sections to the indicated
$\mathbb C^*$-bundle. For a degree-zero line bundle $L_0$
on an elliptic curve, $\mathbb P(\mathcal O\oplus L_0)$
is a product over that curve only if $L_0$ is trivial.
Indeed an isomorphism would give two independent nowhere
vanishing quotient sections of a common twist of
$\mathcal O\oplus L_0$; degree zero forces that twist and
$L_0$ to be trivial. Equivalently, the two distinguished
disjoint sections have normal bundles $L_0,L_0^{-1}$.
A finite isogeny can trivialize $L_0$ only when $L_0$ is
torsion, because the norm of its pullback is a positive
tensor power of $L_0$.

\item The section $P$ of the rational elliptic surface is not
torsion, as is checked by its positive Mordell--Weil height
in \cite[proof of Lemma~2.1, pp.~3--4]{Engel2026}. If its restriction to the root disc
were identically killed by a positive integer, analytic
continuation on the smooth part of $S\to B$ would give the
same identity globally. Each nonzero torsion equation
therefore has a discrete zero set. Outside a countable union
of such sets, $P_s-O_s$ is non-torsion. For each such $s$, the ruled normalization remains nontrivial
after every finite isogeny of the elliptic factor, by the norm
argument above. The torsion isogenies at $s_j=0$ are those
in \cref{fam:src:E.2.1}.
\end{enumerate}
\end{proof}

\begin{theorem}[Finite deck quotients and central cycles]\hypertarget{tag-0301}{}\label{push:src:B.1.2}\label{push:B12}
\StatementRecord{push:src:B.1.2}{proved}{0301}{\SeriesRef{push:src:B.1.1}, \SeriesRef{fam:src:E.2.1}, \SeriesRef{fam:src:A.5.1}}

The twisted finite deck generators extend to the models of
\ref{push:src:B.1.1}. On the twelve component labels their actions are
$i\mapsto i+4$ at $j=1$ and $i\mapsto i-3$ at $j=2$,
with labels in $\mathbb Z/12$. Every nonidentity group
element acts freely, including at every node of the central
cycle. The quotient is a smooth complex fourfold, proper
and flat over $(t,r)=(s_j^{m_j},r)$, with a reduced
normal-crossings $r$-central divisor. Its fibre divisor
over $t=p_j$ has multiplicity exactly $m_j$.

The reduction over $s_j=r=0$ is the central surface of \cref{fam:src:E.2.1}:
a cycle of four ruled surfaces at $j=1$, and a cycle of
three at $j=2$. In the root presentation the normal
character to $s_j=0$ is $e^{-2\pi i/m_j}$ and the
character on $r$ is trivial.
\end{theorem}


\begin{proof}
\begin{enumerate}[label=\arabic*.,leftmargin=*]
\item Write primes for
evaluation at $g_js$, put $\zeta_j=e^{-2\pi i/m_j}$,
and set
\[
 R_1=\begin{pmatrix}-1/\tau&0\\(1-\mu)/\tau&1\end{pmatrix},
 \qquad
 R_2=\begin{pmatrix}1/\tau&0\\-\mu/\tau&1\end{pmatrix}.
\]
The identities $R_j\Pi_c=\Pi'_c A_j$ and $A_jv_j=v_j$
follow by substitution of the displayed matrices. In the coordinates
$(s,\xi,W,r)$ the two affine deck maps become
\begin{align*}
 s'&=\zeta_1s,\qquad
 \xi'=-\xi/\tau+(6\mu'+2\tau'-4)/3,\\
 W'&=r^4 U_1(s,\xi)W,\\
 U_1&=\exp\!\left[-2\pi i
  \left(\frac{1-\mu}{\tau}\xi+
        \frac{\beta^{{\rm part}\prime}+2\mu'}3\right)\right],
 \end{align*}
and
\begin{align*}
 s'&=\zeta_2s,\qquad
 \xi'=\xi/\tau+(-6\mu'-3\tau'+3)/4,\\
 W'&=r^{-3} U_2(s,\xi)W,\\
 U_2&=\exp\!\left[-2\pi i
  \left(-\frac{\mu}{\tau}\xi-
        \frac{\beta^{{\rm part}\prime}+3\mu'}4\right)\right].
 \end{align*}
Each $U_j$ is a nowhere vanishing holomorphic function on
the additive elliptic cover. Period covariance gives the descent rule for $U_j$ and
therefore the required bundle isomorphism after quotient
by $\ker\gamma$. Multiplication by a unit preserves every toric chart.
Multiplication by $r^a$ sends the chart of component index
$i$ to the chart of index $i+a$. Therefore the displayed
maps and their inverses extend holomorphically to the
resolved infinite chain and then to its twelve-component
quotient. The identities $g_1^3=g_2^4=1$ on the quotient
hold on the dense punctured family, and hence everywhere.

\item A nonidentity power has component shift $4,8$ in the first
case and $-3,-6,-9$ in the second. None is zero modulo
twelve. The same shifts act on the edge labels of the
cycle, so no node can be fixed either. At $s\ne0$ the
rotation of $s$ already excludes fixed points. The finite
quotient map is consequently a holomorphic local
isomorphism. Its source is smooth, so its target is smooth.
Properness follows from the proper root model and the
finite base map. Flatness can be checked after the
surjective local-isomorphism cover: there $s^{m_j}$ and
$r$ are a regular sequence, and their fibres have pure
dimension two. The local central equation remains
$r=XY$ or $r=X$. Finally $t=s^{m_j}$ shows that the
Cartier divisor over $p_j$ is exactly $m_j$ times its
reduction, including where this divisor meets a central
node. This also proves the asserted normal characters.

\item At $s=0$ use the explicitly marked isogenies
\[
 L_1(\xi,\eta)=(\eta,\xi-\eta/3),\qquad
 L_2(\xi,\eta)=(\eta,\xi-\eta/2).
\]
Their kernel generators are $(2/3,\rho)$ and $(1/2,i)$,
and the finite deck maps lift to
\[
 (\xi,\eta)\mapsto(\xi+\kappa_1/3,(\rho-1)\eta),
 \qquad
 (\xi,\eta)\mapsto(\xi-\kappa_2/4,-i\eta).
\]
In the Tate coordinate $v=e^{-2\pi i\xi}$ these are the
shifts $4,-3$ just computed. The kernel generators act
by nonzero translations on the smooth elliptic factor
$F_j$, hence freely even at the two fixed points of each
Tate component. Every other element with nonzero deck
power has a nonzero component shift. The full product
groups are therefore free on the product semistable
models. Their component orbits have cardinalities three
and four, giving precisely four and three ruled
components. The original lattice isogenies intertwine
the maps on the dense tori and their unit-interval fans,
so the quotient identifications extend across all nodes.
These quotient identifications give the comparison with \cref{fam:src:E.2.1}.

\item For the punctured family obtained by descending the
$c$-family, the holomorphic correction $\eta_j(s)$ in
\cref{fam:src:A.5.1} causes no additional local obstruction here. At the
shift $c\mapsto c+12$, the reference translations
$-4\Pi_c(\widehat u+\widehat w)$ and
$3\Pi_c(\widehat u+\widehat w)$ are periods.
The remaining correction is equivariant and independent
of $c$. Translation by $-c\eta_j(s)/12$ on the additive
cover conjugates that descent to the period descent.
It commutes with the finite deck action by equivariance.
Thus the constructed finite models fill the local
pullback by $w=r^6$ after a specified holomorphic
identification over $r\ne0$. 
\end{enumerate}
\end{proof}

\subsection{The base-changed cusp fan}\hypertarget{tag-02Z6}{}\label{push:section:B.2}

\begin{proposition}[Period-invariant fans after base change]\hypertarget{tag-0302}{}\label{push:src:B.2.1}\label{push:B21}
\StatementRecord{push:src:B.2.1}{proved}{0302}{\SeriesRef{push:src:A.2.1}, \SeriesRef{push:src:A.1.2}, \SeriesRef{push:src:A.2.2}}

For the plain base change $u=r^n$ of the untwisted cusp
model, the smooth period-invariant refinements which
retain the prescribed bottom triangulation, are jointly
flat over $(q,r)$, and have reduced $r$-central divisor
are indexed by $k=0,\ldots,n$. In a fundamental prism
write its four bottom rays as $A,B,C,D$ as in \ref{push:src:A.2.1},
and put $\rho_i=(0,i,0,1)$ for $0\leq i\leq n$.
The maximal cones for the index $k$ are
\[
 ABC\rho_k,\quad BCD\rho_k,\quad
 AB\rho_i\rho_{i+1}\ (i<k),\quad
 CD\rho_i\rho_{i+1}\ (i\geq k).
\]
There are thirteen such models for $n=12$ and three
for $n=2$. Each has smooth total space and a reduced
normal-crossings central divisor with $n$ components
and no triple locus. Its component labels over the
punctured $q$-disc are globally $\mathbb Z/n$; in
particular their monodromy around $q=0$ is the identity.
\end{proposition}


\begin{proof}
\begin{enumerate}[label=\arabic*.,leftmargin=*]
\item The map of cocharacter lattices for the base change is
\[
 (y_1,y_2,h_q,h_r)\longmapsto
 (y_1,y_2,h_q,nh_r).
\]
The old top ray indexed by $l$ pulls back to
$(0,nl,0,1)$. Smoothness on the two-dimensional top
face forces insertion of every primitive ray
$\rho_{nl+i}$, $0<i<n$: a skipped interval of length
$d>1$ would retain lattice index $d$ on that face.
Joint flatness excludes any new ray with both base
heights positive, and the fixed bottom excludes new
rays of $h_r=0$. Reducedness permits only top rays
of $h_r=1$. These are consequently all the rays.

\item On the left face $AC\rho_0\rho_n$, a triangulation
with these collinear top vertices has a unique triangle
containing the bottom edge $AC$. Write its third
vertex as $\rho_a$. The triangles on either side
of that bridge are then forced. The analogous right
face has bridge $BD\rho_b$. The first period takes
$A,C$ to $B,D$ and fixes every top ray, so periodicity
requires $a=b=k$. The fixed bottom triangles $ABC$
and $BCD$ must then cone to $\rho_k$. The remaining
successive top intervals force exactly the displayed
cones, leaving no choice. Conversely those cones
have determinant $\pm1$ and common-face
intersections. Their normalized volumes sum to
$n+2$, the volume of the prism, so they cover it.
For $n=2$ and $n=12$, the determinant, intersection and
volume calculations use exact rational arithmetic.

\item The period actions are
\[
 G_1:(y_1,y_2)\mapsto(y_1+h_q,y_2),\qquad
 G_2:(y_1,y_2)\mapsto(y_1,y_2+h_q+nh_r).
\]
Thus the choices propagate consistently from the
fundamental prism. The proper toric subdivisions
descend by the covering construction of \ref{push:src:A.1.2}.
Every ray has pure $q$-height or pure $r$-height, so
in a regular affine chart $q$ and $r$ are monomials
in disjoint sets of coordinates. This proves joint
flatness. All nonzero $r$-heights equal one, and
there are at most two top rays in any maximal cone;
this proves reduced normal crossings and absence of
triple points. Properness follows by descent from
the proper toric modification of the base-changed
proper cusp model.

\item The central components and all their intersection
maps are specified by the stars of the $\rho_i$.
Modulo the periods their indices are $i\bmod n$;
adjacent indices meet in the conductor elliptic
surface. Away from the cusp their normalizations
are the ruled bundles of the expanded plumbing.
At the unique bridge class the star has both bottom
triangles, and gives the nodal-curve blow-up and
small resolution described in \ref{push:src:A.2.2}, on the
corresponding ruled bundle. This star description
also fixes the gluing when the bridge is on an
inserted component. The endpoint choices $k=0,n$
are the canonical resolutions of the base changes
of the two original models. Intermediate values
are related to them by the successive side flops.

\item Finally, $G_1$ fixes each index and $G_2$ changes it
by $n$. The units in the cusp periods are
single-valued holomorphic functions of $q$ and do
not change a ray label. The toric charts themselves
are defined over $q$, not a root of $q$. Hence
these labels remain global on the punctured
$q$-disc, proving the monodromy assertion. Further
smooth periodic subdivisions with mixed rays give
infinitely many models flat over $r$ alone, as in
\ref{push:src:A.2.1}; they are excluded by the stated joint-flatness
condition.
\end{enumerate}
\end{proof}

\begin{corollary}[Base-changed cusp refinements]\hypertarget{tag-0303}{}\label{push:src:B.2.2}\label{push:B22}
\hypertarget{tag-0304}{}\label{push:src:B.2.3}\label{push:B23}
\StatementRecord{push:src:B.2.3}{proved}{0304}{\SeriesRef{push:src:A.2.4}, \SeriesRef{push:src:A.2.5}, \SeriesRef{push:src:A.4.7}, \SeriesRef{push:src:A.4.8}, \SeriesRef{push:src:A.4.9}}
\StatementRecord{push:src:B.2.2}{proved}{0303}{\SeriesRef{push:src:A.2.4}, \SeriesRef{push:src:A.2.5}, \SeriesRef{push:src:A.4.7}, \SeriesRef{push:src:A.4.8}, \SeriesRef{push:src:A.4.9}}

For $n=2$ and $n=12$, \ref{push:src:A.2.4}--\ref{push:src:A.2.5} give all
finite-orbit smooth invariant plain-cusp refinements,
their proper flat families, normalization charts,
finite blow-up and blow-down presentations,
multiplicities and all multiple loci. \ref{push:src:A.4.7}--\ref{push:src:A.4.9}
give their complete integral topology and nearby
Euler function. Reducedness is exactly the restriction $\ell=1$ on every positive-height ray.
\end{corollary}

\begin{proof}
\begin{enumerate}[label=\arabic*.,leftmargin=*]
\item The base maps $u=w^2$ and $u=r^{12}$ induce multiplication by
$2$ and $12$ on the integral parameter characteristic. Insert
$n=2$ and $n=12$ respectively into the periodic lattice action
and fan construction of \ref{push:src:A.2.4}. Its converse
recovers the invariant subdivision from each such plain-cusp
model, so the parametrization includes every stated output.

\item For each primitive ray $v$, the local parameter is a monomial
whose exponent on its divisor is its integral height $\ell(v)$.
Thus the central divisor is reduced exactly when every such
height equals one. Smoothness, flatness, properness, the
normalization charts and all intersection maps are those proved
in \ref{push:src:A.2.4}; \ref{push:src:A.2.5} gives the finite
blow-up and blow-down presentation. Finally substitute the
resulting orbit counts and maps into
\ref{push:src:A.4.7}, \ref{push:src:A.4.8}, \ref{push:src:A.4.9} for the
claimed topology.
\end{enumerate}
\end{proof}

\Uses{the hypotheses, notation and definitions stated here; \cref{push:src:A.2.4}; \cref{push:src:A.2.5}; \cref{push:src:A.4.7}; \cref{push:src:A.4.9}}




\Uses{the hypotheses, notation and definitions stated here; \cref{push:src:A.2.4}}


\subsection{Component monodromy and trace obstructions}\hypertarget{tag-02Z7}{}\label{push:section:B.3}

\begin{theorem}[Peripheral component obstruction]\hypertarget{tag-0305}{}\label{push:src:B.3.1}\label{push:B31}
\StatementRecord{push:src:B.3.1}{proved}{0305}{\SeriesRef{push:section:B.1}, \SeriesRef{push:section:B.2}, \SeriesRef{push:src:B.1.2}, \SeriesRef{push:src:B.2.1}}

The finite models of \ref{push:section:B.1} and any of the plain
$u=r^{12}$ cusp models of \ref{push:section:B.2} cannot be glued to a
proper complex space over $B\times\Delta_r$ with
the prescribed central fibres on all three base
neighbourhoods. The obstruction is the peripheral component
monodromy of these prescribed finite and plain cusp models.
\end{theorem}


\begin{proof}
\begin{enumerate}[label=\arabic*.,leftmargin=*]
\item If a further plain cusp refinement changes the top fan,
its central fibre over an ordinary $q\ne0$ has additional
components or different multiplicities and already cannot
agree with the twelve-cycle model of \ref{push:section:B.1} on the overlap.
We may therefore assume its top fan is the standard reduced
twelve-cycle. Additional mixed-height rays lie over $q=0$
and do not change the following argument over the punctured
base.

\item Suppose such a gluing exists and put
$B^\circ=B\setminus\{0,1,\infty\}$. Normalize
its reduced central fibre over $B^\circ$. In
each of the prescribed plumbing charts this is
a proper smooth family whose fibres are the
disjoint union of twelve connected ruled
surfaces. The charts and their transition
isomorphisms therefore give a locally constant
twelve-element set of connected components.
Equivalently, the Stein factor of this
normalization is an unramified degree-twelve
cover of $B^\circ$: the algebra of fibrewise
holomorphic functions is locally $\mathcal O^{12}$,
since each proper connected ruled surface has
only constant holomorphic functions. This description constructs the component cover over
$B^\circ$.

\item The double curves give the adjacency graph of
a twelve-cycle. Thus its transition
permutations lie in the dihedral group
$D_{12}$, the automorphism group of that graph.
We record a peripheral permutation by its
deck-identification map. Holonomy defined by
lifting a clockwise path and returning its
endpoint to the initial root slice uses the
inverse permutations throughout; the relation
and the following contradiction are unchanged.
At $0$ the monodromy is a rotation by four
steps; at $1$ it is a rotation by minus three
steps, by the explicit root deck maps in
\ref{push:src:B.1.2}. Relabelling a cycle can reverse either
orientation, but changes these only to
$\pm4$ and $\pm3$. At the cusp the
monodromy is the identity by \ref{push:src:B.2.1}.

\item The peripheral relation in
$\pi_1(B^\circ)$ is $g_1g_2g_0=1$ with
the consistently chosen clockwise loops.
In any common labelling, the product of the
two finite rotations would therefore have
to be the identity. Its possible rotation
numbers are
\[
 4+3,\quad4-3,\quad-4+3,\quad-4-3
 \equiv7,1,11,5\pmod {12}.
\]
None is zero. This contradicts the peripheral
relation. 
\end{enumerate}
\end{proof}

\begin{proposition}[The marked genus-three component cover]\hypertarget{tag-0306}{}\label{push:src:B.3.2}\label{push:B32}
\StatementRecord{push:src:B.3.2}{proved}{0306}{\SeriesRef{push:src:B.3.1}}

Let $B^\circ=\mathbb P^1\setminus\{0,1,\infty\}$ and use the
clockwise deck-identification convention of \ref{push:src:B.3.1}.
The peripheral rotations $4,-3,-1$ in $\mathbb Z/12$ define a
connected degree-twelve cover of $B^\circ$. Its compact
normalization is the smooth genus-three curve obtained from
\[
 y^{12}=\frac{t^4}{(t-1)^3}.
\]
The ramification indices over $0,1,\infty$ are $3,4,12$,
respectively. The peripheral relation determines the last rotation
from the first two, so this is the marked cyclic component cover
used in the partial filling of \ref{push:src:B.4.19}.
\end{proposition}

\begin{proof}
\begin{enumerate}[label=\arabic*.,leftmargin=*]
\item The presentation $g_1g_2g_0=1$ sends the third peripheral generator
to $-4+3=-1$. The first two images generate $\mathbb Z/12$, so
the resulting action on the twelve sheets is transitive. This
constructs the connected cover by its monodromy representation.

\item The rational function $t^4/(t-1)^3$ has orders $4,-3,-1$ at the
three punctures. Its twelfth-root equation has the same marked
deck representation. The covering-space identification is the map
which sends a chosen twelfth root at one base point to its chosen
sheet and continues both along the same paths. Its equality of
monodromy makes that map single-valued; its local analytic root
coordinates make it holomorphic. Normalization extends this
identification to the compact covers.

\item Locally $y^{12}=z^a u(z)$ has $\gcd(12,a)$ branches of ramification
index $12/\gcd(12,a)$, after absorbing the holomorphic unit in a
local twelfth root. Thus the three contributions to ramification
are $4(3-1)$, $3(4-1)$, and $12-1$. Their sum is $28$.
Riemann--Hurwitz gives $2-2g=2\cdot12-28=-4$, so $g=3$.
Reversing the deck convention is realized by $y\mapsto y^{-1}$,
with the reciprocal defining equation; the compact curve and the
ramification calculation are unchanged.
\end{enumerate}
\end{proof}


\begin{proposition}[Cusp-root multiplicities]\hypertarget{tag-0307}{}\label{push:src:B.3.3}\label{push:B33}
\StatementRecord{push:src:B.3.3}{proved}{0307}{\SeriesRef{push:src:B.3.1}, \SeriesRef{push:src:A.1.2}}

Let $q=s^d$, $u=r^n$,
and suppose a cyclic deck generator has
translation $d^{-1}(aG_1+bG_2)$ in the
additive notation for the two periods.
Assume its $r$-component shift is integral.
Then its $r$-component shift is
$nb/d\bmod n$, whereas the multiplicity
of every generic quotient $q$-component
is a positive multiple of
\[
 \frac{d}{\gcd(d,a,b)}.
\]
A nonzero $r$-component shift therefore
forces a nonreduced $q$-fibre. In
particular a shift $-1\bmod12$ in this
class of constructions forces multiplicity
at least twelve.
\end{proposition}


\begin{proof}
\begin{enumerate}[label=\arabic*.,leftmargin=*]
\item In the coordinates of \ref{push:src:A.1.2} the two
periods have $q$-valuation vectors
$(1,0)$ and $(0,1)$; their $r$-valuation
vectors are $(0,0)$ and $(0,n)$.
After $q=s^d$, the period lattice on
the $q$-central component indices is
$d\mathbb Z^2$. A regular toroidal
resolution has indices in
$(\mathbb Z/d)^2$. The deck
generator translates them by $(a,b)$.
The order of that translation is
$d/\gcd(d,a,b)$, so the subgroup
preserving a component has order
$h=\gcd(d,a,b)$. Its subgroup fixing
a generic point has some order
$h_0$ dividing $h$; unit factors may
make $h_0$ smaller. That isotropy
acts on the normal coordinate $s$
by a primitive $h_0$th root of
unity. The quotient normal
coordinate is $s^{h_0}$, and
$q=(s^{h_0})^{d/h_0}$ has
multiplicity $d/h_0$, a positive
multiple of $d/h$. Unit factors
cannot decrease this lower bound.

\item Reducedness would require $h_0=d$,
and therefore $h=d$,
so both $a,b$ are divisible by $d$.
Then $nb/d$ is zero modulo $n$.
If $nb/d$ is relatively prime to
$n$, its order in $\mathbb Z/n$
is $n$, which divides the order
$d/h$ of the full component
translation; this gives the final lower bound.
\end{enumerate}
\end{proof}

\begin{theorem}[Trace on the split elliptic normalization]\hypertarget{tag-0308}{}\label{push:src:B.3.4}\label{push:B34}
\StatementRecord{push:src:B.3.4}{proved}{0308}{\SeriesRef{push:src:B.3.2}; \cite[Theorems 8.4 and 8.6, table (8.16)]{Shioda90}; \cite[Theorem 4.1]{Shioda90}; \cite[Theorem 1.3]{Shioda90}}

Let $C$ be the genus-three component cover
in \ref{push:src:B.3.2}, and let $\mathscr E\to C$ be
the relatively minimal elliptic surface
obtained by pulling back $S\to B$.
Write $\sigma$ for the order-twelve deck
transformation, with its natural lift to
$\mathscr E$. Then
\[
 \operatorname{Tr}_{\sigma}
   \mathscr E(\mathbb C(C))
 :=\left\{\sum_{i=0}^{11}\sigma^i Q:
                 Q\in\mathscr E(\mathbb C(C))\right\}
 =12\mathbb Z P.
\]
In particular there is no section $Q$
for which $(t_Q\circ\sigma)^{12}=t_{6P}$.
For a normalization over $C$ ruled over this split elliptic
surface, the twelfth iterate of an adjacent conductor map
$t_Q\circ\sigma$ is therefore $t_{12kP}$ for some $k\in\mathbb Z$.

\end{theorem}


\begin{proof}
\begin{enumerate}[label=\arabic*.,leftmargin=*]
\item There is a convenient model of the cover
and its elliptic surface:
\[
 C:\ b^3=a^4+1,\qquad
 t=b^3,\quad y=b/a,\qquad
 \mathscr E:\ Y^2=X^3-3bX+2a^2.
\]
Indeed $y^{12}=t^4/(t-1)^3$.
The substitution from the original
Weierstrass coordinates is
$x=a^2b^4X$, $y_{\rm old}=a^3b^6Y$.
The affine discriminant of the new
equation is $1728$; it has no singular
fibre on the affine curve. The curve
$C$ has one point at infinity, with
pole orders $3,4$ for $a,b$.
Put $z=a/b$ and $c=z^3a=z^4b$ there.
Then $z$ is a uniformizer and
\[
 c^3(1-c)=z^{12},\qquad c(0)=1.
\]
In the integral Weierstrass coordinates
$X_0=z^2X$, $Y_0=z^3Y$, the equation is
\[
 Y_0^2=X_0^3-3cX_0+2c^2.
\]
Its discriminant has order twelve and
its $c_4$ is a unit. The unique singular
fibre of the minimal regular surface
is therefore $I_{12}$. The fundamental
line bundle is $\mathcal O_C(\infty)$,
so $\chi(\mathcal O_{\mathscr E})=1$.
This is a projective elliptic surface.

\item Let $\omega=e^{2\pi i/3}$ and
$\lambda=e^{2\pi i/12}$. The deck lift is
\[
 \sigma(a,b,X,Y)=(ia,\omega b,\lambda^2X,\lambda^3Y).
\]
It fixes the section $P=(0,-\sqrt2a)$.
Two further sections are
\[
 R=\left(\frac{a^2}{b-1},
                 \frac{\sqrt3a}{b-1}\right),
 \qquad
 R'=\sigma R=
   \left(\frac{a^2}{b-\omega},
                 \frac{\sqrt3\omega a}{b-\omega}\right).
\]
Direct substitution verifies their
equations. The three orbit points
$R,\sigma R,\sigma^2R$ lie on the line
\[
 Y=\frac{\sqrt3b}{a}X-\sqrt3a.
\]
Thus $R+\sigma R+\sigma^2R=O$ by the
elliptic chord law, and $\sigma^3R=R$.


\item  The only pole of $R$ on the affine
curve is at $(a,b)=(0,1)$. There $a$ is
a uniformizer and $b-1=a^4/3+O(a^8)$,
so $X=3a^{-2}+O(a^2)$ and
$Y=3\sqrt3a^{-3}+O(a)$.
Consequently $R\cdot O=1$.
The same computation at $(0,\omega)$
gives $R'\cdot O=1$, while $P\cdot O=0$.
At the other two points with $a=0$,
$P$ and $R$ meet simply: the fibre is
smooth at $(X,Y)=(0,0)$, and the
difference of their $Y$ coordinates
has a nonzero linear term in $a$.
Thus $P\cdot R=2$, and likewise
$P\cdot R'=2$. The sections $R,R'$
meet simply at $(a,b)=(0,\omega^2)$.
They cannot meet elsewhere on the
affine curve, since their $X$
coordinates are distinct for $a\ne0$.

\item To check infinity, for the section
$\sigma^jR$ put $k=\omega^j$. Its
integral coordinates are
\[
 X_0=\frac{c^2}{c-kz^4},\qquad
 Y_0=\frac{\sqrt3kc z^4}{c-kz^4}.
\]
At the node use analytic coordinates
\[
 U=Y_0-(X_0-1)\sqrt{X_0+2},\qquad
 V=Y_0+(X_0-1)\sqrt{X_0+2}.
\]
They satisfy
\[
 UV=(1-c)\{3X_0-2(c+1)\},
\]
whose right side is a unit times $z^{12}$.
The section expansions are
\[
 U=-\frac{\sqrt3k^2}{6}z^8+O(z^{12}),
 \qquad V=2\sqrt3kz^4+O(z^8).
\]
Hence $R,R'$ both meet component $4$
of the resolved $I_{12}$ cycle, with
distinct coordinates on that component.
Reversing the cycle numbering replaces
both indices by $8$ and has no effect
on the argument. The section $P$ meets
the identity component, at the smooth
point $(X_0,Y_0)=(0,-\sqrt2)$.
Thus no further intersections occur at
infinity, and $R\cdot R'=1$.

\item The height pairing formula and its
positivity are those of
\cite[Theorems~8.4 and~8.6, table~(8.16), pp.~227--229]{Shioda90}.
For an $I_{12}$ fibre its correction is
$\operatorname{contr}(i,j)=\min(i,j)-ij/12$.
The preceding intersections give, in
the ordered basis $P,R,R'$,
\[
 G=\begin{pmatrix}
 2&0&0\\
 0&4/3&-2/3\\
 0&-2/3&4/3
 \end{pmatrix},\qquad \det G=8/3.
\]
In particular these are three independent
sections. The Leray sequence and
$R^1\pi_*\mathcal O_{\mathscr E}
=\mathcal O_C(-\infty)$ show that its
irregularity equals $g(C)=3$; the
right-hand sheaf has no global section.
This also agrees with the Picard-variety
statement in \cite[Theorem~4.1, pp.~219--220]{Shioda90}.
The equalities $\chi=1$ and
$e(\mathscr E)=12$ give $p_g=3$,
$b_2=22$ and $h^{1,1}=16$.
The zero section, a fibre and the eleven
nonidentity $I_{12}$ components generate
a rank-thirteen trivial lattice.
The identification of sections with
the N\'eron--Severi quotient
\cite[Theorem~1.3, p.~214]{Shioda90}, together
with $\rho\le h^{1,1}$, now shows that
the Mordell--Weil rank is exactly three
and $\rho=16$.

\item There is no nonzero torsion section.
Indeed its height would be zero, so
for some $0\le j<12$ the height formula
would give
\[
 0=2+2(Q\cdot O)-\frac{j(12-j)}{12}.
\]
The correction is at most three and
$Q\cdot O$ is a nonnegative integer.
This forces $Q\cdot O=0$ and
$j(12-j)=24$, which has no integral
solution.

\item Let $L=\mathbb ZP+\mathbb ZR+\mathbb ZR'$.
The integral N\'eron--Severi sublattice
generated by the trivial lattice and
these three sections has determinant
$-32$. This follows either by eliminating
the trivial lattice with the displayed
height matrix, or from the integral
intersection matrix specified by the following entries.
All its entries have just been computed:
sections have square $-1$, their mutual
intersections are $2,2,1$, and $R,R'$
meet precisely component $4$.
If $d=[\mathscr E(\mathbb C(C)):L]$,
the section--N\'eron--Severi quotient
identification gives the same index
for this sublattice in
$\operatorname{NS}(\mathscr E)$.
Its integral determinant implies
$d^2\mid32$, hence $d\in\{1,2,4\}$.
In particular $4Q\in L$ for every section.

\item In the rational basis $P,R,R'$ the
deck action and its trace are
\[
 M=\begin{pmatrix}1&0&0\\0&0&-1\\0&1&-1\end{pmatrix},
 \qquad \sum_{i=0}^{11}M^i=
       \operatorname{diag}(12,0,0).
\]
Suppose first that the trace of $Q$ is
$6P$. By the index bound we may write
\[
 Q=\frac12P+\frac14(uR+vR'),\qquad u,v\in\mathbb Z.
\]
Let $j\in\mathbb Z/12$ be its component
index at infinity. Since both $R,R'$
meet component four, the height formula
modulo integers gives
\[
 \frac{2u-v}{6}\equiv\frac j3\pmod{\mathbb Z},
 \qquad
 \frac{-u+2v}{6}\equiv\frac j3\pmod{\mathbb Z}.
\]
Equivalently $2u-v-2j$ and $-u+2v-2j$
are divisible by six. Reducing modulo
two shows that both $u,v$ are even.
Write $u=2u'$, $v=2v'$. The component
map is a group homomorphism, so
\[
 2j\equiv4(u'+v')\pmod{12}.
\]
Therefore $j$ is even. The geometric
height formula then lies in
$\tfrac13\mathbb Z$. In contrast the
height from $G$ is
\[
 \langle Q,Q\rangle
 =\frac12+\frac{u'^2-u'v'+v'^2}{3},
\]
which lies in the disjoint coset
$\tfrac12+\tfrac13\mathbb Z$.
This is a contradiction.

\item The same argument determines the full
trace image. If $Q=\alpha P+\beta R+\gamma R'$
in the rational span, pairing with the
narrow section $P$ gives
$2\alpha=\langle Q,P\rangle\in\mathbb Z$.
Every odd half-integral value of $\alpha$
leads to the identical congruences and
height contradiction, with $1/2$
replaced by an odd square divided by
two. Thus $\alpha\in\mathbb Z$ and
the trace is in $12\mathbb ZP$.
The section $P$ realizes the generator.

Finally the identity
$(t_Q\circ\sigma)^{12}=t_{\operatorname{Tr}_\sigma Q}$
proves the claimed obstruction.
\end{enumerate}
\end{proof}

\begin{corollary}[Integral Mordell--Weil lattice]\hypertarget{tag-0309}{}\label{push:src:B.3.5}\label{push:B35}
\StatementRecord{push:src:B.3.5}{proved}{0309}{\SeriesRef{push:src:B.3.4}}

In the notation of \ref{push:src:B.3.4},
\[
 \mathscr E(\mathbb C(C))
   =\mathbb ZP\oplus\mathbb ZR\oplus\mathbb ZR',
 \qquad \det\operatorname{NS}(\mathscr E)=-32.
\]
\end{corollary}


\begin{proof}
\begin{enumerate}[label=\arabic*.,leftmargin=*]
\item The affine coordinate ring of $C$ is
$\mathbb C[a,b]/(b^3-a^4-1)$, with
unique representatives of degree at most
two in $b$. The pole orders $3i+4j$ of
$a^ib^j$, $0\le j\le2$, are distinct,
so cancellation cannot create functions
with pole order one or two at infinity.
Consequently
\[
 H^0(C,\mathcal O_C(2\infty))=\mathbb C,
 \qquad
 H^0(C,\mathcal O_C(3\infty))
       =\mathbb C\oplus\mathbb Ca.
\]
A section disjoint from $O$ therefore
has $X=x_0$, $Y=y_0+y_1a$ for constants
$x_0,y_0,y_1$. Substitution in the
Weierstrass equation forces $x_0=y_0=0$
and $y_1^2=2$. Thus the only nonzero
sections with intersection zero with
$O$ are $P,-P$.
Every other nonzero section has
intersection at least one with $O$,
and hence height at least $4-3=1$.

\item \ref{push:src:B.3.4} proves that every section has
integral coefficient along $P$ in the
rational height decomposition. The
remaining lattice can have index only
one, two or four over $\mathbb ZR+\mathbb ZR'$.
Any proper such overlattice contains a
nonzero element of order two modulo
this sublattice. After subtracting an
integral combination of $R,R'$ it is
one of $R/2$, $R'/2$, $(R+R')/2$.
Each has height $1/3$, contradicting
the preceding lower bound. Thus the
index is one. The determinant already
computed in \ref{push:src:B.3.4} is the full
N\'eron--Severi determinant.

\end{enumerate}
\end{proof}

\begin{proposition}[Unramified torsors with prescribed transgression]\hypertarget{tag-030A}{}\label{push:src:B.3.6}\label{push:B36}
\StatementRecord{push:src:B.3.6}{proved}{030A}{\SeriesRef{push:src:B.3.4}, \SeriesRef{push:src:B.4.1}; \cite[Chapter III, Section 4, Theorem 2
and the following remark, pp.129--130]{HS53}; \cite[Lemma 4.7.8]{CDL24}; \cite{CDL24}; \cite[Section 4, equation (4)]{Xiao91}; \cite[Sections 3--4]{Xiao91}}

Put $K=\mathbb C(t)$, $L=\mathbb C(C)$,
$G=\operatorname{Gal}(L/K)=\langle\sigma\rangle$,
and let $E/K$ be the generic elliptic
curve of $S\to B$. There is a nonzero
class
\[
 \xi\in\operatorname{Sha}(E_L/L)^G
 \quad\hbox{whose transgression in}\quad
 H^2(G,E(L))=\mathbb ZP/12\mathbb ZP
 \quad\hbox{is }6P.
\]
Here $\operatorname{Sha}$ denotes the
subgroup of genus-one torsors trivial
over the henselian local fields of all
points of $C$.
Such a class has a relatively
minimal projective genus-one surface
$V\to C$, locally isomorphic to
$\mathscr E\to C$, and an affine lift
$A:V\to V$ of $\sigma$ with
\[
 A^{12}=t_{6P}.
\]
Thus $(V,A)$ is a genus-one normalization base with the
specified affine twelfth power.
\end{proposition}


\begin{proof}
\begin{enumerate}[label=\arabic*.,leftmargin=*]
\item We use Galois cohomology with its
continuous cochains. The five-term
Hochschild--Serre sequence gives
\[
 H^1(K,E)\longrightarrow H^1(L,E)^G
 \xrightarrow{\ d_2\ }H^2(G,E(L))
 \longrightarrow H^2(K,E).
\]
The transgression sequence is
\cite[Chapter III, Section 4, Theorem 2
and the following remark, pp.129--130]{HS53};
its continuous form follows by passing
to finite Galois quotients for the
finitely many values of each continuous
cochain. For this complex function field,
$H^2(K,E)=0$, as in
\cite[Lemma~4.7.8, p.~445]{CDL24}.
One can also see this from
$\operatorname{cd}K=1$: the torsion
submodule of $E(\overline K)$ has no
second cohomology, and its quotient is
a rational vector space, whose positive
continuous cohomology vanishes by
averaging on finite quotients.

\item Cyclic group cohomology gives
$H^2(G,E(L))=E(K)/\operatorname{Tr}_\sigma E(L)$.
By \ref{push:src:B.3.4} this is $\mathbb Z/12$, with
generator the class of $P$. Hence there
is a $G$-invariant torsor class
$\xi_0\in H^1(L,E)$ with $d_2\xi_0=6P$.


\item At a place $v$ of $B$ choose a point
$w$ of $C$ above it and its decomposition
subgroup $G_w$. The only nontrivial
decomposition groups have orders
$3,4,12$, at $0,1,\infty$ respectively.
The component groups of the original
fibres are $\mathbb Z/3$, $\mathbb Z/2$
and zero. Thus $6P$ belongs to the
identity component of the local N\'eron
group at every ramified place.
Multiplication by the ramification
index is surjective on that identity
group over the henselian local field:
on the residue group it is surjective
on $\mathbb G_a(\mathbb C)$ or
$\mathbb G_m(\mathbb C)$, and its
derivative is invertible; Hensel's
lemma lifts a division point.
Choose such a division point over
$K_v^h$. Viewed over $L_w^h$, it is
invariant, so its local norm is $6P$.
The restriction of the transgression
class to $H^2(G_w,E(L_w^h))$ is
therefore zero. For unramified places
the decomposition group is trivial,
so the same assertion is automatic.

\item Naturality and exactness of the local
five-term sequences now show that
$(\xi_0)_w$ is the restriction of some
$\eta_v\in H^1(K_v^h,E)$.
Only finitely many of these classes
need be nonzero. The original elliptic
surface is nontrivial, since
$j=1728t$. Ogg--Shafarevich localization
is therefore surjective:
\[
 H^1(K,E)\longrightarrow
 \bigoplus_{v\in B}H^1(K_v^h,E)
 \longrightarrow0.
\]
This is the combination of exact
sequence~(4.7.4) and Theorem~4.7.9
in \cite[p.~445]{CDL24}; the same surjectivity
and its geometric fibre-twist
interpretation are given in
\cite[Section~4, equation~(4), p.~21]{Xiao91}.
Take a global $\eta$ with these local
values and put
$\xi=\xi_0-\operatorname{res}_{L/K}\eta$.
It is locally trivial at every point
of $C$, remains $G$-invariant, and has
$d_2\xi=6P$. It is nonzero because its
transgression is nonzero.

\item The standard interpretation of the
Weil--Ch\^atelet group represents
$\xi$ by a smooth projective genus-one
curve over $L$; taking its relatively
minimal regular model over $C$ gives
$V$. Local triviality identifies each
henselian model with that of $E_L$.
Over $\mathbb C$ the corresponding
\'etale identifications yield
holomorphic identifications over small
base discs. Equivalently, $V$ is obtained
by gluing copies of $\mathscr E$ over
these discs by translations by local
sections. These translations extend
over the regular $I_{12}$ model.
Thus $V$ is locally isomorphic to
$\mathscr E$, has smooth total surface,
no multiple fibres, and the same single
$I_{12}$ fibre. Properness and flatness
hold locally over the base and hence
globally. The local translation gluing specifies $V$ as in
\cite[Theorem~3.4, pp.~14--15; Section~4, p.~21]{Xiao91}.

\item The $G$-invariance of $\xi$ gives an
affine isomorphism of its generic fibre
covering $\sigma$. It extends to the
relatively minimal regular surface:
in this particular situation the
extension can be checked on a common
resolution. If an exceptional
$(-1)$-curve for one projection is not
contracted by the other, its image is
a vertical rational curve and has
self-intersection at least $-1$,
since blowing down only increases
self-intersection. On either locally
trivial model all vertical rational
curves are components of the $I_{12}$
fibre, of square $-2$. Thus every such
$(-1)$-curve is contracted by both
projections and can be removed from
the common resolution. Repeating this
proves that the map is an isomorphism.
The same argument gives the local
comparison with the Jacobian model
used above. Transporting the base by
$\sigma$ changes none of these facts.
Write $A_0$ for this lift. Its twelfth
power is translation by a section
$N\in E(L)^G=\mathbb ZP$.
Choosing the powers $1,A_0,\ldots,A_0^{11}$
as lifts of $G$ gives the cyclic
factor set
$N\lfloor(i+j)/12\rfloor$.
Its cohomology class is precisely the
descent obstruction $d_2\xi$, so
$N=(6+12m)P$ for an integer $m$.
Replacing $A_0$ by $A=t_{-mP}\circ A_0$
gives $A^{12}=t_{6P}$.


\item To form the threefold filling of \ref{push:src:B.4.1},
the datum $(V,A)$ must be supplemented by its normal line
bundles, cusp charts and specified punctured-family
identification.
\end{enumerate}
\end{proof}

\begin{proposition}[Auxiliary torsor plumbing with cusp multiplicity twelve]\hypertarget{tag-030B}{}\label{push:src:B.3.7}\label{push:B37}
\StatementRecord{push:src:B.3.7}{proved}{030B}{\SeriesRef{push:src:B.3.6}, \SeriesRef{push:src:A.1.2}, \SeriesRef{push:src:A.2.2}; \cite[Proposition 0.9.2 and the opening
of Section 4.2, p.370]{CDL24}; \cite[Proposition 8.12, equation (8.25)]{Shioda90}}

For the torsor and affine lift $(V,A)$
of \ref{push:src:B.3.6}, there is a line bundle $M$ on
$V$ representing the narrow Jacobian
section $P$, with
\[
 M^2=-2,\qquad
 M\otimes(A^*M)^{-1}
       \simeq\pi^*\mathcal O_C(\infty).
\]
These data give a proper flat nodal
plumbing $\mathcal Z^\flat\to\Delta_h$
and two uniform smooth-total-space
semistable modifications
$\mathcal Z^\pm\to\Delta_h$, together
with holomorphic maps to $B$.
For $h\ne0$, the nodal fibre has
twelve ordinary double points and
Euler number twelve; each small-resolved
fibre has Euler number twenty-four.
The cusp fibre of the latter map to
$B$ has multiplicity twelve.
Consequently this auxiliary construction
does not fill the prescribed $S^6$
family. No proper birational modification
over $B$ can turn its multiple cusp
into a reduced cusp.
\end{proposition}


\begin{proof}
\begin{enumerate}[label=\arabic*.,leftmargin=*]
\item The section $P$ determines a relative
Picard class of multidegree zero on
every fibre component of $V\to C$.
It lifts to a line bundle $M$ because
$\operatorname{Br}(C)=0$; this is the
relative Picard consequence of
\cite[Proposition~0.9.2, p.~143; Section~4.2, p.~370]{CDL24}.
In local trivializations of the torsor,
translation acts trivially on the
relative degree-zero Picard class.
Thus the local Poincar\'e lines glue
up to base lines, with precisely that
Brauer obstruction. In particular we
may and do retain multidegree zero
on all twelve fibre components.

\item Choose a finite multisection cover of
degree $d$ which trivializes the torsor
generically, and resolve its base change.
There the pulled-back $M$ becomes the
narrow Poincar\'e line for the pulled-back
$P$, up to a base line. Its square is
$-2d$, by height base change
\cite[Proposition~8.12, equation~(8.25), p.~231]{Shioda90}.
Base lines do not change the square
of a fibre-degree-zero line. The
intersection pullback formula therefore
gives $M^2=-2$.

\item Since $\sigma P=P$, the line
$J=M\otimes(A^*M)^{-1}$ is trivial
on every fibre, including the
$I_{12}$ fibre. Grauert base change
and the evaluation map give
$J=\pi^*J_C$ for a line on $C$.
The line $M\otimes(t_{6P}^*M)^{-1}$
has base degree twelve. To check the
sign and degree, pass to the same
trivializing cover. On a Jacobian,
the degree of the restriction of
$\mathcal O(P-O)$ to a section $Q$
minus its degree on $O$ is
$-\langle P,Q\rangle$, by the height
formula and the narrowness of $P$.
For $Q=6P$ this is $-12d$.
Descending the degree factor $d$
proves the assertion. Iterating $A^*$
twelve times now gives
$12\deg J_C=12$, so $\deg J_C=1$.

\item The forms $da/b^2$, $a\,da/b^2$ and
$da/b$ form a basis of $H^0(C,K_C)$.
Their $\sigma$-eigenvalues are
$\lambda^7,\lambda^{10},\lambda^{11}$.
Thus the characteristic polynomial
on $H^1(C,\mathbb Z)$ is
\[
 p(x)=\Phi_{12}(x)\Phi_6(x)
     =x^6-x^5+x^3-x+1,
 \qquad
 p(x)+(1-x)(x^5+x^2+x)=1.
\]
It follows that $1-\sigma^*$ is an
automorphism of $\operatorname{Pic}^0(C)$,
with inverse $(\sigma^*)^5+(\sigma^*)^2+\sigma^*$.
Twisting $M$ by a base line changes
$J_C$ by $N\otimes(\sigma^*N)^{-1}$.
The displayed inverse therefore supplies
a base twist making
$J_C=\mathcal O_C(\infty)$.
Let $\phi$ be its section with divisor
the reduced $I_{12}$ fibre $F_\infty$.

\item Apply the plumbing maps of \ref{push:src:A.1.2}
to the data $(V,M,A)$.
Set $R_V=\mathbb P_V(\mathcal O\oplus M)$,
and identify its zero and infinity
sections by $s_0(x)\sim s_\infty(Ax)$.
On the node neighbourhood over $V$ use
\[
 a\in M_x,\quad b\in M_{Ax}^{-1},
 \qquad a\otimes b=h\phi_x.
\]
The $a$ collar is glued to the annular
core at $(x,a)$; the $b$ collar is glued
at $(Ax,b^{-1})$. Choose the Hermitian
metric and collar radii as in \ref{push:src:A.1.2}.
Compactness of $V$ and the bounded norm
of $\phi$ give disjoint collars,
Hausdorffness and compact representatives
by exactly the same inequalities.
These are holomorphic gluing maps of
fourfolds over $\Delta_h$.
Their local hypersurface rings are flat
over $h$, and the compact representatives
prove properness. The map to $B$ is
well defined because $t\circ\sigma=t$.

\item At a smooth point of $F_\infty$ the
equation is $ab=hz_C$; at one of its
twelve nodes, $z_C=uv$ and it is
$ab=huv$. For $h\ne0$ the latter
has one ordinary double point.
These twelve centres are distinct
even if $A$ rotates the fibre
components. A centre has $a=b=0$
and is in neither annular collar;
no overlap identifies it with the
centre indexed by $Ax$.
Thus the double locus is parametrized by $V$.
The two uniform prism subdivisions
of \ref{push:src:A.2.2} extend over these charts:
the affine action preserves the
orientation of the $I_{12}$ cycle,
and all remaining changes are units.
They give proper small modifications,
smooth total spaces and reduced
normal-crossings $h$-central fibres.
Equivalently their normalizations are
the blow-up of $R_V$ along $F_\infty$
in one section, followed by the twelve
chosen small resolutions. The normal
product becomes
$\pi^*\mathcal O_C(\infty)(-F_\infty)=\mathcal O_V$.

\item The nearby function of the nodal
plumbing is one on $F_\infty$ and
zero on its complement in the conductor.
Since $e(F_\infty)=12$, its integral
is twelve. Each small resolution
replaces a point by $\mathbb P^1$
and adds one, giving twenty-four.
Independently $e(R_V)=24$; above each
of the twelve nodes the modified
normalization replaces a point by
a chain of two projective lines,
of Euler number three. The open
parts of the fibre are copies of
$\mathbb C^*$ and contribute zero.
The modified normalization therefore
has Euler number $48$, its pinching
has Euler number $36$, and the
semistable nearby count is
$48-2e(V)=24$.

\item Finally $z=a/b$ is a coordinate on
$C$ near infinity and
$q_B=1/t=z^{12}/c^3$, with $c$ a unit.
At a generic point of the auxiliary
cusp fibre this gives multiplicity
twelve. In fact the map $q_B$ is
locally a twelfth power throughout
the smooth nearby threefold; a
holomorphic root of the unit can
be chosen locally. On any common
resolution of a proper birational
comparison over $B$, every divisor
over that fibre consequently has
$q_B$-valuation divisible by twelve.
A reduced fibre on another smooth
model would have a strict transform
of valuation one, a contradiction.

\end{enumerate}
\end{proof}

\begin{proposition}[Hodge theory of analytic unramified torsors]\hypertarget{tag-030C}{}\label{push:src:B.3.8}\label{push:B38}
\StatementRecord{push:src:B.3.8}{proved}{030C}{\SeriesRef{push:src:B.3.4}, \SeriesRef{push:src:B.4.5}, \SeriesRef{push:src:B.3.6}; \cite[Theorem 11, p.297]{Buchdahl1999}; \cite[Section 2.1, p.8]{Chen2026HodgeRoot}}

Let $V\to C$ be any compact smooth analytic genus-one surface
locally isomorphic, as a torsor over every base disc, to the
Jacobian surface $\mathscr E\to C$ of \ref{push:src:B.3.4}.
Then $b_1(V)=6$, and $V$ is K\"ahler. For every holomorphic
line bundle $M$ on $V$, the projective bundle
$\mathbb P_V(\mathcal O\oplus M)$ is K\"ahler. Every smooth
compact threefold bimeromorphic to such a bundle satisfies
the $\partial\bar\partial$-lemma.

Consequently none of these projective bundles or smooth compact
bimeromorphic models can be the non-$\partial\bar\partial$
normalization component required in \ref{push:src:B.4.5}.
\end{proposition}


\begin{proof}
\begin{enumerate}[label=\arabic*.,leftmargin=*]
\item Write $p:\mathscr E\to C$ for the projective Jacobian.
The calculation in \ref{push:src:B.3.4} gives
$R^1p_*\mathcal O_{\mathscr E}=\mathcal O_C(-\infty)$,
and $g(C)=3$. The coherent Leray sequence therefore gives
$h^1(\mathscr E,\mathcal O_{\mathscr E})=3$. Projectivity
and Hodge decomposition imply $b_1(\mathscr E)=6$.
The topological Leray sequence with rational coefficients is
\[
 0\longrightarrow H^1(C;\mathbb Q)
 \longrightarrow H^1(\mathscr E;\mathbb Q)
 \longrightarrow H^0(C,R^1p_*\mathbb Q)
 \longrightarrow H^2(C;\mathbb Q)
 \longrightarrow H^2(\mathscr E;\mathbb Q).
\]
The last arrow is injective because $p$ has a section.
The first arrow is an isomorphism of six-dimensional
vector spaces. Exactness thus implies
$H^0(C,R^1p_*\mathbb Q)=0$.

\item The local torsor identifications differ by translations.
On a smooth elliptic fibre translations act trivially on
$H^1$. On the $I_{12}$ fibre they preserve the orientation
of the component polygon and therefore act trivially on
its one-dimensional $H^1$ as well. These identifications
commute with the local specialization maps. Hence, for
$\pi:V\to C$, they glue to
$R^1\pi_*\mathbb Q\simeq R^1p_*\mathbb Q$ on the whole
curve, including infinity. The same Leray edge sequence
for $V$ now gives
$H^1(V;\mathbb Q)\simeq H^1(C;\mathbb Q)$.
Thus $b_1(V)=6$ is even. Buchdahl's theorem
\cite[Theorem 11, p.297]{Buchdahl1999} proves that $V$
is K\"ahler.

\item Choose a Hermitian metric on $\mathcal O\oplus M$.
The induced metric on the tautological relatively ample
line of its projectivization has a closed real $(1,1)$
curvature form $\theta$ positive in the vertical direction.
If $\omega$ is a K\"ahler form on $V$, compactness implies
that $\theta+A\pi^*\omega$ is positive in every direction
for a sufficiently large constant $A$. It is closed, so
it is a K\"ahler form on the projective bundle. A smooth
compact bimeromorphic model is therefore in Fujiki class
$\mathcal C$ and satisfies the $\partial\bar\partial$-lemma;
the domination statement is recalled, with its precise
reference to the $\partial\bar\partial$ theorem, in
\cite[Section 2.1, p.8]{Chen2026HodgeRoot}. Apply \ref{push:src:B.4.5}
for the final obstruction.
\end{enumerate}
\end{proof}

\subsection{Normal-crossing limits and coherent cohomology}

\begin{remark}[The two logarithms in the finite connector]
\hypertarget{tag-PT09}{}\label{I:finite-connector-two-logarithms}
\StatementRecord{I:finite-connector-two-logarithms}{proved}{PT09}{\cref{fam:src:A.5.1,push:src:B.3.1,night:P2.7}}
The finite-chart comparison of \cref{fam:src:A.3.1} uses the connector
\[
 \frac{\log s_j}{2\pi i}\Pi_c v_j,
 \qquad c=-\frac{12}{2\pi i}\log r.
\]
The identity $\Pi_c v_j=\Pi_0v_j+c\gamma(v_j)e$ separates its vertical term as
\[
 \frac{\log s_j}{2\pi i}\Pi_0v_j
 -\frac{12\gamma(v_j)}{(2\pi i)^2}\log s_j\log r\,e.
\]
This is the exact overlap map on the punctured parameter charts. Proposition~\ref{push:src:B.3.1} computes its component transport, and Theorem~\ref{night:P2.7} gives the resulting obstruction for a proper filling over the full parameter disc.
\end{remark}

\begin{theorem}[Fundamental-group obstruction]\hypertarget{tag-030O}{}\label{push:src:B.4.2}\label{push:B42}
\StatementRecord{push:src:B.4.2}{proved}{030O}{\cite[Theorem 5.1, p.2226]{NO10}}

Let $f:\mathcal Y\to\Delta$ be a
proper holomorphic map from a smooth
complex manifold, with smooth connected
fibres over $\Delta^*$ and reduced
normal-crossings central fibre $D$.
There is a continuous specialization
map from a nearby fibre $Y_t$ to $D$
which induces a surjection
\[
 \pi_1(Y_t)\longrightarrow\pi_1(D).
\]
Hence the central fibre of any such
semistable filling of $S^6$ is simply
connected. In particular it cannot
be a simple pinching of a connected
normalization along one pair of
disjoint nonempty conductor sections,
whether its elliptic normalization
base is split or is a genus-one torsor.
\end{theorem}


\begin{proof}
\begin{enumerate}[label=\arabic*.,leftmargin=*]
\item The surjectivity assertion also follows from
\cite[Theorem~5.1, p.~2226]{NO10}. We use the explicit neighbourhood
retraction needed in \cref{push:src:A.4.7,push:src:B.4.7}. Put
$W=f^{-1}(\overline\Delta_\epsilon)$
and $W^*=W\setminus D$. The negative
gradient of $|f|^2$, for a Hermitian
metric, retracts $W$ onto $D$.
In a product chart
$f=z_1\cdots z_k$ the arithmetic--
geometric mean inequality gives
\[
 \|\nabla |f|^2\|
   \ge c\,|f|^{2-1/k}.
\]
A finite cover of compact $D$ gives
a common exponent $\alpha<1$ in
$\|\nabla\rho\|\ge c\rho^\alpha$,
where $\rho=|f|^2$; for a fourfold
one may take $\alpha=7/8$.
The remaining trajectory length is
bounded by
$c^{-1}\rho^{1-\alpha}/(1-\alpha)$.
Hence trajectories converge uniformly
to $D$, and parametrization by the
level of $\rho$ gives a continuous
strong deformation retraction fixing
$D$. General position moves every
loop of $W$ off the real-codimension-two
normal-crossings strata, so
$\pi_1(W^*)\to\pi_1(W)$ is onto.
The smooth proper fibration over the
punctured disc gives
\[
 \pi_1(Y_t)\longrightarrow\pi_1(W^*)
       \longrightarrow\mathbb Z\longrightarrow0.
\]
A small transverse meridian at a
smooth point of a reduced component
maps to the generator of $\mathbb Z$
and bounds a disc in $W$.
Thus the image of $\pi_1(Y_t)$ alone
generates $\pi_1(W)=\pi_1(D)$. Restricting the retraction
$W\to D$ to $Y_t$ gives the asserted continuous specialization
and its surjection on fundamental groups.

\item For the last assertion, let $R$ be
the connected normalization and
$D_0,D_\infty$ its two disjoint
conductor sections. Choose a
continuous function $h:R\to[0,1]$
equal to zero on $D_0$ and one on
$D_\infty$. The function
$\exp(2\pi i h)$ descends to the
pinching. A path from a point of
$D_0$ to its paired point in
$D_\infty$ becomes a loop whose
image has winding number one.
Consequently the pinching has a
surjection of its fundamental group
onto $\mathbb Z$. This contradicts
the preceding necessary condition
when $Y_t\simeq S^6$.
\end{enumerate}
\end{proof}

\subsection{A Hodge obstruction to the central strata}

\begin{theorem}[The Hodge obstruction to normal-crossing limits]\hypertarget{tag-030Z}{}\label{push:src:B.4.3}
\StatementRecord{push:src:B.4.3}{proved}{030Z}{\SeriesRef{push:src:B.4.33}; \citet[Theorem 1.4, preprint p.~7, and Example 2.1, p.~13]{Nishiguchi1988IIRoot}; \citet[Theorem 7.1]{IstratiEtAl2022Root}}

Let $g:\mathcal Z\to\Delta_s$ be a proper holomorphic map from a
smooth complex fourfold, smooth over $\Delta_s^*$, with reduced
normal-crossings central fibre $D$. Suppose its nearby fibres are
biholomorphic to the six-sphere members of the specified family,
after any finite base change $w=s^N$. Then
\[
 H^1(D;\C)=0,\qquad h^1(D,\OO_D)\geq1.
\]
Consequently the natural map
$H^1(D;\C)\to H^1(D,\OO_D)$ is not surjective. In particular,
$D$ cannot be the analytification of a proper complex algebraic
variety with normal-crossings singularities.

If $D$ has \emph{simple} normal crossings, at least one of its
nonempty closed strata fails the $\partial\bar\partial$-lemma.
In particular, no such simple normal-crossings filling can have
all irreducible components K\"ahler or all irreducible components
Moishezon. These obstructions hold for every $N\geq1$.
\end{theorem}

\begin{proof}
Theorem~\ref{push:src:B.4.2} gives a surjection from the fundamental group of a nearby fibre onto $\pi_1(D)$. Since that fibre is diffeomorphic to $S^6$, it follows that $\pi_1(D)=1$ and $H^1(D;\mathbb C)=0$.

For a compact complex threefold $X$ diffeomorphic to $S^6$, the groups $H^2(X;\mathbb Z)$ and $H^4(X;\mathbb Z)$ vanish, so $c_1(X)=c_2(X)=0$. The analytic Riemann--Roch formula, obtained from \citet[Theorem~4.3, pp.~564--565]{AtiyahSinger1968IIIE}, gives
\[
 \chi(X,\mathcal O_X)=\frac1{24}\langle c_1(X)c_2(X),[X]\rangle=0.
\]
Every holomorphic three-form $\Omega$ is closed and is exact because $H^3(X;\mathbb C)=0$. Stokes' theorem gives $\int_X i\Omega\wedge\overline\Omega=0$. The integrand is pointwise nonnegative for the complex orientation, so $\Omega=0$. Serre duality therefore gives $h^3(X,\mathcal O_X)=0$, and connectedness gives $h^0(X,\mathcal O_X)=1$. Hence
\[
 h^1(X,\mathcal O_X)=1+h^2(X,\mathcal O_X)\geq1.
\]
Locally along the reduced normal-crossing divisor, $g$ has the form $s=z_1\cdots z_k$ after absorbing a holomorphic unit. This makes $\mathcal O_{\mathcal Z}$ torsion-free over the local discrete valuation ring of the disc and hence flat over it. Upper semicontinuity for the proper flat map gives $h^1(D,\mathcal O_D)\geq1$; see \citet[\S7, Satz~5, pp.~63--64]{Grauert60}.

For a proper complex algebraic variety with Du Bois singularities, the canonical map from singular cohomology to coherent cohomology is surjective in every degree; see \citet[\S1.D, p.~120]{GrafKovacs2014Root}. Normal-crossing singularities are Du Bois, including self-intersections of a global component. Thus algebraicity and properness of $D$ would make $H^1(D;\mathbb C)\to H^1(D,\mathcal O_D)$ surjective, contradicting the two groups just computed.

Now suppose $D$ is simple normal crossing and each nonempty closed stratum satisfies the $\partial\bar\partial$-lemma. These are exactly the strata in condition (K1) of \citet[Definition~2.3, p.~9]{Chen2026HodgeRoot}. For a proper semistable degeneration with smooth total space and this condition, \citet[Theorem~2.4(4), p.~10]{Chen2026HodgeRoot} gives $E_1$ degeneration of the Hodge--de Rham spectral sequence on a nearby fibre. In total degree one this yields
\[
 0=b_1(X_s)=h^0(X_s,\Omega^1_{X_s})+h^1(X_s,\mathcal O_{X_s})\geq1,
\]
a contradiction.

A closed smooth submanifold of a compact K\"ahler manifold is K\"ahler. A closed stratum of a Moishezon component is also Moishezon: choose a projective modification, take a component of the inverse image dominating the stratum, and intersect with general relative hyperplanes until the induced proper map is generically finite. This gives a generically finite projective cover and hence maximal algebraic dimension for the stratum. Smooth Moishezon manifolds satisfy the $\partial\bar\partial$-lemma, as recalled in \citet[\S2.1, p.~8]{Chen2026HodgeRoot}. Either stated hypothesis on the irreducible components therefore gives condition (K1), proving the last assertion. A finite parameter cover preserves the biholomorphism type of every nonzero fibre, so the argument applies for every $N\geq1$.
\end{proof}

\begin{theorem}[The obstruction on normalization components]\hypertarget{tag-0315}{}\label{push:src:B.4.5}
\StatementRecord{push:src:B.4.5}{proved}{0315}{\SeriesRef{push:src:B.4.2}, \SeriesRef{push:src:B.4.3}}

Let $D$ be a compact reduced complex space with normal-crossings
singularities, allowing self-intersections of global components.
Suppose that every connected component of its smooth normalization
satisfies the $\partial\bar\partial$-lemma. Then the natural map
\[
 H^1(D;\C)\longrightarrow H^1(D,\OO_D)
\]
is surjective.

Consequently every smooth-total-space reduced normal-crossings
filling of the specified six-sphere family, after any finite base
change, has a normalization component which fails the
$\partial\bar\partial$-lemma. In particular its normalization
cannot consist entirely of compact K\"ahler or Moishezon manifolds.
This conclusion includes normal crossings which are not simple.
\end{theorem}


\begin{proof}
\begin{enumerate}[label=\arabic*.,leftmargin=*]
\item Let
$D^{[p]}$ be the smooth compact branch space whose points are
ordered $(p+1)$-tuples of distinct local branches through one
point of $D$. Equivalently, we normalize the loci where these
branches meet, retaining their ordering. These are distinct
\emph{local branches}, even when they belong to one global
irreducible component. The finite holomorphic maps
$\pi_p:D^{[p]}\to D$ and their face maps are specified by
forgetting one branch. The symmetric group acts by reordering;
we take the alternating summand, with the sign character, in
each degree. Averaging over this finite group is exact over
$\C$. Write $\delta$ for the alternating sum of the face
pullbacks.

\item Both augmented complexes
\[
 \C_D\longrightarrow
 \bigl((\pi_p)_*\C_{D^{[p]}}\bigr)_{p\geq0}^{\mathrm{alt}},
 \qquad
 \OO_D\longrightarrow
 \bigl((\pi_p)_*\OO_{D^{[p]}}\bigr)_{p\geq0}^{\mathrm{alt}}
\]
are resolutions. This can be checked locally without a global
ordering. In a chart $z_1\cdots z_k=0$, the constant complex is
the augmented cochain complex of the simplex on the $k$ branches.
For the holomorphic complex, fix a monomial in the convergent
power series ring. It survives on exactly the intersections
indexed by subsets of those $i$ for which its $z_i$-exponent is
zero. If there are no such indices it already vanishes on $D$;
otherwise its coefficient complex is again an augmented simplex
complex. The usual contraction is a finite combination of
coordinate restrictions, so preserves convergence. This proves exactness in the convergent analytic local rings.
The descriptions agree under permutations of the local branches.

\item Resolve the holomorphic terms by smooth Dolbeault forms and
the constant terms by smooth de Rham forms on the branch spaces.
These compute the respective hypercohomologies. Indeed finite
pushforward is exact here, and partitions of unity from local
ambient smooth functions pull back to make the resulting sheaves
fine. The ordinary de Rham lemma and Dolbeault lemma apply on
each smooth branch space. Projection to the $(0,\bullet)$
part defines the natural comparison between the two total
complexes and induces $H^1(D;\C)\to H^1(D,\OO_D)$.

\item Represent a coherent degree-one class by a total cocycle
$(\alpha,f)$, where $\alpha$ is a smooth $(0,1)$-form on
$D^{[0]}$ and $f$ is an alternating smooth function on
$D^{[1]}$. Choose the total-complex signs so that its equations
are
\[
 \bar\partial\alpha=0,\qquad
 \delta\alpha=\bar\partial f,\qquad \delta f=0.
\]
On a normalization component, $\partial\alpha$ is of type
$(1,1)$, is $\partial$-exact and is $\bar\partial$-closed.
The $\partial\bar\partial$-lemma gives a smooth function $h$
with $\partial\alpha=\partial\bar\partial h$.
Do this on every component, and put
\[
 \alpha'=\alpha-\bar\partial h,\qquad
 f'=f-\delta h.
\]
The pair represents the same coherent class and still satisfies
the displayed cocycle equations; now $d\alpha'=0$.
Holomorphic face maps commute with both differential operators,
so
\[
 \partial\bar\partial f'
       =\partial\delta\alpha'=\delta\partial\alpha'=0.
\]
A smooth complex-valued pluriharmonic function on a compact
connected complex manifold is constant: apply the maximum
principle to its real and imaginary parts in local complex
coordinates. Therefore $f'$ is constant on each connected
component of $D^{[1]}$. It follows that
$df'=0$ and $\delta\alpha'=0$.

\item Thus the same pair $(\alpha',f')$, with $\alpha'$ viewed as
a complex-valued smooth one-form, is a cocycle of the total
de Rham branch complex. Its image is the original coherent
class. This proves surjectivity. The argument on $D^{[1]}$ uses its
compactness through the constancy of pluriharmonic functions.

\item For the asserted application, \ref{push:src:B.4.2} and \ref{push:src:B.4.3} give
$H^1(D;\C)=0$ and $h^1(D,\OO_D)\geq1$.
Surjectivity is impossible, so some normalization component
fails the lemma. Compact K\"ahler and smooth Moishezon
manifolds satisfy it, as in the primary comparison checked
in \ref{push:src:B.4.3}. The same argument applies after every finite base
change, since its smooth fibres are still the specified
six-spheres.
\end{enumerate}
\end{proof}

\begin{proposition}[A restricted non-K\"ahler cap obstruction]\hypertarget{tag-0316}{}\label{push:src:B.4.6}\label{push:B46}
\StatementRecord{push:src:B.4.6}{proved}{0316}{\SeriesRef{push:src:B.3.8}, \SeriesRef{push:src:B.3.7}, \SeriesRef{push:src:B.4.5}, \SeriesRef{push:src:B.4.2}, \SeriesRef{push:src:B.4.3}; \cite[Section 2, p.1352]{IstratiEtAl2022Root}; \cite[Proposition 8.4 and Theorem 8.5]{IstratiEtAl2022Root}; \cite[proof of Theorem 8.8, p.1380]{IstratiEtAl2022Root}; \cite[Theorem 5.2]{IstratiEtAl2022Root}}

Let $D$ be a compact reduced normal-crossings complex
space. Suppose its irreducible components have been
partitioned into closed unions $D_0,\ldots,D_s$, with
their reduced induced structures. Assume that
\begin{equation}
 H^1(D_i,\mathbb C)\longrightarrow H^1(D_i,\mathcal O)
 \quad\hbox{is surjective for every }i,
\label{pub:I:eq:B-46a}
\end{equation}
and that the analogous map is injective for every
pairwise intersection $D_i\cap D_j$. Then
\begin{equation}
 H^1(D,\mathbb C)\longrightarrow H^1(D,\mathcal O_D)
 \quad\hbox{is surjective}.
\label{pub:I:eq:B-46b}
\end{equation}

In particular these hypotheses hold for the following
specified cap assembly. The horizontal block $D_0$ is
a two-section pinching of a ruled bundle over an
analytic locally trivial torsor $V\to C$ of \ref{push:src:B.3.8},
with conductor map $A$ covering $\sigma$ and
$A^{12}=t_{6P}$. The normalization may undergo the
modifications of \ref{push:src:B.3.7}, or any smooth bimeromorphic
modifications for which the two conductor restriction
maps on $H^1$ are the same pullbacks as before.
The other blocks are smooth hyperbolic toric Kato
threefolds or compact smooth spaces satisfying
\eqref{pub:I:eq:B-46a}, for example class-$\mathcal C$ spaces.
Assume every pairwise overlap is either a reduced
union of invariant divisors of a hyperbolic toric
Kato threefold, or a compact normal-crossings surface
whose normalized components and normalized double
curves are rational. No such assembly can be a
semistable central fibre of the $S^6$ family.
\end{proposition}


\begin{proof}
\begin{enumerate}[label=\arabic*.,leftmargin=*]
\item Write $D^{[p]}$ for the disjoint union of intersections
of $p+1$ distinct blocks. The augmented closed-cover
complexes resolve both $\mathbb C_D$ and $\mathcal O_D$.
For the coherent assertion this can be checked in a
normal-crossings coordinate chart. The local equation
of $D$ is a product of distinct coordinates, and each
block groups a disjoint subset of those factors.
The corresponding ideals are generated by products
in disjoint sets of variables. The ordinary alternating
intersection complex is exact, as one sees on each
monomial, where it is the augmented complex of a
simplex. The same contraction acts coefficientwise
on convergent power series. For constant sheaves,
small intersections are contractible unions of
coordinate subspaces, and the same stalkwise simplex
complex is exact. This proves both resolutions also
when the global blocks are self-crossing.

\item Their hypercohomology spectral sequences have
\[
 E_1^{p,q}(\mathcal F)=H^q(D^{[p]},\mathcal F),
 \qquad \mathcal F=\mathbb C\ \hbox{or}\ \mathcal O.
\]
Every connected compact reduced branch intersection
has only constant holomorphic functions. Thus the
comparison maps in the entire $q=0$ row are
isomorphisms. They induce isomorphisms on
$E_2^{1,0}$ and $E_2^{2,0}$. Surjectivity in \eqref{pub:I:eq:B-46a}
and injectivity on pairwise intersections imply
surjectivity on $E_2^{0,1}$: lift a coherent
degree-one class on the blocks, then its topological
restriction difference vanishes because its coherent
image vanishes and the pairwise comparison is
injective.

\item Now use the natural five-term exact sequences
\[
 0\longrightarrow E_2^{1,0}
 \longrightarrow H^1(D,\mathcal F)
 \longrightarrow E_2^{0,1}
 \xrightarrow{d_2} E_2^{2,0}.
\]
The isomorphism on the last term makes a lift of a
coherent $d_2$-closed element topologically
$d_2$-closed. The isomorphism on the first term then
gives \eqref{pub:I:eq:B-46b}. If the block comparisons are
isomorphisms, the same argument also proves
injectivity, provided the map on $E_2^{0,1}$ is
injective.

\item The horizontal-block comparison follows from the cohomology
of its base curve. By \ref{push:src:B.3.8},
$H^1(V,\mathbb C)=H^1(C,\mathbb C)$ and
$H^1(V,\mathcal O)=H^1(C,\mathcal O_C)$.
The first equality is the Leray calculation there;
the second also follows from
$R^1\pi_*\mathcal O_V=\mathcal O_C(-\infty)$,
which has no global section. The action of $A$ on
these groups is $\sigma^*$, since it covers $\sigma$
and the displayed maps are pullbacks from the base.
The characteristic polynomial
\[
 p(x)=\Phi_{12}(x)\Phi_6(x),\qquad p(1)=1,
\]
shows that $1-A^*$ is invertible on both groups.
The normalization sequence of the two-section
pinching, for constants and for holomorphic
functions, therefore gives
\[
 H^1(D_0,\mathbb C)=\mathbb C,
 \qquad H^1(D_0,\mathcal O)=\mathbb C.
\]
Both classes come from the same difference map on
constant functions of the conductor. Their natural
comparison is an isomorphism. Projective-bundle
pullback identifies the normalization groups with
those of $V$. The stipulated modifications preserve
these groups and both restriction maps; in particular
the explicit modifications of \ref{push:src:B.3.7} have this property.
This grouped computation is essential: on an individual
conductor $V$, the comparison has dimensions six and
three and is not injective.

\item For a hyperbolic toric Kato manifold $Z$, its
fundamental group is $\mathbb Z$, as in
\cite[Section 2, p.1352]{IstratiEtAl2022Root};
\cite[Proposition~8.4, p.~1377, and Theorem~8.5, p.~1378]{IstratiEtAl2022Root}
give $H^0(Z,\Omega_Z^1)=0$ and
$H^1(Z,\mathcal O_Z)=\mathbb C$.
The exact sequence
$0\to\mathbb C\to\mathcal O\xrightarrow{d}d\mathcal O\to0$
shows that $H^1(Z,\mathbb C)\to H^1(Z,\mathcal O)$
is injective, because a global section of
$d\mathcal O$ is a closed holomorphic one-form.
The two dimensions are one, so it is an isomorphism.
The same reasoning applies to hyperbolic toric Kato
surfaces. Compact smooth rational varieties have
both groups zero. Class-$\mathcal C$ blocks satisfy
\eqref{pub:I:eq:B-46a} by their $\partial\bar\partial$ property.

\item Finally let $F$ be a reduced union of invariant
divisors in a hyperbolic toric Kato threefold.
Its normalized components are compact toric surfaces
or hyperbolic toric Kato surfaces
\cite[proof of Theorem 8.8, p.1380]{IstratiEtAl2022Root}.
Its normalized double curves are rational, since
every invariant curve is rational in the hyperbolic
case \cite[Theorem~5.2, pp.~1366--1367]{IstratiEtAl2022Root}.
Use the ordered-local-branch resolution from \ref{push:src:B.4.5}.
The comparisons on component $H^1$ are isomorphisms,
and the groups on normalized double curves vanish.
The same five-term argument shows that
$H^1(F,\mathbb C)\to H^1(F,\mathcal O_F)$ is an
isomorphism. For an overlap with rational normalized
components and double curves, the argument is
identical with zero component $H^1$. Thus every
pairwise overlap in the statement has the required
injective comparison, proving \eqref{pub:I:eq:B-46b} for the
specified assembly.

\item For a semistable $S^6$ centre, \ref{push:src:B.4.2} gives
$H^1(D,\mathbb C)=0$, while upper semicontinuity
and the nearby lower bound $h^1(\mathcal O)\geq1$ in \ref{push:src:B.4.3}
give $h^1(D,\mathcal O_D)\geq1$. This contradicts
\eqref{pub:I:eq:B-46b}. The obstruction is the stated surjectivity comparison on the blocks
and injectivity comparison on their pairwise overlaps.
\end{enumerate}
\end{proof}


\begin{remark}[A punctured isomorphism with an essential singularity]
\hypertarget{tag-PT03}{}\label{I:essential-punctured-isomorphism}
\StatementRecord{I:essential-punctured-isomorphism}{proved}{PT03}{}
The automorphism
\[
 (s,[z_0:z_1])\longmapsto(s,[e^{1/s}z_0:z_1])
\]
of $\Delta^*\times\mathbb P^1$ has no meromorphic extension across zero, although the two product families extend properly over the disc. Its restriction to the constant section $[1:1]$ is $[e^{1/s}:1]$, which has an essential singularity. The proper graph of a meromorphic extension, pulled back to this section and normalized in its dominating component, would give a holomorphic extension from the disc, a contradiction.
\end{remark}

\begin{proposition}[A Hopf surface degeneration and its comparison map]\hypertarget{tag-0317}{}\label{push:src:B.4.7}
\StatementRecord{push:src:B.4.7}{proved}{0317}{\SeriesRef{push:src:B.4.6}, \SeriesRef{push:src:B.4.2}, \SeriesRef{push:src:A.4.7}; \citep[Definition 1.1, pp.~85--86]{Fujimoto1990AB}; \citet[Section 2.2]{AbashevaDeev2024AB}}

There is a proper semistable analytic family of compact complex
surfaces $p:\mathcal H\to\Delta$ with smooth total space such that:
\begin{enumerate}
\item its punctured restriction is a holomorphic product with the
scalar primary Hopf surface
\[
 H_-=(\C^2\setminus\{0\})/\langle -e^{-\pi}\id\rangle;
\]
\item its central fibre is $Y=H_1\cup_E H_2$, where both $H_i$
are isomorphic to
\[
 H_+=(\C^2\setminus\{0\})/\langle e^{-2\pi}\id\rangle,
 \qquad E=\C/(\Z+i\Z);
\]
\item $N_{E/H_1}\simeq N_{E/H_2}\simeq\mathcal O_E$, and
\[
 \pi_1(Y)=0,\qquad h^1(Y,\mathcal O_Y)=1,
 \qquad H_*(Y;\Z)=(\Z,0,0,\Z^3,\Z^2);
\]
\item no punctured product isomorphism has a meromorphic extension
whose graph is proper over the two total spaces near the central
fibres.
\end{enumerate}
Consequently $\mathcal H\times\PP^1\to\Delta$ is a
semistable family of compact threefolds with simply connected centre
and $h^1(\mathcal O)=1$, whose normalized components fail the
$\partial\bar\partial$ lemma.  Its smooth fibres have Euler number
zero and fundamental group $\Z$.
\end{proposition}


\begin{proof}
\begin{enumerate}[label=\arabic*.,leftmargin=*]
\item We construct the family by an explicit instance of Fujimoto's
generalized logarithmic transformation
\citep[Definition 1.1, pp.~85--86]{Fujimoto1990AB}.
Let
\[
 b:\mathcal B=\operatorname{Bl}_{(p,0)}(\PP^1\times\Delta)
       \longrightarrow\PP^1\times\Delta.
\]
The central curve is $F_1\cup F_2$, where $F_1$ is the strict
transform of $\PP^1\times\{0\}$ and $F_2$ is the exceptional
projective line.  They meet transversely in one point.  The two
line bundles
\[
 L=b^*\operatorname{pr}_1^*\mathcal O_{\PP^1}(1),
 \qquad A=\mathcal O_{\mathcal B}(F_2)
\]
have degree-restriction matrix
\[
 \begin{array}{c|rr}
       &L&A\\\hline
 F_1&1&1\\
 F_2&0&-1
 \end{array}.
\]
In particular this matrix is unimodular.  Take nonzero holomorphic
transition functions $l_{ab},a_{ab}$ for these two line bundles.
On small intersections choose logarithms and put
\[
 t_{ab}=\left[
       \frac{(1+i)\log l_{ab}-i\log a_{ab}}{2\pi i}
                         \right]\in E.
\]
This is independent of the logarithm choices, since $1+i,-i$
belong to the lattice of $E$.  On triple intersections the sum
is again a lattice element.  The $t_{ab}$ therefore form a
holomorphic cocycle with values in the translation group $E$.
Glue $U_a\times E$ by these translations, obtaining a
holomorphic principal $E$-bundle $\mathcal H\to\mathcal B$.
It is a proper submersion, since it is locally a product with a
compact elliptic curve.  Its total space is smooth and Hausdorff.
Compose with $\mathcal B\to\Delta$.

\item The Chern vector of the bundle, in the lattice $\Z+i\Z$, is
\[
 \kappa=(1+i)c_1(L)-i c_1(A).
\]
Its restrictions to $F_1,F_2$ are respectively $1$ and $i$.
On a nonzero fibre its restriction is $1+i$.  These statements
also follow directly from the displayed transition functions.
Locally at the crossing of $F_1$ and $F_2$ the base map is
$xy=u$, and the principal elliptic bundle is locally trivial.
The composed family is thus locally $(x,y,z)\mapsto xy$,
with $z$ a coordinate on $E$.  It is flat, has smooth total
space and reduced simple normal crossings centre.  It is proper
because both maps in the composition are proper.

\item For a primitive vector
$c\in\Z+i\Z$, choose $d$ so that $(c,d)$ is an oriented
lattice basis.  The homomorphism
\[
 \C^*\longrightarrow E,\qquad
 z\longmapsto\left[\frac{c\log z}{2\pi i}\right]
\]
has kernel generated by
$\alpha_c=\exp(2\pi i d/c)$, with $|\alpha_c|<1$.
The principal $E$-bundle on $\PP^1$ associated to
$\mathcal O(1)^\times$ by this homomorphism is the quotient
of $\mathcal O(1)^\times$ by that cyclic kernel.  Fibrewise
inversion identifies $\mathcal O(1)^\times$ with
$\mathcal O(-1)^\times\simeq\C^2\setminus\{0\}$;
inverting the cyclic generator does not change its group.
The resulting surface is the scalar Hopf surface $H_{\alpha_c}$.
For the three vectors in the construction we may take
\[
 (c,d)=(1,i),\quad(i,-1),\quad(1+i,i),
\]
giving respectively
\[
 \alpha_1=\alpha_i=e^{-2\pi},\qquad
 \alpha_{1+i}=\exp\bigl(2\pi i(1+i)/2\bigr)=-e^{-\pi}.
\]
On $\mathcal B|_{\Delta^*}=\PP^1\times\Delta^*$,
$A$ has the canonical nowhere-zero section of its exceptional
divisor.  Its cocycle is therefore a coboundary there, while
$L$ is the constant $\mathcal O(1)$ bundle.  The cocycle coboundary therefore gives a holomorphic product
isomorphism with $H_-\times\Delta^*$.

\item Each central $H_i\to F_i$ is a principal elliptic bundle.
Its intersection with the other component is a smooth fibre,
so a local trivialization over a base disc gives
$N_{E/H_i}=\mathcal O_E$ exactly.  Thus the product of the
two normal lines is trivial, with the smoothing section supplied
by the coordinate $u=xy$.

\item A scalar Hopf surface is
diffeomorphic to $S^1\times S^3$: divide the radial coordinate
by multiplication by $|\alpha|$ and untwist its constant phase
rotation along the radial interval.  Its fibre inclusion has
fundamental-group map
\[
 \Lambda=\Z^2\longrightarrow\Lambda/\Z c_i.
\]
For $c_1=(1,0)$ and $c_2=(0,1)$, these are the two coordinate
quotients.  The van Kampen pushout kills both generators of
$\Lambda$, and hence $\pi_1(Y)=0$.  In the normalization
Mayer--Vietoris sequence the degree-one map, with the second
inclusion negated, is
\[
 H_1(E;\Z)\longrightarrow H_1(H_1;\Z)\oplus H_1(H_2;\Z),
 \qquad \begin{pmatrix}0&1\\-1&0\end{pmatrix}.
\]
It is unimodular.  The degree-zero map is the primitive column
$(1,-1)^t$, the degree-two target is zero, and $H_3(E)=0$.
The exact sequence now gives
\[
 H_1Y=H_2Y=0,\quad
 0\longrightarrow\Z^2\longrightarrow H_3Y
       \longrightarrow\Z\longrightarrow0,
 \quad H_4Y=\Z^2.
\]
All groups displayed in the proposition are therefore free.

\item For coherent cohomology, translations act trivially on
$H^1(E,\mathcal O_E)$, so for either central principal bundle
the sheaf $R^1\pi_*\mathcal O$ on $\PP^1$ is the constant
line bundle.  Its Leray sequence shows that restriction
\[
 H^1(H_i,\mathcal O_{H_i})\longrightarrow H^1(E,\mathcal O_E)
\]
is an isomorphism.  The normalization exact sequence therefore
has degree-one map $(s,t)\mapsto s-t$ from $\C^2$ to $\C$.
The preceding degree-zero map is surjective.  Its kernel has
dimension one, proving $h^1(Y,\mathcal O_Y)=1$.
This also displays exactly why the overlap-injectivity hypothesis
in \ref{push:src:B.4.6} fails: $H^1(E;\C)\to H^1(E,\mathcal O_E)$ has a
one-dimensional kernel.

\item The two scalar Hopf surfaces used above are
not biholomorphic.  An isomorphism lifts to an automorphism of
$\C^2\setminus\{0\}$ that intertwines the two cyclic generators
or their inverses.  Hartogs extension applies to the lift and
its inverse; both fix zero and have invertible derivative there.
Differentiating the intertwining identity gives
$e^{-2\pi}=(-e^{-\pi})^{\pm1}$, which is false.  They have
no rational curves, since a map $\PP^1\to H_\alpha$ lifts to
$\C^2\setminus\{0\}$ and is constant.  The usual resolution
of a bimeromorphic map of smooth compact surfaces therefore
also shows that these two Hopf surfaces are not bimeromorphic.
This agrees with the explicit scalar case of
\citet[arXiv version~1, Section~2.2, pp.~4--5]{AbashevaDeev2024AB}.

\item The fundamental groups also exclude a meromorphic
extension of the punctured product comparison. A sufficiently
small closed semistable neighbourhood in $\mathcal H$ retracts
onto $Y$, by the monomial-gradient argument of \ref{push:src:B.4.2}.  Its
fundamental group is zero.  The product neighbourhood
$H_-\times\Delta$ has fundamental group $\Z$.  If the
punctured comparison had a proper meromorphic graph, resolving
that graph would give a common smooth proper modification of
the two neighbourhoods.  A proper bimeromorphic map between
smooth complex manifolds preserves $\pi_1$, by the
codimension-two avoidance argument in \ref{push:src:A.4.7}.  This is a
contradiction.  Both models already have the indicated
holomorphic map to the same $\PP^1$; the example therefore
retains an extended base map while the total-space comparison
fails.

\item Finally take the product with $\PP^1$.  Its central integral
Betti numbers are
\[
 (1,0,1,3,2,3,2),
\]
and $h^1(\mathcal O)$ remains one.  The nearby fibre is
$H_-\times\PP^1$, of fundamental group $\Z$ and Euler
number zero.  The resulting fundamental group and Euler number are the stated
invariants of these smooth fibres.
\end{enumerate}
\end{proof}

\begin{remark}[Elliptic-bundle and Hopf degeneration inputs]
\hypertarget{tag-PT10}{}\label{I:elliptic-bundle-primary-comparisons}
\StatementRecord{I:elliptic-bundle-primary-comparisons}{proved}{PT10}{\citep[Definition~1.1 and Remark~1.4, pp.~85--86; Proposition~1.14, pp.~92--93; Example~1.15, pp.~93--94]{Fujimoto1990AB}; \citet[Theorem~1.4, preprint p.~7, and Example~2.1, p.~13]{Nishiguchi1988IIRoot}; \citet[Propositions~3.1--3.4, preprint pp.~6--9]{SchroeerSiebert2002Root}}
The Cartier-divisor construction used in Proposition~\ref{push:src:B.4.7} is the construction of \citet[construction, p.~85; Definition~1.1, p.~86]{Fujimoto1990AB}. His Remark~1.4 attributes non-K\"ahler torus degenerations to Ueno. His Proposition~1.14, pp.~92--93, extends principal elliptic bundles across subsets of codimension at least three; Example~1.15, pp.~93--94, gives a principal-bundle extension failure in codimension two. Hopf components occur in the semistable degenerations of \citet[Theorem~1.4, preprint p.~7, and Example~2.1, p.~13]{Nishiguchi1988IIRoot}. The diagonal Hopf quotients, elliptic-conductor identifications and normal lines in \citet[Propositions~3.1--3.4, preprint pp.~6--9]{SchroeerSiebert2002Root} provide related explicit models. Proposition~\ref{push:src:B.4.7} proves its own proper family using the stated Cartier-divisor cocycle.
\end{remark}

\begin{remark}[Normal degrees along the exceptional ruling]
\hypertarget{tag-PT11}{}\label{I:normal-degree-triple-section}
\StatementRecord{I:normal-degree-triple-section}{proved}{PT11}{\cref{push:src:B.4.7,push:src:B.4.10}}
For the divisor $D=E\times\mathbb P^1$ in $H_i\times\mathbb P^1$, the normal line is pulled back from $E$ and has degree zero on every ruling fibre. For the exceptional divisor obtained by blowing up a smooth elliptic curve in a threefold, the tautological normal restricts to $\mathcal O_{\mathbb P^1}(-1)$ on that fibre. Restricting the normal-crossing identity
\[
 N_{D/A}\otimes N_{D/B}\otimes\mathcal O_D(T)\simeq\mathcal O_D
\]
to the ruling therefore gives $-1+0+\deg(T|_{\mathbb P^1})=0$. In particular an empty triple divisor has the wrong degree, whereas the section $T=E\times\{p\}$ contributes exactly $+1$. Proposition~\ref{push:src:B.4.10} computes the full line bundles, including their elliptic factors.
\end{remark}

\begin{proposition}[Generic simplicity and the distinguished orbit]
\hypertarget{tag-PT12}{}\label{I:generic-torus-simplicity}
\StatementRecord{I:generic-torus-simplicity}{proved}{PT12}{the marked period matrix and the nontorsion elliptic section of the construction}
Fix an ordinary base point at which $\mu\notin\mathbb Q\tau+\mathbb Q$. Such base points exist. Outside a countable set of admissible values of the parameter $\beta$, the corresponding two-dimensional complex torus is simple. For every admissible $\beta$ at this base point, the distinguished line $\mathbb C e$, with $e=\Pi\widehat\delta=(0,1)^t$, meets the period lattice in $\mathbb Z e$. The induced orbit map is the injective holomorphic group map
\[
 \iota:\mathbb C^*\longrightarrow\mathbb C^2/\Pi\Lambda,
 \qquad \iota(e^{2\pi iz})=[ze].
\]
Here $[z]\mapsto e^{2\pi iz}$ identifies $\mathbb C/\mathbb Z$ with the source.
\end{proposition}
\begin{proof}
On an ordinary uniformizing chart the equations $\mu=r\tau+s$, for $(r,s)\in\mathbb Q^2$, define countably many analytic loci. If they covered the chart, Baire's theorem would force one equation on an open set and hence identically. The marked elliptic section would then be torsion, contrary to the construction. Thus the required base point exists.

For integral vectors $v=(a,b,c,d)$ and $v'=(A,B,C,D)$,
\[
 \Pi v=(6a\mu+b\tau+c,\ a\beta+b\mu+d).
\]
The determinant $\det(\Pi v,\Pi v')$ is affine in $\beta$, with coefficient
\[
 (Ab-aB)\tau+(Ac-aC).
\]
If the coefficient vanishes and $a\ne0$, the nonreal number $\tau$ gives $B=Ab/a$ and $C=Ac/a$. The remaining determinant is
$(6a\mu+b\tau+c)(D-Ad/a)$.
Its first factor is nonzero by the rational independence of $1,\tau,\mu$, and its second factor vanishes only when $v,v'$ are rationally dependent. The case $a=0$, $A\ne0$ follows by exchanging the vectors. For $a=A=0$, the determinant is
\[
 (bD-Bd)\tau+(cB-Cb)\mu+(cD-Cd).
\]
Rational independence makes this zero only when all the two-by-two minors vanish, again giving dependence. Therefore an independent integral pair has zero determinant for at most one value of $\beta$. There are countably many such pairs.

A complex line defining an elliptic subtorus contains two real-linearly independent periods. Conversely, any such period pair spans a complex line with a rank-two lattice; admissibility makes the real period map an isomorphism. Avoiding the countable exceptional values therefore gives simplicity. The equation $6a\mu+b\tau+c=0$ forces $a=b=c=0$, which proves $\mathbb C e\cap\Pi\Lambda=\mathbb Z e$ and the claimed orbit map.

For comparison with an elliptic-bundle fibre, every holomorphic map from an elliptic curve to a complex torus lifts to a map $\mathbb C\to\mathbb C^2$ whose derivative is doubly periodic. Compactness makes that derivative constant. A nonconstant such map consequently has image a translate of an elliptic subtorus. Thus a simple torus obtained above has no elliptic-curve fibre of this kind.
\end{proof}

\begin{remark}[The distinguished orbit on an ordinary torus]
\hypertarget{tag-PT04}{}\label{I:noncompactifiable-distinguished-orbit}
\StatementRecord{I:noncompactifiable-distinguished-orbit}{proved}{PT04}{}
The distinguished nonconstant orbit map $\mathbb C^*\to T^2$ on an ordinary torus has no meromorphic extension to $\mathbb P^1$. A meromorphic map from a smooth curve to a compact complex manifold extends holomorphically at its isolated indeterminacies, by normalization of its proper graph. A holomorphic map $\mathbb P^1\to T^2$ lifts to $\mathbb C^2$ and is constant. This contradicts the nonconstant orbit map.
\end{remark}

\begin{proposition}[Normal-line obstruction at good reduction]\hypertarget{tag-0318}{}\label{push:src:B.4.8}
\StatementRecord{push:src:B.4.8}{proved}{0318}{\citet[Section 2.2]{AbashevaDeev2024AB}}

Let $a,b$ be holomorphic units on a disc $\Delta_q$, with
$0<|a(q)|<1$, and let $m\ne0$ be an integer.  Put
\[
 \mathcal E_a=(\C^*\times\Delta_q)/
       \bigl((z,q)\sim(a(q)z,q)\bigr).
\]
This is a smooth proper elliptic family over the full disc.
Over $\Delta_q^*$ define a line bundle $\mathcal L_m$ by
\[
 (z,v,q)\sim\bigl(a(q)z,\ q^m b(q)v,\ q\bigr).
\]
Its restriction to each elliptic fibre has degree zero.
After any finite base change $q=s^N\varepsilon(s)$, where
$N>0$ and $\varepsilon$ is a holomorphic unit, the following
extension is impossible: a proper semistable family
$p:\mathcal C\to\Delta_s$ of connected curves, with smooth
total surface, an isomorphism
$\mathcal C|_{\Delta_s^*}\simeq\mathcal E_a|_{\Delta_s^*}$,
and a line bundle on $\mathcal C$ restricting to
$\mathcal L_m$ under that isomorphism.

The assertion allows any finite collection of rational tails
in the central curve and arbitrary translations in the
punctured isomorphism.  It remains true after tensoring
$\mathcal L_m$ by the restriction of a line bundle extending
on the interface.

In particular the axis $z_1=0$ in the diagonal Hopf surfaces
\[
 H_{q,a(q)}=(\C^2\setminus\{0\})/
       \langle(z_1,z_2)\mapsto(qz_1,a(q)z_2)\rangle
\]
cannot acquire such a semistable completion together with
its normal line, after any finite base change.
\end{proposition}


\begin{proof}
\begin{enumerate}[label=\arabic*.,leftmargin=*]
\item Choose a logarithm of $a$ on the disc and set
$\tau_a=\log a/(2\pi i)$.  Then $\Im\tau_a>0$, and
\[
 \mathcal E_a\simeq
 (\C\times\Delta_q)/(\Z+\tau_a(q)\Z).
\]
The real coordinates $z=\exp(2\pi i(x+\tau_a(q)y))$
give a smooth trivialization by the marked torus
$(x,y)\in\R^2/\Z^2$.  In particular its monodromy on
$H_1$ is the identity.  The line factors in this marking
are $1$ along the $x$ cycle and $q^m b(q)$ along the
$y$ cycle.  On a positive parameter circle, the latter
factor has winding $m$.  Equivalently, the clutching
calculation gives
\[
 \langle c_1(\mathcal L_m),[T_x]\rangle=0,
 \qquad
 \langle c_1(\mathcal L_m),[T_y]\rangle=m,
 \qquad
 \langle c_1(\mathcal L_m),[\mathcal E_{a,q}]\rangle=0.
\]
Here $T_x,T_y$ sweep the indicated cycles around the
parameter circle; the orientation of $T_y$ is chosen to
agree with the clutching convention in the displayed
line factor.  Reversing that convention changes both
sides' sign and has no effect on nonvanishing.  One can
also compute the integer on the square
$[0,1]_y\times S^1_q$: the two $y$ edges are glued by
$q^m b(q)$, whose degree is $m$.

\item Under $q=s^N\varepsilon(s)$ the degree becomes $mN$.
The unit $\varepsilon(s)^m b(q(s))$ has zero winding.
Thus the mixed Chern class is the nonzero covector
\[
  (0,mN)\in
 H^1(\mathcal E_{a,q};\Z)\otimes H^1(S_s^1;\Z).
\]

\item Suppose an extension $\mathcal C$ existed.  Its central
fibre is connected.  Flatness and properness give arithmetic
genus one, and the normalization exact sequence gives
\[
 1=\sum_i g(C_i)+b_1(\Gamma),
\]
where $C_i$ are the smooth normalization components and
$\Gamma$ is the dual graph, including a loop for a
self-node.  Consequently either there is a unique
elliptic component and $\Gamma$ is a tree, or all
components are rational and $\Gamma$ has one cycle.

\item The second case has nontrivial monodromy.  This can be
seen directly in the local plumbing $uv=s$.  Transport
around $s=0$ is a positive Dehn twist on its neck annulus.
An edge outside the unique graph cycle gives a separating
neck and acts trivially on $H_1$.  Each of the $k\ge1$
edges in the graph cycle gives a nonseparating neck.
These disjoint necks represent the same primitive class
up to sign on the smoothed genus-one surface.  The
transvection is unchanged by reversing that class.
Their product therefore has, in a symplectic basis, the
matrix
\[
  \begin{pmatrix}1&k\\0&1\end{pmatrix}
\]
or its inverse under the opposite monodromy convention.
It is not the identity.  The given punctured elliptic
family has identity monodromy, so this case is impossible.

\item There is therefore an elliptic core $C_0$ with rational
trees attached at finitely many points.  Choose two
oriented loops on $C_0$ avoiding those points and with
intersection number one.  The map $p$ is a submersion
near their union.  Compactness of the loops and a smooth
horizontal distribution provide, after shrinking the
disc, families of loops over the full disc.  Each swept
torus on its boundary bounds the resulting solid torus
in $\mathcal C$.  Hence the first Chern class of every
line bundle on $\mathcal C$ evaluates to zero on both
such swept tori.

\item The two transported loops form a basis of the nearby
elliptic fibre's $H_1$: cutting off the rational trees
adds only genus-zero pieces along separating necks,
and their intersection number remains one.  Under the
given punctured isomorphism, their swept tori therefore
project to a basis of the mixed group
$H_1(\mathcal E_{a,q};\Z)\otimes H_1(S_s^1;\Z)$.
They may differ from the marked swept tori by integral
fibre classes.  This accounts in particular for an
arbitrary loop of fibre translations in the isomorphism.
The fibre degree of $\mathcal L_m$ is zero, so these
additional fibre classes do not change either pairing.
The nonzero covector $(0,mN)$ cannot vanish on both
basis vectors.  This contradicts the bounding-solid-torus
calculation and proves the assertion.

\item Finally an extending tensor factor has zero pairing on
both bounding tori by the same argument.  It cannot
cancel the obstruction.  Vertical divisors and any
Cartier triple divisor on such an extended interface
are included among these harmless tensor factors.

\item For the stated Hopf consequence, the axis is precisely
$\mathcal E_{a,q}$.
The derivative of the deck transformation in the $z_1$
direction multiplies a normal vector by $q$; its action
on the axis coordinate is multiplication by $a(q)$.
These are exactly the two factors defining
$\mathcal L_1$.  This is also the explicit diagonal-Hopf
normal-line convention of
\citet[arXiv version~1, Section~2.2, pp.~4--5]{AbashevaDeev2024AB}.
The assertion just proved applies with $m=1$ and $b=1$.
It also applies to the first multiplier $q^m b(q)$ for
any positive $m$.  At the cusp the nonzero unit
$e^{2\pi i\mu(q)}$ has absolute value different from one
at $q=0$.  Indeed the nodal cubic is
$Y^2=(X-1)^2(X+2)$; its normalization
$X=h^2-2$, $Y=h(h^2-3)$ and the multiplicative coordinate
$(h-\sqrt3)/(h+\sqrt3)$ show that the marked section at
$(X,Y)=(0,\pm\sqrt2)$ has multiplier one of
$-(5\pm2\sqrt6)$.
Choose that multiplier or its inverse to obtain
such a contracting $a(q)$.
This proves good reduction for this axis.
\end{enumerate}
\end{proof}

\begin{remark}[The individual mixed Chern class]
\hypertarget{tag-PT13}{}\label{I:individual-normal-winding}
\StatementRecord{I:individual-normal-winding}{proved}{PT13}{\cref{push:src:B.4.7,push:src:B.4.8}}
The evaluation isomorphism $\mathcal L_m\otimes\mathcal L_m^{-1}\simeq\mathcal O$ cancels the two mixed first Chern classes. Individually, Proposition~\ref{push:src:B.4.8} computes the nonzero covector $(0,mN)$ on the two swept tori after a parameter cover of order $N$. Both swept tori bound solid tori on an extending semistable good-reduction interface. Hence the restriction of the first Chern class of every extending line evaluates to zero on both, whereas $c_1(\mathcal L_m)$ does not. This is the explicit restriction map relating the individual obstruction to the tensor-product normal equation. In the constant elliptic family of Proposition~\ref{push:src:B.4.7}, the corresponding mixed class is zero.
\end{remark}

\begin{proposition}[Ordinary torus cocycles and descent classes]\hypertarget{tag-0319}{}\label{push:src:B.4.9}
\StatementRecord{push:src:B.4.9}{proved}{0319}{\SeriesRef{push:src:B.3.1}, \SeriesRef{push:src:C.3.5}; \citep[construction, p.~85; Definition~1.1, p.~86]{Fujimoto1990AB}; \citep[expos\'e~19, \S5, pp.~14--15]{Cartan1952B_AB}}

Let $\mathcal T\to U$, $U=\Delta_q^*\times\Delta_r^*$,
be the Jacobian two-torus family of the marked cusp,
after $w=r^6$.  A holomorphic homomorphism of group
families
\[
 \C^*\times U\longrightarrow\mathcal T
\]
is determined by a flat integral cycle.  For this family
all such cycles are
\[
 c=a\widehat w+b\widehat\delta,\qquad a,b\in\Z.
\]
In the exponential cusp coordinates its homomorphism
is $z\mapsto(z^a,z^b)$.  Thus the ordinary Cartier
logarithm construction along a globally specified
cocharacter has transition maps
\[
 (Z,W)\longmapsto(g_{ij}^{\,a}Z,g_{ij}^{\,b}W).
\]
If these units extend over the prescribed fan charts,
the maps extend there and preserve each toric
boundary component.  These transition units preserve the component permutations
$+4,-3,0$ of \ref{push:src:B.3.1}. For the sheaf $\mathscr T$ of
holomorphic sections of $\mathcal T$,
\[
 H^1(U,\mathscr T)\simeq H^2(U,\Lambda)\simeq\Z^2.
\]
The last group is represented by the coinvariant
classes of $\widehat\gamma,\widehat u$.  The ordinary
fixed-cocharacter Cartier constructions have zero
image in this group.  The cusp family with
the coherent closing translations of \ref{push:src:C.3.5} also
has zero torsor class, by the explicit holomorphic
section below.  A marked torsor isomorphism preserves this cohomology class, so a
torsor with a nonzero class cannot be isomorphic to this marked cusp
torsor.
\end{proposition}


\begin{proof}
\begin{enumerate}[label=\arabic*.,leftmargin=*]
\item First consider any smooth holomorphic family of
complex tori with integral period local system
$\Lambda$ and relative Lie bundle $\mathcal V$.
Lift a group homomorphism from $\C^*$ along the
universal covering $z=\exp(2\pi i t)$.
The lifted homomorphism from the additive group
$\C$ to the universal covering vector space of
each torus is $t\mapsto t v$ for a holomorphic
section $v$ of $\mathcal V$.  Periodicity under
$t\mapsto t+1$ says precisely that $v$ belongs
to the period lattice in every fibre.  In any
local marking its integral coefficients are
continuous, and hence locally constant.  It
therefore defines a flat global section of
$\Lambda$.  Conversely any such section gives
the homomorphism by this formula.  The image
need not be a closed elliptic subgroup; this
argument applies without a polarization and
also to simple two-tori.

\item In the ordered basis
$(\widehat\gamma,\widehat u,\widehat w,\widehat\delta)$,
the commuting linear monodromies have
nilpotent parts
\[
 N_q=A_0-I=
 \begin{pmatrix}
 0&0&0&0\\0&0&0&0\\0&1&0&0\\-1&0&0&0
 \end{pmatrix},\qquad
 N_r=-12\widehat\delta\otimes\gamma.
\]
The sign of $N_r$ is the positive-loop closing
convention of \ref{push:src:C.3.5}; reversing it leaves all
groups below unchanged.  Their common integral
kernel is exactly $\Z\widehat w\oplus\Z\widehat\delta$.
The two corresponding columns of $\Pi_c$ are
$(1,0)^t,(0,1)^t$.  Exponentiating the homomorphism
therefore gives $(z^a,z^b)$, as asserted.

\item Fujimoto's logarithmic cocycle formula
\citep[construction, p.~85; Definition~1.1, p.~86]{Fujimoto1990AB} works in
the same way for a complex torus of any dimension:
apply the displayed homomorphism to the transition
functions $g_{ij}$ of a line bundle.  Equivalently
use $c\log(g_{ij})/(2\pi i)$ modulo the period
lattice.  A different logarithm adds an integral
cycle, and the triple-overlap defect is an
integral cycle as well.  With the present common
invariant cycles this is exactly multiplication
by $(g_{ij}^a,g_{ij}^b)$.  Torus multiplication
extends to every toric variety of a fixed fan
and preserves every torus orbit and boundary
ray.  For an extending unit it is an automorphism
on the full chart.  Consequently it contributes
no permutation of the central component labels.
The product-of-permutations obstruction of \ref{push:src:B.3.1}
is unchanged.  This conclusion keeps the fan,
period action, and local component models fixed.

\item Retract $U$ to its real two-torus.  The cellular
cochain complex with coefficients in the local
system is the two-generator Koszul complex.
Its last differential has image
$N_q\Lambda+N_r\Lambda$, and hence
\[
 H^2(U,\Lambda)
   =\Lambda/(N_q\Lambda+N_r\Lambda)
   =\Lambda/(\Z\widehat w+\Z\widehat\delta)
   \simeq\Z\widehat\gamma\oplus\Z\widehat u.
\]
The relation lattice is primitive: $N_q$ already
has image $\Z\widehat w+\Z\widehat\delta$.

\item The analytic exponential sequence is
\[
 0\longrightarrow\Lambda\longrightarrow
 \mathcal O(\mathcal V)\longrightarrow
 \mathscr T\longrightarrow0.
\]
The product of two punctured discs is Stein,
and $\mathcal O(\mathcal V)$ is coherent.
Cartan's theorem B therefore kills its positive
cohomology
\citep[expos\'e~19, \S5, pp.~14--15]{Cartan1952B_AB}.
The long exact sequence gives the stated
isomorphism $H^1(U,\mathscr T)\simeq H^2(U,\Lambda)$.
The representatives of $H^1(U,\mathscr T)$ give holomorphic torsors
on $U$ by translation gluing.

\item For the last assertion, naturality of the two
exponential sequences sends the first Chern class
of a Cartier line bundle to its multiple of the
flat cycle $c$ in $H^2(U,\Lambda)$.  Both common
invariant cycles lie in $N_q\Lambda$, so this
multiple is zero.

\item We compute the affine class using
the cusp functions of \ref{push:src:C.3.5},
\[
 q=e^{2\pi i s},\qquad
 C_0(q)=\beta^{\rm part}+s,\qquad
 \ell=-s/12+\kappa_0(q),
 \qquad D(q)=C_0(q)-12\kappa_0(q).
\]
The functions $\mu,C_0,\kappa_0$ are single-valued
and holomorphic at $q=0$.  The cusp unit
representative of the even-shift translation is
\[
 \bar b_{h,c}(q)=
 \left(-\frac h2\mu(q),\,
       -\frac h{12}(c+D(q))-\frac{h^2}{24}\right).
\]
On the logarithmic parameter cover define
\[
 \sigma(q,c)=
 \left(-\frac c2\mu(q),\,
       -\frac{c^2}{24}-\frac{cD(q)}{12}\right).
\]
Direct subtraction gives the exact identity
\[
 \sigma(q,c+h)-\sigma(q,c)=\bar b_{h,c}(q).
\]
In particular it is equivariant for the deck
transformation $c\mapsto c-12$ corresponding to
a positive $r$ loop, and it is single-valued in
$q$.  The original cusp has no extra affine
$q$ translation ($l_0=0$).  Thus $\sigma$ descends
to a holomorphic section of the affine
torus family.  Equivalently its $q$ and $r$
affine lifts commute exactly: the $r$ translation
$\bar b_{-12,c}$ is independent of the logarithm
branch $s$.  The resulting integral commutator
representing the Chern class is zero already in
$\Lambda$, before passage to coinvariants.

\item Multiplication by the exponential of $\sigma$
is a torus-unit automorphism of every cusp fan
chart for a fixed nonzero parameter.  Since its
coefficients are holomorphic in the full $q$
disc, this comparison includes the prescribed
Mumford cusp fibre for $r\ne0$.  The displayed comparison is defined on the full $q$ disc for
$r\ne0$, with $c=-(12/(2\pi i))\log r$ on its parameter cover.
Changing coherent closings by a scalar depending
only on the parameter has zero $q$ commutator
and leaves this Chern-class conclusion unchanged.
\end{enumerate}
\end{proof}

\begin{proposition}[Parabolic families and exceptional divisors]\hypertarget{tag-030E}{}\label{push:src:B.4.10}
\StatementRecord{push:src:B.4.10}{proved}{030E}{\SeriesRef{push:src:B.4.8}; \citet[Section 2, pp.1351--1353]{IstratiEtAl2022Root}; \citet[Introduction, pp.1750--1751]{Teleman2010AB}; \citet[Example 9.10, p.1386]{IstratiEtAl2022Root}; \citet[Proposition 5.6, pp.1369--1370]{IstratiEtAl2022Root}; \citet[Section 1.2]{DlousskyTeleman2012AB}}

Fix an integer $n\ge1$ and holomorphic units $a(s),\kappa(s)$ on a
disc, with $0<|a(s)|<1$. After shrinking the disc there is a proper
holomorphic submersion $\mathcal K\to\Delta_s$, with an embedded smooth
proper elliptic family
\[
 \mathcal E=(\C^*\times\Delta_s)/
       ((v,s)\sim(a(s)v,s)).
\]
Its central surface is a parabolic Inoue surface $S_n$, with
$\pi_1(S_n)=\Z$, $b_2(S_n)=n$, and elliptic axis of square $-n$.
For $s\ne0$, the surface $\mathcal K_s$ is obtained from the diagonal
Hopf surface $H_{\kappa(s)s^n,a(s)}$ by $n$ successive blow-ups at one
axis point and its indicated infinitely near strict-axis points.
In a fixed annulus marking of $\mathcal E$, its normal line has factor
\[
 g_s(v)=\kappa(s)(a(s)v+s)^n.
\]
This line extends across $s=0$ and is isomorphic to the holomorphic
family with factor $\kappa(s)a(s)^n v^n$.
For constant $a,\kappa$ it is a constant family of degree $-n$ lines.

The relation with the normal line of the unmodified Hopf axis is
\[
 N_{\mathcal E_s/\mathcal K_s}
   =L_{\kappa(s)s^n}\otimes\OO_{\mathcal E_s}(-n[-s]),
       \qquad s\ne0.
\]
Here $n[-s]$ is the total exceptional correction, with the
multiplicity of each successive blow-up retained. It has no
Cartier extension on $\mathcal E$, even after a finite base change or
the addition of a vertical component. The marked divisor used to change the normal line therefore
fails the extension condition in Proposition~\ref{push:src:B.4.8}.

For $n=8$, the degree $-8$ axis normal can be chosen to be the inverse
of $\OO_E(2)$ on any chosen smooth elliptic complete intersection
$E=(2,2)\subset\PP^3$. This supplies the normal-line matching with
$\operatorname{Bl}_E\PP^3$ described below, including one positive
triple section. The matching is the displayed isomorphism of the two normal lines,
with the stated triple-curve factor.
\end{proposition}


\begin{proof}
\begin{enumerate}[label=\arabic*.,leftmargin=*]
\item Let
$B_R\subset\C^2$ be a ball, and let
$\pi_n:\widehat B_R\to B_R$ be the following modification.
Blow up the origin, then successively the intersection of the newest
exceptional curve with the strict transform of the axis $x=0$.
In the last chart,
\[
 \pi_n(z,w)=(zw^n,w).
\]
Define the parameter-dependent embedding in that chart by
\[
 \sigma_s(z,w)=
       \bigl(\kappa(s)z,\ a(s)w+(1-a(s))s\bigr).
\]
Choose constants $A<\rho<1$ and $K>0$ with
$|a(s)|\le A$ and $|\kappa(s)|\le K$ on a closed smaller disc.
Shrink it further so that
$|a(s)w+(1-a(s))s|\le\rho R$ on $\overline B_R$.
Then
\[
 |\pi_n\sigma_s(z,w)|^2
 \le K^2R^2(\rho R)^{2n}+\rho^2R^2<R^2
\]
when $R$ is small enough. The last-chart coordinates of these images
are bounded, so this is compact containment of
$\sigma_s(\overline B_R)$ in $\widehat B_R$, uniformly in $s$.
For each finite bound on $|\kappa|$, these inequalities give such a radius $R$.

\item Remove the embedded ball and identify the outer and inner collars by
$\sigma_s\pi_n$. The outer collar avoids the exceptional locus, and
$\pi_n$ is a biholomorphism there; thus the collar identification is
a biholomorphism. Carry out the same construction over the parameter
disc. The embeddings and the identifications are holomorphic and
preserve $s$. Uniformly disjoint collars give a Hausdorff smooth
threefold, and the projection is a holomorphic submersion in every
chart. Uniform compact fundamental regions over compact parameter
sets prove properness. This is the usual Kato construction, with
the explicit moving embedding above; see
\citet[Section 2, pp.1351--1353]{IstratiEtAl2022Root}.
The general occurrence of deformations from Kato surfaces to
blown-up Hopf surfaces is classical; see
\citet[Introduction, pp.1750--1751]{Teleman2010AB}.
We use explicit formulas here to retain the elliptic marking and
its exceptional correction.

\item The contracted map $\pi_n\sigma_s$ fixes $(0,s)$.
In the coordinate $v=w-s$ it becomes
\[
 f_s(x,v)=
    \bigl(\kappa(s)x(a(s)v+s)^n,\ a(s)v\bigr).
\]
At $s=0$ this is the parabolic monomial germ with exponent matrix
$\left(\begin{smallmatrix}1&n\\0&1\end{smallmatrix}\right)$.
The chosen modification is exactly the $n$-fold modification in
\citet[Example 9.10, p.1386]{IstratiEtAl2022Root}.
Consequently the central surface is the stated parabolic Inoue
surface; its elliptic axis is disjoint from the rational cycle.
The construction gives $\pi_1=\Z$ and $b_2=n$, or equivalently one
can use the cited description of this surface.

\item The strict axis is preserved in the collar gluing. Its coordinate
$v$ is multiplied by $a(s)$. To see its full quotient, the axis disc
in the blowdown coordinates is
\[
 D_s=\{v:|v+s|<R\}.
\]
For small $s$, it contains $0$ and satisfies $a(s)D_s\Subset D_s$.
Its dilates have union $\C$ and shrinking intersection $\{0\}$.
Thus $D_s\setminus a(s)D_s$, with the two boundary collars glued,
is a fundamental annulus for $\C^*/a(s)^\Z$.
This proves the exact elliptic marking throughout the parameter disc.

\item There are two normal frames, whose distinction is essential.
In the last blow-up chart the collar map $\sigma_s\pi_n$ has factor
\[
 g_z(v)=\kappa(s)(v+s)^n.
\]
On a fixed annulus avoiding $v=0$, shrink $s$ so that both
$v+s$ and $a(s)v+s$ are nonzero. The blowdown normal coordinate is
$x=z(v+s)^n$, and in that frame the factor is
\[
 g_x(v)=\kappa(s)(a(s)v+s)^n.
\]
The change between the frames is a nonzero gauge on these annuli.
The central factor is $\kappa(0)a(0)^n v^n$ and has degree $-n$,
with the factor-of-automorphy convention of
\citet[Proposition 5.6, pp.1369--1370]{IstratiEtAl2022Root}.

\item An explicit regular gauge verifies the whole family. On a fixed
annulus and its overlap with its $a(s)$-translate put
\[
 H_s(v)=\prod_{j\ge0}\left(1+\frac{a(s)^j s}{v}\right)^n.
\]
The product converges locally uniformly in $(v,s)$ and is nonzero
after shrinking the parameter disc. It satisfies
\[
 \frac{H_s(a(s)v)}{H_s(v)}
      =\left(1+\frac{s}{a(s)v}\right)^n.
\]
Replacing $x$ by $x/H_s(v)$ changes $g_x$ to
$\kappa(s)a(s)^n v^n$.
The collar maps are linear in the normal coordinate, so the same
calculation identifies a tubular line neighbourhood.

\item The generic Hopf blowdown also has an explicit
developing map. Fix $s\ne0$ and set
\[
 G_s(v)=\prod_{j\ge1}\left(1+\frac{a(s)^j v}{s}\right)^n,
       \qquad
 \Psi_s(x,v)=(xG_s(v),v).
\]
This product is entire in $v$, converges locally uniformly, and obeys
\[
 G_s(a(s)v)=\frac{G_s(v)}{(1+a(s)v/s)^n},\qquad
 \Psi_s f_s=\operatorname{diag}(\kappa(s)s^n,a(s))\Psi_s.
\]
For sufficiently small nonzero $s$, both diagonal multipliers have
absolute value less than one. Compose $\Psi_s$ with $\pi_n$ on a
Kato fundamental region. The displayed equivariance glues it to a
holomorphic map to $H_{\kappa(s)s^n,a(s)}$.
It never reaches the excluded origin: that would require
$x=v=0$, which is the fixed point $(0,s)$ removed with the inner ball.
All zero-factor curves of $G_s$ have
$v=-s/a(s)^j$, so they lie in the one axis orbit $[-s]$.
The original modification also contracts over $v=-s$.
At that representative,
\[
 G_s(-s)=\prod_{j\ge1}(1-a(s)^j)^n\ne0.
\]
Hence locally the map is precisely $\pi_n$ followed by a coordinate
unit. Away from that orbit it is a local biholomorphism of degree
one. Compactness makes it a proper modification, and the displayed
local model identifies it as the stated $n$ successive blow-ups.
The developing map $\Psi_s$ is defined on the punctured parameter disc $s\ne0$.

\item Each blow-up of a point on the strict elliptic curve subtracts that
point from its normal line. Following all $n$ stages gives
\[
 N_s=L_{\kappa(s)s^n}\otimes\OO_{\mathcal E_s}(-n[-s]).
\]
The multiplicity $n$ records total exceptional transforms.
The reduced final exceptional locus meets the strict axis only once.
One may check the formula directly: the standard theta factor
$\theta_a(av)=-v^{-1}\theta_a(v)$ shows that $\OO([p])$
has factor $-p/v$. At $p=-s$ the tensor product above has factor
\[
 \kappa(s)s^n(v/s)^n=\kappa(s)v^n,
\]
which is equivalent to $\kappa(s)a(s)^n v^n$ by a nonzero monomial
gauge. This checks the signs and degrees independently of the
successive-blow-up formula.

\item The divisor $n[-s]$ cannot extend. For constant $a$, choose two
distinct classes $[v_1],[v_2]\in\C^*/a^\Z$.
The sequences $s_j=-a^jv_i$ tend to zero, whereas the corresponding
divisors are constantly $n[v_i]$ and have different limits in
$\operatorname{Sym}^n(E_a)$.
For varying $a(s)$, choose the holomorphic logarithm
$\tau_a(s)=\log a(s)/(2\pi i)$ on the disc.
The Abel--Jacobi coordinate of $n[-s]$ is
$n\log(-s)/(2\pi i)$ modulo $\Z+\tau_a(s)\Z$.
On a positive parameter loop it winds $n$ times around a primitive
cycle in the smooth elliptic family. An extending section over the
disc would have zero winding. A finite base change of order $N$
multiplies this nonzero integer by $N$.
After removing a vertical Cartier elliptic component, an extending
divisor would still have a proper horizontal part of
degree $n$, finite over the disc, defining a continuous
section of the relative symmetric power, which has just been ruled
out. Nonproper exceptional divisors in deformations of Kato surfaces
are a classical phenomenon, studied by
\citet[arXiv version~2, Section~1.2, PDF pp.~3--5]{DlousskyTeleman2012AB}; the calculation here
identifies the precise weighted divisor in this marked family.

\item This also explains the scope of ordinary elementary modifications.
If a degree-zero line $L_m$ from \ref{push:src:B.4.8} is changed by point blow-ups,
then its new normal is $L_m\otimes\OO(-D_{\rm exc})$.
If both that normal and the weighted Cartier correction extended on
a semistable good-reduction interface, their tensor product would
extend $L_m$, contradicting \ref{push:src:B.4.8}. In this Kato family the correction is exactly $D_{\rm exc}=n[-s]$,
whose nonextension was proved above.

\item Finally let $E=(2,2)\subset\PP^3$ be smooth and set
$L=\OO_E(2)$, of degree eight. Its normal in $\PP^3$ is
$L\oplus L$. For $A=\operatorname{Bl}_E\PP^3$, the exceptional
divisor is $D=E\times\PP^1$ and, in the lines convention,
\[
 N_{D/A}=L\boxtimes\OO_{\PP^1}(-1).
\]
A bare Hopf-product cap has normal $M\boxtimes\OO$ with
$\deg M=0$. The required triple-divisor line would then be
$(L\otimes M)^{-1}\boxtimes\OO(1)$, which has no nonzero section
because its elliptic degree is $-8$. Thus the required triple-divisor line has no effective section
on this interface.

\item For the parabolic surface with $n=8$, changing $\kappa$ varies the
flat factor through all of $\operatorname{Pic}^0(E)$:
constant factors are parametrized by $\C^*/a^\Z$.
Choose its axis normal to equal $L^{-1}$.
In $B=S_8\times\PP^1$ the divisor $D=E\times\PP^1$ then has normal
$L^{-1}\boxtimes\OO$. With triple section $T=E\times\{p\}$,
\[
 N_{D/A}\otimes N_{D/B}\otimes\OO_D(T)
   =(L\boxtimes\OO(-1))\otimes
     (L^{-1}\boxtimes\OO)\otimes(\OO\boxtimes\OO(1))
   \simeq\OO_D.
\]
This proves precisely the first-order matching claimed.
\end{enumerate}
\end{proof}

\begin{proposition}[Globalization at the contracted parabolic elliptic axis]\hypertarget{tag-030F}{}\label{push:src:B.4.11}
\StatementRecord{push:src:B.4.11}{proved}{030F}{\SeriesRef{push:src:B.4.10}; \citep[Section~6, Hauptsatz~IIa, p.~59]{Grauert60}; \citet[Lemma~1, p.~93]{Manetti1991D}; \citet[Section~3.2, Proposition~3.3 and the discussion preceding
Proposition~3.4]{DlousskyTeleman2020D}}

Let $S=S_8(a,\kappa)$ be the central parabolic surface in \ref{push:src:B.4.10},
let $E\subset S$ be its elliptic axis, and put
$N=N_{E/S}$, so $\deg N=-8$.
The contraction $\pi:S\to Y$ of $E$ is a compact normal surface
with a single isolated singular point $p$.
Its germ at $p$ is the vertex of the affine cone over
$(E,N^{-1})$.  In particular,
\[
 H^2(Y,\Theta_Y)=0.
\]
The restriction morphism from a semiuniversal deformation of $Y$
to a semiuniversal deformation of $(Y,p)$ is regular.
Consequently every prescribed analytic one-parameter deformation
of this isolated singularity is induced, after shrinking its
parameter disc, by a proper deformation of the compact surface $Y$.
If the prescribed local deformation smooths the singularity, the
global deformation has smooth fibres away from the central value.
\end{proposition}


\begin{proof}
\begin{enumerate}[label=\arabic*.,leftmargin=*]
\item At the central parameter the Kato gluing of \ref{push:src:B.4.10} is
\[
 (x,v)\longmapsto(\kappa a^8xv^8,av).
\]
It is exactly linear in $x$.  Thus a neighbourhood of $E$ is a
neighbourhood of the zero section in the negative line bundle
$N$, with the indicated factor of automorphy.
Its contraction is the cone associated with the section ring
\[
 R(E,N^{-1})=\bigoplus_{m\ge0}H^0(E,N^{-m}).
\]
A line bundle of degree eight on an elliptic curve is very ample
and normally generated.  The cone is normal, its only singularity
is its vertex, and the blow-up of the vertex maximal ideal is
$\operatorname{Tot}(N)$ near its zero section.
One can see the last assertion directly: the degree-one sections
generate the section ring and have no common zero on $E$;
their ratios on the blow-up recover the projective embedding of
$E$, and their common radial coordinate is the coordinate in $N$.
In particular, with $I=\mathcal O_S(-E)$ and $\mathfrak m_p$ the
vertex maximal ideal,
\[
 \mathfrak m_p\mathcal O_S=I.
\]
The local contraction glues to the identity on $S\setminus E$,
giving the stated proper contraction of the compact surface.

\item Write $C$ for the disjoint rational cycle of the parabolic surface.
The meromorphic two-form
$d\log x\wedge d\log v$ is invariant under the displayed
monomial gluing.  In the toric blow-up charts it has simple poles
along precisely the axis and the rational cycle, and no zeros.
Consequently
\[
 K_S=\mathcal O_S(-E-C).
\]
Also $b_1(S)=1$, by \ref{push:src:B.4.10}.
Put $\mathcal F=\Theta_S(-\log E)$, the locally free sheaf of
vector fields tangent to $E$.
Serre duality and the displayed canonical divisor give
\[
 H^2(S,\mathcal F)^*
 =H^0\bigl(S,\Omega^1_S(\log E)\otimes K_S\bigr)=0.
\]
Indeed there are sheaf injections
\[
 \Omega^1_S(\log E)(-E-C)
       \lhook\joinrel\longrightarrow\Omega^1_S(-C)
       \lhook\joinrel\longrightarrow\Omega^1_S.
\]
Locally the first sends the logarithmic generators
$dx/x,dy$, multiplied by $x$, to the ordinary forms $dx,xdy$.
Finally $H^0(S,\Omega^1_S)=0$ follows already from $b_1(S)=1$.
A holomorphic one-form $\alpha$ on a compact
complex surface is closed: Stokes' theorem applied to
$\alpha\wedge\overline{d\alpha}$ makes the integral of the
positive volume form associated with
$d\alpha\wedge\overline{d\alpha}$ vanish.
If $\alpha\ne0$, its real and imaginary parts give independent
real de Rham classes.  An exact real linear combination would
be the differential of a global real pluriharmonic function,
which is constant on the compact surface.  This would give
$b_1(S)\ge2$, a contradiction.

\item The restriction of $\mathcal F$ to $E$ is the Atiyah extension
of $N$:
\[
 0\longrightarrow\mathcal O_E
 \longrightarrow\mathcal A_N:=\mathcal F|_E
 \longrightarrow\Theta_E\longrightarrow0.
\]
In normal coordinates related by $x_j=g_{ij}(v)x_i$, the
vertical Euler field is $x_i\partial_{x_i}$, and the changes of
horizontal lifts have cocycle $d\log g_{ij}$.
Thus this extension has class $c_1(N)$ in
$H^1(E,\Omega^1_E)$, which is nonzero because $\deg N=-8$.
Since $\Theta_E\simeq\mathcal O_E$, the connecting map
$H^0(E,\Theta_E)\to H^1(E,\mathcal O_E)$ is an isomorphism.
The same cohomology sequence therefore gives the natural
isomorphism
\[
 H^1(E,\mathcal A_N)\xrightarrow{\ \sim\ }
 H^1(E,\Theta_E)\simeq\C.
\]

\item The successive quotients for the formal neighbourhood of $E$ are
\[
 I^m\mathcal F/I^{m+1}\mathcal F
     =\mathcal A_N\otimes N^{-m}\qquad(m\ge0).
\]
For $m\ge1$ both line bundles in the induced extension have
degree $8m>0$ on the elliptic curve, so their first cohomology
vanishes.  It follows that all restriction maps
\[
 H^1\bigl(E,\mathcal F/I^{m+1}\mathcal F\bigr)
 \longrightarrow H^1(E,\mathcal F/I^m\mathcal F)
 \quad(m\ge1)
\]
are isomorphisms.  Their common value is
$H^1(E,\mathcal A_N)\simeq\C$.
The analytic theorem on formal functions
\citep[Section~6, Hauptsatz~IIa, p.~59]{Grauert60}
identifies this inverse limit with the completed stalk of
$R^1\pi_*\mathcal F$ at $p$.
This coherent sheaf is supported at the single point $p$,
so its stalk is a finite-length module and equals its completion.
We have proved
\[
 R^1\pi_*\mathcal F\simeq\C_p,
 \qquad
 H^0(Y,R^1\pi_*\mathcal F)
    \xrightarrow{\ \sim\ }H^1(E,\Theta_E).
\]
The second map is the natural restriction followed by the
Atiyah-extension map just computed.

\item There is also an equality of sheaves
\[
 \pi_*\mathcal F=\Theta_Y.
\]
To check it at $p$, a holomorphic derivation of the cone
integrates locally to a holomorphic flow: embed the cone in a
polydisc, lift its coefficient functions, and use the ordinary
holomorphic vector-field flow, which preserves the defining ideal.
The flow preserves the unique singular point, hence its maximal
ideal.  Functoriality of the blow-up of that ideal lifts the flow
to the resolution, preserving the exceptional curve.
Its derivative is therefore a section of $\mathcal F$.
Conversely a section of $\mathcal F$ differentiates pullbacks
of holomorphic functions on $Y$ and descends to a derivation,
because $\pi_*\mathcal O_S=\mathcal O_Y$.
Both operations agree off $E$ and are inverse there; uniqueness
extends the equality over $E$.

\item The global-to-local cohomology map is surjective for a concrete
reason.  Write $a=e^{2\pi i\tau}$ with $\operatorname{Im}\tau>0$,
and vary $\tau$ to $\tau+u$ in the explicit Kato construction,
keeping its translation parameter zero.
The same uniform collar argument as in \ref{push:src:B.4.10} gives a proper
holomorphic deformation of the pair $(S,E)$.
Its Kodaira--Spencer class belongs to $H^1(S,\mathcal F)$.
Under the natural composite
\[
 H^1(S,\mathcal F)\longrightarrow
 H^0(Y,R^1\pi_*\mathcal F)
 \xrightarrow{\ \sim\ }H^1(E,\Theta_E)
\]
it is the Kodaira--Spencer class of
$\C/(\Z+(\tau+u)\Z)$, which is nonzero.
For example, differentiating its second deck transformation
$z\mapsto z+\tau+u$ gives the usual nonzero period cocycle.
This argument works also at the elliptic curves with extra
automorphisms.
The target is one dimensional, so the first arrow is onto.

\item The degree-two part of the Leray sequence is
\[
 H^1(S,\mathcal F)\longrightarrow
 H^0(Y,R^1\pi_*\mathcal F)\longrightarrow
 H^2(Y,\pi_*\mathcal F)\longrightarrow H^2(S,\mathcal F).
\]
The first map is surjective and the last group is zero.
Using $\pi_*\mathcal F=\Theta_Y$ proves
$H^2(Y,\Theta_Y)=0$.

\item For a compact complex space with isolated singularities,
$H^2(\Theta)=0$ makes the restriction from global to local
deformations formally smooth.
Take a finite Stein cover with one small
neighbourhood of $p$ and all other members smooth.
For a small extension of Artinian bases, prescribe the lifted
deformation on the singular member and take smooth local
lifts on the other members.  Their pairwise identifications
lift on the smooth Stein intersections.  The failure of the
cocycle identity on triple intersections is a \v Cech two-cocycle
in $\Theta_Y$ tensored with the square-zero kernel.
Its class vanishes by the proved $H^2$-vanishing; correcting
the identifications by a one-cochain glues the desired lift.
This proves the relative infinitesimal lifting property.

\item The analytic semiuniversal deformation theory turns this
property into regularity of the induced morphism of analytic
base germs.  The general criterion is
\citet[Lemma~1, p.~93]{Manetti1991D}, stated for reduced compact
complex spaces with isolated singularities and proved by the
local-to-global spectral sequence of the cotangent complex.
It is also used in
\citet[arXiv version~2, Section~3.2, Proposition~3.3 and the following discussion, p.~11]{DlousskyTeleman2020D}.
A regular morphism of analytic
germs is locally a projection from a product with a smooth
germ.  Consequently a specified analytic disc in the local
deformation base lifts to the global base.
The global family is proper by the compact deformation
construction.  Away from $p$ its central fibre is smooth,
and smoothness is open; compactness permits one common
shrinking of the disc.  A local smoothing at $p$ therefore
gives smooth global noncentral fibres, as asserted.
\end{enumerate}
\end{proof}

\begin{proposition}[A proper quadric--parabolic surface smoothing]\hypertarget{tag-030G}{}\label{push:src:B.4.12}
\StatementRecord{push:src:B.4.12}{proved}{030G}{\SeriesRef{push:src:B.4.11}, \SeriesRef{push:src:B.4.10}; \citep[Remark~7.6(iii), p.~46; Section~9, pp.~63--65]{Pinkham1974Root}}

Let $S=S_8(a,\kappa)$ be as in \ref{push:src:B.4.11}, with elliptic axis $E$
and normal line $N$ of degree $-8$.
There is an embedding $E\hookrightarrow Q=\PP^1\times\PP^1$
as a smooth divisor of bidegree $(2,2)$ with
$N_{E/Q}\simeq N^{-1}$.
There is a proper flat map from a smooth complex threefold
\[
 \varpi:\mathcal V\longrightarrow\Delta
\]
with smooth noncentral fibres and reduced simple normal-crossings
central fibre exactly $S\cup_E Q$.  Its two normal lines are
\[
 N_{S/\mathcal V}=\OO_S(-E),\qquad
 N_{Q/\mathcal V}=K_Q=\OO_Q(-2,-2).
\]
The family is obtained by globalizing a specified cone smoothing
and blowing up its vertex over the original parameter disc.

For the elliptic-quartic configuration of \ref{push:src:B.4.10}, the three
pairwise normal-line products, including the triple-curve factors,
can consequently all be made trivial using
\[
 A=\operatorname{Bl}_E\PP^3,\qquad
 B=S\times\PP^1,\qquad C=\mathcal V.
\]
Here $C$ is the constructed proper surface family over an open
disc. Its continuous surjection onto that disc proves that $C$ is
noncompact.
\end{proposition}


\begin{proof}
\begin{enumerate}[label=\arabic*.,leftmargin=*]
\item Put $L=N^{-1}$, of degree eight.  Multiplication by two on the
elliptic Picard variety is surjective, so choose a degree-four
line bundle $M$ with $M^{\otimes2}=L$.
Choose degree-two line bundles $M_1,M_2$ with
$M_1\otimes M_2=M$ and $M_1\not\simeq M_2$.
Each complete linear system is a basepoint-free double cover of
$\PP^1$.  The product of the two maps is an embedding in $Q$.
Indeed two distinct points identified by both maps would give
$M_1\simeq\OO_E(p+q)\simeq M_2$.
A common ramification point would similarly give
$M_1\simeq\OO_E(2p)\simeq M_2$.
Thus the product is injective and immersive; properness makes it a
closed embedding.  Its image has bidegree $(2,2)$, and
\[
 \OO_Q(E)|_E=(M_1\otimes M_2)^{\otimes2}=L.
\]
This also allows us to start with the fixed quadric and elliptic
quartic in \ref{push:src:B.4.10}: the normal line there was chosen to be the
inverse of this same $L$.

\item Use the anticanonical embedding of $Q$ in $\PP^8$ and write
\[
 R_Q=\bigoplus_{d\ge0}H^0(Q,\OO_Q(2d,2d)),\qquad
 Z=\operatorname{Spec}R_Q.
\]
The ring is generated in degree one and is normal.
One direct verification uses the bihomogeneous monomials:
every monomial of bidegree $(2d,2d)$ is a product of $d$
monomials of bidegree $(2,2)$, and the associated monoid is the
intersection of a rational polyhedral cone with its full lattice.
The punctured cone is the complement of the zero section in
$\OO_Q(-2,-2)$, hence is smooth.
Choose the section $e\in H^0(Q,\OO_Q(2,2))$ cutting out $E$.
The degree-one function $t=e$ on $Z$ defines a flat map
$Z\to\mathbb A^1$: its coordinate ring is torsion-free over
$\C[t]$, a principal ideal domain.
For every $d\ge0$, the sequence
\[
 0\longrightarrow\OO_Q(2d-2,2d-2)
 \xrightarrow{\ e\ }\OO_Q(2d,2d)
 \longrightarrow L^{\otimes d}\longrightarrow0
\]
has vanishing first cohomology in its left term.
For $d=0$ this follows from the K\"unneth formula for
$\OO_Q(-2,-2)$; for $d\ge1$ it follows from the same formula
and nonnegative degrees.
Taking sections identifies
\[
 R_Q/(e)=\bigoplus_{d\ge0}H^0(E,L^{\otimes d}).
\]
The central surface is therefore exactly the affine cone over
$(E,L)$, including its scheme structure.
For $t\ne0$, projection from the cone vertex identifies the
fibre with $Q\setminus E$: over each point where $e$ is nonzero
there is exactly one radial coordinate on which $e$ has value $t$.
In particular this is an explicit smoothing of the isolated
singularity of \ref{push:src:B.4.11}.
It is the affine version of Pinkham's cone construction
\citep[Remark~7.6(iii), p.~46; Section~9, pp.~63--65]{Pinkham1974Root}.

\item Blowing up the vertex of $Z$ gives
$\operatorname{Tot}(\OO_Q(-2,-2))=\operatorname{Tot}(K_Q)$.
This can be checked on the blow-up charts given by ratios of the
degree-one homogeneous coordinates; these recover a point of
$Q$, with the remaining common factor the tautological fibre
coordinate $\xi$.
The map to the parameter line is precisely
\[
 t=\xi e.
\]
Its central divisor is the union of the zero section $Q$ and
the restriction $\operatorname{Tot}(K_Q|_E)$, each with
multiplicity one.  Near $E$ the equation is $t=\xi z$, where
$z$ is a local defining function of $E$ in $Q$.
Thus the total space is smooth and the central divisor has
simple normal crossings.

\item Now contract $E$ in the compact surface $S$ as in \ref{push:src:B.4.11}.
The regular global-to-local deformation map in that proposition
lifts the specified analytic germ $Z\to\mathbb A^1$ to a proper
deformation $\mathcal Y\to\Delta$ of $Y=S/E$.
The lifted deformation agrees analytically with this cone
smoothing in a neighbourhood of its singular point.
Away from that neighbourhood the central surface is smooth;
openness of smoothness and properness allow a common shrinking
of $\Delta$ for which the family there is a submersion.
Blow up the single cone vertex in $\mathcal Y$.
The resulting map $\mathcal V\to\Delta$ remains proper.
Near the exceptional divisor its local model is the smooth
space just computed; elsewhere it is unchanged and smooth.
Its central strict transform is exactly $S$, because the
maximal-ideal blow-up of the degree-eight elliptic cone is its
given resolution $\operatorname{Tot}(N)$ near $E$.
The other component is $Q$, and their intersection is the original
marked elliptic curve.  The local equation $t=\xi z$ proves
flatness and reducedness near the intersection; at a smooth
point of either component $t$ is a local coordinate.
The noncentral fibres are the smooth compact fibres of
$\mathcal Y$.
This proves every asserted properness and smoothness property
for the surface family.

\item The exceptional zero section has normal $K_Q$ in its local
total-space model.  Since $S+Q$ is the principal divisor of $t$,
restriction to $S$ gives
\[
 \OO_{\mathcal V}(S)|_S
   =\OO_{\mathcal V}(-Q)|_S=\OO_S(-E).
\]
These are the two asserted normal lines.  They also show directly
that the two branch normals along $E$ multiply to the trivial line.

\item Realize $Q$ as a smooth quadric in $\PP^3$ and $E$ as its
intersection with a second quadric.  Put $L=\OO_E(2)$.
Then $N_{E/\PP^3}=L\oplus L$, so the exceptional divisor
of $A=\operatorname{Bl}_E\PP^3$ is
$D_{AB}=E\times\PP^1$, with normal
$L\boxtimes\OO_{\PP^1}(-1)$.
The strict transform of $Q$ is still $Q$, because the centre is
Cartier on it, and its normal in $A$ is
\[
 N_{Q/A}=\OO_Q(2)\otimes\OO_Q(-E)=\OO_Q.
\]
It meets $D_{AB}$ in the constant section $E\times\{p\}$
corresponding to that member of the pencil of quadrics.
In $B=S\times\PP^1$, use the divisors
$D_{BA}=E\times\PP^1$ and $D_{BC}=S\times\{p\}$.
Their normals are $L^{-1}\boxtimes\OO$ and $\OO_S$.
In $C=\mathcal V$, use its central divisors $Q$ and $S$.
The three normal products with the triple contribution are then
\[
\begin{array}{c|c}
 \text{double surface}&
 N_{D/\text{first branch}}\otimes
 N_{D/\text{second branch}}\otimes\OO_D(\text{triple curve})\\ \hline
 E\times\PP^1&
 (L\boxtimes\OO(-1))\otimes(L^{-1}\boxtimes\OO)
       \otimes(\OO_E\boxtimes\OO(1))=\OO\\
 Q&\OO_Q\otimes K_Q\otimes\OO_Q(E)=\OO_Q\\
 S&\OO_S\otimes\OO_S(-E)\otimes\OO_S(E)=\OO_S.
\end{array}
\]
The three displayed evaluation maps are the required pairwise
normal-line trivializations. The third space maps surjectively and
properly to an open disc, which proves its noncompactness.

\item The affine smooth fibre $Q\setminus E$ has
\[
 H_1(Q\setminus E;\Z)=\Z/2,\qquad
 H_2(Q\setminus E;\Z)=\Z^3.
\]
Indeed the Thom sequence of the pair $(Q,Q\setminus E)$ has
intersection map $H_2(Q;\Z)\to H_0(E;\Z)$ given by
$(a,b)\mapsto2a+2b$.  Its cokernel is $\Z/2$ and its
kernel is $\Z$; the preceding term $H_1(E;\Z)=\Z^2$
injects into $H_2(Q\setminus E;\Z)$.
This exact extension of free groups splits abstractly.
The map $H_4(Q;\Z)\to H_2(E;\Z)$ is an isomorphism,
so $H_3(Q\setminus E;\Z)=0$, and its Euler number is four.
The Euler number is also obtained independently by projection
to either ruling of $Q$: away from the four branch values of
$E\to\PP^1$, the fibre is $\C^*$ and has Euler number zero;
at each branch value it is $\C$ and has Euler number one.
The compact noncentral surface has Euler number
$\chi(S)+\chi(Q)-2\chi(E)=8+4=12$ by the local $\xi z=t$ retraction.
\end{enumerate}
\end{proof}

\begin{remark}[Euler characteristics of the quadric--parabolic smoothing]
\hypertarget{tag-PT14}{}\label{I:quadric-parabolic-euler}
\StatementRecord{I:quadric-parabolic-euler}{proved}{PT14}{\cref{push:src:B.4.10,push:src:B.4.12}}
For the compact noncentral surface $V_t$ of Proposition~\ref{push:src:B.4.12},
\[
 \chi(V_t)=\chi(S_8)+\chi(Q)-2\chi(E)=8+4=12,
 \qquad \chi(V_t,\mathcal O_{V_t})=1.
\]
The first formula is the local $\xi z=t$ smoothing formula proved there. Mayer--Vietoris for the two spherical boundaries in the Kato construction identifies the middle homology and its intersection form with the eight negative exceptional classes. Thus $\chi(S_8)=8$ and the signature is $-8$, so $c_1(S_8)^2=2\chi(S_8)+3\sigma(S_8)=-8$ and analytic Riemann--Roch gives $\chi(S_8,\mathcal O_{S_8})=0$. Since $\chi(Q,\mathcal O_Q)=1$ and $\chi(E,\mathcal O_E)=0$, the normalization sequence for $S_8\cup_E Q$ and flatness give the second formula.
\end{remark}

\begin{proposition}[The obstruction to a marked period lift]\hypertarget{tag-030H}{}\label{push:src:B.4.13}
\StatementRecord{push:src:B.4.13}{proved}{030H}{\SeriesRef{push:src:B.4.10}, \SeriesRef{push:src:A.1.2}, \SeriesRef{push:src:B.4.8}}

Fix $0<|\alpha|,|a|<1$, and put
\[
 H=(\C^2\setminus\{0\})/
       \langle(x,v)\mapsto(\alpha x,av)\rangle,
       \qquad E=\{x=0\}\simeq\C^*/a^\Z.
\]
Let $h$ be the automorphism induced by
$(x,v)\mapsto(c x,bv)$, where $b,c\in\C^*$.
Suppose $p:Y\to H$ is a finite composition of point blow-ups,
allowing infinitely near centres, and at least one point of $E$
is an image of its exceptional locus. If $[b]\in\C^*/a^\Z$ has
infinite order, there is no holomorphic automorphism of $Y$
which agrees, on the dense unmodified core, with the marked
automorphism $h$.

The same obstruction holds for a finite collection of such marked
modifications if the period permutes the collection and
acts as $h$ on each marked core.
More generally, if a compact smooth surface contains a preserved
marked elliptic curve $E$ of nonzero normal degree $d$, every
holomorphic automorphism restricting to translation by $b$ on $E$
must satisfy $d\,b=0$ in $E$.
The latter assertion also excludes a finite cycle of such preserved
marked cap axes when the cycle composite is an infinite-order
translation. These assertions apply to the extension
of the second cusp period to the finite Hopf
modifications in \ref{push:src:B.4.10}, for generic nonzero parameters, after every
finite base change. The period action in these assertions is a holomorphic action on
the full finite modifications with the specified marked core.
\end{proposition}


\begin{proof}
\begin{enumerate}[label=\arabic*.,leftmargin=*]
\item Let $S\subset H$ be the image of the exceptional curves of $p$.
It is finite. Repeated infinitely near blow-ups change the curves
over a point but do not make $S$ infinite.
If $\widetilde h$ were a lift, the maps
\[
 p\widetilde h,\ hp:Y\longrightarrow H
\]
would agree on the dense core and hence everywhere, by the identity
principle for holomorphic maps to the separated complex surface $H$.
The image under $\widetilde h$ of a curve contracted by $p$ is
therefore again contracted by $p$. Applying the same argument to
$\widetilde h^{-1}$ proves that the total exceptional locus is
preserved. Thus $h(S)=S$.

\item The automorphism $h$ preserves $E$, where it is translation by
$[b]$. Consequently the nonempty finite set $S\cap E$ would be
invariant under that translation. The orbit of any one of its
points is finite, so for some $k>0$,
\[
 [b^k v]=[v],\qquad\text{equivalently } b^k\in a^\Z.
\]
This says exactly that $[b]$ is torsion, a contradiction.
In particular, no choice of multiplicities or infinitely near
centres changes the obstruction: only their nonempty finite
support on $E$ is needed.

\item For finitely many modifications $p_i:Y_i\to H$, a deck map has biholomorphisms
$Y_i\to Y_{\sigma(i)}$ satisfying the same core comparison.
The identity principle again gives equivariance everywhere.
The union of the finitely many exceptional image sets is
$h$-invariant; its intersection with $E$ gives the same contradiction.
Equivalently, around a cycle of length $\ell$ the map $h^\ell$
would lift to a single modification, while $[b^\ell]$ still has
infinite order.

\item For the more general assertion, an automorphism preserving $E$ induces through its derivative in the normal
direction an isomorphism
\[
 t_b^*N_{E/S}\simeq N_{E/S}.
\]
For a line bundle $N$ of degree $d$ on an elliptic curve, the
homomorphism $b\mapsto t_b^*N\otimes N^{-1}$ is multiplication
by $-d$ under the usual identification $\operatorname{Pic}^0(E)=E$.
Indeed write $N=\OO_E(\sum_j n_j[p_j])$ with
$\sum_j n_j=d$; the Abel--Jacobi class of its translated divisor
minus the original one is $-d b$.
Thus preservation of $N$ forces $d b=0$.
Taking a power proves the finite-cycle version, since its composite
preserves one axis and its nonzero-degree normal line.
In particular a nonzero-degree parabolic axis cannot carry the
marked generic period translation even if no finite Hopf
modification is specified. The original finite-support argument
also covers modifications where that normal-degree test alone
would give no information.

\item The infinite-order translation at the cusp is computed as follows.
The periods of \ref{push:src:A.1.2} have the form
\[
 g_1:(Z,W)\longmapsto(q c_1 Z,c_2W),\qquad
 g_2:(Z,W)\longmapsto(c_3Z,r^{12}q c_4 W),
\]
where the $c_i$ are nonvanishing holomorphic units near the cusp.
Choose $W$ or $W^{-1}$ so that
$a=c_2^{\varepsilon}$ has modulus less than one,
$\varepsilon\in\{1,-1\}$.
This is possible at the marked cusp because its limiting
unit has modulus different from one, as computed in \ref{push:src:B.4.8}.
For $q\ne0$ small, the partial $g_1$ quotient of $(\C^*)^2$
is the dense core of the diagonal Hopf surface with multipliers
$(q c_1,a)$. The second period induces the above $h$, with
\[
 b(q,r)=(r^{12}q c_4(q,r))^{\varepsilon}.
\]
Fix a sufficiently small nonzero $q$. For each $k>0$ and
$j\in\Z$, the equation $b(q,r)^k=a(q)^j$ has a discrete
solution set in the punctured $r$-disc: its left-hand side has
a zero or pole of nonzero order at $r=0$, while the right-hand
side is a nonzero constant in $r$.
Hence outside a countable union of discrete sets, $[b(q,r)]$
has infinite order. Under a finite base change in $r$ the
nonzero order is multiplied by its ramification index and
the same conclusion holds. A finite base change in $q$ does
not affect this fixed-nonzero-$q$ argument.

\item To compare with \ref{push:src:B.4.10}, write $q c_1=\kappa(s)s^n$ after
the chosen finite cover in the cusp variable. Its generic Kato
fibre is precisely a finite Hopf modification with nonempty
axis centre $[-s]$. A holomorphic extension of the second
period to that whole fibre would therefore contradict the result
just proved. Allowing finitely many translated copies does not
remove the contradiction, by the finite-collection assertion.
\end{enumerate}
\end{proof}

\begin{lemma}[The quotient of indexed cores]
\hypertarget{tag-PT15}{}\label{I:indexed-core-quotient}
\StatementRecord{I:indexed-core-quotient}{proved}{PT15}{the explicit quotient and closure maps in the proof}
Let $h$ be an automorphism of $H$, and let modifications $p_j:Y_j\to H$, indexed by $j\in\mathbb Z$, have isomorphisms $\widetilde h_j:Y_j\to Y_{j+1}$ satisfying $p_{j+1}\widetilde h_j=hp_j$. On an $h$-invariant open set $H^\circ$ over which every $p_j$ is an isomorphism, the shift quotient of the disjoint union of their cores is biholomorphic to $H^\circ$ by
\[
 [(j,y)]\longmapsto h^{-j}p_j(y).
\]
Under this isomorphism, its natural map to the orbit space $H^\circ/\langle h\rangle$ is the ordinary quotient map. Moreover, two embedded compact subspaces of a Hausdorff space that share a subset open and dense in both have the same image.
\end{lemma}
\begin{proof}
The displayed map is invariant under $(j,y)\mapsto(j+1,\widetilde h_jy)$ because $h^{-j-1}p_{j+1}\widetilde h_j=h^{-j}p_j$. Every shift orbit meets the zeroth core exactly once, so the inverse is $x\mapsto[(0,p_0^{-1}x)]$. These maps are holomorphic on the component charts. Composing with $H^\circ\to H^\circ/\langle h\rangle$ gives the asserted orbit-space map. For the last assertion, compact images in a Hausdorff space are closed; each of the two images is consequently the closure of their shared dense subset.
\end{proof}

\begin{proposition}[Tangent representations and associated bundles]\hypertarget{tag-030I}{}\label{push:src:B.4.14}
\StatementRecord{push:src:B.4.14}{proved}{030I}{\SeriesRef{push:src:B.4.11}; \citet[Section~2, equations~(1)--(8), p.~4]{Fujiki2009D}; \citet[Theorem~1.1, pp.~1--2]{Fujiki2009D}; \citet[Theorem~(4.6), p.~566]{AtiyahSinger1968IIIE}}

For the explicit parabolic surface $S=S_8(a,\kappa)$ and its
elliptic contraction $\pi:S\to Y$ of \ref{push:src:B.4.11}, the genuine action
$\lambda\cdot(x,v)=(\lambda x,v)$ induces trivial representations
on the following cohomology groups:
\[
 H^0(S,\Theta_S(-\log E))=\C,\qquad
 H^1(S,\Theta_S(-\log E))=\C^9,
\]
\[
 H^0(Y,\Theta_Y)=\C,\qquad H^1(Y,\Theta_Y)=\C^8.
\]
The two second cohomology groups vanish.

Let $p:Z_L\to\PP^1$ be the compact holomorphic $Y$-bundle
associated to any holomorphic line bundle $L$ and this
$\C^*$ action.  Put
$\mathcal T=\Theta_{Z_L/\PP^1}$.
For every $b\in\PP^1$ one has
\[
 R^0p_*\mathcal T=\mathcal O_{\PP^1},\qquad
 R^1p_*\mathcal T=\mathcal O_{\PP^1}^{\oplus8},\qquad
 R^2p_*\mathcal T=0,
\]
and
\[
 H^2(Z_L,\mathcal T)=0,
 \qquad
 H^2\bigl(Z_L,\mathcal T\otimes
                   p^*\mathcal O_{\PP^1}(-b)\bigr)=0.
\]
The singular locus of the associated bundle is the section induced
by the cone vertex of $Y$; the displayed groups are the relative
tangent cohomology and its specified fibre twist.
\end{proposition}


\begin{proof}
\begin{enumerate}[label=\arabic*.,leftmargin=*]
\item The action commutes with the central Kato gluing and hence
acts on the whole compact surface and its contraction.
Its toric coordinate description is also given in
\citet[Section~2, equations~(1)--(8), p.~4]{Fujiki2009D}.
For any $\lambda\ne1$, its fixed locus consists of the elliptic
curve $E$ and the eight nodes of the rational cycle.
To check that there are no further points, on the dense torus
an equality with a Kato deck iterate would require $a^jv=v$;
since $0<|a|<1$, this forces $j=0$, and then $\lambda=1$.
The action on each rational component is nontrivial and fixes
only its two endpoints.  In the unfolded fan the consecutive
rays are $(j,1)$ and $(j+1,1)$.
Their dual characters evaluate on the scaling cocharacter
$(1,0)$ as $-1,+1$.
Thus the two tangent weights at every cycle node are
$\lambda^{-1},\lambda$.

\item The space of global vector fields is described in
\citet[Theorem~1.1, pp.~1--2]{Fujiki2009D}.
Write the contracted Kato germ as
\[
 F(x,v)=(Kxv^8,av),\qquad K=\kappa a^8.
\]
A global vector field restricts on the outer spherical annulus
to $P(x,v)\partial_x+Q(x,v)\partial_v$.
Hartogs extension gives holomorphic coefficients on the ball.
Compatibility with the collar gluing, followed by analytic
continuation, gives the exact germ identity
\[
 V(F)=DF\cdot V.
\]
Expand $Q(x,v)=\sum_{i\ge0}x^iQ_i(v)$.
The second component says
\[
 K^iv^{8i}Q_i(av)=aQ_i(v).
\]
For $i>0$ the lowest nonzero $v$-order on the two sides would
differ by $8i$, so $Q_i=0$.
For $i=0$, power-series coefficients give $Q_0(v)=cv$.
The constant-in-$x$ coefficient of the first component also
vanishes by the same lowest-order argument.
Its linear coefficient satisfies
$P_1(av)=P_1(v)+8c$.
Evaluation at zero gives $c=0$, and then $P_1$ is constant.
For $i>1$ the remaining equations are
\[
 P_i(v)=K^{i-1}v^{8(i-1)}P_i(av),
\]
so again $P_i=0$.
Consequently $H^0(S,\Theta_S)=\C\,x\partial_x$.
This Euler field is tangent to $E$ and has weight zero.

\item Set $\mathcal F=\Theta_S(-\log E)$.
\ref{push:src:B.4.11} proved $H^2(S,\mathcal F)=0$.
Apply the holomorphic Lefschetz formula to a root of unity
$\lambda\ne1$, using its finite cyclic group action.
The applicable original statement is
\citet[Theorem~(4.6), p.~566]{AtiyahSinger1968IIIE}, which applies
to compact complex manifolds.
On $E$, the Atiyah extension $\mathcal F|_E$ has two
weight-zero terms and Chern character $2$; the normal line
has weight $\lambda$ and degree $-8$.
Since the tangent bundle of the elliptic curve has trivial
Todd class, the contribution of $E$ is
\[
 \int_E\frac{2}{1-\lambda^{-1}\exp(-c_1(N))}
   =-\frac{2\lambda^{-1}\deg N}
           {(1-\lambda^{-1})^2}
   =\frac{16\lambda}{(\lambda-1)^2}.
\]
At a cycle node $\mathcal F$ equals the ordinary tangent
bundle.  Its contribution is
\[
 \frac{\lambda+\lambda^{-1}}
      {(1-\lambda)(1-\lambda^{-1})}
   =-\frac{\lambda^2+1}{(\lambda-1)^2}.
\]
Adding the eight nodes gives the constant character
\[
 \sum_q(-1)^q\operatorname{Tr}
       (\lambda\mid H^q(S,\mathcal F))
 =\frac{16\lambda-8(\lambda^2+1)}{(\lambda-1)^2}
 =-8.
\]
Since $H^0$ is one weight-zero copy and $H^2=0$,
$H^1(S,\mathcal F)$ consists of nine trivial copies.
Equality on infinitely many roots of unity determines the
finite Laurent-polynomial character of the $\C^*$ representation.

\item Every map in the Leray sequence of \ref{push:src:B.4.11} is equivariant.
Its surjection
\[
 H^1(S,\mathcal F)\longrightarrow
 H^0(Y,R^1\pi_*\mathcal F)
       \simeq H^1(E,\Theta_E)=\C
\]
has a trivial one-dimensional target.
The kernel is $H^1(Y,\Theta_Y)$, so it is eight-dimensional
and trivial.  The equality in degree zero follows from
$\pi_*\mathcal F=\Theta_Y$; the second cohomology vanishing
was already proved in \ref{push:src:B.4.11}.

\item Finally form $Z_L$ by choosing line-bundle transition functions
and gluing products $Y\times U_i$ by their action on $Y$.
This is a holomorphic locally trivial bundle with compact
fibre over the compact curve $\PP^1$, hence is proper and
compact.  It has the section of isolated cone points as its
singular locus; that locus is one dimensional in the total
space.  On each product chart, coherent cohomology and base
change identify $R^qp_*\mathcal T$ with
$H^q(Y,\Theta_Y)\otimes\mathcal O_{U_i}$.
The gluing is exactly by the representations computed above.
Their triviality proves the displayed direct-image formulas.
The projection formula and the Leray sequence then give
\[
 H^2\bigl(Z_L,\mathcal T\otimes p^*M\bigr)
      \simeq H^1(\PP^1,M^{\oplus8})
\]
for $M=\mathcal O$ or $\mathcal O(-b)$, both of which have
zero first cohomology on $\PP^1$.
This proves the asserted reduction, including the twist
appropriate to retaining one whole fibre.

\end{enumerate}
\end{proof}

\subsection{A local insertion into the prescribed simple-torus family}

\begin{proposition}[Local insertion of vertical Hopf bundles]\hypertarget{tag-030J}{}\label{push:src:B.4.15}\label{push:B415}
\StatementRecord{push:src:B.4.15}{proved}{030J}{\SeriesRef{push:src:B.3.1}; \citep[expos\'e~19, \S5, pp.~14--15]{Cartan1952B_AB}; \citet[construction, p.~85; Definition~1.1, p.~86]{Fujimoto1990AB}}

Let $q_0$ be a regular point of the elliptic base, and let
$\Delta_z$, $z=q-q_0$, be a simply connected disc avoiding
$0,1,\infty$. After $w=r^6$, take the standard twelve-cycle
model
\[
 \mathcal X\longrightarrow\Delta_z\times\Delta_r
\]
of the prescribed local two-torus family. There is a proper
replacement $\mathcal X'\to\Delta_z\times\Delta_r$, with
a specified isomorphism to $\mathcal X$ off $(z,r)=(0,0)$,
whose total fourfold is smooth and whose divisor $r=0$ is
reduced normal crossings. The composite map to $\Delta_r$
is flat.

The replacement has twelve new compact vertical normalization
components. Six are $\PP^1$-bundles over a scalar Hopf
surface; the other six are the blow-ups of those bundles
along two disjoint elliptic curves. Every new component
has fundamental group $\Z$, and therefore fails the
$\partial\bar\partial$ lemma. The punctured family is the
prescribed family, so its generic simplicity is
preserved.

The map is proper over the bidisc and flat over the degeneration
parameter. Its vertical threefolds have the dimensions stated above.
All twelve component labels, and hence the obstruction of
\ref{push:src:B.3.1}, are unchanged.
\end{proposition}



\begin{proof}
\begin{enumerate}[label=\arabic*.,leftmargin=*]
\item \emph{The action.}
Use the rank-three period subsystem
\[
 \mathbf u(z)=(\tau(z),\mu(z)),\qquad
 \mathbf w=(1,0),\qquad \boldsymbol\delta=(0,1).
\]
It defines the analytic semiabelian group
\[
 G_z=\C^2/(\Z\mathbf u(z)+\Z\mathbf w+
                    \Z\boldsymbol\delta),
 \qquad
 1\longrightarrow\C^*_{\delta}\longrightarrow G_z
       \longrightarrow E_z\longrightarrow0,
 \quad E_z=\C/(\Z+\tau(z)\Z).
\]
This is the punctured degree-zero line bundle determined
by $\mu(z)$. The formula
\begin{equation}
 \chi_z(t)=
 \left[\frac{\log t}{2\pi i}\,\mathbf u(z)\right]
 \quad(t\in\C^*)
\label{pub:I:eq:B-4-15-1}
\end{equation}
is a single-valued holomorphic homomorphism: a logarithm
change adds a lattice period. It varies
holomorphically with $z$. The $r$-monodromy fixes
$\mathbf u$, and the $z$-disc has no period monodromy.

\item The action extends to the entire twelve-cycle model.
Each unfolded normalization
is
\[
 R_z=G_z\mathbin{\times^{\C^*_{\delta}}}\PP^1
       =\PP_{E_z}(\OO\oplus L_z).
\]
Translation in $G_z$ preserves both conductor sections.
More explicitly, a class represented by $(a,b)$ acts
before the closing quotient by
\begin{equation}
  (\xi,W)\longmapsto
            (\xi+a,e^{-2\pi i b}W).
\label{pub:I:eq:B-4-15-2}
\end{equation}
It commutes exactly with the closing period, which is
another translation of the same commutative group and
whose scalar factor contains $r^{12}$. The adjacent line
directions are paired dually in the local equation $xy=r$;
the two scalar factors cancel. Scalar multiplication also
preserves the one-dimensional toric subdivision producing
the twelve components. Thus \eqref{pub:I:eq:B-4-15-2} extends over every
node, preserves each branch, and changes a branch equation
$x$ by a unit times $x$. This checks the required action
using the semiabelian marking of Proposition~\ref{push:src:B.1.1}.

\item The affine family itself admits the necessary local
punctured comparison. On
$U=\Delta_z\times\Delta_r^*$ its lattice local system
$\Lambda$ has $H^2(U,\Lambda)=0$, since $U$ retracts onto
a circle, also for local coefficients. Its relative Lie
algebra $\mathcal V$ is a coherent vector bundle. The
exponential sequence
\[
 0\longrightarrow\Lambda\longrightarrow\mathcal V
       \longrightarrow\mathcal T\longrightarrow0
\]
and Cartan's theorem B on the Stein base give
$H^1(U,\mathcal T)=0$ \citep[expos\'e~19, \S5, pp.~14--15]{Cartan1952B_AB}.
Consequently the affine torsor has a holomorphic
section on $U$, which supplies the stated punctured
comparison with its Jacobian.

\item \emph{The base blow-up and small resolutions.}
Blow up the base origin:
\[
 \beta:\widetilde B=\operatorname{Bl}_{(0,0)}
            (\Delta_z\times\Delta_r)
       \longrightarrow\Delta_z\times\Delta_r,
 \qquad F=\beta^{-1}(0,0)\cong\PP^1.
\]
In the chart $r=zs$, the pulled-back node is
$xy=zs$, with one further smooth elliptic coordinate.
Label the original components by $i\in\Z/12$.
At each node choose $x=0$ as the even-indexed branch,
and blow up $(x,z)$. This is globally the sum of the
pulled-back invertible ideal of the chosen central
component and the ideal $\OO_{\widetilde B}(-F)$.
The ideals are $G$-invariant. Their non-Cartier loci are
the twelve disjoint elliptic node loci, so the choices
at the six even components are compatible.

\item The two strict-transform charts are
\begin{equation}
 \begin{aligned}
 z&=xQ,&y&=Qs,&r&=xQs,\\
 x&=zX,&s&=Xy,&r&=zXy.
 \end{aligned}
\label{pub:I:eq:B-4-15-3}
\end{equation}
They are smooth, and all three factors of $r$ occur once.
In the other base chart $z=r\zeta$, the node is $xy=r$
and the same ideal is $(x,r)=(x)$, so no change occurs.
Away from the indicated node loci it is also invertible.
The resulting total space is smooth and its $r$-central
divisor is reduced normal crossings.

\item For the component identifications, before resolution put
\[
 \begin{array}{ll}
 V_x=\{x=z=0\},&V_y=\{y=z=0\},\\
 H_x=\{x=s=0\},&H_y=\{y=s=0\}.
 \end{array}
\]
The $V$ branches are vertical over $F$, and the $H$
branches lie over the strict transform of $r=0$.
On $V_x$, the two maps in \eqref{pub:I:eq:B-4-15-3} are
\[
 (Q,s)\longmapsto(y,s)=(Qs,s),\qquad
 (X,y)\longmapsto(y,s)=(y,Xy).
\]
These are exactly the blow-up charts of $(y,s)$.
Thus $V_x$ is blown up along its elliptic curve over
$y=s=0$; $V_y$ is unchanged. Similarly $H_y$ is blown
up along the corresponding elliptic curve and $H_x$
is unchanged. Each even vertical component is therefore
blown up at its two disjoint conductor curves, whereas
each odd vertical component is unchanged. The horizontal
parities are reversed.

\item \emph{The Cartier logarithm and properness.}
Set $M=\OO_{\widetilde B}(F)$. On a trivializing cover
write its transition units as $g_{ij}$. Glue the resolved
pullbacks by $\chi_z(g_{ij})$. These are genuine
automorphisms satisfying the cocycle identity.
This is the Cartier-logarithm construction of
\citet[construction, p.~85; Definition~1.1, p.~86]{Fujimoto1990AB}; the explicit
semiabelian action above verifies its applicability
here. The invariant ideals show
that the action lifts to the resolutions; equivalently
one may twist first and then resolve the descended ideals.

\item Over every base chart the glued space is a copy of the
proper resolved pullback. It is Hausdorff: two points
over the same base point are separated in such a chart,
and points over distinct base points are separated in
the base. Properness is local on the target, and
composition with the proper $\beta$ proves properness
over the original bidisc. The canonical divisor section
of $M$ is nonzero off $F$ and trivializes the twist there.
The base blow-up and the resolutions are also
isomorphisms off the origin. This gives the specified
isomorphism with the original family there, in particular
over every $r\ne0$.

\item The construction can be inserted into any larger model
which is standard near $q_0$: the preceding
trivialization supplies the overlap map outside the
origin. The cocharacter $\mathbf u$ is defined on this disc.
Since the local rings of the total space are
regular and $r$ is a nonzero divisor, they are
torsion-free, hence flat, over the one-dimensional
regular parameter ring. This proves flatness over $r$.
The new vertical threefolds rule out joint flatness over
the bidisc. Properness over the $r$-disc alone follows
only when this local replacement is inserted into a
model with compact elliptic base.

\item \emph{The compact vertical components.}
Before resolving, each vertical normalization is the
twist of $R_{q_0}\times F$ by
$M|_F=\OO_{\PP^1}(-1)$, namely
\[
 V_i^0=\OO_{\PP^1}(-1)^\times
                    \mathbin{\times^{\C^*}}R_{q_0}.
\]
Its elliptic quotient is a principal $E_{q_0}$-bundle
$H\to\PP^1$ with primitive Chern cycle
$\pm\mathbf u$. It is explicitly a scalar Hopf surface.
Indeed put
\[
 \tau_0=\tau(q_0),\qquad
 \eta=-1/\tau_0,\qquad a'=e^{2\pi i\eta},
 \qquad |a'|<1.
\]
After dividing the elliptic coordinate by $\tau_0$,
the projection of \eqref{pub:I:eq:B-4-15-1} is
$[\log t/(2\pi i)]$ in $\C/(\Z+\eta\Z)$.
Its kernel is $(a')^\Z$. Consequently
\begin{equation}
 H=\OO_{\PP^1}(-1)^\times/(a')^\Z
   =(\C^2\setminus\{0\})/\langle a'\id\rangle.
\label{pub:I:eq:B-4-15-4}
\end{equation}
The principal $G$-bundle
$\OO(-1)^\times\mathbin{\times^{\C^*}}G_{q_0}$
maps to $H$ with fibre $\C^*_{\delta}$.
Compactifying this principal $\C^*$-bundle gives a
projective bundle
\[
 V_i^0=\PP_H(\OO\oplus M_H)
\]
with its two global conductor sections. On each
elliptic fibre of $H\to\PP^1$, the line $M_H$ retains
the original degree-zero line $L_{q_0}$.

\item The two conductor sections over the distinguished point
$F\cap\operatorname{strict}(r=0)$ are disjoint elliptic
curves. The component calculation in \eqref{pub:I:eq:B-4-15-3} gives
exactly the six unchanged bundles and six two-curve
blow-ups in the statement. A $\PP^1$-bundle has the
fundamental group of its base, and blowing up a smooth
complex codimension-two centre preserves fundamental
group. Hence every new normalization has
$\pi_1=\pi_1(H)=\Z$. A compact
$\partial\bar\partial$ manifold has even first Betti
number, so these normalizations fail that lemma.

\item Finally, the connected translation action fixes every
component label. The rotations $+4,-3,0$ of the
incompatible global assembly are unchanged.
\end{enumerate}
\end{proof}

\begin{remark}[Pullback along a lifted ordinary arc]
\hypertarget{tag-PT16}{}\label{I:lifted-arc-central-cycle}
\StatementRecord{I:lifted-arc-central-cycle}{proved}{PT16}{\cref{push:src:B.4.15}}
Let $\beta:\widetilde B\to\Delta_z\times\Delta_r$ be the blow-up used in Proposition~\ref{push:src:B.4.15}, with exceptional curve $F$. An arc $\iota_k(r)=(kr,r)$ has strict lift $\widetilde\iota_k$ meeting $F$ at $[k:1]$. Choose this point away from the distinguished resolution centres. Pulling the constructed model over $\widetilde B$ back along $\widetilde\iota_k$ gives a central cycle of twelve ruled surfaces over $E_{q_0}$. Indeed the pullback of $\mathcal O_{\widetilde B}(F)$ is a line bundle on a disc and is trivial, and the arc misses the additional resolution centres. The new vertical threefolds of \eqref{pub:I:eq:B-4-15-4} instead lie over the whole exceptional curve $F$ under the map to the original bidisc. The maps being compared are the strict-lift pullback and the composite with $\beta$, respectively.
\end{remark}

\begin{proposition}[The dual complex under local insertions]\hypertarget{tag-030K}{}\label{push:src:B.4.16}
\label{push:B416}
\StatementRecord{push:src:B.4.16}{proved}{030K}{\SeriesRef{push:src:B.4.15}, \SeriesRef{push:src:B.3.1}, \SeriesRef{push:src:B.4.2}}

The local operation \ref{push:src:B.4.15} attaches a triangulated
cylinder to the original twelve-cycle in the dual
complex and admits an elementary collapse back to that
cycle. In a compact normal-crossings model, performing
it at a regular place consequently preserves the
homotopy type of the global dual complex.

The same conclusion holds for finite sequences of these
insertions, twists by a connected holomorphic group
action preserving the prescribed normal-crossings
model, and toroidal subdivisions of its central cone
complex. In particular, starting with a model whose
dual complex has a nonzero first integral cohomology
class, these operations cannot produce a semistable
$S^6$ centre. They also cannot change the component
rotations in \ref{push:src:B.3.1}.

\end{proposition}



\begin{proof}
\begin{enumerate}[label=\arabic*.,leftmargin=*]
\item Write $H_i$, $i\in\Z/12$, for the original component
vertices and $V_i$ for the new vertical ones. At a node
between an even component $i$ and its odd neighbour
$j$, the two charts \eqref{pub:I:eq:B-4-15-3} have triple strata
\begin{equation}
 \{V_i,V_j,H_j\},\qquad
 \{V_i,H_i,H_j\}.
\label{pub:I:eq:B-4-16-1}
\end{equation}
The square with horizontal edges $H_iH_j$ and $V_iV_j$
is thus divided by the diagonal $V_iH_j$. Over all
twelve nodes, \eqref{pub:I:eq:B-4-16-1} is exactly a triangulation of
$C_{12}\times[0,1]$, attached along the old horizontal
cycle $C_{12}\times\{0\}$.

\item There is a literal elementary collapse relative to the
old cycle. First remove each triangle containing a top
edge $V_iV_j$, paired with that free edge. Next remove
each remaining triangle, paired with its now free
diagonal $V_iH_j$. Finally remove each vertical edge
$H_iV_i$, paired with its free vertex $V_i$. What remains
is precisely the original twelve-cycle. These are
$36$ elementary collapses. Equivalently, the retraction
sends $V_i$ to $H_i$ and is the identity on all old
vertices and edges.

\item This description is valid inside a larger dual complex
$K$. The only new strata lie over the base origin,
so the new complex is the mapping-cylinder attachment
\[
 K\ \cup_{\,f:C_{12}\to K}\
           (C_{12}\times[0,1]).
\]
It retracts onto $K$, regardless of possible global
identifications among the local horizontal branches.
For self-crossings the ordered local-branch dual complex records
every local branch and its global identifications.

\item A connected group acts trivially on the finite sets
of component labels and connected components of every
stratum: its permutation is locally constant and is
the identity at the identity element. Cartier twists
by such an action therefore keep this same incidence
complex. Toroidal subdivisions subdivide the
corresponding central polyhedral complex and hence
preserve its underlying homotopy type. These arguments
apply successively, proving the assertion for finite
compositions. Around each regular insertion point the
relative component monodromy is still the identity;
no peripheral rotation is added to \ref{push:src:B.3.1}.

\item A compact normal-crossings space $D$
has a natural injection
\begin{equation}
 H^1(\Delta(D),\C)\ \lhook\joinrel\longrightarrow\
 H^1(D,\C),
\label{pub:I:eq:B-4-16-2}
\end{equation}
where $\Delta(D)$ denotes its local-branch dual
complex with connected strata. Use the ordered finite
branch spaces $D^{[p]}$ and take the alternating
summand for the branch-permutation action. Over $\C$
this is an exact direct-summand operation. Resolve the
constant sheaf by this augmented alternating branch
complex. At a point its row is the augmented simplex
complex, so it is exact, including for self-crossings.
The associated first-quadrant spectral sequence has
\[
 E_1^{p,q}=H^q(D^{[p]},\C)^{\mathrm{alt}},
 \qquad
 E_2^{1,0}=H^1(\Delta(D),\C).
\]
No higher differential enters or leaves $(1,0)$.
Its filtration piece in total degree one gives
\eqref{pub:I:eq:B-4-16-2}. Equivalently one can use the orientation
local system on each unordered branch stratum; an
ordinary constant sheaf there would be incorrect if
branches permute. A nonzero integral first cohomology
class of the finite dual complex remains nonzero over
$\C$, because first integral cohomology is torsion-free.
Thus a surviving circle class forces $H^1(D,\C)\ne0$.
By \ref{push:src:B.4.2} every semistable centre of
a simply connected nearby $S^6$ is simply connected,
which is incompatible with this conclusion.

\item For the twelve-cycle, use the oriented edge-to-vertex and
triangle-to-edge incidence maps of the triangulated cylinder in Step~1.
Their chain ranks are $24,48,24$, and their nonzero Smith entries are
$1^{23}$ and $1^{24}$, respectively. Hence its integral homology is
$(\Z,\Z,0)$. The explicit retraction pulls the integral angular cocycle
back with winding one.
\end{enumerate}
\end{proof}

\begin{proposition}[Trivial monodromy does not force equality of first Betti numbers]\hypertarget{tag-030L}{}\label{push:src:B.4.17}
\StatementRecord{push:src:B.4.17}{proved}{030L}{\SeriesRef{push:src:B.4.7}}

There is a compact connected complex threefold $Z$ and a proper
holomorphic map $f:Z\to\PP^1$ with connected fibres, with precisely
one singular fibre $Y$, such that
\[
 f^{-1}(\PP^1\setminus\{0\})
 \simeq H_-\times(\PP^1\setminus\{0\}),\qquad
 b_1(Y)=0<1=b_1(H_-).
\]
Here $H_-$ is the scalar Hopf surface with multiplier $-e^{-\pi}$,
and $Y=H_+\cup_E H_+$ is a reduced simple-normal-crossings union
with $E=\C/(\Z+i\Z)$ and $H_+$ of multiplier $e^{-2\pi}$.
In particular, trivial local monodromy on $H^1$ does not imply
$b_1(Y)=b_1(H_-)$, even with compact smooth total space.
\end{proposition}


\begin{proof}
\begin{enumerate}[label=\arabic*.,leftmargin=*]
\item This is the compact-base version of the explicit Cartier cocycle
of \ref{push:src:B.4.7}.
Put
\[
 b:\overline B=\operatorname{Bl}_{(p,0)}(\PP^1\times\PP^1)
      \longrightarrow\PP^1\times\PP^1,
 \quad L=b^*\operatorname{pr}_1^*\mathcal O(1),
 \quad A=\mathcal O(F_2),
\]
where $F_2$ is the exceptional curve and $F_1$ the strict transform
of $\PP^1\times\{0\}$.  Given transition units $l_{ab},a_{ab}$,
the formula
\[
 t_{ab}=\left[
 \frac{(1+i)\log l_{ab}-i\log a_{ab}}{2\pi i}
 \right]\in E
\]
is independent of logarithm choices and satisfies the cocycle
identity.  Glue the products $U_a\times E$ by these translations.
This gives a holomorphic principal elliptic bundle
$\pi:Z\to\overline B$.  Its total space is smooth and Hausdorff;
it is compact because $\overline B$ and the fibre are compact and
$\pi$ is locally a proper product.  Let
$f=\operatorname{pr}_2\circ b\circ\pi$.

\item The restrictions of the Chern vector
$(1+i)c_1(L)-i c_1(A)$ to $F_1,F_2$ and a smooth ruling fibre
are $1,i,1+i$, respectively.  As checked directly in \ref{push:src:B.4.7},
the corresponding principal elliptic bundles over $\PP^1$
are the two copies of $H_+$ and the surface $H_-$.  Locally
at $F_1\cap F_2$, the base map is $xy=t$ and the elliptic
bundle is a product.  Thus $f$ has smooth total space and just
the one reduced normal-crossings singular fibre $Y$.
The canonical section of $A$ trivializes $A$ outside $F_2$;
there the blow-up is an isomorphism and $L$ is pulled back
from the first factor.  Consequently the asserted punctured
product is a holomorphic product over the whole
$\PP^1\setminus\{0\}$.  In particular every local monodromy
operator is the identity.

\item The fibre inclusions on fundamental groups
are the two quotient maps
\[
 \Lambda=\Z+i\Z\longrightarrow\Lambda/\Z,
 \qquad
 \Lambda\longrightarrow\Lambda/\Z i.
\]
Their van Kampen pushout kills both lattice generators, so
$\pi_1(Y)=0$.  In Mayer--Vietoris the degree-one map, with
the second inclusion negated, has matrix
\[
 \begin{pmatrix}0&1\\-1&0\end{pmatrix}.
\]
It is unimodular.  Since a scalar Hopf surface is diffeomorphic
to $S^1\times S^3$, the sequence gives
\[
 H_*(Y;\Z)=(\Z,0,0,\Z^3,\Z^2),\qquad b_1(H_-)=1.
\]
This proves the asserted strict inequality despite identity
monodromy.  This compact threefold is not $S^6$: for example,
the principal-bundle spectral sequence gives $b_2(Z)=1$.
Indeed the two Chern classes are independent in
$H^2(\overline B;\Z)\simeq\Z^3$, and the differential on
$\bigwedge^2H^1(E)$ is injective, leaving exactly one
degree-two class.
\end{enumerate}
\end{proof}

\begin{remark}[The invariant-cycle equality in the Hopf degeneration]
\hypertarget{tag-PT17}{}\label{I:Hopf-invariant-cycle-comparison}
\StatementRecord{I:Hopf-invariant-cycle-comparison}{proved}{PT17}{\cref{push:src:B.4.17}; \citet[Theorem~1.1, p.~553; arXiv version~1, Theorems~1.2 and~3.2, pp.~1 and~8, equation~(3.9), p.~9, and Corollary~3.7, p.~11]{HondaViaclovsky2025D}}
The compact fibration of Proposition~\ref{push:src:B.4.17} has identity local monodromy $T$ on $H^1(F;\mathbb R)$, but $b_1(Y)=0$ and $b_1(F)=1$. It therefore refutes the converse equality clause in \citet[Theorem~1.1, p.~553]{HondaViaclovsky2025D}, also stated as Theorems~1.2 and~3.2 in arXiv version~1, pp.~1 and~8. The relative-homology and circle-bundle sequences in that proof give
\[
 b_1(Y)\leq\dim H^1(F;\mathbb R)^T\leq b_1(F)
\]
\citep[arXiv version~1, equation~(3.9), p.~9]{HondaViaclovsky2025D}. Identity of $T$ is equivalent to equality in the second inequality. In the present example the first inequality is $0<1$. Equality of the endpoint Betti numbers still forces $T$ to be the identity by this same chain of inequalities. For an irreducible central fibre, the top-dimensional homology in the displayed relative sequence has rank one, giving $b_1(Y)=\dim H^1(F;\mathbb R)^T$, as proved in \citet[arXiv version~1, Corollary~3.7, p.~11]{HondaViaclovsky2025D}. The example here has two central components, with the inclusion maps on $H_1(E)$ computed in Proposition~\ref{push:src:B.4.17}.
\end{remark}

\subsection{A cusp filling on the punctured bidisc}

\begin{proposition}[A partial cusp filling]\hypertarget{tag-030M}{}\label{push:src:B.4.18}
\label{push:B418}
\StatementRecord{push:src:B.4.18}{proved}{030M}{\SeriesRef{push:src:B.3.1}, \SeriesRef{push:src:B.3.2}, \SeriesRef{push:src:B.4.9}, \SeriesRef{push:src:A.2.1}, \SeriesRef{push:src:A.1.2}, \SeriesRef{push:src:B.3.3}}

Put
\[
 B^\times=(\Delta_q\times\Delta_r)\setminus\{(0,0)\}.
\]
There is a proper Hausdorff analytic family over
$B^\times$ with smooth total fourfold, flat over $r$,
whose restriction to $r\ne0$ is the prescribed
cusp family after $w=r^6$, and whose reduced
normal-crossings $r$-central fibre over $q\ne0$ has
component deck permutation the shift $-1$ on $\Z/12$
for the clockwise cusp generator $g_0$ of \ref{push:src:B.3.1}.
Equivalently, the positive $q$ loop has shift $+1$.
In particular its $q$-fibre for every $r\ne0$ is the
reduced Mumford fibre, while its other axis has
the component transport required by \ref{push:src:B.3.2}.

The overlap is the displayed holomorphic translation with its
affine marking on the two specified base opens.
\end{proposition}



\begin{proof}
\begin{enumerate}[label=\arabic*.,leftmargin=*]
\item \emph{The fractional deck and its exact comparison.}
On the logarithmic cover of
$U=\Delta_q^*\times\Delta_r^*$ write
\[
 s=\frac{\log q}{2\pi i},\qquad
 t=\frac{\log r}{2\pi i},\qquad
 C(q)=\beta^{\rm part}(q)+s.
\]
The functions $\mu(q)$ and $C(q)$ are single-valued
and holomorphic at $q=0$. After the affine gauge of
\ref{push:src:B.4.9} the family is its Jacobian, with period
\begin{equation}
 \boldsymbol\gamma(s,t)
       =(6\mu(q),C(q)-s-12t),\qquad
 \boldsymbol\delta=(0,1).
\label{pub:I:eq:B-4-18-1}
\end{equation}
The other periods are
$\mathbf u=(\tau,\mu)$ and $\mathbf w=(1,0)$.
Under $s\mapsto s+1$, they change by
$\boldsymbol\gamma\mapsto\boldsymbol\gamma-
\boldsymbol\delta$ and $\mathbf u\mapsto\mathbf u+\mathbf w$.
Under $t\mapsto t+1$, the period $\boldsymbol\gamma$
changes by $-12\boldsymbol\delta$.
These are integral changes of the period basis; the
ambient additive $\C^2$ coordinate is unchanged.

\item Set
\[
 b(s,t)=\frac1{12}\boldsymbol\gamma(s,t)
                         -\frac1{24}\boldsymbol\delta.
\]
The two exact identities
\begin{equation}
 \sum_{k=0}^{11}b(s+k,t)
       =\boldsymbol\gamma(s,t)-6\boldsymbol\delta,
 \qquad
 b(s,t+1)-b(s,t)=-\boldsymbol\delta
\label{pub:I:eq:B-4-18-2}
\end{equation}
show that this gives a genuine order-twelve affine
$q$-deck action on the root cover $q=z^{12}$,
well defined around the $r$ loop. The first right-hand
side is a period. The sum of the twelve fractions
$\boldsymbol\gamma(s+k,t)/12$ is
$\boldsymbol\gamma(s,t)-11\boldsymbol\delta/2$.
The twelve constant terms $-\boldsymbol\delta/24$
contribute $-\boldsymbol\delta/2$, giving the integral period
$\boldsymbol\gamma(s,t)-6\boldsymbol\delta$.

\item There is an explicit holomorphic trivialization on $U$.
Put $L=s+12t$ and
\begin{equation}
 F(s,t)=
 \left(\frac{\mu(q)L}{2},\,
       \frac{C(q)L}{12}-\frac{L^2}{24}+6t\right).
\label{pub:I:eq:B-4-18-3}
\end{equation}
Direct subtraction gives
\begin{equation}
 F(s+1,t)-F(s,t)=b(s,t),\qquad
 F(s,t+1)-F(s,t)=\boldsymbol\gamma(s,t).
\label{pub:I:eq:B-4-18-4}
\end{equation}
The second difference is the displayed integral period, so the
translation descends to the torus.
Thus $\zeta\mapsto\zeta+F$ is a holomorphic
isomorphism from the untwisted Jacobian to the
$+b$ affine torus torsor over $U$; its inverse
$\zeta\mapsto\zeta-F$ trivializes that torsor.
Replacing $b,F$ by $-b,-F$ gives the opposite
component rotation. If $\sigma$ denotes the 
affine section of \ref{push:src:B.4.9}, the map from the 
family to the $+b$ torsor is explicitly
$\zeta\mapsto\zeta-\sigma+F$.
This map preserves the affine cusp marking.

\item \emph{A rotating compactification before changing the
$q$-boundary.}
In the root-cover model, the two period shears in cocharacter coordinates
$(y_1,y_2,h_z,h_r)$ are
\[
 \begin{aligned}
 G_1 &: (y_1,y_2)\longmapsto(y_1+12h_z,y_2),\\
 G_2 &: (y_1,y_2)\longmapsto
                         (y_1,y_2+12h_z+12h_r).
 \end{aligned}
\]
The positive fractional deck has the shear
\begin{equation}
 H:(y_1,y_2)\longmapsto(y_1,y_2+h_z+h_r),
 \qquad H^{12}=G_2,
\label{pub:I:eq:B-4-18-5}
\end{equation}
and rotates $z$ by $e^{2\pi i/12}$.
It commutes with $G_1$. Its unit factors are holomorphic
functions of $q$. For example, in the convention
$W=e^{-2\pi i\zeta_2}$, it multiplies $W$ by
\[
 r z\,e^{-2\pi i C(q)/12}e^{\pi i/12}
\]
and the first multiplicative coordinate by
$e^{\pi i\mu(q)}$. The twelve successive root rotations
contribute the factor
$e^{2\pi i(0+\cdots+11)/12}=-1$, and the scalar
correction contributes $e^{\pi i}=-1$.
They cancel. The twelfth iterate is consequently
the second period, including its unit, as
also follows additively from \eqref{pub:I:eq:B-4-18-2}.

\item Use the unit-square cusp triangulation, with twelve
consecutive bottom cells in the $y_1$ direction and
the same bridge choice in every unit prism. Its
translates have bottom cell indices $(i,j)$ and top
interval $[j,j+1]$. Thus $H$ increases the bottom
second index and the top index by one, while $G_1$
shifts the first index by twelve. The fan is invariant
under both, and therefore under $G_2$.
For each bottom cell its three maximal cones are
the determinant-one cones of \ref{push:src:A.2.1}. The usual
compact-core construction of \ref{push:src:A.1.2}, now on the
$I_{12}$ elliptic root model and with twelve
plumbing intervals before the $G_2$ quotient,
gives a proper model on the root cover. Its
unfolded charts are precisely this fan. The compact core and collars therefore prove properness and
proper discontinuity.

\item The affine deck extends to this model because its
unit multiplication preserves every toric chart
and its shear preserves the fan. It is free.
On the $r$-central divisor a point lies on at most
two adjacent component labels in $\Z/12$; no
nonidentity translation can preserve that set.
This covers their intersection over the common
base origin as well. On the $q$-central divisor
with $r\ne0$, the bottom triangle has second
component indices contained in $\{j,j+1\}$,
again not preserved by a nonzero translation
modulo twelve. Off both axes the root-base action
is free. Its finite quotient is therefore smooth,
proper over $(q,r)$, flat over $r$, and has a
reduced normal-crossings $r$-central divisor.

\item Keeping the same positive root-base rotation and using translation
$-b$ gives valuation shear $H^{-1}$ and twelfth iterate
$G_2^{-1}$. This action preserves the same fan and has the same
freeness properties. For the family in the statement we use
translation by $b$ and the clockwise cusp generator
$g_0:s\mapsto s-1$. Thus the
positive fractional deck $b$ has shift $+1$ for
$s\mapsto s+1$ and shift $-1$ for $g_0$, exactly
as required by \ref{push:src:B.3.2}. The finite shifts $4,-3$
and this cusp shift all refer to the same clockwise
deck-identification convention.

\item Denote this rotating model by $\mathcal R$.
Its $q$-fibre for $r\ne0$ has multiplicity twelve.
Indeed the original $z$-central divisor is reduced
and the finite action has trivial generic isotropy
there, while $q=z^{12}$. This is the $q$-multiplicity computed in \ref{push:src:B.3.3}.

\item \emph{The two-open gluing.}
On the base open set $q\ne0$ use $\mathcal R$ with
the positive fractional deck $b$. On the base open set $r\ne0$
use the reduced Mumford cusp family.
Their common base is $U$, where \eqref{pub:I:eq:B-4-18-3}--
\eqref{pub:I:eq:B-4-18-4}, composed with \ref{push:src:B.4.9}, give the specified
holomorphic isomorphism. These two open sets cover
$B^\times$, and there is no additional triple
overlap. Glue by that isomorphism.

\item The result is Hausdorff: points over one base
point are separated in one of the two given
Hausdorff models, and points with different
base images are separated in the base.
Properness is local on the target, so follows
from the properness of those two restrictions.
Smoothness and flatness over $r$ are also local.
The restriction over $r\ne0$ is, by construction,
the prescribed cusp including its reduced
$q$-fibre. Over $r=0,q\ne0$, the construction has
the required twelve-cycle and rotation $-1$.
This proves every asserted property.

\end{enumerate}
\end{proof}

\begin{remark}[The two axial models and their overlap]
\hypertarget{tag-PT18}{}\label{I:two-axial-model-comparison}
\StatementRecord{I:two-axial-model-comparison}{proved}{PT18}{\cref{push:src:B.4.18,push:src:B.3.7,push:src:B.4.30,night:P2.7}}
Proposition~\ref{push:src:B.4.18} uses the rotating model $\mathcal R$ on $q\ne0$ and the reduced Mumford family on $r\ne0$. On their overlap the exact isomorphism is $\zeta\mapsto\zeta-\sigma+F$, with $F$ given by \eqref{pub:I:eq:B-4-18-3}. The larger rotating model has local equation $q=z^{12}$ at its $q$-boundary for $r\ne0$, so every $q$-vertical divisorial valuation there has order in $12\mathbb Z$. The reduced Mumford boundary has a divisor of $q$-order one. A common proper modification over $q$ would identify the function $q$ and the associated divisorial valuations; Proposition~\ref{push:src:B.3.7} gives the resulting contradiction. Thus the overlap isomorphism specifies precisely the two-open gluing on $B^\times$. Its noncompletion at the missing origin is Proposition~\ref{push:src:B.4.30}, and the obstruction for the full global family is Theorem~\ref{night:P2.7}.
\end{remark}

\subsection{The global family with one base point removed}

\begin{proposition}[The global partial filling]\hypertarget{tag-030N}{}\label{push:src:B.4.19}
\label{push:B419}
\StatementRecord{push:src:B.4.19}{proved}{030N}{\SeriesRef{push:src:B.1.2}, \SeriesRef{push:src:B.4.18}, \SeriesRef{push:src:B.4.30}, \SeriesRef{push:src:B.4.33}, \SeriesRef{push:src:C.3.5}, \SeriesRef{push:src:B.1.1}, \SeriesRef{push:src:B.3.2}, \SeriesRef{push:src:B.4.9}, \SeriesRef{push:src:A.1.2}}

Let
\[
 \Omega=(\mathbb P^1\times\Delta_r)
                  \setminus\{(\infty,0)\}.
\]
There is a proper Hausdorff holomorphic map
$\mathcal Y^\circ\to\Omega$ with smooth total fourfold,
flat over $r$, whose restriction to $r\ne0$ is the
family pulled back by $w=r^6$. Its reduced
normal-crossings central divisor has the finite-place
models of \ref{push:src:B.1.2} and the punctured cusp of
\ref{push:src:B.4.18}, with their given markings. Over
$B^\circ=\mathbb P^1\setminus\{0,1,\infty\}$ its
relative component cover is
\[
 y^{12}=K(x),\qquad K(x)=\frac{x^4}{(x-1)^3}.
\]
Thus the global finite-place overlap maps and their
cocycle are constructed. The sole missing base point is $(\infty,0)$.
Proposition~\ref{push:src:B.4.30} proves its noncompletion, and
\ref{push:src:B.4.33} excludes every proper filling of the 
family over the whole $r$ disc.
\end{proposition}



\begin{proof}
\begin{enumerate}[label=\arabic*.,leftmargin=*]
\item We distinguish the ordinary base coordinate $x$ from
the logarithmic parameter
\[
 v=\frac{\log r}{2\pi i},\qquad c=-12v,
 \qquad e=(0,1)^t.
\]
Use the coherent closing of \ref{push:src:C.3.5}. Its additive formula
for a positive $r$ loop, that is $c\mapsto c-12$, is
\begin{equation}
 b_{-12,c}=\Pi_c\widehat\gamma-(k+6)e,
 \qquad k=12\ell=\frac{\log K}{2\pi i}.
\label{pub:I:eq:B-4-19-1}
\end{equation}
The first summand and $6e$ are periods.
Consequently its torus translation has the scalar
representative $-ke$, since addition of those periods
preserves the quotient map.

\item The coherent closing in \ref{push:src:C.3.5} may be used for the
specified punctured family. More explicitly, two
coherent choices differing by a scalar automorphism
give a principal $\mathbb C^*$ torsor over the punctured
parameter disc. That torsor is holomorphically trivial:
the exponential sequence and Cartan's theorem B give
$\operatorname{Pic}(\Delta_r^*)=H^2(\Delta_r^*,\mathbb Z)=0$.
The resulting scalar gauge is a holomorphic
isomorphism of the punctured families.

\item \emph{The ordinary charts and their exact cocycle.}
Choose a good analytic cover of $B^\circ$, refining
period-frame charts, and on each member choose a
holomorphic logarithm $k_i=(2\pi i)^{-1}\log K$.
On the parameter logarithm cover put
\begin{equation}
 \widetilde\zeta_i=\zeta+k_i v e.
\label{pub:I:eq:B-4-19-2}
\end{equation}
Under the positive $r$ deck, $\zeta$ changes by
$-k_i e$ modulo periods and $v$ changes by one.
Thus \eqref{pub:I:eq:B-4-19-2} has the ordinary period descent.
It identifies the local punctured family
with the standard twelve-cycle plumbing of \ref{push:src:B.1.1},
now over an ordinary base chart.

\item On an overlap let $n_{ij}=k_j-k_i\in\mathbb Z$.
The invariant vector $e$ is fixed by every period-frame
change. In the multiplicative coordinate
$W_i=\exp(-2\pi i\widetilde\zeta_{i,2})$, the extra
transition is exactly
\begin{equation}
 W_j=r^{-n_{ij}}W_i.
\label{pub:I:eq:B-4-19-3}
\end{equation}
Every integral power in \eqref{pub:I:eq:B-4-19-3} extends to an
automorphism of the whole twelve-cycle plumbing.
On the unfolded node chart write
\[
 x_j=W/r^j,\qquad y_j=r^{j+1}/W,
 \qquad x_jy_j=r.
\]
Multiplication $W\mapsto r^{-n}W$ sends the chart
index $j$ to $j-n$ and leaves both displayed chart
coordinates unchanged. The same statement holds
in bundle frames, whose additional factors are
holomorphic units. It commutes with the closing
period, which shifts the index by twelve and
translates the elliptic base by $6P$.
The usual semiabelian translations and period-frame
changes preserve the two sections, hence also extend.

\item The integers satisfy
$n_{ik}=n_{ij}+n_{jk}$. Since the scalar translations
commute with each other and with the period-frame
changes, \eqref{pub:I:eq:B-4-19-3} is an exact cocycle on the entire
plumbing. This constructs the family over
$B^\circ\times\Delta_r$, with the original
punctured-family identification \eqref{pub:I:eq:B-4-19-2}.

\item \emph{The finite-place cancellations.}
At the two root charts put
\[
 w_j=\frac{\log s_j}{2\pi i},\qquad
 \epsilon_1=1,\quad\epsilon_2=-1,
 \qquad k_j=12\kappa_j.
\]
The notation is that of \ref{push:src:C.3.5}: $k_j$ is a
holomorphic logarithm of the unit
$K(x)s_j^{-12\epsilon_j}$ on the full root disc.
Its invariance under the finite deck is exact.
Indeed the unit is invariant since $m_j\mid12$;
the difference of two logarithms is an integer
constant, and evaluating at the fixed root centre
makes that constant zero. The scalar representative
of the local positive $r$ closing is $-k_j e$.
Apply the local gauge
$\widetilde\zeta_j=\zeta_j+k_jve$ and use the
smooth twelve-cycle root model of \ref{push:src:B.1.1}, with its
free finite deck action from \ref{push:src:B.1.2}. Because $e$ is
invariant and $k_j$ is invariant, this gauge does
not change that finite action.

\item On a sector of the punctured root disc choose
compatible logarithms. The root-to-ordinary
collar is
\[
 \zeta=\zeta_j+\sigma_{j,c},\qquad
 \sigma_{j,c}=w_j\Pi_c v_j,
 \qquad \gamma(v_j)=\epsilon_j.
\]
Moreover $k-k_j=12\epsilon_jw_j$. Therefore its
expression in the new gauges is
\begin{equation}
 \begin{aligned}
 \widetilde\zeta-\widetilde\zeta_j
 &=w_j\Pi_c v_j+12\epsilon_jw_jve\\
 &=w_j\Pi_{\rm part}v_j,
 \qquad \Pi_{\rm part}=\Pi_c-ce\gamma.
 \end{aligned}
\label{pub:I:eq:B-4-19-4}
\end{equation}
Here $c=-12v$ is essential: the parameter logarithms
cancel exactly. The last translation is independent
of $r$. On each sector it is a holomorphic section
of the semiabelian group and therefore extends to
the ruled normalization and all node charts.

\item Changing a logarithm adds an integral transition
of the form \eqref{pub:I:eq:B-4-19-3}, together with a
period-frame identification. Those maps already
extend. Their deck and overlap identities agree
on the dense locus $r\ne0$ with the collar
identities, hence agree everywhere by the identity
principle. Thus the sector maps descend and glue
exactly to the finite quotients. This constructs
the specified family over
$(\mathbb P^1\setminus\{\infty\})\times\Delta_r$.
It retains the multiplicities three and four and
the reductions already proved in \ref{push:src:B.1.2}.

\item The component cover is determined by the local indices.
If $i_i$ is the
unfolded component index modulo twelve, then
\eqref{pub:I:eq:B-4-19-3} gives $i_j=i_i-n_{ij}$. Hence
\[
 y=\exp(2\pi i k_i/12)\exp(2\pi i i_i/12)
\]
is independent of the chart and satisfies $y^{12}=K$.
The clockwise deck shifts are consequently $4,-3,-1$.
The local finite quotients complete this cover at
$0,1$ exactly as in \ref{push:src:B.3.2}.

\item \emph{The cusp overlap, including $r=0$.}
Retain $s=\log q/(2\pi i)$ and the holomorphic
functions $C,\mu,\kappa_0$ of \ref{push:src:B.4.18} and \ref{push:src:C.3.5}.
Then $k=-s+12\kappa_0$. The cusp-to-rotating comparison of the positive-deck model is
$\zeta\mapsto\zeta-\sigma+F$. Substitution in the
displayed formulas of \ref{push:src:B.4.9} and \ref{push:src:B.4.18} gives
\begin{equation}
 F-\sigma=kve+G(s)+6ve,
 \qquad
 G(s)=\left(\frac{\mu s}{2},\,
                    \frac{Cs}{12}-\frac{s^2}{24}\right).
\label{pub:I:eq:B-4-19-5}
\end{equation}
Thus its comparison with the ordinary gauge
\eqref{pub:I:eq:B-4-19-2} is the translation $G(s)+6ve$.
On a sector of $q\ne0$, $G$ is independent of $r$
and gives a holomorphic semiabelian translation.
The other term is the integral fan shift
$W\mapsto r^{-6}W$. Both extend, with their inverses,
over the full twelve-cycle $r$-central divisor.
The branch identities are
\[
 G(s+1)-G(s)=b(s,v)+ve,
 \qquad k(s+1)-k(s)=-1.
\]
The ordinary coordinate deck contributes $-ve$;
their sum is exactly the positive fractional deck
$b$. A positive $r$ loop changes $6ve$ by the
integral period $6e$. These identities prove
descent of the extended comparison to the full
punctured $q$ disc. In particular the common
clockwise cusp shift is $-1$, as required.

\item Glue the just constructed global model on
$x\ne\infty$ to \ref{push:src:B.4.18} on the cusp neighbourhood
minus its origin, using \eqref{pub:I:eq:B-4-19-5}. The overlap
contains $r=0$ and the map extends there by the
preceding chart calculation. The identities give the exact scalar and translation cocycles
with the prescribed punctured marking.

\item Each chart used is a proper model over its base
open set: this was proved by compact-core plumbing
in \ref{push:src:A.1.2} and by the finite free quotient in \ref{push:src:B.1.2};
\ref{push:src:B.4.18} is proper on its stated punctured bidisc.
Properness is local on the target, so the glued
map to $\Omega$ is proper. The same local argument
gives Hausdorffness, smoothness of the total space,
flatness over $r$, and the reduced normal-crossings
central divisor. For $r\ne0$ the base is the entire
$\mathbb P^1$, and the specified comparison is
precisely the compact family. At $r=0$ the target is the finite base $\mathbb P^1\setminus\{\infty\}$,
as specified in $\Omega$.

\end{enumerate}
\end{proof}

\subsection{An embedded Tate interface at the missing corner}

\begin{proposition}[The other Kato axis]\hypertarget{tag-030P}{}\label{push:src:B.4.20}
\label{push:B420}
\StatementRecord{push:src:B.4.20}{proved}{030P}{\SeriesRef{push:src:B.4.10}, \SeriesRef{push:src:B.4.33}; \cite[Example 9.10]{IstratiEtAl2022Root}}

In the proper smooth Kato family $\mathcal K\to\Delta_s$
of \ref{push:src:B.4.10} there is a second relative effective Cartier
divisor $\mathcal D\subset\mathcal K$, disjoint from
the persistent elliptic divisor $\mathcal E$.
The surface $\mathcal D$ is smooth and proper over
$\Delta_s$. Its general fibre and central fibre are
\[
 \mathcal D_s=\mathbb C^*/(\kappa(s)s^n)^\mathbb Z
 \quad(s\ne0),\qquad
 \mathcal D_0=\hbox{the reduced rational }n\hbox{-cycle}.
\]
At its nodes the relative local equation is $uv=s$.
Its normal line in $\mathcal K$ is the multidegree-zero
Tate line whose factor on the closing collar is $a(s)$.
In particular it has degree zero on every component
of the central cycle. Scaling the first blowdown
coordinate gives a holomorphic $\mathbb C^*$ action
on the pair $(\mathcal K,\mathcal D)$ over the base.

For $n=12$ this constructs an embedded proper copy of
the cusp's elliptic root interface, with its
degree-zero normal line, after choosing the units
$\kappa,a$ to be the Tate and normal multipliers.
In the additive period convention $\mathbf u=(\tau,\mu)$,
$W=e^{-2\pi i\zeta_2}$, these multipliers are
$\kappa(s)s^{12}=e^{2\pi i\tau}$ and
$a(s)=e^{-2\pi i\mu}$. The elliptic translation by
$6P$ is realized by first-coordinate scaling
$a(s)^{-6}$.

The comparison identifies the embedded Cartier divisor and its
normal line. At each central node the ambient base map has the local
form $s=uv-h$ computed below.
\end{proposition}



\begin{proof}
\begin{enumerate}[label=\arabic*.,leftmargin=*]
\item Use the modification, embedding and compact collar
domains of \ref{push:src:B.4.10}. Write $w$ for the second blowdown
coordinate and set, on the modified fundamental region,
\[
 h_s=\pi_n^*w-s.
\]
The collar identification is $\sigma_s\pi_n$, where
$\sigma_s(z,w)=(\kappa(s)z,a(s)w+(1-a(s))s)$ in
the last chart. Consequently
\begin{equation}
 h_s\circ(\sigma_s\pi_n)=a(s)h_s.
\label{pub:I:eq:B-4-20-1}
\end{equation}
Its zeros therefore glue to a relative effective
Cartier divisor, and \eqref{pub:I:eq:B-4-20-1} is simultaneously the
normal-line transition. The divisor is closed in the
proper family $\mathcal K\to\Delta_s$, so it is proper.

\item For $s\ne0$, the line $w=s$ misses every blow-up centre.
In the coordinate $v=w-s$ the contracting germ is
\[
 (x,v)\longmapsto
       (\kappa(s)x(a(s)v+s)^n,a(s)v).
\]
Its restriction to $v=0$ is
$x\mapsto\kappa(s)s^n x$. The generic developing
map of \ref{push:src:B.4.10} restricts to the identity in this
coordinate, since its product $G_s(v)$ has
$G_s(0)=1$. Thus this is precisely the unmodified
other Hopf axis, with the stated Tate multiplier
and normal multiplier $a(s)$.

\item The divisor of $\pi_n^*w$ is reduced. After the first
blow-up, both the strict horizontal axis and the
exceptional curve have coefficient one. Every later
centre is the intersection of the newest exceptional
curve with the strict vertical axis. It lies on no
other component of the divisor of $\pi_n^*w$.
The new exceptional coefficient is therefore again
one. Before the end collars are identified, the
divisor is a chain comprising the strict horizontal
axis and the $n$ exceptional pieces. The end annuli
join the two endpoint pieces, leaving the familiar
rational $n$-cycle of the parabolic surface.
It is disjoint from its elliptic curve, as also
described in the parabolic model of
\cite[Example~9.10, p.~1386]{IstratiEtAl2022Root}.

\item At a crossing of two adjacent pieces, ordinary
blow-up coordinates give $\pi_n^*w=uv$ up to a
holomorphic unit, which can be absorbed in $u$.
Hence $\mathcal D$ has equation $uv=s$. It is
smooth, and its central fibre is reduced nodal.
Away from these nodes its defining equation is
a coordinate minus $s$. This proves smoothness
of the whole divisor and semistability over the disc.
The monomial transition charts are those of the
unit-interval Tate model, and the closing on the
nonzero axis is exactly $x\mapsto\kappa(s)s^n x$.
They identify the marked family with that Tate
model over $\Delta_s$.

\item The same local defining functions trivialize the
normal line. Their only nontrivial closing factor
is \eqref{pub:I:eq:B-4-20-1}, namely the holomorphic unit $a(s)$.
On each rational component the line has degree
zero; equivalently its two adjacent-node degree
corrections cancel its self-intersection $-2$.
For $n=1$ the same statement uses the two branches
of the self-node. The factor $a(s)$ gives the
gluing parameter around the cycle.
This is exactly the multidegree-zero extension
of the line with factor $a(s)$ on the smooth
Tate curves.

\item The first-coordinate scaling commutes with the
modification and the collar maps and preserves
$h_s$. One may equivalently descend the vector
field $x\partial_x$ and its local flow. Properness
over the base makes that flow complete on compact
fibres, uniformly on compact parameter sets;
its $2\pi i$ period yields the stated holomorphic
$\mathbb C^*$ action. On the Tate axis it is
ordinary multiplication of $x$, so a holomorphic
unit $a(s)^{-6}$ realizes the marked translation
by $6P$.

\item For the cusp, the first lattice period
acts in multiplicative coordinates by
\[
 (Z,W)\longmapsto
       (e^{2\pi i\tau}Z,e^{-2\pi i\mu}W).
\]
After the twelfth root change, its first multiplier
is $\kappa(s)s^{12}$ for a holomorphic unit.
The second multiplier tends to the contracting
inverse of $-(5+2\sqrt6)$, so after shrinking
$0<|a(s)|<1$ as required in \ref{push:src:B.4.10}.
The preceding exact Tate charts and the normal
factor identify this interface and its line with
the cusp interface and its normal line. The two section normals in the
ruled model are inverse; one chooses the matching
branch accordingly. More precisely, in the total
line-coordinate convention, this $W$ normal is
$M^{-1}$ for $M=\mathcal O(P-O)$: the theta ratio
with divisor $P-O$ has multiplier $e^{2\pi i\mu}$,
the inverse of $a$. The opposite ruled section
uses $M$. The separate first-coordinate factor
$a^{-6}$ is the translation by $6P$.

\item At a central node, the smooth Kato total space
has coordinates $(u,v,s)$ with divisor
$h=uv-s=0$. In coordinates $(u,v,h)$ its map to
the base is
\begin{equation}
 s=uv-h.
\label{pub:I:eq:B-4-20-2}
\end{equation}
It is a submersion. The total space of the normal
line, with the inherited map from $\mathcal D$,
has $s=uv$, with a critical normal
line over that node. Therefore there is no
relative linear tubular-neighbourhood identification
across the nodes. An isomorphism over the base would identify the critical loci of
the two base maps, contradicting these explicit differentials.

\item With an inverse-normal plumbing variable $f$, the local equation is
\begin{equation}
 r=fh=f(uv-s).
\label{pub:I:eq:B-4-20-3}
\end{equation}
The total fourfold in these coordinates is
smooth and its $r$-central divisor is reduced
normal crossing. For $r\ne0$ one has $f\ne0$
and $s=uv-r/f$, whose derivative in $f$ is
$r/f^2\ne0$. Thus the map to $s$ on this local nonzero-$r$ fibre is a submersion
at the displayed locus.

\item There is also a concrete obstruction to the most
direct global use of this plumbing. For fixed
$s\ne0$, all the $n$ blow-up centres lie on the
first Hopf axis $x=0$, at $[v=-s]$ and its
infinitely near strict-axis points. The other
axis $\mathcal D_s$ is $v=0$, $x\ne0$.
The axes are disjoint, and hence every exceptional
rational curve misses $\mathcal D_s$. A plumbing
that leaves the complement of a small tubular
neighbourhood of $\mathcal D_s$ unchanged retains
these compact curves in its $r\ne0$ fibres over
the same $s$. A complex torus has no rational
curve: a holomorphic map from $\mathbb P^1$
lifts to its vector-space universal cover and
is constant. Thus this fixed-base localized
plumbing of the entire Kato surface cannot
produce the prescribed ordinary two-tori.

\end{enumerate}
\end{proof}

\begin{proposition}[A degree-one Kato model with a linear twelve-gon collar]\hypertarget{tag-030Q}{}\label{push:src:B.4.21}
\StatementRecord{push:src:B.4.21}{proved}{030Q}{\SeriesRef{push:src:B.4.10}, \SeriesRef{push:src:B.4.19}, \SeriesRef{push:src:B.4.33}, \SeriesRef{push:src:B.4.20}}

Start with the $n=1$ Kato family of \ref{push:src:B.4.10} over a disc
$\Delta_q$, with holomorphic units $a(q),\kappa(q)$ and
$0<|a(q)|<1$.  Pull it back by $q=s^{12}$.
There is a proper modification of this pulled-back threefold,
obtained by six point blow-ups and seventy-two smooth curve
blow-ups supported on its central rational cycle, with the
following properties.  Its total space $\mathcal K^*$ is smooth,
it is proper and flat over $\Delta_s$, and its restriction to
$s\ne0$ is the original pulled-back Kato family.

The strict transform $\mathcal D^*$ of the other Cartier axis
is the smooth semistable Tate surface with multiplier
$\kappa(s^{12})s^{12}$ and central fibre $I_{12}$.  A neighbourhood
of $\mathcal D^*$ in $\mathcal K^*$ is biholomorphic over
$\Delta_s$ to a neighbourhood of the zero section in the
multidegree-zero Tate line with closing factor $a(s^{12})$.
The comparison retains the explicit transverse coordinate
\begin{equation}
 H=\frac{h}{q},\qquad h=\pi_1^*w-q,
\label{pub:I:eq:B-4-21-1}
\end{equation}
on the original punctured family.  Thus its normal comparison,
relative to the original punctured normal frame, has factor
$q^{-1}=s^{-12}$; its inverse has factor $q$.

The persistent elliptic axis has an unchanged neighbourhood.
In particular its normal in the strict transform of the original
central Kato surface has degree $-1$.
The displayed function $H$ and the closing factors specify the
linear-collar isomorphism on a neighbourhood of $\mathcal D^*$.
\end{proposition}


\begin{proof}
\begin{enumerate}[label=\arabic*.,leftmargin=*]
\item We distinguish the original parameter $q$ from its twelfth root
$s$.  The pulled-back Kato total space is smooth: the original
map to $q$ is a submersion, so locally its pullback has smooth
coordinates consisting of two fibre coordinates and $s$.
The other axis of \ref{push:src:B.4.20} is Cartier, with closing factor $a(q)$
for its defining function $h$.  At its unique central node its
local equation after this base change is
\begin{equation}
 h=uv-s^{12}.
\label{pub:I:eq:B-4-21-2}
\end{equation}
Away from that point the axis surface is smooth.  The persistent
elliptic axis is disjoint from it throughout the family.

\item \emph{The six point blow-ups.}
Blow up the singular point of the strict transform of the
Cartier surface, six times in succession.  All these are points
of the smooth ambient threefold.  If its current equation is
$uv=s^m$, with $m\ge2$, the $s$ chart of the ambient point
blow-up has
\begin{equation}
 u=su_1,\qquad v=sv_1,
 \qquad u_1v_1=s^{m-2}.
\label{pub:I:eq:B-4-21-3}
\end{equation}
The $u$ chart has $v=uv_1$, $s=us_1$, and strict-transform
equation
\begin{equation}
 v_1=u^{m-2}s_1^m;
\label{pub:I:eq:B-4-21-4}
\end{equation}
the $v$ chart is symmetric.  The latter charts are smooth.
For $m>2$ the only remaining singular point is the origin in
\eqref{pub:I:eq:B-4-21-3}; for $m=2$ that chart has $u_1v_1=1$.
The exponents are $12,10,8,6,4,2$, so precisely six point
blow-ups resolve the surface.

\item Write $D$ for its smooth strict transform.  Its central fibre is
\begin{equation}
 F=C_0+C_1+\cdots+C_{11}=\operatorname{div}_D(s).
\label{pub:I:eq:B-4-21-5}
\end{equation}
Here $C_0$ is the strict transform of the old nodal rational
curve.  Each of the first five point blow-ups contributes two
rational exceptional curves, and the sixth contributes the
smooth conic in its exceptional plane.  In cyclic order the
two chains are $C_1,\ldots,C_5$ and $C_{11},\ldots,C_7$, with
the conic $C_6$ between them.  Equations \eqref{pub:I:eq:B-4-21-3}--
\eqref{pub:I:eq:B-4-21-4} show that every component has multiplicity one,
that consecutive components meet transversely once, and that
there are no other intersections.  Thus the fibre is $I_{12}$.
Since $F$ is principal, $C_i\cdot F=0$ and $C_i^2=-2$.
The same charts identify $D$ with the minimal regular Tate
model of the given multiplier, including the original closing
unit $\kappa(q)$.

\item Let $M$ be the pullback to $D$ of the original Cartier normal
line.  It is the Tate line with factor $a(q)$ and multidegree
zero.  A component contracted to the original node has degree
zero for this pullback, and its restriction to $C_0$ is the
pullback of the degree-zero line on the nodal curve.
In the iterated $s$ chart after $j$ point blow-ups the original
function is
\[
 h=s^{2j}(u_jv_j-s^{12-2j}).
\]
The $j$th exceptional ambient divisor therefore has coefficient
$2j$ in the total transform of $\operatorname{div}(h)$.
For $j<6$ its intersection with $D$ consists of
$C_j+C_{12-j}$, and for $j=6$ it is $C_6$.
Consequently the new normal line is
\begin{equation}
 N_{D/\mathcal K^{(6)}}=M\otimes\mathcal O_D(-Z),
 \quad Z=\sum_{i=0}^{11}z_iC_i,
 \quad z_i=2\min(i,12-i).
\label{pub:I:eq:B-4-21-6}
\end{equation}
The list of coefficients is
$(0,2,4,6,8,10,12,10,8,6,4,2)$.
In particular its degree vector is $-4$ on $C_0$, $+4$ on
$C_6$, and zero on the other ten curves.

\item \emph{The seventy-two curve blow-ups and the function.}
For each $i$, blow up the ambient threefold along the current
copy of $C_i$ exactly
\[
 b_i=12-z_i
\]
times, in increasing order of $i$; the entry $b_6=0$ requires
no blow-up.  At every step $C_i$ is a smooth embedded compact
curve.  Its ideal on $D$ is invertible, so the strict transform
of $D$ is canonically isomorphic to $D$ itself.  The normal
line is tensored by $\mathcal O_D(-C_i)$ at that step.
Thus the final normal is
\begin{equation}
 M\otimes\mathcal O_D(-Z-\textstyle\sum_i b_iC_i)
   =M\otimes\mathcal O_D(-12F)\simeq M.
\label{pub:I:eq:B-4-21-7}
\end{equation}
The sum of the $b_i$ is $72$.
The last isomorphism divides the local generator $s^{12}$
of $\mathcal O_D(-12F)$ by $s^{12}$.  Relative to the
original punctured normal identification it is multiplication
by $s^{-12}$, as asserted.

\item The point-blow-up charts give the following exact local
description near $D$.  At a node $C_i\cap C_{i+1}$ one can
choose ambient coordinates $(x,y,e)$ such that
\begin{equation}
 D=(e=0),\qquad s=xy,\qquad
 h=e x^{z_i}y^{z_{i+1}}\epsilon,
\label{pub:I:eq:B-4-21-8}
\end{equation}
where $\epsilon$ is a holomorphic unit.  At a general point
of $C_i$, omit $y$ and use $s=x$.  These formulas follow
directly from \eqref{pub:I:eq:B-4-21-4}: the strict transform is a coordinate
minus a monomial, and the two factors of $s$ are the two
vertical divisors.  They also show transversality of $D$ to
their intersection.  Absorbing a unit in the second base
coordinate preserves this description if different original
coordinate frames are used.

\item Blowing up $C_i$ is locally the blow-up of $(e,x)$.
In the chart containing the strict transform of $D$, put
$e=xe'$.  Formula \eqref{pub:I:eq:B-4-21-8} keeps $s=xy$ and increases
the exponent of $x$ in $h$ by exactly one.  The old vertical
divisor in this direction no longer meets the strict transform
of $D$.  At an adjacent node the exponent of $y$ is unchanged.
This calculation applies at every repetition and for the
other component as well.  At the end, every vertical ambient
divisor meeting $D$ has $s$-order one and $h$-order twelve.
Hence \eqref{pub:I:eq:B-4-21-8} becomes
\begin{equation}
 h=s^{12}e\epsilon,
 \qquad H=h/q=e\epsilon.
\label{pub:I:eq:B-4-21-9}
\end{equation}
In particular $H$ is holomorphic near all of $D$, has divisor
exactly $D$, and has nonzero normal derivative there.  On the
closing collar its transition is exactly $H\mapsto a(q)H$.
Thus the modification makes the meromorphic function $H$
holomorphic on this neighbourhood.

\item \emph{The relative product charts.}
There is an equivalent useful local target for the preceding
modifications.  The ideal $(h,q)$ is globally defined, since
its changes of generators multiply $h$ by the unit $a(q)$.
Near $D$ the modification maps to its $q$ blow-up chart,
because \eqref{pub:I:eq:B-4-21-9} makes the pullback ideal principal.
At the old node this chart has
\[
 uv=s^{12}(1+H).
\]
Restrict to $|H|<\eta<1$ in a local closing frame and put
$V=v/(1+H)$.  The equation becomes
\[
 uV=s^{12},
\]
times the independent coordinate $H$.
The preceding smooth modification near $D$ is its minimal
resolution times that coordinate.  

\item Label the exceptional chain so that the orders of $u,V,s$
on $C_i$ are $i,12-i,1$, respectively, with the two end
branches at indices $0,12$.  These are the values in the
explicit surface resolution above.  Along any ambient vertical
divisor meeting $C_i$, the orders of $u$ and $V$ sum to twelve,
because $uV=s^{12}$ and $s$ has order one.  Restriction to
$D$ can only increase those orders.  Since the restricted
orders already sum to twelve, the ambient orders are exactly
$i,12-i$.  No other zero or pole of either character occurs
near $D$, because their product is $s^{12}$.

\item On the chart of the two adjacent rays $i,i+1$, the functions
\begin{equation}
 X_i=u/s^i,\qquad Y_i=s^{i+1}/u,
 \qquad X_iY_i=s
\label{pub:I:eq:B-4-21-10}
\end{equation}
therefore extend holomorphically.  A meromorphic function on
the smooth ambient space with no codimension-one pole is
holomorphic.  Their restrictions are the two standard minimal
resolution coordinates on $D$.  Together with $H$, whose
normal derivative is nonzero, they have invertible differential
along $D$.  Thus \eqref{pub:I:eq:B-4-21-10} and $H$ are biholomorphic
product charts near $D$.  Their character identities give the
standard transitions of the minimal resolution of $A_{11}$.

\item \emph{The closing and the uniform linear collar.}
The Kato formula determines the global closing. In its blowdown coordinate $x$
and $h=w-q$, that formula is
\[
 (x,h)\longmapsto
 \bigl(\kappa(q)x(a(q)h+q),a(q)h\bigr).
\]
On the regularized collar this reads
\begin{equation}
 (x,H)\longmapsto
 \bigl(\kappa(q)q x(1+a(q)H),a(q)H\bigr).
\label{pub:I:eq:B-4-21-11}
\end{equation}
Shrink the parameter disc so $|a(q)|\le A<1$.
For $|H|<\eta<1$ the product
\begin{equation}
 G(q,H)=\prod_{j\ge1}(1+a(q)^jH)
\label{pub:I:eq:B-4-21-12}
\end{equation}
converges locally uniformly in both variables, is holomorphic
and nonzero, and satisfies
\[
 G(q,a(q)H)=\frac{G(q,H)}{1+a(q)H}.
\]
Uniform convergence follows from
$\sum_{j\ge1}|a(q)^jH|\le\eta A/(1-A)$; the factors have
no zero in the chosen neighbourhood, and the convergent
logarithm series gives a nonzero product and holomorphic inverse.
In the new character $Z=xG(q,H)$, \eqref{pub:I:eq:B-4-21-11} becomes
\begin{equation}
 (Z,H)\longmapsto
       (\kappa(q)qZ,a(q)H).
\label{pub:I:eq:B-4-21-13}
\end{equation}
Multiplication of a toric character by this unit extends on
the product charts \eqref{pub:I:eq:B-4-21-10}: it multiplies $X_i$ by the
unit and $Y_i$ by its inverse.  It therefore gives 
isomorphisms at all twelve nodes as well as on the annular
core.  The exact closing \eqref{pub:I:eq:B-4-21-13} identifies these charts
with the Tate line of factor $a(q)$, over the Tate family
of multiplier $\kappa(q)q$.  Their inverse maps obey the
same identities.  Compactness of the central divisor and
a finite cover by core and node charts allow one common
shrinking of the base and the collar neighbourhood.
This proves the claimed neighbourhood isomorphism.

\item Every ambient modification used above is a blow-up of a smooth
closed centre and is proper.  The total space remains smooth,
and its map to the parameter disc is flat because a nonzero
base function is a non-zero-divisor in its smooth local rings.
All centres lie over $s=0$, so the punctured family is unchanged.
All centres are disjoint from the persistent elliptic axis,
so its entire neighbourhood is unchanged; in particular its
central normal degree remains $-1$.

\end{enumerate}
\end{proof}

\subsection{A local resolution retaining the cusp factor}

\begin{proposition}[The local plumbing model]\hypertarget{tag-030R}{}\label{push:src:B.4.22}
\label{push:B422}
\StatementRecord{push:src:B.4.22}{proved}{030R}{\SeriesRef{push:src:B.4.21}, \SeriesRef{push:src:B.4.19}, \SeriesRef{push:src:B.4.20}, \SeriesRef{push:src:B.4.30}, \SeriesRef{push:src:B.4.33}}

Let
\[
 X_n=\{fuv=g s^n\}\subset\mathbb C^5,
 \qquad r=g-f,\qquad q=s^n.
\]
There is a proper bimeromorphic local modification
$\widetilde X_n\to X_n$ with smooth total fourfold
and reduced normal-crossings divisor $r=0$.
For fixed $r\ne0$ it resolves the cyclic base change
of a three-branch $q$-normal-crossings chart.
The divisor $s=0$ after the displayed cyclic base change
is reduced; the divisor of the original function $q=s^n$
has multiplicity $n$.
\end{proposition}



\begin{proof}
\begin{enumerate}[label=\arabic*.,leftmargin=*]
\item For $n=1$ first blow up $(u,q)$ and then the pullback
of $(f,g)$. The charts of the resulting modification are
\[
\begin{array}{c|c|c}
 \text{coordinates}&\text{map to }X_1&r\\\hline
 (q,U,v,f)&u=qU,\ g=fUv&f(Uv-1)\\
 (u,Q,G,f)&q=uQ,\ v=GQ,\ g=fG&f(G-1)\\
 (u,v,F,g)&q=uvF,\ f=gF&g(1-F).
\end{array}
\]
To derive them, the $q$ chart of the first blow-up
has $u=qU$, hence $g=fUv$; the ideal $(f,g)$ is
already principal there. Its other chart has
$q=uQ$ and equation $fv=gQ$. Blowing up $(f,g)$
gives the last two charts. These are the 
blow-up charts, with the common factors cancelled
on the closure of the dense torus.
They are smooth. In the first chart $Uv-1$ is a
smooth coordinate along its zero set, and in the
others $G-1$ or $1-F$ is a coordinate. Thus the
displayed products prove the reduced NC assertion
for $r$.

\item For arbitrary $n$, pull back this proper modification
by $q=s^n$. The first chart remains smooth. The
second is $uQ=s^n$ times the two coordinates $G,f$.
Resolve its $A_{n-1}$ factor by the unit-interval fan.
The third chart is
\[
 uvF=s^n
\]
times the independent coordinate $g$. Its cone lattice
may be written as $N=\mathbb Z^3$ with coordinates
$(a,b,d)$, where the valuations of $u,v,F,s$ are
$a,b,nd-a-b,d$. The cone is the cone over
$n\Delta_2$. Insert all rays $(i,j,1)$ with
$i,j\ge0$ and $i+j\le n$, and triangulate by the
unit triangles
\[
\begin{aligned}
 &(i,j),(i+1,j),(i,j+1) &&(i,j\ge0,\ i+j\le n-1),\\
 &(i+1,j),(i,j+1),(i+1,j+1)
                       &&(i,j\ge0,\ i+j\le n-2).
\end{aligned}
\]
Every determinant is one in absolute value; the
triangles partition the large triangle and give a
proper smooth toric subdivision. Their number is
$n^2$. All new rays have $s$-height one.

\item The open set $F\ne0$ is the toric open corresponding
to $nd-a-b=0$. Its fan is the unit subdivision of
the edge $i+j=n$, hence the minimal $A_{n-1}$
resolution times $\mathbb C^*$. This is precisely
the second chart's resolution under $Q=vF$.
Where that chart meets the first, the $A$-type
resolution is an isomorphism. These facts glue
the local proper resolutions to a proper
modification of the whole local $X_n$.

\item The function $F$ is a toric monomial. Above $F=1$
it is submersive: in a regular toric chart a
nonzero monomial has nonzero derivative in a
coordinate with positive exponent. Equivalently,
on $F\ne0$ it is the displayed torus projection.
Thus $r=g(1-F)$ has two smooth transverse
branches in the third chart after resolution.
The second chart has $r=f(G-1)$ independently
of the resolved surface coordinates, and the
first was already checked. This proves smoothness
and reduced normal crossings over $r$ for all $n$.

\item For $r\ne0$ in the third chart,
$g=r/(1-F)$, so the local fibre is the resolved
threefold $uvF=s^n$ with $F\ne1$.
For $n=1$, its $q$ divisor is $uvF=0$, a literal
three-branch normal-crossings chart. For $n=12$
the standard triangulation gives 144 unimodular
charts above this one triple point. This is the
local cyclic-cover count at each cusp triple point. The other
potential singularity, at $f=-r$, is included
in this construction: the first blow-up small-resolves it
for $n=1$, and the second is an isomorphism there
because $f$ is nonzero.

\item The equation arises from the modified neighbourhood of \ref{push:src:B.4.21}. There $h=uv-q$ in the
node chart before modification, and $H=h/q$ is holomorphic
in the modified collar. The unit inverse-normal plumbing
$Hf=r$ is exactly
\[
 f(uv-q)=rq,
 \qquad fuv=(f+r)q.
\]
Putting $g=f+r$ gives $X_n$, including its factor $q$.
The modified collar, where $H$ is small,
identifies with the third chart near $F=1$ by
\[
 F=(1+H)^{-1},\qquad g=f(1+H),\qquad r=fH.
\]
There $uvF=s^n$ is exactly the resolved equation
$uv=s^n(1+H)$ from \ref{push:src:B.4.21}. The charts near $F=0$
are additional local completion charts.


\item The first collar action is $H\mapsto aH$, $f\mapsto a^{-1}f$,
which fixes $r=fH$. The function $G(q,H)$ of \ref{push:src:B.4.21}
conjugates this action on its stated small collar. For fixed $q$,
the product converges locally uniformly for $H\in\C$, because
$\sum_j|a(q)^jH|$ does so on compact sets. At $H=-a^{-j}$ it
vanishes, and the coordinate change $Z=xG(q,H)$ has zero derivative
in $x$. For $b\in\C^*$, on the open set where $r\ne0$, $f\ne0$ and
$W=1+r/f\notin\{1,b^{-1}\}$, the map $f=r/(W-1)$ conjugates
multiplication $W\mapsto bW$ to the rational map
\[
 f\longmapsto\frac{rf}{(b-1)f+br}.
\]
Indeed substituting $b(1+r/f)$ into $r/(W-1)$ gives the displayed
formula. This is the exact coordinate comparison for the test variable $W$.
\end{enumerate}
\end{proof}

\begin{proposition}[Exceptional-curve nonclosure and persistence in a semistable cap]\hypertarget{tag-030S}{}\label{push:src:B.4.23}
\StatementRecord{push:src:B.4.23}{proved}{030S}{\SeriesRef{push:src:B.4.21}, \SeriesRef{push:src:B.4.10}; \citet[Section 1.2]{DlousskyTeleman2012AB}; \citet[Proposition 2.3 and its proof, pp.~6--7 of the
original preprint]{DlousskyTeleman2012AB}; \citet{Merkulov1995D}}

Let $\mathcal K^*\to\Delta_s$ be the modified degree-one
Kato family of \ref{push:src:B.4.21}, with $q=s^{12}$, and let
$\mathcal E\to\Delta_s$ be its unchanged persistent elliptic
axis. For $s\ne0$ let $C_s\simeq\PP^1$ be the exceptional
curve of the single Hopf blow-up.

The union of the $C_s$ has no proper analytic divisor
closure in $\mathcal K^*$. Further finite proper modifications
cannot give these curves a proper analytic divisor closure:
any such closure would push forward to a forbidden one.
In particular, the generic Hopf blowdown does not extend
as a relative proper holomorphic modification, or as a
relative meromorphic contraction represented by a proper
analytic graph, across $s=0$.

For the deformation obstruction, suppose a smooth complex
fourfold $\mathcal Y$ has holomorphic functions $(Q,r)$,
its divisor $r=0$ is reduced normal crossings, and one
component $K$ contains a neighbourhood of some $C_{s_0}$,
$s_0\ne0$, agreeing with the complementary Kato region
and its ordinary base coordinate $Q$. Suppose this curve
is disjoint from every conductor branch of $r=0$.
Then every sufficiently nearby fibre $(Q,r)=(Q_0,r_0)$
contains an embedded rational curve. Consequently such a
cap cannot have the prescribed ordinary two-tori as its
$r\ne0$ surface fibres.
\end{proposition}


\begin{proof}
\begin{enumerate}[label=\arabic*.,leftmargin=*]
\item We first use the exact marked blowdown of \ref{push:src:B.4.10}.
The exceptional curve of the single generic Hopf blow-up
meets the persistent elliptic axis transversely at
\[
 [v=-q]\in E_{a(q)}=\C^*/a(q)^{\Z}.
\]
All modifications in \ref{push:src:B.4.21} are disjoint from that axis.
The intersection section in the modified family is therefore
still
\begin{equation}
 p(s)=[-s^{12}].
\label{pub:I:eq:B-4-23-1}
\end{equation}
The units $a(q)$ obey $0<|a(q)|<1$. Choose a holomorphic
logarithm on the full parameter disc and put
\[
 \tau(s)=\frac{\log a(s^{12})}{2\pi i}.
\]
Then $\mathcal E_s=\C/(\Z+\tau(s)\Z)$ and
$\operatorname{Im}\tau(s)>0$. The maps
$(x,y,s)\mapsto[x+y\tau(s)]$, with $(x,y)\in\R^2/\Z^2$,
give a smooth trivialization of the elliptic family.
A logarithmic lift of \eqref{pub:I:eq:B-4-23-1} is
\[
 z(s)=\frac{\log(-s^{12})}{2\pi i}.
\]
On a positive loop around zero, $\tau$ returns to itself
and $z$ changes by $12$. Thus the loop of the section has
class $(12,0)\ne0$ in the integral first homology of this
trivialized family. A continuous section over the disc
would fill that loop. Hence $p$ has no continuous, and
therefore no holomorphic, extension. A meromorphic section
from a nonsingular curve would extend holomorphically:
its proper graph is generically degree one over the curve,
and normalization gives a finite degree-one map to that
normal base. This is equally impossible. A finite base
change of degree $d>0$ multiplies the winding to $(12d,0)$
and does not change the conclusion.

\item Suppose an analytic divisor $Z$ in $\mathcal K^*$, proper
over the disc, restricted over the punctured disc to the
union of the $C_s$. Its horizontal irreducible part does
not contain $\mathcal E$, since the generic intersection
is the one transverse point \eqref{pub:I:eq:B-4-23-1}. Both are Cartier
locally in the smooth threefold, so their intersection
is an analytic curve, proper over the disc. Its reduced
horizontal part has generic degree one. An irreducible
horizontal analytic curve cannot have a positive-dimensional
fibre without being contained in that fibre. Hence its
normalization is finite of degree one over the disc and
is isomorphic to the disc. It extends $p(s)$, a
contradiction. Vertical components of the intersection
have no effect on this horizontal component.

\item If an additional finite proper modification admitted a
proper analytic closure, its proper image would give
such a $Z$, by the proper mapping theorem; the image has
dimension two and the same generic exceptional curves.
The horizontal exceptional locus of a relative proper
holomorphic blowdown likewise gives $Z$. For a meromorphic
contraction with proper analytic graph, resolve the graph
and apply the same argument to its horizontal exceptional
locus and its proper image in $\mathcal K^*$.
This is a direct marked-intersection proof of a familiar
nonproperness phenomenon for exceptional curves in GSS
surface deformations; see
\citet[arXiv version~2, Section~1.2, PDF pp.~3--5]{DlousskyTeleman2012AB} for the primary
context.

\item Write
$C=C_{s_0}$ and $S=K\cap\{Q=Q_0\}$. Near this curve,
$S$ is a smooth surface and $C$ is its exceptional curve,
so
\[
 N_{C/S}\simeq\OO_{\PP^1}(-1).
\]
The base coordinate $Q$ is submersive on $K$ near $C$.
Because $C$ misses every other central branch, the reduced
function $r$ is a local defining function for $K$ near
$C$ and its normal derivative is nowhere zero.
Consequently the differential of $(Q,r)$ gives the exact
normal sequence
\begin{equation}
 0\longrightarrow\OO_{\PP^1}(-1)
 \longrightarrow N_{C/\mathcal Y}
 \xrightarrow{\ (dQ,dr)\ }
 \OO_{\PP^1}^{\oplus2}\longrightarrow0.
\label{pub:I:eq:B-4-23-2}
\end{equation}
Its extension class vanishes because
$H^1(\PP^1,\OO(-1))=0$. Hence
\begin{equation}
 N_{C/\mathcal Y}\simeq
 \OO(-1)\oplus\OO^{\oplus2},\qquad
 H^1(C,N_{C/\mathcal Y})=0,\qquad
 h^0(C,N_{C/\mathcal Y})=2.
\label{pub:I:eq:B-4-23-3}
\end{equation}
Moreover \eqref{pub:I:eq:B-4-23-2} induces an isomorphism on global
sections, since both $H^0(\OO(-1))$ and $H^1(\OO(-1))$
vanish.

\item Kodaira's embedded-deformation theorem therefore gives a
smooth two-dimensional local parameter space $T$ of
compact embedded deformations of $C$, with tangent space
$H^0(C,N_{C/\mathcal Y})$. In precisely this relative
surface setting one may apply \citet[arXiv version~2, Proposition~2.3 and its proof, PDF pp.~6--7]{DlousskyTeleman2012AB}: after shrinking
the neighbourhood of $C$, $(Q,r)$ is a holomorphic
submersion of relative dimension two, and $C$ is a
compact simply exceptional divisor of its fibre.
The general completeness theorem is due to Kodaira; we use the exact statements in \citet[Section~1, p.~71, and Section~3.1, p.~74]{Merkulov1995D}.
Shrink this family so its curves
remain in the chosen neighbourhood. Their genus remains
zero, so all are isomorphic to $\PP^1$.
The functions $Q,r$ are constant on every compact
connected curve and thus define a holomorphic map
\[
 T\longrightarrow\C^2.
\]
Its derivative at $[C]$ is precisely the isomorphism on
sections induced by \eqref{pub:I:eq:B-4-23-2}. The holomorphic inverse
function theorem shows that the map covers a neighbourhood
of $(Q_0,0)$. In particular the nearby fixed pairs with
$r\ne0$ contain rational curves. A holomorphic map from
$\PP^1$ to a complex torus lifts to its vector-space
universal cover and is constant. Thus those fibres cannot
be the prescribed ordinary two-tori.

\item The conductor condition is essential. For example, if a
further central component meets $K$ along the persistent
elliptic axis, then its intersection number with $C$ is
one and $N_{K/\mathcal Y}|_C\simeq\OO(-1)$.
At the conductor intersection $dr$ vanishes on the normal fibre
of $C$, so $(dQ,dr)$ is not the bundle epimorphism of
\eqref{pub:I:eq:B-4-23-2}. This computes precisely where that
normal-sequence argument uses the conductor-disjoint hypothesis.
\end{enumerate}
\end{proof}

\subsection{The associated logarithmic torsor at the missing origin}

\begin{proposition}[Obstruction in the fixed log-torus category]\hypertarget{tag-030T}{}\label{push:src:B.4.24}
\label{push:B424}
\StatementRecord{push:src:B.4.24}{proved}{030T}{\SeriesRef{push:src:B.4.18}, \SeriesRef{push:src:B.4.19}, \SeriesRef{push:src:B.4.9}, \SeriesRef{push:src:B.4.30}, \SeriesRef{push:src:B.4.33}; \citet{KKN2008Root}; \citep[Definition 1.2.3, pp.~77--78]{KKN2008Root}; \citep[Definition 1.2.6, p.~79]{KKN2008Root}; \citep[Sections 1.3.1--1.3.4, pp.~81--82]{KKN2008Root}; \citet[Sections 7--9]{KKN2008PicardE}}

The local Jacobian admits a nondegenerate logarithmic
complex-torus extension over the full $(q,r)$ bidisc and hence
has proper analytic models in the sense of
\citet{KKN2008Root}. After the normalized base change
$q=Q^{12}$, the two marked axial fan presentations of
\ref{push:src:B.4.18}--\ref{push:src:B.4.19} determine a torsor under its associated
log-torus sheaf on the punctured bidisc. This torsor does
not extend across the origin as a torsor locally trivial
for the usual analytic topology. The obstruction persists
after every further finite coordinate-root base change.

On each axial chart its extending origin identifies the associated sheaf with the Jacobian log-torus sheaf $\mathcal A$. On the overlap the two identifications differ by the section $F$ in \eqref{pub:I:eq:B-4-24-3}. An extension here is a locally trivial $\mathcal A$-torsor carrying these prescribed axial identifications.
\end{proposition}



\begin{proof}
\begin{enumerate}[label=\arabic*.,leftmargin=*]
\item \emph{The Jacobian pairing.}
Retain the period columns and logarithms of \ref{push:src:B.4.18}:
\[
 \gamma=(6\mu,C-s-12v)^t,\quad
 u=(\tau,\mu)^t,\quad w=(1,0)^t,\quad\delta=(0,1)^t,
 \qquad
 s=\frac{\log q}{2\pi i},\quad
 v=\frac{\log r}{2\pi i}.
\]
Here $\tau-s,\mu,C$ are holomorphic at $q=0$.
Use multiplicative characters
$Z=\exp(2\pi i\zeta_1)$ and
$W=\exp(-2\pi i\zeta_2)$, with character lattice
$X=\Z e_1\oplus\Z e_2$ and period lattice
$Y=\Z u\oplus\Z\gamma$. Put
\[
 a_0=e^{2\pi i(\tau-s)},\qquad
 a_1=e^{2\pi i\mu},\qquad a_2=e^{-2\pi iC}.
\]
These are holomorphic units on the full $q$ disc.
The multiplicative pairing is
\begin{equation}
 \begin{array}{c|cc}
       &u&\gamma\\\hline
 e_1&q a_0&a_1^6\\
 e_2&a_1^{-1}&q r^{12}a_2.
 \end{array}
\label{pub:I:eq:B-4-24-1}
\end{equation}
Give the base the fine saturated log structure from
$qr=0$. In its sharp monoid $\N^2$, written additively,
the pairing modulo units is $\operatorname{diag}(q,q+12r)$.
We check the exact admissibility definition
\citep[Definition 1.2.3, pp.~77--78]{KKN2008Root}.
On the full face a nonnegative functional
$N(q)=h,N(r)=k$ gives
$\operatorname{diag}(h,h+12k)$. If $h>0$ it is
nondegenerate. If $h=0<k$, its kernel on each lattice
is the first basis line, exactly the face lattice
for the face generated by $q$; the quotient pairing
is multiplication by $12k$. For the $q$ face the
nonzero face lattices are these first basis lines
and the nontrivial quotient pairing is multiplication
by $h$. The $r$ face and the zero face have zero
face lattices. Zero functionals give zero-dimensional
quotients. These are all the face tests.

\item The residual unit condition is also satisfied
\citep[Definition 1.2.6, p.~79]{KKN2008Root}.
At $q=0$ the residual lattices vanish. At $r=0$,
$q\ne0$, their pairing is $\log|q a_0|\ne0$ after
shrinking the disc. At ordinary points the real
logarithmic matrix has determinant
\[
 \log|q a_0|\log|q r^{12}a_2|
                 +6\bigl(\log|a_1|\bigr)^2>0,
\]
after shrinking so both of the first two logarithms
are negative. Thus \eqref{pub:I:eq:B-4-24-1} defines a log complex
torus $\mathcal A$ by Definition 1.3.4, pp.~81--82.
Theorems 5.2.1 and 5.3.2, pp.~171--172 and 180--181,
respectively supply a complete invariant fan with
finitely many $Y$-orbits and its proper analytic
model.

\item \emph{The two extending origins.}
Set
\[
 B=\Delta_Q\times\Delta_r,\quad q=Q^{12},\qquad
 B_q=\Delta_Q^*\times\Delta_r,\quad
 B_r=\Delta_Q\times\Delta_r^*,\quad U=B_q\cap B_r.
\]
On $B_r$, \ref{push:src:B.4.9} supplies the affine section
\begin{equation}
 \sigma=
 \left(6\mu v,\,-6v^2+v(C-12\kappa_0)\right)^t.
\label{pub:I:eq:B-4-24-2}
\end{equation}
For every local branch of $\log r$, both of its
multiplicative coordinates are holomorphic units
in the full $q$ disc. Its exact $r$-deck difference
is the bounded cusp closing translation
of \ref{push:src:B.4.9}. Hence it descends as a global section
of the affine family, including $q=0$.
Translation by $-\sigma$ identifies the prescribed
cusp fan presentation with its group presentation.
The normalized toric base change $q=Q^{12}$ preserves
this identification and its extending origin.

\item On $B_q$, the rotating model of \ref{push:src:B.4.18} is the
free cyclic quotient of its root fan model. Selecting
$Q$ recovers that model; more generally the normalized
base change recovers it also at the ramification
divisor. The twelfth fractional deck is
$\gamma-6\delta$, a period, so the
multiplicative period action is exactly the second
column of \eqref{pub:I:eq:B-4-24-1}. The point $(Z,W)=(1,1)$
extends in the root fan at $r=0$ and gives the
group origin. Thus both axial presentations have
extending origins in their associated
$\mathcal A$-sheaves. The identity on the character and period lattices identifies the associated quotient sheaves under a fan refinement: both are defined by the same multiplicative pairing \eqref{pub:I:eq:B-4-24-1}.

\item In these two origins the exact comparison
$\zeta\mapsto\zeta-\sigma+F$ of \ref{push:src:B.4.18} becomes
translation by
\begin{equation}
 F(s,v)=\left(
 \frac{\mu(s+12v)}2,
 \frac{C(s+12v)}{12}-\frac{(s+12v)^2}{24}+6v
 \right)^t.
\label{pub:I:eq:B-4-24-3}
\end{equation}
This defines the associated $\mathcal A$-torsor
by gluing trivial torsors on $B_q$ and $B_r$.
The differences are exactly
\begin{equation}
 F(s+12,v)-F(s,v)=\gamma-6\delta,
 \qquad F(s,v+1)-F(s,v)=\gamma.
\label{pub:I:eq:B-4-24-4}
\end{equation}
They are periods, so $F$ is an $\mathcal A$-section
on $U$.

\item \emph{The ghost-valuation obstruction.}
By the quotient-sheaf definition
\citep[Sections 1.3.1--1.3.4, pp.~81--82]{KKN2008Root},
on the underlying analytic site there is an exact sequence
\begin{equation}
 0\longrightarrow Y\longrightarrow
 \mathcal T_{\log}^{(Y)}\longrightarrow
 \mathcal A\longrightarrow0.
\label{pub:I:eq:B-4-24-5}
\end{equation}
The middle term is the sheaf of bounded multiplicative
log characters. On $U$ it is the ordinary
two-dimensional multiplicative torus sheaf.
In particular \eqref{pub:I:eq:B-4-24-4} says that the connecting
class of $F$ in $H^1(U,Y)$ has $r$-loop component
$\gamma$. The terms in $\Z w+\Z\delta$ disappear
on exponentiation, and must not be included in $Y$.

\item Let $i:D_Q=\{Q=0\}\hookrightarrow B_r$.
Taking the $Q$-order of each character gives
\begin{equation}
 \operatorname{ord}_Q:
 \mathcal T_{\log}^{(Y)}\longrightarrow i_*\Z^2,
 \qquad
 u\longmapsto(12,0),\quad
 \gamma\longmapsto(0,12)
\label{pub:I:eq:B-4-24-6}
\end{equation}
on its period subsheaf. If $h$ is any
$\mathcal A$-section on $B_r$, local lifts in
the middle sheaf of \eqref{pub:I:eq:B-4-24-5} represent its
connecting class by a $Y$-cocycle. Applying
\eqref{pub:I:eq:B-4-24-6} makes that cocycle a coboundary,
the difference of the local ghost vectors.
But restriction identifies $H^1(B_r,Y)$ with
$H^1(D_Q,Y)$, since $B_r$ retracts onto $D_Q$;
the induced map to $H^1(D_Q,\Z^2)$ is the
injective diagonal map $(12,12)$. Therefore
the connecting class of $h$ is zero. In
particular its $r$-loop component is zero.

\item For a section on $B_q$ that component is zero
as well, because the $r$ loop contracts in $B_q$.
Consequently $F$ cannot be the difference of
sections from $B_q$ and $B_r$. The glued torsor
is nontrivial on $B\setminus\{0\}$. An extension
locally trivial near the origin would have a
section there; after shrinking the bidisc its
two coordinate expressions would give precisely
the forbidden splitting. This proves nonextension.

\item For $q=Q^{12m}$ and $r=R^n$, with $m,n>0$,
the $R$-loop difference is
\[
 F(s,v+n)-F(s,v)
    =n\gamma-6n(n-1)\delta.
\]
Its $Y$-component is $n\gamma\ne0$,
whereas the $Q$-ghost map on $Y$ is
$\operatorname{diag}(12m,12m)$ and remains
injective. The same proof applies after every
such cover. Finite Kummer covers of the standard
coordinate log structure are locally dominated
by coordinate-root covers, so allowing them
does not remove this obstruction either.

\item \emph{The ordinary gauge.}
The exact identity in Proposition~\ref{push:src:B.4.19} is
\begin{equation}
 F-\sigma-kv\delta=G(s)+6v\delta,
 \quad k=-s+12\kappa_0,
 \quad G(s)=\left(\frac{\mu s}2,
                  \frac{Cs}{12}-\frac{s^2}{24}\right)^t.
\label{pub:I:eq:B-4-24-7}
\end{equation}
On $q\ne0$, this is the raw overlap in the ordinary gauge. Its second multiplicative character is
\[
 e^{-2\pi i kv}=Q^{12v}e^{-24\pi i\kappa_0v},\qquad q=Q^{12}.
\]
The second factor is a holomorphic unit on a fixed branch of $\log r$. A bounded logarithmic character at $Q=0$ has integral $Q$-order, while $12v$ varies with $r$. Multiplication by a period changes that order by an integer. Thus this gauge cannot give an extending log-torus section along the full $Q$-axis. Changes of the two extending axial origins instead lie in the image of
\[
 H^0(B_q,\mathcal A)\oplus H^0(B_r,\mathcal A)
 \longrightarrow H^0(U,\mathcal A),\qquad(h_q,h_r)\longmapsto h_r-h_q.
\]
The ghost calculation shows that $F$ lies outside this image.

\end{enumerate}
\end{proof}

\begin{proposition}[Noncompletion of a one-eigenvalue Hopf boundary]\hypertarget{tag-030U}{}\label{push:src:B.4.25}
\label{push:B425}
\StatementRecord{push:src:B.4.25}{proved}{030U}{\SeriesRef{push:src:B.4.21}, \SeriesRef{push:src:B.4.19}; \citep[Section 6, Hauptsatz I, p.~59]{Grauert60}; \citet[Theorem~1.6 and the following noncompact extension, p.~215]{BM97}}

Let $\alpha(q)=q^m\kappa(q)$, where $m\ge1$ and
$\kappa$ is a holomorphic unit on a disc, and let $\beta$
be holomorphic with $0<|\beta(0)|<1$. Shrink the disc so
$\beta$ is a unit and $0<|\alpha(q)|<|\beta(q)|<1$ for
$q\ne0$. Consider the diagonal primary Hopf surfaces
\begin{equation}
 H_q=(\C^2\setminus\{0\})/
       \langle(z,w)\mapsto(\alpha(q)z,\beta(q)w)\rangle.
\label{pub:I:eq:B-4-25-1}
\end{equation}
Their holomorphic family over the punctured disc has no
proper complex-analytic completion over the full disc,
even allowing singular or nonreduced central fibres and
a singular total space. This remains true after every
finite base change ramified at zero.

In particular one cannot replace the blown-up Hopf
fibres in \ref{push:src:B.4.21} by the unblown family with one multiplier
$\kappa(q)q\to0$ and the other tending to a nonzero
contraction.
\end{proposition}


\begin{proof}
\begin{enumerate}[label=\arabic*.,leftmargin=*]
\item \emph{Anticanonical sections of a diagonal Hopf surface.}
A section of $K_{H_q}^{-1}$ lifts to
$f(z,w)\,\partial_z\wedge\partial_w$ on
$\C^2\setminus\{0\}$. Hartogs extension makes $f$
entire on $\C^2$. Deck equivariance is exactly
\begin{equation}
 f(\alpha z,\beta w)=\alpha\beta f(z,w).
\label{pub:I:eq:B-4-25-2}
\end{equation}
For the Taylor coefficient of $z^i w^j$, with $i,j\ge0$,
this says
\begin{equation}
 \alpha^{i-1}\beta^{j-1}=1.
\label{pub:I:eq:B-4-25-3}
\end{equation}
When $i,j\ge1$, the contraction inequalities force
$i=j=1$. When $j=0$, \eqref{pub:I:eq:B-4-25-3} would give
$\beta=\alpha^{i-1}$, which is impossible under
$|\alpha|<|\beta|<1$ (and also for $i\le1$).
When $i=0$, the only possibility is
\begin{equation}
 \alpha=\beta^N,\qquad j=N+1,\qquad N\ge2.
\label{pub:I:eq:B-4-25-4}
\end{equation}
At a fixed parameter there is at most one such $N$,
since $0<|\beta|<1$. Thus
\begin{equation}
 h^0(H_q,K_{H_q}^{-1})=
 \begin{cases}
 2,&\alpha(q)=\beta(q)^N\text{ for some }N\ge2,\\
 1,&\text{otherwise}.
 \end{cases}
\label{pub:I:eq:B-4-25-5}
\end{equation}
The sections in the first case are represented by
$zw$ and $w^{N+1}$. The computation uses 
the displayed bivector automorphy.

\item \emph{Infinitely many accumulating resonances.}
Put $b=|\beta(0)|\in(0,1)$ and $k=|\kappa(0)|>0$.
For all sufficiently small $q$ we may arrange
$|\kappa(q)|\ge3k/4$ and
\[
 |\beta(q)|\le b\exp(C|q|)
\]
for some fixed $C$, by a holomorphic logarithm of
$\beta(q)/\beta(0)$. Set
\[
 R_N=\left(\frac{2b^N}{k}\right)^{1/m}.
\]
Then $R_N\to0$ and $N R_N\to0$. On $|q|=R_N$, for
all sufficiently large $N$,
\[
 |q^m\kappa(q)|\ge\frac32 b^N,
 \qquad |\beta(q)^N|\le b^N\exp(CN R_N)<\frac54 b^N.
\]
Rouch\'e's theorem shows that
$q^m\kappa(q)-\beta(q)^N$ has $m$ zeros counted with
multiplicity in $|q|<R_N$. None is zero. Choose one,
$q_N$. Different $N$ give different points, since a
nonzero contraction cannot have two equal distinct
positive powers. Hence there are infinitely many
resonances $q_N\to0$ with the first value in
\eqref{pub:I:eq:B-4-25-5}.

\item For each $N$ the resonance function is nonzero at
$q=0$, so is not identically zero. Its zeros in the
punctured disc are discrete. The countable union of
these discrete sets cannot contain an open disc.
There are therefore also nonresonant points arbitrarily
near zero, with the second value in \eqref{pub:I:eq:B-4-25-5}.

\item \emph{Coherent direct images forbid this pattern.}
Suppose first that a proper completion
$\pi:\mathcal X\to\Delta$ has a smooth connected
total threefold and agrees with \eqref{pub:I:eq:B-4-25-1} over
$\Delta^*$. Put $L=K_{\mathcal X}^{-1}$.
On each nonzero fibre, adjunction identifies
$L|_{\mathcal X_q}$ with $K_{H_q}^{-1}$, because
the normal direction is the trivial base direction.
Grauert's proper direct-image theorem gives coherent
sheaves
\[
 M=\pi_*L,\qquad N=R^1\pi_*L
\]
on the disc
\citep[Section 6, Hauptsatz I, p.~59]{Grauert60}.
Multiplication by a nonzero base function is injective
on $L$, and hence on $M$, since the total space is
smooth and the map is nonconstant. Thus $M$ is
locally free of a fixed rank on the disc. The torsion
subsheaf of $N$ is coherent with locally finite
support. After shrinking, it has no support in
$\Delta^*$.

\item For any $a\in\Delta^*$, apply the direct-image long
exact sequence to
\[
 0\longrightarrow L\xrightarrow{\ q-a\ }L
 \longrightarrow L|_{\mathcal X_a}\longrightarrow0.
\]
At $a$ its beginning yields
\begin{equation}
 0\longrightarrow M_a/(q-a)M_a
 \longrightarrow H^0(\mathcal X_a,L|_{\mathcal X_a})
 \longrightarrow \ker(q-a:N_a\to N_a)
 \longrightarrow0.
\label{pub:I:eq:B-4-25-6}
\end{equation}
The last term vanishes. Therefore
$h^0(H_a,K_{H_a}^{-1})=\operatorname{rank}M$ for
\emph{every} sufficiently small $a\ne0$. This
contradicts \eqref{pub:I:eq:B-4-25-5} and the accumulating resonances.
The displayed coherent direct-image exact sequence proves the
required local constancy on the punctured disc.

\item Now suppose an arbitrary proper analytic completion
$\mathcal Z\to\Delta$ exists. Pass to its reduction
and retain the unique irreducible component containing
the connected smooth punctured total space. This is a
closed analytic component, remains proper over the
disc, and dominates the disc. Its restriction over
$\Delta^*$ is exactly the given family. Resolve this
component by a proper analytic resolution chosen to
be an isomorphism on its already nonsingular locus;
see \citet[Theorem~1.6 and the following noncompact extension, p.~215]{BM97}.
All singular points lie over zero, so the resolution
is an isomorphism over the entire punctured family.
Properness of the composition is preserved. Near the
compact central fibre, local finiteness of the
analytic resolution permits one shrinking of the
disc; its dominating smooth component is a smooth
proper total threefold as in the preceding paragraph.
Its map to the disc is nonconstant, so the base
function is a non-zero-divisor on this dominating component.
The preceding contradiction applies. Reduction also
shows that nilpotent structure supported at zero
cannot evade it.

\item Finally, a finite base change replaces $\alpha$ by
a holomorphic function of some positive vanishing
order and replaces $\beta$ by a unit whose value
at zero still has modulus strictly between zero and
one. The same proof applies verbatim.
\end{enumerate}
\end{proof}

\begin{remark}[The toric ray and blow-up comparisons]
\hypertarget{tag-PT19}{}\label{I:Hopf-ray-blowup-comparisons}
\StatementRecord{I:Hopf-ray-blowup-comparisons}{proved}{PT19}{\citet[\S13, pp.~116--118]{Oda1978D}; \cref{push:src:B.4.10,push:src:B.4.25}}
For the toric ray datum of \citet[\S13, retypeset pp.~116--118]{Oda1978D}, put
\[
 v_1=n,\qquad v_2=-n+a n',\qquad a>0.
\]
The original period is $\lambda^{n'}$. The identity $a n'=v_1+v_2$ shows that the base change $\lambda=t^a$ acts in the universal-cover coordinates of $v_1,v_2$ by the two multipliers $(t,t)$, up to the fixed torus units of that construction. Their two vanishing orders are therefore $(1,1)$. The family in Proposition~\ref{push:src:B.4.25} instead has the explicitly specified orders $(m,0)$, with second multiplier tending to the nonzero contraction $\beta(0)$.

The generic fibres of Proposition~\ref{push:src:B.4.10} carry the modification $\pi:\widetilde H_q\to H_q$ formed by successive blow-ups at the marked axis points. At one blow-up, $K_{\widetilde H_q}^{-1}=\pi^*K_{H_q}^{-1}\otimes\mathcal O(-E_{\rm exc})$, so its sections correspond exactly to anticanonical sections on $H_q$ vanishing at that centre. Iterating gives the vanishing condition at each successive infinitely near centre. This is the precise map from the anticanonical spaces of the blown-up fibres to the spaces computed in \eqref{pub:I:eq:B-4-25-5}.
\end{remark}

\subsection{A necessary vertical fibre over the missing origin}

\begin{proposition}[Every origin component is a threefold]\hypertarget{tag-030V}{}\label{push:src:B.4.26}
\label{push:B426}
\StatementRecord{push:src:B.4.26}{proved}{030V}{\SeriesRef{push:src:B.4.19}, \SeriesRef{push:src:B.4.30}, \SeriesRef{push:src:B.4.33}, \SeriesRef{push:src:B.4.24}; \citet[Chapter II, Sections 7--8]{Demailly2012E}}

Suppose the partial family of \ref{push:src:B.4.19} has a
proper analytic completion $f:\mathcal Z\to\Delta_q\times\Delta_r$
with smooth pure four-dimensional total space, with
$r\circ f$ flat, and with its prescribed identification
over the punctured bidisc. The map to this fixed bidisc
is assumed holomorphic. Then every irreducible component
of $f^{-1}(0,0)$ has dimension three. In particular,
there is no equidimensional completion, and no completion
having even one two-dimensional component of the origin
fibre.

The stated map to the bidisc includes a holomorphic extension of
$q$. Proposition~\ref{push:src:B.4.30} excludes completions of this fixed
partial family, including the displayed vertical-component geometry;
Theorem~\ref{push:src:B.4.33} excludes proper fillings of the global family.
\end{proposition}



\begin{proof}
\begin{enumerate}[label=\arabic*.,leftmargin=*]
\item Use the log-torsor obstruction of Proposition~\ref{push:src:B.4.24}
after the normalized finite base change $q=Q^{12}$.
The resulting space $\mathcal Z'$ is normal and
pure four-dimensional. Its restrictions to the two
punctured axes have the explicit normalized fan
presentations used in \ref{push:src:B.4.24}. Write
\[
 B=\Delta_Q\times\Delta_r,\quad
 B_q=\Delta_Q^*\times\Delta_r,\quad
 B_r=\Delta_Q\times\Delta_r^*,\quad
 U=B_q\cap B_r.
\]

\item Assume there is a finite surjective map
$p:V\to B$ of positive generic degree $d$, together
with a holomorphic map $V\to\mathcal Z'$ over $B$.
We may normalize $V$ and retain any collection of
components dominating $B$; its degree remains a
positive integer. A normal surface is Cohen--Macaulay;
finite over the smooth bidisc, it is flat. Thus the
finite holomorphic trace and the degree-$d$ sum with
multiplicities are available also at branch points.
On the ordinary locus, sum the
$d$ points of this multisection, with multiplicity,
in each of the two axial group origins of
\ref{push:src:B.4.24}. This gives two local expressions for the
trace in the Jacobian group. They extend as log-torus sections on $B_q$ and $B_r$.

\item \emph{Ramification away from the axes.}
Over the ordinary locus the torus family is a smooth
holomorphic group family. Near a branch point, choose
small neighbourhoods of the finitely many preimages
in $V$ and local Lie lifts of their maps to the torus.
The trace of these holomorphic functions under the
finite map is holomorphic, and its exponential is
the ordinary sum. Changes of lifts add periods,
which disappear in the group quotient. Thus the
ordinary trace extends across the ordinary branch
locus. Equivalently this is the usual addition map
on the relative symmetric power, applied to the
finite multisection.

\item \emph{The explicit norm at an axis.}
After shrinking $B$, every component of the finite
cover's branch divisor which is not an axis meets
each punctured axis nowhere: its intersection with
either axis is an isolated analytic germ at the
origin. At a point of a punctured axis choose
coordinates $(x,y)$, with $x$ the normal coordinate
to that axis. Locally the normal finite cover is
a disjoint union of covers
\begin{equation}
 x=t^e,\qquad y=y,
\label{pub:I:eq:B-4-26-1}
\end{equation}
including $e=1$ sheets. Indeed the complement of
the axis in a small polydisc has cyclic fundamental
group; each connected unramified finite cover is
the displayed root cover there, and its normal
finite extension is unique.

\item The map from such a root polydisc into the
prescribed fan model gives a section of its
associated log-torus sheaf. It has a local bounded
multiplicative log lift. This can also be seen
directly in a toric chart: the pulled-back characters
are meromorphic functions with zeros and poles
supported on $t=0$, hence are $t^k u(t,y)$ with
$k\in\Z$ and $u$ a holomorphic unit. The quotient
by the period lattice causes only the period
ambiguity already retained in \ref{push:src:B.4.24}.

\item The multiplicative norm of a character on \eqref{pub:I:eq:B-4-26-1}
is explicitly
\begin{equation}
 \prod_{j=0}^{e-1}(\zeta_e^j t)^k
              u(\zeta_e^j t,y)
 =(-1)^{k(e-1)}x^k
              \prod_{j=0}^{e-1}u(\zeta_e^j t,y).
\label{pub:I:eq:B-4-26-2}
\end{equation}
The last product is a holomorphic unit in $(x,y)$.
Thus the norm has an integral ghost vector and
is a log character on the base. It is
bounded in the required sense: on the $r$ axis
the first character has zero normal order and
the second is unrestricted; on the $Q$ axis
both are unrestricted. These are precisely the
bounded-character conditions for \eqref{pub:I:eq:B-4-24-1}.
Multiplication of lifts corresponds to addition
in the log torus. Changing a lift by a period
changes its norm by an integral period, so the
resulting quotient section is independent of
the lift. Summing over the finitely many covers
\eqref{pub:I:eq:B-4-26-1}, whose degrees add to $d$, proves the
claimed extension along that axis.

\item The local extensions agree on overlaps because
they agree on the ordinary locus. Lift locally to log characters:
agreement there is an equality up to a locally constant $Y$-period,
and that meromorphic equality extends across the axis.
We have therefore obtained sections $h_q$ of $\mathcal A$ on
$B_q$ and $h_r$ on $B_r$.
To fix the sign, write $z_r=\zeta-\sigma$ for the cusp group
coordinate. The comparison map of \ref{push:src:B.4.18} gives the
rotating coordinate $z_q=\zeta-\sigma+F=z_r+F$.
Summing over the $d$ points, with multiplicity, therefore gives
\begin{equation}
 h_q-h_r=dF
\label{pub:I:eq:B-4-26-3}
\end{equation}
on $U$. The ghost proof of \ref{push:src:B.4.24} forbids
\eqref{pub:I:eq:B-4-26-3}: its $r$-loop $Y$-component is
$d\gamma\ne0$. Thus no such finite multisection can exist.

\item \emph{A two-dimensional fibre germ would supply one.}
Suppose an origin point of $\mathcal Z$ has local
fibre dimension two. Finite base change and
normalization are finite maps and preserve that
local fibre dimension at points lying above it.
At a resulting point of $\mathcal Z'$, choose
two generic local linear sections in an analytic
embedding. Their intersection with the fibre is
isolated. Intersecting with $\mathcal Z'$ gives
a two-dimensional analytic germ whose map to
$B$ is finite. This is the local finite-projection
criterion from Weierstrass preparation; see also
\citet[Chapter~II, Proposition~8.2 and its proof, pp.~116--117]{Demailly2012E}.
Choose a component through the point. Its finite
image has dimension two and is therefore the
whole base germ. A finite representative of this
germ, after shrinking the bidisc, is finite and
surjective over the entire smaller bidisc. Its
normalization is still finite and maps to
$\mathcal Z'$, giving exactly the forbidden
multisection of positive degree. Properness of
this finite representative prevents loss of
sheets at its boundary.

\item Every fibre component of a dominant map from a
pure fourfold to a smooth surface has dimension
at least two. If a two-dimensional irreducible
component existed, a generic point on it would
have local fibre dimension two, which we have
excluded. Flatness over $r$ makes $r=0$ a
Cartier divisor of pure dimension three, so
the origin fibre has no component of dimension
greater than three. All its irreducible
components consequently have dimension three,
as asserted.
\end{enumerate}
\end{proof}

\subsection{The obstruction survives modifications of the base}

\begin{proposition}[Logarithmic torsors under point blowups]\hypertarget{tag-030W}{}\label{push:src:B.4.27}
\label{push:B427}
\StatementRecord{push:src:B.4.27}{proved}{030W}{\SeriesRef{push:src:B.4.24}, \SeriesRef{push:src:B.4.30}, \SeriesRef{push:src:B.4.33}; \citep[Lemma 1.2.9, pp.~80--81]{KKN2008Root}}

Let $B=\Delta_Q\times\Delta_R$, with
$q=Q^{12m}$ and $r=R^n$, where $m,n$ are positive
integers. Let $\beta:S\to B$ be any finite sequence
of point blowups whose centres lie over the origin.
Give $S$ the divisorial log structure of the reduced
inverse image of $QR=0$, and let $\mathcal A_S$ be
the pullback of the Jacobian log torus of
\ref{push:src:B.4.24}. Then
\begin{equation}
 H^1(S,\mathcal A_S)=0
\label{pub:I:eq:B-4-27-1}
\end{equation}
for the usual analytic topology. In particular the
associated axial log-torus torsor on the
punctured bidisc has no locally trivial
$\mathcal A_S$-torsor extension after these
modifications.

Writing $U_S=\beta^{-1}(B\setminus\{0\})$, the extension datum is an $\mathcal A_S$-torsor whose restriction along the isomorphism $\beta|_{U_S}:U_S\to B\setminus\{0\}$ is the specified axial torsor.
\end{proposition}



\begin{proof}
\begin{enumerate}[label=\arabic*.,leftmargin=*]
\item Write $E_1,\ldots,E_N$ for the exceptional curves,
and $D_Q,D_R$ for the strict transforms of the
two coordinate axes. The reduced total boundary
is normal crossing. Every $E_i$ is a rational
curve, and the exceptional divisor is a connected
tree if $N>0$. Its two orders
\begin{equation}
 a_i=\operatorname{ord}_{E_i}(Q),\qquad
 b_i=\operatorname{ord}_{E_i}(R)
\label{pub:I:eq:B-4-27-2}
\end{equation}
are positive: the image of $E_i$ is the origin.
Set
\[
 D_1=D_Q+\sum_i E_i,\qquad
 D_2=D_Q+D_R+\sum_i E_i.
\]
For a reduced normal-crossing divisor $D$, denote
by $\mathcal M_D^{\mathrm{gp}}$ the sheaf of
meromorphic units whose divisors are supported
on $D$.

\item \emph{The exact bounded-character subsheaf.}
The period pairing remains \eqref{pub:I:eq:B-4-24-1}, pulled
back to $S$. By the base-change lemma
\citep[Lemma 1.2.9, pp.~80--81]{KKN2008Root}
it defines $\mathcal A_S$. More explicitly its
bounded multiplicative-character sheaf is
\begin{equation}
 \mathcal T_{\log}^{(Y)}
   =\mathcal M_{D_1}^{\mathrm{gp}}
          \times\mathcal M_{D_2}^{\mathrm{gp}}.
\label{pub:I:eq:B-4-27-3}
\end{equation}
Indeed the first character is bounded by powers
of $q$, and therefore has zero order on $D_R$.
It may have arbitrary integer order on $D_Q$
and on each $E_i$, where $q$ has positive order.
The second character is bounded by powers of
$q r^{12}$, whose order is positive on every
boundary component, so all its boundary orders
are permitted. These individual conditions
also suffice for every character combination:
if its second coefficient is nonzero, a suitable
signed multiple of the second period bounds
both local normal orders; otherwise powers of
the first period do so on the permitted support.
At a crossing one chooses the bounding powers
simultaneously for its two integer orders.
This proves \eqref{pub:I:eq:B-4-27-3}, including at crossings
and at intersections with $D_R$. In particular the first character has zero order along $D_R$.

\item \emph{Cohomology of the blown-up bidisc.}
For a point blowup $p$ of
a smooth surface,
$p_*\mathcal O=\mathcal O$ and
$R^j p_*\mathcal O=0$ for $j>0$.
The first identity follows from Hartogs.
For the second, the two standard blowup charts
give a local Laurent calculation: on their
intersection a holomorphic function expands
as a convergent sum of $u^k z^\ell$ with
$k\ge0$. Terms with $\ell\ge0$ extend in
the first chart. Terms with $\ell<0$ extend
in the second, because in coordinates
$z'=1/z$, $u'=uz$ they are
$(u')^k(z')^{k-\ell}$. Cauchy estimates on
smaller coordinate polydiscs preserve convergence
of the two sums. The charts and their
intersection are Stein, so this kills the
first cohomology and there is no higher
\v Cech cohomology. Iterating and using the
Leray sequence over the Stein bidisc gives
\begin{equation}
 H^j(S,\mathcal O_S)=0\qquad(j>0).
\label{pub:I:eq:B-4-27-4}
\end{equation}

\item The exceptional tree has one $H^2(\ ,\Z)$
generator per rational component and no positive
cohomology in other degrees. Proper base change
and Leray for $\beta$ therefore give
\begin{equation}
 H^2(S,\Z)\xrightarrow{\ \sim\ }
       \bigoplus_i H^2(E_i,\Z)\simeq\Z^N,
 \qquad H^j(S,\Z)=0\quad(j=1,3,4).
\label{pub:I:eq:B-4-27-5}
\end{equation}
Equivalently these groups follow from the usual
tubular plumbing of the exceptional tree.
The exceptional intersection matrix is
unimodular: at each blowup the total-transform
basis adds a vector of square $-1$ orthogonal
to the preceding basis. Passing to the strict
transform basis is an integral unimodular
change of basis. Thus the Chern classes of
$\mathcal O(E_i)$ form a basis of
$H^2(S,\Z)$. The exponential sequence and
\eqref{pub:I:eq:B-4-27-4}--\eqref{pub:I:eq:B-4-27-5} give
\begin{equation}
 \operatorname{Pic}(S)\simeq\Z^N
       =\langle\mathcal O(E_i)\rangle,
 \qquad
 H^2(S,\mathcal O_S^*)=H^3(S,\mathcal O_S^*)=0.
\label{pub:I:eq:B-4-27-6}
\end{equation}

\item For either $D=D_1$ or $D_2$ the divisor-order
sequence is exact:
\begin{equation}
 0\longrightarrow\mathcal O_S^*
 \longrightarrow\mathcal M_D^{\mathrm{gp}}
 \longrightarrow
   \bigoplus_{C\subset D}i_{C*}\Z
 \longrightarrow0.
\label{pub:I:eq:B-4-27-7}
\end{equation}
The map from global integer orders to the Picard
group sends a component to its divisor line
bundle. It is surjective because both divisors
contain every $E_i$. Every component has
$H^1(C,\Z)=0$: the strict transforms of the
axes are discs and the exceptional curves are
rational. Their only nonzero $H^2$ contributions
are the $N$ exceptional curves. Consequently
the long exact sequence of \eqref{pub:I:eq:B-4-27-7} gives
\begin{equation}
 H^1(S,\mathcal M_{D_j}^{\mathrm{gp}})=0,
 \qquad
 H^2(S,\mathcal M_{D_j}^{\mathrm{gp}})
       \simeq\bigoplus_i H^2(E_i,\Z)
       \simeq\Z^N\quad(j=1,2).
\label{pub:I:eq:B-4-27-8}
\end{equation}
The second isomorphism is induced by the 
divisor-order map, a fact needed for the next
calculation.

\item \emph{Injectivity of the period map in degree two.}
The quotient definition of $\mathcal A_S$ gives
\begin{equation}
 0\longrightarrow Y\longrightarrow
 \mathcal T_{\log}^{(Y)}\longrightarrow
 \mathcal A_S\longrightarrow0,
 \qquad Y=\Z u\oplus\Z\gamma.
\label{pub:I:eq:B-4-27-9}
\end{equation}
By \eqref{pub:I:eq:B-4-27-3} and \eqref{pub:I:eq:B-4-27-8}, its middle
sheaf has zero $H^1$ and $H^2\simeq\Z^{2N}$.
Using restriction to $E_i$ as the coordinates
in \eqref{pub:I:eq:B-4-27-5}, the map
\[
 H^2(S,Y)\longrightarrow
 H^2(S,\mathcal T_{\log}^{(Y)})
\]
is block diagonal. Its block at $E_i$ is the
valuation matrix of the two period columns,
\begin{equation}
 \begin{pmatrix}
 12m a_i&0\\0&12m a_i+12n b_i
 \end{pmatrix}.
\label{pub:I:eq:B-4-27-10}
\end{equation}
Both diagonal entries are positive. Thus this
integral map is injective. The long exact sequence of
\eqref{pub:I:eq:B-4-27-9} now gives \eqref{pub:I:eq:B-4-27-1}.

\item Put $U_S=\beta^{-1}(B\setminus\{0\})$.
The torsor on $U_S\simeq B\setminus\{0\}$
is nontrivial by \ref{push:src:B.4.24}, including after the
stated finite coordinate roots. An extending
torsor on $S$ would be trivial by \eqref{pub:I:eq:B-4-27-1},
and its restriction would be trivial as well,
a contradiction. For $N>0$, $U_S$ is the
complement of the exceptional union; for
$N=0$ it is the punctured bidisc. Thus the
case $N=0$ is included as well.

\end{enumerate}
\end{proof}

\begin{proposition}[Maximal monodromy with trivial central dualizing bundle]\hypertarget{tag-030X}{}\label{push:src:B.4.28}
\StatementRecord{push:src:B.4.28}{proved}{030X}{\citet[Definition 1.1 and Lemmas 1.2, 2.2, 2.7]{Friedman2016AB}; \citet[Corollary 2.25, Lemma 2.26 and Proposition
2.27]{Dloussky2005AB}; \citet[Sections 3.2--3.3]{FujikiPontecorvo2010AB}; \citet[Section 15, Theorem 15.1 and the following
paragraph, pp.~130--131]{Oda1978D}; \citet[Chapter~III, Theorem~4.12(ii), pp.~134--135]{BS76}; \citep[Chapter~III, Theorem~4.12(iii), pp.~134--135]{BS76}}

Let $p:\mathcal X\to\Delta$ be a proper holomorphic map
from a smooth complex threefold, smooth over $\Delta^*$,
whose nonzero fibres are complex two-tori. Suppose that
$D=p^{-1}(0)$ is a reduced normal-crossing divisor and that $\omega_D\simeq\OO_D$.
If its monodromy on $H_1(\mathcal X_t,\Q)$ satisfies
\begin{equation}
 \operatorname{rank}(T-1)=2,
\label{pub:I:eq:B-4-28-1}
\end{equation}
then every normalization component of $D$ is a smooth
rational toric surface, its conductor is its toric
boundary, and the dual complex of $D$ is a topological
two-torus. In particular all normalized components are
K\"ahler.

The trivialization of $\omega_D$ in the hypotheses is the central
section used in the residue and nearby-period arguments below.
\end{proposition}


\begin{proof}
\begin{enumerate}[label=\arabic*.,leftmargin=*]
\item Write $V_i$ for the compact smooth normalization
components and $C_i$ for their reduced conductor divisors.
Adjunction gives
\begin{equation}
 K_{V_i}+C_i\sim0.
\label{pub:I:eq:B-4-28-2}
\end{equation}
If $D$ is smooth, the map is a proper submersion near
its compact central fibre and its monodromy is trivial.
Otherwise every $C_i$ is nonempty, by connectedness.
Adjunction on a conductor branch says that it is either
a smooth elliptic curve disjoint from the other branches,
or a rational curve with two boundary points, counted
on its normalization. The rational branches form cycles,
including the nodal irreducible and two-component cases.

\item \emph{Nonnegative Euler contributions.}
A nonempty effective anticanonical divisor makes every
pluricanonical system empty: multiplying a nonzero
pluricanonical section by the appropriate power of the
anticanonical section would give a nonzero holomorphic
function with a nonempty zero divisor. Thus the surface
has Kodaira dimension $-\infty$. Its minimal model is
rational, ruled, or of class VII.

\item On a rational surface, the conductor is connected, by
$0\to K_V\to\OO_V\to\OO_C\to0$ and $H^1(K_V)=0$.
For a rational cycle, the charge identity and positivity
are
\begin{equation}
 e(V\setminus C)=12-C^2-\#\{\text{cycle components}\}
 \ge0,\qquad
 e(V\setminus C)=0\ \Longleftrightarrow\ (V,C)
 \text{ is toric}.
\label{pub:I:eq:B-4-28-3}
\end{equation}
Here $e$ denotes compactly supported Euler characteristic.
This includes the nodal irreducible convention:
\citet[arXiv version~2, Section~1, p.~3; Definition~1.1 and Lemma~1.2, pp.~4--5; Lemma~2.2, p.~11; Lemma~2.7, p.~14]{Friedman2016AB}.
For a smooth elliptic anticanonical curve,
$e(V\setminus C)=2+b_2(V)>0$ instead.

\item On a nonrational ruled minimal model, a conductor
component dominating the base has degree at most two
and genus at most one. The base is therefore elliptic.
The minimal conductor is two disjoint elliptic sections
or one unramified elliptic bisection. Its Euler
contribution is zero. The two section squares sum to
zero, and the bisection has square zero, also directly
from $K_V+C\sim0$ and $K_V^2=0$.

\item A minimal class-VII surface with $b_2=0$ and a curve is
a Hopf surface. Its elliptic curves have square zero,
and a reduced anticanonical divisor has at most two
elliptic components. For this last assertion the exact
sequence $0\to K_V\to\OO_V\to\OO_C\to0$ gives
$h^0(\OO_C)\le1+h^1(K_V)=2$.
The Euler contribution is zero. This also covers
secondary Hopf surfaces.

\item For a minimal class-VII surface with $b_2>0$, the
reduced anticanonical divisor is a numerically
anticanonical divisor whose coefficients are all one.
\citet[arXiv version~3, Corollary~2.25, pp.~16--17; Lemma~2.26, p.~17; Proposition~2.27, pp.~17--19]{Dloussky2005AB} imply that it is parabolic Inoue
or Inoue--Hirzebruch. Half-Inoue has only a nontrivially
twisted anticanonical divisor and cannot satisfy
\eqref{pub:I:eq:B-4-28-2}. The remaining possibilities are hyperbolic
Inoue with two disjoint rational cycles, or parabolic
Inoue with one rational cycle and a disjoint elliptic
curve. If $n=b_2>0$, there are $n$ rational cycle
components and $e(V)=n$, so the contribution is zero.
In the parabolic case the elliptic curve has square
$-n$. The curve configurations, anticanonical identities
and underlying $(S^1\times S^3)\#n\overline{\PP}^2$
topology are recorded in
\citet[Sections~3.2--3.3, pp.~20--22]{FujikiPontecorvo2010AB}.

\item These assertions extend to nonminimal pairs. An
exceptional curve outside $C$ meets it once, transversely
at a smooth point, since $-K_V\cdot E=1$; its contraction
reverses an interior blow-up. An exceptional curve in
$C$ meets the rest at two points and its contraction
reverses a corner blow-up. The reduced anticanonical
condition is preserved. Interior blow-up increases
$e(V\setminus C)$ by one, whereas corner blow-up leaves
it unchanged. Thus every contribution is nonnegative.
Equality permits only toric rational pairs, minimal
elliptic ruled or Hopf pairs, and corner blow-ups of
hyperbolic or parabolic Inoue pairs. For these Inoue
pairs the cycle count and $b_2$ increase together, and
the fundamental group stays $\Z$. Corner blow-ups do
not touch the parabolic elliptic curve, whose square
remains strictly negative.

\item The local nearby fibre at a multiplicity-$k$ stratum
of a reduced normal-crossing divisor retracts to
$(S^1)^{k-1}$. Euler integration consequently gives
\begin{equation}
 0=e(\mathcal X_t)=\sum_i e(V_i\setminus C_i).
\label{pub:I:eq:B-4-28-4}
\end{equation}
Every summand is nonnegative and therefore zero.
Rational components with a smooth elliptic conductor,
and all interior blow-ups, are excluded.

\item \emph{Elliptic normal-degree balance.}
An elliptic double curve contains no triple point.
The local model $p=xy$ gives opposite degrees to its
two normal lines. Sum these degrees over all elliptic
conductor branches. The total is zero. An elliptic
ruled or Hopf component contributes zero; a parabolic
Inoue component contributes the strictly negative
degree of its sole elliptic branch. No rational
component with elliptic conductor remains. Hence there
are no parabolic Inoue components.

\item By connectedness, all conductors are now elliptic or
all are rational: no remaining component has both
types. In the elliptic case there are no triple points
and every vertex of the dual graph has valence at most
two. The graph is a segment or circle, allowing a loop
or multiple edges. In this case
\begin{equation}
 \operatorname{rank}(T-1)\le b_1(\text{dual graph})
 \le1.
\label{pub:I:eq:B-4-28-5}
\end{equation}
To prove this, remove
small tubes about the disjoint double curves. The
remaining compact cores extend as smooth submersions
over the whole disc and can be trivialized there.
Monodromy is then a product of the annular twists
$xy=t$. A passage of a loop across one neck adds its
oriented normal-circle class, with coefficient the
signed crossing number of that edge. Core loops have
zero change. Crossing numbers of a closed loop form a
cycle of the dual graph. Thus $T-1$ factors through its
rational cycle space. The normal-circle rotation exists
on the circle bundle of the normal line. This proves
\eqref{pub:I:eq:B-4-28-5} and excludes the elliptic case.

\item \emph{Hyperbolic components contradict a nonzero nearby period.}
The remaining possibilities are toric rational pairs
and corner blow-ups of hyperbolic Inoue pairs. Suppose
one component is of the latter kind. A boundary of a
small neighborhood of either rational conductor cycle,
lying in $V_i\setminus C_i$, is a torus mapping torus
$M_A$, whose monodromy on its fibre lattice is
\begin{equation}
 A\in\operatorname{SL}_2(\Z),\qquad
 |\operatorname{tr}A|>2.
\label{pub:I:eq:B-4-28-6}
\end{equation}
In
\citet[Section 15, Theorem 15.1 and the following
paragraph, retypeset pp.~130--131]{Oda1978D}, the universal
toric construction has lattice $N$ and deck action
given by multiplication by a nontrivial totally
positive unit of a real quadratic field. This is
the hyperbolic matrix $A$. The complement of the
two cycles is $\R$ times the compact angular-torus
bundle over a circle. A path in the logarithmic
absolute-value space from a point to its image by
$A$ realizes this mapping torus. The point can be
chosen in a node cone arbitrarily near a conductor
node. Its two adjacent primitive rays form a lattice
basis, so the angular fibre is exactly the Clifford
torus in the corresponding local coordinates
$(x,y)$. Corner blow-ups are toric star subdivisions;
they alter neither the dense torus nor $A$, only
this unimodular coordinate basis. Thus the same
marked mapping torus exists for every corner
blow-up allowed above.

\item Choose a compact such $M_A$ in the smooth open stratum.
A neighborhood of this compact set extends as a
smooth submersion over the full parameter disc.
Transport gives an embedding $M_A\to\mathcal X_t$ for
all small $t$, with the same marked fibre $L_t\simeq
(S^1)^2$. In the abelian group
$H_1(\mathcal X_t,\Q)$, the mapping-torus relations give
\[
 (A-1)H_1(L_t,\Q)\longmapsto0.
\]
Since $A-1$ is invertible over $\Q$, both angular
generators map to zero. The cohomology ring of a
complex two-torus is the exterior algebra on $H^1$.
It follows that the map $H^2(\mathcal X_t,\Q)\to
H^2(L_t,\Q)$ is zero and hence
\begin{equation}
 [L_t]=0\quad\hbox{in }H_2(\mathcal X_t,\Q).
\label{pub:I:eq:B-4-28-7}
\end{equation}

\item On the other hand, $h^0(D,\omega_D)=1$, since $D$ is
compact and connected, and the corresponding dimension
on every nearby two-torus is also one. Grauert's
cohomology and base-change theorem applied to the
invertible relative dualizing sheaf supplies a relative
section extending the chosen nonzero central generator;
see \citet[Chapter~III, Theorem~4.12(ii), pp.~134--135]{BS76}.
Evaluation is
nonzero at every point of $D$. Compactness and shrinking
make this section nonvanishing near the entire fibre.
Write it as $\Omega$. At the conductor node just used,
the total family has coordinates $xyz=t$ and
\begin{equation}
 \Omega_t=U(x,y,z)\frac{dx}{x}\wedge\frac{dy}{y},
 \qquad U(0,0,0)\ne0.
\label{pub:I:eq:B-4-28-8}
\end{equation}
The angular torus can be transported specifically as
$|x|=\epsilon_1$, $|y|=\epsilon_2$, $z=t/(xy)$.
This agrees, up to isotopy, with its compact-core
transport. If $U=\sum a_{ijk}x^iy^jz^k$ in a small
polydisc, direct integration of the convergent series
gives
\begin{equation}
 \int_{L_t}\Omega_t
   =(2\pi i)^2\sum_{k\ge0}a_{kkk}t^k
   \longrightarrow(2\pi i)^2U(0,0,0)\ne0.
\label{pub:I:eq:B-4-28-9}
\end{equation}
For small $t$ this contradicts \eqref{pub:I:eq:B-4-28-7}, because
the holomorphic two-form $\Omega_t$ is closed. This
excludes every hyperbolic Inoue component.

\item \emph{The dual complex.}
All normalization pairs are now rational toric pairs.
Every vertex link is one circle and every edge has
two incident faces, because each rational double
curve has two triple points. Thus the finite dual
complex $K$ is a connected closed topological surface.
It is oriented: the iterated residues of a
trivialization of $\omega_D$ are nonzero at every
triple point, and the residue theorem on each rational
double curve makes its two coefficients opposite,
with the gluing sign for the two branches. These
coefficients form a nowhere-zero two-cycle on $K$.
Constancy of coherent Euler characteristic in the
proper flat analytic family
\citep[Chapter~III, Theorem~4.12(iii), pp.~134--135]{BS76}, and the alternating
normalization resolution, give
\begin{equation}
 0=\chi(\OO_D)=\chi(K).
\label{pub:I:eq:B-4-28-10}
\end{equation}
The surface therefore has genus one.

\item For self-crossings, use the finite local-branch
incidence cell structure; the link of a normalized
toric component is its circular boundary. Over the
normalization $\PP^1$ of a double curve, the two-branch
cover extends unramified over its triple-point ends
by local normal-crossing coordinates and hence splits.
Subdivision
displays the same closed surface and the residue signs.
The alternating normalization resolution uses these
branch signs and gives the same Euler count. The
hyperbolic-link argument took place in the open
normalization stratum, which embeds in $D$, and is
unchanged by self-identification of boundary branches.
\end{enumerate}
\end{proof}


\begin{remark}[Vertical divisors and component adjunction]
\hypertarget{tag-PT05}{}\label{I:vertical-divisor-adjunction}
\StatementRecord{I:vertical-divisor-adjunction}{proved}{PT05}{}
For a proper semistable surface family with smooth total space and reduced simple normal-crossing central divisor $\sum_i V_i$, let a nonzero meromorphic relative two-form have vertical divisor $\sum_i a_iV_i$. Put $C_{ij}=V_i\cap V_j$ and $C_i=\sum_{j\ne i}C_{ij}$. Then
\begin{equation}
 K_{V_i}+C_i\sim\sum_{j\ne i}(a_j-a_i)C_{ij}.
\label{pub:I:eq:B-4-28-11}
\end{equation}
Indeed, adjunction gives $K_{V_i}=(K_{\mathcal X/\Delta}+V_i)|_{V_i}$, while the principal total fibre gives $\mathcal O_{\mathcal X}(V_i)|_{V_i}\sim-C_i$. Restricting the divisor of the form yields the displayed equality.
\end{remark}

\subsection{What invariant forms do not control}

\begin{proposition}[Full-rank periods do not bound an affine comparison]\hypertarget{tag-030Y}{}\label{push:src:B.4.29}
\label{push:B429}
\StatementRecord{push:src:B.4.29}{proved}{030Y}{\SeriesRef{push:src:B.4.30}, \SeriesRef{push:src:B.4.31}, \SeriesRef{push:src:B.4.33}; \citet[Sections 1.1.1--1.1.3, pp.~73--74]{KKN2008Root}}

There is a proper semistable family of two-dimensional
complex tori over a disc, with smooth total threefold
and rank-two unipotent monodromy on $H^1$, and a
fibre-preserving automorphism over the punctured disc
with the following properties. It acts identically
on all relative invariant one- and two-forms and
on the Gauss--Manin variation. Their coherent
extensions supplied by the proper family therefore
give the same lattices under this identification.
An absolute holomorphic canonical form extending
across the centre is also unchanged. Nevertheless,
the automorphism has no meromorphic extension to
the total family.

Thus these invariant-form and connection data do
not alone control an affine origin or a marked
total-space comparison, even in the presence of
a proper smooth semistable cap.
\end{proposition}



\begin{proof}
\begin{enumerate}[label=\arabic*.,leftmargin=*]
\item \emph{A proper full-rank example.}
Let $\mathcal E\to\Delta_t$ be the standard Tate
family with ordinary fibre
\[
 E_t=\C^*/t^{\Z}\qquad(t\ne0)
\]
and reduced nodal rational central fibre. Its
proper analytic model and node chart $xy=t$
are the explicit example in
\citet[Sections 1.1.1--1.1.3, pp.~73--74]{KKN2008Root}.
Form $\mathcal E\times_{\Delta}\mathcal E$.
Its only total-space singularity is the product
of the two nodes, with local equations
\[
 xy=t,\qquad uv=t,\qquad xy=uv.
\]
Blow up the ideal $(x,u)$ in this conifold and
take its strict-transform charts. They are
\[
 u=xA,\quad y=Av,\quad t=xAv
\]
with smooth coordinates $(x,A,v)$, and
\[
 x=uB,\quad v=By,\quad t=uBy
\]
with smooth coordinates $(u,B,y)$. This is the
standard proper small resolution, an isomorphism
away from the conifold point. Glue it to the
identity on the complement of that point. The
resulting $p:\mathcal Y\to\Delta$ is proper,
has smooth total threefold, and has reduced
normal-crossing centre: its displayed equations
are products of distinct coordinates. The
resolution does not alter the ordinary fibres
$E_t\times E_t$.

\item For each elliptic factor choose period cycles
$w_i,u_i$ with $u_i\mapsto u_i+w_i$ around
$t=0$. On the direct sum the logarithm of
monodromy has rank two and square zero.
The same rank holds on dual $H^1$. Thus this
example has the full-rank weight-one
unipotent degeneration stated above.

\item \emph{The nonextending translation.}
On the punctured disc use multiplicative
coordinates $(Z,W)$ for the two factors and put
\begin{equation}
 T_h(Z,W,t)=\bigl(e^{2\pi i h(t)}Z,W,t\bigr),
 \qquad h(t)=1/t.
\label{pub:I:eq:B-4-29-1}
\end{equation}
This is a holomorphic automorphism of the
punctured family. The zero section
$e(t)=([1],[1])$ extends to the smooth torus
locus of the central normalization. For
positive integers $j$ tending to infinity,
\[
 t_j=1/j,\qquad t'_j=1/(j+1/2)
\]
lie in the disc and tend to zero, whereas
\[
 T_h(e(t_j))=([1],[1]),\qquad
 T_h(e(t'_j))=([-1],[1]).
\]
In the Tate model these approach the distinct
central smooth points with multiplicative
coordinates $(1,1)$ and $(-1,1)$.
The small resolution is an isomorphism near
both points. Hence $T_h\circ e$ has no
continuous extension to the centre.

\item This also excludes a meromorphic total extension.
If such an extension existed, let
$\Gamma\to\mathcal Y$ be its proper
meromorphic graph and pull it back along
$e:\Delta\to\mathcal Y$. The resulting
analytic space is proper over the disc and
has a single point in each noncentral fibre.
Its unique component dominating the disc
is a proper analytic curve, generically of
degree one. After normalization it is finite
of degree one over the normal disc, hence
is the disc itself. Its map to the target
$\mathcal Y$ would extend $T_h\circ e$
holomorphically. This contradicts the two
distinct limits just found. Vertical
components in the pulled-back graph do
not affect this dominating curve.

\item \emph{Relative forms and their coherent extensions.}
On the ordinary locus the invariant relative
one-forms are
\[
 \omega_1=dZ/Z,\qquad \omega_2=dW/W
 \quad\pmod{dt}.
\]
They and $\omega_1\wedge\omega_2$ are
unchanged by every $T_h$. The holomorphic
forms of a compact torus are invariant, so
this accounts for all relative holomorphic
one- and two-forms on the ordinary fibres.
The relative log forms extend in the Tate
model. Ordinary coherent relative forms
also give finite-order extensions: for
example $t\omega_1,t\omega_2$, and
$t^2\omega_1\wedge\omega_2$ extend as
ordinary relative forms. In the first
resolution chart the first two are explicitly
\[
 t\omega_1=Av\,dx,\qquad
 t\omega_2=Av\,dx+xv\,dA;
\]
the other chart gives the analogous regular
expressions. At an isolated Tate node the
same identity is $t\,dx/x=y\,dx$.

\item More generally, take the torsion-free image
of each coherent direct image of relative
forms in its bundle on the punctured disc.
Under the two punctured identifications
$\mathrm{id}$ and $T_h$, these are identical
coherent lattices, since the induced map on
the entire ordinary relative-form bundle
is the identity. Coherence therefore gives
no bound on $h$ in this example.

\item \emph{Gauss--Manin and an absolute form.}
Every fibre translation is homotopic to the
identity and acts identically on integral
cohomology. Consequently $T_h$ induces the
identity on the flat Gauss--Manin bundle,
its connection, and its Hodge subbundle.
The logarithmic period growth and regular
singularity of that connection are unchanged.

\item The absolute form
\begin{equation}
 \Omega=dt\wedge\frac{dZ}{Z}\wedge\frac{dW}{W}.
\label{pub:I:eq:B-4-29-2}
\end{equation}
descends under both Tate period actions,
since their extra differentials are multiples
of $dt$. In the first small-resolution chart
it is
\[
 d(xAv)\wedge\frac{dx}{x}
          \wedge\frac{d(xA)}{xA}
       =dv\wedge dx\wedge dA,
\]
and in the other chart it is
$dy\wedge dB\wedge du$.
Thus $\Omega$ extends as a nonvanishing
holomorphic canonical form on the smooth
total space. Translation adds
$2\pi i h'(t)dt$ to $dZ/Z$, which vanishes
in \eqref{pub:I:eq:B-4-29-2}, so $T_h^*\Omega=\Omega$.

\item \emph{The horizontal term of an absolute differential.}
On a local multiplicative chart, an absolute
lift of the first relative form transforms by
\begin{equation}
 T_h^*(dZ/Z)=dZ/Z+2\pi i h'(t)dt.
\label{pub:I:eq:B-4-29-3}
\end{equation}
This horizontal term is erased both on passing
to relative forms and on wedging with the base
differential. It is not controlled by the
linear Gauss--Manin connection. The stronger
choice $h(t)=e^{1/t}$ makes that term
essential, while leaving all the preceding
relative and canonical calculations unchanged.
For this choice the two sequences
$t_j=1/\log j$ and
$t'_j=1/\log(j+1/2)$ give the same two
distinct central limits.

\item Both markings belong to the same proper family. Their relative invariant forms and their absolute canonical form agree by the displayed pullbacks, while the two limits of the translated zero section are distinct. A common translation cancels in a difference of two sections.
\end{enumerate}
\end{proof}

\subsection{Noncompletion of the prescribed partial filling}

\begin{proposition}[Noncompletion of the fixed partial filling]\hypertarget{tag-0310}{}\label{push:src:B.4.30}
\label{push:B430}
\StatementRecord{push:src:B.4.30}{proved}{0310}{\SeriesRef{push:src:B.4.19}, \SeriesRef{push:src:B.4.4}, \SeriesRef{push:src:B.4.31}, \SeriesRef{push:src:B.4.33}, \SeriesRef{push:src:B.4.27}, \SeriesRef{push:src:B.4.29}, \SeriesRef{push:src:B.4.24}, \SeriesRef{push:src:B.4.26}, \SeriesRef{push:src:B.4.32}; \citet[Theorems 6.2 and 6.3,
pp.~32 and 34]{IKN2005E}; \citet{IKN2007ErratumE}; \citet[Definition 7.1(i) and Section 9.2,
pp.~164--165 and 173]{KKN2008PicardE}; \citet[Corollary 1.2 and its proof,
manuscript pp.~3--6]{Parusinski1990E}; \citet[Theorems~1.6 and 1.10, pp.~215--216, including the noncompact extension after Theorem~1.6]{BM97}; \citet[Lemma 2.14, pp.~6--7]{EH2024E}}

The partial filling of \ref{push:src:B.4.19} has no proper
Hausdorff complex-analytic completion over its
missing base point. This remains true after finite
coordinate roots of the two base coordinates and
allows arbitrary additional three-dimensional
components of the origin fibre.

\end{proposition}



\begin{proof}
\begin{enumerate}[label=\arabic*.,leftmargin=*]
\item Every finite coordinate-root cover is dominated by one with $q=Q^{12m}$ and $r=R^n$, where $m,n>0$; proper pullback reduces to this case. Suppose such a completion exists. Retain and normalize its reduced dominating component, resolve it, and principalize the inverse image of the two axes, as in the proof of Theorem~\ref{night:P2.7}. These proper modifications preserve the given family on the complement of the axes and the smooth open parts of the two prescribed axial models. The resulting proper map to the bidisc has connected fibres: its finite Stein factor is generically one sheet and the bidisc is normal.

\item We identify the axial markings with the standard specialization of Theorem~\ref{night:P2.3}, namely $(l_0,l_1,l_2)=(0,1,-1)$ and $a=b=1$. In the notation of Proposition~\ref{push:src:B.4.24}, one has $U=-12\kappa_0$. The two origins in Theorem~\ref{night:P2.3} equal the two origins there plus $6ve$, where $e=\Pi_c\widehat\delta=(0,1)^{\mathsf t}$. Translation by $6ve$ acts on the character $W$ as $W\mapsto r^{-6}W$. It shifts the twelve-chain index by six and extends, with its inverse, across the parameter axis; along the cusp axis $r\ne0$ it is multiplication by a holomorphic unit. Thus this is an explicit comparison of the prescribed axial models. It changes meridian labels only by multiples of $\widehat\delta$ and leaves their $\widehat\gamma$ coefficients unchanged.

\item In the commuting lifts of Theorem~\ref{night:P2.7}, an axis component has orders $(\alpha_i,\beta_i)$ and $\widehat\gamma$ meridian coefficient $\kappa_i$. The same group calculation gives
\[
 (m\alpha_i+n\beta_i)\kappa_j
 =(m\alpha_j+n\beta_j)\kappa_i
\]
for intersecting components. Every denominator $m\alpha_i+n\beta_i$ is positive, including for exceptional components over the origin. Hence the ratio $\kappa_i/(m\alpha_i+n\beta_i)$ is constant on the incidence graph of components meeting the corner.

\item The prescribed parameter-axis model has an extending ordinary origin. Its positive parameter loop has zero lattice label in that origin, so the strict parameter component has ratio zero. The prescribed cusp-axis origin has ratio $a=1$, by the exact covered origin formulas \eqref{eq:P2FQ}--\eqref{eq:P2FR} and the commuting lift $Q_o=T_{-m\widehat\gamma}Q$. Both strict components meet the corner by properness. The compact connected corner fibre makes the incidence graph connected, exactly as in the proof of Theorem~\ref{night:P2.7}. Its ratios would therefore be both zero and one. This contradiction proves noncompletion for every pair of positive root degrees and for all additional corner components.
\end{enumerate}
\end{proof}

\subsection{Global elliptic monodromy constrains every central meridian}

\begin{proposition}[The global meridians have only the prescribed toric direction]\hypertarget{tag-0311}{}\label{push:src:B.4.31}
\StatementRecord{push:src:B.4.31}{proved}{0311}{\SeriesRef{push:src:B.4.4}, \SeriesRef{push:src:B.4.19}, \SeriesRef{push:src:B.4.9}; \citet[Theorems~1.6 and 1.10, pp.~215--216, including the noncompact extension after Theorem~1.6]{BM97}}

Suppose $g:\mathcal Z\to\Delta_r$ is any proper complex-analytic
filling of a finite parameter cover of the six-sphere
family. Retain its reduced main component, normalize and resolve
it, extend its intrinsically marked map to $\PP^1$ by \ref{push:src:B.4.4},
resolve the graph, and resolve the support of its central divisor. Denote
the resulting smooth proper modification by $\widetilde{\mathcal Z}$,
and its holomorphic marked map by $h$. All modifications can be
chosen to be isomorphisms over $r\ne0$. The central divisor has normal-crossings support and positive integral
component multiplicities.

There is a finite subset $\Sigma\subset\PP^1$, containing
$0,1,\infty$, with the following property. For every
$q\in B^\circ:=\PP^1\setminus\Sigma$, every local angular
vanishing circle of the $r$-central normal-crossings divisor,
viewed in the homology of the nearby ordinary two-torus, lies in
\begin{equation}
             \Z\widehat\delta\subset H_1(X_{q,r};\Z).
\label{pub:I:eq:B-4-31-1}
\end{equation}
The same conclusion holds for the image of the entire local
angular torus at a higher crossing over $B^\circ$.
Here $\widehat\delta$ is the primitive cycle in the fixed
period marking. The conclusion is unaffected by any further
finite cover of the degeneration parameter.


\end{proposition}


\begin{proof}
\begin{enumerate}[label=\arabic*.,leftmargin=*]
\item \emph{The global model and the finite exceptional set.}
The punctured total family is connected and smooth, hence
irreducible. Its closure is the unique reduced main component
dominating the parameter disc. Retain this component, normalize
it, and resolve it properly, without changing its smooth
punctured part. The resulting total space is a smooth connected
fourfold, proper over the disc. The parameter is a nonzero
divisor, so its map to the disc is flat. The torsion-free coherent bi-anticanonical pencil and labelled
critical-value argument of \ref{push:src:B.4.4} apply to this smooth model.
Its marked graph is therefore a proper modification.
Further analytic resolution and principalization of the central
ideal give the stated smooth
model and normal-crossings support; the complex-analytic scope,
including locally finite resolution, is explicit in
\citet[Theorems~1.6 and 1.10, pp.~215--216, including the noncompact extension after Theorem~1.6]{BM97}. Properness over a smaller parameter disc
allows us to work over a compact central fibre throughout.
The modifications have their centres over $r=0$, so neither
the punctured family nor its marked map changes.
Flatness over $r$ persists because $r$ is a nonzero divisor in
the smooth local rings; all central multiplicities are kept in the
local monomial equations.

\item For possible self-crossings use the finite closed spaces of
ordered local branch intersections, with each local component
normalized. Locally these are the coordinate intersections of
the central branches. In particular, the double-branch spaces
are compact smooth surfaces and the triple-branch spaces are
compact smooth curves, even when two branches belong to one
global component. There are finitely many such components.

\item Add to $\Sigma$ the images of all branch-intersection components
on which $h$ is constant and the critical values of $h$ on every
remaining branch-intersection component, including the
normalizations of the single branches. These sets are finite.
Indeed the maps are proper; on a component mapping nontrivially
to the curve $\PP^1$ their critical image is a proper analytic
subset of that curve. Zero-dimensional deeper strata contribute
only finitely many further values.

\item Thus, over $B^\circ$, the closed double-branch spaces form
proper smooth families of compact curves. The connected
components of their fibres form finite covering spaces of
$B^\circ$: this follows either from proper submersion
trivializations, or from their finite Stein factorizations.
Removing the finitely many triple points from one connected
smooth curve leaves it connected. Branch orderings, central
multiplicities, and these connected fibre components therefore
constitute finite discrete data over $B^\circ$.

\item \emph{The local angular circles, including multiplicities.}
At a point of an open double-branch stratum, the submersion
condition for $h$ gives coordinates in which $q=h$ is a
tangential coordinate and
\begin{equation}
                 r=x^a y^b,\qquad a,b>0.
\label{pub:I:eq:B-4-31-2}
\end{equation}
A holomorphic unit has been absorbed in one coordinate on this
small chart. Put $d=\gcd(a,b)$. Each connected component of the
local nonzero fibre has the primitive angular circle
\begin{equation}
 \theta\longmapsto
 \bigl(x_0e^{i(b/d)\theta},\,
       y_0e^{-i(a/d)\theta}\bigr),
       \qquad 0\le\theta\le2\pi,
\label{pub:I:eq:B-4-31-3}
\end{equation}
where $x_0^a y_0^b=r$. There are at most $d$ choices of the
connected monomial-fibre component. Interchanging the ordered
branches can reverse the circle's orientation. These are only
finite choices.

\item Move the chosen base point along a path in the same open
double-curve fibre. Local product charts and the normal line
bundles transport the annular circle continuously. Changes of
normal frame rotate its parameter and do not add a tangential
one-cycle. Equivalently, this is the fibre circle of the normal
circle bundle along that path. If a based loop is used, changing
the path changes it by conjugation, which does not affect its
class in $H_1$. Consequently, for fixed small nonzero $r$ and
fixed $q$, each of the finitely many branch-component data
determines only finitely many homology classes, including signs.

\item The same local formula proves monodromy invariance in the
$r$ direction. Continuing around $r=0$ permutes the finitely
many connected components of $x^a y^b=r$; after a positive
power it carries \eqref{pub:I:eq:B-4-31-3} to the same oriented circle class.
In the marking, after a finite parameter cover the
$r$-monodromy is
\begin{equation}
 T_r=I-mE,\qquad
 E=\widehat\delta\otimes\gamma,\qquad m\in\Z_{>0},
 \quad E^2=0.
\label{pub:I:eq:B-4-31-4}
\end{equation}
For example, $w=r^6$ gives $m=12$ with the positive-loop
convention used in \ref{push:src:B.4.9}. If $T_r^k v=v$ for one of the
angular classes $v$, then $kmEv=0$. The homology lattice is
torsion-free, so
\begin{equation}
 v\in\ker E
    =\Z\widehat u+\Z\widehat w+\Z\widehat\delta.
\label{pub:I:eq:B-4-31-5}
\end{equation}
The argument uses the monomial equations with their displayed
positive multiplicities.

\item \emph{Why global continuation gives a finite orbit.}
Choose a base point $q_*\in B^\circ$. Let $\mathscr S$ be the
finite set of angular classes furnished at $q_*$ by all the
connected double-curve fibres, their ordered branches, their
monomial-fibre components, and both orientations.
Take generators of the finitely generated group
$\pi_1(B^\circ,q_*)$. A generator can be lifted to a path in
the appropriate family of double curves; it ends in one of
the finitely many connected fibre components. Avoid the
triple points during this lift. Properness over the compact
generator paths and finitely many product charts give one
sufficiently small $r_*$ for which all their annular circles
can be transported in the smooth fibres
$X_{q,r_*}$.

\item The transported circles are a continuous family of cycles,
so they realize the Gauss--Manin transport of $H_1$.
At the end, moving the base point within the resulting
connected punctured double curve does not change its $H_1$
class, by the previous paragraph. Hence each generator
permutes $\mathscr S$. The same holds for all words in the
generators. This proves a finite orbit in the marked lattice.
Its projection to
\begin{equation}
 \ker E/\Z\widehat\delta
       \simeq \Z\widehat u+\Z\widehat w
\label{pub:I:eq:B-4-31-6}
\end{equation}
therefore has finite orbit as well.

\item \emph{The exact elliptic-monodromy calculation.}
In the fixed marking the two finite peripheral
matrices are
\begin{equation}
 A_1=\begin{pmatrix}
 1&0&0&0\\6&0&1&0\\-6&-1&-1&0\\-2&1&0&1
 \end{pmatrix},\qquad
 A_2=\begin{pmatrix}
 1&0&0&0\\0&0&-1&0\\-6&1&0&0\\3&0&1&1
 \end{pmatrix}.
\label{pub:I:eq:B-4-31-7}
\end{equation}
They fix $\widehat\delta$ and $\gamma$, commute with $E$,
and induce on \eqref{pub:I:eq:B-4-31-6}
\[
 M_1=\begin{pmatrix}0&1\\-1&-1\end{pmatrix},\qquad
 M_2=\begin{pmatrix}0&-1\\1&0\end{pmatrix}.
\]
Deleting additional ordinary values from the elliptic base
does not decrease this monodromy image: the map on fundamental
groups to $\PP^1\setminus\{0,1,\infty\}$ is surjective.
The image therefore contains
\begin{equation}
 U=M_1M_2=\begin{pmatrix}1&0\\-1&1\end{pmatrix},\qquad
 V=M_2UM_2^{-1}=\begin{pmatrix}1&1\\0&1\end{pmatrix}.
\label{pub:I:eq:B-4-31-8}
\end{equation}
If a vector $\bar v$ has finite orbit, some positive powers
$U^k$ and $V^\ell$ both fix it. Their fixed lines are,
respectively, $\Q(0,1)$ and $\Q(1,0)$, whose intersection
is zero. Thus $\bar v=0$. Applying this to the projected
classes in $\mathscr S$ proves that each angular class lies
in $\Q\widehat\delta$. Since $\widehat\delta$ is a member
of an integral basis of $H_1(T^4;\Z)$, it is primitive;
an integral class in this rational line belongs to
$\Z\widehat\delta$. This proves \eqref{pub:I:eq:B-4-31-1}.

\item At a crossing with local equation
$r=\prod_{j=1}^k x_j^{a_j}$, the angular torus has
rational first homology
\[
 \ker\bigl((a_1,\ldots,a_k):\Q^k\longrightarrow\Q\bigr).
\]
It is spanned by the pairwise angular vectors
$a_j e_i-a_i e_j$. Each is a rational multiple of the
primitive circle of the corresponding double branch,
which can be moved to an open double-curve point by
keeping the other coordinates small and nonzero.
The same calculation therefore places the entire angular
image in $\Q\widehat\delta$, and primitivity again gives
the integral conclusion. This also handles self-crossings,
because the ordered local branches were retained throughout.
\end{enumerate}
\end{proof}


\subsection{The radial-loop class of every global cap}

\begin{proposition}[Vanishing of the normalized radial-loop defect]\hypertarget{tag-0312}{}\label{push:src:B.4.32}
\label{push:B432}
\StatementRecord{push:src:B.4.32}{proved}{0312}{\SeriesRef{push:src:B.4.31}, \SeriesRef{push:src:B.4.19}, \SeriesRef{push:src:B.4.4}, \SeriesRef{push:src:B.1.2}, \SeriesRef{push:src:B.4.30}; \citet[Theorem 6.3]{IKN2005E}}

Suppose a proper complex-analytic filling of a finite parameter
cover of the six-sphere family exists. Make the
modifications in \ref{push:src:B.4.31} and delete its finite exceptional set
from the marked base. Use the ordinary local origins of
\ref{push:src:B.4.19} as reference. There is a well-defined normalized
radial-loop defect of a local extending cap section,
\begin{equation}
 \lambda\in
 V:=\bigl(\Lambda\otimes\Q\bigr)/\Q\widehat\delta .
\label{pub:I:eq:B-4-32-1}
\end{equation}
It is independent of the section, of the finite parameter root
needed to obtain that section, and of the modifications used.
It patches to a flat section of the local system $V$
over the remaining ordinary base, and this section is zero.

In particular, let $t$ be the trace of an extending finite
multisection of degree $d$, and let $o$ be the reference ordinary
origin. The radial-loop defect of $t-do$ in $V$ is zero.
The same holds for an origin obtained from multiplied such
traces by bounded logarithmic differences along $r=0$.

\end{proposition}



\begin{proof}
\begin{enumerate}[label=\arabic*.,leftmargin=*]
\item Use the smooth marked model of Theorem~\ref{night:P2.4} and remove the finitely many base values which are images of vertical central components or critical values of the maps on the central branch strata. Every central component above the remaining ordinary base dominates the original $\mathbb P^1$. Its positive normal meridian has zero lattice class in $\Lambda/\mathbb Z\widehat\delta$ by Theorem~\ref{night:P2.6}.

\item The ordinary origin is exactly the one used there. Indeed, write $k=\log K/(2\pi i)$ with $K=t^4/(t-1)^3$ and $v=\log r/(2\pi i)$. The parameter descent of Theorem~\ref{night:P2.1} is $H=-k$, so subtracting the local origin $vHe$, with $e=\Pi_c\widehat\delta$ gives the ordinary gauge $\widetilde\zeta=\zeta+kv e$ of Proposition~\ref{push:src:B.4.19}. On an overlap, $k$ changes by an integer, and the resulting radial label is an integral multiple of $\widehat\delta$. The finite-root cancellation and the cusp comparison in that proposition have the same property. Thus the splitting used in Theorem~\ref{night:P2.6} computes precisely the radial-loop defect relative to the stated ordinary origins.

\item The connecting map is the endpoint period of a continued additive lift. With $T_r=I-mE$, $E=\widehat\delta\otimes\gamma$ and $E^2=0$, changing the lift adds $(T_r-I)\Lambda\subset\mathbb Z\widehat\delta$, so its image $\nu_r$ in $V$ is well defined. Under $r=R^e$ its label changes by
\[
 \sum_{j=0}^{e-1}T_r^jv
 =ev-\frac{me(e-1)}2Ev
 \equiv ev\pmod{\mathbb Q\widehat\delta}.
\]
Consequently the normalized defect of a section $b$ on that root cover is $e^{-1}\nu_R(b-o(R^e))$.

\item Lift an extending section to a common proper resolution. For a fixed ordinary base point the lift extends over the parameter origin: normalize the dominating curve in its proper graph, which has degree one over the disc. At its limit choose a normal-crossing chart in which
\[
 r=\text{unit}\prod_{i=1}^k x_i^{a_i},\qquad a_i>0.
\]
On the section, $x_i(R)=R^{b_i}u_i(R)$ with $u_i$ a holomorphic unit and $b_i\ge0$. Thus $e=\sum_i a_i b_i$, and its positive $R$-loop is the sum of $b_i$ positive component meridians. The unit loops contract in the chart. In the quotient presentation of Theorem~\ref{night:P2.6}, the lattice label is the same sum of the component labels, each zero modulo $\widehat\delta$. Hence the normalized radial defect is zero. This argument applies to every extending section and every root, and is preserved under common proper refinements. In a varying ordinary base chart the labels are locally constant, so the vanishing also holds across the isolated points where a simultaneous section lift has indeterminacy.

\item The connecting homomorphism is additive. Split a finite multisection after a cyclic root at a generic ordinary point, add the zero normalized defects of its branches, and use the preceding root formula. Its trace of degree $d$ therefore satisfies $\nu_r(t-do)=0$ in $V$. These zero local classes patch under the lattice transport, independently of the section and all choices. Finally, by the bounded-character presentation of Proposition~\ref{push:src:B.4.27}, a bounded logarithmic difference along the strict parameter axis has characters in $\mathcal O^*\times\mathcal M_r^{\mathrm{gp}}$: its first character has zero winding and its second contributes only a multiple of $\widehat\delta$. Adding it to a multiplied trace preserves the zero class. This proves every stated assertion.
\end{enumerate}
\end{proof}


\begin{remark}[A zero loop class with a nonextending translation]
\hypertarget{tag-PT07}{}\label{I:zero-loop-essential-translation}
\StatementRecord{I:zero-loop-essential-translation}{proved}{PT07}{\SeriesRef{push:src:B.4.29}}
In Proposition~\ref{push:src:B.4.29}, the additive lift of the translation is $h(t)=1/t$. It is single-valued on the punctured disc, so its radial deck difference and loop class are zero. Its translated zero section has the two distinct central limits computed there, and hence the translation has no meromorphic extension. The loop class and the affine comparison are related by the endpoint-period map used in Proposition~\ref{push:src:B.4.32}; this example has zero image under that map.
\end{remark}

\section{Other logarithmic transformations and finite covers}\label{pub:I:other-multiplicities}

\subsection{Comparison under logarithmic transformation}



\begin{proposition}[The relative cohomology comparison]\hypertarget{tag-01S3}{}\label{inv:B1.1}
\StatementRecord{inv:B1.1}{proved}{01S3}{the definitions and constructions stated in the proof}

Fix the period family and cusp filling. Let $X'$ be obtained by any
admissible translation twists and multiplicities $m_1=3k$, $m_2=4l$,
where $k,l$ are positive integers. Write $D'=W+S'_1+S'_2$ and
\[
\mathcal E'=\Omega^1_{X'}(\log D')/f'^*\Omega^1_B(\log\Sigma).
\]
The identity on the common smooth family induces isomorphisms on $B$
\[
R^qf_*\OO_X\simeq R^qf'_*\OO_{X'},\qquad
R^qf_*\mathcal E\simeq R^qf'_*\mathcal E'\quad(0\le q\le2).
\]
These isomorphisms identify the logarithmic connecting maps $\kappa_q$.
They also hold when the cusp collar is changed by an admissible translation.
\end{proposition}


\begin{proof}
\begin{enumerate}[label=\arabic*.,leftmargin=*]
\item Over $B\setminus\Sigma$
the two families are identified, and a translation acts trivially on
$H^q(F,\OO_F)$ and $H^q(F,\Omega^1_F)$ for a complex torus $F$.
For the second assertion, translation fixes the invariant holomorphic
one-forms, and the first follows either from invariant Dolbeault forms
or the exterior algebra $H^\bullet(F,\OO_F)=\bigwedge^\bullet H^1(F,\OO_F)$.
These are statements about the induced maps on cohomology, including
their holomorphic variation.

\item At a point with original order $d=3$ or $4$, let
$\widetilde f:T\to\Delta_z$ be the smooth proper torus model with
coarse coordinate $q=z^d$. The new model is its pullback by $z=s^e$,
where $m=de$, followed by the free cyclic quotient of order $m$.
Smooth proper base change identifies the higher direct images upstairs
with the pullbacks of those for $T$. The relative logarithmic cotangent
bundle upstairs is precisely $\Omega^1_{T/\Delta_z}$ pulled back:
the base logarithmic direction is killed in the quotient defining $\mathcal E'$.
The derivative of a translation contributes only a multiple of $ds$;
it consequently acts on this quotient as ordinary translation on relative forms.

\item Average over the cyclic action and choose holomorphic eigenframes.
If a frame for $H^q(\OO)$ or $H^q(\Omega^1)$ has character
$\zeta_d^{-a}$, $0\le a<d$, its original invariant extension is
$z^a\OO_{\Delta_q}$. On the order-$m$ cover its character is
$\zeta_m^{-ea}$, and its invariant extension is
$s^{ea}\OO_{\Delta_q}=z^a\OO_{\Delta_q}$.
Taking invariants is exact in characteristic zero; hence the map $s^{ea}v\mapsto z^av$, with $z=s^e$,
identifies the extensions inside the common punctured-disc bundle. A translation
twist induces the identity on the cohomology just used and changes none
of these lattices. This proves the comparison at both finite points,
including the full nonreduced fibres.

\item At the cusp the relative logarithmic forms have the global toric frame
$d\log x_1,d\log x_2$, modulo $d\log q$.
A toric gluing has the form
$x_i\mapsto c_i(q)q^{b_i}x_i$ and therefore changes $d\log x_i$ by
$(q c_i'/c_i+b_i)d\log q$. It induces the identity on the relative
logarithmic quotient. The structure-sheaf comparison is also the identity
on the normalization cochain complex. This verifies the cusp extension
and the assertion about translated collars.

\item The logarithmic cotangent exact sequences are isomorphic on the smooth
base, so their connecting maps agree there under these comparisons.
Their source and target sheaves are locally free by the local calculations
in \cref{inv:A1.0,inv:A1.1} and the eigenframe calculation above. A holomorphic map between
vector bundles on a reduced curve which is zero on a dense open set is
zero everywhere. Thus the connecting maps agree globally.
\end{enumerate}
\end{proof}

\begin{proposition}[All admissible twists have the same Hodge numbers]\hypertarget{tag-01S4}{}\label{inv:B1.2}
\StatementRecord{inv:B1.2}{proved}{01S4}{\SeriesRef{inv:B1.1}}

For every family in \ref{inv:B1.1} whose cusp is the one-component hexagonal filling,
the direct images of the absolute cotangent sheaves are
\[
\begin{array}{c|ccc}
 &R^0f'_*&R^1f'_*&R^2f'_*\\\hline
\Omega^1_{X'}&\OO(-2)&\OO\oplus\OO(-1)^2&\OO\oplus\OO(-1)\\
\Omega^2_{X'}&\OO(-2)\oplus\OO(-1)&\OO(-2)\oplus\OO(-1)^2&\OO
\end{array}
\]
and $R^qf'_*\OO_{X'}=(\OO,\OO\oplus\OO(-1),\OO(-1))$.
In particular every admissible $(3,4)$ member with
$|12l_0-4l_1-3l_2|=1$ has the Hodge diamond of \cref{inv:A3.1}.
\end{proposition}


\begin{proof}
\begin{enumerate}[label=\arabic*.,leftmargin=*]
\item \ref{inv:B1.1} identifies the logarithmic direct images and their connecting maps.
By \cref{inv:A2.1,inv:A2.4}, the cokernel of $\kappa_0$ is
$\C_{p_1}\oplus\C_{p_2}$, and $\kappa_1$ is surjective with kernel
$\OO\oplus\OO(-1)^2$. These statements therefore hold for $X'$.
The residue sheaf of its normalized divisor has higher direct images
\[
R^0f'_*\OO_{\widetilde D'}=\C_{p_0}\oplus\C_{p_1}\oplus\C_{p_2},
\quad R^1f'_*\OO_{\widetilde D'}=\C_{p_1}\oplus\C_{p_2},
\quad R^2=0.
\]
Indeed the finite reductions are bielliptic, and the cusp normalization
is rational. The base logarithmic residue is evaluation with weights
$(1,m_1,m_2)$; its kernel is $\OO(1)(-\Sigma)=\OO(-2)$.
On $H^1$ the residue of the invariant class $\gamma^{0,1}$ is
$m_j\gamma^{0,1}|_{S'_j}\ne0$. The other eigenline has zero residue.
The local rank-one image of $\kappa_0$ is its kernel, as follows either
from the residue diagram of \cref{inv:A1.3} or from the eigenframe calculation
in \ref{inv:B1.1}. Thus the two finite skyscrapers cancel exactly. This proves
the first row, with no surviving torsion.

\item The proper flat map is Gorenstein of relative dimension two. Relative
duality for $F=\Omega^1_{X'}$, using
$F^*\otimes\omega_{X'/B}=\Omega^2_{X'}\otimes f'^*\OO(2)$,
gives the second row because all sheaves in the first row are locally
free. Coherent Leray on $\Pone$ gives every Hodge number. In particular
$h^{p,0}=0$ for $p>0$ follows from these direct images and duality,
as the degree-zero entries and their duals show.
\end{enumerate}
\end{proof}

\subsection{Higher multiplicities and the anticanonical ring}

\begin{proposition}[Canonical bundles for all coprime multiplicities]\hypertarget{tag-01S5}{}\label{inv:B2.1}
\StatementRecord{inv:B2.1}{proved}{01S5}{\SeriesRef{inv:B1.1}, \SeriesRef{inv:B1.2}; \citet[Theorem 1.1 and Proposition 2.3]{ViaclovskyLog2026}}

Let $X=X_{3k,4l}(u)$ with admissible translation data as in
\citet[Theorem~1.1, p.~1, and Proposition~2.3, pp.~5--6]{ViaclovskyLog2026},
$\gcd(3k,4l)=1$, and any nonzero $u$ in its construction domain.
Then
\[
K_X=f^*\OO_B(-1)\otimes
\OO_X\bigl((k-1)S_1+(3l-1)S_2\bigr).
\]
Moreover $(h^{0,q})_{q=0}^3=(1,1,0,0)$, $h^{p,0}=0$ for $p>0$,
and $(h^{1,1},h^{1,2})=(2,1)$ for every such pair.
The relative comparison maps in Proposition~\ref{inv:B1.1}
identify the cohomological calculation for all these multiplicities.
\end{proposition}


\begin{proof}
\begin{enumerate}[label=\arabic*.,leftmargin=*]
\item The cited proposition supplies exactly the free cyclic quotients used
in \ref{inv:B1.1}; its Theorem 1.1 supplies the stated sphere topology for its
admissible choices. Proposition~\ref{inv:B1.2} computes the cohomology by the relative
cotangent comparison.
For the canonical bundle we check the natural evaluation
$f^*f_*\omega_{X/B}\to\omega_{X/B}$.
The invariant relative determinant has weights $a_1=2k$ and $a_2=l$
under the clockwise generators of orders $3k$ and $4l$. The pullback
of a canonical form is $F(s)\,ds\wedge e$, where the character of
$ds\wedge e$ is $\zeta_m^{1+a}$. Therefore its invariant coefficients
are $s^{m-1-a}$ times an invariant function and a nowhere vanishing
unit. Since the quotient of total spaces is \'etale, the evaluation
vanishes to order $m-1-a$ on the reduced fibre downstairs. These orders
are $k-1$ and $3l-1$. At the cusp the form
$d\log x_1\wedge d\log x_2\wedge dq$ is a nowhere vanishing canonical
frame in every smooth toric chart, so no cusp component occurs.

\item By \ref{inv:B1.1} and relative duality,
$f_*\omega_{X/B}=(R^2f_*\OO_X)^*=\OO(1)$.
The evaluation divisor just computed therefore gives
$\omega_{X/B}=f^*\OO(1)\otimes
\OO((k-1)S_1+(3l-1)S_2)$. Tensoring by $f^*K_B=f^*\OO(-2)$
proves the formula.
\end{enumerate}
\end{proof}

\begin{proposition}[The complete graded anticanonical ring]\hypertarget{tag-01S6}{}\label{inv:B2.2}
\StatementRecord{inv:B2.2}{proved}{01S6}{\SeriesRef{inv:B2.1}}

Put $m_1=3k$, $m_2=4l$, and $d=4l+3k-kl$.
Choose integers $r,s$ with $m_2r+m_1s=1$ and set
$\mathcal L=\OO_X(rS_1+sS_2)$.
There is a graded algebra isomorphism
\[
\bigoplus_{n\ge0}H^0(X,-nK_X)
\simeq\bigoplus_{n\ge0}\C[x,y]_{dn},
\qquad\deg x=m_2,\quad\deg y=m_1,
\]
where the right-hand summand is assigned degree $n$.
For negative weights the polynomial space is zero. Explicitly
\[
h^0(X,-nK_X)=
\max\left\{1+n-\left\lceil\frac{n(k-1)}{3k}\right\rceil
-\left\lceil\frac{n(3l-1)}{4l}\right\rceil,\,0\right\}.
\]
The integer $d$ is nonzero for admissible pairs. If $d<0$ the
anticanonical ring is $\C$; if $d>0$ it is the indicated weighted
Veronese ring. For $(k,l)=(1,1)$ it is
$\C[U,V]$, $\deg U=1$, $\deg V=2$.
\end{proposition}


\begin{proof}
\begin{enumerate}[label=\arabic*.,leftmargin=*]
\item The divisor of $f^*(t/(t-1))$ is $m_1S_1-m_2S_2$.
It follows, retaining this rational function as the comparison, that
$\mathcal L^{m_2}\simeq\OO(S_1)$,
$\mathcal L^{m_1}\simeq\OO(S_2)$ and
$\mathcal L^{m_1m_2}\simeq f^*\OO(1)$.
\ref{inv:B2.1} then gives $-K_X=\mathcal L^d$.

\item For every integer
power of $\mathcal L$, a meromorphic function representing a section
of $\OO(aS_1+bS_2)$ is holomorphic on every regular compact fibre
and is constant there. It descends to the base on the smooth locus.
On a cyclic cover $q=s^m$, the finite pole bound gives a Laurent
expansion bounded below, and invariance leaves only powers divisible
by $m$. Thus descent extends meromorphically across the finite point,
with the exact valuation inequality $m\operatorname{ord}_q(h)+a\ge0$.
The identical argument at the reduced cusp has $m=1$.
Conversely each base function satisfying these bounds pulls back to
an allowed section. Consequently, as subsheaves of meromorphic functions,
\[
f_*\OO(aS_1+bS_2)=
\OO_B\left(\left\lfloor a/m_1\right\rfloor p_1+
\left\lfloor b/m_2\right\rfloor p_2\right).
\]
This proves the displayed dimension formula using \ref{inv:B2.1} and
$h^0(\Pone,\OO(e))=\max\{e+1,0\}$.

\item For the algebra, let $x,y$ be the canonical sections of
$\OO(S_1),\OO(S_2)$ under the chosen identifications with powers of
$\mathcal L$. Their powers $x^{m_1},y^{m_2}$ pull back two independent
sections of $\OO_B(1)$, with zeros at the two distinct points.
For any integer $j$, choose the unique residues
$0\le a_0<m_1$, $0\le b_0<m_2$ for which
$j-m_2a_0-m_1b_0=e m_1m_2$ is integral.
The valuation formula identifies $H^0(\mathcal L^j)$ with
$x^{a_0}y^{b_0}H^0(\OO_B(e))$. When $e\ge0$ a basis is
\[
x^{a_0+m_1i}y^{b_0+m_2(e-i)},\qquad0\le i\le e;
\]
when $e<0$ there are no sections. These are exactly the monomials of
weighted degree $j$. Their multiplication is ordinary polynomial
multiplication, proving the algebra assertion in every degree.
When $k=l=1$, the weight-$6n$ monomials are generated by $y^2$ and
$x^3$, of regraded degrees $1,2$.
Finally, $d=0$ would imply $k\mid4l$ and $l\mid3k$; coprimality
forces $k=l=1$, which instead gives $d=6$.
\end{enumerate}
\end{proof}

\begin{remark}[Six multiplicity pairs]
\hypertarget{tag-PI18}{}\label{I:small-multiplicity-pairs}
\StatementRecord{I:small-multiplicity-pairs}{proved}{PI18}{}

The formulas of \cref{inv:B2.1,inv:B2.2} give the following values,
where $r_1,r_2$ are the canonical-divisor coefficients and
$P_n^-=h^0(-nK_X)$:
\[
\begin{array}{c|c|c|c|c|ccc}
(k,l)&(m_1,m_2)&(r_1,r_2)&(\deg x,\deg y)&d&P^-_1&P^-_2&P^-_3\\\hline
(1,1)&(3,4)&(0,2)&(4,3)&6&1&2&2\\
(1,2)&(3,8)&(0,5)&(8,3)&9&1&1&2\\
(3,1)&(9,4)&(2,2)&(4,9)&10&0&1&1\\
(1,4)&(3,16)&(0,11)&(16,3)&15&1&1&1\\
(5,1)&(15,4)&(4,2)&(4,15)&14&0&1&1\\
(1,5)&(3,20)&(0,14)&(20,3)&18&1&1&1
\end{array}
\]
Each entry follows by substitution in the displayed floor formula,
or by counting monomials of weights $(m_2,m_1)$.
\end{remark}

\subsection{Finite fundamental groups and universal covers}

\begin{lemma}[Contraction of the cusp tube]
\hypertarget{tag-PI19}{}\label{I:lem:cusp-tube-contraction}
\StatementRecord{I:lem:cusp-tube-contraction}{proved}{PI19}{The normal-crossing charts of the proper local cusp filling}
The reduced one-component hexagonal cusp fibre $W$ in the proper local
filling has a closed tube $N=\{|t_c|\le\epsilon\}$ that strongly
deformation retracts onto $W$. The retraction fixes $W$ pointwise.
\end{lemma}

\begin{proof}
Choose $\epsilon>0$ small enough that the compact tube is covered by
finitely many normal-crossing charts
$t_c=z_1\cdots z_r$, $1\le r\le3$. Off the divisor, take the real
vector field whose $z_i$ component is
\[
 -\frac{z_i|z_i|^{-2}}{\sum_{j=1}^r|z_j|^{-2}},
\]
with its other coordinate components zero. It satisfies
$\dot t_c=-t_c$. Each component has absolute value at most
$(\sum_j|z_j|^{-2})^{-1/2}\le\min_j|z_j|
\le|t_c|^{1/r}$, so the norm is at most $C|t_c|^{1/r}$.
A finite partition of unity preserves $\dot t_c=-t_c$ and gives a
global field of norm at most $C'|t_c|^{1/3}$ after shrinking the tube
so that $|t_c|\le1$. Extend the field by zero on $W$.

Its flow $\Phi_T$ on $N\setminus W$ satisfies
$t_c(\Phi_T(x))=e^{-T}t_c(x)$. It therefore remains in the compact
tube and exists for every $T\ge0$. Its length after time $T$ is at most
\[
 3C'|t_c(x)|^{1/3}e^{-T/3}
 \le3C'\epsilon^{1/3}e^{-T/3}.
\]
The flow converges uniformly to a continuous map $r:N\to W$,
with $r|_W=\mathrm{id}_W$. The same estimate gives continuity at
$W$ after extending every $\Phi_T$ by the identity there. The map
\[
 H(x,u)=
 \begin{cases}
 \Phi_{-\log(1-u)}(x),&0\le u<1,\\
 r(x),&u=1
 \end{cases}
\]
is continuous at $u=1$ by the uniform tail estimate. It has
$H(x,0)=x$, $H(x,1)\in W$, and $H(w,u)=w$ for $w\in W$,
which proves the strong deformation retraction.
\end{proof}

\begin{proposition}[Connected cyclic covers and the cusp comparison]\hypertarget{tag-01S7}{}\label{inv:B3.1}
\StatementRecord{inv:B3.1}{proved}{01S7}{\cref{I:lem:cusp-tube-contraction}, \SeriesRef{inv:B1.1}}

For an admissible $(3,4)$ member with
$n=|12l_0-4l_1-3l_2|>1$, let $\pi:Y\to X$ be the universal
cover and $g=f\pi$. It is a finite \'etale holomorphic map of degree $n$.
Every fibre of $g$ is connected; its finite reductions are bielliptic,
its multiplicities are $3,4$, and its cusp $W_n$ normalizes to a
disjoint union of $n$ copies of $\mathrm{dP}_6$.
For $\mathcal E_X$ of Definition~\ref{I:def:coarse-log-cotangent}
and the identical quotient $\mathcal E_Y$ on $Y$, pullback induces isomorphisms
\[
R^qf_*\OO_X\simeq R^qg_*\OO_Y,\qquad
R^qf_*\mathcal E_X\simeq R^qg_*\mathcal E_Y.
\]
\end{proposition}


\begin{proof}
\begin{enumerate}[label=\arabic*.,leftmargin=*]
\item The smooth family over the punctured base has zero section, so its
fundamental group is $\Lambda\rtimes\langle x,y\rangle$, where
$x,y$ are the analytic peripheral generators fixed by the root-deck convention and act by $A_1,A_2$.
On a finite filling the lifted affine generator has $m_j$th power
translation by $v_j$: its $m_j$ translations add to $v_j$ because
$A_jv_j=v_j$. Thus filling the normal disc imposes
$x^3=t_{v_1}$ and $y^4=t_{v_2}$, with the signs fixed by the
clockwise convention on the first page.

\item The toric covering space of the cusp tube is simply connected.
Indeed its dense torus has cocharacter lattice $\Z^3$; a loop may
be perturbed into that torus, and the normal discs of the three rays
$(0,1),(e_1,1),(e_2,1)$ kill a basis of this lattice.
This argument applies also to the infinite fan, since a compact loop
and a homotopy meet only finitely many of its locally finite charts.
The same conclusion holds for the small tube: if a loop in
$|t_c|<\epsilon$ bounds a disc in the toric space, the torus action
$(x_1,x_2,t_c)\mapsto(x_1,x_2,a t_c)$ for a sufficiently small
positive $a$ moves that compact disc into the tube, and the scaling
homotopy of its boundary stays in the tube.
The cusp quotient therefore has fundamental group
$\Lambda/\Lambda_{\rm tor}$, and its toric meridian bounds the
normal disc at fixed $x_1,x_2$. A translated collar makes this
meridian $t_\mu y^{-1}x^{-1}$, where $\gamma(\mu)=l_0$.
Van Kampen consequently kills $\Lambda_{\rm tor}$ and imposes
$xy=t_\mu$.

\item The normal closure of $\Lambda_{\rm tor}$ in the fibre lattice is
$\ker\gamma$: it contains $\hat w,\hat\delta$ and
$A_1\hat w+\hat w=\hat u$, and all these relations lie in
$\ker\gamma$. The remaining generator $c$, the image of
$\hat\gamma$, is central because $\gamma A_j=\gamma$.
The presentation is thus
$\langle c,x,y\mid c\text{ central},\ xy=c^{l_0},\
x^3=c^{l_1},\ y^4=c^{l_2}\rangle$.
Eliminating $y$ makes it abelian, with relation matrix
\[
 \begin{pmatrix}3&-l_1\\4&l_2-4l_0\end{pmatrix},
 \qquad\det=-p.
\]
Its first Smith divisor is one because its entries include three
and four; hence the group is $\Z/|p|$.
After imposing $c=1$ the quotient has $x^3=x^4=1$, so it is trivial.
Thus the image of the regular-fibre lattice, containing $c$, generates
the entire group $\Z/n$. The same is true of each reduced special
fibre. At a finite point the cyclic model is a smooth torus family
over a disc, with a free affine finite-group action. An equivariant
smooth trivialization, obtained by averaging a horizontal connection,
contracts the disc to its centre and descends to a deformation
retraction of the local total space onto $S_j$. Thus the image of
the regular-fibre lattice in the local fundamental group is contained
in the image of $\pi_1(S_j)$; the same holds after mapping to
$\pi_1(X)$. At the cusp $\hat\gamma$ survives in
$\Lambda/\Lambda_{\rm tor}$. The explicit normal-crossing contraction
of the tube onto $W$ in Lemma~\ref{I:lem:cusp-tube-contraction} identifies their fundamental groups,
so the same generator is in the image of $\pi_1(W)$.
The covering-space criterion therefore makes the inverse image of
every fibre connected. Admissibility gives $3\nmid l_1$ and odd $l_2$,
so $\gcd(n,12)=1$. The \'etale cover preserves fibre multiplicities.
On a finite reduction it cannot kill the nontrivial order-$3$ or
order-$4$ linear holonomy, since the covering degree is prime to it;
the resulting finite torus quotient is again bielliptic.

\item Normalization commutes with the finite \'etale base change $Y\to X$.
The pullback of the simply connected normalization $\mathrm{dP}_6$
is its trivial degree-$n$ cover. The double curves and triple points
also pull back \'etale, and their incidence maps are the lifts of the
original incidence maps. The lifted incidence maps therefore make the dual complex a
degree-$n$ covering of the original real two-torus. It is connected because
$W_n$ is connected. The normalization resolution for $\OO_{W_n}$
has acyclic components, double curves and point strata. Its global
cochain complex is consequently the cellular cochain complex of this
dual torus, with $n,3n,2n$ cells. Thus its coherent cohomology is
$(1,2,1)$, and pullback from $W$ is an isomorphism over $\C$:
on a finite cover of real tori it is induced by a full-rank map of
their integral first cohomology lattices.

\item Modulo $d\log q$, the two toric forms $d\log x_i$ are invariant
under all cusp gluings as computed in \ref{inv:B1.1}. Hence
$\mathcal E_X|_W=\OO_W^2$ and $\mathcal E_Y|_{W_n}=\OO_{W_n}^2$,
and the same pullback argument applies to their cohomology.
Over a regular fibre, an isogeny of complex tori induces isomorphisms
on $H^q(\OO)$ and $H^q(\Omega^1)$ over $\C$.
At a finite fibre, make the cyclic base change. The map is an isogeny
of the smooth torus families and induces isomorphisms of their
cohomology bundles. Taking the same cyclic invariant eigenspaces
descends these isomorphisms, including all nilpotent layers of the
scheme fibre. Grauert base change at the cusp, and this explicit
calculation at the finite points, prove the asserted bundle maps are
isomorphisms on all of $B$.
\end{enumerate}
\end{proof}

\begin{proposition}[Cotangent direct images and Hodge numbers of the covers]\hypertarget{tag-01S8}{}\label{inv:B3.2}
\StatementRecord{inv:B3.2}{proved}{01S8}{\SeriesRef{inv:B3.1}, \SeriesRef{inv:B1.2}}

Write $T_n=\C_{p_0}^{\,n-1}$. Then
\[
\begin{array}{c|ccc}
 &R^0g_*&R^1g_*&R^2g_*\\\hline
\Omega^1_Y&\OO(-2)&\OO\oplus\OO(-1)^2\oplus T_n&\OO\oplus\OO(-1)\\
\Omega^2_Y&\OO(-2)\oplus\OO(-1)&\OO(-2)\oplus\OO(-1)^2&\OO\oplus T_n
\end{array}
\]
Writing $V$ for the locally free quotient, the splittings form an
affine space under $\operatorname{Hom}(V,T_n)$, by the exact
sequence in the proof.
The only nonzero Hodge numbers of $Y$ are
\[
h^{0,0}=h^{0,1}=h^{3,2}=h^{3,3}=1,\quad
h^{1,2}=h^{2,1}=1,\quad h^{1,1}=h^{2,2}=n+1.
\]
The original finite-fundamental-group threefold $X$ itself has the
same Hodge diamond as the simply connected $(3,4)$ member.
\end{proposition}


\begin{proof}
\begin{enumerate}[label=\arabic*.,leftmargin=*]
\item \ref{inv:B3.1} identifies the logarithmic cotangent sequence on cohomology,
and therefore its connecting maps and direct images, with those of
the original family. The normalized residue sheaf differs only at
the cusp: its degree-zero cohomology there is now $\C^n$.
The residue of a base logarithmic form is the diagonal vector
$(1,\ldots,1)$, since each cusp component has multiplicity one.
Its cokernel is $T_n$. The degree-one residues at the finite fibres
cancel the same two skyscrapers as in \ref{inv:B1.2}. The long exact sequence
therefore gives
\[
0\longrightarrow T_n\longrightarrow R^1g_*\Omega^1_Y
\longrightarrow\OO\oplus\OO(-1)^2\longrightarrow0
\]
and the remaining entries of the first row. The extension splits
because $\operatorname{Ext}^1(V,T_n)=H^1(V^*\otimes T_n)=0$
for a vector bundle $V$ on $B$.

\item For a sheaf $F$ locally free on the total space, relative duality
in its derived form gives
\[
Rg_*(F^*\otimes\omega_{Y/B})
\simeq R\mathcal Hom_B(Rg_*F,\OO_B)[-2].
\]
For $F=\Omega^1_Y$, the torsion in $R^1g_*F$ contributes
$\mathcal Ext^1_B(T_n,\OO_B)$ in degree two of the right side.
This is a length-$(n-1)$ skyscraper at $p_0$. After tensoring by
$\OO(-2)$ it still has that length. The locally free parts dualize
in complementary degrees, giving exactly the second row; the
extension by the locally free quotient splits for the same reason
as above. Coherent Leray computes the displayed Hodge numbers.
For $X$ with one-component cusp, \ref{inv:B1.2} already applies independently
of its fundamental group.
\end{enumerate}
\end{proof}

\begin{proposition}[Canonical rings and rational topology of the covers]\hypertarget{tag-01S9}{}\label{inv:B3.3}
\StatementRecord{inv:B3.3}{proved}{01S9}{\SeriesRef{inv:B3.1}, \SeriesRef{inv:B2.2}, \SeriesRef{inv:B3.2}}

For $X$ and $Y$ as above,
\[
K_Y=\pi^*K_X=\OO_Y(-2S'_2),\quad
\bigoplus_{m\ge0}H^0(Y,-mK_Y)=\C[U,V],\quad
\deg(U,V)=(1,2).
\]
The same ring holds for $X$.
Over $\C$, the Betti numbers of $Y$ are
$(1,0,n-1,0,n-1,0,1)$; its Fr\"olicher sequence degenerates at $E_2$.
For every nontrivial character of the deck group, the character
part of Dolbeault cohomology consists of one class in $(1,1)$ and
one in $(2,2)$ and no others.
\end{proposition}


\begin{proof}
\begin{enumerate}[label=\arabic*.,leftmargin=*]
\item The canonical formula pulls back through an \'etale map, and
$4S'_2=g^*p_2$. \ref{inv:B3.1} gives $g_*\OO_Y=\OO_B$.
The valuation argument in \ref{inv:B2.2} therefore applies unchanged and proves
the ring statement with all multiplication maps.

\item The comparisons of \ref{inv:B3.1} are pullback maps and identify the logarithmic
cohomology bundles with the invariant part; the deck group acts
trivially on those bundles. At the cusp its action on the $n$
normalization components is the regular permutation representation.
The quotient by the diagonal has each nontrivial character once.
Thus the first-row additional skyscraper in \ref{inv:B3.2} has precisely these
characters, and duality gives the conjugate characters in its
second-row additional skyscraper. All other cohomology is invariant
and equal to that of $X$. Coherent Leray now proves the assertion
about the Dolbeault character spaces.

\item The de Rham complex and its Fr\"olicher sequence split under the
idempotents of the finite deck group. A nontrivial character has just
the two stated Dolbeault entries. No differential can connect them
or connect either to another nonzero entry, so both survive.
The invariant summand is the de Rham complex of $X$ by finite
\'etale descent and averaging. Its rational topology is computed by the following attachment maps. In the remainder of this proof
all homology and cohomology coefficients are $\Q$.

\item Let $N$ be a small closed cusp tube, $M=\partial N$, and
$U=X\setminus\operatorname{int}N$. The space $U$ is covered by
two pieces retracting onto $S_1,S_2$, whose intersection retracts
onto a regular fibre $F$. The retractions are obtained from the
equivariant smooth torus trivializations used in \ref{inv:B3.1}. The maps
$F\to S_j$ are the central finite covering maps; translation of
their origins induces the identity on their rational cohomology.
Finite-cover transfer identifies their pullbacks with the invariant
subspaces for $A_j$. Put
\[
 \psi_1=2u+w+3\delta,\qquad \psi_2=u+w+2\delta,
 \qquad \eta=u\wedge w+6\gamma\wedge\delta.
\]
Direct substitution in $A_1,A_2$ shows that the invariant
degree-one spaces are $\langle\gamma,\psi_j\rangle$ and the
invariant degree-two spaces are
$\langle\gamma\wedge\psi_j,\eta\rangle$.
Their dimensions are two: the nontrivial eigenvalues of $A_j$ on
the complementary plane are the inverse primitive third or fourth
roots, so the exterior-square fixed space has exactly the two
displayed generators.
In the transfer-normalized dual homology bases, the two
Mayer--Vietoris maps are therefore the following matrices:
\[
 \alpha_1=\begin{pmatrix}
 1&0&0&0\\0&2&1&3\\-1&0&0&0\\0&-1&-1&-2
 \end{pmatrix},\qquad
 \alpha_2=\begin{pmatrix}
 2&1&3&0&0&0\\0&0&6&1&0&0\\
 -1&-1&-2&0&0&0\\0&0&-6&-1&0&0
 \end{pmatrix}.
\]
The second matrix uses the ordered wedge basis
$(\hat\gamma\hat u,\hat\gamma\hat w,\hat\gamma\hat\delta,
\hat u\hat w,\hat u\hat\delta,\hat w\hat\delta)$.
Both matrices have rank three, and
$\ker\alpha_1=\Q\nu_0$, with
$\nu_0=-\hat u-\hat w+\hat\delta$.
Consequently
\begin{equation}
 H_1(U)=\Q,\qquad
 0\longrightarrow\operatorname{coker}\alpha_2=\Q
 \longrightarrow H_2(U)
 \xrightarrow{\partial}\Q\nu_0\longrightarrow0.
\label{pub:I:eq:B-1}
\end{equation}
The generator $\Theta\in H^1(U)$ is normalized by
$\Theta|_F=\gamma$. The finite affine generators have
$3\bar g_1=l_1c$, $4\bar g_2=l_2c$ in rational homology, where
$c$ is the image of $\hat\gamma$ and $\Theta(c)=1$.
For the toric meridian $a_0=t_\mu y^{-1}x^{-1}$ this gives the
explicit collar value
\begin{equation}
 \Theta(j_*a_0)=l_0-l_1/3-l_2/4=p/12.
\label{pub:I:eq:B-2}
\end{equation}

\item Lemma~\ref{I:lem:cusp-tube-contraction} gives the strong
deformation retraction $N\to W$, fixing $W$ pointwise.

\item The constant-sheaf normalization resolution of $W$ has strata
$\mathrm{dP}_6$, three copies of $\Pone$, and two points.
Its row of degree-zero cohomology is
\[
 \Q\xrightarrow{0}\Q^3
 \xrightarrow{(a,b,c)\mapsto(a-b+c,a-b+c)}\Q^2.
\]
The strata have no degree-one cohomology. This resolution therefore
gives $H_1(N)=H_1(W)=\Q^2$, $H_4(N)=\Q[W]$ and $H_5(N)=0$;
it also gives $e(W)=6-3\cdot2+2=2$.
Poincar\'e--Lefschetz duality for the compact smooth manifold with
boundary $(N,M)$ now shows that $H_1(M)\to H_1(N)$ and
$H_2(M)\to H_2(N)$ are surjective. For degree one the next
relative group is $H_1(N,M)=H^5(N)=0$.
For degree two the map to $H_2(N,M)=H^4(N)$ is the intersection
pairing with $[W]$, which is zero because $W$ is the divisor of
the global function $t_c$ on $N$. The Wang sequence for the
boundary mapping torus, with $\rank(M_0-I)=2$, gives
$\dim H_1(M)=3$. Its kernel to $H_1(N)$ is exactly
$\Q a_0$: this meridian is nonzero under projection to the base
circle and bounds its normal disc in $N$.

\item We exhibit two classes in
$K_2=\ker(H_2(M)\to H_2(N))$.
For $\nu=\hat w,\hat\delta$, let $\xi_\nu$ be the two-torus
obtained by translating the toric meridian around the angular
circle $S^1_\nu$. In the chart $t_c=z_0z_1z_2$ corresponding to
the rays $(0,1),(e_1,1),(e_2,1)$, these are the boundaries of
$D^2_{z_0}\times S^1_{z_i}$ with the other coordinate fixed
and the radii chosen so that $|t_c|=\epsilon$ on the boundary.
They therefore bound in $N$ with the specified toric meridian.
Cutting the meridian into the two arcs of the cover defining
\eqref{pub:I:eq:B-1}, the two cylinders have end circles differing by
$(I-A_1^{-1})\nu$. With this orientation of the Mayer--Vietoris
boundary, matrix multiplication gives
\begin{equation}
 \partial j_*\xi_{\hat w}=-\nu_0,\qquad
 \partial j_*\xi_{\hat\delta}=0.
\label{pub:I:eq:B-3}
\end{equation}
The affine translation of either endpoint has no effect on its
homology class, which explains why this boundary formula is the
same for every translation twist.

\item Choose $\mathcal A\in H^2(U)$ restricting to $\eta$ on $F$;
it exists by Mayer--Vietoris because the two transfer-normalized
classes $\eta$ have equal restrictions. Fibrewise translation by
$S^1_{\hat\delta}$ defines a global smooth action on $U$:
$A_j\hat\delta=\hat\delta$, and all affine translations commute
with this circle. Slant product of its action map with the circle
defines $\iota:H^2(U)\to H^1(U)$. Naturality on $F$ gives
$\iota\mathcal A|_F=\iota_{\hat\delta}\eta=-6\gamma$.
The restriction $H^1(U)\to H^1(F)$ is injective by \eqref{pub:I:eq:B-1}, so
$\iota\mathcal A=-6\Theta$. By the definition of slant product,
\begin{equation}
 \langle\mathcal A,j_*\xi_{\hat\delta}\rangle
 =\langle\iota\mathcal A,j_*a_0\rangle=-p/2.
\label{pub:I:eq:B-4}
\end{equation}
We may add the class pulled back from $(\Q\nu_0)^*$ under
the boundary map in \eqref{pub:I:eq:B-1} so that
$\langle\mathcal A,j_*\xi_{\hat w}\rangle=0$;
this does not change its restriction to $F$ or \eqref{pub:I:eq:B-4}.
In the coordinates on $H_2(U)$ given by $\partial/\nu_0$ and
this normalized pairing, the two attachment cycles have
matrix
\begin{equation}
 \bigl[j_*\xi_{\hat w}\ \ j_*\xi_{\hat\delta}\bigr]
 =\begin{pmatrix}-1&0\\0&-p/2\end{pmatrix}.
\label{pub:I:eq:B-5}
\end{equation}
It is invertible over $\Q$ because admissibility implies
$p\ne0$. Thus the image of $K_2$ already fills $H_2(U)$.
Likewise \eqref{pub:I:eq:B-2} says the image of $\ker(H_1(M)\to H_1(N))$
is nonzero in $H_1(U)$. Using the surjectivity to $N$ just proved,
Mayer--Vietoris for $X=U\cup N$ now gives $H_1(X)=H_2(X)=0$.
Every other fibre has Euler number zero, so additivity over the
three special fibres and the smooth torus bundle gives $e(X)=2$.
Poincar\'e duality then gives $b_4=b_5=0$ and $b_3=0$.
This proves the asserted rational homology of $X$ for all admissible
twists; the only changing determinant in \eqref{pub:I:eq:B-2}--\eqref{pub:I:eq:B-5} is a nonzero
constant multiple of $p$. 

\item Proposition~\ref{inv:C1.1}, applied to this rational topology and
the same Hodge table, gives degeneration of the invariant summand at $E_2$.
Summing it with the nontrivial character spaces proves the Betti
numbers and the degeneration statement for $Y$. In particular
$e(Y)=2n$, as also follows from the covering degree.
\end{enumerate}
\end{proof}



\subsection{CM and torsion}
\begin{cited}[The CM torsion fibres]\hypertarget{tag-000U}{}\label{ari:C.5.1}
\StatementRecord{ari:C.5.1}{proved}{000U}{\citet[Theorem~4.15]{BourdonClarkStankewicz2017}; \citet[Theorem~20.12 and Corollary~20.13]{Sutherland2023CM}}

There is no $t\in\mathbb C\setminus\{0,1\}$ for which $E_t$ has
complex multiplication and $P_t$ is torsion. On the smooth models
at $t=0,1$ after base change of orders $3,4$, respectively, the
limiting points have orders $3,2$. Thus these are the only CM
torsion points on the completed family.
\end{cited}
\begin{proof}
Theorem~4.15(a) of
\cite[arXiv v3, p.~23]{BourdonClarkStankewicz2017} applies to an
elliptic curve with CM by any imaginary quadratic order over
$F=\mathbb Q(j(E))$. Its torsion group is one of
$0,\mathbb Z/2,\mathbb Z/3,\mathbb Z/4,\mathbb Z/6,(\mathbb Z/2)^2$.
Here $F=\mathbb Q(t)$ because $j=1728t$, and the equation
and section of \cref{I:def:theta-height} are defined over $F$.
On the smooth base, $\psi_2(P)$ and $\psi_3(P)$ are nonzero.
Thus only orders four and six remain.

The order-four value has $j=55296/5$. The order-six values have
irreducible monic polynomial
\[
 X^2-\frac{10395648}{11}X+\frac{24461180928}{11}.
\]
None is an algebraic integer. The integrality of singular moduli
for every imaginary quadratic order
\cite[Theorem~20.12 and Corollary~20.13, Lecture~20, p.~6]{Sutherland2023CM}
therefore excludes CM in all these cases.

At zero, the changes $t=s^3$, $x=s^4X$, $y=s^6Y$ give
\[
 Y^2=X^3-\tfrac34s(s^3-1)X+\tfrac14(s^3-1)^2,\qquad
 P_s=(0,(s^3-1)/2).
\]
The central equation is $Y^2=X^3+1/4$, and the limiting point
$(0,-1/2)$ has order three by its horizontal tangent.
At one, $t=1+s^4$, $x=s^2X$, $y=s^3Y$ give
\[
 Y^2=X^3-\tfrac34(1+s^4)^3X+
            \tfrac14(1+s^4)^4s^2,\qquad
 P_s=(0,(1+s^4)^2s/2).
\]
The central equation is $Y^2=X^3-3X/4$, and $(0,0)$ has
exact order two. Both central discriminants are nonzero.
Their CM orders are $\mathbb Z[\rho]$ and $\mathbb Z[i]$,
and their norm values are those of \cref{ari:C.1.2}.
\end{proof}



\section{Global heights and the eleven rational CM fibres}\label{pub:I:cm-height}

\begin{remark}[Arithmetic normalization]
\hypertarget{tag-IA25}{}\label{I:global-height-convention}
\StatementRecord{I:global-height-convention}{proved}{IA25}{}
The arithmetic equation, invariant differential, place weights,
and height normalization are those of \cref{I:def:theta-height}.
\end{remark}

\begin{proposition}[The global height formula]\hypertarget{tag-000V}{}\label{ari:C.6.1}
\StatementRecord{ari:C.6.1}{proved}{000V}{\SeriesRef{ari:C.7.1}; \citet[Section 2.1]{demarco2017}}

For every number field $K$, every $t\in K\setminus\{0,1\}$,
and the section just specified,
\begin{equation}\hypertarget{tag-0039}{}\label{ari:eq:ch-global}
 \sum_{\sigma:K\hookrightarrow\C}\log N(P_{\sigma(t)})
 =\sum_{v\nmid\infty}n_v\lambda_v(P_t)
          -[K:\Q]\widehat h_O(P_t).
\end{equation}
Here the local heights are invariant under change of Weierstrass
model and satisfy, when $nP\ne O$,
\begin{equation}\hypertarget{tag-003B}{}\label{ari:eq:ch-local-mult}
 \lambda_v(nP)=n^2\lambda_v(P)-\log|\psi_n(P)|_v
                   +\frac{n^2-1}{12}\log|\Delta|_v.
\end{equation}
At a complex embedding $\lambda_v(P)=-\log N(P)$; complex
embeddings are counted separately in \eqref{ari:eq:ch-global}.
For rational $t$, put $A_t=\exp(\sum_p\lambda_p(P_t))$.
Then $A_t$ is a positive real algebraic number and
\[
 N(P_t)=A_t e^{-\widehat h_O(P_t)},\qquad
 D_{\rm part}(t)=\frac6\pi
       \left(\log\frac{A_t}{2\sqrt3}-\widehat h_O(P_t)\right).
\]
Thus $N(P_t)$ is algebraic if and only if
$e^{\widehat h_O(P_t)}$ is algebraic.
\end{proposition}


\begin{proof}
\begin{enumerate}[label=\arabic*.,leftmargin=*]
\item The sigma expression for the archimedean local height is
\[
 \lambda_v(P)=-\log|\sigma(u)|
       +\frac12\Re(u\eta(u))-\frac1{12}\log|\Delta|.
\]
The sigma--theta relation and the Legendre relation identify this
with $-\log N$ exactly, including its constant; this is also the
normalization in \cref{ari:C.4.1}. The division identity
$\sigma(nu)=(-1)^{n-1}\psi_n(P)\sigma(u)^{n^2}$ gives
\eqref{ari:eq:ch-local-mult}. At finite places the same identity is
the defining multiplication law of the canonical N\'eron function.
Its normalization is uniquely determined by continuity away from
$O$, its logarithmic pole at $O$, and the multiplication law.
These are the three properties printed in
\citet[Section~2.1, arXiv v2, pp.~5--6]{demarco2017}; their duplication law includes
the discriminant term $\frac14\log|\Delta|_v$.
For the homothety
$(x,y,\Delta,\psi_n)\mapsto
(c^{-2}x,c^{-3}y,c^{-12}\Delta,c^{1-n^2}\psi_n)$,
the change in the two logarithmic terms of
\eqref{ari:eq:ch-local-mult} is
$-(1-n^2)\log|c|_v-(n^2-1)\log|c|_v=0$.
The pole normalization is transported by the same formal-parameter map.

\item Sum the duplication law over places. The product formula cancels
the division polynomial and discriminant terms, so that the sum
$H(P)=[K:\Q]^{-1}\sum_vn_v\lambda_v(P)$ satisfies
$H(2P)=4H(P)$. The pole normalization implies
$H(P)=\frac12h(x(P))+O_E(1)$. Iterating multiplication by two
therefore identifies $H$ with $\widehat h_O$. Moving the
archimedean terms to the other side gives \eqref{ari:eq:ch-global}.
For a finite place, a multiple of $P$ lies in the identity component
of a minimal model. There the local height is a rational multiple
of $\log p$; the multiplication law shows the same for $P$.
Only finitely many places contribute. This proves the assertion
about $A_t$ and both of its consequences.
\end{enumerate}
\end{proof}

\begin{proposition}[Finite local heights]\hypertarget{tag-000W}{}\label{ari:C.6.2}
\StatementRecord{ari:C.6.2}{proved}{000W}{\SeriesRef{ari:eq:ch-local-mult}}

Let $E/\Q$ be a global minimal Weierstrass model. If $mP$ is
in the identity component at $p$, and $mP\ne O$, then
\begin{equation}\hypertarget{tag-003A}{}\label{ari:eq:ch-local-exact}
 \frac{\lambda_p(P)}{\log p}
 =\frac{v_p(\Delta)}{12}
   +\frac{\frac12\max(0,-v_p(x(mP)))-v_p(\psi_m(P))}{m^2}.
\end{equation}
For $m=1$ we read $\psi_1=1$. This formula, with exact rational
coordinates and valuations, gives all the finite contributions
in the tables below.
\end{proposition}


\begin{proof}
\begin{enumerate}[label=\arabic*.,leftmargin=*]
\item On the identity component of a minimal model the canonical
N\'eron function is
$\frac12\log\max(1,|x|_p)-\frac1{12}\log|\Delta|_p$.
This follows either from the formal parameter at $O$ followed
by the group law on the smooth locus, or from the uniqueness
properties used above. Substitute it for $\lambda_p(mP)$ in
\eqref{ari:eq:ch-local-mult} and divide by $m^2\log p$.
The contribution $v_p(\Delta)/(12m^2)$ combines with the
discriminant term to give $v_p(\Delta)/12$.
Membership in the identity component is checked by reduction
to the nonsingular locus of the minimal cubic. Thus the test is
exact also at $2$ and $3$.
\end{enumerate}
\end{proof}

\begin{proposition}[The CM models and their height identities]\hypertarget{tag-000X}{}\label{ari:C.6.3}
\StatementRecord{ari:C.6.3}{proved}{000X}{\SeriesRef{ari:eq:ch-local-exact}}

The eleven rational CM fibres have the global minimal models,
section coordinates, component classes, and finite local heights
in the following tables. On each model $a_1=a_2=0$, and we
write $[a_3,a_4,a_6]$ for
$y^2+a_3y=x^3+a_4x+a_6$. The displayed point is $(0,y_P)$.
For row $i$, put $t_i=j_i/1728$ and
\[
 u_i=-\frac{3t_i(2y_{P,i}+a_{3,i})}{4a_{4,i}}.
\]
The exact pointed isomorphism to the arithmetic equation is
\[
 X=u_i^2x,\qquad Y=u_i^3(y+a_{3,i}/2).
\]
It sends the displayed point to $P_{t_i}$ and pulls back
$dX/(2Y)$ to $u_i^{-1}dx/(2y+a_{3,i})$.
The constants $A_t$ in increasing row order are
\[
 \sqrt6,\quad 2^{2/3}3^{1/4},\quad 2^{3/4},\quad
 2,\quad 2,\quad \sqrt3,\quad 1,\quad 3^{5/12},\quad
 1,\quad1,\quad1.
\]
\begingroup\normalfont\footnotesize\setlength{\tabcolsep}{4pt}
\begin{longtable}{r r r r r r}
\toprule Row & $j$ & $a_3$ & $a_4$ & $a_6$ & $y_P$\\\midrule\endhead
1 & 8000 & 0 & -5880 & 153664 & 392\\
2 & 54000 & 0 & -7260 & 234256 & 484\\
3 & 287496 & 0 & -2156 & 38416 & 196\\
4 & 16581375 & 0 & -7732620 & 8275904784 & 90972\\
5 & -3375 & 0 & -140 & 784 & -28\\
6 & -32768 & 1 & -3234 & 72630 & -270\\
7 & -884736 & 1 & -38 & 90 & -10\\
8 & -12288000 & 1 & -1920270 & 1024288020 & -32005\\
9 & -884736000 & 1 & -379260 & 89898842 & -9482\\
10 & -147197952000 & 1 & -347045930 & 2488447882842 & -1577482\\
11 & -262537412640768000 & 1 & -75065352465446420 & 7916029533155605844266842 & -2813543945482\\
\bottomrule\end{longtable}
\begin{longtable}{r r l c c l}
\toprule Row & $p$ & Type & Component & Order & $\lambda_p/\log p$\\\midrule\endhead
1 & 2 & $III^*$ & $c_P$ & 2 & $\frac{1}{2}$\\
1 & 3 & $I_0^*$ & $0$ & 1 & $\frac{1}{2}$\\
1 & 7 & $I_0^*$ & $c_P$ & 2 & $0$\\
2 & 2 & $I_{4}^*$ & $c_P$ & 2 & $\frac{2}{3}$\\
2 & 3 & $III$ & $0$ & 1 & $\frac{1}{4}$\\
2 & 11 & $I_0^*$ & $c_P$ & 2 & $0$\\
3 & 2 & $I_{5}^*$ & $c_P$ & 2 & $\frac{3}{4}$\\
3 & 7 & $I_0^*$ & $c_P$ & 2 & $0$\\
4 & 2 & $I_{8}^*$ & $c_P$ & 2 & $1$\\
4 & 3 & $I_0^*$ & $c_P$ & 2 & $0$\\
4 & 7 & $III$ & $c_P$ & 2 & $0$\\
4 & 19 & $I_0^*$ & $c_P$ & 2 & $0$\\
5 & 2 & $I_{8}^*$ & $c_P$ & 2 & $1$\\
5 & 7 & $III$ & $c_P$ & 2 & $0$\\
6 & 3 & $I_0^*$ & $0$ & 1 & $\frac{1}{2}$\\
6 & 7 & $I_0^*$ & $c_P$ & 2 & $0$\\
6 & 11 & $III$ & $c_P$ & 2 & $0$\\
7 & 19 & $III$ & $c_P$ & 2 & $0$\\
8 & 3 & $IV$ & $0$ & 1 & $\frac{5}{12}$\\
8 & 11 & $I_0^*$ & $c_P$ & 2 & $0$\\
8 & 23 & $I_0^*$ & $c_P$ & 2 & $0$\\
9 & 3 & $I_0^*$ & $c_P$ & 2 & $0$\\
9 & 7 & $I_0^*$ & $c_P$ & 2 & $0$\\
9 & 43 & $III$ & $c_P$ & 2 & $0$\\
10 & 7 & $I_0^*$ & $c_P$ & 2 & $0$\\
10 & 31 & $I_0^*$ & $c_P$ & 2 & $0$\\
10 & 67 & $III$ & $c_P$ & 2 & $0$\\
11 & 7 & $I_0^*$ & $c_P$ & 2 & $0$\\
11 & 11 & $I_0^*$ & $c_P$ & 2 & $0$\\
11 & 19 & $I_0^*$ & $c_P$ & 2 & $0$\\
11 & 127 & $I_0^*$ & $c_P$ & 2 & $0$\\
11 & 163 & $III$ & $c_P$ & 2 & $0$\\
\bottomrule\end{longtable}
\endgroup
\end{proposition}


\begin{proof}
\begin{enumerate}[label=\arabic*.,leftmargin=*]
\item Substitution verifies the displayed pointed isomorphisms and
$\Delta(E_t)=u_i^{12}\Delta(E_i)$.
Tate's reduction algorithm on each integral equation gives the
minimality, fibre types and component groups in the table.
For each bad prime, rational tangent addition computes $mP$,
the two partial derivatives of the minimal cubic test membership
in its nonsingular locus, and the division recurrence computes $\psi_m$.
Formula
\eqref{ari:eq:ch-local-exact} then proves each entry. Outside the
listed primes the discriminant is a unit and $P$ is integral,
so the local height is zero. Multiplying $p^{\lambda_p/\log p}$
gives the listed constants $A_t$.

\item The component labels are defined as follows. The label $0$
means the identity component, and $c_P$ is the class represented
by the explicitly displayed section $P$. At an odd prime of
type $I_0^*$ this nonidentity class is the outer component
indexed by the root zero of the residual Tate cubic; the
scaling $x=pX$ identifies that root. When there are two
rational components, $c_P$ is their unique nonzero class.
The only remaining ambiguity is row 5 at $2$, where the
component group is $(\Z/2)^2$. Here let $T=(-14,0)$;
we label it with the basis $(c_P,c_T)$. Indeed
$P+T=(18,64)$ after the corresponding choice of sign,
and $P,T,P+T$ all reduce to the singular point, whereas
$2P=(25/4,99/8)$ reduces to the identity component.
Thus the three nonzero classes are explicitly separated.


\end{enumerate}
\end{proof}


\begin{remark}[Numerical sizes of the CM heights]
\hypertarget{tag-IA02}{}\label{I:CM-height-sizes}
\StatementRecord{I:CM-height-sizes}{proved}{IA02}{}

The following approximations have absolute error less than $10^{-8}$,
by the enclosures in \cref{ari:C.6.5}.
\[
\begin{array}{r|rr}
i&\widehat h_O(P_i)&\log N(P_i)\\\hline
1&0.68444759&0.21143215\\
2&0.64179980&0.09495139\\
3&0.53353021&-0.01366983\\
4&0.99406687&-0.30091969\\
5&0.46736568&0.22578150\\
6&0.42860917&0.12069697\\
7&0.09064584&-0.09064584\\
8&0.73662083&-0.27886571\\
9&0.60182962&-0.60182962\\
10&1.00264929&-1.00264929\\
11&2.16548925&-2.16548925\\
\end{array}
\]
\end{remark}

\begin{proposition}[The maximum as a height inequality]\hypertarget{tag-000Y}{}\label{ari:C.6.4}
\StatementRecord{ari:C.6.4}{proved}{000Y}{\SeriesRef{ari:C.8.5}, \SeriesRef{ari:eq:ch-global}; \citet{demarco2017}; \citet[pp.~183--184]{callSilverman1993}; \citet{silverman2006heights}}

For $t\in K\setminus\{0,1\}$,
\begin{equation}\hypertarget{tag-0038}{}\label{ari:eq:ch-bound}
 \widehat h_O(P_t)\geq
 \frac1{[K:\Q]}\sum_{v\nmid\infty}n_v\lambda_v(P_t)
                    -\frac12\log2.
\end{equation}
The inequality is strict on this smooth base. Over $\Q$ it
is $\widehat h_O(P_t)\geq\log A_t-\frac12\log2$.
The generic height is $\widehat h_O(P)=1/12$, and the
variation divisor is $D_P=(\infty)/12$. Consequently
\[
 \widehat h_O(P_t)=\frac1{12}h(t)+O(1)
 \qquad(t\in\overline\Q\setminus\{0,1\}),
\]
where the constant is uniform over algebraic $t$ and is for
the fixed family and section.
\end{proposition}


\begin{proof}
\begin{enumerate}[label=\arabic*.,leftmargin=*]
\item Apply the analytic maximum theorem \ref{ari:C.8.5} at every embedding
and substitute into \eqref{ari:eq:ch-global}. Equality would
require $j(\sigma(t))=1728$ for every embedding, hence $t=1$,
which has been excluded.

\item For the generic local heights at $0$ and $1$, the minimal
discriminant orders are $8$ and $3$, and the section meets
the nonidentity component of types $IV^*$ and $III$.
Their local corrections are $2/3$ and $1/4$, respectively,
so their normalized local heights are zero. At infinity,
put $s=1/t$, $X=s^{-2}x$, $Y=s^{-3}y$. The new equation is
\[
 y^2=x^3-\tfrac34(1-s)x+\tfrac14(1-s)^2,
 \qquad P=(0,(1-s)/2),\qquad\Delta=27s(1-s)^3.
\]
At $s=0$ the node is $(1/2,0)$, whereas $P=(0,1/2)$ is
nonsingular. Its height contribution is $1/12$.
At all other function-field places the point is integral
and the discriminant is a unit. This proves
$D_P=(\infty)/12$ and the generic height.

\item Theorem~1.1 of \citet[arXiv v2, p.~2]{demarco2017} applies with $K=\Q$,
$B=\mathbb P^1$, the displayed elliptic surface and rational
section, since the generic height is positive. It equips
this divisor with continuous adelic semipositive metrics
whose associated height is exactly $\widehat h_O(P_t)$.
Comparing with the usual height for $(\infty)/12$ gives the
bounded difference. This is the specialization here of the
variation theorem of Tate, with the local refinements of
Silverman reproduced in \citet[Theorems~2.2--2.3, arXiv v2, pp.~7--8]{demarco2017}:
outside finitely many places the local correction is zero,
and at each remaining place it extends continuously after
subtracting its prescribed logarithmic singularity.
Our inequality is an explicit archimedean estimate with
constant $\frac12\log2$.

\item The same generic coefficient appears in
\cite[Theorem~4.1 and the following remark, pp.~183--184]{callSilverman1993}:
for $T=\mathbb P^1$ and $[2]^*(O)\sim4(O)$, the specialization
ratio tends to $1/12$. The bounded remainder used here is the
adelic-height conclusion applied in the preceding step.
\end{enumerate}
\end{proof}

\begin{proposition}[Explicit height enclosures]\hypertarget{tag-000Z}{}\label{ari:C.6.5}
\StatementRecord{ari:C.6.5}{proved}{000Z}{\SeriesRef{ari:C.6.3}, \SeriesRef{ari:C.5.1}}

For each of the eleven rational CM parameters in \cref{ari:C.6.3},
the Fourier and product algorithms below give rational interval
enclosures of $\log N(P)$ and $\widehat h_O(P)$ of radius less
than $10^{-236}$.
\end{proposition}


\begin{proof}
\begin{enumerate}[label=\arabic*.,leftmargin=*]
\item The CM discriminants, in the order of \ref{ari:C.6.3}, are
\[
 -8,-12,-16,-28,-7,-11,-19,-27,-43,-67,-163.
\]
Enumeration of primitive reduced forms
$(a,b,c)$, $|b|\le a\le c$, $b^2-4ac=D$, with the usual
nonnegative boundary convention, gives exactly one form of
each discriminant, always with $a=1$.
The range is finite and proved by $a\le\sqrt{|D|/3}$.
Thus the CM theorem used in \ref{ari:C.5.1} makes the corresponding
singular modulus an integer. The Fourier enclosures below
place it within $1/4$ of the specified $j$, proving its
identity. 

\item Write $\tau=iT$ for the first four entries and
$\tau=1/2+iT$ for the other seven, with
\[
 (T_1,T_2,T_3,T_4)=(\sqrt2,\sqrt3,2,\sqrt7),\qquad
 T_k=\sqrt{|D_k|}/2\quad(k\ge5).
\]
Put $Q=e^{-2\pi T}$, $q=Q$ or $-Q$ accordingly,
$M=350$, and $Y=2/5$. On $0\le y\le Y$,
\begin{equation}\hypertarget{tag-003C}{}\label{ari:eq:ch-wp-bracket}
 \frac{\wp(1/2-iy)}{\pi^2}
 =\frac1{\cosh^2\pi y}-\frac13+
 8\sum_{n\ge1}\frac{nq^n}{1-q^n}
                 \{1-(-1)^n\cosh(2\pi ny)\}.
\end{equation}
For $a=Qe^{2\pi Y}$, the absolute tail after $M$ in
$\wp$ is at most
\[
 \frac{16\pi^2}{1-Q}
 \frac{a^{M+1}\{M+1-Ma\}}{(1-a)^2}<10^{-875}.
\]
The proof is $1+\cosh(2\pi ny)\le2e^{2\pi nY}$ and
$\sum_{n>M}na^n=a^{M+1}(M+1-Ma)/(1-a)^2$.
A numerical zero approximation selects a rational centre $y_0$. At the two rational endpoints
$y_0\pm10^{-240}$ the truncated expression plus this
infinite tail is respectively strictly positive and strictly
negative. These strict endpoint signs are evaluated with outward error bounds.
Continuity produces a zero within each bracket. Every zero
of $\wp$ is one of the pair $\pm P$, and $N$ agrees on
that pair, so either zero gives the same value of $N$.

\item The product formula on this bracket has positive real factors:
\begin{align*}
 \log N={}&\log(2\cosh\pi y)-\frac{\pi T}{6}
                      -\frac{\pi y^2}{T}\\
 &+\sum_{n\ge1}\{\log(1+q^ne^{2\pi y})
                         +\log(1+q^ne^{-2\pi y})\}.
\end{align*}
Since $|\log(1+x)|\le |x|/(1-|x|)$ for real $|x|<1$,
the omitted tail is bounded by
\[
 \frac{2e^{2\pi Y}Q^{M+1}}
 {(1-Q)(1-e^{2\pi Y}Q^{M+1})}<10^{-1260}.
\]
The finite sums, elementary functions and tails are evaluated by
outward interval arithmetic. The radius and sign comparisons are strict. The resulting radius is less than $10^{-236}$.
Adding the exact finite local logarithms of \ref{ari:C.6.3} and
subtracting this enclosure gives the claimed height bound.

\item The singular-modulus certificate uses
$E_4=1+240\sum n^3q^n/(1-q^n)$ and
$\Delta_0=q\prod(1-q^n)^{24}$. Their respective error
bounds, after truncation at $M$, are
\[
 \frac{240(M+1)^3Q^{M+1}(1+4Q+Q^2)}{(1-Q)^5},
 \qquad
 \frac{24Q^{M+1}}{(1-Q)(1-Q^{M+1})}
\]
for $E_4$ and for the logarithm of the product. The first
uses $(M+1+k)^3\le(M+1)^3(k+1)^3$ and the differentiated
geometric series; the second uses the same logarithm bound.
Thus $E_4^3/\Delta_0$ is enclosed by the stated finite sums and tails.
\end{enumerate}
\end{proof}


\subsection{Ranks, saturation, and central derivatives}

\begin{remark}[The real period and diagonal height pairing]
\hypertarget{tag-IA26}{}\label{I:rank-period-pairing}
\StatementRecord{I:rank-period-pairing}{proved}{IA26}{}
Write $E_i$ for the eleven minimal models in \cref{ari:C.6.3}.
The full real period and the height relative to $(O)$ are those of
\cref{I:def:theta-height}. Pullback of the normalized Poincar\'e bundle
along the diagonal gives
$[2]^*\mathcal O(O)\otimes\mathcal O(O)^{-2}\simeq\mathcal O(2O)$.
Consequently the diagonal N\'eron--Tate pairing used in the
leading-term formula is $2\widehat h_O$.
\end{remark}

\begin{proposition}[Mordell--Weil ranks and primitive sections]\hypertarget{tag-0010}{}\label{ari:C.6.R1}
\StatementRecord{ari:C.6.R1}{proved}{0010}{\citet[Theorem~1, PDF p.~8]{CrCremonaSiksek2006}}

The algebraic and analytic ranks agree for the eleven curves. They are
two for $j=8000$ and $j=-12288000$, and one for all nine other values.
For every rank-one curve the section $P$ is a primitive generator of
$E(\Q)/E(\Q)_{\rm tors}$. Its index is therefore one. For the two
rank-two curves the index of $\Z P$ in this free quotient is infinite.
\[
\begin{array}{r|r|r|r}
i&j&\operatorname{rank}E(\Q)&[E(\Q)/E(\Q)_{\rm tors}:\Z P]\\\hline
1&8000&2&\infty\\
2&54000&1&1\\
3&287496&1&1\\
4&16581375&1&1\\
5&-3375&1&1\\
6&-32768&1&1\\
7&-884736&1&1\\
8&-12288000&2&\infty\\
9&-884736000&1&1\\
10&-147197952000&1&1\\
11&-262537412640768000&1&1
\end{array}
\]
\end{proposition}


\begin{proof}[Arithmetic verification]
\begin{enumerate}[label=\arabic*.,leftmargin=*]
\item Exact $2$-descent gives matching lower and upper rank bounds
for all eleven equations. For each of the nine rank-one sections,
height bounds restrict its possible index to a finite interval;
reduction maps at good primes test divisibility by every prime in
that interval and every prime dividing a component-group order.
These saturation tests give index one in all nine cases.

\item The integer parts of the bounds for the indices in the subgroup
of points having good reduction everywhere are respectively
\[
 1,1,1,2,2,2,3,5,5
\]
in the rank-one order $i=2,3,4,5,6,7,9,10,11$; the algorithm also
checks the primes dividing component-group orders. The underlying
height lower bound is the one in
\citet[Theorem~1, PDF p.~8]{CrCremonaSiksek2006}, applied to a minimal
Weierstrass equation and the subgroup
$E(\Q)\cap E^0(\R)\cap\bigcap_pE^0(\Q_p)$.
For the largest curve the lower bound on the nonzero height in this
subgroup is $0.651940632595967949\ldots$, its regulator is
$17.323913974368488616\ldots$, and the resulting index bound is
$5.154885040175890631\ldots$. The finite-field saturation maps for
$2,3,5$ already have full rank at residue primes $17,17,59$.
The index bounds and full-rank reduction maps certify primitivity.

\item The analytic ranks follow from the explicit derivative sums and bounds
proved below. The two rank-two bases furnished by descent are
$(-7,441),(105,833)$ for $i=1$, and
$(0,32004),(22218,3305456)$ for $i=8$.
\end{enumerate}
\end{proof}

\begin{lemma}[The CM coefficient series]
\hypertarget{tag-IA08}{}\label{I:CM-coefficient-series}
\StatementRecord{I:CM-coefficient-series}{proved}{IA08}{}

For $s=11,43,67,163$, let $E_s$ be the respective rational
elliptic curve $121\mathrm b1$, $1849\mathrm a1$,
$4489\mathrm a1$, or $26569\mathrm a1$.
The coefficients of its normalized newform are
\[
 b_n=\sum_{a^2=n,\ a>0}\left(\frac as\right)a
 +\sum_{\substack{u,v>0,\ u\equiv v\ (2)\\u^2+sv^2=4n}}
       \left(\frac{(s+1)u/2}{s}\right)u.
\]
Quadratic twisting multiplies $b_n$ by the corresponding
quadratic character.
\end{lemma}
\begin{proof}
For $s=11,43,67,163$, put
$F_s=\Q(\sqrt{-s})$ and $\mathfrak p=(\sqrt{-s})$.
These fields have class number one and units $\{\pm1\}$.
The character
\[
 \psi((\alpha))=
  \left(\frac{\alpha\bmod\mathfrak p}{s}\right)\alpha
\]
is well defined because $(-1/s)=-1$. It has infinity type $(1,0)$
and finite conductor $\mathfrak p$. Hecke's theta construction
places $\sum_{\mathfrak a}\psi(\mathfrak a)q^{N\mathfrak a}$ in
$S_2(\Gamma_0(s^2))$: the nebentypus is
$\chi_{-s}(n)(n/s)=1$. This exact special case is also written
in \citet[\S3, pp.~4--5]{CrDarmon2009}. Grouping conjugate
generators $\alpha=(u+v\sqrt{-s})/2$ gives the integral formula
for the coefficients:
\[
 b_n=\sum_{a^2=n,\ a>0}\left(\frac as\right)a
  +\sum_{\substack{u,v>0,\ u\equiv v\ (2)\\u^2+sv^2=4n}}
       \left(\frac{(s+1)u/2}{s}\right)u.
\]
For the four curves in the statement, exact point counts at good primes, the bad-prime
Euler factors, and multiplicativity give elliptic-curve coefficients
equal to $b_n$ for $1\leq n\leq10000$.

\par The difference of the two weight-two forms has integral
Fourier coefficients. Apply Sturm's congruence theorem
\cite[Theorem~9.18, p.~171]{CgzStein2007} to this difference
modulo each rational prime. Agreement through
$\lfloor kI/12\rfloor$, where
\[
 I=[\mathrm{SL}_2(\mathbb Z):\Gamma_0(N)]
   =N\prod_{p\mid N}(1+1/p),
\]
forces every coefficient to vanish modulo every prime;
the index is also stated in
\cite[Corollary~9.20, p.~174]{CgzStein2007}.
An integer divisible by every prime is zero, so the forms
are equal over $\mathbb C$.
For $k=2,N=s^2$
the four bounds are
\[
 \left\lfloor\frac{s^2+s}{6}\right\rfloor
     =22,315,759,4455,
\]
all below $10000$. Applying the exact quadratic twist character
then identifies every coefficient in both the target and auxiliary
streams, including coefficients past the tested range. 

\end{proof}

\begin{proposition}[Central derivatives with explicit tails]\hypertarget{tag-0011}{}\label{ari:C.6.R2}
\StatementRecord{ari:C.6.R2}{proved}{0011}{\SeriesRef{ari:eq:Cr-tails}; \citet[\S3, pp.~4--5]{CrDarmon2009}; \citet[Theorem~10.1, p.~47]{CrRubin1999}}

The first derivative ratios are enclosed by the following values
with absolute error less than $10^{-8}$. The zeros in rows one
and eight are exact.
\[
\begin{array}{r|r}
i&L'(E_i,1)/\Omega_i\\\hline
1&0\\
2&5.13439843\\
3&2.13412085\\
4&31.81013979\\
5&1.86946271\\
6&3.42887338\\
7&0.36258335\\
8&0\\
9&9.62927392\\
10&16.04238858\\
11&1247.32180615\\
\end{array}
\]
For the same full real periods,
\[
 L''(E_1,1)/\Omega_1=28.71408723\ldots,\qquad
 L''(E_8,1)/\Omega_8=78.03222042\ldots,
\]
again with absolute error less than $10^{-8}$.
The leading Taylor coefficients in these two rows are half
the second derivatives.
\end{proposition}


\begin{proof}[Derivative formulas and error control]

\begin{enumerate}[label=\arabic*.,leftmargin=*]
\item Let $N$ be the conductor, $c=2\pi/\sqrt N$, and
$L(E,s)=\sum_{n\geq1}a_nn^{-s}$. Put
\[
 E_1(x)=\int_x^\infty\frac{e^{-u}}u\,du,
 \qquad G(x)=\int_x^\infty e^{-u}\frac{\log(u/x)}u\,du.
\]
Splitting the Mellin integral at $1/\sqrt N$ and using the functional
equation gives, for root number $-1$,
\begin{equation}\hypertarget{tag-0020}{}\label{ari:eq:Cr-first}
 L'(E,1)=2\sum_{n\geq1}\frac{a_n}{n}E_1(cn).
\end{equation}
For root number $+1$ and $L(E,1)=0$, it gives
\begin{equation}\hypertarget{tag-0022}{}\label{ari:eq:Cr-second}
 L''(E,1)=4\sum_{n\geq1}\frac{a_n}{n}G(cn).
\end{equation}
One can verify the factors directly by writing the completed Mellin
integral as an integral over $[1,\infty)$ of
$(u^{s-1}+wu^{1-s})f(iu/\sqrt N)$: its first derivative has
the factor $2\log u$ when $w=-1$, and its second derivative
the factor $2\log^2u$ when $w=1$. Integration by parts gives
the two formulas above. The vanishing of the lower derivatives
removes all derivatives of the gamma and conductor factors.

\item Hasse's bound and multiplicativity imply
$|a_n|\leq d(n)\sqrt n\leq2n$. Also
$E_1(x)\leq e^{-x}/x$ and $G(x)\leq e^{-x}/x^2$.
The latter follows on writing $u=x+v$ and using
$\log(1+v/x)\leq v/x$. Consequently the absolute tails after
$n=M$ are at most
\begin{equation}\hypertarget{tag-0023}{}\label{ari:eq:Cr-tails}
 \frac{4e^{-c(M+1)}}{c(M+1)(1-e^{-c})},\qquad
 \frac{8e^{-c(M+1)}}{c^2(M+1)^2(1-e^{-c})},
\end{equation}
respectively. For the first ten curves take
$M=\lceil180/c\rceil$ and compute the coefficients exactly.
For the second-derivative kernel use the convergent series
\[
 G(x)=\frac12(\log x+\gamma)^2+\frac{\pi^2}{12}
       +\sum_{k\geq1}\frac{(-x)^k}{k^2k!}.
\]
After $k>x$ its terms decrease in absolute value. The alternating
remainder was cut below $10^{-140}$. At most 740 terms were needed;
the largest intermediate term was less than $1.4\cdot10^{72}$.
Evaluating these finite sums with the rounding allowance below
controls this cancellation. The largest infinite-tail bound in
\eqref{ari:eq:Cr-tails} is $3.46\cdot10^{-77}$.

\item For $i=11$, generate the coefficients by the CM theta series.
The curve is the quadratic twist, by
$d=185801=7\cdot11\cdot19\cdot127$, of the conductor-$163^2$
curve
\[
 y^2+y=x^3-2174420x+1234136692.
\]
Write $D=163$, $\chi_D(a)=(a/D)$, and $\chi_d(n)=(d/n)$.
The exact coefficient $a_n$ is the sum of
\[
 \chi_d(n)\chi_D(a)a \quad(n=a^2,\ a>0)
\]
and
\[
 \chi_d(n)\chi_D((D+1)u/2)u
 \quad\left(n=\frac{u^2+Dv^2}{4},\ u,v>0,
                           \ u\equiv v\pmod2\right).
\]
This is the Hecke-character theta series for
$\mathcal O_{\Q(\sqrt{-D})}$, grouping conjugate algebraic integers
and then applying the quadratic twist; see
\citet[\S3, pp.~4--5]{CrDarmon2009}. The finite modular-form comparison in \cref{I:CM-coefficient-series}
proves the global identification of this coefficient formula.

\item To accelerate the sum, in a block $n=n_0+r$, $|r|\leq R$, form the
exact integer moments $M_k=\sum a_nr^k$. For $x=cn_0$ write
\[
 e_k=e^{-x}(-x)^k/k!,\quad g_0=e^{-x},\quad
 g_k=e_k-g_{k-1},\quad
 \phi_0=E_1(x),\quad\phi_k=-\phi_{k-1}-g_{k-1}/k.
\]
Then the block contribution to $L'/2$ is
$\sum_k\phi_k M_k/n_0^{k+1}$. This follows by Taylor expansion
of $E_1(x(1+s))/(1+s)$. On $|s|=1/2$ its modulus is at most
$2E_1(x/2)$, which is less than $14$ throughout the block range.
We take enough terms that $(2R/n_0)^{k+1}<10^{-85}$.
The Cauchy remainder, summed using the bound
$\sum|a_n|<2\cdot10^{12}$ and $n_0>7700$, is less than
$10^{-75}$. The values of $E_1$ at consecutive block centres
are continued by the analogous Taylor expansion; after at most
1024 such continuations, enclose $E_1$ directly again.
The continuation bound is $10^{-110}$ per normalized expansion.
All moments are exact integers.

\item Taking $R=96$ and cutoff $cM\geq190$ uses
$915818441$ coefficients and gives an infinite tail less than
$3.10\cdot10^{-78}$. In the recurrences,
$|e_k|\leq1$, $|g_k|\leq1$, and $|\phi_k|<20$.
There are fewer than $5\cdot10^6$ blocks and at most 1024
consecutive continuations. Bounding each normalized moment by
$\sum|a_n|(R/n_0)^k$ controls accumulated rounding. The elementary evaluations are carried
out to make that contribution less than $10^{-90}$; adding the
Taylor remainder gives the derivative allowance used below.

\item Enclose the real periods from the exact minimal invariants.  For rows $1,2,3,4$
put respectively $T=\sqrt2,\sqrt3,2,\sqrt7$ and $\tau=iT$;
for rows $5,\ldots,11$ put
$T=\sqrt{7}/2,\sqrt{11}/2,\sqrt{19}/2,\sqrt{27}/2,
\sqrt{43}/2,\sqrt{67}/2,\sqrt{163}/2$ and $\tau=1/2+iT$.
Let $r=e^{-2\pi T}$ and $q=e^{2\pi i\tau}$.  The identity
$E_4=1+240\sum_{n\geq1}\sigma_3(n)q^n$ and
$\sigma_3(n)<2n^3$ give, after $M=100$ terms, the tail bound
\[
480r^{A}\left\{
\frac{A^3}{1-r}+\frac{3A^2r}{(1-r)^2}
+\frac{3Ar(1+r)}{(1-r)^3}
+\frac{r(1+4r+r^2)}{(1-r)^4}\right\},\qquad A=101.
\]
This follows by expanding $(A+j)^3$ and summing each geometric
moment.  All eleven parameters satisfy
$1-1008r/(1-32r)>0$; the bound
$\sigma_5(n)<2n^5$ and $n^5\leq32^{n-1}$ therefore proves
$E_6(\tau)>0$.  The exact Weierstrass invariant relations yield
$\kappa=2\pi(E_4/c_4)^{1/4}$.  For $c_6>0$ the real lattice
generator has length $\kappa$.  For $c_6<0$ the lattice is rotated
by $i$: its real generator has length $T\kappa$ in the rectangular
case and $2T\kappa$ in the other case.  Indeed the real vectors in
$i(\Z+iT\Z)$ are $T\Z$, while those in
$i(\Z+(1/2+iT)\Z)$ are $2T\Z$.  Multiplication by the number
$n_{\R}=2$ for $\Delta>0$ and $n_{\R}=1$ for $\Delta<0$
gives the full real period used in the statement.

Evaluate these operations with outward intervals.
For the first ten derivatives enlarge the explicit tail by one
percent and add
$10^{-100}$ for arithmetic and kernel-series errors.  This last
bound is conservative: fewer than $5\cdot10^5$ coefficients and
740 series terms occur, the intermediate terms are below
$1.4\cdot10^{72}$, and elementary precision chosen to bound accumulated rounding,
together with the $10^{-140}$ kernel remainders, stays below
this allowance.
For the last derivative add $1.1\cdot10^{-75}$, covering the
block Taylor and rounding bounds above. Division by the strictly
positive period intervals gives the stated enclosures.

\item The exact root numbers are $+1$ for $i=1,8$ and $-1$ otherwise.
All points are nontorsion by \cref{ari:C.5.1}. For a CM curve,
the Coates--Wiles theorem, in the form
\citet[Theorem~10.1, p.~47]{CrRubin1999}, says that a nonzero central
Hecke $L$-value forces the Mordell--Weil group over the CM field
to be finite. \citet[Proposition~5.3, p.~15]{CrRubin1999} supplies, for a nonmaximal CM order, an isogenous
curve with maximal order: the quotient is by the Galois-stable subgroup attached to the conductor ideal. The isogeny induces an isomorphism on Mordell--Weil groups tensored with $\Q$ and on their rational Tate modules, preserving ranks and local $L$-factors.
Since $L(E/K,s)=L(E/\Q,s)^2$ for the CM field $K$, nontorsion
therefore forces $L(E,1)=0$. This proves the needed central
vanishing also in the two sign-$+1$ cases. The functional equation
then forces $L'(E,1)=0$ for those two curves. The positive first
or second derivatives just evaluated prove exactly the asserted
analytic ranks.
\end{enumerate}

\end{proof}

\begin{proposition}[Rational leading factors and finite Tate--Shafarevich groups]\hypertarget{tag-0012}{}\label{ari:C.6.R3}
\StatementRecord{ari:C.6.R3}{proved}{0012}{\citep[Chapter~I, Theorem~7.3, preprint p.~12]{CrGrossZagier1986}; \citet[Theorem~1.3, pp.~235--236]{CrGross1991}; \citet[Theorem~1.1, p.~2 of arXiv v1]{CrJetchev2008}; \citet[Introduction, Theorem~(ii), pp.~543--544]{BumpFriedbergHoffstein1990}}

For each of the nine rank-one curves there is a positive rational
number $q_i$ such that
\begin{equation}\hypertarget{tag-0021}{}\label{ari:eq:Cr-rational-factor}
 \frac{L'(E_i,1)}{\Omega_i}=q_i\widehat h(P_i).
\end{equation}
Their Tate--Shafarevich groups over $\Q$ are finite. Consequently
the corresponding $D_{\rm part}$ is an explicit logarithm minus
$(6/(\pi q_i))L'(E_i,1)/\Omega_i$. For the two rank-two curves, $L'(E_i,1)=0$ while
$\widehat h(P_i)>0$.
\end{proposition}


\begin{proof}

\begin{enumerate}[label=\arabic*.,leftmargin=*]
\item Gross--Zagier's theorem applies because each curve is modular and
$L(E_i,1)=0$. CM modularity follows from its Hecke-character
theta series. Their Theorem~I.7.3 \citep[Chapter~I, Theorem~7.3, preprint p.~12]{CrGrossZagier1986} gives a rational point $R$ and
a nonzero rational factor relating $L'(E,1)$ to its canonical
height times the real period. The differentials in the two conventions are rational multiples:
write $\omega'=a\omega_E$ with $a\in\Q^*$.  A primitive real
period for $\omega'$ is $|a|\Omega_i/n_{\R}$, where
$n_{\R}\in\{1,2\}$ is the number of real components.
Thus the change to the full real period multiplies the rational
factor by $|a|/n_{\R}$.
As $L'(E_i,1)>0$, $R$ has infinite order. Primitivity gives
$R=mP_i+T$, $m\ne0$, $T$ torsion; its height is
$m^2\widehat h(P_i)$. This proves
\eqref{ari:eq:Cr-rational-factor}, including positivity.

\item The rank-one Gross--Zagier--Kolyvagin implication asserts rank one
and finite Tate--Shafarevich group when the central zero is simple.
One can use Kolyvagin's Heegner-point theorem as stated in
\citet[Theorem~1.3, pp.~235--236]{CrGross1991}, together with the nonzero
Gross--Zagier Heegner height; a modern primary research statement
of the rank conclusion is \citet[Theorem~1.1, p.~2 of arXiv v1]{CrJetchev2008}.
For each of the nine rows, the associated weight-two newform has
sign $-1$ by \cref{ari:C.6.R2}. Apply the nonvanishing theorem
\citep[Introduction, Theorem~(ii), pp.~543--544]{BumpFriedbergHoffstein1990}
with $S=\{2\}\cup\{p:p\mid N_i\}$, where $N_i$ is the conductor.
It gives a negative fundamental discriminant $D$ such that every
prime in $S$ splits in $K=\Q(\sqrt D)$ and $L(E_i^D,1)\ne0$.
Since $2$ splits, $D\equiv1\pmod8$, so $D\ne-3,-4$ and
$\mathcal O_K^\times=\{\pm1\}$, as required in Gross's setup.
The Gross--Zagier formula
\citep[Chapter~V, Theorem~2.1, preprint p.~127; proof of I.7.3, p.~129]{CrGrossZagier1986}
then gives a nonzero height for the Heegner trace, because
$L'(E_i/K,1)=L'(E_i,1)L(E_i^D,1)\ne0$.
Thus Gross's Theorem~1.3 gives finite $\operatorname{Sha}(E_i/K)$.
The restriction map
\[
\operatorname{Sha}(E_i/\Q)\longrightarrow
\operatorname{Sha}(E_i/K)
\]
has kernel contained in
$H^1(\operatorname{Gal}(K/\Q),E_i(K))$ by restriction--inflation.
This last group is finite: the Mordell--Weil group is finitely
generated, and its degree-one cohomology for this order-two group
is finitely generated and killed by two.  The finite kernel and
finite image prove the claimed finiteness over $\Q$.
For the two other rows the established equations are
$L'(E_i,1)=0$ and $\widehat h_O(P_i)>0$; therefore an equation
$L'(E_i,1)/\Omega_i=q\widehat h_O(P_i)$ with $q>0$ would contradict
these exact values.  Their rank-two regulator, for any free basis
$P,Q$, is the determinant
$\widehat h_O(P)\widehat h_O(Q)-\langle P,Q\rangle^2$.
\end{enumerate}

\end{proof}

\begin{remark}[The analytic leading quotients]
\hypertarget{tag-IA04}{}\label{I:analytic-leading-quotients}
\StatementRecord{I:analytic-leading-quotients}{proved}{IA04}{}

For the nine rank-one rows, define
\[
 Q_i=\frac{L'(E_i,1)/\Omega_i}
 {2\widehat h_O(P_i)\prod_pc_p/\#E_i(\mathbb Q)_{\rm tors}^{\,2}}.
\]
The exact factors in \cref{ari:C.6.R4,ari:C.6.R5} give $Q_i=1$
for $i=2,3,4,5,6,7,9,10$ and $Q_{11}=9$.
In rows $3,5,7$, \cref{ari:C.6.R4} also computes
$\#\operatorname{Sha}(E_i/\mathbb Q)=1$.
The remaining Tate--Shafarevich orders are discussed in
\cref{ari:C.7.5}.
\end{remark}

\begin{cited}[Three exact leading factors]\hypertarget{tag-0013}{}\label{ari:C.6.R4}
\StatementRecord{ari:C.6.R4}{proved}{0013}{\SeriesRef{ari:C.6.R1}, \SeriesRef{ari:C.6.R2}; \citet[Definition~1.1, Theorem~1.2, and Table~9]{miller2011}}

For $i=3,5,7$, namely $j=287496,-3375,-884736$, one has
\[
 \operatorname{Sha}(E_i/\Q)=0,\qquad
 \frac{L'(E_i,1)}{\Omega_i}=4\widehat h_O(P_i).
\]
Thus the factors $q_3=q_5=q_7=4$ are determined exactly by the integer bounds in the cited application. For the smallest of these examples,
$E:y^2+y=x^3-38x+90$ and $P=(0,9)$, this gives the exact
reduction
\[
 N(P)=\exp\left(-\frac{L'(E,1)}{4\Omega_E}\right).
\]
\end{cited}

\begin{proof}
Miller's theorem
\cite[Definition~1.1 and Theorem~1.2, arXiv v3, p.~2;
Table~9, p.~23]{miller2011}
applies to analytic rank at most one and conductor less than $5000$,
outside the listed exceptional curve--prime pairs.
The present conductors are $3136,3136,361$; none occurs in
that exception table, and their analytic ranks are one by
\cref{ari:C.6.R2}. The theorem applies to every curve in each
isogeny class.

For these three rows,
$2\prod_pc_p/\#E(\mathbb Q)_{\rm tors}^2=4$, and
\cref{ari:C.6.R1} proves primitivity. Therefore
$L'(E_i,1)/(4\Omega_i\widehat h_O(P_i))$ is the positive integer
$\#\operatorname{Sha}(E_i/\mathbb Q)$.
To isolate it, complete the square and write the duplication map
as relatively prime integral quartics $[F(X,Z):G(X,Z)]$.
The extended Euclidean algorithm gives integer homogeneous cubics
$A_0,B_0,A_1,B_1$ and $D\ne0$ with
\[
 A_0F+B_0G=DZ^7,\qquad A_1F+B_1G=DX^7.
\]
Let $C$ dominate the coefficient one-norms of $F,G$ and the
sums $\|A_j\|_1+\|B_j\|_1$. For coprime $a,b$, put
$H=\max(|a|,|b|)$ and $d=\gcd(F(a,b),G(a,b))$.
The two identities give $d\mid D$. Using the identity whose
right side contains the coordinate of size $H$ also gives
$\max(|F(a,b)|,|G(a,b)|)\geq |D|H^4/C$.
After division by $d$, the projective size is at least $H^4/C$
and at most $CH^4$. Therefore
\[
 |h(x(2Q))-4h(x(Q))|\le\log C.
\]
Telescoping yields
\[
 \left|\widehat h_O(P)-
       \frac{h(x(2^6P))}{2\cdot4^6}\right|
       \le\frac{\log C}{6\cdot4^6}.
\]
The exact constants for rows $3,5,7$ are
$26610617537591808$, $188716953600$, and $285931133$.
Six rational doublings and outward logarithm enclosures, together
with the derivative tails of \cref{ari:C.6.R2}, give respective
quotient intervals $1\pm0.003$, $1\pm0.003$, and $1\pm0.009$.
Each contains only the integer one.
Finally $A_t=1$ in row seven by \cref{ari:C.6.3}, giving the
displayed formula for $N(P)$.
\end{proof}


\begin{remark}[Ordinary CM leading-term theorems]
\hypertarget{tag-IA05}{}\label{I:ordinary-CM-scope}
\StatementRecord{I:ordinary-CM-scope}{proved}{IA05}{}

The applicable ordinary CM results concern a prime $p$ at which
$E$ has good ordinary reduction, equivalently $p$ splits in the CM
field. One formulation of Rubin's rank-one formula
is \citet[Theorem~1, p.~262]{CrBDP2012}: for a curve over $\Q$ with CM
by the maximal order, an ordinary prime prime to the conductor,
and a simple central zero, an appropriate Katz $p$-adic $L$-value
is, up to a nonzero algebraic factor, the square of a formal
logarithm of a nontorsion rational point divided by a $p$-adic
period. The analytic-rank-one input holds by \cref{ari:C.6.R2}.
The formal logarithm, Katz value and period in this formula
are the $p$-adic quantities specified in that theorem.

For the $p$-part of the real leading-term formula, a modern precise
version is \citet[Theorem~A, p.~2 of the author version]{CrCastella2025}. In addition to analytic
rank one it requires full CM, abelian torsion field over the CM
field, $p\nmid6h_K$, splitting of $p$, conductor prime to $p$,
and that the different of $K$ divide the Hecke conductor exactly.
For the maximal odd-discriminant cases here, the latter follows
from the exponent two of the ramified prime in the rational
conductor. The nonmaximal cases of discriminant $-12$ and $-28$
can be passed through a $2$-isogeny at $p>3$. The discriminant
$-16$ case requires care: the direct application of this version
has a different-conductor hypothesis at $2$ that fails, so this formulation does not apply directly to that case.

For the last curve, $3$ is inert in $\mathbb Q(\sqrt{-163})$
and is excluded by the stated $p>3$ hypotheses.
The exact analytic quotient $Q_{11}=9$ is proved in
\cref{ari:C.6.R5}; the remaining group-order question is
recorded in \cref{ari:C.7.5}.
\end{remark}

\subsection{Exact leading factors from a genus-character denominator bound}

\begin{proposition}[The remaining rational leading factors]\hypertarget{tag-0014}{}\label{ari:C.6.R5}
\StatementRecord{ari:C.6.R5}{proved}{0014}{\SeriesRef{ari:C.6.R4}, \SeriesRef{ari:C.6.R2}, \SeriesRef{ari:C.6.R3}, \SeriesRef{ari:C.6.R1}, \SeriesRef{ari:eq:Cgz-product}, \SeriesRef{ari:eq:Cr-first}, \SeriesRef{ari:C.6.5}, \SeriesRef{ari:eq:Cgz-central-integrality}, \SeriesRef{ari:eq:Cgz-denominator}; \citet[Appendix B, p.~102, unnumbered $26569\mathrm a1/26569\mathrm a2$
block]{CgzBelabas2019}; \citet[Theorem 1.1, pp.~2524--2525]{CgzCai2014}; \citet[\S7.1.1, pp.~218--219, Theorem 7.1.2]{CgzYZZ2011}; \citet[\S3, pp.~4--5]{CrDarmon2009}; \citet[Theorem 9.18]{CgzStein2007}}

Use the full real period and the absolute height of
\cref{I:def:theta-height}.

For the six rank-one rows $i=2,4,6,9,10,11$, one has
\[
 \frac{L'(E_i,1)}{\Omega_i}
   =q_i\widehat h_O(P_i),\qquad
 (q_2,q_4,q_6,q_9,q_{10},q_{11})=(8,32,8,16,16,576).
\]
Together with \ref{ari:C.6.R4}, this determines every rank-one factor:
\[
\begin{array}{c|rrrrrrrrr}
 i&2&3&4&5&6&7&9&10&11\\\hline
 q_i&8&4&32&4&8&4&16&16&576.
\end{array}
\]
The certified auxiliary ratios and denominator bounds used to identify
the six additional factors are as follows. Here $D$ is the auxiliary
negative fundamental discriminant,
$r_i=L(E_i^D,1)/\Omega_{i,D}$, and the denominator of $q_i$ is at most
$B_i$.
\[
\begin{array}{c|r|r|r|r|r}
 i&D&r_i&B_i&q_i&|a_i|\\\hline
 2&-23&64&256&8&32\\
 4&-31&64&256&32&64\\
 6&-19&8&32&8&8\\
 9&-19&16&64&16&16\\
 10&-3&400&1600&16&80\\
 11&-3&64&256&576&192
\end{array}
\]
The last column is the absolute genus-trace integer in the proof. These
identities determine the analytic BSD quotients exactly: they are one
for $i=2,4,6,9,10$ and nine for $i=11$.
\end{proposition}



\begin{proof}
\begin{enumerate}[label=\arabic*.,leftmargin=*]
\item \textbf{A parametrization under quadratic twisting.}
Let $E_0/\Q$ have conductor $N_0$, and suppose a parametrization
$\phi:X_0(N_0)\to E_0$ sends $\infty$ to $O$ and satisfies
\[
 \phi^*\omega_0=c_0f_0(q)\,dq/q,
\]
where $\omega_0$ is a minimal differential and $f_0$ is the normalized
newform. Let $b$ be an odd fundamental discriminant coprime to $N_0$,
and suppose that the minimal twist has conductor $N_0b^2$ and exact
invariants
\[
 c_4(E_0^b)=b^2c_4(E_0),\qquad
 c_6(E_0^b)=b^3c_6(E_0).
\]
Then $E_0^b$ has a parametrization from $X_0(N_0b^2)$ with the same
Manin factor $c_0$.

\item Indeed, put $m=|b|$ and, using addition on $E_0$, define
\[
 F(z)=\sum_{a\bmod m}\chi_b(a)\phi(z+a/m).
\]
For $\gamma=\left(\begin{smallmatrix}A&B\\C&D_0\end{smallmatrix}\right)$
in $\Gamma_0(N_0m^2)$ choose $a'\equiv aD_0^2\pmod m$.
Direct multiplication gives
$T_{a/m}\gamma T_{a'/m}^{-1}\in\Gamma_0(N_0)$, where $T_t$
denotes translation by $t$. Because $\chi_b(D_0^2)=1$, the sum is
invariant. Each summand is an algebraic morphism on a common
congruence cover; invariance therefore descends the sum to
$X_0(N_0m^2)$ over $\Q(\zeta_m)$. Its expansion at $\infty$
gives $F^{\sigma_u}=\chi_b(u)F$ for
$\sigma_u(\zeta_m)=\zeta_m^u$. Quadratic descent gives the required
map to $E_0^b$ over $\Q$, with $\infty$ mapping to $O$.
On differentials, the Gauss-sum identity gives
\[
 F^*\omega_0=c_0\tau(\chi_b)f_b(q)\,dq/q,
 \qquad \tau(\chi_b)^2=b.
\]
For the twist isomorphism $\iota:E_0^b\to E_0$ the exact invariant
identities imply
$\iota^*\omega_0=\sqrt b\,\omega_b$ up to sign. Cancellation proves the assertion.

\item If $c_0=1$, let $t_b=\#E_0^b(\Q)_{\rm tors}$ and
$n_b=\#\pi_0(E_0^b(\R))$. Manin--Drinfeld \citep[Section~1, unnumbered theorem, p.~83 of the Russian original]{Drinfeld1973} says that the image of
the rational cuspidal divisor $[0]-[\infty]$ is torsion. Its elliptic
logarithm, along the imaginary axis, is $\pm L(E_0^b,1)$.
Consequently $t_bL(E_0^b,1)$ is an integral multiple of the positive
real lattice period $\Omega_b^+$, and
\begin{equation}\hypertarget{tag-001X}{}\label{ari:eq:Cgz-central-integrality}
 t_bn_b\,\frac{L(E_0^b,1)}{\Omega_b}\in\Z,
 \qquad \Omega_b=n_b\Omega_b^+.
\end{equation}

\item \textbf{The six small base curves and Heegner hypotheses.}
Use the following data; $E_i=E_0^d$ and $K=\Q(\sqrt D)$.
\[
\begin{array}{c|c|r|r|r}
 i&E_0&N_0&d&D\\\hline
 2&576\mathrm e2&576&-11&-23\\
 4&3136\mathrm r2&3136&57&-31\\
 6&121\mathrm b1&121&21&-19\\
 9&1849\mathrm a1&1849&21&-19\\
 10&4489\mathrm a1&4489&217&-3\\
 11&26569\mathrm a1&26569&185801&-3
\end{array}
\]
All six base Manin factors are one. For the first five bases this follows from the exact
modular-symbol algorithm. For $26569\mathrm a1$, the value one and modular degree
$59368$ are explicitly tabulated by
\citet[Appendix B, p.~102, unnumbered $26569\mathrm a1/26569\mathrm a2$
block]{CgzBelabas2019}; their Algorithm 1 describes the computation.
Exact rational transformations check both $c_4,c_6$ scaling
identities for $b=d,dD$. Tate's algorithm and rational torsion
computations give both conductor identities and torsion orders;
quadratic residue symbols check the splitting conditions.

\item In every row $d,D,dD$ are fundamental discriminants,
$\gcd(d,D)=\gcd(dD,N_0)=1$, and every prime dividing $N_0$ splits
in $K$. The genus character for $K(\sqrt d)/K$ factors through
$\operatorname{Pic}(\mathcal O_{|d|})$ and has conductor exactly
$|d|$. On ideals prime to $d$ it is
\[
 \chi(\mathfrak a)=\chi_d(N\mathfrak a).
\]
For a unit of the order locally at a prime $p\mid d$, its norm is
a square modulo $p$, so this character is trivial on the order's
units. It is nontrivial on the full local unit group: at a split
prime the norm is multiplication on two copies of $\Q_p^\times$,
and at an inert prime the residue-field norm onto $\mathbb F_p^\times$
is surjective. Since $p$ is odd and unramified in $K$, this also
proves that the conductor contains every $p\mid d$ exactly once.
Thus the primitive ring-class-character hypotheses of
\citet[Theorem 1.1, pp.~2524--2525]{CgzCai2014} hold. Their
$\mu(N_0,D)=0$ and
$u=[\mathcal O_{|d|}^{\times}:\Z^{\times}]=1$; the latter remains
true for $D=-3$ because $|d|>1$.

\item Let $P_{\mathfrak a}$ be the corresponding Heegner points and put
\[
 Y=\sum_{[\mathfrak a]}\chi(\mathfrak a)\phi(P_{\mathfrak a}).
\]
The sum is taken in $E_0(\overline{\Q})$ with coefficients
$\chi(\mathfrak a)\in\{\pm1\}$, so $Y\in E_0(\overline{\Q})$. It is fixed
by $\ker\chi$ and anti-invariant under the other automorphism of
$K(\sqrt d)/K$. Twisting therefore gives $Y_d\in E_i(K)$.
Artin factorization gives
\[
 L(s,E_0,\chi)=L(E_i,s)L(E_i^D,s).
\]
The intervals below prove $L(E_i^D,1)>0$, hence the auxiliary
curve has rank zero. The rank-zero implication used here is the
same unconditional modular/CM rank implication used in \cref{ari:C.6.R2}, \cref{ari:C.6.R3}.
Thus $E_i(K)$ has rank one. \ref{ari:C.6.R1} proves that $P_i$ is primitive
over $\Q$. Since $Y_d-\overline{Y_d}$ is torsion and
$Y_d+\overline{Y_d}$ is rational, there is an integer $a_i$ with
\[
 2Y_d=a_iP_i\pmod{E_i(K)_{\rm tors}},\qquad
 \widehat h_O(Y_d)=\frac{a_i^2}{4}\widehat h_O(P_i).
\]

\item \textbf{The exact height and period normalization.}
Write $V_0=\operatorname{area}(\C/\Lambda_0)$ for the period area
of the minimal differential. Integrating $\phi^*\omega_0$ gives
\[
 4\pi^2(f_0,f_0)_{\Gamma_0(N_0)}
   =\frac{\deg\phi}{c_0^2}V_0.
\]
The degree in the denominator of CST's Theorem 1.1 cancels, giving
\begin{equation}\hypertarget{tag-001Z}{}\label{ari:eq:Cgz-product}
 L'(E_i,1)L(E_i^D,1)
  =\frac{2V_0}{c_0^2|d|\sqrt{|D|}}\,\widehat h_K(Y).
\end{equation}
Here $\widehat h_K$ is the diagonal N\'eron--Tate pairing relative
to $K$. This convention is explicit in
\citet[\S7.1.1, pp.~218--219, Theorem 7.1.2]{CgzYZZ2011}:
their underlying height is normalized by $[L:K]^{-1}$ over a
field of definition $L$, and their pairing is
$\widehat h_{\mathcal L}(x+y)-\widehat h_{\mathcal L}(x)
-\widehat h_{\mathcal L}(y)$.

\item One can also see the factor geometrically. For
$\mathcal L=\mathcal O(O)$, pull the normalized Poincar\'e bundle
back under $1\times\lambda_{\mathcal L}$. The result is
$m^*\mathcal L\otimes p_1^*\mathcal L^{-1}\otimes
p_2^*\mathcal L^{-1}$. On the diagonal its canonical height is
that of $[2]^*\mathcal L\otimes\mathcal L^{-2}\simeq
\mathcal L^2$, namely $2\widehat h_O$ in the absolute convention.
Changing from the absolute convention to relative heights over
$K$ multiplies by $[K:\Q]=2$. Therefore
\[
 \widehat h_K(Y)=4\widehat h_O(Y_d)
                =a_i^2\widehat h_O(P_i).
\]
This is the factor from the Poincar\'e pairing and the base field.

\item A real-adapted basis of $\Lambda_0$ gives
\[
 \Omega_i\Omega_{i,D}=
 \begin{cases}
 2V_0/(|d|\sqrt{|D|}),&\Delta(E_0)<0,\\
 4V_0/(|d|\sqrt{|D|}),&\Delta(E_0)>0.
 \end{cases}
\]
Indeed, a triangular real lattice has one real component and its
purely imaginary primitive period has twice the height of its
nonreal basis vector. A rectangular lattice has two real
components. Twisting by the two opposite signs interchanges the
real and imaginary directions, and the exact minimal differential
scalings in step 2 supply $1/(|d|\sqrt{|D|})$.
Substituting into \eqref{ari:eq:Cgz-product} gives
\begin{equation}\hypertarget{tag-001Y}{}\label{ari:eq:Cgz-denominator}
 q_ir_i=
 \begin{cases}
 a_i^2/c_0^2,&\Delta(E_0)<0,\\
 a_i^2/(2c_0^2),&\Delta(E_0)>0.
 \end{cases}
\end{equation}
In particular, if $r_i=A_i/B_i'$ is reduced and $c_0=1$, the
denominator of $q_i$ divides $2|A_i|$. We use the weaker upper
bound $B_i=4|A_i|$ in the interval certificate.

\item \textbf{Global identification of the CM coefficient series.}
Apply \cref{I:CM-coefficient-series} with $s=11,43,67,163$,
and then multiply the $n$th coefficient by the exact quadratic
twist character. This identifies both target and auxiliary series.

\item \textbf{Directed enclosures of both $L$-sums.}
For $c=2\pi/\sqrt N$, the root-number $+1$ formula is
\[
 L(E,1)=2\sum_{n\geq1}a_n\frac{e^{-cn}}n,
\]
while for root number $-1$ the first derivative is
\eqref{ari:eq:Cr-first}. Generate the coefficients in the smaller cases by point counts
and the Euler-product recurrence. The largest auxiliary stream uses
$s=163$, signed twist $-557403$, conductor
$8254937936042721$, and $M=867617471$; the target stream uses
$s=163,d=185801,M=289205824$. Both use the cutoff $cM\geq60$
up to rounding of the integer cutoff, with the displayed integer
$M$ used in every tail bound.

\item Evaluate the finite sums with outward interval arithmetic.
In a block centred at $b$ with radius $R=128$, set
$S_k=\sum a_n(n-b)^k\in\Z$. For $g(z)=e^{-cz}/z$, put
\[
 C_0=1,\quad C_k=C_{k-1}c/k,\qquad
 T_0=1/b,\quad T_k=(T_{k-1}+C_k)/b.
\]
The coefficient of $(z-b)^k$ is $(-1)^ke^{-cb}T_k$.
For $h(z)=E_1(cz)/z$, put
\[
 e_0=E_1(cb),\quad e_k=(-1)^ke^{-cb}T_{k-1}/k\ (k\geq1),
 \qquad h_0=e_0/b,\quad h_k=(e_k-h_{k-1})/b.
\]
The block sums are the finite sums of these interval coefficients
times the exact integers $S_k$.

\item On $|z-b|=b/2$, the two kernels have bounds
\[
 |g(z)|\leq\frac{2e^{-cb/2}}b=:U_g,
 \qquad
 |h(z)|\leq\frac{4e^{-cb/2}}{cb^2}=:U_h.
\]
For the second bound, use
$E_1(cz)=\int_1^\infty e^{-czt}\,dt/t$ when $\Re z>0$.
Thus truncation at degree $k$ contributes at most
\[
 U_g\frac{(2R/b)^{k+1}}{1-2R/b}\sum|a_n|
 \quad\hbox{or}\quad
 U_h\frac{(2R/b)^{k+1}}{1-2R/b}\sum|a_n|.
\]
Select the degree to make the geometric factor small, then evaluate
the displayed remainder bound with upward rounding.
For the derivative stream, the initial $E_1$ in each segment is
evaluated using correctly rounded $-\operatorname{Ei}(-x)$ and
monotonicity. Subsequent centres, spaced by $W=257$, use
\[
 E_1(c(b+W))=E_1(cb)-\int_0^Wg(b+x)\,dx.
\]
Integrating the same Taylor expansion gives the interval for this
advance. Its omitted part is bounded by
$WU_g(2W/b)^{k+1}/(1-2W/b)$ and added to both endpoints.
The first $10280$ terms are summed directly. Segment boundaries
are multiples of $257$, so no coefficient is omitted or repeated.

\item Hasse's bound and multiplicativity give $|a_n|\leq2n$.
The infinite tails, including the factor two, are consequently at
most
\[
 \frac{4e^{-c(M+1)}}{1-e^{-c}},\qquad
 \frac{4e^{-c(M+1)}}{c(M+1)(1-e^{-c})},
\]
respectively. Evaluate both tails with a lower bound for $c$ and upper bounds
for the remaining positive expressions, then add the block
remainders and elementary-evaluation errors. For row eleven the central-value tail is below
$5.07\cdot10^{-19}$, and the derivative tail is below
$2.82\cdot10^{-21}$. These modest precisions already suffice for
rational separation. 

\item \textbf{Periods, heights, and exact rational separation.}
Apply the real-period algorithm in \cref{ari:C.6.R2} separately
to each target and auxiliary minimal equation, using
$\tau=i\sqrt s$ for $i=2,4$ and
$\tau=(1+i\sqrt s)/2$ for $i=6,9,10,11$, where
$s=3,7,11,43,67,163$ respectively.
These parameters and their class-number-one identification are
proved in \cref{ari:C.6.5}. In both equations the exact
$c_4,c_6$ determine the same homothety rule
$\kappa=2\pi(E_4/c_4)^{1/4}$; the signs of $c_6$ and $\Delta$
select the real generator and number of components.
Thus the $E_4$ tail, positive $E_6$ bound and real-lattice
transport of \cref{ari:C.6.R2} give both full-period intervals.
Use the height intervals of \cref{ari:C.6.5}.

\item First, the intervals for
$r_i$ isolate a unique integer after multiplication by the exact
$t_bn_b$ in \eqref{ari:eq:Cgz-central-integrality}. This proves the
six auxiliary rational values from the displayed integrality and interval bounds. Next form the
interval $L'(E_i,1)/(\Omega_i\widehat h_O(P_i))$. Outward enlargement of the intervals gives the following rational
error bounds:
\[
\begin{array}{c|c|c}
 i&|r_i-r_i^{\rm listed}|\text{ is at most}
   &|q_i-q_i^{\rm listed}|\text{ is at most}\\\hline
 2&4\cdot10^{-22}&2\cdot10^{-53}\\
 4&4\cdot10^{-20}&3\cdot10^{-45}\\
 6&4\cdot10^{-22}&2\cdot10^{-53}\\
 9&5\cdot10^{-21}&4\cdot10^{-53}\\
 10&3\cdot10^{-20}&10^{-23}\\
 11&7\cdot10^{-15}&8\cdot10^{-19}
\end{array}
\]
Each last-column bound is strictly less than $1/B_i$. A rational
number with denominator at most $B_i$ differing from the listed
integer would differ by at least $1/B_i$. Equation
\eqref{ari:eq:Cgz-denominator} therefore makes every listed $q_i$
exact. The same equation gives the listed $|a_i|$ by integer
square-root checks. 
\end{enumerate}
\end{proof}


\subsection{The three smooth torsion fibres}

\begin{proposition}[The rational order-four fibre]\hypertarget{tag-0015}{}\label{ari:C.6.T1}
\StatementRecord{ari:C.6.T1}{proved}{0015}{the definitions and constructions stated in the proof}

At $t=32/5$ a global minimal model and the image of the section are
\[
 E:y^2=x^3-2x+1,\qquad P=(0,1),\qquad \Delta=80.
\]
The conductor is $40$, $E(\Q)=\Z/4\Z$ is generated by $P$,
and $\widehat h(P)=0$. The index of $\Z P$ in the quotient
modulo torsion is one. The complete bad-prime data are
\[
\begin{array}{c|c|c|c|c}
p&\operatorname{ord}_p\Delta&\text{type}&\text{component of }P&\lambda_p(P)\\\hline
2&4&\mathrm{III}&\text{unique nonzero class, order }2&\log2/12\\
5&1&\mathrm I_1&\text{identity}&\log5/12.
\end{array}
\]
All other finite local heights vanish. Consequently
\[
 \log N=\frac{\log10}{12},\qquad
 D_{\rm part}=\frac6\pi\log\frac{10^{1/12}}{2\sqrt3}.
\]
The central values satisfy
\[
 \frac{L(E,1)}{\Omega}=\frac18,\qquad
 \frac{L'(E,1)}{\Omega}
   =\frac{\gamma+\log(2\pi/\sqrt{40})}{8},
\]
where $\gamma$ is Euler's constant. In particular the analytic rank is zero.
\end{proposition}


\begin{proof}
\begin{enumerate}[label=\arabic*.,leftmargin=*]
\item From the displayed minimal coordinates to the family coordinates
of \cref{I:def:theta-height}, the map is
$X=(24/5)^2x$, $Y=(24/5)^3y$.
It sends $(0,1)$ to $P_{32/5}$ and pulls the family differential
back to $(5/24)\,dx/(2y)$. Tate's algorithm gives the
listed minimal model, conductor and local data. Since $P$ has
order four, $3P=-P$. The local division identity consequently
gives
\[
 \lambda_p(P)=\left(\frac{v_p(\Delta)}{12}
                         -\frac{v_p(\psi_3(P))}{8}\right)\log p.
\]
Here $\psi_3(P)=-4$, proving both local coefficients exactly.
Exact $2$-descent gives rank zero, and the torsion calculation
gives $\mathbb Z/4\mathbb Z$ generated by $P$.
The plus modular symbol from $0$ to $\infty$ is $-1/4$
in the primitive real-period normalization. There are two real
components, so $L(E,1)/\Omega=1/8$ for the full real period.
Differentiating the completed functional equation at its central
point gives the displayed formula for $L'(E,1)/\Omega$.
\end{enumerate}
\end{proof}

\begin{proposition}[The conjugate order-six fibres]\hypertarget{tag-0016}{}\label{ari:C.6.T2}
\StatementRecord{ari:C.6.T2}{proved}{0016}{\SeriesRef{ari:C.4.4}}

Put $s=\sqrt3$ and
$t_\pm=(3008\pm1728s)/11$. Over $K=\Q(s)$, a global minimal
model for the fibre at $t_-$ is
\[
 E_-:y^2=x^3-(18+10s)x+21+12s,\qquad P_-=(0,3+2s).
\]
The model and point for $t_+$ are its conjugates. The
discriminant of this model is $1032912+596352s$. Define
\[
 \mathfrak p_2=(2,1+s),\qquad\mathfrak p_3=(s),\qquad
 \mathfrak p_{11,-}=(11,s+5),\quad
 \mathfrak p_{11,+}=(11,s-5).
\]
The bad-prime data of $E_-$ are
\[
\begin{array}{c|c|c|c|c|c}
v&e_v&v(\Delta)&\text{type}&\text{component order of }P_-&\lambda_v(P_-)\\\hline
\mathfrak p_2&2&8&\mathrm{IV}^*&3&0\\
\mathfrak p_3&2&3&\mathrm{III}&2&0\\
\mathfrak p_{11,-}&1&1&\mathrm I_1&1&\log11/12.
\end{array}
\]
For $E_+$ replace $\mathfrak p_{11,-}$ by
$\mathfrak p_{11,+}$. The nonzero component at
$\mathfrak p_3$ is unique; the two nonzero classes at
$\mathfrak p_2$ are interchanged by $P\mapsto-P$ and the
class of $P$ is a generator. All omitted local heights vanish.
Both points have order six, and
\begin{align*}
 \log N(P_-)&=\frac1{12}\log(17+10\sqrt3),\\
 \log N(P_+)&=\frac1{12}\log(10\sqrt3-17),\\
 \sum_{\sigma:K\hookrightarrow\R}\log N(P_\sigma)
 &=\frac{\log11}{12}
 =\sum_{v<\infty}n_v\lambda_v(P)-2\widehat h(P).
\end{align*}
In particular $A_\mp=(17+10\sqrt3)^{1/12}$ and
$(10\sqrt3-17)^{1/12}$ respectively give
$D_{\rm part}(t_\mp)=(6/\pi)\log(A_\mp/(2\sqrt3))$.

The two $j$-invariants are $1728t_\pm\in K\setminus\mathbb Q$.
An elliptic curve over $\mathbb Q$ has rational $j$, so these
fibres admit no model over $\mathbb Q$.
Over $K$, both curves have rank one and torsion group $\mathbb Z/6\mathbb Z$. Thus the index of the torsion section in the
free quotient is infinite. For example $(1,-1-s)$ has infinite
order on $E_-$.
\end{proposition}


\begin{proof}
\begin{enumerate}[label=\arabic*.,leftmargin=*]
\item For $t_-$, put $u=48(-64+37s)/11$.
The map from the displayed minimal model to the family equation is
$X=u^2x$, $Y=u^3y$, with differential
$dX/(2Y)=u^{-1}dx/(2y)$.
It is obtained by multiplying the first global minimalization
scale by the unit $(2+s)^{-3}$. Hence it preserves all finite
minimal valuations. The conjugate scale gives $E_+$.
Tate's algorithm yields the displayed table. Its discriminant
ideal is
$\mathfrak p_2^8\mathfrak p_3^3\mathfrak p_{11,-}$ and its
conductor ideal is
$\mathfrak p_2^2\mathfrak p_3^2\mathfrak p_{11,-}$.

\item As $5P=-P$, the normalized local heights are computed exactly by
\[
 \frac{\lambda_v(P)}{\log p}
 =\frac1{e_v}\left(\frac{v(\Delta)}{12}
                         -\frac{v(\psi_5(P))}{24}\right).
\]
The valuations of $\psi_5(P)$ at the three listed primes are
$16,6,0$, respectively. This proves the table, including zero
contribution at the two additive primes. Here $f_v=1$ at all
three primes, so weighting by $n_v=e_vf_v$ gives precisely
$\log11/12$. The exact torsion values proved in Proposition~\ref{ari:C.4.4} have product
$(17+10\sqrt3)(10\sqrt3-17)=11$, proving the archimedean side.

\item For the rank, translate the rational two-torsion point to zero.
The resulting equation is
\[
 y^2=x^3-(9+6s)x^2+(45+26s)x.
\]
The $2$-isogeny descent gives
$[E(K):\widehat\phi E'(K)]=2$ and
$[E'(K):\phi E(K)]=4$. Every surviving covering has a
point; every excluded covering has a local
obstruction at $\mathfrak p_2$ or $\mathfrak p_3$.
The quotient formula $2^{r+2}=2\cdot4$ proves $r=1$.
The torsion calculation
gives $\Z/6\Z$, generated by $P$, independently of the rank
calculation.
\end{enumerate}
\end{proof}

\begin{remark}[The number-field period convention]
\hypertarget{tag-IA06}{}\label{I:number-field-period}
\StatementRecord{I:number-field-period}{proved}{IA06}{}

For $K=\mathbb Q(\sqrt3)$, put
\[
 \Omega_{\mathrm{prod}}=
 \prod_{\sigma:K\hookrightarrow\mathbb R}
       \int_{E_\sigma(\mathbb R)}|\omega_\sigma|,
 \qquad
 \Omega_K=\frac{\Omega_{\mathrm{prod}}}{\sqrt{12}}.
\]
The two conjugate equations give the same Hasse--Weil $L$-function
over $K$. Its analytic conductor is $12^2\cdot396=57024$,
its gamma parameters are $(0,0,1,1)$, and its sign is $-1$.
The exact conversion of derivative ratios is
$L'(E/K,1)/\Omega_K=\sqrt{12}\,L'(E/K,1)/\Omega_{\mathrm{prod}}$.
\end{remark}

\subsection{Integral and Gaussian descriptions}

\begin{proposition}[The Mahler-measure identities]\hypertarget{tag-001T}{}\label{ari:C.8.6}
\StatementRecord{ari:C.8.6}{proved}{001T}{\SeriesRef{ari:C.8.5}}

Let $\Lambda\subset\C$ be a lattice of area $A$, let $v$ be a zero
of $\wp_\Lambda$, and let $d\mu=A^{-1}dA_z$ be Haar probability
measure on $\C/\Lambda$. Put $U(z)=\log N(z)$ and
$L=U(v)$. Then
\begin{align}
 6L&=\int_{\C/\Lambda}\log\left|\frac{\wp'(z)^2}{\wp(z)^3}\right|d\mu(z),
 \hypertarget{tag-002E}{}\label{ari:eq:ca-mahler}\\
 N(v)^2&=|\Delta(\Lambda)|^{1/6}
      \exp\left(-\int_{\C/\Lambda}\log|\wp(z)|d\mu(z)\right).
 \hypertarget{tag-002H}{}\label{ari:eq:ca-wpmean}
\end{align}
If $C(u)=1-(g_2/4)u^2-(g_3/4)u^3$, the pushforward of $d\mu$
under $u=1/\wp(z)$ is
\begin{equation}\hypertarget{tag-002F}{}\label{ari:eq:ca-nu}
 d\nu(u)=\frac1{2A}\frac{dA_u}{|u|\,|C(u)|}.
\end{equation}
In particular its unnormalized integral is exactly $2A$, and
\[
 N(v)^6=4\exp\left(\int_{\C}\log|C(u)|d\nu(u)\right).
\]
Writing $e_1,e_2,e_3$ for the three branch values and
$p(x)=\int\log|x-\wp(z)|d\mu(z)$, the maximum theorem is equivalently
\[
 \int\log|C|d\nu\leq\log2,\qquad
 p(0)\geq\frac{p(e_1)+p(e_2)+p(e_3)}3-\frac{\log2}3.
\]
Equality in either inequality occurs exactly when $g_3=0$.
\end{proposition}


\begin{proof}
\begin{enumerate}[label=\arabic*.,leftmargin=*]
\item The mean of $U$ is zero by the theta product calculation in \cref{ari:C.1.1},
and near zero
$U(z)=\log|z|+\frac1{12}\log|\Delta|+o(1)$.
The divisor of $\wp$ is $(v)+(-v)-2(O)$, with multiplicity
two when $v=-v$. Subtract its logarithmic divisor from the
corresponding combination of $U$'s. The result is harmonic
on the compact torus and hence constant. Its expansion at
zero gives
\[
 \log|\wp(z)|=U(z-v)+U(z+v)-2U(z)
                  +\frac16\log|\Delta|-2L.
\]
Integration proves \eqref{ari:eq:ca-wpmean}.
Similarly the divisor of $\wp'$ is the sum of the three
nonzero half-periods minus $3(O)$. The multiplication-by-two
identity of \cref{ari:C.1.1} gives the product of $N$ at those half-periods
as $2$. Since $\wp'(z)\sim-2z^{-3}$, comparison at zero yields
\[
 \log|\wp'(z)|=\sum_{w\in E[2]\setminus\{O\}}U(z-w)-3U(z)
                         +\frac14\log|\Delta|.
\]
Its mean is $\frac14\log|\Delta|$. Twice this identity minus
three times the preceding mean proves \eqref{ari:eq:ca-mahler}.
All logarithmic singularities are locally integrable, so these
integrations also cover a zero at a half-period.

\item The map $z\mapsto x=\wp(z)$ has degree two. Its two inverse
sheets each have Jacobian $|\wp'(z)|^{-2}$; branch values form
a measure-zero set. Thus its pushforward density is
$2/(A|4x^3-g_2x-g_3|)$. Under $x=1/u$ the Jacobian is
$|u|^{-4}$ and the denominator is $4|u|^{-3}|C(u)|$.
This proves \eqref{ari:eq:ca-nu} and its normalization. The identity
$\wp'^2/\wp^3=4C(u)$ gives the next formula.

\item Replacing $v$ above by a half-period gives
$p(e_j)=\frac16\log|\Delta|-2U(w_j)$; summing and using
$\sum_jU(w_j)=\log2$ proves the centroid formulation.
The equality characterization is that of \ref{ari:C.8.5},
since $j=1728$ is equivalent to $g_3=0$ for a nonsingular lattice.
\end{enumerate}
\end{proof}

\begin{proposition}[The two complex beta integrals]\hypertarget{tag-001U}{}\label{ari:C.8.7}
\StatementRecord{ari:C.8.7}{proved}{001U}{\SeriesRef{ari:C.8.6}, \SeriesRef{ari:eq:ca-mahler}}

For the square and equianharmonic lattices respectively, the
integrals in \ref{ari:C.8.6} are
\[
 \int\log|C|\,d\nu=\log2,\qquad
 \int\log|C|\,d\nu=\frac32\log3-2\log2.
\]
They recover $N^2=2$ and $N^4=3$ with the same normalization.
\end{proposition}


\begin{proof}
\begin{enumerate}[label=\arabic*.,leftmargin=*]
\item By homogeneity it suffices to take $C(u)=1-u^m$, with $m=2$
or $m=3$. For $-1<\Re s<1-1/m$ put
\[
 I_m(s)=\int_{\C} |u|^{-1}|1-u^m|^{s-1}dA_u.
\]
The substitution $w=u^m$, with all $m$ sheets counted, gives
\[
 I_m(s)=\frac1m B_{\C}\left(\frac1{2m},\frac{s+1}2\right),
 \quad
 B_{\C}(a,b)=\pi\frac{\Gamma(a)\Gamma(b)\Gamma(1-a-b)}
                   {\Gamma(1-a)\Gamma(1-b)\Gamma(a+b)}.
\]
The complex beta identity \cite[\S\S1.1--1.2, equations~(1.6)--(1.7),
pp.~38--39]{Neretin2024Beta} applies with parameter pairs $a\mid a$ and $b\mid b$.
In that notation $\Gamma^{\mathbb C}(a\mid a)=\Gamma(a)/\Gamma(1-a)$ and
\[
 B_{\C}(a,b)=\int_{\C}|w|^{2a-2}|1-w|^{2b-2}\,dA_w
           =\pi B^{\mathbb C}(a\mid a,b\mid b).
\]
The source's absolute-convergence conditions are
$\Re a>0$, $\Re b>0$ and $\Re(a+b)<1$. Under
$a=1/(2m)$ and $b=(s+1)/2$, these give exactly the stated strip.
The strict inequalities give integrable power bounds at $u=0$, at each
root of $1-u^m$, and at infinity. On a sufficiently small closed
neighbourhood of $s=0$, their positive margins also absorb the
logarithmic factor introduced by differentiation. Thus differentiation
under the integral at $s=0$ is justified.

\item Consequently, with $a=1/(2m)$ and $\psi=\Gamma'/\Gamma$,
\[
 \frac{I_m'(0)}{I_m(0)}
 =\frac12\left(2\psi(1/2)-\psi(1/2-a)-\psi(1/2+a)\right).
\]
The gamma duplication and triplication formulas, differentiated
logarithmically, evaluate this as $\log2$ for $m=2$ and
$\frac32\log3-2\log2$ for $m=3$. These can also be checked
directly from
$\psi(1/2)=-\gamma-2\log2$,
$\psi(1/4)+\psi(3/4)=-2\gamma-6\log2$ and
$\psi(1/3)+\psi(2/3)=-2\gamma-3\log3$.
Equation \eqref{ari:eq:ca-mahler} gives the stated values of $N$.
\end{enumerate}
\end{proof}

\begin{proposition}[The signed Gaussian identity]\hypertarget{tag-001V}{}\label{ari:C.8.8}
\StatementRecord{ari:C.8.8}{proved}{001V}{the definitions and constructions stated in the proof}

For $\Lambda=\Z+\tau\Z$, $T=\Im\tau$, and every $z\in\C$,
\[
 \sqrt{2T}\,|\eta(\tau)|^2N(z,\tau)^2
 =\sum_{m\in\frac12\Lambda}\epsilon(m)
                         e^{-2\pi|z-m|^2/T},
\]
where $\epsilon=-1$ on $\Lambda$ and $+1$ on the three other
cosets of $\Lambda$ in $\frac12\Lambda$.
Both sides and all their derivatives converge locally uniformly.
\end{proposition}


\begin{proof}
\begin{enumerate}[label=\arabic*.,leftmargin=*]
\item Write $z=b+a\tau$ and expand in the Fourier characters
$e^{2\pi i(kb+la)}$ of $(a,b)\in(\R/\Z)^2$.
For $G(z)=\sum_{\lambda\in\Lambda}e^{-2\pi|z-\lambda|^2/T}$,
Poisson summation gives the coefficient
\[
 \frac12\exp\left(-\frac{\pi|k\tau-l|^2}{2T}\right).
\]
The right side is the sum of the four half-period translates
of $G$, minus $2G$. Its coefficient is therefore
\[
 (-1)^{k+l+kl}
       \exp\left(-\frac{\pi|k\tau-l|^2}{2T}\right).
\]
To check the sign, it is $+1$ when both $k,l$ are even and
$-1$ otherwise.

\item On the left use
\[
 \vartheta_1(z,\tau)
 =-i\sum_{n\in\Z}(-1)^n
        e^{\pi i(n+1/2)^2\tau+2\pi i(n+1/2)z}.
\]
Multiply by its conjugate and by
$\sqrt{2T}e^{-2\pi(\Im z)^2/T}$.
In the double sum put $k=n-m$ and apply one-dimensional
Poisson summation to the remaining Gaussian in $m$.
Completing the square gives precisely the coefficient just
displayed. This proves equality of the Fourier series.
Positive definiteness of $|k\tau-l|^2$ and Gaussian decay
justify the changes of order, Poisson summation, and every
claimed differentiation.
\end{enumerate}
\end{proof}




\subsection{The Betti measure and torsion equidistribution}

\begin{definition}[The current convention]
\hypertarget{tag-IA07}{}\label{I:current-convention}
\StatementRecord{I:current-convention}{proved}{IA07}{}

We take $dd^c=(i/\pi)\partial\bar\partial$, so that
$dd^c\log|w|=\delta_0$. For a local complex coordinate $w$,
$dA_w=(i/2)dw\wedge d\bar w$.
\end{definition}

\begin{proposition}[The positive Betti measure and its current]\hypertarget{tag-001W}{}\label{ari:C.8.9}
\StatementRecord{ari:C.8.9}{proved}{001W}{\SeriesRef{ari:C.3.2}, \SeriesRef{ari:eq:ca-cuspexp}}

On a local period chart write $\mu=a\tau+b$, with $a,b$ real,
and put $T=\Im\tau$. The positive Betti measure is
\[
 m_P=db\wedge da
 =\frac{|\mu_w-a\tau_w|^2}{T}\,dA_w.
\]
Reversing the wedge order gives $da\wedge db=-m_P$.
The globally single-valued function $L(t)=\log N(P_t)$ satisfies
on $\mathbb P^1(\C)$
\begin{equation}\hypertarget{tag-0026}{}\label{ari:eq:ca-current}
 dd^cL=\frac1{12}\delta_\infty-m_P,\qquad
               \int_{\mathbb P^1}m_P=\frac1{12}.
\end{equation}
There are no atoms of $m_P$ at $0$, $1$ or $\infty$.
Let $T_0=\sqrt3/2$, $\rho=(1+i\sqrt3)/2$, and
$p=\wp'_\rho((2-\rho)/3)$. The density in the ordinary
$t$ coordinate satisfies
\begin{align}
 \frac{dm_P}{dA_t}
 &=\frac{|t|^{-2/3}}{64T_0|p|^{2/3}}
                      \bigl(1+O(|t|^{1/3})\bigr),
                         &&t\longrightarrow0,\hypertarget{tag-0029}{}\label{ari:eq:ca-density0}\\
 \frac{dm_P}{dA_t}
 &=\frac{|t-1|^{-3/2}}{24\pi^2\sqrt{E_4(i)}}
                      \bigl(1+O(|t-1|^{1/2})\bigr),
                         &&t\longrightarrow1.\hypertarget{tag-002A}{}\label{ari:eq:ca-density1}
\end{align}
\end{proposition}


\begin{proof}
\begin{enumerate}[label=\arabic*.,leftmargin=*]
\item On a chart avoiding a zero of the section, the theta quotient
is holomorphic and nonzero. Its log modulus is pluriharmonic.
Only $-\pi(\Im\mu)^2/T$ contributes to the curvature.
Direct differentiation of this real quotient gives
\[
 \partial_w\partial_{\bar w}\frac{(\Im\mu)^2}{T}
       =\frac{|\mu_w-a\tau_w|^2}{2T}.
\]
Also
\[
 a_w=\frac{\mu_w-a\tau_w}{2iT},\qquad
 b_w=-\bar\tau\frac{\mu_w-a\tau_w}{2iT}.
\]
Taking the wedge product proves the formula for $m_P$ and
$dd^cL=-m_P$ on the chart. The transformation of $(a,b)$
under an integral change of period basis has determinant one,
so the measure patches. Translation of $\mu$ by a locally
constant integral period changes neither the formula nor
the measure. The formula remains valid on a root chart
at an elliptic modulus.

\item At $t=0$ the expansion proved in Proposition~\ref{ari:C.3.2} is
\[
 L(t)=\frac14\log3+\frac7{192}\Re t
       -B|t|^{4/3}+O(|t|^{5/3}),\qquad
 B=\frac{9\pi}{512T_0|p|^{2/3}}.
\]
The expansion is a convergent real-analytic expansion on a
root chart, so it can be differentiated termwise with the
indicated orders. The Euclidean identity
$\Delta_{\R^2}|t|^\alpha=\alpha^2|t|^{\alpha-2}$ gives
\eqref{ari:eq:ca-density0}.

\item At $t=1$ put $h=\mu-(1-\tau)/2$. The expansion in Proposition~\ref{ari:C.3.1}
and the modular $j$ expansion give
\[
 t-1=A_1h^4(1+O(h^2)),\qquad
 A_1=\frac94\pi^4E_4(i)>0,\qquad
 L=\log\sqrt2-\frac\pi2|h|^2+O(|h|^4).
\]
Therefore
$L=\log\sqrt2-\pi|t-1|^{1/2}/(2\sqrt{A_1})+O(|t-1|)$.
Termwise differentiation on the root chart proves
\eqref{ari:eq:ca-density1}. Both densities are locally integrable,
and the continuous potentials at $0$ and $1$ have no atoms.

\item For infinity take $u=1/t$, and set $\ell=\log(5+2\sqrt6)$.
The cusp branch is
$\mu=1/2-i\ell/(2\pi)+O(q)$, $q=e^{2\pi i\tau}$.
The theta and eta products, with the nonzero factor
$|2\sin\pi\mu|=2\sqrt3+O(q)$, give
\begin{equation}\hypertarget{tag-0027}{}\label{ari:eq:ca-cuspexp}
 L=\frac1{12}\log|u|+\log(2\sqrt3)-\frac1{12}\log1728
          -\frac{\ell^2}{2\log(1728/|u|)}
          +O\left(\frac{|u|}{|\log|u||}+|u|\right).
\end{equation}
Here $q=u/1728+O(u^2)$; the estimate follows directly by
bounding the convergent products. The nonlogarithmic
correction is continuous at $u=0$ and has no point mass.
Equivalently its Betti density is
\[
 \frac{dm_P}{dA_u}
  =\frac{\ell^2}{2\pi|u|^2\log(1728/|u|)^3}
       \bigl(1+o(1)\bigr),
\]
which is integrable. The single logarithmic term in
\eqref{ari:eq:ca-cuspexp} contributes $\delta_\infty/12$.
This proves the global current identity; pairing it with
the constant function one gives the total mass.
\end{enumerate}
\end{proof}

\begin{proposition}[Torsion equidistribution]\hypertarget{tag-001K}{}\label{ari:C.8.10}
\StatementRecord{ari:C.8.10}{proved}{001K}{\SeriesRef{ari:C.6.4}, \SeriesRef{ari:C.8.9}; \citet[Theorem~1.1 and Corollary~1.2, arXiv v2, p.~2]{demarco2017}}

Let $D_n$ be the divisor, with intersection multiplicity, of
parameters $t\in\mathbb P^1\setminus\{0,1,\infty\}$ for which
$nP_t=O$. Then
\[
 \deg D_n=\frac{n^2}{12}+O(1),\qquad
 \frac1{n^2}[D_n]\ \longrightarrow\ m_P,\qquad
 \frac1{\deg D_n}[D_n]\ \longrightarrow\ 12m_P
\]
weakly as $n\to\infty$ through indices for which $D_n\ne0$.
For any nonrepeating sequence of individual algebraic torsion
parameters, their Galois-orbit probability measures likewise
converge to $12m_P$. The factor $12$ normalizes the mass to one.
\end{proposition}


\begin{proof}
\begin{enumerate}[label=\arabic*.,leftmargin=*]
\item We apply \citet[Theorem~1.1 and Corollary~1.2, arXiv v2, p.~2]{demarco2017}.
Their hypotheses are a number field, a smooth projective
base curve, an elliptic surface and a rational section of
nonzero generic canonical height. Here they are $\Q$,
$\mathbb P^1$, the displayed rational elliptic surface,
and $P$, whose generic height for $(O)$ is $1/12$ by \ref{ari:C.6.4}.
At a torsion specialization the canonical height is zero.
Corollary 1.2 consequently gives convergence of every
nonrepeating sequence of Galois-orbit probability measures.

\item The limiting measure is the curvature of the height metric
divided by the degree of its line bundle. In our normalization,
\ref{ari:C.8.9} computes the curvature as $m_P$, and \ref{ari:C.6.4} computes
that degree as $1/12$. Hence the probability measure is
$12m_P$. The total mass is one because
$12\int_{\mathbb P^1}m_P=1$.

\item The surface height
formula, with our half-self-pairing convention, reads
\[
 \widehat h_O(Q)=\chi(\mathcal O_{\mathcal E})+(Q\cdot O)
                -\frac12\sum_v\operatorname{contr}_v(Q).
\]
Here $\chi=1$, and the finitely many fibre corrections are
bounded as $Q=nP$ varies, because they depend only on its
component. Since $\widehat h_O(nP)=n^2/12$, it follows
that $(nP\cdot O)=n^2/12+O(1)$. Removing the singular fibres
changes this by $O(1)$: the local intersection at each
such fibre is bounded as $n$ varies in characteristic zero.
One can see this in the formal group after a fixed
multiple enters it; multiplication by a nonzero integer
is a local analytic isomorphism over the residue field
$\C$. Thus $\deg D_n=n^2/12+O(1)$.

\item At any fixed smooth torsion parameter $t_0$, of exact order
$d$, the multiplicity in $D_n$ is zero unless $d\mid n$,
and is constant when $d\mid n$. Indeed the formal parameter
of $nP=[n/d](dP)$ has first term $(n/d)$ times that of $dP$,
with nonzero constant coefficient. Fix a continuous test
function and $\varepsilon>0$. Corollary 1.2 implies that
all but finitely many distinct Galois orbits of torsion
parameters have average within $\varepsilon$ of the limiting
integral; otherwise a countersequence would violate it.
The total weighted contribution of the finitely many
exceptional orbits is bounded independently of $n$, by
the multiplicity observation. Divide by $\deg D_n$ and
let $n\to\infty$. This proves the probability convergence
for $D_n$, and the degree asymptotic proves the unnormalized
version.
\end{enumerate}
\end{proof}





\section{Critical points and the global level curves}\label{pub:I:critical}

\begin{remark}[Critical-point coordinates]
\hypertarget{tag-IA27}{}\label{I:critical-point-convention}
\StatementRecord{I:critical-point-convention}{proved}{IA27}{}
Use $L=\log N$ and $t=j/1728$ as in \cref{I:def:theta-height}.
\end{remark}

\begin{lemma}[The bounded cusp zero and its differential equation]
\hypertarget{tag-IA09}{}\label{I:cusp-zero-equation}
\StatementRecord{I:cusp-zero-equation}{proved}{IA09}{}
For the cusp-bounded lift of the zero of $\wp_\tau$, one has
\begin{equation}\hypertarget{tag-002I}{}\label{ari:eq:cc-EZ}
 \begin{aligned}
 \mu&=\frac12-i\beta_{\rm cusp}-
       \frac{36i\sqrt6}{\pi}\sum_{n\ge1}\frac{a_nq^n}{n^2},
 &\beta_{\rm cusp}&=\frac{\log(5+2\sqrt6)}{2\pi},\\
 \sum_{n\ge1}a_nq^n&=\frac{\Delta_0}{E_6^{3/2}},
 &\mu''&=144\pi i\sqrt6\frac{\Delta_0}{E_6^{3/2}}.
 \end{aligned}
\end{equation}
The branch of $E_6^{3/2}$ tends to one at the cusp. The series
holds on $\Im\tau>1$ and continues on the chosen logarithm
branch away from half-period branch points.
\end{lemma}
\begin{proof}
Use the Eichler--Zagier formula
\cite[unnumbered theorem, p.~399]{EichlerZagier1982},
in the form of
\cite[Section~2, unnumbered Eichler--Zagier theorem and equation~(6),
author preprint p.~4]{DukeImamoglu2008}.
In equation~(6), take $m=n=0$ and the minus sign, so that the
resulting zero $z_0$ tends to $1/2+i\beta_{\rm cusp}$.
The exact affine map is $\mu=1-z_0$, as in \cref{inv:A2.3};
periodicity and evenness give
$\wp(1-z_0,\tau)=\wp(z_0,\tau)$.
Termwise integration uses
$\int_\tau^{i\infty}(\sigma-\tau)e^{2\pi in\sigma}\,d\sigma
=-q^n/(4\pi^2n^2)$ and gives the displayed series.
Termwise differentiation using $q'=2\pi iq$ gives its
second derivative. At $i$ the continuation uses the square-root
coordinate of \cref{ari:C.3.1}.
\end{proof}

\begin{proposition}[The radial derivative]\hypertarget{tag-001M}{}\label{ari:C.8.12}
\StatementRecord{ari:C.8.12}{proved}{001M}{\SeriesRef{ari:C.8.11}, \SeriesRef{ari:C.3.2}}

Put $p=\wp'_\tau(\mu)$ and
\[
 X=\frac{3Wp}{g_2},\qquad F=1-X,
 \qquad W=\frac{\vartheta_1'(\mu,\tau)}{\vartheta_1(\mu,\tau)}
              +\frac{2\pi i\Im\mu}{\Im\tau}.
\]
These are single-valued modular invariants wherever $g_2\ne0$.
They have continuous extensions at $t=0,1$, with
$X(0)=1/4$ and $X(1)=0$. At every $t\ne0,1$,
\begin{equation}\hypertarget{tag-002W}{}\label{ari:eq:cc-radial}
 (t-1)L_t=\frac{X(X-2)}{24},\qquad
 r\frac{d}{dr}L(1+re^{i\theta})
       =\frac{\Re(X(X-2))}{12}.
\end{equation}
Also, if $a=\Im\mu/\Im\tau$ and $R=\mu'-a$, then
\begin{equation}\hypertarget{tag-002J}{}\label{ari:eq:cc-FR}
 F=\frac{6\pi ip}{g_2}R.
\end{equation}
Consequently $|F|<1$ implies strict decrease along the ray
from $1$ through $t$.
\end{proposition}


\begin{proof}
\begin{enumerate}[label=\arabic*.,leftmargin=*]
\item A lattice translation leaves $W$ unchanged, and replacing $\mu$
by $-\mu$ changes the signs of both $W$ and $p$. Under a modular
change, $W,p,g_2$ have respective weights $1,3,4$, so $X$ is
invariant. The two identities of \ref{ari:C.8.11} give
\[
 L_\tau=\frac{iW(3Wp-2g_2)}{24\pi p},\qquad
 R=\frac{g_2-3Wp}{6\pi ip}.
\]
Ramanujan's identities and $p^2=-g_3$ give
\[
 t'_\tau=-2\pi it\frac{E_6}{E_4},\qquad
 \frac{t-1}{t'_\tau}=-\frac{E_6}{2\pi iE_4^2}.
\]
Substitution, using $g_2=4\pi^4E_4/3$ and
$g_3=8\pi^6E_6/27$, proves \eqref{ari:eq:cc-radial} and
\eqref{ari:eq:cc-FR}. Since $X(X-2)=F^2-1$, $|F|<1$ makes
its real part strictly negative.

\item At $i$, $W=p=0$ and $g_2\ne0$, so $X\to0$. At $\rho$, put $\delta=\tau-\rho$ and $p_0=\wp'_\rho(\mu_0)$.
The expansion in \ref{ari:C.3.2} and Ramanujan's identities give
\[
 R=\frac{3p_0}{8\pi^2}\delta+O(|\delta|^2),\qquad
 g_2=\frac{3ip_0^2}{\pi}\delta+O(\delta^2).
\]
Substitution into \eqref{ari:eq:cc-FR} gives $F(\rho)=3/4$.
Equivalently, the two possible limits furnished by
\eqref{ari:eq:cc-radial} and $L_t(0)=7/384$ are $X=1/4,7/4$;
the expansion of $R$ selects $1/4$. The same expansions show
continuity on the coarse $t$-line.
\end{enumerate}
\end{proof}

\begin{proposition}[A bound on the cusp region]\hypertarget{tag-001N}{}\label{ari:C.8.13}
\StatementRecord{ari:C.8.13}{proved}{001N}{\SeriesRef{ari:eq:cc-EZ}, \SeriesRef{ari:eq:ca-exp}}

If $\Im\tau\ge7/5$, then $|F|<17/20$.
\end{proposition}


\begin{proof}
\begin{enumerate}[label=\arabic*.,leftmargin=*]
\item Here $|q|<1/6600$. On the auxiliary circle $|q|=1/2000$,
the Lambert series and eta product give
\[
 |E_6-1|\le
 504r\frac{1+26r+66r^2+26r^3+r^4}{(1-r)^7},\qquad
 |\Delta_0|\le r\exp\!\frac{24r}{(1-r)^2},
 \quad r=\frac1{2000}.
\]
The right sides imply $|\Delta_0/E_6^{3/2}|<1/1250$.
In particular $E_6$ has no zero on this disc and the specified
branch is holomorphic there. By Cauchy's inequality, with
$s=10/33$,
\[
 \sum_{n\ge1}\frac{|a_nq^n|}{n}
 \le\frac{-\log(1-s)}{1250},\qquad
 \sum_{n\ge1}\frac{|a_nq^n|}{n^2}
 \le\frac1{1250}\left(s+\frac{s^2}{4(1-s)}\right).
\]
Equation \eqref{ari:eq:cc-EZ} therefore yields
\[
 |\mu-(1/2-i\beta_{\rm cusp})|<\frac1{125},\qquad
 |\mu'|<\frac{13}{250},\qquad \beta_{\rm cusp}<\frac{365}{1000}.
\]
For $r=1/6600$, the corresponding weight-four and weight-six
Lambert bounds are
\[
 240r\frac{1+4r+r^2}{(1-r)^5}<\frac{37}{1000},\qquad
 504r\frac{1+26r+66r^2+26r^3+r^4}{(1-r)^7}
       <\frac{77}{1000}.
\]
Thus $|E_4|>963/1000$ and $|E_6|^{1/2}<1038/1000$.
Since $6\pi|p/g_2|=\sqrt6|E_6|^{1/2}/|E_4|$,
\begin{equation}\hypertarget{tag-002O}{}\label{ari:eq:cc-cuspbound}
 |F|<\sqrt6\frac{1038}{963}
 \left\{\frac{365/1000+1/125}{7/5}+\frac{13}{250}\right\}
 <\frac{17}{20}.
\end{equation}
All scalar comparisons here follow from geometric series and
the exponential Taylor bounds \eqref{ari:eq:ca-exp}; for the logarithm
use $-\log(1-s)=\sum s^n/n$ with its geometric remainder.
For reference, the upper quantities before the last rounding are
less than $789/10^6$, $755/10^5$, $51/1000$, and $841/1000$
for $|\Delta_0/E_6^{3/2}|$, the error in $\mu$, the bound on
$\mu'$, and the final bound on $|F|$, respectively.
\end{enumerate}
\end{proof}

\begin{proposition}[A bound on the compact modular region]\hypertarget{tag-001O}{}\label{ari:C.8.14}
\StatementRecord{ari:C.8.14}{proved}{001O}{\SeriesRef{ari:eq:cc-EZ}, \SeriesRef{ari:eq:cc-Fseries}, \SeriesRef{ari:eq:cc-ode}, \SeriesRef{ari:eq:cc-recurrence}, \SeriesRef{ari:eq:cc-power}, \SeriesRef{ari:eq:cc-finite-table}, \SeriesRef{ari:eq:cc-norms}, \SeriesRef{ari:eq:cc-Ffactored}, \SeriesRef{ari:eq:cc-scalar}, \SeriesRef{ari:C.8.12}}

Let $\chi=(\tau-i)/(\tau+i)$. On
$0<|\chi|\le2-\sqrt3$, the invariant $F$ satisfies $|F|<1$.
At the two equianharmonic points on the boundary its continuous
value is $3/4$.
\end{proposition}


\begin{proof}
\begin{enumerate}[label=\arabic*.,leftmargin=*]
\item Set
\[
 \begin{gathered}
 r_0=2-\sqrt3,\qquad A=E_4(i),\qquad k=\pi^2A,\\
 v=\pi\sqrt A\,\chi,\qquad z=v^2=k\chi^2,\\
 z_0=kr_0^2,\qquad c_0=1-\frac1{\sqrt3}.
 \end{gathered}
\]
In this proof $z$ is a coordinate on the modular disc. Define analytic functions $f_2,f_4,f_6$
by
\[
 \begin{gathered}
 \sqrt A f_2(v)=\frac{E_2(\tau)}{(1-\chi)^2}
                         -\frac3{\pi(1-\chi)},\\
 A f_4(v)=\frac{E_4(\tau)}{(1-\chi)^4},\qquad
 A^{3/2} f_6(v)=\frac{E_6(\tau)}{(1-\chi)^6}.
 \end{gathered}
\]
Ramanujan's differential identities become the rational system
\begin{equation}\hypertarget{tag-002T}{}\label{ari:eq:cc-ode}
 \begin{gathered}
 f_2'=\frac{f_4-f_2^2}{3},\qquad
 f_4'=\frac43(f_6-f_2f_4),\qquad
 f_6'=2(f_4^2-f_2f_6),\\
 (f_2,f_4,f_6)(0)=(0,1,0).
 \end{gathered}
\end{equation}
Uniqueness makes $f_4$ even and $f_2,f_6$ odd. Write
\[
 B(z)=f_4(v),\quad J(z)=\frac{f_6(v)}{2v},\quad
 D(z)=\frac{B(z)^3-4zJ(z)^2}{J(z)^{3/2}},\qquad
 H(z)=\sum_{n\ge0}\frac{-[z^n]D(z)}{16n^2-1}z^n.
\]
The branch has $J(0)^{1/2}=1$. The estimates below establish
its regularity on a larger disc, so these definitions are
unambiguous on $|z|\le z_0$.

\item Write $h=(1-\chi)(\mu-(1-\tau)/2)$. The Möbius chain rule and
\eqref{ari:eq:cc-EZ} give
\[
 h''(\chi)=-\frac{4\mu''(\tau)}{(1-\chi)^3}
       =-\frac{i}{2\sqrt{3\pi}}\chi^{-3/2}D(k\chi^2).
\]
The root expansion at $i$ gives the branch
\begin{equation}\hypertarget{tag-002R}{}\label{ari:eq:cc-h}
 h=\frac{2i}{\sqrt{3\pi}}\sqrt \chi\,H(z).
\end{equation}
Indeed differentiating this series twice gives precisely the
coefficient denominator $16n^2-1$ above. The two possible affine
homogeneous terms are absent because the half-period difference
has a Puiseux expansion consisting only of half-integral powers.
Changing the sign of the square root changes the signs of $h,p,R$
and leaves $F$ unchanged.

\item Put $r=|\chi|$. The definitions of $a$ and $R$ give exactly
\[
 R=\frac{1-\chi}{2i}
       \left(h'+\frac{\overline \chi h+\overline h}{1-r^2}\right),
 \qquad
 \frac p{g_2}=-\frac{i\sqrt \chi}{\sqrt{3\pi}(1-\chi)}
                         \frac{\sqrt J}{B}.
\]
Consequently
\begin{equation}\hypertarget{tag-002L}{}\label{ari:eq:cc-Fseries}
 F=\frac{\sqrt J}{B}
 \left\{H+4zH'+\frac{2r}{1-r^2}(rH-\overline H)\right\}.
\end{equation}
All conjugates in this display are values at the same point;
$H$ has real coefficients, so $\overline{H(z)}=H(\overline z)$.

\item The apparent pole at the equianharmonic points must be removed
before estimating this expression. We have $B(-z_0)=0$.
At these points $R=0$, since the expansion at $\rho$ gives
$\mu'=-1/3=a$. Thus, if
\[
 S=\sqrt J H,\qquad T=\sqrt J(H+4zH'),
\]
then $T(-z_0)=c_0S(-z_0)$. Define the analytic quotients
\[
 C(z)=\frac{B(z)}{1+z/z_0},\qquad
 U(z)=\frac{T(z)-c_0S(z)}{(1+z/z_0)C(z)},\qquad
 V(z)=\frac{\sqrt{J(z)}}{C(z)}.
\]
The estimates below prove $C\ne0$ on the needed larger disc.
For $r<r_0$, equation \eqref{ari:eq:cc-Fseries} is now
\begin{equation}\hypertarget{tag-002K}{}\label{ari:eq:cc-Ffactored}
 F=U+V\frac{dH+q(H-\overline H)}{1+z/z_0},
 \qquad
 d=\frac{2(r_0-r)}{(1+r_0)(1+r)},\quad q=\frac{2r}{1-r^2}.
\end{equation}
Here $q$ is only a real abbreviation within this display and
its estimates.

\item For an analytic
power series $K=\sum K_nz^n$, put
$\|K\|_s=\sum|K_n|s^n$. Then, for $0\le s\le z_0$,
\begin{equation}\hypertarget{tag-002S}{}\label{ari:eq:cc-norms}
 \begin{aligned}
 \|U\|_s&\le\frac1{\sqrt3}+\frac{13}{200}s,&
 \|V\|_s&\le1+\frac7{50}s,\\
 \|H\|_s&\le1+\frac{113}{1000}s,&
 \|H'\|_s&\le\frac{13}{100}.
 \end{aligned}
\end{equation}
In
\eqref{ari:eq:cc-ode}, write $f_2=\sum a_nv^n$,
$f_4=\sum b_nv^n$, $f_6=\sum c_nv^n$. Starting with
$(a_0,b_0,c_0)=(0,1,0)$ (only in this recurrence the symbol
$c_0$ denotes the zero initial coefficient), compute
\begin{equation}\hypertarget{tag-002Y}{}\label{ari:eq:cc-recurrence}
 \begin{aligned}
 a_{n+1}&=\frac{b_n-\sum_{j=0}^na_ja_{n-j}}{3(n+1)},\\
 b_{n+1}&=\frac{4(c_n-\sum_{j=0}^na_jb_{n-j})}{3(n+1)},\\
 c_{n+1}&=\frac{2(\sum_{j=0}^nb_jb_{n-j}-
                          \sum_{j=0}^na_jc_{n-j})}{n+1}.
 \end{aligned}
\end{equation}
Thus $B_n=b_{2n}$ and $J_n=c_{2n+1}/2$. For powers of a series
$K$ with $K_0=1$, the equally elementary recurrence is
\begin{equation}\hypertarget{tag-002V}{}\label{ari:eq:cc-power}
 [z^n]K^\alpha=\frac1n\sum_{j=1}^n
                 ((\alpha+1)j-n)K_j[z^{n-j}]K^\alpha.
\end{equation}
Together with $H_n=-D_n/(16n^2-1)$, these determine every
coefficient used. For example,
\[
 \begin{array}{c|rrrr}
 n&0&1&2&3\\ \hline
 B_n&1&10/9&5/27&4/81\\
 J_n&1&14/27&28/135&14/405\\
 H_n&1&13/135&-1999/153090&29299/28146690.
 \end{array}
\]
Division by $1+z/z_0$ uses $C_0=1$ and
$C_n=B_n-C_{n-1}/z_0$; the same recurrence applies to
$(T-c_0S)/(1+z/z_0)$. Division by $C$ uses the usual
convolution recurrence with constant coefficient one.

\item Here are sufficient rational enclosures for the only nonrational
constants in this arithmetic:
\begin{equation}\hypertarget{tag-002N}{}\label{ari:eq:cc-constants}
 \begin{gathered}
 14<k<72/5,\qquad
 1.0315619043558156493568263343177356527388253980203
 <z_0\\
 <1.0315619043558156493568263343177356527388253980205.
 \end{gathered}
\end{equation}
They follow from
$\pi=16\arctan(1/5)-4\arctan(1/239)$, the first 100 terms
of the alternating arctangent series, the exponential series
through degree 160, and
\[
 A=1+240\sum_{n=1}^{49}\frac{n^3e^{-2\pi n}}{1-e^{-2\pi n}}
       +\varepsilon_A,\qquad
 0<\varepsilon_A<
 \frac{240(4e^{-2\pi})^{50}}
 {(1-4e^{-2\pi})(1-e^{-2\pi})}.
\]
The last bound uses $n^3\le4^n$. Bounds for $\sqrt3$ follow
by squaring rational endpoints. All ensuing coefficient
comparisons are rational arithmetic using these enclosures.

\item To bound infinite tails, first consider $|z|=12/5$.
Here $|\chi|<5/12$, $\Im\tau>7/17$ and
$|e^{2\pi i\tau}|<1/13$. Lambert series give
\begin{equation}\hypertarget{tag-002U}{}\label{ari:eq:cc-outer}
 |B|<400,\qquad |J|<3000\qquad (|z|=12/5).
\end{equation}
The respective upper bounds before rounding are
\[
 \left(\frac{12}{7}\right)^4
 \left(1+240u\frac{1+4u+u^2}{(1-u)^5}\right)<400,
 \quad u=1/13,
\]
\[
 \frac13\left(\frac{12}{7}\right)^6
 \left(1+504u\frac{1+26u+66u^2+26u^3+u^4}{(1-u)^7}\right)<3000.
\]
We have used $A>1$ and $2\sqrt{12/5}>3$.
For $R=3/2$, let $\xi=5/8$ and retain coefficients through
degree 59. The tails of $B,J,C$ are respectively at most
\begin{equation}\hypertarget{tag-0031}{}\label{ari:eq:cc-tail60}
 \frac{400\xi^{60}}{1-\xi},\qquad
 \frac{3000\xi^{60}}{1-\xi},\qquad
 \frac{400z_0}{12/5-z_0}\frac{\xi^{60}}{1-\xi}.
\end{equation}
For the last bound, $B(-z_0)=0$ implies
$|C_n|\le400z_0(12/5)^{-n}/(12/5-z_0)$.
The finite sums from \eqref{ari:eq:cc-recurrence}--\eqref{ari:eq:cc-power}
with these tails give
\[
 \begin{array}{c|cccc}
 &\|B\|_{3/2}&\|J\|_{3/2}&\|C-1\|_{3/2}
                   &\min_{|z|\le3/2}\Re J(z)\\ \hline
 \text{bound}&<3276/1000&<2411/1000&<335/1000&>2049/10000.
 \end{array}
\]
For the last column, the lower bound used is exactly
\[
 1-J_2R^2-\frac{J_1^2}{8J_2}
       -\sum_{n=3}^{59}|J_n|R^n-
                   \frac{3000\xi^{60}}{1-\xi}.
\]
This is obtained by minimizing the quadratic
$\Re(1+J_1z+J_2z^2)$ on $|z|=R$; the minimum occurs at
$\cos\arg z=-J_1/(4J_2R)$. In particular, on $|z|\le3/2$,
\[
 |B|<33/10,\quad |J|<5/2,\quad |C|>3/5,\quad
 \Re J>1/5,\quad
 |D|\le\frac{(33/10)^3+6(5/2)^2}{(1/5)^{3/2}}<1000.
\]
Cauchy's inequality and the denominators in $H_n$ now give
\[
 |H|\le1+1000\sum_{n\ge1}\frac1{16n^2-1}
       =1+1000\left(\frac12-\frac\pi8\right)<110
                      \quad (|z|\le3/2).
\]
On $|z|\le7/5$, the same argument gives
\[
 |H+4zH'|\le1+1000\sum_{n\ge1}
                 \frac{(14/15)^n}{4n-1}
       \le1+\frac{1000}{3}\log15<910.
\]
It follows on $|z|=7/5$ that $|U|<10000$ and $|V|<3$;
indeed $|(1+z/z_0)C|>(7/(5z_0)-1)3/5$, and
$|\sqrt J|<8/5$.

\item Finally retain coefficients through degree 99 and put
$\eta=5z_0/7$. The tails in the coefficient norms for
$U,V,H,H'$ at $z_0$ are bounded respectively by
\begin{equation}\hypertarget{tag-0030}{}\label{ari:eq:cc-tail100}
 \frac{10000\eta^{100}}{1-\eta},\quad
 \frac{3\eta^{100}}{1-\eta},\quad
 \frac{110\eta^{100}}{1-\eta},\quad
 \frac{110\eta^{100}}{z_0}
 \left(\frac{100}{1-\eta}+\frac{\eta}{(1-\eta)^2}\right).
\end{equation}
The finite coefficient sums plus precisely these tails give
\begin{equation}\hypertarget{tag-002P}{}\label{ari:eq:cc-finite-table}
 \begin{array}{c|cccc}
 & (\|U\|_{z_0}-1/\sqrt3)/z_0
 & (\|V\|_{z_0}-1)/z_0
 & (\|H\|_{z_0}-1)/z_0 & \|H'\|_{z_0}\\ \hline
 \text{upper bound}&635/10000&1368/10000&1115/10000&1293/10000.
 \end{array}
\end{equation}
The recurrences, cutoffs and remainders above specify each
finite comparison by rational arithmetic.
Equations \eqref{ari:eq:cc-finite-table} imply
\eqref{ari:eq:cc-norms}, with room in each inequality.

\item For $s=|z|=kr^2$,
\[
 |H-\overline H|\le2|\Im z|\,\|H'\|_s,\qquad
 \frac{|\Im z|}{|1+z/z_0|}\le s.
\]
The second inequality follows by writing $z/z_0=ue^{i\phi}$
and using $(\cos\phi+u)^2\ge0$. The useful cancellation is
\[
 |dH+q(H-\overline H)|^2
      =d^2|H|^2+4q(d+q)(\Im H)^2.
\]
Also
\[
 \frac d{|1+z/z_0|}\le
 e(r):=\frac{c_0}{(1+r)(1+r/r_0)},\qquad
 q\le\kappa r,\quad d+q\le c_0+\kappa r^2,
 \quad\kappa=1+\frac2{\sqrt3}.
\]
The last bounds use $(1-r_0^2)^{-1}=1/2+1/\sqrt3$.
It follows from \eqref{ari:eq:cc-Ffactored}--\eqref{ari:eq:cc-norms},
using $k<72/5$, that
\begin{equation}\hypertarget{tag-002Z}{}\label{ari:eq:cc-scalar}
 |F|\le U_0(r)+V_0(r)
 \sqrt{e(r)^2H_0(r)^2+
       4\kappa r(c_0+\kappa r^2)(72/5)^2r^4(13/100)^2},
\end{equation}
where
\[
 U_0=\frac1{\sqrt3}+\frac{13}{200}\frac{72}{5}r^2,
 \quad V_0=1+\frac7{50}\frac{72}{5}r^2,
 \quad H_0=1+\frac{113}{1000}\frac{72}{5}r^2.
\]
To verify that the right side is less than one, let
\[
 M(r)=(1+r)^2(1+r/r_0)^2,
\]
\begin{equation}\hypertarget{tag-002M}{}\label{ari:eq:cc-P}
 \begin{aligned}
 P(r)={}&(1-U_0)^2M(r)\\
 &-V_0^2\left\{c_0^2H_0^2+
 4\kappa r(c_0+\kappa r^2)(72/5)^2r^4(13/100)^2M(r)\right\}.
 \end{aligned}
\end{equation}
This polynomial has $P(0)=0$. Write
\[
 \frac{P(r_0x)}{r_0x}
     =\sum_{j=0}^{14}b_j\binom{14}{j}x^j(1-x)^{14-j}.
\]
All coefficients belong to $\Q(\sqrt3)$. Explicit lower bounds,
obtained by expanding the displayed factored polynomial, are
\[
 \begin{array}{c|rrrrrrrr}
 j&0&1&2&3&4&5&6&7\\ \hline
 100b_j>&169&175&181&187&192&197&202&205\\[2pt]
 j&8&9&10&11&12&13&14&\\ \hline
 100b_j>&207&207&204&197&184&160&118&.
 \end{array}
\]
Explicitly, if $P(r)/r=\sum q_nr^n$,
then
\[
 b_j=\sum_{n=0}^j q_nr_0^n
                  \frac{\binom jn}{\binom{14}n}.
\]
For example $b_0=4-4\sqrt3/3$ and
$b_1=25357/13125-453\sqrt3/4375$. The rational enclosure
\[
 \begin{gathered}
 1.73205080756887729352744634150587236694280525381038<\sqrt3,\\
 \sqrt3<1.73205080756887729352744634150587236694280525381039.
 \end{gathered}
\]
verifies every entry by integer multiplication; its endpoints
are verified by squaring. Every Bernstein basis function is
nonnegative on $[0,1]$, so $P(r)>0$ for $0<r\le r_0$.
Since $U_0<1$ there, squaring \eqref{ari:eq:cc-scalar} is legitimate
and proves $|F|<1$. Formula \eqref{ari:eq:cc-Ffactored} applies
at $r=r_0$ except at $z=-z_0$ by continuity; at that remaining
point \ref{ari:C.8.12} gives $F=3/4$. This proves the assertion.
\end{enumerate}
\end{proof}

\begin{proposition}[The critical points and global level curves]\hypertarget{tag-001P}{}\label{ari:C.8.15}
\StatementRecord{ari:C.8.15}{proved}{001P}{\SeriesRef{ari:C.8.13}, \SeriesRef{ari:C.8.14}, \SeriesRef{ari:eq:cc-radial}, \SeriesRef{ari:eq:cc-EZ}, \SeriesRef{fam:src:C.2.3}}

For every zero $\mu$ of $\wp_\tau$, the zero sets of the two
critical factors are exactly
\[
 \begin{aligned}
 W=0&\quad\Longleftrightarrow\quad
       \tau\in\operatorname{SL}_2(\Z)\{i,\rho\},\\
 3W\wp'_\tau(\mu)-2g_2=0&\quad\Longleftrightarrow\quad
       \tau\in\operatorname{SL}_2(\Z)\rho.
 \end{aligned}
\]
On the coarse $t$-plane the function $L$ has no critical point
away from $t=1$. The point $t=0$ is a noncritical $C^1$ point,
with $L_t(0)=7/384$. The point $t=1$ is the unique global
maximum, of value $\log\sqrt2$, and has the fractional-power
local maximum described in \cref{ari:C.3.1}. In particular there is no
saddle in $0<t<1$. On the root charts, $h=0$ over $t=1$
is a nondegenerate maximum, while $\delta=0$ over $t=0$
is a degenerate critical point with nonzero harmonic cubic
leading term and three alternating ascending and descending
sectors. These are all critical points on the root charts.

For every angle $\theta$, $r\mapsto L(1+re^{i\theta})$ is
strictly decreasing for $r>0$, from $\log\sqrt2$ to $-\infty$.
For each $c<\log\sqrt2$, the level set $\{L=c\}$ is one
Jordan curve, strictly star-shaped about $1$. Its bounded
complement is $\{L>c\}$; its unbounded complement is
$\{L<c\}$. The curve is real analytic except possibly at
$t=0$ when $c=\frac14\log3$, where it is $C^1$ with nonzero
tangent. The closed superlevel sets form nested topological
discs. At the maximum the level set is the single point $1$,
and above it the level set is empty. The corresponding
statements for $D_{\rm part}$ follow by its increasing affine
change from $L$.
\end{proposition}


\begin{proof}
\begin{enumerate}[label=\arabic*.,leftmargin=*]
\item Take a representative in the closed standard modular fundamental
region. If $T\ge7/5$, \ref{ari:C.8.13} applies. Otherwise
$|\Re\tau|\le1/2$ and $\sqrt3/2\le T\le7/5$. For fixed
$\Re\tau$, the function
\[
 |\chi|^2=\frac{(\Re\tau)^2+(T-1)^2}
                    {(\Re\tau)^2+(T+1)^2}
\]
has its maximum on this interval at an endpoint and increases
with $(\Re\tau)^2$. The two endpoint bounds are $r_0^2$
and $41/601<r_0^2$. Thus \ref{ari:C.8.14} covers the entire remaining
region. We conclude $|1-X|<1$ for every $t\ne1$.
Equations \eqref{ari:eq:cc-radial} now prove strict radial
monotonicity, including $t=0$ by continuity. They also exclude
both critical equations $W=0$ and $3Wp-2g_2=0$ on the
ordinary base. The elliptic points must be checked on their
root charts: at $\rho$, the nonzero coarse derivative already
given rules out a critical point; at $i$, $W=p=0$ but $3Wp-2g_2=-2g_2\ne0$, and \cref{ari:C.3.1} gives
the strict maximum. At $\rho$, $W=g_2=0$ and $p\ne0$,
so both displayed factors vanish. This proves their exact zero
sets, including the elliptic orbits. In particular the vanishing of both factors at
$\rho$ in the $\tau$ coordinate is a ramification effect,
not a saddle on the coarse base. More explicitly, on the root
chart at $\rho$,
\[
 L=L(0)+\frac7{192}\Re\left(\frac{ip_0^2}{\pi^3}\delta^3\right)
                   +O(|\delta|^4),
\]
with $p_0\ne0$, whereas on the chart at $i$,
$L=\log\sqrt2-\pi|h|^2/2+O(|h|^4)$.
These expansions give the asserted root-chart critical types.

\item At the cusp, the product formula and \eqref{ari:eq:cc-EZ} give,
uniformly in $\Re\tau$,
\[
 L=-\frac{\pi T}{6}+\log(2\sqrt3)-
                      \frac{\pi\beta_{\rm cusp}^2}{T}+O(e^{-2\pi T}),
 \qquad |t|=\frac{e^{2\pi T}}{1728}(1+o(1)).
\]
Hence $L\to-\infty$ uniformly as $|t|\to\infty$.
Continuity at $t=1$ supplies the other endpoint of each ray.
The intermediate value theorem and strict monotonicity give
one and only one radius $r_c(\theta)>0$ at each level $c$.
Properness and strict monotonicity imply continuity in $\theta$;
the nonvanishing radial derivative gives the stated local
regularity by the implicit-function theorem. Thus
$\theta\mapsto1+r_c(\theta)e^{i\theta}$ is a continuous
injective parametrization of a Jordan curve and has exactly
the asserted complementary regions.

\end{enumerate}
\end{proof}



\begin{remark}[Interval enclosures]
\hypertarget{tag-034N}{}\label{push:section:E.4}
\StatementRecord{push:section:E.4}{proved}{034N}{}
\label{sec:e4}
\hypertarget{tag-033L}{}\label{push:src:E.4.1}
\hypertarget{tag-033M}{}\label{push:src:E.4.2}
The interval bounds for the norm in \cref{fam:src:C.2.3} combine
contraction enclosures for the zeros of $\wp$, explicit theta and eta
tails, and an exact cover of the compact modular region.
\end{remark}

\section{Arithmetic periods and compact subgroups}\label{pub:I:arithmetic}

\begin{definition}[Algebraic period coordinates]
\hypertarget{tag-IA10}{}\label{I:algebraic-period-coordinates}
\StatementRecord{I:algebraic-period-coordinates}{proved}{IA10}{}
Use an algebraic Weierstrass equation and invariant differential
as in \cref{I:def:theta-height}. Write
$\Lambda=\mathbb Z\omega_1+\mathbb Z\omega_2$ with
$\Im(\omega_2/\omega_1)>0$,
$\zeta(z+\omega_j)=\zeta(z)+\eta_j$, and extend
$\eta:\Lambda\to\mathbb C$ real-linearly.
The discriminant is that of the equation in this differential.
For a homothety $\Lambda\mapsto c\Lambda$, $u\mapsto cu$,
\[
 \sigma_{c\Lambda}(cu)=c\sigma_\Lambda(u),\qquad
 \eta_{c\Lambda}(cu)=c^{-1}\eta_\Lambda(u),\qquad
 \Delta(c\Lambda)=c^{-12}\Delta(\Lambda).
\]
These formulas cancel the factors in \eqref{ari:eq:cp-klein},
so $N$ has the same value in the transported coordinates.
\end{definition}

\begin{proposition}[The sigma expression and third-kind period]\hypertarget{tag-0017}{}\label{ari:C.7.1}\hypertarget{tag-IA01}{}\label{ari:C.7.1:lift-transport}
\StatementRecord{ari:C.7.1}{proved}{0017}{\SeriesRef{ari:eq:cp-periods}; \citet[Definitions 2.3--2.4 and
Section 6]{bertolin2022}}

Let $P\in E(\overline\Q)$, with $P\ne O$ and $2P\ne O$, and
let $u$ be an elliptic logarithm of $P$. Then
\begin{equation}\hypertarget{tag-003K}{}\label{ari:eq:cp-sigma}
 2\log N(P)=\frac16\log|\Delta|
       +\Re\{\Lambda_1+u(\zeta(u)-\eta(u))\},
 \qquad \Lambda_1=2\log\sigma(u)-u\zeta(u).
\end{equation}
For a fixed elliptic logarithm $u$, changing the branch of
$\log\sigma(u)$ changes the expression inside the real part by an
integer multiple of $4\pi i$. Put
$B(u)=\Lambda_1+u(\zeta(u)-\eta(u))$. Under a lattice translation,
\[
 B(u+\omega)-B(u)\equiv
 u\eta(\omega)-\omega\eta(u)+2\log\chi(\omega)
 \pmod{4\pi i\Z},\qquad \omega\in\Lambda,
\]
where $\chi(\omega)\in\{1,-1\}$ is the sigma semicharacter.
This difference is purely imaginary, so the displayed real part
is independent of both choices.
There is an extension $G_P$ of $E$ by $\mathbb G_m$, defined
over $\overline\Q$ and corresponding, under the principal
polarization, to $P$, for which $\Lambda_1$, modulo a logarithm
of an algebraic number, is the negative of a relative period
of $[\Z\to G_P]$. Changing the third-kind differential while retaining its simple
poles and residues adds an algebraic multiple of $u$.
Adding an exact algebraic differential to a relative-integral
representative adds algebraic endpoint differences when both
endpoints are regular.
\end{proposition}


\begin{proof}
\begin{enumerate}[label=\arabic*.,leftmargin=*]
\item The sigma--theta product, or its divisor and automorphy law
together with its derivative at zero, gives
\begin{equation}\hypertarget{tag-003F}{}\label{ari:eq:cp-klein}
 N(P)=|\Delta|^{1/12}|\sigma(u)|
                   \exp\!\left(-\frac12\Re(u\eta(u))\right).
\end{equation}
To check the constant, for $\Lambda=\Z+\tau\Z$ one uses
$\sigma(z)=e^{\eta_1z^2/2}\vartheta_1(z,\tau)/(2\pi\eta(\tau)^3)$,
$\Delta=(2\pi)^{12}\eta(\tau)^{24}$, and
$\eta(z)=\eta_1z-2\pi i\Im z/T$. Substitution gives exactly
the defining Gaussian in $N$. Homogeneity gives the general
lattice formula. Taking logarithms proves \eqref{ari:eq:cp-sigma}.
The automorphy law
$\sigma(u+\omega)=\chi(\omega)
 e^{\eta(\omega)(u+\omega/2)}\sigma(u)$ and real-linearity of
$\eta$ give the stated lift transport. Explicitly, writing
$u=a\omega_1+b\omega_2$ with $a,b\in\R$ and
$\omega=m\omega_1+n\omega_2$ with $m,n\in\Z$, Legendre's
relation gives
$u\eta(\omega)-\omega\eta(u)=2\pi i(bm-an)$.
This also proves that the transported difference is purely imaginary.

\item Put
\[
 f_u(z)=\frac{\sigma(z+u)}{\sigma(z)\sigma(u)}e^{-\zeta(u)z}.
\]
Its logarithmic differential is
\[
 \frac{df_u}{f_u}
 =\frac12\frac{\wp'(z)-\wp'(u)}{\wp(z)-\wp(u)}\,dz.
\]
This is the pullback of an algebraic differential with residues
$+1$ at $-P$ and $-1$ at $O$.
Two algebraic differentials with precisely these simple poles
and residues differ by a holomorphic differential, hence by
$c\,dz$ with $c\in\overline{\mathbb Q}$.
Integration gives $cu$; integration of $df$ gives the stated
endpoint difference.
Here the principal polarization identifies $P$ with
$t_P^*\mathcal O_E(O)\otimes\mathcal O_E(O)^{-1}
=\mathcal O_E((-P)-(O))$; this fixes the sign of the extension class. The apparent $0/0$ at $z=u$
has a removable singularity when $P$ is not $2$-torsion.
The rational cocycle $f_u(z+z')/(f_u(z)f_u(z'))$ defines
the algebraic extension; in its rational trivialization the
exponential map has the form
\[
 \exp_{G_P}(z,v)=(\exp_E z,e^v f_u(z)).
\]
Its periods include
\begin{equation}\hypertarget{tag-003G}{}\label{ari:eq:cp-periods}
 (0,2\pi i),\qquad
 (\omega_j,\zeta(u)\omega_j-\eta_j u),\quad j=1,2.
\end{equation}
These signs follow directly from the sigma automorphy law.
The identity
$\sigma(2u)=-\wp'(u)\sigma(u)^4$ implies
\[
 \log f_u(u)=\Lambda_1+\log(-\wp'(u)).
\]
For an algebraic lift $R=(P,r)$, $r\in\overline\Q^{\times}$,
the second coordinate of $\log_{G_P}R$ is
\begin{equation}\hypertarget{tag-003L}{}\label{ari:eq:cp-v}
 v=\log r-\Lambda_1-\log(-\wp'(u)).
\end{equation}
It is a period of the relative homology defining the indicated
$1$-motive.
The description agrees with the generalized logarithm and
Serre function in \citet[Section~2.3, equations~(2.10)--(2.13), pp.~7--8;
Definitions~2.3--2.4, p.~8; Section~6.1, pp.~26--27,
arXiv v2]{bertolin2022}; the signs in
\eqref{ari:eq:cp-periods} use the displayed exponential convention.
\end{enumerate}
\end{proof}

\begin{proposition}[The compact subgroup and real-axis formula]\hypertarget{tag-001B}{}\label{ari:C.7.2}
\StatementRecord{ari:C.7.2}{proved}{001B}{\SeriesRef{ari:C.7.1}, \SeriesRef{ari:eq:cp-periods}, \SeriesRef{ari:eq:cp-v}, \SeriesRef{ari:eq:cp-chi}; \citet[PDF slides~5--8]{bertrand2017}}

For $P$ as in \cref{ari:C.7.1}, $N(P)$ is algebraic if and only
if the fibre of $G_P$ above $P$ contains an algebraic point of
the maximal compact subgroup $G_P(\C)^c$.
For the family in the stated conventions, if $t>1$ is real, choose
$\tau=iT$, $T>1$, and the zero $\mu=1/2-iy$, $0<y<T/2$.
With $x=2y/T$, one has
\[
 N(P)^4=3^{3/4}(t-1)^{-1/4}
       \left|\frac{\vartheta_1(x,i/T)}{\eta(i/T)}\right|.
\]
\end{proposition}


\begin{proof}
\begin{enumerate}[label=\arabic*.,leftmargin=*]
\item On the analytic universal cover of $G_P$ define the real
homomorphism
\[
 \chi(z,v)=\Re\{v-\zeta(u)z+u\eta(z)\}.
\]
It vanishes on the periods \eqref{ari:eq:cp-periods}, hence descends
to $G_P(\C)$. Its restriction to the torus is $\log|\cdot|$.
The kernel is compact: in the universal cover it is a real
three-dimensional vector space modulo the rank-three period
lattice. Every compact subgroup has zero image in $\R$, so
this kernel is exactly $G_P(\C)^c$.
The identity
$\Re(u\eta(z))=\Re(z\eta(u))$ follows from the Legendre relation;
equivalently $\eta(z)=sz+\pi\bar z/\operatorname{area}(\Lambda)$
for a complex constant $s$. Substituting
\eqref{ari:eq:cp-v} at $z=u$ gives
\begin{equation}\hypertarget{tag-003E}{}\label{ari:eq:cp-chi}
 \chi(R)=\log|r|-\log|\wp'(u)|
            -2\log N(P)+\frac16\log|\Delta|.
\end{equation}
If $r$ is algebraic and $\chi(R)=0$, this equation makes
$N(P)^2$ algebraic. Conversely, if $N(P)$ is algebraic,
choose the positive real algebraic number
$r=|\wp'(u)|\,|\Delta|^{-1/6}N(P)^2$.
Then \eqref{ari:eq:cp-chi} vanishes. This proves both directions
with the chosen rational trivialization of the fibre.
The compact interpretation is the one discussed by
\citet[PDF slides~5--8]{bertrand2017}; here the calculation fixes
its normalization on the diagonal $P=P$.

\item For the real formula, the rectangular elliptic curve has its
two zeros of $\wp$ on $1/2+i\R$; along that edge $\wp$ runs
between the positive root $e_1$ and the negative root $e_2$.
Thus the indicated representative exists. Translate $2\mu$
by $-1$ and apply the modular substitution $\tau\mapsto-1/\tau$.
The resulting elliptic argument is $-2y/T$, which is real.
The Gaussian is therefore one. Evenness of $N$ and the
duplication identity prove the displayed formula. This removes
the Gaussian from this particular expression.
\end{enumerate}
\end{proof}

\subsection{Motivic groups from the point and extension data}

\begin{remark}[The motivic dimension formula]
\hypertarget{tag-IA11}{}\label{I:motivic-dimension-input}
\StatementRecord{I:motivic-dimension-input}{proved}{IA11}{}
\citet[Construction~3.4, p.~12, and Lemma~3.5, p.~13, arXiv v2]{bertolin2022} compute the motivic
group from its abelian point and extension data and the central
Weil-pairing bracket. For the $1$-motives used here, let $B$ be the
smallest abelian subvariety generated by these data, and let $Z(1)$
be the central torus obtained from that bracket and the remaining
algebraic lift. Their dimension formula is
\begin{equation}\hypertarget{tag-003D}{}\label{ari:eq:cp-MTdim}
 \dim G_{\mathrm{mot}}(M)
 =\dim G_{\mathrm{mot}}(A)+2\dim B+\dim Z(1).
\end{equation}
The bracket is described in their Section~6.1, pp.~26--27, and
their corrected Table~2, p.~29, with the correction on p.~28, gives unipotent dimension
three in the dependent nondeficient case.
\end{remark}

\begin{proposition}[The motivic groups of the diagonal data]\hypertarget{tag-001C}{}\label{ari:C.7.3}
\StatementRecord{ari:C.7.3}{proved}{001C}{\SeriesRef{ari:eq:cp-MTdim}}

Let $P$ be nontorsion and let $R$ be any algebraic lift above
$P$ in $G_P$. Put $M_R=[\Z\to G_P]$, $1\mapsto R$.
Its connected motivic group is, after a choice of Levi subgroup,
the semidirect product of $G_{\mathrm{mot}}(E)$ and the
three-dimensional Heisenberg group attached to $H_1(E)$ and
its principal polarization. Its dimension is $5$ for CM and
$7$ without CM.

For a non-real embedding put
$\mathcal M=M_R\oplus\overline{M_R}$.
If $E$ and $\bar E$ are not isogenous, its dimension is $13$.
If they are isogenous, choose an isogeny
$\phi:\bar E\to E$ of degree $d$ and write
$Q=\phi(\bar P)$. Let $r$ be the dimension over
$F=\End(E)\otimes\Q$ of the span of $P,Q$, modulo torsion.
Then
\[
\begin{array}{c|c|c|c}
\text{case}&\dim B&\dim Z&\dim G_{\mathrm{mot}}(\mathcal M)\\ \hline
r=2,\ E\text{ CM}&2&2&8\\
r=2,\ E\text{ not CM}&2&2&10\\
r=1,\ E\text{ CM}&1&1+e&5+e\\
r=1,\ E\text{ not CM}&1&1+e&7+e.
\end{array}
\]
Here $e$ is either zero or one: it is zero precisely when the
algebraic lift in the one-dimensional quotient of the two
central tori by the bracket torus is torsion. This criterion is
independent of the choices made modulo isogeny. In the real
case we use the single motive $M_R$.
\end{proposition}


\begin{proof}
\begin{enumerate}[label=\arabic*.,leftmargin=*]
\item All curves, polarizations, points, and lifts are algebraic over
$\overline\Q$, as required by Construction 3.4 and Lemma 3.5.
For $M_R$ the smallest abelian subvariety containing $(P,P)$ is
the diagonal $B=E$, since $P$ is nontorsion. If $\psi$ is the
polarization pairing, the bracket on this diagonal is
\[
 [(x,x),(y,y)]=\psi(x,y)-\psi(y,x)=2\psi(x,y).
\]
It is nonzero and spans $\Q(1)$. Hence $Z'=Z=\mathbb G_m$.
The pure group is $\mathrm{GL}_2$ without CM and
$\operatorname{Res}_{F/\Q}\mathbb G_m$ for the imaginary
quadratic CM field. Their dimensions are four and two.
This proves the single-motive assertion, including its Lie
bracket and Levi action: the action on the centre is the
polarization multiplier.

\item For the direct sum only the two diagonal central coordinates
occur. If the two elliptic curves are not isogenous, they
cannot be CM conjugates. The pure group is
$\{(g_1,g_2)\in\mathrm{GL}_2^2:\det g_1=\det g_2\}$,
of dimension seven. The point data generate $E\times\bar E$,
and the two brackets are independent, giving $7+4+2=13$.

\item If the curves are isogenous, transport the second polarized
factor through $\phi$. Its polarization pairing becomes
$\psi/d$. The smallest abelian subvariety has dimension $r$.
For $r=2$, independent variation in its two factors makes the
two diagonal brackets span both central coordinates; thus
$Z'=Z=\mathbb G_m^2$. The unipotent group is the product
of two Heisenberg groups, with the same pure group acting
on both standard factors.
For $r=1$, after multiplying the points by a common integer
we may write $Q=aP$, $a\in F^\times$. The bracket image is
the one-dimensional line with proportional central coordinates
\begin{equation}\hypertarget{tag-003H}{}\label{ari:eq:cp-positive-centre}
 (1,c),\qquad c=\frac{a a^\dagger}{d}>0,
\end{equation}
where $\dagger$ is the Rosati involution. On an elliptic curve
$aa^\dagger$ is a positive rational number: a square in the
non-CM case and the positive field norm in the CM case.
Construction 3.4(3) adds a further central direction exactly
when the image of the algebraic lift in the quotient torus
is nontorsion. This is the integer $e$ above.
The unipotent Lie algebra is the Heisenberg algebra, with one
additional central Tate line when $e=1$. Formula
\eqref{ari:eq:cp-MTdim} proves every entry of the table.
\end{enumerate}
\end{proof}










\begin{lemma}[The full-period field and the compactness relation]
\hypertarget{tag-IA19}{}\label{I:compact-period-relation}
\StatementRecord{I:compact-period-relation}{proved}{IA19}{}
Let $E/\overline{\mathbb Q}$ and the nontorsion point $P$
have the coordinates of \cref{ari:C.7.1}.
Put $G=G_P\times G_{\bar P}$ and
$K_0=\overline{\mathbb Q}(\Omega_G)$, where $\Omega_G$
is the full de Rham--Betti period matrix, including the
second-kind periods of both factors.
Let $R_N=(P,r_N)$ be the complex compact point specified by
\cref{ari:C.7.2}, with
$r_N=|\wp'(u)|\,|\Delta|^{-1/6}N(P)^2$, and let $(u,v)$
be a logarithm in its algebraic coordinates.
Then $u,\zeta(u),\eta(u)$ and their conjugates belong to
$K_0$, and the compactness equation is
\begin{equation}\hypertarget{tag-003J}{}\label{ari:eq:cp-relation}
 v+\bar v=u\zeta(u)+\bar u\overline{\zeta(u)}
                -u\eta(u)-\bar u\overline{\eta(u)}\in K_0.
\end{equation}
\end{lemma}
\begin{proof}
The full period matrix contains
$\kappa_j=\zeta(u)\omega_j-\eta_j u$.
The coefficient matrix for $(\zeta(u),u)$ is
$\left(\begin{smallmatrix}\omega_1&-\eta_1\\
\omega_2&-\eta_2\end{smallmatrix}\right)$.
Its determinant is the nonzero Legendre determinant, so
$u,\zeta(u)\in K_0$. The same calculation in the conjugate
factor gives their conjugates.
Writing $u=a\omega_1+b\omega_2$ with $a,b\in\mathbb R$
and solving together with the conjugate equation gives
\begin{equation}\hypertarget{tag-003I}{}\label{ari:eq:cp-real-linear}
 \eta(u)=
 \frac{(u\bar\omega_2-\bar u\omega_2)\eta_1
       +(\bar u\omega_1-u\bar\omega_1)\eta_2}
      {\omega_1\bar\omega_2-\bar\omega_1\omega_2}.
\end{equation}
Its denominator is nonzero by real-linear independence.
Thus $\eta(u)$ and its conjugate are also in $K_0$.
Finally \cref{ari:C.7.2} gives
$\Re(v-u\zeta(u)+u\eta(u))=0$.
Adding this expression to its conjugate proves
\eqref{ari:eq:cp-relation}.
\end{proof}

\subsection{CM theta normalization}

\begin{remark}[The theta and metric normalizations]
\hypertarget{tag-IA28}{}\label{I:CM-theta-metric-convention}
\StatementRecord{I:CM-theta-metric-convention}{proved}{IA28}{}
The multiplication law in \cref{ari:C.7.6} retains its finite sign
character. Its exact relation to the metric norm is
\eqref{ari:eq:ct-Gaussian}.
\end{remark}

\begin{proposition}[The CM multiplication law]\hypertarget{tag-001F}{}\label{ari:C.7.6}
\StatementRecord{ari:C.7.6}{proved}{001F}{the definitions and constructions stated in the proof}

Let $E/\overline{\Q}$ have complex multiplication by an order
$\mathcal O$, and use an algebraic invariant differential. Write
$A=\operatorname{area}(\Lambda)$ and
\[
 \eta(z)=s_2z+\frac{\pi}{A}\bar z,\qquad
 \theta(z)=\sigma(z)e^{-s_2z^2/2}.
\]
Then $s_2\in\overline{\Q}$. For every nonzero
$\alpha\in\mathcal O$, with $d=\alpha\bar\alpha$, there is
$F_\alpha\in\overline{\Q}(E)$ such that
\begin{equation}\hypertarget{tag-003Q}{}\label{ari:eq:ct-CMmult}
 \theta(\alpha z)^2=F_\alpha(\exp_E z)\theta(z)^{2d}.
\end{equation}
Its divisor and leading Laurent coefficient are
\[
 \operatorname{div}(F_\alpha)=2[\alpha]^*(O)-2d(O),\qquad
 F_\alpha(\exp_Ez)=\alpha^2z^{2-2d}(1+O(z^2)).
\]
The unsquared quotient is a meromorphic function on the
cover of degree at most two with function field
$\overline{\mathbb Q}(E)(\sqrt{F_\alpha})$.
It descends to $E$ exactly when its residual sign character is trivial. For a nontorsion
$P=\exp_Eu\in E(\overline{\Q})$, algebraicity of $\theta(u)$
would imply algebraicity of $\theta(\alpha u)$ for all
$\alpha\in\mathcal O$. The additive group $\mathcal O u$ has rank two.
\end{proposition}


\begin{proof}
\begin{enumerate}[label=\arabic*.,leftmargin=*]
\item The Legendre relation gives the displayed real-linear decomposition
of $\eta$. The sigma automorphy formula gives
\[
 \theta(z+\omega)=\epsilon(\omega)
  \exp\left(\frac{\pi\bar\omega}{A}
                         (z+\omega/2)\right)\theta(z),
 \qquad \epsilon(\omega)\in\{1,-1\}.
\]
For $\alpha\Lambda\subset\Lambda$, the exponential multipliers
of $\theta(\alpha z)/\theta(z)^d$ cancel because
$\alpha\bar\alpha=d$. Its remaining multiplier is a sign.
Squaring therefore gives a meromorphic function on $E$, with the
stated divisor. This divisor is defined over $\overline{\Q}$.
The existence of a rational function with this divisor over $\C$
implies its existence over $\overline{\Q}$: the relevant
one-dimensional Riemann--Roch space is obtained by extension of
scalars. Normalize its leading Laurent coefficient, in the formal
parameter of the algebraic differential, to $\alpha^2$. This
identifies it with the squared analytic quotient. The action of
$\alpha$ on this differential is algebraic, so $\alpha^2$ is an
allowed normalization.

\item Choose a nonreal $\alpha\in\mathcal O$. Since the normalized
sigma series is $\sigma(z)=z+O(z^5)$ with algebraic Laurent
coefficients, the next coefficient of the squared quotient is
\[
 \alpha^2z^{2-2d}
       \{1-s_2(\alpha^2-d)z^2+O(z^4)\}.
\]
The coefficients of the rational function in this algebraic formal
parameter are algebraic. As $\alpha^2-d\ne0$, this proves
$s_2\in\overline{\Q}$. It also proves that the functions in
\eqref{ari:eq:ct-CMmult} are algebraic over the chosen field.
At nontorsion $u$, none of the arguments $\alpha u$ is a lattice
point; evaluating the identity proves the algebraicity implication.
The injective map $\mathcal O\to\C$, $\alpha\mapsto\alpha u$,
proves the rank assertion.

\item The qualification about rationality is necessary. For a square
lattice, $s_2=0$ and $\alpha=1+i$ give
$F_{1+i}=2i\wp$. The unsquared quotient has a simple zero at
the nonzero point in $\ker(1+i)$ and a simple pole at $O$,
which is not a principal divisor on $E$; its square is the
rational function just displayed.
\end{enumerate}
\end{proof}

\begin{cited}[Two normalized transcendental values]\hypertarget{tag-001G}{}\label{ari:C.7.7}
\StatementRecord{ari:C.7.7}{proved}{001G}{\cref{ari:C.7.6}; \cite[unnumbered theorem, Section~3, p.~113; Corollary~1, Section~4, pp.~113--114]{waldschmidt1979sigma}; \cite[Theorem~6.1, arXiv v8, p.~43]{huber2022}}

Let $E/\overline{\Q}$ have an algebraic invariant differential,
and let $u$ be an elliptic logarithm of an algebraic nontorsion
point $P$. Then
\[
 \sigma(u)\exp\left(-\frac12u\zeta(u)\right)
\]
is transcendental. If $E$ has CM and $s_2$ is normalized as in
\ref{ari:C.7.6}, then $\zeta(u)-s_2u$ is also transcendental. The exact relation of these normalized values to $N(P)$
and $\theta(u)=\sigma(u)e^{-s_2u^2/2}$ is
\eqref{ari:eq:ct-Gaussian}.
\end{cited}
\begin{proof}
The original French statement
\cite[unnumbered theorem, Section~3, p.~113; Corollary~1,
Section~4, pp.~113--114]{waldschmidt1979sigma}
applies to algebraic invariants, algebraic elliptic points
$\exp_Eu,\exp_Ev$, and $\beta\in\overline{\mathbb Q}$.
For nontorsion $\exp_Ev$ and $u,u+v\notin\Lambda$, it proves
transcendence of
\[
 \frac{\sigma(u+v)}{\sigma(u)\sigma(v)}
              e^{(\beta-\zeta(v))u}.
\]
Take $v=u$ and $\beta=0$. Since
$\sigma(2u)=-\wp'(u)\sigma(u)^4$, this value is
$-\wp'(u)\sigma(u)^2e^{-u\zeta(u)}$.
The nonzero factor $-\wp'(u)$ is algebraic, which proves
the first assertion.

For the second, use the algebraic universal vector extension
with analytic periods $(\omega,\eta(\omega))$.
On $E\setminus\{O\}$ an algebraic fibre coordinate is
$t=w-\zeta(z)$: the zeta addition formula makes its group
cocycle rational in the Weierstrass coordinates.
Near $O$, use
$t_O=t-\wp'(z)/(2\wp(z))=w+O(z^3)$.
The overlap change is rational, so $(z,w)$ are the tangent
coordinates of an algebraic chart at the identity.
If $b=\zeta(u)-s_2u$ were algebraic, the logarithm
$(u,\zeta(u)-b)$ would exponentiate to the algebraic point
with fibre coordinate $-b$ and would lie on the algebraic
line $w=s_2z$, since \cref{ari:C.7.6} gives
$s_2\in\overline{\mathbb Q}$.
The analytic subgroup theorem
\cite[Theorem~6.1, arXiv v8, p.~43]{huber2022}
would then give a one-dimensional algebraic subgroup projecting
nontrivially to $E$. Such a subgroup is an elliptic curve
mapping isogenously to $E$: a one-dimensional affine algebraic
group has no nonconstant group morphism to an elliptic curve.
Its inclusion would split the vector extension after that
isogeny. The extension class in $H^1(E,\mathcal O_E)$ is nonzero
and stays nonzero under isogeny, because trace divided by the
degree is a left inverse to pullback. Equivalently, its period
vectors cannot lie on one complex line, by Legendre's relation.
This contradiction proves the second assertion.
\end{proof}

\begin{remark}[Lift transport for the theta expression]
\hypertarget{tag-IA12}{}\label{I:theta-lift-transport}
\StatementRecord{I:theta-lift-transport}{proved}{IA12}{}

Put
$g(u,v)=\log(\sigma(u+v)/(\sigma(u)\sigma(v)))$ where the
three sigma factors are nonzero. Under $v\mapsto v+\omega$,
the change of $g(u,v)-s_2uv$, modulo $2\pi i\mathbb Z$, is
\[
 (\eta(\omega)-s_2\omega)u=\frac{\pi\bar\omega u}{A}.
\]
Both lifts represent the same elliptic point. Thus a relation
$N\exp_Ev=f(\exp_Eu)$ for a Rosati-antisymmetric endomorphism
retains its value, while the theta expression changes by this
explicit term. The pointwise criterion in
\cite[PDF slide~15, comment~(c)]{bertrand2017} uses compatible lifts.

For algebraic invariants and a CM elliptic curve,
\cite[p.~15-03]{waldschmidt1988theta} states, attributing the
result to Bertrand, that $\theta(u)=\sigma(u)e^{-s_2u^2/2}$
is transcendental at an algebraic nontorsion elliptic point.
The $s_2$ there is the regularized lattice sum. Its
quasiperiod identity is
$\eta(z)=s_2z+\pi\bar z/A$, the convention of \cref{ari:C.7.6}.
\end{remark}

\begin{remark}[The Gaussian factor and arithmetic criteria]
\hypertarget{tag-IA13}{}\label{I:Gaussian-comparison}
\StatementRecord{I:Gaussian-comparison}{proved}{IA13}{}
The exact comparison is
\begin{equation}\hypertarget{tag-003R}{}\label{ari:eq:ct-Gaussian}
 N(P)=|\Delta|^{1/12}|\theta(u)|
                         \exp\left(-\frac{\pi|u|^2}{2A}\right).
\end{equation}
It follows from \ref{ari:C.7.1} and the real-linear decomposition of
$\eta$. The normalized value in \cref{ari:C.7.7} gives the exact identity
\[
 N(P)=|\Delta|^{1/12}
 \left|\sigma(u)e^{-u\zeta(u)/2}\right|
 \exp\left(\frac12\Re\{u(\zeta(u)-\eta(u))\}\right).
\]
These identities specify the remaining exponential exactly.
The CM sampling group $\mathcal O u$ has rank two, and
$\theta'(u)/\theta(u)=\zeta(u)-s_2u$ is transcendental by
\cref{ari:C.7.7}. The arithmetic nature of the full product
in \eqref{ari:eq:ct-Gaussian} is the open question in
\cref{ari:C.7.5}.

The related independence theorem of
\citet[Corollary~2, p.~292]{reyssat1982weierstrass} says that two of
$e^{i\pi u/\omega_1}$, $s(u)$ and $s'(u)$ are algebraically
independent, where $s(z)=\sigma(z)e^{-\eta_1z^2/(2\omega_1)}$.
For real $u$, \eqref{ari:eq:cp-klein} gives
$N(P)=|\Delta|^{1/12}|s(u)|$.
This is the exact change of normalization for comparing that theorem
with the metric norm.
The periodic-functional-equation theorem reported in
\citet[p.~15-05]{waldschmidt1988theta} assumes algebraic
$q=e^{i\pi\tau}$, whereas this $q$ is transcendental for
imaginary quadratic $\tau$ by Gelfond--Schneider.  

For vector extensions,
\citet[Theorem~3.1.2, p.~386]{bost2009} proves that an algebraic point
of the universal vector extension of an abelian variety which
lies in its maximal compact subgroup at one complex embedding
is torsion. Its input is an algebraic line bundle with an algebraic connection
having unitary monodromy. The compact-character construction in
\cref{ari:C.7.2} specifies a point in a multiplicative extension.
An application of the vector-extension theorem therefore requires
construction of the stated algebraic connection.

Let $f_E$ be the Latt\`es map defined by
$x(2Q)=f_E(x(Q))$. Its canonical height is
$\widehat h_{f_E}(x(P))=2\widehat h_O(P)$. For rational
parameters \cref{ari:C.6.1} makes algebraicity of $N(P)$
equivalent to algebraicity of
$\exp(\widehat h_{f_E}(x(P)))$.
\citet[Theorem~1.4 and Corollary~1.6, arXiv v1, p.~3]{nguyen2023} settle the
corresponding multiplicative-height problem for nonintegrable
polynomials using B\"ottcher coordinates in the basin of
infinity. An elliptic
duplication Latt\`es map has no exceptional point: the inverse
images of any point contain the images of all successive
division points of an elliptic point, an infinite set. A
polynomial has the exceptional point $\infty$, so no conjugate
of $f_E$ is polynomial. The Julia set is the whole sphere
and the required basin of infinity is absent. 

The arithmetic algebraicity criterion of
\citet[Theorem~1.3, p.~3, and Theorem~2.3, p.~12, arXiv v1]{xie2024algebraicity}
requires sufficiently rapid height-normalized approximation by
points in a fixed number field and an analytic germ formally
defined over that field. This hypothesis fails
for the compact leaf in \ref{ari:C.7.9}: its tangent has the proved
transcendental slope
$c=\zeta(u)-\eta_1u/\omega_1$.
A smooth formal germ defined over $\overline{\Q}$ has an
algebraic tangent; algebraic translations, coordinate changes
and isogenies preserve this obstruction. Assuming $N(P)$
algebraic supplies an algebraic point on the leaf and its
integer multiples, but not an algebraic formal germ. Attempting
to force convergence by successive division also loses the
fixed-number-field hypothesis: if $[2]^nQ_n=Q$ with all
$Q_n\in E(K)$ for one number field, then
$\widehat h_O(Q_n)=4^{-n}\widehat h_O(Q)>0$.
The $Q_n$ are distinct with bounded ample height, contrary to
Northcott finiteness.
\end{remark}

\subsection{Ordinary nonarchimedean heights}
\begin{cited}[Nonzero ordinary heights]\hypertarget{tag-001H}{}\label{ari:C.7.8}
\StatementRecord{ari:C.7.8}{proved}{001H}{\cref{ari:C.6.R1,ari:C.6.R2}; \cite[Appendix~A, Theorem~A.1 and Lemma~A.2, p.~1043]{agboola2006}; \cite[Section~2.9, pp.~587--588; Remark~2.2, p.~588]{MazurSteinTate2006}}

Let $E/\Q$ have complex multiplication by the full ring of integers
of an imaginary quadratic field $K$ of class number one. Assume
that its functional-equation sign is $-1$, and let $p>3$ be a
prime of good ordinary reduction, $p\mathcal O_K=\mathfrak p
\mathfrak p^*$. For a nontorsion $P\in E(\Q)$ the
$\mathfrak p$-adic height and the cyclotomic $p$-adic height
of $P$ over $K$, in the normalization of
\citet[Appendix A]{agboola2006}, are nonzero.
In particular this applies to
\[
 E:y^2+y=x^3-38x+90,\qquad P=(0,9),\qquad p=5.
\]
\end{cited}
\begin{proof}
Use Theorem~A.1 and Lemma~A.2 in Rubin's appendix
\cite[Appendix~A, p.~1043]{agboola2006}, with its standing
CM, ordinary-prime and sign hypotheses. Theorem~A.1, attributed
there to Bertrand, gives nonzero $\mathfrak p$-adic height for
the nontorsion point over $K$.
The Galois-equivariance in
\cite[Section~2.9, pp.~587--588; Remark~2.2, p.~588]{MazurSteinTate2006}
makes the anticyclotomic height of a rational point zero:
complex conjugation fixes the point and acts by $-1$ on that
height space. Lemma~A.2 says that vanishing of two of the
cyclotomic, anticyclotomic and $\mathfrak p$ projections forces
the third to vanish. Hence the cyclotomic height is nonzero.

For the displayed curve, \cref{ari:C.6.R1,ari:C.6.R2} give
nontorsion and analytic rank one, hence sign $-1$.
The CM order has discriminant $-19$ and is maximal.
Its reduced positive binary forms have the single class
$(1,1,5)$, so the class number is one.
The minimal discriminant is $-19^3$.
There are seven points over $\mathbb F_5$, so $a_5=-1$,
proving good ordinary reduction. Equivalently, $-19$ is a
nonzero square modulo $5$. These verify every stated hypothesis.
\end{proof}

\begin{remark}[The compact subgroup and the CM correction]
\hypertarget{tag-IA14}{}\label{I:CM-compact-comparison}
\StatementRecord{I:CM-compact-comparison}{proved}{IA14}{}

The complex holomorphic CM correction is $\theta$ of
\cref{ari:C.7.6}; its exact relation to the real norm includes
$\exp(-\pi|u|^2/(2A))$ as in \eqref{ari:eq:ct-Gaussian}.
The height in \cref{ari:C.7.8} has the ordinary nonarchimedean
normalization specified in its cited appendix.

\citet[Theorem~2.1, p.~147]{bertrand1995} constructs nontrivial
multiplicative extensions with nontorsion algebraic points in
their maximal real compact subgroup. Its data include a
Rosati-antisymmetric endomorphism and a related nontorsion
point over a real number field.
For the diagonal class here, $NP=f(P)$ with $N>0$ and
$f+f^\dagger=0$ forces torsion: the endomorphism $N-f$ is an
isogeny because
$(N-f)^\dagger(N-f)=N^2-f^2$ is positive.
The concrete antisymmetric and diagonal extensions are
compared by the maps in \cref{ari:C.7.10}.
The vector-extension statements
\cite[Theorem~3.1, p.~150]{bertrand1995} and \cite[Theorem~3.1.2, p.~386]{bost2009}
use an algebraic connection, as described in
\cref{I:Gaussian-comparison}.
\end{remark}

\subsection{The fixed semiabelian extension and its compact slope}

\begin{proposition}[Endomorphisms and the compact slope]\hypertarget{tag-001I}{}\label{ari:C.7.9}
\StatementRecord{ari:C.7.9}{proved}{001I}{\SeriesRef{ari:C.7.1}, \SeriesRef{ari:eq:cd-compatibility}, \SeriesRef{ari:C.6.R1}, \SeriesRef{ari:C.7.2}; \citep[Theorem 6.1]{huber2022}}

Let $E/\overline\Q$ be an elliptic curve, let
$q\in E^\vee(\overline\Q)$ be nontorsion, and let
\[
 1\longrightarrow\mathbb G_m\longrightarrow G_q
       \longrightarrow E\longrightarrow0
\]
be the extension of class $q$. Then
$\operatorname{End}_{\overline\Q}(G_q)=\Z$, including when
$E$ has complex multiplication. In particular a noninteger
CM endomorphism of $E$ does not lift to an endomorphism of this
fixed extension.

For
\[
 E:y^2+y=x^3-38x+90,\qquad P=(0,9),\qquad Q=2P=(4,-2),
\]
one has
\[
 N(Q)=19^{1/4}N(P)^4.
\]
Use the differential $dx/(2y+1)$, choose the positive real
period $\omega_1$, and choose the real logarithm
$u\in(0,\omega_1/2)$ of $Q$. Put
\begin{equation}\hypertarget{tag-0033}{}\label{ari:eq:cd-slope}
 c=\zeta(u)-\frac{\eta_1}{\omega_1}u,
 \qquad \eta_1=\eta(\omega_1).
\end{equation}
Then $c$ is a real transcendental number. The point
\[
 R_N=(Q,3\,19^{-1/2}N(Q)^2)\in G_Q(\C)
\]
is nontorsion, belongs to the maximal compact subgroup, and has
the real logarithm $(u,cu)$. Its algebraicity is equivalent to
the algebraicity of $N(P)$. In the algebraic logarithm coordinates
of \ref{ari:C.7.1}, the line $\ell_c=\C(1,c)$ has transcendental
slope.
\end{proposition}


\begin{proof}
\begin{enumerate}[label=\arabic*.,leftmargin=*]
\item The torus is the maximal connected affine subgroup of $G_q$.
Consequently every endomorphism of $G_q$ induces multiplication
by some $n\in\Z$ on $\mathbb G_m$ and an endomorphism
$\alpha$ on $E$. Functoriality of extension classes gives
\begin{equation}\hypertarget{tag-0032}{}\label{ari:eq:cd-compatibility}
 \alpha^*q=nq\quad\text{in }E^\vee.
\end{equation}
Here $\alpha^*$ is the dual endomorphism. To make the
convention explicit, choose a principal polarization
$\lambda:E\to E^\vee$, put $Q_0=\lambda^{-1}q$, and
write $\alpha^\dagger=\lambda^{-1}\alpha^*\lambda$.
Equation \eqref{ari:eq:cd-compatibility} says
$(\alpha^\dagger-[n])Q_0=O$. A nonzero endomorphism of an
elliptic curve is an isogeny and has finite kernel. Since
$Q_0$ is nontorsion, $\alpha^\dagger=[n]$, hence
$\alpha=[n]$ as well.

\item The difference between the given endomorphism and $[n]$ has
image in the torus and is trivial on the torus. It therefore
factors through an algebraic-group homomorphism
$E\to\mathbb G_m$. This homomorphism is zero, since a
morphism from the proper connected curve $E$ to the affine
group $\mathbb G_m$ is constant and a group homomorphism
preserves the identity. Thus the endomorphism is exactly $[n]$. Conversely all integer
multiplications exist. This proves the first assertion.

\item For the concrete curve, the tangent slope at $P$ is $-2$,
giving $2P=(4,-2)$. Its discriminant is $-19^3$ and
$\psi_2(P)=2\cdot9+1=19$. The division identity in \cref{ari:C.4.2}
therefore gives the asserted formula for $N(Q)$, so replacing
$P$ by $Q$ preserves algebraicity of the norm. This doubling
also places the base point in the identity component of the
real elliptic group. In the
chosen differential the Weierstrass coordinates are
$(\wp,\wp')=(x,2y+1)$, with $g_2=152$, $g_3=-361$.
The negative discriminant means that the real cubic has one
real root. Along $(0,\omega_1/2)$, $\wp$ decreases from
$+\infty$ to that root and $\wp'<0$. Since
$\wp'(u)=2(-2)+1=-3$, the indicated real logarithm of $Q$
exists.

\item The period vectors of $G_Q$ in \ref{ari:C.7.1} include
\[
 (\omega_1,\zeta(u)\omega_1-\eta_1u)
       =\omega_1(1,c).
\]
These are algebraic logarithm coordinates: $dz$ and the
invariant differential $d\log r-df_u/f_u$ are algebraic,
because $df_u/f_u$ is the algebraic third-kind differential
displayed in \ref{ari:C.7.1}. Real invariants and real $u,\omega_1$
make $c$ real.

\item Suppose $c$ were algebraic. Then $\ell_c$ would be an
algebraic Lie subspace. Dividing the nonzero period
$\omega_1(1,c)$ by a sufficiently large positive integer
gives a vector in $\ell_c$ whose exponential is a nonzero
torsion point of $G_Q$, hence an algebraic point. The
analytic subgroup theorem
\citep[Theorem~6.1, arXiv v8, p.~43]{huber2022} supplies a positive-dimensional
connected algebraic subgroup with Lie algebra contained in
$\ell_c$. Its dimension is one, and its projection to $E$
is nontrivial. Such a subgroup must be an elliptic curve
$E'$ projecting isogenously to $E$: a connected
one-dimensional affine algebraic group cannot map
nontrivially to an elliptic curve. Its inclusion splits the
pullback of $G_Q$ to $E'$. The dual of that isogeny would
therefore kill the extension class $\lambda(Q)$, forcing
$Q$ torsion. This contradicts \ref{ari:C.6.R1}. We conclude that $c$
is transcendental. In particular the line is not defined over
$\overline\Q$.

\item Define the positive real number
\[
 r=3\,19^{-1/2}N(Q)^2.
\]
Proposition~\ref{ari:C.7.2} places the point $R_N=(Q,r)$
in the maximal compact subgroup of $G_Q$. Since $Q$ is algebraic,
$R_N$ is algebraic exactly when $r$ is algebraic, equivalently
when $N(Q)$, and hence $N(P)$, is algebraic. The chosen real
logarithm has
\[
 f_u(u)=-\wp'(u)\sigma(u)^2e^{-u\zeta(u)}
        =3\sigma(u)^2e^{-u\zeta(u)}>0.
\]
Consequently $v=\log(r/f_u(u))$ is real and
$(u,v)$ is a logarithm of $R_N$. Substituting this real logarithm into the exact
compactness equation $\chi=0$ in \ref{ari:C.7.2} gives
\[
 v=\zeta(u)u-u\eta(u)
   =u\left(\zeta(u)-\frac{\eta_1}{\omega_1}u\right)=cu.
\]
Its projection $Q$ is nontorsion, so $R_N$ is nontorsion.
The line $\ell_c$ containing this compact logarithm has the
transcendental slope proved above. Its CM decomposition is
\[
 c=\zeta(u)-s_2u-\frac{\pi}{A}u.
\]
The last term is the derivative of the real Gaussian
correction.
\end{enumerate}
\end{proof}

\begin{remark}[The Zariski closure of the compact-slope graph]
\hypertarget{tag-IA15}{}\label{I:compact-slope-graph}
\StatementRecord{I:compact-slope-graph}{proved}{IA15}{}

Consider
\[
 f:\C\longrightarrow\operatorname{Lie}(G_Q),\qquad
 f(z)=(z,cz).
\]
No proper connected algebraic subgroup of $G_Q$ projects
onto $E$: as in the proof above, such a subgroup would split
the nontorsion extension after an isogeny. Every proper
connected subgroup is therefore contained in the torus,
whereas $\exp_{G_Q}f$ has nonconstant elliptic projection.
Thus $f$ is $G_Q$-nondegenerate in the subgroup condition (1.2) of \citet[author preprint, p.~2]{noguchi2024}. Their Theorem~1.3 on the same page gives
\[
 \dim_{\C}\overline{\{(\exp_{G_Q}f(z),f(z)):z\in\C\}}^{\rm Zar}
       \ge \dim G_Q+1=3.
\]
The graph is contained in $G_Q\times\ell_c$, of dimension
three, so equality holds. The graph therefore has Zariski closure
$G_Q\times\ell_c$. This equality is the functional statement
used here; the algebraicity of the particular value
$R_N=\exp_{G_Q}(u,cu)$ remains the question in
\cref{ari:C.7.5}.
\end{remark}

\subsection{A CM-stable enlargement and its Gaussian eigencoordinate}

\begin{definition}[The CM enlargement coordinates]
\hypertarget{tag-IA16}{}\label{I:CM-enlargement-coordinates}
\StatementRecord{I:CM-enlargement-coordinates}{proved}{IA16}{}
Retain the curve, point $Q=2P$, and real logarithm $u$ of
\cref{ari:C.7.9}. Put
$\beta=\sqrt{-19}=i\sqrt{19}$, $d=19$,
$D=\zeta(u)-s_2u$, and $k=\pi/A$.
All endomorphisms and extensions below are defined over
$\overline{\mathbb Q}$.
\end{definition}

\begin{proposition}[The CM action and its Gaussian eigencoordinate]\hypertarget{tag-0018}{}\label{ari:C.7.10}
\StatementRecord{ari:C.7.10}{proved}{0018}{\SeriesRef{ari:C.7.9}, \SeriesRef{ari:C.7.6}, \SeriesRef{ari:eq:ce-zeta}}

Let
\[
 \mathcal G=G_Q\mathbin{\times_E}G_{\beta Q},\qquad
 1\longrightarrow\mathbb G_m^2\longrightarrow\mathcal G
       \longrightarrow E\longrightarrow0.
\]
There is an endomorphism $B$ of $\mathcal G$ inducing
$\beta$ on $E$ and
$(t_1,t_2)\mapsto(t_2^{-1},t_1^d)$ on the torus.
It satisfies $B^2=[-d]$. Moreover
\[
 \operatorname{Hom}(\mathcal G,\mathbb G_m)=0,
 \qquad
 \operatorname{End}(\mathcal G)\otimes\Q=\Q(\beta).
\]
Thus $\mathcal G$ has no algebraic $\mathbb G_m$ quotient,
including after isogeny.

There is an algebraic compact point in the second extension
above $Q$ whose logarithm is
\[
 (u,v_2),\qquad v_2=-\beta uD+au,
 \qquad
 a=\frac1{2\beta}
       \frac{d}{du}\log F_\beta(Q)\in\overline\Q.
\]
Here the derivative means differentiation of the rational
function $F_\beta(\exp_E z)$ with respect to the algebraic
invariant differential, evaluated at $z=u$.
Pair this point with the compact point $R_N$ of
\ref{ari:C.7.9}. The resulting compact point of $\mathcal G(\C)$
has logarithm
\[
 U=(u,v_1,v_2),\qquad v_1=uD-ku^2.
\]
Its algebraicity is equivalent to the algebraicity of $N(P)$.
In these algebraic logarithm coordinates the differential of
$B$ and its $-\beta$ eigenprojection are
\[
 dB=\begin{pmatrix}
 \beta&0&0\\ a&0&-1\\ \beta a&d&0
 \end{pmatrix},\qquad
 \Pi_-:=\frac12\left(I-\frac{dB}{\beta}\right),
\]
\[
 \Pi_-U=-\frac{\pi u^2}{2A}(0,1,\beta).
\]
The first torus coordinate of $\exp(\Pi_-U)$ is exactly
the Gaussian $\exp(-\pi u^2/(2A))$. The projector is not
the differential of an algebraic quasi-endomorphism. No
sequence of rational polynomials in $B$, even allowing
inverses of isogenies, can converge to it as complex linear
operators. At least one of the two torus coordinates of
$\exp(\Pi_-U)$ is transcendental.
\end{proposition}


\begin{proof}
\begin{enumerate}[label=\arabic*.,leftmargin=*]
\item Use the principal polarization to identify the two extension
classes with $Q$ and $\beta Q$. They are linearly independent
over $\Q$ modulo torsion: any nonzero $m+n\beta$ is an
isogeny, so $(m+n\beta)Q$ cannot be torsion. A character
of $\mathbb G_m^2$ given by $(m,n)\in\Z^2$ extends to
$\mathcal G$ exactly when the pushout extension has zero
class, namely when $mQ+n\beta Q=O$. Hence the restriction
of every homomorphism $\mathcal G\to\mathbb G_m$ to
the torus is zero. Such a homomorphism then factors through
$E$ and is zero. This proves the asserted absence of
characters. An isogenous semiabelian variety cannot acquire
a nonzero character, because composing it with the isogeny
would give a nonzero character of $\mathcal G$.

\item The Rosati adjoint of $\beta$ is $-\beta$. Pullback by
$\beta$ sends the ordered extension classes to
$(-\beta Q,dQ)$, exactly as does pushout by the displayed
torus matrix. Equality of these extension classes provides
the endomorphism $B$. It is unique because
$\operatorname{Hom}(E,\mathbb G_m^2)=0$. Its square
induces $[-d]$ on the torus and on $E$, so the same
uniqueness gives $B^2=[-d]$. The map
\[
 \operatorname{End}(\mathcal G)\otimes\Q
     \longrightarrow\operatorname{End}(E)\otimes\Q
          =\Q(\beta)
\]
is injective: an endomorphism inducing zero on $E$ has image
in the torus, and the absence of characters makes it zero.
Its image contains $\beta$, hence equals $\Q(\beta)$.

\item The identities in \ref{ari:C.7.6}, applied at $\beta$
and $1+\beta$, give the algebraic nonzero number
\begin{equation}\hypertarget{tag-0036}{}\label{ari:eq:ce-ribet-r}
 r_2=\frac{\theta((1+\beta)u)}{\theta(u)\theta(\beta u)},
 \qquad r_2^2=\frac{F_{1+\beta}(Q)}{F_\beta(Q)}.
\end{equation}
Indeed $(1+\beta)(1+\bar\beta)=1+d$, so the powers
of $\theta(u)$ cancel. None of the arguments is a lattice
point since $Q$ is nontorsion. Written in sigma functions,
\[
 r_2=\frac{\sigma((1+\beta)u)}{\sigma(u)\sigma(\beta u)}
                         e^{-s_2\beta u^2}.
\]
The point $(Q,r_2)$ in the rational trivialization of
$G_{\beta Q}$ therefore has logarithm
\[
 (u,v_2),\qquad
 v_2=u\{\zeta(\beta u)-s_2\beta u\}.
\]
Differentiating \ref{ari:C.7.6} shows that
\begin{equation}\hypertarget{tag-0037}{}\label{ari:eq:ce-zeta}
 \zeta(\beta u)-s_2\beta u=-\beta D+a,
 \qquad a\in\overline\Q,
\end{equation}
with the value of $a$ in the statement. The compactness
character of this extension vanishes, since
\[
 \Re\{v_2-\zeta(\beta u)u+\beta u\eta(u)\}
    =\Re\left(\frac{\beta\pi|u|^2}{A}\right)=0.
\]

\item In the algebraic coordinates, the torus block of $dB$ is the stated integer matrix and its first entry is
$\beta$, so initially write its first column as
$(\beta,a_1,a_2)^t$, with $a_1,a_2$ algebraic.
The identity $B^2=[-d]$ gives $a_2=\beta a_1$.
For $\omega\in\Lambda$, set
\[
 \kappa_1(\omega)=\zeta(u)\omega-\eta(\omega)u,
 \quad
 \kappa_2(\omega)=\zeta(\beta u)\omega-\eta(\omega)\beta u.
\]
Equation \eqref{ari:eq:ce-zeta} and the real-linear formula for
$\eta$ imply
\[
 \kappa_1(\beta\omega)=a\omega-\kappa_2(\omega).
\]
Comparing the second coordinates of the corresponding period
vectors under $B$ gives
$(a_1-a)\omega_j\in2\pi i\Z$ for $j=1,2$.
Since $\omega_2/\omega_1$ is nonreal, both integers must
be zero, so $a_1=a$. This proves the matrix formula.

\item The first compact logarithm in \ref{ari:C.7.9} is
$v_1=uD-ku^2$. Pairing the two compact lifts over the same base
point gives $U$. The second lift
is algebraic, so the paired point is algebraic exactly when its
first lift is algebraic.
Substitution into the displayed matrix yields
\[
 v_2+\beta v_1-au=-\beta ku^2,
 \qquad
 \Pi_-U=-\frac{ku^2}{2}(0,1,\beta),
\]
which proves the exact Gaussian identity.

\item The two eigenspaces of $dB$ have eigenvalues $\beta$ and
$-\beta$. Every algebraic quasi-endomorphism has differential
$a_0I+b_0dB$ for rational $a_0,b_0$, with eigenvalues
$a_0+b_0\beta$ and $a_0-b_0\beta$. The projector
$\Pi_-$ would require these values to be $0$ and $1$,
which is impossible. There is also a uniform obstruction to
approximation. For arbitrary real $a_0,b_0$, put
$x=a_0+b_0\beta$, so that its other eigenvalue is
$\bar x$. Then
\begin{equation}\hypertarget{tag-0035}{}\label{ari:eq:ce-gap}
 \max\{|x|,|\bar x-1|\}\ge\frac12,
\end{equation}
because $1\le|\bar x|+|\bar x-1|=|x|+|\bar x-1|$.
These are exactly the two eigenvalue errors in approximating
$\Pi_-$. The spectral radius is bounded by every induced
operator norm, so no such sequence can converge to the
projector. Compositions and inverses of isogenies still belong
to the same endomorphism algebra and satisfy this bound.

\item Finally the image of $\Pi_-$ is the toric Lie line
$\C(0,1,\beta)$. It is not the Lie algebra of an algebraic
subtorus: an algebraic one-dimensional subtorus of
$\mathbb G_m^2$ has rational slope. Since $u\ne0$,
the projected vector is nonzero. The analytic subgroup
theorem in fact shows that its exponential cannot be an
algebraic point of both torus coordinates. Indeed its nonzero logarithm lies in the algebraic line
$\C(1,\beta)$, whereas a one-dimensional algebraic subtorus
has a rational slope. The analytic subgroup theorem would put
a positive-dimensional subtorus in that line for an algebraic
exponential, contradicting these two slope computations.

\item At the ordinary prime $5$, the CM number $\alpha=(1+\beta)/2$ satisfies
$\alpha^2-\alpha+5=0$. Its two archimedean conjugates
both have absolute value $\sqrt5$. At the two places over
$5$, their normalized valuations are $0$ and $1$, because
the polynomial reduces to $X(X-1)$ and has product $5$.
The separation available in that nonarchimedean linear
algebra is absent in \eqref{ari:eq:ce-gap}.

\item \textbf{An exact compact witness on the conductor-$361$ curve.}
In the preceding normalization the constants and the second
fibre coordinate are explicitly
\begin{equation}\hypertarget{tag-0034}{}\label{ari:eq:ce-explicit-witness}
 s_2=-2,\qquad a=-\frac{21\beta}{703},\qquad
 r_2=-\frac{165}{74}+\frac{1839\beta}{1406}.
\end{equation}
Put $f=4x^3-152x+361$ and
\[
 \begin{aligned}
 h={}&x^9-38x^8+437x^7-1444x^6-7942x^5\\
     &+82308x^4-274360x^3+390963x^2-130321x-130321,\\
 N={}&(-19x+76)h^2+f(h''h-(h')^2)+(6x^2-76)h'h,\\
 K={}&N'h-2Nh'.
 \end{aligned}
\]
Direct expansion gives
\[
 -K^2f=4N^3-152\cdot19^2Nh^4+361\cdot19^3h^6,
 \qquad \gcd(N,h)=1.
\]
Hence, in coordinates $Y=2y+1$, the map
\[
 (x,Y)\longmapsto
 \left(\frac{N}{19h^2},\frac{KY}{19\beta h^3}\right)
\]
is a degree-$19$ endomorphism fixing $O$, with pullback of
$dx/Y$ equal to $\beta\,dx/Y$. It is therefore the CM
endomorphism $\beta$. Divisors and the leading coefficient
at $O$ give $F_\beta=\beta^2h^2$. Its Laurent coefficient
of relative degree two is $-76$. Comparison with \ref{ari:C.7.6}
gives $38s_2=-76$, proving $s_2=-2$.

\item For the other kernel put $h_+=(H_0+\beta H_1)/10$, where
\[
 \begin{aligned}
 H_0={}&10x^8-380x^7+4864x^6-28880x^5+72561x^4\\
       &+54872x^3-795644x^2+1824494x-1433531,\\
 H_1={}&20x^7-304x^6+760x^5+10469x^4-83752x^3\\
       &+246924x^2-315514x+130321.
 \end{aligned}
\]
Writing $\bar h_+$ for replacement of $\beta$ by $-\beta$,
the finite polynomial identities and coprimality checks are
\[
 \begin{gathered}
 19xh^2-N=5f h_+\bar h_+,\qquad
 h_+\mid K+19\beta h^3,\\
 \gcd(h_+,fh)=\gcd(h_+,h_+')=1.
 \end{gathered}
\]
At a root of $h_+$ the displayed CM map sends $(x,Y)$ to
$(x,-Y)$, so these are precisely the eight pairs of
non-$2$-torsion points in $\ker(1+\beta)$. The other three
nonzero kernel points are $E[2]\setminus\{O\}$, since
$1+\beta=2\alpha$ in the CM order. Its norm is $20$,
which verifies the count. Thus
\[
 F_{1+\beta}=\frac{(1+\beta)^2}{4}f h_+^2.
\]
Taking the square ratio in \eqref{ari:eq:ce-ribet-r} and fixing
the sign by its expansion at $u=0$ gives the exact rational
function
\[
 r_2=-\frac{1+\beta}{2\beta}\frac{\wp'(u)h_+(\wp(u))}
                                             {h(\wp(u))}.
\]
At $Q$, the needed evaluations are
\[
 h(4)=-37,\qquad h'(4)=7,\qquad
 h_+(4)=\frac{829-279\beta}{10},\qquad\wp'(u)=-3.
\]
They give exactly \eqref{ari:eq:ce-explicit-witness}. In
particular $(Q,r_2)$ is an explicit nontorsion algebraic
point in the maximal compact subgroup of $G_{\beta Q}$.
Its torus-coordinate norm is
$r_2\bar r_2=974799/26011=4-\wp(\beta u)$.

\end{enumerate}
\end{proof}

\begin{remark}[A real period basis for the CM witness]
\hypertarget{tag-IA17}{}\label{I:CM-real-period-basis}
\StatementRecord{I:CM-real-period-basis}{proved}{IA17}{}

Put $\tau_0=(1+i\sqrt{19})/2$. The integral basis change
$(1,\tau_0)\mapsto(1-2\tau_0,1-\tau_0)$ has determinant one.
Applied to the rotated basis $i(1,\tau_0)$, it gives the ratio
\[
 \tau=\frac{1+i/\sqrt{19}}2,\qquad
 \omega_1=\left(\frac{g_2(\tau)}{152}\right)^{1/4}>0.
\]
The logarithm $u$ is selected by
$0<u<\omega_1/2$, $\wp(u)=4$, and $\wp'(u)=-3$.
This fixes the analytic branch of the algebraic witness in
\cref{ari:C.7.10}.
\end{remark}

\subsection{Complexifying the Gaussian: values, jets, and the rank criterion}

\begin{proposition}[The complexified norm and its second derivatives]\hypertarget{tag-0019}{}\label{ari:C.7.11}
\StatementRecord{ari:C.7.11}{proved}{0019}{\SeriesRef{ari:C.7.9}, \SeriesRef{ari:C.7.1}, \SeriesRef{ari:C.7.6}, \SeriesRef{ari:eq:cc-fixedfield}, \SeriesRef{ari:C.7.7}; \citep[Theorem 6.1]{huber2022}}

Use the curve and conjugation of
\cref{I:CM-enlargement-coordinates}.

Use the curve, point and real logarithm of \ref{ari:C.7.9}:
\[
 E:y^2+y=x^3-38x+90,\qquad Q=(4,-2)=\exp_Eu,
 \qquad 0<u<\omega_1/2.
\]
Put $\beta=i\sqrt{19}$, $k=\pi/A$, $s_2=-2$, and
\[
 \theta(z)=\sigma(z)e^{-s_2z^2/2},\qquad
 D(z)=\zeta(z)-s_2z,\qquad
 \mathcal F(z,w)=\sqrt{19}\,\theta(z)\theta(w)e^{-kzw},
 \qquad H=\mathcal F^2.
\]
Then $\mathcal F(z,\bar z)=N(z)^2$, and $\mathcal F$ is an
entire function of two complex variables of order at most two.
It is periodic under $(z,w)\mapsto(z+\omega,w+\bar\omega)$,
$\omega\in\Lambda$. The subgroup
\[
 \Gamma=\{(\alpha u+\omega,\bar\alpha u+\bar\omega):
                   \alpha\in\Z[\beta],\ \omega\in\Lambda\}
 \subset\C^2
\]
has rank four over $\Z$ and lies in the real plane
$\{(z,\bar z):z\in\C\}$. There is one number field $K$ such
that, with $H_0=H(u,u)=N(Q)^4$, every value of $H$ on
$\Gamma$ belongs to $K[H_0]$. More precisely, the nonzero values
are algebraic multiples of the powers $H_0^{\alpha\bar\alpha}$,
as specified in \eqref{ari:eq:cc-fixedfield}.

For every nonzero $\alpha\in\Z[\beta]$ there is an algebraic
number $a_\alpha$ such that
\[
 D(\alpha u)=\bar\alpha D(u)+a_\alpha.
\]
At $p_\alpha=(\alpha u,\bar\alpha u)$ the algebraic constant
direction $V_\alpha=\alpha\partial_z-\bar\alpha\partial_w$
satisfies
\begin{align*}
 V_\alpha\log\mathcal F(p_\alpha)
    &=\alpha a_\alpha-\bar\alpha\bar a_\alpha
                         \in\overline{\Q},\\
 V_\alpha^2\log\mathcal F(p_\alpha)
    &=-\alpha^2\{\wp(\alpha u)+s_2\}
      -\bar\alpha^2\{\wp(\bar\alpha u)+s_2\}
      +2\alpha\bar\alpha k.
\end{align*}
The number $k$ is transcendental. For every nonzero $\alpha$,
the exact normalized second derivative satisfies
\[
 \overline\Q\!\left(\frac{V_\alpha^2H(p_\alpha)}{H(p_\alpha)}\right)
      =\overline\Q(k).
\] In particular,
\[
 (\partial_z-\partial_w)H(u,u)=0,\qquad
 \frac{(\partial_z-\partial_w)^2H(u,u)}{H(u,u)}=4k-8.
\]
Thus the normalized first derivative is algebraic and the
normalized second derivative generates $\overline\Q(k)$. The rank-four sampling set and growth order give the numerical
comparison with \eqref{ari:eq:cc-rankcriterion} proved below.
\end{proposition}



\begin{proof}
\begin{enumerate}[label=\arabic*.,leftmargin=*]
\item The real model and its algebraic differential give real sigma
coefficients and a conjugation-stable lattice. In particular
$\overline{\theta(z)}=\theta(\bar z)$. The discriminant is
$-19^3$, so \ref{ari:C.7.1} gives the asserted diagonal identity.
The sigma function and the exponential of the quadratic
polynomial $-kzw$ have order at most two. From \ref{ari:C.7.6},
\[
 \theta(z+\omega)=\epsilon(\omega)
         e^{k\bar\omega(z+\omega/2)}\theta(z),
 \qquad \epsilon(\bar\omega)=\epsilon(\omega)\in\{1,-1\}.
\]
Multiplying this formula and its conjugate version, the two
signs multiply to one and the exponent cancels the change of
$-kzw$. This proves periodicity on all of $\C^2$.

\item If $\alpha u+\omega=0$, then $[\alpha]Q=O$.
Nontorsion of $Q$ and the fact that a nonzero elliptic
endomorphism is an isogeny imply $\alpha=0$, and then
$\omega=0$. Thus the rank-two module $\Z[\beta]u$ and the
rank-two period lattice are independent over $\Z$.
This proves the rank assertion.

\item For $d=\alpha\bar\alpha$, the squared CM law \ref{ari:C.7.6} gives
the exact fixed-field identity
\begin{equation}\hypertarget{tag-002Q}{}\label{ari:eq:cc-fixedfield}
 H(\alpha u,\bar\alpha u)
   =19^{1-d}F_\alpha(Q)\overline{F_\alpha(Q)}H(u,u)^d.
\end{equation}
All the endomorphisms and their normalized rational functions
$F_\alpha$ are defined over a single number field: choose
a field containing the coefficients of the finitely many
generators of the CM order, and close it under conjugation.
The uniqueness of the Laurent normalization in \ref{ari:C.7.6} gives
the same field of definition for every $F_\alpha$.
Writing that fixed number field as $K$, formula
\eqref{ari:eq:cc-fixedfield} places all values in $K[H_0]$.
The value for $\alpha=0$ is zero, and period shifts do not
change it. Squaring was included because algebraicity of the
individual positive square roots alone would not establish
their membership in one fixed number field.

\item Logarithmically differentiating the squared CM law at $u$
gives
\[
 a_\alpha=\frac{1}{2\alpha}
     \frac{(d/dz)F_\alpha(\exp_Ez)|_{z=u}}{F_\alpha(Q)}
                  \in\overline{\Q}.
\]
The denominator is nonzero since $Q$ is nontorsion. The
derivation $d/dz$ is the algebraic invariant derivation on
$E$, so the displayed derivative has algebraic value.
Since $D(u)$ is real, conjugation gives the corresponding
formula at $\bar\alpha u$. On a neighborhood where
$\mathcal F\ne0$, its logarithmic gradient and Hessian are
\label{layout:I:3}
\[
\begin{gathered}
(\partial_z\log\mathcal F,\partial_w\log\mathcal F)
       =(D(z)-kw,D(w)-kz),\qquad
 \\
\operatorname{Hess}\log\mathcal F=
 \begin{pmatrix}-\wp(z)-s_2&-k\\-k&-\wp(w)-s_2\end{pmatrix}.
\end{gathered}
\]
Evaluation on $(\alpha,-\bar\alpha)$ proves both identities
in the statement. In particular the Gaussian cancels in the
first derivative and reappears in the second derivative.

\item To prove $k$ transcendental, recall the real period formula
$\eta_1=(s_2+k)\omega_1$. If $k$ were algebraic, the
nonzero period vector $(\omega_1,\eta_1)$ of the algebraic
universal vector extension of $E$ would lie on the algebraic
line of slope $s_2+k$. Dividing this period by a sufficiently
large integer supplies a nonzero algebraic torsion exponential.
The analytic subgroup theorem
\citep[Theorem~6.1, arXiv v8, p.~43]{huber2022} would give a one-dimensional
algebraic subgroup with that Lie line. Its projection to $E$ is nonzero, so it splits the vector
extension after an isogeny. The nonsplitting argument in
\cref{ari:C.7.7} excludes this: the coordinates here are
exactly $(z,w)$ with periods $(\omega,\eta(\omega))$ and
algebraic charts $t=w-\zeta(z)$,
$t_O=t-\wp'(z)/(2\wp(z))$.

\item Write the two logarithmic derivatives as
$V_\alpha\log\mathcal F=B_\alpha$ and
$V_\alpha^2\log\mathcal F=A_\alpha+2dk$, where
$A_\alpha,B_\alpha$ are algebraic. Since $H=\mathcal F^2$,
\[
 \frac{V_\alpha^2H(p_\alpha)}{H(p_\alpha)}
       =4B_\alpha^2+2A_\alpha+4dk.
\]
Here $d>0$, so the displayed affine expression both belongs to
$\overline\Q(k)$ and recovers $k$ by subtraction of its algebraic
constant and division by $4d$. This proves the asserted field
equality; $H(p_\alpha)$ is nonzero because $Q$ is nontorsion.
At $\alpha=1$, $a_1=0$, $\wp(u)=4$ and $s_2=-2$,
which give the last displayed formula in the statement.

\item \textbf{Differential and growth criteria.}
The differential criterion of
\citet[Theorems~3.6--3.7, author proof, pp.~21--22]{waldschmidt2009auxiliary}
requires a finite algebraically independent family whose
partial derivatives belong to the polynomial ring generated
by that family over one number field. Algebraic values of
the full family then give algebraic values of all iterated
partial derivatives. The calculation above gives the exact value-field obstruction:
the field containing $H(p_\alpha)$ and its second directional
derivative contains $k$ as soon as the nonzero value is inverted.
In particular its normalized second jet is transcendental. The issue persists
although the first cancelling direction is algebraic; the
direction also varies with $\alpha$.

\item Theorem~1 of
\cite[p.~24]{waldschmidt1975dimension} uses $d$ meromorphic functions on $\C^n$, growth and
arithmetic-size exponents $\rho_1,\ldots,\rho_d$, and sampling
sets of cardinality $N^{\ell-o(1)}$ in a ball of radius $O(N)$,
with separation at least $N^{-\lambda-o(1)}$. Its strict numerical
requirement is
\begin{equation}\hypertarget{tag-002X}{}\label{ari:eq:cc-rankcriterion}
 (d-1)\ell>\sum_{j=1}^d\rho_j
                     +2d(n-1)(\lambda+1).
\end{equation}
The theorem also requires all function values to belong to
one number field and bounds on denominators in the entire
representations of the meromorphic functions. 

\item The natural independent triple here is
$\wp(z),\wp(w),H(z,w)$, all of order at most two.
The first two are algebraically independent. If $H$ satisfied
a polynomial over $\C(\wp(z),\wp(w))$, fix a generic pair
$(z,w)$ and translate only $z$ by $m\omega$, with
$0\ne\omega\in\Lambda$. The coefficients of the polynomial
remain constant, whereas the absolute values of its proposed
roots $H(z+m\omega,w)$ grow as
$\exp(k|\omega|^2m^2+O(|m|))$. This contradicts the finite
set of roots of a nonzero polynomial. Thus the triple is
independent.

\item The growth-order lower envelope is $\rho_j=2$. Substitution of
this envelope and $n=2$, $d=3$, $\ell=4$ into
\eqref{ari:eq:cc-rankcriterion} gives the required inequality
\[
 8>12+6\lambda,
\]
which fails even for $\lambda=0$. Moreover, the $\asymp N^4$
points obtained by bounding
the four integer coefficients lie in a ball of radius $O(N)$
in a real plane. Packing separated discs in that plane gives
$N^4\ll N^{2+2\lambda+o(1)}$, hence $\lambda\ge1$.
Removing the rank-two set of poles does not change this count.
Larger growth exponents only worsen the inequality. Thus the stated rank and growth envelopes do not satisfy
\eqref{ari:eq:cc-rankcriterion}. The norm question is
recorded in \cref{ari:C.7.5}.
\end{enumerate}
\end{proof}

\subsection{Division values and interpolation bounds}


\begin{proposition}[Division ratios and translated functions]\hypertarget{tag-001A}{}\label{ari:C.7.12}
\StatementRecord{ari:C.7.12}{proved}{001A}{\SeriesRef{ari:C.4.2}, \SeriesRef{ari:C.7.1}, \SeriesRef{ari:C.7.6}, \SeriesRef{ari:C.7.9}, \SeriesRef{ari:C.7.11}}

Let $E:y^2+y=x^3-38x+90$, $Q=(4,-2)$ and $N_0=N(Q)$.
Use the positive real logarithm $u$ and period $\omega_1$ of
\ref{ari:C.7.9}. Put
\[
 \delta=19^{1/4},\qquad b=\eta_1/\omega_1,\qquad
 g(z)=\delta\sigma(z)e^{-bz^2/2}.
\]
Thus $|g(z)|=N(\exp_Ez)$ on the real axis. Let $\mathcal D$
be the CM division hull of $Q$. For $R\in\mathcal D$, the
equations $[n]R=[\alpha]Q$, $n>0$, determine a unique
$\kappa_R=\alpha/n\in\operatorname{End}(E)\otimes\Q$.
Set $w_R=\kappa_R\bar\kappa_R$. For $R\ne O$,
\[
 A_R:=\frac{N(R)}{N_0^{w_R}}\in\overline\Q_{>0}.
\]
In particular, for $[n]R=[m]Q$ with $m\ne0$,
\begin{equation}\hypertarget{tag-003S}{}\label{ari:eq:division-ratio-model}
 A_R^{n^2}=\delta^{n^2-m^2}
                  \frac{|\psi_m(Q)|}{|\psi_n(R)|}.
\end{equation}
For each positive algebraic number $\lambda$, define
\[
 N_\lambda(R)=A_R\lambda^{w_R}
       =N(R)\left(\frac{\lambda}{N_0}\right)^{w_R}.
\]
All these values are algebraic. For $[n]R=[m]Q$, $|m|\le n$,
$R\ne O$, there are constants $D,H$ depending on $\lambda,E,Q$
such that
\[
 h(N_\lambda(R))\le H,\qquad
 [\Q(R,N_\lambda(R)):\Q]\le
 \begin{cases}Dn^4,&R\in E(\R),\\2Dn^6,&R\in E(\C).
 \end{cases}
\]
The transformation is exact for every positive real parameter:
$N_\lambda(R)/N_\mu(R)=(\lambda/\mu)^{w_R}$ and
$N_{N_0}(R)=N(R)$.

For any finite collection of real division-hull logarithms $a$,
the functions $g(z)$ and $g(z+a)/g(a)$, with $\exp_Ea\ne O$,
have function-field transcendence degree at most three over
$\overline\Q$ after adjoining the elliptic coordinates.
The three functions $g(z),g(z+v)/g(v),g(z-v)/g(v)$, $v=u/q$,
attain three. The normalized holomorphic extensions
$F_a(z)/N(a)^2$ of $N(\exp_E(z+a))^2$ for fixed complex
division-hull shifts belong to the same type of field.
The unnormalized translates have the same upper bound, attained by
the displayed triple, over the algebraic closure $L$ of $\overline\Q(N_0)$ in $\C$.

For the scalar models on the real division grids
$z=(mu+k\omega_1)/n$, $|m|\le Cn$, $k$ modulo $n$, with
fixed $C$, every fixed finite family of these values has a common
degree envelope $O(n^5)$ and bounded heights. For the values at $\lambda=N_0$, the same construction
gives a relative degree envelope over a fixed finite extension
of $\Q(N_0)$. The algebraicity of $N_0$ is the open question in
\cref{ari:C.7.5}. 
\end{proposition}

\begin{proof}\leavevmode\par


 
\begin{enumerate}[label=\arabic*.,leftmargin=0pt,labelsep=.5em]
\item \emph{The algebraic ratios and their exact change of parameter.}
Two equations $[n]R=[\alpha]Q$ and $[n']R=[\alpha']Q$
give $[n'\alpha-n\alpha']Q=O$. A nonzero elliptic endomorphism
has finite kernel, whereas $Q$ is nontorsion. Thus
$\alpha/n=\alpha'/n'$, proving uniqueness. The resulting
map $R\mapsto\kappa_R$ is additive, and its kernel is the
torsion subgroup. Every such $R$ is algebraic, since it lies
in a finite isogeny fibre above an algebraic point.


\begin{enumerate}[label=(\alph*),leftmargin=1.3em,labelsep=.5em]
\item  The division identity and the squared CM law give, for
$d=\alpha\bar\alpha$ and $\alpha\ne0$,
\[
 N(R)^{n^2}=\frac{N(nR)\delta^{n^2-1}}{|\psi_n(R)|},
 \qquad
 N(\alpha Q)^2=\delta^{2-2d}|F_\alpha(Q)|N_0^{2d}.
\]
Consequently
\[
 A_R^{2n^2}=\frac{\delta^{2n^2-2d}|F_\alpha(Q)|}
                         {|\psi_n(R)|^2}\in\overline\Q_{>0}.
\]
The positive root proves $A_R\in\overline\Q$. For integral
$\alpha=m$, the usual multiplication identity is
$N(mQ)=N_0^{m^2}|\psi_m(Q)|/\delta^{m^2-1}$;
substitution proves \eqref{ari:eq:division-ratio-model}.
For nonzero $n$-torsion $R$, take the limit at the simple zero
of $\psi_n$ in the uniformizing coordinate. Since
$N(z)=\delta|z|+o(|z|)$ at zero,
\[
 A_R^{n^2}=N(R)^{n^2}
       =\frac{n\delta^{n^2}}{|\partial_z\psi_n(R)|}.
\]
The invariant derivation is algebraic, so this also is algebraic.
All exponents $w_R$ are rational. The asserted scalar models,
their parameter-change identity and the specialization at $N_0$
follow directly. In particular their definition is independent
of every presentation of $R$ by a division equation.


\item  Put $K_\lambda=\Q(\delta,\lambda)$ and
$D=[K_\lambda:\Q]$. For integral $m\ne0$,
\[
 N_\lambda(mQ)=\lambda^{m^2}
                     |\psi_m(Q)|/\delta^{m^2-1}\in K_\lambda,
\]
because $Q$ is rational. The fibre of $[n]$ has $n^2$ points,
so adjoining $R$ costs at most $n^2$. For real $R$ the absolute
value of $\psi_n(R)$ differs from it by a sign. The division
equation for $N_\lambda(R)$ costs at most another $n^2$.
For complex $R$, adjoin $R$ and $\bar R$, at cost at most
$n^4$, and use
\[
 N_\lambda(R)^{2n^2}
 =\frac{N_\lambda(mQ)^2\delta^{2(n^2-1)}}
              {\psi_n(R)\psi_n(\bar R)}.
\]
This gives $2Dn^6$. At nonzero torsion points the derivative
formula proves the same bounds.

 
\end{enumerate}

\item \emph{Uniform heights.}
The following coefficient estimate proves the height bound.  On the fixed integral model, $\psi_n$ is a polynomial
in $x$ for odd $n$, and $(2y+1)$ times a polynomial in $x$
for even $n$.  These polynomials have degree $O(n^2)$.
Choose a fixed circle $|x|=A>1$ avoiding the finite branch values.
On its compact inverse image in $E(\C)$, $N$ has a positive
minimum $m_A$, and $|2y+1|$ has a positive minimum.  Let $M$
be the maximum of $N$ on the whole compact elliptic curve.
The rearranged division identity gives
\[
 |\psi_n(R)|\leq M\delta^{n^2-1}m_A^{-n^2}
       \qquad(|x(R)|=A).
\]
Cauchy's coefficient estimate, and division by $2y+1$ for
even $n$, show that the logarithm of the coefficient
$\ell^1$ norm is $O_E(n^2)$.  The same assertion holds for
$\partial_z\psi_n$, because differentiation and the fixed
Weierstrass equation add only polynomial factors in $n$.


\begin{enumerate}[label=(\alph*),leftmargin=1.3em,labelsep=.5em]
\item  The elementary height estimate for a fixed degree-four
duplication map, followed by telescoping, gives
$h(x(R))=2\widehat h(R)+O_E(1)$ uniformly over
$E(\overline\Q)$.  Since
$\widehat h(R)=(m/n)^2\widehat h(Q)$, the coordinate heights
are uniformly bounded when $|m|\leq n$.
At each nonarchimedean place the ultrametric inequality bounds a
polynomial value by its largest monomial; at each archimedean
place use the sum of the monomials. Summing these logarithmic
bounds by the product formula gives the polynomial evaluation
height inequality, and hence
$h(\psi_n(R))=O_{E,Q}(n^2)$, and the same estimate for its
invariant derivative at torsion points.  Also
\[
 h(N_\lambda(mQ))\leq m^2h(\lambda)+O_{E,Q}(m^2).
\]
Apply the same division identity to $N_\lambda$ and divide its height by $n^2$.
For a complex value one uses
$h(|\alpha|)\leq h(\alpha)$, which follows by squaring and
using $|\alpha|^2=\alpha\bar\alpha$.
This proves a uniform bound for $h(N_\lambda(R))$.

 
\end{enumerate}

\item \emph{Normalized translates and their algebraic scalars.}
For a real logarithm $a$ of a nonzero division point put
\[
 h_a(z)=\frac{\sigma(z+a)}{\sigma(z)\sigma(a)}e^{-baz}.
\]
The exact identities are
\[
 g(z+a)=\delta^{-1}g(z)g(a)h_a(z),\qquad
 h_a(z+\omega_1)=h_a(z),\quad
 h_a(z+\omega_2)=e^{-2\pi ia/\omega_1}h_a(z).
\]
They follow by substitution from sigma automorphy and Legendre's
relation. A real point in the CM division hull lies in the
ordinary division hull: conjugating $[n]R=[\alpha]Q$ gives
$[\alpha-\bar\alpha]Q=O$, hence $\alpha=\bar\alpha\in\Z$.
The real group is connected because the discriminant is negative.
Its division logarithms are therefore exactly
$\widetilde{\mathcal D}=\Q u+\Q\omega_1$.
The coefficient map
\[
 \nu:\widetilde{\mathcal D}\longrightarrow\Q,
 \qquad \nu(ru+s\omega_1)=r
\]
is well defined and additive: a rational relation between $u$
and $\omega_1$ would make $Q$ torsion.


\begin{enumerate}[label=(\alph*),leftmargin=1.3em,labelsep=.5em]
\item  For any positive real $\lambda$ define, on this additive group,
\[
 \widetilde g_\lambda(z)
   =g(z)(\lambda/N_0)^{\nu(z)^2},\qquad
 h_{a,\lambda}(z)
   =h_a(z)(\lambda/N_0)^{2\nu(z)\nu(a)}.
\]
At nonzero division values, the absolute value of
$\widetilde g_\lambda$ is $N_\lambda$ and its sign is the real
sign of $g$. Thus these values are algebraic for algebraic
$\lambda$. The identity
\[
 h_{a,\lambda}(z)=
 \frac{\delta\widetilde g_\lambda(z+a)}
      {\widetilde g_\lambda(z)\widetilde g_\lambda(a)}
\]
is proved by expanding $(\nu(z)+\nu(a))^2$.
It specifies the exact transformation on the sampling group.


\item  Choose a common denominator $q$ for finitely many shifts,
write $a_j=(m_ju+k_j\omega_1)/q$, and put $v=u/q$.
The quotient $h_{a_j}/h_v^{m_j}$ is elliptic for
$\Lambda_q=\Z\omega_1+\Z q\omega_2$. This finite isogeny
of $E$ has a model over $\overline\Q$ with the stated
differential. The quotient has an algebraic divisor, hence
is a rational function over $\overline\Q$ times a scalar.
The scalar is algebraic. Indeed the parameter factors cancel:
\[
 \frac{h_{a_j}(z)}{h_v(z)^{m_j}}
  =\frac{h_{a_j,1}(z)}{h_{v,1}(z)^{m_j}}
       \qquad(z\in\widetilde{\mathcal D}).
\]
Choose a real division point avoiding the finite zeros and poles.
The right side and the value of the chosen rational function
are algebraic and nonzero, so their quotient is algebraic.
This proves that $g(z)$ and all normalized real translates are
algebraic over $\overline\Q(E)(g,h_v)$, after the fixed finite
isogeny. The upper transcendence-degree bound is three.


\item  For the lower bound, work even over $\C$. The multiplier of
$h_v$ under $q\omega_2$ is $e^{-2\pi iu/\omega_1}$,
which has infinite order. A polynomial relation over the elliptic
field, evaluated at a generic $z$ and translated by $mq\omega_2$,
would have infinitely many distinct roots obtained by these
powers. Thus $h_v$ is transcendental over that field. Moreover
\[
 |g(z+mq\omega_2)|=\exp(c m^2+O(|m|)),
 \qquad c=\pi q^2\Im(\omega_2/\omega_1)>0.
\]
In a proposed polynomial relation for $g$ over the elliptic field
and $h_v$, clear denominators. The nonzero leading coefficient
is a polynomial evaluated along an irrational rotation; it is
bounded away from zero on a subsequence, and all other
coefficients remain bounded there. Quadratic growth makes the
highest power of $g$ dominate, a contradiction. Finally
\[
 \wp(z)=\wp(v)-
 \frac{\delta^2g(z+v)g(z-v)}{g(z)^2g(v)^2}
\]
recovers the elliptic coordinate from the three stated translates,
and $h_v=\delta g(z+v)/(g(z)g(v))$. They therefore attain three.
The constants $g(a)$ lie in $L$, by the already proved ratios.
Multiplication by these constants gives the exact map from
normalized to raw translates, proving the assertion over $L$.


\item  For a complex CM division shift $a$, the holomorphic extension
from the real axis is
\[
 F_a(z)=\frac{N(a)^2}{\delta^2}g(z)^2h_a(z)h_{\bar a}(z).
\]
Its multiplier depends on $a+\bar a\in\Q u+\Q\omega_1$.
Choose $q,m$ with $a+\bar a=(mu+k\omega_1)/q$.
Then $h_ah_{\bar a}/h_v^m$ is elliptic on $\C/\Lambda_q$
with algebraic divisor. Its scalar is algebraic by the same
evaluation argument, using the exact identity on real division
points
\[
 h_a(z)h_{\bar a}(z)(\lambda/N_0)^{2\nu(z)m/q}
 =\frac{\delta^2N_\lambda(\exp_E(z+a))^2}
        {N_\lambda(\exp_Ea)^2\widetilde g_\lambda(z)^2}.
\]
Here $w_{\exp_E(z+a)}=|\kappa_{\exp_Ea}+\nu(z)|^2$;
expanding the square proves the identity. Set $\lambda=1$
and divide by $h_{v,1}^m$. This proves the normalized complex
translate assertion with all scalar factors accounted for.

 
\end{enumerate}

\item \emph{Joint degrees and heights on the sampling group.}
In the algebraic meromorphic section coordinate of
\ref{ari:C.7.1},
\[
 z\longmapsto(\exp_Ez,h_v(z))
   =\exp_{G_V}\bigl(z,(\zeta(v)-bv)z\bigr),\qquad
 V=\exp_Ev.
\]
For the scalar model, multiply its torus coordinate by
$(\lambda/N_0)^{2\nu(z)\nu(v)}$. This factor is a character
of the additive sampling group, so the resulting map
$\Phi_\lambda:\widetilde{\mathcal D}\to G_V(\C)$ is an
abstract group homomorphism. Its real period is $\omega_1$,
because $\nu(\omega_1)=0$ and the displayed analytic map
has that period. For algebraic $\lambda$, its values in the displayed section
chart are algebraic by the formula for $h_{v,\lambda}$ above.
For an exceptional sampling point at a zero or pole of that
section, write $z=z_1+z_2$ with both $z_j$ avoiding the finite
exceptional set modulo the elliptic lattice. Such a choice exists
in the infinite division group. Then
$\Phi_\lambda(z)=\Phi_\lambda(z_1)\Phi_\lambda(z_2)$ is
algebraic by the algebraic group law. This also covers the identity.


\begin{enumerate}[label=(\alph*),leftmargin=1.3em,labelsep=.5em]
\item  For $z=(mu+k\omega_1)/n$,
$[n]\Phi_\lambda(z)=[m]\Phi_\lambda(u)$.
Choose a fixed number field containing the coefficients of
$G_V$, the point $\Phi_\lambda(u)$ and the finitely many
fixed constants. Multiplication by $n$ on $G_V$ has degree
$n^3$, since its kernel is an extension of $E[n]$ by
$\mathbb G_m[n]$. Adjoining $\widetilde g_\lambda(z)$ costs
at most $n^2$ by the norm division equation. Passing to the
fixed isogeny and adjoining fixed constants changes only the
constant in $O(n^5)$. The same argument for any positive real
$\lambda$ works over a fixed finite extension of $\Q(\lambda)$,
because all the fixed division values are algebraic over that
field. At $\lambda=N_0$ it gives the claimed relative envelope.


\item  The height proof above applies with $|m|\le Cn$ for any fixed
$C$, changing only its constant. It bounds the model values at
$z$ and every fixed real translate $z+a$. Their quotients bound
$h_{v,\lambda}(z)$. Coordinates on the fixed isogenous elliptic
curve have bounded height by the fixed isogeny height comparison.
The normalized complex translates are fixed rational functions
of those coordinates, $\widetilde g_\lambda^2$ and a fixed power
of $h_{v,\lambda}$, with algebraic coefficients just proved.
The polynomial evaluation inequality bounds their heights.
Any removable zero or pole in these expressions lies in a fixed
finite set and is handled by its specified model value. Taking the
positive square root for a fixed complex-shifted norm divides
height by two and costs at most a factor two in degree.

 
\end{enumerate}

\Needspace{6\baselineskip}
\item \emph{Interpolation inequalities and algebraic tangents.}\par
The sufficient inequality in Gramain--Mignotte--Waldschmidt
\citep[Theorem~2.1, pp.~101--102]{gramain1986values} has the form
\[
 d(N)\prod_{j=1}^k\max\{d(N)\psi_j(N),\theta_j(R(N))\}
 \le c_0\mu(N)^{k-1}\varphi(N)^k,
\]
\[
 \varphi(N)=\log\frac{R(N)^2+r(N)r(N+1)}
                         {R(N)(r(N)+r(N+1))}.
\]
On fixed nested discs the growth and geometric factors are
constant. The scalar-model height envelopes are constant too.
The numerical comparison of these envelopes is therefore
$d^{k+1}\ll\mu^{k-1}$. At denominator $n$ the real group
has $O(n^2)$ sample points, and through denominator $N$ it
has $O(N^3)$. Substitution of the proved envelopes gives
$N^{12}$ against $N^3$ for two functions and $N^{20}$ against
$N^6$ for three. These upper envelopes give the exponent comparisons
$12>3$ and $20>6$. Verifying the sufficient inequality for these grids
requires bounds sharper than these envelopes. For samples at $\lambda=N_0$, the base field is the
displayed extension of $\Q(N_0)$; its number-field status is
exactly the unresolved algebraicity issue. Extra normalized
translates remain in the function field of transcendence degree
three proved above.


\begin{enumerate}[label=(\alph*),leftmargin=1.3em,labelsep=.5em]
\item  The arithmetic-germ criteria
\citep[Theorem~6.1, pp.~31--32]{herblot2012meromorphic}
\citep[Theorem~4.16, author preprint, p.~25]{gasbarri2010analytic}
can be compared with the logarithmic tangent. At a real
nontorsion division point $z$, the exact formula is
\[
 \frac1{g(z)}\frac{dg}{dx}
       =\frac{\zeta(z)-bz}{\wp'(z)}.
\]
Its numerator is transcendental by the period-line proof of
\ref{ari:C.7.9}, applied to this algebraic nontorsion point;
its denominator is nonzero algebraic. A formal graph defined
over $\overline\Q$ would have an algebraic nonzero constant
value and algebraic first derivative, contradicting their
displayed ratio. The scalar change on the sampling group has
not changed this exact tangent of the analytic graph.


\item  For the Mahler system in
\citep[Theorem~1.1 and system~(1.2), p.~56]{adamczewski2017mahler},
the needed algebraic Taylor coefficients and algebraic regular
argument can be tested directly: the coefficient of $z^3$ in
$g$ is $-\delta b/2$, and $b=s_2+k$ is transcendental by
\ref{ari:C.7.11}. The logarithm $u$ is transcendental too:
an algebraic $u\ne0$ would give an algebraic point $(u,Q)$
on the diagonal Lie line of $\mathbb G_a\times E$; the analytic
subgroup theorem would produce an algebraic one-dimensional
subgroup projecting nontrivially to both factors, which is
impossible since an elliptic curve has no nonconstant map to
$\mathbb G_a$ and $\mathbb G_a$ has no nonzero group map to $E$.
The duplication identity acts by the explicit elliptic
multiplication map. The entropy identity in
\citep[Theorem~5.1, author preprint, p.~9]{einsiedler2001entropy} uses the sequence
built from the denominators of $x(2^nQ)$ with scale $4^n$.


\item  Finally the period-refined rank comparison
\citep[Theorem~4.1, pp.~385--386, and Corollary~4.3, p.~386]{waldschmidt1988periodrank}
\citep[main theorem, p.~363]{roy1993periodrank}
has coefficient $\alpha(G)=2\dim G-2d_0(G)-d_1(G)$,
where the last terms count additive and multiplicative
quotients. Both vanish for $G_Q$. An additive homomorphism
vanishes on its torus and then on the proper quotient $E$;
a nonzero torus character extending to $G_Q$ would kill a
nonzero integer multiple of its nontorsion extension class.
Thus $\alpha(G_Q)=4$.


\item  The line $\ell_c$ has period rank one. A general period is
$(\omega,\zeta(u)\omega-\eta(\omega)u+2\pi i j)$.
Writing $\omega=a\omega_1+d\omega_2$, membership in
$\ell_c$ gives $du/\omega_1\in\Z$ by Legendre's relation.
Since $u/\omega_1$ is irrational, $d=j=0$. The numerical
rank bound in the cited theorem is
$(4-1)\dim\ell_c=3$. The period and the logarithm
$(u,cu)$ are independent over $\Q$ and span rank two;
the algebraicity of the latter's exponential remains exactly
the point-algebraicity question in \ref{ari:C.7.9}. \qedhere


\end{enumerate}

\end{enumerate}
\end{proof}


\clearpage\appendix
\section{Parameter equivalence and deformations}\label{pub:I:parameter-appendix}

\begin{remark}[Period and deformation comparison]\hypertarget{tag-02YS}{}\label{night:P5.7}
\StatementRecord{night:P5.7}{proved}{02YS}{\SeriesRef{night:P5.3}, \SeriesRef{night:P5.4}, \SeriesRef{night:P5.5}, \SeriesRef{night:P5.6}; \cite[Theorems 1.3--1.4, p.~3; Theorem 6.12, pp.~81--82; Proposition 6.16, p.~85; Theorem 7.14, p.~95]{DuEtAl2026v1}; \cite[title, abstract, contents and introduction, pp.~1--3]{DuEtAl2026v2}}
For the fixed standard $(3,4,\infty)$ data, the tangent
cohomology $(1,1,1,0)$, smooth one-dimensional Kuranishi germ,
automorphism group $\mathbb C^*$, and equivalence
$X_c\simeq X_{c'}\Longleftrightarrow c'-c\in2\mathbb Z$
agree with the corresponding assertions of
\cite{DuEtAl2026v1}. Version~2, entitled
\emph{Oka Geometry of Complex Six-Spheres}, omits those
tangent-cohomology, miniversality, full-automorphism and exact-period
assertions. Version~1, submitted on
23 September 2026, states the exact $2\mathbb Z$ equivalence as its Theorem~1.4.
\end{remark}
\begin{proof}
\begin{enumerate}[label=\arabic*.]
\item Use the same oriented triangle, ordered integral basis
$(\widehat\gamma,\widehat u,\widehat w,\widehat\delta)$,
finite twists, and untranslated unit-square cusp fan as version~1.
Its chosen function $\beta_0$ and our $\beta^{\rm part}$ obey the
same affine transformation laws, and both $\beta_0+\tau$ and
$\beta^{\rm part}+\tau$ extend over the cusp. Their difference
descends to a holomorphic function on $\mathbb P^1$, hence is the
constant $\delta=\beta^{\rm part}-\beta_0$. The exact parameter
map and inverse are $a=c+\delta$ and $c=a-\delta$: substitution
gives $\Omega_{c+\delta}=\Pi_c$ in the stated ordered basis.
Writing $M_0=\sup(\Im\beta_0-6(\Im\mu)^2/\Im\tau)$,
our corresponding bound is $M_{\rm part}=M_0+\Im\delta$, so this
map also takes our admissible half-plane bijectively to version~1's.
The ordinary quotients, finite logarithmic maps, and cusp charts
then identify by their identical marked periods and specified
gluing data. In particular $a'-a=c'-c$. Its Theorem~1.4 therefore
gives the displayed exact-period equivalence and the full
automorphism group in this standard-data scope.
\item Theorem~6.12 of version~1 gives the tangent-cohomology tuple,
and Proposition~6.16 gives the smooth reduced Kuranishi germ and
the locally miniversal parameter family. These are the
four assertions compared with \cref{night:P5.3,night:P5.4,night:P5.5,night:P5.6}
under that same parameter and marking.
\item The second version retains the period construction and the
Oka conclusions. Its displayed section list and its mathematical
statements omit the first version's logarithmic-deformation and
classification sections and the four assertions just listed.
\end{enumerate}
\end{proof}

\subsection{The integral obstruction}\hypertarget{tag-01PB}{}\label{fam:section:A.2}\hypertarget{tag-03CR}{}\label{wb:pshift:sec:classification}

\begin{proposition}[The integral obstruction to a parameter shift]\hypertarget{tag-01Q1}{}\label{fam:src:A.2.1}\label{fam:st:A.2.1}\hypertarget{tag-03CQ}{}\label{wb:pshift:lem:centralizer}
\StatementRecord{fam:src:A.2.1}{proved}{01Q1}{\SeriesRef{fam:src:A.1.2}}

The rational centralizer is
\begin{equation}
 \operatorname{Cent}_{\mathbb Q}(A_1,A_2)
       =\{aI+bE:a,b\in\mathbb Q\}.
 \hypertarget{tag-03BZ}{}\label{wb:pshift:eq:centralizer}
\end{equation}
A biholomorphism over the identity of $B^\circ$ between the torus
families forces $c'-c\in\mathbb Z$ and has complex-linear part
$I_2$ or $-I_2$ on every fibre. Every biholomorphism $X_c\to X_{c'}$ has $c'-c=n\in2\mathbb Z$, complex-linear fibre part $I_2$, and source-to-target lattice map $I-nE$. In the full fixed lattices,
\[
 \Lambda^{A_1}/N_1\Lambda\simeq\mathbb Z/3,\qquad
 \Lambda^{A_2}/N_2\Lambda\simeq\mathbb Z/4\oplus\mathbb Z/2,
\]
with $d$ generating the second factor of the latter group and vanishing in the former.
\end{proposition}
\begin{proof}
\begin{enumerate}[label=\arabic*.,leftmargin=*]
\item For an integer $n$, put $S_n=I+nE$. Since $S_n^{-1}=S_{-n}$,
\begin{equation}
 \Pi_{c+n}=\Pi_c S_n,\qquad
 \Pi_{c+n}\Lambda=\Pi_c\Lambda.
 \hypertarget{tag-03CF}{}\label{wb:pshift:eq:integer-shear}
\end{equation}
Both monodromy matrices fix $\widehat\delta$ and preserve $\gamma$.
Consequently $A_jE=EA_j=E$, and the same commutation holds for every
monodromy word. The identity of $\mathbb C^2$ therefore identifies
the two full middle torus families with their $R_g$-action.
Its source-to-target lattice map is $S_{-n}$, since
$\Pi_{c+n}S_{-n}=\Pi_c$.
\item Solving the two fixed-vector systems from \eqref{wb:pshift:eq:finite-data} gives
\[
 (\Lambda\otimes\mathbb Q)^{A_1}
   =\mathbb Q(1,2,-4,0)^{\mathsf t}+\mathbb Q\widehat\delta,
 \quad
 (\Lambda\otimes\mathbb Q)^{A_2}
   =\mathbb Q(1,3,-3,0)^{\mathsf t}+\mathbb Q\widehat\delta.
\]
Their intersection is $\mathbb Q\widehat\delta$. The fixed covectors
are respectively
\[
 \mathbb Q\gamma+\mathbb Q(2u+w+3\delta),\qquad
 \mathbb Q\gamma+\mathbb Q(u+w+2\delta).
\]
Their intersection is $\mathbb Q\gamma$: comparing the $u$ and $w$
coefficients of a common multiple gives $2x=y$ and $x=y$, hence
$x=y=0$. Every commuting endomorphism therefore preserves
$\mathbb Q\widehat\delta$ and $\ker\gamma$.

\item On $\ker\gamma/\mathbb Q\widehat\delta$, in the classes of
$(\widehat u,\widehat w)$, the two matrices are
\[
 Q_1=\begin{pmatrix}0&1\\-1&-1\end{pmatrix},\qquad
 Q_2=\begin{pmatrix}0&-1\\1&0\end{pmatrix}.
\]
A matrix commuting with $Q_2$ has the form
$\left(\begin{smallmatrix}b&d\\-d&b\end{smallmatrix}\right)$.
Commutation with $Q_1$ gives $d=0$. Thus a commuting endomorphism $H$
has columns
\[
 H\widehat\gamma=(a,r,s,k)^{\mathsf t},\quad
 H\widehat u=b\widehat u+u_0\widehat\delta,\quad
 H\widehat w=b\widehat w+w_0\widehat\delta,\quad
 H\widehat\delta=d_0\widehat\delta.
\]
Use the columns $A_1\widehat u=-\widehat w+\widehat\delta$,
$A_1\widehat w=\widehat u-\widehat w$, and
$A_2\widehat u=\widehat w$. Commutation on these columns gives
\[
 d_0-w_0=b+u_0,\qquad u_0=2w_0,\qquad w_0=u_0.
\]
Hence $u_0=w_0=0$ and $d_0=b$. Commutation on
$\widehat\gamma$ with $A_1$ gives
\[
 r=2(a-b),\qquad s=-4(a-b).
\]
Commutation on that column with $A_2$ also gives $r=-s$.
Therefore $a=b$, $r=s=0$, and $H=aI+kE$. Conversely these matrices
commute, by the fixed-vector and fixed-covector identities preceding
\eqref{wb:pshift:eq:integer-shear}. This proves \eqref{wb:pshift:eq:centralizer}.

\item For the geometric assertion, a holomorphic map between two compact
complex tori lifts to their universal covers. The derivative of the
lift is periodic, so its entries descend to holomorphic functions on
the compact source torus and are constant. The lift is consequently
affine, with a complex-linear part $L$. In a holomorphic family these
linear parts vary holomorphically. Their lattice maps are integral
and locally constant, hence give an isomorphism of the two marked
local systems. Because the map is over the identity of the base,
that integral map $H$ commutes with both $A_j$. Thus
$H=aI+bE$, with $a,b\in\mathbb Z$. Its determinant is $a^4$;
invertibility over $\mathbb Z$ gives $a=1$ or $a=-1$.

\item The complete period equation is
\begin{equation}
                   L(z)\Pi_c(z)=\Pi_{c'}(z)(aI+bE).
 \hypertarget{tag-03CM}{}\label{wb:pshift:eq:period-morphism}
\end{equation}
Its third and fourth columns give $L(z)=aI_2$. Its first column
then gives $a(c'-c)+b=0$. Therefore $c'-c=-b/a\in\mathbb Z$,
and the asserted two linear parts exhaust the possibilities.
\item \emph{The finite-root obstruction.}
The two cyclic norm matrices are, by multiplication of \eqref{wb:pshift:eq:finite-data},
\[
 N_1=I+A_1+A_1^2=
 \begin{pmatrix}3&0&0&0\\6&0&0&0\\-12&0&0&0\\0&2&1&3\end{pmatrix},
\]
\begin{equation}
 N_2=I+A_2+A_2^2+A_2^3=
 \begin{pmatrix}4&0&0&0\\12&0&0&0\\-12&0&0&0\\0&2&2&4\end{pmatrix}.
 \hypertarget{tag-03CK}{}\label{wb:pshift:eq:norms}
\end{equation}
For example these formulas give
$N_1\widehat w=\widehat\delta$,
$N_2\widehat u=N_2\widehat w=2\widehat\delta$,
$N_1\widehat\delta=3\widehat\delta$, and
$N_2\widehat\delta=4\widehat\delta$.

The ramified base change
$t_j=u^{m_j}$ has local equation
\[
 s_j^{m_j}-u^{m_j}
   =\prod_{a=0}^{m_j-1}(s_j-e^{2\pi ia/m_j}u)=0.
\]
Its normalization separates all these smooth branches, each with
coordinates $(z_1,z_2,u)$. The marked map
$y\mapsto(q_j(y),s_j(y))$ identifies the marked torus family $J'_j$
with this normalization, as proved by the explicit branches in
Proposition~\ref{fam:src:A.1.2}. A biholomorphism over the identity of
$D_j$ base-changes over the identity of the $u$-disc and lifts uniquely
to the normalization: on local rings it carries integral elements in
the total fraction ring to integral elements, and its inverse does
the same. This lift commutes with the marked deck action because
both composites lift the same map $(x,u)\mapsto(F(x),\zeta_j u)$.

Let $c'=c+n$, where $n$ is an integer, and suppose the linear part of
the lifted map is $aI_2$, $a\in\{1,-1\}$. The lifted map on the torus
family is $[z,\zeta]\mapsto[z,a\zeta+B_j(z)]$. Its translating
section has a holomorphic lift $B_j:\Delta_j\to\mathbb C^2$:
$\Delta_j\times\mathbb C^2\to J'_j$ is a covering map and the disc is
simply connected. Local inverses of this covering show that the lift
is holomorphic. The deck-equivariance identity holds modulo a
lattice vector, so, after choosing this lift, it is
\begin{equation}
 B_j(g_jz)-R_{g_j}(z)B_j(z)
  =\Pi_c(g_jz)\left(
       \frac{(1-a)v_j+n\epsilon_j\widehat\delta}{m_j}
          +\ell_j\right),\qquad \ell_j\in\Lambda.
 \hypertarget{tag-03C9}{}\label{wb:pshift:eq:finite-equation}
\end{equation}
The integer column $\ell_j$ is constant: it is locally constant by
discreteness of the marked lattice and the disc is connected.
Iterating $m_j$ times, with $A_jv_j=v_j$ and
$A_j\widehat\delta=\widehat\delta$, yields the necessary equality
\begin{equation}
 N_j\ell_j=-(1-a)v_j-n\epsilon_j\widehat\delta.
 \hypertarget{tag-03CB}{}\label{wb:pshift:eq:finite-norm-obstruction}
\end{equation}
Here the left side of the iterated equation is zero because the
base point and the marked vector coordinate return after $m_j$
steps; real invertibility of $\Pi_c$ makes the lattice equality exact.

If $a=-1$, applying $\gamma$ to \eqref{wb:pshift:eq:finite-norm-obstruction} gives
\[
                   m_j\gamma(\ell_j)=-2\epsilon_j.
\]
This is impossible already at $j=1$, since $3$ does not divide $2$.
Thus a global isomorphism cannot have negative linear part.

For $a=1$, the first row of \eqref{wb:pshift:eq:finite-norm-obstruction} gives $\gamma(\ell_j)=0$.
Equation \eqref{wb:pshift:eq:norms} then proves that the order-three filling imposes no
restriction on $n$, whereas the order-four filling imposes precisely
$n\in2\mathbb Z$. More intrinsically, the finite obstruction is the
class of $-n\epsilon_j\widehat\delta$ in
\begin{equation}
 \Lambda^{A_j}/N_j\Lambda.
 \hypertarget{tag-03C7}{}\label{wb:pshift:eq:finite-class}
\end{equation}
In the marked fixed bases this quotient is
$\mathbb Z/3$ at $j=1$, with $\widehat\delta$ zero, and
$\mathbb Z/4\oplus\mathbb Z/2$ at $j=2$, with
$\widehat\delta$ generating the second summand. Indeed \eqref{wb:pshift:eq:norms} gives
\[
 N_1\Lambda=3\mathbb Z v_1\oplus\mathbb Z\widehat\delta,
 \quad
 N_2\Lambda=4\mathbb Z\varepsilon'\oplus
                      2\mathbb Z\widehat\delta,
 \qquad \varepsilon'=(1,3,-3,0)^{\mathsf t}.
\]
This proves the parity obstruction.
\end{enumerate}
\end{proof}

\subsection{The effective period parameter}
\hypertarget{tag-03CV}{}\label{wb:sec:retained-period-deformation}
\begin{proposition}[Exact period variation and the Kodaira--Spencer class]
\hypertarget{tag-01Q4}{}\label{fam:src:A.4.2}\label{fam:st:A.4.2}
\hypertarget{tag-03DP}{}\label{wb:thm:retained-period-deformation}
\StatementRecord{fam:src:A.4.2}{proved}{01Q4}{\cref{fam:src:A.1.1,fam:src:A.3.1}}
Use the proper family of \cref{fam:src:A.3.1} over a disc
$Q$ with $\overline Q\subset\mathcal P$.
Fix a regular base point and one of its lifts $z$, and fix
$c_*\in Q$. With all functions below evaluated at that $z$, set
\begin{equation}
 T=\operatorname{Im}\tau>0,\quad m=\operatorname{Im}\mu,\quad
 D=\operatorname{Im}\beta_{c_*}-\frac{6m^2}{T}<0,\quad a=\frac mT,
 \qquad P=\begin{pmatrix}0&0\\-a&1\end{pmatrix}.
\label{pub:I:eq:PD3}
\end{equation}
For $\delta=c-c_*$, the exact marked real-linear map from the universal
cover of the fibre at $c_*$ to that at $c$ is
\begin{equation}
 w=A_\delta\zeta+B_\delta\overline\zeta,\qquad
 A_\delta=I+\frac{\delta}{2iD}P,\qquad
 B_\delta=-\frac{\delta}{2iD}P.
\label{pub:I:eq:PD4}
\end{equation}
It maps the marked period lattice at $c_*$ onto that at $c$, and
therefore descends to a real-analytic marked diffeomorphism of the two
tori. In the explicit convention that the transported $(0,1)$ tangent
space is spanned by
$\partial_{\overline\zeta_k}+\sum_j\nu^j_k\partial_{\zeta_j}$,
its exact coefficient matrix and its first derivative are
\begin{equation}
 \nu_\delta=\frac{\delta}{2iD+\delta}P,\qquad
 \dot\nu_z=\frac{1}{2iD}\,
 \partial_{\zeta_2}\otimes
 (d\overline\zeta_2-a\,d\overline\zeta_1).
\label{pub:I:eq:PD5}
\end{equation}
The class of $\dot\nu_z$ in
$H^{0,1}(f_{c_*}^{-1}(t(z)),T^{1,0})$ is nonzero.
With the tangent-exact-sequence convention
$\operatorname{KS}(\partial_c)=[\bar\partial h]$ for a smooth
$(1,0)$ lift $h$ of $\partial_c$, the fibre class is
$\operatorname{KS}_z(\partial_c)=-[\dot\nu_z]$.
Moreover the first-order deformation of the entire $X_{c_*}$ obtained
from \cref{fam:src:A.3.1} is nontrivial: its Kodaira--Spencer class in
$H^1(X_{c_*},T_{X_{c_*}})$ is nonzero. In particular
$h^1(X_{c_*},T_{X_{c_*}})\geq1$.
\end{proposition}
\begin{proof}
\begin{enumerate}[label=\arabic*.,leftmargin=*]
\item \emph{The exact variation on a regular fibre.}
Write $u=(u_1,u_2,u_3,u_4)^t$ in the ordered lattice basis and
$\zeta=\Pi(z,c_*)u$. Taking imaginary parts gives the two exact equations
\begin{equation}
 \operatorname{Im}\zeta_1=6m u_1+T u_2,\qquad
 \operatorname{Im}\zeta_2=(D+6m^2/T)u_1+m u_2.
\label{pub:I:eq:PD13}
\end{equation}
Subtracting $a$ times the first equation from the second yields
\begin{equation}
 u_1=\frac{\operatorname{Im}\zeta_2-a\operatorname{Im}\zeta_1}{D}.
\label{pub:I:eq:PD14}
\end{equation}
In the changed periods the same $u$ maps to
$w=\Pi(z,c)u=\zeta+\delta e_2u_1$. Substitution of
$\operatorname{Im}\zeta=(\zeta-\overline\zeta)/(2i)$ proves \eqref{pub:I:eq:PD4}.
This also proves its asserted lattice map and real invertibility:
it is exactly $\Pi(z,c)\Pi(z,c_*)^{-1}$ as a real-linear map.
Its real determinant is $(D+\operatorname{Im}\delta)/D>0$.


\item Here $P^2=P$, and
\begin{equation}
 \det A_\delta=1+\frac\delta{2iD},\qquad
 A_\delta^{-1}=I-\frac{\delta}{2iD+\delta}P.
\label{pub:I:eq:PD15}
\end{equation}
The denominator can vanish only for $\delta=-2iD$. At that parameter
$D+\operatorname{Im}\delta=-D>0$, which is outside the admissible
domain. Thus this inverse is defined at every $c\in Q$.
Annihilating $w$ by the declared new $(0,1)$ fields gives
$A_\delta\nu_\delta+B_\delta=0$, which proves \eqref{pub:I:eq:PD5}, including
its sign. All its coefficients are constant in the fibre coordinates.
Thus its first derivative is a globally defined $\bar\partial$-closed
tangent-valued $(0,1)$ form on the torus.

\item It cannot be $\bar\partial V$ for a globally defined smooth vector
field $V$. Indeed the invariant holomorphic frame
$\partial_{\zeta_1},\partial_{\zeta_2}$ writes $V$ with periodic
smooth coefficients. The integral of each of their derivatives
$\partial_{\overline\zeta_k}V^j$ against the translation-invariant
probability volume is zero: real constant vector fields preserve this
volume, and integration of their derivatives on the compact torus is
zero. But the coefficient of
$\partial_{\zeta_2}\otimes d\overline\zeta_2$ in $\dot\nu_z$ is
the nonzero constant $1/(2iD)$. Its integral is that same constant.
This proves that the Dolbeault class is nonzero.

\item The conversion to the tangent-exact-sequence convention retains a
minus sign. In the moving coordinates the real-marking lift has
$(1,0)$ part $h=\partial_c+u_1\partial_{w_2}$ at $c_*$.
Since \eqref{pub:I:eq:PD14} gives
\[
 \bar\partial u_1=-\frac{d\overline\zeta_2-a\,d\overline\zeta_1}{2iD},
\]
the connecting class $[\bar\partial h]$ is exactly $-[\dot\nu_z]$.
Thus the two conventions are related by the displayed sign and both give a nonzero class.

\item \emph{The intrinsic pencil in the parameter family.}
Let $\mathcal S_2\subset\mathcal X_Q$ be the reduced finite divisor
over $p_2\times Q$, so $F^*(p_2\times Q)=4\mathcal S_2$.
The proof of \cref{wb:thm:canonical-ring-recovery} works for relative
forms over $Q$, as can be checked on each chart.
On the middle family the frame is $d\zeta_1\wedge d\zeta_2$ relative
to $B\times Q$. At the cusp the relative canonical frame over $Q$ is
$(dx_1/x_1)\wedge(dx_2/x_2)\wedge dt_c$.
Extra $dc$ terms in transition differentials vanish in relative forms.
The finite determinant characters and the holomorphic units in
the canonical calculation of \cref{fam:src:A.1.1} are independent
of $c$, since $R_g,\tau,\mu$ are.
Thus their coefficient identities hold jointly over $Q$.
The evaluation exponents remain $m_j-1-a_j=0,2$.
It follows by these transition functions and divisors that
\begin{equation}
 \omega_{\mathcal X_Q/Q}\simeq\mathcal O_{\mathcal X_Q}(-2\mathcal S_2),
 \qquad
 \omega_{\mathcal X_Q/Q}^{-2}\simeq
          F^*\operatorname{pr}_B^*\mathcal O_B(p_2).
\label{pub:I:eq:PD16}
\end{equation}
Here the base identifications retain $h=1/(t-1)$ and
$\eta_B=dt/(t-1)^2$ from \eqref{pub:I:eq:CR1}; these functions also do not depend on $c$.
The joint transition identities prove \eqref{pub:I:eq:PD16} over the whole parameter disc.

\item For its complete degree-two system over the dual numbers, put
$A=\mathbb C[\epsilon]/(\epsilon^2)$, base-change by
$c=c_*+\epsilon$, write $X_A$ for the resulting deformation and
$L_A=\omega_{X_A/A}^{-2}$. For $L_*=\omega_{X_{c_*}}^{-2}$,
flatness of the deformation and local freeness of $L_A$ give
\begin{equation}
 0\longrightarrow L_*\xrightarrow{\ \epsilon\ }L_A
        \longrightarrow L_*\longrightarrow0.
\label{pub:I:eq:PD17}
\end{equation}
The two sections of $\mathcal O_B(p_2)$ used in \eqref{pub:I:eq:CR4}, namely
$\sigma_{p_2}$ and $\ell=A_0/(t-1)+B_0'$ with $A_0\ne0$,
pull back by \eqref{pub:I:eq:PD16} to sections $s_0,s_1$ of $L_A$ whose reductions
are a basis of $H^0(X_{c_*},L_*)$. This basis assertion is exactly
the proved degree-two case of \eqref{pub:I:eq:CR2}--\eqref{pub:I:eq:CR4}.
Any section of $L_A$ reduces to a complex linear combination of these
two reductions; subtract that lifted combination. Exactness of \eqref{pub:I:eq:PD17}
then writes the remainder as $\epsilon$ times a section of $L_*$,
again a combination of that basis. Conversely, if an $A$-linear
combination of $s_0,s_1$ is zero, reduce modulo $\epsilon$ and then
use the injective first arrow of \eqref{pub:I:eq:PD17} to see that all four complex
coefficients vanish. Therefore
\begin{equation}
 H^0(X_A,L_A)=A s_0\oplus A s_1.
\label{pub:I:eq:PD18}
\end{equation}
The system is basepoint-free and its morphism is $f_A$ followed by
the fixed coordinate map $t\mapsto[t-1:A_0+B_0'(t-1)]$.
This is the base-change isomorphism for the complete degree-two system.

\item \emph{Nontriviality for the complete threefold.}
Suppose its first-order deformation were trivial, with an isomorphism
$X_A\simeq X_{c_*}\times\operatorname{Spec}A$ reducing to the identity.
Intrinsic relative canonical bundles and their complete spaces of
sections are preserved by that isomorphism. By \eqref{pub:I:eq:PD18}, it would
therefore identify $f_A$ with $a_A\circ f_{c_*}$ for a projective
linear change $a_A\in\operatorname{PGL}_2(A)$ reducing to the identity.

\item This change fixes all three distinguished points even to first order.
In the finite covering charts
$f_A$ is $t_j=s_j^{m_j}$, so its relative critical ideal is
$(s_j^{m_j-1})$, and its restriction to that critical scheme has
$t_j=0$. The map from the parameter ring $A$ to this critical
scheme is injective; it contains a constant-parameter section given
by $s_j=0$ and the zero section of the covering torus.
At the cusp, in a chart containing a triple point,
$t_c=z_0z_1z_2$ and the critical ideal is
$(z_1z_2,z_0z_2,z_0z_1)$. Again $t_c=0$ in its ring and $A$ injects,
using $z_0=z_1=z_2=0$. At a double chart the same assertion follows
from $t_c=z_0z_1$; a smooth chart has no critical point.
These critical ideals are intrinsic to the differential of the map
over $A$, and the charts are local biholomorphisms onto $X_A$.
Away from the three distinguished base values the map is a
submersion, so it has no critical point.
Consequently the scheme-theoretic critical-value sections are exactly
the three constant sections $p_1,p_2,p_0$. An automorphism reducing
to the identity cannot permute these three distinct closed points.
Its value at each is therefore the same point over $A$, not a
nonzero infinitesimal displacement: such a displacement would be a
nonzero element of $\epsilon\mathbb C$ in an injective copy of $A$
on the just displayed critical chart.

\item Write $a_A(t)=t+\epsilon(b_0+b_1t+b_2t^2)$.
Fixing $0$ gives $b_0=0$, fixing $1$ gives $b_0+b_1+b_2=0$,
and fixing infinity, in $u=1/t$, gives $b_2=0$.
Hence $a_A$ is the identity. The proposed trivialization thus
preserves $f$, and restricting it to any regular base value would
trivialize the torus deformation \eqref{pub:I:eq:PD4}. Its nonzero class \eqref{pub:I:eq:PD5}
excludes this.

\item Finally, the zero Kodaira--Spencer class is precisely the condition for
such a first-order trivialization. One may see the needed direction
directly: choose local product charts for the holomorphic submersion
$\pi$ and differentiate their transition functions at $c_*$. The
resulting tangent-valued overlap functions form a \v Cech $1$-cocycle.
If its class vanishes, after refinement it is the coboundary of local
holomorphic vector fields. Changing each chart by that vector field
times $\epsilon$ cancels the first-order overlap terms, producing
the identity-reducing trivialization. This is the usual
Kodaira--Spencer connecting class (for the standard deformation-space
interpretation see \cite[Section~67, p.~15]{McMullen1996Complex});
the displayed conversion
from the plus-graph convention in \eqref{pub:I:eq:PD5} changes its sign, not its
vanishing. The contradiction proves the final assertion.
\end{enumerate}
\end{proof}



\section{Tangent cohomology and deformations}\label{pub:I:deformation-appendix}
\subsection{Tangent direct images and local miniversality}
\hypertarget{tag-000D}{}\label{app:rederive:fam-src-A.5.2}

\hypertarget{tag-02ZF}{}\label{push:section:E.2}
\begin{remark}[The standard deformation theorem]
\hypertarget{tag-PI11}{}\label{I:standard-deformation-credit}
\StatementRecord{I:standard-deformation-credit}{proved}{PI11}{}
For the standard data $(\ell_0,\ell_1,\ell_2)=(0,1,-1)$,
the tangent dimensions and smooth reduced Kuranishi germ are given in
\cite[Theorem~6.12, pp.~81--82; Proposition~6.16, p.~85]{DuEtAl2026v1}.
The weighted calculation
also covers the coprime multiplicities stated below, and the
standard case includes the coarse tangent direct images.
\end{remark}

\begin{theorem}[Tangent cohomology, direct images and miniversality]\hypertarget{tag-032U}{}\label{push:src:E.2.1}
\hypertarget{tag-01RH}{}\label{fam:src:A.5.2}\label{fam:st:A.5.2}\label{inv:D1.3}\label{inv:D2.1}\label{inv:D3.1}\label{cdp:st:E9}
\hypertarget{tag-02YO}{}\label{night:P5.3}\label{stmt:P5.3}
\hypertarget{tag-02YP}{}\label{night:P5.4}\label{stmt:P5.4}
\StatementRecord{push:src:E.2.1}{proved}{032U}{\cref{push:section:E.1,push:src:E.3.1,inv:D1.1,inv:D1.2}}
Let $k,l$ be positive integers with $\gcd(3k,4l)=1$, and let $X$
be a sphere member with multiplicities $(3k,4l)$ and the admissible
logarithmic data of \cref{push:section:E.1}. Put $z=kl$ and
$d=3k+4l$. On $\mathfrak B_{k,l}=\mathbb P(3k,4l)$, the vertical
root bundle $E$, the root lift $\bar f$ and its direct images satisfy
\[
 \begin{split}
 E&=\mathcal O\oplus\mathcal O(-z),\\
 R^q\bar f_*\mathcal O_X&=
 \mathcal O,\quad\mathcal O\oplus\mathcal O(-z),\quad\mathcal O(-z),\\
 K_X&=\bar f^*\mathcal O(z-d).
 \end{split}
\]
The tangent cohomology is
$\bigl(h^0,h^1,h^2,h^3\bigr)(X,T_X)=(1,1,1,0)$.
For the transition-arrow convention
$\kappa(\Phi)=(\partial_c\Phi)\Phi^{-1}$, the Kodaira--Spencer
class of $\partial_c$ is $e\otimes[\gamma]$, and the class of
$\partial_u$ is $-(2\pi i u)^{-1}e\otimes[\gamma]$.
Both classes are nonzero. The primary Kuranishi bracket vanishes,
the Kuranishi germ is a smooth reduced disc, and the fixed-data
period family is locally miniversal.

For the standard $(3,4)$ data, $e=\xi$ and the coarse direct
images are
\begin{equation}
 f_*T_X=\mathcal O\oplus\mathcal O(-1),\qquad
 R^1f_*T_X=\mathcal O\oplus\mathcal O(-1)^2,\qquad
 R^2f_*T_X=\mathcal O.
\label{pub:I:eq:D-4}
\end{equation}
\end{theorem}

\begin{proof}
\begin{enumerate}[label=\arabic*.,leftmargin=*]
\item The root lift, its proper flat Gorenstein charts, base change and
relative duality are those of \cref{push:src:E.3.1}.
Write $m=3k$ and $n=4l$. The finite root charts have the specified
free affine cyclic actions, and the cusp charts are $q=z_0z_1z_2$.

\item Use weighted coordinates $x,y$ of weights $m,n$, with coarse
coordinate $t=y^m/(y^m-x^n)$.
The two root divisors have bundles $\mathcal O(n)$ and
$\mathcal O(m)$, and the cusp has bundle $\mathcal O(mn)$.
The root-forgetting morphism
\[
 r:\mathfrak B_{k,l}\longrightarrow\mathfrak B_{1,1},
 \qquad (x,y)\longmapsto(x^l,y^k),
\]
uses the group homomorphism $\lambda\mapsto\lambda^{kl}$ on
the weighted quotient.  It is a holomorphic morphism of analytic
stacks over $B$, with local maps $s\mapsto s^k$ and
$s\mapsto s^l$.  In particular $r^*\mathcal O(1)=\mathcal O(z)$.
The analytic Picard group of this root curve has
the presentation
$\langle H,R_1,R_2\mid mR_1=H,\ nR_2=H\rangle$:
after subtracting the two stabilizer characters, a line bundle
descends to the coarse projective line.  Its only relations are
the two defining root identities.  The Smith form and
$\gcd(m,n)=1$ identify this group with $\mathbb Z$, with
$H=mn\mathcal O(1)$, $R_1=n\mathcal O(1)$, and
$R_2=m\mathcal O(1)$.  This proves the analytic torsion-free
Picard assertion used below.
Translations have identity vertical derivative and act trivially
on structure-sheaf cohomology.  Thus the vertical bundle computed
from the linear period factors is $r^*(\mathcal O\oplus
\mathcal O(-1))$.  The fixed vector gives a subline $\mathcal O$; its extension
splits because $H^1(\mathcal O(z))=0$.

\item Under the marked root-forgetting map in the preceding step,
the constant lattice cocycle $\gamma$ is unchanged. Its finite-chart
representatives and its cusp edge values $(0,-1,-1)$ therefore
give the nowhere-zero cohomology line proved in \cref{inv:D1.2}.
The cusp resolution and cup-product pairing are those of
\cref{push:src:E.3.1}; pulling the linear period factors through $r$
gives the quotient line $\mathcal O(-z)$.
Relative duality gives
$R^2\bar f_*\mathcal O_X=\mathcal O(-z)$ and
$\det R^1\bar f_*\mathcal O_X=\mathcal O(-z)$.
The extension of that quotient by the specified $[\gamma]$-line
splits because $H^1(\mathcal O(z))=0$, as computed in step~5.
Connected fibres give $\bar f_*\mathcal O_X=\mathcal O$.
These maps prove the three structure-sheaf direct images.

\item The logarithmic tangent sequence is
\[
 0\longrightarrow\mathcal V=\bar f^*E
 \longrightarrow\mathcal T=T_X(-\log W)
 \longrightarrow\mathcal H=\bar f^*\mathcal O(d-12z)
 \longrightarrow0.
\]
Its last map is the relative differential on smooth root charts
and sends $\sum a_i z_i\partial_{z_i}$ to
$(\sum a_i)q\partial_q$ at the cusp.  It is surjective there.
Since $K_{\mathfrak B}=\mathcal O(-d)$, determinants give
$K_X=\bar f^*\mathcal O(z-d)$.

\item The two quotient charts $x\ne0$, $y\ne0$ are finite quotients
of affine lines, and their overlap is a finite quotient of a
punctured line.  Finite averaging and Stein acyclicity compute
coherent cohomology by this two-chart complex.  Its homogeneous
Laurent monomials give
\[
 \begin{split}
 h^0\mathcal O(a)&=\#\{(i,j)\in\mathbb Z_{\ge0}^2:mi+nj=a\},\\
 h^1\mathcal O(a)&=h^0\mathcal O(-a-d).
 \end{split}
\]
The second equality identifies the remaining monomials
$x^{-i-1}y^{-j-1}$.
The bundles $\mathcal O(-z),\mathcal O(-2z)$ are acyclic.
Indeed a possible degree-one monomial would give
$m(i+1)+n(j+1)=az$, with $a=1$ or $2$.
Coprimality forces $i+1=lA$, $j+1=kB$, hence
$3A+4B=a$ with $A,B\ge1$, which is impossible.
Positive degree has no degree-one monomials, so
$H^1\mathcal O(z)=0$ as used above.

\item Put $b=d-12z$.  Both $b$ and $b-z$ are negative.
Degree-one monomials for $\mathcal O(b)$ would satisfy
$3A+4B=12$, with $A,B\ge1$, and none exist.
For $\mathcal O(b-z)$ the equation is $3A+4B=13$;
its only positive solution is $(A,B)=(3,1)$, producing
$x^{3l-2}y^{k-2}$ in the dual description.  It is allowed
exactly when $k>1$; write $e_k=1$ in that case and $e_k=0$
otherwise.  Projection formula and the two-column Leray sequence
now give
\[
 h^\bullet(\mathcal V)=(1,1,0,0),\qquad
 h^\bullet(\mathcal H)=(0,0,e_k,e_k),\qquad
 h^\bullet(\mathcal T)=(1,1,e_k,e_k).
\]

\item On the normalization $\nu:P=\mathrm{dP}_6\to W\subset X$,
the branchwise normal line is $K_P$: the invariant cusp form
$d\log x_1\wedge d\log x_2\wedge dq$ trivializes $K_X$
on the entire cusp neighbourhood.  Local normal coefficients
therefore give the exact sequence
\[
 0\longrightarrow\mathcal T\longrightarrow T_X
 \longrightarrow\nu_*K_P\longrightarrow0,
 \qquad h^\bullet(K_P)=(0,0,1).
\]
For $k>1$, the dual logarithmic sequence, tensored with $K_X$, has horizontal
subline $\bar f^*\mathcal O(13z-2d)$ and vertical quotient
$\bar f^*(\mathcal O(z-d)\oplus\mathcal O(2z-d))$.
The latter has no global sections, by the same equations
$3A+4B=1,2$ above.  Consequently
\[
 H^0(\Omega_X^1(\log W)\otimes K_X)
 =H^0(\mathfrak B,\mathcal O(13z-2d)).
\]
For $k>1$ this is generated by the monomial
$x^{3l-2}y^{k-2}$ multiplying the horizontal logarithmic
differential.  Both $x$ and $y$ are nonzero over $p_0$, so its
residue on $P$ is nonzero.  The ordinary logarithmic residue
sequence, tensored with $K_X$, identifies this residue map with
the Serre dual of
$H^2(P,K_P)\to H^3(X,\mathcal T)$.
One can check this identification directly in each chart
$q=z_0z_1z_2$: the residues of $dq/q$ are one on each branch,
and contraction with the invariant three-form is the normal
coefficient pairing.  Thus that connecting map has rank $e_k$.
The long exact sequence proves $h^2(T_X)=1$ and $h^3(T_X)=0$;
degrees zero and one are unchanged.

\item The derivative $\partial_c\Pi=e\otimes\gamma$ gives the constant
groupoid cocycle $\lambda\mapsto\gamma(\lambda)e$ on every
regular torus.  On the enlarged finite roots the derivative of the
affine generator is $(r_j/m_j)e$, and the derivative of its
punctured connector is $(\log s_j)/(2\pi i)\,r_je$.
The identity $m_j(r_j/m_j)=\gamma(w_j)$ verifies the power
relation, while $\gamma A_j=\gamma$ verifies conjugation with
every lattice arrow.  Changing $\log s_j$ by $2\pi i N$
adds the derivative of the integral arrow $Nw_j$; the connector
formula is used only on the punctured overlap.  Its central
root-chart representative is the constant $(r_j/m_j)e$,
which is holomorphic.  At the cusp $\gamma$ kills the toric periods,
so the constant deck cocycle is exactly its extending cocycle
above.  Its image under the two isomorphisms
$H^1(\mathcal V)\simeq H^1(\mathcal T)\simeq H^1(T_X)$
is $e\otimes[\gamma]\ne0$.  The chain rule gives the asserted
$u$-class.  With the tangent-lift convention $[\bar\partial h]$, this
class changes sign, as the explicit lift comparison in
\cref{fam:src:A.4.2} shows. The parameter and the transition-arrow
convention in the statement specify the displayed sign.

\item Let $\alpha$ denote the nonzero tangent class just constructed.
The holomorphic family gives a convergent deformation with this tangent.
Differentiating its integrability equation gives
$\bar\partial\alpha_2=-[\alpha,\alpha]/2$. Thus the primary
bracket vanishes; bilinearity and $h^1(T_X)=1$ give vanishing on
every pair.

By \cite[\S VII.8, principal theorem, pp.~598--599;
\S VIII.1, Theorem~1 and its corollary, p.~600]{Douady1974},
there is a minimal versal analytic germ with tangent space $H^1(T_X)$. The identification of this tangent space
uses tangent-valued \v Cech cocycles of transition maps, with changes
of trivialization acting by coboundaries. Its local algebra is
$\mathbb C\{t\}/I$. The parameter family induces
$t\mapsto g(c)$ with $g'(c_0)\ne0$. By the analytic inverse
function theorem, substitution by $g$ is injective; it annihilates
$I$, so $I=0$. The classifying map is an isomorphism of germs.
Pullback of the versal family realizes every sufficiently small
deformation and proves local miniversality.

\item For the standard data we compute the coarse direct images
by the following Koszul maps. All root bundles in this step are on
$\mathcal B=\mathbb P(3,4)$.
\begin{enumerate}[label=(\alph*),leftmargin=*]
\item Put $L=\OO_{\mathcal B}(-5)$. The first connecting map of \eqref{fam:eq:ab-log-tangent} is
the holomorphic bundle map
$\kappa:L\to E\otimes G$. In the splittings fixed in \ref{inv:D1.2} the
coefficient in $\OO\xi\otimes\OO\gamma$ lies in $H^0(\OO(5))=0$.
The remaining coefficients lie in
$H^0(\OO(4))=\C y$, twice, and $H^0(\OO(3))=\C x$.
Therefore, as a map from $E^*$ to $G$ with its $L$ factor understood,
\begin{equation}
 \kappa=\begin{pmatrix}0&a y\\ b y&c_0x\end{pmatrix}.
\label{pub:I:eq:D-5}
\end{equation}
We use $c_0$ for this coefficient to distinguish it from the deformation
parameter. The logarithmic period derivative at the cusp has leading
term the invertible matrix $B_0/(2\pi i)$. Hence
$\det\kappa=-ab y^2$ is nonzero there, so $ab\ne0$.
At the order-three point $\tau-\rho$ has order one on the root chart:
$j=1728t$ and the elliptic ramification order of $j$ at $\rho$ is
three. The marked torus period derivative is consequently nonzero
there: the two constant identity columns in the normalized period
matrix exclude an infinitesimal linear-frame change that could
otherwise conceal a nonzero derivative. In \eqref{pub:I:eq:D-5}, $y=0$ and
$x\ne0$ at that point, so $c_0\ne0$.
For the following matrices only, write the frames of $E,G$ as
$(e_1,e_2)$ and $(g_1,g_2)$ and replace them by
\[
 e'_1=(a/c_0)e_1,\quad e'_2=e_2,\qquad
 g'_1=b g_1,\quad g'_2=c_0g_2.
\]
The coefficients of $e'_1\otimes g'_2$, $e'_2\otimes g'_1$ and
$e'_2\otimes g'_2$ are then $y,y,x$.
The geometric cocycle $\xi\otimes\gamma$ keeps its original
meaning; $e'_1\otimes g'_1=(ab/c_0)\xi\otimes\gamma$ in these
matrix frames.

\item The next connecting map is cup product with $\kappa$. In suitable
bases it is
\begin{equation}
 L\otimes G=\OO(-5)\oplus\OO(-6)
 \xrightarrow{\ \delta_1\ }
 E\otimes\det G=\OO(-1)\oplus\OO(-2),\qquad
 \delta_1=\begin{pmatrix}-y&0\\-x&y\end{pmatrix}.
\label{pub:I:eq:D-6}
\end{equation}
The sign follows by wedging first with $\gamma$ and then with the
second generator of $G$; reversing both sign conventions leaves the
cokernels unchanged. As its determinant is $-y^2$, it is injective as
a morphism of sheaves. The first connecting map is
$\delta_0=(0,y,y,x)^t$ and is a subbundle inclusion, since $x,y$ have
no common zero. Its cokernel is
\begin{equation}
 \OO\oplus\OO(-1)\oplus\OO(2).
\label{pub:I:eq:D-7}
\end{equation}
Use the ordered frames
$u=e'_1\otimes g'_2+e'_2\otimes g'_1$ and
$v=e'_1\otimes g'_2$ for the two $\OO(-1)$ summands.
The change matrix has determinant $-1$, the image is $yu+xe'_2g'_2$,
and $v$ gives the unused summand. The remaining quotient is the exact sequence
$0\to\OO(-5)\xrightarrow{(y,x)}\OO(-1)\oplus\OO(-2)
\xrightarrow{(x,-y)}\OO(2)\to0$.

\item The cokernel of \eqref{pub:I:eq:D-6} is supported at the order-three root point. There $x$ is a unit. For target generators $q_1,q_2$, the
relations are $-yq_1-xq_2=0$ and $yq_2=0$. Eliminating
$q_2=-(y/x)q_1$ gives the single relation $y^2q_1=0$ and hence
\begin{equation}
 Q=\OO(-1)\big|_{2p_1},
\label{pub:I:eq:D-8}
\end{equation}
where $2p_1$ is the length-two divisor in the root coordinate $y$. The generator is the first target component,
and its annihilator is $y^2$. The long direct-image sequence now gives
\[
 R^0\bar f_*\mathcal T=E,\quad
 R^1\bar f_*\mathcal T=\OO\oplus\OO(-1)\oplus\OO(2),
\]
\begin{equation}
 0\longrightarrow Q\longrightarrow R^2\bar f_*\mathcal T
 \longrightarrow\OO(-6)\longrightarrow0.
\label{pub:I:eq:D-9}
\end{equation}
This last extension splits, since $\operatorname{Ext}^1(\OO(-6),Q)
=H^1(Q(6))=0$.

\item Let $\nu:P=\mathrm dP_6\to W\hookrightarrow X$ be the normalization
immersion. There is an exact sequence
\begin{equation}
 \hypertarget{tag-01P1}{}\label{fam:eq:ab-normal-quotient}
0\longrightarrow\mathcal T\longrightarrow T_X
 \longrightarrow\nu_*K_P\longrightarrow0.             
\end{equation}
To verify the quotient, in a chart $W=(z_1\cdots z_r=0)$ the map takes
a vector field to its normal coefficient on each of the $r$ separate
branches. The coefficients lift independently; its kernel is precisely
the vector fields tangent to every branch. Thus the normal-coefficient map identifies the quotient with
the pushforward of the normal line of $\nu$.
The form $d\log x_1\wedge d\log x_2\wedge dq$ becomes a unit times
$dz_1\wedge dz_2\wedge dz_3$ in every unimodular chart and is invariant
under every cusp deck transformation. It trivializes $\nu^*K_X$.
Adjunction for the immersion gives its normal line $K_P$, proving
\eqref{fam:eq:ab-normal-quotient}. Since $P$ is rational, its $K_P$ cohomology is $(0,0,1)$ by
Serre duality. Finite pushforward preserves that cohomology.

\item For the degree-two extension at $p_0$, the kernel of
\[
 \Omega_X^1|_W\longrightarrow\nu_*\Omega_P^1
\]
is $\nu_*\OO_P\otimes N_W^*$. At $xy=0$ its sections are
$a(y,z)y\,dx+b(x,z)x\,dy$; at $xyz=0$ they are
$a(y,z)yz\,dx+b(x,z)xz\,dy+c(x,y)xy\,dz$.
These identifications glue by the common conormal $dq$.
Since $N_W$ and $K_X|_W$ are trivial and $H^0(P,\Omega_P^1)=0$, the
Serre-dual space of $H^2(W,T_X|_W)$ is one-dimensional.
Top-degree base change identifies this group with the residue fibre
of $(R^2\bar f_*T_X)_{p_0}$. Its generic rank is one by \eqref{pub:I:eq:D-9}.
A finite module over the discrete valuation ring $\C\{q\}$ of rank
one and residue dimension one is free of rank one. Thus the extension
by $\C_{p_0}$ increases the degree of its line part by $p_0$, without
adding torsion. We obtain the complete root direct images
\[
 R^0\bar f_*T_X=\OO\oplus\OO(-1),\quad
 R^1\bar f_*T_X=\OO\oplus\OO(-1)\oplus\OO(2),
\]
\begin{equation}
 R^2\bar f_*T_X=\OO(6)\oplus Q.
\label{pub:I:eq:D-11}
\end{equation}
On the order-three chart, an invariant coefficient of the generator
of $Q$ must have exponent $2$ modulo $3$, whereas $Q$ only retains
exponents $0,1$. Therefore $\rho_*Q=0$. Applying the exact $\rho_*$
to \eqref{pub:I:eq:D-11} and using
$d(0)=0,d(-1)=d(2)=-1,d(6)=0$ proves \eqref{pub:I:eq:D-4}.
The coarse direct images are the ones displayed in the statement;
their Leray cohomology agrees with the tuple obtained in step~7.
\end{enumerate}
\end{enumerate}
\end{proof}


\section{The integral parameter comparison}\label{pub:I:torsor-appendix}

\subsection{Ordinary tangent sheaves and the conductor}
\begin{definition}[The double locus and its tangent quotient]
\hypertarget{tag-PI12}{}\label{I:ordinary-tangent-sheaves}
\StatementRecord{I:ordinary-tangent-sheaves}{proved}{PI12}{}

For the $(3,4)$ family, let
$D_{\mathrm{dbl}}=(\operatorname{Sing}W)_{\mathrm{red}}$.
It consists of three rational curves joining the two triple points.
Put
\[
 L=f^*\mathcal O_B(-1)\otimes\mathcal O_X(2S_1+S_2),\qquad
 F=\mathcal O_X(S_1+S_2),\qquad
 \mathcal V=\ker(df),\qquad
 J=\mathcal I_{D_{\mathrm{dbl}}}\otimes F.
\]
The root conventions give
$\rho^*\mathcal O_B(1)=\mathcal O_{\mathcal B}(12)$,
$\mathcal O_X(S_1)=\bar f^*\mathcal O_{\mathcal B}(4)$ and
$\mathcal O_X(S_2)=\bar f^*\mathcal O_{\mathcal B}(3)$.
The tensor maps on their canonical divisor sections therefore give
\[
 L=\bar f^*\widehat L,\quad \widehat L=\mathcal O_{\mathcal B}(-1),
 \qquad F=\bar f^*\widehat F,\quad
 \widehat F=\mathcal O_{\mathcal B}(7).
\]
The weights are $-12+2\cdot4+3=-1$ and $4+3=7$.
The vertical root bundle is $E$ of \cref{inv:D1.2}.
\end{definition}

\begin{theorem}[The ordinary tangent sequence]\hypertarget{tag-02YM}{}\label{night:P5.1}\label{stmt:P5.1}
\StatementRecord{night:P5.1}{proved}{02YM}{\cref{I:ordinary-tangent-sheaves,inv:D1.2}}
The sheaf $\mathcal V$ is locally free of rank two and equals
$\bar f^*E$. The differential induces an exact sequence
\begin{equation}\hypertarget{tag-01OK}{}\label{eq:P5tangent}
 0\longrightarrow\mathcal V\longrightarrow T_X
 \longrightarrow J\longrightarrow0.
\end{equation}
\end{theorem}
\begin{proof}
On a finite root chart $t=s^m$, the differential has image
$(s^{m-1})f^*T_B$. The free cyclic quotient makes the root chart
\'etale over $X$, so the orders along $S_1,S_2$ are $2,3$.
At a cusp double point $q=xy$, its image is $(x,y)$; at a triple
point $q=xyz$, its image is $(xy,xz,yz)$, the reduced ideal of the
three axes. At a smooth point of $W$, it is surjective.
Thus the global image is
\[
 \mathcal I_{D_{\mathrm{dbl}}}\otimes f^*T_B
 \otimes\mathcal O_X(-2S_1-3S_2)=J.
\]
Here $t(t-1)\partial_t$ trivializes $T_B(-p_1-p_2)$,
and $f^*(p_1)=3S_1$, $f^*(p_2)=4S_2$ identify the remaining
line factor with $\mathcal O_X(S_1+S_2)$.
The kernel and its finite, toric and period-transition maps were
computed in \cref{inv:D1.2}; those evaluation maps identify it
with $\bar f^*E$, including at the double and triple points.
\end{proof}


\begin{theorem}[The ordinary direct images]\hypertarget{tag-02YN}{}\label{night:P5.2}\label{stmt:P5.2}
\StatementRecord{night:P5.2}{proved}{02YN}{\cref{inv:D1.1,inv:D1.2,push:src:E.3.1,night:P5.1}}
For the sheaves of \cref{I:ordinary-tangent-sheaves}, the coarse
direct images are
\[
\begin{array}{c|ccc}
 & f_* & R^1f_* &R^2f_*\\ \hline
 \mathcal O_X&\mathcal O&\mathcal O\oplus\mathcal O(-1)&\mathcal O(-1)\\
 L&\mathcal O(-1)&\mathcal O(-1)^{\oplus2}&\mathcal O(-1)\\
 F&\mathcal O&\mathcal O^{\oplus2}&\mathcal O\\
 \mathcal V&\mathcal O\oplus\mathcal O(-1)&\mathcal O\oplus\mathcal O(-1)^{\oplus3}
       &\mathcal O(-1)^{\oplus2}\\
 J&\mathcal O(-1)&\mathcal O(-1)^{\oplus2}&\mathcal O .
\end{array}
\]
\end{theorem}
\begin{proof}
\begin{enumerate}[label=\arabic*.,leftmargin=*]
\item The root splitting and structure-sheaf direct images are given in
\cref{inv:D1.2,push:src:E.3.1}; in the standard data these are also computed in
\cite[Proposition~6.6, pp.~77--78; Lemma~6.7, p.~78;
Proposition~6.10, pp.~79--80]{DuEtAl2026v1}.
The determinant line in these conventions is
\begin{equation}\hypertarget{tag-01OG}{}\label{eq:P5L}
 \det E=\widehat L
 =\rho^*\mathcal O_B(-1)\otimes
   \mathcal O_{\mathcal B}(2p_1/3+p_2/4)
 =\mathcal O_{\mathcal B}(-1).
\end{equation}
The equality is the tensor product of the divisor-section maps
in \cref{I:ordinary-tangent-sheaves}; its degree is $-12+8+3$.
For integers $d,a,b$, the invariant root-line formula is
\begin{equation}\hypertarget{tag-01OI}{}\label{eq:P5floor}
 \rho_*\bigl(\rho^*\mathcal O_B(d)\otimes
 \mathcal O_{\mathcal B}(ap_1/3+bp_2/4)\bigr)
 =\mathcal O_B\bigl(d+\lfloor a/3\rfloor+\lfloor b/4\rfloor\bigr).
\end{equation}
Indeed writing $a=3q+r$, $b=4q'+r'$ with
$0\leq r<3$, $0\leq r'<4$ separates the coarse divisor
$(d+q+q')$ from the two root characters. This is the same
invariant-section map as \cref{inv:D1.1}, and applies to all
integer $a,b$.

\item For a root line bundle $A$, projection and exact finite
invariants give
\[
 R^qf_*\bar f^*A=\rho_*\bigl(A\otimes
              \bigwedge^q(\mathcal O\oplus\widehat L)\bigr).
\]
For $A=\mathcal O,\widehat L,\widehat F$, and for
$E=\mathcal O\oplus\widehat L$, the formula
\eqref{eq:P5floor} gives the first four rows. The root-to-coarse
map, including the full scheme fibres, is \cref{I:root-coarse-comparison}.

\item For the last row one must compute restriction to $D_{\mathrm{dbl}}$.  The conductor
square gives
\[
0\to\mathcal O_W\to\nu_*\mathcal O_S\oplus\mathcal O_{D_{\mathrm{dbl}}}
\to\nu_*\mathcal O_{\partial S}\to0,
\]
where $\partial S$ is the hexagon.  Here $h^1(\mathcal O_{\partial S})=1$
and $h^1(\mathcal O_{D_{\mathrm{dbl}}})=2$.  The equality conditions at the two triple
points give incidence matrix
\[
 \begin{pmatrix}1&-1&0\\0&1&-1\\1&-1&0\\0&1&-1\end{pmatrix}.
\]
Orient the three theta-graph edges from $P$ to $Q$.  The hexagon
traversal maps to $e_1-e_2+e_3-e_1+e_2-e_3=0$.  The constant-to-coherent
branch comparison is an isomorphism in degree one, so the induced map
$H^1(D_{\mathrm{dbl}},\mathcal O_{D_{\mathrm{dbl}}})\to H^1(\partial S,\mathcal O_{\partial S})$ is
zero.  The conductor sequence then makes
$H^1(W,\mathcal O_W)\to H^1(D_{\mathrm{dbl}},\mathcal O_{D_{\mathrm{dbl}}})$ an isomorphism.

\item The canonical section of $F$ is nonzero near $W$, so $F|_{D_{\mathrm{dbl}}}=\mathcal O_{D_{\mathrm{dbl}}}$.
In $0\to J\to F\to\mathcal O_{D_{\mathrm{dbl}}}\to0$, the maps of direct images
$\mathcal O\to\mathbb C_{p_0}$ and
$\mathcal O^{\oplus2}\to\mathbb C_{p_0}^{\oplus2}$ are evaluation and
evaluation followed by this isomorphism.  They are surjective, with
kernels $\mathcal O(-p_0)$ in the indicated ranks.  Finally
$H^2(D_{\mathrm{dbl}},\mathcal O_{D_{\mathrm{dbl}}})=0$ and $R^2f_*F=\mathcal O$.  This gives the last row.
\end{enumerate}
\end{proof}






\subsection{The full automorphism group}\hypertarget{tag-000E}{}\label{app:rederive:fam-src-B.4.1}

\begin{theorem}[The full automorphism group]\hypertarget{tag-01QF}{}\label{fam:src:B.4.1}\label{fam:st:B.4.1}\hypertarget{tag-02YQ}{}\label{night:P5.5}\label{stmt:P5.5}
\StatementRecord{night:P5.5}{proved}{02YQ}{the definitions and constructions stated in the proof}

For every admissible twist of the $(3,4)$ construction,
\[
 \operatorname{Aut}(X_c)=\mathbb C/\mathbb Z\simeq\mathbb C^*.
\]
The parameter $a\bmod\mathbb Z$ acts by $\zeta\mapsto\zeta+ae$,
where $e=\Pi_c\widehat\delta$. The exponential group isomorphism
is $a\bmod\mathbb Z\mapsto e^{2\pi ia}$.
\end{theorem}
\begin{proof}
\begin{enumerate}[label=\arabic*.,leftmargin=*]
\item Proposition~\ref{fam:src:A.1.1} identifies the complete
anticanonical pencil with $f$ and proves that every automorphism
induces the identity on its base.

\item Apply the affine-lift and centralizer calculation of
\cref{fam:src:A.2.1} with $c'=c$. The period equation
\eqref{wb:pshift:eq:period-morphism} gives the integral linear part
$P=aI+bE_{41}$, $a=\pm1$, and the complex linear part $H=aI$;
its first column gives $b=0$.

\item The finite-root equation
\eqref{wb:pshift:eq:finite-norm-obstruction}, with $n=0$,
gives $N_1\ell=-(1-a)v_1$. For $a=-1$, applying the invariant
covector $\gamma$ gives $3\gamma(\ell)=-2l_1$, where
$l_1=\gamma(v_1)$. Freeness gives $3\nmid l_1$, so this
integer equality is impossible. Thus $a=1$, and every
automorphism is fibrewise translation.

\item Let $j$ be the
inclusion of the cusp-punctured orbifold and set
$\mathcal H=j_*\Lambda$.  Its cusp stalk is
$\ker(M_0-I)=\mathbb Z\widehat w\oplus\mathbb Z\widehat\delta$.
Let $\mathcal A$ be the sheaf of fibrewise translations extending to
the local compactification.  There is an exact sequence
\begin{equation}\hypertarget{tag-01OH}{}\label{eq:P5exponential}
0\longrightarrow\mathcal H\longrightarrow E\longrightarrow\mathcal A
\longrightarrow0.
\end{equation}
For the local check at the cusp, use the toroidal covering whose
component lattice is $\Gamma=\Lambda/\Lambda_{\rm tor}$.  A translation
acts identically on the nearby fibre lattice and on $\Gamma$, and
therefore lifts.  Its meridian remains bounding a disc through a smooth
point of $W$, which verifies the same assertion for the full local
cover.  On the punctured torus the lift has form $x\mapsto a(q)x$.
It sends the component $E_0$ to $E_v$; the index-one map
$B_0:\Gamma\to\mathbb Z^2$ is an isomorphism, so composing with the
unique corresponding deck shear makes it preserve $E_0$.  It then
preserves the torus $E_0\setminus\partial E_0$.  Evaluation at $(1,1)$
in the chart $(\mathbb C^*)^2\times\Delta_q$ shows that $a_i(q)$ are
holomorphic units at zero, and hence have logarithms.  Conversely their
multiplication extends on every toric chart and commutes with the deck
shears.  The logarithm kernel consists precisely of those periods with
no logarithmic term, namely $\ker(M_0-I)$.  On a finite root disc, a holomorphic lift has lattice defect in
$\ker N_j=(I-A_j)\Lambda$, by the saturated-image calculation
of \cref{fam:src:A.1.2}. Subtracting a period kills the defect
and preserves holomorphy. This gives the same exact sequence
with its finite equivariance.

\item   Put
\[
 M_0=(A_1A_2)^{-1}=\begin{pmatrix}1&0&0&0\\0&1&0&0\\0&1&1&0\\-1&0&0&1\end{pmatrix}.
\]
A cocycle on the punctured orbifold is a pair $(x_1,x_2)$ with
$N_1x_1=N_2x_2=0$, for the norm matrices in
\eqref{wb:pshift:eq:norms}.
Extension over the cusp says that $x_1+A_1x_2$ vanishes in
$\operatorname{coker}(M_0-I)$, so its first two coordinates vanish.
The complete integral solution is
\[
 x_1=(0,-b,2b-3c,c),\qquad x_2=(0,-2a-b,b,a),\quad a,b,c\in\mathbb Z.
\]
The principal cocycle of $(g,u,w,d)$ has parameters
\[
 \begin{pmatrix}a\\b\\c\end{pmatrix}
 =\begin{pmatrix}3&0&1\\-6&1&-1\\-2&1&0\end{pmatrix}
  \begin{pmatrix}g\\u\\w\end{pmatrix}.
\]
This matrix has determinant $-1$; its inverse is explicitly
$g=c-a-b$, $u=3c-2a-2b$, $w=4a+3b-3c$.  Hence
$H^1(\mathcal B,\mathcal H)=0$ integrally.

\item Sequence \eqref{eq:P5exponential} now makes every global translation a
global logarithm in $H^0(\mathcal B,E)=\mathbb C$.  Its kernel is
$\Lambda^{A_1,A_2}=\mathbb Z\widehat\delta$.  The flow extends by
commutation with the affine finite actions and by torus multiplication
on the cusp charts.  This proves existence, faithfulness, and
completeness of the asserted group.
\end{enumerate}
\end{proof}

\begin{theorem}[Integral torsor classes and the parameter period]\hypertarget{tag-02YR}{}\label{night:P5.6}\label{stmt:P5.6}
\StatementRecord{night:P5.6}{proved}{02YR}{the definitions and constructions stated in the proof}

For fixed admissible twist data, $X_c\simeq X_{c'}$ precisely when
$c-c'\in2\mathbb Z$.  This fixes the integral period of $c$; the
construction still permits an additive choice of $\beta^{\rm part}$.
\end{theorem}
\begin{proof}
\begin{enumerate}[label=\arabic*.,leftmargin=*]
\item The analytically computed complete linear system
$|\omega_{X_c}^{-2}|$ is the marked map $f_c$, by
\cref{wb:thm:canonical-ring-recovery}. A biholomorphism preserves its
canonical bundle, the complete space of sections, and its evaluation
map. Consequently it induces $\alpha\in\operatorname{PGL}_2(\mathbb C)$
with $f_{c'}F=\alpha f_c$. The three special fibres have distinct
intrinsic types: the cusp is reduced and nonnormal, while the other
two have multiplicities three and four. All remaining fibres are
smooth reduced tori. Thus $\alpha$ fixes $p_0,p_1,p_2$ individually.
In the marked coordinate, a fractional-linear map fixing
$\infty,0,1$ has zero lower-left and upper-right entries and equal
diagonal entries; hence $\alpha=\mathrm{id}_B$.

The centralizer computation of \cref{wb:pshift:lem:centralizer}
then gives $P=aI+bE_{41}$ with $a=\pm1$ and $b=a(c-c')$.
The root-fibre character calculation in that proposition excludes
$a=-1$. Thus $b=c-c'$ is an integer.

\item Compute $H^2(\mathcal B,\mathcal H)$ with a pair of pants, two root discs,
and a cusp disc.  The pair of pants has cochains $\Lambda\to\Lambda^2$
with differential $((A_1-I),(A_2-I))$.  A root disc has the periodic
cyclic resolution with alternating differentials $A_j-I,N_j$; the
cusp disc has $\ker(M_0-I)$ in degree zero.  An annulus has
$\Lambda\xrightarrow{A_j-I}\Lambda$.  The degree-one restriction from
the pair of pants is
\[
 (x_1,x_2)\longmapsto(x_1,x_2,-M_0(x_1+A_1x_2)).
\]
The total-complex ranks in degrees zero through three are $14,28,20,8$.
The resulting integer differentials have sizes
$28\times14$, $20\times28$ and $8\times20$, and their consecutive
products vanish by the restriction identities.
Smith decompositions with unimodular changes of basis identify each
integral kernel before quotienting by the preceding image, giving
$H^0=\mathbb Z$, $H^1=0$, $H^2=\mathbb Z\oplus\mathbb Z/2$.


\item The degree-two presentation also reduces by hand.  The invariant
lattices have bases $(\widehat\delta,\varepsilon_j)$, where
$\varepsilon_1=(1,2,-4,0)$ and $\varepsilon_2=(1,3,-3,0)$.
The five generators $(\delta_1,\varepsilon_1,\delta_2,\varepsilon_2,h_c)$
have relations
\begin{equation}\hypertarget{tag-01OF}{}\label{eq:P5H2}
 \delta_1=0,\qquad 2\delta_2=0,\qquad
 3\varepsilon_1=h_c=4\varepsilon_2.
\end{equation}
For the constant sheaf $\mathbb Z$, the same computation has generators
$z_1,z_2,h_c$ and relations $3z_1=h_c=4z_2$; its primitive generator is
$z_1-z_2$.  Inclusion $\mathbb Z\widehat\delta\to\mathcal H$ sends it
to the nonzero order-two class $\delta_2$.  The character $\gamma$
sends $\varepsilon_j$ to $z_j$ and kills $\delta_j$.

\item Both $H^1(\mathcal B,E)$ and $H^2(\mathcal B,E)$ vanish, by
$\rho_*E=\mathcal O\oplus\mathcal O(-1)$ and the dimension of the coarse
curve.  Consequently \eqref{eq:P5exponential} identifies compactified
translation torsors $H^1(\mathcal B,\mathcal A)$ with
$H^2(\mathcal B,\mathcal H)$.  This identification takes logarithms of
local translation transitions and their lattice-valued \v Cech
coboundaries; the two vanishings also reverse that construction.

\item The cyclic norm of the translation $v_j/m_j$ is $v_j$, while the cusp
collar has the opposite boundary orientation.  Thus the torsor
character is
\[
 l_1z_1+l_2z_2-l_0h_c
 =(4l_1+3l_2-12l_0)(z_1-z_2)=-p(z_1-z_2).
\]
The shear $S_b-I=b\widehat\delta\otimes\gamma$ changes its class by
$-bp\delta_2$.  Freeness forces $l_2$ odd and thus $p$ odd.  The class
is preserved precisely for even $b$.

\item For an integral shear the Jacobian charts themselves identify: the
period lattices differ by $S_b$, and the bounded cusp matrix changes by
an integer in its $\beta$ entry.  Every toric unit
$\exp(2\pi iC\lambda)$ is therefore unchanged.  The fan, lattice, and
toric gluings are the same, and on root discs the basis change
identifies the smooth Jacobian and its cyclic action.  Equality of the
torsor classes now gives a holomorphic isomorphism for the fixed sheaf
$\mathcal A$, which consists precisely of translations extending over
the cusp.  It extends on all three special charts.  Conversely every
isomorphism preserves the class, so forces $b$ even.
\end{enumerate}
\end{proof}

\clearpage
\part{Results not yet established}
\section{Open arithmetic questions for the normalized norm}
\hypertarget{tag-003T}{}\label{ari:open:norm}
\hypertarget{tag-001E}{}\label{ari:C.7.5}

The arithmetic nature of $N(P)$ remains open for algebraic
nontorsion points in this family. An explicit instance is
\[
 E:y^2=x^3-140x+784,\qquad P=(0,28),\qquad j=-3375,
 \qquad N(P)=2e^{-\widehat h_O(P)}.
\]
The point is primitive modulo rational two-torsion, and
$\widehat h_O(P)=0.46736568\ldots$.
The unresolved assertion is transcendence of
$2e^{-\widehat h_O(P)}$.
For the conductor-$361$ example
$y^2+y=x^3-38x+90$, $P=(0,9)$, the corresponding exact
formula is
\[
 N(P)=e^{-\widehat h_O(P)}
     =\exp\!\left(-\frac{L'(E,1)}{4\Omega_E}\right).
\]
Here every finite local contribution vanishes.
The compact-point equivalence in \cref{ari:C.7.2} and the
explicit constructions in \cref{ari:C.7.9,ari:C.7.10}
give precise group-theoretic forms of this same unresolved value.

The exact analytic leading quotients in
\cref{I:analytic-leading-quotients} also leave group-order
questions: the orders of the Tate--Shafarevich groups are
identified here for rows $3,5,7$ only.
In particular the analytic quotient of row eleven is exactly
$9$; its identification with
$\#\operatorname{Sha}(E_{11}/\mathbb Q)$, including its
$3$-primary part, remains open in the results of this paper.

\subsection{The unresolved relative-period bound}
\hypertarget{tag-001D}{}\label{ari:C.7.4}

The open lower-bound question is the Relative Semi-abelian Conjecture
\cite[Conjecture~4.7, arXiv v2, p.~18]{bertolin2022} which asks whether algebraic lifts in the indicated semiabelian
varieties satisfy the bound
\[
 \operatorname{trdeg}_{\overline{\mathbb Q}(\Omega_G)}
 \overline{\mathbb Q}\bigl(\Omega_G,R,\bar R,
                   \widetilde{\log}_G(R,\bar R)\bigr)
 \geq 2\dim B_Q+\dim Z(1).
\]
The full matrix $\Omega_G$ includes the second-kind periods.
This lower bound remains open for the compact diagonal
configuration at issue here. Its relation to Grothendieck's
period conjecture for algebraic $1$-motives is given in
\cite[Proposition~4.12, arXiv v2, pp.~19--20]{bertolin2022}.

For the algebraic motives of \cref{ari:C.7.3}, projection
from the point-extension variety to the extension-coordinate
variety is an isogeny: the point coordinates agree with the
extension data under the stated polarizations.
Equivalently, $(P,\bar P)=(P,0)+(0,\bar P)$ in
$E\times\bar E$. Formula~(4.4), p.~19 of the same source
therefore gives $\dim B_Q=0$.
The remaining kernel of connected motivic groups is the
weight-$-2$ vector group $\operatorname{Lie}Z$.
After fixing the base periods its fibre is an affine torsor
of dimension $\dim Z$.

In positive Tate coordinates the central action changes
$(v,\bar v)$ by $2\pi i(c_1,c_2)$.
For an isogeny $\phi:\bar E\to E$ of degree $d$,
$\phi^{\mathsf t}J_E\phi=dJ_{\bar E}$.
Thus the dependent bracket line of \cref{ari:C.7.3}
has coordinates $(1,aa^\dagger/d)$ with positive second
coordinate; in the independent case the bracket spans both
central coordinates.
The functional $(c_1,c_2)\mapsto c_1+c_2$ is nonzero
on this central group, including when there is a further
quotient Tate direction.
Accordingly, the compactness equation
\eqref{ari:eq:cp-relation} cuts a proper affine hyperplane.
The unconditional period-field calculation is
\cref{I:compact-period-relation}; the proposed lower bound
needed to exclude an algebraic compact lift is the open step.

\clearpage

\section*{Programs}
\addcontentsline{toc}{section}{Programs}
The following files contain the indicated computations at commit
\begin{center}\small\ttfamily\SsixRepositoryCommit\end{center}
Each link identifies the complete repository path.

{\footnotesize
\renewcommand{\arraystretch}{.9}
\setlength{\tabcolsep}{4pt}
\begin{longtable}{@{}>{\raggedright\arraybackslash}p{.57\textwidth}>{\raggedright\arraybackslash}p{.39\textwidth}@{}}
\caption{Programs and their mathematical scope.}\label{tab:I:programs}\\
\toprule
Computation and statement & Source file \\
\midrule
\endfirsthead
\toprule
Computation and statement & Source file \\
\midrule
\endhead
\bottomrule
\endfoot
\hypertarget{prog:I:integral-algebra}{}Period determinants, integral parameter equations and tangent ranks; cusp deck identities and finite actions; initial node blowup with saturated equation $bw=xy$\par \cref{inv:A2.1,night:P5.1,night:P5.2,night:P5.6,night:P2.1,night:P2.2,night:P2.5,night:P2.8,night:P2.9,night:P2.10} &
\href{https://raw.githubusercontent.com/KokunoYumeto/s6-complex-structure-workbench/\SsixRepositoryCommit/pro_returns/S6_SERIES_GRIND_checkpoint_10_20261120T0211Z/code/current_audits/audit_paper_i/code/priority_rows_exact.py}{\nolinkurl{priority_rows_exact.py}}\par \href{https://raw.githubusercontent.com/KokunoYumeto/s6-complex-structure-workbench/\SsixRepositoryCommit/pro_returns/S6_SERIES_GRIND_checkpoint_10_20261120T0211Z/code/current_audits/audit_paper_i/code/adjacent_exact.py}{\nolinkurl{adjacent_exact.py}} \\[3pt]
\hypertarget{prog:I:double-complex}{}Finite double-complex matrices and their cohomology\par \cref{inv:C1.1,inv:C2.1} &
\href{https://raw.githubusercontent.com/KokunoYumeto/s6-complex-structure-workbench/\SsixRepositoryCommit/pro_returns/S6_SERIES_GRIND_checkpoint_10_20261120T0211Z/code/current_audits/audit_paper_i/code/main_D_double_complex.py}{\nolinkurl{main_D_double_complex.py}} \\[3pt]
\hypertarget{prog:I:coefficient-atlas}{}Regular-chart identities, cusp Cartier factors and finite-collar first jets\par \cref{inv:C3.8} &
\href{https://raw.githubusercontent.com/KokunoYumeto/s6-complex-structure-workbench/\SsixRepositoryCommit/pro_returns/S6_SERIES_GRIND_checkpoint_10_20261120T0211Z/code/current_audits/audit_paper_i/open_atlas/exact_regular_atlas.py}{\nolinkurl{exact_regular_atlas.py}}\par \href{https://raw.githubusercontent.com/KokunoYumeto/s6-complex-structure-workbench/\SsixRepositoryCommit/pro_returns/S6_SERIES_GRIND_checkpoint_10_20261120T0211Z/code/current_audits/audit_paper_i/open_atlas/cusp_and_finite_exact.py}{\nolinkurl{cusp_and_finite_exact.py}}\par \href{https://raw.githubusercontent.com/KokunoYumeto/s6-complex-structure-workbench/\SsixRepositoryCommit/pro_returns/S6_SERIES_GRIND_checkpoint_10_20261120T0211Z/code/current_audits/audit_paper_i/open_atlas/finite_collar_jet_verifier.py}{\nolinkurl{finite_collar_jet_verifier.py}} \\[3pt]
\hypertarget{prog:I:koszul}{}Tangent-twist Smith modules and stabilizer-character jet ranks\par \cref{inv:D4.4} &
\href{https://raw.githubusercontent.com/KokunoYumeto/s6-complex-structure-workbench/\SsixRepositoryCommit/pro_returns/S6_10_PAPERS_checkpoint_15_20261119T1045Z/S6_10_PAPERS/code/I/replay_invariants/d_koszul.py}{\nolinkurl{d_koszul.py}} \\[3pt]
\hypertarget{prog:I:covers}{}Weighted section counts and rational boundary maps of finite covers\par \cref{inv:B2.2,inv:B3.1,inv:B3.3} &
\href{https://raw.githubusercontent.com/KokunoYumeto/s6-complex-structure-workbench/\SsixRepositoryCommit/pro_returns/S6_01_INVARIANTS/return/code/b_families.py}{\nolinkurl{b_families.py}}\par \href{https://raw.githubusercontent.com/KokunoYumeto/s6-complex-structure-workbench/\SsixRepositoryCommit/pro_returns/S6_01_INVARIANTS/return/code/b_topology.py}{\nolinkurl{b_topology.py}} \\[3pt]
\hypertarget{prog:I:resonance}{}Division-polynomial factors at the marked elliptic section\par \cref{inv:D4.3} &
\href{https://raw.githubusercontent.com/KokunoYumeto/s6-complex-structure-workbench/\SsixRepositoryCommit/pro_returns/S6_01_INVARIANTS/return/code/d_division_polynomials.py}{\nolinkurl{d_division_polynomials.py}} \\[3pt]
% Analytic/arithmetic rows for the one final Paper I program table.
\hypertarget{prog:I:division}{}Exact division recurrences and resultants\par \cref{ari:C.4.3,ari:C.4.4,ari:C.5.1,ari:C.8.4} &
\href{https://raw.githubusercontent.com/KokunoYumeto/s6-complex-structure-workbench/\SsixRepositoryCommit/pro_returns/S6_10_PAPERS_checkpoint_15_20261119T1045Z/S6_10_PAPERS/code/I/replay_special/Cspecial_division.py}{\nolinkurl{Cspecial_division.py}}\par \href{https://raw.githubusercontent.com/KokunoYumeto/s6-complex-structure-workbench/\SsixRepositoryCommit/pro_returns/S6_10_PAPERS_checkpoint_15_20261119T1045Z/S6_10_PAPERS/code/I/replay_special/Cspecial_relations.py}{\nolinkurl{Cspecial_relations.py}} \\[3pt]
\hypertarget{prog:I:cm_models}{}Rational CM models, local reductions and finite heights\par \cref{ari:C.6.3} &
\href{https://raw.githubusercontent.com/KokunoYumeto/s6-complex-structure-workbench/\SsixRepositoryCommit/pro_returns/S6_10_PAPERS_checkpoint_15_20261119T1045Z/S6_10_PAPERS/code/I/replay_arithmetic/C_cm_models.gp}{\nolinkurl{C_cm_models.gp}}\par \href{https://raw.githubusercontent.com/KokunoYumeto/s6-complex-structure-workbench/\SsixRepositoryCommit/pro_returns/S6_10_PAPERS_checkpoint_15_20261119T1045Z/S6_10_PAPERS/code/I/replay_arithmetic/C_cm_local_heights.gp}{\nolinkurl{C_cm_local_heights.gp}} \\[3pt]
\hypertarget{prog:I:cm_heights}{}Outward interval enclosures for CM heights and singular moduli\par \cref{ari:C.6.5} &
\href{https://raw.githubusercontent.com/KokunoYumeto/s6-complex-structure-workbench/\SsixRepositoryCommit/pro_returns/S6_10_PAPERS_checkpoint_15_20261119T1045Z/S6_10_PAPERS/code/I/replay_arithmetic/C_cm_height_arb.py}{\nolinkurl{C_cm_height_arb.py}} \\[3pt]
\hypertarget{prog:I:ranks}{}Exact descent, local rank data and saturation\par \cref{ari:C.6.R1} &
\href{https://raw.githubusercontent.com/KokunoYumeto/s6-complex-structure-workbench/\SsixRepositoryCommit/pro_returns/S6_10_PAPERS_checkpoint_15_20261119T1045Z/S6_10_PAPERS/code/I/replay_arithmetic/C_cm_rank_lvalues.gp}{\nolinkurl{C_cm_rank_lvalues.gp}}\par \href{https://raw.githubusercontent.com/KokunoYumeto/s6-complex-structure-workbench/\SsixRepositoryCommit/pro_returns/S6_10_PAPERS_checkpoint_15_20261119T1045Z/S6_10_PAPERS/code/I/replay_arithmetic/Crank_saturation.cpp}{\nolinkurl{Crank_saturation.cpp}} \\[3pt]
\hypertarget{prog:I:derivatives}{}Integer coefficient streams, derivative sums and full-real-period enclosures\par \cref{I:CM-coefficient-series,ari:C.6.R2} &
\href{https://raw.githubusercontent.com/KokunoYumeto/s6-complex-structure-workbench/\SsixRepositoryCommit/pro_returns/S6_10_PAPERS_checkpoint_15_20261119T1045Z/S6_10_PAPERS/code/I/replay_arithmetic/Crank_coefficients.gp}{\nolinkurl{Crank_coefficients.gp}}\par \href{https://raw.githubusercontent.com/KokunoYumeto/s6-complex-structure-workbench/\SsixRepositoryCommit/pro_returns/S6_10_PAPERS_checkpoint_15_20261119T1045Z/S6_10_PAPERS/code/I/replay_arithmetic/Crank_derivatives.cpp}{\nolinkurl{Crank_derivatives.cpp}}\par \href{https://raw.githubusercontent.com/KokunoYumeto/s6-complex-structure-workbench/\SsixRepositoryCommit/pro_returns/S6_10_PAPERS_checkpoint_15_20261119T1045Z/S6_10_PAPERS/code/I/replay_arithmetic/Crank_cm_stream.cpp}{\nolinkurl{Crank_cm_stream.cpp}}\par \href{https://raw.githubusercontent.com/KokunoYumeto/s6-complex-structure-workbench/\SsixRepositoryCommit/pro_returns/S6_10_PAPERS_checkpoint_15_20261119T1045Z/S6_10_PAPERS/code/I/certify_all_period_ratios.py}{\nolinkurl{certify_all_period_ratios.py}} \\[3pt]
\hypertarget{prog:I:coarse_height}{}Exact Bezout height bounds, six rational doublings and integer isolation\par \cref{ari:C.6.R4} &
\href{https://raw.githubusercontent.com/KokunoYumeto/s6-complex-structure-workbench/\SsixRepositoryCommit/pro_returns/S6_10_PAPERS_checkpoint_15_20261119T1045Z/S6_10_PAPERS/code/I/replay_arithmetic/C_height_coarse_certificate.py}{\nolinkurl{C_height_coarse_certificate.py}} \\[3pt]
\hypertarget{prog:I:leading_factors}{}Genus data, modular symbols, central sums and rational denominator isolation\par \cref{ari:C.6.R5} &
\href{https://raw.githubusercontent.com/KokunoYumeto/s6-complex-structure-workbench/\SsixRepositoryCommit/pro_returns/S6_10_PAPERS_checkpoint_15_20261119T1045Z/S6_10_PAPERS/code/I/replay_arithmetic/B_C6_gz_exact_metadata.gp}{\nolinkurl{B_C6_gz_exact_metadata.gp}}\par \href{https://raw.githubusercontent.com/KokunoYumeto/s6-complex-structure-workbench/\SsixRepositoryCommit/pro_returns/S6_10_PAPERS_checkpoint_15_20261119T1045Z/S6_10_PAPERS/code/I/replay_arithmetic/B_C6_manin_2.gp}{\nolinkurl{B_C6_manin_2.gp}}\par \href{https://raw.githubusercontent.com/KokunoYumeto/s6-complex-structure-workbench/\SsixRepositoryCommit/pro_returns/S6_10_PAPERS_checkpoint_15_20261119T1045Z/S6_10_PAPERS/code/I/replay_arithmetic/B_C6_manin_4.gp}{\nolinkurl{B_C6_manin_4.gp}}\par \href{https://raw.githubusercontent.com/KokunoYumeto/s6-complex-structure-workbench/\SsixRepositoryCommit/pro_returns/S6_10_PAPERS_checkpoint_15_20261119T1045Z/S6_10_PAPERS/code/I/replay_arithmetic/B_C6_manin_6.gp}{\nolinkurl{B_C6_manin_6.gp}}\par \href{https://raw.githubusercontent.com/KokunoYumeto/s6-complex-structure-workbench/\SsixRepositoryCommit/pro_returns/S6_10_PAPERS_checkpoint_15_20261119T1045Z/S6_10_PAPERS/code/I/replay_arithmetic/B_C6_manin_9.gp}{\nolinkurl{B_C6_manin_9.gp}}\par \href{https://raw.githubusercontent.com/KokunoYumeto/s6-complex-structure-workbench/\SsixRepositoryCommit/pro_returns/S6_10_PAPERS_checkpoint_15_20261119T1045Z/S6_10_PAPERS/code/I/replay_arithmetic/B_C6_manin_10.gp}{\nolinkurl{B_C6_manin_10.gp}}\par \href{https://raw.githubusercontent.com/KokunoYumeto/s6-complex-structure-workbench/\SsixRepositoryCommit/pro_returns/S6_10_PAPERS_checkpoint_15_20261119T1045Z/S6_10_PAPERS/code/I/replay_arithmetic/B_C6_central_stream.cpp}{\nolinkurl{B_C6_central_stream.cpp}}\par \href{https://raw.githubusercontent.com/KokunoYumeto/s6-complex-structure-workbench/\SsixRepositoryCommit/pro_returns/S6_10_PAPERS_checkpoint_15_20261119T1045Z/S6_10_PAPERS/code/I/replay_arithmetic/B_C6_exact_factors.py}{\nolinkurl{B_C6_exact_factors.py}}\par \href{https://raw.githubusercontent.com/KokunoYumeto/s6-complex-structure-workbench/\SsixRepositoryCommit/pro_returns/S6_10_PAPERS_checkpoint_15_20261119T1045Z/S6_10_PAPERS/code/I/replay_arithmetic/auxiliary_coefficients.gp}{\nolinkurl{auxiliary_coefficients.gp}} \\[3pt]
\hypertarget{prog:I:torsion}{}Exact local data, torsion and two-isogeny descent over $\mathbb Q(\sqrt3)$\par \cref{ari:C.6.T1,ari:C.6.T2} &
\href{https://raw.githubusercontent.com/KokunoYumeto/s6-complex-structure-workbench/\SsixRepositoryCommit/pro_returns/S6_10_PAPERS_checkpoint_15_20261119T1045Z/S6_10_PAPERS/code/I/replay_torsion/Crank_torsion_local.gp}{\nolinkurl{Crank_torsion_local.gp}}\par \href{https://raw.githubusercontent.com/KokunoYumeto/s6-complex-structure-workbench/\SsixRepositoryCommit/pro_returns/S6_10_PAPERS_checkpoint_15_20261119T1045Z/S6_10_PAPERS/code/I/replay_torsion/Crank_torsion_final.gp}{\nolinkurl{Crank_torsion_final.gp}}\par \href{https://raw.githubusercontent.com/KokunoYumeto/s6-complex-structure-workbench/\SsixRepositoryCommit/pro_returns/S6_10_PAPERS_checkpoint_15_20261119T1045Z/S6_10_PAPERS/code/I/replay_torsion/Crank_torsion_descent.gp}{\nolinkurl{Crank_torsion_descent.gp}}\par \href{https://raw.githubusercontent.com/KokunoYumeto/s6-complex-structure-workbench/\SsixRepositoryCommit/pro_returns/S6_10_PAPERS_checkpoint_15_20261119T1045Z/S6_10_PAPERS/code/I/replay_torsion/Crank_Sage_ell.gp}{\nolinkurl{Crank_Sage_ell.gp}}\par \href{https://raw.githubusercontent.com/KokunoYumeto/s6-complex-structure-workbench/\SsixRepositoryCommit/pro_returns/S6_10_PAPERS_checkpoint_15_20261119T1045Z/S6_10_PAPERS/code/I/replay_torsion/Crank_Sage_ellcommon.gp}{\nolinkurl{Crank_Sage_ellcommon.gp}} \\[3pt]
\hypertarget{prog:I:critical}{}Rational power-series recurrence and Bernstein coefficient inequalities\par \cref{ari:C.8.14} &
\href{https://raw.githubusercontent.com/KokunoYumeto/s6-complex-structure-workbench/\SsixRepositoryCommit/pro_returns/S6_10_PAPERS_checkpoint_15_20261119T1045Z/S6_10_PAPERS/code/I/replay_arithmetic/Ccritical_rational_certificate.py}{\nolinkurl{Ccritical_rational_certificate.py}}\par \href{https://raw.githubusercontent.com/KokunoYumeto/s6-complex-structure-workbench/\SsixRepositoryCommit/pro_returns/S6_10_PAPERS_checkpoint_15_20261119T1045Z/S6_10_PAPERS/code/I/replay_arithmetic/Ccritical_scalar.py}{\nolinkurl{Ccritical_scalar.py}} \\[3pt]
\hypertarget{prog:I:cm_witness}{}Exact CM endomorphism, kernel polynomial and compact-point witness identities\par \cref{ari:C.7.10} &
\href{https://raw.githubusercontent.com/KokunoYumeto/s6-complex-structure-workbench/\SsixRepositoryCommit/pro_returns/S6_SERIES_GRIND_checkpoint_10_20261120T0211Z/code/current_audits/audit_paper_i/arithmetic_06/cm_witness_exact.py}{\nolinkurl{cm_witness_exact.py}} \\[3pt]
Exact sixth- and fourth-division recurrences for the exceptional-fibre calculations\par \cref{fam:src:D.1.1,fam:src:D.1.3} &
\href{https://raw.githubusercontent.com/KokunoYumeto/s6-complex-structure-workbench/\SsixRepositoryCommit/pro_returns/S6_10_PAPERS_checkpoint_15_20261119T1045Z/S6_10_PAPERS/code/I/replay_family_applications/cd_torsion_independent.py}{\nolinkurl{cd_torsion_independent.py}} \\[3pt]
Integral and affine fixed-set Smith/Hermite computations for the three real parameter arcs\par \cref{mapsI:src:A.3.2} &
\href{https://raw.githubusercontent.com/KokunoYumeto/s6-complex-structure-workbench/\SsixRepositoryCommit/pro_returns/S6_10_PAPERS_checkpoint_15_20261119T1045Z/S6_10_PAPERS/code/I/replay_maps/real_involutions.py}{\nolinkurl{real_involutions.py}} \\[3pt]
Denominator-24 finite-root fixed sets and deck permutations\par \cref{mapsI:src:A.3.2} &
\href{https://raw.githubusercontent.com/KokunoYumeto/s6-complex-structure-workbench/\SsixRepositoryCommit/pro_returns/S6_10_PAPERS_checkpoint_15_20261119T1045Z/S6_10_PAPERS/code/I/replay_maps/real_finite_four.py}{\nolinkurl{real_finite_four.py}} \\[3pt]
Thirty-six primitive-ray pairs, twelve fan maps and six curve actions\par \cref{fam:src:C.4.1} &
\href{https://raw.githubusercontent.com/KokunoYumeto/s6-complex-structure-workbench/\SsixRepositoryCommit/pro_returns/S6_SERIES_GRIND_checkpoint_10_20261120T0211Z/code/current_audits/audit_paper_i/code/cusp_hexagon_exact.py}{\nolinkurl{cusp_hexagon_exact.py}} \\[3pt]
Two-period determinant, square-zero monodromies, root differences and exceptional-valuation identities\par \cref{push:src:B.4.24,push:src:B.4.19} &
\href{https://raw.githubusercontent.com/KokunoYumeto/s6-complex-structure-workbench/\SsixRepositoryCommit/pro_returns/S6_02_PUSH_final_20261115T2217Z/S6_02_PUSH/code/e_push_completion_diagnostics.py}{\nolinkurl{e_push_completion_diagnostics.py}} \\[3pt]
Exceptional valuation matrices and intersection determinants for labelled point-blowup sequences through four exceptional curves\par \cref{push:src:B.4.27} &
\href{https://raw.githubusercontent.com/KokunoYumeto/s6-complex-structure-workbench/\SsixRepositoryCommit/pro_returns/S6_02_PUSH_final_20261115T2217Z/S6_02_PUSH/code/e_push_log_blowup.py}{\nolinkurl{e_push_log_blowup.py}} \\[3pt]
Canonical-form coefficients in the small charts, multiplicative limits and horizontal derivatives\par \cref{push:src:B.4.29} &
\href{https://raw.githubusercontent.com/KokunoYumeto/s6-complex-structure-workbench/\SsixRepositoryCommit/pro_returns/S6_02_PUSH_final_20261115T2217Z/S6_02_PUSH/code/e_push_forms_translation.py}{\nolinkurl{e_push_forms_translation.py}} \\[3pt]
Normalization-lattice Gram and monodromy identities; integral Smith forms for the order-three and order-four cycle matrices\par \cref{fam:src:E.1.3,fam:src:E.1.4,fam:src:E.2.1} &
\href{https://raw.githubusercontent.com/KokunoYumeto/s6-complex-structure-workbench/\SsixRepositoryCommit/pro_returns/S6_02_FAMILY/return/code/e_degeneration.py}{\nolinkurl{e_degeneration.py}} \\[3pt]
Character--ray pairings, monoid relations, three affine chart maps, section doubling, normalization Smith forms and virtual Chern pushforward\par \cref{push:src:A.1.1,push:src:A.2.2,push:src:A.3.2,push:src:A.4.1} &
\href{https://raw.githubusercontent.com/KokunoYumeto/s6-complex-structure-workbench/\SsixRepositoryCommit/pro_returns/S6_02_PUSH_final_20261115T2217Z/S6_02_PUSH/code/ab_push_local_checks.py}{\nolinkurl{ab_push_local_checks.py}} \\[3pt]
Breadth-first star examples and local substitutions for exponents $0,1,2$ and centre codimensions $2,3,4$\par \cref{push:src:A.2.5} &
\href{https://raw.githubusercontent.com/KokunoYumeto/s6-complex-structure-workbench/\SsixRepositoryCommit/pro_returns/S6_02_PUSH_final_20261115T2217Z/S6_02_PUSH/code/ab_push_normalization_algorithm.py}{\nolinkurl{ab_push_normalization_algorithm.py}} \\[3pt]
Real-period determinant, covector equation, nonseparation vector and rational-circle generators with denominators at most twelve\par \cref{push:src:D.2.1,push:src:D.2.2,push:src:D.3.1} &
\href{https://raw.githubusercontent.com/KokunoYumeto/s6-complex-structure-workbench/\SsixRepositoryCommit/pro_returns/S6_02_PUSH_final_20261115T2217Z/S6_02_PUSH/code/d_push_exact.py}{\nolinkurl{d_push_exact.py}} \\[3pt]
Finite-series numerical comparison of the theta derivative and symmetric finite differences at three points\par \cref{push:src:D.4.4} &
\href{https://raw.githubusercontent.com/KokunoYumeto/s6-complex-structure-workbench/\SsixRepositoryCommit/pro_returns/S6_02_PUSH_final_20261115T2217Z/S6_02_PUSH/code/d_push_gradient.py}{\nolinkurl{d_push_gradient.py}} \\[3pt]
Exact degree-three endomorphism and kernel congruences; formal logarithmic coefficients and rational derivative coefficient\par \cref{push:src:D.4.5} &
\href{https://raw.githubusercontent.com/KokunoYumeto/s6-complex-structure-workbench/\SsixRepositoryCommit/pro_returns/S6_02_PUSH_final_20261115T2217Z/S6_02_PUSH/code/d_push_equianharmonic.py}{\nolinkurl{d_push_equianharmonic.py}} \\[3pt]
Finite-product numerical derivative comparisons at three displacements near the equianharmonic point\par \cref{push:src:D.4.5} &
\href{https://raw.githubusercontent.com/KokunoYumeto/s6-complex-structure-workbench/\SsixRepositoryCommit/pro_returns/S6_02_PUSH_final_20261115T2217Z/S6_02_PUSH/code/d_push_equianharmonic_crossing.py}{\nolinkurl{d_push_equianharmonic_crossing.py}} \\[3pt]
Integral monodromy and common-fixed equations; root-norm examples for degrees $1$--$64$ and connector cancellation\par \cref{push:src:B.4.32,night:P2.5} &
\href{https://raw.githubusercontent.com/KokunoYumeto/s6-complex-structure-workbench/\SsixRepositoryCommit/pro_returns/S6_02_PUSH_final_20261115T2217Z/S6_02_PUSH/code/d_push_global_loop_defect.py}{\nolinkurl{d_push_global_loop_defect.py}} \\[3pt]
Affine defects of the initial parameter lift and its twelve-shift comparison terms\par \cref{push:src:C.3.5} &
\href{https://raw.githubusercontent.com/KokunoYumeto/s6-complex-structure-workbench/\SsixRepositoryCommit/pro_returns/S6_02_PUSH_final_20261115T2217Z/S6_02_PUSH/code/d_push_monodromy.py}{\nolinkurl{d_push_monodromy.py}} \\[3pt]
Coherent affine composition, finite and cusp defects, period/scalar comparison and twelve-shift factor\par \cref{push:src:C.3.5} &
\href{https://raw.githubusercontent.com/KokunoYumeto/s6-complex-structure-workbench/\SsixRepositoryCommit/pro_returns/S6_02_PUSH_final_20261115T2217Z/S6_02_PUSH/code/d_push_closing.py}{\nolinkurl{d_push_closing.py}} \\[3pt]
Primitive-ray lists at heights $1,2,3$ for $n=1,2,12$; bounded determinant tests and eight successive point stars\par \cref{push:src:A.2.4} &
\href{https://raw.githubusercontent.com/KokunoYumeto/s6-complex-structure-workbench/\SsixRepositoryCommit/pro_returns/S6_02_PUSH_final_20261115T2217Z/S6_02_PUSH/code/ab_push_all_refinements.py}{\nolinkurl{ab_push_all_refinements.py}} \\[3pt]
Integral $E_8$ and neck matrices, Smith form and the square-zero rank-six neck operator\par \cref{push:src:A.4.2} &
\href{https://raw.githubusercontent.com/KokunoYumeto/s6-complex-structure-workbench/\SsixRepositoryCommit/pro_returns/S6_02_PUSH_final_20261115T2217Z/S6_02_PUSH/code/ab_push_monodromy.py}{\nolinkurl{ab_push_monodromy.py}} \\[3pt]
Cone and orbit counts for $n=1,2,12$, eight point stars and five conductor stars\par \cref{push:src:A.4.6,push:src:A.4.7} &
\href{https://raw.githubusercontent.com/KokunoYumeto/s6-complex-structure-workbench/\SsixRepositoryCommit/pro_returns/S6_02_PUSH_final_20261115T2217Z/S6_02_PUSH/code/ab_push_all_fan_topology.py}{\nolinkurl{ab_push_all_fan_topology.py}} \\[3pt]
Normalization-chain Smith forms for $n=1,2,3,4,6,12$ and blowup-centre increments\par \cref{push:src:A.4.8} &
\href{https://raw.githubusercontent.com/KokunoYumeto/s6-complex-structure-workbench/\SsixRepositoryCommit/pro_returns/S6_02_PUSH_final_20261115T2217Z/S6_02_PUSH/code/ab_push_all_fan_homology.py}{\nolinkurl{ab_push_all_fan_homology.py}} \\[3pt]
Local exceptional complexes, unit Smith factors and Gysin maps for centre codimensions $2,3,4$\par \cref{push:src:A.4.9} &
\href{https://raw.githubusercontent.com/KokunoYumeto/s6-complex-structure-workbench/\SsixRepositoryCommit/pro_returns/S6_02_PUSH_final_20261115T2217Z/S6_02_PUSH/code/ab_push_central_pushforward.py}{\nolinkurl{ab_push_central_pushforward.py}} \\[3pt]
Augmented-simplex differentials and signed permutation chain maps for at most four branches\par \cref{push:src:B.4.5} &
\href{https://raw.githubusercontent.com/KokunoYumeto/s6-complex-structure-workbench/\SsixRepositoryCommit/pro_returns/S6_02_PUSH_final_20261115T2217Z/S6_02_PUSH/code/b_nc_degree_one.py}{\nolinkurl{b_nc_degree_one.py}} \\[3pt]
Unipotent winding, common fixed cycles, coinvariant Smith form and the exact affine closing identity\par \cref{push:src:B.4.9} &
\href{https://raw.githubusercontent.com/KokunoYumeto/s6-complex-structure-workbench/\SsixRepositoryCommit/pro_returns/S6_02_PUSH_final_20261115T2217Z/S6_02_PUSH/code/ab_push_hopf_winding.py}{\nolinkurl{ab_push_hopf_winding.py}} \\[3pt]
Four affine smoothness calculations, Hilbert difference, normal-degree cancellations and intersection Smith form\par \cref{push:src:B.4.12} &
\href{https://raw.githubusercontent.com/KokunoYumeto/s6-complex-structure-workbench/\SsixRepositoryCommit/pro_returns/S6_02_PUSH_final_20261115T2217Z/S6_02_PUSH/code/root_quadric_smoothing.py}{\nolinkurl{root_quadric_smoothing.py}} \\[3pt]
Fixed-curve and node character identities, tangent weights and vector-field identities for $n=1,\ldots,12$\par \cref{push:src:B.4.14} &
\href{https://raw.githubusercontent.com/KokunoYumeto/s6-complex-structure-workbench/\SsixRepositoryCommit/pro_returns/S6_02_PUSH_final_20261115T2217Z/S6_02_PUSH/code/ab_push_kato_weights.py}{\nolinkurl{ab_push_kato_weights.py}} \\[3pt]
Integral chain maps, elementary collapses and angular cocycles for $n=4,6,12,24$\par \cref{push:src:B.4.16} &
\href{https://raw.githubusercontent.com/KokunoYumeto/s6-complex-structure-workbench/\SsixRepositoryCommit/pro_returns/S6_02_PUSH_final_20261115T2217Z/S6_02_PUSH/code/ab_push_regular_cap_dual.py}{\nolinkurl{ab_push_regular_cap_dual.py}} \\[3pt]
Deck differences, twelve-step affine norm, commuting shears and thirty-six primitive-cone determinants\par \cref{push:src:B.4.18} &
\href{https://raw.githubusercontent.com/KokunoYumeto/s6-complex-structure-workbench/\SsixRepositoryCommit/pro_returns/S6_02_PUSH_final_20261115T2217Z/S6_02_PUSH/code/ab_push_corrected_cusp_overlap.py}{\nolinkurl{ab_push_corrected_cusp_overlap.py}} \\[3pt]
Finite and cusp gauge identities; node substitutions and the additive fan cocycle on the punctured parameter domain\par \cref{push:src:B.4.19} &
\href{https://raw.githubusercontent.com/KokunoYumeto/s6-complex-structure-workbench/\SsixRepositoryCommit/pro_returns/S6_02_PUSH_final_20261115T2217Z/S6_02_PUSH/code/ab_push_global_partial.py}{\nolinkurl{ab_push_global_partial.py}} \\[3pt]
Cartier and Tate factors for $n=1,2,3,4,8,12$, local base-map differentials and the closing factor\par \cref{push:src:B.4.20} &
\href{https://raw.githubusercontent.com/KokunoYumeto/s6-complex-structure-workbench/\SsixRepositoryCommit/pro_returns/S6_02_PUSH_final_20261115T2217Z/S6_02_PUSH/code/ab_push_tate_interface.py}{\nolinkurl{ab_push_tate_interface.py}} \\[3pt]
Point and curve blowup substitutions, node exponents, intersection degrees and finite collar-product identities\par \cref{push:src:B.4.21} &
\href{https://raw.githubusercontent.com/KokunoYumeto/s6-complex-structure-workbench/\SsixRepositoryCommit/pro_returns/S6_02_PUSH_final_20261115T2217Z/S6_02_PUSH/code/d_push_degree_one_repair.py}{\nolinkurl{d_push_degree_one_repair.py}} \\[3pt]
Three local chart maps, triangular fans for $n=1,2,3,4,6,12$ and collar identities\par \cref{push:src:B.4.22} &
\href{https://raw.githubusercontent.com/KokunoYumeto/s6-complex-structure-workbench/\SsixRepositoryCommit/pro_returns/S6_02_PUSH_final_20261115T2217Z/S6_02_PUSH/code/ab_push_corner_resolution.py}{\nolinkurl{ab_push_corner_resolution.py}} \\[3pt]
Exact cone, ray and fan calculations\par \cref{push:src:A.2.1} &
\href{https://raw.githubusercontent.com/KokunoYumeto/s6-complex-structure-workbench/\SsixRepositoryCommit/pro_returns/S6_10_PAPERS_checkpoint_15_20261119T1045Z/S6_10_PAPERS/code/I/replay_family_applications/ab_push_fan.py}{\nolinkurl{ab_push_fan.py}} \\[3pt]
Characteristic polynomial, integral inverse identity and both Euler calculations (\cref{push:src:B.3.7}) & \href{https://raw.githubusercontent.com/KokunoYumeto/s6-complex-structure-workbench/\SsixRepositoryCommit/pro_returns/S6_02_PUSH_final_20261115T2217Z/S6_02_PUSH/code/ab_push_auxiliary_plumbing.py}{\texttt{ab\_push\_auxiliary\_plumbing.py}} \\
Rational fan certificates for the base changes $n=2$ and $n=12$ (\cref{push:src:B.2.1}) & \href{https://raw.githubusercontent.com/KokunoYumeto/s6-complex-structure-workbench/\SsixRepositoryCommit/pro_returns/S6_02_PUSH_final_20261115T2217Z/S6_02_PUSH/code/ab_push_basechange_fan.py}{\texttt{ab\_push\_basechange\_fan.py}} \\
Finite period covariance, affine deck formulas, component cycles and peripheral residues (\cref{push:src:B.1.2,push:src:B.3.1}) & \href{https://raw.githubusercontent.com/KokunoYumeto/s6-complex-structure-workbench/\SsixRepositoryCommit/pro_returns/S6_02_PUSH_final_20261115T2217Z/S6_02_PUSH/code/ab_push_finite_decks.py}{\texttt{ab\_push\_finite\_decks.py}} \\
Finite cyclic carry cocycles and local component arithmetic (\cref{push:src:B.3.6}) & \href{https://raw.githubusercontent.com/KokunoYumeto/s6-complex-structure-workbench/\SsixRepositoryCommit/pro_returns/S6_02_PUSH_final_20261115T2217Z/S6_02_PUSH/code/ab_push_torsor_cocycles.py}{\texttt{ab\_push\_torsor\_cocycles.py}} \\
Node series, height and intersection lattices, trace congruences and integral saturation (\cref{push:src:B.3.4,push:src:B.3.5}) & \href{https://raw.githubusercontent.com/KokunoYumeto/s6-complex-structure-workbench/\SsixRepositoryCommit/pro_returns/S6_02_PUSH_final_20261115T2217Z/S6_02_PUSH/code/ab_push_trace_obstruction.py}{\texttt{ab\_push\_trace\_obstruction.py}} \\
Section equations and collinearity identities on $b^3=a^4+1$ (\cref{push:src:B.3.4}) & \href{https://raw.githubusercontent.com/KokunoYumeto/s6-complex-structure-workbench/\SsixRepositoryCommit/pro_returns/S6_02_PUSH_final_20261115T2217Z/S6_02_PUSH/code/ab_push_trace_sections.py}{\texttt{ab\_push\_trace\_sections.py}} \\
\end{longtable}
}

\clearpage

\begin{thebibliography}{99}
\providecommand{\natexlab}[1]{#1}
\providecommand{\doi}[1]{\url{https://doi.org/#1}}
\bibitem[Abasheva and D{\'e}ev(2024)]{AbashevaDeev2024AB}
Anna Abasheva and Rodion D{\'e}ev.

\newblock Complex surfaces with many algebraic structures.

\newblock \emph{International Mathematics Research Notices}, 2024\penalty0
  (9):\penalty0 7379--7400, 2024.

\newblock \doi{10.1093/imrn/rnad190}.

\newblock URL \url{https://arxiv.org/pdf/2303.10764v1}.

\newblock Cited preprint edition: arXiv:2303.10764v1, 19 March 2023; preprint page and statement references use that edition.

\bibitem[Abate et~al.(2013)Abate, Bracci, Suwa, and Tovena]{AbateBracciSuwaTovena2010Atiyah}
Marco Abate, Filippo Bracci, Tatsuo Suwa, and Francesca Tovena.

\newblock Localization of Atiyah classes.

\newblock \emph{Revista Matem\'atica Iberoamericana}, 29\penalty0 (2):\penalty0
  547--578, 2013.

\newblock \doi{10.4171/RMI/730}.

\newblock The cited edition is the author manuscript dated 7 May 2010;
  page and statement references use that edition.

\newblock URL \url{https://pagine.dm.unipi.it/abate/articoli/artric/files/LocalAtiClasses.pdf}.


\bibitem[Abramovich et~al.(2002)Abramovich, Karu, Matsuki, and
  W{\l}odarczyk]{AKMW02AB}
Dan Abramovich, Kalle Karu, Kenji Matsuki, and Jaros{\l}aw W{\l}odarczyk.

\newblock Torification and factorization of birational maps.

\newblock \emph{Journal of the American Mathematical Society}, 15\penalty0
  (3):\penalty0 531--572, 2002.

\newblock \doi{10.1090/S0894-0347-02-00396-X}.

\newblock URL \url{https://personal.math.ubc.ca/~karu/papers/semist.pdf}.

\bibitem[Adamczewski and Faverjon(2017)]{adamczewski2017mahler}
Boris Adamczewski and Colin Faverjon.

\newblock M\'ethode de {Mahler}: relations lin\'eaires, transcendance et
  applications aux nombres automatiques.

\newblock \emph{Proceedings of the London Mathematical Society}, 115:\penalty0
  55--90, 2017.

\newblock \doi{10.1112/plms.12038}.

\newblock URL
  \url{https://faverjon.perso.math.cnrs.fr/AdamczewskiFaverjon_mahler1.pdf}.

\bibitem[Agboola and Howard(2006)]{agboola2006}
Adebisi Agboola and Benjamin Howard.

\newblock Anticyclotomic {Iwasawa} theory of {CM} elliptic curves.

\newblock \emph{Annales de l'Institut Fourier}, 56\penalty0 (4):\penalty0
  1001--1048, 2006.

\newblock \doi{10.5802/aif.2206}.

\newblock URL
  \url{https://web.math.ucsb.edu/~agboola/papers/bisi_aif_2006.pdf}.

\bibitem[Angella(2018)]{Angella2018}
Daniele Angella.

\newblock Hodge numbers of a hypothetical complex structure on {$S^6$}.

\newblock \emph{Differential Geometry and its Applications}, 57:\penalty0
  105--120, 2018.

\newblock \doi{10.1016/j.difgeo.2017.10.012}.

\newblock URL \url{https://arxiv.org/abs/1705.10518v2}.

\newblock Cited preprint edition: arXiv:1705.10518v2, 10 November 2017; preprint page and statement references use that edition.

\bibitem[Atiyah and Singer(1968)]{AtiyahSinger1968IIIE}
Michael~F. Atiyah and Isadore~M. Singer.

\newblock The index of elliptic operators. III.

\newblock \emph{Annals of Mathematics}, 87\penalty0 (3):\penalty0 546--604,
  1968.

\newblock \doi{10.2307/1970717}.

\bibitem[B{\u a}nic{\u a} and St{\u a}n{\u a}{\c s}il{\u a}(1976)]{BS76}
Constantin B{\u a}nic{\u a} and Octavian St{\u a}n{\u a}{\c s}il{\u a}.

\newblock \emph{Algebraic Methods in the Global Theory of Complex Spaces}.

\newblock Editura Academiei and John Wiley \& Sons, Bucharest and London, 1976.

\bibitem[Belabas et~al.(2019)Belabas, Bernardi, and
  Perrin-Riou]{CgzBelabas2019}
Karim Belabas, Dominique Bernardi, and Bernadette Perrin-Riou.

\newblock La constante de {Manin} et le degr\'e modulaire d'une courbe
  elliptique.

\newblock \emph{Publications math\'ematiques de Besan\c con}, 2019/2:\penalty0
  81--103, 2019.

\newblock \doi{10.5802/pmb.37}.

\newblock URL \url{https://www.numdam.org/item/10.5802/pmb.37.pdf}.

\bibitem[Bertolin et~al.(2022)Bertolin, Philippon, Saha, and
  Saha]{bertolin2022}
Cristiana Bertolin, Patrice Philippon, Biswajyoti Saha, and Ekata Saha.

\newblock Semi-abelian analogues of {Schanuel} conjecture and applications.

\newblock \emph{Journal of Algebra}, 596:\penalty0 250--288, 2022.

\newblock \doi{10.1016/j.jalgebra.2021.12.040}.

\newblock URL \url{https://arxiv.org/pdf/2010.15170v2}.

\newblock The cited version is arXiv:2010.15170v2, 15 November 2022; page and statement references use that version.

\bibitem[Bertolini et~al.(2012)Bertolini, Darmon, and Prasanna]{CrBDP2012}
Massimo Bertolini, Henri Darmon, and Kartik Prasanna.

\newblock {$p$}-adic {Rankin} {L}-series and rational points on {CM} elliptic
  curves.

\newblock \emph{Pacific Journal of Mathematics}, 260\penalty0 (2):\penalty0
  261--303, 2012.

\newblock \doi{10.2140/pjm.2012.260.261}.

\newblock URL
  \url{https://www.math.mcgill.ca/darmon/pub/Articles/Research/52.BDP2/bdp-pjm.pdf}.

\bibitem[Bertrand(1995)]{bertrand1995}
Daniel Bertrand.

\newblock Points rationnels sur les sous-groupes compacts des groupes
  alg{\'e}briques.

\newblock \emph{Experimental Mathematics}, 4\penalty0 (2):\penalty0 145--151,
  1995.

\newblock \doi{10.1080/10586458.1995.10504315}.

\newblock URL
  \url{https://www.maths.tcd.ie/EMIS/journals/EM/restricted/4/4.2/bertrand.ps}.

\bibitem[Bertrand(2017)]{bertrand2017}
Daniel Bertrand.

\newblock Transcendence problems related to heights.

\newblock Fields Institute, Heights and Applications to Unlikely Intersections,
  February 13--17, 2017.

\newblock URL
  \url{https://webusers.imj-prg.fr/~daniel.bertrand/Recherche/rpdf/DB_Fieldspost.pdf}.

\bibitem[Bierstone and Milman(1997)]{BM97}
Edward Bierstone and Pierre~D. Milman.

\newblock Canonical desingularization in characteristic zero by blowing up the
  maximum strata of a local invariant.

\newblock \emph{Inventiones Mathematicae}, 128:\penalty0 207--302, 1997.

\newblock \doi{10.1007/s002220050141}.

\newblock Earlier Chapter~I preprint: arXiv:alg-geom/9508005v2 (1995), with separate pagination. \newblock \url{https://arxiv.org/abs/alg-geom/9508005v2}.

\bibitem[Bost and K{\"u}nnemann(2009)]{bost2009}
Jean-Beno{\^i}t Bost and Klaus K{\"u}nnemann.

\newblock Hermitian vector bundles and extension groups on arithmetic schemes
  {II}. {The} arithmetic {Atiyah} extension.

\newblock In \emph{From probability to geometry (I)---Volume in honor of the 60th birthday of Jean-Michel Bismut}, number 327 of \emph{Ast{\'e}risque}, pages 361--424. Soci{\'e}t{\'e} math{\'e}matique de France, 2009.

\newblock URL \url{https://www.numdam.org/article/AST_2009__327__361_0.pdf}.

\bibitem[Bourdon et~al.(2017)Bourdon, Clark, and
  Stankewicz]{BourdonClarkStankewicz2017}
Abbey Bourdon, Pete~L. Clark, and James Stankewicz.

\newblock Torsion points on {CM} elliptic curves over real number fields.

\newblock \emph{Transactions of the American Mathematical Society},
  369\penalty0 (12):\penalty0 8457--8496, 2017.

\newblock \doi{10.1090/tran/6905}.

\newblock URL \url{https://arxiv.org/pdf/1411.2742v3}.

\newblock Cited preprint edition: arXiv:1411.2742v3, 25 September 2015; preprint page and statement references use that edition.

\bibitem[Buchdahl(1999)]{Buchdahl1999}
Nicholas Buchdahl.

\newblock On compact Kähler surfaces.

\newblock \emph{Annales de l'Institut Fourier}, 49\penalty0 (1):\penalty0
  287--302, 1999.

\newblock \doi{10.5802/aif.1674}.

\newblock URL \url{https://www.numdam.org/article/AIF_1999__49_1_287_0.pdf}.

\bibitem[Bump et~al.(1990)Bump, Friedberg, and Hoffstein]{BumpFriedbergHoffstein1990}
Daniel Bump, Solomon Friedberg, and Jeffrey Hoffstein.

\newblock Nonvanishing theorems for {$L$}-functions of modular forms and their derivatives.

\newblock \emph{Inventiones Mathematicae}, 102:\penalty0 543--618, 1990.

\newblock \doi{10.1007/BF01233440}.

\bibitem[Cai et~al.(2014)Cai, Shu, and Tian]{CgzCai2014}
Li~Cai, Jie Shu, and Ye~Tian.

\newblock Explicit {Gross--Zagier} and {Waldspurger} formulae.

\newblock \emph{Algebra \& Number Theory}, 8\penalty0 (10):\penalty0
  2523--2572, 2014.

\newblock \doi{10.2140/ant.2014.8.2523}.

\newblock URL \url{https://msp.org/ant/2014/8-10/ant-v8-n10-p05-p.pdf}.

\bibitem[Call and Silverman(1993)]{callSilverman1993}
Gregory~S. Call and Joseph~H. Silverman.

\newblock Canonical heights on varieties with morphisms.

\newblock \emph{Compositio Mathematica}, 89\penalty0 (2):\penalty0 163--205,
  1993.

\newblock URL \url{https://www.numdam.org/item/CM_1993__89_2_163_0.pdf}.

\bibitem[Cartan(1952)]{Cartan1952B_AB}
Henri Cartan.

\newblock Faisceaux analytiques sur les vari{\'e}t{\'e}s de {Stein}:
  d{\'e}monstration des th{\'e}or{\`e}mes fondamentaux.

\newblock \emph{S{\'e}minaire Henri Cartan}, volume 4 (1951--1952), expos{\'e} 19, pages 1--15. Secr{\'e}tariat math{\'e}matique, Paris, 1952.

\newblock URL \url{https://www.numdam.org/article/SHC_1951-1952__4__A19_0.pdf}.

\bibitem[Castella(2025)]{CrCastella2025}
Francesc Castella.

\newblock Tamagawa number conjecture for {CM} modular forms and
  {R}ankin--{S}elberg convolutions.

\newblock \emph{Proceedings of the London Mathematical Society}, 131\penalty0
  (4):\penalty0 e70089, 2025.

\newblock \doi{10.1112/plms.70089}.

\newblock URL \url{https://web.math.ucsb.edu/~castella/Katz.pdf}.

\newblock Page references use the linked author version dated 14 February 2025.

\bibitem[Chen(2026)]{Chen2026HodgeRoot}
Kuan-Wen Chen.

\newblock On the {Hodge} structures of global smoothings of normal crossing
  varieties, 2026.

\newblock URL \url{https://arxiv.org/abs/2404.19229v2}.

\bibitem[{Claude} and {L. Alp{\"o}ge}(2026)]{AC2026}
{Claude} and {L. Alp{\"o}ge}.

\newblock The $(3,4,\infty)$ modular family of 2-tori, completed at its three
  special points, is a complex structure on {$S^6$}.

\newblock 108-page manuscript, 2026.

\newblock URL \url{https://alpo.ge/s6.pdf}.

\bibitem[Conrad(n.d.)]{ConradSL2Z}
Keith Conrad.

\newblock $\operatorname{SL}_2(\mathbb Z)$.

\newblock Author's notes, undated, 19-page version retrieved
  21 November 2026.

\newblock URL
  \url{https://kconrad.math.uconn.edu/blurbs/grouptheory/SL(2,Z).pdf}.

\bibitem[Cornulier(2011)]{Cornulier2011}
Yves Cornulier.

\newblock On the {Chabauty} space of locally compact abelian groups.

\newblock \emph{Algebraic \& Geometric Topology}, 11(4):2007--2035, 2011.

\newblock \doi{10.2140/agt.2011.11.2007}.

\bibitem[Cossec et~al.(2024)Cossec, Dolgachev, and Liedtke]{CDL24}
Fran\c{c}ois Cossec, Igor Dolgachev, and Christian Liedtke.

\newblock \emph{Enriques surfaces {I}}. Author manuscript, 19 April 2024
  (the cited pagination).

\newblock Second edition: Springer Singapore, 2025.
  \doi{10.1007/978-981-96-1214-7}.

\newblock Manuscript URL
  \url{https://sites.lsa.umich.edu/idolga/wp-content/uploads/sites/1334/2024/08/EnriquesOne.pdf}.

\bibitem[Cremona and Siksek(2006)]{CrCremonaSiksek2006}
John~E. Cremona and Samir Siksek.

\newblock Computing a lower bound for the canonical height on elliptic curves
  over {$\mathbb Q$}.

\newblock In \emph{Algorithmic Number Theory}, volume 4076 of \emph{Lecture
  Notes in Computer Science}, pages 275--286. Springer, 2006.

\newblock \doi{10.1007/11792086_20}.

\newblock URL \url{https://johncremona.github.io/papers/HeightLB.pdf}.

\newblock The linked author version has unnumbered pages; its locators use PDF page numbers.

\bibitem[Darmon(2009)]{CrDarmon2009}
Henri Darmon.

\newblock Cycles on modular varieties and rational points on elliptic curves,
  2009.

\newblock In \emph{Explicit Methods in Number Theory}, \emph{Oberwolfach
  Reports}, 6\penalty0 (3):\penalty0 1871--1885, 2009.

\newblock Report DOI: \doi{10.4171/OWR/2009/33}.

\newblock URL
  \url{https://www.math.mcgill.ca/darmon/pub/Articles/Expository/13.Oberwolfach/paper.pdf}.

\newblock Author lecture notes dated 31 July 2009; the cited pagination is that of this version.

\bibitem[de~Cataldo et~al.(2018)de~Cataldo, Migliorini, and
  Musta\c{t}\u{a}]{DMM16}
Mark Andrea~A. de~Cataldo, Luca Migliorini, and Mircea Musta\c{t}\u{a}.

\newblock Combinatorics and topology of proper toric maps.

\newblock {\em Journal f{\"u}r die reine und angewandte Mathematik}, 744:133--163, 2018.
\newblock \doi{10.1515/crelle-2015-0104}.
\newblock Cited edition (page and statement numbering): arXiv:1407.3497v2, 18 January 2016.
\newblock URL \url{https://arxiv.org/pdf/1407.3497v2}.

\bibitem[Demailly(2012)]{Demailly2012E}
Jean-Pierre Demailly.

\newblock Complex analytic and differential geometry, 2012.

\newblock URL
  \url{https://www-fourier.univ-grenoble-alpes.fr/~demailly/manuscripts/agbook.pdf}.

\newblock Cited author edition dated 21 June 2012; page and statement references use that edition.

\bibitem[DeMarco and Mavraki(2020)]{demarco2017}
Laura DeMarco and Niki~Myrto Mavraki.

\newblock Variation of canonical height and equidistribution.

\newblock \emph{American Journal of Mathematics}, 142\penalty0 (2):\penalty0
  443--473, 2020.

\newblock \doi{10.1353/ajm.2020.0012}.

\newblock URL \url{https://arxiv.org/pdf/1701.07947v2}.

\newblock The cited version is arXiv:1701.07947v2, 2 March 2017; page and statement references use that version.

\bibitem[Dloussky(2006)]{Dloussky2005AB}
Georges Dloussky.

\newblock On surfaces of class {VII}$_0^+$ with numerically anticanonical
  divisor.

\newblock {\em American Journal of Mathematics}, 128(3):639--670, 2006.
\newblock \doi{10.1353/ajm.2006.0021}.
\newblock Cited edition (page and statement numbering): arXiv:math/0406387v3, 13 April 2005.
\newblock URL \url{https://arxiv.org/pdf/math/0406387v3}.

\bibitem[Dloussky and Teleman(2012)]{DlousskyTeleman2012AB}
Georges Dloussky and Andrei Teleman.

\newblock Infinite bubbling in non-{K}{\"a}hlerian geometry.

\newblock \emph{Mathematische Annalen}, 353\penalty0 (4):\penalty0 1283--1314,
  2012.

\newblock \doi{10.1007/s00208-011-0713-9}.

\newblock URL \url{https://arxiv.org/pdf/1012.5247v2}.

\newblock Cited preprint edition: arXiv:1012.5247v2, 5 October 2011; preprint page and statement references use that edition.

\bibitem[Dloussky and Teleman(2020)]{DlousskyTeleman2020D}
Georges Dloussky and Andrei Teleman.

\newblock Smooth deformations of singular contractions of class {VII} surfaces.

\newblock \emph{Mathematische Zeitschrift}, 296:\penalty0 1521--1537, 2020.

\newblock \doi{10.1007/s00209-020-02481-0}.

\newblock URL \url{https://arxiv.org/pdf/1803.07631v2}.

\newblock Cited preprint edition: arXiv:1803.07631v2, 28 February 2020; preprint page and statement references use that edition.

\bibitem[Douady(1974)]{Douady1974}
Adrien Douady.

\newblock Le probl\`eme des modules locaux pour les espaces
  $\mathbf C$-analytiques compacts.

\newblock \emph{Annales scientifiques de l'\'Ecole Normale Sup\'erieure},
  (4) 7(4):569--602, 1974.

\newblock \doi{10.24033/asens.1280}.
\newblock URL
  \url{https://www.numdam.org/article/ASENS_1974_4_7_4_569_0.pdf}.

\bibitem[Drinfeld(1973)]{Drinfeld1973}
V.~G. Drinfeld.

\newblock Two theorems on modular curves.

\newblock \emph{Funktsional. Anal. i Prilozhen.}, 7(2):83--84, 1973 (Russian).

\newblock English translation: \emph{Functional Analysis and Its Applications}, 7(2):155--156, 1973.

\newblock \doi{10.1007/BF01078890}.

\newblock \url{https://www.mathnet.ru/eng/faa2422}.

\bibitem[Du et~al.(2026{\natexlab{a}})Du, Guo, Kusakabe, and Xie]{DuEtAl2026v1}
Yun-Heng Du, Bin Guo, Yuta Kusakabe, and Song-Yan Xie.

\newblock Oka geometry and deformations of torus-fibred complex six-spheres.

\newblock arXiv:2610.00250v1, 2026{\natexlab{a}}.

\newblock URL \url{https://arxiv.org/abs/2610.00250v1}.

\bibitem[Du et~al.(2026{\natexlab{b}})Du, Guo, Kusakabe, and Xie]{DuEtAl2026v2}
Yun-Heng Du, Bin Guo, Yuta Kusakabe, and Song-Yan Xie.

\newblock Oka geometry of complex six-spheres.

\newblock arXiv:2610.00250v2, 2026{\natexlab{b}}.

\newblock URL \url{https://arxiv.org/abs/2610.00250v2}.

\bibitem[Duke and Imamo{\u g}lu(2008)]{DukeImamoglu2008}
William Duke and {\"O}zlem Imamo{\u g}lu.

\newblock The zeros of the {Weierstrass} {$\wp$}-function and hypergeometric
  series.

\newblock \emph{Mathematische Annalen}, 340\penalty0 (4):\penalty0 897--905,
  2008.

\newblock \doi{10.1007/s00208-007-0174-3}.

\newblock URL \url{https://www.math.ucla.edu/~wdduke/preprints/zeros.pdf}.
\newblock The cited pages are those of the author preprint: Theorem~1 and equation~(5), p.~2; equation~(6), p.~4.

\bibitem[Eichler and Zagier(1982)]{EichlerZagier1982}
Martin Eichler and Don Zagier.

\newblock On the zeros of the {Weierstrass} {$\wp$}-function.

\newblock \emph{Mathematische Annalen}, 258\penalty0 (4):\penalty0 399--407,
  1982.

\newblock \doi{10.1007/BF01453974}.

\newblock URL
  \url{https://people.mpim-bonn.mpg.de/zagier/files/doi/10.1007/BF01453974/fulltext.pdf}.

\bibitem[Einsiedler et~al.(2001)Einsiedler, Everest, and
  Ward]{einsiedler2001entropy}
Manfred Einsiedler, Graham Everest, and Thomas Ward.

\newblock Entropy and the canonical height.

\newblock \emph{Journal of Number Theory}, 91\penalty0 (2):\penalty0 256--273,
  2001.

\newblock \doi{10.1006/jnth.2001.2682}.

\newblock URL \url{https://people.math.ethz.ch/~einsiedl/height.pdf}.

\bibitem[Engel et~al.(2025)Engel, de~Gaay~Fortman, and Schreieder]{EGFS25}
Philip Engel, Olivier de~Gaay~Fortman, and Stefan Schreieder.

\newblock Combinatorics and {H}odge theory of degenerations of abelian
  varieties: A survey of the {M}umford construction, 2025.

\newblock URL \url{https://arxiv.org/html/2507.15695v3}.

\newblock Cited edition: arXiv:2507.15695v3, 27 March 2026; the first version appeared in 2025.

\bibitem[Engel(2026)]{Engel2026}
Philip Engel.

\newblock Complex structures on {$S^6$}.

\newblock arXiv:2609.38442v1, 2026.

\newblock URL \url{https://arxiv.org/abs/2609.38442v1}.

\bibitem[Enokizono and Hashizume(2025)]{EH2024E}
Makoto Enokizono and Kenta Hashizume.

\newblock Semistable reduction for complex analytic spaces.

\newblock {\em Transactions of the American Mathematical Society}, 378(11):7667--7688, 2025.
\newblock \doi{10.1090/tran/9365}.
\newblock Cited edition (page and statement numbering): arXiv:2404.05124v3, 12 October 2024.
\newblock URL \url{https://arxiv.org/pdf/2404.05124v3}.

\bibitem[Erban and Van~Gorder(2021)]{ErbanVanGorder2021}
Radek Erban and Robert~A. Van~Gorder.

\newblock On the sum of positive divisors functions.

\newblock \emph{Research in Number Theory}, 7, Article~25, 2021.

\newblock \doi{10.1007/s40993-021-00240-6}.

\newblock URL \url{https://people.maths.ox.ac.uk/erban/papers/paperRNT.pdf}.


\bibitem[Fantechi and G{\"o}ttsche(2010)]{FG10}
Barbara Fantechi and Lothar G{\"o}ttsche.

\newblock Riemann--Roch theorems and elliptic genus for virtually smooth
  schemes.

\newblock \emph{Geometry \& Topology}, 14\penalty0 (1):\penalty0 83--115, 2010.

\newblock \doi{10.2140/gt.2010.14.83}.

\newblock URL \url{https://msp.org/gt/2010/14-1/gt-v14-n1-p02-p.pdf}.

\bibitem[Friedman(2016)]{Friedman2016AB}
Robert Friedman.

\newblock On the geometry of anticanonical pairs.

\newblock Cited edition (page and statement numbering): arXiv:1502.02560v2, 22 July 2016.
\newblock URL \url{https://arxiv.org/pdf/1502.02560v2}.

\bibitem[Fujiki(2009)]{Fujiki2009D}
Akira Fujiki.

\newblock Automorphisms of parabolic {Inoue} surfaces.

\newblock Cited edition (page and statement numbering): arXiv:0903.5374v1, 31 March 2009.
\newblock URL \url{https://arxiv.org/pdf/0903.5374v1}.

\bibitem[Fujiki and Pontecorvo(2010)]{FujikiPontecorvo2010AB}
Akira Fujiki and Massimiliano Pontecorvo.

\newblock Anti-self-dual bihermitian structures on {Inoue} surfaces.

\newblock \emph{Journal of Differential Geometry}, 85\penalty0 (1):\penalty0
  15--71, 2010.

\newblock \doi{10.4310/jdg/1284557925}.

\newblock URL
  \url{https://archive.ymsc.tsinghua.edu.cn/pacm_download/21/834-JDG-85-1-A2-fujiki.pdf}.

\bibitem[Fujimoto(1990)]{Fujimoto1990AB}
Yoshio Fujimoto.

\newblock Logarithmic transformations on elliptic fiber spaces. II.---elliptic bundle case.

\newblock In K.~Konno, M.-H. Saito, and S.~Usui, editors, \emph{Proceedings of
  the Meeting and the Workshop ``Algebraic Geometry and Hodge Theory'', Vol.
  I}, volume~16 of \emph{Hokkaido University Technical Report Series in
  Mathematics}, pages 81--102. Hokkaido University, 1990.

\newblock Collection DOI: \doi{10.14943/5135}.

\newblock Journal publication under the same title: \emph{Mathematische
  Annalen}, 288:\penalty0 559--570, 1990; \doi{10.1007/BF01444550}.

\newblock Page and statement references use the conference contribution
  on pp.~81--102 in the linked collection.

\newblock URL
  \url{https://eprints.lib.hokudai.ac.jp/repo/huscap/all/5450/16.pdf}.

\bibitem[Gasbarri(2010)]{gasbarri2010analytic}
Carlo Gasbarri.

\newblock Analytic subvarieties with many rational points.

\newblock \emph{Mathematische Annalen}, 346\penalty0 (1):\penalty0 199--243,
  2010.

\newblock \doi{10.1007/s00208-009-0394-9}.

\newblock URL \url{https://preprints.ihes.fr/storage/newscheiderlang.pdf}.

\bibitem[Gauduchon(1985)]{Gauduchon1985}
Paul Gauduchon.

\newblock Les m\'etriques standard d'une surface complexe compacte \`a premier
  nombre de {Betti} pair.

\newblock In \emph{G{\'e}om{\'e}trie des surfaces K3: modules et p{\'e}riodes---S{\'e}minaire Palaiseau}, number 126 of \emph{Ast\'erisque}, pages 129--135. Soci{\'e}t{\'e} math{\'e}matique de France, 1985.

\newblock URL \url{https://numdam.org/item/AST_1985__126__129_0/}.

\bibitem[Graf and Kov{\'a}cs(2014)]{GrafKovacs2014Root}
Patrick Graf and S{\'a}ndor~J. Kov{\'a}cs.

\newblock Potentially {Du Bois} spaces.

\newblock \emph{Journal of Singularities}, 8:\penalty0 117--134, 2014.

\newblock \doi{10.5427/jsing.2014.8i}.

\newblock URL
  \url{https://sites.math.washington.edu/~kovacs/current/papers/Graf_Kovacs__Potentially_Du_Bois_spaces.pdf}.

\bibitem[Gramain et~al.(1986)Gramain, Mignotte, and
  Waldschmidt]{gramain1986values}
Fran\c{c}ois Gramain, Maurice Mignotte, and Michel Waldschmidt.

\newblock Valeurs alg\'ebriques de fonctions analytiques.

\newblock \emph{Acta Arithmetica}, 47\penalty0 (2):\penalty0 97--121, 1986.

\newblock \doi{10.4064/aa-47-2-97-121}.

\newblock URL \url{https://matwbn.icm.edu.pl/ksiazki/aa/aa47/aa4721.pdf}.

\bibitem[Grauert(1960)]{Grauert60}
Hans Grauert.

\newblock Ein Theorem der analytischen Garbentheorie und die Modulräume
  komplexer Strukturen.

\newblock \emph{Publications Mathématiques de l'IHÉS}, 5:\penalty0 5--64,
  1960.

\newblock \doi{10.1007/BF02684746}.

\newblock URL \url{https://www.numdam.org/article/PMIHES_1960__5__5_0.pdf}.

\bibitem[Grauert(1963)]{Grauert1963ErratumD}
Hans Grauert.

\newblock Berichtigung: Ein Theorem der analytischen Garbentheorie und die
  Modulräume komplexer Strukturen.

\newblock \emph{Publications Mathématiques de l'IHÉS}, 16:\penalty0 35--36,
  1963.

\newblock \doi{10.1007/BF02684291}.

\newblock URL \url{https://www.numdam.org/item/PMIHES_1963__16__35_0.pdf}.

\bibitem[Gross and Zagier(1986)]{CrGrossZagier1986}
Benedict~H. Gross and Don~B. Zagier.

\newblock Heegner points and derivatives of {L}-series.

\newblock \emph{Inventiones Mathematicae}, 84\penalty0 (2):\penalty0 225--320,
  1986.

\newblock \doi{10.1007/BF01388809}.

\newblock URL
  \url{https://archive.mpim-bonn.mpg.de/id/eprint/3578/1/preprint_1985_30.pdf}.

\newblock The linked text is the MPI/SFB 85-30 preprint (1985); citations marked as preprint pages use its pagination.

\bibitem[Gross(1991)]{CrGross1991}
Benedict~H. Gross.

\newblock Kolyvagin's work on modular elliptic curves.

\newblock In \emph{L-functions and Arithmetic}, volume 153 of \emph{London
  Mathematical Society Lecture Note Series}, pages 235--256. Cambridge
  University Press, 1991.

\newblock \doi{10.1017/CBO9780511526053.009}.


\bibitem[Herblot(2012)]{herblot2012meromorphic}
Mathilde Herblot.

\newblock Algebraic points on meromorphic curves, 2012.

\newblock URL \url{https://webusers.imj-prg.fr/~mathilde.herblot/publi/SL.pdf}.

\newblock Cited author manuscript dated 18 April 2012; page and statement references use that edition.

\bibitem[Hochschild and Serre(1953)]{HS53}
Gerhard Hochschild and Jean-Pierre Serre.

\newblock Cohomology of group extensions.

\newblock \emph{Transactions of the American Mathematical Society}, 74\penalty0
  (1):\penalty0 110--134, 1953.

\newblock \doi{10.1090/S0002-9947-1953-0052438-8}.

\bibitem[Honda and Viaclovsky(2025)]{HondaViaclovsky2025D}
Nobuhiro Honda and Jeff Viaclovsky.

\newblock Some remarks on fibrations in complex geometry.

\newblock \emph{Journal of Mathematical Study}, 58\penalty0 (4):\penalty0
  552--574, 2025.

\newblock \doi{10.4208/jms.v58n4.25.11}.

\newblock URL \url{https://arxiv.org/abs/2504.14420v1}.

\newblock Theorem~1.1 on p.~553 is cited from the journal edition; references marked arXiv version~1 use arXiv:2504.14420v1, 19 April 2025.

\bibitem[Huber and W{\"u}stholz(2022)]{huber2022}
Annette Huber and Gisbert W{\"u}stholz.

\newblock \emph{Transcendence and linear relations of 1-periods}, volume 227 of
  \emph{Cambridge Tracts in Mathematics}.

\newblock Cambridge University Press, 2022.

\newblock \doi{10.1017/9781009019729}.
\newblock The analytic subgroup theorem is cited from the separately paginated
  \href{https://arxiv.org/abs/1805.10104v8}{arXiv:1805.10104v8}, Theorem~6.1, p.~43.

\bibitem[Illman(1972)]{Illman1972}
S{\"o}ren Illman.

\newblock \emph{Equivariant Algebraic Topology}.

\newblock PhD thesis, Princeton University, 1972.

\bibitem[Illusie et~al.(2005)Illusie, Kato, and Nakayama]{IKN2005E}
Luc Illusie, Kazuya Kato, and Chikara Nakayama.

\newblock Quasi-unipotent logarithmic {Riemann--Hilbert} correspondences.

\newblock \emph{Journal of Mathematical Sciences, the University of Tokyo},
  12\penalty0 (1):\penalty0 1--66, 2005.

\newblock URL \url{https://www.ms.u-tokyo.ac.jp/journal/pdf/jms120101.pdf}.

\bibitem[Illusie et~al.(2007)Illusie, Kato, and Nakayama]{IKN2007ErratumE}
Luc Illusie, Kazuya Kato, and Chikara Nakayama.

\newblock Erratum to ``quasi-unipotent logarithmic {Riemann--Hilbert}
  correspondences''.

\newblock \emph{Journal of Mathematical Sciences, the University of Tokyo},
  14\penalty0 (1):\penalty0 113--116, 2007.

\newblock URL \url{https://www.ms.u-tokyo.ac.jp/journal/pdf/jms140105.pdf}.

\bibitem[Istrati et~al.(2022)Istrati, Otiman, Pontecorvo, and
  Ruggiero]{IstratiEtAl2022Root}
Nicolina Istrati, Alexandra Otiman, Massimiliano Pontecorvo, and Matteo
  Ruggiero.

\newblock Toric {Kato} manifolds.

\newblock \emph{Journal de l'{\'E}cole polytechnique---Math{\'e}matiques},
  9:\penalty0 1347--1395, 2022.

\newblock \doi{10.5802/jep.208}.

\newblock URL \url{https://www.numdam.org/item/JEP_2022__9__1347_0.pdf}.

\bibitem[Jetchev et~al.(2009)Jetchev, Lauter, and Stein]{CrJetchev2008}
Dimitar Jetchev, Kristin Lauter, and William Stein.

\newblock Explicit {H}eegner points: {K}olyvagin's conjecture and non-trivial
  elements in the {S}hafarevich--{T}ate group.

\newblock \emph{Journal of Number Theory}, 129\penalty0 (2):\penalty0 284--302, 2009.

\newblock \doi{10.1016/j.jnt.2008.05.007}.

\newblock URL \url{https://arxiv.org/pdf/0707.0032v1}.

\newblock The cited version is arXiv:0707.0032v1, 30 June 2007; page references use that version.

\bibitem[Kajiwara et~al.(2008{\natexlab{a}})Kajiwara, Kato, and
  Nakayama]{KKN2008PicardE}
Takeshi Kajiwara, Kazuya Kato, and Chikara Nakayama.

\newblock Analytic log {Picard} varieties.

\newblock \emph{Nagoya Mathematical Journal}, 191:\penalty0 149--180,
  2008{\natexlab{a}}.

\newblock \doi{10.1017/S0027763000025940}.

\newblock URL
  \url{https://www.cambridge.org/core/services/aop-cambridge-core/content/view/E8B336E3252DE1F448583CC5865DB1DA/S0027763000025940a.pdf/analytic_log_picard_varieties.pdf}.

\bibitem[Kajiwara et~al.(2008{\natexlab{b}})Kajiwara, Kato, and
  Nakayama]{KKN2008Root}
Takeshi Kajiwara, Kazuya Kato, and Chikara Nakayama.

\newblock Logarithmic abelian varieties, {Part I}: Complex analytic theory.

\newblock \emph{Journal of Mathematical Sciences, the University of Tokyo},
  15\penalty0 (1):\penalty0 69--193, 2008{\natexlab{b}}.

\newblock URL \url{https://www.ms.u-tokyo.ac.jp/journal/pdf/jms150104.pdf}.

\bibitem[Knudson(2004)]{Knudson2004}
Kevin~P. Knudson.

\newblock Invariant chains and the homology of quotient spaces.

\newblock Cited edition (page and statement numbering): arXiv:math/0409133v1, 8 September 2004.
\newblock URL \url{https://arxiv.org/pdf/math/0409133v1}.

\bibitem[Levy(2008)]{Levy08}
Roni~N. Levy.

\newblock Riemann--Roch theorem for higher bivariant {K}-functors.

\newblock \emph{Annales de l'Institut Fourier}, 58\penalty0 (2):\penalty0
  571--601, 2008.

\newblock \doi{10.5802/aif.2361}.

\newblock URL \url{https://www.numdam.org/article/AIF_2008__58_2_571_0.pdf}.

\bibitem[Lin and Wang(2010)]{LinWang2010}
Chang-Shou Lin and Chin-Lung Wang.

\newblock Elliptic functions, Green functions and the mean field equations on tori.

\newblock \emph{Annals of Mathematics}, 172(2):911--954, 2010.

\newblock \doi{10.4007/annals.2010.172.911}.

\newblock URL
\url{https://annals.math.princeton.edu/wp-content/uploads/annals-v172-n2-p03-p.pdf}.

\bibitem[\L{}ojasiewicz(1964)]{Loj64}
Stanis{\l}aw \L{}ojasiewicz.

\newblock Triangulation of semi-analytic sets.

\newblock \emph{Annali della Scuola Normale Superiore di Pisa, Classe di
  Scienze}, Series~3, 18\penalty0 (4):\penalty0 449--474, 1964.

\newblock URL \url{https://www.numdam.org/item/ASNSP_1964_3_18_4_449_0.pdf}.

\bibitem[Manetti(1991)]{Manetti1991D}
Marco Manetti.

\newblock Normal degenerations of the complex projective plane.

\newblock \emph{Journal f{\"u}r die reine und angewandte Mathematik},
  419:\penalty0 89--118, 1991.

\newblock \doi{10.1515/crll.1991.419.89}.

\newblock URL \url{https://gdz.sub.uni-goettingen.de/id/PPN243919689_0419}.

\bibitem[Matsuki(2000)]{Matsuki00ErrAB}
Kenji Matsuki.

\newblock Correction: ``A note on the factorization theorem of toric
  birational maps after {Morelli} and its toroidal extension''.

\newblock {\em Tohoku Mathematical Journal}, 52(4):629--631, 2000.
\newblock \doi{10.2748/tmj/1178207758}.
\newblock Cited edition (page and statement numbering): arXiv:math/0006038v1, 5 June 2000.
\newblock URL \url{https://arxiv.org/abs/math/0006038v1}.

\bibitem[Mazur et~al.(2006)Mazur, Stein, and Tate]{MazurSteinTate2006}
Barry Mazur, William Stein, and John Tate.

\newblock Computation of $p$-adic heights and log convergence.

\newblock \emph{Documenta Mathematica}, Extra Volume Coates:577--614, 2006.

\newblock URL \url{https://www.maths.tcd.ie/EMIS/journals/DMJDMV/vol-coates/mazur_stein_tate.pdf}.
\newblock \doi{10.4171/DMS/4/17}.

\bibitem[McHugh(2019)]{McHugh2018}
Andrew McHugh.

\newblock Bott--Chern--Aeppli, Dolbeault, and Fr{\"o}licher on compact
  complex $3$-folds.

\newblock \emph{Rivista di Matematica della Universit\`a di Parma},
  10(1):25--62, 2019.

\newblock URL
  \url{https://www.rivmat.unipr.it/fulltext/2019-10-1/Riv_Parma_10-1_2019_02.pdf}.
\newblock The locators in this paper refer to the separately paginated
  arXiv:1708.03251v4, 15 November 2018.
\newblock Cited edition: \url{https://arxiv.org/pdf/1708.03251v4}.

\bibitem[McMullen(1996)]{McMullen1996Complex}
Curtis~T. McMullen.

\newblock Complex manifolds: Math 241 course outline.

\newblock University of California, Berkeley, course notes, 1996.

\newblock URL
  \url{https://people.math.harvard.edu/~ctm/home/text/class/berkeley/241/96/course/course.pdf}.

\bibitem[Merkulov(1995)]{Merkulov1995D}
Sergey~A. Merkulov.

\newblock Moduli spaces of compact complex submanifolds of complex fibered
  manifolds.

\newblock \emph{Mathematical Proceedings of the Cambridge Philosophical
  Society}, 118\penalty0 (1):\penalty0 71--91, 1995.

\newblock \doi{10.1017/S0305004100073473}.


\bibitem[Miller(2011)]{miller2011}
Robert~L. Miller.

\newblock Proving the {Birch} and {Swinnerton-Dyer} conjecture for specific
  elliptic curves of analytic rank zero and one.

\newblock \emph{LMS Journal of Computation and Mathematics}, 14:\penalty0
  327--350, 2011.

\newblock \doi{10.1112/S1461157011000180}.

\newblock URL \url{https://arxiv.org/pdf/1010.2431v3}.

\newblock Cited preprint edition: arXiv:1010.2431v3, 21 December 2011; preprint page and statement references use that edition.

\bibitem[Movasati(2024)]{Movasati24}
Hossein Movasati.

\newblock A course in complex geometry and holomorphic foliations.

\newblock Author version dated 13 May 2024.

\newblock URL \url{https://w3.impa.br/~hossein/myarticles/ComplexGeometry.pdf}.

\bibitem[Nakayama and Ogus(2010)]{NO10}
Chikara Nakayama and Arthur Ogus.

\newblock Relative rounding in toric and logarithmic geometry.

\newblock \emph{Geometry \& Topology}, 14\penalty0 (4):\penalty0 2189--2241,
  2010.

\newblock \doi{10.2140/gt.2010.14.2189}.

\newblock URL \url{https://msp.org/gt/2010/14-4/gt-v14-n4-p09-s.pdf}.

\bibitem[{National Institute of Standards and Technology}(2026)]{DLMFTheta2026}
{National Institute of Standards and Technology}.

\newblock {NIST Digital Library of Mathematical Functions}, chapters 20 and 23,
  Version~1.2.8, released 15 September 2026. Accessed 21 November 2026.

\newblock URL \url{https://dlmf.nist.gov/}.

\bibitem[Neretin(2024)]{Neretin2024Beta}
Yury A. Neretin.

\newblock On the Dotsenko--Fateev complex twin of the Selberg integral and its extensions.

\newblock \emph{The Ramanujan Journal}, 64\penalty0 (1):\penalty0 37--55, 2024.

\newblock \doi{10.1007/s11139-023-00804-3}.

\newblock URL \url{https://link.springer.com/content/pdf/10.1007/s11139-023-00804-3.pdf}.

\bibitem[Nguyen(2023)]{nguyen2023}
Khoa~D. Nguyen.

\newblock Transcendence of polynomial canonical heights.

\newblock \emph{Mathematische Annalen}, 387(1--2):\penalty0 1--15, 2023.

\newblock URL \url{https://arxiv.org/pdf/2112.14937v1}.
\newblock The cited edition is arXiv:2112.14937v1, 30 December 2021.
\newblock \doi{10.1007/s00208-022-02465-x}.

\bibitem[Nishiguchi(1988)]{Nishiguchi1988IIRoot}
Kenji Nishiguchi.

\newblock Degeneration of {K3} surfaces, {II}.

\newblock Max-Planck-Institut f{\"u}r Mathematik preprint MPI/88-48, 1988.

\newblock URL
  \url{https://archive.mpim-bonn.mpg.de/id/eprint/3390/1/preprint_1988_48.pdf}.

\bibitem[Noguchi(2024)]{noguchi2024}
Junjiro Noguchi.

\newblock Analytic {Ax--Schanuel} for semi-abelian varieties and {Nevanlinna}
  theory.

\newblock \emph{Journal of the Mathematical Society of Japan}, 76\penalty0
  (1):\penalty0 1--22, 2024.

\newblock \doi{10.2969/jmsj/89588958}.

\newblock URL \url{https://www.ms.u-tokyo.ac.jp/~noguchi/An-axsch.pdf}.

\bibitem[Oda(1978)]{Oda1978D}
Tadao Oda.

\newblock \emph{Lectures on Torus Embeddings and Applications}.

\newblock Number~58 in Tata Institute of Fundamental Research Lectures on
  Mathematics and Physics. Springer-Verlag, Berlin, Heidelberg, New York, 1978.

\newblock URL
  \url{https://mathweb.tifr.res.in/Documents/Publications/Lectures/tifr58.pdf}.
\newblock Citations use the page numbers of this official TIFR retypeset edition.

\bibitem[Parusi{\'n}ski(1994)]{Parusinski1990E}
Adam Parusi{\'n}ski.

\newblock Subanalytic functions.

\newblock \emph{Transactions of the American Mathematical Society},
  344\penalty0 (2):\penalty0 583--595, 1994.

\newblock \doi{10.1090/S0002-9947-1994-1160156-4}.

\newblock The cited edition is Technical Report MPI/90-58,
  Max-Planck-Institut f{\"u}r Mathematik, 1990; manuscript page and statement
  references use that edition.

\newblock URL
  \url{https://archive.mpim-bonn.mpg.de/id/eprint/1622/1/preprint_1990_58.pdf}.

\bibitem[Pinkham(1974)]{Pinkham1974Root}
Henry~C. Pinkham.

\newblock \emph{Deformations of algebraic varieties with {$G_m$} action}.

\newblock Number~20 in Ast\'erisque. Soci\'et\'e math\'ematique de France,
  1974.

\newblock URL \url{https://www.numdam.org/item/AST_1974__20__1_0/}.

\bibitem[Ramis and Ruget, 1970]{RR70push}
J.-P. Ramis and G. Ruget.

\newblock Complexe dualisant et th{\'e}or{\`e}mes de dualit{\'e} en
  g{\'e}om{\'e}trie analytique complexe.

\newblock {\em Publications Math{\'e}matiques de l'IH{\'E}S}, 38:77--91, 1970.

\newblock \doi{10.1007/BF02684652}.

\bibitem[Ramis et~al., 1971]{RRV71}
J.-P. Ramis, G. Ruget, and J.-L. Verdier.

\newblock Dualit{\'e} relative en g{\'e}om{\'e}trie analytique complexe.

\newblock {\em Inventiones Mathematicae}, 13:261--283, 1971.

\newblock \doi{10.1007/BF01406078}.

\bibitem[Reyssat(1982)]{reyssat1982weierstrass}
Eric Reyssat.

\newblock Propri\'et\'es d'ind\'ependance alg\'ebrique de nombres li\'es aux
  fonctions de Weierstrass.

\newblock \emph{Acta Arithmetica}, 41\penalty0 (3):\penalty0 291--310, 1982.

\newblock \doi{10.4064/aa-41-3-291-310}.

\newblock URL \url{https://matwbn.icm.edu.pl/ksiazki/aa/aa41/aa4137.pdf}.

\bibitem[Roy and Waldschmidt(1993)]{roy1993periodrank}
Damien Roy and Michel Waldschmidt.

\newblock Autour du th\'eor\`eme du sous-groupe alg\'ebrique.

\newblock \emph{Canadian Mathematical Bulletin}, 36\penalty0 (3):\penalty0
  358--367, 1993.

\newblock \doi{10.4153/CMB-1993-049-8}.

\newblock URL
  \url{https://www.cambridge.org/core/journals/canadian-mathematical-bulletin/article/autour-du-theoreme-du-sousgroupe-algebrique/A1E35F3798690D084664ACD83D1557E8}.

\bibitem[Rubin(1999)]{CrRubin1999}
Karl Rubin.

\newblock Elliptic curves with complex multiplication and the conjecture of
  {B}irch and {S}winnerton-{D}yer.

\newblock In \emph{Arithmetic Theory of Elliptic Curves}, volume 1716 of
  \emph{Lecture Notes in Mathematics}, pages 167--234. Springer, 1999.

\newblock URL \url{https://swc-math.github.io/aws/1999/99RubinCM.pdf}.

\newblock Page references use the linked lecture-note version, paginated 1--68.
\newblock \doi{10.1007/BFb0093455}.

\bibitem[Schnell(2019)]{Schnell19}
Christian Schnell.

\newblock Lecture 4: The {Appel--Humbert} theorem.

\newblock MAT 615, Stony Brook University, lecture of 6 February, 2019.

\newblock URL
  \url{https://www.math.stonybrook.edu/~cschnell/mat615/lectures/lecture4.pdf}.

\bibitem[Schr{\"o}er and Siebert(2002)]{SchroeerSiebert2002Root}
Stefan Schr{\"o}er and Bernd Siebert.

\newblock Irreducible degenerations of primary {Kodaira} surfaces.

\newblock In \emph{Complex geometry ({G}{\"o}ttingen, 2000)}, pages 193--222.
  Springer, Berlin, 2002.

\newblock URL
  \url{https://www.math.uni-duesseldorf.de/~schroeer/publications_pdf/kodaira.pdf}.
\newblock \doi{10.1007/978-3-642-56202-0_14}.

\bibitem[Shin(2015)]{Shin2015}
Dong~Hwa Shin.

\newblock Generation of class fields by {Siegel--Ramachandra} invariants.

\newblock \emph{Journal of the Korean Mathematical Society}, 52\penalty0
  (5):\penalty0 907--928, 2015.

\newblock URL \url{https://arxiv.org/pdf/1009.2253v6}.
\newblock Page references use arXiv:1009.2253v6, 18 October 2014.
\newblock \doi{10.4134/JKMS.2015.52.5.907}.

\bibitem[Shioda(1990)]{Shioda90}
Tetsuji Shioda.

\newblock On the {M}ordell--{W}eil lattices.

\newblock \emph{Commentarii Mathematici Universitatis Sancti Pauli},
  39\penalty0 (2):\penalty0 211--240, 1990.

\newblock \doi{10.14992/00009994}.

\newblock URL
  \url{https://rikkyo.repo.nii.ac.jp/record/10027/files/AA00610867_39-02_05.pdf}.

\bibitem[Silverman(2006)]{silverman2006heights}
Joseph~H. Silverman.

\newblock An introduction to height functions.

\newblock MSRI workshop lecture slides, January 2006.

\newblock URL
  \url{https://slmath-production-assets.sfo3.digitaloceanspaces.com/attachments/workshops/301/HtSurveyMSRIJan06.pdf}.

\bibitem[Stein(2007)]{CgzStein2007}
William Stein.

\newblock \emph{Modular Forms: A Computational Approach}, volume~79 of
  \emph{Graduate Studies in Mathematics}.

\newblock American Mathematical Society, 2007.

\newblock URL \url{https://www.wstein.org/books/modform/stein-modform.pdf}.
\newblock \doi{10.1090/gsm/079}.

\bibitem[Stelzig(2021)]{Stelzig2021}
Jonas Stelzig.

\newblock On the structure of double complexes.

\newblock \emph{Journal of the London Mathematical Society}, 104\penalty0
  (2):\penalty0 956--988, 2021.

\newblock \doi{10.1112/jlms.12453}.

\bibitem[Sutherland(2023)]{Sutherland2023CM}
Andrew~V. Sutherland.

\newblock The {Hilbert} class polynomial.

\newblock MIT 18.783 Elliptic Curves, Lecture 20, 21 November 2023.

\newblock URL
  \url{https://math.mit.edu/classes/18.783/2023/LectureNotes20.pdf}.

\bibitem[Sz\'ekelyhidi et~al.(2017)Sz\'ekelyhidi, Tosatti, and
  Weinkove]{STW2017}
G\'abor Sz\'ekelyhidi, Valentino Tosatti, and Ben Weinkove.

\newblock Gauduchon metrics with prescribed volume form.

\newblock \emph{Acta Mathematica}, 219(1):181--211, 2017.

\newblock \doi{10.4310/ACTA.2017.v219.n1.a6}.

\newblock Earlier manuscript: arXiv:1503.04491v2, 8 December 2017, with separate pagination.

\newblock URL \url{https://arxiv.org/abs/1503.04491v2}.

\bibitem[Teleman(2010)]{Teleman2010AB}
Andrei Teleman.

\newblock Instantons and curves on class {VII} surfaces.

\newblock \emph{Annals of Mathematics}, 172\penalty0 (3):\penalty0 1749--1804,
  2010.

\newblock \doi{10.4007/annals.2010.172.1749}.

\newblock URL
  \url{https://annals.math.princeton.edu/wp-content/uploads/annals-v172-n3-p07-p.pdf}.

\bibitem[Viaclovsky(2026{\natexlab{a}})]{ViaclovskyLog2026}
Jeff Viaclovsky.

\newblock Logarithmic transformations on the six-sphere.

\newblock arXiv:2610.39102v1, 2026{\natexlab{a}}.

\newblock URL \url{https://arxiv.org/abs/2610.39102v1}.

\bibitem[Viaclovsky(2026{\natexlab{b}})]{ViaclovskyTwo2026}
Jeff Viaclovsky.

\newblock A two-parameter family of complex structures on {$S^6$}.

\newblock arXiv:2609.33785v1, 2026{\natexlab{b}}.

\newblock URL \url{https://arxiv.org/abs/2609.33785v1}.

\bibitem[Waldschmidt(1975)]{waldschmidt1975dimension}
Michel Waldschmidt.

\newblock Dimension alg{\'e}brique de sous-groupes analytiques de
  vari{\'e}t{\'e}s de groupe.

\newblock \emph{Annales de l'Institut Fourier}, 25\penalty0 (1):\penalty0
  23--33, 1975.

\newblock \doi{10.5802/aif.543}.

\newblock URL
  \url{https://webusers.imj-prg.fr/~michel.waldschmidt/articles/pdf/Fourier.pdf}.

\bibitem[Waldschmidt(1979)]{waldschmidt1979sigma}
Michel Waldschmidt.

\newblock Nombres transcendants et fonctions sigma de {Weierstrass}.

\newblock \emph{Comptes Rendus Math\'ematiques de l'Acad\'emie des Sciences,
  Canada}, 1\penalty0 (2):\penalty0 111--114, 1979.

\newblock URL
  \url{https://webusers.imj-prg.fr/~michel.waldschmidt/articles/pdf/GDL/SigmaNoteCanada1979-fr.pdf}.

\bibitem[Waldschmidt(1988{\natexlab{a}})]{waldschmidt1988periodrank}
Michel Waldschmidt.

\newblock On the transcendence methods of {Gelfond} and {Schneider} in several
  variables.

\newblock In Alan Baker, editor, \emph{New Advances in Transcendence Theory},
  chapter~24, pages 375--398. Cambridge University Press, 1988{\natexlab{a}}.

\newblock URL
  \url{https://webusers.imj-prg.fr/~michel.waldschmidt/articles/pdf/Durham1986TMGSSV.pdf}.
\newblock \doi{10.1017/CBO9780511897184.025}.

\bibitem[Waldschmidt(1988{\natexlab{b}})]{waldschmidt1988theta}
Michel Waldschmidt.

\newblock Fonctions th\^eta et nombres transcendants.

\newblock S\'eminaire de Th\'eorie des Nombres de Bordeaux, 1987--1988,
  expos\'e 15, pp. 1--7, 1988{\natexlab{b}}.

\newblock URL
  \url{https://webusers.imj-prg.fr/~michel.waldschmidt/articles/pdf/PPN320141322_0016__0_0.pdf}.

\bibitem[Waldschmidt(2009)]{waldschmidt2009auxiliary}
Michel Waldschmidt.

\newblock Auxiliary functions in transcendental number theory.

\newblock \emph{The Ramanujan Journal}, 20\penalty0 (3):\penalty0 341--373,
  2009.

\newblock \doi{10.1007/s11139-009-9204-y}.

\newblock URL
  \url{https://webusers.imj-prg.fr/~michel.waldschmidt/articles/pdf/SASTRAaftp.pdf}.

\bibitem[Xiao(1991)]{Xiao91}
Gang Xiao.

\newblock Elliptic surfaces with fixed Jacobian.

\newblock Technical Report MPI/91-30, Max-Planck-Institut f{\"u}r Mathematik,
  1991.

\newblock URL
  \url{https://archive.mpim-bonn.mpg.de/id/eprint/2953/1/preprint_1991_30.pdf}.

\bibitem[Xie(2024)]{xie2024algebraicity}
Junyi Xie.

\newblock Algebraicity criteria, invariant subvarieties and transcendence
  problems from arithmetic dynamics.

\newblock \emph{Peking Mathematical Journal}, 7:\penalty0 345--398, 2024.

\newblock \doi{10.1007/s42543-022-00059-9}.

\newblock URL \url{https://arxiv.org/pdf/2202.09032v1}.

\newblock Cited preprint edition: arXiv:2202.09032v1, 18 February 2022; preprint page and statement references use that edition.

\bibitem[Yang(2026)]{Yang2026}
Xiufan Yang.

\newblock Compact complex threefolds fibred by two-dimensional complex tori.

\newblock arXiv:2609.29666v1, 2026.

\newblock URL \url{https://arxiv.org/abs/2609.29666v1}.

\bibitem[Yuan et~al.(2013)Yuan, Zhang, and Zhang]{CgzYZZ2011}
Xinyi Yuan, Shou-Wu Zhang, and Wei Zhang.

\newblock \emph{The Gross--Zagier Formula on Shimura Curves}, volume~184 of
  \emph{Annals of Mathematics Studies}.

\newblock Princeton University Press, Princeton, 2013.

\newblock \doi{10.1515/9781400845644}.

\newblock The cited edition is the preprint dated 6 November 2011;
  page and statement references use that edition.

\newblock Cited preprint: \url{https://www.researchgate.net/publication/267551160_Gross-Zagier_Formula_On_Shimura_Curves}.

\end{thebibliography}

\end{document}
